
% =================================================================
% REGULATIONS FOR PREPARING THE GRADUATION INTERNSHIP REPORT
% USING LATEX
% =================================================================
%
% This file is intended to be included from the main document with:
%
%     
% =================================================================
% REGULATIONS FOR PREPARING THE GRADUATION INTERNSHIP REPORT
% USING LATEX
% =================================================================
%
% This file is intended to be included from the main document with:
%
%     
% =================================================================
% REGULATIONS FOR PREPARING THE GRADUATION INTERNSHIP REPORT
% USING LATEX
% =================================================================
%
% This file is intended to be included from the main document with:
%
%     
% =================================================================
% REGULATIONS FOR PREPARING THE GRADUATION INTERNSHIP REPORT
% USING LATEX
% =================================================================
%
% This file is intended to be included from the main document with:
%
%     \input{Regulations.tex}
%
% IMPORTANT:
% This section is provided for guidance only. Students shall remove
% or comment out the \input{Regulations.tex} command before submitting
% the final internship report.
%
% =================================================================

\clearpage
\thispagestyle{empty}

\begin{center}
    {\fontsize{16pt}{20pt}\selectfont\bfseries
    SPECIFIC REGULATIONS FOR PREPARING\\
    GRADUATION INTERNSHIP REPORTS USING \LaTeX}
\end{center}

\vspace{0.5cm}

{\fontsize{13pt}{17pt}\selectfont
\setstretch{1.3}
\justifying

\noindent
\textbf{Note:}
This section is provided for guidance only and is included in the main
document using the command \verb|\input{Regulations.tex}|.
Students shall remove or comment out this command before submitting the
final internship report.

\vspace{0.4cm}

% ================================================================
% I. GENERAL REQUIREMENTS
% ================================================================

\noindent
\textbf{I.\quad GENERAL REQUIREMENTS}

\vspace{0.15cm}

The Graduation Internship Report shall be prepared using the official
\LaTeX{} template provided by the Faculty. Students shall preserve the
predefined page layout, typography, line spacing, heading styles,
numbering system, figure and table formatting, and bibliography settings
unless otherwise instructed by the Faculty or the supervisor.

The report shall be written in English.

Students are responsible for the academic and technical content of the
report. The official \LaTeX{} template is responsible for standardized
document formatting.

\vspace{0.3cm}

% ================================================================
% II. REPORT STRUCTURE
% ================================================================

\noindent
\textbf{II.\quad REPORT STRUCTURE}

\vspace{0.15cm}

The report shall normally be arranged in the following order:

\begin{enumerate}[label=\arabic*., leftmargin=1.2cm, itemsep=2pt]
    \item Outer Cover;
    \item Comments of the Internship Host Organization;
    \item Supervisor's Evaluation and Graduation Internship Defense Results;
    \item Table of Contents;
    \item List of Symbols and Abbreviations, if applicable;
    \item List of Figures;
    \item List of Tables;
    \item Abstract;
    \item Main Chapters of the Internship Report;
    \item References;
    \item Appendices, if applicable.
\end{enumerate}

The preliminary pages shall not be numbered as chapters. The main report
content shall be organized into consecutively numbered chapters beginning
with Chapter 1.

\vspace{0.3cm}

% ================================================================
% III. PAGE SIZE AND MARGINS
% ================================================================

\noindent
\textbf{III.\quad PAGE SIZE AND MARGINS}

\vspace{0.15cm}

The report shall be formatted on A4 paper with the following margins:

\begin{center}
\begin{tabular}{ll}
    Left margin:   & 3.5 cm \\
    Right margin:  & 2.0 cm \\
    Top margin:    & 2.5 cm \\
    Bottom margin: & 2.5 cm
\end{tabular}
\end{center}

These settings are predefined in the official template and shall not be
modified by students.

\vspace{0.3cm}

% ================================================================
% IV. FONT, FONT SIZE, AND LINE SPACING
% ================================================================

\noindent
\textbf{IV.\quad FONT, FONT SIZE, AND LINE SPACING}

\vspace{0.15cm}

The main text shall use a 13 pt font size with 1.3-line spacing.
A Times-family typeface shall be used in accordance with the official
template.

The text shall be fully justified.

Students shall not use repeated spaces, repeated manual line breaks, or
multiple blank lines to position text. Paragraph and heading spacing shall
be controlled by the template.

\vspace{0.3cm}

% ================================================================
% V. OUTER COVER
% ================================================================

\noindent
\textbf{V.\quad OUTER COVER}

\vspace{0.15cm}

The outer cover shall follow the official format provided in the template
and shall contain the following information:

\begin{itemize}[leftmargin=0.8cm, itemsep=2pt]
    \item Thai Nguyen University of Technology;
    \item Faculty of International Training;
    \item Engineering and Technology Department;
    \item University logo and Faculty logo;
    \item the title \textbf{INTERNSHIP REPORT};
    \item Internship Organization;
    \item Author;
    \item Class;
    \item Supervisor;
    \item place and year of completion.
\end{itemize}

Students shall only replace the required information fields and shall not
modify the cover-page layout unless instructed otherwise.

\vspace{0.3cm}

% ================================================================
% VI. COMMENTS OF THE INTERNSHIP HOST ORGANIZATION
% ================================================================

\noindent
\textbf{VI.\quad COMMENTS OF THE INTERNSHIP HOST ORGANIZATION}

\vspace{0.15cm}

The page entitled \textbf{COMMENTS OF THE INTERNSHIP HOST ORGANIZATION}
is reserved for the internship host organization.

It includes:

\begin{itemize}[leftmargin=0.8cm, itemsep=2pt]
    \item the student's full name;
    \item class;
    \item comments from the host organization;
    \item score;
    \item date;
    \item signature and official seal.
\end{itemize}

Students shall not alter or complete official evaluation fields unless
specifically authorized.

\vspace{0.3cm}

% ================================================================
% VII. SUPERVISOR'S EVALUATION AND DEFENSE RESULTS
% ================================================================

\noindent
\textbf{VII.\quad SUPERVISOR'S EVALUATION AND DEFENSE RESULTS}

\vspace{0.15cm}

The section entitled
\textbf{SUPERVISOR'S EVALUATION AND GRADUATION INTERNSHIP DEFENSE RESULTS}
is reserved for the supervisor and defense committee.

The evaluation may include:

\begin{itemize}[leftmargin=0.8cm, itemsep=2pt]
    \item Internship Progress and Student's Attitude;
    \item Report Content;
    \item Report Presentation;
    \item Other Comments;
    \item Supervisor's Opinion;
    \item Defense Result;
    \item Score.
\end{itemize}

Students shall not enter evaluation scores or defense results before they
are officially determined.

\vspace{0.3cm}

% ================================================================
% VIII. CHAPTERS, SECTIONS, AND SUBSECTIONS
% ================================================================

\noindent
\textbf{VIII.\quad CHAPTERS, SECTIONS, AND SUBSECTIONS}

\vspace{0.15cm}

The report structure shall be generated using the standard \LaTeX{}
commands:

\begin{verbatim}
\chapter{...}
\section{...}
\subsection{...}
\end{verbatim}

Chapter, section, and subsection numbers shall be generated automatically
by \LaTeX{}.

Students shall not manually type chapter or section numbers and shall not
reproduce heading styles using commands such as \verb|\textbf|,
\verb|\fontsize|, or \verb|\centering| instead of the appropriate
structural command.

Excessive subdivision should be avoided. In most cases, the hierarchy
Chapter--Section--Subsection is sufficient.

\vspace{0.3cm}

% ================================================================
% IX. ORGANIZATION OF THE MAIN CONTENT
% ================================================================

\noindent
\textbf{IX.\quad ORGANIZATION OF THE MAIN CONTENT}

\vspace{0.15cm}

The specific content of the report shall depend on the internship
organization and the assigned tasks. A typical structure may include:

\begin{itemize}[leftmargin=0.8cm, itemsep=2pt]
    \item \textbf{Chapter 1. Overview of the Internship Organization};
    \item \textbf{Chapter 2. Internship Activities and Assigned Tasks};
    \item \textbf{Chapter 3. Results, Experience, and Evaluation}.
\end{itemize}

These chapter titles are recommendations only and may be adapted to the
actual internship content.

The report shall reflect work actually performed, observed, analyzed, or
learned during the internship. Unrelated theoretical material should not
be included merely to increase report length.

\vspace{0.3cm}

% ================================================================
% X. TABLE OF CONTENTS
% ================================================================

\noindent
\textbf{X.\quad TABLE OF CONTENTS}

\vspace{0.15cm}

The Table of Contents shall be generated automatically using:

\begin{verbatim}
\tableofcontents
\end{verbatim}

Students shall not create or manually edit page numbers in the Table of
Contents.

The document shall be compiled sufficiently to ensure that all page
numbers and entries are correctly updated.

\vspace{0.3cm}

% ================================================================
% XI. LIST OF SYMBOLS AND ABBREVIATIONS
% ================================================================

\noindent
\textbf{XI.\quad LIST OF SYMBOLS AND ABBREVIATIONS}

\vspace{0.15cm}

Symbols and abbreviations that are used repeatedly in the report shall be
listed in the \textbf{List of Symbols and Abbreviations}.

The list may contain the following columns:

\begin{center}
\begin{tabular}{llll}
    No. & Symbol / Abbreviation & Unit & Definition
\end{tabular}
\end{center}

Only symbols and abbreviations actually used in the final report shall be
retained. Example entries contained in the template shall be removed if
they are not relevant to the student's report.

Each abbreviation should normally be defined when it first appears in
the main text.

\vspace{0.3cm}

% ================================================================
% XII. FIGURES, DIAGRAMS, AND PHOTOGRAPHS
% ================================================================

\noindent
\textbf{XII.\quad FIGURES, DIAGRAMS, AND PHOTOGRAPHS}

\vspace{0.15cm}

Figures shall be clear, legible, and directly relevant to the report
content.

Figures shall be inserted using the standard \verb|figure| environment.

Example:

\begin{verbatim}
\begin{figure}[H]
    \centering
    \includegraphics[width=0.8\textwidth]
                    {figures/company_structure.pdf}
    \caption{Organizational structure of the company}
    \label{fig:company-structure}
\end{figure}
\end{verbatim}

Figure captions shall be placed below the figure.

Each referenced figure shall have a unique \verb|\label| and shall be
referenced in the text using, for example:

\begin{verbatim}
Figure~\ref{fig:company-structure}
\end{verbatim}

Students shall not manually type figure numbers.

Figures reproduced or adapted from external sources shall include an
appropriate citation.

Photographs taken during the internship shall be relevant and sufficiently
clear. Students shall not include confidential, restricted, or proprietary
information without permission from the internship organization.

Vector graphics or high-resolution images are recommended whenever
possible.

\vspace{0.3cm}

% ================================================================
% XIII. TABLES
% ================================================================

\noindent
\textbf{XIII.\quad TABLES}

\vspace{0.15cm}

Tables shall be created using suitable \LaTeX{} environments such as
\verb|table|, \verb|tabular|, \verb|tabularx|, or \verb|longtable|.

Table captions shall be placed above the table.

Example:

\begin{verbatim}
\begin{table}[H]
    \centering
    \caption{Internship schedule}
    \label{tab:internship-schedule}

    \begin{tabular}{lll}
        \toprule
        Week & Task & Result \\
        \midrule
        1 & Orientation & Completed \\
        2 & System study & Completed \\
        \bottomrule
    \end{tabular}
\end{table}
\end{verbatim}

Tables shall be referenced using:

\begin{verbatim}
Table~\ref{tab:internship-schedule}
\end{verbatim}

Students shall not manually type table numbers.

Tables should be created directly in \LaTeX{} rather than inserted as
screenshots whenever practical.

\vspace{0.3cm}

% ================================================================
% XIV. EQUATIONS AND TECHNICAL QUANTITIES
% ================================================================

\noindent
\textbf{XIV.\quad EQUATIONS AND TECHNICAL QUANTITIES}

\vspace{0.15cm}

Mathematical expressions shall be typeset directly using \LaTeX{}.
Images or screenshots of equations shall not be used.

Short expressions may be written in inline mathematical mode, for example:

\begin{verbatim}
$R = 0.7\,\Omega$
\end{verbatim}

Important equations shall use a numbered equation environment.

Example:

\begin{verbatim}
\begin{equation}
    P = UI
    \label{eq:power}
\end{equation}
\end{verbatim}

Equations shall be referenced using:

\begin{verbatim}
Equation~\eqref{eq:power}
\end{verbatim}

Equation numbers shall not be manually entered.

SI units shall be used wherever applicable. The \verb|siunitx| package
should be used where appropriate.

Examples:

\begin{verbatim}
\SI{220}{\volt}
\SI{50}{\hertz}
\SI{2.5}{\ampere}
\end{verbatim}

Notation and units shall be consistent throughout the report.

\vspace{0.3cm}

% ================================================================
% XV. SOURCE CODE
% ================================================================

\noindent
\textbf{XV.\quad SOURCE CODE}

\vspace{0.15cm}

If the internship involves programming, PLCs, microcontrollers,
simulation, or software development, relevant source code may be included
in the report.

Source code shall be presented using the \verb|listings| environment or
another environment defined by the official template.

Example:

\begin{verbatim}
\begin{lstlisting}[language=st]
IF Start = TRUE THEN
    Motor := TRUE;
ELSE
    Motor := FALSE;
END_IF;
\end{lstlisting}
\end{verbatim}

Only code necessary for explaining the work should appear in the main
report. Long source-code listings should be moved to an appendix or
submitted as supplementary files if required.

Source code should not normally be inserted as screenshots.

\vspace{0.3cm}

% ================================================================
% XVI. ABSTRACT
% ================================================================

\noindent
\textbf{XVI.\quad ABSTRACT}

\vspace{0.15cm}

The Abstract shall briefly summarize:

\begin{itemize}[leftmargin=0.8cm, itemsep=2pt]
    \item the internship organization;
    \item the objectives of the internship;
    \item the main tasks performed;
    \item the main results, knowledge, and skills obtained.
\end{itemize}

The Abstract should be concise and should not reproduce the full content
of the report chapter by chapter.

\vspace{0.3cm}

% ================================================================
% XVII. REFERENCES AND CITATIONS
% ================================================================

\noindent
\textbf{XVII.\quad REFERENCES AND CITATIONS}

\vspace{0.15cm}

All books, journal papers, conference papers, standards, websites,
datasheets, manuals, technical documents, figures, tables, or other
external sources used in preparing the report shall be properly cited.

References shall be managed through the bibliography system provided with
the official template.

Bibliographic entries shall normally be stored in:

\begin{verbatim}
References.bib
\end{verbatim}

References shall be cited in the report using:

\begin{verbatim}
\cite{reference-key}
\end{verbatim}

Students shall not manually type reference numbers such as ``[1]'',
``[2]'', etc.

Titles of foreign-language publications shall remain in their original
language in the bibliography.

All borrowed ideas, data, figures, tables, equations, methods, or results
shall be properly acknowledged.

\vspace{0.3cm}

% ================================================================
% XVIII. CROSS-REFERENCING
% ================================================================

\noindent
\textbf{XVIII.\quad CROSS-REFERENCING}

\vspace{0.15cm}

Numbered elements shall be cross-referenced using the standard \LaTeX{}
mechanism.

Use the following commands where appropriate:

\begin{verbatim}
\label{...}
\ref{...}
\eqref{...}
\end{verbatim}

For example:

\begin{verbatim}
As shown in Figure~\ref{fig:system}, ...
\end{verbatim}

Students shall not manually type numbers such as ``Figure 2.3'' when
those numbers can be generated automatically.

\vspace{0.3cm}

% ================================================================
% XIX. PAGE NUMBERING
% ================================================================

\noindent
\textbf{XIX.\quad PAGE NUMBERING}

\vspace{0.15cm}

Cover pages and official evaluation forms at the beginning of the report
shall not display page numbers.

The main body of the report shall use Arabic page numbers beginning from
Chapter 1.

Page numbering shall be generated automatically by the official template.
Students shall not manually enter page numbers.

\vspace{0.3cm}

% ================================================================
% XX. PROJECT FILE ORGANIZATION
% ================================================================

\noindent
\textbf{XX.\quad PROJECT FILE ORGANIZATION}

\vspace{0.15cm}

Students should maintain a clear and organized \LaTeX{} project
structure. A recommended structure is:

\begin{verbatim}
InternshipReport/
|
|-- main.tex
|-- References.bib
|
|-- figures/
|   |-- company.pdf
|   |-- system.png
|   `-- ...
|
`-- chapters/
    |-- chapter1.tex
    |-- chapter2.tex
    `-- chapter3.tex
\end{verbatim}

Figures should normally be stored in the \verb|figures| directory
specified by the official template.

File and directory names should preferably contain Latin letters,
numbers, hyphens, or underscores.

Recommended examples include:

\begin{verbatim}
company_structure.pdf
motor_control.png
chapter_01.tex
\end{verbatim}

Spaces and unnecessary special characters in file names should be
avoided where possible.

\vspace{0.3cm}

% ================================================================
% XXI. DOCUMENT COMPILATION AND FINAL CHECK
% ================================================================

\noindent
\textbf{XXI.\quad DOCUMENT COMPILATION AND FINAL CHECK}

\vspace{0.15cm}

Before submission, the complete \LaTeX{} project shall compile
successfully and produce the final PDF report.

The final document shall not contain unresolved references such as:

\begin{verbatim}
??
[?]
\end{verbatim}

Before generating the final PDF, students shall check:

\begin{itemize}[leftmargin=0.8cm, itemsep=2pt]
    \item the Table of Contents;
    \item page numbering;
    \item chapter, section, and subsection numbering;
    \item figure numbering;
    \item table numbering;
    \item equation numbering;
    \item cross-references;
    \item citations;
    \item bibliography;
    \item font display;
    \item mathematical symbols;
    \item image quality;
    \item completeness of all required forms.
\end{itemize}

The document may be compiled using TeX Live, MiKTeX, or another
compatible \LaTeX{} environment. Regardless of the compilation
environment, the resulting PDF shall conform to the official Faculty
template.

\vspace{0.3cm}

% ================================================================
% XXII. CONTENT THAT STUDENTS MAY EDIT
% ================================================================

\noindent
\textbf{XXII.\quad CONTENT THAT STUDENTS MAY EDIT}

\vspace{0.15cm}

Students may edit or replace:

\begin{itemize}[leftmargin=0.8cm, itemsep=2pt]
    \item student name;
    \item class;
    \item internship organization;
    \item supervisor name;
    \item Abstract;
    \item List of Symbols and Abbreviations;
    \item chapter titles and chapter content;
    \item figures and photographs;
    \item tables;
    \item equations;
    \item citations and references.
\end{itemize}

Students shall not modify the main document formatting settings in the
preamble, including the document class, page geometry, line spacing,
heading formatting, table-of-contents formatting, font configuration,
page layout, or cover-page macros unless specifically instructed.

\vspace{0.3cm}

% ================================================================
% XXIII. FINAL SUBMISSION REQUIREMENTS
% ================================================================

\noindent
\textbf{XXIII.\quad FINAL SUBMISSION REQUIREMENTS}

\vspace{0.15cm}

Before submission, the final report shall:

\begin{itemize}[leftmargin=0.8cm, itemsep=2pt]
    \item compile without errors that affect the report content or presentation;
    \item contain all required forms and report sections;
    \item contain correct student, class, internship organization, and supervisor information;
    \item contain an updated Table of Contents, List of Figures, and List of Tables;
    \item contain no unused example content from the template;
    \item contain no placeholder text such as \textit{Title of Chapter 1},
    \textit{Section 1.1}, or other sample entries;
    \item contain all required citations and references;
    \item contain no unresolved cross-references.
\end{itemize}

If the Faculty requires submission of the \LaTeX{} source, the complete
project shall include all \verb|.tex| files, the bibliography database,
figures, and any other files required to reproduce the submitted PDF.

\vspace{0.4cm}

\noindent
\textbf{Important:}
Before submitting the final internship report, students shall remove or
comment out the following command in the main document:

\begin{verbatim}
\input{Regulations.tex}
\end{verbatim}

For example:

\begin{verbatim}
% \input{Regulations.tex}
\end{verbatim}

This removes the guidance pages from the final submitted report while
preserving the regulations file in the \LaTeX{} project for reference.

} % End 13pt content

\clearpage

%
% IMPORTANT:
% This section is provided for guidance only. Students shall remove
% or comment out the 
% =================================================================
% REGULATIONS FOR PREPARING THE GRADUATION INTERNSHIP REPORT
% USING LATEX
% =================================================================
%
% This file is intended to be included from the main document with:
%
%     \input{Regulations.tex}
%
% IMPORTANT:
% This section is provided for guidance only. Students shall remove
% or comment out the \input{Regulations.tex} command before submitting
% the final internship report.
%
% =================================================================

\clearpage
\thispagestyle{empty}

\begin{center}
    {\fontsize{16pt}{20pt}\selectfont\bfseries
    SPECIFIC REGULATIONS FOR PREPARING\\
    GRADUATION INTERNSHIP REPORTS USING \LaTeX}
\end{center}

\vspace{0.5cm}

{\fontsize{13pt}{17pt}\selectfont
\setstretch{1.3}
\justifying

\noindent
\textbf{Note:}
This section is provided for guidance only and is included in the main
document using the command \verb|\input{Regulations.tex}|.
Students shall remove or comment out this command before submitting the
final internship report.

\vspace{0.4cm}

% ================================================================
% I. GENERAL REQUIREMENTS
% ================================================================

\noindent
\textbf{I.\quad GENERAL REQUIREMENTS}

\vspace{0.15cm}

The Graduation Internship Report shall be prepared using the official
\LaTeX{} template provided by the Faculty. Students shall preserve the
predefined page layout, typography, line spacing, heading styles,
numbering system, figure and table formatting, and bibliography settings
unless otherwise instructed by the Faculty or the supervisor.

The report shall be written in English.

Students are responsible for the academic and technical content of the
report. The official \LaTeX{} template is responsible for standardized
document formatting.

\vspace{0.3cm}

% ================================================================
% II. REPORT STRUCTURE
% ================================================================

\noindent
\textbf{II.\quad REPORT STRUCTURE}

\vspace{0.15cm}

The report shall normally be arranged in the following order:

\begin{enumerate}[label=\arabic*., leftmargin=1.2cm, itemsep=2pt]
    \item Outer Cover;
    \item Comments of the Internship Host Organization;
    \item Supervisor's Evaluation and Graduation Internship Defense Results;
    \item Table of Contents;
    \item List of Symbols and Abbreviations, if applicable;
    \item List of Figures;
    \item List of Tables;
    \item Abstract;
    \item Main Chapters of the Internship Report;
    \item References;
    \item Appendices, if applicable.
\end{enumerate}

The preliminary pages shall not be numbered as chapters. The main report
content shall be organized into consecutively numbered chapters beginning
with Chapter 1.

\vspace{0.3cm}

% ================================================================
% III. PAGE SIZE AND MARGINS
% ================================================================

\noindent
\textbf{III.\quad PAGE SIZE AND MARGINS}

\vspace{0.15cm}

The report shall be formatted on A4 paper with the following margins:

\begin{center}
\begin{tabular}{ll}
    Left margin:   & 3.5 cm \\
    Right margin:  & 2.0 cm \\
    Top margin:    & 2.5 cm \\
    Bottom margin: & 2.5 cm
\end{tabular}
\end{center}

These settings are predefined in the official template and shall not be
modified by students.

\vspace{0.3cm}

% ================================================================
% IV. FONT, FONT SIZE, AND LINE SPACING
% ================================================================

\noindent
\textbf{IV.\quad FONT, FONT SIZE, AND LINE SPACING}

\vspace{0.15cm}

The main text shall use a 13 pt font size with 1.3-line spacing.
A Times-family typeface shall be used in accordance with the official
template.

The text shall be fully justified.

Students shall not use repeated spaces, repeated manual line breaks, or
multiple blank lines to position text. Paragraph and heading spacing shall
be controlled by the template.

\vspace{0.3cm}

% ================================================================
% V. OUTER COVER
% ================================================================

\noindent
\textbf{V.\quad OUTER COVER}

\vspace{0.15cm}

The outer cover shall follow the official format provided in the template
and shall contain the following information:

\begin{itemize}[leftmargin=0.8cm, itemsep=2pt]
    \item Thai Nguyen University of Technology;
    \item Faculty of International Training;
    \item Engineering and Technology Department;
    \item University logo and Faculty logo;
    \item the title \textbf{INTERNSHIP REPORT};
    \item Internship Organization;
    \item Author;
    \item Class;
    \item Supervisor;
    \item place and year of completion.
\end{itemize}

Students shall only replace the required information fields and shall not
modify the cover-page layout unless instructed otherwise.

\vspace{0.3cm}

% ================================================================
% VI. COMMENTS OF THE INTERNSHIP HOST ORGANIZATION
% ================================================================

\noindent
\textbf{VI.\quad COMMENTS OF THE INTERNSHIP HOST ORGANIZATION}

\vspace{0.15cm}

The page entitled \textbf{COMMENTS OF THE INTERNSHIP HOST ORGANIZATION}
is reserved for the internship host organization.

It includes:

\begin{itemize}[leftmargin=0.8cm, itemsep=2pt]
    \item the student's full name;
    \item class;
    \item comments from the host organization;
    \item score;
    \item date;
    \item signature and official seal.
\end{itemize}

Students shall not alter or complete official evaluation fields unless
specifically authorized.

\vspace{0.3cm}

% ================================================================
% VII. SUPERVISOR'S EVALUATION AND DEFENSE RESULTS
% ================================================================

\noindent
\textbf{VII.\quad SUPERVISOR'S EVALUATION AND DEFENSE RESULTS}

\vspace{0.15cm}

The section entitled
\textbf{SUPERVISOR'S EVALUATION AND GRADUATION INTERNSHIP DEFENSE RESULTS}
is reserved for the supervisor and defense committee.

The evaluation may include:

\begin{itemize}[leftmargin=0.8cm, itemsep=2pt]
    \item Internship Progress and Student's Attitude;
    \item Report Content;
    \item Report Presentation;
    \item Other Comments;
    \item Supervisor's Opinion;
    \item Defense Result;
    \item Score.
\end{itemize}

Students shall not enter evaluation scores or defense results before they
are officially determined.

\vspace{0.3cm}

% ================================================================
% VIII. CHAPTERS, SECTIONS, AND SUBSECTIONS
% ================================================================

\noindent
\textbf{VIII.\quad CHAPTERS, SECTIONS, AND SUBSECTIONS}

\vspace{0.15cm}

The report structure shall be generated using the standard \LaTeX{}
commands:

\begin{verbatim}
\chapter{...}
\section{...}
\subsection{...}
\end{verbatim}

Chapter, section, and subsection numbers shall be generated automatically
by \LaTeX{}.

Students shall not manually type chapter or section numbers and shall not
reproduce heading styles using commands such as \verb|\textbf|,
\verb|\fontsize|, or \verb|\centering| instead of the appropriate
structural command.

Excessive subdivision should be avoided. In most cases, the hierarchy
Chapter--Section--Subsection is sufficient.

\vspace{0.3cm}

% ================================================================
% IX. ORGANIZATION OF THE MAIN CONTENT
% ================================================================

\noindent
\textbf{IX.\quad ORGANIZATION OF THE MAIN CONTENT}

\vspace{0.15cm}

The specific content of the report shall depend on the internship
organization and the assigned tasks. A typical structure may include:

\begin{itemize}[leftmargin=0.8cm, itemsep=2pt]
    \item \textbf{Chapter 1. Overview of the Internship Organization};
    \item \textbf{Chapter 2. Internship Activities and Assigned Tasks};
    \item \textbf{Chapter 3. Results, Experience, and Evaluation}.
\end{itemize}

These chapter titles are recommendations only and may be adapted to the
actual internship content.

The report shall reflect work actually performed, observed, analyzed, or
learned during the internship. Unrelated theoretical material should not
be included merely to increase report length.

\vspace{0.3cm}

% ================================================================
% X. TABLE OF CONTENTS
% ================================================================

\noindent
\textbf{X.\quad TABLE OF CONTENTS}

\vspace{0.15cm}

The Table of Contents shall be generated automatically using:

\begin{verbatim}
\tableofcontents
\end{verbatim}

Students shall not create or manually edit page numbers in the Table of
Contents.

The document shall be compiled sufficiently to ensure that all page
numbers and entries are correctly updated.

\vspace{0.3cm}

% ================================================================
% XI. LIST OF SYMBOLS AND ABBREVIATIONS
% ================================================================

\noindent
\textbf{XI.\quad LIST OF SYMBOLS AND ABBREVIATIONS}

\vspace{0.15cm}

Symbols and abbreviations that are used repeatedly in the report shall be
listed in the \textbf{List of Symbols and Abbreviations}.

The list may contain the following columns:

\begin{center}
\begin{tabular}{llll}
    No. & Symbol / Abbreviation & Unit & Definition
\end{tabular}
\end{center}

Only symbols and abbreviations actually used in the final report shall be
retained. Example entries contained in the template shall be removed if
they are not relevant to the student's report.

Each abbreviation should normally be defined when it first appears in
the main text.

\vspace{0.3cm}

% ================================================================
% XII. FIGURES, DIAGRAMS, AND PHOTOGRAPHS
% ================================================================

\noindent
\textbf{XII.\quad FIGURES, DIAGRAMS, AND PHOTOGRAPHS}

\vspace{0.15cm}

Figures shall be clear, legible, and directly relevant to the report
content.

Figures shall be inserted using the standard \verb|figure| environment.

Example:

\begin{verbatim}
\begin{figure}[H]
    \centering
    \includegraphics[width=0.8\textwidth]
                    {figures/company_structure.pdf}
    \caption{Organizational structure of the company}
    \label{fig:company-structure}
\end{figure}
\end{verbatim}

Figure captions shall be placed below the figure.

Each referenced figure shall have a unique \verb|\label| and shall be
referenced in the text using, for example:

\begin{verbatim}
Figure~\ref{fig:company-structure}
\end{verbatim}

Students shall not manually type figure numbers.

Figures reproduced or adapted from external sources shall include an
appropriate citation.

Photographs taken during the internship shall be relevant and sufficiently
clear. Students shall not include confidential, restricted, or proprietary
information without permission from the internship organization.

Vector graphics or high-resolution images are recommended whenever
possible.

\vspace{0.3cm}

% ================================================================
% XIII. TABLES
% ================================================================

\noindent
\textbf{XIII.\quad TABLES}

\vspace{0.15cm}

Tables shall be created using suitable \LaTeX{} environments such as
\verb|table|, \verb|tabular|, \verb|tabularx|, or \verb|longtable|.

Table captions shall be placed above the table.

Example:

\begin{verbatim}
\begin{table}[H]
    \centering
    \caption{Internship schedule}
    \label{tab:internship-schedule}

    \begin{tabular}{lll}
        \toprule
        Week & Task & Result \\
        \midrule
        1 & Orientation & Completed \\
        2 & System study & Completed \\
        \bottomrule
    \end{tabular}
\end{table}
\end{verbatim}

Tables shall be referenced using:

\begin{verbatim}
Table~\ref{tab:internship-schedule}
\end{verbatim}

Students shall not manually type table numbers.

Tables should be created directly in \LaTeX{} rather than inserted as
screenshots whenever practical.

\vspace{0.3cm}

% ================================================================
% XIV. EQUATIONS AND TECHNICAL QUANTITIES
% ================================================================

\noindent
\textbf{XIV.\quad EQUATIONS AND TECHNICAL QUANTITIES}

\vspace{0.15cm}

Mathematical expressions shall be typeset directly using \LaTeX{}.
Images or screenshots of equations shall not be used.

Short expressions may be written in inline mathematical mode, for example:

\begin{verbatim}
$R = 0.7\,\Omega$
\end{verbatim}

Important equations shall use a numbered equation environment.

Example:

\begin{verbatim}
\begin{equation}
    P = UI
    \label{eq:power}
\end{equation}
\end{verbatim}

Equations shall be referenced using:

\begin{verbatim}
Equation~\eqref{eq:power}
\end{verbatim}

Equation numbers shall not be manually entered.

SI units shall be used wherever applicable. The \verb|siunitx| package
should be used where appropriate.

Examples:

\begin{verbatim}
\SI{220}{\volt}
\SI{50}{\hertz}
\SI{2.5}{\ampere}
\end{verbatim}

Notation and units shall be consistent throughout the report.

\vspace{0.3cm}

% ================================================================
% XV. SOURCE CODE
% ================================================================

\noindent
\textbf{XV.\quad SOURCE CODE}

\vspace{0.15cm}

If the internship involves programming, PLCs, microcontrollers,
simulation, or software development, relevant source code may be included
in the report.

Source code shall be presented using the \verb|listings| environment or
another environment defined by the official template.

Example:

\begin{verbatim}
\begin{lstlisting}[language=st]
IF Start = TRUE THEN
    Motor := TRUE;
ELSE
    Motor := FALSE;
END_IF;
\end{lstlisting}
\end{verbatim}

Only code necessary for explaining the work should appear in the main
report. Long source-code listings should be moved to an appendix or
submitted as supplementary files if required.

Source code should not normally be inserted as screenshots.

\vspace{0.3cm}

% ================================================================
% XVI. ABSTRACT
% ================================================================

\noindent
\textbf{XVI.\quad ABSTRACT}

\vspace{0.15cm}

The Abstract shall briefly summarize:

\begin{itemize}[leftmargin=0.8cm, itemsep=2pt]
    \item the internship organization;
    \item the objectives of the internship;
    \item the main tasks performed;
    \item the main results, knowledge, and skills obtained.
\end{itemize}

The Abstract should be concise and should not reproduce the full content
of the report chapter by chapter.

\vspace{0.3cm}

% ================================================================
% XVII. REFERENCES AND CITATIONS
% ================================================================

\noindent
\textbf{XVII.\quad REFERENCES AND CITATIONS}

\vspace{0.15cm}

All books, journal papers, conference papers, standards, websites,
datasheets, manuals, technical documents, figures, tables, or other
external sources used in preparing the report shall be properly cited.

References shall be managed through the bibliography system provided with
the official template.

Bibliographic entries shall normally be stored in:

\begin{verbatim}
References.bib
\end{verbatim}

References shall be cited in the report using:

\begin{verbatim}
\cite{reference-key}
\end{verbatim}

Students shall not manually type reference numbers such as ``[1]'',
``[2]'', etc.

Titles of foreign-language publications shall remain in their original
language in the bibliography.

All borrowed ideas, data, figures, tables, equations, methods, or results
shall be properly acknowledged.

\vspace{0.3cm}

% ================================================================
% XVIII. CROSS-REFERENCING
% ================================================================

\noindent
\textbf{XVIII.\quad CROSS-REFERENCING}

\vspace{0.15cm}

Numbered elements shall be cross-referenced using the standard \LaTeX{}
mechanism.

Use the following commands where appropriate:

\begin{verbatim}
\label{...}
\ref{...}
\eqref{...}
\end{verbatim}

For example:

\begin{verbatim}
As shown in Figure~\ref{fig:system}, ...
\end{verbatim}

Students shall not manually type numbers such as ``Figure 2.3'' when
those numbers can be generated automatically.

\vspace{0.3cm}

% ================================================================
% XIX. PAGE NUMBERING
% ================================================================

\noindent
\textbf{XIX.\quad PAGE NUMBERING}

\vspace{0.15cm}

Cover pages and official evaluation forms at the beginning of the report
shall not display page numbers.

The main body of the report shall use Arabic page numbers beginning from
Chapter 1.

Page numbering shall be generated automatically by the official template.
Students shall not manually enter page numbers.

\vspace{0.3cm}

% ================================================================
% XX. PROJECT FILE ORGANIZATION
% ================================================================

\noindent
\textbf{XX.\quad PROJECT FILE ORGANIZATION}

\vspace{0.15cm}

Students should maintain a clear and organized \LaTeX{} project
structure. A recommended structure is:

\begin{verbatim}
InternshipReport/
|
|-- main.tex
|-- References.bib
|
|-- figures/
|   |-- company.pdf
|   |-- system.png
|   `-- ...
|
`-- chapters/
    |-- chapter1.tex
    |-- chapter2.tex
    `-- chapter3.tex
\end{verbatim}

Figures should normally be stored in the \verb|figures| directory
specified by the official template.

File and directory names should preferably contain Latin letters,
numbers, hyphens, or underscores.

Recommended examples include:

\begin{verbatim}
company_structure.pdf
motor_control.png
chapter_01.tex
\end{verbatim}

Spaces and unnecessary special characters in file names should be
avoided where possible.

\vspace{0.3cm}

% ================================================================
% XXI. DOCUMENT COMPILATION AND FINAL CHECK
% ================================================================

\noindent
\textbf{XXI.\quad DOCUMENT COMPILATION AND FINAL CHECK}

\vspace{0.15cm}

Before submission, the complete \LaTeX{} project shall compile
successfully and produce the final PDF report.

The final document shall not contain unresolved references such as:

\begin{verbatim}
??
[?]
\end{verbatim}

Before generating the final PDF, students shall check:

\begin{itemize}[leftmargin=0.8cm, itemsep=2pt]
    \item the Table of Contents;
    \item page numbering;
    \item chapter, section, and subsection numbering;
    \item figure numbering;
    \item table numbering;
    \item equation numbering;
    \item cross-references;
    \item citations;
    \item bibliography;
    \item font display;
    \item mathematical symbols;
    \item image quality;
    \item completeness of all required forms.
\end{itemize}

The document may be compiled using TeX Live, MiKTeX, or another
compatible \LaTeX{} environment. Regardless of the compilation
environment, the resulting PDF shall conform to the official Faculty
template.

\vspace{0.3cm}

% ================================================================
% XXII. CONTENT THAT STUDENTS MAY EDIT
% ================================================================

\noindent
\textbf{XXII.\quad CONTENT THAT STUDENTS MAY EDIT}

\vspace{0.15cm}

Students may edit or replace:

\begin{itemize}[leftmargin=0.8cm, itemsep=2pt]
    \item student name;
    \item class;
    \item internship organization;
    \item supervisor name;
    \item Abstract;
    \item List of Symbols and Abbreviations;
    \item chapter titles and chapter content;
    \item figures and photographs;
    \item tables;
    \item equations;
    \item citations and references.
\end{itemize}

Students shall not modify the main document formatting settings in the
preamble, including the document class, page geometry, line spacing,
heading formatting, table-of-contents formatting, font configuration,
page layout, or cover-page macros unless specifically instructed.

\vspace{0.3cm}

% ================================================================
% XXIII. FINAL SUBMISSION REQUIREMENTS
% ================================================================

\noindent
\textbf{XXIII.\quad FINAL SUBMISSION REQUIREMENTS}

\vspace{0.15cm}

Before submission, the final report shall:

\begin{itemize}[leftmargin=0.8cm, itemsep=2pt]
    \item compile without errors that affect the report content or presentation;
    \item contain all required forms and report sections;
    \item contain correct student, class, internship organization, and supervisor information;
    \item contain an updated Table of Contents, List of Figures, and List of Tables;
    \item contain no unused example content from the template;
    \item contain no placeholder text such as \textit{Title of Chapter 1},
    \textit{Section 1.1}, or other sample entries;
    \item contain all required citations and references;
    \item contain no unresolved cross-references.
\end{itemize}

If the Faculty requires submission of the \LaTeX{} source, the complete
project shall include all \verb|.tex| files, the bibliography database,
figures, and any other files required to reproduce the submitted PDF.

\vspace{0.4cm}

\noindent
\textbf{Important:}
Before submitting the final internship report, students shall remove or
comment out the following command in the main document:

\begin{verbatim}
\input{Regulations.tex}
\end{verbatim}

For example:

\begin{verbatim}
% \input{Regulations.tex}
\end{verbatim}

This removes the guidance pages from the final submitted report while
preserving the regulations file in the \LaTeX{} project for reference.

} % End 13pt content

\clearpage
 command before submitting
% the final internship report.
%
% =================================================================

\clearpage
\thispagestyle{empty}

\begin{center}
    {\fontsize{16pt}{20pt}\selectfont\bfseries
    SPECIFIC REGULATIONS FOR PREPARING\\
    GRADUATION INTERNSHIP REPORTS USING \LaTeX}
\end{center}

\vspace{0.5cm}

{\fontsize{13pt}{17pt}\selectfont
\setstretch{1.3}
\justifying

\noindent
\textbf{Note:}
This section is provided for guidance only and is included in the main
document using the command \verb|
% =================================================================
% REGULATIONS FOR PREPARING THE GRADUATION INTERNSHIP REPORT
% USING LATEX
% =================================================================
%
% This file is intended to be included from the main document with:
%
%     \input{Regulations.tex}
%
% IMPORTANT:
% This section is provided for guidance only. Students shall remove
% or comment out the \input{Regulations.tex} command before submitting
% the final internship report.
%
% =================================================================

\clearpage
\thispagestyle{empty}

\begin{center}
    {\fontsize{16pt}{20pt}\selectfont\bfseries
    SPECIFIC REGULATIONS FOR PREPARING\\
    GRADUATION INTERNSHIP REPORTS USING \LaTeX}
\end{center}

\vspace{0.5cm}

{\fontsize{13pt}{17pt}\selectfont
\setstretch{1.3}
\justifying

\noindent
\textbf{Note:}
This section is provided for guidance only and is included in the main
document using the command \verb|\input{Regulations.tex}|.
Students shall remove or comment out this command before submitting the
final internship report.

\vspace{0.4cm}

% ================================================================
% I. GENERAL REQUIREMENTS
% ================================================================

\noindent
\textbf{I.\quad GENERAL REQUIREMENTS}

\vspace{0.15cm}

The Graduation Internship Report shall be prepared using the official
\LaTeX{} template provided by the Faculty. Students shall preserve the
predefined page layout, typography, line spacing, heading styles,
numbering system, figure and table formatting, and bibliography settings
unless otherwise instructed by the Faculty or the supervisor.

The report shall be written in English.

Students are responsible for the academic and technical content of the
report. The official \LaTeX{} template is responsible for standardized
document formatting.

\vspace{0.3cm}

% ================================================================
% II. REPORT STRUCTURE
% ================================================================

\noindent
\textbf{II.\quad REPORT STRUCTURE}

\vspace{0.15cm}

The report shall normally be arranged in the following order:

\begin{enumerate}[label=\arabic*., leftmargin=1.2cm, itemsep=2pt]
    \item Outer Cover;
    \item Comments of the Internship Host Organization;
    \item Supervisor's Evaluation and Graduation Internship Defense Results;
    \item Table of Contents;
    \item List of Symbols and Abbreviations, if applicable;
    \item List of Figures;
    \item List of Tables;
    \item Abstract;
    \item Main Chapters of the Internship Report;
    \item References;
    \item Appendices, if applicable.
\end{enumerate}

The preliminary pages shall not be numbered as chapters. The main report
content shall be organized into consecutively numbered chapters beginning
with Chapter 1.

\vspace{0.3cm}

% ================================================================
% III. PAGE SIZE AND MARGINS
% ================================================================

\noindent
\textbf{III.\quad PAGE SIZE AND MARGINS}

\vspace{0.15cm}

The report shall be formatted on A4 paper with the following margins:

\begin{center}
\begin{tabular}{ll}
    Left margin:   & 3.5 cm \\
    Right margin:  & 2.0 cm \\
    Top margin:    & 2.5 cm \\
    Bottom margin: & 2.5 cm
\end{tabular}
\end{center}

These settings are predefined in the official template and shall not be
modified by students.

\vspace{0.3cm}

% ================================================================
% IV. FONT, FONT SIZE, AND LINE SPACING
% ================================================================

\noindent
\textbf{IV.\quad FONT, FONT SIZE, AND LINE SPACING}

\vspace{0.15cm}

The main text shall use a 13 pt font size with 1.3-line spacing.
A Times-family typeface shall be used in accordance with the official
template.

The text shall be fully justified.

Students shall not use repeated spaces, repeated manual line breaks, or
multiple blank lines to position text. Paragraph and heading spacing shall
be controlled by the template.

\vspace{0.3cm}

% ================================================================
% V. OUTER COVER
% ================================================================

\noindent
\textbf{V.\quad OUTER COVER}

\vspace{0.15cm}

The outer cover shall follow the official format provided in the template
and shall contain the following information:

\begin{itemize}[leftmargin=0.8cm, itemsep=2pt]
    \item Thai Nguyen University of Technology;
    \item Faculty of International Training;
    \item Engineering and Technology Department;
    \item University logo and Faculty logo;
    \item the title \textbf{INTERNSHIP REPORT};
    \item Internship Organization;
    \item Author;
    \item Class;
    \item Supervisor;
    \item place and year of completion.
\end{itemize}

Students shall only replace the required information fields and shall not
modify the cover-page layout unless instructed otherwise.

\vspace{0.3cm}

% ================================================================
% VI. COMMENTS OF THE INTERNSHIP HOST ORGANIZATION
% ================================================================

\noindent
\textbf{VI.\quad COMMENTS OF THE INTERNSHIP HOST ORGANIZATION}

\vspace{0.15cm}

The page entitled \textbf{COMMENTS OF THE INTERNSHIP HOST ORGANIZATION}
is reserved for the internship host organization.

It includes:

\begin{itemize}[leftmargin=0.8cm, itemsep=2pt]
    \item the student's full name;
    \item class;
    \item comments from the host organization;
    \item score;
    \item date;
    \item signature and official seal.
\end{itemize}

Students shall not alter or complete official evaluation fields unless
specifically authorized.

\vspace{0.3cm}

% ================================================================
% VII. SUPERVISOR'S EVALUATION AND DEFENSE RESULTS
% ================================================================

\noindent
\textbf{VII.\quad SUPERVISOR'S EVALUATION AND DEFENSE RESULTS}

\vspace{0.15cm}

The section entitled
\textbf{SUPERVISOR'S EVALUATION AND GRADUATION INTERNSHIP DEFENSE RESULTS}
is reserved for the supervisor and defense committee.

The evaluation may include:

\begin{itemize}[leftmargin=0.8cm, itemsep=2pt]
    \item Internship Progress and Student's Attitude;
    \item Report Content;
    \item Report Presentation;
    \item Other Comments;
    \item Supervisor's Opinion;
    \item Defense Result;
    \item Score.
\end{itemize}

Students shall not enter evaluation scores or defense results before they
are officially determined.

\vspace{0.3cm}

% ================================================================
% VIII. CHAPTERS, SECTIONS, AND SUBSECTIONS
% ================================================================

\noindent
\textbf{VIII.\quad CHAPTERS, SECTIONS, AND SUBSECTIONS}

\vspace{0.15cm}

The report structure shall be generated using the standard \LaTeX{}
commands:

\begin{verbatim}
\chapter{...}
\section{...}
\subsection{...}
\end{verbatim}

Chapter, section, and subsection numbers shall be generated automatically
by \LaTeX{}.

Students shall not manually type chapter or section numbers and shall not
reproduce heading styles using commands such as \verb|\textbf|,
\verb|\fontsize|, or \verb|\centering| instead of the appropriate
structural command.

Excessive subdivision should be avoided. In most cases, the hierarchy
Chapter--Section--Subsection is sufficient.

\vspace{0.3cm}

% ================================================================
% IX. ORGANIZATION OF THE MAIN CONTENT
% ================================================================

\noindent
\textbf{IX.\quad ORGANIZATION OF THE MAIN CONTENT}

\vspace{0.15cm}

The specific content of the report shall depend on the internship
organization and the assigned tasks. A typical structure may include:

\begin{itemize}[leftmargin=0.8cm, itemsep=2pt]
    \item \textbf{Chapter 1. Overview of the Internship Organization};
    \item \textbf{Chapter 2. Internship Activities and Assigned Tasks};
    \item \textbf{Chapter 3. Results, Experience, and Evaluation}.
\end{itemize}

These chapter titles are recommendations only and may be adapted to the
actual internship content.

The report shall reflect work actually performed, observed, analyzed, or
learned during the internship. Unrelated theoretical material should not
be included merely to increase report length.

\vspace{0.3cm}

% ================================================================
% X. TABLE OF CONTENTS
% ================================================================

\noindent
\textbf{X.\quad TABLE OF CONTENTS}

\vspace{0.15cm}

The Table of Contents shall be generated automatically using:

\begin{verbatim}
\tableofcontents
\end{verbatim}

Students shall not create or manually edit page numbers in the Table of
Contents.

The document shall be compiled sufficiently to ensure that all page
numbers and entries are correctly updated.

\vspace{0.3cm}

% ================================================================
% XI. LIST OF SYMBOLS AND ABBREVIATIONS
% ================================================================

\noindent
\textbf{XI.\quad LIST OF SYMBOLS AND ABBREVIATIONS}

\vspace{0.15cm}

Symbols and abbreviations that are used repeatedly in the report shall be
listed in the \textbf{List of Symbols and Abbreviations}.

The list may contain the following columns:

\begin{center}
\begin{tabular}{llll}
    No. & Symbol / Abbreviation & Unit & Definition
\end{tabular}
\end{center}

Only symbols and abbreviations actually used in the final report shall be
retained. Example entries contained in the template shall be removed if
they are not relevant to the student's report.

Each abbreviation should normally be defined when it first appears in
the main text.

\vspace{0.3cm}

% ================================================================
% XII. FIGURES, DIAGRAMS, AND PHOTOGRAPHS
% ================================================================

\noindent
\textbf{XII.\quad FIGURES, DIAGRAMS, AND PHOTOGRAPHS}

\vspace{0.15cm}

Figures shall be clear, legible, and directly relevant to the report
content.

Figures shall be inserted using the standard \verb|figure| environment.

Example:

\begin{verbatim}
\begin{figure}[H]
    \centering
    \includegraphics[width=0.8\textwidth]
                    {figures/company_structure.pdf}
    \caption{Organizational structure of the company}
    \label{fig:company-structure}
\end{figure}
\end{verbatim}

Figure captions shall be placed below the figure.

Each referenced figure shall have a unique \verb|\label| and shall be
referenced in the text using, for example:

\begin{verbatim}
Figure~\ref{fig:company-structure}
\end{verbatim}

Students shall not manually type figure numbers.

Figures reproduced or adapted from external sources shall include an
appropriate citation.

Photographs taken during the internship shall be relevant and sufficiently
clear. Students shall not include confidential, restricted, or proprietary
information without permission from the internship organization.

Vector graphics or high-resolution images are recommended whenever
possible.

\vspace{0.3cm}

% ================================================================
% XIII. TABLES
% ================================================================

\noindent
\textbf{XIII.\quad TABLES}

\vspace{0.15cm}

Tables shall be created using suitable \LaTeX{} environments such as
\verb|table|, \verb|tabular|, \verb|tabularx|, or \verb|longtable|.

Table captions shall be placed above the table.

Example:

\begin{verbatim}
\begin{table}[H]
    \centering
    \caption{Internship schedule}
    \label{tab:internship-schedule}

    \begin{tabular}{lll}
        \toprule
        Week & Task & Result \\
        \midrule
        1 & Orientation & Completed \\
        2 & System study & Completed \\
        \bottomrule
    \end{tabular}
\end{table}
\end{verbatim}

Tables shall be referenced using:

\begin{verbatim}
Table~\ref{tab:internship-schedule}
\end{verbatim}

Students shall not manually type table numbers.

Tables should be created directly in \LaTeX{} rather than inserted as
screenshots whenever practical.

\vspace{0.3cm}

% ================================================================
% XIV. EQUATIONS AND TECHNICAL QUANTITIES
% ================================================================

\noindent
\textbf{XIV.\quad EQUATIONS AND TECHNICAL QUANTITIES}

\vspace{0.15cm}

Mathematical expressions shall be typeset directly using \LaTeX{}.
Images or screenshots of equations shall not be used.

Short expressions may be written in inline mathematical mode, for example:

\begin{verbatim}
$R = 0.7\,\Omega$
\end{verbatim}

Important equations shall use a numbered equation environment.

Example:

\begin{verbatim}
\begin{equation}
    P = UI
    \label{eq:power}
\end{equation}
\end{verbatim}

Equations shall be referenced using:

\begin{verbatim}
Equation~\eqref{eq:power}
\end{verbatim}

Equation numbers shall not be manually entered.

SI units shall be used wherever applicable. The \verb|siunitx| package
should be used where appropriate.

Examples:

\begin{verbatim}
\SI{220}{\volt}
\SI{50}{\hertz}
\SI{2.5}{\ampere}
\end{verbatim}

Notation and units shall be consistent throughout the report.

\vspace{0.3cm}

% ================================================================
% XV. SOURCE CODE
% ================================================================

\noindent
\textbf{XV.\quad SOURCE CODE}

\vspace{0.15cm}

If the internship involves programming, PLCs, microcontrollers,
simulation, or software development, relevant source code may be included
in the report.

Source code shall be presented using the \verb|listings| environment or
another environment defined by the official template.

Example:

\begin{verbatim}
\begin{lstlisting}[language=st]
IF Start = TRUE THEN
    Motor := TRUE;
ELSE
    Motor := FALSE;
END_IF;
\end{lstlisting}
\end{verbatim}

Only code necessary for explaining the work should appear in the main
report. Long source-code listings should be moved to an appendix or
submitted as supplementary files if required.

Source code should not normally be inserted as screenshots.

\vspace{0.3cm}

% ================================================================
% XVI. ABSTRACT
% ================================================================

\noindent
\textbf{XVI.\quad ABSTRACT}

\vspace{0.15cm}

The Abstract shall briefly summarize:

\begin{itemize}[leftmargin=0.8cm, itemsep=2pt]
    \item the internship organization;
    \item the objectives of the internship;
    \item the main tasks performed;
    \item the main results, knowledge, and skills obtained.
\end{itemize}

The Abstract should be concise and should not reproduce the full content
of the report chapter by chapter.

\vspace{0.3cm}

% ================================================================
% XVII. REFERENCES AND CITATIONS
% ================================================================

\noindent
\textbf{XVII.\quad REFERENCES AND CITATIONS}

\vspace{0.15cm}

All books, journal papers, conference papers, standards, websites,
datasheets, manuals, technical documents, figures, tables, or other
external sources used in preparing the report shall be properly cited.

References shall be managed through the bibliography system provided with
the official template.

Bibliographic entries shall normally be stored in:

\begin{verbatim}
References.bib
\end{verbatim}

References shall be cited in the report using:

\begin{verbatim}
\cite{reference-key}
\end{verbatim}

Students shall not manually type reference numbers such as ``[1]'',
``[2]'', etc.

Titles of foreign-language publications shall remain in their original
language in the bibliography.

All borrowed ideas, data, figures, tables, equations, methods, or results
shall be properly acknowledged.

\vspace{0.3cm}

% ================================================================
% XVIII. CROSS-REFERENCING
% ================================================================

\noindent
\textbf{XVIII.\quad CROSS-REFERENCING}

\vspace{0.15cm}

Numbered elements shall be cross-referenced using the standard \LaTeX{}
mechanism.

Use the following commands where appropriate:

\begin{verbatim}
\label{...}
\ref{...}
\eqref{...}
\end{verbatim}

For example:

\begin{verbatim}
As shown in Figure~\ref{fig:system}, ...
\end{verbatim}

Students shall not manually type numbers such as ``Figure 2.3'' when
those numbers can be generated automatically.

\vspace{0.3cm}

% ================================================================
% XIX. PAGE NUMBERING
% ================================================================

\noindent
\textbf{XIX.\quad PAGE NUMBERING}

\vspace{0.15cm}

Cover pages and official evaluation forms at the beginning of the report
shall not display page numbers.

The main body of the report shall use Arabic page numbers beginning from
Chapter 1.

Page numbering shall be generated automatically by the official template.
Students shall not manually enter page numbers.

\vspace{0.3cm}

% ================================================================
% XX. PROJECT FILE ORGANIZATION
% ================================================================

\noindent
\textbf{XX.\quad PROJECT FILE ORGANIZATION}

\vspace{0.15cm}

Students should maintain a clear and organized \LaTeX{} project
structure. A recommended structure is:

\begin{verbatim}
InternshipReport/
|
|-- main.tex
|-- References.bib
|
|-- figures/
|   |-- company.pdf
|   |-- system.png
|   `-- ...
|
`-- chapters/
    |-- chapter1.tex
    |-- chapter2.tex
    `-- chapter3.tex
\end{verbatim}

Figures should normally be stored in the \verb|figures| directory
specified by the official template.

File and directory names should preferably contain Latin letters,
numbers, hyphens, or underscores.

Recommended examples include:

\begin{verbatim}
company_structure.pdf
motor_control.png
chapter_01.tex
\end{verbatim}

Spaces and unnecessary special characters in file names should be
avoided where possible.

\vspace{0.3cm}

% ================================================================
% XXI. DOCUMENT COMPILATION AND FINAL CHECK
% ================================================================

\noindent
\textbf{XXI.\quad DOCUMENT COMPILATION AND FINAL CHECK}

\vspace{0.15cm}

Before submission, the complete \LaTeX{} project shall compile
successfully and produce the final PDF report.

The final document shall not contain unresolved references such as:

\begin{verbatim}
??
[?]
\end{verbatim}

Before generating the final PDF, students shall check:

\begin{itemize}[leftmargin=0.8cm, itemsep=2pt]
    \item the Table of Contents;
    \item page numbering;
    \item chapter, section, and subsection numbering;
    \item figure numbering;
    \item table numbering;
    \item equation numbering;
    \item cross-references;
    \item citations;
    \item bibliography;
    \item font display;
    \item mathematical symbols;
    \item image quality;
    \item completeness of all required forms.
\end{itemize}

The document may be compiled using TeX Live, MiKTeX, or another
compatible \LaTeX{} environment. Regardless of the compilation
environment, the resulting PDF shall conform to the official Faculty
template.

\vspace{0.3cm}

% ================================================================
% XXII. CONTENT THAT STUDENTS MAY EDIT
% ================================================================

\noindent
\textbf{XXII.\quad CONTENT THAT STUDENTS MAY EDIT}

\vspace{0.15cm}

Students may edit or replace:

\begin{itemize}[leftmargin=0.8cm, itemsep=2pt]
    \item student name;
    \item class;
    \item internship organization;
    \item supervisor name;
    \item Abstract;
    \item List of Symbols and Abbreviations;
    \item chapter titles and chapter content;
    \item figures and photographs;
    \item tables;
    \item equations;
    \item citations and references.
\end{itemize}

Students shall not modify the main document formatting settings in the
preamble, including the document class, page geometry, line spacing,
heading formatting, table-of-contents formatting, font configuration,
page layout, or cover-page macros unless specifically instructed.

\vspace{0.3cm}

% ================================================================
% XXIII. FINAL SUBMISSION REQUIREMENTS
% ================================================================

\noindent
\textbf{XXIII.\quad FINAL SUBMISSION REQUIREMENTS}

\vspace{0.15cm}

Before submission, the final report shall:

\begin{itemize}[leftmargin=0.8cm, itemsep=2pt]
    \item compile without errors that affect the report content or presentation;
    \item contain all required forms and report sections;
    \item contain correct student, class, internship organization, and supervisor information;
    \item contain an updated Table of Contents, List of Figures, and List of Tables;
    \item contain no unused example content from the template;
    \item contain no placeholder text such as \textit{Title of Chapter 1},
    \textit{Section 1.1}, or other sample entries;
    \item contain all required citations and references;
    \item contain no unresolved cross-references.
\end{itemize}

If the Faculty requires submission of the \LaTeX{} source, the complete
project shall include all \verb|.tex| files, the bibliography database,
figures, and any other files required to reproduce the submitted PDF.

\vspace{0.4cm}

\noindent
\textbf{Important:}
Before submitting the final internship report, students shall remove or
comment out the following command in the main document:

\begin{verbatim}
\input{Regulations.tex}
\end{verbatim}

For example:

\begin{verbatim}
% \input{Regulations.tex}
\end{verbatim}

This removes the guidance pages from the final submitted report while
preserving the regulations file in the \LaTeX{} project for reference.

} % End 13pt content

\clearpage
|.
Students shall remove or comment out this command before submitting the
final internship report.

\vspace{0.4cm}

% ================================================================
% I. GENERAL REQUIREMENTS
% ================================================================

\noindent
\textbf{I.\quad GENERAL REQUIREMENTS}

\vspace{0.15cm}

The Graduation Internship Report shall be prepared using the official
\LaTeX{} template provided by the Faculty. Students shall preserve the
predefined page layout, typography, line spacing, heading styles,
numbering system, figure and table formatting, and bibliography settings
unless otherwise instructed by the Faculty or the supervisor.

The report shall be written in English.

Students are responsible for the academic and technical content of the
report. The official \LaTeX{} template is responsible for standardized
document formatting.

\vspace{0.3cm}

% ================================================================
% II. REPORT STRUCTURE
% ================================================================

\noindent
\textbf{II.\quad REPORT STRUCTURE}

\vspace{0.15cm}

The report shall normally be arranged in the following order:

\begin{enumerate}[label=\arabic*., leftmargin=1.2cm, itemsep=2pt]
    \item Outer Cover;
    \item Comments of the Internship Host Organization;
    \item Supervisor's Evaluation and Graduation Internship Defense Results;
    \item Table of Contents;
    \item List of Symbols and Abbreviations, if applicable;
    \item List of Figures;
    \item List of Tables;
    \item Abstract;
    \item Main Chapters of the Internship Report;
    \item References;
    \item Appendices, if applicable.
\end{enumerate}

The preliminary pages shall not be numbered as chapters. The main report
content shall be organized into consecutively numbered chapters beginning
with Chapter 1.

\vspace{0.3cm}

% ================================================================
% III. PAGE SIZE AND MARGINS
% ================================================================

\noindent
\textbf{III.\quad PAGE SIZE AND MARGINS}

\vspace{0.15cm}

The report shall be formatted on A4 paper with the following margins:

\begin{center}
\begin{tabular}{ll}
    Left margin:   & 3.5 cm \\
    Right margin:  & 2.0 cm \\
    Top margin:    & 2.5 cm \\
    Bottom margin: & 2.5 cm
\end{tabular}
\end{center}

These settings are predefined in the official template and shall not be
modified by students.

\vspace{0.3cm}

% ================================================================
% IV. FONT, FONT SIZE, AND LINE SPACING
% ================================================================

\noindent
\textbf{IV.\quad FONT, FONT SIZE, AND LINE SPACING}

\vspace{0.15cm}

The main text shall use a 13 pt font size with 1.3-line spacing.
A Times-family typeface shall be used in accordance with the official
template.

The text shall be fully justified.

Students shall not use repeated spaces, repeated manual line breaks, or
multiple blank lines to position text. Paragraph and heading spacing shall
be controlled by the template.

\vspace{0.3cm}

% ================================================================
% V. OUTER COVER
% ================================================================

\noindent
\textbf{V.\quad OUTER COVER}

\vspace{0.15cm}

The outer cover shall follow the official format provided in the template
and shall contain the following information:

\begin{itemize}[leftmargin=0.8cm, itemsep=2pt]
    \item Thai Nguyen University of Technology;
    \item Faculty of International Training;
    \item Engineering and Technology Department;
    \item University logo and Faculty logo;
    \item the title \textbf{INTERNSHIP REPORT};
    \item Internship Organization;
    \item Author;
    \item Class;
    \item Supervisor;
    \item place and year of completion.
\end{itemize}

Students shall only replace the required information fields and shall not
modify the cover-page layout unless instructed otherwise.

\vspace{0.3cm}

% ================================================================
% VI. COMMENTS OF THE INTERNSHIP HOST ORGANIZATION
% ================================================================

\noindent
\textbf{VI.\quad COMMENTS OF THE INTERNSHIP HOST ORGANIZATION}

\vspace{0.15cm}

The page entitled \textbf{COMMENTS OF THE INTERNSHIP HOST ORGANIZATION}
is reserved for the internship host organization.

It includes:

\begin{itemize}[leftmargin=0.8cm, itemsep=2pt]
    \item the student's full name;
    \item class;
    \item comments from the host organization;
    \item score;
    \item date;
    \item signature and official seal.
\end{itemize}

Students shall not alter or complete official evaluation fields unless
specifically authorized.

\vspace{0.3cm}

% ================================================================
% VII. SUPERVISOR'S EVALUATION AND DEFENSE RESULTS
% ================================================================

\noindent
\textbf{VII.\quad SUPERVISOR'S EVALUATION AND DEFENSE RESULTS}

\vspace{0.15cm}

The section entitled
\textbf{SUPERVISOR'S EVALUATION AND GRADUATION INTERNSHIP DEFENSE RESULTS}
is reserved for the supervisor and defense committee.

The evaluation may include:

\begin{itemize}[leftmargin=0.8cm, itemsep=2pt]
    \item Internship Progress and Student's Attitude;
    \item Report Content;
    \item Report Presentation;
    \item Other Comments;
    \item Supervisor's Opinion;
    \item Defense Result;
    \item Score.
\end{itemize}

Students shall not enter evaluation scores or defense results before they
are officially determined.

\vspace{0.3cm}

% ================================================================
% VIII. CHAPTERS, SECTIONS, AND SUBSECTIONS
% ================================================================

\noindent
\textbf{VIII.\quad CHAPTERS, SECTIONS, AND SUBSECTIONS}

\vspace{0.15cm}

The report structure shall be generated using the standard \LaTeX{}
commands:

\begin{verbatim}
\chapter{...}
\section{...}
\subsection{...}
\end{verbatim}

Chapter, section, and subsection numbers shall be generated automatically
by \LaTeX{}.

Students shall not manually type chapter or section numbers and shall not
reproduce heading styles using commands such as \verb|\textbf|,
\verb|\fontsize|, or \verb|\centering| instead of the appropriate
structural command.

Excessive subdivision should be avoided. In most cases, the hierarchy
Chapter--Section--Subsection is sufficient.

\vspace{0.3cm}

% ================================================================
% IX. ORGANIZATION OF THE MAIN CONTENT
% ================================================================

\noindent
\textbf{IX.\quad ORGANIZATION OF THE MAIN CONTENT}

\vspace{0.15cm}

The specific content of the report shall depend on the internship
organization and the assigned tasks. A typical structure may include:

\begin{itemize}[leftmargin=0.8cm, itemsep=2pt]
    \item \textbf{Chapter 1. Overview of the Internship Organization};
    \item \textbf{Chapter 2. Internship Activities and Assigned Tasks};
    \item \textbf{Chapter 3. Results, Experience, and Evaluation}.
\end{itemize}

These chapter titles are recommendations only and may be adapted to the
actual internship content.

The report shall reflect work actually performed, observed, analyzed, or
learned during the internship. Unrelated theoretical material should not
be included merely to increase report length.

\vspace{0.3cm}

% ================================================================
% X. TABLE OF CONTENTS
% ================================================================

\noindent
\textbf{X.\quad TABLE OF CONTENTS}

\vspace{0.15cm}

The Table of Contents shall be generated automatically using:

\begin{verbatim}
\tableofcontents
\end{verbatim}

Students shall not create or manually edit page numbers in the Table of
Contents.

The document shall be compiled sufficiently to ensure that all page
numbers and entries are correctly updated.

\vspace{0.3cm}

% ================================================================
% XI. LIST OF SYMBOLS AND ABBREVIATIONS
% ================================================================

\noindent
\textbf{XI.\quad LIST OF SYMBOLS AND ABBREVIATIONS}

\vspace{0.15cm}

Symbols and abbreviations that are used repeatedly in the report shall be
listed in the \textbf{List of Symbols and Abbreviations}.

The list may contain the following columns:

\begin{center}
\begin{tabular}{llll}
    No. & Symbol / Abbreviation & Unit & Definition
\end{tabular}
\end{center}

Only symbols and abbreviations actually used in the final report shall be
retained. Example entries contained in the template shall be removed if
they are not relevant to the student's report.

Each abbreviation should normally be defined when it first appears in
the main text.

\vspace{0.3cm}

% ================================================================
% XII. FIGURES, DIAGRAMS, AND PHOTOGRAPHS
% ================================================================

\noindent
\textbf{XII.\quad FIGURES, DIAGRAMS, AND PHOTOGRAPHS}

\vspace{0.15cm}

Figures shall be clear, legible, and directly relevant to the report
content.

Figures shall be inserted using the standard \verb|figure| environment.

Example:

\begin{verbatim}
\begin{figure}[H]
    \centering
    \includegraphics[width=0.8\textwidth]
                    {figures/company_structure.pdf}
    \caption{Organizational structure of the company}
    \label{fig:company-structure}
\end{figure}
\end{verbatim}

Figure captions shall be placed below the figure.

Each referenced figure shall have a unique \verb|\label| and shall be
referenced in the text using, for example:

\begin{verbatim}
Figure~\ref{fig:company-structure}
\end{verbatim}

Students shall not manually type figure numbers.

Figures reproduced or adapted from external sources shall include an
appropriate citation.

Photographs taken during the internship shall be relevant and sufficiently
clear. Students shall not include confidential, restricted, or proprietary
information without permission from the internship organization.

Vector graphics or high-resolution images are recommended whenever
possible.

\vspace{0.3cm}

% ================================================================
% XIII. TABLES
% ================================================================

\noindent
\textbf{XIII.\quad TABLES}

\vspace{0.15cm}

Tables shall be created using suitable \LaTeX{} environments such as
\verb|table|, \verb|tabular|, \verb|tabularx|, or \verb|longtable|.

Table captions shall be placed above the table.

Example:

\begin{verbatim}
\begin{table}[H]
    \centering
    \caption{Internship schedule}
    \label{tab:internship-schedule}

    \begin{tabular}{lll}
        \toprule
        Week & Task & Result \\
        \midrule
        1 & Orientation & Completed \\
        2 & System study & Completed \\
        \bottomrule
    \end{tabular}
\end{table}
\end{verbatim}

Tables shall be referenced using:

\begin{verbatim}
Table~\ref{tab:internship-schedule}
\end{verbatim}

Students shall not manually type table numbers.

Tables should be created directly in \LaTeX{} rather than inserted as
screenshots whenever practical.

\vspace{0.3cm}

% ================================================================
% XIV. EQUATIONS AND TECHNICAL QUANTITIES
% ================================================================

\noindent
\textbf{XIV.\quad EQUATIONS AND TECHNICAL QUANTITIES}

\vspace{0.15cm}

Mathematical expressions shall be typeset directly using \LaTeX{}.
Images or screenshots of equations shall not be used.

Short expressions may be written in inline mathematical mode, for example:

\begin{verbatim}
$R = 0.7\,\Omega$
\end{verbatim}

Important equations shall use a numbered equation environment.

Example:

\begin{verbatim}
\begin{equation}
    P = UI
    \label{eq:power}
\end{equation}
\end{verbatim}

Equations shall be referenced using:

\begin{verbatim}
Equation~\eqref{eq:power}
\end{verbatim}

Equation numbers shall not be manually entered.

SI units shall be used wherever applicable. The \verb|siunitx| package
should be used where appropriate.

Examples:

\begin{verbatim}
\SI{220}{\volt}
\SI{50}{\hertz}
\SI{2.5}{\ampere}
\end{verbatim}

Notation and units shall be consistent throughout the report.

\vspace{0.3cm}

% ================================================================
% XV. SOURCE CODE
% ================================================================

\noindent
\textbf{XV.\quad SOURCE CODE}

\vspace{0.15cm}

If the internship involves programming, PLCs, microcontrollers,
simulation, or software development, relevant source code may be included
in the report.

Source code shall be presented using the \verb|listings| environment or
another environment defined by the official template.

Example:

\begin{verbatim}
\begin{lstlisting}[language=st]
IF Start = TRUE THEN
    Motor := TRUE;
ELSE
    Motor := FALSE;
END_IF;
\end{lstlisting}
\end{verbatim}

Only code necessary for explaining the work should appear in the main
report. Long source-code listings should be moved to an appendix or
submitted as supplementary files if required.

Source code should not normally be inserted as screenshots.

\vspace{0.3cm}

% ================================================================
% XVI. ABSTRACT
% ================================================================

\noindent
\textbf{XVI.\quad ABSTRACT}

\vspace{0.15cm}

The Abstract shall briefly summarize:

\begin{itemize}[leftmargin=0.8cm, itemsep=2pt]
    \item the internship organization;
    \item the objectives of the internship;
    \item the main tasks performed;
    \item the main results, knowledge, and skills obtained.
\end{itemize}

The Abstract should be concise and should not reproduce the full content
of the report chapter by chapter.

\vspace{0.3cm}

% ================================================================
% XVII. REFERENCES AND CITATIONS
% ================================================================

\noindent
\textbf{XVII.\quad REFERENCES AND CITATIONS}

\vspace{0.15cm}

All books, journal papers, conference papers, standards, websites,
datasheets, manuals, technical documents, figures, tables, or other
external sources used in preparing the report shall be properly cited.

References shall be managed through the bibliography system provided with
the official template.

Bibliographic entries shall normally be stored in:

\begin{verbatim}
References.bib
\end{verbatim}

References shall be cited in the report using:

\begin{verbatim}
\cite{reference-key}
\end{verbatim}

Students shall not manually type reference numbers such as ``[1]'',
``[2]'', etc.

Titles of foreign-language publications shall remain in their original
language in the bibliography.

All borrowed ideas, data, figures, tables, equations, methods, or results
shall be properly acknowledged.

\vspace{0.3cm}

% ================================================================
% XVIII. CROSS-REFERENCING
% ================================================================

\noindent
\textbf{XVIII.\quad CROSS-REFERENCING}

\vspace{0.15cm}

Numbered elements shall be cross-referenced using the standard \LaTeX{}
mechanism.

Use the following commands where appropriate:

\begin{verbatim}
\label{...}
\ref{...}
\eqref{...}
\end{verbatim}

For example:

\begin{verbatim}
As shown in Figure~\ref{fig:system}, ...
\end{verbatim}

Students shall not manually type numbers such as ``Figure 2.3'' when
those numbers can be generated automatically.

\vspace{0.3cm}

% ================================================================
% XIX. PAGE NUMBERING
% ================================================================

\noindent
\textbf{XIX.\quad PAGE NUMBERING}

\vspace{0.15cm}

Cover pages and official evaluation forms at the beginning of the report
shall not display page numbers.

The main body of the report shall use Arabic page numbers beginning from
Chapter 1.

Page numbering shall be generated automatically by the official template.
Students shall not manually enter page numbers.

\vspace{0.3cm}

% ================================================================
% XX. PROJECT FILE ORGANIZATION
% ================================================================

\noindent
\textbf{XX.\quad PROJECT FILE ORGANIZATION}

\vspace{0.15cm}

Students should maintain a clear and organized \LaTeX{} project
structure. A recommended structure is:

\begin{verbatim}
InternshipReport/
|
|-- main.tex
|-- References.bib
|
|-- figures/
|   |-- company.pdf
|   |-- system.png
|   `-- ...
|
`-- chapters/
    |-- chapter1.tex
    |-- chapter2.tex
    `-- chapter3.tex
\end{verbatim}

Figures should normally be stored in the \verb|figures| directory
specified by the official template.

File and directory names should preferably contain Latin letters,
numbers, hyphens, or underscores.

Recommended examples include:

\begin{verbatim}
company_structure.pdf
motor_control.png
chapter_01.tex
\end{verbatim}

Spaces and unnecessary special characters in file names should be
avoided where possible.

\vspace{0.3cm}

% ================================================================
% XXI. DOCUMENT COMPILATION AND FINAL CHECK
% ================================================================

\noindent
\textbf{XXI.\quad DOCUMENT COMPILATION AND FINAL CHECK}

\vspace{0.15cm}

Before submission, the complete \LaTeX{} project shall compile
successfully and produce the final PDF report.

The final document shall not contain unresolved references such as:

\begin{verbatim}
??
[?]
\end{verbatim}

Before generating the final PDF, students shall check:

\begin{itemize}[leftmargin=0.8cm, itemsep=2pt]
    \item the Table of Contents;
    \item page numbering;
    \item chapter, section, and subsection numbering;
    \item figure numbering;
    \item table numbering;
    \item equation numbering;
    \item cross-references;
    \item citations;
    \item bibliography;
    \item font display;
    \item mathematical symbols;
    \item image quality;
    \item completeness of all required forms.
\end{itemize}

The document may be compiled using TeX Live, MiKTeX, or another
compatible \LaTeX{} environment. Regardless of the compilation
environment, the resulting PDF shall conform to the official Faculty
template.

\vspace{0.3cm}

% ================================================================
% XXII. CONTENT THAT STUDENTS MAY EDIT
% ================================================================

\noindent
\textbf{XXII.\quad CONTENT THAT STUDENTS MAY EDIT}

\vspace{0.15cm}

Students may edit or replace:

\begin{itemize}[leftmargin=0.8cm, itemsep=2pt]
    \item student name;
    \item class;
    \item internship organization;
    \item supervisor name;
    \item Abstract;
    \item List of Symbols and Abbreviations;
    \item chapter titles and chapter content;
    \item figures and photographs;
    \item tables;
    \item equations;
    \item citations and references.
\end{itemize}

Students shall not modify the main document formatting settings in the
preamble, including the document class, page geometry, line spacing,
heading formatting, table-of-contents formatting, font configuration,
page layout, or cover-page macros unless specifically instructed.

\vspace{0.3cm}

% ================================================================
% XXIII. FINAL SUBMISSION REQUIREMENTS
% ================================================================

\noindent
\textbf{XXIII.\quad FINAL SUBMISSION REQUIREMENTS}

\vspace{0.15cm}

Before submission, the final report shall:

\begin{itemize}[leftmargin=0.8cm, itemsep=2pt]
    \item compile without errors that affect the report content or presentation;
    \item contain all required forms and report sections;
    \item contain correct student, class, internship organization, and supervisor information;
    \item contain an updated Table of Contents, List of Figures, and List of Tables;
    \item contain no unused example content from the template;
    \item contain no placeholder text such as \textit{Title of Chapter 1},
    \textit{Section 1.1}, or other sample entries;
    \item contain all required citations and references;
    \item contain no unresolved cross-references.
\end{itemize}

If the Faculty requires submission of the \LaTeX{} source, the complete
project shall include all \verb|.tex| files, the bibliography database,
figures, and any other files required to reproduce the submitted PDF.

\vspace{0.4cm}

\noindent
\textbf{Important:}
Before submitting the final internship report, students shall remove or
comment out the following command in the main document:

\begin{verbatim}

% =================================================================
% REGULATIONS FOR PREPARING THE GRADUATION INTERNSHIP REPORT
% USING LATEX
% =================================================================
%
% This file is intended to be included from the main document with:
%
%     \input{Regulations.tex}
%
% IMPORTANT:
% This section is provided for guidance only. Students shall remove
% or comment out the \input{Regulations.tex} command before submitting
% the final internship report.
%
% =================================================================

\clearpage
\thispagestyle{empty}

\begin{center}
    {\fontsize{16pt}{20pt}\selectfont\bfseries
    SPECIFIC REGULATIONS FOR PREPARING\\
    GRADUATION INTERNSHIP REPORTS USING \LaTeX}
\end{center}

\vspace{0.5cm}

{\fontsize{13pt}{17pt}\selectfont
\setstretch{1.3}
\justifying

\noindent
\textbf{Note:}
This section is provided for guidance only and is included in the main
document using the command \verb|\input{Regulations.tex}|.
Students shall remove or comment out this command before submitting the
final internship report.

\vspace{0.4cm}

% ================================================================
% I. GENERAL REQUIREMENTS
% ================================================================

\noindent
\textbf{I.\quad GENERAL REQUIREMENTS}

\vspace{0.15cm}

The Graduation Internship Report shall be prepared using the official
\LaTeX{} template provided by the Faculty. Students shall preserve the
predefined page layout, typography, line spacing, heading styles,
numbering system, figure and table formatting, and bibliography settings
unless otherwise instructed by the Faculty or the supervisor.

The report shall be written in English.

Students are responsible for the academic and technical content of the
report. The official \LaTeX{} template is responsible for standardized
document formatting.

\vspace{0.3cm}

% ================================================================
% II. REPORT STRUCTURE
% ================================================================

\noindent
\textbf{II.\quad REPORT STRUCTURE}

\vspace{0.15cm}

The report shall normally be arranged in the following order:

\begin{enumerate}[label=\arabic*., leftmargin=1.2cm, itemsep=2pt]
    \item Outer Cover;
    \item Comments of the Internship Host Organization;
    \item Supervisor's Evaluation and Graduation Internship Defense Results;
    \item Table of Contents;
    \item List of Symbols and Abbreviations, if applicable;
    \item List of Figures;
    \item List of Tables;
    \item Abstract;
    \item Main Chapters of the Internship Report;
    \item References;
    \item Appendices, if applicable.
\end{enumerate}

The preliminary pages shall not be numbered as chapters. The main report
content shall be organized into consecutively numbered chapters beginning
with Chapter 1.

\vspace{0.3cm}

% ================================================================
% III. PAGE SIZE AND MARGINS
% ================================================================

\noindent
\textbf{III.\quad PAGE SIZE AND MARGINS}

\vspace{0.15cm}

The report shall be formatted on A4 paper with the following margins:

\begin{center}
\begin{tabular}{ll}
    Left margin:   & 3.5 cm \\
    Right margin:  & 2.0 cm \\
    Top margin:    & 2.5 cm \\
    Bottom margin: & 2.5 cm
\end{tabular}
\end{center}

These settings are predefined in the official template and shall not be
modified by students.

\vspace{0.3cm}

% ================================================================
% IV. FONT, FONT SIZE, AND LINE SPACING
% ================================================================

\noindent
\textbf{IV.\quad FONT, FONT SIZE, AND LINE SPACING}

\vspace{0.15cm}

The main text shall use a 13 pt font size with 1.3-line spacing.
A Times-family typeface shall be used in accordance with the official
template.

The text shall be fully justified.

Students shall not use repeated spaces, repeated manual line breaks, or
multiple blank lines to position text. Paragraph and heading spacing shall
be controlled by the template.

\vspace{0.3cm}

% ================================================================
% V. OUTER COVER
% ================================================================

\noindent
\textbf{V.\quad OUTER COVER}

\vspace{0.15cm}

The outer cover shall follow the official format provided in the template
and shall contain the following information:

\begin{itemize}[leftmargin=0.8cm, itemsep=2pt]
    \item Thai Nguyen University of Technology;
    \item Faculty of International Training;
    \item Engineering and Technology Department;
    \item University logo and Faculty logo;
    \item the title \textbf{INTERNSHIP REPORT};
    \item Internship Organization;
    \item Author;
    \item Class;
    \item Supervisor;
    \item place and year of completion.
\end{itemize}

Students shall only replace the required information fields and shall not
modify the cover-page layout unless instructed otherwise.

\vspace{0.3cm}

% ================================================================
% VI. COMMENTS OF THE INTERNSHIP HOST ORGANIZATION
% ================================================================

\noindent
\textbf{VI.\quad COMMENTS OF THE INTERNSHIP HOST ORGANIZATION}

\vspace{0.15cm}

The page entitled \textbf{COMMENTS OF THE INTERNSHIP HOST ORGANIZATION}
is reserved for the internship host organization.

It includes:

\begin{itemize}[leftmargin=0.8cm, itemsep=2pt]
    \item the student's full name;
    \item class;
    \item comments from the host organization;
    \item score;
    \item date;
    \item signature and official seal.
\end{itemize}

Students shall not alter or complete official evaluation fields unless
specifically authorized.

\vspace{0.3cm}

% ================================================================
% VII. SUPERVISOR'S EVALUATION AND DEFENSE RESULTS
% ================================================================

\noindent
\textbf{VII.\quad SUPERVISOR'S EVALUATION AND DEFENSE RESULTS}

\vspace{0.15cm}

The section entitled
\textbf{SUPERVISOR'S EVALUATION AND GRADUATION INTERNSHIP DEFENSE RESULTS}
is reserved for the supervisor and defense committee.

The evaluation may include:

\begin{itemize}[leftmargin=0.8cm, itemsep=2pt]
    \item Internship Progress and Student's Attitude;
    \item Report Content;
    \item Report Presentation;
    \item Other Comments;
    \item Supervisor's Opinion;
    \item Defense Result;
    \item Score.
\end{itemize}

Students shall not enter evaluation scores or defense results before they
are officially determined.

\vspace{0.3cm}

% ================================================================
% VIII. CHAPTERS, SECTIONS, AND SUBSECTIONS
% ================================================================

\noindent
\textbf{VIII.\quad CHAPTERS, SECTIONS, AND SUBSECTIONS}

\vspace{0.15cm}

The report structure shall be generated using the standard \LaTeX{}
commands:

\begin{verbatim}
\chapter{...}
\section{...}
\subsection{...}
\end{verbatim}

Chapter, section, and subsection numbers shall be generated automatically
by \LaTeX{}.

Students shall not manually type chapter or section numbers and shall not
reproduce heading styles using commands such as \verb|\textbf|,
\verb|\fontsize|, or \verb|\centering| instead of the appropriate
structural command.

Excessive subdivision should be avoided. In most cases, the hierarchy
Chapter--Section--Subsection is sufficient.

\vspace{0.3cm}

% ================================================================
% IX. ORGANIZATION OF THE MAIN CONTENT
% ================================================================

\noindent
\textbf{IX.\quad ORGANIZATION OF THE MAIN CONTENT}

\vspace{0.15cm}

The specific content of the report shall depend on the internship
organization and the assigned tasks. A typical structure may include:

\begin{itemize}[leftmargin=0.8cm, itemsep=2pt]
    \item \textbf{Chapter 1. Overview of the Internship Organization};
    \item \textbf{Chapter 2. Internship Activities and Assigned Tasks};
    \item \textbf{Chapter 3. Results, Experience, and Evaluation}.
\end{itemize}

These chapter titles are recommendations only and may be adapted to the
actual internship content.

The report shall reflect work actually performed, observed, analyzed, or
learned during the internship. Unrelated theoretical material should not
be included merely to increase report length.

\vspace{0.3cm}

% ================================================================
% X. TABLE OF CONTENTS
% ================================================================

\noindent
\textbf{X.\quad TABLE OF CONTENTS}

\vspace{0.15cm}

The Table of Contents shall be generated automatically using:

\begin{verbatim}
\tableofcontents
\end{verbatim}

Students shall not create or manually edit page numbers in the Table of
Contents.

The document shall be compiled sufficiently to ensure that all page
numbers and entries are correctly updated.

\vspace{0.3cm}

% ================================================================
% XI. LIST OF SYMBOLS AND ABBREVIATIONS
% ================================================================

\noindent
\textbf{XI.\quad LIST OF SYMBOLS AND ABBREVIATIONS}

\vspace{0.15cm}

Symbols and abbreviations that are used repeatedly in the report shall be
listed in the \textbf{List of Symbols and Abbreviations}.

The list may contain the following columns:

\begin{center}
\begin{tabular}{llll}
    No. & Symbol / Abbreviation & Unit & Definition
\end{tabular}
\end{center}

Only symbols and abbreviations actually used in the final report shall be
retained. Example entries contained in the template shall be removed if
they are not relevant to the student's report.

Each abbreviation should normally be defined when it first appears in
the main text.

\vspace{0.3cm}

% ================================================================
% XII. FIGURES, DIAGRAMS, AND PHOTOGRAPHS
% ================================================================

\noindent
\textbf{XII.\quad FIGURES, DIAGRAMS, AND PHOTOGRAPHS}

\vspace{0.15cm}

Figures shall be clear, legible, and directly relevant to the report
content.

Figures shall be inserted using the standard \verb|figure| environment.

Example:

\begin{verbatim}
\begin{figure}[H]
    \centering
    \includegraphics[width=0.8\textwidth]
                    {figures/company_structure.pdf}
    \caption{Organizational structure of the company}
    \label{fig:company-structure}
\end{figure}
\end{verbatim}

Figure captions shall be placed below the figure.

Each referenced figure shall have a unique \verb|\label| and shall be
referenced in the text using, for example:

\begin{verbatim}
Figure~\ref{fig:company-structure}
\end{verbatim}

Students shall not manually type figure numbers.

Figures reproduced or adapted from external sources shall include an
appropriate citation.

Photographs taken during the internship shall be relevant and sufficiently
clear. Students shall not include confidential, restricted, or proprietary
information without permission from the internship organization.

Vector graphics or high-resolution images are recommended whenever
possible.

\vspace{0.3cm}

% ================================================================
% XIII. TABLES
% ================================================================

\noindent
\textbf{XIII.\quad TABLES}

\vspace{0.15cm}

Tables shall be created using suitable \LaTeX{} environments such as
\verb|table|, \verb|tabular|, \verb|tabularx|, or \verb|longtable|.

Table captions shall be placed above the table.

Example:

\begin{verbatim}
\begin{table}[H]
    \centering
    \caption{Internship schedule}
    \label{tab:internship-schedule}

    \begin{tabular}{lll}
        \toprule
        Week & Task & Result \\
        \midrule
        1 & Orientation & Completed \\
        2 & System study & Completed \\
        \bottomrule
    \end{tabular}
\end{table}
\end{verbatim}

Tables shall be referenced using:

\begin{verbatim}
Table~\ref{tab:internship-schedule}
\end{verbatim}

Students shall not manually type table numbers.

Tables should be created directly in \LaTeX{} rather than inserted as
screenshots whenever practical.

\vspace{0.3cm}

% ================================================================
% XIV. EQUATIONS AND TECHNICAL QUANTITIES
% ================================================================

\noindent
\textbf{XIV.\quad EQUATIONS AND TECHNICAL QUANTITIES}

\vspace{0.15cm}

Mathematical expressions shall be typeset directly using \LaTeX{}.
Images or screenshots of equations shall not be used.

Short expressions may be written in inline mathematical mode, for example:

\begin{verbatim}
$R = 0.7\,\Omega$
\end{verbatim}

Important equations shall use a numbered equation environment.

Example:

\begin{verbatim}
\begin{equation}
    P = UI
    \label{eq:power}
\end{equation}
\end{verbatim}

Equations shall be referenced using:

\begin{verbatim}
Equation~\eqref{eq:power}
\end{verbatim}

Equation numbers shall not be manually entered.

SI units shall be used wherever applicable. The \verb|siunitx| package
should be used where appropriate.

Examples:

\begin{verbatim}
\SI{220}{\volt}
\SI{50}{\hertz}
\SI{2.5}{\ampere}
\end{verbatim}

Notation and units shall be consistent throughout the report.

\vspace{0.3cm}

% ================================================================
% XV. SOURCE CODE
% ================================================================

\noindent
\textbf{XV.\quad SOURCE CODE}

\vspace{0.15cm}

If the internship involves programming, PLCs, microcontrollers,
simulation, or software development, relevant source code may be included
in the report.

Source code shall be presented using the \verb|listings| environment or
another environment defined by the official template.

Example:

\begin{verbatim}
\begin{lstlisting}[language=st]
IF Start = TRUE THEN
    Motor := TRUE;
ELSE
    Motor := FALSE;
END_IF;
\end{lstlisting}
\end{verbatim}

Only code necessary for explaining the work should appear in the main
report. Long source-code listings should be moved to an appendix or
submitted as supplementary files if required.

Source code should not normally be inserted as screenshots.

\vspace{0.3cm}

% ================================================================
% XVI. ABSTRACT
% ================================================================

\noindent
\textbf{XVI.\quad ABSTRACT}

\vspace{0.15cm}

The Abstract shall briefly summarize:

\begin{itemize}[leftmargin=0.8cm, itemsep=2pt]
    \item the internship organization;
    \item the objectives of the internship;
    \item the main tasks performed;
    \item the main results, knowledge, and skills obtained.
\end{itemize}

The Abstract should be concise and should not reproduce the full content
of the report chapter by chapter.

\vspace{0.3cm}

% ================================================================
% XVII. REFERENCES AND CITATIONS
% ================================================================

\noindent
\textbf{XVII.\quad REFERENCES AND CITATIONS}

\vspace{0.15cm}

All books, journal papers, conference papers, standards, websites,
datasheets, manuals, technical documents, figures, tables, or other
external sources used in preparing the report shall be properly cited.

References shall be managed through the bibliography system provided with
the official template.

Bibliographic entries shall normally be stored in:

\begin{verbatim}
References.bib
\end{verbatim}

References shall be cited in the report using:

\begin{verbatim}
\cite{reference-key}
\end{verbatim}

Students shall not manually type reference numbers such as ``[1]'',
``[2]'', etc.

Titles of foreign-language publications shall remain in their original
language in the bibliography.

All borrowed ideas, data, figures, tables, equations, methods, or results
shall be properly acknowledged.

\vspace{0.3cm}

% ================================================================
% XVIII. CROSS-REFERENCING
% ================================================================

\noindent
\textbf{XVIII.\quad CROSS-REFERENCING}

\vspace{0.15cm}

Numbered elements shall be cross-referenced using the standard \LaTeX{}
mechanism.

Use the following commands where appropriate:

\begin{verbatim}
\label{...}
\ref{...}
\eqref{...}
\end{verbatim}

For example:

\begin{verbatim}
As shown in Figure~\ref{fig:system}, ...
\end{verbatim}

Students shall not manually type numbers such as ``Figure 2.3'' when
those numbers can be generated automatically.

\vspace{0.3cm}

% ================================================================
% XIX. PAGE NUMBERING
% ================================================================

\noindent
\textbf{XIX.\quad PAGE NUMBERING}

\vspace{0.15cm}

Cover pages and official evaluation forms at the beginning of the report
shall not display page numbers.

The main body of the report shall use Arabic page numbers beginning from
Chapter 1.

Page numbering shall be generated automatically by the official template.
Students shall not manually enter page numbers.

\vspace{0.3cm}

% ================================================================
% XX. PROJECT FILE ORGANIZATION
% ================================================================

\noindent
\textbf{XX.\quad PROJECT FILE ORGANIZATION}

\vspace{0.15cm}

Students should maintain a clear and organized \LaTeX{} project
structure. A recommended structure is:

\begin{verbatim}
InternshipReport/
|
|-- main.tex
|-- References.bib
|
|-- figures/
|   |-- company.pdf
|   |-- system.png
|   `-- ...
|
`-- chapters/
    |-- chapter1.tex
    |-- chapter2.tex
    `-- chapter3.tex
\end{verbatim}

Figures should normally be stored in the \verb|figures| directory
specified by the official template.

File and directory names should preferably contain Latin letters,
numbers, hyphens, or underscores.

Recommended examples include:

\begin{verbatim}
company_structure.pdf
motor_control.png
chapter_01.tex
\end{verbatim}

Spaces and unnecessary special characters in file names should be
avoided where possible.

\vspace{0.3cm}

% ================================================================
% XXI. DOCUMENT COMPILATION AND FINAL CHECK
% ================================================================

\noindent
\textbf{XXI.\quad DOCUMENT COMPILATION AND FINAL CHECK}

\vspace{0.15cm}

Before submission, the complete \LaTeX{} project shall compile
successfully and produce the final PDF report.

The final document shall not contain unresolved references such as:

\begin{verbatim}
??
[?]
\end{verbatim}

Before generating the final PDF, students shall check:

\begin{itemize}[leftmargin=0.8cm, itemsep=2pt]
    \item the Table of Contents;
    \item page numbering;
    \item chapter, section, and subsection numbering;
    \item figure numbering;
    \item table numbering;
    \item equation numbering;
    \item cross-references;
    \item citations;
    \item bibliography;
    \item font display;
    \item mathematical symbols;
    \item image quality;
    \item completeness of all required forms.
\end{itemize}

The document may be compiled using TeX Live, MiKTeX, or another
compatible \LaTeX{} environment. Regardless of the compilation
environment, the resulting PDF shall conform to the official Faculty
template.

\vspace{0.3cm}

% ================================================================
% XXII. CONTENT THAT STUDENTS MAY EDIT
% ================================================================

\noindent
\textbf{XXII.\quad CONTENT THAT STUDENTS MAY EDIT}

\vspace{0.15cm}

Students may edit or replace:

\begin{itemize}[leftmargin=0.8cm, itemsep=2pt]
    \item student name;
    \item class;
    \item internship organization;
    \item supervisor name;
    \item Abstract;
    \item List of Symbols and Abbreviations;
    \item chapter titles and chapter content;
    \item figures and photographs;
    \item tables;
    \item equations;
    \item citations and references.
\end{itemize}

Students shall not modify the main document formatting settings in the
preamble, including the document class, page geometry, line spacing,
heading formatting, table-of-contents formatting, font configuration,
page layout, or cover-page macros unless specifically instructed.

\vspace{0.3cm}

% ================================================================
% XXIII. FINAL SUBMISSION REQUIREMENTS
% ================================================================

\noindent
\textbf{XXIII.\quad FINAL SUBMISSION REQUIREMENTS}

\vspace{0.15cm}

Before submission, the final report shall:

\begin{itemize}[leftmargin=0.8cm, itemsep=2pt]
    \item compile without errors that affect the report content or presentation;
    \item contain all required forms and report sections;
    \item contain correct student, class, internship organization, and supervisor information;
    \item contain an updated Table of Contents, List of Figures, and List of Tables;
    \item contain no unused example content from the template;
    \item contain no placeholder text such as \textit{Title of Chapter 1},
    \textit{Section 1.1}, or other sample entries;
    \item contain all required citations and references;
    \item contain no unresolved cross-references.
\end{itemize}

If the Faculty requires submission of the \LaTeX{} source, the complete
project shall include all \verb|.tex| files, the bibliography database,
figures, and any other files required to reproduce the submitted PDF.

\vspace{0.4cm}

\noindent
\textbf{Important:}
Before submitting the final internship report, students shall remove or
comment out the following command in the main document:

\begin{verbatim}
\input{Regulations.tex}
\end{verbatim}

For example:

\begin{verbatim}
% \input{Regulations.tex}
\end{verbatim}

This removes the guidance pages from the final submitted report while
preserving the regulations file in the \LaTeX{} project for reference.

} % End 13pt content

\clearpage

\end{verbatim}

For example:

\begin{verbatim}
% 
% =================================================================
% REGULATIONS FOR PREPARING THE GRADUATION INTERNSHIP REPORT
% USING LATEX
% =================================================================
%
% This file is intended to be included from the main document with:
%
%     \input{Regulations.tex}
%
% IMPORTANT:
% This section is provided for guidance only. Students shall remove
% or comment out the \input{Regulations.tex} command before submitting
% the final internship report.
%
% =================================================================

\clearpage
\thispagestyle{empty}

\begin{center}
    {\fontsize{16pt}{20pt}\selectfont\bfseries
    SPECIFIC REGULATIONS FOR PREPARING\\
    GRADUATION INTERNSHIP REPORTS USING \LaTeX}
\end{center}

\vspace{0.5cm}

{\fontsize{13pt}{17pt}\selectfont
\setstretch{1.3}
\justifying

\noindent
\textbf{Note:}
This section is provided for guidance only and is included in the main
document using the command \verb|\input{Regulations.tex}|.
Students shall remove or comment out this command before submitting the
final internship report.

\vspace{0.4cm}

% ================================================================
% I. GENERAL REQUIREMENTS
% ================================================================

\noindent
\textbf{I.\quad GENERAL REQUIREMENTS}

\vspace{0.15cm}

The Graduation Internship Report shall be prepared using the official
\LaTeX{} template provided by the Faculty. Students shall preserve the
predefined page layout, typography, line spacing, heading styles,
numbering system, figure and table formatting, and bibliography settings
unless otherwise instructed by the Faculty or the supervisor.

The report shall be written in English.

Students are responsible for the academic and technical content of the
report. The official \LaTeX{} template is responsible for standardized
document formatting.

\vspace{0.3cm}

% ================================================================
% II. REPORT STRUCTURE
% ================================================================

\noindent
\textbf{II.\quad REPORT STRUCTURE}

\vspace{0.15cm}

The report shall normally be arranged in the following order:

\begin{enumerate}[label=\arabic*., leftmargin=1.2cm, itemsep=2pt]
    \item Outer Cover;
    \item Comments of the Internship Host Organization;
    \item Supervisor's Evaluation and Graduation Internship Defense Results;
    \item Table of Contents;
    \item List of Symbols and Abbreviations, if applicable;
    \item List of Figures;
    \item List of Tables;
    \item Abstract;
    \item Main Chapters of the Internship Report;
    \item References;
    \item Appendices, if applicable.
\end{enumerate}

The preliminary pages shall not be numbered as chapters. The main report
content shall be organized into consecutively numbered chapters beginning
with Chapter 1.

\vspace{0.3cm}

% ================================================================
% III. PAGE SIZE AND MARGINS
% ================================================================

\noindent
\textbf{III.\quad PAGE SIZE AND MARGINS}

\vspace{0.15cm}

The report shall be formatted on A4 paper with the following margins:

\begin{center}
\begin{tabular}{ll}
    Left margin:   & 3.5 cm \\
    Right margin:  & 2.0 cm \\
    Top margin:    & 2.5 cm \\
    Bottom margin: & 2.5 cm
\end{tabular}
\end{center}

These settings are predefined in the official template and shall not be
modified by students.

\vspace{0.3cm}

% ================================================================
% IV. FONT, FONT SIZE, AND LINE SPACING
% ================================================================

\noindent
\textbf{IV.\quad FONT, FONT SIZE, AND LINE SPACING}

\vspace{0.15cm}

The main text shall use a 13 pt font size with 1.3-line spacing.
A Times-family typeface shall be used in accordance with the official
template.

The text shall be fully justified.

Students shall not use repeated spaces, repeated manual line breaks, or
multiple blank lines to position text. Paragraph and heading spacing shall
be controlled by the template.

\vspace{0.3cm}

% ================================================================
% V. OUTER COVER
% ================================================================

\noindent
\textbf{V.\quad OUTER COVER}

\vspace{0.15cm}

The outer cover shall follow the official format provided in the template
and shall contain the following information:

\begin{itemize}[leftmargin=0.8cm, itemsep=2pt]
    \item Thai Nguyen University of Technology;
    \item Faculty of International Training;
    \item Engineering and Technology Department;
    \item University logo and Faculty logo;
    \item the title \textbf{INTERNSHIP REPORT};
    \item Internship Organization;
    \item Author;
    \item Class;
    \item Supervisor;
    \item place and year of completion.
\end{itemize}

Students shall only replace the required information fields and shall not
modify the cover-page layout unless instructed otherwise.

\vspace{0.3cm}

% ================================================================
% VI. COMMENTS OF THE INTERNSHIP HOST ORGANIZATION
% ================================================================

\noindent
\textbf{VI.\quad COMMENTS OF THE INTERNSHIP HOST ORGANIZATION}

\vspace{0.15cm}

The page entitled \textbf{COMMENTS OF THE INTERNSHIP HOST ORGANIZATION}
is reserved for the internship host organization.

It includes:

\begin{itemize}[leftmargin=0.8cm, itemsep=2pt]
    \item the student's full name;
    \item class;
    \item comments from the host organization;
    \item score;
    \item date;
    \item signature and official seal.
\end{itemize}

Students shall not alter or complete official evaluation fields unless
specifically authorized.

\vspace{0.3cm}

% ================================================================
% VII. SUPERVISOR'S EVALUATION AND DEFENSE RESULTS
% ================================================================

\noindent
\textbf{VII.\quad SUPERVISOR'S EVALUATION AND DEFENSE RESULTS}

\vspace{0.15cm}

The section entitled
\textbf{SUPERVISOR'S EVALUATION AND GRADUATION INTERNSHIP DEFENSE RESULTS}
is reserved for the supervisor and defense committee.

The evaluation may include:

\begin{itemize}[leftmargin=0.8cm, itemsep=2pt]
    \item Internship Progress and Student's Attitude;
    \item Report Content;
    \item Report Presentation;
    \item Other Comments;
    \item Supervisor's Opinion;
    \item Defense Result;
    \item Score.
\end{itemize}

Students shall not enter evaluation scores or defense results before they
are officially determined.

\vspace{0.3cm}

% ================================================================
% VIII. CHAPTERS, SECTIONS, AND SUBSECTIONS
% ================================================================

\noindent
\textbf{VIII.\quad CHAPTERS, SECTIONS, AND SUBSECTIONS}

\vspace{0.15cm}

The report structure shall be generated using the standard \LaTeX{}
commands:

\begin{verbatim}
\chapter{...}
\section{...}
\subsection{...}
\end{verbatim}

Chapter, section, and subsection numbers shall be generated automatically
by \LaTeX{}.

Students shall not manually type chapter or section numbers and shall not
reproduce heading styles using commands such as \verb|\textbf|,
\verb|\fontsize|, or \verb|\centering| instead of the appropriate
structural command.

Excessive subdivision should be avoided. In most cases, the hierarchy
Chapter--Section--Subsection is sufficient.

\vspace{0.3cm}

% ================================================================
% IX. ORGANIZATION OF THE MAIN CONTENT
% ================================================================

\noindent
\textbf{IX.\quad ORGANIZATION OF THE MAIN CONTENT}

\vspace{0.15cm}

The specific content of the report shall depend on the internship
organization and the assigned tasks. A typical structure may include:

\begin{itemize}[leftmargin=0.8cm, itemsep=2pt]
    \item \textbf{Chapter 1. Overview of the Internship Organization};
    \item \textbf{Chapter 2. Internship Activities and Assigned Tasks};
    \item \textbf{Chapter 3. Results, Experience, and Evaluation}.
\end{itemize}

These chapter titles are recommendations only and may be adapted to the
actual internship content.

The report shall reflect work actually performed, observed, analyzed, or
learned during the internship. Unrelated theoretical material should not
be included merely to increase report length.

\vspace{0.3cm}

% ================================================================
% X. TABLE OF CONTENTS
% ================================================================

\noindent
\textbf{X.\quad TABLE OF CONTENTS}

\vspace{0.15cm}

The Table of Contents shall be generated automatically using:

\begin{verbatim}
\tableofcontents
\end{verbatim}

Students shall not create or manually edit page numbers in the Table of
Contents.

The document shall be compiled sufficiently to ensure that all page
numbers and entries are correctly updated.

\vspace{0.3cm}

% ================================================================
% XI. LIST OF SYMBOLS AND ABBREVIATIONS
% ================================================================

\noindent
\textbf{XI.\quad LIST OF SYMBOLS AND ABBREVIATIONS}

\vspace{0.15cm}

Symbols and abbreviations that are used repeatedly in the report shall be
listed in the \textbf{List of Symbols and Abbreviations}.

The list may contain the following columns:

\begin{center}
\begin{tabular}{llll}
    No. & Symbol / Abbreviation & Unit & Definition
\end{tabular}
\end{center}

Only symbols and abbreviations actually used in the final report shall be
retained. Example entries contained in the template shall be removed if
they are not relevant to the student's report.

Each abbreviation should normally be defined when it first appears in
the main text.

\vspace{0.3cm}

% ================================================================
% XII. FIGURES, DIAGRAMS, AND PHOTOGRAPHS
% ================================================================

\noindent
\textbf{XII.\quad FIGURES, DIAGRAMS, AND PHOTOGRAPHS}

\vspace{0.15cm}

Figures shall be clear, legible, and directly relevant to the report
content.

Figures shall be inserted using the standard \verb|figure| environment.

Example:

\begin{verbatim}
\begin{figure}[H]
    \centering
    \includegraphics[width=0.8\textwidth]
                    {figures/company_structure.pdf}
    \caption{Organizational structure of the company}
    \label{fig:company-structure}
\end{figure}
\end{verbatim}

Figure captions shall be placed below the figure.

Each referenced figure shall have a unique \verb|\label| and shall be
referenced in the text using, for example:

\begin{verbatim}
Figure~\ref{fig:company-structure}
\end{verbatim}

Students shall not manually type figure numbers.

Figures reproduced or adapted from external sources shall include an
appropriate citation.

Photographs taken during the internship shall be relevant and sufficiently
clear. Students shall not include confidential, restricted, or proprietary
information without permission from the internship organization.

Vector graphics or high-resolution images are recommended whenever
possible.

\vspace{0.3cm}

% ================================================================
% XIII. TABLES
% ================================================================

\noindent
\textbf{XIII.\quad TABLES}

\vspace{0.15cm}

Tables shall be created using suitable \LaTeX{} environments such as
\verb|table|, \verb|tabular|, \verb|tabularx|, or \verb|longtable|.

Table captions shall be placed above the table.

Example:

\begin{verbatim}
\begin{table}[H]
    \centering
    \caption{Internship schedule}
    \label{tab:internship-schedule}

    \begin{tabular}{lll}
        \toprule
        Week & Task & Result \\
        \midrule
        1 & Orientation & Completed \\
        2 & System study & Completed \\
        \bottomrule
    \end{tabular}
\end{table}
\end{verbatim}

Tables shall be referenced using:

\begin{verbatim}
Table~\ref{tab:internship-schedule}
\end{verbatim}

Students shall not manually type table numbers.

Tables should be created directly in \LaTeX{} rather than inserted as
screenshots whenever practical.

\vspace{0.3cm}

% ================================================================
% XIV. EQUATIONS AND TECHNICAL QUANTITIES
% ================================================================

\noindent
\textbf{XIV.\quad EQUATIONS AND TECHNICAL QUANTITIES}

\vspace{0.15cm}

Mathematical expressions shall be typeset directly using \LaTeX{}.
Images or screenshots of equations shall not be used.

Short expressions may be written in inline mathematical mode, for example:

\begin{verbatim}
$R = 0.7\,\Omega$
\end{verbatim}

Important equations shall use a numbered equation environment.

Example:

\begin{verbatim}
\begin{equation}
    P = UI
    \label{eq:power}
\end{equation}
\end{verbatim}

Equations shall be referenced using:

\begin{verbatim}
Equation~\eqref{eq:power}
\end{verbatim}

Equation numbers shall not be manually entered.

SI units shall be used wherever applicable. The \verb|siunitx| package
should be used where appropriate.

Examples:

\begin{verbatim}
\SI{220}{\volt}
\SI{50}{\hertz}
\SI{2.5}{\ampere}
\end{verbatim}

Notation and units shall be consistent throughout the report.

\vspace{0.3cm}

% ================================================================
% XV. SOURCE CODE
% ================================================================

\noindent
\textbf{XV.\quad SOURCE CODE}

\vspace{0.15cm}

If the internship involves programming, PLCs, microcontrollers,
simulation, or software development, relevant source code may be included
in the report.

Source code shall be presented using the \verb|listings| environment or
another environment defined by the official template.

Example:

\begin{verbatim}
\begin{lstlisting}[language=st]
IF Start = TRUE THEN
    Motor := TRUE;
ELSE
    Motor := FALSE;
END_IF;
\end{lstlisting}
\end{verbatim}

Only code necessary for explaining the work should appear in the main
report. Long source-code listings should be moved to an appendix or
submitted as supplementary files if required.

Source code should not normally be inserted as screenshots.

\vspace{0.3cm}

% ================================================================
% XVI. ABSTRACT
% ================================================================

\noindent
\textbf{XVI.\quad ABSTRACT}

\vspace{0.15cm}

The Abstract shall briefly summarize:

\begin{itemize}[leftmargin=0.8cm, itemsep=2pt]
    \item the internship organization;
    \item the objectives of the internship;
    \item the main tasks performed;
    \item the main results, knowledge, and skills obtained.
\end{itemize}

The Abstract should be concise and should not reproduce the full content
of the report chapter by chapter.

\vspace{0.3cm}

% ================================================================
% XVII. REFERENCES AND CITATIONS
% ================================================================

\noindent
\textbf{XVII.\quad REFERENCES AND CITATIONS}

\vspace{0.15cm}

All books, journal papers, conference papers, standards, websites,
datasheets, manuals, technical documents, figures, tables, or other
external sources used in preparing the report shall be properly cited.

References shall be managed through the bibliography system provided with
the official template.

Bibliographic entries shall normally be stored in:

\begin{verbatim}
References.bib
\end{verbatim}

References shall be cited in the report using:

\begin{verbatim}
\cite{reference-key}
\end{verbatim}

Students shall not manually type reference numbers such as ``[1]'',
``[2]'', etc.

Titles of foreign-language publications shall remain in their original
language in the bibliography.

All borrowed ideas, data, figures, tables, equations, methods, or results
shall be properly acknowledged.

\vspace{0.3cm}

% ================================================================
% XVIII. CROSS-REFERENCING
% ================================================================

\noindent
\textbf{XVIII.\quad CROSS-REFERENCING}

\vspace{0.15cm}

Numbered elements shall be cross-referenced using the standard \LaTeX{}
mechanism.

Use the following commands where appropriate:

\begin{verbatim}
\label{...}
\ref{...}
\eqref{...}
\end{verbatim}

For example:

\begin{verbatim}
As shown in Figure~\ref{fig:system}, ...
\end{verbatim}

Students shall not manually type numbers such as ``Figure 2.3'' when
those numbers can be generated automatically.

\vspace{0.3cm}

% ================================================================
% XIX. PAGE NUMBERING
% ================================================================

\noindent
\textbf{XIX.\quad PAGE NUMBERING}

\vspace{0.15cm}

Cover pages and official evaluation forms at the beginning of the report
shall not display page numbers.

The main body of the report shall use Arabic page numbers beginning from
Chapter 1.

Page numbering shall be generated automatically by the official template.
Students shall not manually enter page numbers.

\vspace{0.3cm}

% ================================================================
% XX. PROJECT FILE ORGANIZATION
% ================================================================

\noindent
\textbf{XX.\quad PROJECT FILE ORGANIZATION}

\vspace{0.15cm}

Students should maintain a clear and organized \LaTeX{} project
structure. A recommended structure is:

\begin{verbatim}
InternshipReport/
|
|-- main.tex
|-- References.bib
|
|-- figures/
|   |-- company.pdf
|   |-- system.png
|   `-- ...
|
`-- chapters/
    |-- chapter1.tex
    |-- chapter2.tex
    `-- chapter3.tex
\end{verbatim}

Figures should normally be stored in the \verb|figures| directory
specified by the official template.

File and directory names should preferably contain Latin letters,
numbers, hyphens, or underscores.

Recommended examples include:

\begin{verbatim}
company_structure.pdf
motor_control.png
chapter_01.tex
\end{verbatim}

Spaces and unnecessary special characters in file names should be
avoided where possible.

\vspace{0.3cm}

% ================================================================
% XXI. DOCUMENT COMPILATION AND FINAL CHECK
% ================================================================

\noindent
\textbf{XXI.\quad DOCUMENT COMPILATION AND FINAL CHECK}

\vspace{0.15cm}

Before submission, the complete \LaTeX{} project shall compile
successfully and produce the final PDF report.

The final document shall not contain unresolved references such as:

\begin{verbatim}
??
[?]
\end{verbatim}

Before generating the final PDF, students shall check:

\begin{itemize}[leftmargin=0.8cm, itemsep=2pt]
    \item the Table of Contents;
    \item page numbering;
    \item chapter, section, and subsection numbering;
    \item figure numbering;
    \item table numbering;
    \item equation numbering;
    \item cross-references;
    \item citations;
    \item bibliography;
    \item font display;
    \item mathematical symbols;
    \item image quality;
    \item completeness of all required forms.
\end{itemize}

The document may be compiled using TeX Live, MiKTeX, or another
compatible \LaTeX{} environment. Regardless of the compilation
environment, the resulting PDF shall conform to the official Faculty
template.

\vspace{0.3cm}

% ================================================================
% XXII. CONTENT THAT STUDENTS MAY EDIT
% ================================================================

\noindent
\textbf{XXII.\quad CONTENT THAT STUDENTS MAY EDIT}

\vspace{0.15cm}

Students may edit or replace:

\begin{itemize}[leftmargin=0.8cm, itemsep=2pt]
    \item student name;
    \item class;
    \item internship organization;
    \item supervisor name;
    \item Abstract;
    \item List of Symbols and Abbreviations;
    \item chapter titles and chapter content;
    \item figures and photographs;
    \item tables;
    \item equations;
    \item citations and references.
\end{itemize}

Students shall not modify the main document formatting settings in the
preamble, including the document class, page geometry, line spacing,
heading formatting, table-of-contents formatting, font configuration,
page layout, or cover-page macros unless specifically instructed.

\vspace{0.3cm}

% ================================================================
% XXIII. FINAL SUBMISSION REQUIREMENTS
% ================================================================

\noindent
\textbf{XXIII.\quad FINAL SUBMISSION REQUIREMENTS}

\vspace{0.15cm}

Before submission, the final report shall:

\begin{itemize}[leftmargin=0.8cm, itemsep=2pt]
    \item compile without errors that affect the report content or presentation;
    \item contain all required forms and report sections;
    \item contain correct student, class, internship organization, and supervisor information;
    \item contain an updated Table of Contents, List of Figures, and List of Tables;
    \item contain no unused example content from the template;
    \item contain no placeholder text such as \textit{Title of Chapter 1},
    \textit{Section 1.1}, or other sample entries;
    \item contain all required citations and references;
    \item contain no unresolved cross-references.
\end{itemize}

If the Faculty requires submission of the \LaTeX{} source, the complete
project shall include all \verb|.tex| files, the bibliography database,
figures, and any other files required to reproduce the submitted PDF.

\vspace{0.4cm}

\noindent
\textbf{Important:}
Before submitting the final internship report, students shall remove or
comment out the following command in the main document:

\begin{verbatim}
\input{Regulations.tex}
\end{verbatim}

For example:

\begin{verbatim}
% \input{Regulations.tex}
\end{verbatim}

This removes the guidance pages from the final submitted report while
preserving the regulations file in the \LaTeX{} project for reference.

} % End 13pt content

\clearpage

\end{verbatim}

This removes the guidance pages from the final submitted report while
preserving the regulations file in the \LaTeX{} project for reference.

} % End 13pt content

\clearpage

%
% IMPORTANT:
% This section is provided for guidance only. Students shall remove
% or comment out the 
% =================================================================
% REGULATIONS FOR PREPARING THE GRADUATION INTERNSHIP REPORT
% USING LATEX
% =================================================================
%
% This file is intended to be included from the main document with:
%
%     
% =================================================================
% REGULATIONS FOR PREPARING THE GRADUATION INTERNSHIP REPORT
% USING LATEX
% =================================================================
%
% This file is intended to be included from the main document with:
%
%     \input{Regulations.tex}
%
% IMPORTANT:
% This section is provided for guidance only. Students shall remove
% or comment out the \input{Regulations.tex} command before submitting
% the final internship report.
%
% =================================================================

\clearpage
\thispagestyle{empty}

\begin{center}
    {\fontsize{16pt}{20pt}\selectfont\bfseries
    SPECIFIC REGULATIONS FOR PREPARING\\
    GRADUATION INTERNSHIP REPORTS USING \LaTeX}
\end{center}

\vspace{0.5cm}

{\fontsize{13pt}{17pt}\selectfont
\setstretch{1.3}
\justifying

\noindent
\textbf{Note:}
This section is provided for guidance only and is included in the main
document using the command \verb|\input{Regulations.tex}|.
Students shall remove or comment out this command before submitting the
final internship report.

\vspace{0.4cm}

% ================================================================
% I. GENERAL REQUIREMENTS
% ================================================================

\noindent
\textbf{I.\quad GENERAL REQUIREMENTS}

\vspace{0.15cm}

The Graduation Internship Report shall be prepared using the official
\LaTeX{} template provided by the Faculty. Students shall preserve the
predefined page layout, typography, line spacing, heading styles,
numbering system, figure and table formatting, and bibliography settings
unless otherwise instructed by the Faculty or the supervisor.

The report shall be written in English.

Students are responsible for the academic and technical content of the
report. The official \LaTeX{} template is responsible for standardized
document formatting.

\vspace{0.3cm}

% ================================================================
% II. REPORT STRUCTURE
% ================================================================

\noindent
\textbf{II.\quad REPORT STRUCTURE}

\vspace{0.15cm}

The report shall normally be arranged in the following order:

\begin{enumerate}[label=\arabic*., leftmargin=1.2cm, itemsep=2pt]
    \item Outer Cover;
    \item Comments of the Internship Host Organization;
    \item Supervisor's Evaluation and Graduation Internship Defense Results;
    \item Table of Contents;
    \item List of Symbols and Abbreviations, if applicable;
    \item List of Figures;
    \item List of Tables;
    \item Abstract;
    \item Main Chapters of the Internship Report;
    \item References;
    \item Appendices, if applicable.
\end{enumerate}

The preliminary pages shall not be numbered as chapters. The main report
content shall be organized into consecutively numbered chapters beginning
with Chapter 1.

\vspace{0.3cm}

% ================================================================
% III. PAGE SIZE AND MARGINS
% ================================================================

\noindent
\textbf{III.\quad PAGE SIZE AND MARGINS}

\vspace{0.15cm}

The report shall be formatted on A4 paper with the following margins:

\begin{center}
\begin{tabular}{ll}
    Left margin:   & 3.5 cm \\
    Right margin:  & 2.0 cm \\
    Top margin:    & 2.5 cm \\
    Bottom margin: & 2.5 cm
\end{tabular}
\end{center}

These settings are predefined in the official template and shall not be
modified by students.

\vspace{0.3cm}

% ================================================================
% IV. FONT, FONT SIZE, AND LINE SPACING
% ================================================================

\noindent
\textbf{IV.\quad FONT, FONT SIZE, AND LINE SPACING}

\vspace{0.15cm}

The main text shall use a 13 pt font size with 1.3-line spacing.
A Times-family typeface shall be used in accordance with the official
template.

The text shall be fully justified.

Students shall not use repeated spaces, repeated manual line breaks, or
multiple blank lines to position text. Paragraph and heading spacing shall
be controlled by the template.

\vspace{0.3cm}

% ================================================================
% V. OUTER COVER
% ================================================================

\noindent
\textbf{V.\quad OUTER COVER}

\vspace{0.15cm}

The outer cover shall follow the official format provided in the template
and shall contain the following information:

\begin{itemize}[leftmargin=0.8cm, itemsep=2pt]
    \item Thai Nguyen University of Technology;
    \item Faculty of International Training;
    \item Engineering and Technology Department;
    \item University logo and Faculty logo;
    \item the title \textbf{INTERNSHIP REPORT};
    \item Internship Organization;
    \item Author;
    \item Class;
    \item Supervisor;
    \item place and year of completion.
\end{itemize}

Students shall only replace the required information fields and shall not
modify the cover-page layout unless instructed otherwise.

\vspace{0.3cm}

% ================================================================
% VI. COMMENTS OF THE INTERNSHIP HOST ORGANIZATION
% ================================================================

\noindent
\textbf{VI.\quad COMMENTS OF THE INTERNSHIP HOST ORGANIZATION}

\vspace{0.15cm}

The page entitled \textbf{COMMENTS OF THE INTERNSHIP HOST ORGANIZATION}
is reserved for the internship host organization.

It includes:

\begin{itemize}[leftmargin=0.8cm, itemsep=2pt]
    \item the student's full name;
    \item class;
    \item comments from the host organization;
    \item score;
    \item date;
    \item signature and official seal.
\end{itemize}

Students shall not alter or complete official evaluation fields unless
specifically authorized.

\vspace{0.3cm}

% ================================================================
% VII. SUPERVISOR'S EVALUATION AND DEFENSE RESULTS
% ================================================================

\noindent
\textbf{VII.\quad SUPERVISOR'S EVALUATION AND DEFENSE RESULTS}

\vspace{0.15cm}

The section entitled
\textbf{SUPERVISOR'S EVALUATION AND GRADUATION INTERNSHIP DEFENSE RESULTS}
is reserved for the supervisor and defense committee.

The evaluation may include:

\begin{itemize}[leftmargin=0.8cm, itemsep=2pt]
    \item Internship Progress and Student's Attitude;
    \item Report Content;
    \item Report Presentation;
    \item Other Comments;
    \item Supervisor's Opinion;
    \item Defense Result;
    \item Score.
\end{itemize}

Students shall not enter evaluation scores or defense results before they
are officially determined.

\vspace{0.3cm}

% ================================================================
% VIII. CHAPTERS, SECTIONS, AND SUBSECTIONS
% ================================================================

\noindent
\textbf{VIII.\quad CHAPTERS, SECTIONS, AND SUBSECTIONS}

\vspace{0.15cm}

The report structure shall be generated using the standard \LaTeX{}
commands:

\begin{verbatim}
\chapter{...}
\section{...}
\subsection{...}
\end{verbatim}

Chapter, section, and subsection numbers shall be generated automatically
by \LaTeX{}.

Students shall not manually type chapter or section numbers and shall not
reproduce heading styles using commands such as \verb|\textbf|,
\verb|\fontsize|, or \verb|\centering| instead of the appropriate
structural command.

Excessive subdivision should be avoided. In most cases, the hierarchy
Chapter--Section--Subsection is sufficient.

\vspace{0.3cm}

% ================================================================
% IX. ORGANIZATION OF THE MAIN CONTENT
% ================================================================

\noindent
\textbf{IX.\quad ORGANIZATION OF THE MAIN CONTENT}

\vspace{0.15cm}

The specific content of the report shall depend on the internship
organization and the assigned tasks. A typical structure may include:

\begin{itemize}[leftmargin=0.8cm, itemsep=2pt]
    \item \textbf{Chapter 1. Overview of the Internship Organization};
    \item \textbf{Chapter 2. Internship Activities and Assigned Tasks};
    \item \textbf{Chapter 3. Results, Experience, and Evaluation}.
\end{itemize}

These chapter titles are recommendations only and may be adapted to the
actual internship content.

The report shall reflect work actually performed, observed, analyzed, or
learned during the internship. Unrelated theoretical material should not
be included merely to increase report length.

\vspace{0.3cm}

% ================================================================
% X. TABLE OF CONTENTS
% ================================================================

\noindent
\textbf{X.\quad TABLE OF CONTENTS}

\vspace{0.15cm}

The Table of Contents shall be generated automatically using:

\begin{verbatim}
\tableofcontents
\end{verbatim}

Students shall not create or manually edit page numbers in the Table of
Contents.

The document shall be compiled sufficiently to ensure that all page
numbers and entries are correctly updated.

\vspace{0.3cm}

% ================================================================
% XI. LIST OF SYMBOLS AND ABBREVIATIONS
% ================================================================

\noindent
\textbf{XI.\quad LIST OF SYMBOLS AND ABBREVIATIONS}

\vspace{0.15cm}

Symbols and abbreviations that are used repeatedly in the report shall be
listed in the \textbf{List of Symbols and Abbreviations}.

The list may contain the following columns:

\begin{center}
\begin{tabular}{llll}
    No. & Symbol / Abbreviation & Unit & Definition
\end{tabular}
\end{center}

Only symbols and abbreviations actually used in the final report shall be
retained. Example entries contained in the template shall be removed if
they are not relevant to the student's report.

Each abbreviation should normally be defined when it first appears in
the main text.

\vspace{0.3cm}

% ================================================================
% XII. FIGURES, DIAGRAMS, AND PHOTOGRAPHS
% ================================================================

\noindent
\textbf{XII.\quad FIGURES, DIAGRAMS, AND PHOTOGRAPHS}

\vspace{0.15cm}

Figures shall be clear, legible, and directly relevant to the report
content.

Figures shall be inserted using the standard \verb|figure| environment.

Example:

\begin{verbatim}
\begin{figure}[H]
    \centering
    \includegraphics[width=0.8\textwidth]
                    {figures/company_structure.pdf}
    \caption{Organizational structure of the company}
    \label{fig:company-structure}
\end{figure}
\end{verbatim}

Figure captions shall be placed below the figure.

Each referenced figure shall have a unique \verb|\label| and shall be
referenced in the text using, for example:

\begin{verbatim}
Figure~\ref{fig:company-structure}
\end{verbatim}

Students shall not manually type figure numbers.

Figures reproduced or adapted from external sources shall include an
appropriate citation.

Photographs taken during the internship shall be relevant and sufficiently
clear. Students shall not include confidential, restricted, or proprietary
information without permission from the internship organization.

Vector graphics or high-resolution images are recommended whenever
possible.

\vspace{0.3cm}

% ================================================================
% XIII. TABLES
% ================================================================

\noindent
\textbf{XIII.\quad TABLES}

\vspace{0.15cm}

Tables shall be created using suitable \LaTeX{} environments such as
\verb|table|, \verb|tabular|, \verb|tabularx|, or \verb|longtable|.

Table captions shall be placed above the table.

Example:

\begin{verbatim}
\begin{table}[H]
    \centering
    \caption{Internship schedule}
    \label{tab:internship-schedule}

    \begin{tabular}{lll}
        \toprule
        Week & Task & Result \\
        \midrule
        1 & Orientation & Completed \\
        2 & System study & Completed \\
        \bottomrule
    \end{tabular}
\end{table}
\end{verbatim}

Tables shall be referenced using:

\begin{verbatim}
Table~\ref{tab:internship-schedule}
\end{verbatim}

Students shall not manually type table numbers.

Tables should be created directly in \LaTeX{} rather than inserted as
screenshots whenever practical.

\vspace{0.3cm}

% ================================================================
% XIV. EQUATIONS AND TECHNICAL QUANTITIES
% ================================================================

\noindent
\textbf{XIV.\quad EQUATIONS AND TECHNICAL QUANTITIES}

\vspace{0.15cm}

Mathematical expressions shall be typeset directly using \LaTeX{}.
Images or screenshots of equations shall not be used.

Short expressions may be written in inline mathematical mode, for example:

\begin{verbatim}
$R = 0.7\,\Omega$
\end{verbatim}

Important equations shall use a numbered equation environment.

Example:

\begin{verbatim}
\begin{equation}
    P = UI
    \label{eq:power}
\end{equation}
\end{verbatim}

Equations shall be referenced using:

\begin{verbatim}
Equation~\eqref{eq:power}
\end{verbatim}

Equation numbers shall not be manually entered.

SI units shall be used wherever applicable. The \verb|siunitx| package
should be used where appropriate.

Examples:

\begin{verbatim}
\SI{220}{\volt}
\SI{50}{\hertz}
\SI{2.5}{\ampere}
\end{verbatim}

Notation and units shall be consistent throughout the report.

\vspace{0.3cm}

% ================================================================
% XV. SOURCE CODE
% ================================================================

\noindent
\textbf{XV.\quad SOURCE CODE}

\vspace{0.15cm}

If the internship involves programming, PLCs, microcontrollers,
simulation, or software development, relevant source code may be included
in the report.

Source code shall be presented using the \verb|listings| environment or
another environment defined by the official template.

Example:

\begin{verbatim}
\begin{lstlisting}[language=st]
IF Start = TRUE THEN
    Motor := TRUE;
ELSE
    Motor := FALSE;
END_IF;
\end{lstlisting}
\end{verbatim}

Only code necessary for explaining the work should appear in the main
report. Long source-code listings should be moved to an appendix or
submitted as supplementary files if required.

Source code should not normally be inserted as screenshots.

\vspace{0.3cm}

% ================================================================
% XVI. ABSTRACT
% ================================================================

\noindent
\textbf{XVI.\quad ABSTRACT}

\vspace{0.15cm}

The Abstract shall briefly summarize:

\begin{itemize}[leftmargin=0.8cm, itemsep=2pt]
    \item the internship organization;
    \item the objectives of the internship;
    \item the main tasks performed;
    \item the main results, knowledge, and skills obtained.
\end{itemize}

The Abstract should be concise and should not reproduce the full content
of the report chapter by chapter.

\vspace{0.3cm}

% ================================================================
% XVII. REFERENCES AND CITATIONS
% ================================================================

\noindent
\textbf{XVII.\quad REFERENCES AND CITATIONS}

\vspace{0.15cm}

All books, journal papers, conference papers, standards, websites,
datasheets, manuals, technical documents, figures, tables, or other
external sources used in preparing the report shall be properly cited.

References shall be managed through the bibliography system provided with
the official template.

Bibliographic entries shall normally be stored in:

\begin{verbatim}
References.bib
\end{verbatim}

References shall be cited in the report using:

\begin{verbatim}
\cite{reference-key}
\end{verbatim}

Students shall not manually type reference numbers such as ``[1]'',
``[2]'', etc.

Titles of foreign-language publications shall remain in their original
language in the bibliography.

All borrowed ideas, data, figures, tables, equations, methods, or results
shall be properly acknowledged.

\vspace{0.3cm}

% ================================================================
% XVIII. CROSS-REFERENCING
% ================================================================

\noindent
\textbf{XVIII.\quad CROSS-REFERENCING}

\vspace{0.15cm}

Numbered elements shall be cross-referenced using the standard \LaTeX{}
mechanism.

Use the following commands where appropriate:

\begin{verbatim}
\label{...}
\ref{...}
\eqref{...}
\end{verbatim}

For example:

\begin{verbatim}
As shown in Figure~\ref{fig:system}, ...
\end{verbatim}

Students shall not manually type numbers such as ``Figure 2.3'' when
those numbers can be generated automatically.

\vspace{0.3cm}

% ================================================================
% XIX. PAGE NUMBERING
% ================================================================

\noindent
\textbf{XIX.\quad PAGE NUMBERING}

\vspace{0.15cm}

Cover pages and official evaluation forms at the beginning of the report
shall not display page numbers.

The main body of the report shall use Arabic page numbers beginning from
Chapter 1.

Page numbering shall be generated automatically by the official template.
Students shall not manually enter page numbers.

\vspace{0.3cm}

% ================================================================
% XX. PROJECT FILE ORGANIZATION
% ================================================================

\noindent
\textbf{XX.\quad PROJECT FILE ORGANIZATION}

\vspace{0.15cm}

Students should maintain a clear and organized \LaTeX{} project
structure. A recommended structure is:

\begin{verbatim}
InternshipReport/
|
|-- main.tex
|-- References.bib
|
|-- figures/
|   |-- company.pdf
|   |-- system.png
|   `-- ...
|
`-- chapters/
    |-- chapter1.tex
    |-- chapter2.tex
    `-- chapter3.tex
\end{verbatim}

Figures should normally be stored in the \verb|figures| directory
specified by the official template.

File and directory names should preferably contain Latin letters,
numbers, hyphens, or underscores.

Recommended examples include:

\begin{verbatim}
company_structure.pdf
motor_control.png
chapter_01.tex
\end{verbatim}

Spaces and unnecessary special characters in file names should be
avoided where possible.

\vspace{0.3cm}

% ================================================================
% XXI. DOCUMENT COMPILATION AND FINAL CHECK
% ================================================================

\noindent
\textbf{XXI.\quad DOCUMENT COMPILATION AND FINAL CHECK}

\vspace{0.15cm}

Before submission, the complete \LaTeX{} project shall compile
successfully and produce the final PDF report.

The final document shall not contain unresolved references such as:

\begin{verbatim}
??
[?]
\end{verbatim}

Before generating the final PDF, students shall check:

\begin{itemize}[leftmargin=0.8cm, itemsep=2pt]
    \item the Table of Contents;
    \item page numbering;
    \item chapter, section, and subsection numbering;
    \item figure numbering;
    \item table numbering;
    \item equation numbering;
    \item cross-references;
    \item citations;
    \item bibliography;
    \item font display;
    \item mathematical symbols;
    \item image quality;
    \item completeness of all required forms.
\end{itemize}

The document may be compiled using TeX Live, MiKTeX, or another
compatible \LaTeX{} environment. Regardless of the compilation
environment, the resulting PDF shall conform to the official Faculty
template.

\vspace{0.3cm}

% ================================================================
% XXII. CONTENT THAT STUDENTS MAY EDIT
% ================================================================

\noindent
\textbf{XXII.\quad CONTENT THAT STUDENTS MAY EDIT}

\vspace{0.15cm}

Students may edit or replace:

\begin{itemize}[leftmargin=0.8cm, itemsep=2pt]
    \item student name;
    \item class;
    \item internship organization;
    \item supervisor name;
    \item Abstract;
    \item List of Symbols and Abbreviations;
    \item chapter titles and chapter content;
    \item figures and photographs;
    \item tables;
    \item equations;
    \item citations and references.
\end{itemize}

Students shall not modify the main document formatting settings in the
preamble, including the document class, page geometry, line spacing,
heading formatting, table-of-contents formatting, font configuration,
page layout, or cover-page macros unless specifically instructed.

\vspace{0.3cm}

% ================================================================
% XXIII. FINAL SUBMISSION REQUIREMENTS
% ================================================================

\noindent
\textbf{XXIII.\quad FINAL SUBMISSION REQUIREMENTS}

\vspace{0.15cm}

Before submission, the final report shall:

\begin{itemize}[leftmargin=0.8cm, itemsep=2pt]
    \item compile without errors that affect the report content or presentation;
    \item contain all required forms and report sections;
    \item contain correct student, class, internship organization, and supervisor information;
    \item contain an updated Table of Contents, List of Figures, and List of Tables;
    \item contain no unused example content from the template;
    \item contain no placeholder text such as \textit{Title of Chapter 1},
    \textit{Section 1.1}, or other sample entries;
    \item contain all required citations and references;
    \item contain no unresolved cross-references.
\end{itemize}

If the Faculty requires submission of the \LaTeX{} source, the complete
project shall include all \verb|.tex| files, the bibliography database,
figures, and any other files required to reproduce the submitted PDF.

\vspace{0.4cm}

\noindent
\textbf{Important:}
Before submitting the final internship report, students shall remove or
comment out the following command in the main document:

\begin{verbatim}
\input{Regulations.tex}
\end{verbatim}

For example:

\begin{verbatim}
% \input{Regulations.tex}
\end{verbatim}

This removes the guidance pages from the final submitted report while
preserving the regulations file in the \LaTeX{} project for reference.

} % End 13pt content

\clearpage

%
% IMPORTANT:
% This section is provided for guidance only. Students shall remove
% or comment out the 
% =================================================================
% REGULATIONS FOR PREPARING THE GRADUATION INTERNSHIP REPORT
% USING LATEX
% =================================================================
%
% This file is intended to be included from the main document with:
%
%     \input{Regulations.tex}
%
% IMPORTANT:
% This section is provided for guidance only. Students shall remove
% or comment out the \input{Regulations.tex} command before submitting
% the final internship report.
%
% =================================================================

\clearpage
\thispagestyle{empty}

\begin{center}
    {\fontsize{16pt}{20pt}\selectfont\bfseries
    SPECIFIC REGULATIONS FOR PREPARING\\
    GRADUATION INTERNSHIP REPORTS USING \LaTeX}
\end{center}

\vspace{0.5cm}

{\fontsize{13pt}{17pt}\selectfont
\setstretch{1.3}
\justifying

\noindent
\textbf{Note:}
This section is provided for guidance only and is included in the main
document using the command \verb|\input{Regulations.tex}|.
Students shall remove or comment out this command before submitting the
final internship report.

\vspace{0.4cm}

% ================================================================
% I. GENERAL REQUIREMENTS
% ================================================================

\noindent
\textbf{I.\quad GENERAL REQUIREMENTS}

\vspace{0.15cm}

The Graduation Internship Report shall be prepared using the official
\LaTeX{} template provided by the Faculty. Students shall preserve the
predefined page layout, typography, line spacing, heading styles,
numbering system, figure and table formatting, and bibliography settings
unless otherwise instructed by the Faculty or the supervisor.

The report shall be written in English.

Students are responsible for the academic and technical content of the
report. The official \LaTeX{} template is responsible for standardized
document formatting.

\vspace{0.3cm}

% ================================================================
% II. REPORT STRUCTURE
% ================================================================

\noindent
\textbf{II.\quad REPORT STRUCTURE}

\vspace{0.15cm}

The report shall normally be arranged in the following order:

\begin{enumerate}[label=\arabic*., leftmargin=1.2cm, itemsep=2pt]
    \item Outer Cover;
    \item Comments of the Internship Host Organization;
    \item Supervisor's Evaluation and Graduation Internship Defense Results;
    \item Table of Contents;
    \item List of Symbols and Abbreviations, if applicable;
    \item List of Figures;
    \item List of Tables;
    \item Abstract;
    \item Main Chapters of the Internship Report;
    \item References;
    \item Appendices, if applicable.
\end{enumerate}

The preliminary pages shall not be numbered as chapters. The main report
content shall be organized into consecutively numbered chapters beginning
with Chapter 1.

\vspace{0.3cm}

% ================================================================
% III. PAGE SIZE AND MARGINS
% ================================================================

\noindent
\textbf{III.\quad PAGE SIZE AND MARGINS}

\vspace{0.15cm}

The report shall be formatted on A4 paper with the following margins:

\begin{center}
\begin{tabular}{ll}
    Left margin:   & 3.5 cm \\
    Right margin:  & 2.0 cm \\
    Top margin:    & 2.5 cm \\
    Bottom margin: & 2.5 cm
\end{tabular}
\end{center}

These settings are predefined in the official template and shall not be
modified by students.

\vspace{0.3cm}

% ================================================================
% IV. FONT, FONT SIZE, AND LINE SPACING
% ================================================================

\noindent
\textbf{IV.\quad FONT, FONT SIZE, AND LINE SPACING}

\vspace{0.15cm}

The main text shall use a 13 pt font size with 1.3-line spacing.
A Times-family typeface shall be used in accordance with the official
template.

The text shall be fully justified.

Students shall not use repeated spaces, repeated manual line breaks, or
multiple blank lines to position text. Paragraph and heading spacing shall
be controlled by the template.

\vspace{0.3cm}

% ================================================================
% V. OUTER COVER
% ================================================================

\noindent
\textbf{V.\quad OUTER COVER}

\vspace{0.15cm}

The outer cover shall follow the official format provided in the template
and shall contain the following information:

\begin{itemize}[leftmargin=0.8cm, itemsep=2pt]
    \item Thai Nguyen University of Technology;
    \item Faculty of International Training;
    \item Engineering and Technology Department;
    \item University logo and Faculty logo;
    \item the title \textbf{INTERNSHIP REPORT};
    \item Internship Organization;
    \item Author;
    \item Class;
    \item Supervisor;
    \item place and year of completion.
\end{itemize}

Students shall only replace the required information fields and shall not
modify the cover-page layout unless instructed otherwise.

\vspace{0.3cm}

% ================================================================
% VI. COMMENTS OF THE INTERNSHIP HOST ORGANIZATION
% ================================================================

\noindent
\textbf{VI.\quad COMMENTS OF THE INTERNSHIP HOST ORGANIZATION}

\vspace{0.15cm}

The page entitled \textbf{COMMENTS OF THE INTERNSHIP HOST ORGANIZATION}
is reserved for the internship host organization.

It includes:

\begin{itemize}[leftmargin=0.8cm, itemsep=2pt]
    \item the student's full name;
    \item class;
    \item comments from the host organization;
    \item score;
    \item date;
    \item signature and official seal.
\end{itemize}

Students shall not alter or complete official evaluation fields unless
specifically authorized.

\vspace{0.3cm}

% ================================================================
% VII. SUPERVISOR'S EVALUATION AND DEFENSE RESULTS
% ================================================================

\noindent
\textbf{VII.\quad SUPERVISOR'S EVALUATION AND DEFENSE RESULTS}

\vspace{0.15cm}

The section entitled
\textbf{SUPERVISOR'S EVALUATION AND GRADUATION INTERNSHIP DEFENSE RESULTS}
is reserved for the supervisor and defense committee.

The evaluation may include:

\begin{itemize}[leftmargin=0.8cm, itemsep=2pt]
    \item Internship Progress and Student's Attitude;
    \item Report Content;
    \item Report Presentation;
    \item Other Comments;
    \item Supervisor's Opinion;
    \item Defense Result;
    \item Score.
\end{itemize}

Students shall not enter evaluation scores or defense results before they
are officially determined.

\vspace{0.3cm}

% ================================================================
% VIII. CHAPTERS, SECTIONS, AND SUBSECTIONS
% ================================================================

\noindent
\textbf{VIII.\quad CHAPTERS, SECTIONS, AND SUBSECTIONS}

\vspace{0.15cm}

The report structure shall be generated using the standard \LaTeX{}
commands:

\begin{verbatim}
\chapter{...}
\section{...}
\subsection{...}
\end{verbatim}

Chapter, section, and subsection numbers shall be generated automatically
by \LaTeX{}.

Students shall not manually type chapter or section numbers and shall not
reproduce heading styles using commands such as \verb|\textbf|,
\verb|\fontsize|, or \verb|\centering| instead of the appropriate
structural command.

Excessive subdivision should be avoided. In most cases, the hierarchy
Chapter--Section--Subsection is sufficient.

\vspace{0.3cm}

% ================================================================
% IX. ORGANIZATION OF THE MAIN CONTENT
% ================================================================

\noindent
\textbf{IX.\quad ORGANIZATION OF THE MAIN CONTENT}

\vspace{0.15cm}

The specific content of the report shall depend on the internship
organization and the assigned tasks. A typical structure may include:

\begin{itemize}[leftmargin=0.8cm, itemsep=2pt]
    \item \textbf{Chapter 1. Overview of the Internship Organization};
    \item \textbf{Chapter 2. Internship Activities and Assigned Tasks};
    \item \textbf{Chapter 3. Results, Experience, and Evaluation}.
\end{itemize}

These chapter titles are recommendations only and may be adapted to the
actual internship content.

The report shall reflect work actually performed, observed, analyzed, or
learned during the internship. Unrelated theoretical material should not
be included merely to increase report length.

\vspace{0.3cm}

% ================================================================
% X. TABLE OF CONTENTS
% ================================================================

\noindent
\textbf{X.\quad TABLE OF CONTENTS}

\vspace{0.15cm}

The Table of Contents shall be generated automatically using:

\begin{verbatim}
\tableofcontents
\end{verbatim}

Students shall not create or manually edit page numbers in the Table of
Contents.

The document shall be compiled sufficiently to ensure that all page
numbers and entries are correctly updated.

\vspace{0.3cm}

% ================================================================
% XI. LIST OF SYMBOLS AND ABBREVIATIONS
% ================================================================

\noindent
\textbf{XI.\quad LIST OF SYMBOLS AND ABBREVIATIONS}

\vspace{0.15cm}

Symbols and abbreviations that are used repeatedly in the report shall be
listed in the \textbf{List of Symbols and Abbreviations}.

The list may contain the following columns:

\begin{center}
\begin{tabular}{llll}
    No. & Symbol / Abbreviation & Unit & Definition
\end{tabular}
\end{center}

Only symbols and abbreviations actually used in the final report shall be
retained. Example entries contained in the template shall be removed if
they are not relevant to the student's report.

Each abbreviation should normally be defined when it first appears in
the main text.

\vspace{0.3cm}

% ================================================================
% XII. FIGURES, DIAGRAMS, AND PHOTOGRAPHS
% ================================================================

\noindent
\textbf{XII.\quad FIGURES, DIAGRAMS, AND PHOTOGRAPHS}

\vspace{0.15cm}

Figures shall be clear, legible, and directly relevant to the report
content.

Figures shall be inserted using the standard \verb|figure| environment.

Example:

\begin{verbatim}
\begin{figure}[H]
    \centering
    \includegraphics[width=0.8\textwidth]
                    {figures/company_structure.pdf}
    \caption{Organizational structure of the company}
    \label{fig:company-structure}
\end{figure}
\end{verbatim}

Figure captions shall be placed below the figure.

Each referenced figure shall have a unique \verb|\label| and shall be
referenced in the text using, for example:

\begin{verbatim}
Figure~\ref{fig:company-structure}
\end{verbatim}

Students shall not manually type figure numbers.

Figures reproduced or adapted from external sources shall include an
appropriate citation.

Photographs taken during the internship shall be relevant and sufficiently
clear. Students shall not include confidential, restricted, or proprietary
information without permission from the internship organization.

Vector graphics or high-resolution images are recommended whenever
possible.

\vspace{0.3cm}

% ================================================================
% XIII. TABLES
% ================================================================

\noindent
\textbf{XIII.\quad TABLES}

\vspace{0.15cm}

Tables shall be created using suitable \LaTeX{} environments such as
\verb|table|, \verb|tabular|, \verb|tabularx|, or \verb|longtable|.

Table captions shall be placed above the table.

Example:

\begin{verbatim}
\begin{table}[H]
    \centering
    \caption{Internship schedule}
    \label{tab:internship-schedule}

    \begin{tabular}{lll}
        \toprule
        Week & Task & Result \\
        \midrule
        1 & Orientation & Completed \\
        2 & System study & Completed \\
        \bottomrule
    \end{tabular}
\end{table}
\end{verbatim}

Tables shall be referenced using:

\begin{verbatim}
Table~\ref{tab:internship-schedule}
\end{verbatim}

Students shall not manually type table numbers.

Tables should be created directly in \LaTeX{} rather than inserted as
screenshots whenever practical.

\vspace{0.3cm}

% ================================================================
% XIV. EQUATIONS AND TECHNICAL QUANTITIES
% ================================================================

\noindent
\textbf{XIV.\quad EQUATIONS AND TECHNICAL QUANTITIES}

\vspace{0.15cm}

Mathematical expressions shall be typeset directly using \LaTeX{}.
Images or screenshots of equations shall not be used.

Short expressions may be written in inline mathematical mode, for example:

\begin{verbatim}
$R = 0.7\,\Omega$
\end{verbatim}

Important equations shall use a numbered equation environment.

Example:

\begin{verbatim}
\begin{equation}
    P = UI
    \label{eq:power}
\end{equation}
\end{verbatim}

Equations shall be referenced using:

\begin{verbatim}
Equation~\eqref{eq:power}
\end{verbatim}

Equation numbers shall not be manually entered.

SI units shall be used wherever applicable. The \verb|siunitx| package
should be used where appropriate.

Examples:

\begin{verbatim}
\SI{220}{\volt}
\SI{50}{\hertz}
\SI{2.5}{\ampere}
\end{verbatim}

Notation and units shall be consistent throughout the report.

\vspace{0.3cm}

% ================================================================
% XV. SOURCE CODE
% ================================================================

\noindent
\textbf{XV.\quad SOURCE CODE}

\vspace{0.15cm}

If the internship involves programming, PLCs, microcontrollers,
simulation, or software development, relevant source code may be included
in the report.

Source code shall be presented using the \verb|listings| environment or
another environment defined by the official template.

Example:

\begin{verbatim}
\begin{lstlisting}[language=st]
IF Start = TRUE THEN
    Motor := TRUE;
ELSE
    Motor := FALSE;
END_IF;
\end{lstlisting}
\end{verbatim}

Only code necessary for explaining the work should appear in the main
report. Long source-code listings should be moved to an appendix or
submitted as supplementary files if required.

Source code should not normally be inserted as screenshots.

\vspace{0.3cm}

% ================================================================
% XVI. ABSTRACT
% ================================================================

\noindent
\textbf{XVI.\quad ABSTRACT}

\vspace{0.15cm}

The Abstract shall briefly summarize:

\begin{itemize}[leftmargin=0.8cm, itemsep=2pt]
    \item the internship organization;
    \item the objectives of the internship;
    \item the main tasks performed;
    \item the main results, knowledge, and skills obtained.
\end{itemize}

The Abstract should be concise and should not reproduce the full content
of the report chapter by chapter.

\vspace{0.3cm}

% ================================================================
% XVII. REFERENCES AND CITATIONS
% ================================================================

\noindent
\textbf{XVII.\quad REFERENCES AND CITATIONS}

\vspace{0.15cm}

All books, journal papers, conference papers, standards, websites,
datasheets, manuals, technical documents, figures, tables, or other
external sources used in preparing the report shall be properly cited.

References shall be managed through the bibliography system provided with
the official template.

Bibliographic entries shall normally be stored in:

\begin{verbatim}
References.bib
\end{verbatim}

References shall be cited in the report using:

\begin{verbatim}
\cite{reference-key}
\end{verbatim}

Students shall not manually type reference numbers such as ``[1]'',
``[2]'', etc.

Titles of foreign-language publications shall remain in their original
language in the bibliography.

All borrowed ideas, data, figures, tables, equations, methods, or results
shall be properly acknowledged.

\vspace{0.3cm}

% ================================================================
% XVIII. CROSS-REFERENCING
% ================================================================

\noindent
\textbf{XVIII.\quad CROSS-REFERENCING}

\vspace{0.15cm}

Numbered elements shall be cross-referenced using the standard \LaTeX{}
mechanism.

Use the following commands where appropriate:

\begin{verbatim}
\label{...}
\ref{...}
\eqref{...}
\end{verbatim}

For example:

\begin{verbatim}
As shown in Figure~\ref{fig:system}, ...
\end{verbatim}

Students shall not manually type numbers such as ``Figure 2.3'' when
those numbers can be generated automatically.

\vspace{0.3cm}

% ================================================================
% XIX. PAGE NUMBERING
% ================================================================

\noindent
\textbf{XIX.\quad PAGE NUMBERING}

\vspace{0.15cm}

Cover pages and official evaluation forms at the beginning of the report
shall not display page numbers.

The main body of the report shall use Arabic page numbers beginning from
Chapter 1.

Page numbering shall be generated automatically by the official template.
Students shall not manually enter page numbers.

\vspace{0.3cm}

% ================================================================
% XX. PROJECT FILE ORGANIZATION
% ================================================================

\noindent
\textbf{XX.\quad PROJECT FILE ORGANIZATION}

\vspace{0.15cm}

Students should maintain a clear and organized \LaTeX{} project
structure. A recommended structure is:

\begin{verbatim}
InternshipReport/
|
|-- main.tex
|-- References.bib
|
|-- figures/
|   |-- company.pdf
|   |-- system.png
|   `-- ...
|
`-- chapters/
    |-- chapter1.tex
    |-- chapter2.tex
    `-- chapter3.tex
\end{verbatim}

Figures should normally be stored in the \verb|figures| directory
specified by the official template.

File and directory names should preferably contain Latin letters,
numbers, hyphens, or underscores.

Recommended examples include:

\begin{verbatim}
company_structure.pdf
motor_control.png
chapter_01.tex
\end{verbatim}

Spaces and unnecessary special characters in file names should be
avoided where possible.

\vspace{0.3cm}

% ================================================================
% XXI. DOCUMENT COMPILATION AND FINAL CHECK
% ================================================================

\noindent
\textbf{XXI.\quad DOCUMENT COMPILATION AND FINAL CHECK}

\vspace{0.15cm}

Before submission, the complete \LaTeX{} project shall compile
successfully and produce the final PDF report.

The final document shall not contain unresolved references such as:

\begin{verbatim}
??
[?]
\end{verbatim}

Before generating the final PDF, students shall check:

\begin{itemize}[leftmargin=0.8cm, itemsep=2pt]
    \item the Table of Contents;
    \item page numbering;
    \item chapter, section, and subsection numbering;
    \item figure numbering;
    \item table numbering;
    \item equation numbering;
    \item cross-references;
    \item citations;
    \item bibliography;
    \item font display;
    \item mathematical symbols;
    \item image quality;
    \item completeness of all required forms.
\end{itemize}

The document may be compiled using TeX Live, MiKTeX, or another
compatible \LaTeX{} environment. Regardless of the compilation
environment, the resulting PDF shall conform to the official Faculty
template.

\vspace{0.3cm}

% ================================================================
% XXII. CONTENT THAT STUDENTS MAY EDIT
% ================================================================

\noindent
\textbf{XXII.\quad CONTENT THAT STUDENTS MAY EDIT}

\vspace{0.15cm}

Students may edit or replace:

\begin{itemize}[leftmargin=0.8cm, itemsep=2pt]
    \item student name;
    \item class;
    \item internship organization;
    \item supervisor name;
    \item Abstract;
    \item List of Symbols and Abbreviations;
    \item chapter titles and chapter content;
    \item figures and photographs;
    \item tables;
    \item equations;
    \item citations and references.
\end{itemize}

Students shall not modify the main document formatting settings in the
preamble, including the document class, page geometry, line spacing,
heading formatting, table-of-contents formatting, font configuration,
page layout, or cover-page macros unless specifically instructed.

\vspace{0.3cm}

% ================================================================
% XXIII. FINAL SUBMISSION REQUIREMENTS
% ================================================================

\noindent
\textbf{XXIII.\quad FINAL SUBMISSION REQUIREMENTS}

\vspace{0.15cm}

Before submission, the final report shall:

\begin{itemize}[leftmargin=0.8cm, itemsep=2pt]
    \item compile without errors that affect the report content or presentation;
    \item contain all required forms and report sections;
    \item contain correct student, class, internship organization, and supervisor information;
    \item contain an updated Table of Contents, List of Figures, and List of Tables;
    \item contain no unused example content from the template;
    \item contain no placeholder text such as \textit{Title of Chapter 1},
    \textit{Section 1.1}, or other sample entries;
    \item contain all required citations and references;
    \item contain no unresolved cross-references.
\end{itemize}

If the Faculty requires submission of the \LaTeX{} source, the complete
project shall include all \verb|.tex| files, the bibliography database,
figures, and any other files required to reproduce the submitted PDF.

\vspace{0.4cm}

\noindent
\textbf{Important:}
Before submitting the final internship report, students shall remove or
comment out the following command in the main document:

\begin{verbatim}
\input{Regulations.tex}
\end{verbatim}

For example:

\begin{verbatim}
% \input{Regulations.tex}
\end{verbatim}

This removes the guidance pages from the final submitted report while
preserving the regulations file in the \LaTeX{} project for reference.

} % End 13pt content

\clearpage
 command before submitting
% the final internship report.
%
% =================================================================

\clearpage
\thispagestyle{empty}

\begin{center}
    {\fontsize{16pt}{20pt}\selectfont\bfseries
    SPECIFIC REGULATIONS FOR PREPARING\\
    GRADUATION INTERNSHIP REPORTS USING \LaTeX}
\end{center}

\vspace{0.5cm}

{\fontsize{13pt}{17pt}\selectfont
\setstretch{1.3}
\justifying

\noindent
\textbf{Note:}
This section is provided for guidance only and is included in the main
document using the command \verb|
% =================================================================
% REGULATIONS FOR PREPARING THE GRADUATION INTERNSHIP REPORT
% USING LATEX
% =================================================================
%
% This file is intended to be included from the main document with:
%
%     \input{Regulations.tex}
%
% IMPORTANT:
% This section is provided for guidance only. Students shall remove
% or comment out the \input{Regulations.tex} command before submitting
% the final internship report.
%
% =================================================================

\clearpage
\thispagestyle{empty}

\begin{center}
    {\fontsize{16pt}{20pt}\selectfont\bfseries
    SPECIFIC REGULATIONS FOR PREPARING\\
    GRADUATION INTERNSHIP REPORTS USING \LaTeX}
\end{center}

\vspace{0.5cm}

{\fontsize{13pt}{17pt}\selectfont
\setstretch{1.3}
\justifying

\noindent
\textbf{Note:}
This section is provided for guidance only and is included in the main
document using the command \verb|\input{Regulations.tex}|.
Students shall remove or comment out this command before submitting the
final internship report.

\vspace{0.4cm}

% ================================================================
% I. GENERAL REQUIREMENTS
% ================================================================

\noindent
\textbf{I.\quad GENERAL REQUIREMENTS}

\vspace{0.15cm}

The Graduation Internship Report shall be prepared using the official
\LaTeX{} template provided by the Faculty. Students shall preserve the
predefined page layout, typography, line spacing, heading styles,
numbering system, figure and table formatting, and bibliography settings
unless otherwise instructed by the Faculty or the supervisor.

The report shall be written in English.

Students are responsible for the academic and technical content of the
report. The official \LaTeX{} template is responsible for standardized
document formatting.

\vspace{0.3cm}

% ================================================================
% II. REPORT STRUCTURE
% ================================================================

\noindent
\textbf{II.\quad REPORT STRUCTURE}

\vspace{0.15cm}

The report shall normally be arranged in the following order:

\begin{enumerate}[label=\arabic*., leftmargin=1.2cm, itemsep=2pt]
    \item Outer Cover;
    \item Comments of the Internship Host Organization;
    \item Supervisor's Evaluation and Graduation Internship Defense Results;
    \item Table of Contents;
    \item List of Symbols and Abbreviations, if applicable;
    \item List of Figures;
    \item List of Tables;
    \item Abstract;
    \item Main Chapters of the Internship Report;
    \item References;
    \item Appendices, if applicable.
\end{enumerate}

The preliminary pages shall not be numbered as chapters. The main report
content shall be organized into consecutively numbered chapters beginning
with Chapter 1.

\vspace{0.3cm}

% ================================================================
% III. PAGE SIZE AND MARGINS
% ================================================================

\noindent
\textbf{III.\quad PAGE SIZE AND MARGINS}

\vspace{0.15cm}

The report shall be formatted on A4 paper with the following margins:

\begin{center}
\begin{tabular}{ll}
    Left margin:   & 3.5 cm \\
    Right margin:  & 2.0 cm \\
    Top margin:    & 2.5 cm \\
    Bottom margin: & 2.5 cm
\end{tabular}
\end{center}

These settings are predefined in the official template and shall not be
modified by students.

\vspace{0.3cm}

% ================================================================
% IV. FONT, FONT SIZE, AND LINE SPACING
% ================================================================

\noindent
\textbf{IV.\quad FONT, FONT SIZE, AND LINE SPACING}

\vspace{0.15cm}

The main text shall use a 13 pt font size with 1.3-line spacing.
A Times-family typeface shall be used in accordance with the official
template.

The text shall be fully justified.

Students shall not use repeated spaces, repeated manual line breaks, or
multiple blank lines to position text. Paragraph and heading spacing shall
be controlled by the template.

\vspace{0.3cm}

% ================================================================
% V. OUTER COVER
% ================================================================

\noindent
\textbf{V.\quad OUTER COVER}

\vspace{0.15cm}

The outer cover shall follow the official format provided in the template
and shall contain the following information:

\begin{itemize}[leftmargin=0.8cm, itemsep=2pt]
    \item Thai Nguyen University of Technology;
    \item Faculty of International Training;
    \item Engineering and Technology Department;
    \item University logo and Faculty logo;
    \item the title \textbf{INTERNSHIP REPORT};
    \item Internship Organization;
    \item Author;
    \item Class;
    \item Supervisor;
    \item place and year of completion.
\end{itemize}

Students shall only replace the required information fields and shall not
modify the cover-page layout unless instructed otherwise.

\vspace{0.3cm}

% ================================================================
% VI. COMMENTS OF THE INTERNSHIP HOST ORGANIZATION
% ================================================================

\noindent
\textbf{VI.\quad COMMENTS OF THE INTERNSHIP HOST ORGANIZATION}

\vspace{0.15cm}

The page entitled \textbf{COMMENTS OF THE INTERNSHIP HOST ORGANIZATION}
is reserved for the internship host organization.

It includes:

\begin{itemize}[leftmargin=0.8cm, itemsep=2pt]
    \item the student's full name;
    \item class;
    \item comments from the host organization;
    \item score;
    \item date;
    \item signature and official seal.
\end{itemize}

Students shall not alter or complete official evaluation fields unless
specifically authorized.

\vspace{0.3cm}

% ================================================================
% VII. SUPERVISOR'S EVALUATION AND DEFENSE RESULTS
% ================================================================

\noindent
\textbf{VII.\quad SUPERVISOR'S EVALUATION AND DEFENSE RESULTS}

\vspace{0.15cm}

The section entitled
\textbf{SUPERVISOR'S EVALUATION AND GRADUATION INTERNSHIP DEFENSE RESULTS}
is reserved for the supervisor and defense committee.

The evaluation may include:

\begin{itemize}[leftmargin=0.8cm, itemsep=2pt]
    \item Internship Progress and Student's Attitude;
    \item Report Content;
    \item Report Presentation;
    \item Other Comments;
    \item Supervisor's Opinion;
    \item Defense Result;
    \item Score.
\end{itemize}

Students shall not enter evaluation scores or defense results before they
are officially determined.

\vspace{0.3cm}

% ================================================================
% VIII. CHAPTERS, SECTIONS, AND SUBSECTIONS
% ================================================================

\noindent
\textbf{VIII.\quad CHAPTERS, SECTIONS, AND SUBSECTIONS}

\vspace{0.15cm}

The report structure shall be generated using the standard \LaTeX{}
commands:

\begin{verbatim}
\chapter{...}
\section{...}
\subsection{...}
\end{verbatim}

Chapter, section, and subsection numbers shall be generated automatically
by \LaTeX{}.

Students shall not manually type chapter or section numbers and shall not
reproduce heading styles using commands such as \verb|\textbf|,
\verb|\fontsize|, or \verb|\centering| instead of the appropriate
structural command.

Excessive subdivision should be avoided. In most cases, the hierarchy
Chapter--Section--Subsection is sufficient.

\vspace{0.3cm}

% ================================================================
% IX. ORGANIZATION OF THE MAIN CONTENT
% ================================================================

\noindent
\textbf{IX.\quad ORGANIZATION OF THE MAIN CONTENT}

\vspace{0.15cm}

The specific content of the report shall depend on the internship
organization and the assigned tasks. A typical structure may include:

\begin{itemize}[leftmargin=0.8cm, itemsep=2pt]
    \item \textbf{Chapter 1. Overview of the Internship Organization};
    \item \textbf{Chapter 2. Internship Activities and Assigned Tasks};
    \item \textbf{Chapter 3. Results, Experience, and Evaluation}.
\end{itemize}

These chapter titles are recommendations only and may be adapted to the
actual internship content.

The report shall reflect work actually performed, observed, analyzed, or
learned during the internship. Unrelated theoretical material should not
be included merely to increase report length.

\vspace{0.3cm}

% ================================================================
% X. TABLE OF CONTENTS
% ================================================================

\noindent
\textbf{X.\quad TABLE OF CONTENTS}

\vspace{0.15cm}

The Table of Contents shall be generated automatically using:

\begin{verbatim}
\tableofcontents
\end{verbatim}

Students shall not create or manually edit page numbers in the Table of
Contents.

The document shall be compiled sufficiently to ensure that all page
numbers and entries are correctly updated.

\vspace{0.3cm}

% ================================================================
% XI. LIST OF SYMBOLS AND ABBREVIATIONS
% ================================================================

\noindent
\textbf{XI.\quad LIST OF SYMBOLS AND ABBREVIATIONS}

\vspace{0.15cm}

Symbols and abbreviations that are used repeatedly in the report shall be
listed in the \textbf{List of Symbols and Abbreviations}.

The list may contain the following columns:

\begin{center}
\begin{tabular}{llll}
    No. & Symbol / Abbreviation & Unit & Definition
\end{tabular}
\end{center}

Only symbols and abbreviations actually used in the final report shall be
retained. Example entries contained in the template shall be removed if
they are not relevant to the student's report.

Each abbreviation should normally be defined when it first appears in
the main text.

\vspace{0.3cm}

% ================================================================
% XII. FIGURES, DIAGRAMS, AND PHOTOGRAPHS
% ================================================================

\noindent
\textbf{XII.\quad FIGURES, DIAGRAMS, AND PHOTOGRAPHS}

\vspace{0.15cm}

Figures shall be clear, legible, and directly relevant to the report
content.

Figures shall be inserted using the standard \verb|figure| environment.

Example:

\begin{verbatim}
\begin{figure}[H]
    \centering
    \includegraphics[width=0.8\textwidth]
                    {figures/company_structure.pdf}
    \caption{Organizational structure of the company}
    \label{fig:company-structure}
\end{figure}
\end{verbatim}

Figure captions shall be placed below the figure.

Each referenced figure shall have a unique \verb|\label| and shall be
referenced in the text using, for example:

\begin{verbatim}
Figure~\ref{fig:company-structure}
\end{verbatim}

Students shall not manually type figure numbers.

Figures reproduced or adapted from external sources shall include an
appropriate citation.

Photographs taken during the internship shall be relevant and sufficiently
clear. Students shall not include confidential, restricted, or proprietary
information without permission from the internship organization.

Vector graphics or high-resolution images are recommended whenever
possible.

\vspace{0.3cm}

% ================================================================
% XIII. TABLES
% ================================================================

\noindent
\textbf{XIII.\quad TABLES}

\vspace{0.15cm}

Tables shall be created using suitable \LaTeX{} environments such as
\verb|table|, \verb|tabular|, \verb|tabularx|, or \verb|longtable|.

Table captions shall be placed above the table.

Example:

\begin{verbatim}
\begin{table}[H]
    \centering
    \caption{Internship schedule}
    \label{tab:internship-schedule}

    \begin{tabular}{lll}
        \toprule
        Week & Task & Result \\
        \midrule
        1 & Orientation & Completed \\
        2 & System study & Completed \\
        \bottomrule
    \end{tabular}
\end{table}
\end{verbatim}

Tables shall be referenced using:

\begin{verbatim}
Table~\ref{tab:internship-schedule}
\end{verbatim}

Students shall not manually type table numbers.

Tables should be created directly in \LaTeX{} rather than inserted as
screenshots whenever practical.

\vspace{0.3cm}

% ================================================================
% XIV. EQUATIONS AND TECHNICAL QUANTITIES
% ================================================================

\noindent
\textbf{XIV.\quad EQUATIONS AND TECHNICAL QUANTITIES}

\vspace{0.15cm}

Mathematical expressions shall be typeset directly using \LaTeX{}.
Images or screenshots of equations shall not be used.

Short expressions may be written in inline mathematical mode, for example:

\begin{verbatim}
$R = 0.7\,\Omega$
\end{verbatim}

Important equations shall use a numbered equation environment.

Example:

\begin{verbatim}
\begin{equation}
    P = UI
    \label{eq:power}
\end{equation}
\end{verbatim}

Equations shall be referenced using:

\begin{verbatim}
Equation~\eqref{eq:power}
\end{verbatim}

Equation numbers shall not be manually entered.

SI units shall be used wherever applicable. The \verb|siunitx| package
should be used where appropriate.

Examples:

\begin{verbatim}
\SI{220}{\volt}
\SI{50}{\hertz}
\SI{2.5}{\ampere}
\end{verbatim}

Notation and units shall be consistent throughout the report.

\vspace{0.3cm}

% ================================================================
% XV. SOURCE CODE
% ================================================================

\noindent
\textbf{XV.\quad SOURCE CODE}

\vspace{0.15cm}

If the internship involves programming, PLCs, microcontrollers,
simulation, or software development, relevant source code may be included
in the report.

Source code shall be presented using the \verb|listings| environment or
another environment defined by the official template.

Example:

\begin{verbatim}
\begin{lstlisting}[language=st]
IF Start = TRUE THEN
    Motor := TRUE;
ELSE
    Motor := FALSE;
END_IF;
\end{lstlisting}
\end{verbatim}

Only code necessary for explaining the work should appear in the main
report. Long source-code listings should be moved to an appendix or
submitted as supplementary files if required.

Source code should not normally be inserted as screenshots.

\vspace{0.3cm}

% ================================================================
% XVI. ABSTRACT
% ================================================================

\noindent
\textbf{XVI.\quad ABSTRACT}

\vspace{0.15cm}

The Abstract shall briefly summarize:

\begin{itemize}[leftmargin=0.8cm, itemsep=2pt]
    \item the internship organization;
    \item the objectives of the internship;
    \item the main tasks performed;
    \item the main results, knowledge, and skills obtained.
\end{itemize}

The Abstract should be concise and should not reproduce the full content
of the report chapter by chapter.

\vspace{0.3cm}

% ================================================================
% XVII. REFERENCES AND CITATIONS
% ================================================================

\noindent
\textbf{XVII.\quad REFERENCES AND CITATIONS}

\vspace{0.15cm}

All books, journal papers, conference papers, standards, websites,
datasheets, manuals, technical documents, figures, tables, or other
external sources used in preparing the report shall be properly cited.

References shall be managed through the bibliography system provided with
the official template.

Bibliographic entries shall normally be stored in:

\begin{verbatim}
References.bib
\end{verbatim}

References shall be cited in the report using:

\begin{verbatim}
\cite{reference-key}
\end{verbatim}

Students shall not manually type reference numbers such as ``[1]'',
``[2]'', etc.

Titles of foreign-language publications shall remain in their original
language in the bibliography.

All borrowed ideas, data, figures, tables, equations, methods, or results
shall be properly acknowledged.

\vspace{0.3cm}

% ================================================================
% XVIII. CROSS-REFERENCING
% ================================================================

\noindent
\textbf{XVIII.\quad CROSS-REFERENCING}

\vspace{0.15cm}

Numbered elements shall be cross-referenced using the standard \LaTeX{}
mechanism.

Use the following commands where appropriate:

\begin{verbatim}
\label{...}
\ref{...}
\eqref{...}
\end{verbatim}

For example:

\begin{verbatim}
As shown in Figure~\ref{fig:system}, ...
\end{verbatim}

Students shall not manually type numbers such as ``Figure 2.3'' when
those numbers can be generated automatically.

\vspace{0.3cm}

% ================================================================
% XIX. PAGE NUMBERING
% ================================================================

\noindent
\textbf{XIX.\quad PAGE NUMBERING}

\vspace{0.15cm}

Cover pages and official evaluation forms at the beginning of the report
shall not display page numbers.

The main body of the report shall use Arabic page numbers beginning from
Chapter 1.

Page numbering shall be generated automatically by the official template.
Students shall not manually enter page numbers.

\vspace{0.3cm}

% ================================================================
% XX. PROJECT FILE ORGANIZATION
% ================================================================

\noindent
\textbf{XX.\quad PROJECT FILE ORGANIZATION}

\vspace{0.15cm}

Students should maintain a clear and organized \LaTeX{} project
structure. A recommended structure is:

\begin{verbatim}
InternshipReport/
|
|-- main.tex
|-- References.bib
|
|-- figures/
|   |-- company.pdf
|   |-- system.png
|   `-- ...
|
`-- chapters/
    |-- chapter1.tex
    |-- chapter2.tex
    `-- chapter3.tex
\end{verbatim}

Figures should normally be stored in the \verb|figures| directory
specified by the official template.

File and directory names should preferably contain Latin letters,
numbers, hyphens, or underscores.

Recommended examples include:

\begin{verbatim}
company_structure.pdf
motor_control.png
chapter_01.tex
\end{verbatim}

Spaces and unnecessary special characters in file names should be
avoided where possible.

\vspace{0.3cm}

% ================================================================
% XXI. DOCUMENT COMPILATION AND FINAL CHECK
% ================================================================

\noindent
\textbf{XXI.\quad DOCUMENT COMPILATION AND FINAL CHECK}

\vspace{0.15cm}

Before submission, the complete \LaTeX{} project shall compile
successfully and produce the final PDF report.

The final document shall not contain unresolved references such as:

\begin{verbatim}
??
[?]
\end{verbatim}

Before generating the final PDF, students shall check:

\begin{itemize}[leftmargin=0.8cm, itemsep=2pt]
    \item the Table of Contents;
    \item page numbering;
    \item chapter, section, and subsection numbering;
    \item figure numbering;
    \item table numbering;
    \item equation numbering;
    \item cross-references;
    \item citations;
    \item bibliography;
    \item font display;
    \item mathematical symbols;
    \item image quality;
    \item completeness of all required forms.
\end{itemize}

The document may be compiled using TeX Live, MiKTeX, or another
compatible \LaTeX{} environment. Regardless of the compilation
environment, the resulting PDF shall conform to the official Faculty
template.

\vspace{0.3cm}

% ================================================================
% XXII. CONTENT THAT STUDENTS MAY EDIT
% ================================================================

\noindent
\textbf{XXII.\quad CONTENT THAT STUDENTS MAY EDIT}

\vspace{0.15cm}

Students may edit or replace:

\begin{itemize}[leftmargin=0.8cm, itemsep=2pt]
    \item student name;
    \item class;
    \item internship organization;
    \item supervisor name;
    \item Abstract;
    \item List of Symbols and Abbreviations;
    \item chapter titles and chapter content;
    \item figures and photographs;
    \item tables;
    \item equations;
    \item citations and references.
\end{itemize}

Students shall not modify the main document formatting settings in the
preamble, including the document class, page geometry, line spacing,
heading formatting, table-of-contents formatting, font configuration,
page layout, or cover-page macros unless specifically instructed.

\vspace{0.3cm}

% ================================================================
% XXIII. FINAL SUBMISSION REQUIREMENTS
% ================================================================

\noindent
\textbf{XXIII.\quad FINAL SUBMISSION REQUIREMENTS}

\vspace{0.15cm}

Before submission, the final report shall:

\begin{itemize}[leftmargin=0.8cm, itemsep=2pt]
    \item compile without errors that affect the report content or presentation;
    \item contain all required forms and report sections;
    \item contain correct student, class, internship organization, and supervisor information;
    \item contain an updated Table of Contents, List of Figures, and List of Tables;
    \item contain no unused example content from the template;
    \item contain no placeholder text such as \textit{Title of Chapter 1},
    \textit{Section 1.1}, or other sample entries;
    \item contain all required citations and references;
    \item contain no unresolved cross-references.
\end{itemize}

If the Faculty requires submission of the \LaTeX{} source, the complete
project shall include all \verb|.tex| files, the bibliography database,
figures, and any other files required to reproduce the submitted PDF.

\vspace{0.4cm}

\noindent
\textbf{Important:}
Before submitting the final internship report, students shall remove or
comment out the following command in the main document:

\begin{verbatim}
\input{Regulations.tex}
\end{verbatim}

For example:

\begin{verbatim}
% \input{Regulations.tex}
\end{verbatim}

This removes the guidance pages from the final submitted report while
preserving the regulations file in the \LaTeX{} project for reference.

} % End 13pt content

\clearpage
|.
Students shall remove or comment out this command before submitting the
final internship report.

\vspace{0.4cm}

% ================================================================
% I. GENERAL REQUIREMENTS
% ================================================================

\noindent
\textbf{I.\quad GENERAL REQUIREMENTS}

\vspace{0.15cm}

The Graduation Internship Report shall be prepared using the official
\LaTeX{} template provided by the Faculty. Students shall preserve the
predefined page layout, typography, line spacing, heading styles,
numbering system, figure and table formatting, and bibliography settings
unless otherwise instructed by the Faculty or the supervisor.

The report shall be written in English.

Students are responsible for the academic and technical content of the
report. The official \LaTeX{} template is responsible for standardized
document formatting.

\vspace{0.3cm}

% ================================================================
% II. REPORT STRUCTURE
% ================================================================

\noindent
\textbf{II.\quad REPORT STRUCTURE}

\vspace{0.15cm}

The report shall normally be arranged in the following order:

\begin{enumerate}[label=\arabic*., leftmargin=1.2cm, itemsep=2pt]
    \item Outer Cover;
    \item Comments of the Internship Host Organization;
    \item Supervisor's Evaluation and Graduation Internship Defense Results;
    \item Table of Contents;
    \item List of Symbols and Abbreviations, if applicable;
    \item List of Figures;
    \item List of Tables;
    \item Abstract;
    \item Main Chapters of the Internship Report;
    \item References;
    \item Appendices, if applicable.
\end{enumerate}

The preliminary pages shall not be numbered as chapters. The main report
content shall be organized into consecutively numbered chapters beginning
with Chapter 1.

\vspace{0.3cm}

% ================================================================
% III. PAGE SIZE AND MARGINS
% ================================================================

\noindent
\textbf{III.\quad PAGE SIZE AND MARGINS}

\vspace{0.15cm}

The report shall be formatted on A4 paper with the following margins:

\begin{center}
\begin{tabular}{ll}
    Left margin:   & 3.5 cm \\
    Right margin:  & 2.0 cm \\
    Top margin:    & 2.5 cm \\
    Bottom margin: & 2.5 cm
\end{tabular}
\end{center}

These settings are predefined in the official template and shall not be
modified by students.

\vspace{0.3cm}

% ================================================================
% IV. FONT, FONT SIZE, AND LINE SPACING
% ================================================================

\noindent
\textbf{IV.\quad FONT, FONT SIZE, AND LINE SPACING}

\vspace{0.15cm}

The main text shall use a 13 pt font size with 1.3-line spacing.
A Times-family typeface shall be used in accordance with the official
template.

The text shall be fully justified.

Students shall not use repeated spaces, repeated manual line breaks, or
multiple blank lines to position text. Paragraph and heading spacing shall
be controlled by the template.

\vspace{0.3cm}

% ================================================================
% V. OUTER COVER
% ================================================================

\noindent
\textbf{V.\quad OUTER COVER}

\vspace{0.15cm}

The outer cover shall follow the official format provided in the template
and shall contain the following information:

\begin{itemize}[leftmargin=0.8cm, itemsep=2pt]
    \item Thai Nguyen University of Technology;
    \item Faculty of International Training;
    \item Engineering and Technology Department;
    \item University logo and Faculty logo;
    \item the title \textbf{INTERNSHIP REPORT};
    \item Internship Organization;
    \item Author;
    \item Class;
    \item Supervisor;
    \item place and year of completion.
\end{itemize}

Students shall only replace the required information fields and shall not
modify the cover-page layout unless instructed otherwise.

\vspace{0.3cm}

% ================================================================
% VI. COMMENTS OF THE INTERNSHIP HOST ORGANIZATION
% ================================================================

\noindent
\textbf{VI.\quad COMMENTS OF THE INTERNSHIP HOST ORGANIZATION}

\vspace{0.15cm}

The page entitled \textbf{COMMENTS OF THE INTERNSHIP HOST ORGANIZATION}
is reserved for the internship host organization.

It includes:

\begin{itemize}[leftmargin=0.8cm, itemsep=2pt]
    \item the student's full name;
    \item class;
    \item comments from the host organization;
    \item score;
    \item date;
    \item signature and official seal.
\end{itemize}

Students shall not alter or complete official evaluation fields unless
specifically authorized.

\vspace{0.3cm}

% ================================================================
% VII. SUPERVISOR'S EVALUATION AND DEFENSE RESULTS
% ================================================================

\noindent
\textbf{VII.\quad SUPERVISOR'S EVALUATION AND DEFENSE RESULTS}

\vspace{0.15cm}

The section entitled
\textbf{SUPERVISOR'S EVALUATION AND GRADUATION INTERNSHIP DEFENSE RESULTS}
is reserved for the supervisor and defense committee.

The evaluation may include:

\begin{itemize}[leftmargin=0.8cm, itemsep=2pt]
    \item Internship Progress and Student's Attitude;
    \item Report Content;
    \item Report Presentation;
    \item Other Comments;
    \item Supervisor's Opinion;
    \item Defense Result;
    \item Score.
\end{itemize}

Students shall not enter evaluation scores or defense results before they
are officially determined.

\vspace{0.3cm}

% ================================================================
% VIII. CHAPTERS, SECTIONS, AND SUBSECTIONS
% ================================================================

\noindent
\textbf{VIII.\quad CHAPTERS, SECTIONS, AND SUBSECTIONS}

\vspace{0.15cm}

The report structure shall be generated using the standard \LaTeX{}
commands:

\begin{verbatim}
\chapter{...}
\section{...}
\subsection{...}
\end{verbatim}

Chapter, section, and subsection numbers shall be generated automatically
by \LaTeX{}.

Students shall not manually type chapter or section numbers and shall not
reproduce heading styles using commands such as \verb|\textbf|,
\verb|\fontsize|, or \verb|\centering| instead of the appropriate
structural command.

Excessive subdivision should be avoided. In most cases, the hierarchy
Chapter--Section--Subsection is sufficient.

\vspace{0.3cm}

% ================================================================
% IX. ORGANIZATION OF THE MAIN CONTENT
% ================================================================

\noindent
\textbf{IX.\quad ORGANIZATION OF THE MAIN CONTENT}

\vspace{0.15cm}

The specific content of the report shall depend on the internship
organization and the assigned tasks. A typical structure may include:

\begin{itemize}[leftmargin=0.8cm, itemsep=2pt]
    \item \textbf{Chapter 1. Overview of the Internship Organization};
    \item \textbf{Chapter 2. Internship Activities and Assigned Tasks};
    \item \textbf{Chapter 3. Results, Experience, and Evaluation}.
\end{itemize}

These chapter titles are recommendations only and may be adapted to the
actual internship content.

The report shall reflect work actually performed, observed, analyzed, or
learned during the internship. Unrelated theoretical material should not
be included merely to increase report length.

\vspace{0.3cm}

% ================================================================
% X. TABLE OF CONTENTS
% ================================================================

\noindent
\textbf{X.\quad TABLE OF CONTENTS}

\vspace{0.15cm}

The Table of Contents shall be generated automatically using:

\begin{verbatim}
\tableofcontents
\end{verbatim}

Students shall not create or manually edit page numbers in the Table of
Contents.

The document shall be compiled sufficiently to ensure that all page
numbers and entries are correctly updated.

\vspace{0.3cm}

% ================================================================
% XI. LIST OF SYMBOLS AND ABBREVIATIONS
% ================================================================

\noindent
\textbf{XI.\quad LIST OF SYMBOLS AND ABBREVIATIONS}

\vspace{0.15cm}

Symbols and abbreviations that are used repeatedly in the report shall be
listed in the \textbf{List of Symbols and Abbreviations}.

The list may contain the following columns:

\begin{center}
\begin{tabular}{llll}
    No. & Symbol / Abbreviation & Unit & Definition
\end{tabular}
\end{center}

Only symbols and abbreviations actually used in the final report shall be
retained. Example entries contained in the template shall be removed if
they are not relevant to the student's report.

Each abbreviation should normally be defined when it first appears in
the main text.

\vspace{0.3cm}

% ================================================================
% XII. FIGURES, DIAGRAMS, AND PHOTOGRAPHS
% ================================================================

\noindent
\textbf{XII.\quad FIGURES, DIAGRAMS, AND PHOTOGRAPHS}

\vspace{0.15cm}

Figures shall be clear, legible, and directly relevant to the report
content.

Figures shall be inserted using the standard \verb|figure| environment.

Example:

\begin{verbatim}
\begin{figure}[H]
    \centering
    \includegraphics[width=0.8\textwidth]
                    {figures/company_structure.pdf}
    \caption{Organizational structure of the company}
    \label{fig:company-structure}
\end{figure}
\end{verbatim}

Figure captions shall be placed below the figure.

Each referenced figure shall have a unique \verb|\label| and shall be
referenced in the text using, for example:

\begin{verbatim}
Figure~\ref{fig:company-structure}
\end{verbatim}

Students shall not manually type figure numbers.

Figures reproduced or adapted from external sources shall include an
appropriate citation.

Photographs taken during the internship shall be relevant and sufficiently
clear. Students shall not include confidential, restricted, or proprietary
information without permission from the internship organization.

Vector graphics or high-resolution images are recommended whenever
possible.

\vspace{0.3cm}

% ================================================================
% XIII. TABLES
% ================================================================

\noindent
\textbf{XIII.\quad TABLES}

\vspace{0.15cm}

Tables shall be created using suitable \LaTeX{} environments such as
\verb|table|, \verb|tabular|, \verb|tabularx|, or \verb|longtable|.

Table captions shall be placed above the table.

Example:

\begin{verbatim}
\begin{table}[H]
    \centering
    \caption{Internship schedule}
    \label{tab:internship-schedule}

    \begin{tabular}{lll}
        \toprule
        Week & Task & Result \\
        \midrule
        1 & Orientation & Completed \\
        2 & System study & Completed \\
        \bottomrule
    \end{tabular}
\end{table}
\end{verbatim}

Tables shall be referenced using:

\begin{verbatim}
Table~\ref{tab:internship-schedule}
\end{verbatim}

Students shall not manually type table numbers.

Tables should be created directly in \LaTeX{} rather than inserted as
screenshots whenever practical.

\vspace{0.3cm}

% ================================================================
% XIV. EQUATIONS AND TECHNICAL QUANTITIES
% ================================================================

\noindent
\textbf{XIV.\quad EQUATIONS AND TECHNICAL QUANTITIES}

\vspace{0.15cm}

Mathematical expressions shall be typeset directly using \LaTeX{}.
Images or screenshots of equations shall not be used.

Short expressions may be written in inline mathematical mode, for example:

\begin{verbatim}
$R = 0.7\,\Omega$
\end{verbatim}

Important equations shall use a numbered equation environment.

Example:

\begin{verbatim}
\begin{equation}
    P = UI
    \label{eq:power}
\end{equation}
\end{verbatim}

Equations shall be referenced using:

\begin{verbatim}
Equation~\eqref{eq:power}
\end{verbatim}

Equation numbers shall not be manually entered.

SI units shall be used wherever applicable. The \verb|siunitx| package
should be used where appropriate.

Examples:

\begin{verbatim}
\SI{220}{\volt}
\SI{50}{\hertz}
\SI{2.5}{\ampere}
\end{verbatim}

Notation and units shall be consistent throughout the report.

\vspace{0.3cm}

% ================================================================
% XV. SOURCE CODE
% ================================================================

\noindent
\textbf{XV.\quad SOURCE CODE}

\vspace{0.15cm}

If the internship involves programming, PLCs, microcontrollers,
simulation, or software development, relevant source code may be included
in the report.

Source code shall be presented using the \verb|listings| environment or
another environment defined by the official template.

Example:

\begin{verbatim}
\begin{lstlisting}[language=st]
IF Start = TRUE THEN
    Motor := TRUE;
ELSE
    Motor := FALSE;
END_IF;
\end{lstlisting}
\end{verbatim}

Only code necessary for explaining the work should appear in the main
report. Long source-code listings should be moved to an appendix or
submitted as supplementary files if required.

Source code should not normally be inserted as screenshots.

\vspace{0.3cm}

% ================================================================
% XVI. ABSTRACT
% ================================================================

\noindent
\textbf{XVI.\quad ABSTRACT}

\vspace{0.15cm}

The Abstract shall briefly summarize:

\begin{itemize}[leftmargin=0.8cm, itemsep=2pt]
    \item the internship organization;
    \item the objectives of the internship;
    \item the main tasks performed;
    \item the main results, knowledge, and skills obtained.
\end{itemize}

The Abstract should be concise and should not reproduce the full content
of the report chapter by chapter.

\vspace{0.3cm}

% ================================================================
% XVII. REFERENCES AND CITATIONS
% ================================================================

\noindent
\textbf{XVII.\quad REFERENCES AND CITATIONS}

\vspace{0.15cm}

All books, journal papers, conference papers, standards, websites,
datasheets, manuals, technical documents, figures, tables, or other
external sources used in preparing the report shall be properly cited.

References shall be managed through the bibliography system provided with
the official template.

Bibliographic entries shall normally be stored in:

\begin{verbatim}
References.bib
\end{verbatim}

References shall be cited in the report using:

\begin{verbatim}
\cite{reference-key}
\end{verbatim}

Students shall not manually type reference numbers such as ``[1]'',
``[2]'', etc.

Titles of foreign-language publications shall remain in their original
language in the bibliography.

All borrowed ideas, data, figures, tables, equations, methods, or results
shall be properly acknowledged.

\vspace{0.3cm}

% ================================================================
% XVIII. CROSS-REFERENCING
% ================================================================

\noindent
\textbf{XVIII.\quad CROSS-REFERENCING}

\vspace{0.15cm}

Numbered elements shall be cross-referenced using the standard \LaTeX{}
mechanism.

Use the following commands where appropriate:

\begin{verbatim}
\label{...}
\ref{...}
\eqref{...}
\end{verbatim}

For example:

\begin{verbatim}
As shown in Figure~\ref{fig:system}, ...
\end{verbatim}

Students shall not manually type numbers such as ``Figure 2.3'' when
those numbers can be generated automatically.

\vspace{0.3cm}

% ================================================================
% XIX. PAGE NUMBERING
% ================================================================

\noindent
\textbf{XIX.\quad PAGE NUMBERING}

\vspace{0.15cm}

Cover pages and official evaluation forms at the beginning of the report
shall not display page numbers.

The main body of the report shall use Arabic page numbers beginning from
Chapter 1.

Page numbering shall be generated automatically by the official template.
Students shall not manually enter page numbers.

\vspace{0.3cm}

% ================================================================
% XX. PROJECT FILE ORGANIZATION
% ================================================================

\noindent
\textbf{XX.\quad PROJECT FILE ORGANIZATION}

\vspace{0.15cm}

Students should maintain a clear and organized \LaTeX{} project
structure. A recommended structure is:

\begin{verbatim}
InternshipReport/
|
|-- main.tex
|-- References.bib
|
|-- figures/
|   |-- company.pdf
|   |-- system.png
|   `-- ...
|
`-- chapters/
    |-- chapter1.tex
    |-- chapter2.tex
    `-- chapter3.tex
\end{verbatim}

Figures should normally be stored in the \verb|figures| directory
specified by the official template.

File and directory names should preferably contain Latin letters,
numbers, hyphens, or underscores.

Recommended examples include:

\begin{verbatim}
company_structure.pdf
motor_control.png
chapter_01.tex
\end{verbatim}

Spaces and unnecessary special characters in file names should be
avoided where possible.

\vspace{0.3cm}

% ================================================================
% XXI. DOCUMENT COMPILATION AND FINAL CHECK
% ================================================================

\noindent
\textbf{XXI.\quad DOCUMENT COMPILATION AND FINAL CHECK}

\vspace{0.15cm}

Before submission, the complete \LaTeX{} project shall compile
successfully and produce the final PDF report.

The final document shall not contain unresolved references such as:

\begin{verbatim}
??
[?]
\end{verbatim}

Before generating the final PDF, students shall check:

\begin{itemize}[leftmargin=0.8cm, itemsep=2pt]
    \item the Table of Contents;
    \item page numbering;
    \item chapter, section, and subsection numbering;
    \item figure numbering;
    \item table numbering;
    \item equation numbering;
    \item cross-references;
    \item citations;
    \item bibliography;
    \item font display;
    \item mathematical symbols;
    \item image quality;
    \item completeness of all required forms.
\end{itemize}

The document may be compiled using TeX Live, MiKTeX, or another
compatible \LaTeX{} environment. Regardless of the compilation
environment, the resulting PDF shall conform to the official Faculty
template.

\vspace{0.3cm}

% ================================================================
% XXII. CONTENT THAT STUDENTS MAY EDIT
% ================================================================

\noindent
\textbf{XXII.\quad CONTENT THAT STUDENTS MAY EDIT}

\vspace{0.15cm}

Students may edit or replace:

\begin{itemize}[leftmargin=0.8cm, itemsep=2pt]
    \item student name;
    \item class;
    \item internship organization;
    \item supervisor name;
    \item Abstract;
    \item List of Symbols and Abbreviations;
    \item chapter titles and chapter content;
    \item figures and photographs;
    \item tables;
    \item equations;
    \item citations and references.
\end{itemize}

Students shall not modify the main document formatting settings in the
preamble, including the document class, page geometry, line spacing,
heading formatting, table-of-contents formatting, font configuration,
page layout, or cover-page macros unless specifically instructed.

\vspace{0.3cm}

% ================================================================
% XXIII. FINAL SUBMISSION REQUIREMENTS
% ================================================================

\noindent
\textbf{XXIII.\quad FINAL SUBMISSION REQUIREMENTS}

\vspace{0.15cm}

Before submission, the final report shall:

\begin{itemize}[leftmargin=0.8cm, itemsep=2pt]
    \item compile without errors that affect the report content or presentation;
    \item contain all required forms and report sections;
    \item contain correct student, class, internship organization, and supervisor information;
    \item contain an updated Table of Contents, List of Figures, and List of Tables;
    \item contain no unused example content from the template;
    \item contain no placeholder text such as \textit{Title of Chapter 1},
    \textit{Section 1.1}, or other sample entries;
    \item contain all required citations and references;
    \item contain no unresolved cross-references.
\end{itemize}

If the Faculty requires submission of the \LaTeX{} source, the complete
project shall include all \verb|.tex| files, the bibliography database,
figures, and any other files required to reproduce the submitted PDF.

\vspace{0.4cm}

\noindent
\textbf{Important:}
Before submitting the final internship report, students shall remove or
comment out the following command in the main document:

\begin{verbatim}

% =================================================================
% REGULATIONS FOR PREPARING THE GRADUATION INTERNSHIP REPORT
% USING LATEX
% =================================================================
%
% This file is intended to be included from the main document with:
%
%     \input{Regulations.tex}
%
% IMPORTANT:
% This section is provided for guidance only. Students shall remove
% or comment out the \input{Regulations.tex} command before submitting
% the final internship report.
%
% =================================================================

\clearpage
\thispagestyle{empty}

\begin{center}
    {\fontsize{16pt}{20pt}\selectfont\bfseries
    SPECIFIC REGULATIONS FOR PREPARING\\
    GRADUATION INTERNSHIP REPORTS USING \LaTeX}
\end{center}

\vspace{0.5cm}

{\fontsize{13pt}{17pt}\selectfont
\setstretch{1.3}
\justifying

\noindent
\textbf{Note:}
This section is provided for guidance only and is included in the main
document using the command \verb|\input{Regulations.tex}|.
Students shall remove or comment out this command before submitting the
final internship report.

\vspace{0.4cm}

% ================================================================
% I. GENERAL REQUIREMENTS
% ================================================================

\noindent
\textbf{I.\quad GENERAL REQUIREMENTS}

\vspace{0.15cm}

The Graduation Internship Report shall be prepared using the official
\LaTeX{} template provided by the Faculty. Students shall preserve the
predefined page layout, typography, line spacing, heading styles,
numbering system, figure and table formatting, and bibliography settings
unless otherwise instructed by the Faculty or the supervisor.

The report shall be written in English.

Students are responsible for the academic and technical content of the
report. The official \LaTeX{} template is responsible for standardized
document formatting.

\vspace{0.3cm}

% ================================================================
% II. REPORT STRUCTURE
% ================================================================

\noindent
\textbf{II.\quad REPORT STRUCTURE}

\vspace{0.15cm}

The report shall normally be arranged in the following order:

\begin{enumerate}[label=\arabic*., leftmargin=1.2cm, itemsep=2pt]
    \item Outer Cover;
    \item Comments of the Internship Host Organization;
    \item Supervisor's Evaluation and Graduation Internship Defense Results;
    \item Table of Contents;
    \item List of Symbols and Abbreviations, if applicable;
    \item List of Figures;
    \item List of Tables;
    \item Abstract;
    \item Main Chapters of the Internship Report;
    \item References;
    \item Appendices, if applicable.
\end{enumerate}

The preliminary pages shall not be numbered as chapters. The main report
content shall be organized into consecutively numbered chapters beginning
with Chapter 1.

\vspace{0.3cm}

% ================================================================
% III. PAGE SIZE AND MARGINS
% ================================================================

\noindent
\textbf{III.\quad PAGE SIZE AND MARGINS}

\vspace{0.15cm}

The report shall be formatted on A4 paper with the following margins:

\begin{center}
\begin{tabular}{ll}
    Left margin:   & 3.5 cm \\
    Right margin:  & 2.0 cm \\
    Top margin:    & 2.5 cm \\
    Bottom margin: & 2.5 cm
\end{tabular}
\end{center}

These settings are predefined in the official template and shall not be
modified by students.

\vspace{0.3cm}

% ================================================================
% IV. FONT, FONT SIZE, AND LINE SPACING
% ================================================================

\noindent
\textbf{IV.\quad FONT, FONT SIZE, AND LINE SPACING}

\vspace{0.15cm}

The main text shall use a 13 pt font size with 1.3-line spacing.
A Times-family typeface shall be used in accordance with the official
template.

The text shall be fully justified.

Students shall not use repeated spaces, repeated manual line breaks, or
multiple blank lines to position text. Paragraph and heading spacing shall
be controlled by the template.

\vspace{0.3cm}

% ================================================================
% V. OUTER COVER
% ================================================================

\noindent
\textbf{V.\quad OUTER COVER}

\vspace{0.15cm}

The outer cover shall follow the official format provided in the template
and shall contain the following information:

\begin{itemize}[leftmargin=0.8cm, itemsep=2pt]
    \item Thai Nguyen University of Technology;
    \item Faculty of International Training;
    \item Engineering and Technology Department;
    \item University logo and Faculty logo;
    \item the title \textbf{INTERNSHIP REPORT};
    \item Internship Organization;
    \item Author;
    \item Class;
    \item Supervisor;
    \item place and year of completion.
\end{itemize}

Students shall only replace the required information fields and shall not
modify the cover-page layout unless instructed otherwise.

\vspace{0.3cm}

% ================================================================
% VI. COMMENTS OF THE INTERNSHIP HOST ORGANIZATION
% ================================================================

\noindent
\textbf{VI.\quad COMMENTS OF THE INTERNSHIP HOST ORGANIZATION}

\vspace{0.15cm}

The page entitled \textbf{COMMENTS OF THE INTERNSHIP HOST ORGANIZATION}
is reserved for the internship host organization.

It includes:

\begin{itemize}[leftmargin=0.8cm, itemsep=2pt]
    \item the student's full name;
    \item class;
    \item comments from the host organization;
    \item score;
    \item date;
    \item signature and official seal.
\end{itemize}

Students shall not alter or complete official evaluation fields unless
specifically authorized.

\vspace{0.3cm}

% ================================================================
% VII. SUPERVISOR'S EVALUATION AND DEFENSE RESULTS
% ================================================================

\noindent
\textbf{VII.\quad SUPERVISOR'S EVALUATION AND DEFENSE RESULTS}

\vspace{0.15cm}

The section entitled
\textbf{SUPERVISOR'S EVALUATION AND GRADUATION INTERNSHIP DEFENSE RESULTS}
is reserved for the supervisor and defense committee.

The evaluation may include:

\begin{itemize}[leftmargin=0.8cm, itemsep=2pt]
    \item Internship Progress and Student's Attitude;
    \item Report Content;
    \item Report Presentation;
    \item Other Comments;
    \item Supervisor's Opinion;
    \item Defense Result;
    \item Score.
\end{itemize}

Students shall not enter evaluation scores or defense results before they
are officially determined.

\vspace{0.3cm}

% ================================================================
% VIII. CHAPTERS, SECTIONS, AND SUBSECTIONS
% ================================================================

\noindent
\textbf{VIII.\quad CHAPTERS, SECTIONS, AND SUBSECTIONS}

\vspace{0.15cm}

The report structure shall be generated using the standard \LaTeX{}
commands:

\begin{verbatim}
\chapter{...}
\section{...}
\subsection{...}
\end{verbatim}

Chapter, section, and subsection numbers shall be generated automatically
by \LaTeX{}.

Students shall not manually type chapter or section numbers and shall not
reproduce heading styles using commands such as \verb|\textbf|,
\verb|\fontsize|, or \verb|\centering| instead of the appropriate
structural command.

Excessive subdivision should be avoided. In most cases, the hierarchy
Chapter--Section--Subsection is sufficient.

\vspace{0.3cm}

% ================================================================
% IX. ORGANIZATION OF THE MAIN CONTENT
% ================================================================

\noindent
\textbf{IX.\quad ORGANIZATION OF THE MAIN CONTENT}

\vspace{0.15cm}

The specific content of the report shall depend on the internship
organization and the assigned tasks. A typical structure may include:

\begin{itemize}[leftmargin=0.8cm, itemsep=2pt]
    \item \textbf{Chapter 1. Overview of the Internship Organization};
    \item \textbf{Chapter 2. Internship Activities and Assigned Tasks};
    \item \textbf{Chapter 3. Results, Experience, and Evaluation}.
\end{itemize}

These chapter titles are recommendations only and may be adapted to the
actual internship content.

The report shall reflect work actually performed, observed, analyzed, or
learned during the internship. Unrelated theoretical material should not
be included merely to increase report length.

\vspace{0.3cm}

% ================================================================
% X. TABLE OF CONTENTS
% ================================================================

\noindent
\textbf{X.\quad TABLE OF CONTENTS}

\vspace{0.15cm}

The Table of Contents shall be generated automatically using:

\begin{verbatim}
\tableofcontents
\end{verbatim}

Students shall not create or manually edit page numbers in the Table of
Contents.

The document shall be compiled sufficiently to ensure that all page
numbers and entries are correctly updated.

\vspace{0.3cm}

% ================================================================
% XI. LIST OF SYMBOLS AND ABBREVIATIONS
% ================================================================

\noindent
\textbf{XI.\quad LIST OF SYMBOLS AND ABBREVIATIONS}

\vspace{0.15cm}

Symbols and abbreviations that are used repeatedly in the report shall be
listed in the \textbf{List of Symbols and Abbreviations}.

The list may contain the following columns:

\begin{center}
\begin{tabular}{llll}
    No. & Symbol / Abbreviation & Unit & Definition
\end{tabular}
\end{center}

Only symbols and abbreviations actually used in the final report shall be
retained. Example entries contained in the template shall be removed if
they are not relevant to the student's report.

Each abbreviation should normally be defined when it first appears in
the main text.

\vspace{0.3cm}

% ================================================================
% XII. FIGURES, DIAGRAMS, AND PHOTOGRAPHS
% ================================================================

\noindent
\textbf{XII.\quad FIGURES, DIAGRAMS, AND PHOTOGRAPHS}

\vspace{0.15cm}

Figures shall be clear, legible, and directly relevant to the report
content.

Figures shall be inserted using the standard \verb|figure| environment.

Example:

\begin{verbatim}
\begin{figure}[H]
    \centering
    \includegraphics[width=0.8\textwidth]
                    {figures/company_structure.pdf}
    \caption{Organizational structure of the company}
    \label{fig:company-structure}
\end{figure}
\end{verbatim}

Figure captions shall be placed below the figure.

Each referenced figure shall have a unique \verb|\label| and shall be
referenced in the text using, for example:

\begin{verbatim}
Figure~\ref{fig:company-structure}
\end{verbatim}

Students shall not manually type figure numbers.

Figures reproduced or adapted from external sources shall include an
appropriate citation.

Photographs taken during the internship shall be relevant and sufficiently
clear. Students shall not include confidential, restricted, or proprietary
information without permission from the internship organization.

Vector graphics or high-resolution images are recommended whenever
possible.

\vspace{0.3cm}

% ================================================================
% XIII. TABLES
% ================================================================

\noindent
\textbf{XIII.\quad TABLES}

\vspace{0.15cm}

Tables shall be created using suitable \LaTeX{} environments such as
\verb|table|, \verb|tabular|, \verb|tabularx|, or \verb|longtable|.

Table captions shall be placed above the table.

Example:

\begin{verbatim}
\begin{table}[H]
    \centering
    \caption{Internship schedule}
    \label{tab:internship-schedule}

    \begin{tabular}{lll}
        \toprule
        Week & Task & Result \\
        \midrule
        1 & Orientation & Completed \\
        2 & System study & Completed \\
        \bottomrule
    \end{tabular}
\end{table}
\end{verbatim}

Tables shall be referenced using:

\begin{verbatim}
Table~\ref{tab:internship-schedule}
\end{verbatim}

Students shall not manually type table numbers.

Tables should be created directly in \LaTeX{} rather than inserted as
screenshots whenever practical.

\vspace{0.3cm}

% ================================================================
% XIV. EQUATIONS AND TECHNICAL QUANTITIES
% ================================================================

\noindent
\textbf{XIV.\quad EQUATIONS AND TECHNICAL QUANTITIES}

\vspace{0.15cm}

Mathematical expressions shall be typeset directly using \LaTeX{}.
Images or screenshots of equations shall not be used.

Short expressions may be written in inline mathematical mode, for example:

\begin{verbatim}
$R = 0.7\,\Omega$
\end{verbatim}

Important equations shall use a numbered equation environment.

Example:

\begin{verbatim}
\begin{equation}
    P = UI
    \label{eq:power}
\end{equation}
\end{verbatim}

Equations shall be referenced using:

\begin{verbatim}
Equation~\eqref{eq:power}
\end{verbatim}

Equation numbers shall not be manually entered.

SI units shall be used wherever applicable. The \verb|siunitx| package
should be used where appropriate.

Examples:

\begin{verbatim}
\SI{220}{\volt}
\SI{50}{\hertz}
\SI{2.5}{\ampere}
\end{verbatim}

Notation and units shall be consistent throughout the report.

\vspace{0.3cm}

% ================================================================
% XV. SOURCE CODE
% ================================================================

\noindent
\textbf{XV.\quad SOURCE CODE}

\vspace{0.15cm}

If the internship involves programming, PLCs, microcontrollers,
simulation, or software development, relevant source code may be included
in the report.

Source code shall be presented using the \verb|listings| environment or
another environment defined by the official template.

Example:

\begin{verbatim}
\begin{lstlisting}[language=st]
IF Start = TRUE THEN
    Motor := TRUE;
ELSE
    Motor := FALSE;
END_IF;
\end{lstlisting}
\end{verbatim}

Only code necessary for explaining the work should appear in the main
report. Long source-code listings should be moved to an appendix or
submitted as supplementary files if required.

Source code should not normally be inserted as screenshots.

\vspace{0.3cm}

% ================================================================
% XVI. ABSTRACT
% ================================================================

\noindent
\textbf{XVI.\quad ABSTRACT}

\vspace{0.15cm}

The Abstract shall briefly summarize:

\begin{itemize}[leftmargin=0.8cm, itemsep=2pt]
    \item the internship organization;
    \item the objectives of the internship;
    \item the main tasks performed;
    \item the main results, knowledge, and skills obtained.
\end{itemize}

The Abstract should be concise and should not reproduce the full content
of the report chapter by chapter.

\vspace{0.3cm}

% ================================================================
% XVII. REFERENCES AND CITATIONS
% ================================================================

\noindent
\textbf{XVII.\quad REFERENCES AND CITATIONS}

\vspace{0.15cm}

All books, journal papers, conference papers, standards, websites,
datasheets, manuals, technical documents, figures, tables, or other
external sources used in preparing the report shall be properly cited.

References shall be managed through the bibliography system provided with
the official template.

Bibliographic entries shall normally be stored in:

\begin{verbatim}
References.bib
\end{verbatim}

References shall be cited in the report using:

\begin{verbatim}
\cite{reference-key}
\end{verbatim}

Students shall not manually type reference numbers such as ``[1]'',
``[2]'', etc.

Titles of foreign-language publications shall remain in their original
language in the bibliography.

All borrowed ideas, data, figures, tables, equations, methods, or results
shall be properly acknowledged.

\vspace{0.3cm}

% ================================================================
% XVIII. CROSS-REFERENCING
% ================================================================

\noindent
\textbf{XVIII.\quad CROSS-REFERENCING}

\vspace{0.15cm}

Numbered elements shall be cross-referenced using the standard \LaTeX{}
mechanism.

Use the following commands where appropriate:

\begin{verbatim}
\label{...}
\ref{...}
\eqref{...}
\end{verbatim}

For example:

\begin{verbatim}
As shown in Figure~\ref{fig:system}, ...
\end{verbatim}

Students shall not manually type numbers such as ``Figure 2.3'' when
those numbers can be generated automatically.

\vspace{0.3cm}

% ================================================================
% XIX. PAGE NUMBERING
% ================================================================

\noindent
\textbf{XIX.\quad PAGE NUMBERING}

\vspace{0.15cm}

Cover pages and official evaluation forms at the beginning of the report
shall not display page numbers.

The main body of the report shall use Arabic page numbers beginning from
Chapter 1.

Page numbering shall be generated automatically by the official template.
Students shall not manually enter page numbers.

\vspace{0.3cm}

% ================================================================
% XX. PROJECT FILE ORGANIZATION
% ================================================================

\noindent
\textbf{XX.\quad PROJECT FILE ORGANIZATION}

\vspace{0.15cm}

Students should maintain a clear and organized \LaTeX{} project
structure. A recommended structure is:

\begin{verbatim}
InternshipReport/
|
|-- main.tex
|-- References.bib
|
|-- figures/
|   |-- company.pdf
|   |-- system.png
|   `-- ...
|
`-- chapters/
    |-- chapter1.tex
    |-- chapter2.tex
    `-- chapter3.tex
\end{verbatim}

Figures should normally be stored in the \verb|figures| directory
specified by the official template.

File and directory names should preferably contain Latin letters,
numbers, hyphens, or underscores.

Recommended examples include:

\begin{verbatim}
company_structure.pdf
motor_control.png
chapter_01.tex
\end{verbatim}

Spaces and unnecessary special characters in file names should be
avoided where possible.

\vspace{0.3cm}

% ================================================================
% XXI. DOCUMENT COMPILATION AND FINAL CHECK
% ================================================================

\noindent
\textbf{XXI.\quad DOCUMENT COMPILATION AND FINAL CHECK}

\vspace{0.15cm}

Before submission, the complete \LaTeX{} project shall compile
successfully and produce the final PDF report.

The final document shall not contain unresolved references such as:

\begin{verbatim}
??
[?]
\end{verbatim}

Before generating the final PDF, students shall check:

\begin{itemize}[leftmargin=0.8cm, itemsep=2pt]
    \item the Table of Contents;
    \item page numbering;
    \item chapter, section, and subsection numbering;
    \item figure numbering;
    \item table numbering;
    \item equation numbering;
    \item cross-references;
    \item citations;
    \item bibliography;
    \item font display;
    \item mathematical symbols;
    \item image quality;
    \item completeness of all required forms.
\end{itemize}

The document may be compiled using TeX Live, MiKTeX, or another
compatible \LaTeX{} environment. Regardless of the compilation
environment, the resulting PDF shall conform to the official Faculty
template.

\vspace{0.3cm}

% ================================================================
% XXII. CONTENT THAT STUDENTS MAY EDIT
% ================================================================

\noindent
\textbf{XXII.\quad CONTENT THAT STUDENTS MAY EDIT}

\vspace{0.15cm}

Students may edit or replace:

\begin{itemize}[leftmargin=0.8cm, itemsep=2pt]
    \item student name;
    \item class;
    \item internship organization;
    \item supervisor name;
    \item Abstract;
    \item List of Symbols and Abbreviations;
    \item chapter titles and chapter content;
    \item figures and photographs;
    \item tables;
    \item equations;
    \item citations and references.
\end{itemize}

Students shall not modify the main document formatting settings in the
preamble, including the document class, page geometry, line spacing,
heading formatting, table-of-contents formatting, font configuration,
page layout, or cover-page macros unless specifically instructed.

\vspace{0.3cm}

% ================================================================
% XXIII. FINAL SUBMISSION REQUIREMENTS
% ================================================================

\noindent
\textbf{XXIII.\quad FINAL SUBMISSION REQUIREMENTS}

\vspace{0.15cm}

Before submission, the final report shall:

\begin{itemize}[leftmargin=0.8cm, itemsep=2pt]
    \item compile without errors that affect the report content or presentation;
    \item contain all required forms and report sections;
    \item contain correct student, class, internship organization, and supervisor information;
    \item contain an updated Table of Contents, List of Figures, and List of Tables;
    \item contain no unused example content from the template;
    \item contain no placeholder text such as \textit{Title of Chapter 1},
    \textit{Section 1.1}, or other sample entries;
    \item contain all required citations and references;
    \item contain no unresolved cross-references.
\end{itemize}

If the Faculty requires submission of the \LaTeX{} source, the complete
project shall include all \verb|.tex| files, the bibliography database,
figures, and any other files required to reproduce the submitted PDF.

\vspace{0.4cm}

\noindent
\textbf{Important:}
Before submitting the final internship report, students shall remove or
comment out the following command in the main document:

\begin{verbatim}
\input{Regulations.tex}
\end{verbatim}

For example:

\begin{verbatim}
% \input{Regulations.tex}
\end{verbatim}

This removes the guidance pages from the final submitted report while
preserving the regulations file in the \LaTeX{} project for reference.

} % End 13pt content

\clearpage

\end{verbatim}

For example:

\begin{verbatim}
% 
% =================================================================
% REGULATIONS FOR PREPARING THE GRADUATION INTERNSHIP REPORT
% USING LATEX
% =================================================================
%
% This file is intended to be included from the main document with:
%
%     \input{Regulations.tex}
%
% IMPORTANT:
% This section is provided for guidance only. Students shall remove
% or comment out the \input{Regulations.tex} command before submitting
% the final internship report.
%
% =================================================================

\clearpage
\thispagestyle{empty}

\begin{center}
    {\fontsize{16pt}{20pt}\selectfont\bfseries
    SPECIFIC REGULATIONS FOR PREPARING\\
    GRADUATION INTERNSHIP REPORTS USING \LaTeX}
\end{center}

\vspace{0.5cm}

{\fontsize{13pt}{17pt}\selectfont
\setstretch{1.3}
\justifying

\noindent
\textbf{Note:}
This section is provided for guidance only and is included in the main
document using the command \verb|\input{Regulations.tex}|.
Students shall remove or comment out this command before submitting the
final internship report.

\vspace{0.4cm}

% ================================================================
% I. GENERAL REQUIREMENTS
% ================================================================

\noindent
\textbf{I.\quad GENERAL REQUIREMENTS}

\vspace{0.15cm}

The Graduation Internship Report shall be prepared using the official
\LaTeX{} template provided by the Faculty. Students shall preserve the
predefined page layout, typography, line spacing, heading styles,
numbering system, figure and table formatting, and bibliography settings
unless otherwise instructed by the Faculty or the supervisor.

The report shall be written in English.

Students are responsible for the academic and technical content of the
report. The official \LaTeX{} template is responsible for standardized
document formatting.

\vspace{0.3cm}

% ================================================================
% II. REPORT STRUCTURE
% ================================================================

\noindent
\textbf{II.\quad REPORT STRUCTURE}

\vspace{0.15cm}

The report shall normally be arranged in the following order:

\begin{enumerate}[label=\arabic*., leftmargin=1.2cm, itemsep=2pt]
    \item Outer Cover;
    \item Comments of the Internship Host Organization;
    \item Supervisor's Evaluation and Graduation Internship Defense Results;
    \item Table of Contents;
    \item List of Symbols and Abbreviations, if applicable;
    \item List of Figures;
    \item List of Tables;
    \item Abstract;
    \item Main Chapters of the Internship Report;
    \item References;
    \item Appendices, if applicable.
\end{enumerate}

The preliminary pages shall not be numbered as chapters. The main report
content shall be organized into consecutively numbered chapters beginning
with Chapter 1.

\vspace{0.3cm}

% ================================================================
% III. PAGE SIZE AND MARGINS
% ================================================================

\noindent
\textbf{III.\quad PAGE SIZE AND MARGINS}

\vspace{0.15cm}

The report shall be formatted on A4 paper with the following margins:

\begin{center}
\begin{tabular}{ll}
    Left margin:   & 3.5 cm \\
    Right margin:  & 2.0 cm \\
    Top margin:    & 2.5 cm \\
    Bottom margin: & 2.5 cm
\end{tabular}
\end{center}

These settings are predefined in the official template and shall not be
modified by students.

\vspace{0.3cm}

% ================================================================
% IV. FONT, FONT SIZE, AND LINE SPACING
% ================================================================

\noindent
\textbf{IV.\quad FONT, FONT SIZE, AND LINE SPACING}

\vspace{0.15cm}

The main text shall use a 13 pt font size with 1.3-line spacing.
A Times-family typeface shall be used in accordance with the official
template.

The text shall be fully justified.

Students shall not use repeated spaces, repeated manual line breaks, or
multiple blank lines to position text. Paragraph and heading spacing shall
be controlled by the template.

\vspace{0.3cm}

% ================================================================
% V. OUTER COVER
% ================================================================

\noindent
\textbf{V.\quad OUTER COVER}

\vspace{0.15cm}

The outer cover shall follow the official format provided in the template
and shall contain the following information:

\begin{itemize}[leftmargin=0.8cm, itemsep=2pt]
    \item Thai Nguyen University of Technology;
    \item Faculty of International Training;
    \item Engineering and Technology Department;
    \item University logo and Faculty logo;
    \item the title \textbf{INTERNSHIP REPORT};
    \item Internship Organization;
    \item Author;
    \item Class;
    \item Supervisor;
    \item place and year of completion.
\end{itemize}

Students shall only replace the required information fields and shall not
modify the cover-page layout unless instructed otherwise.

\vspace{0.3cm}

% ================================================================
% VI. COMMENTS OF THE INTERNSHIP HOST ORGANIZATION
% ================================================================

\noindent
\textbf{VI.\quad COMMENTS OF THE INTERNSHIP HOST ORGANIZATION}

\vspace{0.15cm}

The page entitled \textbf{COMMENTS OF THE INTERNSHIP HOST ORGANIZATION}
is reserved for the internship host organization.

It includes:

\begin{itemize}[leftmargin=0.8cm, itemsep=2pt]
    \item the student's full name;
    \item class;
    \item comments from the host organization;
    \item score;
    \item date;
    \item signature and official seal.
\end{itemize}

Students shall not alter or complete official evaluation fields unless
specifically authorized.

\vspace{0.3cm}

% ================================================================
% VII. SUPERVISOR'S EVALUATION AND DEFENSE RESULTS
% ================================================================

\noindent
\textbf{VII.\quad SUPERVISOR'S EVALUATION AND DEFENSE RESULTS}

\vspace{0.15cm}

The section entitled
\textbf{SUPERVISOR'S EVALUATION AND GRADUATION INTERNSHIP DEFENSE RESULTS}
is reserved for the supervisor and defense committee.

The evaluation may include:

\begin{itemize}[leftmargin=0.8cm, itemsep=2pt]
    \item Internship Progress and Student's Attitude;
    \item Report Content;
    \item Report Presentation;
    \item Other Comments;
    \item Supervisor's Opinion;
    \item Defense Result;
    \item Score.
\end{itemize}

Students shall not enter evaluation scores or defense results before they
are officially determined.

\vspace{0.3cm}

% ================================================================
% VIII. CHAPTERS, SECTIONS, AND SUBSECTIONS
% ================================================================

\noindent
\textbf{VIII.\quad CHAPTERS, SECTIONS, AND SUBSECTIONS}

\vspace{0.15cm}

The report structure shall be generated using the standard \LaTeX{}
commands:

\begin{verbatim}
\chapter{...}
\section{...}
\subsection{...}
\end{verbatim}

Chapter, section, and subsection numbers shall be generated automatically
by \LaTeX{}.

Students shall not manually type chapter or section numbers and shall not
reproduce heading styles using commands such as \verb|\textbf|,
\verb|\fontsize|, or \verb|\centering| instead of the appropriate
structural command.

Excessive subdivision should be avoided. In most cases, the hierarchy
Chapter--Section--Subsection is sufficient.

\vspace{0.3cm}

% ================================================================
% IX. ORGANIZATION OF THE MAIN CONTENT
% ================================================================

\noindent
\textbf{IX.\quad ORGANIZATION OF THE MAIN CONTENT}

\vspace{0.15cm}

The specific content of the report shall depend on the internship
organization and the assigned tasks. A typical structure may include:

\begin{itemize}[leftmargin=0.8cm, itemsep=2pt]
    \item \textbf{Chapter 1. Overview of the Internship Organization};
    \item \textbf{Chapter 2. Internship Activities and Assigned Tasks};
    \item \textbf{Chapter 3. Results, Experience, and Evaluation}.
\end{itemize}

These chapter titles are recommendations only and may be adapted to the
actual internship content.

The report shall reflect work actually performed, observed, analyzed, or
learned during the internship. Unrelated theoretical material should not
be included merely to increase report length.

\vspace{0.3cm}

% ================================================================
% X. TABLE OF CONTENTS
% ================================================================

\noindent
\textbf{X.\quad TABLE OF CONTENTS}

\vspace{0.15cm}

The Table of Contents shall be generated automatically using:

\begin{verbatim}
\tableofcontents
\end{verbatim}

Students shall not create or manually edit page numbers in the Table of
Contents.

The document shall be compiled sufficiently to ensure that all page
numbers and entries are correctly updated.

\vspace{0.3cm}

% ================================================================
% XI. LIST OF SYMBOLS AND ABBREVIATIONS
% ================================================================

\noindent
\textbf{XI.\quad LIST OF SYMBOLS AND ABBREVIATIONS}

\vspace{0.15cm}

Symbols and abbreviations that are used repeatedly in the report shall be
listed in the \textbf{List of Symbols and Abbreviations}.

The list may contain the following columns:

\begin{center}
\begin{tabular}{llll}
    No. & Symbol / Abbreviation & Unit & Definition
\end{tabular}
\end{center}

Only symbols and abbreviations actually used in the final report shall be
retained. Example entries contained in the template shall be removed if
they are not relevant to the student's report.

Each abbreviation should normally be defined when it first appears in
the main text.

\vspace{0.3cm}

% ================================================================
% XII. FIGURES, DIAGRAMS, AND PHOTOGRAPHS
% ================================================================

\noindent
\textbf{XII.\quad FIGURES, DIAGRAMS, AND PHOTOGRAPHS}

\vspace{0.15cm}

Figures shall be clear, legible, and directly relevant to the report
content.

Figures shall be inserted using the standard \verb|figure| environment.

Example:

\begin{verbatim}
\begin{figure}[H]
    \centering
    \includegraphics[width=0.8\textwidth]
                    {figures/company_structure.pdf}
    \caption{Organizational structure of the company}
    \label{fig:company-structure}
\end{figure}
\end{verbatim}

Figure captions shall be placed below the figure.

Each referenced figure shall have a unique \verb|\label| and shall be
referenced in the text using, for example:

\begin{verbatim}
Figure~\ref{fig:company-structure}
\end{verbatim}

Students shall not manually type figure numbers.

Figures reproduced or adapted from external sources shall include an
appropriate citation.

Photographs taken during the internship shall be relevant and sufficiently
clear. Students shall not include confidential, restricted, or proprietary
information without permission from the internship organization.

Vector graphics or high-resolution images are recommended whenever
possible.

\vspace{0.3cm}

% ================================================================
% XIII. TABLES
% ================================================================

\noindent
\textbf{XIII.\quad TABLES}

\vspace{0.15cm}

Tables shall be created using suitable \LaTeX{} environments such as
\verb|table|, \verb|tabular|, \verb|tabularx|, or \verb|longtable|.

Table captions shall be placed above the table.

Example:

\begin{verbatim}
\begin{table}[H]
    \centering
    \caption{Internship schedule}
    \label{tab:internship-schedule}

    \begin{tabular}{lll}
        \toprule
        Week & Task & Result \\
        \midrule
        1 & Orientation & Completed \\
        2 & System study & Completed \\
        \bottomrule
    \end{tabular}
\end{table}
\end{verbatim}

Tables shall be referenced using:

\begin{verbatim}
Table~\ref{tab:internship-schedule}
\end{verbatim}

Students shall not manually type table numbers.

Tables should be created directly in \LaTeX{} rather than inserted as
screenshots whenever practical.

\vspace{0.3cm}

% ================================================================
% XIV. EQUATIONS AND TECHNICAL QUANTITIES
% ================================================================

\noindent
\textbf{XIV.\quad EQUATIONS AND TECHNICAL QUANTITIES}

\vspace{0.15cm}

Mathematical expressions shall be typeset directly using \LaTeX{}.
Images or screenshots of equations shall not be used.

Short expressions may be written in inline mathematical mode, for example:

\begin{verbatim}
$R = 0.7\,\Omega$
\end{verbatim}

Important equations shall use a numbered equation environment.

Example:

\begin{verbatim}
\begin{equation}
    P = UI
    \label{eq:power}
\end{equation}
\end{verbatim}

Equations shall be referenced using:

\begin{verbatim}
Equation~\eqref{eq:power}
\end{verbatim}

Equation numbers shall not be manually entered.

SI units shall be used wherever applicable. The \verb|siunitx| package
should be used where appropriate.

Examples:

\begin{verbatim}
\SI{220}{\volt}
\SI{50}{\hertz}
\SI{2.5}{\ampere}
\end{verbatim}

Notation and units shall be consistent throughout the report.

\vspace{0.3cm}

% ================================================================
% XV. SOURCE CODE
% ================================================================

\noindent
\textbf{XV.\quad SOURCE CODE}

\vspace{0.15cm}

If the internship involves programming, PLCs, microcontrollers,
simulation, or software development, relevant source code may be included
in the report.

Source code shall be presented using the \verb|listings| environment or
another environment defined by the official template.

Example:

\begin{verbatim}
\begin{lstlisting}[language=st]
IF Start = TRUE THEN
    Motor := TRUE;
ELSE
    Motor := FALSE;
END_IF;
\end{lstlisting}
\end{verbatim}

Only code necessary for explaining the work should appear in the main
report. Long source-code listings should be moved to an appendix or
submitted as supplementary files if required.

Source code should not normally be inserted as screenshots.

\vspace{0.3cm}

% ================================================================
% XVI. ABSTRACT
% ================================================================

\noindent
\textbf{XVI.\quad ABSTRACT}

\vspace{0.15cm}

The Abstract shall briefly summarize:

\begin{itemize}[leftmargin=0.8cm, itemsep=2pt]
    \item the internship organization;
    \item the objectives of the internship;
    \item the main tasks performed;
    \item the main results, knowledge, and skills obtained.
\end{itemize}

The Abstract should be concise and should not reproduce the full content
of the report chapter by chapter.

\vspace{0.3cm}

% ================================================================
% XVII. REFERENCES AND CITATIONS
% ================================================================

\noindent
\textbf{XVII.\quad REFERENCES AND CITATIONS}

\vspace{0.15cm}

All books, journal papers, conference papers, standards, websites,
datasheets, manuals, technical documents, figures, tables, or other
external sources used in preparing the report shall be properly cited.

References shall be managed through the bibliography system provided with
the official template.

Bibliographic entries shall normally be stored in:

\begin{verbatim}
References.bib
\end{verbatim}

References shall be cited in the report using:

\begin{verbatim}
\cite{reference-key}
\end{verbatim}

Students shall not manually type reference numbers such as ``[1]'',
``[2]'', etc.

Titles of foreign-language publications shall remain in their original
language in the bibliography.

All borrowed ideas, data, figures, tables, equations, methods, or results
shall be properly acknowledged.

\vspace{0.3cm}

% ================================================================
% XVIII. CROSS-REFERENCING
% ================================================================

\noindent
\textbf{XVIII.\quad CROSS-REFERENCING}

\vspace{0.15cm}

Numbered elements shall be cross-referenced using the standard \LaTeX{}
mechanism.

Use the following commands where appropriate:

\begin{verbatim}
\label{...}
\ref{...}
\eqref{...}
\end{verbatim}

For example:

\begin{verbatim}
As shown in Figure~\ref{fig:system}, ...
\end{verbatim}

Students shall not manually type numbers such as ``Figure 2.3'' when
those numbers can be generated automatically.

\vspace{0.3cm}

% ================================================================
% XIX. PAGE NUMBERING
% ================================================================

\noindent
\textbf{XIX.\quad PAGE NUMBERING}

\vspace{0.15cm}

Cover pages and official evaluation forms at the beginning of the report
shall not display page numbers.

The main body of the report shall use Arabic page numbers beginning from
Chapter 1.

Page numbering shall be generated automatically by the official template.
Students shall not manually enter page numbers.

\vspace{0.3cm}

% ================================================================
% XX. PROJECT FILE ORGANIZATION
% ================================================================

\noindent
\textbf{XX.\quad PROJECT FILE ORGANIZATION}

\vspace{0.15cm}

Students should maintain a clear and organized \LaTeX{} project
structure. A recommended structure is:

\begin{verbatim}
InternshipReport/
|
|-- main.tex
|-- References.bib
|
|-- figures/
|   |-- company.pdf
|   |-- system.png
|   `-- ...
|
`-- chapters/
    |-- chapter1.tex
    |-- chapter2.tex
    `-- chapter3.tex
\end{verbatim}

Figures should normally be stored in the \verb|figures| directory
specified by the official template.

File and directory names should preferably contain Latin letters,
numbers, hyphens, or underscores.

Recommended examples include:

\begin{verbatim}
company_structure.pdf
motor_control.png
chapter_01.tex
\end{verbatim}

Spaces and unnecessary special characters in file names should be
avoided where possible.

\vspace{0.3cm}

% ================================================================
% XXI. DOCUMENT COMPILATION AND FINAL CHECK
% ================================================================

\noindent
\textbf{XXI.\quad DOCUMENT COMPILATION AND FINAL CHECK}

\vspace{0.15cm}

Before submission, the complete \LaTeX{} project shall compile
successfully and produce the final PDF report.

The final document shall not contain unresolved references such as:

\begin{verbatim}
??
[?]
\end{verbatim}

Before generating the final PDF, students shall check:

\begin{itemize}[leftmargin=0.8cm, itemsep=2pt]
    \item the Table of Contents;
    \item page numbering;
    \item chapter, section, and subsection numbering;
    \item figure numbering;
    \item table numbering;
    \item equation numbering;
    \item cross-references;
    \item citations;
    \item bibliography;
    \item font display;
    \item mathematical symbols;
    \item image quality;
    \item completeness of all required forms.
\end{itemize}

The document may be compiled using TeX Live, MiKTeX, or another
compatible \LaTeX{} environment. Regardless of the compilation
environment, the resulting PDF shall conform to the official Faculty
template.

\vspace{0.3cm}

% ================================================================
% XXII. CONTENT THAT STUDENTS MAY EDIT
% ================================================================

\noindent
\textbf{XXII.\quad CONTENT THAT STUDENTS MAY EDIT}

\vspace{0.15cm}

Students may edit or replace:

\begin{itemize}[leftmargin=0.8cm, itemsep=2pt]
    \item student name;
    \item class;
    \item internship organization;
    \item supervisor name;
    \item Abstract;
    \item List of Symbols and Abbreviations;
    \item chapter titles and chapter content;
    \item figures and photographs;
    \item tables;
    \item equations;
    \item citations and references.
\end{itemize}

Students shall not modify the main document formatting settings in the
preamble, including the document class, page geometry, line spacing,
heading formatting, table-of-contents formatting, font configuration,
page layout, or cover-page macros unless specifically instructed.

\vspace{0.3cm}

% ================================================================
% XXIII. FINAL SUBMISSION REQUIREMENTS
% ================================================================

\noindent
\textbf{XXIII.\quad FINAL SUBMISSION REQUIREMENTS}

\vspace{0.15cm}

Before submission, the final report shall:

\begin{itemize}[leftmargin=0.8cm, itemsep=2pt]
    \item compile without errors that affect the report content or presentation;
    \item contain all required forms and report sections;
    \item contain correct student, class, internship organization, and supervisor information;
    \item contain an updated Table of Contents, List of Figures, and List of Tables;
    \item contain no unused example content from the template;
    \item contain no placeholder text such as \textit{Title of Chapter 1},
    \textit{Section 1.1}, or other sample entries;
    \item contain all required citations and references;
    \item contain no unresolved cross-references.
\end{itemize}

If the Faculty requires submission of the \LaTeX{} source, the complete
project shall include all \verb|.tex| files, the bibliography database,
figures, and any other files required to reproduce the submitted PDF.

\vspace{0.4cm}

\noindent
\textbf{Important:}
Before submitting the final internship report, students shall remove or
comment out the following command in the main document:

\begin{verbatim}
\input{Regulations.tex}
\end{verbatim}

For example:

\begin{verbatim}
% \input{Regulations.tex}
\end{verbatim}

This removes the guidance pages from the final submitted report while
preserving the regulations file in the \LaTeX{} project for reference.

} % End 13pt content

\clearpage

\end{verbatim}

This removes the guidance pages from the final submitted report while
preserving the regulations file in the \LaTeX{} project for reference.

} % End 13pt content

\clearpage
 command before submitting
% the final internship report.
%
% =================================================================

\clearpage
\thispagestyle{empty}

\begin{center}
    {\fontsize{16pt}{20pt}\selectfont\bfseries
    SPECIFIC REGULATIONS FOR PREPARING\\
    GRADUATION INTERNSHIP REPORTS USING \LaTeX}
\end{center}

\vspace{0.5cm}

{\fontsize{13pt}{17pt}\selectfont
\setstretch{1.3}
\justifying

\noindent
\textbf{Note:}
This section is provided for guidance only and is included in the main
document using the command \verb|
% =================================================================
% REGULATIONS FOR PREPARING THE GRADUATION INTERNSHIP REPORT
% USING LATEX
% =================================================================
%
% This file is intended to be included from the main document with:
%
%     
% =================================================================
% REGULATIONS FOR PREPARING THE GRADUATION INTERNSHIP REPORT
% USING LATEX
% =================================================================
%
% This file is intended to be included from the main document with:
%
%     \input{Regulations.tex}
%
% IMPORTANT:
% This section is provided for guidance only. Students shall remove
% or comment out the \input{Regulations.tex} command before submitting
% the final internship report.
%
% =================================================================

\clearpage
\thispagestyle{empty}

\begin{center}
    {\fontsize{16pt}{20pt}\selectfont\bfseries
    SPECIFIC REGULATIONS FOR PREPARING\\
    GRADUATION INTERNSHIP REPORTS USING \LaTeX}
\end{center}

\vspace{0.5cm}

{\fontsize{13pt}{17pt}\selectfont
\setstretch{1.3}
\justifying

\noindent
\textbf{Note:}
This section is provided for guidance only and is included in the main
document using the command \verb|\input{Regulations.tex}|.
Students shall remove or comment out this command before submitting the
final internship report.

\vspace{0.4cm}

% ================================================================
% I. GENERAL REQUIREMENTS
% ================================================================

\noindent
\textbf{I.\quad GENERAL REQUIREMENTS}

\vspace{0.15cm}

The Graduation Internship Report shall be prepared using the official
\LaTeX{} template provided by the Faculty. Students shall preserve the
predefined page layout, typography, line spacing, heading styles,
numbering system, figure and table formatting, and bibliography settings
unless otherwise instructed by the Faculty or the supervisor.

The report shall be written in English.

Students are responsible for the academic and technical content of the
report. The official \LaTeX{} template is responsible for standardized
document formatting.

\vspace{0.3cm}

% ================================================================
% II. REPORT STRUCTURE
% ================================================================

\noindent
\textbf{II.\quad REPORT STRUCTURE}

\vspace{0.15cm}

The report shall normally be arranged in the following order:

\begin{enumerate}[label=\arabic*., leftmargin=1.2cm, itemsep=2pt]
    \item Outer Cover;
    \item Comments of the Internship Host Organization;
    \item Supervisor's Evaluation and Graduation Internship Defense Results;
    \item Table of Contents;
    \item List of Symbols and Abbreviations, if applicable;
    \item List of Figures;
    \item List of Tables;
    \item Abstract;
    \item Main Chapters of the Internship Report;
    \item References;
    \item Appendices, if applicable.
\end{enumerate}

The preliminary pages shall not be numbered as chapters. The main report
content shall be organized into consecutively numbered chapters beginning
with Chapter 1.

\vspace{0.3cm}

% ================================================================
% III. PAGE SIZE AND MARGINS
% ================================================================

\noindent
\textbf{III.\quad PAGE SIZE AND MARGINS}

\vspace{0.15cm}

The report shall be formatted on A4 paper with the following margins:

\begin{center}
\begin{tabular}{ll}
    Left margin:   & 3.5 cm \\
    Right margin:  & 2.0 cm \\
    Top margin:    & 2.5 cm \\
    Bottom margin: & 2.5 cm
\end{tabular}
\end{center}

These settings are predefined in the official template and shall not be
modified by students.

\vspace{0.3cm}

% ================================================================
% IV. FONT, FONT SIZE, AND LINE SPACING
% ================================================================

\noindent
\textbf{IV.\quad FONT, FONT SIZE, AND LINE SPACING}

\vspace{0.15cm}

The main text shall use a 13 pt font size with 1.3-line spacing.
A Times-family typeface shall be used in accordance with the official
template.

The text shall be fully justified.

Students shall not use repeated spaces, repeated manual line breaks, or
multiple blank lines to position text. Paragraph and heading spacing shall
be controlled by the template.

\vspace{0.3cm}

% ================================================================
% V. OUTER COVER
% ================================================================

\noindent
\textbf{V.\quad OUTER COVER}

\vspace{0.15cm}

The outer cover shall follow the official format provided in the template
and shall contain the following information:

\begin{itemize}[leftmargin=0.8cm, itemsep=2pt]
    \item Thai Nguyen University of Technology;
    \item Faculty of International Training;
    \item Engineering and Technology Department;
    \item University logo and Faculty logo;
    \item the title \textbf{INTERNSHIP REPORT};
    \item Internship Organization;
    \item Author;
    \item Class;
    \item Supervisor;
    \item place and year of completion.
\end{itemize}

Students shall only replace the required information fields and shall not
modify the cover-page layout unless instructed otherwise.

\vspace{0.3cm}

% ================================================================
% VI. COMMENTS OF THE INTERNSHIP HOST ORGANIZATION
% ================================================================

\noindent
\textbf{VI.\quad COMMENTS OF THE INTERNSHIP HOST ORGANIZATION}

\vspace{0.15cm}

The page entitled \textbf{COMMENTS OF THE INTERNSHIP HOST ORGANIZATION}
is reserved for the internship host organization.

It includes:

\begin{itemize}[leftmargin=0.8cm, itemsep=2pt]
    \item the student's full name;
    \item class;
    \item comments from the host organization;
    \item score;
    \item date;
    \item signature and official seal.
\end{itemize}

Students shall not alter or complete official evaluation fields unless
specifically authorized.

\vspace{0.3cm}

% ================================================================
% VII. SUPERVISOR'S EVALUATION AND DEFENSE RESULTS
% ================================================================

\noindent
\textbf{VII.\quad SUPERVISOR'S EVALUATION AND DEFENSE RESULTS}

\vspace{0.15cm}

The section entitled
\textbf{SUPERVISOR'S EVALUATION AND GRADUATION INTERNSHIP DEFENSE RESULTS}
is reserved for the supervisor and defense committee.

The evaluation may include:

\begin{itemize}[leftmargin=0.8cm, itemsep=2pt]
    \item Internship Progress and Student's Attitude;
    \item Report Content;
    \item Report Presentation;
    \item Other Comments;
    \item Supervisor's Opinion;
    \item Defense Result;
    \item Score.
\end{itemize}

Students shall not enter evaluation scores or defense results before they
are officially determined.

\vspace{0.3cm}

% ================================================================
% VIII. CHAPTERS, SECTIONS, AND SUBSECTIONS
% ================================================================

\noindent
\textbf{VIII.\quad CHAPTERS, SECTIONS, AND SUBSECTIONS}

\vspace{0.15cm}

The report structure shall be generated using the standard \LaTeX{}
commands:

\begin{verbatim}
\chapter{...}
\section{...}
\subsection{...}
\end{verbatim}

Chapter, section, and subsection numbers shall be generated automatically
by \LaTeX{}.

Students shall not manually type chapter or section numbers and shall not
reproduce heading styles using commands such as \verb|\textbf|,
\verb|\fontsize|, or \verb|\centering| instead of the appropriate
structural command.

Excessive subdivision should be avoided. In most cases, the hierarchy
Chapter--Section--Subsection is sufficient.

\vspace{0.3cm}

% ================================================================
% IX. ORGANIZATION OF THE MAIN CONTENT
% ================================================================

\noindent
\textbf{IX.\quad ORGANIZATION OF THE MAIN CONTENT}

\vspace{0.15cm}

The specific content of the report shall depend on the internship
organization and the assigned tasks. A typical structure may include:

\begin{itemize}[leftmargin=0.8cm, itemsep=2pt]
    \item \textbf{Chapter 1. Overview of the Internship Organization};
    \item \textbf{Chapter 2. Internship Activities and Assigned Tasks};
    \item \textbf{Chapter 3. Results, Experience, and Evaluation}.
\end{itemize}

These chapter titles are recommendations only and may be adapted to the
actual internship content.

The report shall reflect work actually performed, observed, analyzed, or
learned during the internship. Unrelated theoretical material should not
be included merely to increase report length.

\vspace{0.3cm}

% ================================================================
% X. TABLE OF CONTENTS
% ================================================================

\noindent
\textbf{X.\quad TABLE OF CONTENTS}

\vspace{0.15cm}

The Table of Contents shall be generated automatically using:

\begin{verbatim}
\tableofcontents
\end{verbatim}

Students shall not create or manually edit page numbers in the Table of
Contents.

The document shall be compiled sufficiently to ensure that all page
numbers and entries are correctly updated.

\vspace{0.3cm}

% ================================================================
% XI. LIST OF SYMBOLS AND ABBREVIATIONS
% ================================================================

\noindent
\textbf{XI.\quad LIST OF SYMBOLS AND ABBREVIATIONS}

\vspace{0.15cm}

Symbols and abbreviations that are used repeatedly in the report shall be
listed in the \textbf{List of Symbols and Abbreviations}.

The list may contain the following columns:

\begin{center}
\begin{tabular}{llll}
    No. & Symbol / Abbreviation & Unit & Definition
\end{tabular}
\end{center}

Only symbols and abbreviations actually used in the final report shall be
retained. Example entries contained in the template shall be removed if
they are not relevant to the student's report.

Each abbreviation should normally be defined when it first appears in
the main text.

\vspace{0.3cm}

% ================================================================
% XII. FIGURES, DIAGRAMS, AND PHOTOGRAPHS
% ================================================================

\noindent
\textbf{XII.\quad FIGURES, DIAGRAMS, AND PHOTOGRAPHS}

\vspace{0.15cm}

Figures shall be clear, legible, and directly relevant to the report
content.

Figures shall be inserted using the standard \verb|figure| environment.

Example:

\begin{verbatim}
\begin{figure}[H]
    \centering
    \includegraphics[width=0.8\textwidth]
                    {figures/company_structure.pdf}
    \caption{Organizational structure of the company}
    \label{fig:company-structure}
\end{figure}
\end{verbatim}

Figure captions shall be placed below the figure.

Each referenced figure shall have a unique \verb|\label| and shall be
referenced in the text using, for example:

\begin{verbatim}
Figure~\ref{fig:company-structure}
\end{verbatim}

Students shall not manually type figure numbers.

Figures reproduced or adapted from external sources shall include an
appropriate citation.

Photographs taken during the internship shall be relevant and sufficiently
clear. Students shall not include confidential, restricted, or proprietary
information without permission from the internship organization.

Vector graphics or high-resolution images are recommended whenever
possible.

\vspace{0.3cm}

% ================================================================
% XIII. TABLES
% ================================================================

\noindent
\textbf{XIII.\quad TABLES}

\vspace{0.15cm}

Tables shall be created using suitable \LaTeX{} environments such as
\verb|table|, \verb|tabular|, \verb|tabularx|, or \verb|longtable|.

Table captions shall be placed above the table.

Example:

\begin{verbatim}
\begin{table}[H]
    \centering
    \caption{Internship schedule}
    \label{tab:internship-schedule}

    \begin{tabular}{lll}
        \toprule
        Week & Task & Result \\
        \midrule
        1 & Orientation & Completed \\
        2 & System study & Completed \\
        \bottomrule
    \end{tabular}
\end{table}
\end{verbatim}

Tables shall be referenced using:

\begin{verbatim}
Table~\ref{tab:internship-schedule}
\end{verbatim}

Students shall not manually type table numbers.

Tables should be created directly in \LaTeX{} rather than inserted as
screenshots whenever practical.

\vspace{0.3cm}

% ================================================================
% XIV. EQUATIONS AND TECHNICAL QUANTITIES
% ================================================================

\noindent
\textbf{XIV.\quad EQUATIONS AND TECHNICAL QUANTITIES}

\vspace{0.15cm}

Mathematical expressions shall be typeset directly using \LaTeX{}.
Images or screenshots of equations shall not be used.

Short expressions may be written in inline mathematical mode, for example:

\begin{verbatim}
$R = 0.7\,\Omega$
\end{verbatim}

Important equations shall use a numbered equation environment.

Example:

\begin{verbatim}
\begin{equation}
    P = UI
    \label{eq:power}
\end{equation}
\end{verbatim}

Equations shall be referenced using:

\begin{verbatim}
Equation~\eqref{eq:power}
\end{verbatim}

Equation numbers shall not be manually entered.

SI units shall be used wherever applicable. The \verb|siunitx| package
should be used where appropriate.

Examples:

\begin{verbatim}
\SI{220}{\volt}
\SI{50}{\hertz}
\SI{2.5}{\ampere}
\end{verbatim}

Notation and units shall be consistent throughout the report.

\vspace{0.3cm}

% ================================================================
% XV. SOURCE CODE
% ================================================================

\noindent
\textbf{XV.\quad SOURCE CODE}

\vspace{0.15cm}

If the internship involves programming, PLCs, microcontrollers,
simulation, or software development, relevant source code may be included
in the report.

Source code shall be presented using the \verb|listings| environment or
another environment defined by the official template.

Example:

\begin{verbatim}
\begin{lstlisting}[language=st]
IF Start = TRUE THEN
    Motor := TRUE;
ELSE
    Motor := FALSE;
END_IF;
\end{lstlisting}
\end{verbatim}

Only code necessary for explaining the work should appear in the main
report. Long source-code listings should be moved to an appendix or
submitted as supplementary files if required.

Source code should not normally be inserted as screenshots.

\vspace{0.3cm}

% ================================================================
% XVI. ABSTRACT
% ================================================================

\noindent
\textbf{XVI.\quad ABSTRACT}

\vspace{0.15cm}

The Abstract shall briefly summarize:

\begin{itemize}[leftmargin=0.8cm, itemsep=2pt]
    \item the internship organization;
    \item the objectives of the internship;
    \item the main tasks performed;
    \item the main results, knowledge, and skills obtained.
\end{itemize}

The Abstract should be concise and should not reproduce the full content
of the report chapter by chapter.

\vspace{0.3cm}

% ================================================================
% XVII. REFERENCES AND CITATIONS
% ================================================================

\noindent
\textbf{XVII.\quad REFERENCES AND CITATIONS}

\vspace{0.15cm}

All books, journal papers, conference papers, standards, websites,
datasheets, manuals, technical documents, figures, tables, or other
external sources used in preparing the report shall be properly cited.

References shall be managed through the bibliography system provided with
the official template.

Bibliographic entries shall normally be stored in:

\begin{verbatim}
References.bib
\end{verbatim}

References shall be cited in the report using:

\begin{verbatim}
\cite{reference-key}
\end{verbatim}

Students shall not manually type reference numbers such as ``[1]'',
``[2]'', etc.

Titles of foreign-language publications shall remain in their original
language in the bibliography.

All borrowed ideas, data, figures, tables, equations, methods, or results
shall be properly acknowledged.

\vspace{0.3cm}

% ================================================================
% XVIII. CROSS-REFERENCING
% ================================================================

\noindent
\textbf{XVIII.\quad CROSS-REFERENCING}

\vspace{0.15cm}

Numbered elements shall be cross-referenced using the standard \LaTeX{}
mechanism.

Use the following commands where appropriate:

\begin{verbatim}
\label{...}
\ref{...}
\eqref{...}
\end{verbatim}

For example:

\begin{verbatim}
As shown in Figure~\ref{fig:system}, ...
\end{verbatim}

Students shall not manually type numbers such as ``Figure 2.3'' when
those numbers can be generated automatically.

\vspace{0.3cm}

% ================================================================
% XIX. PAGE NUMBERING
% ================================================================

\noindent
\textbf{XIX.\quad PAGE NUMBERING}

\vspace{0.15cm}

Cover pages and official evaluation forms at the beginning of the report
shall not display page numbers.

The main body of the report shall use Arabic page numbers beginning from
Chapter 1.

Page numbering shall be generated automatically by the official template.
Students shall not manually enter page numbers.

\vspace{0.3cm}

% ================================================================
% XX. PROJECT FILE ORGANIZATION
% ================================================================

\noindent
\textbf{XX.\quad PROJECT FILE ORGANIZATION}

\vspace{0.15cm}

Students should maintain a clear and organized \LaTeX{} project
structure. A recommended structure is:

\begin{verbatim}
InternshipReport/
|
|-- main.tex
|-- References.bib
|
|-- figures/
|   |-- company.pdf
|   |-- system.png
|   `-- ...
|
`-- chapters/
    |-- chapter1.tex
    |-- chapter2.tex
    `-- chapter3.tex
\end{verbatim}

Figures should normally be stored in the \verb|figures| directory
specified by the official template.

File and directory names should preferably contain Latin letters,
numbers, hyphens, or underscores.

Recommended examples include:

\begin{verbatim}
company_structure.pdf
motor_control.png
chapter_01.tex
\end{verbatim}

Spaces and unnecessary special characters in file names should be
avoided where possible.

\vspace{0.3cm}

% ================================================================
% XXI. DOCUMENT COMPILATION AND FINAL CHECK
% ================================================================

\noindent
\textbf{XXI.\quad DOCUMENT COMPILATION AND FINAL CHECK}

\vspace{0.15cm}

Before submission, the complete \LaTeX{} project shall compile
successfully and produce the final PDF report.

The final document shall not contain unresolved references such as:

\begin{verbatim}
??
[?]
\end{verbatim}

Before generating the final PDF, students shall check:

\begin{itemize}[leftmargin=0.8cm, itemsep=2pt]
    \item the Table of Contents;
    \item page numbering;
    \item chapter, section, and subsection numbering;
    \item figure numbering;
    \item table numbering;
    \item equation numbering;
    \item cross-references;
    \item citations;
    \item bibliography;
    \item font display;
    \item mathematical symbols;
    \item image quality;
    \item completeness of all required forms.
\end{itemize}

The document may be compiled using TeX Live, MiKTeX, or another
compatible \LaTeX{} environment. Regardless of the compilation
environment, the resulting PDF shall conform to the official Faculty
template.

\vspace{0.3cm}

% ================================================================
% XXII. CONTENT THAT STUDENTS MAY EDIT
% ================================================================

\noindent
\textbf{XXII.\quad CONTENT THAT STUDENTS MAY EDIT}

\vspace{0.15cm}

Students may edit or replace:

\begin{itemize}[leftmargin=0.8cm, itemsep=2pt]
    \item student name;
    \item class;
    \item internship organization;
    \item supervisor name;
    \item Abstract;
    \item List of Symbols and Abbreviations;
    \item chapter titles and chapter content;
    \item figures and photographs;
    \item tables;
    \item equations;
    \item citations and references.
\end{itemize}

Students shall not modify the main document formatting settings in the
preamble, including the document class, page geometry, line spacing,
heading formatting, table-of-contents formatting, font configuration,
page layout, or cover-page macros unless specifically instructed.

\vspace{0.3cm}

% ================================================================
% XXIII. FINAL SUBMISSION REQUIREMENTS
% ================================================================

\noindent
\textbf{XXIII.\quad FINAL SUBMISSION REQUIREMENTS}

\vspace{0.15cm}

Before submission, the final report shall:

\begin{itemize}[leftmargin=0.8cm, itemsep=2pt]
    \item compile without errors that affect the report content or presentation;
    \item contain all required forms and report sections;
    \item contain correct student, class, internship organization, and supervisor information;
    \item contain an updated Table of Contents, List of Figures, and List of Tables;
    \item contain no unused example content from the template;
    \item contain no placeholder text such as \textit{Title of Chapter 1},
    \textit{Section 1.1}, or other sample entries;
    \item contain all required citations and references;
    \item contain no unresolved cross-references.
\end{itemize}

If the Faculty requires submission of the \LaTeX{} source, the complete
project shall include all \verb|.tex| files, the bibliography database,
figures, and any other files required to reproduce the submitted PDF.

\vspace{0.4cm}

\noindent
\textbf{Important:}
Before submitting the final internship report, students shall remove or
comment out the following command in the main document:

\begin{verbatim}
\input{Regulations.tex}
\end{verbatim}

For example:

\begin{verbatim}
% \input{Regulations.tex}
\end{verbatim}

This removes the guidance pages from the final submitted report while
preserving the regulations file in the \LaTeX{} project for reference.

} % End 13pt content

\clearpage

%
% IMPORTANT:
% This section is provided for guidance only. Students shall remove
% or comment out the 
% =================================================================
% REGULATIONS FOR PREPARING THE GRADUATION INTERNSHIP REPORT
% USING LATEX
% =================================================================
%
% This file is intended to be included from the main document with:
%
%     \input{Regulations.tex}
%
% IMPORTANT:
% This section is provided for guidance only. Students shall remove
% or comment out the \input{Regulations.tex} command before submitting
% the final internship report.
%
% =================================================================

\clearpage
\thispagestyle{empty}

\begin{center}
    {\fontsize{16pt}{20pt}\selectfont\bfseries
    SPECIFIC REGULATIONS FOR PREPARING\\
    GRADUATION INTERNSHIP REPORTS USING \LaTeX}
\end{center}

\vspace{0.5cm}

{\fontsize{13pt}{17pt}\selectfont
\setstretch{1.3}
\justifying

\noindent
\textbf{Note:}
This section is provided for guidance only and is included in the main
document using the command \verb|\input{Regulations.tex}|.
Students shall remove or comment out this command before submitting the
final internship report.

\vspace{0.4cm}

% ================================================================
% I. GENERAL REQUIREMENTS
% ================================================================

\noindent
\textbf{I.\quad GENERAL REQUIREMENTS}

\vspace{0.15cm}

The Graduation Internship Report shall be prepared using the official
\LaTeX{} template provided by the Faculty. Students shall preserve the
predefined page layout, typography, line spacing, heading styles,
numbering system, figure and table formatting, and bibliography settings
unless otherwise instructed by the Faculty or the supervisor.

The report shall be written in English.

Students are responsible for the academic and technical content of the
report. The official \LaTeX{} template is responsible for standardized
document formatting.

\vspace{0.3cm}

% ================================================================
% II. REPORT STRUCTURE
% ================================================================

\noindent
\textbf{II.\quad REPORT STRUCTURE}

\vspace{0.15cm}

The report shall normally be arranged in the following order:

\begin{enumerate}[label=\arabic*., leftmargin=1.2cm, itemsep=2pt]
    \item Outer Cover;
    \item Comments of the Internship Host Organization;
    \item Supervisor's Evaluation and Graduation Internship Defense Results;
    \item Table of Contents;
    \item List of Symbols and Abbreviations, if applicable;
    \item List of Figures;
    \item List of Tables;
    \item Abstract;
    \item Main Chapters of the Internship Report;
    \item References;
    \item Appendices, if applicable.
\end{enumerate}

The preliminary pages shall not be numbered as chapters. The main report
content shall be organized into consecutively numbered chapters beginning
with Chapter 1.

\vspace{0.3cm}

% ================================================================
% III. PAGE SIZE AND MARGINS
% ================================================================

\noindent
\textbf{III.\quad PAGE SIZE AND MARGINS}

\vspace{0.15cm}

The report shall be formatted on A4 paper with the following margins:

\begin{center}
\begin{tabular}{ll}
    Left margin:   & 3.5 cm \\
    Right margin:  & 2.0 cm \\
    Top margin:    & 2.5 cm \\
    Bottom margin: & 2.5 cm
\end{tabular}
\end{center}

These settings are predefined in the official template and shall not be
modified by students.

\vspace{0.3cm}

% ================================================================
% IV. FONT, FONT SIZE, AND LINE SPACING
% ================================================================

\noindent
\textbf{IV.\quad FONT, FONT SIZE, AND LINE SPACING}

\vspace{0.15cm}

The main text shall use a 13 pt font size with 1.3-line spacing.
A Times-family typeface shall be used in accordance with the official
template.

The text shall be fully justified.

Students shall not use repeated spaces, repeated manual line breaks, or
multiple blank lines to position text. Paragraph and heading spacing shall
be controlled by the template.

\vspace{0.3cm}

% ================================================================
% V. OUTER COVER
% ================================================================

\noindent
\textbf{V.\quad OUTER COVER}

\vspace{0.15cm}

The outer cover shall follow the official format provided in the template
and shall contain the following information:

\begin{itemize}[leftmargin=0.8cm, itemsep=2pt]
    \item Thai Nguyen University of Technology;
    \item Faculty of International Training;
    \item Engineering and Technology Department;
    \item University logo and Faculty logo;
    \item the title \textbf{INTERNSHIP REPORT};
    \item Internship Organization;
    \item Author;
    \item Class;
    \item Supervisor;
    \item place and year of completion.
\end{itemize}

Students shall only replace the required information fields and shall not
modify the cover-page layout unless instructed otherwise.

\vspace{0.3cm}

% ================================================================
% VI. COMMENTS OF THE INTERNSHIP HOST ORGANIZATION
% ================================================================

\noindent
\textbf{VI.\quad COMMENTS OF THE INTERNSHIP HOST ORGANIZATION}

\vspace{0.15cm}

The page entitled \textbf{COMMENTS OF THE INTERNSHIP HOST ORGANIZATION}
is reserved for the internship host organization.

It includes:

\begin{itemize}[leftmargin=0.8cm, itemsep=2pt]
    \item the student's full name;
    \item class;
    \item comments from the host organization;
    \item score;
    \item date;
    \item signature and official seal.
\end{itemize}

Students shall not alter or complete official evaluation fields unless
specifically authorized.

\vspace{0.3cm}

% ================================================================
% VII. SUPERVISOR'S EVALUATION AND DEFENSE RESULTS
% ================================================================

\noindent
\textbf{VII.\quad SUPERVISOR'S EVALUATION AND DEFENSE RESULTS}

\vspace{0.15cm}

The section entitled
\textbf{SUPERVISOR'S EVALUATION AND GRADUATION INTERNSHIP DEFENSE RESULTS}
is reserved for the supervisor and defense committee.

The evaluation may include:

\begin{itemize}[leftmargin=0.8cm, itemsep=2pt]
    \item Internship Progress and Student's Attitude;
    \item Report Content;
    \item Report Presentation;
    \item Other Comments;
    \item Supervisor's Opinion;
    \item Defense Result;
    \item Score.
\end{itemize}

Students shall not enter evaluation scores or defense results before they
are officially determined.

\vspace{0.3cm}

% ================================================================
% VIII. CHAPTERS, SECTIONS, AND SUBSECTIONS
% ================================================================

\noindent
\textbf{VIII.\quad CHAPTERS, SECTIONS, AND SUBSECTIONS}

\vspace{0.15cm}

The report structure shall be generated using the standard \LaTeX{}
commands:

\begin{verbatim}
\chapter{...}
\section{...}
\subsection{...}
\end{verbatim}

Chapter, section, and subsection numbers shall be generated automatically
by \LaTeX{}.

Students shall not manually type chapter or section numbers and shall not
reproduce heading styles using commands such as \verb|\textbf|,
\verb|\fontsize|, or \verb|\centering| instead of the appropriate
structural command.

Excessive subdivision should be avoided. In most cases, the hierarchy
Chapter--Section--Subsection is sufficient.

\vspace{0.3cm}

% ================================================================
% IX. ORGANIZATION OF THE MAIN CONTENT
% ================================================================

\noindent
\textbf{IX.\quad ORGANIZATION OF THE MAIN CONTENT}

\vspace{0.15cm}

The specific content of the report shall depend on the internship
organization and the assigned tasks. A typical structure may include:

\begin{itemize}[leftmargin=0.8cm, itemsep=2pt]
    \item \textbf{Chapter 1. Overview of the Internship Organization};
    \item \textbf{Chapter 2. Internship Activities and Assigned Tasks};
    \item \textbf{Chapter 3. Results, Experience, and Evaluation}.
\end{itemize}

These chapter titles are recommendations only and may be adapted to the
actual internship content.

The report shall reflect work actually performed, observed, analyzed, or
learned during the internship. Unrelated theoretical material should not
be included merely to increase report length.

\vspace{0.3cm}

% ================================================================
% X. TABLE OF CONTENTS
% ================================================================

\noindent
\textbf{X.\quad TABLE OF CONTENTS}

\vspace{0.15cm}

The Table of Contents shall be generated automatically using:

\begin{verbatim}
\tableofcontents
\end{verbatim}

Students shall not create or manually edit page numbers in the Table of
Contents.

The document shall be compiled sufficiently to ensure that all page
numbers and entries are correctly updated.

\vspace{0.3cm}

% ================================================================
% XI. LIST OF SYMBOLS AND ABBREVIATIONS
% ================================================================

\noindent
\textbf{XI.\quad LIST OF SYMBOLS AND ABBREVIATIONS}

\vspace{0.15cm}

Symbols and abbreviations that are used repeatedly in the report shall be
listed in the \textbf{List of Symbols and Abbreviations}.

The list may contain the following columns:

\begin{center}
\begin{tabular}{llll}
    No. & Symbol / Abbreviation & Unit & Definition
\end{tabular}
\end{center}

Only symbols and abbreviations actually used in the final report shall be
retained. Example entries contained in the template shall be removed if
they are not relevant to the student's report.

Each abbreviation should normally be defined when it first appears in
the main text.

\vspace{0.3cm}

% ================================================================
% XII. FIGURES, DIAGRAMS, AND PHOTOGRAPHS
% ================================================================

\noindent
\textbf{XII.\quad FIGURES, DIAGRAMS, AND PHOTOGRAPHS}

\vspace{0.15cm}

Figures shall be clear, legible, and directly relevant to the report
content.

Figures shall be inserted using the standard \verb|figure| environment.

Example:

\begin{verbatim}
\begin{figure}[H]
    \centering
    \includegraphics[width=0.8\textwidth]
                    {figures/company_structure.pdf}
    \caption{Organizational structure of the company}
    \label{fig:company-structure}
\end{figure}
\end{verbatim}

Figure captions shall be placed below the figure.

Each referenced figure shall have a unique \verb|\label| and shall be
referenced in the text using, for example:

\begin{verbatim}
Figure~\ref{fig:company-structure}
\end{verbatim}

Students shall not manually type figure numbers.

Figures reproduced or adapted from external sources shall include an
appropriate citation.

Photographs taken during the internship shall be relevant and sufficiently
clear. Students shall not include confidential, restricted, or proprietary
information without permission from the internship organization.

Vector graphics or high-resolution images are recommended whenever
possible.

\vspace{0.3cm}

% ================================================================
% XIII. TABLES
% ================================================================

\noindent
\textbf{XIII.\quad TABLES}

\vspace{0.15cm}

Tables shall be created using suitable \LaTeX{} environments such as
\verb|table|, \verb|tabular|, \verb|tabularx|, or \verb|longtable|.

Table captions shall be placed above the table.

Example:

\begin{verbatim}
\begin{table}[H]
    \centering
    \caption{Internship schedule}
    \label{tab:internship-schedule}

    \begin{tabular}{lll}
        \toprule
        Week & Task & Result \\
        \midrule
        1 & Orientation & Completed \\
        2 & System study & Completed \\
        \bottomrule
    \end{tabular}
\end{table}
\end{verbatim}

Tables shall be referenced using:

\begin{verbatim}
Table~\ref{tab:internship-schedule}
\end{verbatim}

Students shall not manually type table numbers.

Tables should be created directly in \LaTeX{} rather than inserted as
screenshots whenever practical.

\vspace{0.3cm}

% ================================================================
% XIV. EQUATIONS AND TECHNICAL QUANTITIES
% ================================================================

\noindent
\textbf{XIV.\quad EQUATIONS AND TECHNICAL QUANTITIES}

\vspace{0.15cm}

Mathematical expressions shall be typeset directly using \LaTeX{}.
Images or screenshots of equations shall not be used.

Short expressions may be written in inline mathematical mode, for example:

\begin{verbatim}
$R = 0.7\,\Omega$
\end{verbatim}

Important equations shall use a numbered equation environment.

Example:

\begin{verbatim}
\begin{equation}
    P = UI
    \label{eq:power}
\end{equation}
\end{verbatim}

Equations shall be referenced using:

\begin{verbatim}
Equation~\eqref{eq:power}
\end{verbatim}

Equation numbers shall not be manually entered.

SI units shall be used wherever applicable. The \verb|siunitx| package
should be used where appropriate.

Examples:

\begin{verbatim}
\SI{220}{\volt}
\SI{50}{\hertz}
\SI{2.5}{\ampere}
\end{verbatim}

Notation and units shall be consistent throughout the report.

\vspace{0.3cm}

% ================================================================
% XV. SOURCE CODE
% ================================================================

\noindent
\textbf{XV.\quad SOURCE CODE}

\vspace{0.15cm}

If the internship involves programming, PLCs, microcontrollers,
simulation, or software development, relevant source code may be included
in the report.

Source code shall be presented using the \verb|listings| environment or
another environment defined by the official template.

Example:

\begin{verbatim}
\begin{lstlisting}[language=st]
IF Start = TRUE THEN
    Motor := TRUE;
ELSE
    Motor := FALSE;
END_IF;
\end{lstlisting}
\end{verbatim}

Only code necessary for explaining the work should appear in the main
report. Long source-code listings should be moved to an appendix or
submitted as supplementary files if required.

Source code should not normally be inserted as screenshots.

\vspace{0.3cm}

% ================================================================
% XVI. ABSTRACT
% ================================================================

\noindent
\textbf{XVI.\quad ABSTRACT}

\vspace{0.15cm}

The Abstract shall briefly summarize:

\begin{itemize}[leftmargin=0.8cm, itemsep=2pt]
    \item the internship organization;
    \item the objectives of the internship;
    \item the main tasks performed;
    \item the main results, knowledge, and skills obtained.
\end{itemize}

The Abstract should be concise and should not reproduce the full content
of the report chapter by chapter.

\vspace{0.3cm}

% ================================================================
% XVII. REFERENCES AND CITATIONS
% ================================================================

\noindent
\textbf{XVII.\quad REFERENCES AND CITATIONS}

\vspace{0.15cm}

All books, journal papers, conference papers, standards, websites,
datasheets, manuals, technical documents, figures, tables, or other
external sources used in preparing the report shall be properly cited.

References shall be managed through the bibliography system provided with
the official template.

Bibliographic entries shall normally be stored in:

\begin{verbatim}
References.bib
\end{verbatim}

References shall be cited in the report using:

\begin{verbatim}
\cite{reference-key}
\end{verbatim}

Students shall not manually type reference numbers such as ``[1]'',
``[2]'', etc.

Titles of foreign-language publications shall remain in their original
language in the bibliography.

All borrowed ideas, data, figures, tables, equations, methods, or results
shall be properly acknowledged.

\vspace{0.3cm}

% ================================================================
% XVIII. CROSS-REFERENCING
% ================================================================

\noindent
\textbf{XVIII.\quad CROSS-REFERENCING}

\vspace{0.15cm}

Numbered elements shall be cross-referenced using the standard \LaTeX{}
mechanism.

Use the following commands where appropriate:

\begin{verbatim}
\label{...}
\ref{...}
\eqref{...}
\end{verbatim}

For example:

\begin{verbatim}
As shown in Figure~\ref{fig:system}, ...
\end{verbatim}

Students shall not manually type numbers such as ``Figure 2.3'' when
those numbers can be generated automatically.

\vspace{0.3cm}

% ================================================================
% XIX. PAGE NUMBERING
% ================================================================

\noindent
\textbf{XIX.\quad PAGE NUMBERING}

\vspace{0.15cm}

Cover pages and official evaluation forms at the beginning of the report
shall not display page numbers.

The main body of the report shall use Arabic page numbers beginning from
Chapter 1.

Page numbering shall be generated automatically by the official template.
Students shall not manually enter page numbers.

\vspace{0.3cm}

% ================================================================
% XX. PROJECT FILE ORGANIZATION
% ================================================================

\noindent
\textbf{XX.\quad PROJECT FILE ORGANIZATION}

\vspace{0.15cm}

Students should maintain a clear and organized \LaTeX{} project
structure. A recommended structure is:

\begin{verbatim}
InternshipReport/
|
|-- main.tex
|-- References.bib
|
|-- figures/
|   |-- company.pdf
|   |-- system.png
|   `-- ...
|
`-- chapters/
    |-- chapter1.tex
    |-- chapter2.tex
    `-- chapter3.tex
\end{verbatim}

Figures should normally be stored in the \verb|figures| directory
specified by the official template.

File and directory names should preferably contain Latin letters,
numbers, hyphens, or underscores.

Recommended examples include:

\begin{verbatim}
company_structure.pdf
motor_control.png
chapter_01.tex
\end{verbatim}

Spaces and unnecessary special characters in file names should be
avoided where possible.

\vspace{0.3cm}

% ================================================================
% XXI. DOCUMENT COMPILATION AND FINAL CHECK
% ================================================================

\noindent
\textbf{XXI.\quad DOCUMENT COMPILATION AND FINAL CHECK}

\vspace{0.15cm}

Before submission, the complete \LaTeX{} project shall compile
successfully and produce the final PDF report.

The final document shall not contain unresolved references such as:

\begin{verbatim}
??
[?]
\end{verbatim}

Before generating the final PDF, students shall check:

\begin{itemize}[leftmargin=0.8cm, itemsep=2pt]
    \item the Table of Contents;
    \item page numbering;
    \item chapter, section, and subsection numbering;
    \item figure numbering;
    \item table numbering;
    \item equation numbering;
    \item cross-references;
    \item citations;
    \item bibliography;
    \item font display;
    \item mathematical symbols;
    \item image quality;
    \item completeness of all required forms.
\end{itemize}

The document may be compiled using TeX Live, MiKTeX, or another
compatible \LaTeX{} environment. Regardless of the compilation
environment, the resulting PDF shall conform to the official Faculty
template.

\vspace{0.3cm}

% ================================================================
% XXII. CONTENT THAT STUDENTS MAY EDIT
% ================================================================

\noindent
\textbf{XXII.\quad CONTENT THAT STUDENTS MAY EDIT}

\vspace{0.15cm}

Students may edit or replace:

\begin{itemize}[leftmargin=0.8cm, itemsep=2pt]
    \item student name;
    \item class;
    \item internship organization;
    \item supervisor name;
    \item Abstract;
    \item List of Symbols and Abbreviations;
    \item chapter titles and chapter content;
    \item figures and photographs;
    \item tables;
    \item equations;
    \item citations and references.
\end{itemize}

Students shall not modify the main document formatting settings in the
preamble, including the document class, page geometry, line spacing,
heading formatting, table-of-contents formatting, font configuration,
page layout, or cover-page macros unless specifically instructed.

\vspace{0.3cm}

% ================================================================
% XXIII. FINAL SUBMISSION REQUIREMENTS
% ================================================================

\noindent
\textbf{XXIII.\quad FINAL SUBMISSION REQUIREMENTS}

\vspace{0.15cm}

Before submission, the final report shall:

\begin{itemize}[leftmargin=0.8cm, itemsep=2pt]
    \item compile without errors that affect the report content or presentation;
    \item contain all required forms and report sections;
    \item contain correct student, class, internship organization, and supervisor information;
    \item contain an updated Table of Contents, List of Figures, and List of Tables;
    \item contain no unused example content from the template;
    \item contain no placeholder text such as \textit{Title of Chapter 1},
    \textit{Section 1.1}, or other sample entries;
    \item contain all required citations and references;
    \item contain no unresolved cross-references.
\end{itemize}

If the Faculty requires submission of the \LaTeX{} source, the complete
project shall include all \verb|.tex| files, the bibliography database,
figures, and any other files required to reproduce the submitted PDF.

\vspace{0.4cm}

\noindent
\textbf{Important:}
Before submitting the final internship report, students shall remove or
comment out the following command in the main document:

\begin{verbatim}
\input{Regulations.tex}
\end{verbatim}

For example:

\begin{verbatim}
% \input{Regulations.tex}
\end{verbatim}

This removes the guidance pages from the final submitted report while
preserving the regulations file in the \LaTeX{} project for reference.

} % End 13pt content

\clearpage
 command before submitting
% the final internship report.
%
% =================================================================

\clearpage
\thispagestyle{empty}

\begin{center}
    {\fontsize{16pt}{20pt}\selectfont\bfseries
    SPECIFIC REGULATIONS FOR PREPARING\\
    GRADUATION INTERNSHIP REPORTS USING \LaTeX}
\end{center}

\vspace{0.5cm}

{\fontsize{13pt}{17pt}\selectfont
\setstretch{1.3}
\justifying

\noindent
\textbf{Note:}
This section is provided for guidance only and is included in the main
document using the command \verb|
% =================================================================
% REGULATIONS FOR PREPARING THE GRADUATION INTERNSHIP REPORT
% USING LATEX
% =================================================================
%
% This file is intended to be included from the main document with:
%
%     \input{Regulations.tex}
%
% IMPORTANT:
% This section is provided for guidance only. Students shall remove
% or comment out the \input{Regulations.tex} command before submitting
% the final internship report.
%
% =================================================================

\clearpage
\thispagestyle{empty}

\begin{center}
    {\fontsize{16pt}{20pt}\selectfont\bfseries
    SPECIFIC REGULATIONS FOR PREPARING\\
    GRADUATION INTERNSHIP REPORTS USING \LaTeX}
\end{center}

\vspace{0.5cm}

{\fontsize{13pt}{17pt}\selectfont
\setstretch{1.3}
\justifying

\noindent
\textbf{Note:}
This section is provided for guidance only and is included in the main
document using the command \verb|\input{Regulations.tex}|.
Students shall remove or comment out this command before submitting the
final internship report.

\vspace{0.4cm}

% ================================================================
% I. GENERAL REQUIREMENTS
% ================================================================

\noindent
\textbf{I.\quad GENERAL REQUIREMENTS}

\vspace{0.15cm}

The Graduation Internship Report shall be prepared using the official
\LaTeX{} template provided by the Faculty. Students shall preserve the
predefined page layout, typography, line spacing, heading styles,
numbering system, figure and table formatting, and bibliography settings
unless otherwise instructed by the Faculty or the supervisor.

The report shall be written in English.

Students are responsible for the academic and technical content of the
report. The official \LaTeX{} template is responsible for standardized
document formatting.

\vspace{0.3cm}

% ================================================================
% II. REPORT STRUCTURE
% ================================================================

\noindent
\textbf{II.\quad REPORT STRUCTURE}

\vspace{0.15cm}

The report shall normally be arranged in the following order:

\begin{enumerate}[label=\arabic*., leftmargin=1.2cm, itemsep=2pt]
    \item Outer Cover;
    \item Comments of the Internship Host Organization;
    \item Supervisor's Evaluation and Graduation Internship Defense Results;
    \item Table of Contents;
    \item List of Symbols and Abbreviations, if applicable;
    \item List of Figures;
    \item List of Tables;
    \item Abstract;
    \item Main Chapters of the Internship Report;
    \item References;
    \item Appendices, if applicable.
\end{enumerate}

The preliminary pages shall not be numbered as chapters. The main report
content shall be organized into consecutively numbered chapters beginning
with Chapter 1.

\vspace{0.3cm}

% ================================================================
% III. PAGE SIZE AND MARGINS
% ================================================================

\noindent
\textbf{III.\quad PAGE SIZE AND MARGINS}

\vspace{0.15cm}

The report shall be formatted on A4 paper with the following margins:

\begin{center}
\begin{tabular}{ll}
    Left margin:   & 3.5 cm \\
    Right margin:  & 2.0 cm \\
    Top margin:    & 2.5 cm \\
    Bottom margin: & 2.5 cm
\end{tabular}
\end{center}

These settings are predefined in the official template and shall not be
modified by students.

\vspace{0.3cm}

% ================================================================
% IV. FONT, FONT SIZE, AND LINE SPACING
% ================================================================

\noindent
\textbf{IV.\quad FONT, FONT SIZE, AND LINE SPACING}

\vspace{0.15cm}

The main text shall use a 13 pt font size with 1.3-line spacing.
A Times-family typeface shall be used in accordance with the official
template.

The text shall be fully justified.

Students shall not use repeated spaces, repeated manual line breaks, or
multiple blank lines to position text. Paragraph and heading spacing shall
be controlled by the template.

\vspace{0.3cm}

% ================================================================
% V. OUTER COVER
% ================================================================

\noindent
\textbf{V.\quad OUTER COVER}

\vspace{0.15cm}

The outer cover shall follow the official format provided in the template
and shall contain the following information:

\begin{itemize}[leftmargin=0.8cm, itemsep=2pt]
    \item Thai Nguyen University of Technology;
    \item Faculty of International Training;
    \item Engineering and Technology Department;
    \item University logo and Faculty logo;
    \item the title \textbf{INTERNSHIP REPORT};
    \item Internship Organization;
    \item Author;
    \item Class;
    \item Supervisor;
    \item place and year of completion.
\end{itemize}

Students shall only replace the required information fields and shall not
modify the cover-page layout unless instructed otherwise.

\vspace{0.3cm}

% ================================================================
% VI. COMMENTS OF THE INTERNSHIP HOST ORGANIZATION
% ================================================================

\noindent
\textbf{VI.\quad COMMENTS OF THE INTERNSHIP HOST ORGANIZATION}

\vspace{0.15cm}

The page entitled \textbf{COMMENTS OF THE INTERNSHIP HOST ORGANIZATION}
is reserved for the internship host organization.

It includes:

\begin{itemize}[leftmargin=0.8cm, itemsep=2pt]
    \item the student's full name;
    \item class;
    \item comments from the host organization;
    \item score;
    \item date;
    \item signature and official seal.
\end{itemize}

Students shall not alter or complete official evaluation fields unless
specifically authorized.

\vspace{0.3cm}

% ================================================================
% VII. SUPERVISOR'S EVALUATION AND DEFENSE RESULTS
% ================================================================

\noindent
\textbf{VII.\quad SUPERVISOR'S EVALUATION AND DEFENSE RESULTS}

\vspace{0.15cm}

The section entitled
\textbf{SUPERVISOR'S EVALUATION AND GRADUATION INTERNSHIP DEFENSE RESULTS}
is reserved for the supervisor and defense committee.

The evaluation may include:

\begin{itemize}[leftmargin=0.8cm, itemsep=2pt]
    \item Internship Progress and Student's Attitude;
    \item Report Content;
    \item Report Presentation;
    \item Other Comments;
    \item Supervisor's Opinion;
    \item Defense Result;
    \item Score.
\end{itemize}

Students shall not enter evaluation scores or defense results before they
are officially determined.

\vspace{0.3cm}

% ================================================================
% VIII. CHAPTERS, SECTIONS, AND SUBSECTIONS
% ================================================================

\noindent
\textbf{VIII.\quad CHAPTERS, SECTIONS, AND SUBSECTIONS}

\vspace{0.15cm}

The report structure shall be generated using the standard \LaTeX{}
commands:

\begin{verbatim}
\chapter{...}
\section{...}
\subsection{...}
\end{verbatim}

Chapter, section, and subsection numbers shall be generated automatically
by \LaTeX{}.

Students shall not manually type chapter or section numbers and shall not
reproduce heading styles using commands such as \verb|\textbf|,
\verb|\fontsize|, or \verb|\centering| instead of the appropriate
structural command.

Excessive subdivision should be avoided. In most cases, the hierarchy
Chapter--Section--Subsection is sufficient.

\vspace{0.3cm}

% ================================================================
% IX. ORGANIZATION OF THE MAIN CONTENT
% ================================================================

\noindent
\textbf{IX.\quad ORGANIZATION OF THE MAIN CONTENT}

\vspace{0.15cm}

The specific content of the report shall depend on the internship
organization and the assigned tasks. A typical structure may include:

\begin{itemize}[leftmargin=0.8cm, itemsep=2pt]
    \item \textbf{Chapter 1. Overview of the Internship Organization};
    \item \textbf{Chapter 2. Internship Activities and Assigned Tasks};
    \item \textbf{Chapter 3. Results, Experience, and Evaluation}.
\end{itemize}

These chapter titles are recommendations only and may be adapted to the
actual internship content.

The report shall reflect work actually performed, observed, analyzed, or
learned during the internship. Unrelated theoretical material should not
be included merely to increase report length.

\vspace{0.3cm}

% ================================================================
% X. TABLE OF CONTENTS
% ================================================================

\noindent
\textbf{X.\quad TABLE OF CONTENTS}

\vspace{0.15cm}

The Table of Contents shall be generated automatically using:

\begin{verbatim}
\tableofcontents
\end{verbatim}

Students shall not create or manually edit page numbers in the Table of
Contents.

The document shall be compiled sufficiently to ensure that all page
numbers and entries are correctly updated.

\vspace{0.3cm}

% ================================================================
% XI. LIST OF SYMBOLS AND ABBREVIATIONS
% ================================================================

\noindent
\textbf{XI.\quad LIST OF SYMBOLS AND ABBREVIATIONS}

\vspace{0.15cm}

Symbols and abbreviations that are used repeatedly in the report shall be
listed in the \textbf{List of Symbols and Abbreviations}.

The list may contain the following columns:

\begin{center}
\begin{tabular}{llll}
    No. & Symbol / Abbreviation & Unit & Definition
\end{tabular}
\end{center}

Only symbols and abbreviations actually used in the final report shall be
retained. Example entries contained in the template shall be removed if
they are not relevant to the student's report.

Each abbreviation should normally be defined when it first appears in
the main text.

\vspace{0.3cm}

% ================================================================
% XII. FIGURES, DIAGRAMS, AND PHOTOGRAPHS
% ================================================================

\noindent
\textbf{XII.\quad FIGURES, DIAGRAMS, AND PHOTOGRAPHS}

\vspace{0.15cm}

Figures shall be clear, legible, and directly relevant to the report
content.

Figures shall be inserted using the standard \verb|figure| environment.

Example:

\begin{verbatim}
\begin{figure}[H]
    \centering
    \includegraphics[width=0.8\textwidth]
                    {figures/company_structure.pdf}
    \caption{Organizational structure of the company}
    \label{fig:company-structure}
\end{figure}
\end{verbatim}

Figure captions shall be placed below the figure.

Each referenced figure shall have a unique \verb|\label| and shall be
referenced in the text using, for example:

\begin{verbatim}
Figure~\ref{fig:company-structure}
\end{verbatim}

Students shall not manually type figure numbers.

Figures reproduced or adapted from external sources shall include an
appropriate citation.

Photographs taken during the internship shall be relevant and sufficiently
clear. Students shall not include confidential, restricted, or proprietary
information without permission from the internship organization.

Vector graphics or high-resolution images are recommended whenever
possible.

\vspace{0.3cm}

% ================================================================
% XIII. TABLES
% ================================================================

\noindent
\textbf{XIII.\quad TABLES}

\vspace{0.15cm}

Tables shall be created using suitable \LaTeX{} environments such as
\verb|table|, \verb|tabular|, \verb|tabularx|, or \verb|longtable|.

Table captions shall be placed above the table.

Example:

\begin{verbatim}
\begin{table}[H]
    \centering
    \caption{Internship schedule}
    \label{tab:internship-schedule}

    \begin{tabular}{lll}
        \toprule
        Week & Task & Result \\
        \midrule
        1 & Orientation & Completed \\
        2 & System study & Completed \\
        \bottomrule
    \end{tabular}
\end{table}
\end{verbatim}

Tables shall be referenced using:

\begin{verbatim}
Table~\ref{tab:internship-schedule}
\end{verbatim}

Students shall not manually type table numbers.

Tables should be created directly in \LaTeX{} rather than inserted as
screenshots whenever practical.

\vspace{0.3cm}

% ================================================================
% XIV. EQUATIONS AND TECHNICAL QUANTITIES
% ================================================================

\noindent
\textbf{XIV.\quad EQUATIONS AND TECHNICAL QUANTITIES}

\vspace{0.15cm}

Mathematical expressions shall be typeset directly using \LaTeX{}.
Images or screenshots of equations shall not be used.

Short expressions may be written in inline mathematical mode, for example:

\begin{verbatim}
$R = 0.7\,\Omega$
\end{verbatim}

Important equations shall use a numbered equation environment.

Example:

\begin{verbatim}
\begin{equation}
    P = UI
    \label{eq:power}
\end{equation}
\end{verbatim}

Equations shall be referenced using:

\begin{verbatim}
Equation~\eqref{eq:power}
\end{verbatim}

Equation numbers shall not be manually entered.

SI units shall be used wherever applicable. The \verb|siunitx| package
should be used where appropriate.

Examples:

\begin{verbatim}
\SI{220}{\volt}
\SI{50}{\hertz}
\SI{2.5}{\ampere}
\end{verbatim}

Notation and units shall be consistent throughout the report.

\vspace{0.3cm}

% ================================================================
% XV. SOURCE CODE
% ================================================================

\noindent
\textbf{XV.\quad SOURCE CODE}

\vspace{0.15cm}

If the internship involves programming, PLCs, microcontrollers,
simulation, or software development, relevant source code may be included
in the report.

Source code shall be presented using the \verb|listings| environment or
another environment defined by the official template.

Example:

\begin{verbatim}
\begin{lstlisting}[language=st]
IF Start = TRUE THEN
    Motor := TRUE;
ELSE
    Motor := FALSE;
END_IF;
\end{lstlisting}
\end{verbatim}

Only code necessary for explaining the work should appear in the main
report. Long source-code listings should be moved to an appendix or
submitted as supplementary files if required.

Source code should not normally be inserted as screenshots.

\vspace{0.3cm}

% ================================================================
% XVI. ABSTRACT
% ================================================================

\noindent
\textbf{XVI.\quad ABSTRACT}

\vspace{0.15cm}

The Abstract shall briefly summarize:

\begin{itemize}[leftmargin=0.8cm, itemsep=2pt]
    \item the internship organization;
    \item the objectives of the internship;
    \item the main tasks performed;
    \item the main results, knowledge, and skills obtained.
\end{itemize}

The Abstract should be concise and should not reproduce the full content
of the report chapter by chapter.

\vspace{0.3cm}

% ================================================================
% XVII. REFERENCES AND CITATIONS
% ================================================================

\noindent
\textbf{XVII.\quad REFERENCES AND CITATIONS}

\vspace{0.15cm}

All books, journal papers, conference papers, standards, websites,
datasheets, manuals, technical documents, figures, tables, or other
external sources used in preparing the report shall be properly cited.

References shall be managed through the bibliography system provided with
the official template.

Bibliographic entries shall normally be stored in:

\begin{verbatim}
References.bib
\end{verbatim}

References shall be cited in the report using:

\begin{verbatim}
\cite{reference-key}
\end{verbatim}

Students shall not manually type reference numbers such as ``[1]'',
``[2]'', etc.

Titles of foreign-language publications shall remain in their original
language in the bibliography.

All borrowed ideas, data, figures, tables, equations, methods, or results
shall be properly acknowledged.

\vspace{0.3cm}

% ================================================================
% XVIII. CROSS-REFERENCING
% ================================================================

\noindent
\textbf{XVIII.\quad CROSS-REFERENCING}

\vspace{0.15cm}

Numbered elements shall be cross-referenced using the standard \LaTeX{}
mechanism.

Use the following commands where appropriate:

\begin{verbatim}
\label{...}
\ref{...}
\eqref{...}
\end{verbatim}

For example:

\begin{verbatim}
As shown in Figure~\ref{fig:system}, ...
\end{verbatim}

Students shall not manually type numbers such as ``Figure 2.3'' when
those numbers can be generated automatically.

\vspace{0.3cm}

% ================================================================
% XIX. PAGE NUMBERING
% ================================================================

\noindent
\textbf{XIX.\quad PAGE NUMBERING}

\vspace{0.15cm}

Cover pages and official evaluation forms at the beginning of the report
shall not display page numbers.

The main body of the report shall use Arabic page numbers beginning from
Chapter 1.

Page numbering shall be generated automatically by the official template.
Students shall not manually enter page numbers.

\vspace{0.3cm}

% ================================================================
% XX. PROJECT FILE ORGANIZATION
% ================================================================

\noindent
\textbf{XX.\quad PROJECT FILE ORGANIZATION}

\vspace{0.15cm}

Students should maintain a clear and organized \LaTeX{} project
structure. A recommended structure is:

\begin{verbatim}
InternshipReport/
|
|-- main.tex
|-- References.bib
|
|-- figures/
|   |-- company.pdf
|   |-- system.png
|   `-- ...
|
`-- chapters/
    |-- chapter1.tex
    |-- chapter2.tex
    `-- chapter3.tex
\end{verbatim}

Figures should normally be stored in the \verb|figures| directory
specified by the official template.

File and directory names should preferably contain Latin letters,
numbers, hyphens, or underscores.

Recommended examples include:

\begin{verbatim}
company_structure.pdf
motor_control.png
chapter_01.tex
\end{verbatim}

Spaces and unnecessary special characters in file names should be
avoided where possible.

\vspace{0.3cm}

% ================================================================
% XXI. DOCUMENT COMPILATION AND FINAL CHECK
% ================================================================

\noindent
\textbf{XXI.\quad DOCUMENT COMPILATION AND FINAL CHECK}

\vspace{0.15cm}

Before submission, the complete \LaTeX{} project shall compile
successfully and produce the final PDF report.

The final document shall not contain unresolved references such as:

\begin{verbatim}
??
[?]
\end{verbatim}

Before generating the final PDF, students shall check:

\begin{itemize}[leftmargin=0.8cm, itemsep=2pt]
    \item the Table of Contents;
    \item page numbering;
    \item chapter, section, and subsection numbering;
    \item figure numbering;
    \item table numbering;
    \item equation numbering;
    \item cross-references;
    \item citations;
    \item bibliography;
    \item font display;
    \item mathematical symbols;
    \item image quality;
    \item completeness of all required forms.
\end{itemize}

The document may be compiled using TeX Live, MiKTeX, or another
compatible \LaTeX{} environment. Regardless of the compilation
environment, the resulting PDF shall conform to the official Faculty
template.

\vspace{0.3cm}

% ================================================================
% XXII. CONTENT THAT STUDENTS MAY EDIT
% ================================================================

\noindent
\textbf{XXII.\quad CONTENT THAT STUDENTS MAY EDIT}

\vspace{0.15cm}

Students may edit or replace:

\begin{itemize}[leftmargin=0.8cm, itemsep=2pt]
    \item student name;
    \item class;
    \item internship organization;
    \item supervisor name;
    \item Abstract;
    \item List of Symbols and Abbreviations;
    \item chapter titles and chapter content;
    \item figures and photographs;
    \item tables;
    \item equations;
    \item citations and references.
\end{itemize}

Students shall not modify the main document formatting settings in the
preamble, including the document class, page geometry, line spacing,
heading formatting, table-of-contents formatting, font configuration,
page layout, or cover-page macros unless specifically instructed.

\vspace{0.3cm}

% ================================================================
% XXIII. FINAL SUBMISSION REQUIREMENTS
% ================================================================

\noindent
\textbf{XXIII.\quad FINAL SUBMISSION REQUIREMENTS}

\vspace{0.15cm}

Before submission, the final report shall:

\begin{itemize}[leftmargin=0.8cm, itemsep=2pt]
    \item compile without errors that affect the report content or presentation;
    \item contain all required forms and report sections;
    \item contain correct student, class, internship organization, and supervisor information;
    \item contain an updated Table of Contents, List of Figures, and List of Tables;
    \item contain no unused example content from the template;
    \item contain no placeholder text such as \textit{Title of Chapter 1},
    \textit{Section 1.1}, or other sample entries;
    \item contain all required citations and references;
    \item contain no unresolved cross-references.
\end{itemize}

If the Faculty requires submission of the \LaTeX{} source, the complete
project shall include all \verb|.tex| files, the bibliography database,
figures, and any other files required to reproduce the submitted PDF.

\vspace{0.4cm}

\noindent
\textbf{Important:}
Before submitting the final internship report, students shall remove or
comment out the following command in the main document:

\begin{verbatim}
\input{Regulations.tex}
\end{verbatim}

For example:

\begin{verbatim}
% \input{Regulations.tex}
\end{verbatim}

This removes the guidance pages from the final submitted report while
preserving the regulations file in the \LaTeX{} project for reference.

} % End 13pt content

\clearpage
|.
Students shall remove or comment out this command before submitting the
final internship report.

\vspace{0.4cm}

% ================================================================
% I. GENERAL REQUIREMENTS
% ================================================================

\noindent
\textbf{I.\quad GENERAL REQUIREMENTS}

\vspace{0.15cm}

The Graduation Internship Report shall be prepared using the official
\LaTeX{} template provided by the Faculty. Students shall preserve the
predefined page layout, typography, line spacing, heading styles,
numbering system, figure and table formatting, and bibliography settings
unless otherwise instructed by the Faculty or the supervisor.

The report shall be written in English.

Students are responsible for the academic and technical content of the
report. The official \LaTeX{} template is responsible for standardized
document formatting.

\vspace{0.3cm}

% ================================================================
% II. REPORT STRUCTURE
% ================================================================

\noindent
\textbf{II.\quad REPORT STRUCTURE}

\vspace{0.15cm}

The report shall normally be arranged in the following order:

\begin{enumerate}[label=\arabic*., leftmargin=1.2cm, itemsep=2pt]
    \item Outer Cover;
    \item Comments of the Internship Host Organization;
    \item Supervisor's Evaluation and Graduation Internship Defense Results;
    \item Table of Contents;
    \item List of Symbols and Abbreviations, if applicable;
    \item List of Figures;
    \item List of Tables;
    \item Abstract;
    \item Main Chapters of the Internship Report;
    \item References;
    \item Appendices, if applicable.
\end{enumerate}

The preliminary pages shall not be numbered as chapters. The main report
content shall be organized into consecutively numbered chapters beginning
with Chapter 1.

\vspace{0.3cm}

% ================================================================
% III. PAGE SIZE AND MARGINS
% ================================================================

\noindent
\textbf{III.\quad PAGE SIZE AND MARGINS}

\vspace{0.15cm}

The report shall be formatted on A4 paper with the following margins:

\begin{center}
\begin{tabular}{ll}
    Left margin:   & 3.5 cm \\
    Right margin:  & 2.0 cm \\
    Top margin:    & 2.5 cm \\
    Bottom margin: & 2.5 cm
\end{tabular}
\end{center}

These settings are predefined in the official template and shall not be
modified by students.

\vspace{0.3cm}

% ================================================================
% IV. FONT, FONT SIZE, AND LINE SPACING
% ================================================================

\noindent
\textbf{IV.\quad FONT, FONT SIZE, AND LINE SPACING}

\vspace{0.15cm}

The main text shall use a 13 pt font size with 1.3-line spacing.
A Times-family typeface shall be used in accordance with the official
template.

The text shall be fully justified.

Students shall not use repeated spaces, repeated manual line breaks, or
multiple blank lines to position text. Paragraph and heading spacing shall
be controlled by the template.

\vspace{0.3cm}

% ================================================================
% V. OUTER COVER
% ================================================================

\noindent
\textbf{V.\quad OUTER COVER}

\vspace{0.15cm}

The outer cover shall follow the official format provided in the template
and shall contain the following information:

\begin{itemize}[leftmargin=0.8cm, itemsep=2pt]
    \item Thai Nguyen University of Technology;
    \item Faculty of International Training;
    \item Engineering and Technology Department;
    \item University logo and Faculty logo;
    \item the title \textbf{INTERNSHIP REPORT};
    \item Internship Organization;
    \item Author;
    \item Class;
    \item Supervisor;
    \item place and year of completion.
\end{itemize}

Students shall only replace the required information fields and shall not
modify the cover-page layout unless instructed otherwise.

\vspace{0.3cm}

% ================================================================
% VI. COMMENTS OF THE INTERNSHIP HOST ORGANIZATION
% ================================================================

\noindent
\textbf{VI.\quad COMMENTS OF THE INTERNSHIP HOST ORGANIZATION}

\vspace{0.15cm}

The page entitled \textbf{COMMENTS OF THE INTERNSHIP HOST ORGANIZATION}
is reserved for the internship host organization.

It includes:

\begin{itemize}[leftmargin=0.8cm, itemsep=2pt]
    \item the student's full name;
    \item class;
    \item comments from the host organization;
    \item score;
    \item date;
    \item signature and official seal.
\end{itemize}

Students shall not alter or complete official evaluation fields unless
specifically authorized.

\vspace{0.3cm}

% ================================================================
% VII. SUPERVISOR'S EVALUATION AND DEFENSE RESULTS
% ================================================================

\noindent
\textbf{VII.\quad SUPERVISOR'S EVALUATION AND DEFENSE RESULTS}

\vspace{0.15cm}

The section entitled
\textbf{SUPERVISOR'S EVALUATION AND GRADUATION INTERNSHIP DEFENSE RESULTS}
is reserved for the supervisor and defense committee.

The evaluation may include:

\begin{itemize}[leftmargin=0.8cm, itemsep=2pt]
    \item Internship Progress and Student's Attitude;
    \item Report Content;
    \item Report Presentation;
    \item Other Comments;
    \item Supervisor's Opinion;
    \item Defense Result;
    \item Score.
\end{itemize}

Students shall not enter evaluation scores or defense results before they
are officially determined.

\vspace{0.3cm}

% ================================================================
% VIII. CHAPTERS, SECTIONS, AND SUBSECTIONS
% ================================================================

\noindent
\textbf{VIII.\quad CHAPTERS, SECTIONS, AND SUBSECTIONS}

\vspace{0.15cm}

The report structure shall be generated using the standard \LaTeX{}
commands:

\begin{verbatim}
\chapter{...}
\section{...}
\subsection{...}
\end{verbatim}

Chapter, section, and subsection numbers shall be generated automatically
by \LaTeX{}.

Students shall not manually type chapter or section numbers and shall not
reproduce heading styles using commands such as \verb|\textbf|,
\verb|\fontsize|, or \verb|\centering| instead of the appropriate
structural command.

Excessive subdivision should be avoided. In most cases, the hierarchy
Chapter--Section--Subsection is sufficient.

\vspace{0.3cm}

% ================================================================
% IX. ORGANIZATION OF THE MAIN CONTENT
% ================================================================

\noindent
\textbf{IX.\quad ORGANIZATION OF THE MAIN CONTENT}

\vspace{0.15cm}

The specific content of the report shall depend on the internship
organization and the assigned tasks. A typical structure may include:

\begin{itemize}[leftmargin=0.8cm, itemsep=2pt]
    \item \textbf{Chapter 1. Overview of the Internship Organization};
    \item \textbf{Chapter 2. Internship Activities and Assigned Tasks};
    \item \textbf{Chapter 3. Results, Experience, and Evaluation}.
\end{itemize}

These chapter titles are recommendations only and may be adapted to the
actual internship content.

The report shall reflect work actually performed, observed, analyzed, or
learned during the internship. Unrelated theoretical material should not
be included merely to increase report length.

\vspace{0.3cm}

% ================================================================
% X. TABLE OF CONTENTS
% ================================================================

\noindent
\textbf{X.\quad TABLE OF CONTENTS}

\vspace{0.15cm}

The Table of Contents shall be generated automatically using:

\begin{verbatim}
\tableofcontents
\end{verbatim}

Students shall not create or manually edit page numbers in the Table of
Contents.

The document shall be compiled sufficiently to ensure that all page
numbers and entries are correctly updated.

\vspace{0.3cm}

% ================================================================
% XI. LIST OF SYMBOLS AND ABBREVIATIONS
% ================================================================

\noindent
\textbf{XI.\quad LIST OF SYMBOLS AND ABBREVIATIONS}

\vspace{0.15cm}

Symbols and abbreviations that are used repeatedly in the report shall be
listed in the \textbf{List of Symbols and Abbreviations}.

The list may contain the following columns:

\begin{center}
\begin{tabular}{llll}
    No. & Symbol / Abbreviation & Unit & Definition
\end{tabular}
\end{center}

Only symbols and abbreviations actually used in the final report shall be
retained. Example entries contained in the template shall be removed if
they are not relevant to the student's report.

Each abbreviation should normally be defined when it first appears in
the main text.

\vspace{0.3cm}

% ================================================================
% XII. FIGURES, DIAGRAMS, AND PHOTOGRAPHS
% ================================================================

\noindent
\textbf{XII.\quad FIGURES, DIAGRAMS, AND PHOTOGRAPHS}

\vspace{0.15cm}

Figures shall be clear, legible, and directly relevant to the report
content.

Figures shall be inserted using the standard \verb|figure| environment.

Example:

\begin{verbatim}
\begin{figure}[H]
    \centering
    \includegraphics[width=0.8\textwidth]
                    {figures/company_structure.pdf}
    \caption{Organizational structure of the company}
    \label{fig:company-structure}
\end{figure}
\end{verbatim}

Figure captions shall be placed below the figure.

Each referenced figure shall have a unique \verb|\label| and shall be
referenced in the text using, for example:

\begin{verbatim}
Figure~\ref{fig:company-structure}
\end{verbatim}

Students shall not manually type figure numbers.

Figures reproduced or adapted from external sources shall include an
appropriate citation.

Photographs taken during the internship shall be relevant and sufficiently
clear. Students shall not include confidential, restricted, or proprietary
information without permission from the internship organization.

Vector graphics or high-resolution images are recommended whenever
possible.

\vspace{0.3cm}

% ================================================================
% XIII. TABLES
% ================================================================

\noindent
\textbf{XIII.\quad TABLES}

\vspace{0.15cm}

Tables shall be created using suitable \LaTeX{} environments such as
\verb|table|, \verb|tabular|, \verb|tabularx|, or \verb|longtable|.

Table captions shall be placed above the table.

Example:

\begin{verbatim}
\begin{table}[H]
    \centering
    \caption{Internship schedule}
    \label{tab:internship-schedule}

    \begin{tabular}{lll}
        \toprule
        Week & Task & Result \\
        \midrule
        1 & Orientation & Completed \\
        2 & System study & Completed \\
        \bottomrule
    \end{tabular}
\end{table}
\end{verbatim}

Tables shall be referenced using:

\begin{verbatim}
Table~\ref{tab:internship-schedule}
\end{verbatim}

Students shall not manually type table numbers.

Tables should be created directly in \LaTeX{} rather than inserted as
screenshots whenever practical.

\vspace{0.3cm}

% ================================================================
% XIV. EQUATIONS AND TECHNICAL QUANTITIES
% ================================================================

\noindent
\textbf{XIV.\quad EQUATIONS AND TECHNICAL QUANTITIES}

\vspace{0.15cm}

Mathematical expressions shall be typeset directly using \LaTeX{}.
Images or screenshots of equations shall not be used.

Short expressions may be written in inline mathematical mode, for example:

\begin{verbatim}
$R = 0.7\,\Omega$
\end{verbatim}

Important equations shall use a numbered equation environment.

Example:

\begin{verbatim}
\begin{equation}
    P = UI
    \label{eq:power}
\end{equation}
\end{verbatim}

Equations shall be referenced using:

\begin{verbatim}
Equation~\eqref{eq:power}
\end{verbatim}

Equation numbers shall not be manually entered.

SI units shall be used wherever applicable. The \verb|siunitx| package
should be used where appropriate.

Examples:

\begin{verbatim}
\SI{220}{\volt}
\SI{50}{\hertz}
\SI{2.5}{\ampere}
\end{verbatim}

Notation and units shall be consistent throughout the report.

\vspace{0.3cm}

% ================================================================
% XV. SOURCE CODE
% ================================================================

\noindent
\textbf{XV.\quad SOURCE CODE}

\vspace{0.15cm}

If the internship involves programming, PLCs, microcontrollers,
simulation, or software development, relevant source code may be included
in the report.

Source code shall be presented using the \verb|listings| environment or
another environment defined by the official template.

Example:

\begin{verbatim}
\begin{lstlisting}[language=st]
IF Start = TRUE THEN
    Motor := TRUE;
ELSE
    Motor := FALSE;
END_IF;
\end{lstlisting}
\end{verbatim}

Only code necessary for explaining the work should appear in the main
report. Long source-code listings should be moved to an appendix or
submitted as supplementary files if required.

Source code should not normally be inserted as screenshots.

\vspace{0.3cm}

% ================================================================
% XVI. ABSTRACT
% ================================================================

\noindent
\textbf{XVI.\quad ABSTRACT}

\vspace{0.15cm}

The Abstract shall briefly summarize:

\begin{itemize}[leftmargin=0.8cm, itemsep=2pt]
    \item the internship organization;
    \item the objectives of the internship;
    \item the main tasks performed;
    \item the main results, knowledge, and skills obtained.
\end{itemize}

The Abstract should be concise and should not reproduce the full content
of the report chapter by chapter.

\vspace{0.3cm}

% ================================================================
% XVII. REFERENCES AND CITATIONS
% ================================================================

\noindent
\textbf{XVII.\quad REFERENCES AND CITATIONS}

\vspace{0.15cm}

All books, journal papers, conference papers, standards, websites,
datasheets, manuals, technical documents, figures, tables, or other
external sources used in preparing the report shall be properly cited.

References shall be managed through the bibliography system provided with
the official template.

Bibliographic entries shall normally be stored in:

\begin{verbatim}
References.bib
\end{verbatim}

References shall be cited in the report using:

\begin{verbatim}
\cite{reference-key}
\end{verbatim}

Students shall not manually type reference numbers such as ``[1]'',
``[2]'', etc.

Titles of foreign-language publications shall remain in their original
language in the bibliography.

All borrowed ideas, data, figures, tables, equations, methods, or results
shall be properly acknowledged.

\vspace{0.3cm}

% ================================================================
% XVIII. CROSS-REFERENCING
% ================================================================

\noindent
\textbf{XVIII.\quad CROSS-REFERENCING}

\vspace{0.15cm}

Numbered elements shall be cross-referenced using the standard \LaTeX{}
mechanism.

Use the following commands where appropriate:

\begin{verbatim}
\label{...}
\ref{...}
\eqref{...}
\end{verbatim}

For example:

\begin{verbatim}
As shown in Figure~\ref{fig:system}, ...
\end{verbatim}

Students shall not manually type numbers such as ``Figure 2.3'' when
those numbers can be generated automatically.

\vspace{0.3cm}

% ================================================================
% XIX. PAGE NUMBERING
% ================================================================

\noindent
\textbf{XIX.\quad PAGE NUMBERING}

\vspace{0.15cm}

Cover pages and official evaluation forms at the beginning of the report
shall not display page numbers.

The main body of the report shall use Arabic page numbers beginning from
Chapter 1.

Page numbering shall be generated automatically by the official template.
Students shall not manually enter page numbers.

\vspace{0.3cm}

% ================================================================
% XX. PROJECT FILE ORGANIZATION
% ================================================================

\noindent
\textbf{XX.\quad PROJECT FILE ORGANIZATION}

\vspace{0.15cm}

Students should maintain a clear and organized \LaTeX{} project
structure. A recommended structure is:

\begin{verbatim}
InternshipReport/
|
|-- main.tex
|-- References.bib
|
|-- figures/
|   |-- company.pdf
|   |-- system.png
|   `-- ...
|
`-- chapters/
    |-- chapter1.tex
    |-- chapter2.tex
    `-- chapter3.tex
\end{verbatim}

Figures should normally be stored in the \verb|figures| directory
specified by the official template.

File and directory names should preferably contain Latin letters,
numbers, hyphens, or underscores.

Recommended examples include:

\begin{verbatim}
company_structure.pdf
motor_control.png
chapter_01.tex
\end{verbatim}

Spaces and unnecessary special characters in file names should be
avoided where possible.

\vspace{0.3cm}

% ================================================================
% XXI. DOCUMENT COMPILATION AND FINAL CHECK
% ================================================================

\noindent
\textbf{XXI.\quad DOCUMENT COMPILATION AND FINAL CHECK}

\vspace{0.15cm}

Before submission, the complete \LaTeX{} project shall compile
successfully and produce the final PDF report.

The final document shall not contain unresolved references such as:

\begin{verbatim}
??
[?]
\end{verbatim}

Before generating the final PDF, students shall check:

\begin{itemize}[leftmargin=0.8cm, itemsep=2pt]
    \item the Table of Contents;
    \item page numbering;
    \item chapter, section, and subsection numbering;
    \item figure numbering;
    \item table numbering;
    \item equation numbering;
    \item cross-references;
    \item citations;
    \item bibliography;
    \item font display;
    \item mathematical symbols;
    \item image quality;
    \item completeness of all required forms.
\end{itemize}

The document may be compiled using TeX Live, MiKTeX, or another
compatible \LaTeX{} environment. Regardless of the compilation
environment, the resulting PDF shall conform to the official Faculty
template.

\vspace{0.3cm}

% ================================================================
% XXII. CONTENT THAT STUDENTS MAY EDIT
% ================================================================

\noindent
\textbf{XXII.\quad CONTENT THAT STUDENTS MAY EDIT}

\vspace{0.15cm}

Students may edit or replace:

\begin{itemize}[leftmargin=0.8cm, itemsep=2pt]
    \item student name;
    \item class;
    \item internship organization;
    \item supervisor name;
    \item Abstract;
    \item List of Symbols and Abbreviations;
    \item chapter titles and chapter content;
    \item figures and photographs;
    \item tables;
    \item equations;
    \item citations and references.
\end{itemize}

Students shall not modify the main document formatting settings in the
preamble, including the document class, page geometry, line spacing,
heading formatting, table-of-contents formatting, font configuration,
page layout, or cover-page macros unless specifically instructed.

\vspace{0.3cm}

% ================================================================
% XXIII. FINAL SUBMISSION REQUIREMENTS
% ================================================================

\noindent
\textbf{XXIII.\quad FINAL SUBMISSION REQUIREMENTS}

\vspace{0.15cm}

Before submission, the final report shall:

\begin{itemize}[leftmargin=0.8cm, itemsep=2pt]
    \item compile without errors that affect the report content or presentation;
    \item contain all required forms and report sections;
    \item contain correct student, class, internship organization, and supervisor information;
    \item contain an updated Table of Contents, List of Figures, and List of Tables;
    \item contain no unused example content from the template;
    \item contain no placeholder text such as \textit{Title of Chapter 1},
    \textit{Section 1.1}, or other sample entries;
    \item contain all required citations and references;
    \item contain no unresolved cross-references.
\end{itemize}

If the Faculty requires submission of the \LaTeX{} source, the complete
project shall include all \verb|.tex| files, the bibliography database,
figures, and any other files required to reproduce the submitted PDF.

\vspace{0.4cm}

\noindent
\textbf{Important:}
Before submitting the final internship report, students shall remove or
comment out the following command in the main document:

\begin{verbatim}

% =================================================================
% REGULATIONS FOR PREPARING THE GRADUATION INTERNSHIP REPORT
% USING LATEX
% =================================================================
%
% This file is intended to be included from the main document with:
%
%     \input{Regulations.tex}
%
% IMPORTANT:
% This section is provided for guidance only. Students shall remove
% or comment out the \input{Regulations.tex} command before submitting
% the final internship report.
%
% =================================================================

\clearpage
\thispagestyle{empty}

\begin{center}
    {\fontsize{16pt}{20pt}\selectfont\bfseries
    SPECIFIC REGULATIONS FOR PREPARING\\
    GRADUATION INTERNSHIP REPORTS USING \LaTeX}
\end{center}

\vspace{0.5cm}

{\fontsize{13pt}{17pt}\selectfont
\setstretch{1.3}
\justifying

\noindent
\textbf{Note:}
This section is provided for guidance only and is included in the main
document using the command \verb|\input{Regulations.tex}|.
Students shall remove or comment out this command before submitting the
final internship report.

\vspace{0.4cm}

% ================================================================
% I. GENERAL REQUIREMENTS
% ================================================================

\noindent
\textbf{I.\quad GENERAL REQUIREMENTS}

\vspace{0.15cm}

The Graduation Internship Report shall be prepared using the official
\LaTeX{} template provided by the Faculty. Students shall preserve the
predefined page layout, typography, line spacing, heading styles,
numbering system, figure and table formatting, and bibliography settings
unless otherwise instructed by the Faculty or the supervisor.

The report shall be written in English.

Students are responsible for the academic and technical content of the
report. The official \LaTeX{} template is responsible for standardized
document formatting.

\vspace{0.3cm}

% ================================================================
% II. REPORT STRUCTURE
% ================================================================

\noindent
\textbf{II.\quad REPORT STRUCTURE}

\vspace{0.15cm}

The report shall normally be arranged in the following order:

\begin{enumerate}[label=\arabic*., leftmargin=1.2cm, itemsep=2pt]
    \item Outer Cover;
    \item Comments of the Internship Host Organization;
    \item Supervisor's Evaluation and Graduation Internship Defense Results;
    \item Table of Contents;
    \item List of Symbols and Abbreviations, if applicable;
    \item List of Figures;
    \item List of Tables;
    \item Abstract;
    \item Main Chapters of the Internship Report;
    \item References;
    \item Appendices, if applicable.
\end{enumerate}

The preliminary pages shall not be numbered as chapters. The main report
content shall be organized into consecutively numbered chapters beginning
with Chapter 1.

\vspace{0.3cm}

% ================================================================
% III. PAGE SIZE AND MARGINS
% ================================================================

\noindent
\textbf{III.\quad PAGE SIZE AND MARGINS}

\vspace{0.15cm}

The report shall be formatted on A4 paper with the following margins:

\begin{center}
\begin{tabular}{ll}
    Left margin:   & 3.5 cm \\
    Right margin:  & 2.0 cm \\
    Top margin:    & 2.5 cm \\
    Bottom margin: & 2.5 cm
\end{tabular}
\end{center}

These settings are predefined in the official template and shall not be
modified by students.

\vspace{0.3cm}

% ================================================================
% IV. FONT, FONT SIZE, AND LINE SPACING
% ================================================================

\noindent
\textbf{IV.\quad FONT, FONT SIZE, AND LINE SPACING}

\vspace{0.15cm}

The main text shall use a 13 pt font size with 1.3-line spacing.
A Times-family typeface shall be used in accordance with the official
template.

The text shall be fully justified.

Students shall not use repeated spaces, repeated manual line breaks, or
multiple blank lines to position text. Paragraph and heading spacing shall
be controlled by the template.

\vspace{0.3cm}

% ================================================================
% V. OUTER COVER
% ================================================================

\noindent
\textbf{V.\quad OUTER COVER}

\vspace{0.15cm}

The outer cover shall follow the official format provided in the template
and shall contain the following information:

\begin{itemize}[leftmargin=0.8cm, itemsep=2pt]
    \item Thai Nguyen University of Technology;
    \item Faculty of International Training;
    \item Engineering and Technology Department;
    \item University logo and Faculty logo;
    \item the title \textbf{INTERNSHIP REPORT};
    \item Internship Organization;
    \item Author;
    \item Class;
    \item Supervisor;
    \item place and year of completion.
\end{itemize}

Students shall only replace the required information fields and shall not
modify the cover-page layout unless instructed otherwise.

\vspace{0.3cm}

% ================================================================
% VI. COMMENTS OF THE INTERNSHIP HOST ORGANIZATION
% ================================================================

\noindent
\textbf{VI.\quad COMMENTS OF THE INTERNSHIP HOST ORGANIZATION}

\vspace{0.15cm}

The page entitled \textbf{COMMENTS OF THE INTERNSHIP HOST ORGANIZATION}
is reserved for the internship host organization.

It includes:

\begin{itemize}[leftmargin=0.8cm, itemsep=2pt]
    \item the student's full name;
    \item class;
    \item comments from the host organization;
    \item score;
    \item date;
    \item signature and official seal.
\end{itemize}

Students shall not alter or complete official evaluation fields unless
specifically authorized.

\vspace{0.3cm}

% ================================================================
% VII. SUPERVISOR'S EVALUATION AND DEFENSE RESULTS
% ================================================================

\noindent
\textbf{VII.\quad SUPERVISOR'S EVALUATION AND DEFENSE RESULTS}

\vspace{0.15cm}

The section entitled
\textbf{SUPERVISOR'S EVALUATION AND GRADUATION INTERNSHIP DEFENSE RESULTS}
is reserved for the supervisor and defense committee.

The evaluation may include:

\begin{itemize}[leftmargin=0.8cm, itemsep=2pt]
    \item Internship Progress and Student's Attitude;
    \item Report Content;
    \item Report Presentation;
    \item Other Comments;
    \item Supervisor's Opinion;
    \item Defense Result;
    \item Score.
\end{itemize}

Students shall not enter evaluation scores or defense results before they
are officially determined.

\vspace{0.3cm}

% ================================================================
% VIII. CHAPTERS, SECTIONS, AND SUBSECTIONS
% ================================================================

\noindent
\textbf{VIII.\quad CHAPTERS, SECTIONS, AND SUBSECTIONS}

\vspace{0.15cm}

The report structure shall be generated using the standard \LaTeX{}
commands:

\begin{verbatim}
\chapter{...}
\section{...}
\subsection{...}
\end{verbatim}

Chapter, section, and subsection numbers shall be generated automatically
by \LaTeX{}.

Students shall not manually type chapter or section numbers and shall not
reproduce heading styles using commands such as \verb|\textbf|,
\verb|\fontsize|, or \verb|\centering| instead of the appropriate
structural command.

Excessive subdivision should be avoided. In most cases, the hierarchy
Chapter--Section--Subsection is sufficient.

\vspace{0.3cm}

% ================================================================
% IX. ORGANIZATION OF THE MAIN CONTENT
% ================================================================

\noindent
\textbf{IX.\quad ORGANIZATION OF THE MAIN CONTENT}

\vspace{0.15cm}

The specific content of the report shall depend on the internship
organization and the assigned tasks. A typical structure may include:

\begin{itemize}[leftmargin=0.8cm, itemsep=2pt]
    \item \textbf{Chapter 1. Overview of the Internship Organization};
    \item \textbf{Chapter 2. Internship Activities and Assigned Tasks};
    \item \textbf{Chapter 3. Results, Experience, and Evaluation}.
\end{itemize}

These chapter titles are recommendations only and may be adapted to the
actual internship content.

The report shall reflect work actually performed, observed, analyzed, or
learned during the internship. Unrelated theoretical material should not
be included merely to increase report length.

\vspace{0.3cm}

% ================================================================
% X. TABLE OF CONTENTS
% ================================================================

\noindent
\textbf{X.\quad TABLE OF CONTENTS}

\vspace{0.15cm}

The Table of Contents shall be generated automatically using:

\begin{verbatim}
\tableofcontents
\end{verbatim}

Students shall not create or manually edit page numbers in the Table of
Contents.

The document shall be compiled sufficiently to ensure that all page
numbers and entries are correctly updated.

\vspace{0.3cm}

% ================================================================
% XI. LIST OF SYMBOLS AND ABBREVIATIONS
% ================================================================

\noindent
\textbf{XI.\quad LIST OF SYMBOLS AND ABBREVIATIONS}

\vspace{0.15cm}

Symbols and abbreviations that are used repeatedly in the report shall be
listed in the \textbf{List of Symbols and Abbreviations}.

The list may contain the following columns:

\begin{center}
\begin{tabular}{llll}
    No. & Symbol / Abbreviation & Unit & Definition
\end{tabular}
\end{center}

Only symbols and abbreviations actually used in the final report shall be
retained. Example entries contained in the template shall be removed if
they are not relevant to the student's report.

Each abbreviation should normally be defined when it first appears in
the main text.

\vspace{0.3cm}

% ================================================================
% XII. FIGURES, DIAGRAMS, AND PHOTOGRAPHS
% ================================================================

\noindent
\textbf{XII.\quad FIGURES, DIAGRAMS, AND PHOTOGRAPHS}

\vspace{0.15cm}

Figures shall be clear, legible, and directly relevant to the report
content.

Figures shall be inserted using the standard \verb|figure| environment.

Example:

\begin{verbatim}
\begin{figure}[H]
    \centering
    \includegraphics[width=0.8\textwidth]
                    {figures/company_structure.pdf}
    \caption{Organizational structure of the company}
    \label{fig:company-structure}
\end{figure}
\end{verbatim}

Figure captions shall be placed below the figure.

Each referenced figure shall have a unique \verb|\label| and shall be
referenced in the text using, for example:

\begin{verbatim}
Figure~\ref{fig:company-structure}
\end{verbatim}

Students shall not manually type figure numbers.

Figures reproduced or adapted from external sources shall include an
appropriate citation.

Photographs taken during the internship shall be relevant and sufficiently
clear. Students shall not include confidential, restricted, or proprietary
information without permission from the internship organization.

Vector graphics or high-resolution images are recommended whenever
possible.

\vspace{0.3cm}

% ================================================================
% XIII. TABLES
% ================================================================

\noindent
\textbf{XIII.\quad TABLES}

\vspace{0.15cm}

Tables shall be created using suitable \LaTeX{} environments such as
\verb|table|, \verb|tabular|, \verb|tabularx|, or \verb|longtable|.

Table captions shall be placed above the table.

Example:

\begin{verbatim}
\begin{table}[H]
    \centering
    \caption{Internship schedule}
    \label{tab:internship-schedule}

    \begin{tabular}{lll}
        \toprule
        Week & Task & Result \\
        \midrule
        1 & Orientation & Completed \\
        2 & System study & Completed \\
        \bottomrule
    \end{tabular}
\end{table}
\end{verbatim}

Tables shall be referenced using:

\begin{verbatim}
Table~\ref{tab:internship-schedule}
\end{verbatim}

Students shall not manually type table numbers.

Tables should be created directly in \LaTeX{} rather than inserted as
screenshots whenever practical.

\vspace{0.3cm}

% ================================================================
% XIV. EQUATIONS AND TECHNICAL QUANTITIES
% ================================================================

\noindent
\textbf{XIV.\quad EQUATIONS AND TECHNICAL QUANTITIES}

\vspace{0.15cm}

Mathematical expressions shall be typeset directly using \LaTeX{}.
Images or screenshots of equations shall not be used.

Short expressions may be written in inline mathematical mode, for example:

\begin{verbatim}
$R = 0.7\,\Omega$
\end{verbatim}

Important equations shall use a numbered equation environment.

Example:

\begin{verbatim}
\begin{equation}
    P = UI
    \label{eq:power}
\end{equation}
\end{verbatim}

Equations shall be referenced using:

\begin{verbatim}
Equation~\eqref{eq:power}
\end{verbatim}

Equation numbers shall not be manually entered.

SI units shall be used wherever applicable. The \verb|siunitx| package
should be used where appropriate.

Examples:

\begin{verbatim}
\SI{220}{\volt}
\SI{50}{\hertz}
\SI{2.5}{\ampere}
\end{verbatim}

Notation and units shall be consistent throughout the report.

\vspace{0.3cm}

% ================================================================
% XV. SOURCE CODE
% ================================================================

\noindent
\textbf{XV.\quad SOURCE CODE}

\vspace{0.15cm}

If the internship involves programming, PLCs, microcontrollers,
simulation, or software development, relevant source code may be included
in the report.

Source code shall be presented using the \verb|listings| environment or
another environment defined by the official template.

Example:

\begin{verbatim}
\begin{lstlisting}[language=st]
IF Start = TRUE THEN
    Motor := TRUE;
ELSE
    Motor := FALSE;
END_IF;
\end{lstlisting}
\end{verbatim}

Only code necessary for explaining the work should appear in the main
report. Long source-code listings should be moved to an appendix or
submitted as supplementary files if required.

Source code should not normally be inserted as screenshots.

\vspace{0.3cm}

% ================================================================
% XVI. ABSTRACT
% ================================================================

\noindent
\textbf{XVI.\quad ABSTRACT}

\vspace{0.15cm}

The Abstract shall briefly summarize:

\begin{itemize}[leftmargin=0.8cm, itemsep=2pt]
    \item the internship organization;
    \item the objectives of the internship;
    \item the main tasks performed;
    \item the main results, knowledge, and skills obtained.
\end{itemize}

The Abstract should be concise and should not reproduce the full content
of the report chapter by chapter.

\vspace{0.3cm}

% ================================================================
% XVII. REFERENCES AND CITATIONS
% ================================================================

\noindent
\textbf{XVII.\quad REFERENCES AND CITATIONS}

\vspace{0.15cm}

All books, journal papers, conference papers, standards, websites,
datasheets, manuals, technical documents, figures, tables, or other
external sources used in preparing the report shall be properly cited.

References shall be managed through the bibliography system provided with
the official template.

Bibliographic entries shall normally be stored in:

\begin{verbatim}
References.bib
\end{verbatim}

References shall be cited in the report using:

\begin{verbatim}
\cite{reference-key}
\end{verbatim}

Students shall not manually type reference numbers such as ``[1]'',
``[2]'', etc.

Titles of foreign-language publications shall remain in their original
language in the bibliography.

All borrowed ideas, data, figures, tables, equations, methods, or results
shall be properly acknowledged.

\vspace{0.3cm}

% ================================================================
% XVIII. CROSS-REFERENCING
% ================================================================

\noindent
\textbf{XVIII.\quad CROSS-REFERENCING}

\vspace{0.15cm}

Numbered elements shall be cross-referenced using the standard \LaTeX{}
mechanism.

Use the following commands where appropriate:

\begin{verbatim}
\label{...}
\ref{...}
\eqref{...}
\end{verbatim}

For example:

\begin{verbatim}
As shown in Figure~\ref{fig:system}, ...
\end{verbatim}

Students shall not manually type numbers such as ``Figure 2.3'' when
those numbers can be generated automatically.

\vspace{0.3cm}

% ================================================================
% XIX. PAGE NUMBERING
% ================================================================

\noindent
\textbf{XIX.\quad PAGE NUMBERING}

\vspace{0.15cm}

Cover pages and official evaluation forms at the beginning of the report
shall not display page numbers.

The main body of the report shall use Arabic page numbers beginning from
Chapter 1.

Page numbering shall be generated automatically by the official template.
Students shall not manually enter page numbers.

\vspace{0.3cm}

% ================================================================
% XX. PROJECT FILE ORGANIZATION
% ================================================================

\noindent
\textbf{XX.\quad PROJECT FILE ORGANIZATION}

\vspace{0.15cm}

Students should maintain a clear and organized \LaTeX{} project
structure. A recommended structure is:

\begin{verbatim}
InternshipReport/
|
|-- main.tex
|-- References.bib
|
|-- figures/
|   |-- company.pdf
|   |-- system.png
|   `-- ...
|
`-- chapters/
    |-- chapter1.tex
    |-- chapter2.tex
    `-- chapter3.tex
\end{verbatim}

Figures should normally be stored in the \verb|figures| directory
specified by the official template.

File and directory names should preferably contain Latin letters,
numbers, hyphens, or underscores.

Recommended examples include:

\begin{verbatim}
company_structure.pdf
motor_control.png
chapter_01.tex
\end{verbatim}

Spaces and unnecessary special characters in file names should be
avoided where possible.

\vspace{0.3cm}

% ================================================================
% XXI. DOCUMENT COMPILATION AND FINAL CHECK
% ================================================================

\noindent
\textbf{XXI.\quad DOCUMENT COMPILATION AND FINAL CHECK}

\vspace{0.15cm}

Before submission, the complete \LaTeX{} project shall compile
successfully and produce the final PDF report.

The final document shall not contain unresolved references such as:

\begin{verbatim}
??
[?]
\end{verbatim}

Before generating the final PDF, students shall check:

\begin{itemize}[leftmargin=0.8cm, itemsep=2pt]
    \item the Table of Contents;
    \item page numbering;
    \item chapter, section, and subsection numbering;
    \item figure numbering;
    \item table numbering;
    \item equation numbering;
    \item cross-references;
    \item citations;
    \item bibliography;
    \item font display;
    \item mathematical symbols;
    \item image quality;
    \item completeness of all required forms.
\end{itemize}

The document may be compiled using TeX Live, MiKTeX, or another
compatible \LaTeX{} environment. Regardless of the compilation
environment, the resulting PDF shall conform to the official Faculty
template.

\vspace{0.3cm}

% ================================================================
% XXII. CONTENT THAT STUDENTS MAY EDIT
% ================================================================

\noindent
\textbf{XXII.\quad CONTENT THAT STUDENTS MAY EDIT}

\vspace{0.15cm}

Students may edit or replace:

\begin{itemize}[leftmargin=0.8cm, itemsep=2pt]
    \item student name;
    \item class;
    \item internship organization;
    \item supervisor name;
    \item Abstract;
    \item List of Symbols and Abbreviations;
    \item chapter titles and chapter content;
    \item figures and photographs;
    \item tables;
    \item equations;
    \item citations and references.
\end{itemize}

Students shall not modify the main document formatting settings in the
preamble, including the document class, page geometry, line spacing,
heading formatting, table-of-contents formatting, font configuration,
page layout, or cover-page macros unless specifically instructed.

\vspace{0.3cm}

% ================================================================
% XXIII. FINAL SUBMISSION REQUIREMENTS
% ================================================================

\noindent
\textbf{XXIII.\quad FINAL SUBMISSION REQUIREMENTS}

\vspace{0.15cm}

Before submission, the final report shall:

\begin{itemize}[leftmargin=0.8cm, itemsep=2pt]
    \item compile without errors that affect the report content or presentation;
    \item contain all required forms and report sections;
    \item contain correct student, class, internship organization, and supervisor information;
    \item contain an updated Table of Contents, List of Figures, and List of Tables;
    \item contain no unused example content from the template;
    \item contain no placeholder text such as \textit{Title of Chapter 1},
    \textit{Section 1.1}, or other sample entries;
    \item contain all required citations and references;
    \item contain no unresolved cross-references.
\end{itemize}

If the Faculty requires submission of the \LaTeX{} source, the complete
project shall include all \verb|.tex| files, the bibliography database,
figures, and any other files required to reproduce the submitted PDF.

\vspace{0.4cm}

\noindent
\textbf{Important:}
Before submitting the final internship report, students shall remove or
comment out the following command in the main document:

\begin{verbatim}
\input{Regulations.tex}
\end{verbatim}

For example:

\begin{verbatim}
% \input{Regulations.tex}
\end{verbatim}

This removes the guidance pages from the final submitted report while
preserving the regulations file in the \LaTeX{} project for reference.

} % End 13pt content

\clearpage

\end{verbatim}

For example:

\begin{verbatim}
% 
% =================================================================
% REGULATIONS FOR PREPARING THE GRADUATION INTERNSHIP REPORT
% USING LATEX
% =================================================================
%
% This file is intended to be included from the main document with:
%
%     \input{Regulations.tex}
%
% IMPORTANT:
% This section is provided for guidance only. Students shall remove
% or comment out the \input{Regulations.tex} command before submitting
% the final internship report.
%
% =================================================================

\clearpage
\thispagestyle{empty}

\begin{center}
    {\fontsize{16pt}{20pt}\selectfont\bfseries
    SPECIFIC REGULATIONS FOR PREPARING\\
    GRADUATION INTERNSHIP REPORTS USING \LaTeX}
\end{center}

\vspace{0.5cm}

{\fontsize{13pt}{17pt}\selectfont
\setstretch{1.3}
\justifying

\noindent
\textbf{Note:}
This section is provided for guidance only and is included in the main
document using the command \verb|\input{Regulations.tex}|.
Students shall remove or comment out this command before submitting the
final internship report.

\vspace{0.4cm}

% ================================================================
% I. GENERAL REQUIREMENTS
% ================================================================

\noindent
\textbf{I.\quad GENERAL REQUIREMENTS}

\vspace{0.15cm}

The Graduation Internship Report shall be prepared using the official
\LaTeX{} template provided by the Faculty. Students shall preserve the
predefined page layout, typography, line spacing, heading styles,
numbering system, figure and table formatting, and bibliography settings
unless otherwise instructed by the Faculty or the supervisor.

The report shall be written in English.

Students are responsible for the academic and technical content of the
report. The official \LaTeX{} template is responsible for standardized
document formatting.

\vspace{0.3cm}

% ================================================================
% II. REPORT STRUCTURE
% ================================================================

\noindent
\textbf{II.\quad REPORT STRUCTURE}

\vspace{0.15cm}

The report shall normally be arranged in the following order:

\begin{enumerate}[label=\arabic*., leftmargin=1.2cm, itemsep=2pt]
    \item Outer Cover;
    \item Comments of the Internship Host Organization;
    \item Supervisor's Evaluation and Graduation Internship Defense Results;
    \item Table of Contents;
    \item List of Symbols and Abbreviations, if applicable;
    \item List of Figures;
    \item List of Tables;
    \item Abstract;
    \item Main Chapters of the Internship Report;
    \item References;
    \item Appendices, if applicable.
\end{enumerate}

The preliminary pages shall not be numbered as chapters. The main report
content shall be organized into consecutively numbered chapters beginning
with Chapter 1.

\vspace{0.3cm}

% ================================================================
% III. PAGE SIZE AND MARGINS
% ================================================================

\noindent
\textbf{III.\quad PAGE SIZE AND MARGINS}

\vspace{0.15cm}

The report shall be formatted on A4 paper with the following margins:

\begin{center}
\begin{tabular}{ll}
    Left margin:   & 3.5 cm \\
    Right margin:  & 2.0 cm \\
    Top margin:    & 2.5 cm \\
    Bottom margin: & 2.5 cm
\end{tabular}
\end{center}

These settings are predefined in the official template and shall not be
modified by students.

\vspace{0.3cm}

% ================================================================
% IV. FONT, FONT SIZE, AND LINE SPACING
% ================================================================

\noindent
\textbf{IV.\quad FONT, FONT SIZE, AND LINE SPACING}

\vspace{0.15cm}

The main text shall use a 13 pt font size with 1.3-line spacing.
A Times-family typeface shall be used in accordance with the official
template.

The text shall be fully justified.

Students shall not use repeated spaces, repeated manual line breaks, or
multiple blank lines to position text. Paragraph and heading spacing shall
be controlled by the template.

\vspace{0.3cm}

% ================================================================
% V. OUTER COVER
% ================================================================

\noindent
\textbf{V.\quad OUTER COVER}

\vspace{0.15cm}

The outer cover shall follow the official format provided in the template
and shall contain the following information:

\begin{itemize}[leftmargin=0.8cm, itemsep=2pt]
    \item Thai Nguyen University of Technology;
    \item Faculty of International Training;
    \item Engineering and Technology Department;
    \item University logo and Faculty logo;
    \item the title \textbf{INTERNSHIP REPORT};
    \item Internship Organization;
    \item Author;
    \item Class;
    \item Supervisor;
    \item place and year of completion.
\end{itemize}

Students shall only replace the required information fields and shall not
modify the cover-page layout unless instructed otherwise.

\vspace{0.3cm}

% ================================================================
% VI. COMMENTS OF THE INTERNSHIP HOST ORGANIZATION
% ================================================================

\noindent
\textbf{VI.\quad COMMENTS OF THE INTERNSHIP HOST ORGANIZATION}

\vspace{0.15cm}

The page entitled \textbf{COMMENTS OF THE INTERNSHIP HOST ORGANIZATION}
is reserved for the internship host organization.

It includes:

\begin{itemize}[leftmargin=0.8cm, itemsep=2pt]
    \item the student's full name;
    \item class;
    \item comments from the host organization;
    \item score;
    \item date;
    \item signature and official seal.
\end{itemize}

Students shall not alter or complete official evaluation fields unless
specifically authorized.

\vspace{0.3cm}

% ================================================================
% VII. SUPERVISOR'S EVALUATION AND DEFENSE RESULTS
% ================================================================

\noindent
\textbf{VII.\quad SUPERVISOR'S EVALUATION AND DEFENSE RESULTS}

\vspace{0.15cm}

The section entitled
\textbf{SUPERVISOR'S EVALUATION AND GRADUATION INTERNSHIP DEFENSE RESULTS}
is reserved for the supervisor and defense committee.

The evaluation may include:

\begin{itemize}[leftmargin=0.8cm, itemsep=2pt]
    \item Internship Progress and Student's Attitude;
    \item Report Content;
    \item Report Presentation;
    \item Other Comments;
    \item Supervisor's Opinion;
    \item Defense Result;
    \item Score.
\end{itemize}

Students shall not enter evaluation scores or defense results before they
are officially determined.

\vspace{0.3cm}

% ================================================================
% VIII. CHAPTERS, SECTIONS, AND SUBSECTIONS
% ================================================================

\noindent
\textbf{VIII.\quad CHAPTERS, SECTIONS, AND SUBSECTIONS}

\vspace{0.15cm}

The report structure shall be generated using the standard \LaTeX{}
commands:

\begin{verbatim}
\chapter{...}
\section{...}
\subsection{...}
\end{verbatim}

Chapter, section, and subsection numbers shall be generated automatically
by \LaTeX{}.

Students shall not manually type chapter or section numbers and shall not
reproduce heading styles using commands such as \verb|\textbf|,
\verb|\fontsize|, or \verb|\centering| instead of the appropriate
structural command.

Excessive subdivision should be avoided. In most cases, the hierarchy
Chapter--Section--Subsection is sufficient.

\vspace{0.3cm}

% ================================================================
% IX. ORGANIZATION OF THE MAIN CONTENT
% ================================================================

\noindent
\textbf{IX.\quad ORGANIZATION OF THE MAIN CONTENT}

\vspace{0.15cm}

The specific content of the report shall depend on the internship
organization and the assigned tasks. A typical structure may include:

\begin{itemize}[leftmargin=0.8cm, itemsep=2pt]
    \item \textbf{Chapter 1. Overview of the Internship Organization};
    \item \textbf{Chapter 2. Internship Activities and Assigned Tasks};
    \item \textbf{Chapter 3. Results, Experience, and Evaluation}.
\end{itemize}

These chapter titles are recommendations only and may be adapted to the
actual internship content.

The report shall reflect work actually performed, observed, analyzed, or
learned during the internship. Unrelated theoretical material should not
be included merely to increase report length.

\vspace{0.3cm}

% ================================================================
% X. TABLE OF CONTENTS
% ================================================================

\noindent
\textbf{X.\quad TABLE OF CONTENTS}

\vspace{0.15cm}

The Table of Contents shall be generated automatically using:

\begin{verbatim}
\tableofcontents
\end{verbatim}

Students shall not create or manually edit page numbers in the Table of
Contents.

The document shall be compiled sufficiently to ensure that all page
numbers and entries are correctly updated.

\vspace{0.3cm}

% ================================================================
% XI. LIST OF SYMBOLS AND ABBREVIATIONS
% ================================================================

\noindent
\textbf{XI.\quad LIST OF SYMBOLS AND ABBREVIATIONS}

\vspace{0.15cm}

Symbols and abbreviations that are used repeatedly in the report shall be
listed in the \textbf{List of Symbols and Abbreviations}.

The list may contain the following columns:

\begin{center}
\begin{tabular}{llll}
    No. & Symbol / Abbreviation & Unit & Definition
\end{tabular}
\end{center}

Only symbols and abbreviations actually used in the final report shall be
retained. Example entries contained in the template shall be removed if
they are not relevant to the student's report.

Each abbreviation should normally be defined when it first appears in
the main text.

\vspace{0.3cm}

% ================================================================
% XII. FIGURES, DIAGRAMS, AND PHOTOGRAPHS
% ================================================================

\noindent
\textbf{XII.\quad FIGURES, DIAGRAMS, AND PHOTOGRAPHS}

\vspace{0.15cm}

Figures shall be clear, legible, and directly relevant to the report
content.

Figures shall be inserted using the standard \verb|figure| environment.

Example:

\begin{verbatim}
\begin{figure}[H]
    \centering
    \includegraphics[width=0.8\textwidth]
                    {figures/company_structure.pdf}
    \caption{Organizational structure of the company}
    \label{fig:company-structure}
\end{figure}
\end{verbatim}

Figure captions shall be placed below the figure.

Each referenced figure shall have a unique \verb|\label| and shall be
referenced in the text using, for example:

\begin{verbatim}
Figure~\ref{fig:company-structure}
\end{verbatim}

Students shall not manually type figure numbers.

Figures reproduced or adapted from external sources shall include an
appropriate citation.

Photographs taken during the internship shall be relevant and sufficiently
clear. Students shall not include confidential, restricted, or proprietary
information without permission from the internship organization.

Vector graphics or high-resolution images are recommended whenever
possible.

\vspace{0.3cm}

% ================================================================
% XIII. TABLES
% ================================================================

\noindent
\textbf{XIII.\quad TABLES}

\vspace{0.15cm}

Tables shall be created using suitable \LaTeX{} environments such as
\verb|table|, \verb|tabular|, \verb|tabularx|, or \verb|longtable|.

Table captions shall be placed above the table.

Example:

\begin{verbatim}
\begin{table}[H]
    \centering
    \caption{Internship schedule}
    \label{tab:internship-schedule}

    \begin{tabular}{lll}
        \toprule
        Week & Task & Result \\
        \midrule
        1 & Orientation & Completed \\
        2 & System study & Completed \\
        \bottomrule
    \end{tabular}
\end{table}
\end{verbatim}

Tables shall be referenced using:

\begin{verbatim}
Table~\ref{tab:internship-schedule}
\end{verbatim}

Students shall not manually type table numbers.

Tables should be created directly in \LaTeX{} rather than inserted as
screenshots whenever practical.

\vspace{0.3cm}

% ================================================================
% XIV. EQUATIONS AND TECHNICAL QUANTITIES
% ================================================================

\noindent
\textbf{XIV.\quad EQUATIONS AND TECHNICAL QUANTITIES}

\vspace{0.15cm}

Mathematical expressions shall be typeset directly using \LaTeX{}.
Images or screenshots of equations shall not be used.

Short expressions may be written in inline mathematical mode, for example:

\begin{verbatim}
$R = 0.7\,\Omega$
\end{verbatim}

Important equations shall use a numbered equation environment.

Example:

\begin{verbatim}
\begin{equation}
    P = UI
    \label{eq:power}
\end{equation}
\end{verbatim}

Equations shall be referenced using:

\begin{verbatim}
Equation~\eqref{eq:power}
\end{verbatim}

Equation numbers shall not be manually entered.

SI units shall be used wherever applicable. The \verb|siunitx| package
should be used where appropriate.

Examples:

\begin{verbatim}
\SI{220}{\volt}
\SI{50}{\hertz}
\SI{2.5}{\ampere}
\end{verbatim}

Notation and units shall be consistent throughout the report.

\vspace{0.3cm}

% ================================================================
% XV. SOURCE CODE
% ================================================================

\noindent
\textbf{XV.\quad SOURCE CODE}

\vspace{0.15cm}

If the internship involves programming, PLCs, microcontrollers,
simulation, or software development, relevant source code may be included
in the report.

Source code shall be presented using the \verb|listings| environment or
another environment defined by the official template.

Example:

\begin{verbatim}
\begin{lstlisting}[language=st]
IF Start = TRUE THEN
    Motor := TRUE;
ELSE
    Motor := FALSE;
END_IF;
\end{lstlisting}
\end{verbatim}

Only code necessary for explaining the work should appear in the main
report. Long source-code listings should be moved to an appendix or
submitted as supplementary files if required.

Source code should not normally be inserted as screenshots.

\vspace{0.3cm}

% ================================================================
% XVI. ABSTRACT
% ================================================================

\noindent
\textbf{XVI.\quad ABSTRACT}

\vspace{0.15cm}

The Abstract shall briefly summarize:

\begin{itemize}[leftmargin=0.8cm, itemsep=2pt]
    \item the internship organization;
    \item the objectives of the internship;
    \item the main tasks performed;
    \item the main results, knowledge, and skills obtained.
\end{itemize}

The Abstract should be concise and should not reproduce the full content
of the report chapter by chapter.

\vspace{0.3cm}

% ================================================================
% XVII. REFERENCES AND CITATIONS
% ================================================================

\noindent
\textbf{XVII.\quad REFERENCES AND CITATIONS}

\vspace{0.15cm}

All books, journal papers, conference papers, standards, websites,
datasheets, manuals, technical documents, figures, tables, or other
external sources used in preparing the report shall be properly cited.

References shall be managed through the bibliography system provided with
the official template.

Bibliographic entries shall normally be stored in:

\begin{verbatim}
References.bib
\end{verbatim}

References shall be cited in the report using:

\begin{verbatim}
\cite{reference-key}
\end{verbatim}

Students shall not manually type reference numbers such as ``[1]'',
``[2]'', etc.

Titles of foreign-language publications shall remain in their original
language in the bibliography.

All borrowed ideas, data, figures, tables, equations, methods, or results
shall be properly acknowledged.

\vspace{0.3cm}

% ================================================================
% XVIII. CROSS-REFERENCING
% ================================================================

\noindent
\textbf{XVIII.\quad CROSS-REFERENCING}

\vspace{0.15cm}

Numbered elements shall be cross-referenced using the standard \LaTeX{}
mechanism.

Use the following commands where appropriate:

\begin{verbatim}
\label{...}
\ref{...}
\eqref{...}
\end{verbatim}

For example:

\begin{verbatim}
As shown in Figure~\ref{fig:system}, ...
\end{verbatim}

Students shall not manually type numbers such as ``Figure 2.3'' when
those numbers can be generated automatically.

\vspace{0.3cm}

% ================================================================
% XIX. PAGE NUMBERING
% ================================================================

\noindent
\textbf{XIX.\quad PAGE NUMBERING}

\vspace{0.15cm}

Cover pages and official evaluation forms at the beginning of the report
shall not display page numbers.

The main body of the report shall use Arabic page numbers beginning from
Chapter 1.

Page numbering shall be generated automatically by the official template.
Students shall not manually enter page numbers.

\vspace{0.3cm}

% ================================================================
% XX. PROJECT FILE ORGANIZATION
% ================================================================

\noindent
\textbf{XX.\quad PROJECT FILE ORGANIZATION}

\vspace{0.15cm}

Students should maintain a clear and organized \LaTeX{} project
structure. A recommended structure is:

\begin{verbatim}
InternshipReport/
|
|-- main.tex
|-- References.bib
|
|-- figures/
|   |-- company.pdf
|   |-- system.png
|   `-- ...
|
`-- chapters/
    |-- chapter1.tex
    |-- chapter2.tex
    `-- chapter3.tex
\end{verbatim}

Figures should normally be stored in the \verb|figures| directory
specified by the official template.

File and directory names should preferably contain Latin letters,
numbers, hyphens, or underscores.

Recommended examples include:

\begin{verbatim}
company_structure.pdf
motor_control.png
chapter_01.tex
\end{verbatim}

Spaces and unnecessary special characters in file names should be
avoided where possible.

\vspace{0.3cm}

% ================================================================
% XXI. DOCUMENT COMPILATION AND FINAL CHECK
% ================================================================

\noindent
\textbf{XXI.\quad DOCUMENT COMPILATION AND FINAL CHECK}

\vspace{0.15cm}

Before submission, the complete \LaTeX{} project shall compile
successfully and produce the final PDF report.

The final document shall not contain unresolved references such as:

\begin{verbatim}
??
[?]
\end{verbatim}

Before generating the final PDF, students shall check:

\begin{itemize}[leftmargin=0.8cm, itemsep=2pt]
    \item the Table of Contents;
    \item page numbering;
    \item chapter, section, and subsection numbering;
    \item figure numbering;
    \item table numbering;
    \item equation numbering;
    \item cross-references;
    \item citations;
    \item bibliography;
    \item font display;
    \item mathematical symbols;
    \item image quality;
    \item completeness of all required forms.
\end{itemize}

The document may be compiled using TeX Live, MiKTeX, or another
compatible \LaTeX{} environment. Regardless of the compilation
environment, the resulting PDF shall conform to the official Faculty
template.

\vspace{0.3cm}

% ================================================================
% XXII. CONTENT THAT STUDENTS MAY EDIT
% ================================================================

\noindent
\textbf{XXII.\quad CONTENT THAT STUDENTS MAY EDIT}

\vspace{0.15cm}

Students may edit or replace:

\begin{itemize}[leftmargin=0.8cm, itemsep=2pt]
    \item student name;
    \item class;
    \item internship organization;
    \item supervisor name;
    \item Abstract;
    \item List of Symbols and Abbreviations;
    \item chapter titles and chapter content;
    \item figures and photographs;
    \item tables;
    \item equations;
    \item citations and references.
\end{itemize}

Students shall not modify the main document formatting settings in the
preamble, including the document class, page geometry, line spacing,
heading formatting, table-of-contents formatting, font configuration,
page layout, or cover-page macros unless specifically instructed.

\vspace{0.3cm}

% ================================================================
% XXIII. FINAL SUBMISSION REQUIREMENTS
% ================================================================

\noindent
\textbf{XXIII.\quad FINAL SUBMISSION REQUIREMENTS}

\vspace{0.15cm}

Before submission, the final report shall:

\begin{itemize}[leftmargin=0.8cm, itemsep=2pt]
    \item compile without errors that affect the report content or presentation;
    \item contain all required forms and report sections;
    \item contain correct student, class, internship organization, and supervisor information;
    \item contain an updated Table of Contents, List of Figures, and List of Tables;
    \item contain no unused example content from the template;
    \item contain no placeholder text such as \textit{Title of Chapter 1},
    \textit{Section 1.1}, or other sample entries;
    \item contain all required citations and references;
    \item contain no unresolved cross-references.
\end{itemize}

If the Faculty requires submission of the \LaTeX{} source, the complete
project shall include all \verb|.tex| files, the bibliography database,
figures, and any other files required to reproduce the submitted PDF.

\vspace{0.4cm}

\noindent
\textbf{Important:}
Before submitting the final internship report, students shall remove or
comment out the following command in the main document:

\begin{verbatim}
\input{Regulations.tex}
\end{verbatim}

For example:

\begin{verbatim}
% \input{Regulations.tex}
\end{verbatim}

This removes the guidance pages from the final submitted report while
preserving the regulations file in the \LaTeX{} project for reference.

} % End 13pt content

\clearpage

\end{verbatim}

This removes the guidance pages from the final submitted report while
preserving the regulations file in the \LaTeX{} project for reference.

} % End 13pt content

\clearpage
|.
Students shall remove or comment out this command before submitting the
final internship report.

\vspace{0.4cm}

% ================================================================
% I. GENERAL REQUIREMENTS
% ================================================================

\noindent
\textbf{I.\quad GENERAL REQUIREMENTS}

\vspace{0.15cm}

The Graduation Internship Report shall be prepared using the official
\LaTeX{} template provided by the Faculty. Students shall preserve the
predefined page layout, typography, line spacing, heading styles,
numbering system, figure and table formatting, and bibliography settings
unless otherwise instructed by the Faculty or the supervisor.

The report shall be written in English.

Students are responsible for the academic and technical content of the
report. The official \LaTeX{} template is responsible for standardized
document formatting.

\vspace{0.3cm}

% ================================================================
% II. REPORT STRUCTURE
% ================================================================

\noindent
\textbf{II.\quad REPORT STRUCTURE}

\vspace{0.15cm}

The report shall normally be arranged in the following order:

\begin{enumerate}[label=\arabic*., leftmargin=1.2cm, itemsep=2pt]
    \item Outer Cover;
    \item Comments of the Internship Host Organization;
    \item Supervisor's Evaluation and Graduation Internship Defense Results;
    \item Table of Contents;
    \item List of Symbols and Abbreviations, if applicable;
    \item List of Figures;
    \item List of Tables;
    \item Abstract;
    \item Main Chapters of the Internship Report;
    \item References;
    \item Appendices, if applicable.
\end{enumerate}

The preliminary pages shall not be numbered as chapters. The main report
content shall be organized into consecutively numbered chapters beginning
with Chapter 1.

\vspace{0.3cm}

% ================================================================
% III. PAGE SIZE AND MARGINS
% ================================================================

\noindent
\textbf{III.\quad PAGE SIZE AND MARGINS}

\vspace{0.15cm}

The report shall be formatted on A4 paper with the following margins:

\begin{center}
\begin{tabular}{ll}
    Left margin:   & 3.5 cm \\
    Right margin:  & 2.0 cm \\
    Top margin:    & 2.5 cm \\
    Bottom margin: & 2.5 cm
\end{tabular}
\end{center}

These settings are predefined in the official template and shall not be
modified by students.

\vspace{0.3cm}

% ================================================================
% IV. FONT, FONT SIZE, AND LINE SPACING
% ================================================================

\noindent
\textbf{IV.\quad FONT, FONT SIZE, AND LINE SPACING}

\vspace{0.15cm}

The main text shall use a 13 pt font size with 1.3-line spacing.
A Times-family typeface shall be used in accordance with the official
template.

The text shall be fully justified.

Students shall not use repeated spaces, repeated manual line breaks, or
multiple blank lines to position text. Paragraph and heading spacing shall
be controlled by the template.

\vspace{0.3cm}

% ================================================================
% V. OUTER COVER
% ================================================================

\noindent
\textbf{V.\quad OUTER COVER}

\vspace{0.15cm}

The outer cover shall follow the official format provided in the template
and shall contain the following information:

\begin{itemize}[leftmargin=0.8cm, itemsep=2pt]
    \item Thai Nguyen University of Technology;
    \item Faculty of International Training;
    \item Engineering and Technology Department;
    \item University logo and Faculty logo;
    \item the title \textbf{INTERNSHIP REPORT};
    \item Internship Organization;
    \item Author;
    \item Class;
    \item Supervisor;
    \item place and year of completion.
\end{itemize}

Students shall only replace the required information fields and shall not
modify the cover-page layout unless instructed otherwise.

\vspace{0.3cm}

% ================================================================
% VI. COMMENTS OF THE INTERNSHIP HOST ORGANIZATION
% ================================================================

\noindent
\textbf{VI.\quad COMMENTS OF THE INTERNSHIP HOST ORGANIZATION}

\vspace{0.15cm}

The page entitled \textbf{COMMENTS OF THE INTERNSHIP HOST ORGANIZATION}
is reserved for the internship host organization.

It includes:

\begin{itemize}[leftmargin=0.8cm, itemsep=2pt]
    \item the student's full name;
    \item class;
    \item comments from the host organization;
    \item score;
    \item date;
    \item signature and official seal.
\end{itemize}

Students shall not alter or complete official evaluation fields unless
specifically authorized.

\vspace{0.3cm}

% ================================================================
% VII. SUPERVISOR'S EVALUATION AND DEFENSE RESULTS
% ================================================================

\noindent
\textbf{VII.\quad SUPERVISOR'S EVALUATION AND DEFENSE RESULTS}

\vspace{0.15cm}

The section entitled
\textbf{SUPERVISOR'S EVALUATION AND GRADUATION INTERNSHIP DEFENSE RESULTS}
is reserved for the supervisor and defense committee.

The evaluation may include:

\begin{itemize}[leftmargin=0.8cm, itemsep=2pt]
    \item Internship Progress and Student's Attitude;
    \item Report Content;
    \item Report Presentation;
    \item Other Comments;
    \item Supervisor's Opinion;
    \item Defense Result;
    \item Score.
\end{itemize}

Students shall not enter evaluation scores or defense results before they
are officially determined.

\vspace{0.3cm}

% ================================================================
% VIII. CHAPTERS, SECTIONS, AND SUBSECTIONS
% ================================================================

\noindent
\textbf{VIII.\quad CHAPTERS, SECTIONS, AND SUBSECTIONS}

\vspace{0.15cm}

The report structure shall be generated using the standard \LaTeX{}
commands:

\begin{verbatim}
\chapter{...}
\section{...}
\subsection{...}
\end{verbatim}

Chapter, section, and subsection numbers shall be generated automatically
by \LaTeX{}.

Students shall not manually type chapter or section numbers and shall not
reproduce heading styles using commands such as \verb|\textbf|,
\verb|\fontsize|, or \verb|\centering| instead of the appropriate
structural command.

Excessive subdivision should be avoided. In most cases, the hierarchy
Chapter--Section--Subsection is sufficient.

\vspace{0.3cm}

% ================================================================
% IX. ORGANIZATION OF THE MAIN CONTENT
% ================================================================

\noindent
\textbf{IX.\quad ORGANIZATION OF THE MAIN CONTENT}

\vspace{0.15cm}

The specific content of the report shall depend on the internship
organization and the assigned tasks. A typical structure may include:

\begin{itemize}[leftmargin=0.8cm, itemsep=2pt]
    \item \textbf{Chapter 1. Overview of the Internship Organization};
    \item \textbf{Chapter 2. Internship Activities and Assigned Tasks};
    \item \textbf{Chapter 3. Results, Experience, and Evaluation}.
\end{itemize}

These chapter titles are recommendations only and may be adapted to the
actual internship content.

The report shall reflect work actually performed, observed, analyzed, or
learned during the internship. Unrelated theoretical material should not
be included merely to increase report length.

\vspace{0.3cm}

% ================================================================
% X. TABLE OF CONTENTS
% ================================================================

\noindent
\textbf{X.\quad TABLE OF CONTENTS}

\vspace{0.15cm}

The Table of Contents shall be generated automatically using:

\begin{verbatim}
\tableofcontents
\end{verbatim}

Students shall not create or manually edit page numbers in the Table of
Contents.

The document shall be compiled sufficiently to ensure that all page
numbers and entries are correctly updated.

\vspace{0.3cm}

% ================================================================
% XI. LIST OF SYMBOLS AND ABBREVIATIONS
% ================================================================

\noindent
\textbf{XI.\quad LIST OF SYMBOLS AND ABBREVIATIONS}

\vspace{0.15cm}

Symbols and abbreviations that are used repeatedly in the report shall be
listed in the \textbf{List of Symbols and Abbreviations}.

The list may contain the following columns:

\begin{center}
\begin{tabular}{llll}
    No. & Symbol / Abbreviation & Unit & Definition
\end{tabular}
\end{center}

Only symbols and abbreviations actually used in the final report shall be
retained. Example entries contained in the template shall be removed if
they are not relevant to the student's report.

Each abbreviation should normally be defined when it first appears in
the main text.

\vspace{0.3cm}

% ================================================================
% XII. FIGURES, DIAGRAMS, AND PHOTOGRAPHS
% ================================================================

\noindent
\textbf{XII.\quad FIGURES, DIAGRAMS, AND PHOTOGRAPHS}

\vspace{0.15cm}

Figures shall be clear, legible, and directly relevant to the report
content.

Figures shall be inserted using the standard \verb|figure| environment.

Example:

\begin{verbatim}
\begin{figure}[H]
    \centering
    \includegraphics[width=0.8\textwidth]
                    {figures/company_structure.pdf}
    \caption{Organizational structure of the company}
    \label{fig:company-structure}
\end{figure}
\end{verbatim}

Figure captions shall be placed below the figure.

Each referenced figure shall have a unique \verb|\label| and shall be
referenced in the text using, for example:

\begin{verbatim}
Figure~\ref{fig:company-structure}
\end{verbatim}

Students shall not manually type figure numbers.

Figures reproduced or adapted from external sources shall include an
appropriate citation.

Photographs taken during the internship shall be relevant and sufficiently
clear. Students shall not include confidential, restricted, or proprietary
information without permission from the internship organization.

Vector graphics or high-resolution images are recommended whenever
possible.

\vspace{0.3cm}

% ================================================================
% XIII. TABLES
% ================================================================

\noindent
\textbf{XIII.\quad TABLES}

\vspace{0.15cm}

Tables shall be created using suitable \LaTeX{} environments such as
\verb|table|, \verb|tabular|, \verb|tabularx|, or \verb|longtable|.

Table captions shall be placed above the table.

Example:

\begin{verbatim}
\begin{table}[H]
    \centering
    \caption{Internship schedule}
    \label{tab:internship-schedule}

    \begin{tabular}{lll}
        \toprule
        Week & Task & Result \\
        \midrule
        1 & Orientation & Completed \\
        2 & System study & Completed \\
        \bottomrule
    \end{tabular}
\end{table}
\end{verbatim}

Tables shall be referenced using:

\begin{verbatim}
Table~\ref{tab:internship-schedule}
\end{verbatim}

Students shall not manually type table numbers.

Tables should be created directly in \LaTeX{} rather than inserted as
screenshots whenever practical.

\vspace{0.3cm}

% ================================================================
% XIV. EQUATIONS AND TECHNICAL QUANTITIES
% ================================================================

\noindent
\textbf{XIV.\quad EQUATIONS AND TECHNICAL QUANTITIES}

\vspace{0.15cm}

Mathematical expressions shall be typeset directly using \LaTeX{}.
Images or screenshots of equations shall not be used.

Short expressions may be written in inline mathematical mode, for example:

\begin{verbatim}
$R = 0.7\,\Omega$
\end{verbatim}

Important equations shall use a numbered equation environment.

Example:

\begin{verbatim}
\begin{equation}
    P = UI
    \label{eq:power}
\end{equation}
\end{verbatim}

Equations shall be referenced using:

\begin{verbatim}
Equation~\eqref{eq:power}
\end{verbatim}

Equation numbers shall not be manually entered.

SI units shall be used wherever applicable. The \verb|siunitx| package
should be used where appropriate.

Examples:

\begin{verbatim}
\SI{220}{\volt}
\SI{50}{\hertz}
\SI{2.5}{\ampere}
\end{verbatim}

Notation and units shall be consistent throughout the report.

\vspace{0.3cm}

% ================================================================
% XV. SOURCE CODE
% ================================================================

\noindent
\textbf{XV.\quad SOURCE CODE}

\vspace{0.15cm}

If the internship involves programming, PLCs, microcontrollers,
simulation, or software development, relevant source code may be included
in the report.

Source code shall be presented using the \verb|listings| environment or
another environment defined by the official template.

Example:

\begin{verbatim}
\begin{lstlisting}[language=st]
IF Start = TRUE THEN
    Motor := TRUE;
ELSE
    Motor := FALSE;
END_IF;
\end{lstlisting}
\end{verbatim}

Only code necessary for explaining the work should appear in the main
report. Long source-code listings should be moved to an appendix or
submitted as supplementary files if required.

Source code should not normally be inserted as screenshots.

\vspace{0.3cm}

% ================================================================
% XVI. ABSTRACT
% ================================================================

\noindent
\textbf{XVI.\quad ABSTRACT}

\vspace{0.15cm}

The Abstract shall briefly summarize:

\begin{itemize}[leftmargin=0.8cm, itemsep=2pt]
    \item the internship organization;
    \item the objectives of the internship;
    \item the main tasks performed;
    \item the main results, knowledge, and skills obtained.
\end{itemize}

The Abstract should be concise and should not reproduce the full content
of the report chapter by chapter.

\vspace{0.3cm}

% ================================================================
% XVII. REFERENCES AND CITATIONS
% ================================================================

\noindent
\textbf{XVII.\quad REFERENCES AND CITATIONS}

\vspace{0.15cm}

All books, journal papers, conference papers, standards, websites,
datasheets, manuals, technical documents, figures, tables, or other
external sources used in preparing the report shall be properly cited.

References shall be managed through the bibliography system provided with
the official template.

Bibliographic entries shall normally be stored in:

\begin{verbatim}
References.bib
\end{verbatim}

References shall be cited in the report using:

\begin{verbatim}
\cite{reference-key}
\end{verbatim}

Students shall not manually type reference numbers such as ``[1]'',
``[2]'', etc.

Titles of foreign-language publications shall remain in their original
language in the bibliography.

All borrowed ideas, data, figures, tables, equations, methods, or results
shall be properly acknowledged.

\vspace{0.3cm}

% ================================================================
% XVIII. CROSS-REFERENCING
% ================================================================

\noindent
\textbf{XVIII.\quad CROSS-REFERENCING}

\vspace{0.15cm}

Numbered elements shall be cross-referenced using the standard \LaTeX{}
mechanism.

Use the following commands where appropriate:

\begin{verbatim}
\label{...}
\ref{...}
\eqref{...}
\end{verbatim}

For example:

\begin{verbatim}
As shown in Figure~\ref{fig:system}, ...
\end{verbatim}

Students shall not manually type numbers such as ``Figure 2.3'' when
those numbers can be generated automatically.

\vspace{0.3cm}

% ================================================================
% XIX. PAGE NUMBERING
% ================================================================

\noindent
\textbf{XIX.\quad PAGE NUMBERING}

\vspace{0.15cm}

Cover pages and official evaluation forms at the beginning of the report
shall not display page numbers.

The main body of the report shall use Arabic page numbers beginning from
Chapter 1.

Page numbering shall be generated automatically by the official template.
Students shall not manually enter page numbers.

\vspace{0.3cm}

% ================================================================
% XX. PROJECT FILE ORGANIZATION
% ================================================================

\noindent
\textbf{XX.\quad PROJECT FILE ORGANIZATION}

\vspace{0.15cm}

Students should maintain a clear and organized \LaTeX{} project
structure. A recommended structure is:

\begin{verbatim}
InternshipReport/
|
|-- main.tex
|-- References.bib
|
|-- figures/
|   |-- company.pdf
|   |-- system.png
|   `-- ...
|
`-- chapters/
    |-- chapter1.tex
    |-- chapter2.tex
    `-- chapter3.tex
\end{verbatim}

Figures should normally be stored in the \verb|figures| directory
specified by the official template.

File and directory names should preferably contain Latin letters,
numbers, hyphens, or underscores.

Recommended examples include:

\begin{verbatim}
company_structure.pdf
motor_control.png
chapter_01.tex
\end{verbatim}

Spaces and unnecessary special characters in file names should be
avoided where possible.

\vspace{0.3cm}

% ================================================================
% XXI. DOCUMENT COMPILATION AND FINAL CHECK
% ================================================================

\noindent
\textbf{XXI.\quad DOCUMENT COMPILATION AND FINAL CHECK}

\vspace{0.15cm}

Before submission, the complete \LaTeX{} project shall compile
successfully and produce the final PDF report.

The final document shall not contain unresolved references such as:

\begin{verbatim}
??
[?]
\end{verbatim}

Before generating the final PDF, students shall check:

\begin{itemize}[leftmargin=0.8cm, itemsep=2pt]
    \item the Table of Contents;
    \item page numbering;
    \item chapter, section, and subsection numbering;
    \item figure numbering;
    \item table numbering;
    \item equation numbering;
    \item cross-references;
    \item citations;
    \item bibliography;
    \item font display;
    \item mathematical symbols;
    \item image quality;
    \item completeness of all required forms.
\end{itemize}

The document may be compiled using TeX Live, MiKTeX, or another
compatible \LaTeX{} environment. Regardless of the compilation
environment, the resulting PDF shall conform to the official Faculty
template.

\vspace{0.3cm}

% ================================================================
% XXII. CONTENT THAT STUDENTS MAY EDIT
% ================================================================

\noindent
\textbf{XXII.\quad CONTENT THAT STUDENTS MAY EDIT}

\vspace{0.15cm}

Students may edit or replace:

\begin{itemize}[leftmargin=0.8cm, itemsep=2pt]
    \item student name;
    \item class;
    \item internship organization;
    \item supervisor name;
    \item Abstract;
    \item List of Symbols and Abbreviations;
    \item chapter titles and chapter content;
    \item figures and photographs;
    \item tables;
    \item equations;
    \item citations and references.
\end{itemize}

Students shall not modify the main document formatting settings in the
preamble, including the document class, page geometry, line spacing,
heading formatting, table-of-contents formatting, font configuration,
page layout, or cover-page macros unless specifically instructed.

\vspace{0.3cm}

% ================================================================
% XXIII. FINAL SUBMISSION REQUIREMENTS
% ================================================================

\noindent
\textbf{XXIII.\quad FINAL SUBMISSION REQUIREMENTS}

\vspace{0.15cm}

Before submission, the final report shall:

\begin{itemize}[leftmargin=0.8cm, itemsep=2pt]
    \item compile without errors that affect the report content or presentation;
    \item contain all required forms and report sections;
    \item contain correct student, class, internship organization, and supervisor information;
    \item contain an updated Table of Contents, List of Figures, and List of Tables;
    \item contain no unused example content from the template;
    \item contain no placeholder text such as \textit{Title of Chapter 1},
    \textit{Section 1.1}, or other sample entries;
    \item contain all required citations and references;
    \item contain no unresolved cross-references.
\end{itemize}

If the Faculty requires submission of the \LaTeX{} source, the complete
project shall include all \verb|.tex| files, the bibliography database,
figures, and any other files required to reproduce the submitted PDF.

\vspace{0.4cm}

\noindent
\textbf{Important:}
Before submitting the final internship report, students shall remove or
comment out the following command in the main document:

\begin{verbatim}

% =================================================================
% REGULATIONS FOR PREPARING THE GRADUATION INTERNSHIP REPORT
% USING LATEX
% =================================================================
%
% This file is intended to be included from the main document with:
%
%     
% =================================================================
% REGULATIONS FOR PREPARING THE GRADUATION INTERNSHIP REPORT
% USING LATEX
% =================================================================
%
% This file is intended to be included from the main document with:
%
%     \input{Regulations.tex}
%
% IMPORTANT:
% This section is provided for guidance only. Students shall remove
% or comment out the \input{Regulations.tex} command before submitting
% the final internship report.
%
% =================================================================

\clearpage
\thispagestyle{empty}

\begin{center}
    {\fontsize{16pt}{20pt}\selectfont\bfseries
    SPECIFIC REGULATIONS FOR PREPARING\\
    GRADUATION INTERNSHIP REPORTS USING \LaTeX}
\end{center}

\vspace{0.5cm}

{\fontsize{13pt}{17pt}\selectfont
\setstretch{1.3}
\justifying

\noindent
\textbf{Note:}
This section is provided for guidance only and is included in the main
document using the command \verb|\input{Regulations.tex}|.
Students shall remove or comment out this command before submitting the
final internship report.

\vspace{0.4cm}

% ================================================================
% I. GENERAL REQUIREMENTS
% ================================================================

\noindent
\textbf{I.\quad GENERAL REQUIREMENTS}

\vspace{0.15cm}

The Graduation Internship Report shall be prepared using the official
\LaTeX{} template provided by the Faculty. Students shall preserve the
predefined page layout, typography, line spacing, heading styles,
numbering system, figure and table formatting, and bibliography settings
unless otherwise instructed by the Faculty or the supervisor.

The report shall be written in English.

Students are responsible for the academic and technical content of the
report. The official \LaTeX{} template is responsible for standardized
document formatting.

\vspace{0.3cm}

% ================================================================
% II. REPORT STRUCTURE
% ================================================================

\noindent
\textbf{II.\quad REPORT STRUCTURE}

\vspace{0.15cm}

The report shall normally be arranged in the following order:

\begin{enumerate}[label=\arabic*., leftmargin=1.2cm, itemsep=2pt]
    \item Outer Cover;
    \item Comments of the Internship Host Organization;
    \item Supervisor's Evaluation and Graduation Internship Defense Results;
    \item Table of Contents;
    \item List of Symbols and Abbreviations, if applicable;
    \item List of Figures;
    \item List of Tables;
    \item Abstract;
    \item Main Chapters of the Internship Report;
    \item References;
    \item Appendices, if applicable.
\end{enumerate}

The preliminary pages shall not be numbered as chapters. The main report
content shall be organized into consecutively numbered chapters beginning
with Chapter 1.

\vspace{0.3cm}

% ================================================================
% III. PAGE SIZE AND MARGINS
% ================================================================

\noindent
\textbf{III.\quad PAGE SIZE AND MARGINS}

\vspace{0.15cm}

The report shall be formatted on A4 paper with the following margins:

\begin{center}
\begin{tabular}{ll}
    Left margin:   & 3.5 cm \\
    Right margin:  & 2.0 cm \\
    Top margin:    & 2.5 cm \\
    Bottom margin: & 2.5 cm
\end{tabular}
\end{center}

These settings are predefined in the official template and shall not be
modified by students.

\vspace{0.3cm}

% ================================================================
% IV. FONT, FONT SIZE, AND LINE SPACING
% ================================================================

\noindent
\textbf{IV.\quad FONT, FONT SIZE, AND LINE SPACING}

\vspace{0.15cm}

The main text shall use a 13 pt font size with 1.3-line spacing.
A Times-family typeface shall be used in accordance with the official
template.

The text shall be fully justified.

Students shall not use repeated spaces, repeated manual line breaks, or
multiple blank lines to position text. Paragraph and heading spacing shall
be controlled by the template.

\vspace{0.3cm}

% ================================================================
% V. OUTER COVER
% ================================================================

\noindent
\textbf{V.\quad OUTER COVER}

\vspace{0.15cm}

The outer cover shall follow the official format provided in the template
and shall contain the following information:

\begin{itemize}[leftmargin=0.8cm, itemsep=2pt]
    \item Thai Nguyen University of Technology;
    \item Faculty of International Training;
    \item Engineering and Technology Department;
    \item University logo and Faculty logo;
    \item the title \textbf{INTERNSHIP REPORT};
    \item Internship Organization;
    \item Author;
    \item Class;
    \item Supervisor;
    \item place and year of completion.
\end{itemize}

Students shall only replace the required information fields and shall not
modify the cover-page layout unless instructed otherwise.

\vspace{0.3cm}

% ================================================================
% VI. COMMENTS OF THE INTERNSHIP HOST ORGANIZATION
% ================================================================

\noindent
\textbf{VI.\quad COMMENTS OF THE INTERNSHIP HOST ORGANIZATION}

\vspace{0.15cm}

The page entitled \textbf{COMMENTS OF THE INTERNSHIP HOST ORGANIZATION}
is reserved for the internship host organization.

It includes:

\begin{itemize}[leftmargin=0.8cm, itemsep=2pt]
    \item the student's full name;
    \item class;
    \item comments from the host organization;
    \item score;
    \item date;
    \item signature and official seal.
\end{itemize}

Students shall not alter or complete official evaluation fields unless
specifically authorized.

\vspace{0.3cm}

% ================================================================
% VII. SUPERVISOR'S EVALUATION AND DEFENSE RESULTS
% ================================================================

\noindent
\textbf{VII.\quad SUPERVISOR'S EVALUATION AND DEFENSE RESULTS}

\vspace{0.15cm}

The section entitled
\textbf{SUPERVISOR'S EVALUATION AND GRADUATION INTERNSHIP DEFENSE RESULTS}
is reserved for the supervisor and defense committee.

The evaluation may include:

\begin{itemize}[leftmargin=0.8cm, itemsep=2pt]
    \item Internship Progress and Student's Attitude;
    \item Report Content;
    \item Report Presentation;
    \item Other Comments;
    \item Supervisor's Opinion;
    \item Defense Result;
    \item Score.
\end{itemize}

Students shall not enter evaluation scores or defense results before they
are officially determined.

\vspace{0.3cm}

% ================================================================
% VIII. CHAPTERS, SECTIONS, AND SUBSECTIONS
% ================================================================

\noindent
\textbf{VIII.\quad CHAPTERS, SECTIONS, AND SUBSECTIONS}

\vspace{0.15cm}

The report structure shall be generated using the standard \LaTeX{}
commands:

\begin{verbatim}
\chapter{...}
\section{...}
\subsection{...}
\end{verbatim}

Chapter, section, and subsection numbers shall be generated automatically
by \LaTeX{}.

Students shall not manually type chapter or section numbers and shall not
reproduce heading styles using commands such as \verb|\textbf|,
\verb|\fontsize|, or \verb|\centering| instead of the appropriate
structural command.

Excessive subdivision should be avoided. In most cases, the hierarchy
Chapter--Section--Subsection is sufficient.

\vspace{0.3cm}

% ================================================================
% IX. ORGANIZATION OF THE MAIN CONTENT
% ================================================================

\noindent
\textbf{IX.\quad ORGANIZATION OF THE MAIN CONTENT}

\vspace{0.15cm}

The specific content of the report shall depend on the internship
organization and the assigned tasks. A typical structure may include:

\begin{itemize}[leftmargin=0.8cm, itemsep=2pt]
    \item \textbf{Chapter 1. Overview of the Internship Organization};
    \item \textbf{Chapter 2. Internship Activities and Assigned Tasks};
    \item \textbf{Chapter 3. Results, Experience, and Evaluation}.
\end{itemize}

These chapter titles are recommendations only and may be adapted to the
actual internship content.

The report shall reflect work actually performed, observed, analyzed, or
learned during the internship. Unrelated theoretical material should not
be included merely to increase report length.

\vspace{0.3cm}

% ================================================================
% X. TABLE OF CONTENTS
% ================================================================

\noindent
\textbf{X.\quad TABLE OF CONTENTS}

\vspace{0.15cm}

The Table of Contents shall be generated automatically using:

\begin{verbatim}
\tableofcontents
\end{verbatim}

Students shall not create or manually edit page numbers in the Table of
Contents.

The document shall be compiled sufficiently to ensure that all page
numbers and entries are correctly updated.

\vspace{0.3cm}

% ================================================================
% XI. LIST OF SYMBOLS AND ABBREVIATIONS
% ================================================================

\noindent
\textbf{XI.\quad LIST OF SYMBOLS AND ABBREVIATIONS}

\vspace{0.15cm}

Symbols and abbreviations that are used repeatedly in the report shall be
listed in the \textbf{List of Symbols and Abbreviations}.

The list may contain the following columns:

\begin{center}
\begin{tabular}{llll}
    No. & Symbol / Abbreviation & Unit & Definition
\end{tabular}
\end{center}

Only symbols and abbreviations actually used in the final report shall be
retained. Example entries contained in the template shall be removed if
they are not relevant to the student's report.

Each abbreviation should normally be defined when it first appears in
the main text.

\vspace{0.3cm}

% ================================================================
% XII. FIGURES, DIAGRAMS, AND PHOTOGRAPHS
% ================================================================

\noindent
\textbf{XII.\quad FIGURES, DIAGRAMS, AND PHOTOGRAPHS}

\vspace{0.15cm}

Figures shall be clear, legible, and directly relevant to the report
content.

Figures shall be inserted using the standard \verb|figure| environment.

Example:

\begin{verbatim}
\begin{figure}[H]
    \centering
    \includegraphics[width=0.8\textwidth]
                    {figures/company_structure.pdf}
    \caption{Organizational structure of the company}
    \label{fig:company-structure}
\end{figure}
\end{verbatim}

Figure captions shall be placed below the figure.

Each referenced figure shall have a unique \verb|\label| and shall be
referenced in the text using, for example:

\begin{verbatim}
Figure~\ref{fig:company-structure}
\end{verbatim}

Students shall not manually type figure numbers.

Figures reproduced or adapted from external sources shall include an
appropriate citation.

Photographs taken during the internship shall be relevant and sufficiently
clear. Students shall not include confidential, restricted, or proprietary
information without permission from the internship organization.

Vector graphics or high-resolution images are recommended whenever
possible.

\vspace{0.3cm}

% ================================================================
% XIII. TABLES
% ================================================================

\noindent
\textbf{XIII.\quad TABLES}

\vspace{0.15cm}

Tables shall be created using suitable \LaTeX{} environments such as
\verb|table|, \verb|tabular|, \verb|tabularx|, or \verb|longtable|.

Table captions shall be placed above the table.

Example:

\begin{verbatim}
\begin{table}[H]
    \centering
    \caption{Internship schedule}
    \label{tab:internship-schedule}

    \begin{tabular}{lll}
        \toprule
        Week & Task & Result \\
        \midrule
        1 & Orientation & Completed \\
        2 & System study & Completed \\
        \bottomrule
    \end{tabular}
\end{table}
\end{verbatim}

Tables shall be referenced using:

\begin{verbatim}
Table~\ref{tab:internship-schedule}
\end{verbatim}

Students shall not manually type table numbers.

Tables should be created directly in \LaTeX{} rather than inserted as
screenshots whenever practical.

\vspace{0.3cm}

% ================================================================
% XIV. EQUATIONS AND TECHNICAL QUANTITIES
% ================================================================

\noindent
\textbf{XIV.\quad EQUATIONS AND TECHNICAL QUANTITIES}

\vspace{0.15cm}

Mathematical expressions shall be typeset directly using \LaTeX{}.
Images or screenshots of equations shall not be used.

Short expressions may be written in inline mathematical mode, for example:

\begin{verbatim}
$R = 0.7\,\Omega$
\end{verbatim}

Important equations shall use a numbered equation environment.

Example:

\begin{verbatim}
\begin{equation}
    P = UI
    \label{eq:power}
\end{equation}
\end{verbatim}

Equations shall be referenced using:

\begin{verbatim}
Equation~\eqref{eq:power}
\end{verbatim}

Equation numbers shall not be manually entered.

SI units shall be used wherever applicable. The \verb|siunitx| package
should be used where appropriate.

Examples:

\begin{verbatim}
\SI{220}{\volt}
\SI{50}{\hertz}
\SI{2.5}{\ampere}
\end{verbatim}

Notation and units shall be consistent throughout the report.

\vspace{0.3cm}

% ================================================================
% XV. SOURCE CODE
% ================================================================

\noindent
\textbf{XV.\quad SOURCE CODE}

\vspace{0.15cm}

If the internship involves programming, PLCs, microcontrollers,
simulation, or software development, relevant source code may be included
in the report.

Source code shall be presented using the \verb|listings| environment or
another environment defined by the official template.

Example:

\begin{verbatim}
\begin{lstlisting}[language=st]
IF Start = TRUE THEN
    Motor := TRUE;
ELSE
    Motor := FALSE;
END_IF;
\end{lstlisting}
\end{verbatim}

Only code necessary for explaining the work should appear in the main
report. Long source-code listings should be moved to an appendix or
submitted as supplementary files if required.

Source code should not normally be inserted as screenshots.

\vspace{0.3cm}

% ================================================================
% XVI. ABSTRACT
% ================================================================

\noindent
\textbf{XVI.\quad ABSTRACT}

\vspace{0.15cm}

The Abstract shall briefly summarize:

\begin{itemize}[leftmargin=0.8cm, itemsep=2pt]
    \item the internship organization;
    \item the objectives of the internship;
    \item the main tasks performed;
    \item the main results, knowledge, and skills obtained.
\end{itemize}

The Abstract should be concise and should not reproduce the full content
of the report chapter by chapter.

\vspace{0.3cm}

% ================================================================
% XVII. REFERENCES AND CITATIONS
% ================================================================

\noindent
\textbf{XVII.\quad REFERENCES AND CITATIONS}

\vspace{0.15cm}

All books, journal papers, conference papers, standards, websites,
datasheets, manuals, technical documents, figures, tables, or other
external sources used in preparing the report shall be properly cited.

References shall be managed through the bibliography system provided with
the official template.

Bibliographic entries shall normally be stored in:

\begin{verbatim}
References.bib
\end{verbatim}

References shall be cited in the report using:

\begin{verbatim}
\cite{reference-key}
\end{verbatim}

Students shall not manually type reference numbers such as ``[1]'',
``[2]'', etc.

Titles of foreign-language publications shall remain in their original
language in the bibliography.

All borrowed ideas, data, figures, tables, equations, methods, or results
shall be properly acknowledged.

\vspace{0.3cm}

% ================================================================
% XVIII. CROSS-REFERENCING
% ================================================================

\noindent
\textbf{XVIII.\quad CROSS-REFERENCING}

\vspace{0.15cm}

Numbered elements shall be cross-referenced using the standard \LaTeX{}
mechanism.

Use the following commands where appropriate:

\begin{verbatim}
\label{...}
\ref{...}
\eqref{...}
\end{verbatim}

For example:

\begin{verbatim}
As shown in Figure~\ref{fig:system}, ...
\end{verbatim}

Students shall not manually type numbers such as ``Figure 2.3'' when
those numbers can be generated automatically.

\vspace{0.3cm}

% ================================================================
% XIX. PAGE NUMBERING
% ================================================================

\noindent
\textbf{XIX.\quad PAGE NUMBERING}

\vspace{0.15cm}

Cover pages and official evaluation forms at the beginning of the report
shall not display page numbers.

The main body of the report shall use Arabic page numbers beginning from
Chapter 1.

Page numbering shall be generated automatically by the official template.
Students shall not manually enter page numbers.

\vspace{0.3cm}

% ================================================================
% XX. PROJECT FILE ORGANIZATION
% ================================================================

\noindent
\textbf{XX.\quad PROJECT FILE ORGANIZATION}

\vspace{0.15cm}

Students should maintain a clear and organized \LaTeX{} project
structure. A recommended structure is:

\begin{verbatim}
InternshipReport/
|
|-- main.tex
|-- References.bib
|
|-- figures/
|   |-- company.pdf
|   |-- system.png
|   `-- ...
|
`-- chapters/
    |-- chapter1.tex
    |-- chapter2.tex
    `-- chapter3.tex
\end{verbatim}

Figures should normally be stored in the \verb|figures| directory
specified by the official template.

File and directory names should preferably contain Latin letters,
numbers, hyphens, or underscores.

Recommended examples include:

\begin{verbatim}
company_structure.pdf
motor_control.png
chapter_01.tex
\end{verbatim}

Spaces and unnecessary special characters in file names should be
avoided where possible.

\vspace{0.3cm}

% ================================================================
% XXI. DOCUMENT COMPILATION AND FINAL CHECK
% ================================================================

\noindent
\textbf{XXI.\quad DOCUMENT COMPILATION AND FINAL CHECK}

\vspace{0.15cm}

Before submission, the complete \LaTeX{} project shall compile
successfully and produce the final PDF report.

The final document shall not contain unresolved references such as:

\begin{verbatim}
??
[?]
\end{verbatim}

Before generating the final PDF, students shall check:

\begin{itemize}[leftmargin=0.8cm, itemsep=2pt]
    \item the Table of Contents;
    \item page numbering;
    \item chapter, section, and subsection numbering;
    \item figure numbering;
    \item table numbering;
    \item equation numbering;
    \item cross-references;
    \item citations;
    \item bibliography;
    \item font display;
    \item mathematical symbols;
    \item image quality;
    \item completeness of all required forms.
\end{itemize}

The document may be compiled using TeX Live, MiKTeX, or another
compatible \LaTeX{} environment. Regardless of the compilation
environment, the resulting PDF shall conform to the official Faculty
template.

\vspace{0.3cm}

% ================================================================
% XXII. CONTENT THAT STUDENTS MAY EDIT
% ================================================================

\noindent
\textbf{XXII.\quad CONTENT THAT STUDENTS MAY EDIT}

\vspace{0.15cm}

Students may edit or replace:

\begin{itemize}[leftmargin=0.8cm, itemsep=2pt]
    \item student name;
    \item class;
    \item internship organization;
    \item supervisor name;
    \item Abstract;
    \item List of Symbols and Abbreviations;
    \item chapter titles and chapter content;
    \item figures and photographs;
    \item tables;
    \item equations;
    \item citations and references.
\end{itemize}

Students shall not modify the main document formatting settings in the
preamble, including the document class, page geometry, line spacing,
heading formatting, table-of-contents formatting, font configuration,
page layout, or cover-page macros unless specifically instructed.

\vspace{0.3cm}

% ================================================================
% XXIII. FINAL SUBMISSION REQUIREMENTS
% ================================================================

\noindent
\textbf{XXIII.\quad FINAL SUBMISSION REQUIREMENTS}

\vspace{0.15cm}

Before submission, the final report shall:

\begin{itemize}[leftmargin=0.8cm, itemsep=2pt]
    \item compile without errors that affect the report content or presentation;
    \item contain all required forms and report sections;
    \item contain correct student, class, internship organization, and supervisor information;
    \item contain an updated Table of Contents, List of Figures, and List of Tables;
    \item contain no unused example content from the template;
    \item contain no placeholder text such as \textit{Title of Chapter 1},
    \textit{Section 1.1}, or other sample entries;
    \item contain all required citations and references;
    \item contain no unresolved cross-references.
\end{itemize}

If the Faculty requires submission of the \LaTeX{} source, the complete
project shall include all \verb|.tex| files, the bibliography database,
figures, and any other files required to reproduce the submitted PDF.

\vspace{0.4cm}

\noindent
\textbf{Important:}
Before submitting the final internship report, students shall remove or
comment out the following command in the main document:

\begin{verbatim}
\input{Regulations.tex}
\end{verbatim}

For example:

\begin{verbatim}
% \input{Regulations.tex}
\end{verbatim}

This removes the guidance pages from the final submitted report while
preserving the regulations file in the \LaTeX{} project for reference.

} % End 13pt content

\clearpage

%
% IMPORTANT:
% This section is provided for guidance only. Students shall remove
% or comment out the 
% =================================================================
% REGULATIONS FOR PREPARING THE GRADUATION INTERNSHIP REPORT
% USING LATEX
% =================================================================
%
% This file is intended to be included from the main document with:
%
%     \input{Regulations.tex}
%
% IMPORTANT:
% This section is provided for guidance only. Students shall remove
% or comment out the \input{Regulations.tex} command before submitting
% the final internship report.
%
% =================================================================

\clearpage
\thispagestyle{empty}

\begin{center}
    {\fontsize{16pt}{20pt}\selectfont\bfseries
    SPECIFIC REGULATIONS FOR PREPARING\\
    GRADUATION INTERNSHIP REPORTS USING \LaTeX}
\end{center}

\vspace{0.5cm}

{\fontsize{13pt}{17pt}\selectfont
\setstretch{1.3}
\justifying

\noindent
\textbf{Note:}
This section is provided for guidance only and is included in the main
document using the command \verb|\input{Regulations.tex}|.
Students shall remove or comment out this command before submitting the
final internship report.

\vspace{0.4cm}

% ================================================================
% I. GENERAL REQUIREMENTS
% ================================================================

\noindent
\textbf{I.\quad GENERAL REQUIREMENTS}

\vspace{0.15cm}

The Graduation Internship Report shall be prepared using the official
\LaTeX{} template provided by the Faculty. Students shall preserve the
predefined page layout, typography, line spacing, heading styles,
numbering system, figure and table formatting, and bibliography settings
unless otherwise instructed by the Faculty or the supervisor.

The report shall be written in English.

Students are responsible for the academic and technical content of the
report. The official \LaTeX{} template is responsible for standardized
document formatting.

\vspace{0.3cm}

% ================================================================
% II. REPORT STRUCTURE
% ================================================================

\noindent
\textbf{II.\quad REPORT STRUCTURE}

\vspace{0.15cm}

The report shall normally be arranged in the following order:

\begin{enumerate}[label=\arabic*., leftmargin=1.2cm, itemsep=2pt]
    \item Outer Cover;
    \item Comments of the Internship Host Organization;
    \item Supervisor's Evaluation and Graduation Internship Defense Results;
    \item Table of Contents;
    \item List of Symbols and Abbreviations, if applicable;
    \item List of Figures;
    \item List of Tables;
    \item Abstract;
    \item Main Chapters of the Internship Report;
    \item References;
    \item Appendices, if applicable.
\end{enumerate}

The preliminary pages shall not be numbered as chapters. The main report
content shall be organized into consecutively numbered chapters beginning
with Chapter 1.

\vspace{0.3cm}

% ================================================================
% III. PAGE SIZE AND MARGINS
% ================================================================

\noindent
\textbf{III.\quad PAGE SIZE AND MARGINS}

\vspace{0.15cm}

The report shall be formatted on A4 paper with the following margins:

\begin{center}
\begin{tabular}{ll}
    Left margin:   & 3.5 cm \\
    Right margin:  & 2.0 cm \\
    Top margin:    & 2.5 cm \\
    Bottom margin: & 2.5 cm
\end{tabular}
\end{center}

These settings are predefined in the official template and shall not be
modified by students.

\vspace{0.3cm}

% ================================================================
% IV. FONT, FONT SIZE, AND LINE SPACING
% ================================================================

\noindent
\textbf{IV.\quad FONT, FONT SIZE, AND LINE SPACING}

\vspace{0.15cm}

The main text shall use a 13 pt font size with 1.3-line spacing.
A Times-family typeface shall be used in accordance with the official
template.

The text shall be fully justified.

Students shall not use repeated spaces, repeated manual line breaks, or
multiple blank lines to position text. Paragraph and heading spacing shall
be controlled by the template.

\vspace{0.3cm}

% ================================================================
% V. OUTER COVER
% ================================================================

\noindent
\textbf{V.\quad OUTER COVER}

\vspace{0.15cm}

The outer cover shall follow the official format provided in the template
and shall contain the following information:

\begin{itemize}[leftmargin=0.8cm, itemsep=2pt]
    \item Thai Nguyen University of Technology;
    \item Faculty of International Training;
    \item Engineering and Technology Department;
    \item University logo and Faculty logo;
    \item the title \textbf{INTERNSHIP REPORT};
    \item Internship Organization;
    \item Author;
    \item Class;
    \item Supervisor;
    \item place and year of completion.
\end{itemize}

Students shall only replace the required information fields and shall not
modify the cover-page layout unless instructed otherwise.

\vspace{0.3cm}

% ================================================================
% VI. COMMENTS OF THE INTERNSHIP HOST ORGANIZATION
% ================================================================

\noindent
\textbf{VI.\quad COMMENTS OF THE INTERNSHIP HOST ORGANIZATION}

\vspace{0.15cm}

The page entitled \textbf{COMMENTS OF THE INTERNSHIP HOST ORGANIZATION}
is reserved for the internship host organization.

It includes:

\begin{itemize}[leftmargin=0.8cm, itemsep=2pt]
    \item the student's full name;
    \item class;
    \item comments from the host organization;
    \item score;
    \item date;
    \item signature and official seal.
\end{itemize}

Students shall not alter or complete official evaluation fields unless
specifically authorized.

\vspace{0.3cm}

% ================================================================
% VII. SUPERVISOR'S EVALUATION AND DEFENSE RESULTS
% ================================================================

\noindent
\textbf{VII.\quad SUPERVISOR'S EVALUATION AND DEFENSE RESULTS}

\vspace{0.15cm}

The section entitled
\textbf{SUPERVISOR'S EVALUATION AND GRADUATION INTERNSHIP DEFENSE RESULTS}
is reserved for the supervisor and defense committee.

The evaluation may include:

\begin{itemize}[leftmargin=0.8cm, itemsep=2pt]
    \item Internship Progress and Student's Attitude;
    \item Report Content;
    \item Report Presentation;
    \item Other Comments;
    \item Supervisor's Opinion;
    \item Defense Result;
    \item Score.
\end{itemize}

Students shall not enter evaluation scores or defense results before they
are officially determined.

\vspace{0.3cm}

% ================================================================
% VIII. CHAPTERS, SECTIONS, AND SUBSECTIONS
% ================================================================

\noindent
\textbf{VIII.\quad CHAPTERS, SECTIONS, AND SUBSECTIONS}

\vspace{0.15cm}

The report structure shall be generated using the standard \LaTeX{}
commands:

\begin{verbatim}
\chapter{...}
\section{...}
\subsection{...}
\end{verbatim}

Chapter, section, and subsection numbers shall be generated automatically
by \LaTeX{}.

Students shall not manually type chapter or section numbers and shall not
reproduce heading styles using commands such as \verb|\textbf|,
\verb|\fontsize|, or \verb|\centering| instead of the appropriate
structural command.

Excessive subdivision should be avoided. In most cases, the hierarchy
Chapter--Section--Subsection is sufficient.

\vspace{0.3cm}

% ================================================================
% IX. ORGANIZATION OF THE MAIN CONTENT
% ================================================================

\noindent
\textbf{IX.\quad ORGANIZATION OF THE MAIN CONTENT}

\vspace{0.15cm}

The specific content of the report shall depend on the internship
organization and the assigned tasks. A typical structure may include:

\begin{itemize}[leftmargin=0.8cm, itemsep=2pt]
    \item \textbf{Chapter 1. Overview of the Internship Organization};
    \item \textbf{Chapter 2. Internship Activities and Assigned Tasks};
    \item \textbf{Chapter 3. Results, Experience, and Evaluation}.
\end{itemize}

These chapter titles are recommendations only and may be adapted to the
actual internship content.

The report shall reflect work actually performed, observed, analyzed, or
learned during the internship. Unrelated theoretical material should not
be included merely to increase report length.

\vspace{0.3cm}

% ================================================================
% X. TABLE OF CONTENTS
% ================================================================

\noindent
\textbf{X.\quad TABLE OF CONTENTS}

\vspace{0.15cm}

The Table of Contents shall be generated automatically using:

\begin{verbatim}
\tableofcontents
\end{verbatim}

Students shall not create or manually edit page numbers in the Table of
Contents.

The document shall be compiled sufficiently to ensure that all page
numbers and entries are correctly updated.

\vspace{0.3cm}

% ================================================================
% XI. LIST OF SYMBOLS AND ABBREVIATIONS
% ================================================================

\noindent
\textbf{XI.\quad LIST OF SYMBOLS AND ABBREVIATIONS}

\vspace{0.15cm}

Symbols and abbreviations that are used repeatedly in the report shall be
listed in the \textbf{List of Symbols and Abbreviations}.

The list may contain the following columns:

\begin{center}
\begin{tabular}{llll}
    No. & Symbol / Abbreviation & Unit & Definition
\end{tabular}
\end{center}

Only symbols and abbreviations actually used in the final report shall be
retained. Example entries contained in the template shall be removed if
they are not relevant to the student's report.

Each abbreviation should normally be defined when it first appears in
the main text.

\vspace{0.3cm}

% ================================================================
% XII. FIGURES, DIAGRAMS, AND PHOTOGRAPHS
% ================================================================

\noindent
\textbf{XII.\quad FIGURES, DIAGRAMS, AND PHOTOGRAPHS}

\vspace{0.15cm}

Figures shall be clear, legible, and directly relevant to the report
content.

Figures shall be inserted using the standard \verb|figure| environment.

Example:

\begin{verbatim}
\begin{figure}[H]
    \centering
    \includegraphics[width=0.8\textwidth]
                    {figures/company_structure.pdf}
    \caption{Organizational structure of the company}
    \label{fig:company-structure}
\end{figure}
\end{verbatim}

Figure captions shall be placed below the figure.

Each referenced figure shall have a unique \verb|\label| and shall be
referenced in the text using, for example:

\begin{verbatim}
Figure~\ref{fig:company-structure}
\end{verbatim}

Students shall not manually type figure numbers.

Figures reproduced or adapted from external sources shall include an
appropriate citation.

Photographs taken during the internship shall be relevant and sufficiently
clear. Students shall not include confidential, restricted, or proprietary
information without permission from the internship organization.

Vector graphics or high-resolution images are recommended whenever
possible.

\vspace{0.3cm}

% ================================================================
% XIII. TABLES
% ================================================================

\noindent
\textbf{XIII.\quad TABLES}

\vspace{0.15cm}

Tables shall be created using suitable \LaTeX{} environments such as
\verb|table|, \verb|tabular|, \verb|tabularx|, or \verb|longtable|.

Table captions shall be placed above the table.

Example:

\begin{verbatim}
\begin{table}[H]
    \centering
    \caption{Internship schedule}
    \label{tab:internship-schedule}

    \begin{tabular}{lll}
        \toprule
        Week & Task & Result \\
        \midrule
        1 & Orientation & Completed \\
        2 & System study & Completed \\
        \bottomrule
    \end{tabular}
\end{table}
\end{verbatim}

Tables shall be referenced using:

\begin{verbatim}
Table~\ref{tab:internship-schedule}
\end{verbatim}

Students shall not manually type table numbers.

Tables should be created directly in \LaTeX{} rather than inserted as
screenshots whenever practical.

\vspace{0.3cm}

% ================================================================
% XIV. EQUATIONS AND TECHNICAL QUANTITIES
% ================================================================

\noindent
\textbf{XIV.\quad EQUATIONS AND TECHNICAL QUANTITIES}

\vspace{0.15cm}

Mathematical expressions shall be typeset directly using \LaTeX{}.
Images or screenshots of equations shall not be used.

Short expressions may be written in inline mathematical mode, for example:

\begin{verbatim}
$R = 0.7\,\Omega$
\end{verbatim}

Important equations shall use a numbered equation environment.

Example:

\begin{verbatim}
\begin{equation}
    P = UI
    \label{eq:power}
\end{equation}
\end{verbatim}

Equations shall be referenced using:

\begin{verbatim}
Equation~\eqref{eq:power}
\end{verbatim}

Equation numbers shall not be manually entered.

SI units shall be used wherever applicable. The \verb|siunitx| package
should be used where appropriate.

Examples:

\begin{verbatim}
\SI{220}{\volt}
\SI{50}{\hertz}
\SI{2.5}{\ampere}
\end{verbatim}

Notation and units shall be consistent throughout the report.

\vspace{0.3cm}

% ================================================================
% XV. SOURCE CODE
% ================================================================

\noindent
\textbf{XV.\quad SOURCE CODE}

\vspace{0.15cm}

If the internship involves programming, PLCs, microcontrollers,
simulation, or software development, relevant source code may be included
in the report.

Source code shall be presented using the \verb|listings| environment or
another environment defined by the official template.

Example:

\begin{verbatim}
\begin{lstlisting}[language=st]
IF Start = TRUE THEN
    Motor := TRUE;
ELSE
    Motor := FALSE;
END_IF;
\end{lstlisting}
\end{verbatim}

Only code necessary for explaining the work should appear in the main
report. Long source-code listings should be moved to an appendix or
submitted as supplementary files if required.

Source code should not normally be inserted as screenshots.

\vspace{0.3cm}

% ================================================================
% XVI. ABSTRACT
% ================================================================

\noindent
\textbf{XVI.\quad ABSTRACT}

\vspace{0.15cm}

The Abstract shall briefly summarize:

\begin{itemize}[leftmargin=0.8cm, itemsep=2pt]
    \item the internship organization;
    \item the objectives of the internship;
    \item the main tasks performed;
    \item the main results, knowledge, and skills obtained.
\end{itemize}

The Abstract should be concise and should not reproduce the full content
of the report chapter by chapter.

\vspace{0.3cm}

% ================================================================
% XVII. REFERENCES AND CITATIONS
% ================================================================

\noindent
\textbf{XVII.\quad REFERENCES AND CITATIONS}

\vspace{0.15cm}

All books, journal papers, conference papers, standards, websites,
datasheets, manuals, technical documents, figures, tables, or other
external sources used in preparing the report shall be properly cited.

References shall be managed through the bibliography system provided with
the official template.

Bibliographic entries shall normally be stored in:

\begin{verbatim}
References.bib
\end{verbatim}

References shall be cited in the report using:

\begin{verbatim}
\cite{reference-key}
\end{verbatim}

Students shall not manually type reference numbers such as ``[1]'',
``[2]'', etc.

Titles of foreign-language publications shall remain in their original
language in the bibliography.

All borrowed ideas, data, figures, tables, equations, methods, or results
shall be properly acknowledged.

\vspace{0.3cm}

% ================================================================
% XVIII. CROSS-REFERENCING
% ================================================================

\noindent
\textbf{XVIII.\quad CROSS-REFERENCING}

\vspace{0.15cm}

Numbered elements shall be cross-referenced using the standard \LaTeX{}
mechanism.

Use the following commands where appropriate:

\begin{verbatim}
\label{...}
\ref{...}
\eqref{...}
\end{verbatim}

For example:

\begin{verbatim}
As shown in Figure~\ref{fig:system}, ...
\end{verbatim}

Students shall not manually type numbers such as ``Figure 2.3'' when
those numbers can be generated automatically.

\vspace{0.3cm}

% ================================================================
% XIX. PAGE NUMBERING
% ================================================================

\noindent
\textbf{XIX.\quad PAGE NUMBERING}

\vspace{0.15cm}

Cover pages and official evaluation forms at the beginning of the report
shall not display page numbers.

The main body of the report shall use Arabic page numbers beginning from
Chapter 1.

Page numbering shall be generated automatically by the official template.
Students shall not manually enter page numbers.

\vspace{0.3cm}

% ================================================================
% XX. PROJECT FILE ORGANIZATION
% ================================================================

\noindent
\textbf{XX.\quad PROJECT FILE ORGANIZATION}

\vspace{0.15cm}

Students should maintain a clear and organized \LaTeX{} project
structure. A recommended structure is:

\begin{verbatim}
InternshipReport/
|
|-- main.tex
|-- References.bib
|
|-- figures/
|   |-- company.pdf
|   |-- system.png
|   `-- ...
|
`-- chapters/
    |-- chapter1.tex
    |-- chapter2.tex
    `-- chapter3.tex
\end{verbatim}

Figures should normally be stored in the \verb|figures| directory
specified by the official template.

File and directory names should preferably contain Latin letters,
numbers, hyphens, or underscores.

Recommended examples include:

\begin{verbatim}
company_structure.pdf
motor_control.png
chapter_01.tex
\end{verbatim}

Spaces and unnecessary special characters in file names should be
avoided where possible.

\vspace{0.3cm}

% ================================================================
% XXI. DOCUMENT COMPILATION AND FINAL CHECK
% ================================================================

\noindent
\textbf{XXI.\quad DOCUMENT COMPILATION AND FINAL CHECK}

\vspace{0.15cm}

Before submission, the complete \LaTeX{} project shall compile
successfully and produce the final PDF report.

The final document shall not contain unresolved references such as:

\begin{verbatim}
??
[?]
\end{verbatim}

Before generating the final PDF, students shall check:

\begin{itemize}[leftmargin=0.8cm, itemsep=2pt]
    \item the Table of Contents;
    \item page numbering;
    \item chapter, section, and subsection numbering;
    \item figure numbering;
    \item table numbering;
    \item equation numbering;
    \item cross-references;
    \item citations;
    \item bibliography;
    \item font display;
    \item mathematical symbols;
    \item image quality;
    \item completeness of all required forms.
\end{itemize}

The document may be compiled using TeX Live, MiKTeX, or another
compatible \LaTeX{} environment. Regardless of the compilation
environment, the resulting PDF shall conform to the official Faculty
template.

\vspace{0.3cm}

% ================================================================
% XXII. CONTENT THAT STUDENTS MAY EDIT
% ================================================================

\noindent
\textbf{XXII.\quad CONTENT THAT STUDENTS MAY EDIT}

\vspace{0.15cm}

Students may edit or replace:

\begin{itemize}[leftmargin=0.8cm, itemsep=2pt]
    \item student name;
    \item class;
    \item internship organization;
    \item supervisor name;
    \item Abstract;
    \item List of Symbols and Abbreviations;
    \item chapter titles and chapter content;
    \item figures and photographs;
    \item tables;
    \item equations;
    \item citations and references.
\end{itemize}

Students shall not modify the main document formatting settings in the
preamble, including the document class, page geometry, line spacing,
heading formatting, table-of-contents formatting, font configuration,
page layout, or cover-page macros unless specifically instructed.

\vspace{0.3cm}

% ================================================================
% XXIII. FINAL SUBMISSION REQUIREMENTS
% ================================================================

\noindent
\textbf{XXIII.\quad FINAL SUBMISSION REQUIREMENTS}

\vspace{0.15cm}

Before submission, the final report shall:

\begin{itemize}[leftmargin=0.8cm, itemsep=2pt]
    \item compile without errors that affect the report content or presentation;
    \item contain all required forms and report sections;
    \item contain correct student, class, internship organization, and supervisor information;
    \item contain an updated Table of Contents, List of Figures, and List of Tables;
    \item contain no unused example content from the template;
    \item contain no placeholder text such as \textit{Title of Chapter 1},
    \textit{Section 1.1}, or other sample entries;
    \item contain all required citations and references;
    \item contain no unresolved cross-references.
\end{itemize}

If the Faculty requires submission of the \LaTeX{} source, the complete
project shall include all \verb|.tex| files, the bibliography database,
figures, and any other files required to reproduce the submitted PDF.

\vspace{0.4cm}

\noindent
\textbf{Important:}
Before submitting the final internship report, students shall remove or
comment out the following command in the main document:

\begin{verbatim}
\input{Regulations.tex}
\end{verbatim}

For example:

\begin{verbatim}
% \input{Regulations.tex}
\end{verbatim}

This removes the guidance pages from the final submitted report while
preserving the regulations file in the \LaTeX{} project for reference.

} % End 13pt content

\clearpage
 command before submitting
% the final internship report.
%
% =================================================================

\clearpage
\thispagestyle{empty}

\begin{center}
    {\fontsize{16pt}{20pt}\selectfont\bfseries
    SPECIFIC REGULATIONS FOR PREPARING\\
    GRADUATION INTERNSHIP REPORTS USING \LaTeX}
\end{center}

\vspace{0.5cm}

{\fontsize{13pt}{17pt}\selectfont
\setstretch{1.3}
\justifying

\noindent
\textbf{Note:}
This section is provided for guidance only and is included in the main
document using the command \verb|
% =================================================================
% REGULATIONS FOR PREPARING THE GRADUATION INTERNSHIP REPORT
% USING LATEX
% =================================================================
%
% This file is intended to be included from the main document with:
%
%     \input{Regulations.tex}
%
% IMPORTANT:
% This section is provided for guidance only. Students shall remove
% or comment out the \input{Regulations.tex} command before submitting
% the final internship report.
%
% =================================================================

\clearpage
\thispagestyle{empty}

\begin{center}
    {\fontsize{16pt}{20pt}\selectfont\bfseries
    SPECIFIC REGULATIONS FOR PREPARING\\
    GRADUATION INTERNSHIP REPORTS USING \LaTeX}
\end{center}

\vspace{0.5cm}

{\fontsize{13pt}{17pt}\selectfont
\setstretch{1.3}
\justifying

\noindent
\textbf{Note:}
This section is provided for guidance only and is included in the main
document using the command \verb|\input{Regulations.tex}|.
Students shall remove or comment out this command before submitting the
final internship report.

\vspace{0.4cm}

% ================================================================
% I. GENERAL REQUIREMENTS
% ================================================================

\noindent
\textbf{I.\quad GENERAL REQUIREMENTS}

\vspace{0.15cm}

The Graduation Internship Report shall be prepared using the official
\LaTeX{} template provided by the Faculty. Students shall preserve the
predefined page layout, typography, line spacing, heading styles,
numbering system, figure and table formatting, and bibliography settings
unless otherwise instructed by the Faculty or the supervisor.

The report shall be written in English.

Students are responsible for the academic and technical content of the
report. The official \LaTeX{} template is responsible for standardized
document formatting.

\vspace{0.3cm}

% ================================================================
% II. REPORT STRUCTURE
% ================================================================

\noindent
\textbf{II.\quad REPORT STRUCTURE}

\vspace{0.15cm}

The report shall normally be arranged in the following order:

\begin{enumerate}[label=\arabic*., leftmargin=1.2cm, itemsep=2pt]
    \item Outer Cover;
    \item Comments of the Internship Host Organization;
    \item Supervisor's Evaluation and Graduation Internship Defense Results;
    \item Table of Contents;
    \item List of Symbols and Abbreviations, if applicable;
    \item List of Figures;
    \item List of Tables;
    \item Abstract;
    \item Main Chapters of the Internship Report;
    \item References;
    \item Appendices, if applicable.
\end{enumerate}

The preliminary pages shall not be numbered as chapters. The main report
content shall be organized into consecutively numbered chapters beginning
with Chapter 1.

\vspace{0.3cm}

% ================================================================
% III. PAGE SIZE AND MARGINS
% ================================================================

\noindent
\textbf{III.\quad PAGE SIZE AND MARGINS}

\vspace{0.15cm}

The report shall be formatted on A4 paper with the following margins:

\begin{center}
\begin{tabular}{ll}
    Left margin:   & 3.5 cm \\
    Right margin:  & 2.0 cm \\
    Top margin:    & 2.5 cm \\
    Bottom margin: & 2.5 cm
\end{tabular}
\end{center}

These settings are predefined in the official template and shall not be
modified by students.

\vspace{0.3cm}

% ================================================================
% IV. FONT, FONT SIZE, AND LINE SPACING
% ================================================================

\noindent
\textbf{IV.\quad FONT, FONT SIZE, AND LINE SPACING}

\vspace{0.15cm}

The main text shall use a 13 pt font size with 1.3-line spacing.
A Times-family typeface shall be used in accordance with the official
template.

The text shall be fully justified.

Students shall not use repeated spaces, repeated manual line breaks, or
multiple blank lines to position text. Paragraph and heading spacing shall
be controlled by the template.

\vspace{0.3cm}

% ================================================================
% V. OUTER COVER
% ================================================================

\noindent
\textbf{V.\quad OUTER COVER}

\vspace{0.15cm}

The outer cover shall follow the official format provided in the template
and shall contain the following information:

\begin{itemize}[leftmargin=0.8cm, itemsep=2pt]
    \item Thai Nguyen University of Technology;
    \item Faculty of International Training;
    \item Engineering and Technology Department;
    \item University logo and Faculty logo;
    \item the title \textbf{INTERNSHIP REPORT};
    \item Internship Organization;
    \item Author;
    \item Class;
    \item Supervisor;
    \item place and year of completion.
\end{itemize}

Students shall only replace the required information fields and shall not
modify the cover-page layout unless instructed otherwise.

\vspace{0.3cm}

% ================================================================
% VI. COMMENTS OF THE INTERNSHIP HOST ORGANIZATION
% ================================================================

\noindent
\textbf{VI.\quad COMMENTS OF THE INTERNSHIP HOST ORGANIZATION}

\vspace{0.15cm}

The page entitled \textbf{COMMENTS OF THE INTERNSHIP HOST ORGANIZATION}
is reserved for the internship host organization.

It includes:

\begin{itemize}[leftmargin=0.8cm, itemsep=2pt]
    \item the student's full name;
    \item class;
    \item comments from the host organization;
    \item score;
    \item date;
    \item signature and official seal.
\end{itemize}

Students shall not alter or complete official evaluation fields unless
specifically authorized.

\vspace{0.3cm}

% ================================================================
% VII. SUPERVISOR'S EVALUATION AND DEFENSE RESULTS
% ================================================================

\noindent
\textbf{VII.\quad SUPERVISOR'S EVALUATION AND DEFENSE RESULTS}

\vspace{0.15cm}

The section entitled
\textbf{SUPERVISOR'S EVALUATION AND GRADUATION INTERNSHIP DEFENSE RESULTS}
is reserved for the supervisor and defense committee.

The evaluation may include:

\begin{itemize}[leftmargin=0.8cm, itemsep=2pt]
    \item Internship Progress and Student's Attitude;
    \item Report Content;
    \item Report Presentation;
    \item Other Comments;
    \item Supervisor's Opinion;
    \item Defense Result;
    \item Score.
\end{itemize}

Students shall not enter evaluation scores or defense results before they
are officially determined.

\vspace{0.3cm}

% ================================================================
% VIII. CHAPTERS, SECTIONS, AND SUBSECTIONS
% ================================================================

\noindent
\textbf{VIII.\quad CHAPTERS, SECTIONS, AND SUBSECTIONS}

\vspace{0.15cm}

The report structure shall be generated using the standard \LaTeX{}
commands:

\begin{verbatim}
\chapter{...}
\section{...}
\subsection{...}
\end{verbatim}

Chapter, section, and subsection numbers shall be generated automatically
by \LaTeX{}.

Students shall not manually type chapter or section numbers and shall not
reproduce heading styles using commands such as \verb|\textbf|,
\verb|\fontsize|, or \verb|\centering| instead of the appropriate
structural command.

Excessive subdivision should be avoided. In most cases, the hierarchy
Chapter--Section--Subsection is sufficient.

\vspace{0.3cm}

% ================================================================
% IX. ORGANIZATION OF THE MAIN CONTENT
% ================================================================

\noindent
\textbf{IX.\quad ORGANIZATION OF THE MAIN CONTENT}

\vspace{0.15cm}

The specific content of the report shall depend on the internship
organization and the assigned tasks. A typical structure may include:

\begin{itemize}[leftmargin=0.8cm, itemsep=2pt]
    \item \textbf{Chapter 1. Overview of the Internship Organization};
    \item \textbf{Chapter 2. Internship Activities and Assigned Tasks};
    \item \textbf{Chapter 3. Results, Experience, and Evaluation}.
\end{itemize}

These chapter titles are recommendations only and may be adapted to the
actual internship content.

The report shall reflect work actually performed, observed, analyzed, or
learned during the internship. Unrelated theoretical material should not
be included merely to increase report length.

\vspace{0.3cm}

% ================================================================
% X. TABLE OF CONTENTS
% ================================================================

\noindent
\textbf{X.\quad TABLE OF CONTENTS}

\vspace{0.15cm}

The Table of Contents shall be generated automatically using:

\begin{verbatim}
\tableofcontents
\end{verbatim}

Students shall not create or manually edit page numbers in the Table of
Contents.

The document shall be compiled sufficiently to ensure that all page
numbers and entries are correctly updated.

\vspace{0.3cm}

% ================================================================
% XI. LIST OF SYMBOLS AND ABBREVIATIONS
% ================================================================

\noindent
\textbf{XI.\quad LIST OF SYMBOLS AND ABBREVIATIONS}

\vspace{0.15cm}

Symbols and abbreviations that are used repeatedly in the report shall be
listed in the \textbf{List of Symbols and Abbreviations}.

The list may contain the following columns:

\begin{center}
\begin{tabular}{llll}
    No. & Symbol / Abbreviation & Unit & Definition
\end{tabular}
\end{center}

Only symbols and abbreviations actually used in the final report shall be
retained. Example entries contained in the template shall be removed if
they are not relevant to the student's report.

Each abbreviation should normally be defined when it first appears in
the main text.

\vspace{0.3cm}

% ================================================================
% XII. FIGURES, DIAGRAMS, AND PHOTOGRAPHS
% ================================================================

\noindent
\textbf{XII.\quad FIGURES, DIAGRAMS, AND PHOTOGRAPHS}

\vspace{0.15cm}

Figures shall be clear, legible, and directly relevant to the report
content.

Figures shall be inserted using the standard \verb|figure| environment.

Example:

\begin{verbatim}
\begin{figure}[H]
    \centering
    \includegraphics[width=0.8\textwidth]
                    {figures/company_structure.pdf}
    \caption{Organizational structure of the company}
    \label{fig:company-structure}
\end{figure}
\end{verbatim}

Figure captions shall be placed below the figure.

Each referenced figure shall have a unique \verb|\label| and shall be
referenced in the text using, for example:

\begin{verbatim}
Figure~\ref{fig:company-structure}
\end{verbatim}

Students shall not manually type figure numbers.

Figures reproduced or adapted from external sources shall include an
appropriate citation.

Photographs taken during the internship shall be relevant and sufficiently
clear. Students shall not include confidential, restricted, or proprietary
information without permission from the internship organization.

Vector graphics or high-resolution images are recommended whenever
possible.

\vspace{0.3cm}

% ================================================================
% XIII. TABLES
% ================================================================

\noindent
\textbf{XIII.\quad TABLES}

\vspace{0.15cm}

Tables shall be created using suitable \LaTeX{} environments such as
\verb|table|, \verb|tabular|, \verb|tabularx|, or \verb|longtable|.

Table captions shall be placed above the table.

Example:

\begin{verbatim}
\begin{table}[H]
    \centering
    \caption{Internship schedule}
    \label{tab:internship-schedule}

    \begin{tabular}{lll}
        \toprule
        Week & Task & Result \\
        \midrule
        1 & Orientation & Completed \\
        2 & System study & Completed \\
        \bottomrule
    \end{tabular}
\end{table}
\end{verbatim}

Tables shall be referenced using:

\begin{verbatim}
Table~\ref{tab:internship-schedule}
\end{verbatim}

Students shall not manually type table numbers.

Tables should be created directly in \LaTeX{} rather than inserted as
screenshots whenever practical.

\vspace{0.3cm}

% ================================================================
% XIV. EQUATIONS AND TECHNICAL QUANTITIES
% ================================================================

\noindent
\textbf{XIV.\quad EQUATIONS AND TECHNICAL QUANTITIES}

\vspace{0.15cm}

Mathematical expressions shall be typeset directly using \LaTeX{}.
Images or screenshots of equations shall not be used.

Short expressions may be written in inline mathematical mode, for example:

\begin{verbatim}
$R = 0.7\,\Omega$
\end{verbatim}

Important equations shall use a numbered equation environment.

Example:

\begin{verbatim}
\begin{equation}
    P = UI
    \label{eq:power}
\end{equation}
\end{verbatim}

Equations shall be referenced using:

\begin{verbatim}
Equation~\eqref{eq:power}
\end{verbatim}

Equation numbers shall not be manually entered.

SI units shall be used wherever applicable. The \verb|siunitx| package
should be used where appropriate.

Examples:

\begin{verbatim}
\SI{220}{\volt}
\SI{50}{\hertz}
\SI{2.5}{\ampere}
\end{verbatim}

Notation and units shall be consistent throughout the report.

\vspace{0.3cm}

% ================================================================
% XV. SOURCE CODE
% ================================================================

\noindent
\textbf{XV.\quad SOURCE CODE}

\vspace{0.15cm}

If the internship involves programming, PLCs, microcontrollers,
simulation, or software development, relevant source code may be included
in the report.

Source code shall be presented using the \verb|listings| environment or
another environment defined by the official template.

Example:

\begin{verbatim}
\begin{lstlisting}[language=st]
IF Start = TRUE THEN
    Motor := TRUE;
ELSE
    Motor := FALSE;
END_IF;
\end{lstlisting}
\end{verbatim}

Only code necessary for explaining the work should appear in the main
report. Long source-code listings should be moved to an appendix or
submitted as supplementary files if required.

Source code should not normally be inserted as screenshots.

\vspace{0.3cm}

% ================================================================
% XVI. ABSTRACT
% ================================================================

\noindent
\textbf{XVI.\quad ABSTRACT}

\vspace{0.15cm}

The Abstract shall briefly summarize:

\begin{itemize}[leftmargin=0.8cm, itemsep=2pt]
    \item the internship organization;
    \item the objectives of the internship;
    \item the main tasks performed;
    \item the main results, knowledge, and skills obtained.
\end{itemize}

The Abstract should be concise and should not reproduce the full content
of the report chapter by chapter.

\vspace{0.3cm}

% ================================================================
% XVII. REFERENCES AND CITATIONS
% ================================================================

\noindent
\textbf{XVII.\quad REFERENCES AND CITATIONS}

\vspace{0.15cm}

All books, journal papers, conference papers, standards, websites,
datasheets, manuals, technical documents, figures, tables, or other
external sources used in preparing the report shall be properly cited.

References shall be managed through the bibliography system provided with
the official template.

Bibliographic entries shall normally be stored in:

\begin{verbatim}
References.bib
\end{verbatim}

References shall be cited in the report using:

\begin{verbatim}
\cite{reference-key}
\end{verbatim}

Students shall not manually type reference numbers such as ``[1]'',
``[2]'', etc.

Titles of foreign-language publications shall remain in their original
language in the bibliography.

All borrowed ideas, data, figures, tables, equations, methods, or results
shall be properly acknowledged.

\vspace{0.3cm}

% ================================================================
% XVIII. CROSS-REFERENCING
% ================================================================

\noindent
\textbf{XVIII.\quad CROSS-REFERENCING}

\vspace{0.15cm}

Numbered elements shall be cross-referenced using the standard \LaTeX{}
mechanism.

Use the following commands where appropriate:

\begin{verbatim}
\label{...}
\ref{...}
\eqref{...}
\end{verbatim}

For example:

\begin{verbatim}
As shown in Figure~\ref{fig:system}, ...
\end{verbatim}

Students shall not manually type numbers such as ``Figure 2.3'' when
those numbers can be generated automatically.

\vspace{0.3cm}

% ================================================================
% XIX. PAGE NUMBERING
% ================================================================

\noindent
\textbf{XIX.\quad PAGE NUMBERING}

\vspace{0.15cm}

Cover pages and official evaluation forms at the beginning of the report
shall not display page numbers.

The main body of the report shall use Arabic page numbers beginning from
Chapter 1.

Page numbering shall be generated automatically by the official template.
Students shall not manually enter page numbers.

\vspace{0.3cm}

% ================================================================
% XX. PROJECT FILE ORGANIZATION
% ================================================================

\noindent
\textbf{XX.\quad PROJECT FILE ORGANIZATION}

\vspace{0.15cm}

Students should maintain a clear and organized \LaTeX{} project
structure. A recommended structure is:

\begin{verbatim}
InternshipReport/
|
|-- main.tex
|-- References.bib
|
|-- figures/
|   |-- company.pdf
|   |-- system.png
|   `-- ...
|
`-- chapters/
    |-- chapter1.tex
    |-- chapter2.tex
    `-- chapter3.tex
\end{verbatim}

Figures should normally be stored in the \verb|figures| directory
specified by the official template.

File and directory names should preferably contain Latin letters,
numbers, hyphens, or underscores.

Recommended examples include:

\begin{verbatim}
company_structure.pdf
motor_control.png
chapter_01.tex
\end{verbatim}

Spaces and unnecessary special characters in file names should be
avoided where possible.

\vspace{0.3cm}

% ================================================================
% XXI. DOCUMENT COMPILATION AND FINAL CHECK
% ================================================================

\noindent
\textbf{XXI.\quad DOCUMENT COMPILATION AND FINAL CHECK}

\vspace{0.15cm}

Before submission, the complete \LaTeX{} project shall compile
successfully and produce the final PDF report.

The final document shall not contain unresolved references such as:

\begin{verbatim}
??
[?]
\end{verbatim}

Before generating the final PDF, students shall check:

\begin{itemize}[leftmargin=0.8cm, itemsep=2pt]
    \item the Table of Contents;
    \item page numbering;
    \item chapter, section, and subsection numbering;
    \item figure numbering;
    \item table numbering;
    \item equation numbering;
    \item cross-references;
    \item citations;
    \item bibliography;
    \item font display;
    \item mathematical symbols;
    \item image quality;
    \item completeness of all required forms.
\end{itemize}

The document may be compiled using TeX Live, MiKTeX, or another
compatible \LaTeX{} environment. Regardless of the compilation
environment, the resulting PDF shall conform to the official Faculty
template.

\vspace{0.3cm}

% ================================================================
% XXII. CONTENT THAT STUDENTS MAY EDIT
% ================================================================

\noindent
\textbf{XXII.\quad CONTENT THAT STUDENTS MAY EDIT}

\vspace{0.15cm}

Students may edit or replace:

\begin{itemize}[leftmargin=0.8cm, itemsep=2pt]
    \item student name;
    \item class;
    \item internship organization;
    \item supervisor name;
    \item Abstract;
    \item List of Symbols and Abbreviations;
    \item chapter titles and chapter content;
    \item figures and photographs;
    \item tables;
    \item equations;
    \item citations and references.
\end{itemize}

Students shall not modify the main document formatting settings in the
preamble, including the document class, page geometry, line spacing,
heading formatting, table-of-contents formatting, font configuration,
page layout, or cover-page macros unless specifically instructed.

\vspace{0.3cm}

% ================================================================
% XXIII. FINAL SUBMISSION REQUIREMENTS
% ================================================================

\noindent
\textbf{XXIII.\quad FINAL SUBMISSION REQUIREMENTS}

\vspace{0.15cm}

Before submission, the final report shall:

\begin{itemize}[leftmargin=0.8cm, itemsep=2pt]
    \item compile without errors that affect the report content or presentation;
    \item contain all required forms and report sections;
    \item contain correct student, class, internship organization, and supervisor information;
    \item contain an updated Table of Contents, List of Figures, and List of Tables;
    \item contain no unused example content from the template;
    \item contain no placeholder text such as \textit{Title of Chapter 1},
    \textit{Section 1.1}, or other sample entries;
    \item contain all required citations and references;
    \item contain no unresolved cross-references.
\end{itemize}

If the Faculty requires submission of the \LaTeX{} source, the complete
project shall include all \verb|.tex| files, the bibliography database,
figures, and any other files required to reproduce the submitted PDF.

\vspace{0.4cm}

\noindent
\textbf{Important:}
Before submitting the final internship report, students shall remove or
comment out the following command in the main document:

\begin{verbatim}
\input{Regulations.tex}
\end{verbatim}

For example:

\begin{verbatim}
% \input{Regulations.tex}
\end{verbatim}

This removes the guidance pages from the final submitted report while
preserving the regulations file in the \LaTeX{} project for reference.

} % End 13pt content

\clearpage
|.
Students shall remove or comment out this command before submitting the
final internship report.

\vspace{0.4cm}

% ================================================================
% I. GENERAL REQUIREMENTS
% ================================================================

\noindent
\textbf{I.\quad GENERAL REQUIREMENTS}

\vspace{0.15cm}

The Graduation Internship Report shall be prepared using the official
\LaTeX{} template provided by the Faculty. Students shall preserve the
predefined page layout, typography, line spacing, heading styles,
numbering system, figure and table formatting, and bibliography settings
unless otherwise instructed by the Faculty or the supervisor.

The report shall be written in English.

Students are responsible for the academic and technical content of the
report. The official \LaTeX{} template is responsible for standardized
document formatting.

\vspace{0.3cm}

% ================================================================
% II. REPORT STRUCTURE
% ================================================================

\noindent
\textbf{II.\quad REPORT STRUCTURE}

\vspace{0.15cm}

The report shall normally be arranged in the following order:

\begin{enumerate}[label=\arabic*., leftmargin=1.2cm, itemsep=2pt]
    \item Outer Cover;
    \item Comments of the Internship Host Organization;
    \item Supervisor's Evaluation and Graduation Internship Defense Results;
    \item Table of Contents;
    \item List of Symbols and Abbreviations, if applicable;
    \item List of Figures;
    \item List of Tables;
    \item Abstract;
    \item Main Chapters of the Internship Report;
    \item References;
    \item Appendices, if applicable.
\end{enumerate}

The preliminary pages shall not be numbered as chapters. The main report
content shall be organized into consecutively numbered chapters beginning
with Chapter 1.

\vspace{0.3cm}

% ================================================================
% III. PAGE SIZE AND MARGINS
% ================================================================

\noindent
\textbf{III.\quad PAGE SIZE AND MARGINS}

\vspace{0.15cm}

The report shall be formatted on A4 paper with the following margins:

\begin{center}
\begin{tabular}{ll}
    Left margin:   & 3.5 cm \\
    Right margin:  & 2.0 cm \\
    Top margin:    & 2.5 cm \\
    Bottom margin: & 2.5 cm
\end{tabular}
\end{center}

These settings are predefined in the official template and shall not be
modified by students.

\vspace{0.3cm}

% ================================================================
% IV. FONT, FONT SIZE, AND LINE SPACING
% ================================================================

\noindent
\textbf{IV.\quad FONT, FONT SIZE, AND LINE SPACING}

\vspace{0.15cm}

The main text shall use a 13 pt font size with 1.3-line spacing.
A Times-family typeface shall be used in accordance with the official
template.

The text shall be fully justified.

Students shall not use repeated spaces, repeated manual line breaks, or
multiple blank lines to position text. Paragraph and heading spacing shall
be controlled by the template.

\vspace{0.3cm}

% ================================================================
% V. OUTER COVER
% ================================================================

\noindent
\textbf{V.\quad OUTER COVER}

\vspace{0.15cm}

The outer cover shall follow the official format provided in the template
and shall contain the following information:

\begin{itemize}[leftmargin=0.8cm, itemsep=2pt]
    \item Thai Nguyen University of Technology;
    \item Faculty of International Training;
    \item Engineering and Technology Department;
    \item University logo and Faculty logo;
    \item the title \textbf{INTERNSHIP REPORT};
    \item Internship Organization;
    \item Author;
    \item Class;
    \item Supervisor;
    \item place and year of completion.
\end{itemize}

Students shall only replace the required information fields and shall not
modify the cover-page layout unless instructed otherwise.

\vspace{0.3cm}

% ================================================================
% VI. COMMENTS OF THE INTERNSHIP HOST ORGANIZATION
% ================================================================

\noindent
\textbf{VI.\quad COMMENTS OF THE INTERNSHIP HOST ORGANIZATION}

\vspace{0.15cm}

The page entitled \textbf{COMMENTS OF THE INTERNSHIP HOST ORGANIZATION}
is reserved for the internship host organization.

It includes:

\begin{itemize}[leftmargin=0.8cm, itemsep=2pt]
    \item the student's full name;
    \item class;
    \item comments from the host organization;
    \item score;
    \item date;
    \item signature and official seal.
\end{itemize}

Students shall not alter or complete official evaluation fields unless
specifically authorized.

\vspace{0.3cm}

% ================================================================
% VII. SUPERVISOR'S EVALUATION AND DEFENSE RESULTS
% ================================================================

\noindent
\textbf{VII.\quad SUPERVISOR'S EVALUATION AND DEFENSE RESULTS}

\vspace{0.15cm}

The section entitled
\textbf{SUPERVISOR'S EVALUATION AND GRADUATION INTERNSHIP DEFENSE RESULTS}
is reserved for the supervisor and defense committee.

The evaluation may include:

\begin{itemize}[leftmargin=0.8cm, itemsep=2pt]
    \item Internship Progress and Student's Attitude;
    \item Report Content;
    \item Report Presentation;
    \item Other Comments;
    \item Supervisor's Opinion;
    \item Defense Result;
    \item Score.
\end{itemize}

Students shall not enter evaluation scores or defense results before they
are officially determined.

\vspace{0.3cm}

% ================================================================
% VIII. CHAPTERS, SECTIONS, AND SUBSECTIONS
% ================================================================

\noindent
\textbf{VIII.\quad CHAPTERS, SECTIONS, AND SUBSECTIONS}

\vspace{0.15cm}

The report structure shall be generated using the standard \LaTeX{}
commands:

\begin{verbatim}
\chapter{...}
\section{...}
\subsection{...}
\end{verbatim}

Chapter, section, and subsection numbers shall be generated automatically
by \LaTeX{}.

Students shall not manually type chapter or section numbers and shall not
reproduce heading styles using commands such as \verb|\textbf|,
\verb|\fontsize|, or \verb|\centering| instead of the appropriate
structural command.

Excessive subdivision should be avoided. In most cases, the hierarchy
Chapter--Section--Subsection is sufficient.

\vspace{0.3cm}

% ================================================================
% IX. ORGANIZATION OF THE MAIN CONTENT
% ================================================================

\noindent
\textbf{IX.\quad ORGANIZATION OF THE MAIN CONTENT}

\vspace{0.15cm}

The specific content of the report shall depend on the internship
organization and the assigned tasks. A typical structure may include:

\begin{itemize}[leftmargin=0.8cm, itemsep=2pt]
    \item \textbf{Chapter 1. Overview of the Internship Organization};
    \item \textbf{Chapter 2. Internship Activities and Assigned Tasks};
    \item \textbf{Chapter 3. Results, Experience, and Evaluation}.
\end{itemize}

These chapter titles are recommendations only and may be adapted to the
actual internship content.

The report shall reflect work actually performed, observed, analyzed, or
learned during the internship. Unrelated theoretical material should not
be included merely to increase report length.

\vspace{0.3cm}

% ================================================================
% X. TABLE OF CONTENTS
% ================================================================

\noindent
\textbf{X.\quad TABLE OF CONTENTS}

\vspace{0.15cm}

The Table of Contents shall be generated automatically using:

\begin{verbatim}
\tableofcontents
\end{verbatim}

Students shall not create or manually edit page numbers in the Table of
Contents.

The document shall be compiled sufficiently to ensure that all page
numbers and entries are correctly updated.

\vspace{0.3cm}

% ================================================================
% XI. LIST OF SYMBOLS AND ABBREVIATIONS
% ================================================================

\noindent
\textbf{XI.\quad LIST OF SYMBOLS AND ABBREVIATIONS}

\vspace{0.15cm}

Symbols and abbreviations that are used repeatedly in the report shall be
listed in the \textbf{List of Symbols and Abbreviations}.

The list may contain the following columns:

\begin{center}
\begin{tabular}{llll}
    No. & Symbol / Abbreviation & Unit & Definition
\end{tabular}
\end{center}

Only symbols and abbreviations actually used in the final report shall be
retained. Example entries contained in the template shall be removed if
they are not relevant to the student's report.

Each abbreviation should normally be defined when it first appears in
the main text.

\vspace{0.3cm}

% ================================================================
% XII. FIGURES, DIAGRAMS, AND PHOTOGRAPHS
% ================================================================

\noindent
\textbf{XII.\quad FIGURES, DIAGRAMS, AND PHOTOGRAPHS}

\vspace{0.15cm}

Figures shall be clear, legible, and directly relevant to the report
content.

Figures shall be inserted using the standard \verb|figure| environment.

Example:

\begin{verbatim}
\begin{figure}[H]
    \centering
    \includegraphics[width=0.8\textwidth]
                    {figures/company_structure.pdf}
    \caption{Organizational structure of the company}
    \label{fig:company-structure}
\end{figure}
\end{verbatim}

Figure captions shall be placed below the figure.

Each referenced figure shall have a unique \verb|\label| and shall be
referenced in the text using, for example:

\begin{verbatim}
Figure~\ref{fig:company-structure}
\end{verbatim}

Students shall not manually type figure numbers.

Figures reproduced or adapted from external sources shall include an
appropriate citation.

Photographs taken during the internship shall be relevant and sufficiently
clear. Students shall not include confidential, restricted, or proprietary
information without permission from the internship organization.

Vector graphics or high-resolution images are recommended whenever
possible.

\vspace{0.3cm}

% ================================================================
% XIII. TABLES
% ================================================================

\noindent
\textbf{XIII.\quad TABLES}

\vspace{0.15cm}

Tables shall be created using suitable \LaTeX{} environments such as
\verb|table|, \verb|tabular|, \verb|tabularx|, or \verb|longtable|.

Table captions shall be placed above the table.

Example:

\begin{verbatim}
\begin{table}[H]
    \centering
    \caption{Internship schedule}
    \label{tab:internship-schedule}

    \begin{tabular}{lll}
        \toprule
        Week & Task & Result \\
        \midrule
        1 & Orientation & Completed \\
        2 & System study & Completed \\
        \bottomrule
    \end{tabular}
\end{table}
\end{verbatim}

Tables shall be referenced using:

\begin{verbatim}
Table~\ref{tab:internship-schedule}
\end{verbatim}

Students shall not manually type table numbers.

Tables should be created directly in \LaTeX{} rather than inserted as
screenshots whenever practical.

\vspace{0.3cm}

% ================================================================
% XIV. EQUATIONS AND TECHNICAL QUANTITIES
% ================================================================

\noindent
\textbf{XIV.\quad EQUATIONS AND TECHNICAL QUANTITIES}

\vspace{0.15cm}

Mathematical expressions shall be typeset directly using \LaTeX{}.
Images or screenshots of equations shall not be used.

Short expressions may be written in inline mathematical mode, for example:

\begin{verbatim}
$R = 0.7\,\Omega$
\end{verbatim}

Important equations shall use a numbered equation environment.

Example:

\begin{verbatim}
\begin{equation}
    P = UI
    \label{eq:power}
\end{equation}
\end{verbatim}

Equations shall be referenced using:

\begin{verbatim}
Equation~\eqref{eq:power}
\end{verbatim}

Equation numbers shall not be manually entered.

SI units shall be used wherever applicable. The \verb|siunitx| package
should be used where appropriate.

Examples:

\begin{verbatim}
\SI{220}{\volt}
\SI{50}{\hertz}
\SI{2.5}{\ampere}
\end{verbatim}

Notation and units shall be consistent throughout the report.

\vspace{0.3cm}

% ================================================================
% XV. SOURCE CODE
% ================================================================

\noindent
\textbf{XV.\quad SOURCE CODE}

\vspace{0.15cm}

If the internship involves programming, PLCs, microcontrollers,
simulation, or software development, relevant source code may be included
in the report.

Source code shall be presented using the \verb|listings| environment or
another environment defined by the official template.

Example:

\begin{verbatim}
\begin{lstlisting}[language=st]
IF Start = TRUE THEN
    Motor := TRUE;
ELSE
    Motor := FALSE;
END_IF;
\end{lstlisting}
\end{verbatim}

Only code necessary for explaining the work should appear in the main
report. Long source-code listings should be moved to an appendix or
submitted as supplementary files if required.

Source code should not normally be inserted as screenshots.

\vspace{0.3cm}

% ================================================================
% XVI. ABSTRACT
% ================================================================

\noindent
\textbf{XVI.\quad ABSTRACT}

\vspace{0.15cm}

The Abstract shall briefly summarize:

\begin{itemize}[leftmargin=0.8cm, itemsep=2pt]
    \item the internship organization;
    \item the objectives of the internship;
    \item the main tasks performed;
    \item the main results, knowledge, and skills obtained.
\end{itemize}

The Abstract should be concise and should not reproduce the full content
of the report chapter by chapter.

\vspace{0.3cm}

% ================================================================
% XVII. REFERENCES AND CITATIONS
% ================================================================

\noindent
\textbf{XVII.\quad REFERENCES AND CITATIONS}

\vspace{0.15cm}

All books, journal papers, conference papers, standards, websites,
datasheets, manuals, technical documents, figures, tables, or other
external sources used in preparing the report shall be properly cited.

References shall be managed through the bibliography system provided with
the official template.

Bibliographic entries shall normally be stored in:

\begin{verbatim}
References.bib
\end{verbatim}

References shall be cited in the report using:

\begin{verbatim}
\cite{reference-key}
\end{verbatim}

Students shall not manually type reference numbers such as ``[1]'',
``[2]'', etc.

Titles of foreign-language publications shall remain in their original
language in the bibliography.

All borrowed ideas, data, figures, tables, equations, methods, or results
shall be properly acknowledged.

\vspace{0.3cm}

% ================================================================
% XVIII. CROSS-REFERENCING
% ================================================================

\noindent
\textbf{XVIII.\quad CROSS-REFERENCING}

\vspace{0.15cm}

Numbered elements shall be cross-referenced using the standard \LaTeX{}
mechanism.

Use the following commands where appropriate:

\begin{verbatim}
\label{...}
\ref{...}
\eqref{...}
\end{verbatim}

For example:

\begin{verbatim}
As shown in Figure~\ref{fig:system}, ...
\end{verbatim}

Students shall not manually type numbers such as ``Figure 2.3'' when
those numbers can be generated automatically.

\vspace{0.3cm}

% ================================================================
% XIX. PAGE NUMBERING
% ================================================================

\noindent
\textbf{XIX.\quad PAGE NUMBERING}

\vspace{0.15cm}

Cover pages and official evaluation forms at the beginning of the report
shall not display page numbers.

The main body of the report shall use Arabic page numbers beginning from
Chapter 1.

Page numbering shall be generated automatically by the official template.
Students shall not manually enter page numbers.

\vspace{0.3cm}

% ================================================================
% XX. PROJECT FILE ORGANIZATION
% ================================================================

\noindent
\textbf{XX.\quad PROJECT FILE ORGANIZATION}

\vspace{0.15cm}

Students should maintain a clear and organized \LaTeX{} project
structure. A recommended structure is:

\begin{verbatim}
InternshipReport/
|
|-- main.tex
|-- References.bib
|
|-- figures/
|   |-- company.pdf
|   |-- system.png
|   `-- ...
|
`-- chapters/
    |-- chapter1.tex
    |-- chapter2.tex
    `-- chapter3.tex
\end{verbatim}

Figures should normally be stored in the \verb|figures| directory
specified by the official template.

File and directory names should preferably contain Latin letters,
numbers, hyphens, or underscores.

Recommended examples include:

\begin{verbatim}
company_structure.pdf
motor_control.png
chapter_01.tex
\end{verbatim}

Spaces and unnecessary special characters in file names should be
avoided where possible.

\vspace{0.3cm}

% ================================================================
% XXI. DOCUMENT COMPILATION AND FINAL CHECK
% ================================================================

\noindent
\textbf{XXI.\quad DOCUMENT COMPILATION AND FINAL CHECK}

\vspace{0.15cm}

Before submission, the complete \LaTeX{} project shall compile
successfully and produce the final PDF report.

The final document shall not contain unresolved references such as:

\begin{verbatim}
??
[?]
\end{verbatim}

Before generating the final PDF, students shall check:

\begin{itemize}[leftmargin=0.8cm, itemsep=2pt]
    \item the Table of Contents;
    \item page numbering;
    \item chapter, section, and subsection numbering;
    \item figure numbering;
    \item table numbering;
    \item equation numbering;
    \item cross-references;
    \item citations;
    \item bibliography;
    \item font display;
    \item mathematical symbols;
    \item image quality;
    \item completeness of all required forms.
\end{itemize}

The document may be compiled using TeX Live, MiKTeX, or another
compatible \LaTeX{} environment. Regardless of the compilation
environment, the resulting PDF shall conform to the official Faculty
template.

\vspace{0.3cm}

% ================================================================
% XXII. CONTENT THAT STUDENTS MAY EDIT
% ================================================================

\noindent
\textbf{XXII.\quad CONTENT THAT STUDENTS MAY EDIT}

\vspace{0.15cm}

Students may edit or replace:

\begin{itemize}[leftmargin=0.8cm, itemsep=2pt]
    \item student name;
    \item class;
    \item internship organization;
    \item supervisor name;
    \item Abstract;
    \item List of Symbols and Abbreviations;
    \item chapter titles and chapter content;
    \item figures and photographs;
    \item tables;
    \item equations;
    \item citations and references.
\end{itemize}

Students shall not modify the main document formatting settings in the
preamble, including the document class, page geometry, line spacing,
heading formatting, table-of-contents formatting, font configuration,
page layout, or cover-page macros unless specifically instructed.

\vspace{0.3cm}

% ================================================================
% XXIII. FINAL SUBMISSION REQUIREMENTS
% ================================================================

\noindent
\textbf{XXIII.\quad FINAL SUBMISSION REQUIREMENTS}

\vspace{0.15cm}

Before submission, the final report shall:

\begin{itemize}[leftmargin=0.8cm, itemsep=2pt]
    \item compile without errors that affect the report content or presentation;
    \item contain all required forms and report sections;
    \item contain correct student, class, internship organization, and supervisor information;
    \item contain an updated Table of Contents, List of Figures, and List of Tables;
    \item contain no unused example content from the template;
    \item contain no placeholder text such as \textit{Title of Chapter 1},
    \textit{Section 1.1}, or other sample entries;
    \item contain all required citations and references;
    \item contain no unresolved cross-references.
\end{itemize}

If the Faculty requires submission of the \LaTeX{} source, the complete
project shall include all \verb|.tex| files, the bibliography database,
figures, and any other files required to reproduce the submitted PDF.

\vspace{0.4cm}

\noindent
\textbf{Important:}
Before submitting the final internship report, students shall remove or
comment out the following command in the main document:

\begin{verbatim}

% =================================================================
% REGULATIONS FOR PREPARING THE GRADUATION INTERNSHIP REPORT
% USING LATEX
% =================================================================
%
% This file is intended to be included from the main document with:
%
%     \input{Regulations.tex}
%
% IMPORTANT:
% This section is provided for guidance only. Students shall remove
% or comment out the \input{Regulations.tex} command before submitting
% the final internship report.
%
% =================================================================

\clearpage
\thispagestyle{empty}

\begin{center}
    {\fontsize{16pt}{20pt}\selectfont\bfseries
    SPECIFIC REGULATIONS FOR PREPARING\\
    GRADUATION INTERNSHIP REPORTS USING \LaTeX}
\end{center}

\vspace{0.5cm}

{\fontsize{13pt}{17pt}\selectfont
\setstretch{1.3}
\justifying

\noindent
\textbf{Note:}
This section is provided for guidance only and is included in the main
document using the command \verb|\input{Regulations.tex}|.
Students shall remove or comment out this command before submitting the
final internship report.

\vspace{0.4cm}

% ================================================================
% I. GENERAL REQUIREMENTS
% ================================================================

\noindent
\textbf{I.\quad GENERAL REQUIREMENTS}

\vspace{0.15cm}

The Graduation Internship Report shall be prepared using the official
\LaTeX{} template provided by the Faculty. Students shall preserve the
predefined page layout, typography, line spacing, heading styles,
numbering system, figure and table formatting, and bibliography settings
unless otherwise instructed by the Faculty or the supervisor.

The report shall be written in English.

Students are responsible for the academic and technical content of the
report. The official \LaTeX{} template is responsible for standardized
document formatting.

\vspace{0.3cm}

% ================================================================
% II. REPORT STRUCTURE
% ================================================================

\noindent
\textbf{II.\quad REPORT STRUCTURE}

\vspace{0.15cm}

The report shall normally be arranged in the following order:

\begin{enumerate}[label=\arabic*., leftmargin=1.2cm, itemsep=2pt]
    \item Outer Cover;
    \item Comments of the Internship Host Organization;
    \item Supervisor's Evaluation and Graduation Internship Defense Results;
    \item Table of Contents;
    \item List of Symbols and Abbreviations, if applicable;
    \item List of Figures;
    \item List of Tables;
    \item Abstract;
    \item Main Chapters of the Internship Report;
    \item References;
    \item Appendices, if applicable.
\end{enumerate}

The preliminary pages shall not be numbered as chapters. The main report
content shall be organized into consecutively numbered chapters beginning
with Chapter 1.

\vspace{0.3cm}

% ================================================================
% III. PAGE SIZE AND MARGINS
% ================================================================

\noindent
\textbf{III.\quad PAGE SIZE AND MARGINS}

\vspace{0.15cm}

The report shall be formatted on A4 paper with the following margins:

\begin{center}
\begin{tabular}{ll}
    Left margin:   & 3.5 cm \\
    Right margin:  & 2.0 cm \\
    Top margin:    & 2.5 cm \\
    Bottom margin: & 2.5 cm
\end{tabular}
\end{center}

These settings are predefined in the official template and shall not be
modified by students.

\vspace{0.3cm}

% ================================================================
% IV. FONT, FONT SIZE, AND LINE SPACING
% ================================================================

\noindent
\textbf{IV.\quad FONT, FONT SIZE, AND LINE SPACING}

\vspace{0.15cm}

The main text shall use a 13 pt font size with 1.3-line spacing.
A Times-family typeface shall be used in accordance with the official
template.

The text shall be fully justified.

Students shall not use repeated spaces, repeated manual line breaks, or
multiple blank lines to position text. Paragraph and heading spacing shall
be controlled by the template.

\vspace{0.3cm}

% ================================================================
% V. OUTER COVER
% ================================================================

\noindent
\textbf{V.\quad OUTER COVER}

\vspace{0.15cm}

The outer cover shall follow the official format provided in the template
and shall contain the following information:

\begin{itemize}[leftmargin=0.8cm, itemsep=2pt]
    \item Thai Nguyen University of Technology;
    \item Faculty of International Training;
    \item Engineering and Technology Department;
    \item University logo and Faculty logo;
    \item the title \textbf{INTERNSHIP REPORT};
    \item Internship Organization;
    \item Author;
    \item Class;
    \item Supervisor;
    \item place and year of completion.
\end{itemize}

Students shall only replace the required information fields and shall not
modify the cover-page layout unless instructed otherwise.

\vspace{0.3cm}

% ================================================================
% VI. COMMENTS OF THE INTERNSHIP HOST ORGANIZATION
% ================================================================

\noindent
\textbf{VI.\quad COMMENTS OF THE INTERNSHIP HOST ORGANIZATION}

\vspace{0.15cm}

The page entitled \textbf{COMMENTS OF THE INTERNSHIP HOST ORGANIZATION}
is reserved for the internship host organization.

It includes:

\begin{itemize}[leftmargin=0.8cm, itemsep=2pt]
    \item the student's full name;
    \item class;
    \item comments from the host organization;
    \item score;
    \item date;
    \item signature and official seal.
\end{itemize}

Students shall not alter or complete official evaluation fields unless
specifically authorized.

\vspace{0.3cm}

% ================================================================
% VII. SUPERVISOR'S EVALUATION AND DEFENSE RESULTS
% ================================================================

\noindent
\textbf{VII.\quad SUPERVISOR'S EVALUATION AND DEFENSE RESULTS}

\vspace{0.15cm}

The section entitled
\textbf{SUPERVISOR'S EVALUATION AND GRADUATION INTERNSHIP DEFENSE RESULTS}
is reserved for the supervisor and defense committee.

The evaluation may include:

\begin{itemize}[leftmargin=0.8cm, itemsep=2pt]
    \item Internship Progress and Student's Attitude;
    \item Report Content;
    \item Report Presentation;
    \item Other Comments;
    \item Supervisor's Opinion;
    \item Defense Result;
    \item Score.
\end{itemize}

Students shall not enter evaluation scores or defense results before they
are officially determined.

\vspace{0.3cm}

% ================================================================
% VIII. CHAPTERS, SECTIONS, AND SUBSECTIONS
% ================================================================

\noindent
\textbf{VIII.\quad CHAPTERS, SECTIONS, AND SUBSECTIONS}

\vspace{0.15cm}

The report structure shall be generated using the standard \LaTeX{}
commands:

\begin{verbatim}
\chapter{...}
\section{...}
\subsection{...}
\end{verbatim}

Chapter, section, and subsection numbers shall be generated automatically
by \LaTeX{}.

Students shall not manually type chapter or section numbers and shall not
reproduce heading styles using commands such as \verb|\textbf|,
\verb|\fontsize|, or \verb|\centering| instead of the appropriate
structural command.

Excessive subdivision should be avoided. In most cases, the hierarchy
Chapter--Section--Subsection is sufficient.

\vspace{0.3cm}

% ================================================================
% IX. ORGANIZATION OF THE MAIN CONTENT
% ================================================================

\noindent
\textbf{IX.\quad ORGANIZATION OF THE MAIN CONTENT}

\vspace{0.15cm}

The specific content of the report shall depend on the internship
organization and the assigned tasks. A typical structure may include:

\begin{itemize}[leftmargin=0.8cm, itemsep=2pt]
    \item \textbf{Chapter 1. Overview of the Internship Organization};
    \item \textbf{Chapter 2. Internship Activities and Assigned Tasks};
    \item \textbf{Chapter 3. Results, Experience, and Evaluation}.
\end{itemize}

These chapter titles are recommendations only and may be adapted to the
actual internship content.

The report shall reflect work actually performed, observed, analyzed, or
learned during the internship. Unrelated theoretical material should not
be included merely to increase report length.

\vspace{0.3cm}

% ================================================================
% X. TABLE OF CONTENTS
% ================================================================

\noindent
\textbf{X.\quad TABLE OF CONTENTS}

\vspace{0.15cm}

The Table of Contents shall be generated automatically using:

\begin{verbatim}
\tableofcontents
\end{verbatim}

Students shall not create or manually edit page numbers in the Table of
Contents.

The document shall be compiled sufficiently to ensure that all page
numbers and entries are correctly updated.

\vspace{0.3cm}

% ================================================================
% XI. LIST OF SYMBOLS AND ABBREVIATIONS
% ================================================================

\noindent
\textbf{XI.\quad LIST OF SYMBOLS AND ABBREVIATIONS}

\vspace{0.15cm}

Symbols and abbreviations that are used repeatedly in the report shall be
listed in the \textbf{List of Symbols and Abbreviations}.

The list may contain the following columns:

\begin{center}
\begin{tabular}{llll}
    No. & Symbol / Abbreviation & Unit & Definition
\end{tabular}
\end{center}

Only symbols and abbreviations actually used in the final report shall be
retained. Example entries contained in the template shall be removed if
they are not relevant to the student's report.

Each abbreviation should normally be defined when it first appears in
the main text.

\vspace{0.3cm}

% ================================================================
% XII. FIGURES, DIAGRAMS, AND PHOTOGRAPHS
% ================================================================

\noindent
\textbf{XII.\quad FIGURES, DIAGRAMS, AND PHOTOGRAPHS}

\vspace{0.15cm}

Figures shall be clear, legible, and directly relevant to the report
content.

Figures shall be inserted using the standard \verb|figure| environment.

Example:

\begin{verbatim}
\begin{figure}[H]
    \centering
    \includegraphics[width=0.8\textwidth]
                    {figures/company_structure.pdf}
    \caption{Organizational structure of the company}
    \label{fig:company-structure}
\end{figure}
\end{verbatim}

Figure captions shall be placed below the figure.

Each referenced figure shall have a unique \verb|\label| and shall be
referenced in the text using, for example:

\begin{verbatim}
Figure~\ref{fig:company-structure}
\end{verbatim}

Students shall not manually type figure numbers.

Figures reproduced or adapted from external sources shall include an
appropriate citation.

Photographs taken during the internship shall be relevant and sufficiently
clear. Students shall not include confidential, restricted, or proprietary
information without permission from the internship organization.

Vector graphics or high-resolution images are recommended whenever
possible.

\vspace{0.3cm}

% ================================================================
% XIII. TABLES
% ================================================================

\noindent
\textbf{XIII.\quad TABLES}

\vspace{0.15cm}

Tables shall be created using suitable \LaTeX{} environments such as
\verb|table|, \verb|tabular|, \verb|tabularx|, or \verb|longtable|.

Table captions shall be placed above the table.

Example:

\begin{verbatim}
\begin{table}[H]
    \centering
    \caption{Internship schedule}
    \label{tab:internship-schedule}

    \begin{tabular}{lll}
        \toprule
        Week & Task & Result \\
        \midrule
        1 & Orientation & Completed \\
        2 & System study & Completed \\
        \bottomrule
    \end{tabular}
\end{table}
\end{verbatim}

Tables shall be referenced using:

\begin{verbatim}
Table~\ref{tab:internship-schedule}
\end{verbatim}

Students shall not manually type table numbers.

Tables should be created directly in \LaTeX{} rather than inserted as
screenshots whenever practical.

\vspace{0.3cm}

% ================================================================
% XIV. EQUATIONS AND TECHNICAL QUANTITIES
% ================================================================

\noindent
\textbf{XIV.\quad EQUATIONS AND TECHNICAL QUANTITIES}

\vspace{0.15cm}

Mathematical expressions shall be typeset directly using \LaTeX{}.
Images or screenshots of equations shall not be used.

Short expressions may be written in inline mathematical mode, for example:

\begin{verbatim}
$R = 0.7\,\Omega$
\end{verbatim}

Important equations shall use a numbered equation environment.

Example:

\begin{verbatim}
\begin{equation}
    P = UI
    \label{eq:power}
\end{equation}
\end{verbatim}

Equations shall be referenced using:

\begin{verbatim}
Equation~\eqref{eq:power}
\end{verbatim}

Equation numbers shall not be manually entered.

SI units shall be used wherever applicable. The \verb|siunitx| package
should be used where appropriate.

Examples:

\begin{verbatim}
\SI{220}{\volt}
\SI{50}{\hertz}
\SI{2.5}{\ampere}
\end{verbatim}

Notation and units shall be consistent throughout the report.

\vspace{0.3cm}

% ================================================================
% XV. SOURCE CODE
% ================================================================

\noindent
\textbf{XV.\quad SOURCE CODE}

\vspace{0.15cm}

If the internship involves programming, PLCs, microcontrollers,
simulation, or software development, relevant source code may be included
in the report.

Source code shall be presented using the \verb|listings| environment or
another environment defined by the official template.

Example:

\begin{verbatim}
\begin{lstlisting}[language=st]
IF Start = TRUE THEN
    Motor := TRUE;
ELSE
    Motor := FALSE;
END_IF;
\end{lstlisting}
\end{verbatim}

Only code necessary for explaining the work should appear in the main
report. Long source-code listings should be moved to an appendix or
submitted as supplementary files if required.

Source code should not normally be inserted as screenshots.

\vspace{0.3cm}

% ================================================================
% XVI. ABSTRACT
% ================================================================

\noindent
\textbf{XVI.\quad ABSTRACT}

\vspace{0.15cm}

The Abstract shall briefly summarize:

\begin{itemize}[leftmargin=0.8cm, itemsep=2pt]
    \item the internship organization;
    \item the objectives of the internship;
    \item the main tasks performed;
    \item the main results, knowledge, and skills obtained.
\end{itemize}

The Abstract should be concise and should not reproduce the full content
of the report chapter by chapter.

\vspace{0.3cm}

% ================================================================
% XVII. REFERENCES AND CITATIONS
% ================================================================

\noindent
\textbf{XVII.\quad REFERENCES AND CITATIONS}

\vspace{0.15cm}

All books, journal papers, conference papers, standards, websites,
datasheets, manuals, technical documents, figures, tables, or other
external sources used in preparing the report shall be properly cited.

References shall be managed through the bibliography system provided with
the official template.

Bibliographic entries shall normally be stored in:

\begin{verbatim}
References.bib
\end{verbatim}

References shall be cited in the report using:

\begin{verbatim}
\cite{reference-key}
\end{verbatim}

Students shall not manually type reference numbers such as ``[1]'',
``[2]'', etc.

Titles of foreign-language publications shall remain in their original
language in the bibliography.

All borrowed ideas, data, figures, tables, equations, methods, or results
shall be properly acknowledged.

\vspace{0.3cm}

% ================================================================
% XVIII. CROSS-REFERENCING
% ================================================================

\noindent
\textbf{XVIII.\quad CROSS-REFERENCING}

\vspace{0.15cm}

Numbered elements shall be cross-referenced using the standard \LaTeX{}
mechanism.

Use the following commands where appropriate:

\begin{verbatim}
\label{...}
\ref{...}
\eqref{...}
\end{verbatim}

For example:

\begin{verbatim}
As shown in Figure~\ref{fig:system}, ...
\end{verbatim}

Students shall not manually type numbers such as ``Figure 2.3'' when
those numbers can be generated automatically.

\vspace{0.3cm}

% ================================================================
% XIX. PAGE NUMBERING
% ================================================================

\noindent
\textbf{XIX.\quad PAGE NUMBERING}

\vspace{0.15cm}

Cover pages and official evaluation forms at the beginning of the report
shall not display page numbers.

The main body of the report shall use Arabic page numbers beginning from
Chapter 1.

Page numbering shall be generated automatically by the official template.
Students shall not manually enter page numbers.

\vspace{0.3cm}

% ================================================================
% XX. PROJECT FILE ORGANIZATION
% ================================================================

\noindent
\textbf{XX.\quad PROJECT FILE ORGANIZATION}

\vspace{0.15cm}

Students should maintain a clear and organized \LaTeX{} project
structure. A recommended structure is:

\begin{verbatim}
InternshipReport/
|
|-- main.tex
|-- References.bib
|
|-- figures/
|   |-- company.pdf
|   |-- system.png
|   `-- ...
|
`-- chapters/
    |-- chapter1.tex
    |-- chapter2.tex
    `-- chapter3.tex
\end{verbatim}

Figures should normally be stored in the \verb|figures| directory
specified by the official template.

File and directory names should preferably contain Latin letters,
numbers, hyphens, or underscores.

Recommended examples include:

\begin{verbatim}
company_structure.pdf
motor_control.png
chapter_01.tex
\end{verbatim}

Spaces and unnecessary special characters in file names should be
avoided where possible.

\vspace{0.3cm}

% ================================================================
% XXI. DOCUMENT COMPILATION AND FINAL CHECK
% ================================================================

\noindent
\textbf{XXI.\quad DOCUMENT COMPILATION AND FINAL CHECK}

\vspace{0.15cm}

Before submission, the complete \LaTeX{} project shall compile
successfully and produce the final PDF report.

The final document shall not contain unresolved references such as:

\begin{verbatim}
??
[?]
\end{verbatim}

Before generating the final PDF, students shall check:

\begin{itemize}[leftmargin=0.8cm, itemsep=2pt]
    \item the Table of Contents;
    \item page numbering;
    \item chapter, section, and subsection numbering;
    \item figure numbering;
    \item table numbering;
    \item equation numbering;
    \item cross-references;
    \item citations;
    \item bibliography;
    \item font display;
    \item mathematical symbols;
    \item image quality;
    \item completeness of all required forms.
\end{itemize}

The document may be compiled using TeX Live, MiKTeX, or another
compatible \LaTeX{} environment. Regardless of the compilation
environment, the resulting PDF shall conform to the official Faculty
template.

\vspace{0.3cm}

% ================================================================
% XXII. CONTENT THAT STUDENTS MAY EDIT
% ================================================================

\noindent
\textbf{XXII.\quad CONTENT THAT STUDENTS MAY EDIT}

\vspace{0.15cm}

Students may edit or replace:

\begin{itemize}[leftmargin=0.8cm, itemsep=2pt]
    \item student name;
    \item class;
    \item internship organization;
    \item supervisor name;
    \item Abstract;
    \item List of Symbols and Abbreviations;
    \item chapter titles and chapter content;
    \item figures and photographs;
    \item tables;
    \item equations;
    \item citations and references.
\end{itemize}

Students shall not modify the main document formatting settings in the
preamble, including the document class, page geometry, line spacing,
heading formatting, table-of-contents formatting, font configuration,
page layout, or cover-page macros unless specifically instructed.

\vspace{0.3cm}

% ================================================================
% XXIII. FINAL SUBMISSION REQUIREMENTS
% ================================================================

\noindent
\textbf{XXIII.\quad FINAL SUBMISSION REQUIREMENTS}

\vspace{0.15cm}

Before submission, the final report shall:

\begin{itemize}[leftmargin=0.8cm, itemsep=2pt]
    \item compile without errors that affect the report content or presentation;
    \item contain all required forms and report sections;
    \item contain correct student, class, internship organization, and supervisor information;
    \item contain an updated Table of Contents, List of Figures, and List of Tables;
    \item contain no unused example content from the template;
    \item contain no placeholder text such as \textit{Title of Chapter 1},
    \textit{Section 1.1}, or other sample entries;
    \item contain all required citations and references;
    \item contain no unresolved cross-references.
\end{itemize}

If the Faculty requires submission of the \LaTeX{} source, the complete
project shall include all \verb|.tex| files, the bibliography database,
figures, and any other files required to reproduce the submitted PDF.

\vspace{0.4cm}

\noindent
\textbf{Important:}
Before submitting the final internship report, students shall remove or
comment out the following command in the main document:

\begin{verbatim}
\input{Regulations.tex}
\end{verbatim}

For example:

\begin{verbatim}
% \input{Regulations.tex}
\end{verbatim}

This removes the guidance pages from the final submitted report while
preserving the regulations file in the \LaTeX{} project for reference.

} % End 13pt content

\clearpage

\end{verbatim}

For example:

\begin{verbatim}
% 
% =================================================================
% REGULATIONS FOR PREPARING THE GRADUATION INTERNSHIP REPORT
% USING LATEX
% =================================================================
%
% This file is intended to be included from the main document with:
%
%     \input{Regulations.tex}
%
% IMPORTANT:
% This section is provided for guidance only. Students shall remove
% or comment out the \input{Regulations.tex} command before submitting
% the final internship report.
%
% =================================================================

\clearpage
\thispagestyle{empty}

\begin{center}
    {\fontsize{16pt}{20pt}\selectfont\bfseries
    SPECIFIC REGULATIONS FOR PREPARING\\
    GRADUATION INTERNSHIP REPORTS USING \LaTeX}
\end{center}

\vspace{0.5cm}

{\fontsize{13pt}{17pt}\selectfont
\setstretch{1.3}
\justifying

\noindent
\textbf{Note:}
This section is provided for guidance only and is included in the main
document using the command \verb|\input{Regulations.tex}|.
Students shall remove or comment out this command before submitting the
final internship report.

\vspace{0.4cm}

% ================================================================
% I. GENERAL REQUIREMENTS
% ================================================================

\noindent
\textbf{I.\quad GENERAL REQUIREMENTS}

\vspace{0.15cm}

The Graduation Internship Report shall be prepared using the official
\LaTeX{} template provided by the Faculty. Students shall preserve the
predefined page layout, typography, line spacing, heading styles,
numbering system, figure and table formatting, and bibliography settings
unless otherwise instructed by the Faculty or the supervisor.

The report shall be written in English.

Students are responsible for the academic and technical content of the
report. The official \LaTeX{} template is responsible for standardized
document formatting.

\vspace{0.3cm}

% ================================================================
% II. REPORT STRUCTURE
% ================================================================

\noindent
\textbf{II.\quad REPORT STRUCTURE}

\vspace{0.15cm}

The report shall normally be arranged in the following order:

\begin{enumerate}[label=\arabic*., leftmargin=1.2cm, itemsep=2pt]
    \item Outer Cover;
    \item Comments of the Internship Host Organization;
    \item Supervisor's Evaluation and Graduation Internship Defense Results;
    \item Table of Contents;
    \item List of Symbols and Abbreviations, if applicable;
    \item List of Figures;
    \item List of Tables;
    \item Abstract;
    \item Main Chapters of the Internship Report;
    \item References;
    \item Appendices, if applicable.
\end{enumerate}

The preliminary pages shall not be numbered as chapters. The main report
content shall be organized into consecutively numbered chapters beginning
with Chapter 1.

\vspace{0.3cm}

% ================================================================
% III. PAGE SIZE AND MARGINS
% ================================================================

\noindent
\textbf{III.\quad PAGE SIZE AND MARGINS}

\vspace{0.15cm}

The report shall be formatted on A4 paper with the following margins:

\begin{center}
\begin{tabular}{ll}
    Left margin:   & 3.5 cm \\
    Right margin:  & 2.0 cm \\
    Top margin:    & 2.5 cm \\
    Bottom margin: & 2.5 cm
\end{tabular}
\end{center}

These settings are predefined in the official template and shall not be
modified by students.

\vspace{0.3cm}

% ================================================================
% IV. FONT, FONT SIZE, AND LINE SPACING
% ================================================================

\noindent
\textbf{IV.\quad FONT, FONT SIZE, AND LINE SPACING}

\vspace{0.15cm}

The main text shall use a 13 pt font size with 1.3-line spacing.
A Times-family typeface shall be used in accordance with the official
template.

The text shall be fully justified.

Students shall not use repeated spaces, repeated manual line breaks, or
multiple blank lines to position text. Paragraph and heading spacing shall
be controlled by the template.

\vspace{0.3cm}

% ================================================================
% V. OUTER COVER
% ================================================================

\noindent
\textbf{V.\quad OUTER COVER}

\vspace{0.15cm}

The outer cover shall follow the official format provided in the template
and shall contain the following information:

\begin{itemize}[leftmargin=0.8cm, itemsep=2pt]
    \item Thai Nguyen University of Technology;
    \item Faculty of International Training;
    \item Engineering and Technology Department;
    \item University logo and Faculty logo;
    \item the title \textbf{INTERNSHIP REPORT};
    \item Internship Organization;
    \item Author;
    \item Class;
    \item Supervisor;
    \item place and year of completion.
\end{itemize}

Students shall only replace the required information fields and shall not
modify the cover-page layout unless instructed otherwise.

\vspace{0.3cm}

% ================================================================
% VI. COMMENTS OF THE INTERNSHIP HOST ORGANIZATION
% ================================================================

\noindent
\textbf{VI.\quad COMMENTS OF THE INTERNSHIP HOST ORGANIZATION}

\vspace{0.15cm}

The page entitled \textbf{COMMENTS OF THE INTERNSHIP HOST ORGANIZATION}
is reserved for the internship host organization.

It includes:

\begin{itemize}[leftmargin=0.8cm, itemsep=2pt]
    \item the student's full name;
    \item class;
    \item comments from the host organization;
    \item score;
    \item date;
    \item signature and official seal.
\end{itemize}

Students shall not alter or complete official evaluation fields unless
specifically authorized.

\vspace{0.3cm}

% ================================================================
% VII. SUPERVISOR'S EVALUATION AND DEFENSE RESULTS
% ================================================================

\noindent
\textbf{VII.\quad SUPERVISOR'S EVALUATION AND DEFENSE RESULTS}

\vspace{0.15cm}

The section entitled
\textbf{SUPERVISOR'S EVALUATION AND GRADUATION INTERNSHIP DEFENSE RESULTS}
is reserved for the supervisor and defense committee.

The evaluation may include:

\begin{itemize}[leftmargin=0.8cm, itemsep=2pt]
    \item Internship Progress and Student's Attitude;
    \item Report Content;
    \item Report Presentation;
    \item Other Comments;
    \item Supervisor's Opinion;
    \item Defense Result;
    \item Score.
\end{itemize}

Students shall not enter evaluation scores or defense results before they
are officially determined.

\vspace{0.3cm}

% ================================================================
% VIII. CHAPTERS, SECTIONS, AND SUBSECTIONS
% ================================================================

\noindent
\textbf{VIII.\quad CHAPTERS, SECTIONS, AND SUBSECTIONS}

\vspace{0.15cm}

The report structure shall be generated using the standard \LaTeX{}
commands:

\begin{verbatim}
\chapter{...}
\section{...}
\subsection{...}
\end{verbatim}

Chapter, section, and subsection numbers shall be generated automatically
by \LaTeX{}.

Students shall not manually type chapter or section numbers and shall not
reproduce heading styles using commands such as \verb|\textbf|,
\verb|\fontsize|, or \verb|\centering| instead of the appropriate
structural command.

Excessive subdivision should be avoided. In most cases, the hierarchy
Chapter--Section--Subsection is sufficient.

\vspace{0.3cm}

% ================================================================
% IX. ORGANIZATION OF THE MAIN CONTENT
% ================================================================

\noindent
\textbf{IX.\quad ORGANIZATION OF THE MAIN CONTENT}

\vspace{0.15cm}

The specific content of the report shall depend on the internship
organization and the assigned tasks. A typical structure may include:

\begin{itemize}[leftmargin=0.8cm, itemsep=2pt]
    \item \textbf{Chapter 1. Overview of the Internship Organization};
    \item \textbf{Chapter 2. Internship Activities and Assigned Tasks};
    \item \textbf{Chapter 3. Results, Experience, and Evaluation}.
\end{itemize}

These chapter titles are recommendations only and may be adapted to the
actual internship content.

The report shall reflect work actually performed, observed, analyzed, or
learned during the internship. Unrelated theoretical material should not
be included merely to increase report length.

\vspace{0.3cm}

% ================================================================
% X. TABLE OF CONTENTS
% ================================================================

\noindent
\textbf{X.\quad TABLE OF CONTENTS}

\vspace{0.15cm}

The Table of Contents shall be generated automatically using:

\begin{verbatim}
\tableofcontents
\end{verbatim}

Students shall not create or manually edit page numbers in the Table of
Contents.

The document shall be compiled sufficiently to ensure that all page
numbers and entries are correctly updated.

\vspace{0.3cm}

% ================================================================
% XI. LIST OF SYMBOLS AND ABBREVIATIONS
% ================================================================

\noindent
\textbf{XI.\quad LIST OF SYMBOLS AND ABBREVIATIONS}

\vspace{0.15cm}

Symbols and abbreviations that are used repeatedly in the report shall be
listed in the \textbf{List of Symbols and Abbreviations}.

The list may contain the following columns:

\begin{center}
\begin{tabular}{llll}
    No. & Symbol / Abbreviation & Unit & Definition
\end{tabular}
\end{center}

Only symbols and abbreviations actually used in the final report shall be
retained. Example entries contained in the template shall be removed if
they are not relevant to the student's report.

Each abbreviation should normally be defined when it first appears in
the main text.

\vspace{0.3cm}

% ================================================================
% XII. FIGURES, DIAGRAMS, AND PHOTOGRAPHS
% ================================================================

\noindent
\textbf{XII.\quad FIGURES, DIAGRAMS, AND PHOTOGRAPHS}

\vspace{0.15cm}

Figures shall be clear, legible, and directly relevant to the report
content.

Figures shall be inserted using the standard \verb|figure| environment.

Example:

\begin{verbatim}
\begin{figure}[H]
    \centering
    \includegraphics[width=0.8\textwidth]
                    {figures/company_structure.pdf}
    \caption{Organizational structure of the company}
    \label{fig:company-structure}
\end{figure}
\end{verbatim}

Figure captions shall be placed below the figure.

Each referenced figure shall have a unique \verb|\label| and shall be
referenced in the text using, for example:

\begin{verbatim}
Figure~\ref{fig:company-structure}
\end{verbatim}

Students shall not manually type figure numbers.

Figures reproduced or adapted from external sources shall include an
appropriate citation.

Photographs taken during the internship shall be relevant and sufficiently
clear. Students shall not include confidential, restricted, or proprietary
information without permission from the internship organization.

Vector graphics or high-resolution images are recommended whenever
possible.

\vspace{0.3cm}

% ================================================================
% XIII. TABLES
% ================================================================

\noindent
\textbf{XIII.\quad TABLES}

\vspace{0.15cm}

Tables shall be created using suitable \LaTeX{} environments such as
\verb|table|, \verb|tabular|, \verb|tabularx|, or \verb|longtable|.

Table captions shall be placed above the table.

Example:

\begin{verbatim}
\begin{table}[H]
    \centering
    \caption{Internship schedule}
    \label{tab:internship-schedule}

    \begin{tabular}{lll}
        \toprule
        Week & Task & Result \\
        \midrule
        1 & Orientation & Completed \\
        2 & System study & Completed \\
        \bottomrule
    \end{tabular}
\end{table}
\end{verbatim}

Tables shall be referenced using:

\begin{verbatim}
Table~\ref{tab:internship-schedule}
\end{verbatim}

Students shall not manually type table numbers.

Tables should be created directly in \LaTeX{} rather than inserted as
screenshots whenever practical.

\vspace{0.3cm}

% ================================================================
% XIV. EQUATIONS AND TECHNICAL QUANTITIES
% ================================================================

\noindent
\textbf{XIV.\quad EQUATIONS AND TECHNICAL QUANTITIES}

\vspace{0.15cm}

Mathematical expressions shall be typeset directly using \LaTeX{}.
Images or screenshots of equations shall not be used.

Short expressions may be written in inline mathematical mode, for example:

\begin{verbatim}
$R = 0.7\,\Omega$
\end{verbatim}

Important equations shall use a numbered equation environment.

Example:

\begin{verbatim}
\begin{equation}
    P = UI
    \label{eq:power}
\end{equation}
\end{verbatim}

Equations shall be referenced using:

\begin{verbatim}
Equation~\eqref{eq:power}
\end{verbatim}

Equation numbers shall not be manually entered.

SI units shall be used wherever applicable. The \verb|siunitx| package
should be used where appropriate.

Examples:

\begin{verbatim}
\SI{220}{\volt}
\SI{50}{\hertz}
\SI{2.5}{\ampere}
\end{verbatim}

Notation and units shall be consistent throughout the report.

\vspace{0.3cm}

% ================================================================
% XV. SOURCE CODE
% ================================================================

\noindent
\textbf{XV.\quad SOURCE CODE}

\vspace{0.15cm}

If the internship involves programming, PLCs, microcontrollers,
simulation, or software development, relevant source code may be included
in the report.

Source code shall be presented using the \verb|listings| environment or
another environment defined by the official template.

Example:

\begin{verbatim}
\begin{lstlisting}[language=st]
IF Start = TRUE THEN
    Motor := TRUE;
ELSE
    Motor := FALSE;
END_IF;
\end{lstlisting}
\end{verbatim}

Only code necessary for explaining the work should appear in the main
report. Long source-code listings should be moved to an appendix or
submitted as supplementary files if required.

Source code should not normally be inserted as screenshots.

\vspace{0.3cm}

% ================================================================
% XVI. ABSTRACT
% ================================================================

\noindent
\textbf{XVI.\quad ABSTRACT}

\vspace{0.15cm}

The Abstract shall briefly summarize:

\begin{itemize}[leftmargin=0.8cm, itemsep=2pt]
    \item the internship organization;
    \item the objectives of the internship;
    \item the main tasks performed;
    \item the main results, knowledge, and skills obtained.
\end{itemize}

The Abstract should be concise and should not reproduce the full content
of the report chapter by chapter.

\vspace{0.3cm}

% ================================================================
% XVII. REFERENCES AND CITATIONS
% ================================================================

\noindent
\textbf{XVII.\quad REFERENCES AND CITATIONS}

\vspace{0.15cm}

All books, journal papers, conference papers, standards, websites,
datasheets, manuals, technical documents, figures, tables, or other
external sources used in preparing the report shall be properly cited.

References shall be managed through the bibliography system provided with
the official template.

Bibliographic entries shall normally be stored in:

\begin{verbatim}
References.bib
\end{verbatim}

References shall be cited in the report using:

\begin{verbatim}
\cite{reference-key}
\end{verbatim}

Students shall not manually type reference numbers such as ``[1]'',
``[2]'', etc.

Titles of foreign-language publications shall remain in their original
language in the bibliography.

All borrowed ideas, data, figures, tables, equations, methods, or results
shall be properly acknowledged.

\vspace{0.3cm}

% ================================================================
% XVIII. CROSS-REFERENCING
% ================================================================

\noindent
\textbf{XVIII.\quad CROSS-REFERENCING}

\vspace{0.15cm}

Numbered elements shall be cross-referenced using the standard \LaTeX{}
mechanism.

Use the following commands where appropriate:

\begin{verbatim}
\label{...}
\ref{...}
\eqref{...}
\end{verbatim}

For example:

\begin{verbatim}
As shown in Figure~\ref{fig:system}, ...
\end{verbatim}

Students shall not manually type numbers such as ``Figure 2.3'' when
those numbers can be generated automatically.

\vspace{0.3cm}

% ================================================================
% XIX. PAGE NUMBERING
% ================================================================

\noindent
\textbf{XIX.\quad PAGE NUMBERING}

\vspace{0.15cm}

Cover pages and official evaluation forms at the beginning of the report
shall not display page numbers.

The main body of the report shall use Arabic page numbers beginning from
Chapter 1.

Page numbering shall be generated automatically by the official template.
Students shall not manually enter page numbers.

\vspace{0.3cm}

% ================================================================
% XX. PROJECT FILE ORGANIZATION
% ================================================================

\noindent
\textbf{XX.\quad PROJECT FILE ORGANIZATION}

\vspace{0.15cm}

Students should maintain a clear and organized \LaTeX{} project
structure. A recommended structure is:

\begin{verbatim}
InternshipReport/
|
|-- main.tex
|-- References.bib
|
|-- figures/
|   |-- company.pdf
|   |-- system.png
|   `-- ...
|
`-- chapters/
    |-- chapter1.tex
    |-- chapter2.tex
    `-- chapter3.tex
\end{verbatim}

Figures should normally be stored in the \verb|figures| directory
specified by the official template.

File and directory names should preferably contain Latin letters,
numbers, hyphens, or underscores.

Recommended examples include:

\begin{verbatim}
company_structure.pdf
motor_control.png
chapter_01.tex
\end{verbatim}

Spaces and unnecessary special characters in file names should be
avoided where possible.

\vspace{0.3cm}

% ================================================================
% XXI. DOCUMENT COMPILATION AND FINAL CHECK
% ================================================================

\noindent
\textbf{XXI.\quad DOCUMENT COMPILATION AND FINAL CHECK}

\vspace{0.15cm}

Before submission, the complete \LaTeX{} project shall compile
successfully and produce the final PDF report.

The final document shall not contain unresolved references such as:

\begin{verbatim}
??
[?]
\end{verbatim}

Before generating the final PDF, students shall check:

\begin{itemize}[leftmargin=0.8cm, itemsep=2pt]
    \item the Table of Contents;
    \item page numbering;
    \item chapter, section, and subsection numbering;
    \item figure numbering;
    \item table numbering;
    \item equation numbering;
    \item cross-references;
    \item citations;
    \item bibliography;
    \item font display;
    \item mathematical symbols;
    \item image quality;
    \item completeness of all required forms.
\end{itemize}

The document may be compiled using TeX Live, MiKTeX, or another
compatible \LaTeX{} environment. Regardless of the compilation
environment, the resulting PDF shall conform to the official Faculty
template.

\vspace{0.3cm}

% ================================================================
% XXII. CONTENT THAT STUDENTS MAY EDIT
% ================================================================

\noindent
\textbf{XXII.\quad CONTENT THAT STUDENTS MAY EDIT}

\vspace{0.15cm}

Students may edit or replace:

\begin{itemize}[leftmargin=0.8cm, itemsep=2pt]
    \item student name;
    \item class;
    \item internship organization;
    \item supervisor name;
    \item Abstract;
    \item List of Symbols and Abbreviations;
    \item chapter titles and chapter content;
    \item figures and photographs;
    \item tables;
    \item equations;
    \item citations and references.
\end{itemize}

Students shall not modify the main document formatting settings in the
preamble, including the document class, page geometry, line spacing,
heading formatting, table-of-contents formatting, font configuration,
page layout, or cover-page macros unless specifically instructed.

\vspace{0.3cm}

% ================================================================
% XXIII. FINAL SUBMISSION REQUIREMENTS
% ================================================================

\noindent
\textbf{XXIII.\quad FINAL SUBMISSION REQUIREMENTS}

\vspace{0.15cm}

Before submission, the final report shall:

\begin{itemize}[leftmargin=0.8cm, itemsep=2pt]
    \item compile without errors that affect the report content or presentation;
    \item contain all required forms and report sections;
    \item contain correct student, class, internship organization, and supervisor information;
    \item contain an updated Table of Contents, List of Figures, and List of Tables;
    \item contain no unused example content from the template;
    \item contain no placeholder text such as \textit{Title of Chapter 1},
    \textit{Section 1.1}, or other sample entries;
    \item contain all required citations and references;
    \item contain no unresolved cross-references.
\end{itemize}

If the Faculty requires submission of the \LaTeX{} source, the complete
project shall include all \verb|.tex| files, the bibliography database,
figures, and any other files required to reproduce the submitted PDF.

\vspace{0.4cm}

\noindent
\textbf{Important:}
Before submitting the final internship report, students shall remove or
comment out the following command in the main document:

\begin{verbatim}
\input{Regulations.tex}
\end{verbatim}

For example:

\begin{verbatim}
% \input{Regulations.tex}
\end{verbatim}

This removes the guidance pages from the final submitted report while
preserving the regulations file in the \LaTeX{} project for reference.

} % End 13pt content

\clearpage

\end{verbatim}

This removes the guidance pages from the final submitted report while
preserving the regulations file in the \LaTeX{} project for reference.

} % End 13pt content

\clearpage

\end{verbatim}

For example:

\begin{verbatim}
% 
% =================================================================
% REGULATIONS FOR PREPARING THE GRADUATION INTERNSHIP REPORT
% USING LATEX
% =================================================================
%
% This file is intended to be included from the main document with:
%
%     
% =================================================================
% REGULATIONS FOR PREPARING THE GRADUATION INTERNSHIP REPORT
% USING LATEX
% =================================================================
%
% This file is intended to be included from the main document with:
%
%     \input{Regulations.tex}
%
% IMPORTANT:
% This section is provided for guidance only. Students shall remove
% or comment out the \input{Regulations.tex} command before submitting
% the final internship report.
%
% =================================================================

\clearpage
\thispagestyle{empty}

\begin{center}
    {\fontsize{16pt}{20pt}\selectfont\bfseries
    SPECIFIC REGULATIONS FOR PREPARING\\
    GRADUATION INTERNSHIP REPORTS USING \LaTeX}
\end{center}

\vspace{0.5cm}

{\fontsize{13pt}{17pt}\selectfont
\setstretch{1.3}
\justifying

\noindent
\textbf{Note:}
This section is provided for guidance only and is included in the main
document using the command \verb|\input{Regulations.tex}|.
Students shall remove or comment out this command before submitting the
final internship report.

\vspace{0.4cm}

% ================================================================
% I. GENERAL REQUIREMENTS
% ================================================================

\noindent
\textbf{I.\quad GENERAL REQUIREMENTS}

\vspace{0.15cm}

The Graduation Internship Report shall be prepared using the official
\LaTeX{} template provided by the Faculty. Students shall preserve the
predefined page layout, typography, line spacing, heading styles,
numbering system, figure and table formatting, and bibliography settings
unless otherwise instructed by the Faculty or the supervisor.

The report shall be written in English.

Students are responsible for the academic and technical content of the
report. The official \LaTeX{} template is responsible for standardized
document formatting.

\vspace{0.3cm}

% ================================================================
% II. REPORT STRUCTURE
% ================================================================

\noindent
\textbf{II.\quad REPORT STRUCTURE}

\vspace{0.15cm}

The report shall normally be arranged in the following order:

\begin{enumerate}[label=\arabic*., leftmargin=1.2cm, itemsep=2pt]
    \item Outer Cover;
    \item Comments of the Internship Host Organization;
    \item Supervisor's Evaluation and Graduation Internship Defense Results;
    \item Table of Contents;
    \item List of Symbols and Abbreviations, if applicable;
    \item List of Figures;
    \item List of Tables;
    \item Abstract;
    \item Main Chapters of the Internship Report;
    \item References;
    \item Appendices, if applicable.
\end{enumerate}

The preliminary pages shall not be numbered as chapters. The main report
content shall be organized into consecutively numbered chapters beginning
with Chapter 1.

\vspace{0.3cm}

% ================================================================
% III. PAGE SIZE AND MARGINS
% ================================================================

\noindent
\textbf{III.\quad PAGE SIZE AND MARGINS}

\vspace{0.15cm}

The report shall be formatted on A4 paper with the following margins:

\begin{center}
\begin{tabular}{ll}
    Left margin:   & 3.5 cm \\
    Right margin:  & 2.0 cm \\
    Top margin:    & 2.5 cm \\
    Bottom margin: & 2.5 cm
\end{tabular}
\end{center}

These settings are predefined in the official template and shall not be
modified by students.

\vspace{0.3cm}

% ================================================================
% IV. FONT, FONT SIZE, AND LINE SPACING
% ================================================================

\noindent
\textbf{IV.\quad FONT, FONT SIZE, AND LINE SPACING}

\vspace{0.15cm}

The main text shall use a 13 pt font size with 1.3-line spacing.
A Times-family typeface shall be used in accordance with the official
template.

The text shall be fully justified.

Students shall not use repeated spaces, repeated manual line breaks, or
multiple blank lines to position text. Paragraph and heading spacing shall
be controlled by the template.

\vspace{0.3cm}

% ================================================================
% V. OUTER COVER
% ================================================================

\noindent
\textbf{V.\quad OUTER COVER}

\vspace{0.15cm}

The outer cover shall follow the official format provided in the template
and shall contain the following information:

\begin{itemize}[leftmargin=0.8cm, itemsep=2pt]
    \item Thai Nguyen University of Technology;
    \item Faculty of International Training;
    \item Engineering and Technology Department;
    \item University logo and Faculty logo;
    \item the title \textbf{INTERNSHIP REPORT};
    \item Internship Organization;
    \item Author;
    \item Class;
    \item Supervisor;
    \item place and year of completion.
\end{itemize}

Students shall only replace the required information fields and shall not
modify the cover-page layout unless instructed otherwise.

\vspace{0.3cm}

% ================================================================
% VI. COMMENTS OF THE INTERNSHIP HOST ORGANIZATION
% ================================================================

\noindent
\textbf{VI.\quad COMMENTS OF THE INTERNSHIP HOST ORGANIZATION}

\vspace{0.15cm}

The page entitled \textbf{COMMENTS OF THE INTERNSHIP HOST ORGANIZATION}
is reserved for the internship host organization.

It includes:

\begin{itemize}[leftmargin=0.8cm, itemsep=2pt]
    \item the student's full name;
    \item class;
    \item comments from the host organization;
    \item score;
    \item date;
    \item signature and official seal.
\end{itemize}

Students shall not alter or complete official evaluation fields unless
specifically authorized.

\vspace{0.3cm}

% ================================================================
% VII. SUPERVISOR'S EVALUATION AND DEFENSE RESULTS
% ================================================================

\noindent
\textbf{VII.\quad SUPERVISOR'S EVALUATION AND DEFENSE RESULTS}

\vspace{0.15cm}

The section entitled
\textbf{SUPERVISOR'S EVALUATION AND GRADUATION INTERNSHIP DEFENSE RESULTS}
is reserved for the supervisor and defense committee.

The evaluation may include:

\begin{itemize}[leftmargin=0.8cm, itemsep=2pt]
    \item Internship Progress and Student's Attitude;
    \item Report Content;
    \item Report Presentation;
    \item Other Comments;
    \item Supervisor's Opinion;
    \item Defense Result;
    \item Score.
\end{itemize}

Students shall not enter evaluation scores or defense results before they
are officially determined.

\vspace{0.3cm}

% ================================================================
% VIII. CHAPTERS, SECTIONS, AND SUBSECTIONS
% ================================================================

\noindent
\textbf{VIII.\quad CHAPTERS, SECTIONS, AND SUBSECTIONS}

\vspace{0.15cm}

The report structure shall be generated using the standard \LaTeX{}
commands:

\begin{verbatim}
\chapter{...}
\section{...}
\subsection{...}
\end{verbatim}

Chapter, section, and subsection numbers shall be generated automatically
by \LaTeX{}.

Students shall not manually type chapter or section numbers and shall not
reproduce heading styles using commands such as \verb|\textbf|,
\verb|\fontsize|, or \verb|\centering| instead of the appropriate
structural command.

Excessive subdivision should be avoided. In most cases, the hierarchy
Chapter--Section--Subsection is sufficient.

\vspace{0.3cm}

% ================================================================
% IX. ORGANIZATION OF THE MAIN CONTENT
% ================================================================

\noindent
\textbf{IX.\quad ORGANIZATION OF THE MAIN CONTENT}

\vspace{0.15cm}

The specific content of the report shall depend on the internship
organization and the assigned tasks. A typical structure may include:

\begin{itemize}[leftmargin=0.8cm, itemsep=2pt]
    \item \textbf{Chapter 1. Overview of the Internship Organization};
    \item \textbf{Chapter 2. Internship Activities and Assigned Tasks};
    \item \textbf{Chapter 3. Results, Experience, and Evaluation}.
\end{itemize}

These chapter titles are recommendations only and may be adapted to the
actual internship content.

The report shall reflect work actually performed, observed, analyzed, or
learned during the internship. Unrelated theoretical material should not
be included merely to increase report length.

\vspace{0.3cm}

% ================================================================
% X. TABLE OF CONTENTS
% ================================================================

\noindent
\textbf{X.\quad TABLE OF CONTENTS}

\vspace{0.15cm}

The Table of Contents shall be generated automatically using:

\begin{verbatim}
\tableofcontents
\end{verbatim}

Students shall not create or manually edit page numbers in the Table of
Contents.

The document shall be compiled sufficiently to ensure that all page
numbers and entries are correctly updated.

\vspace{0.3cm}

% ================================================================
% XI. LIST OF SYMBOLS AND ABBREVIATIONS
% ================================================================

\noindent
\textbf{XI.\quad LIST OF SYMBOLS AND ABBREVIATIONS}

\vspace{0.15cm}

Symbols and abbreviations that are used repeatedly in the report shall be
listed in the \textbf{List of Symbols and Abbreviations}.

The list may contain the following columns:

\begin{center}
\begin{tabular}{llll}
    No. & Symbol / Abbreviation & Unit & Definition
\end{tabular}
\end{center}

Only symbols and abbreviations actually used in the final report shall be
retained. Example entries contained in the template shall be removed if
they are not relevant to the student's report.

Each abbreviation should normally be defined when it first appears in
the main text.

\vspace{0.3cm}

% ================================================================
% XII. FIGURES, DIAGRAMS, AND PHOTOGRAPHS
% ================================================================

\noindent
\textbf{XII.\quad FIGURES, DIAGRAMS, AND PHOTOGRAPHS}

\vspace{0.15cm}

Figures shall be clear, legible, and directly relevant to the report
content.

Figures shall be inserted using the standard \verb|figure| environment.

Example:

\begin{verbatim}
\begin{figure}[H]
    \centering
    \includegraphics[width=0.8\textwidth]
                    {figures/company_structure.pdf}
    \caption{Organizational structure of the company}
    \label{fig:company-structure}
\end{figure}
\end{verbatim}

Figure captions shall be placed below the figure.

Each referenced figure shall have a unique \verb|\label| and shall be
referenced in the text using, for example:

\begin{verbatim}
Figure~\ref{fig:company-structure}
\end{verbatim}

Students shall not manually type figure numbers.

Figures reproduced or adapted from external sources shall include an
appropriate citation.

Photographs taken during the internship shall be relevant and sufficiently
clear. Students shall not include confidential, restricted, or proprietary
information without permission from the internship organization.

Vector graphics or high-resolution images are recommended whenever
possible.

\vspace{0.3cm}

% ================================================================
% XIII. TABLES
% ================================================================

\noindent
\textbf{XIII.\quad TABLES}

\vspace{0.15cm}

Tables shall be created using suitable \LaTeX{} environments such as
\verb|table|, \verb|tabular|, \verb|tabularx|, or \verb|longtable|.

Table captions shall be placed above the table.

Example:

\begin{verbatim}
\begin{table}[H]
    \centering
    \caption{Internship schedule}
    \label{tab:internship-schedule}

    \begin{tabular}{lll}
        \toprule
        Week & Task & Result \\
        \midrule
        1 & Orientation & Completed \\
        2 & System study & Completed \\
        \bottomrule
    \end{tabular}
\end{table}
\end{verbatim}

Tables shall be referenced using:

\begin{verbatim}
Table~\ref{tab:internship-schedule}
\end{verbatim}

Students shall not manually type table numbers.

Tables should be created directly in \LaTeX{} rather than inserted as
screenshots whenever practical.

\vspace{0.3cm}

% ================================================================
% XIV. EQUATIONS AND TECHNICAL QUANTITIES
% ================================================================

\noindent
\textbf{XIV.\quad EQUATIONS AND TECHNICAL QUANTITIES}

\vspace{0.15cm}

Mathematical expressions shall be typeset directly using \LaTeX{}.
Images or screenshots of equations shall not be used.

Short expressions may be written in inline mathematical mode, for example:

\begin{verbatim}
$R = 0.7\,\Omega$
\end{verbatim}

Important equations shall use a numbered equation environment.

Example:

\begin{verbatim}
\begin{equation}
    P = UI
    \label{eq:power}
\end{equation}
\end{verbatim}

Equations shall be referenced using:

\begin{verbatim}
Equation~\eqref{eq:power}
\end{verbatim}

Equation numbers shall not be manually entered.

SI units shall be used wherever applicable. The \verb|siunitx| package
should be used where appropriate.

Examples:

\begin{verbatim}
\SI{220}{\volt}
\SI{50}{\hertz}
\SI{2.5}{\ampere}
\end{verbatim}

Notation and units shall be consistent throughout the report.

\vspace{0.3cm}

% ================================================================
% XV. SOURCE CODE
% ================================================================

\noindent
\textbf{XV.\quad SOURCE CODE}

\vspace{0.15cm}

If the internship involves programming, PLCs, microcontrollers,
simulation, or software development, relevant source code may be included
in the report.

Source code shall be presented using the \verb|listings| environment or
another environment defined by the official template.

Example:

\begin{verbatim}
\begin{lstlisting}[language=st]
IF Start = TRUE THEN
    Motor := TRUE;
ELSE
    Motor := FALSE;
END_IF;
\end{lstlisting}
\end{verbatim}

Only code necessary for explaining the work should appear in the main
report. Long source-code listings should be moved to an appendix or
submitted as supplementary files if required.

Source code should not normally be inserted as screenshots.

\vspace{0.3cm}

% ================================================================
% XVI. ABSTRACT
% ================================================================

\noindent
\textbf{XVI.\quad ABSTRACT}

\vspace{0.15cm}

The Abstract shall briefly summarize:

\begin{itemize}[leftmargin=0.8cm, itemsep=2pt]
    \item the internship organization;
    \item the objectives of the internship;
    \item the main tasks performed;
    \item the main results, knowledge, and skills obtained.
\end{itemize}

The Abstract should be concise and should not reproduce the full content
of the report chapter by chapter.

\vspace{0.3cm}

% ================================================================
% XVII. REFERENCES AND CITATIONS
% ================================================================

\noindent
\textbf{XVII.\quad REFERENCES AND CITATIONS}

\vspace{0.15cm}

All books, journal papers, conference papers, standards, websites,
datasheets, manuals, technical documents, figures, tables, or other
external sources used in preparing the report shall be properly cited.

References shall be managed through the bibliography system provided with
the official template.

Bibliographic entries shall normally be stored in:

\begin{verbatim}
References.bib
\end{verbatim}

References shall be cited in the report using:

\begin{verbatim}
\cite{reference-key}
\end{verbatim}

Students shall not manually type reference numbers such as ``[1]'',
``[2]'', etc.

Titles of foreign-language publications shall remain in their original
language in the bibliography.

All borrowed ideas, data, figures, tables, equations, methods, or results
shall be properly acknowledged.

\vspace{0.3cm}

% ================================================================
% XVIII. CROSS-REFERENCING
% ================================================================

\noindent
\textbf{XVIII.\quad CROSS-REFERENCING}

\vspace{0.15cm}

Numbered elements shall be cross-referenced using the standard \LaTeX{}
mechanism.

Use the following commands where appropriate:

\begin{verbatim}
\label{...}
\ref{...}
\eqref{...}
\end{verbatim}

For example:

\begin{verbatim}
As shown in Figure~\ref{fig:system}, ...
\end{verbatim}

Students shall not manually type numbers such as ``Figure 2.3'' when
those numbers can be generated automatically.

\vspace{0.3cm}

% ================================================================
% XIX. PAGE NUMBERING
% ================================================================

\noindent
\textbf{XIX.\quad PAGE NUMBERING}

\vspace{0.15cm}

Cover pages and official evaluation forms at the beginning of the report
shall not display page numbers.

The main body of the report shall use Arabic page numbers beginning from
Chapter 1.

Page numbering shall be generated automatically by the official template.
Students shall not manually enter page numbers.

\vspace{0.3cm}

% ================================================================
% XX. PROJECT FILE ORGANIZATION
% ================================================================

\noindent
\textbf{XX.\quad PROJECT FILE ORGANIZATION}

\vspace{0.15cm}

Students should maintain a clear and organized \LaTeX{} project
structure. A recommended structure is:

\begin{verbatim}
InternshipReport/
|
|-- main.tex
|-- References.bib
|
|-- figures/
|   |-- company.pdf
|   |-- system.png
|   `-- ...
|
`-- chapters/
    |-- chapter1.tex
    |-- chapter2.tex
    `-- chapter3.tex
\end{verbatim}

Figures should normally be stored in the \verb|figures| directory
specified by the official template.

File and directory names should preferably contain Latin letters,
numbers, hyphens, or underscores.

Recommended examples include:

\begin{verbatim}
company_structure.pdf
motor_control.png
chapter_01.tex
\end{verbatim}

Spaces and unnecessary special characters in file names should be
avoided where possible.

\vspace{0.3cm}

% ================================================================
% XXI. DOCUMENT COMPILATION AND FINAL CHECK
% ================================================================

\noindent
\textbf{XXI.\quad DOCUMENT COMPILATION AND FINAL CHECK}

\vspace{0.15cm}

Before submission, the complete \LaTeX{} project shall compile
successfully and produce the final PDF report.

The final document shall not contain unresolved references such as:

\begin{verbatim}
??
[?]
\end{verbatim}

Before generating the final PDF, students shall check:

\begin{itemize}[leftmargin=0.8cm, itemsep=2pt]
    \item the Table of Contents;
    \item page numbering;
    \item chapter, section, and subsection numbering;
    \item figure numbering;
    \item table numbering;
    \item equation numbering;
    \item cross-references;
    \item citations;
    \item bibliography;
    \item font display;
    \item mathematical symbols;
    \item image quality;
    \item completeness of all required forms.
\end{itemize}

The document may be compiled using TeX Live, MiKTeX, or another
compatible \LaTeX{} environment. Regardless of the compilation
environment, the resulting PDF shall conform to the official Faculty
template.

\vspace{0.3cm}

% ================================================================
% XXII. CONTENT THAT STUDENTS MAY EDIT
% ================================================================

\noindent
\textbf{XXII.\quad CONTENT THAT STUDENTS MAY EDIT}

\vspace{0.15cm}

Students may edit or replace:

\begin{itemize}[leftmargin=0.8cm, itemsep=2pt]
    \item student name;
    \item class;
    \item internship organization;
    \item supervisor name;
    \item Abstract;
    \item List of Symbols and Abbreviations;
    \item chapter titles and chapter content;
    \item figures and photographs;
    \item tables;
    \item equations;
    \item citations and references.
\end{itemize}

Students shall not modify the main document formatting settings in the
preamble, including the document class, page geometry, line spacing,
heading formatting, table-of-contents formatting, font configuration,
page layout, or cover-page macros unless specifically instructed.

\vspace{0.3cm}

% ================================================================
% XXIII. FINAL SUBMISSION REQUIREMENTS
% ================================================================

\noindent
\textbf{XXIII.\quad FINAL SUBMISSION REQUIREMENTS}

\vspace{0.15cm}

Before submission, the final report shall:

\begin{itemize}[leftmargin=0.8cm, itemsep=2pt]
    \item compile without errors that affect the report content or presentation;
    \item contain all required forms and report sections;
    \item contain correct student, class, internship organization, and supervisor information;
    \item contain an updated Table of Contents, List of Figures, and List of Tables;
    \item contain no unused example content from the template;
    \item contain no placeholder text such as \textit{Title of Chapter 1},
    \textit{Section 1.1}, or other sample entries;
    \item contain all required citations and references;
    \item contain no unresolved cross-references.
\end{itemize}

If the Faculty requires submission of the \LaTeX{} source, the complete
project shall include all \verb|.tex| files, the bibliography database,
figures, and any other files required to reproduce the submitted PDF.

\vspace{0.4cm}

\noindent
\textbf{Important:}
Before submitting the final internship report, students shall remove or
comment out the following command in the main document:

\begin{verbatim}
\input{Regulations.tex}
\end{verbatim}

For example:

\begin{verbatim}
% \input{Regulations.tex}
\end{verbatim}

This removes the guidance pages from the final submitted report while
preserving the regulations file in the \LaTeX{} project for reference.

} % End 13pt content

\clearpage

%
% IMPORTANT:
% This section is provided for guidance only. Students shall remove
% or comment out the 
% =================================================================
% REGULATIONS FOR PREPARING THE GRADUATION INTERNSHIP REPORT
% USING LATEX
% =================================================================
%
% This file is intended to be included from the main document with:
%
%     \input{Regulations.tex}
%
% IMPORTANT:
% This section is provided for guidance only. Students shall remove
% or comment out the \input{Regulations.tex} command before submitting
% the final internship report.
%
% =================================================================

\clearpage
\thispagestyle{empty}

\begin{center}
    {\fontsize{16pt}{20pt}\selectfont\bfseries
    SPECIFIC REGULATIONS FOR PREPARING\\
    GRADUATION INTERNSHIP REPORTS USING \LaTeX}
\end{center}

\vspace{0.5cm}

{\fontsize{13pt}{17pt}\selectfont
\setstretch{1.3}
\justifying

\noindent
\textbf{Note:}
This section is provided for guidance only and is included in the main
document using the command \verb|\input{Regulations.tex}|.
Students shall remove or comment out this command before submitting the
final internship report.

\vspace{0.4cm}

% ================================================================
% I. GENERAL REQUIREMENTS
% ================================================================

\noindent
\textbf{I.\quad GENERAL REQUIREMENTS}

\vspace{0.15cm}

The Graduation Internship Report shall be prepared using the official
\LaTeX{} template provided by the Faculty. Students shall preserve the
predefined page layout, typography, line spacing, heading styles,
numbering system, figure and table formatting, and bibliography settings
unless otherwise instructed by the Faculty or the supervisor.

The report shall be written in English.

Students are responsible for the academic and technical content of the
report. The official \LaTeX{} template is responsible for standardized
document formatting.

\vspace{0.3cm}

% ================================================================
% II. REPORT STRUCTURE
% ================================================================

\noindent
\textbf{II.\quad REPORT STRUCTURE}

\vspace{0.15cm}

The report shall normally be arranged in the following order:

\begin{enumerate}[label=\arabic*., leftmargin=1.2cm, itemsep=2pt]
    \item Outer Cover;
    \item Comments of the Internship Host Organization;
    \item Supervisor's Evaluation and Graduation Internship Defense Results;
    \item Table of Contents;
    \item List of Symbols and Abbreviations, if applicable;
    \item List of Figures;
    \item List of Tables;
    \item Abstract;
    \item Main Chapters of the Internship Report;
    \item References;
    \item Appendices, if applicable.
\end{enumerate}

The preliminary pages shall not be numbered as chapters. The main report
content shall be organized into consecutively numbered chapters beginning
with Chapter 1.

\vspace{0.3cm}

% ================================================================
% III. PAGE SIZE AND MARGINS
% ================================================================

\noindent
\textbf{III.\quad PAGE SIZE AND MARGINS}

\vspace{0.15cm}

The report shall be formatted on A4 paper with the following margins:

\begin{center}
\begin{tabular}{ll}
    Left margin:   & 3.5 cm \\
    Right margin:  & 2.0 cm \\
    Top margin:    & 2.5 cm \\
    Bottom margin: & 2.5 cm
\end{tabular}
\end{center}

These settings are predefined in the official template and shall not be
modified by students.

\vspace{0.3cm}

% ================================================================
% IV. FONT, FONT SIZE, AND LINE SPACING
% ================================================================

\noindent
\textbf{IV.\quad FONT, FONT SIZE, AND LINE SPACING}

\vspace{0.15cm}

The main text shall use a 13 pt font size with 1.3-line spacing.
A Times-family typeface shall be used in accordance with the official
template.

The text shall be fully justified.

Students shall not use repeated spaces, repeated manual line breaks, or
multiple blank lines to position text. Paragraph and heading spacing shall
be controlled by the template.

\vspace{0.3cm}

% ================================================================
% V. OUTER COVER
% ================================================================

\noindent
\textbf{V.\quad OUTER COVER}

\vspace{0.15cm}

The outer cover shall follow the official format provided in the template
and shall contain the following information:

\begin{itemize}[leftmargin=0.8cm, itemsep=2pt]
    \item Thai Nguyen University of Technology;
    \item Faculty of International Training;
    \item Engineering and Technology Department;
    \item University logo and Faculty logo;
    \item the title \textbf{INTERNSHIP REPORT};
    \item Internship Organization;
    \item Author;
    \item Class;
    \item Supervisor;
    \item place and year of completion.
\end{itemize}

Students shall only replace the required information fields and shall not
modify the cover-page layout unless instructed otherwise.

\vspace{0.3cm}

% ================================================================
% VI. COMMENTS OF THE INTERNSHIP HOST ORGANIZATION
% ================================================================

\noindent
\textbf{VI.\quad COMMENTS OF THE INTERNSHIP HOST ORGANIZATION}

\vspace{0.15cm}

The page entitled \textbf{COMMENTS OF THE INTERNSHIP HOST ORGANIZATION}
is reserved for the internship host organization.

It includes:

\begin{itemize}[leftmargin=0.8cm, itemsep=2pt]
    \item the student's full name;
    \item class;
    \item comments from the host organization;
    \item score;
    \item date;
    \item signature and official seal.
\end{itemize}

Students shall not alter or complete official evaluation fields unless
specifically authorized.

\vspace{0.3cm}

% ================================================================
% VII. SUPERVISOR'S EVALUATION AND DEFENSE RESULTS
% ================================================================

\noindent
\textbf{VII.\quad SUPERVISOR'S EVALUATION AND DEFENSE RESULTS}

\vspace{0.15cm}

The section entitled
\textbf{SUPERVISOR'S EVALUATION AND GRADUATION INTERNSHIP DEFENSE RESULTS}
is reserved for the supervisor and defense committee.

The evaluation may include:

\begin{itemize}[leftmargin=0.8cm, itemsep=2pt]
    \item Internship Progress and Student's Attitude;
    \item Report Content;
    \item Report Presentation;
    \item Other Comments;
    \item Supervisor's Opinion;
    \item Defense Result;
    \item Score.
\end{itemize}

Students shall not enter evaluation scores or defense results before they
are officially determined.

\vspace{0.3cm}

% ================================================================
% VIII. CHAPTERS, SECTIONS, AND SUBSECTIONS
% ================================================================

\noindent
\textbf{VIII.\quad CHAPTERS, SECTIONS, AND SUBSECTIONS}

\vspace{0.15cm}

The report structure shall be generated using the standard \LaTeX{}
commands:

\begin{verbatim}
\chapter{...}
\section{...}
\subsection{...}
\end{verbatim}

Chapter, section, and subsection numbers shall be generated automatically
by \LaTeX{}.

Students shall not manually type chapter or section numbers and shall not
reproduce heading styles using commands such as \verb|\textbf|,
\verb|\fontsize|, or \verb|\centering| instead of the appropriate
structural command.

Excessive subdivision should be avoided. In most cases, the hierarchy
Chapter--Section--Subsection is sufficient.

\vspace{0.3cm}

% ================================================================
% IX. ORGANIZATION OF THE MAIN CONTENT
% ================================================================

\noindent
\textbf{IX.\quad ORGANIZATION OF THE MAIN CONTENT}

\vspace{0.15cm}

The specific content of the report shall depend on the internship
organization and the assigned tasks. A typical structure may include:

\begin{itemize}[leftmargin=0.8cm, itemsep=2pt]
    \item \textbf{Chapter 1. Overview of the Internship Organization};
    \item \textbf{Chapter 2. Internship Activities and Assigned Tasks};
    \item \textbf{Chapter 3. Results, Experience, and Evaluation}.
\end{itemize}

These chapter titles are recommendations only and may be adapted to the
actual internship content.

The report shall reflect work actually performed, observed, analyzed, or
learned during the internship. Unrelated theoretical material should not
be included merely to increase report length.

\vspace{0.3cm}

% ================================================================
% X. TABLE OF CONTENTS
% ================================================================

\noindent
\textbf{X.\quad TABLE OF CONTENTS}

\vspace{0.15cm}

The Table of Contents shall be generated automatically using:

\begin{verbatim}
\tableofcontents
\end{verbatim}

Students shall not create or manually edit page numbers in the Table of
Contents.

The document shall be compiled sufficiently to ensure that all page
numbers and entries are correctly updated.

\vspace{0.3cm}

% ================================================================
% XI. LIST OF SYMBOLS AND ABBREVIATIONS
% ================================================================

\noindent
\textbf{XI.\quad LIST OF SYMBOLS AND ABBREVIATIONS}

\vspace{0.15cm}

Symbols and abbreviations that are used repeatedly in the report shall be
listed in the \textbf{List of Symbols and Abbreviations}.

The list may contain the following columns:

\begin{center}
\begin{tabular}{llll}
    No. & Symbol / Abbreviation & Unit & Definition
\end{tabular}
\end{center}

Only symbols and abbreviations actually used in the final report shall be
retained. Example entries contained in the template shall be removed if
they are not relevant to the student's report.

Each abbreviation should normally be defined when it first appears in
the main text.

\vspace{0.3cm}

% ================================================================
% XII. FIGURES, DIAGRAMS, AND PHOTOGRAPHS
% ================================================================

\noindent
\textbf{XII.\quad FIGURES, DIAGRAMS, AND PHOTOGRAPHS}

\vspace{0.15cm}

Figures shall be clear, legible, and directly relevant to the report
content.

Figures shall be inserted using the standard \verb|figure| environment.

Example:

\begin{verbatim}
\begin{figure}[H]
    \centering
    \includegraphics[width=0.8\textwidth]
                    {figures/company_structure.pdf}
    \caption{Organizational structure of the company}
    \label{fig:company-structure}
\end{figure}
\end{verbatim}

Figure captions shall be placed below the figure.

Each referenced figure shall have a unique \verb|\label| and shall be
referenced in the text using, for example:

\begin{verbatim}
Figure~\ref{fig:company-structure}
\end{verbatim}

Students shall not manually type figure numbers.

Figures reproduced or adapted from external sources shall include an
appropriate citation.

Photographs taken during the internship shall be relevant and sufficiently
clear. Students shall not include confidential, restricted, or proprietary
information without permission from the internship organization.

Vector graphics or high-resolution images are recommended whenever
possible.

\vspace{0.3cm}

% ================================================================
% XIII. TABLES
% ================================================================

\noindent
\textbf{XIII.\quad TABLES}

\vspace{0.15cm}

Tables shall be created using suitable \LaTeX{} environments such as
\verb|table|, \verb|tabular|, \verb|tabularx|, or \verb|longtable|.

Table captions shall be placed above the table.

Example:

\begin{verbatim}
\begin{table}[H]
    \centering
    \caption{Internship schedule}
    \label{tab:internship-schedule}

    \begin{tabular}{lll}
        \toprule
        Week & Task & Result \\
        \midrule
        1 & Orientation & Completed \\
        2 & System study & Completed \\
        \bottomrule
    \end{tabular}
\end{table}
\end{verbatim}

Tables shall be referenced using:

\begin{verbatim}
Table~\ref{tab:internship-schedule}
\end{verbatim}

Students shall not manually type table numbers.

Tables should be created directly in \LaTeX{} rather than inserted as
screenshots whenever practical.

\vspace{0.3cm}

% ================================================================
% XIV. EQUATIONS AND TECHNICAL QUANTITIES
% ================================================================

\noindent
\textbf{XIV.\quad EQUATIONS AND TECHNICAL QUANTITIES}

\vspace{0.15cm}

Mathematical expressions shall be typeset directly using \LaTeX{}.
Images or screenshots of equations shall not be used.

Short expressions may be written in inline mathematical mode, for example:

\begin{verbatim}
$R = 0.7\,\Omega$
\end{verbatim}

Important equations shall use a numbered equation environment.

Example:

\begin{verbatim}
\begin{equation}
    P = UI
    \label{eq:power}
\end{equation}
\end{verbatim}

Equations shall be referenced using:

\begin{verbatim}
Equation~\eqref{eq:power}
\end{verbatim}

Equation numbers shall not be manually entered.

SI units shall be used wherever applicable. The \verb|siunitx| package
should be used where appropriate.

Examples:

\begin{verbatim}
\SI{220}{\volt}
\SI{50}{\hertz}
\SI{2.5}{\ampere}
\end{verbatim}

Notation and units shall be consistent throughout the report.

\vspace{0.3cm}

% ================================================================
% XV. SOURCE CODE
% ================================================================

\noindent
\textbf{XV.\quad SOURCE CODE}

\vspace{0.15cm}

If the internship involves programming, PLCs, microcontrollers,
simulation, or software development, relevant source code may be included
in the report.

Source code shall be presented using the \verb|listings| environment or
another environment defined by the official template.

Example:

\begin{verbatim}
\begin{lstlisting}[language=st]
IF Start = TRUE THEN
    Motor := TRUE;
ELSE
    Motor := FALSE;
END_IF;
\end{lstlisting}
\end{verbatim}

Only code necessary for explaining the work should appear in the main
report. Long source-code listings should be moved to an appendix or
submitted as supplementary files if required.

Source code should not normally be inserted as screenshots.

\vspace{0.3cm}

% ================================================================
% XVI. ABSTRACT
% ================================================================

\noindent
\textbf{XVI.\quad ABSTRACT}

\vspace{0.15cm}

The Abstract shall briefly summarize:

\begin{itemize}[leftmargin=0.8cm, itemsep=2pt]
    \item the internship organization;
    \item the objectives of the internship;
    \item the main tasks performed;
    \item the main results, knowledge, and skills obtained.
\end{itemize}

The Abstract should be concise and should not reproduce the full content
of the report chapter by chapter.

\vspace{0.3cm}

% ================================================================
% XVII. REFERENCES AND CITATIONS
% ================================================================

\noindent
\textbf{XVII.\quad REFERENCES AND CITATIONS}

\vspace{0.15cm}

All books, journal papers, conference papers, standards, websites,
datasheets, manuals, technical documents, figures, tables, or other
external sources used in preparing the report shall be properly cited.

References shall be managed through the bibliography system provided with
the official template.

Bibliographic entries shall normally be stored in:

\begin{verbatim}
References.bib
\end{verbatim}

References shall be cited in the report using:

\begin{verbatim}
\cite{reference-key}
\end{verbatim}

Students shall not manually type reference numbers such as ``[1]'',
``[2]'', etc.

Titles of foreign-language publications shall remain in their original
language in the bibliography.

All borrowed ideas, data, figures, tables, equations, methods, or results
shall be properly acknowledged.

\vspace{0.3cm}

% ================================================================
% XVIII. CROSS-REFERENCING
% ================================================================

\noindent
\textbf{XVIII.\quad CROSS-REFERENCING}

\vspace{0.15cm}

Numbered elements shall be cross-referenced using the standard \LaTeX{}
mechanism.

Use the following commands where appropriate:

\begin{verbatim}
\label{...}
\ref{...}
\eqref{...}
\end{verbatim}

For example:

\begin{verbatim}
As shown in Figure~\ref{fig:system}, ...
\end{verbatim}

Students shall not manually type numbers such as ``Figure 2.3'' when
those numbers can be generated automatically.

\vspace{0.3cm}

% ================================================================
% XIX. PAGE NUMBERING
% ================================================================

\noindent
\textbf{XIX.\quad PAGE NUMBERING}

\vspace{0.15cm}

Cover pages and official evaluation forms at the beginning of the report
shall not display page numbers.

The main body of the report shall use Arabic page numbers beginning from
Chapter 1.

Page numbering shall be generated automatically by the official template.
Students shall not manually enter page numbers.

\vspace{0.3cm}

% ================================================================
% XX. PROJECT FILE ORGANIZATION
% ================================================================

\noindent
\textbf{XX.\quad PROJECT FILE ORGANIZATION}

\vspace{0.15cm}

Students should maintain a clear and organized \LaTeX{} project
structure. A recommended structure is:

\begin{verbatim}
InternshipReport/
|
|-- main.tex
|-- References.bib
|
|-- figures/
|   |-- company.pdf
|   |-- system.png
|   `-- ...
|
`-- chapters/
    |-- chapter1.tex
    |-- chapter2.tex
    `-- chapter3.tex
\end{verbatim}

Figures should normally be stored in the \verb|figures| directory
specified by the official template.

File and directory names should preferably contain Latin letters,
numbers, hyphens, or underscores.

Recommended examples include:

\begin{verbatim}
company_structure.pdf
motor_control.png
chapter_01.tex
\end{verbatim}

Spaces and unnecessary special characters in file names should be
avoided where possible.

\vspace{0.3cm}

% ================================================================
% XXI. DOCUMENT COMPILATION AND FINAL CHECK
% ================================================================

\noindent
\textbf{XXI.\quad DOCUMENT COMPILATION AND FINAL CHECK}

\vspace{0.15cm}

Before submission, the complete \LaTeX{} project shall compile
successfully and produce the final PDF report.

The final document shall not contain unresolved references such as:

\begin{verbatim}
??
[?]
\end{verbatim}

Before generating the final PDF, students shall check:

\begin{itemize}[leftmargin=0.8cm, itemsep=2pt]
    \item the Table of Contents;
    \item page numbering;
    \item chapter, section, and subsection numbering;
    \item figure numbering;
    \item table numbering;
    \item equation numbering;
    \item cross-references;
    \item citations;
    \item bibliography;
    \item font display;
    \item mathematical symbols;
    \item image quality;
    \item completeness of all required forms.
\end{itemize}

The document may be compiled using TeX Live, MiKTeX, or another
compatible \LaTeX{} environment. Regardless of the compilation
environment, the resulting PDF shall conform to the official Faculty
template.

\vspace{0.3cm}

% ================================================================
% XXII. CONTENT THAT STUDENTS MAY EDIT
% ================================================================

\noindent
\textbf{XXII.\quad CONTENT THAT STUDENTS MAY EDIT}

\vspace{0.15cm}

Students may edit or replace:

\begin{itemize}[leftmargin=0.8cm, itemsep=2pt]
    \item student name;
    \item class;
    \item internship organization;
    \item supervisor name;
    \item Abstract;
    \item List of Symbols and Abbreviations;
    \item chapter titles and chapter content;
    \item figures and photographs;
    \item tables;
    \item equations;
    \item citations and references.
\end{itemize}

Students shall not modify the main document formatting settings in the
preamble, including the document class, page geometry, line spacing,
heading formatting, table-of-contents formatting, font configuration,
page layout, or cover-page macros unless specifically instructed.

\vspace{0.3cm}

% ================================================================
% XXIII. FINAL SUBMISSION REQUIREMENTS
% ================================================================

\noindent
\textbf{XXIII.\quad FINAL SUBMISSION REQUIREMENTS}

\vspace{0.15cm}

Before submission, the final report shall:

\begin{itemize}[leftmargin=0.8cm, itemsep=2pt]
    \item compile without errors that affect the report content or presentation;
    \item contain all required forms and report sections;
    \item contain correct student, class, internship organization, and supervisor information;
    \item contain an updated Table of Contents, List of Figures, and List of Tables;
    \item contain no unused example content from the template;
    \item contain no placeholder text such as \textit{Title of Chapter 1},
    \textit{Section 1.1}, or other sample entries;
    \item contain all required citations and references;
    \item contain no unresolved cross-references.
\end{itemize}

If the Faculty requires submission of the \LaTeX{} source, the complete
project shall include all \verb|.tex| files, the bibliography database,
figures, and any other files required to reproduce the submitted PDF.

\vspace{0.4cm}

\noindent
\textbf{Important:}
Before submitting the final internship report, students shall remove or
comment out the following command in the main document:

\begin{verbatim}
\input{Regulations.tex}
\end{verbatim}

For example:

\begin{verbatim}
% \input{Regulations.tex}
\end{verbatim}

This removes the guidance pages from the final submitted report while
preserving the regulations file in the \LaTeX{} project for reference.

} % End 13pt content

\clearpage
 command before submitting
% the final internship report.
%
% =================================================================

\clearpage
\thispagestyle{empty}

\begin{center}
    {\fontsize{16pt}{20pt}\selectfont\bfseries
    SPECIFIC REGULATIONS FOR PREPARING\\
    GRADUATION INTERNSHIP REPORTS USING \LaTeX}
\end{center}

\vspace{0.5cm}

{\fontsize{13pt}{17pt}\selectfont
\setstretch{1.3}
\justifying

\noindent
\textbf{Note:}
This section is provided for guidance only and is included in the main
document using the command \verb|
% =================================================================
% REGULATIONS FOR PREPARING THE GRADUATION INTERNSHIP REPORT
% USING LATEX
% =================================================================
%
% This file is intended to be included from the main document with:
%
%     \input{Regulations.tex}
%
% IMPORTANT:
% This section is provided for guidance only. Students shall remove
% or comment out the \input{Regulations.tex} command before submitting
% the final internship report.
%
% =================================================================

\clearpage
\thispagestyle{empty}

\begin{center}
    {\fontsize{16pt}{20pt}\selectfont\bfseries
    SPECIFIC REGULATIONS FOR PREPARING\\
    GRADUATION INTERNSHIP REPORTS USING \LaTeX}
\end{center}

\vspace{0.5cm}

{\fontsize{13pt}{17pt}\selectfont
\setstretch{1.3}
\justifying

\noindent
\textbf{Note:}
This section is provided for guidance only and is included in the main
document using the command \verb|\input{Regulations.tex}|.
Students shall remove or comment out this command before submitting the
final internship report.

\vspace{0.4cm}

% ================================================================
% I. GENERAL REQUIREMENTS
% ================================================================

\noindent
\textbf{I.\quad GENERAL REQUIREMENTS}

\vspace{0.15cm}

The Graduation Internship Report shall be prepared using the official
\LaTeX{} template provided by the Faculty. Students shall preserve the
predefined page layout, typography, line spacing, heading styles,
numbering system, figure and table formatting, and bibliography settings
unless otherwise instructed by the Faculty or the supervisor.

The report shall be written in English.

Students are responsible for the academic and technical content of the
report. The official \LaTeX{} template is responsible for standardized
document formatting.

\vspace{0.3cm}

% ================================================================
% II. REPORT STRUCTURE
% ================================================================

\noindent
\textbf{II.\quad REPORT STRUCTURE}

\vspace{0.15cm}

The report shall normally be arranged in the following order:

\begin{enumerate}[label=\arabic*., leftmargin=1.2cm, itemsep=2pt]
    \item Outer Cover;
    \item Comments of the Internship Host Organization;
    \item Supervisor's Evaluation and Graduation Internship Defense Results;
    \item Table of Contents;
    \item List of Symbols and Abbreviations, if applicable;
    \item List of Figures;
    \item List of Tables;
    \item Abstract;
    \item Main Chapters of the Internship Report;
    \item References;
    \item Appendices, if applicable.
\end{enumerate}

The preliminary pages shall not be numbered as chapters. The main report
content shall be organized into consecutively numbered chapters beginning
with Chapter 1.

\vspace{0.3cm}

% ================================================================
% III. PAGE SIZE AND MARGINS
% ================================================================

\noindent
\textbf{III.\quad PAGE SIZE AND MARGINS}

\vspace{0.15cm}

The report shall be formatted on A4 paper with the following margins:

\begin{center}
\begin{tabular}{ll}
    Left margin:   & 3.5 cm \\
    Right margin:  & 2.0 cm \\
    Top margin:    & 2.5 cm \\
    Bottom margin: & 2.5 cm
\end{tabular}
\end{center}

These settings are predefined in the official template and shall not be
modified by students.

\vspace{0.3cm}

% ================================================================
% IV. FONT, FONT SIZE, AND LINE SPACING
% ================================================================

\noindent
\textbf{IV.\quad FONT, FONT SIZE, AND LINE SPACING}

\vspace{0.15cm}

The main text shall use a 13 pt font size with 1.3-line spacing.
A Times-family typeface shall be used in accordance with the official
template.

The text shall be fully justified.

Students shall not use repeated spaces, repeated manual line breaks, or
multiple blank lines to position text. Paragraph and heading spacing shall
be controlled by the template.

\vspace{0.3cm}

% ================================================================
% V. OUTER COVER
% ================================================================

\noindent
\textbf{V.\quad OUTER COVER}

\vspace{0.15cm}

The outer cover shall follow the official format provided in the template
and shall contain the following information:

\begin{itemize}[leftmargin=0.8cm, itemsep=2pt]
    \item Thai Nguyen University of Technology;
    \item Faculty of International Training;
    \item Engineering and Technology Department;
    \item University logo and Faculty logo;
    \item the title \textbf{INTERNSHIP REPORT};
    \item Internship Organization;
    \item Author;
    \item Class;
    \item Supervisor;
    \item place and year of completion.
\end{itemize}

Students shall only replace the required information fields and shall not
modify the cover-page layout unless instructed otherwise.

\vspace{0.3cm}

% ================================================================
% VI. COMMENTS OF THE INTERNSHIP HOST ORGANIZATION
% ================================================================

\noindent
\textbf{VI.\quad COMMENTS OF THE INTERNSHIP HOST ORGANIZATION}

\vspace{0.15cm}

The page entitled \textbf{COMMENTS OF THE INTERNSHIP HOST ORGANIZATION}
is reserved for the internship host organization.

It includes:

\begin{itemize}[leftmargin=0.8cm, itemsep=2pt]
    \item the student's full name;
    \item class;
    \item comments from the host organization;
    \item score;
    \item date;
    \item signature and official seal.
\end{itemize}

Students shall not alter or complete official evaluation fields unless
specifically authorized.

\vspace{0.3cm}

% ================================================================
% VII. SUPERVISOR'S EVALUATION AND DEFENSE RESULTS
% ================================================================

\noindent
\textbf{VII.\quad SUPERVISOR'S EVALUATION AND DEFENSE RESULTS}

\vspace{0.15cm}

The section entitled
\textbf{SUPERVISOR'S EVALUATION AND GRADUATION INTERNSHIP DEFENSE RESULTS}
is reserved for the supervisor and defense committee.

The evaluation may include:

\begin{itemize}[leftmargin=0.8cm, itemsep=2pt]
    \item Internship Progress and Student's Attitude;
    \item Report Content;
    \item Report Presentation;
    \item Other Comments;
    \item Supervisor's Opinion;
    \item Defense Result;
    \item Score.
\end{itemize}

Students shall not enter evaluation scores or defense results before they
are officially determined.

\vspace{0.3cm}

% ================================================================
% VIII. CHAPTERS, SECTIONS, AND SUBSECTIONS
% ================================================================

\noindent
\textbf{VIII.\quad CHAPTERS, SECTIONS, AND SUBSECTIONS}

\vspace{0.15cm}

The report structure shall be generated using the standard \LaTeX{}
commands:

\begin{verbatim}
\chapter{...}
\section{...}
\subsection{...}
\end{verbatim}

Chapter, section, and subsection numbers shall be generated automatically
by \LaTeX{}.

Students shall not manually type chapter or section numbers and shall not
reproduce heading styles using commands such as \verb|\textbf|,
\verb|\fontsize|, or \verb|\centering| instead of the appropriate
structural command.

Excessive subdivision should be avoided. In most cases, the hierarchy
Chapter--Section--Subsection is sufficient.

\vspace{0.3cm}

% ================================================================
% IX. ORGANIZATION OF THE MAIN CONTENT
% ================================================================

\noindent
\textbf{IX.\quad ORGANIZATION OF THE MAIN CONTENT}

\vspace{0.15cm}

The specific content of the report shall depend on the internship
organization and the assigned tasks. A typical structure may include:

\begin{itemize}[leftmargin=0.8cm, itemsep=2pt]
    \item \textbf{Chapter 1. Overview of the Internship Organization};
    \item \textbf{Chapter 2. Internship Activities and Assigned Tasks};
    \item \textbf{Chapter 3. Results, Experience, and Evaluation}.
\end{itemize}

These chapter titles are recommendations only and may be adapted to the
actual internship content.

The report shall reflect work actually performed, observed, analyzed, or
learned during the internship. Unrelated theoretical material should not
be included merely to increase report length.

\vspace{0.3cm}

% ================================================================
% X. TABLE OF CONTENTS
% ================================================================

\noindent
\textbf{X.\quad TABLE OF CONTENTS}

\vspace{0.15cm}

The Table of Contents shall be generated automatically using:

\begin{verbatim}
\tableofcontents
\end{verbatim}

Students shall not create or manually edit page numbers in the Table of
Contents.

The document shall be compiled sufficiently to ensure that all page
numbers and entries are correctly updated.

\vspace{0.3cm}

% ================================================================
% XI. LIST OF SYMBOLS AND ABBREVIATIONS
% ================================================================

\noindent
\textbf{XI.\quad LIST OF SYMBOLS AND ABBREVIATIONS}

\vspace{0.15cm}

Symbols and abbreviations that are used repeatedly in the report shall be
listed in the \textbf{List of Symbols and Abbreviations}.

The list may contain the following columns:

\begin{center}
\begin{tabular}{llll}
    No. & Symbol / Abbreviation & Unit & Definition
\end{tabular}
\end{center}

Only symbols and abbreviations actually used in the final report shall be
retained. Example entries contained in the template shall be removed if
they are not relevant to the student's report.

Each abbreviation should normally be defined when it first appears in
the main text.

\vspace{0.3cm}

% ================================================================
% XII. FIGURES, DIAGRAMS, AND PHOTOGRAPHS
% ================================================================

\noindent
\textbf{XII.\quad FIGURES, DIAGRAMS, AND PHOTOGRAPHS}

\vspace{0.15cm}

Figures shall be clear, legible, and directly relevant to the report
content.

Figures shall be inserted using the standard \verb|figure| environment.

Example:

\begin{verbatim}
\begin{figure}[H]
    \centering
    \includegraphics[width=0.8\textwidth]
                    {figures/company_structure.pdf}
    \caption{Organizational structure of the company}
    \label{fig:company-structure}
\end{figure}
\end{verbatim}

Figure captions shall be placed below the figure.

Each referenced figure shall have a unique \verb|\label| and shall be
referenced in the text using, for example:

\begin{verbatim}
Figure~\ref{fig:company-structure}
\end{verbatim}

Students shall not manually type figure numbers.

Figures reproduced or adapted from external sources shall include an
appropriate citation.

Photographs taken during the internship shall be relevant and sufficiently
clear. Students shall not include confidential, restricted, or proprietary
information without permission from the internship organization.

Vector graphics or high-resolution images are recommended whenever
possible.

\vspace{0.3cm}

% ================================================================
% XIII. TABLES
% ================================================================

\noindent
\textbf{XIII.\quad TABLES}

\vspace{0.15cm}

Tables shall be created using suitable \LaTeX{} environments such as
\verb|table|, \verb|tabular|, \verb|tabularx|, or \verb|longtable|.

Table captions shall be placed above the table.

Example:

\begin{verbatim}
\begin{table}[H]
    \centering
    \caption{Internship schedule}
    \label{tab:internship-schedule}

    \begin{tabular}{lll}
        \toprule
        Week & Task & Result \\
        \midrule
        1 & Orientation & Completed \\
        2 & System study & Completed \\
        \bottomrule
    \end{tabular}
\end{table}
\end{verbatim}

Tables shall be referenced using:

\begin{verbatim}
Table~\ref{tab:internship-schedule}
\end{verbatim}

Students shall not manually type table numbers.

Tables should be created directly in \LaTeX{} rather than inserted as
screenshots whenever practical.

\vspace{0.3cm}

% ================================================================
% XIV. EQUATIONS AND TECHNICAL QUANTITIES
% ================================================================

\noindent
\textbf{XIV.\quad EQUATIONS AND TECHNICAL QUANTITIES}

\vspace{0.15cm}

Mathematical expressions shall be typeset directly using \LaTeX{}.
Images or screenshots of equations shall not be used.

Short expressions may be written in inline mathematical mode, for example:

\begin{verbatim}
$R = 0.7\,\Omega$
\end{verbatim}

Important equations shall use a numbered equation environment.

Example:

\begin{verbatim}
\begin{equation}
    P = UI
    \label{eq:power}
\end{equation}
\end{verbatim}

Equations shall be referenced using:

\begin{verbatim}
Equation~\eqref{eq:power}
\end{verbatim}

Equation numbers shall not be manually entered.

SI units shall be used wherever applicable. The \verb|siunitx| package
should be used where appropriate.

Examples:

\begin{verbatim}
\SI{220}{\volt}
\SI{50}{\hertz}
\SI{2.5}{\ampere}
\end{verbatim}

Notation and units shall be consistent throughout the report.

\vspace{0.3cm}

% ================================================================
% XV. SOURCE CODE
% ================================================================

\noindent
\textbf{XV.\quad SOURCE CODE}

\vspace{0.15cm}

If the internship involves programming, PLCs, microcontrollers,
simulation, or software development, relevant source code may be included
in the report.

Source code shall be presented using the \verb|listings| environment or
another environment defined by the official template.

Example:

\begin{verbatim}
\begin{lstlisting}[language=st]
IF Start = TRUE THEN
    Motor := TRUE;
ELSE
    Motor := FALSE;
END_IF;
\end{lstlisting}
\end{verbatim}

Only code necessary for explaining the work should appear in the main
report. Long source-code listings should be moved to an appendix or
submitted as supplementary files if required.

Source code should not normally be inserted as screenshots.

\vspace{0.3cm}

% ================================================================
% XVI. ABSTRACT
% ================================================================

\noindent
\textbf{XVI.\quad ABSTRACT}

\vspace{0.15cm}

The Abstract shall briefly summarize:

\begin{itemize}[leftmargin=0.8cm, itemsep=2pt]
    \item the internship organization;
    \item the objectives of the internship;
    \item the main tasks performed;
    \item the main results, knowledge, and skills obtained.
\end{itemize}

The Abstract should be concise and should not reproduce the full content
of the report chapter by chapter.

\vspace{0.3cm}

% ================================================================
% XVII. REFERENCES AND CITATIONS
% ================================================================

\noindent
\textbf{XVII.\quad REFERENCES AND CITATIONS}

\vspace{0.15cm}

All books, journal papers, conference papers, standards, websites,
datasheets, manuals, technical documents, figures, tables, or other
external sources used in preparing the report shall be properly cited.

References shall be managed through the bibliography system provided with
the official template.

Bibliographic entries shall normally be stored in:

\begin{verbatim}
References.bib
\end{verbatim}

References shall be cited in the report using:

\begin{verbatim}
\cite{reference-key}
\end{verbatim}

Students shall not manually type reference numbers such as ``[1]'',
``[2]'', etc.

Titles of foreign-language publications shall remain in their original
language in the bibliography.

All borrowed ideas, data, figures, tables, equations, methods, or results
shall be properly acknowledged.

\vspace{0.3cm}

% ================================================================
% XVIII. CROSS-REFERENCING
% ================================================================

\noindent
\textbf{XVIII.\quad CROSS-REFERENCING}

\vspace{0.15cm}

Numbered elements shall be cross-referenced using the standard \LaTeX{}
mechanism.

Use the following commands where appropriate:

\begin{verbatim}
\label{...}
\ref{...}
\eqref{...}
\end{verbatim}

For example:

\begin{verbatim}
As shown in Figure~\ref{fig:system}, ...
\end{verbatim}

Students shall not manually type numbers such as ``Figure 2.3'' when
those numbers can be generated automatically.

\vspace{0.3cm}

% ================================================================
% XIX. PAGE NUMBERING
% ================================================================

\noindent
\textbf{XIX.\quad PAGE NUMBERING}

\vspace{0.15cm}

Cover pages and official evaluation forms at the beginning of the report
shall not display page numbers.

The main body of the report shall use Arabic page numbers beginning from
Chapter 1.

Page numbering shall be generated automatically by the official template.
Students shall not manually enter page numbers.

\vspace{0.3cm}

% ================================================================
% XX. PROJECT FILE ORGANIZATION
% ================================================================

\noindent
\textbf{XX.\quad PROJECT FILE ORGANIZATION}

\vspace{0.15cm}

Students should maintain a clear and organized \LaTeX{} project
structure. A recommended structure is:

\begin{verbatim}
InternshipReport/
|
|-- main.tex
|-- References.bib
|
|-- figures/
|   |-- company.pdf
|   |-- system.png
|   `-- ...
|
`-- chapters/
    |-- chapter1.tex
    |-- chapter2.tex
    `-- chapter3.tex
\end{verbatim}

Figures should normally be stored in the \verb|figures| directory
specified by the official template.

File and directory names should preferably contain Latin letters,
numbers, hyphens, or underscores.

Recommended examples include:

\begin{verbatim}
company_structure.pdf
motor_control.png
chapter_01.tex
\end{verbatim}

Spaces and unnecessary special characters in file names should be
avoided where possible.

\vspace{0.3cm}

% ================================================================
% XXI. DOCUMENT COMPILATION AND FINAL CHECK
% ================================================================

\noindent
\textbf{XXI.\quad DOCUMENT COMPILATION AND FINAL CHECK}

\vspace{0.15cm}

Before submission, the complete \LaTeX{} project shall compile
successfully and produce the final PDF report.

The final document shall not contain unresolved references such as:

\begin{verbatim}
??
[?]
\end{verbatim}

Before generating the final PDF, students shall check:

\begin{itemize}[leftmargin=0.8cm, itemsep=2pt]
    \item the Table of Contents;
    \item page numbering;
    \item chapter, section, and subsection numbering;
    \item figure numbering;
    \item table numbering;
    \item equation numbering;
    \item cross-references;
    \item citations;
    \item bibliography;
    \item font display;
    \item mathematical symbols;
    \item image quality;
    \item completeness of all required forms.
\end{itemize}

The document may be compiled using TeX Live, MiKTeX, or another
compatible \LaTeX{} environment. Regardless of the compilation
environment, the resulting PDF shall conform to the official Faculty
template.

\vspace{0.3cm}

% ================================================================
% XXII. CONTENT THAT STUDENTS MAY EDIT
% ================================================================

\noindent
\textbf{XXII.\quad CONTENT THAT STUDENTS MAY EDIT}

\vspace{0.15cm}

Students may edit or replace:

\begin{itemize}[leftmargin=0.8cm, itemsep=2pt]
    \item student name;
    \item class;
    \item internship organization;
    \item supervisor name;
    \item Abstract;
    \item List of Symbols and Abbreviations;
    \item chapter titles and chapter content;
    \item figures and photographs;
    \item tables;
    \item equations;
    \item citations and references.
\end{itemize}

Students shall not modify the main document formatting settings in the
preamble, including the document class, page geometry, line spacing,
heading formatting, table-of-contents formatting, font configuration,
page layout, or cover-page macros unless specifically instructed.

\vspace{0.3cm}

% ================================================================
% XXIII. FINAL SUBMISSION REQUIREMENTS
% ================================================================

\noindent
\textbf{XXIII.\quad FINAL SUBMISSION REQUIREMENTS}

\vspace{0.15cm}

Before submission, the final report shall:

\begin{itemize}[leftmargin=0.8cm, itemsep=2pt]
    \item compile without errors that affect the report content or presentation;
    \item contain all required forms and report sections;
    \item contain correct student, class, internship organization, and supervisor information;
    \item contain an updated Table of Contents, List of Figures, and List of Tables;
    \item contain no unused example content from the template;
    \item contain no placeholder text such as \textit{Title of Chapter 1},
    \textit{Section 1.1}, or other sample entries;
    \item contain all required citations and references;
    \item contain no unresolved cross-references.
\end{itemize}

If the Faculty requires submission of the \LaTeX{} source, the complete
project shall include all \verb|.tex| files, the bibliography database,
figures, and any other files required to reproduce the submitted PDF.

\vspace{0.4cm}

\noindent
\textbf{Important:}
Before submitting the final internship report, students shall remove or
comment out the following command in the main document:

\begin{verbatim}
\input{Regulations.tex}
\end{verbatim}

For example:

\begin{verbatim}
% \input{Regulations.tex}
\end{verbatim}

This removes the guidance pages from the final submitted report while
preserving the regulations file in the \LaTeX{} project for reference.

} % End 13pt content

\clearpage
|.
Students shall remove or comment out this command before submitting the
final internship report.

\vspace{0.4cm}

% ================================================================
% I. GENERAL REQUIREMENTS
% ================================================================

\noindent
\textbf{I.\quad GENERAL REQUIREMENTS}

\vspace{0.15cm}

The Graduation Internship Report shall be prepared using the official
\LaTeX{} template provided by the Faculty. Students shall preserve the
predefined page layout, typography, line spacing, heading styles,
numbering system, figure and table formatting, and bibliography settings
unless otherwise instructed by the Faculty or the supervisor.

The report shall be written in English.

Students are responsible for the academic and technical content of the
report. The official \LaTeX{} template is responsible for standardized
document formatting.

\vspace{0.3cm}

% ================================================================
% II. REPORT STRUCTURE
% ================================================================

\noindent
\textbf{II.\quad REPORT STRUCTURE}

\vspace{0.15cm}

The report shall normally be arranged in the following order:

\begin{enumerate}[label=\arabic*., leftmargin=1.2cm, itemsep=2pt]
    \item Outer Cover;
    \item Comments of the Internship Host Organization;
    \item Supervisor's Evaluation and Graduation Internship Defense Results;
    \item Table of Contents;
    \item List of Symbols and Abbreviations, if applicable;
    \item List of Figures;
    \item List of Tables;
    \item Abstract;
    \item Main Chapters of the Internship Report;
    \item References;
    \item Appendices, if applicable.
\end{enumerate}

The preliminary pages shall not be numbered as chapters. The main report
content shall be organized into consecutively numbered chapters beginning
with Chapter 1.

\vspace{0.3cm}

% ================================================================
% III. PAGE SIZE AND MARGINS
% ================================================================

\noindent
\textbf{III.\quad PAGE SIZE AND MARGINS}

\vspace{0.15cm}

The report shall be formatted on A4 paper with the following margins:

\begin{center}
\begin{tabular}{ll}
    Left margin:   & 3.5 cm \\
    Right margin:  & 2.0 cm \\
    Top margin:    & 2.5 cm \\
    Bottom margin: & 2.5 cm
\end{tabular}
\end{center}

These settings are predefined in the official template and shall not be
modified by students.

\vspace{0.3cm}

% ================================================================
% IV. FONT, FONT SIZE, AND LINE SPACING
% ================================================================

\noindent
\textbf{IV.\quad FONT, FONT SIZE, AND LINE SPACING}

\vspace{0.15cm}

The main text shall use a 13 pt font size with 1.3-line spacing.
A Times-family typeface shall be used in accordance with the official
template.

The text shall be fully justified.

Students shall not use repeated spaces, repeated manual line breaks, or
multiple blank lines to position text. Paragraph and heading spacing shall
be controlled by the template.

\vspace{0.3cm}

% ================================================================
% V. OUTER COVER
% ================================================================

\noindent
\textbf{V.\quad OUTER COVER}

\vspace{0.15cm}

The outer cover shall follow the official format provided in the template
and shall contain the following information:

\begin{itemize}[leftmargin=0.8cm, itemsep=2pt]
    \item Thai Nguyen University of Technology;
    \item Faculty of International Training;
    \item Engineering and Technology Department;
    \item University logo and Faculty logo;
    \item the title \textbf{INTERNSHIP REPORT};
    \item Internship Organization;
    \item Author;
    \item Class;
    \item Supervisor;
    \item place and year of completion.
\end{itemize}

Students shall only replace the required information fields and shall not
modify the cover-page layout unless instructed otherwise.

\vspace{0.3cm}

% ================================================================
% VI. COMMENTS OF THE INTERNSHIP HOST ORGANIZATION
% ================================================================

\noindent
\textbf{VI.\quad COMMENTS OF THE INTERNSHIP HOST ORGANIZATION}

\vspace{0.15cm}

The page entitled \textbf{COMMENTS OF THE INTERNSHIP HOST ORGANIZATION}
is reserved for the internship host organization.

It includes:

\begin{itemize}[leftmargin=0.8cm, itemsep=2pt]
    \item the student's full name;
    \item class;
    \item comments from the host organization;
    \item score;
    \item date;
    \item signature and official seal.
\end{itemize}

Students shall not alter or complete official evaluation fields unless
specifically authorized.

\vspace{0.3cm}

% ================================================================
% VII. SUPERVISOR'S EVALUATION AND DEFENSE RESULTS
% ================================================================

\noindent
\textbf{VII.\quad SUPERVISOR'S EVALUATION AND DEFENSE RESULTS}

\vspace{0.15cm}

The section entitled
\textbf{SUPERVISOR'S EVALUATION AND GRADUATION INTERNSHIP DEFENSE RESULTS}
is reserved for the supervisor and defense committee.

The evaluation may include:

\begin{itemize}[leftmargin=0.8cm, itemsep=2pt]
    \item Internship Progress and Student's Attitude;
    \item Report Content;
    \item Report Presentation;
    \item Other Comments;
    \item Supervisor's Opinion;
    \item Defense Result;
    \item Score.
\end{itemize}

Students shall not enter evaluation scores or defense results before they
are officially determined.

\vspace{0.3cm}

% ================================================================
% VIII. CHAPTERS, SECTIONS, AND SUBSECTIONS
% ================================================================

\noindent
\textbf{VIII.\quad CHAPTERS, SECTIONS, AND SUBSECTIONS}

\vspace{0.15cm}

The report structure shall be generated using the standard \LaTeX{}
commands:

\begin{verbatim}
\chapter{...}
\section{...}
\subsection{...}
\end{verbatim}

Chapter, section, and subsection numbers shall be generated automatically
by \LaTeX{}.

Students shall not manually type chapter or section numbers and shall not
reproduce heading styles using commands such as \verb|\textbf|,
\verb|\fontsize|, or \verb|\centering| instead of the appropriate
structural command.

Excessive subdivision should be avoided. In most cases, the hierarchy
Chapter--Section--Subsection is sufficient.

\vspace{0.3cm}

% ================================================================
% IX. ORGANIZATION OF THE MAIN CONTENT
% ================================================================

\noindent
\textbf{IX.\quad ORGANIZATION OF THE MAIN CONTENT}

\vspace{0.15cm}

The specific content of the report shall depend on the internship
organization and the assigned tasks. A typical structure may include:

\begin{itemize}[leftmargin=0.8cm, itemsep=2pt]
    \item \textbf{Chapter 1. Overview of the Internship Organization};
    \item \textbf{Chapter 2. Internship Activities and Assigned Tasks};
    \item \textbf{Chapter 3. Results, Experience, and Evaluation}.
\end{itemize}

These chapter titles are recommendations only and may be adapted to the
actual internship content.

The report shall reflect work actually performed, observed, analyzed, or
learned during the internship. Unrelated theoretical material should not
be included merely to increase report length.

\vspace{0.3cm}

% ================================================================
% X. TABLE OF CONTENTS
% ================================================================

\noindent
\textbf{X.\quad TABLE OF CONTENTS}

\vspace{0.15cm}

The Table of Contents shall be generated automatically using:

\begin{verbatim}
\tableofcontents
\end{verbatim}

Students shall not create or manually edit page numbers in the Table of
Contents.

The document shall be compiled sufficiently to ensure that all page
numbers and entries are correctly updated.

\vspace{0.3cm}

% ================================================================
% XI. LIST OF SYMBOLS AND ABBREVIATIONS
% ================================================================

\noindent
\textbf{XI.\quad LIST OF SYMBOLS AND ABBREVIATIONS}

\vspace{0.15cm}

Symbols and abbreviations that are used repeatedly in the report shall be
listed in the \textbf{List of Symbols and Abbreviations}.

The list may contain the following columns:

\begin{center}
\begin{tabular}{llll}
    No. & Symbol / Abbreviation & Unit & Definition
\end{tabular}
\end{center}

Only symbols and abbreviations actually used in the final report shall be
retained. Example entries contained in the template shall be removed if
they are not relevant to the student's report.

Each abbreviation should normally be defined when it first appears in
the main text.

\vspace{0.3cm}

% ================================================================
% XII. FIGURES, DIAGRAMS, AND PHOTOGRAPHS
% ================================================================

\noindent
\textbf{XII.\quad FIGURES, DIAGRAMS, AND PHOTOGRAPHS}

\vspace{0.15cm}

Figures shall be clear, legible, and directly relevant to the report
content.

Figures shall be inserted using the standard \verb|figure| environment.

Example:

\begin{verbatim}
\begin{figure}[H]
    \centering
    \includegraphics[width=0.8\textwidth]
                    {figures/company_structure.pdf}
    \caption{Organizational structure of the company}
    \label{fig:company-structure}
\end{figure}
\end{verbatim}

Figure captions shall be placed below the figure.

Each referenced figure shall have a unique \verb|\label| and shall be
referenced in the text using, for example:

\begin{verbatim}
Figure~\ref{fig:company-structure}
\end{verbatim}

Students shall not manually type figure numbers.

Figures reproduced or adapted from external sources shall include an
appropriate citation.

Photographs taken during the internship shall be relevant and sufficiently
clear. Students shall not include confidential, restricted, or proprietary
information without permission from the internship organization.

Vector graphics or high-resolution images are recommended whenever
possible.

\vspace{0.3cm}

% ================================================================
% XIII. TABLES
% ================================================================

\noindent
\textbf{XIII.\quad TABLES}

\vspace{0.15cm}

Tables shall be created using suitable \LaTeX{} environments such as
\verb|table|, \verb|tabular|, \verb|tabularx|, or \verb|longtable|.

Table captions shall be placed above the table.

Example:

\begin{verbatim}
\begin{table}[H]
    \centering
    \caption{Internship schedule}
    \label{tab:internship-schedule}

    \begin{tabular}{lll}
        \toprule
        Week & Task & Result \\
        \midrule
        1 & Orientation & Completed \\
        2 & System study & Completed \\
        \bottomrule
    \end{tabular}
\end{table}
\end{verbatim}

Tables shall be referenced using:

\begin{verbatim}
Table~\ref{tab:internship-schedule}
\end{verbatim}

Students shall not manually type table numbers.

Tables should be created directly in \LaTeX{} rather than inserted as
screenshots whenever practical.

\vspace{0.3cm}

% ================================================================
% XIV. EQUATIONS AND TECHNICAL QUANTITIES
% ================================================================

\noindent
\textbf{XIV.\quad EQUATIONS AND TECHNICAL QUANTITIES}

\vspace{0.15cm}

Mathematical expressions shall be typeset directly using \LaTeX{}.
Images or screenshots of equations shall not be used.

Short expressions may be written in inline mathematical mode, for example:

\begin{verbatim}
$R = 0.7\,\Omega$
\end{verbatim}

Important equations shall use a numbered equation environment.

Example:

\begin{verbatim}
\begin{equation}
    P = UI
    \label{eq:power}
\end{equation}
\end{verbatim}

Equations shall be referenced using:

\begin{verbatim}
Equation~\eqref{eq:power}
\end{verbatim}

Equation numbers shall not be manually entered.

SI units shall be used wherever applicable. The \verb|siunitx| package
should be used where appropriate.

Examples:

\begin{verbatim}
\SI{220}{\volt}
\SI{50}{\hertz}
\SI{2.5}{\ampere}
\end{verbatim}

Notation and units shall be consistent throughout the report.

\vspace{0.3cm}

% ================================================================
% XV. SOURCE CODE
% ================================================================

\noindent
\textbf{XV.\quad SOURCE CODE}

\vspace{0.15cm}

If the internship involves programming, PLCs, microcontrollers,
simulation, or software development, relevant source code may be included
in the report.

Source code shall be presented using the \verb|listings| environment or
another environment defined by the official template.

Example:

\begin{verbatim}
\begin{lstlisting}[language=st]
IF Start = TRUE THEN
    Motor := TRUE;
ELSE
    Motor := FALSE;
END_IF;
\end{lstlisting}
\end{verbatim}

Only code necessary for explaining the work should appear in the main
report. Long source-code listings should be moved to an appendix or
submitted as supplementary files if required.

Source code should not normally be inserted as screenshots.

\vspace{0.3cm}

% ================================================================
% XVI. ABSTRACT
% ================================================================

\noindent
\textbf{XVI.\quad ABSTRACT}

\vspace{0.15cm}

The Abstract shall briefly summarize:

\begin{itemize}[leftmargin=0.8cm, itemsep=2pt]
    \item the internship organization;
    \item the objectives of the internship;
    \item the main tasks performed;
    \item the main results, knowledge, and skills obtained.
\end{itemize}

The Abstract should be concise and should not reproduce the full content
of the report chapter by chapter.

\vspace{0.3cm}

% ================================================================
% XVII. REFERENCES AND CITATIONS
% ================================================================

\noindent
\textbf{XVII.\quad REFERENCES AND CITATIONS}

\vspace{0.15cm}

All books, journal papers, conference papers, standards, websites,
datasheets, manuals, technical documents, figures, tables, or other
external sources used in preparing the report shall be properly cited.

References shall be managed through the bibliography system provided with
the official template.

Bibliographic entries shall normally be stored in:

\begin{verbatim}
References.bib
\end{verbatim}

References shall be cited in the report using:

\begin{verbatim}
\cite{reference-key}
\end{verbatim}

Students shall not manually type reference numbers such as ``[1]'',
``[2]'', etc.

Titles of foreign-language publications shall remain in their original
language in the bibliography.

All borrowed ideas, data, figures, tables, equations, methods, or results
shall be properly acknowledged.

\vspace{0.3cm}

% ================================================================
% XVIII. CROSS-REFERENCING
% ================================================================

\noindent
\textbf{XVIII.\quad CROSS-REFERENCING}

\vspace{0.15cm}

Numbered elements shall be cross-referenced using the standard \LaTeX{}
mechanism.

Use the following commands where appropriate:

\begin{verbatim}
\label{...}
\ref{...}
\eqref{...}
\end{verbatim}

For example:

\begin{verbatim}
As shown in Figure~\ref{fig:system}, ...
\end{verbatim}

Students shall not manually type numbers such as ``Figure 2.3'' when
those numbers can be generated automatically.

\vspace{0.3cm}

% ================================================================
% XIX. PAGE NUMBERING
% ================================================================

\noindent
\textbf{XIX.\quad PAGE NUMBERING}

\vspace{0.15cm}

Cover pages and official evaluation forms at the beginning of the report
shall not display page numbers.

The main body of the report shall use Arabic page numbers beginning from
Chapter 1.

Page numbering shall be generated automatically by the official template.
Students shall not manually enter page numbers.

\vspace{0.3cm}

% ================================================================
% XX. PROJECT FILE ORGANIZATION
% ================================================================

\noindent
\textbf{XX.\quad PROJECT FILE ORGANIZATION}

\vspace{0.15cm}

Students should maintain a clear and organized \LaTeX{} project
structure. A recommended structure is:

\begin{verbatim}
InternshipReport/
|
|-- main.tex
|-- References.bib
|
|-- figures/
|   |-- company.pdf
|   |-- system.png
|   `-- ...
|
`-- chapters/
    |-- chapter1.tex
    |-- chapter2.tex
    `-- chapter3.tex
\end{verbatim}

Figures should normally be stored in the \verb|figures| directory
specified by the official template.

File and directory names should preferably contain Latin letters,
numbers, hyphens, or underscores.

Recommended examples include:

\begin{verbatim}
company_structure.pdf
motor_control.png
chapter_01.tex
\end{verbatim}

Spaces and unnecessary special characters in file names should be
avoided where possible.

\vspace{0.3cm}

% ================================================================
% XXI. DOCUMENT COMPILATION AND FINAL CHECK
% ================================================================

\noindent
\textbf{XXI.\quad DOCUMENT COMPILATION AND FINAL CHECK}

\vspace{0.15cm}

Before submission, the complete \LaTeX{} project shall compile
successfully and produce the final PDF report.

The final document shall not contain unresolved references such as:

\begin{verbatim}
??
[?]
\end{verbatim}

Before generating the final PDF, students shall check:

\begin{itemize}[leftmargin=0.8cm, itemsep=2pt]
    \item the Table of Contents;
    \item page numbering;
    \item chapter, section, and subsection numbering;
    \item figure numbering;
    \item table numbering;
    \item equation numbering;
    \item cross-references;
    \item citations;
    \item bibliography;
    \item font display;
    \item mathematical symbols;
    \item image quality;
    \item completeness of all required forms.
\end{itemize}

The document may be compiled using TeX Live, MiKTeX, or another
compatible \LaTeX{} environment. Regardless of the compilation
environment, the resulting PDF shall conform to the official Faculty
template.

\vspace{0.3cm}

% ================================================================
% XXII. CONTENT THAT STUDENTS MAY EDIT
% ================================================================

\noindent
\textbf{XXII.\quad CONTENT THAT STUDENTS MAY EDIT}

\vspace{0.15cm}

Students may edit or replace:

\begin{itemize}[leftmargin=0.8cm, itemsep=2pt]
    \item student name;
    \item class;
    \item internship organization;
    \item supervisor name;
    \item Abstract;
    \item List of Symbols and Abbreviations;
    \item chapter titles and chapter content;
    \item figures and photographs;
    \item tables;
    \item equations;
    \item citations and references.
\end{itemize}

Students shall not modify the main document formatting settings in the
preamble, including the document class, page geometry, line spacing,
heading formatting, table-of-contents formatting, font configuration,
page layout, or cover-page macros unless specifically instructed.

\vspace{0.3cm}

% ================================================================
% XXIII. FINAL SUBMISSION REQUIREMENTS
% ================================================================

\noindent
\textbf{XXIII.\quad FINAL SUBMISSION REQUIREMENTS}

\vspace{0.15cm}

Before submission, the final report shall:

\begin{itemize}[leftmargin=0.8cm, itemsep=2pt]
    \item compile without errors that affect the report content or presentation;
    \item contain all required forms and report sections;
    \item contain correct student, class, internship organization, and supervisor information;
    \item contain an updated Table of Contents, List of Figures, and List of Tables;
    \item contain no unused example content from the template;
    \item contain no placeholder text such as \textit{Title of Chapter 1},
    \textit{Section 1.1}, or other sample entries;
    \item contain all required citations and references;
    \item contain no unresolved cross-references.
\end{itemize}

If the Faculty requires submission of the \LaTeX{} source, the complete
project shall include all \verb|.tex| files, the bibliography database,
figures, and any other files required to reproduce the submitted PDF.

\vspace{0.4cm}

\noindent
\textbf{Important:}
Before submitting the final internship report, students shall remove or
comment out the following command in the main document:

\begin{verbatim}

% =================================================================
% REGULATIONS FOR PREPARING THE GRADUATION INTERNSHIP REPORT
% USING LATEX
% =================================================================
%
% This file is intended to be included from the main document with:
%
%     \input{Regulations.tex}
%
% IMPORTANT:
% This section is provided for guidance only. Students shall remove
% or comment out the \input{Regulations.tex} command before submitting
% the final internship report.
%
% =================================================================

\clearpage
\thispagestyle{empty}

\begin{center}
    {\fontsize{16pt}{20pt}\selectfont\bfseries
    SPECIFIC REGULATIONS FOR PREPARING\\
    GRADUATION INTERNSHIP REPORTS USING \LaTeX}
\end{center}

\vspace{0.5cm}

{\fontsize{13pt}{17pt}\selectfont
\setstretch{1.3}
\justifying

\noindent
\textbf{Note:}
This section is provided for guidance only and is included in the main
document using the command \verb|\input{Regulations.tex}|.
Students shall remove or comment out this command before submitting the
final internship report.

\vspace{0.4cm}

% ================================================================
% I. GENERAL REQUIREMENTS
% ================================================================

\noindent
\textbf{I.\quad GENERAL REQUIREMENTS}

\vspace{0.15cm}

The Graduation Internship Report shall be prepared using the official
\LaTeX{} template provided by the Faculty. Students shall preserve the
predefined page layout, typography, line spacing, heading styles,
numbering system, figure and table formatting, and bibliography settings
unless otherwise instructed by the Faculty or the supervisor.

The report shall be written in English.

Students are responsible for the academic and technical content of the
report. The official \LaTeX{} template is responsible for standardized
document formatting.

\vspace{0.3cm}

% ================================================================
% II. REPORT STRUCTURE
% ================================================================

\noindent
\textbf{II.\quad REPORT STRUCTURE}

\vspace{0.15cm}

The report shall normally be arranged in the following order:

\begin{enumerate}[label=\arabic*., leftmargin=1.2cm, itemsep=2pt]
    \item Outer Cover;
    \item Comments of the Internship Host Organization;
    \item Supervisor's Evaluation and Graduation Internship Defense Results;
    \item Table of Contents;
    \item List of Symbols and Abbreviations, if applicable;
    \item List of Figures;
    \item List of Tables;
    \item Abstract;
    \item Main Chapters of the Internship Report;
    \item References;
    \item Appendices, if applicable.
\end{enumerate}

The preliminary pages shall not be numbered as chapters. The main report
content shall be organized into consecutively numbered chapters beginning
with Chapter 1.

\vspace{0.3cm}

% ================================================================
% III. PAGE SIZE AND MARGINS
% ================================================================

\noindent
\textbf{III.\quad PAGE SIZE AND MARGINS}

\vspace{0.15cm}

The report shall be formatted on A4 paper with the following margins:

\begin{center}
\begin{tabular}{ll}
    Left margin:   & 3.5 cm \\
    Right margin:  & 2.0 cm \\
    Top margin:    & 2.5 cm \\
    Bottom margin: & 2.5 cm
\end{tabular}
\end{center}

These settings are predefined in the official template and shall not be
modified by students.

\vspace{0.3cm}

% ================================================================
% IV. FONT, FONT SIZE, AND LINE SPACING
% ================================================================

\noindent
\textbf{IV.\quad FONT, FONT SIZE, AND LINE SPACING}

\vspace{0.15cm}

The main text shall use a 13 pt font size with 1.3-line spacing.
A Times-family typeface shall be used in accordance with the official
template.

The text shall be fully justified.

Students shall not use repeated spaces, repeated manual line breaks, or
multiple blank lines to position text. Paragraph and heading spacing shall
be controlled by the template.

\vspace{0.3cm}

% ================================================================
% V. OUTER COVER
% ================================================================

\noindent
\textbf{V.\quad OUTER COVER}

\vspace{0.15cm}

The outer cover shall follow the official format provided in the template
and shall contain the following information:

\begin{itemize}[leftmargin=0.8cm, itemsep=2pt]
    \item Thai Nguyen University of Technology;
    \item Faculty of International Training;
    \item Engineering and Technology Department;
    \item University logo and Faculty logo;
    \item the title \textbf{INTERNSHIP REPORT};
    \item Internship Organization;
    \item Author;
    \item Class;
    \item Supervisor;
    \item place and year of completion.
\end{itemize}

Students shall only replace the required information fields and shall not
modify the cover-page layout unless instructed otherwise.

\vspace{0.3cm}

% ================================================================
% VI. COMMENTS OF THE INTERNSHIP HOST ORGANIZATION
% ================================================================

\noindent
\textbf{VI.\quad COMMENTS OF THE INTERNSHIP HOST ORGANIZATION}

\vspace{0.15cm}

The page entitled \textbf{COMMENTS OF THE INTERNSHIP HOST ORGANIZATION}
is reserved for the internship host organization.

It includes:

\begin{itemize}[leftmargin=0.8cm, itemsep=2pt]
    \item the student's full name;
    \item class;
    \item comments from the host organization;
    \item score;
    \item date;
    \item signature and official seal.
\end{itemize}

Students shall not alter or complete official evaluation fields unless
specifically authorized.

\vspace{0.3cm}

% ================================================================
% VII. SUPERVISOR'S EVALUATION AND DEFENSE RESULTS
% ================================================================

\noindent
\textbf{VII.\quad SUPERVISOR'S EVALUATION AND DEFENSE RESULTS}

\vspace{0.15cm}

The section entitled
\textbf{SUPERVISOR'S EVALUATION AND GRADUATION INTERNSHIP DEFENSE RESULTS}
is reserved for the supervisor and defense committee.

The evaluation may include:

\begin{itemize}[leftmargin=0.8cm, itemsep=2pt]
    \item Internship Progress and Student's Attitude;
    \item Report Content;
    \item Report Presentation;
    \item Other Comments;
    \item Supervisor's Opinion;
    \item Defense Result;
    \item Score.
\end{itemize}

Students shall not enter evaluation scores or defense results before they
are officially determined.

\vspace{0.3cm}

% ================================================================
% VIII. CHAPTERS, SECTIONS, AND SUBSECTIONS
% ================================================================

\noindent
\textbf{VIII.\quad CHAPTERS, SECTIONS, AND SUBSECTIONS}

\vspace{0.15cm}

The report structure shall be generated using the standard \LaTeX{}
commands:

\begin{verbatim}
\chapter{...}
\section{...}
\subsection{...}
\end{verbatim}

Chapter, section, and subsection numbers shall be generated automatically
by \LaTeX{}.

Students shall not manually type chapter or section numbers and shall not
reproduce heading styles using commands such as \verb|\textbf|,
\verb|\fontsize|, or \verb|\centering| instead of the appropriate
structural command.

Excessive subdivision should be avoided. In most cases, the hierarchy
Chapter--Section--Subsection is sufficient.

\vspace{0.3cm}

% ================================================================
% IX. ORGANIZATION OF THE MAIN CONTENT
% ================================================================

\noindent
\textbf{IX.\quad ORGANIZATION OF THE MAIN CONTENT}

\vspace{0.15cm}

The specific content of the report shall depend on the internship
organization and the assigned tasks. A typical structure may include:

\begin{itemize}[leftmargin=0.8cm, itemsep=2pt]
    \item \textbf{Chapter 1. Overview of the Internship Organization};
    \item \textbf{Chapter 2. Internship Activities and Assigned Tasks};
    \item \textbf{Chapter 3. Results, Experience, and Evaluation}.
\end{itemize}

These chapter titles are recommendations only and may be adapted to the
actual internship content.

The report shall reflect work actually performed, observed, analyzed, or
learned during the internship. Unrelated theoretical material should not
be included merely to increase report length.

\vspace{0.3cm}

% ================================================================
% X. TABLE OF CONTENTS
% ================================================================

\noindent
\textbf{X.\quad TABLE OF CONTENTS}

\vspace{0.15cm}

The Table of Contents shall be generated automatically using:

\begin{verbatim}
\tableofcontents
\end{verbatim}

Students shall not create or manually edit page numbers in the Table of
Contents.

The document shall be compiled sufficiently to ensure that all page
numbers and entries are correctly updated.

\vspace{0.3cm}

% ================================================================
% XI. LIST OF SYMBOLS AND ABBREVIATIONS
% ================================================================

\noindent
\textbf{XI.\quad LIST OF SYMBOLS AND ABBREVIATIONS}

\vspace{0.15cm}

Symbols and abbreviations that are used repeatedly in the report shall be
listed in the \textbf{List of Symbols and Abbreviations}.

The list may contain the following columns:

\begin{center}
\begin{tabular}{llll}
    No. & Symbol / Abbreviation & Unit & Definition
\end{tabular}
\end{center}

Only symbols and abbreviations actually used in the final report shall be
retained. Example entries contained in the template shall be removed if
they are not relevant to the student's report.

Each abbreviation should normally be defined when it first appears in
the main text.

\vspace{0.3cm}

% ================================================================
% XII. FIGURES, DIAGRAMS, AND PHOTOGRAPHS
% ================================================================

\noindent
\textbf{XII.\quad FIGURES, DIAGRAMS, AND PHOTOGRAPHS}

\vspace{0.15cm}

Figures shall be clear, legible, and directly relevant to the report
content.

Figures shall be inserted using the standard \verb|figure| environment.

Example:

\begin{verbatim}
\begin{figure}[H]
    \centering
    \includegraphics[width=0.8\textwidth]
                    {figures/company_structure.pdf}
    \caption{Organizational structure of the company}
    \label{fig:company-structure}
\end{figure}
\end{verbatim}

Figure captions shall be placed below the figure.

Each referenced figure shall have a unique \verb|\label| and shall be
referenced in the text using, for example:

\begin{verbatim}
Figure~\ref{fig:company-structure}
\end{verbatim}

Students shall not manually type figure numbers.

Figures reproduced or adapted from external sources shall include an
appropriate citation.

Photographs taken during the internship shall be relevant and sufficiently
clear. Students shall not include confidential, restricted, or proprietary
information without permission from the internship organization.

Vector graphics or high-resolution images are recommended whenever
possible.

\vspace{0.3cm}

% ================================================================
% XIII. TABLES
% ================================================================

\noindent
\textbf{XIII.\quad TABLES}

\vspace{0.15cm}

Tables shall be created using suitable \LaTeX{} environments such as
\verb|table|, \verb|tabular|, \verb|tabularx|, or \verb|longtable|.

Table captions shall be placed above the table.

Example:

\begin{verbatim}
\begin{table}[H]
    \centering
    \caption{Internship schedule}
    \label{tab:internship-schedule}

    \begin{tabular}{lll}
        \toprule
        Week & Task & Result \\
        \midrule
        1 & Orientation & Completed \\
        2 & System study & Completed \\
        \bottomrule
    \end{tabular}
\end{table}
\end{verbatim}

Tables shall be referenced using:

\begin{verbatim}
Table~\ref{tab:internship-schedule}
\end{verbatim}

Students shall not manually type table numbers.

Tables should be created directly in \LaTeX{} rather than inserted as
screenshots whenever practical.

\vspace{0.3cm}

% ================================================================
% XIV. EQUATIONS AND TECHNICAL QUANTITIES
% ================================================================

\noindent
\textbf{XIV.\quad EQUATIONS AND TECHNICAL QUANTITIES}

\vspace{0.15cm}

Mathematical expressions shall be typeset directly using \LaTeX{}.
Images or screenshots of equations shall not be used.

Short expressions may be written in inline mathematical mode, for example:

\begin{verbatim}
$R = 0.7\,\Omega$
\end{verbatim}

Important equations shall use a numbered equation environment.

Example:

\begin{verbatim}
\begin{equation}
    P = UI
    \label{eq:power}
\end{equation}
\end{verbatim}

Equations shall be referenced using:

\begin{verbatim}
Equation~\eqref{eq:power}
\end{verbatim}

Equation numbers shall not be manually entered.

SI units shall be used wherever applicable. The \verb|siunitx| package
should be used where appropriate.

Examples:

\begin{verbatim}
\SI{220}{\volt}
\SI{50}{\hertz}
\SI{2.5}{\ampere}
\end{verbatim}

Notation and units shall be consistent throughout the report.

\vspace{0.3cm}

% ================================================================
% XV. SOURCE CODE
% ================================================================

\noindent
\textbf{XV.\quad SOURCE CODE}

\vspace{0.15cm}

If the internship involves programming, PLCs, microcontrollers,
simulation, or software development, relevant source code may be included
in the report.

Source code shall be presented using the \verb|listings| environment or
another environment defined by the official template.

Example:

\begin{verbatim}
\begin{lstlisting}[language=st]
IF Start = TRUE THEN
    Motor := TRUE;
ELSE
    Motor := FALSE;
END_IF;
\end{lstlisting}
\end{verbatim}

Only code necessary for explaining the work should appear in the main
report. Long source-code listings should be moved to an appendix or
submitted as supplementary files if required.

Source code should not normally be inserted as screenshots.

\vspace{0.3cm}

% ================================================================
% XVI. ABSTRACT
% ================================================================

\noindent
\textbf{XVI.\quad ABSTRACT}

\vspace{0.15cm}

The Abstract shall briefly summarize:

\begin{itemize}[leftmargin=0.8cm, itemsep=2pt]
    \item the internship organization;
    \item the objectives of the internship;
    \item the main tasks performed;
    \item the main results, knowledge, and skills obtained.
\end{itemize}

The Abstract should be concise and should not reproduce the full content
of the report chapter by chapter.

\vspace{0.3cm}

% ================================================================
% XVII. REFERENCES AND CITATIONS
% ================================================================

\noindent
\textbf{XVII.\quad REFERENCES AND CITATIONS}

\vspace{0.15cm}

All books, journal papers, conference papers, standards, websites,
datasheets, manuals, technical documents, figures, tables, or other
external sources used in preparing the report shall be properly cited.

References shall be managed through the bibliography system provided with
the official template.

Bibliographic entries shall normally be stored in:

\begin{verbatim}
References.bib
\end{verbatim}

References shall be cited in the report using:

\begin{verbatim}
\cite{reference-key}
\end{verbatim}

Students shall not manually type reference numbers such as ``[1]'',
``[2]'', etc.

Titles of foreign-language publications shall remain in their original
language in the bibliography.

All borrowed ideas, data, figures, tables, equations, methods, or results
shall be properly acknowledged.

\vspace{0.3cm}

% ================================================================
% XVIII. CROSS-REFERENCING
% ================================================================

\noindent
\textbf{XVIII.\quad CROSS-REFERENCING}

\vspace{0.15cm}

Numbered elements shall be cross-referenced using the standard \LaTeX{}
mechanism.

Use the following commands where appropriate:

\begin{verbatim}
\label{...}
\ref{...}
\eqref{...}
\end{verbatim}

For example:

\begin{verbatim}
As shown in Figure~\ref{fig:system}, ...
\end{verbatim}

Students shall not manually type numbers such as ``Figure 2.3'' when
those numbers can be generated automatically.

\vspace{0.3cm}

% ================================================================
% XIX. PAGE NUMBERING
% ================================================================

\noindent
\textbf{XIX.\quad PAGE NUMBERING}

\vspace{0.15cm}

Cover pages and official evaluation forms at the beginning of the report
shall not display page numbers.

The main body of the report shall use Arabic page numbers beginning from
Chapter 1.

Page numbering shall be generated automatically by the official template.
Students shall not manually enter page numbers.

\vspace{0.3cm}

% ================================================================
% XX. PROJECT FILE ORGANIZATION
% ================================================================

\noindent
\textbf{XX.\quad PROJECT FILE ORGANIZATION}

\vspace{0.15cm}

Students should maintain a clear and organized \LaTeX{} project
structure. A recommended structure is:

\begin{verbatim}
InternshipReport/
|
|-- main.tex
|-- References.bib
|
|-- figures/
|   |-- company.pdf
|   |-- system.png
|   `-- ...
|
`-- chapters/
    |-- chapter1.tex
    |-- chapter2.tex
    `-- chapter3.tex
\end{verbatim}

Figures should normally be stored in the \verb|figures| directory
specified by the official template.

File and directory names should preferably contain Latin letters,
numbers, hyphens, or underscores.

Recommended examples include:

\begin{verbatim}
company_structure.pdf
motor_control.png
chapter_01.tex
\end{verbatim}

Spaces and unnecessary special characters in file names should be
avoided where possible.

\vspace{0.3cm}

% ================================================================
% XXI. DOCUMENT COMPILATION AND FINAL CHECK
% ================================================================

\noindent
\textbf{XXI.\quad DOCUMENT COMPILATION AND FINAL CHECK}

\vspace{0.15cm}

Before submission, the complete \LaTeX{} project shall compile
successfully and produce the final PDF report.

The final document shall not contain unresolved references such as:

\begin{verbatim}
??
[?]
\end{verbatim}

Before generating the final PDF, students shall check:

\begin{itemize}[leftmargin=0.8cm, itemsep=2pt]
    \item the Table of Contents;
    \item page numbering;
    \item chapter, section, and subsection numbering;
    \item figure numbering;
    \item table numbering;
    \item equation numbering;
    \item cross-references;
    \item citations;
    \item bibliography;
    \item font display;
    \item mathematical symbols;
    \item image quality;
    \item completeness of all required forms.
\end{itemize}

The document may be compiled using TeX Live, MiKTeX, or another
compatible \LaTeX{} environment. Regardless of the compilation
environment, the resulting PDF shall conform to the official Faculty
template.

\vspace{0.3cm}

% ================================================================
% XXII. CONTENT THAT STUDENTS MAY EDIT
% ================================================================

\noindent
\textbf{XXII.\quad CONTENT THAT STUDENTS MAY EDIT}

\vspace{0.15cm}

Students may edit or replace:

\begin{itemize}[leftmargin=0.8cm, itemsep=2pt]
    \item student name;
    \item class;
    \item internship organization;
    \item supervisor name;
    \item Abstract;
    \item List of Symbols and Abbreviations;
    \item chapter titles and chapter content;
    \item figures and photographs;
    \item tables;
    \item equations;
    \item citations and references.
\end{itemize}

Students shall not modify the main document formatting settings in the
preamble, including the document class, page geometry, line spacing,
heading formatting, table-of-contents formatting, font configuration,
page layout, or cover-page macros unless specifically instructed.

\vspace{0.3cm}

% ================================================================
% XXIII. FINAL SUBMISSION REQUIREMENTS
% ================================================================

\noindent
\textbf{XXIII.\quad FINAL SUBMISSION REQUIREMENTS}

\vspace{0.15cm}

Before submission, the final report shall:

\begin{itemize}[leftmargin=0.8cm, itemsep=2pt]
    \item compile without errors that affect the report content or presentation;
    \item contain all required forms and report sections;
    \item contain correct student, class, internship organization, and supervisor information;
    \item contain an updated Table of Contents, List of Figures, and List of Tables;
    \item contain no unused example content from the template;
    \item contain no placeholder text such as \textit{Title of Chapter 1},
    \textit{Section 1.1}, or other sample entries;
    \item contain all required citations and references;
    \item contain no unresolved cross-references.
\end{itemize}

If the Faculty requires submission of the \LaTeX{} source, the complete
project shall include all \verb|.tex| files, the bibliography database,
figures, and any other files required to reproduce the submitted PDF.

\vspace{0.4cm}

\noindent
\textbf{Important:}
Before submitting the final internship report, students shall remove or
comment out the following command in the main document:

\begin{verbatim}
\input{Regulations.tex}
\end{verbatim}

For example:

\begin{verbatim}
% \input{Regulations.tex}
\end{verbatim}

This removes the guidance pages from the final submitted report while
preserving the regulations file in the \LaTeX{} project for reference.

} % End 13pt content

\clearpage

\end{verbatim}

For example:

\begin{verbatim}
% 
% =================================================================
% REGULATIONS FOR PREPARING THE GRADUATION INTERNSHIP REPORT
% USING LATEX
% =================================================================
%
% This file is intended to be included from the main document with:
%
%     \input{Regulations.tex}
%
% IMPORTANT:
% This section is provided for guidance only. Students shall remove
% or comment out the \input{Regulations.tex} command before submitting
% the final internship report.
%
% =================================================================

\clearpage
\thispagestyle{empty}

\begin{center}
    {\fontsize{16pt}{20pt}\selectfont\bfseries
    SPECIFIC REGULATIONS FOR PREPARING\\
    GRADUATION INTERNSHIP REPORTS USING \LaTeX}
\end{center}

\vspace{0.5cm}

{\fontsize{13pt}{17pt}\selectfont
\setstretch{1.3}
\justifying

\noindent
\textbf{Note:}
This section is provided for guidance only and is included in the main
document using the command \verb|\input{Regulations.tex}|.
Students shall remove or comment out this command before submitting the
final internship report.

\vspace{0.4cm}

% ================================================================
% I. GENERAL REQUIREMENTS
% ================================================================

\noindent
\textbf{I.\quad GENERAL REQUIREMENTS}

\vspace{0.15cm}

The Graduation Internship Report shall be prepared using the official
\LaTeX{} template provided by the Faculty. Students shall preserve the
predefined page layout, typography, line spacing, heading styles,
numbering system, figure and table formatting, and bibliography settings
unless otherwise instructed by the Faculty or the supervisor.

The report shall be written in English.

Students are responsible for the academic and technical content of the
report. The official \LaTeX{} template is responsible for standardized
document formatting.

\vspace{0.3cm}

% ================================================================
% II. REPORT STRUCTURE
% ================================================================

\noindent
\textbf{II.\quad REPORT STRUCTURE}

\vspace{0.15cm}

The report shall normally be arranged in the following order:

\begin{enumerate}[label=\arabic*., leftmargin=1.2cm, itemsep=2pt]
    \item Outer Cover;
    \item Comments of the Internship Host Organization;
    \item Supervisor's Evaluation and Graduation Internship Defense Results;
    \item Table of Contents;
    \item List of Symbols and Abbreviations, if applicable;
    \item List of Figures;
    \item List of Tables;
    \item Abstract;
    \item Main Chapters of the Internship Report;
    \item References;
    \item Appendices, if applicable.
\end{enumerate}

The preliminary pages shall not be numbered as chapters. The main report
content shall be organized into consecutively numbered chapters beginning
with Chapter 1.

\vspace{0.3cm}

% ================================================================
% III. PAGE SIZE AND MARGINS
% ================================================================

\noindent
\textbf{III.\quad PAGE SIZE AND MARGINS}

\vspace{0.15cm}

The report shall be formatted on A4 paper with the following margins:

\begin{center}
\begin{tabular}{ll}
    Left margin:   & 3.5 cm \\
    Right margin:  & 2.0 cm \\
    Top margin:    & 2.5 cm \\
    Bottom margin: & 2.5 cm
\end{tabular}
\end{center}

These settings are predefined in the official template and shall not be
modified by students.

\vspace{0.3cm}

% ================================================================
% IV. FONT, FONT SIZE, AND LINE SPACING
% ================================================================

\noindent
\textbf{IV.\quad FONT, FONT SIZE, AND LINE SPACING}

\vspace{0.15cm}

The main text shall use a 13 pt font size with 1.3-line spacing.
A Times-family typeface shall be used in accordance with the official
template.

The text shall be fully justified.

Students shall not use repeated spaces, repeated manual line breaks, or
multiple blank lines to position text. Paragraph and heading spacing shall
be controlled by the template.

\vspace{0.3cm}

% ================================================================
% V. OUTER COVER
% ================================================================

\noindent
\textbf{V.\quad OUTER COVER}

\vspace{0.15cm}

The outer cover shall follow the official format provided in the template
and shall contain the following information:

\begin{itemize}[leftmargin=0.8cm, itemsep=2pt]
    \item Thai Nguyen University of Technology;
    \item Faculty of International Training;
    \item Engineering and Technology Department;
    \item University logo and Faculty logo;
    \item the title \textbf{INTERNSHIP REPORT};
    \item Internship Organization;
    \item Author;
    \item Class;
    \item Supervisor;
    \item place and year of completion.
\end{itemize}

Students shall only replace the required information fields and shall not
modify the cover-page layout unless instructed otherwise.

\vspace{0.3cm}

% ================================================================
% VI. COMMENTS OF THE INTERNSHIP HOST ORGANIZATION
% ================================================================

\noindent
\textbf{VI.\quad COMMENTS OF THE INTERNSHIP HOST ORGANIZATION}

\vspace{0.15cm}

The page entitled \textbf{COMMENTS OF THE INTERNSHIP HOST ORGANIZATION}
is reserved for the internship host organization.

It includes:

\begin{itemize}[leftmargin=0.8cm, itemsep=2pt]
    \item the student's full name;
    \item class;
    \item comments from the host organization;
    \item score;
    \item date;
    \item signature and official seal.
\end{itemize}

Students shall not alter or complete official evaluation fields unless
specifically authorized.

\vspace{0.3cm}

% ================================================================
% VII. SUPERVISOR'S EVALUATION AND DEFENSE RESULTS
% ================================================================

\noindent
\textbf{VII.\quad SUPERVISOR'S EVALUATION AND DEFENSE RESULTS}

\vspace{0.15cm}

The section entitled
\textbf{SUPERVISOR'S EVALUATION AND GRADUATION INTERNSHIP DEFENSE RESULTS}
is reserved for the supervisor and defense committee.

The evaluation may include:

\begin{itemize}[leftmargin=0.8cm, itemsep=2pt]
    \item Internship Progress and Student's Attitude;
    \item Report Content;
    \item Report Presentation;
    \item Other Comments;
    \item Supervisor's Opinion;
    \item Defense Result;
    \item Score.
\end{itemize}

Students shall not enter evaluation scores or defense results before they
are officially determined.

\vspace{0.3cm}

% ================================================================
% VIII. CHAPTERS, SECTIONS, AND SUBSECTIONS
% ================================================================

\noindent
\textbf{VIII.\quad CHAPTERS, SECTIONS, AND SUBSECTIONS}

\vspace{0.15cm}

The report structure shall be generated using the standard \LaTeX{}
commands:

\begin{verbatim}
\chapter{...}
\section{...}
\subsection{...}
\end{verbatim}

Chapter, section, and subsection numbers shall be generated automatically
by \LaTeX{}.

Students shall not manually type chapter or section numbers and shall not
reproduce heading styles using commands such as \verb|\textbf|,
\verb|\fontsize|, or \verb|\centering| instead of the appropriate
structural command.

Excessive subdivision should be avoided. In most cases, the hierarchy
Chapter--Section--Subsection is sufficient.

\vspace{0.3cm}

% ================================================================
% IX. ORGANIZATION OF THE MAIN CONTENT
% ================================================================

\noindent
\textbf{IX.\quad ORGANIZATION OF THE MAIN CONTENT}

\vspace{0.15cm}

The specific content of the report shall depend on the internship
organization and the assigned tasks. A typical structure may include:

\begin{itemize}[leftmargin=0.8cm, itemsep=2pt]
    \item \textbf{Chapter 1. Overview of the Internship Organization};
    \item \textbf{Chapter 2. Internship Activities and Assigned Tasks};
    \item \textbf{Chapter 3. Results, Experience, and Evaluation}.
\end{itemize}

These chapter titles are recommendations only and may be adapted to the
actual internship content.

The report shall reflect work actually performed, observed, analyzed, or
learned during the internship. Unrelated theoretical material should not
be included merely to increase report length.

\vspace{0.3cm}

% ================================================================
% X. TABLE OF CONTENTS
% ================================================================

\noindent
\textbf{X.\quad TABLE OF CONTENTS}

\vspace{0.15cm}

The Table of Contents shall be generated automatically using:

\begin{verbatim}
\tableofcontents
\end{verbatim}

Students shall not create or manually edit page numbers in the Table of
Contents.

The document shall be compiled sufficiently to ensure that all page
numbers and entries are correctly updated.

\vspace{0.3cm}

% ================================================================
% XI. LIST OF SYMBOLS AND ABBREVIATIONS
% ================================================================

\noindent
\textbf{XI.\quad LIST OF SYMBOLS AND ABBREVIATIONS}

\vspace{0.15cm}

Symbols and abbreviations that are used repeatedly in the report shall be
listed in the \textbf{List of Symbols and Abbreviations}.

The list may contain the following columns:

\begin{center}
\begin{tabular}{llll}
    No. & Symbol / Abbreviation & Unit & Definition
\end{tabular}
\end{center}

Only symbols and abbreviations actually used in the final report shall be
retained. Example entries contained in the template shall be removed if
they are not relevant to the student's report.

Each abbreviation should normally be defined when it first appears in
the main text.

\vspace{0.3cm}

% ================================================================
% XII. FIGURES, DIAGRAMS, AND PHOTOGRAPHS
% ================================================================

\noindent
\textbf{XII.\quad FIGURES, DIAGRAMS, AND PHOTOGRAPHS}

\vspace{0.15cm}

Figures shall be clear, legible, and directly relevant to the report
content.

Figures shall be inserted using the standard \verb|figure| environment.

Example:

\begin{verbatim}
\begin{figure}[H]
    \centering
    \includegraphics[width=0.8\textwidth]
                    {figures/company_structure.pdf}
    \caption{Organizational structure of the company}
    \label{fig:company-structure}
\end{figure}
\end{verbatim}

Figure captions shall be placed below the figure.

Each referenced figure shall have a unique \verb|\label| and shall be
referenced in the text using, for example:

\begin{verbatim}
Figure~\ref{fig:company-structure}
\end{verbatim}

Students shall not manually type figure numbers.

Figures reproduced or adapted from external sources shall include an
appropriate citation.

Photographs taken during the internship shall be relevant and sufficiently
clear. Students shall not include confidential, restricted, or proprietary
information without permission from the internship organization.

Vector graphics or high-resolution images are recommended whenever
possible.

\vspace{0.3cm}

% ================================================================
% XIII. TABLES
% ================================================================

\noindent
\textbf{XIII.\quad TABLES}

\vspace{0.15cm}

Tables shall be created using suitable \LaTeX{} environments such as
\verb|table|, \verb|tabular|, \verb|tabularx|, or \verb|longtable|.

Table captions shall be placed above the table.

Example:

\begin{verbatim}
\begin{table}[H]
    \centering
    \caption{Internship schedule}
    \label{tab:internship-schedule}

    \begin{tabular}{lll}
        \toprule
        Week & Task & Result \\
        \midrule
        1 & Orientation & Completed \\
        2 & System study & Completed \\
        \bottomrule
    \end{tabular}
\end{table}
\end{verbatim}

Tables shall be referenced using:

\begin{verbatim}
Table~\ref{tab:internship-schedule}
\end{verbatim}

Students shall not manually type table numbers.

Tables should be created directly in \LaTeX{} rather than inserted as
screenshots whenever practical.

\vspace{0.3cm}

% ================================================================
% XIV. EQUATIONS AND TECHNICAL QUANTITIES
% ================================================================

\noindent
\textbf{XIV.\quad EQUATIONS AND TECHNICAL QUANTITIES}

\vspace{0.15cm}

Mathematical expressions shall be typeset directly using \LaTeX{}.
Images or screenshots of equations shall not be used.

Short expressions may be written in inline mathematical mode, for example:

\begin{verbatim}
$R = 0.7\,\Omega$
\end{verbatim}

Important equations shall use a numbered equation environment.

Example:

\begin{verbatim}
\begin{equation}
    P = UI
    \label{eq:power}
\end{equation}
\end{verbatim}

Equations shall be referenced using:

\begin{verbatim}
Equation~\eqref{eq:power}
\end{verbatim}

Equation numbers shall not be manually entered.

SI units shall be used wherever applicable. The \verb|siunitx| package
should be used where appropriate.

Examples:

\begin{verbatim}
\SI{220}{\volt}
\SI{50}{\hertz}
\SI{2.5}{\ampere}
\end{verbatim}

Notation and units shall be consistent throughout the report.

\vspace{0.3cm}

% ================================================================
% XV. SOURCE CODE
% ================================================================

\noindent
\textbf{XV.\quad SOURCE CODE}

\vspace{0.15cm}

If the internship involves programming, PLCs, microcontrollers,
simulation, or software development, relevant source code may be included
in the report.

Source code shall be presented using the \verb|listings| environment or
another environment defined by the official template.

Example:

\begin{verbatim}
\begin{lstlisting}[language=st]
IF Start = TRUE THEN
    Motor := TRUE;
ELSE
    Motor := FALSE;
END_IF;
\end{lstlisting}
\end{verbatim}

Only code necessary for explaining the work should appear in the main
report. Long source-code listings should be moved to an appendix or
submitted as supplementary files if required.

Source code should not normally be inserted as screenshots.

\vspace{0.3cm}

% ================================================================
% XVI. ABSTRACT
% ================================================================

\noindent
\textbf{XVI.\quad ABSTRACT}

\vspace{0.15cm}

The Abstract shall briefly summarize:

\begin{itemize}[leftmargin=0.8cm, itemsep=2pt]
    \item the internship organization;
    \item the objectives of the internship;
    \item the main tasks performed;
    \item the main results, knowledge, and skills obtained.
\end{itemize}

The Abstract should be concise and should not reproduce the full content
of the report chapter by chapter.

\vspace{0.3cm}

% ================================================================
% XVII. REFERENCES AND CITATIONS
% ================================================================

\noindent
\textbf{XVII.\quad REFERENCES AND CITATIONS}

\vspace{0.15cm}

All books, journal papers, conference papers, standards, websites,
datasheets, manuals, technical documents, figures, tables, or other
external sources used in preparing the report shall be properly cited.

References shall be managed through the bibliography system provided with
the official template.

Bibliographic entries shall normally be stored in:

\begin{verbatim}
References.bib
\end{verbatim}

References shall be cited in the report using:

\begin{verbatim}
\cite{reference-key}
\end{verbatim}

Students shall not manually type reference numbers such as ``[1]'',
``[2]'', etc.

Titles of foreign-language publications shall remain in their original
language in the bibliography.

All borrowed ideas, data, figures, tables, equations, methods, or results
shall be properly acknowledged.

\vspace{0.3cm}

% ================================================================
% XVIII. CROSS-REFERENCING
% ================================================================

\noindent
\textbf{XVIII.\quad CROSS-REFERENCING}

\vspace{0.15cm}

Numbered elements shall be cross-referenced using the standard \LaTeX{}
mechanism.

Use the following commands where appropriate:

\begin{verbatim}
\label{...}
\ref{...}
\eqref{...}
\end{verbatim}

For example:

\begin{verbatim}
As shown in Figure~\ref{fig:system}, ...
\end{verbatim}

Students shall not manually type numbers such as ``Figure 2.3'' when
those numbers can be generated automatically.

\vspace{0.3cm}

% ================================================================
% XIX. PAGE NUMBERING
% ================================================================

\noindent
\textbf{XIX.\quad PAGE NUMBERING}

\vspace{0.15cm}

Cover pages and official evaluation forms at the beginning of the report
shall not display page numbers.

The main body of the report shall use Arabic page numbers beginning from
Chapter 1.

Page numbering shall be generated automatically by the official template.
Students shall not manually enter page numbers.

\vspace{0.3cm}

% ================================================================
% XX. PROJECT FILE ORGANIZATION
% ================================================================

\noindent
\textbf{XX.\quad PROJECT FILE ORGANIZATION}

\vspace{0.15cm}

Students should maintain a clear and organized \LaTeX{} project
structure. A recommended structure is:

\begin{verbatim}
InternshipReport/
|
|-- main.tex
|-- References.bib
|
|-- figures/
|   |-- company.pdf
|   |-- system.png
|   `-- ...
|
`-- chapters/
    |-- chapter1.tex
    |-- chapter2.tex
    `-- chapter3.tex
\end{verbatim}

Figures should normally be stored in the \verb|figures| directory
specified by the official template.

File and directory names should preferably contain Latin letters,
numbers, hyphens, or underscores.

Recommended examples include:

\begin{verbatim}
company_structure.pdf
motor_control.png
chapter_01.tex
\end{verbatim}

Spaces and unnecessary special characters in file names should be
avoided where possible.

\vspace{0.3cm}

% ================================================================
% XXI. DOCUMENT COMPILATION AND FINAL CHECK
% ================================================================

\noindent
\textbf{XXI.\quad DOCUMENT COMPILATION AND FINAL CHECK}

\vspace{0.15cm}

Before submission, the complete \LaTeX{} project shall compile
successfully and produce the final PDF report.

The final document shall not contain unresolved references such as:

\begin{verbatim}
??
[?]
\end{verbatim}

Before generating the final PDF, students shall check:

\begin{itemize}[leftmargin=0.8cm, itemsep=2pt]
    \item the Table of Contents;
    \item page numbering;
    \item chapter, section, and subsection numbering;
    \item figure numbering;
    \item table numbering;
    \item equation numbering;
    \item cross-references;
    \item citations;
    \item bibliography;
    \item font display;
    \item mathematical symbols;
    \item image quality;
    \item completeness of all required forms.
\end{itemize}

The document may be compiled using TeX Live, MiKTeX, or another
compatible \LaTeX{} environment. Regardless of the compilation
environment, the resulting PDF shall conform to the official Faculty
template.

\vspace{0.3cm}

% ================================================================
% XXII. CONTENT THAT STUDENTS MAY EDIT
% ================================================================

\noindent
\textbf{XXII.\quad CONTENT THAT STUDENTS MAY EDIT}

\vspace{0.15cm}

Students may edit or replace:

\begin{itemize}[leftmargin=0.8cm, itemsep=2pt]
    \item student name;
    \item class;
    \item internship organization;
    \item supervisor name;
    \item Abstract;
    \item List of Symbols and Abbreviations;
    \item chapter titles and chapter content;
    \item figures and photographs;
    \item tables;
    \item equations;
    \item citations and references.
\end{itemize}

Students shall not modify the main document formatting settings in the
preamble, including the document class, page geometry, line spacing,
heading formatting, table-of-contents formatting, font configuration,
page layout, or cover-page macros unless specifically instructed.

\vspace{0.3cm}

% ================================================================
% XXIII. FINAL SUBMISSION REQUIREMENTS
% ================================================================

\noindent
\textbf{XXIII.\quad FINAL SUBMISSION REQUIREMENTS}

\vspace{0.15cm}

Before submission, the final report shall:

\begin{itemize}[leftmargin=0.8cm, itemsep=2pt]
    \item compile without errors that affect the report content or presentation;
    \item contain all required forms and report sections;
    \item contain correct student, class, internship organization, and supervisor information;
    \item contain an updated Table of Contents, List of Figures, and List of Tables;
    \item contain no unused example content from the template;
    \item contain no placeholder text such as \textit{Title of Chapter 1},
    \textit{Section 1.1}, or other sample entries;
    \item contain all required citations and references;
    \item contain no unresolved cross-references.
\end{itemize}

If the Faculty requires submission of the \LaTeX{} source, the complete
project shall include all \verb|.tex| files, the bibliography database,
figures, and any other files required to reproduce the submitted PDF.

\vspace{0.4cm}

\noindent
\textbf{Important:}
Before submitting the final internship report, students shall remove or
comment out the following command in the main document:

\begin{verbatim}
\input{Regulations.tex}
\end{verbatim}

For example:

\begin{verbatim}
% \input{Regulations.tex}
\end{verbatim}

This removes the guidance pages from the final submitted report while
preserving the regulations file in the \LaTeX{} project for reference.

} % End 13pt content

\clearpage

\end{verbatim}

This removes the guidance pages from the final submitted report while
preserving the regulations file in the \LaTeX{} project for reference.

} % End 13pt content

\clearpage

\end{verbatim}

This removes the guidance pages from the final submitted report while
preserving the regulations file in the \LaTeX{} project for reference.

} % End 13pt content

\clearpage

%
% IMPORTANT:
% This section is provided for guidance only. Students shall remove
% or comment out the 
% =================================================================
% REGULATIONS FOR PREPARING THE GRADUATION INTERNSHIP REPORT
% USING LATEX
% =================================================================
%
% This file is intended to be included from the main document with:
%
%     
% =================================================================
% REGULATIONS FOR PREPARING THE GRADUATION INTERNSHIP REPORT
% USING LATEX
% =================================================================
%
% This file is intended to be included from the main document with:
%
%     
% =================================================================
% REGULATIONS FOR PREPARING THE GRADUATION INTERNSHIP REPORT
% USING LATEX
% =================================================================
%
% This file is intended to be included from the main document with:
%
%     \input{Regulations.tex}
%
% IMPORTANT:
% This section is provided for guidance only. Students shall remove
% or comment out the \input{Regulations.tex} command before submitting
% the final internship report.
%
% =================================================================

\clearpage
\thispagestyle{empty}

\begin{center}
    {\fontsize{16pt}{20pt}\selectfont\bfseries
    SPECIFIC REGULATIONS FOR PREPARING\\
    GRADUATION INTERNSHIP REPORTS USING \LaTeX}
\end{center}

\vspace{0.5cm}

{\fontsize{13pt}{17pt}\selectfont
\setstretch{1.3}
\justifying

\noindent
\textbf{Note:}
This section is provided for guidance only and is included in the main
document using the command \verb|\input{Regulations.tex}|.
Students shall remove or comment out this command before submitting the
final internship report.

\vspace{0.4cm}

% ================================================================
% I. GENERAL REQUIREMENTS
% ================================================================

\noindent
\textbf{I.\quad GENERAL REQUIREMENTS}

\vspace{0.15cm}

The Graduation Internship Report shall be prepared using the official
\LaTeX{} template provided by the Faculty. Students shall preserve the
predefined page layout, typography, line spacing, heading styles,
numbering system, figure and table formatting, and bibliography settings
unless otherwise instructed by the Faculty or the supervisor.

The report shall be written in English.

Students are responsible for the academic and technical content of the
report. The official \LaTeX{} template is responsible for standardized
document formatting.

\vspace{0.3cm}

% ================================================================
% II. REPORT STRUCTURE
% ================================================================

\noindent
\textbf{II.\quad REPORT STRUCTURE}

\vspace{0.15cm}

The report shall normally be arranged in the following order:

\begin{enumerate}[label=\arabic*., leftmargin=1.2cm, itemsep=2pt]
    \item Outer Cover;
    \item Comments of the Internship Host Organization;
    \item Supervisor's Evaluation and Graduation Internship Defense Results;
    \item Table of Contents;
    \item List of Symbols and Abbreviations, if applicable;
    \item List of Figures;
    \item List of Tables;
    \item Abstract;
    \item Main Chapters of the Internship Report;
    \item References;
    \item Appendices, if applicable.
\end{enumerate}

The preliminary pages shall not be numbered as chapters. The main report
content shall be organized into consecutively numbered chapters beginning
with Chapter 1.

\vspace{0.3cm}

% ================================================================
% III. PAGE SIZE AND MARGINS
% ================================================================

\noindent
\textbf{III.\quad PAGE SIZE AND MARGINS}

\vspace{0.15cm}

The report shall be formatted on A4 paper with the following margins:

\begin{center}
\begin{tabular}{ll}
    Left margin:   & 3.5 cm \\
    Right margin:  & 2.0 cm \\
    Top margin:    & 2.5 cm \\
    Bottom margin: & 2.5 cm
\end{tabular}
\end{center}

These settings are predefined in the official template and shall not be
modified by students.

\vspace{0.3cm}

% ================================================================
% IV. FONT, FONT SIZE, AND LINE SPACING
% ================================================================

\noindent
\textbf{IV.\quad FONT, FONT SIZE, AND LINE SPACING}

\vspace{0.15cm}

The main text shall use a 13 pt font size with 1.3-line spacing.
A Times-family typeface shall be used in accordance with the official
template.

The text shall be fully justified.

Students shall not use repeated spaces, repeated manual line breaks, or
multiple blank lines to position text. Paragraph and heading spacing shall
be controlled by the template.

\vspace{0.3cm}

% ================================================================
% V. OUTER COVER
% ================================================================

\noindent
\textbf{V.\quad OUTER COVER}

\vspace{0.15cm}

The outer cover shall follow the official format provided in the template
and shall contain the following information:

\begin{itemize}[leftmargin=0.8cm, itemsep=2pt]
    \item Thai Nguyen University of Technology;
    \item Faculty of International Training;
    \item Engineering and Technology Department;
    \item University logo and Faculty logo;
    \item the title \textbf{INTERNSHIP REPORT};
    \item Internship Organization;
    \item Author;
    \item Class;
    \item Supervisor;
    \item place and year of completion.
\end{itemize}

Students shall only replace the required information fields and shall not
modify the cover-page layout unless instructed otherwise.

\vspace{0.3cm}

% ================================================================
% VI. COMMENTS OF THE INTERNSHIP HOST ORGANIZATION
% ================================================================

\noindent
\textbf{VI.\quad COMMENTS OF THE INTERNSHIP HOST ORGANIZATION}

\vspace{0.15cm}

The page entitled \textbf{COMMENTS OF THE INTERNSHIP HOST ORGANIZATION}
is reserved for the internship host organization.

It includes:

\begin{itemize}[leftmargin=0.8cm, itemsep=2pt]
    \item the student's full name;
    \item class;
    \item comments from the host organization;
    \item score;
    \item date;
    \item signature and official seal.
\end{itemize}

Students shall not alter or complete official evaluation fields unless
specifically authorized.

\vspace{0.3cm}

% ================================================================
% VII. SUPERVISOR'S EVALUATION AND DEFENSE RESULTS
% ================================================================

\noindent
\textbf{VII.\quad SUPERVISOR'S EVALUATION AND DEFENSE RESULTS}

\vspace{0.15cm}

The section entitled
\textbf{SUPERVISOR'S EVALUATION AND GRADUATION INTERNSHIP DEFENSE RESULTS}
is reserved for the supervisor and defense committee.

The evaluation may include:

\begin{itemize}[leftmargin=0.8cm, itemsep=2pt]
    \item Internship Progress and Student's Attitude;
    \item Report Content;
    \item Report Presentation;
    \item Other Comments;
    \item Supervisor's Opinion;
    \item Defense Result;
    \item Score.
\end{itemize}

Students shall not enter evaluation scores or defense results before they
are officially determined.

\vspace{0.3cm}

% ================================================================
% VIII. CHAPTERS, SECTIONS, AND SUBSECTIONS
% ================================================================

\noindent
\textbf{VIII.\quad CHAPTERS, SECTIONS, AND SUBSECTIONS}

\vspace{0.15cm}

The report structure shall be generated using the standard \LaTeX{}
commands:

\begin{verbatim}
\chapter{...}
\section{...}
\subsection{...}
\end{verbatim}

Chapter, section, and subsection numbers shall be generated automatically
by \LaTeX{}.

Students shall not manually type chapter or section numbers and shall not
reproduce heading styles using commands such as \verb|\textbf|,
\verb|\fontsize|, or \verb|\centering| instead of the appropriate
structural command.

Excessive subdivision should be avoided. In most cases, the hierarchy
Chapter--Section--Subsection is sufficient.

\vspace{0.3cm}

% ================================================================
% IX. ORGANIZATION OF THE MAIN CONTENT
% ================================================================

\noindent
\textbf{IX.\quad ORGANIZATION OF THE MAIN CONTENT}

\vspace{0.15cm}

The specific content of the report shall depend on the internship
organization and the assigned tasks. A typical structure may include:

\begin{itemize}[leftmargin=0.8cm, itemsep=2pt]
    \item \textbf{Chapter 1. Overview of the Internship Organization};
    \item \textbf{Chapter 2. Internship Activities and Assigned Tasks};
    \item \textbf{Chapter 3. Results, Experience, and Evaluation}.
\end{itemize}

These chapter titles are recommendations only and may be adapted to the
actual internship content.

The report shall reflect work actually performed, observed, analyzed, or
learned during the internship. Unrelated theoretical material should not
be included merely to increase report length.

\vspace{0.3cm}

% ================================================================
% X. TABLE OF CONTENTS
% ================================================================

\noindent
\textbf{X.\quad TABLE OF CONTENTS}

\vspace{0.15cm}

The Table of Contents shall be generated automatically using:

\begin{verbatim}
\tableofcontents
\end{verbatim}

Students shall not create or manually edit page numbers in the Table of
Contents.

The document shall be compiled sufficiently to ensure that all page
numbers and entries are correctly updated.

\vspace{0.3cm}

% ================================================================
% XI. LIST OF SYMBOLS AND ABBREVIATIONS
% ================================================================

\noindent
\textbf{XI.\quad LIST OF SYMBOLS AND ABBREVIATIONS}

\vspace{0.15cm}

Symbols and abbreviations that are used repeatedly in the report shall be
listed in the \textbf{List of Symbols and Abbreviations}.

The list may contain the following columns:

\begin{center}
\begin{tabular}{llll}
    No. & Symbol / Abbreviation & Unit & Definition
\end{tabular}
\end{center}

Only symbols and abbreviations actually used in the final report shall be
retained. Example entries contained in the template shall be removed if
they are not relevant to the student's report.

Each abbreviation should normally be defined when it first appears in
the main text.

\vspace{0.3cm}

% ================================================================
% XII. FIGURES, DIAGRAMS, AND PHOTOGRAPHS
% ================================================================

\noindent
\textbf{XII.\quad FIGURES, DIAGRAMS, AND PHOTOGRAPHS}

\vspace{0.15cm}

Figures shall be clear, legible, and directly relevant to the report
content.

Figures shall be inserted using the standard \verb|figure| environment.

Example:

\begin{verbatim}
\begin{figure}[H]
    \centering
    \includegraphics[width=0.8\textwidth]
                    {figures/company_structure.pdf}
    \caption{Organizational structure of the company}
    \label{fig:company-structure}
\end{figure}
\end{verbatim}

Figure captions shall be placed below the figure.

Each referenced figure shall have a unique \verb|\label| and shall be
referenced in the text using, for example:

\begin{verbatim}
Figure~\ref{fig:company-structure}
\end{verbatim}

Students shall not manually type figure numbers.

Figures reproduced or adapted from external sources shall include an
appropriate citation.

Photographs taken during the internship shall be relevant and sufficiently
clear. Students shall not include confidential, restricted, or proprietary
information without permission from the internship organization.

Vector graphics or high-resolution images are recommended whenever
possible.

\vspace{0.3cm}

% ================================================================
% XIII. TABLES
% ================================================================

\noindent
\textbf{XIII.\quad TABLES}

\vspace{0.15cm}

Tables shall be created using suitable \LaTeX{} environments such as
\verb|table|, \verb|tabular|, \verb|tabularx|, or \verb|longtable|.

Table captions shall be placed above the table.

Example:

\begin{verbatim}
\begin{table}[H]
    \centering
    \caption{Internship schedule}
    \label{tab:internship-schedule}

    \begin{tabular}{lll}
        \toprule
        Week & Task & Result \\
        \midrule
        1 & Orientation & Completed \\
        2 & System study & Completed \\
        \bottomrule
    \end{tabular}
\end{table}
\end{verbatim}

Tables shall be referenced using:

\begin{verbatim}
Table~\ref{tab:internship-schedule}
\end{verbatim}

Students shall not manually type table numbers.

Tables should be created directly in \LaTeX{} rather than inserted as
screenshots whenever practical.

\vspace{0.3cm}

% ================================================================
% XIV. EQUATIONS AND TECHNICAL QUANTITIES
% ================================================================

\noindent
\textbf{XIV.\quad EQUATIONS AND TECHNICAL QUANTITIES}

\vspace{0.15cm}

Mathematical expressions shall be typeset directly using \LaTeX{}.
Images or screenshots of equations shall not be used.

Short expressions may be written in inline mathematical mode, for example:

\begin{verbatim}
$R = 0.7\,\Omega$
\end{verbatim}

Important equations shall use a numbered equation environment.

Example:

\begin{verbatim}
\begin{equation}
    P = UI
    \label{eq:power}
\end{equation}
\end{verbatim}

Equations shall be referenced using:

\begin{verbatim}
Equation~\eqref{eq:power}
\end{verbatim}

Equation numbers shall not be manually entered.

SI units shall be used wherever applicable. The \verb|siunitx| package
should be used where appropriate.

Examples:

\begin{verbatim}
\SI{220}{\volt}
\SI{50}{\hertz}
\SI{2.5}{\ampere}
\end{verbatim}

Notation and units shall be consistent throughout the report.

\vspace{0.3cm}

% ================================================================
% XV. SOURCE CODE
% ================================================================

\noindent
\textbf{XV.\quad SOURCE CODE}

\vspace{0.15cm}

If the internship involves programming, PLCs, microcontrollers,
simulation, or software development, relevant source code may be included
in the report.

Source code shall be presented using the \verb|listings| environment or
another environment defined by the official template.

Example:

\begin{verbatim}
\begin{lstlisting}[language=st]
IF Start = TRUE THEN
    Motor := TRUE;
ELSE
    Motor := FALSE;
END_IF;
\end{lstlisting}
\end{verbatim}

Only code necessary for explaining the work should appear in the main
report. Long source-code listings should be moved to an appendix or
submitted as supplementary files if required.

Source code should not normally be inserted as screenshots.

\vspace{0.3cm}

% ================================================================
% XVI. ABSTRACT
% ================================================================

\noindent
\textbf{XVI.\quad ABSTRACT}

\vspace{0.15cm}

The Abstract shall briefly summarize:

\begin{itemize}[leftmargin=0.8cm, itemsep=2pt]
    \item the internship organization;
    \item the objectives of the internship;
    \item the main tasks performed;
    \item the main results, knowledge, and skills obtained.
\end{itemize}

The Abstract should be concise and should not reproduce the full content
of the report chapter by chapter.

\vspace{0.3cm}

% ================================================================
% XVII. REFERENCES AND CITATIONS
% ================================================================

\noindent
\textbf{XVII.\quad REFERENCES AND CITATIONS}

\vspace{0.15cm}

All books, journal papers, conference papers, standards, websites,
datasheets, manuals, technical documents, figures, tables, or other
external sources used in preparing the report shall be properly cited.

References shall be managed through the bibliography system provided with
the official template.

Bibliographic entries shall normally be stored in:

\begin{verbatim}
References.bib
\end{verbatim}

References shall be cited in the report using:

\begin{verbatim}
\cite{reference-key}
\end{verbatim}

Students shall not manually type reference numbers such as ``[1]'',
``[2]'', etc.

Titles of foreign-language publications shall remain in their original
language in the bibliography.

All borrowed ideas, data, figures, tables, equations, methods, or results
shall be properly acknowledged.

\vspace{0.3cm}

% ================================================================
% XVIII. CROSS-REFERENCING
% ================================================================

\noindent
\textbf{XVIII.\quad CROSS-REFERENCING}

\vspace{0.15cm}

Numbered elements shall be cross-referenced using the standard \LaTeX{}
mechanism.

Use the following commands where appropriate:

\begin{verbatim}
\label{...}
\ref{...}
\eqref{...}
\end{verbatim}

For example:

\begin{verbatim}
As shown in Figure~\ref{fig:system}, ...
\end{verbatim}

Students shall not manually type numbers such as ``Figure 2.3'' when
those numbers can be generated automatically.

\vspace{0.3cm}

% ================================================================
% XIX. PAGE NUMBERING
% ================================================================

\noindent
\textbf{XIX.\quad PAGE NUMBERING}

\vspace{0.15cm}

Cover pages and official evaluation forms at the beginning of the report
shall not display page numbers.

The main body of the report shall use Arabic page numbers beginning from
Chapter 1.

Page numbering shall be generated automatically by the official template.
Students shall not manually enter page numbers.

\vspace{0.3cm}

% ================================================================
% XX. PROJECT FILE ORGANIZATION
% ================================================================

\noindent
\textbf{XX.\quad PROJECT FILE ORGANIZATION}

\vspace{0.15cm}

Students should maintain a clear and organized \LaTeX{} project
structure. A recommended structure is:

\begin{verbatim}
InternshipReport/
|
|-- main.tex
|-- References.bib
|
|-- figures/
|   |-- company.pdf
|   |-- system.png
|   `-- ...
|
`-- chapters/
    |-- chapter1.tex
    |-- chapter2.tex
    `-- chapter3.tex
\end{verbatim}

Figures should normally be stored in the \verb|figures| directory
specified by the official template.

File and directory names should preferably contain Latin letters,
numbers, hyphens, or underscores.

Recommended examples include:

\begin{verbatim}
company_structure.pdf
motor_control.png
chapter_01.tex
\end{verbatim}

Spaces and unnecessary special characters in file names should be
avoided where possible.

\vspace{0.3cm}

% ================================================================
% XXI. DOCUMENT COMPILATION AND FINAL CHECK
% ================================================================

\noindent
\textbf{XXI.\quad DOCUMENT COMPILATION AND FINAL CHECK}

\vspace{0.15cm}

Before submission, the complete \LaTeX{} project shall compile
successfully and produce the final PDF report.

The final document shall not contain unresolved references such as:

\begin{verbatim}
??
[?]
\end{verbatim}

Before generating the final PDF, students shall check:

\begin{itemize}[leftmargin=0.8cm, itemsep=2pt]
    \item the Table of Contents;
    \item page numbering;
    \item chapter, section, and subsection numbering;
    \item figure numbering;
    \item table numbering;
    \item equation numbering;
    \item cross-references;
    \item citations;
    \item bibliography;
    \item font display;
    \item mathematical symbols;
    \item image quality;
    \item completeness of all required forms.
\end{itemize}

The document may be compiled using TeX Live, MiKTeX, or another
compatible \LaTeX{} environment. Regardless of the compilation
environment, the resulting PDF shall conform to the official Faculty
template.

\vspace{0.3cm}

% ================================================================
% XXII. CONTENT THAT STUDENTS MAY EDIT
% ================================================================

\noindent
\textbf{XXII.\quad CONTENT THAT STUDENTS MAY EDIT}

\vspace{0.15cm}

Students may edit or replace:

\begin{itemize}[leftmargin=0.8cm, itemsep=2pt]
    \item student name;
    \item class;
    \item internship organization;
    \item supervisor name;
    \item Abstract;
    \item List of Symbols and Abbreviations;
    \item chapter titles and chapter content;
    \item figures and photographs;
    \item tables;
    \item equations;
    \item citations and references.
\end{itemize}

Students shall not modify the main document formatting settings in the
preamble, including the document class, page geometry, line spacing,
heading formatting, table-of-contents formatting, font configuration,
page layout, or cover-page macros unless specifically instructed.

\vspace{0.3cm}

% ================================================================
% XXIII. FINAL SUBMISSION REQUIREMENTS
% ================================================================

\noindent
\textbf{XXIII.\quad FINAL SUBMISSION REQUIREMENTS}

\vspace{0.15cm}

Before submission, the final report shall:

\begin{itemize}[leftmargin=0.8cm, itemsep=2pt]
    \item compile without errors that affect the report content or presentation;
    \item contain all required forms and report sections;
    \item contain correct student, class, internship organization, and supervisor information;
    \item contain an updated Table of Contents, List of Figures, and List of Tables;
    \item contain no unused example content from the template;
    \item contain no placeholder text such as \textit{Title of Chapter 1},
    \textit{Section 1.1}, or other sample entries;
    \item contain all required citations and references;
    \item contain no unresolved cross-references.
\end{itemize}

If the Faculty requires submission of the \LaTeX{} source, the complete
project shall include all \verb|.tex| files, the bibliography database,
figures, and any other files required to reproduce the submitted PDF.

\vspace{0.4cm}

\noindent
\textbf{Important:}
Before submitting the final internship report, students shall remove or
comment out the following command in the main document:

\begin{verbatim}
\input{Regulations.tex}
\end{verbatim}

For example:

\begin{verbatim}
% \input{Regulations.tex}
\end{verbatim}

This removes the guidance pages from the final submitted report while
preserving the regulations file in the \LaTeX{} project for reference.

} % End 13pt content

\clearpage

%
% IMPORTANT:
% This section is provided for guidance only. Students shall remove
% or comment out the 
% =================================================================
% REGULATIONS FOR PREPARING THE GRADUATION INTERNSHIP REPORT
% USING LATEX
% =================================================================
%
% This file is intended to be included from the main document with:
%
%     \input{Regulations.tex}
%
% IMPORTANT:
% This section is provided for guidance only. Students shall remove
% or comment out the \input{Regulations.tex} command before submitting
% the final internship report.
%
% =================================================================

\clearpage
\thispagestyle{empty}

\begin{center}
    {\fontsize{16pt}{20pt}\selectfont\bfseries
    SPECIFIC REGULATIONS FOR PREPARING\\
    GRADUATION INTERNSHIP REPORTS USING \LaTeX}
\end{center}

\vspace{0.5cm}

{\fontsize{13pt}{17pt}\selectfont
\setstretch{1.3}
\justifying

\noindent
\textbf{Note:}
This section is provided for guidance only and is included in the main
document using the command \verb|\input{Regulations.tex}|.
Students shall remove or comment out this command before submitting the
final internship report.

\vspace{0.4cm}

% ================================================================
% I. GENERAL REQUIREMENTS
% ================================================================

\noindent
\textbf{I.\quad GENERAL REQUIREMENTS}

\vspace{0.15cm}

The Graduation Internship Report shall be prepared using the official
\LaTeX{} template provided by the Faculty. Students shall preserve the
predefined page layout, typography, line spacing, heading styles,
numbering system, figure and table formatting, and bibliography settings
unless otherwise instructed by the Faculty or the supervisor.

The report shall be written in English.

Students are responsible for the academic and technical content of the
report. The official \LaTeX{} template is responsible for standardized
document formatting.

\vspace{0.3cm}

% ================================================================
% II. REPORT STRUCTURE
% ================================================================

\noindent
\textbf{II.\quad REPORT STRUCTURE}

\vspace{0.15cm}

The report shall normally be arranged in the following order:

\begin{enumerate}[label=\arabic*., leftmargin=1.2cm, itemsep=2pt]
    \item Outer Cover;
    \item Comments of the Internship Host Organization;
    \item Supervisor's Evaluation and Graduation Internship Defense Results;
    \item Table of Contents;
    \item List of Symbols and Abbreviations, if applicable;
    \item List of Figures;
    \item List of Tables;
    \item Abstract;
    \item Main Chapters of the Internship Report;
    \item References;
    \item Appendices, if applicable.
\end{enumerate}

The preliminary pages shall not be numbered as chapters. The main report
content shall be organized into consecutively numbered chapters beginning
with Chapter 1.

\vspace{0.3cm}

% ================================================================
% III. PAGE SIZE AND MARGINS
% ================================================================

\noindent
\textbf{III.\quad PAGE SIZE AND MARGINS}

\vspace{0.15cm}

The report shall be formatted on A4 paper with the following margins:

\begin{center}
\begin{tabular}{ll}
    Left margin:   & 3.5 cm \\
    Right margin:  & 2.0 cm \\
    Top margin:    & 2.5 cm \\
    Bottom margin: & 2.5 cm
\end{tabular}
\end{center}

These settings are predefined in the official template and shall not be
modified by students.

\vspace{0.3cm}

% ================================================================
% IV. FONT, FONT SIZE, AND LINE SPACING
% ================================================================

\noindent
\textbf{IV.\quad FONT, FONT SIZE, AND LINE SPACING}

\vspace{0.15cm}

The main text shall use a 13 pt font size with 1.3-line spacing.
A Times-family typeface shall be used in accordance with the official
template.

The text shall be fully justified.

Students shall not use repeated spaces, repeated manual line breaks, or
multiple blank lines to position text. Paragraph and heading spacing shall
be controlled by the template.

\vspace{0.3cm}

% ================================================================
% V. OUTER COVER
% ================================================================

\noindent
\textbf{V.\quad OUTER COVER}

\vspace{0.15cm}

The outer cover shall follow the official format provided in the template
and shall contain the following information:

\begin{itemize}[leftmargin=0.8cm, itemsep=2pt]
    \item Thai Nguyen University of Technology;
    \item Faculty of International Training;
    \item Engineering and Technology Department;
    \item University logo and Faculty logo;
    \item the title \textbf{INTERNSHIP REPORT};
    \item Internship Organization;
    \item Author;
    \item Class;
    \item Supervisor;
    \item place and year of completion.
\end{itemize}

Students shall only replace the required information fields and shall not
modify the cover-page layout unless instructed otherwise.

\vspace{0.3cm}

% ================================================================
% VI. COMMENTS OF THE INTERNSHIP HOST ORGANIZATION
% ================================================================

\noindent
\textbf{VI.\quad COMMENTS OF THE INTERNSHIP HOST ORGANIZATION}

\vspace{0.15cm}

The page entitled \textbf{COMMENTS OF THE INTERNSHIP HOST ORGANIZATION}
is reserved for the internship host organization.

It includes:

\begin{itemize}[leftmargin=0.8cm, itemsep=2pt]
    \item the student's full name;
    \item class;
    \item comments from the host organization;
    \item score;
    \item date;
    \item signature and official seal.
\end{itemize}

Students shall not alter or complete official evaluation fields unless
specifically authorized.

\vspace{0.3cm}

% ================================================================
% VII. SUPERVISOR'S EVALUATION AND DEFENSE RESULTS
% ================================================================

\noindent
\textbf{VII.\quad SUPERVISOR'S EVALUATION AND DEFENSE RESULTS}

\vspace{0.15cm}

The section entitled
\textbf{SUPERVISOR'S EVALUATION AND GRADUATION INTERNSHIP DEFENSE RESULTS}
is reserved for the supervisor and defense committee.

The evaluation may include:

\begin{itemize}[leftmargin=0.8cm, itemsep=2pt]
    \item Internship Progress and Student's Attitude;
    \item Report Content;
    \item Report Presentation;
    \item Other Comments;
    \item Supervisor's Opinion;
    \item Defense Result;
    \item Score.
\end{itemize}

Students shall not enter evaluation scores or defense results before they
are officially determined.

\vspace{0.3cm}

% ================================================================
% VIII. CHAPTERS, SECTIONS, AND SUBSECTIONS
% ================================================================

\noindent
\textbf{VIII.\quad CHAPTERS, SECTIONS, AND SUBSECTIONS}

\vspace{0.15cm}

The report structure shall be generated using the standard \LaTeX{}
commands:

\begin{verbatim}
\chapter{...}
\section{...}
\subsection{...}
\end{verbatim}

Chapter, section, and subsection numbers shall be generated automatically
by \LaTeX{}.

Students shall not manually type chapter or section numbers and shall not
reproduce heading styles using commands such as \verb|\textbf|,
\verb|\fontsize|, or \verb|\centering| instead of the appropriate
structural command.

Excessive subdivision should be avoided. In most cases, the hierarchy
Chapter--Section--Subsection is sufficient.

\vspace{0.3cm}

% ================================================================
% IX. ORGANIZATION OF THE MAIN CONTENT
% ================================================================

\noindent
\textbf{IX.\quad ORGANIZATION OF THE MAIN CONTENT}

\vspace{0.15cm}

The specific content of the report shall depend on the internship
organization and the assigned tasks. A typical structure may include:

\begin{itemize}[leftmargin=0.8cm, itemsep=2pt]
    \item \textbf{Chapter 1. Overview of the Internship Organization};
    \item \textbf{Chapter 2. Internship Activities and Assigned Tasks};
    \item \textbf{Chapter 3. Results, Experience, and Evaluation}.
\end{itemize}

These chapter titles are recommendations only and may be adapted to the
actual internship content.

The report shall reflect work actually performed, observed, analyzed, or
learned during the internship. Unrelated theoretical material should not
be included merely to increase report length.

\vspace{0.3cm}

% ================================================================
% X. TABLE OF CONTENTS
% ================================================================

\noindent
\textbf{X.\quad TABLE OF CONTENTS}

\vspace{0.15cm}

The Table of Contents shall be generated automatically using:

\begin{verbatim}
\tableofcontents
\end{verbatim}

Students shall not create or manually edit page numbers in the Table of
Contents.

The document shall be compiled sufficiently to ensure that all page
numbers and entries are correctly updated.

\vspace{0.3cm}

% ================================================================
% XI. LIST OF SYMBOLS AND ABBREVIATIONS
% ================================================================

\noindent
\textbf{XI.\quad LIST OF SYMBOLS AND ABBREVIATIONS}

\vspace{0.15cm}

Symbols and abbreviations that are used repeatedly in the report shall be
listed in the \textbf{List of Symbols and Abbreviations}.

The list may contain the following columns:

\begin{center}
\begin{tabular}{llll}
    No. & Symbol / Abbreviation & Unit & Definition
\end{tabular}
\end{center}

Only symbols and abbreviations actually used in the final report shall be
retained. Example entries contained in the template shall be removed if
they are not relevant to the student's report.

Each abbreviation should normally be defined when it first appears in
the main text.

\vspace{0.3cm}

% ================================================================
% XII. FIGURES, DIAGRAMS, AND PHOTOGRAPHS
% ================================================================

\noindent
\textbf{XII.\quad FIGURES, DIAGRAMS, AND PHOTOGRAPHS}

\vspace{0.15cm}

Figures shall be clear, legible, and directly relevant to the report
content.

Figures shall be inserted using the standard \verb|figure| environment.

Example:

\begin{verbatim}
\begin{figure}[H]
    \centering
    \includegraphics[width=0.8\textwidth]
                    {figures/company_structure.pdf}
    \caption{Organizational structure of the company}
    \label{fig:company-structure}
\end{figure}
\end{verbatim}

Figure captions shall be placed below the figure.

Each referenced figure shall have a unique \verb|\label| and shall be
referenced in the text using, for example:

\begin{verbatim}
Figure~\ref{fig:company-structure}
\end{verbatim}

Students shall not manually type figure numbers.

Figures reproduced or adapted from external sources shall include an
appropriate citation.

Photographs taken during the internship shall be relevant and sufficiently
clear. Students shall not include confidential, restricted, or proprietary
information without permission from the internship organization.

Vector graphics or high-resolution images are recommended whenever
possible.

\vspace{0.3cm}

% ================================================================
% XIII. TABLES
% ================================================================

\noindent
\textbf{XIII.\quad TABLES}

\vspace{0.15cm}

Tables shall be created using suitable \LaTeX{} environments such as
\verb|table|, \verb|tabular|, \verb|tabularx|, or \verb|longtable|.

Table captions shall be placed above the table.

Example:

\begin{verbatim}
\begin{table}[H]
    \centering
    \caption{Internship schedule}
    \label{tab:internship-schedule}

    \begin{tabular}{lll}
        \toprule
        Week & Task & Result \\
        \midrule
        1 & Orientation & Completed \\
        2 & System study & Completed \\
        \bottomrule
    \end{tabular}
\end{table}
\end{verbatim}

Tables shall be referenced using:

\begin{verbatim}
Table~\ref{tab:internship-schedule}
\end{verbatim}

Students shall not manually type table numbers.

Tables should be created directly in \LaTeX{} rather than inserted as
screenshots whenever practical.

\vspace{0.3cm}

% ================================================================
% XIV. EQUATIONS AND TECHNICAL QUANTITIES
% ================================================================

\noindent
\textbf{XIV.\quad EQUATIONS AND TECHNICAL QUANTITIES}

\vspace{0.15cm}

Mathematical expressions shall be typeset directly using \LaTeX{}.
Images or screenshots of equations shall not be used.

Short expressions may be written in inline mathematical mode, for example:

\begin{verbatim}
$R = 0.7\,\Omega$
\end{verbatim}

Important equations shall use a numbered equation environment.

Example:

\begin{verbatim}
\begin{equation}
    P = UI
    \label{eq:power}
\end{equation}
\end{verbatim}

Equations shall be referenced using:

\begin{verbatim}
Equation~\eqref{eq:power}
\end{verbatim}

Equation numbers shall not be manually entered.

SI units shall be used wherever applicable. The \verb|siunitx| package
should be used where appropriate.

Examples:

\begin{verbatim}
\SI{220}{\volt}
\SI{50}{\hertz}
\SI{2.5}{\ampere}
\end{verbatim}

Notation and units shall be consistent throughout the report.

\vspace{0.3cm}

% ================================================================
% XV. SOURCE CODE
% ================================================================

\noindent
\textbf{XV.\quad SOURCE CODE}

\vspace{0.15cm}

If the internship involves programming, PLCs, microcontrollers,
simulation, or software development, relevant source code may be included
in the report.

Source code shall be presented using the \verb|listings| environment or
another environment defined by the official template.

Example:

\begin{verbatim}
\begin{lstlisting}[language=st]
IF Start = TRUE THEN
    Motor := TRUE;
ELSE
    Motor := FALSE;
END_IF;
\end{lstlisting}
\end{verbatim}

Only code necessary for explaining the work should appear in the main
report. Long source-code listings should be moved to an appendix or
submitted as supplementary files if required.

Source code should not normally be inserted as screenshots.

\vspace{0.3cm}

% ================================================================
% XVI. ABSTRACT
% ================================================================

\noindent
\textbf{XVI.\quad ABSTRACT}

\vspace{0.15cm}

The Abstract shall briefly summarize:

\begin{itemize}[leftmargin=0.8cm, itemsep=2pt]
    \item the internship organization;
    \item the objectives of the internship;
    \item the main tasks performed;
    \item the main results, knowledge, and skills obtained.
\end{itemize}

The Abstract should be concise and should not reproduce the full content
of the report chapter by chapter.

\vspace{0.3cm}

% ================================================================
% XVII. REFERENCES AND CITATIONS
% ================================================================

\noindent
\textbf{XVII.\quad REFERENCES AND CITATIONS}

\vspace{0.15cm}

All books, journal papers, conference papers, standards, websites,
datasheets, manuals, technical documents, figures, tables, or other
external sources used in preparing the report shall be properly cited.

References shall be managed through the bibliography system provided with
the official template.

Bibliographic entries shall normally be stored in:

\begin{verbatim}
References.bib
\end{verbatim}

References shall be cited in the report using:

\begin{verbatim}
\cite{reference-key}
\end{verbatim}

Students shall not manually type reference numbers such as ``[1]'',
``[2]'', etc.

Titles of foreign-language publications shall remain in their original
language in the bibliography.

All borrowed ideas, data, figures, tables, equations, methods, or results
shall be properly acknowledged.

\vspace{0.3cm}

% ================================================================
% XVIII. CROSS-REFERENCING
% ================================================================

\noindent
\textbf{XVIII.\quad CROSS-REFERENCING}

\vspace{0.15cm}

Numbered elements shall be cross-referenced using the standard \LaTeX{}
mechanism.

Use the following commands where appropriate:

\begin{verbatim}
\label{...}
\ref{...}
\eqref{...}
\end{verbatim}

For example:

\begin{verbatim}
As shown in Figure~\ref{fig:system}, ...
\end{verbatim}

Students shall not manually type numbers such as ``Figure 2.3'' when
those numbers can be generated automatically.

\vspace{0.3cm}

% ================================================================
% XIX. PAGE NUMBERING
% ================================================================

\noindent
\textbf{XIX.\quad PAGE NUMBERING}

\vspace{0.15cm}

Cover pages and official evaluation forms at the beginning of the report
shall not display page numbers.

The main body of the report shall use Arabic page numbers beginning from
Chapter 1.

Page numbering shall be generated automatically by the official template.
Students shall not manually enter page numbers.

\vspace{0.3cm}

% ================================================================
% XX. PROJECT FILE ORGANIZATION
% ================================================================

\noindent
\textbf{XX.\quad PROJECT FILE ORGANIZATION}

\vspace{0.15cm}

Students should maintain a clear and organized \LaTeX{} project
structure. A recommended structure is:

\begin{verbatim}
InternshipReport/
|
|-- main.tex
|-- References.bib
|
|-- figures/
|   |-- company.pdf
|   |-- system.png
|   `-- ...
|
`-- chapters/
    |-- chapter1.tex
    |-- chapter2.tex
    `-- chapter3.tex
\end{verbatim}

Figures should normally be stored in the \verb|figures| directory
specified by the official template.

File and directory names should preferably contain Latin letters,
numbers, hyphens, or underscores.

Recommended examples include:

\begin{verbatim}
company_structure.pdf
motor_control.png
chapter_01.tex
\end{verbatim}

Spaces and unnecessary special characters in file names should be
avoided where possible.

\vspace{0.3cm}

% ================================================================
% XXI. DOCUMENT COMPILATION AND FINAL CHECK
% ================================================================

\noindent
\textbf{XXI.\quad DOCUMENT COMPILATION AND FINAL CHECK}

\vspace{0.15cm}

Before submission, the complete \LaTeX{} project shall compile
successfully and produce the final PDF report.

The final document shall not contain unresolved references such as:

\begin{verbatim}
??
[?]
\end{verbatim}

Before generating the final PDF, students shall check:

\begin{itemize}[leftmargin=0.8cm, itemsep=2pt]
    \item the Table of Contents;
    \item page numbering;
    \item chapter, section, and subsection numbering;
    \item figure numbering;
    \item table numbering;
    \item equation numbering;
    \item cross-references;
    \item citations;
    \item bibliography;
    \item font display;
    \item mathematical symbols;
    \item image quality;
    \item completeness of all required forms.
\end{itemize}

The document may be compiled using TeX Live, MiKTeX, or another
compatible \LaTeX{} environment. Regardless of the compilation
environment, the resulting PDF shall conform to the official Faculty
template.

\vspace{0.3cm}

% ================================================================
% XXII. CONTENT THAT STUDENTS MAY EDIT
% ================================================================

\noindent
\textbf{XXII.\quad CONTENT THAT STUDENTS MAY EDIT}

\vspace{0.15cm}

Students may edit or replace:

\begin{itemize}[leftmargin=0.8cm, itemsep=2pt]
    \item student name;
    \item class;
    \item internship organization;
    \item supervisor name;
    \item Abstract;
    \item List of Symbols and Abbreviations;
    \item chapter titles and chapter content;
    \item figures and photographs;
    \item tables;
    \item equations;
    \item citations and references.
\end{itemize}

Students shall not modify the main document formatting settings in the
preamble, including the document class, page geometry, line spacing,
heading formatting, table-of-contents formatting, font configuration,
page layout, or cover-page macros unless specifically instructed.

\vspace{0.3cm}

% ================================================================
% XXIII. FINAL SUBMISSION REQUIREMENTS
% ================================================================

\noindent
\textbf{XXIII.\quad FINAL SUBMISSION REQUIREMENTS}

\vspace{0.15cm}

Before submission, the final report shall:

\begin{itemize}[leftmargin=0.8cm, itemsep=2pt]
    \item compile without errors that affect the report content or presentation;
    \item contain all required forms and report sections;
    \item contain correct student, class, internship organization, and supervisor information;
    \item contain an updated Table of Contents, List of Figures, and List of Tables;
    \item contain no unused example content from the template;
    \item contain no placeholder text such as \textit{Title of Chapter 1},
    \textit{Section 1.1}, or other sample entries;
    \item contain all required citations and references;
    \item contain no unresolved cross-references.
\end{itemize}

If the Faculty requires submission of the \LaTeX{} source, the complete
project shall include all \verb|.tex| files, the bibliography database,
figures, and any other files required to reproduce the submitted PDF.

\vspace{0.4cm}

\noindent
\textbf{Important:}
Before submitting the final internship report, students shall remove or
comment out the following command in the main document:

\begin{verbatim}
\input{Regulations.tex}
\end{verbatim}

For example:

\begin{verbatim}
% \input{Regulations.tex}
\end{verbatim}

This removes the guidance pages from the final submitted report while
preserving the regulations file in the \LaTeX{} project for reference.

} % End 13pt content

\clearpage
 command before submitting
% the final internship report.
%
% =================================================================

\clearpage
\thispagestyle{empty}

\begin{center}
    {\fontsize{16pt}{20pt}\selectfont\bfseries
    SPECIFIC REGULATIONS FOR PREPARING\\
    GRADUATION INTERNSHIP REPORTS USING \LaTeX}
\end{center}

\vspace{0.5cm}

{\fontsize{13pt}{17pt}\selectfont
\setstretch{1.3}
\justifying

\noindent
\textbf{Note:}
This section is provided for guidance only and is included in the main
document using the command \verb|
% =================================================================
% REGULATIONS FOR PREPARING THE GRADUATION INTERNSHIP REPORT
% USING LATEX
% =================================================================
%
% This file is intended to be included from the main document with:
%
%     \input{Regulations.tex}
%
% IMPORTANT:
% This section is provided for guidance only. Students shall remove
% or comment out the \input{Regulations.tex} command before submitting
% the final internship report.
%
% =================================================================

\clearpage
\thispagestyle{empty}

\begin{center}
    {\fontsize{16pt}{20pt}\selectfont\bfseries
    SPECIFIC REGULATIONS FOR PREPARING\\
    GRADUATION INTERNSHIP REPORTS USING \LaTeX}
\end{center}

\vspace{0.5cm}

{\fontsize{13pt}{17pt}\selectfont
\setstretch{1.3}
\justifying

\noindent
\textbf{Note:}
This section is provided for guidance only and is included in the main
document using the command \verb|\input{Regulations.tex}|.
Students shall remove or comment out this command before submitting the
final internship report.

\vspace{0.4cm}

% ================================================================
% I. GENERAL REQUIREMENTS
% ================================================================

\noindent
\textbf{I.\quad GENERAL REQUIREMENTS}

\vspace{0.15cm}

The Graduation Internship Report shall be prepared using the official
\LaTeX{} template provided by the Faculty. Students shall preserve the
predefined page layout, typography, line spacing, heading styles,
numbering system, figure and table formatting, and bibliography settings
unless otherwise instructed by the Faculty or the supervisor.

The report shall be written in English.

Students are responsible for the academic and technical content of the
report. The official \LaTeX{} template is responsible for standardized
document formatting.

\vspace{0.3cm}

% ================================================================
% II. REPORT STRUCTURE
% ================================================================

\noindent
\textbf{II.\quad REPORT STRUCTURE}

\vspace{0.15cm}

The report shall normally be arranged in the following order:

\begin{enumerate}[label=\arabic*., leftmargin=1.2cm, itemsep=2pt]
    \item Outer Cover;
    \item Comments of the Internship Host Organization;
    \item Supervisor's Evaluation and Graduation Internship Defense Results;
    \item Table of Contents;
    \item List of Symbols and Abbreviations, if applicable;
    \item List of Figures;
    \item List of Tables;
    \item Abstract;
    \item Main Chapters of the Internship Report;
    \item References;
    \item Appendices, if applicable.
\end{enumerate}

The preliminary pages shall not be numbered as chapters. The main report
content shall be organized into consecutively numbered chapters beginning
with Chapter 1.

\vspace{0.3cm}

% ================================================================
% III. PAGE SIZE AND MARGINS
% ================================================================

\noindent
\textbf{III.\quad PAGE SIZE AND MARGINS}

\vspace{0.15cm}

The report shall be formatted on A4 paper with the following margins:

\begin{center}
\begin{tabular}{ll}
    Left margin:   & 3.5 cm \\
    Right margin:  & 2.0 cm \\
    Top margin:    & 2.5 cm \\
    Bottom margin: & 2.5 cm
\end{tabular}
\end{center}

These settings are predefined in the official template and shall not be
modified by students.

\vspace{0.3cm}

% ================================================================
% IV. FONT, FONT SIZE, AND LINE SPACING
% ================================================================

\noindent
\textbf{IV.\quad FONT, FONT SIZE, AND LINE SPACING}

\vspace{0.15cm}

The main text shall use a 13 pt font size with 1.3-line spacing.
A Times-family typeface shall be used in accordance with the official
template.

The text shall be fully justified.

Students shall not use repeated spaces, repeated manual line breaks, or
multiple blank lines to position text. Paragraph and heading spacing shall
be controlled by the template.

\vspace{0.3cm}

% ================================================================
% V. OUTER COVER
% ================================================================

\noindent
\textbf{V.\quad OUTER COVER}

\vspace{0.15cm}

The outer cover shall follow the official format provided in the template
and shall contain the following information:

\begin{itemize}[leftmargin=0.8cm, itemsep=2pt]
    \item Thai Nguyen University of Technology;
    \item Faculty of International Training;
    \item Engineering and Technology Department;
    \item University logo and Faculty logo;
    \item the title \textbf{INTERNSHIP REPORT};
    \item Internship Organization;
    \item Author;
    \item Class;
    \item Supervisor;
    \item place and year of completion.
\end{itemize}

Students shall only replace the required information fields and shall not
modify the cover-page layout unless instructed otherwise.

\vspace{0.3cm}

% ================================================================
% VI. COMMENTS OF THE INTERNSHIP HOST ORGANIZATION
% ================================================================

\noindent
\textbf{VI.\quad COMMENTS OF THE INTERNSHIP HOST ORGANIZATION}

\vspace{0.15cm}

The page entitled \textbf{COMMENTS OF THE INTERNSHIP HOST ORGANIZATION}
is reserved for the internship host organization.

It includes:

\begin{itemize}[leftmargin=0.8cm, itemsep=2pt]
    \item the student's full name;
    \item class;
    \item comments from the host organization;
    \item score;
    \item date;
    \item signature and official seal.
\end{itemize}

Students shall not alter or complete official evaluation fields unless
specifically authorized.

\vspace{0.3cm}

% ================================================================
% VII. SUPERVISOR'S EVALUATION AND DEFENSE RESULTS
% ================================================================

\noindent
\textbf{VII.\quad SUPERVISOR'S EVALUATION AND DEFENSE RESULTS}

\vspace{0.15cm}

The section entitled
\textbf{SUPERVISOR'S EVALUATION AND GRADUATION INTERNSHIP DEFENSE RESULTS}
is reserved for the supervisor and defense committee.

The evaluation may include:

\begin{itemize}[leftmargin=0.8cm, itemsep=2pt]
    \item Internship Progress and Student's Attitude;
    \item Report Content;
    \item Report Presentation;
    \item Other Comments;
    \item Supervisor's Opinion;
    \item Defense Result;
    \item Score.
\end{itemize}

Students shall not enter evaluation scores or defense results before they
are officially determined.

\vspace{0.3cm}

% ================================================================
% VIII. CHAPTERS, SECTIONS, AND SUBSECTIONS
% ================================================================

\noindent
\textbf{VIII.\quad CHAPTERS, SECTIONS, AND SUBSECTIONS}

\vspace{0.15cm}

The report structure shall be generated using the standard \LaTeX{}
commands:

\begin{verbatim}
\chapter{...}
\section{...}
\subsection{...}
\end{verbatim}

Chapter, section, and subsection numbers shall be generated automatically
by \LaTeX{}.

Students shall not manually type chapter or section numbers and shall not
reproduce heading styles using commands such as \verb|\textbf|,
\verb|\fontsize|, or \verb|\centering| instead of the appropriate
structural command.

Excessive subdivision should be avoided. In most cases, the hierarchy
Chapter--Section--Subsection is sufficient.

\vspace{0.3cm}

% ================================================================
% IX. ORGANIZATION OF THE MAIN CONTENT
% ================================================================

\noindent
\textbf{IX.\quad ORGANIZATION OF THE MAIN CONTENT}

\vspace{0.15cm}

The specific content of the report shall depend on the internship
organization and the assigned tasks. A typical structure may include:

\begin{itemize}[leftmargin=0.8cm, itemsep=2pt]
    \item \textbf{Chapter 1. Overview of the Internship Organization};
    \item \textbf{Chapter 2. Internship Activities and Assigned Tasks};
    \item \textbf{Chapter 3. Results, Experience, and Evaluation}.
\end{itemize}

These chapter titles are recommendations only and may be adapted to the
actual internship content.

The report shall reflect work actually performed, observed, analyzed, or
learned during the internship. Unrelated theoretical material should not
be included merely to increase report length.

\vspace{0.3cm}

% ================================================================
% X. TABLE OF CONTENTS
% ================================================================

\noindent
\textbf{X.\quad TABLE OF CONTENTS}

\vspace{0.15cm}

The Table of Contents shall be generated automatically using:

\begin{verbatim}
\tableofcontents
\end{verbatim}

Students shall not create or manually edit page numbers in the Table of
Contents.

The document shall be compiled sufficiently to ensure that all page
numbers and entries are correctly updated.

\vspace{0.3cm}

% ================================================================
% XI. LIST OF SYMBOLS AND ABBREVIATIONS
% ================================================================

\noindent
\textbf{XI.\quad LIST OF SYMBOLS AND ABBREVIATIONS}

\vspace{0.15cm}

Symbols and abbreviations that are used repeatedly in the report shall be
listed in the \textbf{List of Symbols and Abbreviations}.

The list may contain the following columns:

\begin{center}
\begin{tabular}{llll}
    No. & Symbol / Abbreviation & Unit & Definition
\end{tabular}
\end{center}

Only symbols and abbreviations actually used in the final report shall be
retained. Example entries contained in the template shall be removed if
they are not relevant to the student's report.

Each abbreviation should normally be defined when it first appears in
the main text.

\vspace{0.3cm}

% ================================================================
% XII. FIGURES, DIAGRAMS, AND PHOTOGRAPHS
% ================================================================

\noindent
\textbf{XII.\quad FIGURES, DIAGRAMS, AND PHOTOGRAPHS}

\vspace{0.15cm}

Figures shall be clear, legible, and directly relevant to the report
content.

Figures shall be inserted using the standard \verb|figure| environment.

Example:

\begin{verbatim}
\begin{figure}[H]
    \centering
    \includegraphics[width=0.8\textwidth]
                    {figures/company_structure.pdf}
    \caption{Organizational structure of the company}
    \label{fig:company-structure}
\end{figure}
\end{verbatim}

Figure captions shall be placed below the figure.

Each referenced figure shall have a unique \verb|\label| and shall be
referenced in the text using, for example:

\begin{verbatim}
Figure~\ref{fig:company-structure}
\end{verbatim}

Students shall not manually type figure numbers.

Figures reproduced or adapted from external sources shall include an
appropriate citation.

Photographs taken during the internship shall be relevant and sufficiently
clear. Students shall not include confidential, restricted, or proprietary
information without permission from the internship organization.

Vector graphics or high-resolution images are recommended whenever
possible.

\vspace{0.3cm}

% ================================================================
% XIII. TABLES
% ================================================================

\noindent
\textbf{XIII.\quad TABLES}

\vspace{0.15cm}

Tables shall be created using suitable \LaTeX{} environments such as
\verb|table|, \verb|tabular|, \verb|tabularx|, or \verb|longtable|.

Table captions shall be placed above the table.

Example:

\begin{verbatim}
\begin{table}[H]
    \centering
    \caption{Internship schedule}
    \label{tab:internship-schedule}

    \begin{tabular}{lll}
        \toprule
        Week & Task & Result \\
        \midrule
        1 & Orientation & Completed \\
        2 & System study & Completed \\
        \bottomrule
    \end{tabular}
\end{table}
\end{verbatim}

Tables shall be referenced using:

\begin{verbatim}
Table~\ref{tab:internship-schedule}
\end{verbatim}

Students shall not manually type table numbers.

Tables should be created directly in \LaTeX{} rather than inserted as
screenshots whenever practical.

\vspace{0.3cm}

% ================================================================
% XIV. EQUATIONS AND TECHNICAL QUANTITIES
% ================================================================

\noindent
\textbf{XIV.\quad EQUATIONS AND TECHNICAL QUANTITIES}

\vspace{0.15cm}

Mathematical expressions shall be typeset directly using \LaTeX{}.
Images or screenshots of equations shall not be used.

Short expressions may be written in inline mathematical mode, for example:

\begin{verbatim}
$R = 0.7\,\Omega$
\end{verbatim}

Important equations shall use a numbered equation environment.

Example:

\begin{verbatim}
\begin{equation}
    P = UI
    \label{eq:power}
\end{equation}
\end{verbatim}

Equations shall be referenced using:

\begin{verbatim}
Equation~\eqref{eq:power}
\end{verbatim}

Equation numbers shall not be manually entered.

SI units shall be used wherever applicable. The \verb|siunitx| package
should be used where appropriate.

Examples:

\begin{verbatim}
\SI{220}{\volt}
\SI{50}{\hertz}
\SI{2.5}{\ampere}
\end{verbatim}

Notation and units shall be consistent throughout the report.

\vspace{0.3cm}

% ================================================================
% XV. SOURCE CODE
% ================================================================

\noindent
\textbf{XV.\quad SOURCE CODE}

\vspace{0.15cm}

If the internship involves programming, PLCs, microcontrollers,
simulation, or software development, relevant source code may be included
in the report.

Source code shall be presented using the \verb|listings| environment or
another environment defined by the official template.

Example:

\begin{verbatim}
\begin{lstlisting}[language=st]
IF Start = TRUE THEN
    Motor := TRUE;
ELSE
    Motor := FALSE;
END_IF;
\end{lstlisting}
\end{verbatim}

Only code necessary for explaining the work should appear in the main
report. Long source-code listings should be moved to an appendix or
submitted as supplementary files if required.

Source code should not normally be inserted as screenshots.

\vspace{0.3cm}

% ================================================================
% XVI. ABSTRACT
% ================================================================

\noindent
\textbf{XVI.\quad ABSTRACT}

\vspace{0.15cm}

The Abstract shall briefly summarize:

\begin{itemize}[leftmargin=0.8cm, itemsep=2pt]
    \item the internship organization;
    \item the objectives of the internship;
    \item the main tasks performed;
    \item the main results, knowledge, and skills obtained.
\end{itemize}

The Abstract should be concise and should not reproduce the full content
of the report chapter by chapter.

\vspace{0.3cm}

% ================================================================
% XVII. REFERENCES AND CITATIONS
% ================================================================

\noindent
\textbf{XVII.\quad REFERENCES AND CITATIONS}

\vspace{0.15cm}

All books, journal papers, conference papers, standards, websites,
datasheets, manuals, technical documents, figures, tables, or other
external sources used in preparing the report shall be properly cited.

References shall be managed through the bibliography system provided with
the official template.

Bibliographic entries shall normally be stored in:

\begin{verbatim}
References.bib
\end{verbatim}

References shall be cited in the report using:

\begin{verbatim}
\cite{reference-key}
\end{verbatim}

Students shall not manually type reference numbers such as ``[1]'',
``[2]'', etc.

Titles of foreign-language publications shall remain in their original
language in the bibliography.

All borrowed ideas, data, figures, tables, equations, methods, or results
shall be properly acknowledged.

\vspace{0.3cm}

% ================================================================
% XVIII. CROSS-REFERENCING
% ================================================================

\noindent
\textbf{XVIII.\quad CROSS-REFERENCING}

\vspace{0.15cm}

Numbered elements shall be cross-referenced using the standard \LaTeX{}
mechanism.

Use the following commands where appropriate:

\begin{verbatim}
\label{...}
\ref{...}
\eqref{...}
\end{verbatim}

For example:

\begin{verbatim}
As shown in Figure~\ref{fig:system}, ...
\end{verbatim}

Students shall not manually type numbers such as ``Figure 2.3'' when
those numbers can be generated automatically.

\vspace{0.3cm}

% ================================================================
% XIX. PAGE NUMBERING
% ================================================================

\noindent
\textbf{XIX.\quad PAGE NUMBERING}

\vspace{0.15cm}

Cover pages and official evaluation forms at the beginning of the report
shall not display page numbers.

The main body of the report shall use Arabic page numbers beginning from
Chapter 1.

Page numbering shall be generated automatically by the official template.
Students shall not manually enter page numbers.

\vspace{0.3cm}

% ================================================================
% XX. PROJECT FILE ORGANIZATION
% ================================================================

\noindent
\textbf{XX.\quad PROJECT FILE ORGANIZATION}

\vspace{0.15cm}

Students should maintain a clear and organized \LaTeX{} project
structure. A recommended structure is:

\begin{verbatim}
InternshipReport/
|
|-- main.tex
|-- References.bib
|
|-- figures/
|   |-- company.pdf
|   |-- system.png
|   `-- ...
|
`-- chapters/
    |-- chapter1.tex
    |-- chapter2.tex
    `-- chapter3.tex
\end{verbatim}

Figures should normally be stored in the \verb|figures| directory
specified by the official template.

File and directory names should preferably contain Latin letters,
numbers, hyphens, or underscores.

Recommended examples include:

\begin{verbatim}
company_structure.pdf
motor_control.png
chapter_01.tex
\end{verbatim}

Spaces and unnecessary special characters in file names should be
avoided where possible.

\vspace{0.3cm}

% ================================================================
% XXI. DOCUMENT COMPILATION AND FINAL CHECK
% ================================================================

\noindent
\textbf{XXI.\quad DOCUMENT COMPILATION AND FINAL CHECK}

\vspace{0.15cm}

Before submission, the complete \LaTeX{} project shall compile
successfully and produce the final PDF report.

The final document shall not contain unresolved references such as:

\begin{verbatim}
??
[?]
\end{verbatim}

Before generating the final PDF, students shall check:

\begin{itemize}[leftmargin=0.8cm, itemsep=2pt]
    \item the Table of Contents;
    \item page numbering;
    \item chapter, section, and subsection numbering;
    \item figure numbering;
    \item table numbering;
    \item equation numbering;
    \item cross-references;
    \item citations;
    \item bibliography;
    \item font display;
    \item mathematical symbols;
    \item image quality;
    \item completeness of all required forms.
\end{itemize}

The document may be compiled using TeX Live, MiKTeX, or another
compatible \LaTeX{} environment. Regardless of the compilation
environment, the resulting PDF shall conform to the official Faculty
template.

\vspace{0.3cm}

% ================================================================
% XXII. CONTENT THAT STUDENTS MAY EDIT
% ================================================================

\noindent
\textbf{XXII.\quad CONTENT THAT STUDENTS MAY EDIT}

\vspace{0.15cm}

Students may edit or replace:

\begin{itemize}[leftmargin=0.8cm, itemsep=2pt]
    \item student name;
    \item class;
    \item internship organization;
    \item supervisor name;
    \item Abstract;
    \item List of Symbols and Abbreviations;
    \item chapter titles and chapter content;
    \item figures and photographs;
    \item tables;
    \item equations;
    \item citations and references.
\end{itemize}

Students shall not modify the main document formatting settings in the
preamble, including the document class, page geometry, line spacing,
heading formatting, table-of-contents formatting, font configuration,
page layout, or cover-page macros unless specifically instructed.

\vspace{0.3cm}

% ================================================================
% XXIII. FINAL SUBMISSION REQUIREMENTS
% ================================================================

\noindent
\textbf{XXIII.\quad FINAL SUBMISSION REQUIREMENTS}

\vspace{0.15cm}

Before submission, the final report shall:

\begin{itemize}[leftmargin=0.8cm, itemsep=2pt]
    \item compile without errors that affect the report content or presentation;
    \item contain all required forms and report sections;
    \item contain correct student, class, internship organization, and supervisor information;
    \item contain an updated Table of Contents, List of Figures, and List of Tables;
    \item contain no unused example content from the template;
    \item contain no placeholder text such as \textit{Title of Chapter 1},
    \textit{Section 1.1}, or other sample entries;
    \item contain all required citations and references;
    \item contain no unresolved cross-references.
\end{itemize}

If the Faculty requires submission of the \LaTeX{} source, the complete
project shall include all \verb|.tex| files, the bibliography database,
figures, and any other files required to reproduce the submitted PDF.

\vspace{0.4cm}

\noindent
\textbf{Important:}
Before submitting the final internship report, students shall remove or
comment out the following command in the main document:

\begin{verbatim}
\input{Regulations.tex}
\end{verbatim}

For example:

\begin{verbatim}
% \input{Regulations.tex}
\end{verbatim}

This removes the guidance pages from the final submitted report while
preserving the regulations file in the \LaTeX{} project for reference.

} % End 13pt content

\clearpage
|.
Students shall remove or comment out this command before submitting the
final internship report.

\vspace{0.4cm}

% ================================================================
% I. GENERAL REQUIREMENTS
% ================================================================

\noindent
\textbf{I.\quad GENERAL REQUIREMENTS}

\vspace{0.15cm}

The Graduation Internship Report shall be prepared using the official
\LaTeX{} template provided by the Faculty. Students shall preserve the
predefined page layout, typography, line spacing, heading styles,
numbering system, figure and table formatting, and bibliography settings
unless otherwise instructed by the Faculty or the supervisor.

The report shall be written in English.

Students are responsible for the academic and technical content of the
report. The official \LaTeX{} template is responsible for standardized
document formatting.

\vspace{0.3cm}

% ================================================================
% II. REPORT STRUCTURE
% ================================================================

\noindent
\textbf{II.\quad REPORT STRUCTURE}

\vspace{0.15cm}

The report shall normally be arranged in the following order:

\begin{enumerate}[label=\arabic*., leftmargin=1.2cm, itemsep=2pt]
    \item Outer Cover;
    \item Comments of the Internship Host Organization;
    \item Supervisor's Evaluation and Graduation Internship Defense Results;
    \item Table of Contents;
    \item List of Symbols and Abbreviations, if applicable;
    \item List of Figures;
    \item List of Tables;
    \item Abstract;
    \item Main Chapters of the Internship Report;
    \item References;
    \item Appendices, if applicable.
\end{enumerate}

The preliminary pages shall not be numbered as chapters. The main report
content shall be organized into consecutively numbered chapters beginning
with Chapter 1.

\vspace{0.3cm}

% ================================================================
% III. PAGE SIZE AND MARGINS
% ================================================================

\noindent
\textbf{III.\quad PAGE SIZE AND MARGINS}

\vspace{0.15cm}

The report shall be formatted on A4 paper with the following margins:

\begin{center}
\begin{tabular}{ll}
    Left margin:   & 3.5 cm \\
    Right margin:  & 2.0 cm \\
    Top margin:    & 2.5 cm \\
    Bottom margin: & 2.5 cm
\end{tabular}
\end{center}

These settings are predefined in the official template and shall not be
modified by students.

\vspace{0.3cm}

% ================================================================
% IV. FONT, FONT SIZE, AND LINE SPACING
% ================================================================

\noindent
\textbf{IV.\quad FONT, FONT SIZE, AND LINE SPACING}

\vspace{0.15cm}

The main text shall use a 13 pt font size with 1.3-line spacing.
A Times-family typeface shall be used in accordance with the official
template.

The text shall be fully justified.

Students shall not use repeated spaces, repeated manual line breaks, or
multiple blank lines to position text. Paragraph and heading spacing shall
be controlled by the template.

\vspace{0.3cm}

% ================================================================
% V. OUTER COVER
% ================================================================

\noindent
\textbf{V.\quad OUTER COVER}

\vspace{0.15cm}

The outer cover shall follow the official format provided in the template
and shall contain the following information:

\begin{itemize}[leftmargin=0.8cm, itemsep=2pt]
    \item Thai Nguyen University of Technology;
    \item Faculty of International Training;
    \item Engineering and Technology Department;
    \item University logo and Faculty logo;
    \item the title \textbf{INTERNSHIP REPORT};
    \item Internship Organization;
    \item Author;
    \item Class;
    \item Supervisor;
    \item place and year of completion.
\end{itemize}

Students shall only replace the required information fields and shall not
modify the cover-page layout unless instructed otherwise.

\vspace{0.3cm}

% ================================================================
% VI. COMMENTS OF THE INTERNSHIP HOST ORGANIZATION
% ================================================================

\noindent
\textbf{VI.\quad COMMENTS OF THE INTERNSHIP HOST ORGANIZATION}

\vspace{0.15cm}

The page entitled \textbf{COMMENTS OF THE INTERNSHIP HOST ORGANIZATION}
is reserved for the internship host organization.

It includes:

\begin{itemize}[leftmargin=0.8cm, itemsep=2pt]
    \item the student's full name;
    \item class;
    \item comments from the host organization;
    \item score;
    \item date;
    \item signature and official seal.
\end{itemize}

Students shall not alter or complete official evaluation fields unless
specifically authorized.

\vspace{0.3cm}

% ================================================================
% VII. SUPERVISOR'S EVALUATION AND DEFENSE RESULTS
% ================================================================

\noindent
\textbf{VII.\quad SUPERVISOR'S EVALUATION AND DEFENSE RESULTS}

\vspace{0.15cm}

The section entitled
\textbf{SUPERVISOR'S EVALUATION AND GRADUATION INTERNSHIP DEFENSE RESULTS}
is reserved for the supervisor and defense committee.

The evaluation may include:

\begin{itemize}[leftmargin=0.8cm, itemsep=2pt]
    \item Internship Progress and Student's Attitude;
    \item Report Content;
    \item Report Presentation;
    \item Other Comments;
    \item Supervisor's Opinion;
    \item Defense Result;
    \item Score.
\end{itemize}

Students shall not enter evaluation scores or defense results before they
are officially determined.

\vspace{0.3cm}

% ================================================================
% VIII. CHAPTERS, SECTIONS, AND SUBSECTIONS
% ================================================================

\noindent
\textbf{VIII.\quad CHAPTERS, SECTIONS, AND SUBSECTIONS}

\vspace{0.15cm}

The report structure shall be generated using the standard \LaTeX{}
commands:

\begin{verbatim}
\chapter{...}
\section{...}
\subsection{...}
\end{verbatim}

Chapter, section, and subsection numbers shall be generated automatically
by \LaTeX{}.

Students shall not manually type chapter or section numbers and shall not
reproduce heading styles using commands such as \verb|\textbf|,
\verb|\fontsize|, or \verb|\centering| instead of the appropriate
structural command.

Excessive subdivision should be avoided. In most cases, the hierarchy
Chapter--Section--Subsection is sufficient.

\vspace{0.3cm}

% ================================================================
% IX. ORGANIZATION OF THE MAIN CONTENT
% ================================================================

\noindent
\textbf{IX.\quad ORGANIZATION OF THE MAIN CONTENT}

\vspace{0.15cm}

The specific content of the report shall depend on the internship
organization and the assigned tasks. A typical structure may include:

\begin{itemize}[leftmargin=0.8cm, itemsep=2pt]
    \item \textbf{Chapter 1. Overview of the Internship Organization};
    \item \textbf{Chapter 2. Internship Activities and Assigned Tasks};
    \item \textbf{Chapter 3. Results, Experience, and Evaluation}.
\end{itemize}

These chapter titles are recommendations only and may be adapted to the
actual internship content.

The report shall reflect work actually performed, observed, analyzed, or
learned during the internship. Unrelated theoretical material should not
be included merely to increase report length.

\vspace{0.3cm}

% ================================================================
% X. TABLE OF CONTENTS
% ================================================================

\noindent
\textbf{X.\quad TABLE OF CONTENTS}

\vspace{0.15cm}

The Table of Contents shall be generated automatically using:

\begin{verbatim}
\tableofcontents
\end{verbatim}

Students shall not create or manually edit page numbers in the Table of
Contents.

The document shall be compiled sufficiently to ensure that all page
numbers and entries are correctly updated.

\vspace{0.3cm}

% ================================================================
% XI. LIST OF SYMBOLS AND ABBREVIATIONS
% ================================================================

\noindent
\textbf{XI.\quad LIST OF SYMBOLS AND ABBREVIATIONS}

\vspace{0.15cm}

Symbols and abbreviations that are used repeatedly in the report shall be
listed in the \textbf{List of Symbols and Abbreviations}.

The list may contain the following columns:

\begin{center}
\begin{tabular}{llll}
    No. & Symbol / Abbreviation & Unit & Definition
\end{tabular}
\end{center}

Only symbols and abbreviations actually used in the final report shall be
retained. Example entries contained in the template shall be removed if
they are not relevant to the student's report.

Each abbreviation should normally be defined when it first appears in
the main text.

\vspace{0.3cm}

% ================================================================
% XII. FIGURES, DIAGRAMS, AND PHOTOGRAPHS
% ================================================================

\noindent
\textbf{XII.\quad FIGURES, DIAGRAMS, AND PHOTOGRAPHS}

\vspace{0.15cm}

Figures shall be clear, legible, and directly relevant to the report
content.

Figures shall be inserted using the standard \verb|figure| environment.

Example:

\begin{verbatim}
\begin{figure}[H]
    \centering
    \includegraphics[width=0.8\textwidth]
                    {figures/company_structure.pdf}
    \caption{Organizational structure of the company}
    \label{fig:company-structure}
\end{figure}
\end{verbatim}

Figure captions shall be placed below the figure.

Each referenced figure shall have a unique \verb|\label| and shall be
referenced in the text using, for example:

\begin{verbatim}
Figure~\ref{fig:company-structure}
\end{verbatim}

Students shall not manually type figure numbers.

Figures reproduced or adapted from external sources shall include an
appropriate citation.

Photographs taken during the internship shall be relevant and sufficiently
clear. Students shall not include confidential, restricted, or proprietary
information without permission from the internship organization.

Vector graphics or high-resolution images are recommended whenever
possible.

\vspace{0.3cm}

% ================================================================
% XIII. TABLES
% ================================================================

\noindent
\textbf{XIII.\quad TABLES}

\vspace{0.15cm}

Tables shall be created using suitable \LaTeX{} environments such as
\verb|table|, \verb|tabular|, \verb|tabularx|, or \verb|longtable|.

Table captions shall be placed above the table.

Example:

\begin{verbatim}
\begin{table}[H]
    \centering
    \caption{Internship schedule}
    \label{tab:internship-schedule}

    \begin{tabular}{lll}
        \toprule
        Week & Task & Result \\
        \midrule
        1 & Orientation & Completed \\
        2 & System study & Completed \\
        \bottomrule
    \end{tabular}
\end{table}
\end{verbatim}

Tables shall be referenced using:

\begin{verbatim}
Table~\ref{tab:internship-schedule}
\end{verbatim}

Students shall not manually type table numbers.

Tables should be created directly in \LaTeX{} rather than inserted as
screenshots whenever practical.

\vspace{0.3cm}

% ================================================================
% XIV. EQUATIONS AND TECHNICAL QUANTITIES
% ================================================================

\noindent
\textbf{XIV.\quad EQUATIONS AND TECHNICAL QUANTITIES}

\vspace{0.15cm}

Mathematical expressions shall be typeset directly using \LaTeX{}.
Images or screenshots of equations shall not be used.

Short expressions may be written in inline mathematical mode, for example:

\begin{verbatim}
$R = 0.7\,\Omega$
\end{verbatim}

Important equations shall use a numbered equation environment.

Example:

\begin{verbatim}
\begin{equation}
    P = UI
    \label{eq:power}
\end{equation}
\end{verbatim}

Equations shall be referenced using:

\begin{verbatim}
Equation~\eqref{eq:power}
\end{verbatim}

Equation numbers shall not be manually entered.

SI units shall be used wherever applicable. The \verb|siunitx| package
should be used where appropriate.

Examples:

\begin{verbatim}
\SI{220}{\volt}
\SI{50}{\hertz}
\SI{2.5}{\ampere}
\end{verbatim}

Notation and units shall be consistent throughout the report.

\vspace{0.3cm}

% ================================================================
% XV. SOURCE CODE
% ================================================================

\noindent
\textbf{XV.\quad SOURCE CODE}

\vspace{0.15cm}

If the internship involves programming, PLCs, microcontrollers,
simulation, or software development, relevant source code may be included
in the report.

Source code shall be presented using the \verb|listings| environment or
another environment defined by the official template.

Example:

\begin{verbatim}
\begin{lstlisting}[language=st]
IF Start = TRUE THEN
    Motor := TRUE;
ELSE
    Motor := FALSE;
END_IF;
\end{lstlisting}
\end{verbatim}

Only code necessary for explaining the work should appear in the main
report. Long source-code listings should be moved to an appendix or
submitted as supplementary files if required.

Source code should not normally be inserted as screenshots.

\vspace{0.3cm}

% ================================================================
% XVI. ABSTRACT
% ================================================================

\noindent
\textbf{XVI.\quad ABSTRACT}

\vspace{0.15cm}

The Abstract shall briefly summarize:

\begin{itemize}[leftmargin=0.8cm, itemsep=2pt]
    \item the internship organization;
    \item the objectives of the internship;
    \item the main tasks performed;
    \item the main results, knowledge, and skills obtained.
\end{itemize}

The Abstract should be concise and should not reproduce the full content
of the report chapter by chapter.

\vspace{0.3cm}

% ================================================================
% XVII. REFERENCES AND CITATIONS
% ================================================================

\noindent
\textbf{XVII.\quad REFERENCES AND CITATIONS}

\vspace{0.15cm}

All books, journal papers, conference papers, standards, websites,
datasheets, manuals, technical documents, figures, tables, or other
external sources used in preparing the report shall be properly cited.

References shall be managed through the bibliography system provided with
the official template.

Bibliographic entries shall normally be stored in:

\begin{verbatim}
References.bib
\end{verbatim}

References shall be cited in the report using:

\begin{verbatim}
\cite{reference-key}
\end{verbatim}

Students shall not manually type reference numbers such as ``[1]'',
``[2]'', etc.

Titles of foreign-language publications shall remain in their original
language in the bibliography.

All borrowed ideas, data, figures, tables, equations, methods, or results
shall be properly acknowledged.

\vspace{0.3cm}

% ================================================================
% XVIII. CROSS-REFERENCING
% ================================================================

\noindent
\textbf{XVIII.\quad CROSS-REFERENCING}

\vspace{0.15cm}

Numbered elements shall be cross-referenced using the standard \LaTeX{}
mechanism.

Use the following commands where appropriate:

\begin{verbatim}
\label{...}
\ref{...}
\eqref{...}
\end{verbatim}

For example:

\begin{verbatim}
As shown in Figure~\ref{fig:system}, ...
\end{verbatim}

Students shall not manually type numbers such as ``Figure 2.3'' when
those numbers can be generated automatically.

\vspace{0.3cm}

% ================================================================
% XIX. PAGE NUMBERING
% ================================================================

\noindent
\textbf{XIX.\quad PAGE NUMBERING}

\vspace{0.15cm}

Cover pages and official evaluation forms at the beginning of the report
shall not display page numbers.

The main body of the report shall use Arabic page numbers beginning from
Chapter 1.

Page numbering shall be generated automatically by the official template.
Students shall not manually enter page numbers.

\vspace{0.3cm}

% ================================================================
% XX. PROJECT FILE ORGANIZATION
% ================================================================

\noindent
\textbf{XX.\quad PROJECT FILE ORGANIZATION}

\vspace{0.15cm}

Students should maintain a clear and organized \LaTeX{} project
structure. A recommended structure is:

\begin{verbatim}
InternshipReport/
|
|-- main.tex
|-- References.bib
|
|-- figures/
|   |-- company.pdf
|   |-- system.png
|   `-- ...
|
`-- chapters/
    |-- chapter1.tex
    |-- chapter2.tex
    `-- chapter3.tex
\end{verbatim}

Figures should normally be stored in the \verb|figures| directory
specified by the official template.

File and directory names should preferably contain Latin letters,
numbers, hyphens, or underscores.

Recommended examples include:

\begin{verbatim}
company_structure.pdf
motor_control.png
chapter_01.tex
\end{verbatim}

Spaces and unnecessary special characters in file names should be
avoided where possible.

\vspace{0.3cm}

% ================================================================
% XXI. DOCUMENT COMPILATION AND FINAL CHECK
% ================================================================

\noindent
\textbf{XXI.\quad DOCUMENT COMPILATION AND FINAL CHECK}

\vspace{0.15cm}

Before submission, the complete \LaTeX{} project shall compile
successfully and produce the final PDF report.

The final document shall not contain unresolved references such as:

\begin{verbatim}
??
[?]
\end{verbatim}

Before generating the final PDF, students shall check:

\begin{itemize}[leftmargin=0.8cm, itemsep=2pt]
    \item the Table of Contents;
    \item page numbering;
    \item chapter, section, and subsection numbering;
    \item figure numbering;
    \item table numbering;
    \item equation numbering;
    \item cross-references;
    \item citations;
    \item bibliography;
    \item font display;
    \item mathematical symbols;
    \item image quality;
    \item completeness of all required forms.
\end{itemize}

The document may be compiled using TeX Live, MiKTeX, or another
compatible \LaTeX{} environment. Regardless of the compilation
environment, the resulting PDF shall conform to the official Faculty
template.

\vspace{0.3cm}

% ================================================================
% XXII. CONTENT THAT STUDENTS MAY EDIT
% ================================================================

\noindent
\textbf{XXII.\quad CONTENT THAT STUDENTS MAY EDIT}

\vspace{0.15cm}

Students may edit or replace:

\begin{itemize}[leftmargin=0.8cm, itemsep=2pt]
    \item student name;
    \item class;
    \item internship organization;
    \item supervisor name;
    \item Abstract;
    \item List of Symbols and Abbreviations;
    \item chapter titles and chapter content;
    \item figures and photographs;
    \item tables;
    \item equations;
    \item citations and references.
\end{itemize}

Students shall not modify the main document formatting settings in the
preamble, including the document class, page geometry, line spacing,
heading formatting, table-of-contents formatting, font configuration,
page layout, or cover-page macros unless specifically instructed.

\vspace{0.3cm}

% ================================================================
% XXIII. FINAL SUBMISSION REQUIREMENTS
% ================================================================

\noindent
\textbf{XXIII.\quad FINAL SUBMISSION REQUIREMENTS}

\vspace{0.15cm}

Before submission, the final report shall:

\begin{itemize}[leftmargin=0.8cm, itemsep=2pt]
    \item compile without errors that affect the report content or presentation;
    \item contain all required forms and report sections;
    \item contain correct student, class, internship organization, and supervisor information;
    \item contain an updated Table of Contents, List of Figures, and List of Tables;
    \item contain no unused example content from the template;
    \item contain no placeholder text such as \textit{Title of Chapter 1},
    \textit{Section 1.1}, or other sample entries;
    \item contain all required citations and references;
    \item contain no unresolved cross-references.
\end{itemize}

If the Faculty requires submission of the \LaTeX{} source, the complete
project shall include all \verb|.tex| files, the bibliography database,
figures, and any other files required to reproduce the submitted PDF.

\vspace{0.4cm}

\noindent
\textbf{Important:}
Before submitting the final internship report, students shall remove or
comment out the following command in the main document:

\begin{verbatim}

% =================================================================
% REGULATIONS FOR PREPARING THE GRADUATION INTERNSHIP REPORT
% USING LATEX
% =================================================================
%
% This file is intended to be included from the main document with:
%
%     \input{Regulations.tex}
%
% IMPORTANT:
% This section is provided for guidance only. Students shall remove
% or comment out the \input{Regulations.tex} command before submitting
% the final internship report.
%
% =================================================================

\clearpage
\thispagestyle{empty}

\begin{center}
    {\fontsize{16pt}{20pt}\selectfont\bfseries
    SPECIFIC REGULATIONS FOR PREPARING\\
    GRADUATION INTERNSHIP REPORTS USING \LaTeX}
\end{center}

\vspace{0.5cm}

{\fontsize{13pt}{17pt}\selectfont
\setstretch{1.3}
\justifying

\noindent
\textbf{Note:}
This section is provided for guidance only and is included in the main
document using the command \verb|\input{Regulations.tex}|.
Students shall remove or comment out this command before submitting the
final internship report.

\vspace{0.4cm}

% ================================================================
% I. GENERAL REQUIREMENTS
% ================================================================

\noindent
\textbf{I.\quad GENERAL REQUIREMENTS}

\vspace{0.15cm}

The Graduation Internship Report shall be prepared using the official
\LaTeX{} template provided by the Faculty. Students shall preserve the
predefined page layout, typography, line spacing, heading styles,
numbering system, figure and table formatting, and bibliography settings
unless otherwise instructed by the Faculty or the supervisor.

The report shall be written in English.

Students are responsible for the academic and technical content of the
report. The official \LaTeX{} template is responsible for standardized
document formatting.

\vspace{0.3cm}

% ================================================================
% II. REPORT STRUCTURE
% ================================================================

\noindent
\textbf{II.\quad REPORT STRUCTURE}

\vspace{0.15cm}

The report shall normally be arranged in the following order:

\begin{enumerate}[label=\arabic*., leftmargin=1.2cm, itemsep=2pt]
    \item Outer Cover;
    \item Comments of the Internship Host Organization;
    \item Supervisor's Evaluation and Graduation Internship Defense Results;
    \item Table of Contents;
    \item List of Symbols and Abbreviations, if applicable;
    \item List of Figures;
    \item List of Tables;
    \item Abstract;
    \item Main Chapters of the Internship Report;
    \item References;
    \item Appendices, if applicable.
\end{enumerate}

The preliminary pages shall not be numbered as chapters. The main report
content shall be organized into consecutively numbered chapters beginning
with Chapter 1.

\vspace{0.3cm}

% ================================================================
% III. PAGE SIZE AND MARGINS
% ================================================================

\noindent
\textbf{III.\quad PAGE SIZE AND MARGINS}

\vspace{0.15cm}

The report shall be formatted on A4 paper with the following margins:

\begin{center}
\begin{tabular}{ll}
    Left margin:   & 3.5 cm \\
    Right margin:  & 2.0 cm \\
    Top margin:    & 2.5 cm \\
    Bottom margin: & 2.5 cm
\end{tabular}
\end{center}

These settings are predefined in the official template and shall not be
modified by students.

\vspace{0.3cm}

% ================================================================
% IV. FONT, FONT SIZE, AND LINE SPACING
% ================================================================

\noindent
\textbf{IV.\quad FONT, FONT SIZE, AND LINE SPACING}

\vspace{0.15cm}

The main text shall use a 13 pt font size with 1.3-line spacing.
A Times-family typeface shall be used in accordance with the official
template.

The text shall be fully justified.

Students shall not use repeated spaces, repeated manual line breaks, or
multiple blank lines to position text. Paragraph and heading spacing shall
be controlled by the template.

\vspace{0.3cm}

% ================================================================
% V. OUTER COVER
% ================================================================

\noindent
\textbf{V.\quad OUTER COVER}

\vspace{0.15cm}

The outer cover shall follow the official format provided in the template
and shall contain the following information:

\begin{itemize}[leftmargin=0.8cm, itemsep=2pt]
    \item Thai Nguyen University of Technology;
    \item Faculty of International Training;
    \item Engineering and Technology Department;
    \item University logo and Faculty logo;
    \item the title \textbf{INTERNSHIP REPORT};
    \item Internship Organization;
    \item Author;
    \item Class;
    \item Supervisor;
    \item place and year of completion.
\end{itemize}

Students shall only replace the required information fields and shall not
modify the cover-page layout unless instructed otherwise.

\vspace{0.3cm}

% ================================================================
% VI. COMMENTS OF THE INTERNSHIP HOST ORGANIZATION
% ================================================================

\noindent
\textbf{VI.\quad COMMENTS OF THE INTERNSHIP HOST ORGANIZATION}

\vspace{0.15cm}

The page entitled \textbf{COMMENTS OF THE INTERNSHIP HOST ORGANIZATION}
is reserved for the internship host organization.

It includes:

\begin{itemize}[leftmargin=0.8cm, itemsep=2pt]
    \item the student's full name;
    \item class;
    \item comments from the host organization;
    \item score;
    \item date;
    \item signature and official seal.
\end{itemize}

Students shall not alter or complete official evaluation fields unless
specifically authorized.

\vspace{0.3cm}

% ================================================================
% VII. SUPERVISOR'S EVALUATION AND DEFENSE RESULTS
% ================================================================

\noindent
\textbf{VII.\quad SUPERVISOR'S EVALUATION AND DEFENSE RESULTS}

\vspace{0.15cm}

The section entitled
\textbf{SUPERVISOR'S EVALUATION AND GRADUATION INTERNSHIP DEFENSE RESULTS}
is reserved for the supervisor and defense committee.

The evaluation may include:

\begin{itemize}[leftmargin=0.8cm, itemsep=2pt]
    \item Internship Progress and Student's Attitude;
    \item Report Content;
    \item Report Presentation;
    \item Other Comments;
    \item Supervisor's Opinion;
    \item Defense Result;
    \item Score.
\end{itemize}

Students shall not enter evaluation scores or defense results before they
are officially determined.

\vspace{0.3cm}

% ================================================================
% VIII. CHAPTERS, SECTIONS, AND SUBSECTIONS
% ================================================================

\noindent
\textbf{VIII.\quad CHAPTERS, SECTIONS, AND SUBSECTIONS}

\vspace{0.15cm}

The report structure shall be generated using the standard \LaTeX{}
commands:

\begin{verbatim}
\chapter{...}
\section{...}
\subsection{...}
\end{verbatim}

Chapter, section, and subsection numbers shall be generated automatically
by \LaTeX{}.

Students shall not manually type chapter or section numbers and shall not
reproduce heading styles using commands such as \verb|\textbf|,
\verb|\fontsize|, or \verb|\centering| instead of the appropriate
structural command.

Excessive subdivision should be avoided. In most cases, the hierarchy
Chapter--Section--Subsection is sufficient.

\vspace{0.3cm}

% ================================================================
% IX. ORGANIZATION OF THE MAIN CONTENT
% ================================================================

\noindent
\textbf{IX.\quad ORGANIZATION OF THE MAIN CONTENT}

\vspace{0.15cm}

The specific content of the report shall depend on the internship
organization and the assigned tasks. A typical structure may include:

\begin{itemize}[leftmargin=0.8cm, itemsep=2pt]
    \item \textbf{Chapter 1. Overview of the Internship Organization};
    \item \textbf{Chapter 2. Internship Activities and Assigned Tasks};
    \item \textbf{Chapter 3. Results, Experience, and Evaluation}.
\end{itemize}

These chapter titles are recommendations only and may be adapted to the
actual internship content.

The report shall reflect work actually performed, observed, analyzed, or
learned during the internship. Unrelated theoretical material should not
be included merely to increase report length.

\vspace{0.3cm}

% ================================================================
% X. TABLE OF CONTENTS
% ================================================================

\noindent
\textbf{X.\quad TABLE OF CONTENTS}

\vspace{0.15cm}

The Table of Contents shall be generated automatically using:

\begin{verbatim}
\tableofcontents
\end{verbatim}

Students shall not create or manually edit page numbers in the Table of
Contents.

The document shall be compiled sufficiently to ensure that all page
numbers and entries are correctly updated.

\vspace{0.3cm}

% ================================================================
% XI. LIST OF SYMBOLS AND ABBREVIATIONS
% ================================================================

\noindent
\textbf{XI.\quad LIST OF SYMBOLS AND ABBREVIATIONS}

\vspace{0.15cm}

Symbols and abbreviations that are used repeatedly in the report shall be
listed in the \textbf{List of Symbols and Abbreviations}.

The list may contain the following columns:

\begin{center}
\begin{tabular}{llll}
    No. & Symbol / Abbreviation & Unit & Definition
\end{tabular}
\end{center}

Only symbols and abbreviations actually used in the final report shall be
retained. Example entries contained in the template shall be removed if
they are not relevant to the student's report.

Each abbreviation should normally be defined when it first appears in
the main text.

\vspace{0.3cm}

% ================================================================
% XII. FIGURES, DIAGRAMS, AND PHOTOGRAPHS
% ================================================================

\noindent
\textbf{XII.\quad FIGURES, DIAGRAMS, AND PHOTOGRAPHS}

\vspace{0.15cm}

Figures shall be clear, legible, and directly relevant to the report
content.

Figures shall be inserted using the standard \verb|figure| environment.

Example:

\begin{verbatim}
\begin{figure}[H]
    \centering
    \includegraphics[width=0.8\textwidth]
                    {figures/company_structure.pdf}
    \caption{Organizational structure of the company}
    \label{fig:company-structure}
\end{figure}
\end{verbatim}

Figure captions shall be placed below the figure.

Each referenced figure shall have a unique \verb|\label| and shall be
referenced in the text using, for example:

\begin{verbatim}
Figure~\ref{fig:company-structure}
\end{verbatim}

Students shall not manually type figure numbers.

Figures reproduced or adapted from external sources shall include an
appropriate citation.

Photographs taken during the internship shall be relevant and sufficiently
clear. Students shall not include confidential, restricted, or proprietary
information without permission from the internship organization.

Vector graphics or high-resolution images are recommended whenever
possible.

\vspace{0.3cm}

% ================================================================
% XIII. TABLES
% ================================================================

\noindent
\textbf{XIII.\quad TABLES}

\vspace{0.15cm}

Tables shall be created using suitable \LaTeX{} environments such as
\verb|table|, \verb|tabular|, \verb|tabularx|, or \verb|longtable|.

Table captions shall be placed above the table.

Example:

\begin{verbatim}
\begin{table}[H]
    \centering
    \caption{Internship schedule}
    \label{tab:internship-schedule}

    \begin{tabular}{lll}
        \toprule
        Week & Task & Result \\
        \midrule
        1 & Orientation & Completed \\
        2 & System study & Completed \\
        \bottomrule
    \end{tabular}
\end{table}
\end{verbatim}

Tables shall be referenced using:

\begin{verbatim}
Table~\ref{tab:internship-schedule}
\end{verbatim}

Students shall not manually type table numbers.

Tables should be created directly in \LaTeX{} rather than inserted as
screenshots whenever practical.

\vspace{0.3cm}

% ================================================================
% XIV. EQUATIONS AND TECHNICAL QUANTITIES
% ================================================================

\noindent
\textbf{XIV.\quad EQUATIONS AND TECHNICAL QUANTITIES}

\vspace{0.15cm}

Mathematical expressions shall be typeset directly using \LaTeX{}.
Images or screenshots of equations shall not be used.

Short expressions may be written in inline mathematical mode, for example:

\begin{verbatim}
$R = 0.7\,\Omega$
\end{verbatim}

Important equations shall use a numbered equation environment.

Example:

\begin{verbatim}
\begin{equation}
    P = UI
    \label{eq:power}
\end{equation}
\end{verbatim}

Equations shall be referenced using:

\begin{verbatim}
Equation~\eqref{eq:power}
\end{verbatim}

Equation numbers shall not be manually entered.

SI units shall be used wherever applicable. The \verb|siunitx| package
should be used where appropriate.

Examples:

\begin{verbatim}
\SI{220}{\volt}
\SI{50}{\hertz}
\SI{2.5}{\ampere}
\end{verbatim}

Notation and units shall be consistent throughout the report.

\vspace{0.3cm}

% ================================================================
% XV. SOURCE CODE
% ================================================================

\noindent
\textbf{XV.\quad SOURCE CODE}

\vspace{0.15cm}

If the internship involves programming, PLCs, microcontrollers,
simulation, or software development, relevant source code may be included
in the report.

Source code shall be presented using the \verb|listings| environment or
another environment defined by the official template.

Example:

\begin{verbatim}
\begin{lstlisting}[language=st]
IF Start = TRUE THEN
    Motor := TRUE;
ELSE
    Motor := FALSE;
END_IF;
\end{lstlisting}
\end{verbatim}

Only code necessary for explaining the work should appear in the main
report. Long source-code listings should be moved to an appendix or
submitted as supplementary files if required.

Source code should not normally be inserted as screenshots.

\vspace{0.3cm}

% ================================================================
% XVI. ABSTRACT
% ================================================================

\noindent
\textbf{XVI.\quad ABSTRACT}

\vspace{0.15cm}

The Abstract shall briefly summarize:

\begin{itemize}[leftmargin=0.8cm, itemsep=2pt]
    \item the internship organization;
    \item the objectives of the internship;
    \item the main tasks performed;
    \item the main results, knowledge, and skills obtained.
\end{itemize}

The Abstract should be concise and should not reproduce the full content
of the report chapter by chapter.

\vspace{0.3cm}

% ================================================================
% XVII. REFERENCES AND CITATIONS
% ================================================================

\noindent
\textbf{XVII.\quad REFERENCES AND CITATIONS}

\vspace{0.15cm}

All books, journal papers, conference papers, standards, websites,
datasheets, manuals, technical documents, figures, tables, or other
external sources used in preparing the report shall be properly cited.

References shall be managed through the bibliography system provided with
the official template.

Bibliographic entries shall normally be stored in:

\begin{verbatim}
References.bib
\end{verbatim}

References shall be cited in the report using:

\begin{verbatim}
\cite{reference-key}
\end{verbatim}

Students shall not manually type reference numbers such as ``[1]'',
``[2]'', etc.

Titles of foreign-language publications shall remain in their original
language in the bibliography.

All borrowed ideas, data, figures, tables, equations, methods, or results
shall be properly acknowledged.

\vspace{0.3cm}

% ================================================================
% XVIII. CROSS-REFERENCING
% ================================================================

\noindent
\textbf{XVIII.\quad CROSS-REFERENCING}

\vspace{0.15cm}

Numbered elements shall be cross-referenced using the standard \LaTeX{}
mechanism.

Use the following commands where appropriate:

\begin{verbatim}
\label{...}
\ref{...}
\eqref{...}
\end{verbatim}

For example:

\begin{verbatim}
As shown in Figure~\ref{fig:system}, ...
\end{verbatim}

Students shall not manually type numbers such as ``Figure 2.3'' when
those numbers can be generated automatically.

\vspace{0.3cm}

% ================================================================
% XIX. PAGE NUMBERING
% ================================================================

\noindent
\textbf{XIX.\quad PAGE NUMBERING}

\vspace{0.15cm}

Cover pages and official evaluation forms at the beginning of the report
shall not display page numbers.

The main body of the report shall use Arabic page numbers beginning from
Chapter 1.

Page numbering shall be generated automatically by the official template.
Students shall not manually enter page numbers.

\vspace{0.3cm}

% ================================================================
% XX. PROJECT FILE ORGANIZATION
% ================================================================

\noindent
\textbf{XX.\quad PROJECT FILE ORGANIZATION}

\vspace{0.15cm}

Students should maintain a clear and organized \LaTeX{} project
structure. A recommended structure is:

\begin{verbatim}
InternshipReport/
|
|-- main.tex
|-- References.bib
|
|-- figures/
|   |-- company.pdf
|   |-- system.png
|   `-- ...
|
`-- chapters/
    |-- chapter1.tex
    |-- chapter2.tex
    `-- chapter3.tex
\end{verbatim}

Figures should normally be stored in the \verb|figures| directory
specified by the official template.

File and directory names should preferably contain Latin letters,
numbers, hyphens, or underscores.

Recommended examples include:

\begin{verbatim}
company_structure.pdf
motor_control.png
chapter_01.tex
\end{verbatim}

Spaces and unnecessary special characters in file names should be
avoided where possible.

\vspace{0.3cm}

% ================================================================
% XXI. DOCUMENT COMPILATION AND FINAL CHECK
% ================================================================

\noindent
\textbf{XXI.\quad DOCUMENT COMPILATION AND FINAL CHECK}

\vspace{0.15cm}

Before submission, the complete \LaTeX{} project shall compile
successfully and produce the final PDF report.

The final document shall not contain unresolved references such as:

\begin{verbatim}
??
[?]
\end{verbatim}

Before generating the final PDF, students shall check:

\begin{itemize}[leftmargin=0.8cm, itemsep=2pt]
    \item the Table of Contents;
    \item page numbering;
    \item chapter, section, and subsection numbering;
    \item figure numbering;
    \item table numbering;
    \item equation numbering;
    \item cross-references;
    \item citations;
    \item bibliography;
    \item font display;
    \item mathematical symbols;
    \item image quality;
    \item completeness of all required forms.
\end{itemize}

The document may be compiled using TeX Live, MiKTeX, or another
compatible \LaTeX{} environment. Regardless of the compilation
environment, the resulting PDF shall conform to the official Faculty
template.

\vspace{0.3cm}

% ================================================================
% XXII. CONTENT THAT STUDENTS MAY EDIT
% ================================================================

\noindent
\textbf{XXII.\quad CONTENT THAT STUDENTS MAY EDIT}

\vspace{0.15cm}

Students may edit or replace:

\begin{itemize}[leftmargin=0.8cm, itemsep=2pt]
    \item student name;
    \item class;
    \item internship organization;
    \item supervisor name;
    \item Abstract;
    \item List of Symbols and Abbreviations;
    \item chapter titles and chapter content;
    \item figures and photographs;
    \item tables;
    \item equations;
    \item citations and references.
\end{itemize}

Students shall not modify the main document formatting settings in the
preamble, including the document class, page geometry, line spacing,
heading formatting, table-of-contents formatting, font configuration,
page layout, or cover-page macros unless specifically instructed.

\vspace{0.3cm}

% ================================================================
% XXIII. FINAL SUBMISSION REQUIREMENTS
% ================================================================

\noindent
\textbf{XXIII.\quad FINAL SUBMISSION REQUIREMENTS}

\vspace{0.15cm}

Before submission, the final report shall:

\begin{itemize}[leftmargin=0.8cm, itemsep=2pt]
    \item compile without errors that affect the report content or presentation;
    \item contain all required forms and report sections;
    \item contain correct student, class, internship organization, and supervisor information;
    \item contain an updated Table of Contents, List of Figures, and List of Tables;
    \item contain no unused example content from the template;
    \item contain no placeholder text such as \textit{Title of Chapter 1},
    \textit{Section 1.1}, or other sample entries;
    \item contain all required citations and references;
    \item contain no unresolved cross-references.
\end{itemize}

If the Faculty requires submission of the \LaTeX{} source, the complete
project shall include all \verb|.tex| files, the bibliography database,
figures, and any other files required to reproduce the submitted PDF.

\vspace{0.4cm}

\noindent
\textbf{Important:}
Before submitting the final internship report, students shall remove or
comment out the following command in the main document:

\begin{verbatim}
\input{Regulations.tex}
\end{verbatim}

For example:

\begin{verbatim}
% \input{Regulations.tex}
\end{verbatim}

This removes the guidance pages from the final submitted report while
preserving the regulations file in the \LaTeX{} project for reference.

} % End 13pt content

\clearpage

\end{verbatim}

For example:

\begin{verbatim}
% 
% =================================================================
% REGULATIONS FOR PREPARING THE GRADUATION INTERNSHIP REPORT
% USING LATEX
% =================================================================
%
% This file is intended to be included from the main document with:
%
%     \input{Regulations.tex}
%
% IMPORTANT:
% This section is provided for guidance only. Students shall remove
% or comment out the \input{Regulations.tex} command before submitting
% the final internship report.
%
% =================================================================

\clearpage
\thispagestyle{empty}

\begin{center}
    {\fontsize{16pt}{20pt}\selectfont\bfseries
    SPECIFIC REGULATIONS FOR PREPARING\\
    GRADUATION INTERNSHIP REPORTS USING \LaTeX}
\end{center}

\vspace{0.5cm}

{\fontsize{13pt}{17pt}\selectfont
\setstretch{1.3}
\justifying

\noindent
\textbf{Note:}
This section is provided for guidance only and is included in the main
document using the command \verb|\input{Regulations.tex}|.
Students shall remove or comment out this command before submitting the
final internship report.

\vspace{0.4cm}

% ================================================================
% I. GENERAL REQUIREMENTS
% ================================================================

\noindent
\textbf{I.\quad GENERAL REQUIREMENTS}

\vspace{0.15cm}

The Graduation Internship Report shall be prepared using the official
\LaTeX{} template provided by the Faculty. Students shall preserve the
predefined page layout, typography, line spacing, heading styles,
numbering system, figure and table formatting, and bibliography settings
unless otherwise instructed by the Faculty or the supervisor.

The report shall be written in English.

Students are responsible for the academic and technical content of the
report. The official \LaTeX{} template is responsible for standardized
document formatting.

\vspace{0.3cm}

% ================================================================
% II. REPORT STRUCTURE
% ================================================================

\noindent
\textbf{II.\quad REPORT STRUCTURE}

\vspace{0.15cm}

The report shall normally be arranged in the following order:

\begin{enumerate}[label=\arabic*., leftmargin=1.2cm, itemsep=2pt]
    \item Outer Cover;
    \item Comments of the Internship Host Organization;
    \item Supervisor's Evaluation and Graduation Internship Defense Results;
    \item Table of Contents;
    \item List of Symbols and Abbreviations, if applicable;
    \item List of Figures;
    \item List of Tables;
    \item Abstract;
    \item Main Chapters of the Internship Report;
    \item References;
    \item Appendices, if applicable.
\end{enumerate}

The preliminary pages shall not be numbered as chapters. The main report
content shall be organized into consecutively numbered chapters beginning
with Chapter 1.

\vspace{0.3cm}

% ================================================================
% III. PAGE SIZE AND MARGINS
% ================================================================

\noindent
\textbf{III.\quad PAGE SIZE AND MARGINS}

\vspace{0.15cm}

The report shall be formatted on A4 paper with the following margins:

\begin{center}
\begin{tabular}{ll}
    Left margin:   & 3.5 cm \\
    Right margin:  & 2.0 cm \\
    Top margin:    & 2.5 cm \\
    Bottom margin: & 2.5 cm
\end{tabular}
\end{center}

These settings are predefined in the official template and shall not be
modified by students.

\vspace{0.3cm}

% ================================================================
% IV. FONT, FONT SIZE, AND LINE SPACING
% ================================================================

\noindent
\textbf{IV.\quad FONT, FONT SIZE, AND LINE SPACING}

\vspace{0.15cm}

The main text shall use a 13 pt font size with 1.3-line spacing.
A Times-family typeface shall be used in accordance with the official
template.

The text shall be fully justified.

Students shall not use repeated spaces, repeated manual line breaks, or
multiple blank lines to position text. Paragraph and heading spacing shall
be controlled by the template.

\vspace{0.3cm}

% ================================================================
% V. OUTER COVER
% ================================================================

\noindent
\textbf{V.\quad OUTER COVER}

\vspace{0.15cm}

The outer cover shall follow the official format provided in the template
and shall contain the following information:

\begin{itemize}[leftmargin=0.8cm, itemsep=2pt]
    \item Thai Nguyen University of Technology;
    \item Faculty of International Training;
    \item Engineering and Technology Department;
    \item University logo and Faculty logo;
    \item the title \textbf{INTERNSHIP REPORT};
    \item Internship Organization;
    \item Author;
    \item Class;
    \item Supervisor;
    \item place and year of completion.
\end{itemize}

Students shall only replace the required information fields and shall not
modify the cover-page layout unless instructed otherwise.

\vspace{0.3cm}

% ================================================================
% VI. COMMENTS OF THE INTERNSHIP HOST ORGANIZATION
% ================================================================

\noindent
\textbf{VI.\quad COMMENTS OF THE INTERNSHIP HOST ORGANIZATION}

\vspace{0.15cm}

The page entitled \textbf{COMMENTS OF THE INTERNSHIP HOST ORGANIZATION}
is reserved for the internship host organization.

It includes:

\begin{itemize}[leftmargin=0.8cm, itemsep=2pt]
    \item the student's full name;
    \item class;
    \item comments from the host organization;
    \item score;
    \item date;
    \item signature and official seal.
\end{itemize}

Students shall not alter or complete official evaluation fields unless
specifically authorized.

\vspace{0.3cm}

% ================================================================
% VII. SUPERVISOR'S EVALUATION AND DEFENSE RESULTS
% ================================================================

\noindent
\textbf{VII.\quad SUPERVISOR'S EVALUATION AND DEFENSE RESULTS}

\vspace{0.15cm}

The section entitled
\textbf{SUPERVISOR'S EVALUATION AND GRADUATION INTERNSHIP DEFENSE RESULTS}
is reserved for the supervisor and defense committee.

The evaluation may include:

\begin{itemize}[leftmargin=0.8cm, itemsep=2pt]
    \item Internship Progress and Student's Attitude;
    \item Report Content;
    \item Report Presentation;
    \item Other Comments;
    \item Supervisor's Opinion;
    \item Defense Result;
    \item Score.
\end{itemize}

Students shall not enter evaluation scores or defense results before they
are officially determined.

\vspace{0.3cm}

% ================================================================
% VIII. CHAPTERS, SECTIONS, AND SUBSECTIONS
% ================================================================

\noindent
\textbf{VIII.\quad CHAPTERS, SECTIONS, AND SUBSECTIONS}

\vspace{0.15cm}

The report structure shall be generated using the standard \LaTeX{}
commands:

\begin{verbatim}
\chapter{...}
\section{...}
\subsection{...}
\end{verbatim}

Chapter, section, and subsection numbers shall be generated automatically
by \LaTeX{}.

Students shall not manually type chapter or section numbers and shall not
reproduce heading styles using commands such as \verb|\textbf|,
\verb|\fontsize|, or \verb|\centering| instead of the appropriate
structural command.

Excessive subdivision should be avoided. In most cases, the hierarchy
Chapter--Section--Subsection is sufficient.

\vspace{0.3cm}

% ================================================================
% IX. ORGANIZATION OF THE MAIN CONTENT
% ================================================================

\noindent
\textbf{IX.\quad ORGANIZATION OF THE MAIN CONTENT}

\vspace{0.15cm}

The specific content of the report shall depend on the internship
organization and the assigned tasks. A typical structure may include:

\begin{itemize}[leftmargin=0.8cm, itemsep=2pt]
    \item \textbf{Chapter 1. Overview of the Internship Organization};
    \item \textbf{Chapter 2. Internship Activities and Assigned Tasks};
    \item \textbf{Chapter 3. Results, Experience, and Evaluation}.
\end{itemize}

These chapter titles are recommendations only and may be adapted to the
actual internship content.

The report shall reflect work actually performed, observed, analyzed, or
learned during the internship. Unrelated theoretical material should not
be included merely to increase report length.

\vspace{0.3cm}

% ================================================================
% X. TABLE OF CONTENTS
% ================================================================

\noindent
\textbf{X.\quad TABLE OF CONTENTS}

\vspace{0.15cm}

The Table of Contents shall be generated automatically using:

\begin{verbatim}
\tableofcontents
\end{verbatim}

Students shall not create or manually edit page numbers in the Table of
Contents.

The document shall be compiled sufficiently to ensure that all page
numbers and entries are correctly updated.

\vspace{0.3cm}

% ================================================================
% XI. LIST OF SYMBOLS AND ABBREVIATIONS
% ================================================================

\noindent
\textbf{XI.\quad LIST OF SYMBOLS AND ABBREVIATIONS}

\vspace{0.15cm}

Symbols and abbreviations that are used repeatedly in the report shall be
listed in the \textbf{List of Symbols and Abbreviations}.

The list may contain the following columns:

\begin{center}
\begin{tabular}{llll}
    No. & Symbol / Abbreviation & Unit & Definition
\end{tabular}
\end{center}

Only symbols and abbreviations actually used in the final report shall be
retained. Example entries contained in the template shall be removed if
they are not relevant to the student's report.

Each abbreviation should normally be defined when it first appears in
the main text.

\vspace{0.3cm}

% ================================================================
% XII. FIGURES, DIAGRAMS, AND PHOTOGRAPHS
% ================================================================

\noindent
\textbf{XII.\quad FIGURES, DIAGRAMS, AND PHOTOGRAPHS}

\vspace{0.15cm}

Figures shall be clear, legible, and directly relevant to the report
content.

Figures shall be inserted using the standard \verb|figure| environment.

Example:

\begin{verbatim}
\begin{figure}[H]
    \centering
    \includegraphics[width=0.8\textwidth]
                    {figures/company_structure.pdf}
    \caption{Organizational structure of the company}
    \label{fig:company-structure}
\end{figure}
\end{verbatim}

Figure captions shall be placed below the figure.

Each referenced figure shall have a unique \verb|\label| and shall be
referenced in the text using, for example:

\begin{verbatim}
Figure~\ref{fig:company-structure}
\end{verbatim}

Students shall not manually type figure numbers.

Figures reproduced or adapted from external sources shall include an
appropriate citation.

Photographs taken during the internship shall be relevant and sufficiently
clear. Students shall not include confidential, restricted, or proprietary
information without permission from the internship organization.

Vector graphics or high-resolution images are recommended whenever
possible.

\vspace{0.3cm}

% ================================================================
% XIII. TABLES
% ================================================================

\noindent
\textbf{XIII.\quad TABLES}

\vspace{0.15cm}

Tables shall be created using suitable \LaTeX{} environments such as
\verb|table|, \verb|tabular|, \verb|tabularx|, or \verb|longtable|.

Table captions shall be placed above the table.

Example:

\begin{verbatim}
\begin{table}[H]
    \centering
    \caption{Internship schedule}
    \label{tab:internship-schedule}

    \begin{tabular}{lll}
        \toprule
        Week & Task & Result \\
        \midrule
        1 & Orientation & Completed \\
        2 & System study & Completed \\
        \bottomrule
    \end{tabular}
\end{table}
\end{verbatim}

Tables shall be referenced using:

\begin{verbatim}
Table~\ref{tab:internship-schedule}
\end{verbatim}

Students shall not manually type table numbers.

Tables should be created directly in \LaTeX{} rather than inserted as
screenshots whenever practical.

\vspace{0.3cm}

% ================================================================
% XIV. EQUATIONS AND TECHNICAL QUANTITIES
% ================================================================

\noindent
\textbf{XIV.\quad EQUATIONS AND TECHNICAL QUANTITIES}

\vspace{0.15cm}

Mathematical expressions shall be typeset directly using \LaTeX{}.
Images or screenshots of equations shall not be used.

Short expressions may be written in inline mathematical mode, for example:

\begin{verbatim}
$R = 0.7\,\Omega$
\end{verbatim}

Important equations shall use a numbered equation environment.

Example:

\begin{verbatim}
\begin{equation}
    P = UI
    \label{eq:power}
\end{equation}
\end{verbatim}

Equations shall be referenced using:

\begin{verbatim}
Equation~\eqref{eq:power}
\end{verbatim}

Equation numbers shall not be manually entered.

SI units shall be used wherever applicable. The \verb|siunitx| package
should be used where appropriate.

Examples:

\begin{verbatim}
\SI{220}{\volt}
\SI{50}{\hertz}
\SI{2.5}{\ampere}
\end{verbatim}

Notation and units shall be consistent throughout the report.

\vspace{0.3cm}

% ================================================================
% XV. SOURCE CODE
% ================================================================

\noindent
\textbf{XV.\quad SOURCE CODE}

\vspace{0.15cm}

If the internship involves programming, PLCs, microcontrollers,
simulation, or software development, relevant source code may be included
in the report.

Source code shall be presented using the \verb|listings| environment or
another environment defined by the official template.

Example:

\begin{verbatim}
\begin{lstlisting}[language=st]
IF Start = TRUE THEN
    Motor := TRUE;
ELSE
    Motor := FALSE;
END_IF;
\end{lstlisting}
\end{verbatim}

Only code necessary for explaining the work should appear in the main
report. Long source-code listings should be moved to an appendix or
submitted as supplementary files if required.

Source code should not normally be inserted as screenshots.

\vspace{0.3cm}

% ================================================================
% XVI. ABSTRACT
% ================================================================

\noindent
\textbf{XVI.\quad ABSTRACT}

\vspace{0.15cm}

The Abstract shall briefly summarize:

\begin{itemize}[leftmargin=0.8cm, itemsep=2pt]
    \item the internship organization;
    \item the objectives of the internship;
    \item the main tasks performed;
    \item the main results, knowledge, and skills obtained.
\end{itemize}

The Abstract should be concise and should not reproduce the full content
of the report chapter by chapter.

\vspace{0.3cm}

% ================================================================
% XVII. REFERENCES AND CITATIONS
% ================================================================

\noindent
\textbf{XVII.\quad REFERENCES AND CITATIONS}

\vspace{0.15cm}

All books, journal papers, conference papers, standards, websites,
datasheets, manuals, technical documents, figures, tables, or other
external sources used in preparing the report shall be properly cited.

References shall be managed through the bibliography system provided with
the official template.

Bibliographic entries shall normally be stored in:

\begin{verbatim}
References.bib
\end{verbatim}

References shall be cited in the report using:

\begin{verbatim}
\cite{reference-key}
\end{verbatim}

Students shall not manually type reference numbers such as ``[1]'',
``[2]'', etc.

Titles of foreign-language publications shall remain in their original
language in the bibliography.

All borrowed ideas, data, figures, tables, equations, methods, or results
shall be properly acknowledged.

\vspace{0.3cm}

% ================================================================
% XVIII. CROSS-REFERENCING
% ================================================================

\noindent
\textbf{XVIII.\quad CROSS-REFERENCING}

\vspace{0.15cm}

Numbered elements shall be cross-referenced using the standard \LaTeX{}
mechanism.

Use the following commands where appropriate:

\begin{verbatim}
\label{...}
\ref{...}
\eqref{...}
\end{verbatim}

For example:

\begin{verbatim}
As shown in Figure~\ref{fig:system}, ...
\end{verbatim}

Students shall not manually type numbers such as ``Figure 2.3'' when
those numbers can be generated automatically.

\vspace{0.3cm}

% ================================================================
% XIX. PAGE NUMBERING
% ================================================================

\noindent
\textbf{XIX.\quad PAGE NUMBERING}

\vspace{0.15cm}

Cover pages and official evaluation forms at the beginning of the report
shall not display page numbers.

The main body of the report shall use Arabic page numbers beginning from
Chapter 1.

Page numbering shall be generated automatically by the official template.
Students shall not manually enter page numbers.

\vspace{0.3cm}

% ================================================================
% XX. PROJECT FILE ORGANIZATION
% ================================================================

\noindent
\textbf{XX.\quad PROJECT FILE ORGANIZATION}

\vspace{0.15cm}

Students should maintain a clear and organized \LaTeX{} project
structure. A recommended structure is:

\begin{verbatim}
InternshipReport/
|
|-- main.tex
|-- References.bib
|
|-- figures/
|   |-- company.pdf
|   |-- system.png
|   `-- ...
|
`-- chapters/
    |-- chapter1.tex
    |-- chapter2.tex
    `-- chapter3.tex
\end{verbatim}

Figures should normally be stored in the \verb|figures| directory
specified by the official template.

File and directory names should preferably contain Latin letters,
numbers, hyphens, or underscores.

Recommended examples include:

\begin{verbatim}
company_structure.pdf
motor_control.png
chapter_01.tex
\end{verbatim}

Spaces and unnecessary special characters in file names should be
avoided where possible.

\vspace{0.3cm}

% ================================================================
% XXI. DOCUMENT COMPILATION AND FINAL CHECK
% ================================================================

\noindent
\textbf{XXI.\quad DOCUMENT COMPILATION AND FINAL CHECK}

\vspace{0.15cm}

Before submission, the complete \LaTeX{} project shall compile
successfully and produce the final PDF report.

The final document shall not contain unresolved references such as:

\begin{verbatim}
??
[?]
\end{verbatim}

Before generating the final PDF, students shall check:

\begin{itemize}[leftmargin=0.8cm, itemsep=2pt]
    \item the Table of Contents;
    \item page numbering;
    \item chapter, section, and subsection numbering;
    \item figure numbering;
    \item table numbering;
    \item equation numbering;
    \item cross-references;
    \item citations;
    \item bibliography;
    \item font display;
    \item mathematical symbols;
    \item image quality;
    \item completeness of all required forms.
\end{itemize}

The document may be compiled using TeX Live, MiKTeX, or another
compatible \LaTeX{} environment. Regardless of the compilation
environment, the resulting PDF shall conform to the official Faculty
template.

\vspace{0.3cm}

% ================================================================
% XXII. CONTENT THAT STUDENTS MAY EDIT
% ================================================================

\noindent
\textbf{XXII.\quad CONTENT THAT STUDENTS MAY EDIT}

\vspace{0.15cm}

Students may edit or replace:

\begin{itemize}[leftmargin=0.8cm, itemsep=2pt]
    \item student name;
    \item class;
    \item internship organization;
    \item supervisor name;
    \item Abstract;
    \item List of Symbols and Abbreviations;
    \item chapter titles and chapter content;
    \item figures and photographs;
    \item tables;
    \item equations;
    \item citations and references.
\end{itemize}

Students shall not modify the main document formatting settings in the
preamble, including the document class, page geometry, line spacing,
heading formatting, table-of-contents formatting, font configuration,
page layout, or cover-page macros unless specifically instructed.

\vspace{0.3cm}

% ================================================================
% XXIII. FINAL SUBMISSION REQUIREMENTS
% ================================================================

\noindent
\textbf{XXIII.\quad FINAL SUBMISSION REQUIREMENTS}

\vspace{0.15cm}

Before submission, the final report shall:

\begin{itemize}[leftmargin=0.8cm, itemsep=2pt]
    \item compile without errors that affect the report content or presentation;
    \item contain all required forms and report sections;
    \item contain correct student, class, internship organization, and supervisor information;
    \item contain an updated Table of Contents, List of Figures, and List of Tables;
    \item contain no unused example content from the template;
    \item contain no placeholder text such as \textit{Title of Chapter 1},
    \textit{Section 1.1}, or other sample entries;
    \item contain all required citations and references;
    \item contain no unresolved cross-references.
\end{itemize}

If the Faculty requires submission of the \LaTeX{} source, the complete
project shall include all \verb|.tex| files, the bibliography database,
figures, and any other files required to reproduce the submitted PDF.

\vspace{0.4cm}

\noindent
\textbf{Important:}
Before submitting the final internship report, students shall remove or
comment out the following command in the main document:

\begin{verbatim}
\input{Regulations.tex}
\end{verbatim}

For example:

\begin{verbatim}
% \input{Regulations.tex}
\end{verbatim}

This removes the guidance pages from the final submitted report while
preserving the regulations file in the \LaTeX{} project for reference.

} % End 13pt content

\clearpage

\end{verbatim}

This removes the guidance pages from the final submitted report while
preserving the regulations file in the \LaTeX{} project for reference.

} % End 13pt content

\clearpage

%
% IMPORTANT:
% This section is provided for guidance only. Students shall remove
% or comment out the 
% =================================================================
% REGULATIONS FOR PREPARING THE GRADUATION INTERNSHIP REPORT
% USING LATEX
% =================================================================
%
% This file is intended to be included from the main document with:
%
%     
% =================================================================
% REGULATIONS FOR PREPARING THE GRADUATION INTERNSHIP REPORT
% USING LATEX
% =================================================================
%
% This file is intended to be included from the main document with:
%
%     \input{Regulations.tex}
%
% IMPORTANT:
% This section is provided for guidance only. Students shall remove
% or comment out the \input{Regulations.tex} command before submitting
% the final internship report.
%
% =================================================================

\clearpage
\thispagestyle{empty}

\begin{center}
    {\fontsize{16pt}{20pt}\selectfont\bfseries
    SPECIFIC REGULATIONS FOR PREPARING\\
    GRADUATION INTERNSHIP REPORTS USING \LaTeX}
\end{center}

\vspace{0.5cm}

{\fontsize{13pt}{17pt}\selectfont
\setstretch{1.3}
\justifying

\noindent
\textbf{Note:}
This section is provided for guidance only and is included in the main
document using the command \verb|\input{Regulations.tex}|.
Students shall remove or comment out this command before submitting the
final internship report.

\vspace{0.4cm}

% ================================================================
% I. GENERAL REQUIREMENTS
% ================================================================

\noindent
\textbf{I.\quad GENERAL REQUIREMENTS}

\vspace{0.15cm}

The Graduation Internship Report shall be prepared using the official
\LaTeX{} template provided by the Faculty. Students shall preserve the
predefined page layout, typography, line spacing, heading styles,
numbering system, figure and table formatting, and bibliography settings
unless otherwise instructed by the Faculty or the supervisor.

The report shall be written in English.

Students are responsible for the academic and technical content of the
report. The official \LaTeX{} template is responsible for standardized
document formatting.

\vspace{0.3cm}

% ================================================================
% II. REPORT STRUCTURE
% ================================================================

\noindent
\textbf{II.\quad REPORT STRUCTURE}

\vspace{0.15cm}

The report shall normally be arranged in the following order:

\begin{enumerate}[label=\arabic*., leftmargin=1.2cm, itemsep=2pt]
    \item Outer Cover;
    \item Comments of the Internship Host Organization;
    \item Supervisor's Evaluation and Graduation Internship Defense Results;
    \item Table of Contents;
    \item List of Symbols and Abbreviations, if applicable;
    \item List of Figures;
    \item List of Tables;
    \item Abstract;
    \item Main Chapters of the Internship Report;
    \item References;
    \item Appendices, if applicable.
\end{enumerate}

The preliminary pages shall not be numbered as chapters. The main report
content shall be organized into consecutively numbered chapters beginning
with Chapter 1.

\vspace{0.3cm}

% ================================================================
% III. PAGE SIZE AND MARGINS
% ================================================================

\noindent
\textbf{III.\quad PAGE SIZE AND MARGINS}

\vspace{0.15cm}

The report shall be formatted on A4 paper with the following margins:

\begin{center}
\begin{tabular}{ll}
    Left margin:   & 3.5 cm \\
    Right margin:  & 2.0 cm \\
    Top margin:    & 2.5 cm \\
    Bottom margin: & 2.5 cm
\end{tabular}
\end{center}

These settings are predefined in the official template and shall not be
modified by students.

\vspace{0.3cm}

% ================================================================
% IV. FONT, FONT SIZE, AND LINE SPACING
% ================================================================

\noindent
\textbf{IV.\quad FONT, FONT SIZE, AND LINE SPACING}

\vspace{0.15cm}

The main text shall use a 13 pt font size with 1.3-line spacing.
A Times-family typeface shall be used in accordance with the official
template.

The text shall be fully justified.

Students shall not use repeated spaces, repeated manual line breaks, or
multiple blank lines to position text. Paragraph and heading spacing shall
be controlled by the template.

\vspace{0.3cm}

% ================================================================
% V. OUTER COVER
% ================================================================

\noindent
\textbf{V.\quad OUTER COVER}

\vspace{0.15cm}

The outer cover shall follow the official format provided in the template
and shall contain the following information:

\begin{itemize}[leftmargin=0.8cm, itemsep=2pt]
    \item Thai Nguyen University of Technology;
    \item Faculty of International Training;
    \item Engineering and Technology Department;
    \item University logo and Faculty logo;
    \item the title \textbf{INTERNSHIP REPORT};
    \item Internship Organization;
    \item Author;
    \item Class;
    \item Supervisor;
    \item place and year of completion.
\end{itemize}

Students shall only replace the required information fields and shall not
modify the cover-page layout unless instructed otherwise.

\vspace{0.3cm}

% ================================================================
% VI. COMMENTS OF THE INTERNSHIP HOST ORGANIZATION
% ================================================================

\noindent
\textbf{VI.\quad COMMENTS OF THE INTERNSHIP HOST ORGANIZATION}

\vspace{0.15cm}

The page entitled \textbf{COMMENTS OF THE INTERNSHIP HOST ORGANIZATION}
is reserved for the internship host organization.

It includes:

\begin{itemize}[leftmargin=0.8cm, itemsep=2pt]
    \item the student's full name;
    \item class;
    \item comments from the host organization;
    \item score;
    \item date;
    \item signature and official seal.
\end{itemize}

Students shall not alter or complete official evaluation fields unless
specifically authorized.

\vspace{0.3cm}

% ================================================================
% VII. SUPERVISOR'S EVALUATION AND DEFENSE RESULTS
% ================================================================

\noindent
\textbf{VII.\quad SUPERVISOR'S EVALUATION AND DEFENSE RESULTS}

\vspace{0.15cm}

The section entitled
\textbf{SUPERVISOR'S EVALUATION AND GRADUATION INTERNSHIP DEFENSE RESULTS}
is reserved for the supervisor and defense committee.

The evaluation may include:

\begin{itemize}[leftmargin=0.8cm, itemsep=2pt]
    \item Internship Progress and Student's Attitude;
    \item Report Content;
    \item Report Presentation;
    \item Other Comments;
    \item Supervisor's Opinion;
    \item Defense Result;
    \item Score.
\end{itemize}

Students shall not enter evaluation scores or defense results before they
are officially determined.

\vspace{0.3cm}

% ================================================================
% VIII. CHAPTERS, SECTIONS, AND SUBSECTIONS
% ================================================================

\noindent
\textbf{VIII.\quad CHAPTERS, SECTIONS, AND SUBSECTIONS}

\vspace{0.15cm}

The report structure shall be generated using the standard \LaTeX{}
commands:

\begin{verbatim}
\chapter{...}
\section{...}
\subsection{...}
\end{verbatim}

Chapter, section, and subsection numbers shall be generated automatically
by \LaTeX{}.

Students shall not manually type chapter or section numbers and shall not
reproduce heading styles using commands such as \verb|\textbf|,
\verb|\fontsize|, or \verb|\centering| instead of the appropriate
structural command.

Excessive subdivision should be avoided. In most cases, the hierarchy
Chapter--Section--Subsection is sufficient.

\vspace{0.3cm}

% ================================================================
% IX. ORGANIZATION OF THE MAIN CONTENT
% ================================================================

\noindent
\textbf{IX.\quad ORGANIZATION OF THE MAIN CONTENT}

\vspace{0.15cm}

The specific content of the report shall depend on the internship
organization and the assigned tasks. A typical structure may include:

\begin{itemize}[leftmargin=0.8cm, itemsep=2pt]
    \item \textbf{Chapter 1. Overview of the Internship Organization};
    \item \textbf{Chapter 2. Internship Activities and Assigned Tasks};
    \item \textbf{Chapter 3. Results, Experience, and Evaluation}.
\end{itemize}

These chapter titles are recommendations only and may be adapted to the
actual internship content.

The report shall reflect work actually performed, observed, analyzed, or
learned during the internship. Unrelated theoretical material should not
be included merely to increase report length.

\vspace{0.3cm}

% ================================================================
% X. TABLE OF CONTENTS
% ================================================================

\noindent
\textbf{X.\quad TABLE OF CONTENTS}

\vspace{0.15cm}

The Table of Contents shall be generated automatically using:

\begin{verbatim}
\tableofcontents
\end{verbatim}

Students shall not create or manually edit page numbers in the Table of
Contents.

The document shall be compiled sufficiently to ensure that all page
numbers and entries are correctly updated.

\vspace{0.3cm}

% ================================================================
% XI. LIST OF SYMBOLS AND ABBREVIATIONS
% ================================================================

\noindent
\textbf{XI.\quad LIST OF SYMBOLS AND ABBREVIATIONS}

\vspace{0.15cm}

Symbols and abbreviations that are used repeatedly in the report shall be
listed in the \textbf{List of Symbols and Abbreviations}.

The list may contain the following columns:

\begin{center}
\begin{tabular}{llll}
    No. & Symbol / Abbreviation & Unit & Definition
\end{tabular}
\end{center}

Only symbols and abbreviations actually used in the final report shall be
retained. Example entries contained in the template shall be removed if
they are not relevant to the student's report.

Each abbreviation should normally be defined when it first appears in
the main text.

\vspace{0.3cm}

% ================================================================
% XII. FIGURES, DIAGRAMS, AND PHOTOGRAPHS
% ================================================================

\noindent
\textbf{XII.\quad FIGURES, DIAGRAMS, AND PHOTOGRAPHS}

\vspace{0.15cm}

Figures shall be clear, legible, and directly relevant to the report
content.

Figures shall be inserted using the standard \verb|figure| environment.

Example:

\begin{verbatim}
\begin{figure}[H]
    \centering
    \includegraphics[width=0.8\textwidth]
                    {figures/company_structure.pdf}
    \caption{Organizational structure of the company}
    \label{fig:company-structure}
\end{figure}
\end{verbatim}

Figure captions shall be placed below the figure.

Each referenced figure shall have a unique \verb|\label| and shall be
referenced in the text using, for example:

\begin{verbatim}
Figure~\ref{fig:company-structure}
\end{verbatim}

Students shall not manually type figure numbers.

Figures reproduced or adapted from external sources shall include an
appropriate citation.

Photographs taken during the internship shall be relevant and sufficiently
clear. Students shall not include confidential, restricted, or proprietary
information without permission from the internship organization.

Vector graphics or high-resolution images are recommended whenever
possible.

\vspace{0.3cm}

% ================================================================
% XIII. TABLES
% ================================================================

\noindent
\textbf{XIII.\quad TABLES}

\vspace{0.15cm}

Tables shall be created using suitable \LaTeX{} environments such as
\verb|table|, \verb|tabular|, \verb|tabularx|, or \verb|longtable|.

Table captions shall be placed above the table.

Example:

\begin{verbatim}
\begin{table}[H]
    \centering
    \caption{Internship schedule}
    \label{tab:internship-schedule}

    \begin{tabular}{lll}
        \toprule
        Week & Task & Result \\
        \midrule
        1 & Orientation & Completed \\
        2 & System study & Completed \\
        \bottomrule
    \end{tabular}
\end{table}
\end{verbatim}

Tables shall be referenced using:

\begin{verbatim}
Table~\ref{tab:internship-schedule}
\end{verbatim}

Students shall not manually type table numbers.

Tables should be created directly in \LaTeX{} rather than inserted as
screenshots whenever practical.

\vspace{0.3cm}

% ================================================================
% XIV. EQUATIONS AND TECHNICAL QUANTITIES
% ================================================================

\noindent
\textbf{XIV.\quad EQUATIONS AND TECHNICAL QUANTITIES}

\vspace{0.15cm}

Mathematical expressions shall be typeset directly using \LaTeX{}.
Images or screenshots of equations shall not be used.

Short expressions may be written in inline mathematical mode, for example:

\begin{verbatim}
$R = 0.7\,\Omega$
\end{verbatim}

Important equations shall use a numbered equation environment.

Example:

\begin{verbatim}
\begin{equation}
    P = UI
    \label{eq:power}
\end{equation}
\end{verbatim}

Equations shall be referenced using:

\begin{verbatim}
Equation~\eqref{eq:power}
\end{verbatim}

Equation numbers shall not be manually entered.

SI units shall be used wherever applicable. The \verb|siunitx| package
should be used where appropriate.

Examples:

\begin{verbatim}
\SI{220}{\volt}
\SI{50}{\hertz}
\SI{2.5}{\ampere}
\end{verbatim}

Notation and units shall be consistent throughout the report.

\vspace{0.3cm}

% ================================================================
% XV. SOURCE CODE
% ================================================================

\noindent
\textbf{XV.\quad SOURCE CODE}

\vspace{0.15cm}

If the internship involves programming, PLCs, microcontrollers,
simulation, or software development, relevant source code may be included
in the report.

Source code shall be presented using the \verb|listings| environment or
another environment defined by the official template.

Example:

\begin{verbatim}
\begin{lstlisting}[language=st]
IF Start = TRUE THEN
    Motor := TRUE;
ELSE
    Motor := FALSE;
END_IF;
\end{lstlisting}
\end{verbatim}

Only code necessary for explaining the work should appear in the main
report. Long source-code listings should be moved to an appendix or
submitted as supplementary files if required.

Source code should not normally be inserted as screenshots.

\vspace{0.3cm}

% ================================================================
% XVI. ABSTRACT
% ================================================================

\noindent
\textbf{XVI.\quad ABSTRACT}

\vspace{0.15cm}

The Abstract shall briefly summarize:

\begin{itemize}[leftmargin=0.8cm, itemsep=2pt]
    \item the internship organization;
    \item the objectives of the internship;
    \item the main tasks performed;
    \item the main results, knowledge, and skills obtained.
\end{itemize}

The Abstract should be concise and should not reproduce the full content
of the report chapter by chapter.

\vspace{0.3cm}

% ================================================================
% XVII. REFERENCES AND CITATIONS
% ================================================================

\noindent
\textbf{XVII.\quad REFERENCES AND CITATIONS}

\vspace{0.15cm}

All books, journal papers, conference papers, standards, websites,
datasheets, manuals, technical documents, figures, tables, or other
external sources used in preparing the report shall be properly cited.

References shall be managed through the bibliography system provided with
the official template.

Bibliographic entries shall normally be stored in:

\begin{verbatim}
References.bib
\end{verbatim}

References shall be cited in the report using:

\begin{verbatim}
\cite{reference-key}
\end{verbatim}

Students shall not manually type reference numbers such as ``[1]'',
``[2]'', etc.

Titles of foreign-language publications shall remain in their original
language in the bibliography.

All borrowed ideas, data, figures, tables, equations, methods, or results
shall be properly acknowledged.

\vspace{0.3cm}

% ================================================================
% XVIII. CROSS-REFERENCING
% ================================================================

\noindent
\textbf{XVIII.\quad CROSS-REFERENCING}

\vspace{0.15cm}

Numbered elements shall be cross-referenced using the standard \LaTeX{}
mechanism.

Use the following commands where appropriate:

\begin{verbatim}
\label{...}
\ref{...}
\eqref{...}
\end{verbatim}

For example:

\begin{verbatim}
As shown in Figure~\ref{fig:system}, ...
\end{verbatim}

Students shall not manually type numbers such as ``Figure 2.3'' when
those numbers can be generated automatically.

\vspace{0.3cm}

% ================================================================
% XIX. PAGE NUMBERING
% ================================================================

\noindent
\textbf{XIX.\quad PAGE NUMBERING}

\vspace{0.15cm}

Cover pages and official evaluation forms at the beginning of the report
shall not display page numbers.

The main body of the report shall use Arabic page numbers beginning from
Chapter 1.

Page numbering shall be generated automatically by the official template.
Students shall not manually enter page numbers.

\vspace{0.3cm}

% ================================================================
% XX. PROJECT FILE ORGANIZATION
% ================================================================

\noindent
\textbf{XX.\quad PROJECT FILE ORGANIZATION}

\vspace{0.15cm}

Students should maintain a clear and organized \LaTeX{} project
structure. A recommended structure is:

\begin{verbatim}
InternshipReport/
|
|-- main.tex
|-- References.bib
|
|-- figures/
|   |-- company.pdf
|   |-- system.png
|   `-- ...
|
`-- chapters/
    |-- chapter1.tex
    |-- chapter2.tex
    `-- chapter3.tex
\end{verbatim}

Figures should normally be stored in the \verb|figures| directory
specified by the official template.

File and directory names should preferably contain Latin letters,
numbers, hyphens, or underscores.

Recommended examples include:

\begin{verbatim}
company_structure.pdf
motor_control.png
chapter_01.tex
\end{verbatim}

Spaces and unnecessary special characters in file names should be
avoided where possible.

\vspace{0.3cm}

% ================================================================
% XXI. DOCUMENT COMPILATION AND FINAL CHECK
% ================================================================

\noindent
\textbf{XXI.\quad DOCUMENT COMPILATION AND FINAL CHECK}

\vspace{0.15cm}

Before submission, the complete \LaTeX{} project shall compile
successfully and produce the final PDF report.

The final document shall not contain unresolved references such as:

\begin{verbatim}
??
[?]
\end{verbatim}

Before generating the final PDF, students shall check:

\begin{itemize}[leftmargin=0.8cm, itemsep=2pt]
    \item the Table of Contents;
    \item page numbering;
    \item chapter, section, and subsection numbering;
    \item figure numbering;
    \item table numbering;
    \item equation numbering;
    \item cross-references;
    \item citations;
    \item bibliography;
    \item font display;
    \item mathematical symbols;
    \item image quality;
    \item completeness of all required forms.
\end{itemize}

The document may be compiled using TeX Live, MiKTeX, or another
compatible \LaTeX{} environment. Regardless of the compilation
environment, the resulting PDF shall conform to the official Faculty
template.

\vspace{0.3cm}

% ================================================================
% XXII. CONTENT THAT STUDENTS MAY EDIT
% ================================================================

\noindent
\textbf{XXII.\quad CONTENT THAT STUDENTS MAY EDIT}

\vspace{0.15cm}

Students may edit or replace:

\begin{itemize}[leftmargin=0.8cm, itemsep=2pt]
    \item student name;
    \item class;
    \item internship organization;
    \item supervisor name;
    \item Abstract;
    \item List of Symbols and Abbreviations;
    \item chapter titles and chapter content;
    \item figures and photographs;
    \item tables;
    \item equations;
    \item citations and references.
\end{itemize}

Students shall not modify the main document formatting settings in the
preamble, including the document class, page geometry, line spacing,
heading formatting, table-of-contents formatting, font configuration,
page layout, or cover-page macros unless specifically instructed.

\vspace{0.3cm}

% ================================================================
% XXIII. FINAL SUBMISSION REQUIREMENTS
% ================================================================

\noindent
\textbf{XXIII.\quad FINAL SUBMISSION REQUIREMENTS}

\vspace{0.15cm}

Before submission, the final report shall:

\begin{itemize}[leftmargin=0.8cm, itemsep=2pt]
    \item compile without errors that affect the report content or presentation;
    \item contain all required forms and report sections;
    \item contain correct student, class, internship organization, and supervisor information;
    \item contain an updated Table of Contents, List of Figures, and List of Tables;
    \item contain no unused example content from the template;
    \item contain no placeholder text such as \textit{Title of Chapter 1},
    \textit{Section 1.1}, or other sample entries;
    \item contain all required citations and references;
    \item contain no unresolved cross-references.
\end{itemize}

If the Faculty requires submission of the \LaTeX{} source, the complete
project shall include all \verb|.tex| files, the bibliography database,
figures, and any other files required to reproduce the submitted PDF.

\vspace{0.4cm}

\noindent
\textbf{Important:}
Before submitting the final internship report, students shall remove or
comment out the following command in the main document:

\begin{verbatim}
\input{Regulations.tex}
\end{verbatim}

For example:

\begin{verbatim}
% \input{Regulations.tex}
\end{verbatim}

This removes the guidance pages from the final submitted report while
preserving the regulations file in the \LaTeX{} project for reference.

} % End 13pt content

\clearpage

%
% IMPORTANT:
% This section is provided for guidance only. Students shall remove
% or comment out the 
% =================================================================
% REGULATIONS FOR PREPARING THE GRADUATION INTERNSHIP REPORT
% USING LATEX
% =================================================================
%
% This file is intended to be included from the main document with:
%
%     \input{Regulations.tex}
%
% IMPORTANT:
% This section is provided for guidance only. Students shall remove
% or comment out the \input{Regulations.tex} command before submitting
% the final internship report.
%
% =================================================================

\clearpage
\thispagestyle{empty}

\begin{center}
    {\fontsize{16pt}{20pt}\selectfont\bfseries
    SPECIFIC REGULATIONS FOR PREPARING\\
    GRADUATION INTERNSHIP REPORTS USING \LaTeX}
\end{center}

\vspace{0.5cm}

{\fontsize{13pt}{17pt}\selectfont
\setstretch{1.3}
\justifying

\noindent
\textbf{Note:}
This section is provided for guidance only and is included in the main
document using the command \verb|\input{Regulations.tex}|.
Students shall remove or comment out this command before submitting the
final internship report.

\vspace{0.4cm}

% ================================================================
% I. GENERAL REQUIREMENTS
% ================================================================

\noindent
\textbf{I.\quad GENERAL REQUIREMENTS}

\vspace{0.15cm}

The Graduation Internship Report shall be prepared using the official
\LaTeX{} template provided by the Faculty. Students shall preserve the
predefined page layout, typography, line spacing, heading styles,
numbering system, figure and table formatting, and bibliography settings
unless otherwise instructed by the Faculty or the supervisor.

The report shall be written in English.

Students are responsible for the academic and technical content of the
report. The official \LaTeX{} template is responsible for standardized
document formatting.

\vspace{0.3cm}

% ================================================================
% II. REPORT STRUCTURE
% ================================================================

\noindent
\textbf{II.\quad REPORT STRUCTURE}

\vspace{0.15cm}

The report shall normally be arranged in the following order:

\begin{enumerate}[label=\arabic*., leftmargin=1.2cm, itemsep=2pt]
    \item Outer Cover;
    \item Comments of the Internship Host Organization;
    \item Supervisor's Evaluation and Graduation Internship Defense Results;
    \item Table of Contents;
    \item List of Symbols and Abbreviations, if applicable;
    \item List of Figures;
    \item List of Tables;
    \item Abstract;
    \item Main Chapters of the Internship Report;
    \item References;
    \item Appendices, if applicable.
\end{enumerate}

The preliminary pages shall not be numbered as chapters. The main report
content shall be organized into consecutively numbered chapters beginning
with Chapter 1.

\vspace{0.3cm}

% ================================================================
% III. PAGE SIZE AND MARGINS
% ================================================================

\noindent
\textbf{III.\quad PAGE SIZE AND MARGINS}

\vspace{0.15cm}

The report shall be formatted on A4 paper with the following margins:

\begin{center}
\begin{tabular}{ll}
    Left margin:   & 3.5 cm \\
    Right margin:  & 2.0 cm \\
    Top margin:    & 2.5 cm \\
    Bottom margin: & 2.5 cm
\end{tabular}
\end{center}

These settings are predefined in the official template and shall not be
modified by students.

\vspace{0.3cm}

% ================================================================
% IV. FONT, FONT SIZE, AND LINE SPACING
% ================================================================

\noindent
\textbf{IV.\quad FONT, FONT SIZE, AND LINE SPACING}

\vspace{0.15cm}

The main text shall use a 13 pt font size with 1.3-line spacing.
A Times-family typeface shall be used in accordance with the official
template.

The text shall be fully justified.

Students shall not use repeated spaces, repeated manual line breaks, or
multiple blank lines to position text. Paragraph and heading spacing shall
be controlled by the template.

\vspace{0.3cm}

% ================================================================
% V. OUTER COVER
% ================================================================

\noindent
\textbf{V.\quad OUTER COVER}

\vspace{0.15cm}

The outer cover shall follow the official format provided in the template
and shall contain the following information:

\begin{itemize}[leftmargin=0.8cm, itemsep=2pt]
    \item Thai Nguyen University of Technology;
    \item Faculty of International Training;
    \item Engineering and Technology Department;
    \item University logo and Faculty logo;
    \item the title \textbf{INTERNSHIP REPORT};
    \item Internship Organization;
    \item Author;
    \item Class;
    \item Supervisor;
    \item place and year of completion.
\end{itemize}

Students shall only replace the required information fields and shall not
modify the cover-page layout unless instructed otherwise.

\vspace{0.3cm}

% ================================================================
% VI. COMMENTS OF THE INTERNSHIP HOST ORGANIZATION
% ================================================================

\noindent
\textbf{VI.\quad COMMENTS OF THE INTERNSHIP HOST ORGANIZATION}

\vspace{0.15cm}

The page entitled \textbf{COMMENTS OF THE INTERNSHIP HOST ORGANIZATION}
is reserved for the internship host organization.

It includes:

\begin{itemize}[leftmargin=0.8cm, itemsep=2pt]
    \item the student's full name;
    \item class;
    \item comments from the host organization;
    \item score;
    \item date;
    \item signature and official seal.
\end{itemize}

Students shall not alter or complete official evaluation fields unless
specifically authorized.

\vspace{0.3cm}

% ================================================================
% VII. SUPERVISOR'S EVALUATION AND DEFENSE RESULTS
% ================================================================

\noindent
\textbf{VII.\quad SUPERVISOR'S EVALUATION AND DEFENSE RESULTS}

\vspace{0.15cm}

The section entitled
\textbf{SUPERVISOR'S EVALUATION AND GRADUATION INTERNSHIP DEFENSE RESULTS}
is reserved for the supervisor and defense committee.

The evaluation may include:

\begin{itemize}[leftmargin=0.8cm, itemsep=2pt]
    \item Internship Progress and Student's Attitude;
    \item Report Content;
    \item Report Presentation;
    \item Other Comments;
    \item Supervisor's Opinion;
    \item Defense Result;
    \item Score.
\end{itemize}

Students shall not enter evaluation scores or defense results before they
are officially determined.

\vspace{0.3cm}

% ================================================================
% VIII. CHAPTERS, SECTIONS, AND SUBSECTIONS
% ================================================================

\noindent
\textbf{VIII.\quad CHAPTERS, SECTIONS, AND SUBSECTIONS}

\vspace{0.15cm}

The report structure shall be generated using the standard \LaTeX{}
commands:

\begin{verbatim}
\chapter{...}
\section{...}
\subsection{...}
\end{verbatim}

Chapter, section, and subsection numbers shall be generated automatically
by \LaTeX{}.

Students shall not manually type chapter or section numbers and shall not
reproduce heading styles using commands such as \verb|\textbf|,
\verb|\fontsize|, or \verb|\centering| instead of the appropriate
structural command.

Excessive subdivision should be avoided. In most cases, the hierarchy
Chapter--Section--Subsection is sufficient.

\vspace{0.3cm}

% ================================================================
% IX. ORGANIZATION OF THE MAIN CONTENT
% ================================================================

\noindent
\textbf{IX.\quad ORGANIZATION OF THE MAIN CONTENT}

\vspace{0.15cm}

The specific content of the report shall depend on the internship
organization and the assigned tasks. A typical structure may include:

\begin{itemize}[leftmargin=0.8cm, itemsep=2pt]
    \item \textbf{Chapter 1. Overview of the Internship Organization};
    \item \textbf{Chapter 2. Internship Activities and Assigned Tasks};
    \item \textbf{Chapter 3. Results, Experience, and Evaluation}.
\end{itemize}

These chapter titles are recommendations only and may be adapted to the
actual internship content.

The report shall reflect work actually performed, observed, analyzed, or
learned during the internship. Unrelated theoretical material should not
be included merely to increase report length.

\vspace{0.3cm}

% ================================================================
% X. TABLE OF CONTENTS
% ================================================================

\noindent
\textbf{X.\quad TABLE OF CONTENTS}

\vspace{0.15cm}

The Table of Contents shall be generated automatically using:

\begin{verbatim}
\tableofcontents
\end{verbatim}

Students shall not create or manually edit page numbers in the Table of
Contents.

The document shall be compiled sufficiently to ensure that all page
numbers and entries are correctly updated.

\vspace{0.3cm}

% ================================================================
% XI. LIST OF SYMBOLS AND ABBREVIATIONS
% ================================================================

\noindent
\textbf{XI.\quad LIST OF SYMBOLS AND ABBREVIATIONS}

\vspace{0.15cm}

Symbols and abbreviations that are used repeatedly in the report shall be
listed in the \textbf{List of Symbols and Abbreviations}.

The list may contain the following columns:

\begin{center}
\begin{tabular}{llll}
    No. & Symbol / Abbreviation & Unit & Definition
\end{tabular}
\end{center}

Only symbols and abbreviations actually used in the final report shall be
retained. Example entries contained in the template shall be removed if
they are not relevant to the student's report.

Each abbreviation should normally be defined when it first appears in
the main text.

\vspace{0.3cm}

% ================================================================
% XII. FIGURES, DIAGRAMS, AND PHOTOGRAPHS
% ================================================================

\noindent
\textbf{XII.\quad FIGURES, DIAGRAMS, AND PHOTOGRAPHS}

\vspace{0.15cm}

Figures shall be clear, legible, and directly relevant to the report
content.

Figures shall be inserted using the standard \verb|figure| environment.

Example:

\begin{verbatim}
\begin{figure}[H]
    \centering
    \includegraphics[width=0.8\textwidth]
                    {figures/company_structure.pdf}
    \caption{Organizational structure of the company}
    \label{fig:company-structure}
\end{figure}
\end{verbatim}

Figure captions shall be placed below the figure.

Each referenced figure shall have a unique \verb|\label| and shall be
referenced in the text using, for example:

\begin{verbatim}
Figure~\ref{fig:company-structure}
\end{verbatim}

Students shall not manually type figure numbers.

Figures reproduced or adapted from external sources shall include an
appropriate citation.

Photographs taken during the internship shall be relevant and sufficiently
clear. Students shall not include confidential, restricted, or proprietary
information without permission from the internship organization.

Vector graphics or high-resolution images are recommended whenever
possible.

\vspace{0.3cm}

% ================================================================
% XIII. TABLES
% ================================================================

\noindent
\textbf{XIII.\quad TABLES}

\vspace{0.15cm}

Tables shall be created using suitable \LaTeX{} environments such as
\verb|table|, \verb|tabular|, \verb|tabularx|, or \verb|longtable|.

Table captions shall be placed above the table.

Example:

\begin{verbatim}
\begin{table}[H]
    \centering
    \caption{Internship schedule}
    \label{tab:internship-schedule}

    \begin{tabular}{lll}
        \toprule
        Week & Task & Result \\
        \midrule
        1 & Orientation & Completed \\
        2 & System study & Completed \\
        \bottomrule
    \end{tabular}
\end{table}
\end{verbatim}

Tables shall be referenced using:

\begin{verbatim}
Table~\ref{tab:internship-schedule}
\end{verbatim}

Students shall not manually type table numbers.

Tables should be created directly in \LaTeX{} rather than inserted as
screenshots whenever practical.

\vspace{0.3cm}

% ================================================================
% XIV. EQUATIONS AND TECHNICAL QUANTITIES
% ================================================================

\noindent
\textbf{XIV.\quad EQUATIONS AND TECHNICAL QUANTITIES}

\vspace{0.15cm}

Mathematical expressions shall be typeset directly using \LaTeX{}.
Images or screenshots of equations shall not be used.

Short expressions may be written in inline mathematical mode, for example:

\begin{verbatim}
$R = 0.7\,\Omega$
\end{verbatim}

Important equations shall use a numbered equation environment.

Example:

\begin{verbatim}
\begin{equation}
    P = UI
    \label{eq:power}
\end{equation}
\end{verbatim}

Equations shall be referenced using:

\begin{verbatim}
Equation~\eqref{eq:power}
\end{verbatim}

Equation numbers shall not be manually entered.

SI units shall be used wherever applicable. The \verb|siunitx| package
should be used where appropriate.

Examples:

\begin{verbatim}
\SI{220}{\volt}
\SI{50}{\hertz}
\SI{2.5}{\ampere}
\end{verbatim}

Notation and units shall be consistent throughout the report.

\vspace{0.3cm}

% ================================================================
% XV. SOURCE CODE
% ================================================================

\noindent
\textbf{XV.\quad SOURCE CODE}

\vspace{0.15cm}

If the internship involves programming, PLCs, microcontrollers,
simulation, or software development, relevant source code may be included
in the report.

Source code shall be presented using the \verb|listings| environment or
another environment defined by the official template.

Example:

\begin{verbatim}
\begin{lstlisting}[language=st]
IF Start = TRUE THEN
    Motor := TRUE;
ELSE
    Motor := FALSE;
END_IF;
\end{lstlisting}
\end{verbatim}

Only code necessary for explaining the work should appear in the main
report. Long source-code listings should be moved to an appendix or
submitted as supplementary files if required.

Source code should not normally be inserted as screenshots.

\vspace{0.3cm}

% ================================================================
% XVI. ABSTRACT
% ================================================================

\noindent
\textbf{XVI.\quad ABSTRACT}

\vspace{0.15cm}

The Abstract shall briefly summarize:

\begin{itemize}[leftmargin=0.8cm, itemsep=2pt]
    \item the internship organization;
    \item the objectives of the internship;
    \item the main tasks performed;
    \item the main results, knowledge, and skills obtained.
\end{itemize}

The Abstract should be concise and should not reproduce the full content
of the report chapter by chapter.

\vspace{0.3cm}

% ================================================================
% XVII. REFERENCES AND CITATIONS
% ================================================================

\noindent
\textbf{XVII.\quad REFERENCES AND CITATIONS}

\vspace{0.15cm}

All books, journal papers, conference papers, standards, websites,
datasheets, manuals, technical documents, figures, tables, or other
external sources used in preparing the report shall be properly cited.

References shall be managed through the bibliography system provided with
the official template.

Bibliographic entries shall normally be stored in:

\begin{verbatim}
References.bib
\end{verbatim}

References shall be cited in the report using:

\begin{verbatim}
\cite{reference-key}
\end{verbatim}

Students shall not manually type reference numbers such as ``[1]'',
``[2]'', etc.

Titles of foreign-language publications shall remain in their original
language in the bibliography.

All borrowed ideas, data, figures, tables, equations, methods, or results
shall be properly acknowledged.

\vspace{0.3cm}

% ================================================================
% XVIII. CROSS-REFERENCING
% ================================================================

\noindent
\textbf{XVIII.\quad CROSS-REFERENCING}

\vspace{0.15cm}

Numbered elements shall be cross-referenced using the standard \LaTeX{}
mechanism.

Use the following commands where appropriate:

\begin{verbatim}
\label{...}
\ref{...}
\eqref{...}
\end{verbatim}

For example:

\begin{verbatim}
As shown in Figure~\ref{fig:system}, ...
\end{verbatim}

Students shall not manually type numbers such as ``Figure 2.3'' when
those numbers can be generated automatically.

\vspace{0.3cm}

% ================================================================
% XIX. PAGE NUMBERING
% ================================================================

\noindent
\textbf{XIX.\quad PAGE NUMBERING}

\vspace{0.15cm}

Cover pages and official evaluation forms at the beginning of the report
shall not display page numbers.

The main body of the report shall use Arabic page numbers beginning from
Chapter 1.

Page numbering shall be generated automatically by the official template.
Students shall not manually enter page numbers.

\vspace{0.3cm}

% ================================================================
% XX. PROJECT FILE ORGANIZATION
% ================================================================

\noindent
\textbf{XX.\quad PROJECT FILE ORGANIZATION}

\vspace{0.15cm}

Students should maintain a clear and organized \LaTeX{} project
structure. A recommended structure is:

\begin{verbatim}
InternshipReport/
|
|-- main.tex
|-- References.bib
|
|-- figures/
|   |-- company.pdf
|   |-- system.png
|   `-- ...
|
`-- chapters/
    |-- chapter1.tex
    |-- chapter2.tex
    `-- chapter3.tex
\end{verbatim}

Figures should normally be stored in the \verb|figures| directory
specified by the official template.

File and directory names should preferably contain Latin letters,
numbers, hyphens, or underscores.

Recommended examples include:

\begin{verbatim}
company_structure.pdf
motor_control.png
chapter_01.tex
\end{verbatim}

Spaces and unnecessary special characters in file names should be
avoided where possible.

\vspace{0.3cm}

% ================================================================
% XXI. DOCUMENT COMPILATION AND FINAL CHECK
% ================================================================

\noindent
\textbf{XXI.\quad DOCUMENT COMPILATION AND FINAL CHECK}

\vspace{0.15cm}

Before submission, the complete \LaTeX{} project shall compile
successfully and produce the final PDF report.

The final document shall not contain unresolved references such as:

\begin{verbatim}
??
[?]
\end{verbatim}

Before generating the final PDF, students shall check:

\begin{itemize}[leftmargin=0.8cm, itemsep=2pt]
    \item the Table of Contents;
    \item page numbering;
    \item chapter, section, and subsection numbering;
    \item figure numbering;
    \item table numbering;
    \item equation numbering;
    \item cross-references;
    \item citations;
    \item bibliography;
    \item font display;
    \item mathematical symbols;
    \item image quality;
    \item completeness of all required forms.
\end{itemize}

The document may be compiled using TeX Live, MiKTeX, or another
compatible \LaTeX{} environment. Regardless of the compilation
environment, the resulting PDF shall conform to the official Faculty
template.

\vspace{0.3cm}

% ================================================================
% XXII. CONTENT THAT STUDENTS MAY EDIT
% ================================================================

\noindent
\textbf{XXII.\quad CONTENT THAT STUDENTS MAY EDIT}

\vspace{0.15cm}

Students may edit or replace:

\begin{itemize}[leftmargin=0.8cm, itemsep=2pt]
    \item student name;
    \item class;
    \item internship organization;
    \item supervisor name;
    \item Abstract;
    \item List of Symbols and Abbreviations;
    \item chapter titles and chapter content;
    \item figures and photographs;
    \item tables;
    \item equations;
    \item citations and references.
\end{itemize}

Students shall not modify the main document formatting settings in the
preamble, including the document class, page geometry, line spacing,
heading formatting, table-of-contents formatting, font configuration,
page layout, or cover-page macros unless specifically instructed.

\vspace{0.3cm}

% ================================================================
% XXIII. FINAL SUBMISSION REQUIREMENTS
% ================================================================

\noindent
\textbf{XXIII.\quad FINAL SUBMISSION REQUIREMENTS}

\vspace{0.15cm}

Before submission, the final report shall:

\begin{itemize}[leftmargin=0.8cm, itemsep=2pt]
    \item compile without errors that affect the report content or presentation;
    \item contain all required forms and report sections;
    \item contain correct student, class, internship organization, and supervisor information;
    \item contain an updated Table of Contents, List of Figures, and List of Tables;
    \item contain no unused example content from the template;
    \item contain no placeholder text such as \textit{Title of Chapter 1},
    \textit{Section 1.1}, or other sample entries;
    \item contain all required citations and references;
    \item contain no unresolved cross-references.
\end{itemize}

If the Faculty requires submission of the \LaTeX{} source, the complete
project shall include all \verb|.tex| files, the bibliography database,
figures, and any other files required to reproduce the submitted PDF.

\vspace{0.4cm}

\noindent
\textbf{Important:}
Before submitting the final internship report, students shall remove or
comment out the following command in the main document:

\begin{verbatim}
\input{Regulations.tex}
\end{verbatim}

For example:

\begin{verbatim}
% \input{Regulations.tex}
\end{verbatim}

This removes the guidance pages from the final submitted report while
preserving the regulations file in the \LaTeX{} project for reference.

} % End 13pt content

\clearpage
 command before submitting
% the final internship report.
%
% =================================================================

\clearpage
\thispagestyle{empty}

\begin{center}
    {\fontsize{16pt}{20pt}\selectfont\bfseries
    SPECIFIC REGULATIONS FOR PREPARING\\
    GRADUATION INTERNSHIP REPORTS USING \LaTeX}
\end{center}

\vspace{0.5cm}

{\fontsize{13pt}{17pt}\selectfont
\setstretch{1.3}
\justifying

\noindent
\textbf{Note:}
This section is provided for guidance only and is included in the main
document using the command \verb|
% =================================================================
% REGULATIONS FOR PREPARING THE GRADUATION INTERNSHIP REPORT
% USING LATEX
% =================================================================
%
% This file is intended to be included from the main document with:
%
%     \input{Regulations.tex}
%
% IMPORTANT:
% This section is provided for guidance only. Students shall remove
% or comment out the \input{Regulations.tex} command before submitting
% the final internship report.
%
% =================================================================

\clearpage
\thispagestyle{empty}

\begin{center}
    {\fontsize{16pt}{20pt}\selectfont\bfseries
    SPECIFIC REGULATIONS FOR PREPARING\\
    GRADUATION INTERNSHIP REPORTS USING \LaTeX}
\end{center}

\vspace{0.5cm}

{\fontsize{13pt}{17pt}\selectfont
\setstretch{1.3}
\justifying

\noindent
\textbf{Note:}
This section is provided for guidance only and is included in the main
document using the command \verb|\input{Regulations.tex}|.
Students shall remove or comment out this command before submitting the
final internship report.

\vspace{0.4cm}

% ================================================================
% I. GENERAL REQUIREMENTS
% ================================================================

\noindent
\textbf{I.\quad GENERAL REQUIREMENTS}

\vspace{0.15cm}

The Graduation Internship Report shall be prepared using the official
\LaTeX{} template provided by the Faculty. Students shall preserve the
predefined page layout, typography, line spacing, heading styles,
numbering system, figure and table formatting, and bibliography settings
unless otherwise instructed by the Faculty or the supervisor.

The report shall be written in English.

Students are responsible for the academic and technical content of the
report. The official \LaTeX{} template is responsible for standardized
document formatting.

\vspace{0.3cm}

% ================================================================
% II. REPORT STRUCTURE
% ================================================================

\noindent
\textbf{II.\quad REPORT STRUCTURE}

\vspace{0.15cm}

The report shall normally be arranged in the following order:

\begin{enumerate}[label=\arabic*., leftmargin=1.2cm, itemsep=2pt]
    \item Outer Cover;
    \item Comments of the Internship Host Organization;
    \item Supervisor's Evaluation and Graduation Internship Defense Results;
    \item Table of Contents;
    \item List of Symbols and Abbreviations, if applicable;
    \item List of Figures;
    \item List of Tables;
    \item Abstract;
    \item Main Chapters of the Internship Report;
    \item References;
    \item Appendices, if applicable.
\end{enumerate}

The preliminary pages shall not be numbered as chapters. The main report
content shall be organized into consecutively numbered chapters beginning
with Chapter 1.

\vspace{0.3cm}

% ================================================================
% III. PAGE SIZE AND MARGINS
% ================================================================

\noindent
\textbf{III.\quad PAGE SIZE AND MARGINS}

\vspace{0.15cm}

The report shall be formatted on A4 paper with the following margins:

\begin{center}
\begin{tabular}{ll}
    Left margin:   & 3.5 cm \\
    Right margin:  & 2.0 cm \\
    Top margin:    & 2.5 cm \\
    Bottom margin: & 2.5 cm
\end{tabular}
\end{center}

These settings are predefined in the official template and shall not be
modified by students.

\vspace{0.3cm}

% ================================================================
% IV. FONT, FONT SIZE, AND LINE SPACING
% ================================================================

\noindent
\textbf{IV.\quad FONT, FONT SIZE, AND LINE SPACING}

\vspace{0.15cm}

The main text shall use a 13 pt font size with 1.3-line spacing.
A Times-family typeface shall be used in accordance with the official
template.

The text shall be fully justified.

Students shall not use repeated spaces, repeated manual line breaks, or
multiple blank lines to position text. Paragraph and heading spacing shall
be controlled by the template.

\vspace{0.3cm}

% ================================================================
% V. OUTER COVER
% ================================================================

\noindent
\textbf{V.\quad OUTER COVER}

\vspace{0.15cm}

The outer cover shall follow the official format provided in the template
and shall contain the following information:

\begin{itemize}[leftmargin=0.8cm, itemsep=2pt]
    \item Thai Nguyen University of Technology;
    \item Faculty of International Training;
    \item Engineering and Technology Department;
    \item University logo and Faculty logo;
    \item the title \textbf{INTERNSHIP REPORT};
    \item Internship Organization;
    \item Author;
    \item Class;
    \item Supervisor;
    \item place and year of completion.
\end{itemize}

Students shall only replace the required information fields and shall not
modify the cover-page layout unless instructed otherwise.

\vspace{0.3cm}

% ================================================================
% VI. COMMENTS OF THE INTERNSHIP HOST ORGANIZATION
% ================================================================

\noindent
\textbf{VI.\quad COMMENTS OF THE INTERNSHIP HOST ORGANIZATION}

\vspace{0.15cm}

The page entitled \textbf{COMMENTS OF THE INTERNSHIP HOST ORGANIZATION}
is reserved for the internship host organization.

It includes:

\begin{itemize}[leftmargin=0.8cm, itemsep=2pt]
    \item the student's full name;
    \item class;
    \item comments from the host organization;
    \item score;
    \item date;
    \item signature and official seal.
\end{itemize}

Students shall not alter or complete official evaluation fields unless
specifically authorized.

\vspace{0.3cm}

% ================================================================
% VII. SUPERVISOR'S EVALUATION AND DEFENSE RESULTS
% ================================================================

\noindent
\textbf{VII.\quad SUPERVISOR'S EVALUATION AND DEFENSE RESULTS}

\vspace{0.15cm}

The section entitled
\textbf{SUPERVISOR'S EVALUATION AND GRADUATION INTERNSHIP DEFENSE RESULTS}
is reserved for the supervisor and defense committee.

The evaluation may include:

\begin{itemize}[leftmargin=0.8cm, itemsep=2pt]
    \item Internship Progress and Student's Attitude;
    \item Report Content;
    \item Report Presentation;
    \item Other Comments;
    \item Supervisor's Opinion;
    \item Defense Result;
    \item Score.
\end{itemize}

Students shall not enter evaluation scores or defense results before they
are officially determined.

\vspace{0.3cm}

% ================================================================
% VIII. CHAPTERS, SECTIONS, AND SUBSECTIONS
% ================================================================

\noindent
\textbf{VIII.\quad CHAPTERS, SECTIONS, AND SUBSECTIONS}

\vspace{0.15cm}

The report structure shall be generated using the standard \LaTeX{}
commands:

\begin{verbatim}
\chapter{...}
\section{...}
\subsection{...}
\end{verbatim}

Chapter, section, and subsection numbers shall be generated automatically
by \LaTeX{}.

Students shall not manually type chapter or section numbers and shall not
reproduce heading styles using commands such as \verb|\textbf|,
\verb|\fontsize|, or \verb|\centering| instead of the appropriate
structural command.

Excessive subdivision should be avoided. In most cases, the hierarchy
Chapter--Section--Subsection is sufficient.

\vspace{0.3cm}

% ================================================================
% IX. ORGANIZATION OF THE MAIN CONTENT
% ================================================================

\noindent
\textbf{IX.\quad ORGANIZATION OF THE MAIN CONTENT}

\vspace{0.15cm}

The specific content of the report shall depend on the internship
organization and the assigned tasks. A typical structure may include:

\begin{itemize}[leftmargin=0.8cm, itemsep=2pt]
    \item \textbf{Chapter 1. Overview of the Internship Organization};
    \item \textbf{Chapter 2. Internship Activities and Assigned Tasks};
    \item \textbf{Chapter 3. Results, Experience, and Evaluation}.
\end{itemize}

These chapter titles are recommendations only and may be adapted to the
actual internship content.

The report shall reflect work actually performed, observed, analyzed, or
learned during the internship. Unrelated theoretical material should not
be included merely to increase report length.

\vspace{0.3cm}

% ================================================================
% X. TABLE OF CONTENTS
% ================================================================

\noindent
\textbf{X.\quad TABLE OF CONTENTS}

\vspace{0.15cm}

The Table of Contents shall be generated automatically using:

\begin{verbatim}
\tableofcontents
\end{verbatim}

Students shall not create or manually edit page numbers in the Table of
Contents.

The document shall be compiled sufficiently to ensure that all page
numbers and entries are correctly updated.

\vspace{0.3cm}

% ================================================================
% XI. LIST OF SYMBOLS AND ABBREVIATIONS
% ================================================================

\noindent
\textbf{XI.\quad LIST OF SYMBOLS AND ABBREVIATIONS}

\vspace{0.15cm}

Symbols and abbreviations that are used repeatedly in the report shall be
listed in the \textbf{List of Symbols and Abbreviations}.

The list may contain the following columns:

\begin{center}
\begin{tabular}{llll}
    No. & Symbol / Abbreviation & Unit & Definition
\end{tabular}
\end{center}

Only symbols and abbreviations actually used in the final report shall be
retained. Example entries contained in the template shall be removed if
they are not relevant to the student's report.

Each abbreviation should normally be defined when it first appears in
the main text.

\vspace{0.3cm}

% ================================================================
% XII. FIGURES, DIAGRAMS, AND PHOTOGRAPHS
% ================================================================

\noindent
\textbf{XII.\quad FIGURES, DIAGRAMS, AND PHOTOGRAPHS}

\vspace{0.15cm}

Figures shall be clear, legible, and directly relevant to the report
content.

Figures shall be inserted using the standard \verb|figure| environment.

Example:

\begin{verbatim}
\begin{figure}[H]
    \centering
    \includegraphics[width=0.8\textwidth]
                    {figures/company_structure.pdf}
    \caption{Organizational structure of the company}
    \label{fig:company-structure}
\end{figure}
\end{verbatim}

Figure captions shall be placed below the figure.

Each referenced figure shall have a unique \verb|\label| and shall be
referenced in the text using, for example:

\begin{verbatim}
Figure~\ref{fig:company-structure}
\end{verbatim}

Students shall not manually type figure numbers.

Figures reproduced or adapted from external sources shall include an
appropriate citation.

Photographs taken during the internship shall be relevant and sufficiently
clear. Students shall not include confidential, restricted, or proprietary
information without permission from the internship organization.

Vector graphics or high-resolution images are recommended whenever
possible.

\vspace{0.3cm}

% ================================================================
% XIII. TABLES
% ================================================================

\noindent
\textbf{XIII.\quad TABLES}

\vspace{0.15cm}

Tables shall be created using suitable \LaTeX{} environments such as
\verb|table|, \verb|tabular|, \verb|tabularx|, or \verb|longtable|.

Table captions shall be placed above the table.

Example:

\begin{verbatim}
\begin{table}[H]
    \centering
    \caption{Internship schedule}
    \label{tab:internship-schedule}

    \begin{tabular}{lll}
        \toprule
        Week & Task & Result \\
        \midrule
        1 & Orientation & Completed \\
        2 & System study & Completed \\
        \bottomrule
    \end{tabular}
\end{table}
\end{verbatim}

Tables shall be referenced using:

\begin{verbatim}
Table~\ref{tab:internship-schedule}
\end{verbatim}

Students shall not manually type table numbers.

Tables should be created directly in \LaTeX{} rather than inserted as
screenshots whenever practical.

\vspace{0.3cm}

% ================================================================
% XIV. EQUATIONS AND TECHNICAL QUANTITIES
% ================================================================

\noindent
\textbf{XIV.\quad EQUATIONS AND TECHNICAL QUANTITIES}

\vspace{0.15cm}

Mathematical expressions shall be typeset directly using \LaTeX{}.
Images or screenshots of equations shall not be used.

Short expressions may be written in inline mathematical mode, for example:

\begin{verbatim}
$R = 0.7\,\Omega$
\end{verbatim}

Important equations shall use a numbered equation environment.

Example:

\begin{verbatim}
\begin{equation}
    P = UI
    \label{eq:power}
\end{equation}
\end{verbatim}

Equations shall be referenced using:

\begin{verbatim}
Equation~\eqref{eq:power}
\end{verbatim}

Equation numbers shall not be manually entered.

SI units shall be used wherever applicable. The \verb|siunitx| package
should be used where appropriate.

Examples:

\begin{verbatim}
\SI{220}{\volt}
\SI{50}{\hertz}
\SI{2.5}{\ampere}
\end{verbatim}

Notation and units shall be consistent throughout the report.

\vspace{0.3cm}

% ================================================================
% XV. SOURCE CODE
% ================================================================

\noindent
\textbf{XV.\quad SOURCE CODE}

\vspace{0.15cm}

If the internship involves programming, PLCs, microcontrollers,
simulation, or software development, relevant source code may be included
in the report.

Source code shall be presented using the \verb|listings| environment or
another environment defined by the official template.

Example:

\begin{verbatim}
\begin{lstlisting}[language=st]
IF Start = TRUE THEN
    Motor := TRUE;
ELSE
    Motor := FALSE;
END_IF;
\end{lstlisting}
\end{verbatim}

Only code necessary for explaining the work should appear in the main
report. Long source-code listings should be moved to an appendix or
submitted as supplementary files if required.

Source code should not normally be inserted as screenshots.

\vspace{0.3cm}

% ================================================================
% XVI. ABSTRACT
% ================================================================

\noindent
\textbf{XVI.\quad ABSTRACT}

\vspace{0.15cm}

The Abstract shall briefly summarize:

\begin{itemize}[leftmargin=0.8cm, itemsep=2pt]
    \item the internship organization;
    \item the objectives of the internship;
    \item the main tasks performed;
    \item the main results, knowledge, and skills obtained.
\end{itemize}

The Abstract should be concise and should not reproduce the full content
of the report chapter by chapter.

\vspace{0.3cm}

% ================================================================
% XVII. REFERENCES AND CITATIONS
% ================================================================

\noindent
\textbf{XVII.\quad REFERENCES AND CITATIONS}

\vspace{0.15cm}

All books, journal papers, conference papers, standards, websites,
datasheets, manuals, technical documents, figures, tables, or other
external sources used in preparing the report shall be properly cited.

References shall be managed through the bibliography system provided with
the official template.

Bibliographic entries shall normally be stored in:

\begin{verbatim}
References.bib
\end{verbatim}

References shall be cited in the report using:

\begin{verbatim}
\cite{reference-key}
\end{verbatim}

Students shall not manually type reference numbers such as ``[1]'',
``[2]'', etc.

Titles of foreign-language publications shall remain in their original
language in the bibliography.

All borrowed ideas, data, figures, tables, equations, methods, or results
shall be properly acknowledged.

\vspace{0.3cm}

% ================================================================
% XVIII. CROSS-REFERENCING
% ================================================================

\noindent
\textbf{XVIII.\quad CROSS-REFERENCING}

\vspace{0.15cm}

Numbered elements shall be cross-referenced using the standard \LaTeX{}
mechanism.

Use the following commands where appropriate:

\begin{verbatim}
\label{...}
\ref{...}
\eqref{...}
\end{verbatim}

For example:

\begin{verbatim}
As shown in Figure~\ref{fig:system}, ...
\end{verbatim}

Students shall not manually type numbers such as ``Figure 2.3'' when
those numbers can be generated automatically.

\vspace{0.3cm}

% ================================================================
% XIX. PAGE NUMBERING
% ================================================================

\noindent
\textbf{XIX.\quad PAGE NUMBERING}

\vspace{0.15cm}

Cover pages and official evaluation forms at the beginning of the report
shall not display page numbers.

The main body of the report shall use Arabic page numbers beginning from
Chapter 1.

Page numbering shall be generated automatically by the official template.
Students shall not manually enter page numbers.

\vspace{0.3cm}

% ================================================================
% XX. PROJECT FILE ORGANIZATION
% ================================================================

\noindent
\textbf{XX.\quad PROJECT FILE ORGANIZATION}

\vspace{0.15cm}

Students should maintain a clear and organized \LaTeX{} project
structure. A recommended structure is:

\begin{verbatim}
InternshipReport/
|
|-- main.tex
|-- References.bib
|
|-- figures/
|   |-- company.pdf
|   |-- system.png
|   `-- ...
|
`-- chapters/
    |-- chapter1.tex
    |-- chapter2.tex
    `-- chapter3.tex
\end{verbatim}

Figures should normally be stored in the \verb|figures| directory
specified by the official template.

File and directory names should preferably contain Latin letters,
numbers, hyphens, or underscores.

Recommended examples include:

\begin{verbatim}
company_structure.pdf
motor_control.png
chapter_01.tex
\end{verbatim}

Spaces and unnecessary special characters in file names should be
avoided where possible.

\vspace{0.3cm}

% ================================================================
% XXI. DOCUMENT COMPILATION AND FINAL CHECK
% ================================================================

\noindent
\textbf{XXI.\quad DOCUMENT COMPILATION AND FINAL CHECK}

\vspace{0.15cm}

Before submission, the complete \LaTeX{} project shall compile
successfully and produce the final PDF report.

The final document shall not contain unresolved references such as:

\begin{verbatim}
??
[?]
\end{verbatim}

Before generating the final PDF, students shall check:

\begin{itemize}[leftmargin=0.8cm, itemsep=2pt]
    \item the Table of Contents;
    \item page numbering;
    \item chapter, section, and subsection numbering;
    \item figure numbering;
    \item table numbering;
    \item equation numbering;
    \item cross-references;
    \item citations;
    \item bibliography;
    \item font display;
    \item mathematical symbols;
    \item image quality;
    \item completeness of all required forms.
\end{itemize}

The document may be compiled using TeX Live, MiKTeX, or another
compatible \LaTeX{} environment. Regardless of the compilation
environment, the resulting PDF shall conform to the official Faculty
template.

\vspace{0.3cm}

% ================================================================
% XXII. CONTENT THAT STUDENTS MAY EDIT
% ================================================================

\noindent
\textbf{XXII.\quad CONTENT THAT STUDENTS MAY EDIT}

\vspace{0.15cm}

Students may edit or replace:

\begin{itemize}[leftmargin=0.8cm, itemsep=2pt]
    \item student name;
    \item class;
    \item internship organization;
    \item supervisor name;
    \item Abstract;
    \item List of Symbols and Abbreviations;
    \item chapter titles and chapter content;
    \item figures and photographs;
    \item tables;
    \item equations;
    \item citations and references.
\end{itemize}

Students shall not modify the main document formatting settings in the
preamble, including the document class, page geometry, line spacing,
heading formatting, table-of-contents formatting, font configuration,
page layout, or cover-page macros unless specifically instructed.

\vspace{0.3cm}

% ================================================================
% XXIII. FINAL SUBMISSION REQUIREMENTS
% ================================================================

\noindent
\textbf{XXIII.\quad FINAL SUBMISSION REQUIREMENTS}

\vspace{0.15cm}

Before submission, the final report shall:

\begin{itemize}[leftmargin=0.8cm, itemsep=2pt]
    \item compile without errors that affect the report content or presentation;
    \item contain all required forms and report sections;
    \item contain correct student, class, internship organization, and supervisor information;
    \item contain an updated Table of Contents, List of Figures, and List of Tables;
    \item contain no unused example content from the template;
    \item contain no placeholder text such as \textit{Title of Chapter 1},
    \textit{Section 1.1}, or other sample entries;
    \item contain all required citations and references;
    \item contain no unresolved cross-references.
\end{itemize}

If the Faculty requires submission of the \LaTeX{} source, the complete
project shall include all \verb|.tex| files, the bibliography database,
figures, and any other files required to reproduce the submitted PDF.

\vspace{0.4cm}

\noindent
\textbf{Important:}
Before submitting the final internship report, students shall remove or
comment out the following command in the main document:

\begin{verbatim}
\input{Regulations.tex}
\end{verbatim}

For example:

\begin{verbatim}
% \input{Regulations.tex}
\end{verbatim}

This removes the guidance pages from the final submitted report while
preserving the regulations file in the \LaTeX{} project for reference.

} % End 13pt content

\clearpage
|.
Students shall remove or comment out this command before submitting the
final internship report.

\vspace{0.4cm}

% ================================================================
% I. GENERAL REQUIREMENTS
% ================================================================

\noindent
\textbf{I.\quad GENERAL REQUIREMENTS}

\vspace{0.15cm}

The Graduation Internship Report shall be prepared using the official
\LaTeX{} template provided by the Faculty. Students shall preserve the
predefined page layout, typography, line spacing, heading styles,
numbering system, figure and table formatting, and bibliography settings
unless otherwise instructed by the Faculty or the supervisor.

The report shall be written in English.

Students are responsible for the academic and technical content of the
report. The official \LaTeX{} template is responsible for standardized
document formatting.

\vspace{0.3cm}

% ================================================================
% II. REPORT STRUCTURE
% ================================================================

\noindent
\textbf{II.\quad REPORT STRUCTURE}

\vspace{0.15cm}

The report shall normally be arranged in the following order:

\begin{enumerate}[label=\arabic*., leftmargin=1.2cm, itemsep=2pt]
    \item Outer Cover;
    \item Comments of the Internship Host Organization;
    \item Supervisor's Evaluation and Graduation Internship Defense Results;
    \item Table of Contents;
    \item List of Symbols and Abbreviations, if applicable;
    \item List of Figures;
    \item List of Tables;
    \item Abstract;
    \item Main Chapters of the Internship Report;
    \item References;
    \item Appendices, if applicable.
\end{enumerate}

The preliminary pages shall not be numbered as chapters. The main report
content shall be organized into consecutively numbered chapters beginning
with Chapter 1.

\vspace{0.3cm}

% ================================================================
% III. PAGE SIZE AND MARGINS
% ================================================================

\noindent
\textbf{III.\quad PAGE SIZE AND MARGINS}

\vspace{0.15cm}

The report shall be formatted on A4 paper with the following margins:

\begin{center}
\begin{tabular}{ll}
    Left margin:   & 3.5 cm \\
    Right margin:  & 2.0 cm \\
    Top margin:    & 2.5 cm \\
    Bottom margin: & 2.5 cm
\end{tabular}
\end{center}

These settings are predefined in the official template and shall not be
modified by students.

\vspace{0.3cm}

% ================================================================
% IV. FONT, FONT SIZE, AND LINE SPACING
% ================================================================

\noindent
\textbf{IV.\quad FONT, FONT SIZE, AND LINE SPACING}

\vspace{0.15cm}

The main text shall use a 13 pt font size with 1.3-line spacing.
A Times-family typeface shall be used in accordance with the official
template.

The text shall be fully justified.

Students shall not use repeated spaces, repeated manual line breaks, or
multiple blank lines to position text. Paragraph and heading spacing shall
be controlled by the template.

\vspace{0.3cm}

% ================================================================
% V. OUTER COVER
% ================================================================

\noindent
\textbf{V.\quad OUTER COVER}

\vspace{0.15cm}

The outer cover shall follow the official format provided in the template
and shall contain the following information:

\begin{itemize}[leftmargin=0.8cm, itemsep=2pt]
    \item Thai Nguyen University of Technology;
    \item Faculty of International Training;
    \item Engineering and Technology Department;
    \item University logo and Faculty logo;
    \item the title \textbf{INTERNSHIP REPORT};
    \item Internship Organization;
    \item Author;
    \item Class;
    \item Supervisor;
    \item place and year of completion.
\end{itemize}

Students shall only replace the required information fields and shall not
modify the cover-page layout unless instructed otherwise.

\vspace{0.3cm}

% ================================================================
% VI. COMMENTS OF THE INTERNSHIP HOST ORGANIZATION
% ================================================================

\noindent
\textbf{VI.\quad COMMENTS OF THE INTERNSHIP HOST ORGANIZATION}

\vspace{0.15cm}

The page entitled \textbf{COMMENTS OF THE INTERNSHIP HOST ORGANIZATION}
is reserved for the internship host organization.

It includes:

\begin{itemize}[leftmargin=0.8cm, itemsep=2pt]
    \item the student's full name;
    \item class;
    \item comments from the host organization;
    \item score;
    \item date;
    \item signature and official seal.
\end{itemize}

Students shall not alter or complete official evaluation fields unless
specifically authorized.

\vspace{0.3cm}

% ================================================================
% VII. SUPERVISOR'S EVALUATION AND DEFENSE RESULTS
% ================================================================

\noindent
\textbf{VII.\quad SUPERVISOR'S EVALUATION AND DEFENSE RESULTS}

\vspace{0.15cm}

The section entitled
\textbf{SUPERVISOR'S EVALUATION AND GRADUATION INTERNSHIP DEFENSE RESULTS}
is reserved for the supervisor and defense committee.

The evaluation may include:

\begin{itemize}[leftmargin=0.8cm, itemsep=2pt]
    \item Internship Progress and Student's Attitude;
    \item Report Content;
    \item Report Presentation;
    \item Other Comments;
    \item Supervisor's Opinion;
    \item Defense Result;
    \item Score.
\end{itemize}

Students shall not enter evaluation scores or defense results before they
are officially determined.

\vspace{0.3cm}

% ================================================================
% VIII. CHAPTERS, SECTIONS, AND SUBSECTIONS
% ================================================================

\noindent
\textbf{VIII.\quad CHAPTERS, SECTIONS, AND SUBSECTIONS}

\vspace{0.15cm}

The report structure shall be generated using the standard \LaTeX{}
commands:

\begin{verbatim}
\chapter{...}
\section{...}
\subsection{...}
\end{verbatim}

Chapter, section, and subsection numbers shall be generated automatically
by \LaTeX{}.

Students shall not manually type chapter or section numbers and shall not
reproduce heading styles using commands such as \verb|\textbf|,
\verb|\fontsize|, or \verb|\centering| instead of the appropriate
structural command.

Excessive subdivision should be avoided. In most cases, the hierarchy
Chapter--Section--Subsection is sufficient.

\vspace{0.3cm}

% ================================================================
% IX. ORGANIZATION OF THE MAIN CONTENT
% ================================================================

\noindent
\textbf{IX.\quad ORGANIZATION OF THE MAIN CONTENT}

\vspace{0.15cm}

The specific content of the report shall depend on the internship
organization and the assigned tasks. A typical structure may include:

\begin{itemize}[leftmargin=0.8cm, itemsep=2pt]
    \item \textbf{Chapter 1. Overview of the Internship Organization};
    \item \textbf{Chapter 2. Internship Activities and Assigned Tasks};
    \item \textbf{Chapter 3. Results, Experience, and Evaluation}.
\end{itemize}

These chapter titles are recommendations only and may be adapted to the
actual internship content.

The report shall reflect work actually performed, observed, analyzed, or
learned during the internship. Unrelated theoretical material should not
be included merely to increase report length.

\vspace{0.3cm}

% ================================================================
% X. TABLE OF CONTENTS
% ================================================================

\noindent
\textbf{X.\quad TABLE OF CONTENTS}

\vspace{0.15cm}

The Table of Contents shall be generated automatically using:

\begin{verbatim}
\tableofcontents
\end{verbatim}

Students shall not create or manually edit page numbers in the Table of
Contents.

The document shall be compiled sufficiently to ensure that all page
numbers and entries are correctly updated.

\vspace{0.3cm}

% ================================================================
% XI. LIST OF SYMBOLS AND ABBREVIATIONS
% ================================================================

\noindent
\textbf{XI.\quad LIST OF SYMBOLS AND ABBREVIATIONS}

\vspace{0.15cm}

Symbols and abbreviations that are used repeatedly in the report shall be
listed in the \textbf{List of Symbols and Abbreviations}.

The list may contain the following columns:

\begin{center}
\begin{tabular}{llll}
    No. & Symbol / Abbreviation & Unit & Definition
\end{tabular}
\end{center}

Only symbols and abbreviations actually used in the final report shall be
retained. Example entries contained in the template shall be removed if
they are not relevant to the student's report.

Each abbreviation should normally be defined when it first appears in
the main text.

\vspace{0.3cm}

% ================================================================
% XII. FIGURES, DIAGRAMS, AND PHOTOGRAPHS
% ================================================================

\noindent
\textbf{XII.\quad FIGURES, DIAGRAMS, AND PHOTOGRAPHS}

\vspace{0.15cm}

Figures shall be clear, legible, and directly relevant to the report
content.

Figures shall be inserted using the standard \verb|figure| environment.

Example:

\begin{verbatim}
\begin{figure}[H]
    \centering
    \includegraphics[width=0.8\textwidth]
                    {figures/company_structure.pdf}
    \caption{Organizational structure of the company}
    \label{fig:company-structure}
\end{figure}
\end{verbatim}

Figure captions shall be placed below the figure.

Each referenced figure shall have a unique \verb|\label| and shall be
referenced in the text using, for example:

\begin{verbatim}
Figure~\ref{fig:company-structure}
\end{verbatim}

Students shall not manually type figure numbers.

Figures reproduced or adapted from external sources shall include an
appropriate citation.

Photographs taken during the internship shall be relevant and sufficiently
clear. Students shall not include confidential, restricted, or proprietary
information without permission from the internship organization.

Vector graphics or high-resolution images are recommended whenever
possible.

\vspace{0.3cm}

% ================================================================
% XIII. TABLES
% ================================================================

\noindent
\textbf{XIII.\quad TABLES}

\vspace{0.15cm}

Tables shall be created using suitable \LaTeX{} environments such as
\verb|table|, \verb|tabular|, \verb|tabularx|, or \verb|longtable|.

Table captions shall be placed above the table.

Example:

\begin{verbatim}
\begin{table}[H]
    \centering
    \caption{Internship schedule}
    \label{tab:internship-schedule}

    \begin{tabular}{lll}
        \toprule
        Week & Task & Result \\
        \midrule
        1 & Orientation & Completed \\
        2 & System study & Completed \\
        \bottomrule
    \end{tabular}
\end{table}
\end{verbatim}

Tables shall be referenced using:

\begin{verbatim}
Table~\ref{tab:internship-schedule}
\end{verbatim}

Students shall not manually type table numbers.

Tables should be created directly in \LaTeX{} rather than inserted as
screenshots whenever practical.

\vspace{0.3cm}

% ================================================================
% XIV. EQUATIONS AND TECHNICAL QUANTITIES
% ================================================================

\noindent
\textbf{XIV.\quad EQUATIONS AND TECHNICAL QUANTITIES}

\vspace{0.15cm}

Mathematical expressions shall be typeset directly using \LaTeX{}.
Images or screenshots of equations shall not be used.

Short expressions may be written in inline mathematical mode, for example:

\begin{verbatim}
$R = 0.7\,\Omega$
\end{verbatim}

Important equations shall use a numbered equation environment.

Example:

\begin{verbatim}
\begin{equation}
    P = UI
    \label{eq:power}
\end{equation}
\end{verbatim}

Equations shall be referenced using:

\begin{verbatim}
Equation~\eqref{eq:power}
\end{verbatim}

Equation numbers shall not be manually entered.

SI units shall be used wherever applicable. The \verb|siunitx| package
should be used where appropriate.

Examples:

\begin{verbatim}
\SI{220}{\volt}
\SI{50}{\hertz}
\SI{2.5}{\ampere}
\end{verbatim}

Notation and units shall be consistent throughout the report.

\vspace{0.3cm}

% ================================================================
% XV. SOURCE CODE
% ================================================================

\noindent
\textbf{XV.\quad SOURCE CODE}

\vspace{0.15cm}

If the internship involves programming, PLCs, microcontrollers,
simulation, or software development, relevant source code may be included
in the report.

Source code shall be presented using the \verb|listings| environment or
another environment defined by the official template.

Example:

\begin{verbatim}
\begin{lstlisting}[language=st]
IF Start = TRUE THEN
    Motor := TRUE;
ELSE
    Motor := FALSE;
END_IF;
\end{lstlisting}
\end{verbatim}

Only code necessary for explaining the work should appear in the main
report. Long source-code listings should be moved to an appendix or
submitted as supplementary files if required.

Source code should not normally be inserted as screenshots.

\vspace{0.3cm}

% ================================================================
% XVI. ABSTRACT
% ================================================================

\noindent
\textbf{XVI.\quad ABSTRACT}

\vspace{0.15cm}

The Abstract shall briefly summarize:

\begin{itemize}[leftmargin=0.8cm, itemsep=2pt]
    \item the internship organization;
    \item the objectives of the internship;
    \item the main tasks performed;
    \item the main results, knowledge, and skills obtained.
\end{itemize}

The Abstract should be concise and should not reproduce the full content
of the report chapter by chapter.

\vspace{0.3cm}

% ================================================================
% XVII. REFERENCES AND CITATIONS
% ================================================================

\noindent
\textbf{XVII.\quad REFERENCES AND CITATIONS}

\vspace{0.15cm}

All books, journal papers, conference papers, standards, websites,
datasheets, manuals, technical documents, figures, tables, or other
external sources used in preparing the report shall be properly cited.

References shall be managed through the bibliography system provided with
the official template.

Bibliographic entries shall normally be stored in:

\begin{verbatim}
References.bib
\end{verbatim}

References shall be cited in the report using:

\begin{verbatim}
\cite{reference-key}
\end{verbatim}

Students shall not manually type reference numbers such as ``[1]'',
``[2]'', etc.

Titles of foreign-language publications shall remain in their original
language in the bibliography.

All borrowed ideas, data, figures, tables, equations, methods, or results
shall be properly acknowledged.

\vspace{0.3cm}

% ================================================================
% XVIII. CROSS-REFERENCING
% ================================================================

\noindent
\textbf{XVIII.\quad CROSS-REFERENCING}

\vspace{0.15cm}

Numbered elements shall be cross-referenced using the standard \LaTeX{}
mechanism.

Use the following commands where appropriate:

\begin{verbatim}
\label{...}
\ref{...}
\eqref{...}
\end{verbatim}

For example:

\begin{verbatim}
As shown in Figure~\ref{fig:system}, ...
\end{verbatim}

Students shall not manually type numbers such as ``Figure 2.3'' when
those numbers can be generated automatically.

\vspace{0.3cm}

% ================================================================
% XIX. PAGE NUMBERING
% ================================================================

\noindent
\textbf{XIX.\quad PAGE NUMBERING}

\vspace{0.15cm}

Cover pages and official evaluation forms at the beginning of the report
shall not display page numbers.

The main body of the report shall use Arabic page numbers beginning from
Chapter 1.

Page numbering shall be generated automatically by the official template.
Students shall not manually enter page numbers.

\vspace{0.3cm}

% ================================================================
% XX. PROJECT FILE ORGANIZATION
% ================================================================

\noindent
\textbf{XX.\quad PROJECT FILE ORGANIZATION}

\vspace{0.15cm}

Students should maintain a clear and organized \LaTeX{} project
structure. A recommended structure is:

\begin{verbatim}
InternshipReport/
|
|-- main.tex
|-- References.bib
|
|-- figures/
|   |-- company.pdf
|   |-- system.png
|   `-- ...
|
`-- chapters/
    |-- chapter1.tex
    |-- chapter2.tex
    `-- chapter3.tex
\end{verbatim}

Figures should normally be stored in the \verb|figures| directory
specified by the official template.

File and directory names should preferably contain Latin letters,
numbers, hyphens, or underscores.

Recommended examples include:

\begin{verbatim}
company_structure.pdf
motor_control.png
chapter_01.tex
\end{verbatim}

Spaces and unnecessary special characters in file names should be
avoided where possible.

\vspace{0.3cm}

% ================================================================
% XXI. DOCUMENT COMPILATION AND FINAL CHECK
% ================================================================

\noindent
\textbf{XXI.\quad DOCUMENT COMPILATION AND FINAL CHECK}

\vspace{0.15cm}

Before submission, the complete \LaTeX{} project shall compile
successfully and produce the final PDF report.

The final document shall not contain unresolved references such as:

\begin{verbatim}
??
[?]
\end{verbatim}

Before generating the final PDF, students shall check:

\begin{itemize}[leftmargin=0.8cm, itemsep=2pt]
    \item the Table of Contents;
    \item page numbering;
    \item chapter, section, and subsection numbering;
    \item figure numbering;
    \item table numbering;
    \item equation numbering;
    \item cross-references;
    \item citations;
    \item bibliography;
    \item font display;
    \item mathematical symbols;
    \item image quality;
    \item completeness of all required forms.
\end{itemize}

The document may be compiled using TeX Live, MiKTeX, or another
compatible \LaTeX{} environment. Regardless of the compilation
environment, the resulting PDF shall conform to the official Faculty
template.

\vspace{0.3cm}

% ================================================================
% XXII. CONTENT THAT STUDENTS MAY EDIT
% ================================================================

\noindent
\textbf{XXII.\quad CONTENT THAT STUDENTS MAY EDIT}

\vspace{0.15cm}

Students may edit or replace:

\begin{itemize}[leftmargin=0.8cm, itemsep=2pt]
    \item student name;
    \item class;
    \item internship organization;
    \item supervisor name;
    \item Abstract;
    \item List of Symbols and Abbreviations;
    \item chapter titles and chapter content;
    \item figures and photographs;
    \item tables;
    \item equations;
    \item citations and references.
\end{itemize}

Students shall not modify the main document formatting settings in the
preamble, including the document class, page geometry, line spacing,
heading formatting, table-of-contents formatting, font configuration,
page layout, or cover-page macros unless specifically instructed.

\vspace{0.3cm}

% ================================================================
% XXIII. FINAL SUBMISSION REQUIREMENTS
% ================================================================

\noindent
\textbf{XXIII.\quad FINAL SUBMISSION REQUIREMENTS}

\vspace{0.15cm}

Before submission, the final report shall:

\begin{itemize}[leftmargin=0.8cm, itemsep=2pt]
    \item compile without errors that affect the report content or presentation;
    \item contain all required forms and report sections;
    \item contain correct student, class, internship organization, and supervisor information;
    \item contain an updated Table of Contents, List of Figures, and List of Tables;
    \item contain no unused example content from the template;
    \item contain no placeholder text such as \textit{Title of Chapter 1},
    \textit{Section 1.1}, or other sample entries;
    \item contain all required citations and references;
    \item contain no unresolved cross-references.
\end{itemize}

If the Faculty requires submission of the \LaTeX{} source, the complete
project shall include all \verb|.tex| files, the bibliography database,
figures, and any other files required to reproduce the submitted PDF.

\vspace{0.4cm}

\noindent
\textbf{Important:}
Before submitting the final internship report, students shall remove or
comment out the following command in the main document:

\begin{verbatim}

% =================================================================
% REGULATIONS FOR PREPARING THE GRADUATION INTERNSHIP REPORT
% USING LATEX
% =================================================================
%
% This file is intended to be included from the main document with:
%
%     \input{Regulations.tex}
%
% IMPORTANT:
% This section is provided for guidance only. Students shall remove
% or comment out the \input{Regulations.tex} command before submitting
% the final internship report.
%
% =================================================================

\clearpage
\thispagestyle{empty}

\begin{center}
    {\fontsize{16pt}{20pt}\selectfont\bfseries
    SPECIFIC REGULATIONS FOR PREPARING\\
    GRADUATION INTERNSHIP REPORTS USING \LaTeX}
\end{center}

\vspace{0.5cm}

{\fontsize{13pt}{17pt}\selectfont
\setstretch{1.3}
\justifying

\noindent
\textbf{Note:}
This section is provided for guidance only and is included in the main
document using the command \verb|\input{Regulations.tex}|.
Students shall remove or comment out this command before submitting the
final internship report.

\vspace{0.4cm}

% ================================================================
% I. GENERAL REQUIREMENTS
% ================================================================

\noindent
\textbf{I.\quad GENERAL REQUIREMENTS}

\vspace{0.15cm}

The Graduation Internship Report shall be prepared using the official
\LaTeX{} template provided by the Faculty. Students shall preserve the
predefined page layout, typography, line spacing, heading styles,
numbering system, figure and table formatting, and bibliography settings
unless otherwise instructed by the Faculty or the supervisor.

The report shall be written in English.

Students are responsible for the academic and technical content of the
report. The official \LaTeX{} template is responsible for standardized
document formatting.

\vspace{0.3cm}

% ================================================================
% II. REPORT STRUCTURE
% ================================================================

\noindent
\textbf{II.\quad REPORT STRUCTURE}

\vspace{0.15cm}

The report shall normally be arranged in the following order:

\begin{enumerate}[label=\arabic*., leftmargin=1.2cm, itemsep=2pt]
    \item Outer Cover;
    \item Comments of the Internship Host Organization;
    \item Supervisor's Evaluation and Graduation Internship Defense Results;
    \item Table of Contents;
    \item List of Symbols and Abbreviations, if applicable;
    \item List of Figures;
    \item List of Tables;
    \item Abstract;
    \item Main Chapters of the Internship Report;
    \item References;
    \item Appendices, if applicable.
\end{enumerate}

The preliminary pages shall not be numbered as chapters. The main report
content shall be organized into consecutively numbered chapters beginning
with Chapter 1.

\vspace{0.3cm}

% ================================================================
% III. PAGE SIZE AND MARGINS
% ================================================================

\noindent
\textbf{III.\quad PAGE SIZE AND MARGINS}

\vspace{0.15cm}

The report shall be formatted on A4 paper with the following margins:

\begin{center}
\begin{tabular}{ll}
    Left margin:   & 3.5 cm \\
    Right margin:  & 2.0 cm \\
    Top margin:    & 2.5 cm \\
    Bottom margin: & 2.5 cm
\end{tabular}
\end{center}

These settings are predefined in the official template and shall not be
modified by students.

\vspace{0.3cm}

% ================================================================
% IV. FONT, FONT SIZE, AND LINE SPACING
% ================================================================

\noindent
\textbf{IV.\quad FONT, FONT SIZE, AND LINE SPACING}

\vspace{0.15cm}

The main text shall use a 13 pt font size with 1.3-line spacing.
A Times-family typeface shall be used in accordance with the official
template.

The text shall be fully justified.

Students shall not use repeated spaces, repeated manual line breaks, or
multiple blank lines to position text. Paragraph and heading spacing shall
be controlled by the template.

\vspace{0.3cm}

% ================================================================
% V. OUTER COVER
% ================================================================

\noindent
\textbf{V.\quad OUTER COVER}

\vspace{0.15cm}

The outer cover shall follow the official format provided in the template
and shall contain the following information:

\begin{itemize}[leftmargin=0.8cm, itemsep=2pt]
    \item Thai Nguyen University of Technology;
    \item Faculty of International Training;
    \item Engineering and Technology Department;
    \item University logo and Faculty logo;
    \item the title \textbf{INTERNSHIP REPORT};
    \item Internship Organization;
    \item Author;
    \item Class;
    \item Supervisor;
    \item place and year of completion.
\end{itemize}

Students shall only replace the required information fields and shall not
modify the cover-page layout unless instructed otherwise.

\vspace{0.3cm}

% ================================================================
% VI. COMMENTS OF THE INTERNSHIP HOST ORGANIZATION
% ================================================================

\noindent
\textbf{VI.\quad COMMENTS OF THE INTERNSHIP HOST ORGANIZATION}

\vspace{0.15cm}

The page entitled \textbf{COMMENTS OF THE INTERNSHIP HOST ORGANIZATION}
is reserved for the internship host organization.

It includes:

\begin{itemize}[leftmargin=0.8cm, itemsep=2pt]
    \item the student's full name;
    \item class;
    \item comments from the host organization;
    \item score;
    \item date;
    \item signature and official seal.
\end{itemize}

Students shall not alter or complete official evaluation fields unless
specifically authorized.

\vspace{0.3cm}

% ================================================================
% VII. SUPERVISOR'S EVALUATION AND DEFENSE RESULTS
% ================================================================

\noindent
\textbf{VII.\quad SUPERVISOR'S EVALUATION AND DEFENSE RESULTS}

\vspace{0.15cm}

The section entitled
\textbf{SUPERVISOR'S EVALUATION AND GRADUATION INTERNSHIP DEFENSE RESULTS}
is reserved for the supervisor and defense committee.

The evaluation may include:

\begin{itemize}[leftmargin=0.8cm, itemsep=2pt]
    \item Internship Progress and Student's Attitude;
    \item Report Content;
    \item Report Presentation;
    \item Other Comments;
    \item Supervisor's Opinion;
    \item Defense Result;
    \item Score.
\end{itemize}

Students shall not enter evaluation scores or defense results before they
are officially determined.

\vspace{0.3cm}

% ================================================================
% VIII. CHAPTERS, SECTIONS, AND SUBSECTIONS
% ================================================================

\noindent
\textbf{VIII.\quad CHAPTERS, SECTIONS, AND SUBSECTIONS}

\vspace{0.15cm}

The report structure shall be generated using the standard \LaTeX{}
commands:

\begin{verbatim}
\chapter{...}
\section{...}
\subsection{...}
\end{verbatim}

Chapter, section, and subsection numbers shall be generated automatically
by \LaTeX{}.

Students shall not manually type chapter or section numbers and shall not
reproduce heading styles using commands such as \verb|\textbf|,
\verb|\fontsize|, or \verb|\centering| instead of the appropriate
structural command.

Excessive subdivision should be avoided. In most cases, the hierarchy
Chapter--Section--Subsection is sufficient.

\vspace{0.3cm}

% ================================================================
% IX. ORGANIZATION OF THE MAIN CONTENT
% ================================================================

\noindent
\textbf{IX.\quad ORGANIZATION OF THE MAIN CONTENT}

\vspace{0.15cm}

The specific content of the report shall depend on the internship
organization and the assigned tasks. A typical structure may include:

\begin{itemize}[leftmargin=0.8cm, itemsep=2pt]
    \item \textbf{Chapter 1. Overview of the Internship Organization};
    \item \textbf{Chapter 2. Internship Activities and Assigned Tasks};
    \item \textbf{Chapter 3. Results, Experience, and Evaluation}.
\end{itemize}

These chapter titles are recommendations only and may be adapted to the
actual internship content.

The report shall reflect work actually performed, observed, analyzed, or
learned during the internship. Unrelated theoretical material should not
be included merely to increase report length.

\vspace{0.3cm}

% ================================================================
% X. TABLE OF CONTENTS
% ================================================================

\noindent
\textbf{X.\quad TABLE OF CONTENTS}

\vspace{0.15cm}

The Table of Contents shall be generated automatically using:

\begin{verbatim}
\tableofcontents
\end{verbatim}

Students shall not create or manually edit page numbers in the Table of
Contents.

The document shall be compiled sufficiently to ensure that all page
numbers and entries are correctly updated.

\vspace{0.3cm}

% ================================================================
% XI. LIST OF SYMBOLS AND ABBREVIATIONS
% ================================================================

\noindent
\textbf{XI.\quad LIST OF SYMBOLS AND ABBREVIATIONS}

\vspace{0.15cm}

Symbols and abbreviations that are used repeatedly in the report shall be
listed in the \textbf{List of Symbols and Abbreviations}.

The list may contain the following columns:

\begin{center}
\begin{tabular}{llll}
    No. & Symbol / Abbreviation & Unit & Definition
\end{tabular}
\end{center}

Only symbols and abbreviations actually used in the final report shall be
retained. Example entries contained in the template shall be removed if
they are not relevant to the student's report.

Each abbreviation should normally be defined when it first appears in
the main text.

\vspace{0.3cm}

% ================================================================
% XII. FIGURES, DIAGRAMS, AND PHOTOGRAPHS
% ================================================================

\noindent
\textbf{XII.\quad FIGURES, DIAGRAMS, AND PHOTOGRAPHS}

\vspace{0.15cm}

Figures shall be clear, legible, and directly relevant to the report
content.

Figures shall be inserted using the standard \verb|figure| environment.

Example:

\begin{verbatim}
\begin{figure}[H]
    \centering
    \includegraphics[width=0.8\textwidth]
                    {figures/company_structure.pdf}
    \caption{Organizational structure of the company}
    \label{fig:company-structure}
\end{figure}
\end{verbatim}

Figure captions shall be placed below the figure.

Each referenced figure shall have a unique \verb|\label| and shall be
referenced in the text using, for example:

\begin{verbatim}
Figure~\ref{fig:company-structure}
\end{verbatim}

Students shall not manually type figure numbers.

Figures reproduced or adapted from external sources shall include an
appropriate citation.

Photographs taken during the internship shall be relevant and sufficiently
clear. Students shall not include confidential, restricted, or proprietary
information without permission from the internship organization.

Vector graphics or high-resolution images are recommended whenever
possible.

\vspace{0.3cm}

% ================================================================
% XIII. TABLES
% ================================================================

\noindent
\textbf{XIII.\quad TABLES}

\vspace{0.15cm}

Tables shall be created using suitable \LaTeX{} environments such as
\verb|table|, \verb|tabular|, \verb|tabularx|, or \verb|longtable|.

Table captions shall be placed above the table.

Example:

\begin{verbatim}
\begin{table}[H]
    \centering
    \caption{Internship schedule}
    \label{tab:internship-schedule}

    \begin{tabular}{lll}
        \toprule
        Week & Task & Result \\
        \midrule
        1 & Orientation & Completed \\
        2 & System study & Completed \\
        \bottomrule
    \end{tabular}
\end{table}
\end{verbatim}

Tables shall be referenced using:

\begin{verbatim}
Table~\ref{tab:internship-schedule}
\end{verbatim}

Students shall not manually type table numbers.

Tables should be created directly in \LaTeX{} rather than inserted as
screenshots whenever practical.

\vspace{0.3cm}

% ================================================================
% XIV. EQUATIONS AND TECHNICAL QUANTITIES
% ================================================================

\noindent
\textbf{XIV.\quad EQUATIONS AND TECHNICAL QUANTITIES}

\vspace{0.15cm}

Mathematical expressions shall be typeset directly using \LaTeX{}.
Images or screenshots of equations shall not be used.

Short expressions may be written in inline mathematical mode, for example:

\begin{verbatim}
$R = 0.7\,\Omega$
\end{verbatim}

Important equations shall use a numbered equation environment.

Example:

\begin{verbatim}
\begin{equation}
    P = UI
    \label{eq:power}
\end{equation}
\end{verbatim}

Equations shall be referenced using:

\begin{verbatim}
Equation~\eqref{eq:power}
\end{verbatim}

Equation numbers shall not be manually entered.

SI units shall be used wherever applicable. The \verb|siunitx| package
should be used where appropriate.

Examples:

\begin{verbatim}
\SI{220}{\volt}
\SI{50}{\hertz}
\SI{2.5}{\ampere}
\end{verbatim}

Notation and units shall be consistent throughout the report.

\vspace{0.3cm}

% ================================================================
% XV. SOURCE CODE
% ================================================================

\noindent
\textbf{XV.\quad SOURCE CODE}

\vspace{0.15cm}

If the internship involves programming, PLCs, microcontrollers,
simulation, or software development, relevant source code may be included
in the report.

Source code shall be presented using the \verb|listings| environment or
another environment defined by the official template.

Example:

\begin{verbatim}
\begin{lstlisting}[language=st]
IF Start = TRUE THEN
    Motor := TRUE;
ELSE
    Motor := FALSE;
END_IF;
\end{lstlisting}
\end{verbatim}

Only code necessary for explaining the work should appear in the main
report. Long source-code listings should be moved to an appendix or
submitted as supplementary files if required.

Source code should not normally be inserted as screenshots.

\vspace{0.3cm}

% ================================================================
% XVI. ABSTRACT
% ================================================================

\noindent
\textbf{XVI.\quad ABSTRACT}

\vspace{0.15cm}

The Abstract shall briefly summarize:

\begin{itemize}[leftmargin=0.8cm, itemsep=2pt]
    \item the internship organization;
    \item the objectives of the internship;
    \item the main tasks performed;
    \item the main results, knowledge, and skills obtained.
\end{itemize}

The Abstract should be concise and should not reproduce the full content
of the report chapter by chapter.

\vspace{0.3cm}

% ================================================================
% XVII. REFERENCES AND CITATIONS
% ================================================================

\noindent
\textbf{XVII.\quad REFERENCES AND CITATIONS}

\vspace{0.15cm}

All books, journal papers, conference papers, standards, websites,
datasheets, manuals, technical documents, figures, tables, or other
external sources used in preparing the report shall be properly cited.

References shall be managed through the bibliography system provided with
the official template.

Bibliographic entries shall normally be stored in:

\begin{verbatim}
References.bib
\end{verbatim}

References shall be cited in the report using:

\begin{verbatim}
\cite{reference-key}
\end{verbatim}

Students shall not manually type reference numbers such as ``[1]'',
``[2]'', etc.

Titles of foreign-language publications shall remain in their original
language in the bibliography.

All borrowed ideas, data, figures, tables, equations, methods, or results
shall be properly acknowledged.

\vspace{0.3cm}

% ================================================================
% XVIII. CROSS-REFERENCING
% ================================================================

\noindent
\textbf{XVIII.\quad CROSS-REFERENCING}

\vspace{0.15cm}

Numbered elements shall be cross-referenced using the standard \LaTeX{}
mechanism.

Use the following commands where appropriate:

\begin{verbatim}
\label{...}
\ref{...}
\eqref{...}
\end{verbatim}

For example:

\begin{verbatim}
As shown in Figure~\ref{fig:system}, ...
\end{verbatim}

Students shall not manually type numbers such as ``Figure 2.3'' when
those numbers can be generated automatically.

\vspace{0.3cm}

% ================================================================
% XIX. PAGE NUMBERING
% ================================================================

\noindent
\textbf{XIX.\quad PAGE NUMBERING}

\vspace{0.15cm}

Cover pages and official evaluation forms at the beginning of the report
shall not display page numbers.

The main body of the report shall use Arabic page numbers beginning from
Chapter 1.

Page numbering shall be generated automatically by the official template.
Students shall not manually enter page numbers.

\vspace{0.3cm}

% ================================================================
% XX. PROJECT FILE ORGANIZATION
% ================================================================

\noindent
\textbf{XX.\quad PROJECT FILE ORGANIZATION}

\vspace{0.15cm}

Students should maintain a clear and organized \LaTeX{} project
structure. A recommended structure is:

\begin{verbatim}
InternshipReport/
|
|-- main.tex
|-- References.bib
|
|-- figures/
|   |-- company.pdf
|   |-- system.png
|   `-- ...
|
`-- chapters/
    |-- chapter1.tex
    |-- chapter2.tex
    `-- chapter3.tex
\end{verbatim}

Figures should normally be stored in the \verb|figures| directory
specified by the official template.

File and directory names should preferably contain Latin letters,
numbers, hyphens, or underscores.

Recommended examples include:

\begin{verbatim}
company_structure.pdf
motor_control.png
chapter_01.tex
\end{verbatim}

Spaces and unnecessary special characters in file names should be
avoided where possible.

\vspace{0.3cm}

% ================================================================
% XXI. DOCUMENT COMPILATION AND FINAL CHECK
% ================================================================

\noindent
\textbf{XXI.\quad DOCUMENT COMPILATION AND FINAL CHECK}

\vspace{0.15cm}

Before submission, the complete \LaTeX{} project shall compile
successfully and produce the final PDF report.

The final document shall not contain unresolved references such as:

\begin{verbatim}
??
[?]
\end{verbatim}

Before generating the final PDF, students shall check:

\begin{itemize}[leftmargin=0.8cm, itemsep=2pt]
    \item the Table of Contents;
    \item page numbering;
    \item chapter, section, and subsection numbering;
    \item figure numbering;
    \item table numbering;
    \item equation numbering;
    \item cross-references;
    \item citations;
    \item bibliography;
    \item font display;
    \item mathematical symbols;
    \item image quality;
    \item completeness of all required forms.
\end{itemize}

The document may be compiled using TeX Live, MiKTeX, or another
compatible \LaTeX{} environment. Regardless of the compilation
environment, the resulting PDF shall conform to the official Faculty
template.

\vspace{0.3cm}

% ================================================================
% XXII. CONTENT THAT STUDENTS MAY EDIT
% ================================================================

\noindent
\textbf{XXII.\quad CONTENT THAT STUDENTS MAY EDIT}

\vspace{0.15cm}

Students may edit or replace:

\begin{itemize}[leftmargin=0.8cm, itemsep=2pt]
    \item student name;
    \item class;
    \item internship organization;
    \item supervisor name;
    \item Abstract;
    \item List of Symbols and Abbreviations;
    \item chapter titles and chapter content;
    \item figures and photographs;
    \item tables;
    \item equations;
    \item citations and references.
\end{itemize}

Students shall not modify the main document formatting settings in the
preamble, including the document class, page geometry, line spacing,
heading formatting, table-of-contents formatting, font configuration,
page layout, or cover-page macros unless specifically instructed.

\vspace{0.3cm}

% ================================================================
% XXIII. FINAL SUBMISSION REQUIREMENTS
% ================================================================

\noindent
\textbf{XXIII.\quad FINAL SUBMISSION REQUIREMENTS}

\vspace{0.15cm}

Before submission, the final report shall:

\begin{itemize}[leftmargin=0.8cm, itemsep=2pt]
    \item compile without errors that affect the report content or presentation;
    \item contain all required forms and report sections;
    \item contain correct student, class, internship organization, and supervisor information;
    \item contain an updated Table of Contents, List of Figures, and List of Tables;
    \item contain no unused example content from the template;
    \item contain no placeholder text such as \textit{Title of Chapter 1},
    \textit{Section 1.1}, or other sample entries;
    \item contain all required citations and references;
    \item contain no unresolved cross-references.
\end{itemize}

If the Faculty requires submission of the \LaTeX{} source, the complete
project shall include all \verb|.tex| files, the bibliography database,
figures, and any other files required to reproduce the submitted PDF.

\vspace{0.4cm}

\noindent
\textbf{Important:}
Before submitting the final internship report, students shall remove or
comment out the following command in the main document:

\begin{verbatim}
\input{Regulations.tex}
\end{verbatim}

For example:

\begin{verbatim}
% \input{Regulations.tex}
\end{verbatim}

This removes the guidance pages from the final submitted report while
preserving the regulations file in the \LaTeX{} project for reference.

} % End 13pt content

\clearpage

\end{verbatim}

For example:

\begin{verbatim}
% 
% =================================================================
% REGULATIONS FOR PREPARING THE GRADUATION INTERNSHIP REPORT
% USING LATEX
% =================================================================
%
% This file is intended to be included from the main document with:
%
%     \input{Regulations.tex}
%
% IMPORTANT:
% This section is provided for guidance only. Students shall remove
% or comment out the \input{Regulations.tex} command before submitting
% the final internship report.
%
% =================================================================

\clearpage
\thispagestyle{empty}

\begin{center}
    {\fontsize{16pt}{20pt}\selectfont\bfseries
    SPECIFIC REGULATIONS FOR PREPARING\\
    GRADUATION INTERNSHIP REPORTS USING \LaTeX}
\end{center}

\vspace{0.5cm}

{\fontsize{13pt}{17pt}\selectfont
\setstretch{1.3}
\justifying

\noindent
\textbf{Note:}
This section is provided for guidance only and is included in the main
document using the command \verb|\input{Regulations.tex}|.
Students shall remove or comment out this command before submitting the
final internship report.

\vspace{0.4cm}

% ================================================================
% I. GENERAL REQUIREMENTS
% ================================================================

\noindent
\textbf{I.\quad GENERAL REQUIREMENTS}

\vspace{0.15cm}

The Graduation Internship Report shall be prepared using the official
\LaTeX{} template provided by the Faculty. Students shall preserve the
predefined page layout, typography, line spacing, heading styles,
numbering system, figure and table formatting, and bibliography settings
unless otherwise instructed by the Faculty or the supervisor.

The report shall be written in English.

Students are responsible for the academic and technical content of the
report. The official \LaTeX{} template is responsible for standardized
document formatting.

\vspace{0.3cm}

% ================================================================
% II. REPORT STRUCTURE
% ================================================================

\noindent
\textbf{II.\quad REPORT STRUCTURE}

\vspace{0.15cm}

The report shall normally be arranged in the following order:

\begin{enumerate}[label=\arabic*., leftmargin=1.2cm, itemsep=2pt]
    \item Outer Cover;
    \item Comments of the Internship Host Organization;
    \item Supervisor's Evaluation and Graduation Internship Defense Results;
    \item Table of Contents;
    \item List of Symbols and Abbreviations, if applicable;
    \item List of Figures;
    \item List of Tables;
    \item Abstract;
    \item Main Chapters of the Internship Report;
    \item References;
    \item Appendices, if applicable.
\end{enumerate}

The preliminary pages shall not be numbered as chapters. The main report
content shall be organized into consecutively numbered chapters beginning
with Chapter 1.

\vspace{0.3cm}

% ================================================================
% III. PAGE SIZE AND MARGINS
% ================================================================

\noindent
\textbf{III.\quad PAGE SIZE AND MARGINS}

\vspace{0.15cm}

The report shall be formatted on A4 paper with the following margins:

\begin{center}
\begin{tabular}{ll}
    Left margin:   & 3.5 cm \\
    Right margin:  & 2.0 cm \\
    Top margin:    & 2.5 cm \\
    Bottom margin: & 2.5 cm
\end{tabular}
\end{center}

These settings are predefined in the official template and shall not be
modified by students.

\vspace{0.3cm}

% ================================================================
% IV. FONT, FONT SIZE, AND LINE SPACING
% ================================================================

\noindent
\textbf{IV.\quad FONT, FONT SIZE, AND LINE SPACING}

\vspace{0.15cm}

The main text shall use a 13 pt font size with 1.3-line spacing.
A Times-family typeface shall be used in accordance with the official
template.

The text shall be fully justified.

Students shall not use repeated spaces, repeated manual line breaks, or
multiple blank lines to position text. Paragraph and heading spacing shall
be controlled by the template.

\vspace{0.3cm}

% ================================================================
% V. OUTER COVER
% ================================================================

\noindent
\textbf{V.\quad OUTER COVER}

\vspace{0.15cm}

The outer cover shall follow the official format provided in the template
and shall contain the following information:

\begin{itemize}[leftmargin=0.8cm, itemsep=2pt]
    \item Thai Nguyen University of Technology;
    \item Faculty of International Training;
    \item Engineering and Technology Department;
    \item University logo and Faculty logo;
    \item the title \textbf{INTERNSHIP REPORT};
    \item Internship Organization;
    \item Author;
    \item Class;
    \item Supervisor;
    \item place and year of completion.
\end{itemize}

Students shall only replace the required information fields and shall not
modify the cover-page layout unless instructed otherwise.

\vspace{0.3cm}

% ================================================================
% VI. COMMENTS OF THE INTERNSHIP HOST ORGANIZATION
% ================================================================

\noindent
\textbf{VI.\quad COMMENTS OF THE INTERNSHIP HOST ORGANIZATION}

\vspace{0.15cm}

The page entitled \textbf{COMMENTS OF THE INTERNSHIP HOST ORGANIZATION}
is reserved for the internship host organization.

It includes:

\begin{itemize}[leftmargin=0.8cm, itemsep=2pt]
    \item the student's full name;
    \item class;
    \item comments from the host organization;
    \item score;
    \item date;
    \item signature and official seal.
\end{itemize}

Students shall not alter or complete official evaluation fields unless
specifically authorized.

\vspace{0.3cm}

% ================================================================
% VII. SUPERVISOR'S EVALUATION AND DEFENSE RESULTS
% ================================================================

\noindent
\textbf{VII.\quad SUPERVISOR'S EVALUATION AND DEFENSE RESULTS}

\vspace{0.15cm}

The section entitled
\textbf{SUPERVISOR'S EVALUATION AND GRADUATION INTERNSHIP DEFENSE RESULTS}
is reserved for the supervisor and defense committee.

The evaluation may include:

\begin{itemize}[leftmargin=0.8cm, itemsep=2pt]
    \item Internship Progress and Student's Attitude;
    \item Report Content;
    \item Report Presentation;
    \item Other Comments;
    \item Supervisor's Opinion;
    \item Defense Result;
    \item Score.
\end{itemize}

Students shall not enter evaluation scores or defense results before they
are officially determined.

\vspace{0.3cm}

% ================================================================
% VIII. CHAPTERS, SECTIONS, AND SUBSECTIONS
% ================================================================

\noindent
\textbf{VIII.\quad CHAPTERS, SECTIONS, AND SUBSECTIONS}

\vspace{0.15cm}

The report structure shall be generated using the standard \LaTeX{}
commands:

\begin{verbatim}
\chapter{...}
\section{...}
\subsection{...}
\end{verbatim}

Chapter, section, and subsection numbers shall be generated automatically
by \LaTeX{}.

Students shall not manually type chapter or section numbers and shall not
reproduce heading styles using commands such as \verb|\textbf|,
\verb|\fontsize|, or \verb|\centering| instead of the appropriate
structural command.

Excessive subdivision should be avoided. In most cases, the hierarchy
Chapter--Section--Subsection is sufficient.

\vspace{0.3cm}

% ================================================================
% IX. ORGANIZATION OF THE MAIN CONTENT
% ================================================================

\noindent
\textbf{IX.\quad ORGANIZATION OF THE MAIN CONTENT}

\vspace{0.15cm}

The specific content of the report shall depend on the internship
organization and the assigned tasks. A typical structure may include:

\begin{itemize}[leftmargin=0.8cm, itemsep=2pt]
    \item \textbf{Chapter 1. Overview of the Internship Organization};
    \item \textbf{Chapter 2. Internship Activities and Assigned Tasks};
    \item \textbf{Chapter 3. Results, Experience, and Evaluation}.
\end{itemize}

These chapter titles are recommendations only and may be adapted to the
actual internship content.

The report shall reflect work actually performed, observed, analyzed, or
learned during the internship. Unrelated theoretical material should not
be included merely to increase report length.

\vspace{0.3cm}

% ================================================================
% X. TABLE OF CONTENTS
% ================================================================

\noindent
\textbf{X.\quad TABLE OF CONTENTS}

\vspace{0.15cm}

The Table of Contents shall be generated automatically using:

\begin{verbatim}
\tableofcontents
\end{verbatim}

Students shall not create or manually edit page numbers in the Table of
Contents.

The document shall be compiled sufficiently to ensure that all page
numbers and entries are correctly updated.

\vspace{0.3cm}

% ================================================================
% XI. LIST OF SYMBOLS AND ABBREVIATIONS
% ================================================================

\noindent
\textbf{XI.\quad LIST OF SYMBOLS AND ABBREVIATIONS}

\vspace{0.15cm}

Symbols and abbreviations that are used repeatedly in the report shall be
listed in the \textbf{List of Symbols and Abbreviations}.

The list may contain the following columns:

\begin{center}
\begin{tabular}{llll}
    No. & Symbol / Abbreviation & Unit & Definition
\end{tabular}
\end{center}

Only symbols and abbreviations actually used in the final report shall be
retained. Example entries contained in the template shall be removed if
they are not relevant to the student's report.

Each abbreviation should normally be defined when it first appears in
the main text.

\vspace{0.3cm}

% ================================================================
% XII. FIGURES, DIAGRAMS, AND PHOTOGRAPHS
% ================================================================

\noindent
\textbf{XII.\quad FIGURES, DIAGRAMS, AND PHOTOGRAPHS}

\vspace{0.15cm}

Figures shall be clear, legible, and directly relevant to the report
content.

Figures shall be inserted using the standard \verb|figure| environment.

Example:

\begin{verbatim}
\begin{figure}[H]
    \centering
    \includegraphics[width=0.8\textwidth]
                    {figures/company_structure.pdf}
    \caption{Organizational structure of the company}
    \label{fig:company-structure}
\end{figure}
\end{verbatim}

Figure captions shall be placed below the figure.

Each referenced figure shall have a unique \verb|\label| and shall be
referenced in the text using, for example:

\begin{verbatim}
Figure~\ref{fig:company-structure}
\end{verbatim}

Students shall not manually type figure numbers.

Figures reproduced or adapted from external sources shall include an
appropriate citation.

Photographs taken during the internship shall be relevant and sufficiently
clear. Students shall not include confidential, restricted, or proprietary
information without permission from the internship organization.

Vector graphics or high-resolution images are recommended whenever
possible.

\vspace{0.3cm}

% ================================================================
% XIII. TABLES
% ================================================================

\noindent
\textbf{XIII.\quad TABLES}

\vspace{0.15cm}

Tables shall be created using suitable \LaTeX{} environments such as
\verb|table|, \verb|tabular|, \verb|tabularx|, or \verb|longtable|.

Table captions shall be placed above the table.

Example:

\begin{verbatim}
\begin{table}[H]
    \centering
    \caption{Internship schedule}
    \label{tab:internship-schedule}

    \begin{tabular}{lll}
        \toprule
        Week & Task & Result \\
        \midrule
        1 & Orientation & Completed \\
        2 & System study & Completed \\
        \bottomrule
    \end{tabular}
\end{table}
\end{verbatim}

Tables shall be referenced using:

\begin{verbatim}
Table~\ref{tab:internship-schedule}
\end{verbatim}

Students shall not manually type table numbers.

Tables should be created directly in \LaTeX{} rather than inserted as
screenshots whenever practical.

\vspace{0.3cm}

% ================================================================
% XIV. EQUATIONS AND TECHNICAL QUANTITIES
% ================================================================

\noindent
\textbf{XIV.\quad EQUATIONS AND TECHNICAL QUANTITIES}

\vspace{0.15cm}

Mathematical expressions shall be typeset directly using \LaTeX{}.
Images or screenshots of equations shall not be used.

Short expressions may be written in inline mathematical mode, for example:

\begin{verbatim}
$R = 0.7\,\Omega$
\end{verbatim}

Important equations shall use a numbered equation environment.

Example:

\begin{verbatim}
\begin{equation}
    P = UI
    \label{eq:power}
\end{equation}
\end{verbatim}

Equations shall be referenced using:

\begin{verbatim}
Equation~\eqref{eq:power}
\end{verbatim}

Equation numbers shall not be manually entered.

SI units shall be used wherever applicable. The \verb|siunitx| package
should be used where appropriate.

Examples:

\begin{verbatim}
\SI{220}{\volt}
\SI{50}{\hertz}
\SI{2.5}{\ampere}
\end{verbatim}

Notation and units shall be consistent throughout the report.

\vspace{0.3cm}

% ================================================================
% XV. SOURCE CODE
% ================================================================

\noindent
\textbf{XV.\quad SOURCE CODE}

\vspace{0.15cm}

If the internship involves programming, PLCs, microcontrollers,
simulation, or software development, relevant source code may be included
in the report.

Source code shall be presented using the \verb|listings| environment or
another environment defined by the official template.

Example:

\begin{verbatim}
\begin{lstlisting}[language=st]
IF Start = TRUE THEN
    Motor := TRUE;
ELSE
    Motor := FALSE;
END_IF;
\end{lstlisting}
\end{verbatim}

Only code necessary for explaining the work should appear in the main
report. Long source-code listings should be moved to an appendix or
submitted as supplementary files if required.

Source code should not normally be inserted as screenshots.

\vspace{0.3cm}

% ================================================================
% XVI. ABSTRACT
% ================================================================

\noindent
\textbf{XVI.\quad ABSTRACT}

\vspace{0.15cm}

The Abstract shall briefly summarize:

\begin{itemize}[leftmargin=0.8cm, itemsep=2pt]
    \item the internship organization;
    \item the objectives of the internship;
    \item the main tasks performed;
    \item the main results, knowledge, and skills obtained.
\end{itemize}

The Abstract should be concise and should not reproduce the full content
of the report chapter by chapter.

\vspace{0.3cm}

% ================================================================
% XVII. REFERENCES AND CITATIONS
% ================================================================

\noindent
\textbf{XVII.\quad REFERENCES AND CITATIONS}

\vspace{0.15cm}

All books, journal papers, conference papers, standards, websites,
datasheets, manuals, technical documents, figures, tables, or other
external sources used in preparing the report shall be properly cited.

References shall be managed through the bibliography system provided with
the official template.

Bibliographic entries shall normally be stored in:

\begin{verbatim}
References.bib
\end{verbatim}

References shall be cited in the report using:

\begin{verbatim}
\cite{reference-key}
\end{verbatim}

Students shall not manually type reference numbers such as ``[1]'',
``[2]'', etc.

Titles of foreign-language publications shall remain in their original
language in the bibliography.

All borrowed ideas, data, figures, tables, equations, methods, or results
shall be properly acknowledged.

\vspace{0.3cm}

% ================================================================
% XVIII. CROSS-REFERENCING
% ================================================================

\noindent
\textbf{XVIII.\quad CROSS-REFERENCING}

\vspace{0.15cm}

Numbered elements shall be cross-referenced using the standard \LaTeX{}
mechanism.

Use the following commands where appropriate:

\begin{verbatim}
\label{...}
\ref{...}
\eqref{...}
\end{verbatim}

For example:

\begin{verbatim}
As shown in Figure~\ref{fig:system}, ...
\end{verbatim}

Students shall not manually type numbers such as ``Figure 2.3'' when
those numbers can be generated automatically.

\vspace{0.3cm}

% ================================================================
% XIX. PAGE NUMBERING
% ================================================================

\noindent
\textbf{XIX.\quad PAGE NUMBERING}

\vspace{0.15cm}

Cover pages and official evaluation forms at the beginning of the report
shall not display page numbers.

The main body of the report shall use Arabic page numbers beginning from
Chapter 1.

Page numbering shall be generated automatically by the official template.
Students shall not manually enter page numbers.

\vspace{0.3cm}

% ================================================================
% XX. PROJECT FILE ORGANIZATION
% ================================================================

\noindent
\textbf{XX.\quad PROJECT FILE ORGANIZATION}

\vspace{0.15cm}

Students should maintain a clear and organized \LaTeX{} project
structure. A recommended structure is:

\begin{verbatim}
InternshipReport/
|
|-- main.tex
|-- References.bib
|
|-- figures/
|   |-- company.pdf
|   |-- system.png
|   `-- ...
|
`-- chapters/
    |-- chapter1.tex
    |-- chapter2.tex
    `-- chapter3.tex
\end{verbatim}

Figures should normally be stored in the \verb|figures| directory
specified by the official template.

File and directory names should preferably contain Latin letters,
numbers, hyphens, or underscores.

Recommended examples include:

\begin{verbatim}
company_structure.pdf
motor_control.png
chapter_01.tex
\end{verbatim}

Spaces and unnecessary special characters in file names should be
avoided where possible.

\vspace{0.3cm}

% ================================================================
% XXI. DOCUMENT COMPILATION AND FINAL CHECK
% ================================================================

\noindent
\textbf{XXI.\quad DOCUMENT COMPILATION AND FINAL CHECK}

\vspace{0.15cm}

Before submission, the complete \LaTeX{} project shall compile
successfully and produce the final PDF report.

The final document shall not contain unresolved references such as:

\begin{verbatim}
??
[?]
\end{verbatim}

Before generating the final PDF, students shall check:

\begin{itemize}[leftmargin=0.8cm, itemsep=2pt]
    \item the Table of Contents;
    \item page numbering;
    \item chapter, section, and subsection numbering;
    \item figure numbering;
    \item table numbering;
    \item equation numbering;
    \item cross-references;
    \item citations;
    \item bibliography;
    \item font display;
    \item mathematical symbols;
    \item image quality;
    \item completeness of all required forms.
\end{itemize}

The document may be compiled using TeX Live, MiKTeX, or another
compatible \LaTeX{} environment. Regardless of the compilation
environment, the resulting PDF shall conform to the official Faculty
template.

\vspace{0.3cm}

% ================================================================
% XXII. CONTENT THAT STUDENTS MAY EDIT
% ================================================================

\noindent
\textbf{XXII.\quad CONTENT THAT STUDENTS MAY EDIT}

\vspace{0.15cm}

Students may edit or replace:

\begin{itemize}[leftmargin=0.8cm, itemsep=2pt]
    \item student name;
    \item class;
    \item internship organization;
    \item supervisor name;
    \item Abstract;
    \item List of Symbols and Abbreviations;
    \item chapter titles and chapter content;
    \item figures and photographs;
    \item tables;
    \item equations;
    \item citations and references.
\end{itemize}

Students shall not modify the main document formatting settings in the
preamble, including the document class, page geometry, line spacing,
heading formatting, table-of-contents formatting, font configuration,
page layout, or cover-page macros unless specifically instructed.

\vspace{0.3cm}

% ================================================================
% XXIII. FINAL SUBMISSION REQUIREMENTS
% ================================================================

\noindent
\textbf{XXIII.\quad FINAL SUBMISSION REQUIREMENTS}

\vspace{0.15cm}

Before submission, the final report shall:

\begin{itemize}[leftmargin=0.8cm, itemsep=2pt]
    \item compile without errors that affect the report content or presentation;
    \item contain all required forms and report sections;
    \item contain correct student, class, internship organization, and supervisor information;
    \item contain an updated Table of Contents, List of Figures, and List of Tables;
    \item contain no unused example content from the template;
    \item contain no placeholder text such as \textit{Title of Chapter 1},
    \textit{Section 1.1}, or other sample entries;
    \item contain all required citations and references;
    \item contain no unresolved cross-references.
\end{itemize}

If the Faculty requires submission of the \LaTeX{} source, the complete
project shall include all \verb|.tex| files, the bibliography database,
figures, and any other files required to reproduce the submitted PDF.

\vspace{0.4cm}

\noindent
\textbf{Important:}
Before submitting the final internship report, students shall remove or
comment out the following command in the main document:

\begin{verbatim}
\input{Regulations.tex}
\end{verbatim}

For example:

\begin{verbatim}
% \input{Regulations.tex}
\end{verbatim}

This removes the guidance pages from the final submitted report while
preserving the regulations file in the \LaTeX{} project for reference.

} % End 13pt content

\clearpage

\end{verbatim}

This removes the guidance pages from the final submitted report while
preserving the regulations file in the \LaTeX{} project for reference.

} % End 13pt content

\clearpage
 command before submitting
% the final internship report.
%
% =================================================================

\clearpage
\thispagestyle{empty}

\begin{center}
    {\fontsize{16pt}{20pt}\selectfont\bfseries
    SPECIFIC REGULATIONS FOR PREPARING\\
    GRADUATION INTERNSHIP REPORTS USING \LaTeX}
\end{center}

\vspace{0.5cm}

{\fontsize{13pt}{17pt}\selectfont
\setstretch{1.3}
\justifying

\noindent
\textbf{Note:}
This section is provided for guidance only and is included in the main
document using the command \verb|
% =================================================================
% REGULATIONS FOR PREPARING THE GRADUATION INTERNSHIP REPORT
% USING LATEX
% =================================================================
%
% This file is intended to be included from the main document with:
%
%     
% =================================================================
% REGULATIONS FOR PREPARING THE GRADUATION INTERNSHIP REPORT
% USING LATEX
% =================================================================
%
% This file is intended to be included from the main document with:
%
%     \input{Regulations.tex}
%
% IMPORTANT:
% This section is provided for guidance only. Students shall remove
% or comment out the \input{Regulations.tex} command before submitting
% the final internship report.
%
% =================================================================

\clearpage
\thispagestyle{empty}

\begin{center}
    {\fontsize{16pt}{20pt}\selectfont\bfseries
    SPECIFIC REGULATIONS FOR PREPARING\\
    GRADUATION INTERNSHIP REPORTS USING \LaTeX}
\end{center}

\vspace{0.5cm}

{\fontsize{13pt}{17pt}\selectfont
\setstretch{1.3}
\justifying

\noindent
\textbf{Note:}
This section is provided for guidance only and is included in the main
document using the command \verb|\input{Regulations.tex}|.
Students shall remove or comment out this command before submitting the
final internship report.

\vspace{0.4cm}

% ================================================================
% I. GENERAL REQUIREMENTS
% ================================================================

\noindent
\textbf{I.\quad GENERAL REQUIREMENTS}

\vspace{0.15cm}

The Graduation Internship Report shall be prepared using the official
\LaTeX{} template provided by the Faculty. Students shall preserve the
predefined page layout, typography, line spacing, heading styles,
numbering system, figure and table formatting, and bibliography settings
unless otherwise instructed by the Faculty or the supervisor.

The report shall be written in English.

Students are responsible for the academic and technical content of the
report. The official \LaTeX{} template is responsible for standardized
document formatting.

\vspace{0.3cm}

% ================================================================
% II. REPORT STRUCTURE
% ================================================================

\noindent
\textbf{II.\quad REPORT STRUCTURE}

\vspace{0.15cm}

The report shall normally be arranged in the following order:

\begin{enumerate}[label=\arabic*., leftmargin=1.2cm, itemsep=2pt]
    \item Outer Cover;
    \item Comments of the Internship Host Organization;
    \item Supervisor's Evaluation and Graduation Internship Defense Results;
    \item Table of Contents;
    \item List of Symbols and Abbreviations, if applicable;
    \item List of Figures;
    \item List of Tables;
    \item Abstract;
    \item Main Chapters of the Internship Report;
    \item References;
    \item Appendices, if applicable.
\end{enumerate}

The preliminary pages shall not be numbered as chapters. The main report
content shall be organized into consecutively numbered chapters beginning
with Chapter 1.

\vspace{0.3cm}

% ================================================================
% III. PAGE SIZE AND MARGINS
% ================================================================

\noindent
\textbf{III.\quad PAGE SIZE AND MARGINS}

\vspace{0.15cm}

The report shall be formatted on A4 paper with the following margins:

\begin{center}
\begin{tabular}{ll}
    Left margin:   & 3.5 cm \\
    Right margin:  & 2.0 cm \\
    Top margin:    & 2.5 cm \\
    Bottom margin: & 2.5 cm
\end{tabular}
\end{center}

These settings are predefined in the official template and shall not be
modified by students.

\vspace{0.3cm}

% ================================================================
% IV. FONT, FONT SIZE, AND LINE SPACING
% ================================================================

\noindent
\textbf{IV.\quad FONT, FONT SIZE, AND LINE SPACING}

\vspace{0.15cm}

The main text shall use a 13 pt font size with 1.3-line spacing.
A Times-family typeface shall be used in accordance with the official
template.

The text shall be fully justified.

Students shall not use repeated spaces, repeated manual line breaks, or
multiple blank lines to position text. Paragraph and heading spacing shall
be controlled by the template.

\vspace{0.3cm}

% ================================================================
% V. OUTER COVER
% ================================================================

\noindent
\textbf{V.\quad OUTER COVER}

\vspace{0.15cm}

The outer cover shall follow the official format provided in the template
and shall contain the following information:

\begin{itemize}[leftmargin=0.8cm, itemsep=2pt]
    \item Thai Nguyen University of Technology;
    \item Faculty of International Training;
    \item Engineering and Technology Department;
    \item University logo and Faculty logo;
    \item the title \textbf{INTERNSHIP REPORT};
    \item Internship Organization;
    \item Author;
    \item Class;
    \item Supervisor;
    \item place and year of completion.
\end{itemize}

Students shall only replace the required information fields and shall not
modify the cover-page layout unless instructed otherwise.

\vspace{0.3cm}

% ================================================================
% VI. COMMENTS OF THE INTERNSHIP HOST ORGANIZATION
% ================================================================

\noindent
\textbf{VI.\quad COMMENTS OF THE INTERNSHIP HOST ORGANIZATION}

\vspace{0.15cm}

The page entitled \textbf{COMMENTS OF THE INTERNSHIP HOST ORGANIZATION}
is reserved for the internship host organization.

It includes:

\begin{itemize}[leftmargin=0.8cm, itemsep=2pt]
    \item the student's full name;
    \item class;
    \item comments from the host organization;
    \item score;
    \item date;
    \item signature and official seal.
\end{itemize}

Students shall not alter or complete official evaluation fields unless
specifically authorized.

\vspace{0.3cm}

% ================================================================
% VII. SUPERVISOR'S EVALUATION AND DEFENSE RESULTS
% ================================================================

\noindent
\textbf{VII.\quad SUPERVISOR'S EVALUATION AND DEFENSE RESULTS}

\vspace{0.15cm}

The section entitled
\textbf{SUPERVISOR'S EVALUATION AND GRADUATION INTERNSHIP DEFENSE RESULTS}
is reserved for the supervisor and defense committee.

The evaluation may include:

\begin{itemize}[leftmargin=0.8cm, itemsep=2pt]
    \item Internship Progress and Student's Attitude;
    \item Report Content;
    \item Report Presentation;
    \item Other Comments;
    \item Supervisor's Opinion;
    \item Defense Result;
    \item Score.
\end{itemize}

Students shall not enter evaluation scores or defense results before they
are officially determined.

\vspace{0.3cm}

% ================================================================
% VIII. CHAPTERS, SECTIONS, AND SUBSECTIONS
% ================================================================

\noindent
\textbf{VIII.\quad CHAPTERS, SECTIONS, AND SUBSECTIONS}

\vspace{0.15cm}

The report structure shall be generated using the standard \LaTeX{}
commands:

\begin{verbatim}
\chapter{...}
\section{...}
\subsection{...}
\end{verbatim}

Chapter, section, and subsection numbers shall be generated automatically
by \LaTeX{}.

Students shall not manually type chapter or section numbers and shall not
reproduce heading styles using commands such as \verb|\textbf|,
\verb|\fontsize|, or \verb|\centering| instead of the appropriate
structural command.

Excessive subdivision should be avoided. In most cases, the hierarchy
Chapter--Section--Subsection is sufficient.

\vspace{0.3cm}

% ================================================================
% IX. ORGANIZATION OF THE MAIN CONTENT
% ================================================================

\noindent
\textbf{IX.\quad ORGANIZATION OF THE MAIN CONTENT}

\vspace{0.15cm}

The specific content of the report shall depend on the internship
organization and the assigned tasks. A typical structure may include:

\begin{itemize}[leftmargin=0.8cm, itemsep=2pt]
    \item \textbf{Chapter 1. Overview of the Internship Organization};
    \item \textbf{Chapter 2. Internship Activities and Assigned Tasks};
    \item \textbf{Chapter 3. Results, Experience, and Evaluation}.
\end{itemize}

These chapter titles are recommendations only and may be adapted to the
actual internship content.

The report shall reflect work actually performed, observed, analyzed, or
learned during the internship. Unrelated theoretical material should not
be included merely to increase report length.

\vspace{0.3cm}

% ================================================================
% X. TABLE OF CONTENTS
% ================================================================

\noindent
\textbf{X.\quad TABLE OF CONTENTS}

\vspace{0.15cm}

The Table of Contents shall be generated automatically using:

\begin{verbatim}
\tableofcontents
\end{verbatim}

Students shall not create or manually edit page numbers in the Table of
Contents.

The document shall be compiled sufficiently to ensure that all page
numbers and entries are correctly updated.

\vspace{0.3cm}

% ================================================================
% XI. LIST OF SYMBOLS AND ABBREVIATIONS
% ================================================================

\noindent
\textbf{XI.\quad LIST OF SYMBOLS AND ABBREVIATIONS}

\vspace{0.15cm}

Symbols and abbreviations that are used repeatedly in the report shall be
listed in the \textbf{List of Symbols and Abbreviations}.

The list may contain the following columns:

\begin{center}
\begin{tabular}{llll}
    No. & Symbol / Abbreviation & Unit & Definition
\end{tabular}
\end{center}

Only symbols and abbreviations actually used in the final report shall be
retained. Example entries contained in the template shall be removed if
they are not relevant to the student's report.

Each abbreviation should normally be defined when it first appears in
the main text.

\vspace{0.3cm}

% ================================================================
% XII. FIGURES, DIAGRAMS, AND PHOTOGRAPHS
% ================================================================

\noindent
\textbf{XII.\quad FIGURES, DIAGRAMS, AND PHOTOGRAPHS}

\vspace{0.15cm}

Figures shall be clear, legible, and directly relevant to the report
content.

Figures shall be inserted using the standard \verb|figure| environment.

Example:

\begin{verbatim}
\begin{figure}[H]
    \centering
    \includegraphics[width=0.8\textwidth]
                    {figures/company_structure.pdf}
    \caption{Organizational structure of the company}
    \label{fig:company-structure}
\end{figure}
\end{verbatim}

Figure captions shall be placed below the figure.

Each referenced figure shall have a unique \verb|\label| and shall be
referenced in the text using, for example:

\begin{verbatim}
Figure~\ref{fig:company-structure}
\end{verbatim}

Students shall not manually type figure numbers.

Figures reproduced or adapted from external sources shall include an
appropriate citation.

Photographs taken during the internship shall be relevant and sufficiently
clear. Students shall not include confidential, restricted, or proprietary
information without permission from the internship organization.

Vector graphics or high-resolution images are recommended whenever
possible.

\vspace{0.3cm}

% ================================================================
% XIII. TABLES
% ================================================================

\noindent
\textbf{XIII.\quad TABLES}

\vspace{0.15cm}

Tables shall be created using suitable \LaTeX{} environments such as
\verb|table|, \verb|tabular|, \verb|tabularx|, or \verb|longtable|.

Table captions shall be placed above the table.

Example:

\begin{verbatim}
\begin{table}[H]
    \centering
    \caption{Internship schedule}
    \label{tab:internship-schedule}

    \begin{tabular}{lll}
        \toprule
        Week & Task & Result \\
        \midrule
        1 & Orientation & Completed \\
        2 & System study & Completed \\
        \bottomrule
    \end{tabular}
\end{table}
\end{verbatim}

Tables shall be referenced using:

\begin{verbatim}
Table~\ref{tab:internship-schedule}
\end{verbatim}

Students shall not manually type table numbers.

Tables should be created directly in \LaTeX{} rather than inserted as
screenshots whenever practical.

\vspace{0.3cm}

% ================================================================
% XIV. EQUATIONS AND TECHNICAL QUANTITIES
% ================================================================

\noindent
\textbf{XIV.\quad EQUATIONS AND TECHNICAL QUANTITIES}

\vspace{0.15cm}

Mathematical expressions shall be typeset directly using \LaTeX{}.
Images or screenshots of equations shall not be used.

Short expressions may be written in inline mathematical mode, for example:

\begin{verbatim}
$R = 0.7\,\Omega$
\end{verbatim}

Important equations shall use a numbered equation environment.

Example:

\begin{verbatim}
\begin{equation}
    P = UI
    \label{eq:power}
\end{equation}
\end{verbatim}

Equations shall be referenced using:

\begin{verbatim}
Equation~\eqref{eq:power}
\end{verbatim}

Equation numbers shall not be manually entered.

SI units shall be used wherever applicable. The \verb|siunitx| package
should be used where appropriate.

Examples:

\begin{verbatim}
\SI{220}{\volt}
\SI{50}{\hertz}
\SI{2.5}{\ampere}
\end{verbatim}

Notation and units shall be consistent throughout the report.

\vspace{0.3cm}

% ================================================================
% XV. SOURCE CODE
% ================================================================

\noindent
\textbf{XV.\quad SOURCE CODE}

\vspace{0.15cm}

If the internship involves programming, PLCs, microcontrollers,
simulation, or software development, relevant source code may be included
in the report.

Source code shall be presented using the \verb|listings| environment or
another environment defined by the official template.

Example:

\begin{verbatim}
\begin{lstlisting}[language=st]
IF Start = TRUE THEN
    Motor := TRUE;
ELSE
    Motor := FALSE;
END_IF;
\end{lstlisting}
\end{verbatim}

Only code necessary for explaining the work should appear in the main
report. Long source-code listings should be moved to an appendix or
submitted as supplementary files if required.

Source code should not normally be inserted as screenshots.

\vspace{0.3cm}

% ================================================================
% XVI. ABSTRACT
% ================================================================

\noindent
\textbf{XVI.\quad ABSTRACT}

\vspace{0.15cm}

The Abstract shall briefly summarize:

\begin{itemize}[leftmargin=0.8cm, itemsep=2pt]
    \item the internship organization;
    \item the objectives of the internship;
    \item the main tasks performed;
    \item the main results, knowledge, and skills obtained.
\end{itemize}

The Abstract should be concise and should not reproduce the full content
of the report chapter by chapter.

\vspace{0.3cm}

% ================================================================
% XVII. REFERENCES AND CITATIONS
% ================================================================

\noindent
\textbf{XVII.\quad REFERENCES AND CITATIONS}

\vspace{0.15cm}

All books, journal papers, conference papers, standards, websites,
datasheets, manuals, technical documents, figures, tables, or other
external sources used in preparing the report shall be properly cited.

References shall be managed through the bibliography system provided with
the official template.

Bibliographic entries shall normally be stored in:

\begin{verbatim}
References.bib
\end{verbatim}

References shall be cited in the report using:

\begin{verbatim}
\cite{reference-key}
\end{verbatim}

Students shall not manually type reference numbers such as ``[1]'',
``[2]'', etc.

Titles of foreign-language publications shall remain in their original
language in the bibliography.

All borrowed ideas, data, figures, tables, equations, methods, or results
shall be properly acknowledged.

\vspace{0.3cm}

% ================================================================
% XVIII. CROSS-REFERENCING
% ================================================================

\noindent
\textbf{XVIII.\quad CROSS-REFERENCING}

\vspace{0.15cm}

Numbered elements shall be cross-referenced using the standard \LaTeX{}
mechanism.

Use the following commands where appropriate:

\begin{verbatim}
\label{...}
\ref{...}
\eqref{...}
\end{verbatim}

For example:

\begin{verbatim}
As shown in Figure~\ref{fig:system}, ...
\end{verbatim}

Students shall not manually type numbers such as ``Figure 2.3'' when
those numbers can be generated automatically.

\vspace{0.3cm}

% ================================================================
% XIX. PAGE NUMBERING
% ================================================================

\noindent
\textbf{XIX.\quad PAGE NUMBERING}

\vspace{0.15cm}

Cover pages and official evaluation forms at the beginning of the report
shall not display page numbers.

The main body of the report shall use Arabic page numbers beginning from
Chapter 1.

Page numbering shall be generated automatically by the official template.
Students shall not manually enter page numbers.

\vspace{0.3cm}

% ================================================================
% XX. PROJECT FILE ORGANIZATION
% ================================================================

\noindent
\textbf{XX.\quad PROJECT FILE ORGANIZATION}

\vspace{0.15cm}

Students should maintain a clear and organized \LaTeX{} project
structure. A recommended structure is:

\begin{verbatim}
InternshipReport/
|
|-- main.tex
|-- References.bib
|
|-- figures/
|   |-- company.pdf
|   |-- system.png
|   `-- ...
|
`-- chapters/
    |-- chapter1.tex
    |-- chapter2.tex
    `-- chapter3.tex
\end{verbatim}

Figures should normally be stored in the \verb|figures| directory
specified by the official template.

File and directory names should preferably contain Latin letters,
numbers, hyphens, or underscores.

Recommended examples include:

\begin{verbatim}
company_structure.pdf
motor_control.png
chapter_01.tex
\end{verbatim}

Spaces and unnecessary special characters in file names should be
avoided where possible.

\vspace{0.3cm}

% ================================================================
% XXI. DOCUMENT COMPILATION AND FINAL CHECK
% ================================================================

\noindent
\textbf{XXI.\quad DOCUMENT COMPILATION AND FINAL CHECK}

\vspace{0.15cm}

Before submission, the complete \LaTeX{} project shall compile
successfully and produce the final PDF report.

The final document shall not contain unresolved references such as:

\begin{verbatim}
??
[?]
\end{verbatim}

Before generating the final PDF, students shall check:

\begin{itemize}[leftmargin=0.8cm, itemsep=2pt]
    \item the Table of Contents;
    \item page numbering;
    \item chapter, section, and subsection numbering;
    \item figure numbering;
    \item table numbering;
    \item equation numbering;
    \item cross-references;
    \item citations;
    \item bibliography;
    \item font display;
    \item mathematical symbols;
    \item image quality;
    \item completeness of all required forms.
\end{itemize}

The document may be compiled using TeX Live, MiKTeX, or another
compatible \LaTeX{} environment. Regardless of the compilation
environment, the resulting PDF shall conform to the official Faculty
template.

\vspace{0.3cm}

% ================================================================
% XXII. CONTENT THAT STUDENTS MAY EDIT
% ================================================================

\noindent
\textbf{XXII.\quad CONTENT THAT STUDENTS MAY EDIT}

\vspace{0.15cm}

Students may edit or replace:

\begin{itemize}[leftmargin=0.8cm, itemsep=2pt]
    \item student name;
    \item class;
    \item internship organization;
    \item supervisor name;
    \item Abstract;
    \item List of Symbols and Abbreviations;
    \item chapter titles and chapter content;
    \item figures and photographs;
    \item tables;
    \item equations;
    \item citations and references.
\end{itemize}

Students shall not modify the main document formatting settings in the
preamble, including the document class, page geometry, line spacing,
heading formatting, table-of-contents formatting, font configuration,
page layout, or cover-page macros unless specifically instructed.

\vspace{0.3cm}

% ================================================================
% XXIII. FINAL SUBMISSION REQUIREMENTS
% ================================================================

\noindent
\textbf{XXIII.\quad FINAL SUBMISSION REQUIREMENTS}

\vspace{0.15cm}

Before submission, the final report shall:

\begin{itemize}[leftmargin=0.8cm, itemsep=2pt]
    \item compile without errors that affect the report content or presentation;
    \item contain all required forms and report sections;
    \item contain correct student, class, internship organization, and supervisor information;
    \item contain an updated Table of Contents, List of Figures, and List of Tables;
    \item contain no unused example content from the template;
    \item contain no placeholder text such as \textit{Title of Chapter 1},
    \textit{Section 1.1}, or other sample entries;
    \item contain all required citations and references;
    \item contain no unresolved cross-references.
\end{itemize}

If the Faculty requires submission of the \LaTeX{} source, the complete
project shall include all \verb|.tex| files, the bibliography database,
figures, and any other files required to reproduce the submitted PDF.

\vspace{0.4cm}

\noindent
\textbf{Important:}
Before submitting the final internship report, students shall remove or
comment out the following command in the main document:

\begin{verbatim}
\input{Regulations.tex}
\end{verbatim}

For example:

\begin{verbatim}
% \input{Regulations.tex}
\end{verbatim}

This removes the guidance pages from the final submitted report while
preserving the regulations file in the \LaTeX{} project for reference.

} % End 13pt content

\clearpage

%
% IMPORTANT:
% This section is provided for guidance only. Students shall remove
% or comment out the 
% =================================================================
% REGULATIONS FOR PREPARING THE GRADUATION INTERNSHIP REPORT
% USING LATEX
% =================================================================
%
% This file is intended to be included from the main document with:
%
%     \input{Regulations.tex}
%
% IMPORTANT:
% This section is provided for guidance only. Students shall remove
% or comment out the \input{Regulations.tex} command before submitting
% the final internship report.
%
% =================================================================

\clearpage
\thispagestyle{empty}

\begin{center}
    {\fontsize{16pt}{20pt}\selectfont\bfseries
    SPECIFIC REGULATIONS FOR PREPARING\\
    GRADUATION INTERNSHIP REPORTS USING \LaTeX}
\end{center}

\vspace{0.5cm}

{\fontsize{13pt}{17pt}\selectfont
\setstretch{1.3}
\justifying

\noindent
\textbf{Note:}
This section is provided for guidance only and is included in the main
document using the command \verb|\input{Regulations.tex}|.
Students shall remove or comment out this command before submitting the
final internship report.

\vspace{0.4cm}

% ================================================================
% I. GENERAL REQUIREMENTS
% ================================================================

\noindent
\textbf{I.\quad GENERAL REQUIREMENTS}

\vspace{0.15cm}

The Graduation Internship Report shall be prepared using the official
\LaTeX{} template provided by the Faculty. Students shall preserve the
predefined page layout, typography, line spacing, heading styles,
numbering system, figure and table formatting, and bibliography settings
unless otherwise instructed by the Faculty or the supervisor.

The report shall be written in English.

Students are responsible for the academic and technical content of the
report. The official \LaTeX{} template is responsible for standardized
document formatting.

\vspace{0.3cm}

% ================================================================
% II. REPORT STRUCTURE
% ================================================================

\noindent
\textbf{II.\quad REPORT STRUCTURE}

\vspace{0.15cm}

The report shall normally be arranged in the following order:

\begin{enumerate}[label=\arabic*., leftmargin=1.2cm, itemsep=2pt]
    \item Outer Cover;
    \item Comments of the Internship Host Organization;
    \item Supervisor's Evaluation and Graduation Internship Defense Results;
    \item Table of Contents;
    \item List of Symbols and Abbreviations, if applicable;
    \item List of Figures;
    \item List of Tables;
    \item Abstract;
    \item Main Chapters of the Internship Report;
    \item References;
    \item Appendices, if applicable.
\end{enumerate}

The preliminary pages shall not be numbered as chapters. The main report
content shall be organized into consecutively numbered chapters beginning
with Chapter 1.

\vspace{0.3cm}

% ================================================================
% III. PAGE SIZE AND MARGINS
% ================================================================

\noindent
\textbf{III.\quad PAGE SIZE AND MARGINS}

\vspace{0.15cm}

The report shall be formatted on A4 paper with the following margins:

\begin{center}
\begin{tabular}{ll}
    Left margin:   & 3.5 cm \\
    Right margin:  & 2.0 cm \\
    Top margin:    & 2.5 cm \\
    Bottom margin: & 2.5 cm
\end{tabular}
\end{center}

These settings are predefined in the official template and shall not be
modified by students.

\vspace{0.3cm}

% ================================================================
% IV. FONT, FONT SIZE, AND LINE SPACING
% ================================================================

\noindent
\textbf{IV.\quad FONT, FONT SIZE, AND LINE SPACING}

\vspace{0.15cm}

The main text shall use a 13 pt font size with 1.3-line spacing.
A Times-family typeface shall be used in accordance with the official
template.

The text shall be fully justified.

Students shall not use repeated spaces, repeated manual line breaks, or
multiple blank lines to position text. Paragraph and heading spacing shall
be controlled by the template.

\vspace{0.3cm}

% ================================================================
% V. OUTER COVER
% ================================================================

\noindent
\textbf{V.\quad OUTER COVER}

\vspace{0.15cm}

The outer cover shall follow the official format provided in the template
and shall contain the following information:

\begin{itemize}[leftmargin=0.8cm, itemsep=2pt]
    \item Thai Nguyen University of Technology;
    \item Faculty of International Training;
    \item Engineering and Technology Department;
    \item University logo and Faculty logo;
    \item the title \textbf{INTERNSHIP REPORT};
    \item Internship Organization;
    \item Author;
    \item Class;
    \item Supervisor;
    \item place and year of completion.
\end{itemize}

Students shall only replace the required information fields and shall not
modify the cover-page layout unless instructed otherwise.

\vspace{0.3cm}

% ================================================================
% VI. COMMENTS OF THE INTERNSHIP HOST ORGANIZATION
% ================================================================

\noindent
\textbf{VI.\quad COMMENTS OF THE INTERNSHIP HOST ORGANIZATION}

\vspace{0.15cm}

The page entitled \textbf{COMMENTS OF THE INTERNSHIP HOST ORGANIZATION}
is reserved for the internship host organization.

It includes:

\begin{itemize}[leftmargin=0.8cm, itemsep=2pt]
    \item the student's full name;
    \item class;
    \item comments from the host organization;
    \item score;
    \item date;
    \item signature and official seal.
\end{itemize}

Students shall not alter or complete official evaluation fields unless
specifically authorized.

\vspace{0.3cm}

% ================================================================
% VII. SUPERVISOR'S EVALUATION AND DEFENSE RESULTS
% ================================================================

\noindent
\textbf{VII.\quad SUPERVISOR'S EVALUATION AND DEFENSE RESULTS}

\vspace{0.15cm}

The section entitled
\textbf{SUPERVISOR'S EVALUATION AND GRADUATION INTERNSHIP DEFENSE RESULTS}
is reserved for the supervisor and defense committee.

The evaluation may include:

\begin{itemize}[leftmargin=0.8cm, itemsep=2pt]
    \item Internship Progress and Student's Attitude;
    \item Report Content;
    \item Report Presentation;
    \item Other Comments;
    \item Supervisor's Opinion;
    \item Defense Result;
    \item Score.
\end{itemize}

Students shall not enter evaluation scores or defense results before they
are officially determined.

\vspace{0.3cm}

% ================================================================
% VIII. CHAPTERS, SECTIONS, AND SUBSECTIONS
% ================================================================

\noindent
\textbf{VIII.\quad CHAPTERS, SECTIONS, AND SUBSECTIONS}

\vspace{0.15cm}

The report structure shall be generated using the standard \LaTeX{}
commands:

\begin{verbatim}
\chapter{...}
\section{...}
\subsection{...}
\end{verbatim}

Chapter, section, and subsection numbers shall be generated automatically
by \LaTeX{}.

Students shall not manually type chapter or section numbers and shall not
reproduce heading styles using commands such as \verb|\textbf|,
\verb|\fontsize|, or \verb|\centering| instead of the appropriate
structural command.

Excessive subdivision should be avoided. In most cases, the hierarchy
Chapter--Section--Subsection is sufficient.

\vspace{0.3cm}

% ================================================================
% IX. ORGANIZATION OF THE MAIN CONTENT
% ================================================================

\noindent
\textbf{IX.\quad ORGANIZATION OF THE MAIN CONTENT}

\vspace{0.15cm}

The specific content of the report shall depend on the internship
organization and the assigned tasks. A typical structure may include:

\begin{itemize}[leftmargin=0.8cm, itemsep=2pt]
    \item \textbf{Chapter 1. Overview of the Internship Organization};
    \item \textbf{Chapter 2. Internship Activities and Assigned Tasks};
    \item \textbf{Chapter 3. Results, Experience, and Evaluation}.
\end{itemize}

These chapter titles are recommendations only and may be adapted to the
actual internship content.

The report shall reflect work actually performed, observed, analyzed, or
learned during the internship. Unrelated theoretical material should not
be included merely to increase report length.

\vspace{0.3cm}

% ================================================================
% X. TABLE OF CONTENTS
% ================================================================

\noindent
\textbf{X.\quad TABLE OF CONTENTS}

\vspace{0.15cm}

The Table of Contents shall be generated automatically using:

\begin{verbatim}
\tableofcontents
\end{verbatim}

Students shall not create or manually edit page numbers in the Table of
Contents.

The document shall be compiled sufficiently to ensure that all page
numbers and entries are correctly updated.

\vspace{0.3cm}

% ================================================================
% XI. LIST OF SYMBOLS AND ABBREVIATIONS
% ================================================================

\noindent
\textbf{XI.\quad LIST OF SYMBOLS AND ABBREVIATIONS}

\vspace{0.15cm}

Symbols and abbreviations that are used repeatedly in the report shall be
listed in the \textbf{List of Symbols and Abbreviations}.

The list may contain the following columns:

\begin{center}
\begin{tabular}{llll}
    No. & Symbol / Abbreviation & Unit & Definition
\end{tabular}
\end{center}

Only symbols and abbreviations actually used in the final report shall be
retained. Example entries contained in the template shall be removed if
they are not relevant to the student's report.

Each abbreviation should normally be defined when it first appears in
the main text.

\vspace{0.3cm}

% ================================================================
% XII. FIGURES, DIAGRAMS, AND PHOTOGRAPHS
% ================================================================

\noindent
\textbf{XII.\quad FIGURES, DIAGRAMS, AND PHOTOGRAPHS}

\vspace{0.15cm}

Figures shall be clear, legible, and directly relevant to the report
content.

Figures shall be inserted using the standard \verb|figure| environment.

Example:

\begin{verbatim}
\begin{figure}[H]
    \centering
    \includegraphics[width=0.8\textwidth]
                    {figures/company_structure.pdf}
    \caption{Organizational structure of the company}
    \label{fig:company-structure}
\end{figure}
\end{verbatim}

Figure captions shall be placed below the figure.

Each referenced figure shall have a unique \verb|\label| and shall be
referenced in the text using, for example:

\begin{verbatim}
Figure~\ref{fig:company-structure}
\end{verbatim}

Students shall not manually type figure numbers.

Figures reproduced or adapted from external sources shall include an
appropriate citation.

Photographs taken during the internship shall be relevant and sufficiently
clear. Students shall not include confidential, restricted, or proprietary
information without permission from the internship organization.

Vector graphics or high-resolution images are recommended whenever
possible.

\vspace{0.3cm}

% ================================================================
% XIII. TABLES
% ================================================================

\noindent
\textbf{XIII.\quad TABLES}

\vspace{0.15cm}

Tables shall be created using suitable \LaTeX{} environments such as
\verb|table|, \verb|tabular|, \verb|tabularx|, or \verb|longtable|.

Table captions shall be placed above the table.

Example:

\begin{verbatim}
\begin{table}[H]
    \centering
    \caption{Internship schedule}
    \label{tab:internship-schedule}

    \begin{tabular}{lll}
        \toprule
        Week & Task & Result \\
        \midrule
        1 & Orientation & Completed \\
        2 & System study & Completed \\
        \bottomrule
    \end{tabular}
\end{table}
\end{verbatim}

Tables shall be referenced using:

\begin{verbatim}
Table~\ref{tab:internship-schedule}
\end{verbatim}

Students shall not manually type table numbers.

Tables should be created directly in \LaTeX{} rather than inserted as
screenshots whenever practical.

\vspace{0.3cm}

% ================================================================
% XIV. EQUATIONS AND TECHNICAL QUANTITIES
% ================================================================

\noindent
\textbf{XIV.\quad EQUATIONS AND TECHNICAL QUANTITIES}

\vspace{0.15cm}

Mathematical expressions shall be typeset directly using \LaTeX{}.
Images or screenshots of equations shall not be used.

Short expressions may be written in inline mathematical mode, for example:

\begin{verbatim}
$R = 0.7\,\Omega$
\end{verbatim}

Important equations shall use a numbered equation environment.

Example:

\begin{verbatim}
\begin{equation}
    P = UI
    \label{eq:power}
\end{equation}
\end{verbatim}

Equations shall be referenced using:

\begin{verbatim}
Equation~\eqref{eq:power}
\end{verbatim}

Equation numbers shall not be manually entered.

SI units shall be used wherever applicable. The \verb|siunitx| package
should be used where appropriate.

Examples:

\begin{verbatim}
\SI{220}{\volt}
\SI{50}{\hertz}
\SI{2.5}{\ampere}
\end{verbatim}

Notation and units shall be consistent throughout the report.

\vspace{0.3cm}

% ================================================================
% XV. SOURCE CODE
% ================================================================

\noindent
\textbf{XV.\quad SOURCE CODE}

\vspace{0.15cm}

If the internship involves programming, PLCs, microcontrollers,
simulation, or software development, relevant source code may be included
in the report.

Source code shall be presented using the \verb|listings| environment or
another environment defined by the official template.

Example:

\begin{verbatim}
\begin{lstlisting}[language=st]
IF Start = TRUE THEN
    Motor := TRUE;
ELSE
    Motor := FALSE;
END_IF;
\end{lstlisting}
\end{verbatim}

Only code necessary for explaining the work should appear in the main
report. Long source-code listings should be moved to an appendix or
submitted as supplementary files if required.

Source code should not normally be inserted as screenshots.

\vspace{0.3cm}

% ================================================================
% XVI. ABSTRACT
% ================================================================

\noindent
\textbf{XVI.\quad ABSTRACT}

\vspace{0.15cm}

The Abstract shall briefly summarize:

\begin{itemize}[leftmargin=0.8cm, itemsep=2pt]
    \item the internship organization;
    \item the objectives of the internship;
    \item the main tasks performed;
    \item the main results, knowledge, and skills obtained.
\end{itemize}

The Abstract should be concise and should not reproduce the full content
of the report chapter by chapter.

\vspace{0.3cm}

% ================================================================
% XVII. REFERENCES AND CITATIONS
% ================================================================

\noindent
\textbf{XVII.\quad REFERENCES AND CITATIONS}

\vspace{0.15cm}

All books, journal papers, conference papers, standards, websites,
datasheets, manuals, technical documents, figures, tables, or other
external sources used in preparing the report shall be properly cited.

References shall be managed through the bibliography system provided with
the official template.

Bibliographic entries shall normally be stored in:

\begin{verbatim}
References.bib
\end{verbatim}

References shall be cited in the report using:

\begin{verbatim}
\cite{reference-key}
\end{verbatim}

Students shall not manually type reference numbers such as ``[1]'',
``[2]'', etc.

Titles of foreign-language publications shall remain in their original
language in the bibliography.

All borrowed ideas, data, figures, tables, equations, methods, or results
shall be properly acknowledged.

\vspace{0.3cm}

% ================================================================
% XVIII. CROSS-REFERENCING
% ================================================================

\noindent
\textbf{XVIII.\quad CROSS-REFERENCING}

\vspace{0.15cm}

Numbered elements shall be cross-referenced using the standard \LaTeX{}
mechanism.

Use the following commands where appropriate:

\begin{verbatim}
\label{...}
\ref{...}
\eqref{...}
\end{verbatim}

For example:

\begin{verbatim}
As shown in Figure~\ref{fig:system}, ...
\end{verbatim}

Students shall not manually type numbers such as ``Figure 2.3'' when
those numbers can be generated automatically.

\vspace{0.3cm}

% ================================================================
% XIX. PAGE NUMBERING
% ================================================================

\noindent
\textbf{XIX.\quad PAGE NUMBERING}

\vspace{0.15cm}

Cover pages and official evaluation forms at the beginning of the report
shall not display page numbers.

The main body of the report shall use Arabic page numbers beginning from
Chapter 1.

Page numbering shall be generated automatically by the official template.
Students shall not manually enter page numbers.

\vspace{0.3cm}

% ================================================================
% XX. PROJECT FILE ORGANIZATION
% ================================================================

\noindent
\textbf{XX.\quad PROJECT FILE ORGANIZATION}

\vspace{0.15cm}

Students should maintain a clear and organized \LaTeX{} project
structure. A recommended structure is:

\begin{verbatim}
InternshipReport/
|
|-- main.tex
|-- References.bib
|
|-- figures/
|   |-- company.pdf
|   |-- system.png
|   `-- ...
|
`-- chapters/
    |-- chapter1.tex
    |-- chapter2.tex
    `-- chapter3.tex
\end{verbatim}

Figures should normally be stored in the \verb|figures| directory
specified by the official template.

File and directory names should preferably contain Latin letters,
numbers, hyphens, or underscores.

Recommended examples include:

\begin{verbatim}
company_structure.pdf
motor_control.png
chapter_01.tex
\end{verbatim}

Spaces and unnecessary special characters in file names should be
avoided where possible.

\vspace{0.3cm}

% ================================================================
% XXI. DOCUMENT COMPILATION AND FINAL CHECK
% ================================================================

\noindent
\textbf{XXI.\quad DOCUMENT COMPILATION AND FINAL CHECK}

\vspace{0.15cm}

Before submission, the complete \LaTeX{} project shall compile
successfully and produce the final PDF report.

The final document shall not contain unresolved references such as:

\begin{verbatim}
??
[?]
\end{verbatim}

Before generating the final PDF, students shall check:

\begin{itemize}[leftmargin=0.8cm, itemsep=2pt]
    \item the Table of Contents;
    \item page numbering;
    \item chapter, section, and subsection numbering;
    \item figure numbering;
    \item table numbering;
    \item equation numbering;
    \item cross-references;
    \item citations;
    \item bibliography;
    \item font display;
    \item mathematical symbols;
    \item image quality;
    \item completeness of all required forms.
\end{itemize}

The document may be compiled using TeX Live, MiKTeX, or another
compatible \LaTeX{} environment. Regardless of the compilation
environment, the resulting PDF shall conform to the official Faculty
template.

\vspace{0.3cm}

% ================================================================
% XXII. CONTENT THAT STUDENTS MAY EDIT
% ================================================================

\noindent
\textbf{XXII.\quad CONTENT THAT STUDENTS MAY EDIT}

\vspace{0.15cm}

Students may edit or replace:

\begin{itemize}[leftmargin=0.8cm, itemsep=2pt]
    \item student name;
    \item class;
    \item internship organization;
    \item supervisor name;
    \item Abstract;
    \item List of Symbols and Abbreviations;
    \item chapter titles and chapter content;
    \item figures and photographs;
    \item tables;
    \item equations;
    \item citations and references.
\end{itemize}

Students shall not modify the main document formatting settings in the
preamble, including the document class, page geometry, line spacing,
heading formatting, table-of-contents formatting, font configuration,
page layout, or cover-page macros unless specifically instructed.

\vspace{0.3cm}

% ================================================================
% XXIII. FINAL SUBMISSION REQUIREMENTS
% ================================================================

\noindent
\textbf{XXIII.\quad FINAL SUBMISSION REQUIREMENTS}

\vspace{0.15cm}

Before submission, the final report shall:

\begin{itemize}[leftmargin=0.8cm, itemsep=2pt]
    \item compile without errors that affect the report content or presentation;
    \item contain all required forms and report sections;
    \item contain correct student, class, internship organization, and supervisor information;
    \item contain an updated Table of Contents, List of Figures, and List of Tables;
    \item contain no unused example content from the template;
    \item contain no placeholder text such as \textit{Title of Chapter 1},
    \textit{Section 1.1}, or other sample entries;
    \item contain all required citations and references;
    \item contain no unresolved cross-references.
\end{itemize}

If the Faculty requires submission of the \LaTeX{} source, the complete
project shall include all \verb|.tex| files, the bibliography database,
figures, and any other files required to reproduce the submitted PDF.

\vspace{0.4cm}

\noindent
\textbf{Important:}
Before submitting the final internship report, students shall remove or
comment out the following command in the main document:

\begin{verbatim}
\input{Regulations.tex}
\end{verbatim}

For example:

\begin{verbatim}
% \input{Regulations.tex}
\end{verbatim}

This removes the guidance pages from the final submitted report while
preserving the regulations file in the \LaTeX{} project for reference.

} % End 13pt content

\clearpage
 command before submitting
% the final internship report.
%
% =================================================================

\clearpage
\thispagestyle{empty}

\begin{center}
    {\fontsize{16pt}{20pt}\selectfont\bfseries
    SPECIFIC REGULATIONS FOR PREPARING\\
    GRADUATION INTERNSHIP REPORTS USING \LaTeX}
\end{center}

\vspace{0.5cm}

{\fontsize{13pt}{17pt}\selectfont
\setstretch{1.3}
\justifying

\noindent
\textbf{Note:}
This section is provided for guidance only and is included in the main
document using the command \verb|
% =================================================================
% REGULATIONS FOR PREPARING THE GRADUATION INTERNSHIP REPORT
% USING LATEX
% =================================================================
%
% This file is intended to be included from the main document with:
%
%     \input{Regulations.tex}
%
% IMPORTANT:
% This section is provided for guidance only. Students shall remove
% or comment out the \input{Regulations.tex} command before submitting
% the final internship report.
%
% =================================================================

\clearpage
\thispagestyle{empty}

\begin{center}
    {\fontsize{16pt}{20pt}\selectfont\bfseries
    SPECIFIC REGULATIONS FOR PREPARING\\
    GRADUATION INTERNSHIP REPORTS USING \LaTeX}
\end{center}

\vspace{0.5cm}

{\fontsize{13pt}{17pt}\selectfont
\setstretch{1.3}
\justifying

\noindent
\textbf{Note:}
This section is provided for guidance only and is included in the main
document using the command \verb|\input{Regulations.tex}|.
Students shall remove or comment out this command before submitting the
final internship report.

\vspace{0.4cm}

% ================================================================
% I. GENERAL REQUIREMENTS
% ================================================================

\noindent
\textbf{I.\quad GENERAL REQUIREMENTS}

\vspace{0.15cm}

The Graduation Internship Report shall be prepared using the official
\LaTeX{} template provided by the Faculty. Students shall preserve the
predefined page layout, typography, line spacing, heading styles,
numbering system, figure and table formatting, and bibliography settings
unless otherwise instructed by the Faculty or the supervisor.

The report shall be written in English.

Students are responsible for the academic and technical content of the
report. The official \LaTeX{} template is responsible for standardized
document formatting.

\vspace{0.3cm}

% ================================================================
% II. REPORT STRUCTURE
% ================================================================

\noindent
\textbf{II.\quad REPORT STRUCTURE}

\vspace{0.15cm}

The report shall normally be arranged in the following order:

\begin{enumerate}[label=\arabic*., leftmargin=1.2cm, itemsep=2pt]
    \item Outer Cover;
    \item Comments of the Internship Host Organization;
    \item Supervisor's Evaluation and Graduation Internship Defense Results;
    \item Table of Contents;
    \item List of Symbols and Abbreviations, if applicable;
    \item List of Figures;
    \item List of Tables;
    \item Abstract;
    \item Main Chapters of the Internship Report;
    \item References;
    \item Appendices, if applicable.
\end{enumerate}

The preliminary pages shall not be numbered as chapters. The main report
content shall be organized into consecutively numbered chapters beginning
with Chapter 1.

\vspace{0.3cm}

% ================================================================
% III. PAGE SIZE AND MARGINS
% ================================================================

\noindent
\textbf{III.\quad PAGE SIZE AND MARGINS}

\vspace{0.15cm}

The report shall be formatted on A4 paper with the following margins:

\begin{center}
\begin{tabular}{ll}
    Left margin:   & 3.5 cm \\
    Right margin:  & 2.0 cm \\
    Top margin:    & 2.5 cm \\
    Bottom margin: & 2.5 cm
\end{tabular}
\end{center}

These settings are predefined in the official template and shall not be
modified by students.

\vspace{0.3cm}

% ================================================================
% IV. FONT, FONT SIZE, AND LINE SPACING
% ================================================================

\noindent
\textbf{IV.\quad FONT, FONT SIZE, AND LINE SPACING}

\vspace{0.15cm}

The main text shall use a 13 pt font size with 1.3-line spacing.
A Times-family typeface shall be used in accordance with the official
template.

The text shall be fully justified.

Students shall not use repeated spaces, repeated manual line breaks, or
multiple blank lines to position text. Paragraph and heading spacing shall
be controlled by the template.

\vspace{0.3cm}

% ================================================================
% V. OUTER COVER
% ================================================================

\noindent
\textbf{V.\quad OUTER COVER}

\vspace{0.15cm}

The outer cover shall follow the official format provided in the template
and shall contain the following information:

\begin{itemize}[leftmargin=0.8cm, itemsep=2pt]
    \item Thai Nguyen University of Technology;
    \item Faculty of International Training;
    \item Engineering and Technology Department;
    \item University logo and Faculty logo;
    \item the title \textbf{INTERNSHIP REPORT};
    \item Internship Organization;
    \item Author;
    \item Class;
    \item Supervisor;
    \item place and year of completion.
\end{itemize}

Students shall only replace the required information fields and shall not
modify the cover-page layout unless instructed otherwise.

\vspace{0.3cm}

% ================================================================
% VI. COMMENTS OF THE INTERNSHIP HOST ORGANIZATION
% ================================================================

\noindent
\textbf{VI.\quad COMMENTS OF THE INTERNSHIP HOST ORGANIZATION}

\vspace{0.15cm}

The page entitled \textbf{COMMENTS OF THE INTERNSHIP HOST ORGANIZATION}
is reserved for the internship host organization.

It includes:

\begin{itemize}[leftmargin=0.8cm, itemsep=2pt]
    \item the student's full name;
    \item class;
    \item comments from the host organization;
    \item score;
    \item date;
    \item signature and official seal.
\end{itemize}

Students shall not alter or complete official evaluation fields unless
specifically authorized.

\vspace{0.3cm}

% ================================================================
% VII. SUPERVISOR'S EVALUATION AND DEFENSE RESULTS
% ================================================================

\noindent
\textbf{VII.\quad SUPERVISOR'S EVALUATION AND DEFENSE RESULTS}

\vspace{0.15cm}

The section entitled
\textbf{SUPERVISOR'S EVALUATION AND GRADUATION INTERNSHIP DEFENSE RESULTS}
is reserved for the supervisor and defense committee.

The evaluation may include:

\begin{itemize}[leftmargin=0.8cm, itemsep=2pt]
    \item Internship Progress and Student's Attitude;
    \item Report Content;
    \item Report Presentation;
    \item Other Comments;
    \item Supervisor's Opinion;
    \item Defense Result;
    \item Score.
\end{itemize}

Students shall not enter evaluation scores or defense results before they
are officially determined.

\vspace{0.3cm}

% ================================================================
% VIII. CHAPTERS, SECTIONS, AND SUBSECTIONS
% ================================================================

\noindent
\textbf{VIII.\quad CHAPTERS, SECTIONS, AND SUBSECTIONS}

\vspace{0.15cm}

The report structure shall be generated using the standard \LaTeX{}
commands:

\begin{verbatim}
\chapter{...}
\section{...}
\subsection{...}
\end{verbatim}

Chapter, section, and subsection numbers shall be generated automatically
by \LaTeX{}.

Students shall not manually type chapter or section numbers and shall not
reproduce heading styles using commands such as \verb|\textbf|,
\verb|\fontsize|, or \verb|\centering| instead of the appropriate
structural command.

Excessive subdivision should be avoided. In most cases, the hierarchy
Chapter--Section--Subsection is sufficient.

\vspace{0.3cm}

% ================================================================
% IX. ORGANIZATION OF THE MAIN CONTENT
% ================================================================

\noindent
\textbf{IX.\quad ORGANIZATION OF THE MAIN CONTENT}

\vspace{0.15cm}

The specific content of the report shall depend on the internship
organization and the assigned tasks. A typical structure may include:

\begin{itemize}[leftmargin=0.8cm, itemsep=2pt]
    \item \textbf{Chapter 1. Overview of the Internship Organization};
    \item \textbf{Chapter 2. Internship Activities and Assigned Tasks};
    \item \textbf{Chapter 3. Results, Experience, and Evaluation}.
\end{itemize}

These chapter titles are recommendations only and may be adapted to the
actual internship content.

The report shall reflect work actually performed, observed, analyzed, or
learned during the internship. Unrelated theoretical material should not
be included merely to increase report length.

\vspace{0.3cm}

% ================================================================
% X. TABLE OF CONTENTS
% ================================================================

\noindent
\textbf{X.\quad TABLE OF CONTENTS}

\vspace{0.15cm}

The Table of Contents shall be generated automatically using:

\begin{verbatim}
\tableofcontents
\end{verbatim}

Students shall not create or manually edit page numbers in the Table of
Contents.

The document shall be compiled sufficiently to ensure that all page
numbers and entries are correctly updated.

\vspace{0.3cm}

% ================================================================
% XI. LIST OF SYMBOLS AND ABBREVIATIONS
% ================================================================

\noindent
\textbf{XI.\quad LIST OF SYMBOLS AND ABBREVIATIONS}

\vspace{0.15cm}

Symbols and abbreviations that are used repeatedly in the report shall be
listed in the \textbf{List of Symbols and Abbreviations}.

The list may contain the following columns:

\begin{center}
\begin{tabular}{llll}
    No. & Symbol / Abbreviation & Unit & Definition
\end{tabular}
\end{center}

Only symbols and abbreviations actually used in the final report shall be
retained. Example entries contained in the template shall be removed if
they are not relevant to the student's report.

Each abbreviation should normally be defined when it first appears in
the main text.

\vspace{0.3cm}

% ================================================================
% XII. FIGURES, DIAGRAMS, AND PHOTOGRAPHS
% ================================================================

\noindent
\textbf{XII.\quad FIGURES, DIAGRAMS, AND PHOTOGRAPHS}

\vspace{0.15cm}

Figures shall be clear, legible, and directly relevant to the report
content.

Figures shall be inserted using the standard \verb|figure| environment.

Example:

\begin{verbatim}
\begin{figure}[H]
    \centering
    \includegraphics[width=0.8\textwidth]
                    {figures/company_structure.pdf}
    \caption{Organizational structure of the company}
    \label{fig:company-structure}
\end{figure}
\end{verbatim}

Figure captions shall be placed below the figure.

Each referenced figure shall have a unique \verb|\label| and shall be
referenced in the text using, for example:

\begin{verbatim}
Figure~\ref{fig:company-structure}
\end{verbatim}

Students shall not manually type figure numbers.

Figures reproduced or adapted from external sources shall include an
appropriate citation.

Photographs taken during the internship shall be relevant and sufficiently
clear. Students shall not include confidential, restricted, or proprietary
information without permission from the internship organization.

Vector graphics or high-resolution images are recommended whenever
possible.

\vspace{0.3cm}

% ================================================================
% XIII. TABLES
% ================================================================

\noindent
\textbf{XIII.\quad TABLES}

\vspace{0.15cm}

Tables shall be created using suitable \LaTeX{} environments such as
\verb|table|, \verb|tabular|, \verb|tabularx|, or \verb|longtable|.

Table captions shall be placed above the table.

Example:

\begin{verbatim}
\begin{table}[H]
    \centering
    \caption{Internship schedule}
    \label{tab:internship-schedule}

    \begin{tabular}{lll}
        \toprule
        Week & Task & Result \\
        \midrule
        1 & Orientation & Completed \\
        2 & System study & Completed \\
        \bottomrule
    \end{tabular}
\end{table}
\end{verbatim}

Tables shall be referenced using:

\begin{verbatim}
Table~\ref{tab:internship-schedule}
\end{verbatim}

Students shall not manually type table numbers.

Tables should be created directly in \LaTeX{} rather than inserted as
screenshots whenever practical.

\vspace{0.3cm}

% ================================================================
% XIV. EQUATIONS AND TECHNICAL QUANTITIES
% ================================================================

\noindent
\textbf{XIV.\quad EQUATIONS AND TECHNICAL QUANTITIES}

\vspace{0.15cm}

Mathematical expressions shall be typeset directly using \LaTeX{}.
Images or screenshots of equations shall not be used.

Short expressions may be written in inline mathematical mode, for example:

\begin{verbatim}
$R = 0.7\,\Omega$
\end{verbatim}

Important equations shall use a numbered equation environment.

Example:

\begin{verbatim}
\begin{equation}
    P = UI
    \label{eq:power}
\end{equation}
\end{verbatim}

Equations shall be referenced using:

\begin{verbatim}
Equation~\eqref{eq:power}
\end{verbatim}

Equation numbers shall not be manually entered.

SI units shall be used wherever applicable. The \verb|siunitx| package
should be used where appropriate.

Examples:

\begin{verbatim}
\SI{220}{\volt}
\SI{50}{\hertz}
\SI{2.5}{\ampere}
\end{verbatim}

Notation and units shall be consistent throughout the report.

\vspace{0.3cm}

% ================================================================
% XV. SOURCE CODE
% ================================================================

\noindent
\textbf{XV.\quad SOURCE CODE}

\vspace{0.15cm}

If the internship involves programming, PLCs, microcontrollers,
simulation, or software development, relevant source code may be included
in the report.

Source code shall be presented using the \verb|listings| environment or
another environment defined by the official template.

Example:

\begin{verbatim}
\begin{lstlisting}[language=st]
IF Start = TRUE THEN
    Motor := TRUE;
ELSE
    Motor := FALSE;
END_IF;
\end{lstlisting}
\end{verbatim}

Only code necessary for explaining the work should appear in the main
report. Long source-code listings should be moved to an appendix or
submitted as supplementary files if required.

Source code should not normally be inserted as screenshots.

\vspace{0.3cm}

% ================================================================
% XVI. ABSTRACT
% ================================================================

\noindent
\textbf{XVI.\quad ABSTRACT}

\vspace{0.15cm}

The Abstract shall briefly summarize:

\begin{itemize}[leftmargin=0.8cm, itemsep=2pt]
    \item the internship organization;
    \item the objectives of the internship;
    \item the main tasks performed;
    \item the main results, knowledge, and skills obtained.
\end{itemize}

The Abstract should be concise and should not reproduce the full content
of the report chapter by chapter.

\vspace{0.3cm}

% ================================================================
% XVII. REFERENCES AND CITATIONS
% ================================================================

\noindent
\textbf{XVII.\quad REFERENCES AND CITATIONS}

\vspace{0.15cm}

All books, journal papers, conference papers, standards, websites,
datasheets, manuals, technical documents, figures, tables, or other
external sources used in preparing the report shall be properly cited.

References shall be managed through the bibliography system provided with
the official template.

Bibliographic entries shall normally be stored in:

\begin{verbatim}
References.bib
\end{verbatim}

References shall be cited in the report using:

\begin{verbatim}
\cite{reference-key}
\end{verbatim}

Students shall not manually type reference numbers such as ``[1]'',
``[2]'', etc.

Titles of foreign-language publications shall remain in their original
language in the bibliography.

All borrowed ideas, data, figures, tables, equations, methods, or results
shall be properly acknowledged.

\vspace{0.3cm}

% ================================================================
% XVIII. CROSS-REFERENCING
% ================================================================

\noindent
\textbf{XVIII.\quad CROSS-REFERENCING}

\vspace{0.15cm}

Numbered elements shall be cross-referenced using the standard \LaTeX{}
mechanism.

Use the following commands where appropriate:

\begin{verbatim}
\label{...}
\ref{...}
\eqref{...}
\end{verbatim}

For example:

\begin{verbatim}
As shown in Figure~\ref{fig:system}, ...
\end{verbatim}

Students shall not manually type numbers such as ``Figure 2.3'' when
those numbers can be generated automatically.

\vspace{0.3cm}

% ================================================================
% XIX. PAGE NUMBERING
% ================================================================

\noindent
\textbf{XIX.\quad PAGE NUMBERING}

\vspace{0.15cm}

Cover pages and official evaluation forms at the beginning of the report
shall not display page numbers.

The main body of the report shall use Arabic page numbers beginning from
Chapter 1.

Page numbering shall be generated automatically by the official template.
Students shall not manually enter page numbers.

\vspace{0.3cm}

% ================================================================
% XX. PROJECT FILE ORGANIZATION
% ================================================================

\noindent
\textbf{XX.\quad PROJECT FILE ORGANIZATION}

\vspace{0.15cm}

Students should maintain a clear and organized \LaTeX{} project
structure. A recommended structure is:

\begin{verbatim}
InternshipReport/
|
|-- main.tex
|-- References.bib
|
|-- figures/
|   |-- company.pdf
|   |-- system.png
|   `-- ...
|
`-- chapters/
    |-- chapter1.tex
    |-- chapter2.tex
    `-- chapter3.tex
\end{verbatim}

Figures should normally be stored in the \verb|figures| directory
specified by the official template.

File and directory names should preferably contain Latin letters,
numbers, hyphens, or underscores.

Recommended examples include:

\begin{verbatim}
company_structure.pdf
motor_control.png
chapter_01.tex
\end{verbatim}

Spaces and unnecessary special characters in file names should be
avoided where possible.

\vspace{0.3cm}

% ================================================================
% XXI. DOCUMENT COMPILATION AND FINAL CHECK
% ================================================================

\noindent
\textbf{XXI.\quad DOCUMENT COMPILATION AND FINAL CHECK}

\vspace{0.15cm}

Before submission, the complete \LaTeX{} project shall compile
successfully and produce the final PDF report.

The final document shall not contain unresolved references such as:

\begin{verbatim}
??
[?]
\end{verbatim}

Before generating the final PDF, students shall check:

\begin{itemize}[leftmargin=0.8cm, itemsep=2pt]
    \item the Table of Contents;
    \item page numbering;
    \item chapter, section, and subsection numbering;
    \item figure numbering;
    \item table numbering;
    \item equation numbering;
    \item cross-references;
    \item citations;
    \item bibliography;
    \item font display;
    \item mathematical symbols;
    \item image quality;
    \item completeness of all required forms.
\end{itemize}

The document may be compiled using TeX Live, MiKTeX, or another
compatible \LaTeX{} environment. Regardless of the compilation
environment, the resulting PDF shall conform to the official Faculty
template.

\vspace{0.3cm}

% ================================================================
% XXII. CONTENT THAT STUDENTS MAY EDIT
% ================================================================

\noindent
\textbf{XXII.\quad CONTENT THAT STUDENTS MAY EDIT}

\vspace{0.15cm}

Students may edit or replace:

\begin{itemize}[leftmargin=0.8cm, itemsep=2pt]
    \item student name;
    \item class;
    \item internship organization;
    \item supervisor name;
    \item Abstract;
    \item List of Symbols and Abbreviations;
    \item chapter titles and chapter content;
    \item figures and photographs;
    \item tables;
    \item equations;
    \item citations and references.
\end{itemize}

Students shall not modify the main document formatting settings in the
preamble, including the document class, page geometry, line spacing,
heading formatting, table-of-contents formatting, font configuration,
page layout, or cover-page macros unless specifically instructed.

\vspace{0.3cm}

% ================================================================
% XXIII. FINAL SUBMISSION REQUIREMENTS
% ================================================================

\noindent
\textbf{XXIII.\quad FINAL SUBMISSION REQUIREMENTS}

\vspace{0.15cm}

Before submission, the final report shall:

\begin{itemize}[leftmargin=0.8cm, itemsep=2pt]
    \item compile without errors that affect the report content or presentation;
    \item contain all required forms and report sections;
    \item contain correct student, class, internship organization, and supervisor information;
    \item contain an updated Table of Contents, List of Figures, and List of Tables;
    \item contain no unused example content from the template;
    \item contain no placeholder text such as \textit{Title of Chapter 1},
    \textit{Section 1.1}, or other sample entries;
    \item contain all required citations and references;
    \item contain no unresolved cross-references.
\end{itemize}

If the Faculty requires submission of the \LaTeX{} source, the complete
project shall include all \verb|.tex| files, the bibliography database,
figures, and any other files required to reproduce the submitted PDF.

\vspace{0.4cm}

\noindent
\textbf{Important:}
Before submitting the final internship report, students shall remove or
comment out the following command in the main document:

\begin{verbatim}
\input{Regulations.tex}
\end{verbatim}

For example:

\begin{verbatim}
% \input{Regulations.tex}
\end{verbatim}

This removes the guidance pages from the final submitted report while
preserving the regulations file in the \LaTeX{} project for reference.

} % End 13pt content

\clearpage
|.
Students shall remove or comment out this command before submitting the
final internship report.

\vspace{0.4cm}

% ================================================================
% I. GENERAL REQUIREMENTS
% ================================================================

\noindent
\textbf{I.\quad GENERAL REQUIREMENTS}

\vspace{0.15cm}

The Graduation Internship Report shall be prepared using the official
\LaTeX{} template provided by the Faculty. Students shall preserve the
predefined page layout, typography, line spacing, heading styles,
numbering system, figure and table formatting, and bibliography settings
unless otherwise instructed by the Faculty or the supervisor.

The report shall be written in English.

Students are responsible for the academic and technical content of the
report. The official \LaTeX{} template is responsible for standardized
document formatting.

\vspace{0.3cm}

% ================================================================
% II. REPORT STRUCTURE
% ================================================================

\noindent
\textbf{II.\quad REPORT STRUCTURE}

\vspace{0.15cm}

The report shall normally be arranged in the following order:

\begin{enumerate}[label=\arabic*., leftmargin=1.2cm, itemsep=2pt]
    \item Outer Cover;
    \item Comments of the Internship Host Organization;
    \item Supervisor's Evaluation and Graduation Internship Defense Results;
    \item Table of Contents;
    \item List of Symbols and Abbreviations, if applicable;
    \item List of Figures;
    \item List of Tables;
    \item Abstract;
    \item Main Chapters of the Internship Report;
    \item References;
    \item Appendices, if applicable.
\end{enumerate}

The preliminary pages shall not be numbered as chapters. The main report
content shall be organized into consecutively numbered chapters beginning
with Chapter 1.

\vspace{0.3cm}

% ================================================================
% III. PAGE SIZE AND MARGINS
% ================================================================

\noindent
\textbf{III.\quad PAGE SIZE AND MARGINS}

\vspace{0.15cm}

The report shall be formatted on A4 paper with the following margins:

\begin{center}
\begin{tabular}{ll}
    Left margin:   & 3.5 cm \\
    Right margin:  & 2.0 cm \\
    Top margin:    & 2.5 cm \\
    Bottom margin: & 2.5 cm
\end{tabular}
\end{center}

These settings are predefined in the official template and shall not be
modified by students.

\vspace{0.3cm}

% ================================================================
% IV. FONT, FONT SIZE, AND LINE SPACING
% ================================================================

\noindent
\textbf{IV.\quad FONT, FONT SIZE, AND LINE SPACING}

\vspace{0.15cm}

The main text shall use a 13 pt font size with 1.3-line spacing.
A Times-family typeface shall be used in accordance with the official
template.

The text shall be fully justified.

Students shall not use repeated spaces, repeated manual line breaks, or
multiple blank lines to position text. Paragraph and heading spacing shall
be controlled by the template.

\vspace{0.3cm}

% ================================================================
% V. OUTER COVER
% ================================================================

\noindent
\textbf{V.\quad OUTER COVER}

\vspace{0.15cm}

The outer cover shall follow the official format provided in the template
and shall contain the following information:

\begin{itemize}[leftmargin=0.8cm, itemsep=2pt]
    \item Thai Nguyen University of Technology;
    \item Faculty of International Training;
    \item Engineering and Technology Department;
    \item University logo and Faculty logo;
    \item the title \textbf{INTERNSHIP REPORT};
    \item Internship Organization;
    \item Author;
    \item Class;
    \item Supervisor;
    \item place and year of completion.
\end{itemize}

Students shall only replace the required information fields and shall not
modify the cover-page layout unless instructed otherwise.

\vspace{0.3cm}

% ================================================================
% VI. COMMENTS OF THE INTERNSHIP HOST ORGANIZATION
% ================================================================

\noindent
\textbf{VI.\quad COMMENTS OF THE INTERNSHIP HOST ORGANIZATION}

\vspace{0.15cm}

The page entitled \textbf{COMMENTS OF THE INTERNSHIP HOST ORGANIZATION}
is reserved for the internship host organization.

It includes:

\begin{itemize}[leftmargin=0.8cm, itemsep=2pt]
    \item the student's full name;
    \item class;
    \item comments from the host organization;
    \item score;
    \item date;
    \item signature and official seal.
\end{itemize}

Students shall not alter or complete official evaluation fields unless
specifically authorized.

\vspace{0.3cm}

% ================================================================
% VII. SUPERVISOR'S EVALUATION AND DEFENSE RESULTS
% ================================================================

\noindent
\textbf{VII.\quad SUPERVISOR'S EVALUATION AND DEFENSE RESULTS}

\vspace{0.15cm}

The section entitled
\textbf{SUPERVISOR'S EVALUATION AND GRADUATION INTERNSHIP DEFENSE RESULTS}
is reserved for the supervisor and defense committee.

The evaluation may include:

\begin{itemize}[leftmargin=0.8cm, itemsep=2pt]
    \item Internship Progress and Student's Attitude;
    \item Report Content;
    \item Report Presentation;
    \item Other Comments;
    \item Supervisor's Opinion;
    \item Defense Result;
    \item Score.
\end{itemize}

Students shall not enter evaluation scores or defense results before they
are officially determined.

\vspace{0.3cm}

% ================================================================
% VIII. CHAPTERS, SECTIONS, AND SUBSECTIONS
% ================================================================

\noindent
\textbf{VIII.\quad CHAPTERS, SECTIONS, AND SUBSECTIONS}

\vspace{0.15cm}

The report structure shall be generated using the standard \LaTeX{}
commands:

\begin{verbatim}
\chapter{...}
\section{...}
\subsection{...}
\end{verbatim}

Chapter, section, and subsection numbers shall be generated automatically
by \LaTeX{}.

Students shall not manually type chapter or section numbers and shall not
reproduce heading styles using commands such as \verb|\textbf|,
\verb|\fontsize|, or \verb|\centering| instead of the appropriate
structural command.

Excessive subdivision should be avoided. In most cases, the hierarchy
Chapter--Section--Subsection is sufficient.

\vspace{0.3cm}

% ================================================================
% IX. ORGANIZATION OF THE MAIN CONTENT
% ================================================================

\noindent
\textbf{IX.\quad ORGANIZATION OF THE MAIN CONTENT}

\vspace{0.15cm}

The specific content of the report shall depend on the internship
organization and the assigned tasks. A typical structure may include:

\begin{itemize}[leftmargin=0.8cm, itemsep=2pt]
    \item \textbf{Chapter 1. Overview of the Internship Organization};
    \item \textbf{Chapter 2. Internship Activities and Assigned Tasks};
    \item \textbf{Chapter 3. Results, Experience, and Evaluation}.
\end{itemize}

These chapter titles are recommendations only and may be adapted to the
actual internship content.

The report shall reflect work actually performed, observed, analyzed, or
learned during the internship. Unrelated theoretical material should not
be included merely to increase report length.

\vspace{0.3cm}

% ================================================================
% X. TABLE OF CONTENTS
% ================================================================

\noindent
\textbf{X.\quad TABLE OF CONTENTS}

\vspace{0.15cm}

The Table of Contents shall be generated automatically using:

\begin{verbatim}
\tableofcontents
\end{verbatim}

Students shall not create or manually edit page numbers in the Table of
Contents.

The document shall be compiled sufficiently to ensure that all page
numbers and entries are correctly updated.

\vspace{0.3cm}

% ================================================================
% XI. LIST OF SYMBOLS AND ABBREVIATIONS
% ================================================================

\noindent
\textbf{XI.\quad LIST OF SYMBOLS AND ABBREVIATIONS}

\vspace{0.15cm}

Symbols and abbreviations that are used repeatedly in the report shall be
listed in the \textbf{List of Symbols and Abbreviations}.

The list may contain the following columns:

\begin{center}
\begin{tabular}{llll}
    No. & Symbol / Abbreviation & Unit & Definition
\end{tabular}
\end{center}

Only symbols and abbreviations actually used in the final report shall be
retained. Example entries contained in the template shall be removed if
they are not relevant to the student's report.

Each abbreviation should normally be defined when it first appears in
the main text.

\vspace{0.3cm}

% ================================================================
% XII. FIGURES, DIAGRAMS, AND PHOTOGRAPHS
% ================================================================

\noindent
\textbf{XII.\quad FIGURES, DIAGRAMS, AND PHOTOGRAPHS}

\vspace{0.15cm}

Figures shall be clear, legible, and directly relevant to the report
content.

Figures shall be inserted using the standard \verb|figure| environment.

Example:

\begin{verbatim}
\begin{figure}[H]
    \centering
    \includegraphics[width=0.8\textwidth]
                    {figures/company_structure.pdf}
    \caption{Organizational structure of the company}
    \label{fig:company-structure}
\end{figure}
\end{verbatim}

Figure captions shall be placed below the figure.

Each referenced figure shall have a unique \verb|\label| and shall be
referenced in the text using, for example:

\begin{verbatim}
Figure~\ref{fig:company-structure}
\end{verbatim}

Students shall not manually type figure numbers.

Figures reproduced or adapted from external sources shall include an
appropriate citation.

Photographs taken during the internship shall be relevant and sufficiently
clear. Students shall not include confidential, restricted, or proprietary
information without permission from the internship organization.

Vector graphics or high-resolution images are recommended whenever
possible.

\vspace{0.3cm}

% ================================================================
% XIII. TABLES
% ================================================================

\noindent
\textbf{XIII.\quad TABLES}

\vspace{0.15cm}

Tables shall be created using suitable \LaTeX{} environments such as
\verb|table|, \verb|tabular|, \verb|tabularx|, or \verb|longtable|.

Table captions shall be placed above the table.

Example:

\begin{verbatim}
\begin{table}[H]
    \centering
    \caption{Internship schedule}
    \label{tab:internship-schedule}

    \begin{tabular}{lll}
        \toprule
        Week & Task & Result \\
        \midrule
        1 & Orientation & Completed \\
        2 & System study & Completed \\
        \bottomrule
    \end{tabular}
\end{table}
\end{verbatim}

Tables shall be referenced using:

\begin{verbatim}
Table~\ref{tab:internship-schedule}
\end{verbatim}

Students shall not manually type table numbers.

Tables should be created directly in \LaTeX{} rather than inserted as
screenshots whenever practical.

\vspace{0.3cm}

% ================================================================
% XIV. EQUATIONS AND TECHNICAL QUANTITIES
% ================================================================

\noindent
\textbf{XIV.\quad EQUATIONS AND TECHNICAL QUANTITIES}

\vspace{0.15cm}

Mathematical expressions shall be typeset directly using \LaTeX{}.
Images or screenshots of equations shall not be used.

Short expressions may be written in inline mathematical mode, for example:

\begin{verbatim}
$R = 0.7\,\Omega$
\end{verbatim}

Important equations shall use a numbered equation environment.

Example:

\begin{verbatim}
\begin{equation}
    P = UI
    \label{eq:power}
\end{equation}
\end{verbatim}

Equations shall be referenced using:

\begin{verbatim}
Equation~\eqref{eq:power}
\end{verbatim}

Equation numbers shall not be manually entered.

SI units shall be used wherever applicable. The \verb|siunitx| package
should be used where appropriate.

Examples:

\begin{verbatim}
\SI{220}{\volt}
\SI{50}{\hertz}
\SI{2.5}{\ampere}
\end{verbatim}

Notation and units shall be consistent throughout the report.

\vspace{0.3cm}

% ================================================================
% XV. SOURCE CODE
% ================================================================

\noindent
\textbf{XV.\quad SOURCE CODE}

\vspace{0.15cm}

If the internship involves programming, PLCs, microcontrollers,
simulation, or software development, relevant source code may be included
in the report.

Source code shall be presented using the \verb|listings| environment or
another environment defined by the official template.

Example:

\begin{verbatim}
\begin{lstlisting}[language=st]
IF Start = TRUE THEN
    Motor := TRUE;
ELSE
    Motor := FALSE;
END_IF;
\end{lstlisting}
\end{verbatim}

Only code necessary for explaining the work should appear in the main
report. Long source-code listings should be moved to an appendix or
submitted as supplementary files if required.

Source code should not normally be inserted as screenshots.

\vspace{0.3cm}

% ================================================================
% XVI. ABSTRACT
% ================================================================

\noindent
\textbf{XVI.\quad ABSTRACT}

\vspace{0.15cm}

The Abstract shall briefly summarize:

\begin{itemize}[leftmargin=0.8cm, itemsep=2pt]
    \item the internship organization;
    \item the objectives of the internship;
    \item the main tasks performed;
    \item the main results, knowledge, and skills obtained.
\end{itemize}

The Abstract should be concise and should not reproduce the full content
of the report chapter by chapter.

\vspace{0.3cm}

% ================================================================
% XVII. REFERENCES AND CITATIONS
% ================================================================

\noindent
\textbf{XVII.\quad REFERENCES AND CITATIONS}

\vspace{0.15cm}

All books, journal papers, conference papers, standards, websites,
datasheets, manuals, technical documents, figures, tables, or other
external sources used in preparing the report shall be properly cited.

References shall be managed through the bibliography system provided with
the official template.

Bibliographic entries shall normally be stored in:

\begin{verbatim}
References.bib
\end{verbatim}

References shall be cited in the report using:

\begin{verbatim}
\cite{reference-key}
\end{verbatim}

Students shall not manually type reference numbers such as ``[1]'',
``[2]'', etc.

Titles of foreign-language publications shall remain in their original
language in the bibliography.

All borrowed ideas, data, figures, tables, equations, methods, or results
shall be properly acknowledged.

\vspace{0.3cm}

% ================================================================
% XVIII. CROSS-REFERENCING
% ================================================================

\noindent
\textbf{XVIII.\quad CROSS-REFERENCING}

\vspace{0.15cm}

Numbered elements shall be cross-referenced using the standard \LaTeX{}
mechanism.

Use the following commands where appropriate:

\begin{verbatim}
\label{...}
\ref{...}
\eqref{...}
\end{verbatim}

For example:

\begin{verbatim}
As shown in Figure~\ref{fig:system}, ...
\end{verbatim}

Students shall not manually type numbers such as ``Figure 2.3'' when
those numbers can be generated automatically.

\vspace{0.3cm}

% ================================================================
% XIX. PAGE NUMBERING
% ================================================================

\noindent
\textbf{XIX.\quad PAGE NUMBERING}

\vspace{0.15cm}

Cover pages and official evaluation forms at the beginning of the report
shall not display page numbers.

The main body of the report shall use Arabic page numbers beginning from
Chapter 1.

Page numbering shall be generated automatically by the official template.
Students shall not manually enter page numbers.

\vspace{0.3cm}

% ================================================================
% XX. PROJECT FILE ORGANIZATION
% ================================================================

\noindent
\textbf{XX.\quad PROJECT FILE ORGANIZATION}

\vspace{0.15cm}

Students should maintain a clear and organized \LaTeX{} project
structure. A recommended structure is:

\begin{verbatim}
InternshipReport/
|
|-- main.tex
|-- References.bib
|
|-- figures/
|   |-- company.pdf
|   |-- system.png
|   `-- ...
|
`-- chapters/
    |-- chapter1.tex
    |-- chapter2.tex
    `-- chapter3.tex
\end{verbatim}

Figures should normally be stored in the \verb|figures| directory
specified by the official template.

File and directory names should preferably contain Latin letters,
numbers, hyphens, or underscores.

Recommended examples include:

\begin{verbatim}
company_structure.pdf
motor_control.png
chapter_01.tex
\end{verbatim}

Spaces and unnecessary special characters in file names should be
avoided where possible.

\vspace{0.3cm}

% ================================================================
% XXI. DOCUMENT COMPILATION AND FINAL CHECK
% ================================================================

\noindent
\textbf{XXI.\quad DOCUMENT COMPILATION AND FINAL CHECK}

\vspace{0.15cm}

Before submission, the complete \LaTeX{} project shall compile
successfully and produce the final PDF report.

The final document shall not contain unresolved references such as:

\begin{verbatim}
??
[?]
\end{verbatim}

Before generating the final PDF, students shall check:

\begin{itemize}[leftmargin=0.8cm, itemsep=2pt]
    \item the Table of Contents;
    \item page numbering;
    \item chapter, section, and subsection numbering;
    \item figure numbering;
    \item table numbering;
    \item equation numbering;
    \item cross-references;
    \item citations;
    \item bibliography;
    \item font display;
    \item mathematical symbols;
    \item image quality;
    \item completeness of all required forms.
\end{itemize}

The document may be compiled using TeX Live, MiKTeX, or another
compatible \LaTeX{} environment. Regardless of the compilation
environment, the resulting PDF shall conform to the official Faculty
template.

\vspace{0.3cm}

% ================================================================
% XXII. CONTENT THAT STUDENTS MAY EDIT
% ================================================================

\noindent
\textbf{XXII.\quad CONTENT THAT STUDENTS MAY EDIT}

\vspace{0.15cm}

Students may edit or replace:

\begin{itemize}[leftmargin=0.8cm, itemsep=2pt]
    \item student name;
    \item class;
    \item internship organization;
    \item supervisor name;
    \item Abstract;
    \item List of Symbols and Abbreviations;
    \item chapter titles and chapter content;
    \item figures and photographs;
    \item tables;
    \item equations;
    \item citations and references.
\end{itemize}

Students shall not modify the main document formatting settings in the
preamble, including the document class, page geometry, line spacing,
heading formatting, table-of-contents formatting, font configuration,
page layout, or cover-page macros unless specifically instructed.

\vspace{0.3cm}

% ================================================================
% XXIII. FINAL SUBMISSION REQUIREMENTS
% ================================================================

\noindent
\textbf{XXIII.\quad FINAL SUBMISSION REQUIREMENTS}

\vspace{0.15cm}

Before submission, the final report shall:

\begin{itemize}[leftmargin=0.8cm, itemsep=2pt]
    \item compile without errors that affect the report content or presentation;
    \item contain all required forms and report sections;
    \item contain correct student, class, internship organization, and supervisor information;
    \item contain an updated Table of Contents, List of Figures, and List of Tables;
    \item contain no unused example content from the template;
    \item contain no placeholder text such as \textit{Title of Chapter 1},
    \textit{Section 1.1}, or other sample entries;
    \item contain all required citations and references;
    \item contain no unresolved cross-references.
\end{itemize}

If the Faculty requires submission of the \LaTeX{} source, the complete
project shall include all \verb|.tex| files, the bibliography database,
figures, and any other files required to reproduce the submitted PDF.

\vspace{0.4cm}

\noindent
\textbf{Important:}
Before submitting the final internship report, students shall remove or
comment out the following command in the main document:

\begin{verbatim}

% =================================================================
% REGULATIONS FOR PREPARING THE GRADUATION INTERNSHIP REPORT
% USING LATEX
% =================================================================
%
% This file is intended to be included from the main document with:
%
%     \input{Regulations.tex}
%
% IMPORTANT:
% This section is provided for guidance only. Students shall remove
% or comment out the \input{Regulations.tex} command before submitting
% the final internship report.
%
% =================================================================

\clearpage
\thispagestyle{empty}

\begin{center}
    {\fontsize{16pt}{20pt}\selectfont\bfseries
    SPECIFIC REGULATIONS FOR PREPARING\\
    GRADUATION INTERNSHIP REPORTS USING \LaTeX}
\end{center}

\vspace{0.5cm}

{\fontsize{13pt}{17pt}\selectfont
\setstretch{1.3}
\justifying

\noindent
\textbf{Note:}
This section is provided for guidance only and is included in the main
document using the command \verb|\input{Regulations.tex}|.
Students shall remove or comment out this command before submitting the
final internship report.

\vspace{0.4cm}

% ================================================================
% I. GENERAL REQUIREMENTS
% ================================================================

\noindent
\textbf{I.\quad GENERAL REQUIREMENTS}

\vspace{0.15cm}

The Graduation Internship Report shall be prepared using the official
\LaTeX{} template provided by the Faculty. Students shall preserve the
predefined page layout, typography, line spacing, heading styles,
numbering system, figure and table formatting, and bibliography settings
unless otherwise instructed by the Faculty or the supervisor.

The report shall be written in English.

Students are responsible for the academic and technical content of the
report. The official \LaTeX{} template is responsible for standardized
document formatting.

\vspace{0.3cm}

% ================================================================
% II. REPORT STRUCTURE
% ================================================================

\noindent
\textbf{II.\quad REPORT STRUCTURE}

\vspace{0.15cm}

The report shall normally be arranged in the following order:

\begin{enumerate}[label=\arabic*., leftmargin=1.2cm, itemsep=2pt]
    \item Outer Cover;
    \item Comments of the Internship Host Organization;
    \item Supervisor's Evaluation and Graduation Internship Defense Results;
    \item Table of Contents;
    \item List of Symbols and Abbreviations, if applicable;
    \item List of Figures;
    \item List of Tables;
    \item Abstract;
    \item Main Chapters of the Internship Report;
    \item References;
    \item Appendices, if applicable.
\end{enumerate}

The preliminary pages shall not be numbered as chapters. The main report
content shall be organized into consecutively numbered chapters beginning
with Chapter 1.

\vspace{0.3cm}

% ================================================================
% III. PAGE SIZE AND MARGINS
% ================================================================

\noindent
\textbf{III.\quad PAGE SIZE AND MARGINS}

\vspace{0.15cm}

The report shall be formatted on A4 paper with the following margins:

\begin{center}
\begin{tabular}{ll}
    Left margin:   & 3.5 cm \\
    Right margin:  & 2.0 cm \\
    Top margin:    & 2.5 cm \\
    Bottom margin: & 2.5 cm
\end{tabular}
\end{center}

These settings are predefined in the official template and shall not be
modified by students.

\vspace{0.3cm}

% ================================================================
% IV. FONT, FONT SIZE, AND LINE SPACING
% ================================================================

\noindent
\textbf{IV.\quad FONT, FONT SIZE, AND LINE SPACING}

\vspace{0.15cm}

The main text shall use a 13 pt font size with 1.3-line spacing.
A Times-family typeface shall be used in accordance with the official
template.

The text shall be fully justified.

Students shall not use repeated spaces, repeated manual line breaks, or
multiple blank lines to position text. Paragraph and heading spacing shall
be controlled by the template.

\vspace{0.3cm}

% ================================================================
% V. OUTER COVER
% ================================================================

\noindent
\textbf{V.\quad OUTER COVER}

\vspace{0.15cm}

The outer cover shall follow the official format provided in the template
and shall contain the following information:

\begin{itemize}[leftmargin=0.8cm, itemsep=2pt]
    \item Thai Nguyen University of Technology;
    \item Faculty of International Training;
    \item Engineering and Technology Department;
    \item University logo and Faculty logo;
    \item the title \textbf{INTERNSHIP REPORT};
    \item Internship Organization;
    \item Author;
    \item Class;
    \item Supervisor;
    \item place and year of completion.
\end{itemize}

Students shall only replace the required information fields and shall not
modify the cover-page layout unless instructed otherwise.

\vspace{0.3cm}

% ================================================================
% VI. COMMENTS OF THE INTERNSHIP HOST ORGANIZATION
% ================================================================

\noindent
\textbf{VI.\quad COMMENTS OF THE INTERNSHIP HOST ORGANIZATION}

\vspace{0.15cm}

The page entitled \textbf{COMMENTS OF THE INTERNSHIP HOST ORGANIZATION}
is reserved for the internship host organization.

It includes:

\begin{itemize}[leftmargin=0.8cm, itemsep=2pt]
    \item the student's full name;
    \item class;
    \item comments from the host organization;
    \item score;
    \item date;
    \item signature and official seal.
\end{itemize}

Students shall not alter or complete official evaluation fields unless
specifically authorized.

\vspace{0.3cm}

% ================================================================
% VII. SUPERVISOR'S EVALUATION AND DEFENSE RESULTS
% ================================================================

\noindent
\textbf{VII.\quad SUPERVISOR'S EVALUATION AND DEFENSE RESULTS}

\vspace{0.15cm}

The section entitled
\textbf{SUPERVISOR'S EVALUATION AND GRADUATION INTERNSHIP DEFENSE RESULTS}
is reserved for the supervisor and defense committee.

The evaluation may include:

\begin{itemize}[leftmargin=0.8cm, itemsep=2pt]
    \item Internship Progress and Student's Attitude;
    \item Report Content;
    \item Report Presentation;
    \item Other Comments;
    \item Supervisor's Opinion;
    \item Defense Result;
    \item Score.
\end{itemize}

Students shall not enter evaluation scores or defense results before they
are officially determined.

\vspace{0.3cm}

% ================================================================
% VIII. CHAPTERS, SECTIONS, AND SUBSECTIONS
% ================================================================

\noindent
\textbf{VIII.\quad CHAPTERS, SECTIONS, AND SUBSECTIONS}

\vspace{0.15cm}

The report structure shall be generated using the standard \LaTeX{}
commands:

\begin{verbatim}
\chapter{...}
\section{...}
\subsection{...}
\end{verbatim}

Chapter, section, and subsection numbers shall be generated automatically
by \LaTeX{}.

Students shall not manually type chapter or section numbers and shall not
reproduce heading styles using commands such as \verb|\textbf|,
\verb|\fontsize|, or \verb|\centering| instead of the appropriate
structural command.

Excessive subdivision should be avoided. In most cases, the hierarchy
Chapter--Section--Subsection is sufficient.

\vspace{0.3cm}

% ================================================================
% IX. ORGANIZATION OF THE MAIN CONTENT
% ================================================================

\noindent
\textbf{IX.\quad ORGANIZATION OF THE MAIN CONTENT}

\vspace{0.15cm}

The specific content of the report shall depend on the internship
organization and the assigned tasks. A typical structure may include:

\begin{itemize}[leftmargin=0.8cm, itemsep=2pt]
    \item \textbf{Chapter 1. Overview of the Internship Organization};
    \item \textbf{Chapter 2. Internship Activities and Assigned Tasks};
    \item \textbf{Chapter 3. Results, Experience, and Evaluation}.
\end{itemize}

These chapter titles are recommendations only and may be adapted to the
actual internship content.

The report shall reflect work actually performed, observed, analyzed, or
learned during the internship. Unrelated theoretical material should not
be included merely to increase report length.

\vspace{0.3cm}

% ================================================================
% X. TABLE OF CONTENTS
% ================================================================

\noindent
\textbf{X.\quad TABLE OF CONTENTS}

\vspace{0.15cm}

The Table of Contents shall be generated automatically using:

\begin{verbatim}
\tableofcontents
\end{verbatim}

Students shall not create or manually edit page numbers in the Table of
Contents.

The document shall be compiled sufficiently to ensure that all page
numbers and entries are correctly updated.

\vspace{0.3cm}

% ================================================================
% XI. LIST OF SYMBOLS AND ABBREVIATIONS
% ================================================================

\noindent
\textbf{XI.\quad LIST OF SYMBOLS AND ABBREVIATIONS}

\vspace{0.15cm}

Symbols and abbreviations that are used repeatedly in the report shall be
listed in the \textbf{List of Symbols and Abbreviations}.

The list may contain the following columns:

\begin{center}
\begin{tabular}{llll}
    No. & Symbol / Abbreviation & Unit & Definition
\end{tabular}
\end{center}

Only symbols and abbreviations actually used in the final report shall be
retained. Example entries contained in the template shall be removed if
they are not relevant to the student's report.

Each abbreviation should normally be defined when it first appears in
the main text.

\vspace{0.3cm}

% ================================================================
% XII. FIGURES, DIAGRAMS, AND PHOTOGRAPHS
% ================================================================

\noindent
\textbf{XII.\quad FIGURES, DIAGRAMS, AND PHOTOGRAPHS}

\vspace{0.15cm}

Figures shall be clear, legible, and directly relevant to the report
content.

Figures shall be inserted using the standard \verb|figure| environment.

Example:

\begin{verbatim}
\begin{figure}[H]
    \centering
    \includegraphics[width=0.8\textwidth]
                    {figures/company_structure.pdf}
    \caption{Organizational structure of the company}
    \label{fig:company-structure}
\end{figure}
\end{verbatim}

Figure captions shall be placed below the figure.

Each referenced figure shall have a unique \verb|\label| and shall be
referenced in the text using, for example:

\begin{verbatim}
Figure~\ref{fig:company-structure}
\end{verbatim}

Students shall not manually type figure numbers.

Figures reproduced or adapted from external sources shall include an
appropriate citation.

Photographs taken during the internship shall be relevant and sufficiently
clear. Students shall not include confidential, restricted, or proprietary
information without permission from the internship organization.

Vector graphics or high-resolution images are recommended whenever
possible.

\vspace{0.3cm}

% ================================================================
% XIII. TABLES
% ================================================================

\noindent
\textbf{XIII.\quad TABLES}

\vspace{0.15cm}

Tables shall be created using suitable \LaTeX{} environments such as
\verb|table|, \verb|tabular|, \verb|tabularx|, or \verb|longtable|.

Table captions shall be placed above the table.

Example:

\begin{verbatim}
\begin{table}[H]
    \centering
    \caption{Internship schedule}
    \label{tab:internship-schedule}

    \begin{tabular}{lll}
        \toprule
        Week & Task & Result \\
        \midrule
        1 & Orientation & Completed \\
        2 & System study & Completed \\
        \bottomrule
    \end{tabular}
\end{table}
\end{verbatim}

Tables shall be referenced using:

\begin{verbatim}
Table~\ref{tab:internship-schedule}
\end{verbatim}

Students shall not manually type table numbers.

Tables should be created directly in \LaTeX{} rather than inserted as
screenshots whenever practical.

\vspace{0.3cm}

% ================================================================
% XIV. EQUATIONS AND TECHNICAL QUANTITIES
% ================================================================

\noindent
\textbf{XIV.\quad EQUATIONS AND TECHNICAL QUANTITIES}

\vspace{0.15cm}

Mathematical expressions shall be typeset directly using \LaTeX{}.
Images or screenshots of equations shall not be used.

Short expressions may be written in inline mathematical mode, for example:

\begin{verbatim}
$R = 0.7\,\Omega$
\end{verbatim}

Important equations shall use a numbered equation environment.

Example:

\begin{verbatim}
\begin{equation}
    P = UI
    \label{eq:power}
\end{equation}
\end{verbatim}

Equations shall be referenced using:

\begin{verbatim}
Equation~\eqref{eq:power}
\end{verbatim}

Equation numbers shall not be manually entered.

SI units shall be used wherever applicable. The \verb|siunitx| package
should be used where appropriate.

Examples:

\begin{verbatim}
\SI{220}{\volt}
\SI{50}{\hertz}
\SI{2.5}{\ampere}
\end{verbatim}

Notation and units shall be consistent throughout the report.

\vspace{0.3cm}

% ================================================================
% XV. SOURCE CODE
% ================================================================

\noindent
\textbf{XV.\quad SOURCE CODE}

\vspace{0.15cm}

If the internship involves programming, PLCs, microcontrollers,
simulation, or software development, relevant source code may be included
in the report.

Source code shall be presented using the \verb|listings| environment or
another environment defined by the official template.

Example:

\begin{verbatim}
\begin{lstlisting}[language=st]
IF Start = TRUE THEN
    Motor := TRUE;
ELSE
    Motor := FALSE;
END_IF;
\end{lstlisting}
\end{verbatim}

Only code necessary for explaining the work should appear in the main
report. Long source-code listings should be moved to an appendix or
submitted as supplementary files if required.

Source code should not normally be inserted as screenshots.

\vspace{0.3cm}

% ================================================================
% XVI. ABSTRACT
% ================================================================

\noindent
\textbf{XVI.\quad ABSTRACT}

\vspace{0.15cm}

The Abstract shall briefly summarize:

\begin{itemize}[leftmargin=0.8cm, itemsep=2pt]
    \item the internship organization;
    \item the objectives of the internship;
    \item the main tasks performed;
    \item the main results, knowledge, and skills obtained.
\end{itemize}

The Abstract should be concise and should not reproduce the full content
of the report chapter by chapter.

\vspace{0.3cm}

% ================================================================
% XVII. REFERENCES AND CITATIONS
% ================================================================

\noindent
\textbf{XVII.\quad REFERENCES AND CITATIONS}

\vspace{0.15cm}

All books, journal papers, conference papers, standards, websites,
datasheets, manuals, technical documents, figures, tables, or other
external sources used in preparing the report shall be properly cited.

References shall be managed through the bibliography system provided with
the official template.

Bibliographic entries shall normally be stored in:

\begin{verbatim}
References.bib
\end{verbatim}

References shall be cited in the report using:

\begin{verbatim}
\cite{reference-key}
\end{verbatim}

Students shall not manually type reference numbers such as ``[1]'',
``[2]'', etc.

Titles of foreign-language publications shall remain in their original
language in the bibliography.

All borrowed ideas, data, figures, tables, equations, methods, or results
shall be properly acknowledged.

\vspace{0.3cm}

% ================================================================
% XVIII. CROSS-REFERENCING
% ================================================================

\noindent
\textbf{XVIII.\quad CROSS-REFERENCING}

\vspace{0.15cm}

Numbered elements shall be cross-referenced using the standard \LaTeX{}
mechanism.

Use the following commands where appropriate:

\begin{verbatim}
\label{...}
\ref{...}
\eqref{...}
\end{verbatim}

For example:

\begin{verbatim}
As shown in Figure~\ref{fig:system}, ...
\end{verbatim}

Students shall not manually type numbers such as ``Figure 2.3'' when
those numbers can be generated automatically.

\vspace{0.3cm}

% ================================================================
% XIX. PAGE NUMBERING
% ================================================================

\noindent
\textbf{XIX.\quad PAGE NUMBERING}

\vspace{0.15cm}

Cover pages and official evaluation forms at the beginning of the report
shall not display page numbers.

The main body of the report shall use Arabic page numbers beginning from
Chapter 1.

Page numbering shall be generated automatically by the official template.
Students shall not manually enter page numbers.

\vspace{0.3cm}

% ================================================================
% XX. PROJECT FILE ORGANIZATION
% ================================================================

\noindent
\textbf{XX.\quad PROJECT FILE ORGANIZATION}

\vspace{0.15cm}

Students should maintain a clear and organized \LaTeX{} project
structure. A recommended structure is:

\begin{verbatim}
InternshipReport/
|
|-- main.tex
|-- References.bib
|
|-- figures/
|   |-- company.pdf
|   |-- system.png
|   `-- ...
|
`-- chapters/
    |-- chapter1.tex
    |-- chapter2.tex
    `-- chapter3.tex
\end{verbatim}

Figures should normally be stored in the \verb|figures| directory
specified by the official template.

File and directory names should preferably contain Latin letters,
numbers, hyphens, or underscores.

Recommended examples include:

\begin{verbatim}
company_structure.pdf
motor_control.png
chapter_01.tex
\end{verbatim}

Spaces and unnecessary special characters in file names should be
avoided where possible.

\vspace{0.3cm}

% ================================================================
% XXI. DOCUMENT COMPILATION AND FINAL CHECK
% ================================================================

\noindent
\textbf{XXI.\quad DOCUMENT COMPILATION AND FINAL CHECK}

\vspace{0.15cm}

Before submission, the complete \LaTeX{} project shall compile
successfully and produce the final PDF report.

The final document shall not contain unresolved references such as:

\begin{verbatim}
??
[?]
\end{verbatim}

Before generating the final PDF, students shall check:

\begin{itemize}[leftmargin=0.8cm, itemsep=2pt]
    \item the Table of Contents;
    \item page numbering;
    \item chapter, section, and subsection numbering;
    \item figure numbering;
    \item table numbering;
    \item equation numbering;
    \item cross-references;
    \item citations;
    \item bibliography;
    \item font display;
    \item mathematical symbols;
    \item image quality;
    \item completeness of all required forms.
\end{itemize}

The document may be compiled using TeX Live, MiKTeX, or another
compatible \LaTeX{} environment. Regardless of the compilation
environment, the resulting PDF shall conform to the official Faculty
template.

\vspace{0.3cm}

% ================================================================
% XXII. CONTENT THAT STUDENTS MAY EDIT
% ================================================================

\noindent
\textbf{XXII.\quad CONTENT THAT STUDENTS MAY EDIT}

\vspace{0.15cm}

Students may edit or replace:

\begin{itemize}[leftmargin=0.8cm, itemsep=2pt]
    \item student name;
    \item class;
    \item internship organization;
    \item supervisor name;
    \item Abstract;
    \item List of Symbols and Abbreviations;
    \item chapter titles and chapter content;
    \item figures and photographs;
    \item tables;
    \item equations;
    \item citations and references.
\end{itemize}

Students shall not modify the main document formatting settings in the
preamble, including the document class, page geometry, line spacing,
heading formatting, table-of-contents formatting, font configuration,
page layout, or cover-page macros unless specifically instructed.

\vspace{0.3cm}

% ================================================================
% XXIII. FINAL SUBMISSION REQUIREMENTS
% ================================================================

\noindent
\textbf{XXIII.\quad FINAL SUBMISSION REQUIREMENTS}

\vspace{0.15cm}

Before submission, the final report shall:

\begin{itemize}[leftmargin=0.8cm, itemsep=2pt]
    \item compile without errors that affect the report content or presentation;
    \item contain all required forms and report sections;
    \item contain correct student, class, internship organization, and supervisor information;
    \item contain an updated Table of Contents, List of Figures, and List of Tables;
    \item contain no unused example content from the template;
    \item contain no placeholder text such as \textit{Title of Chapter 1},
    \textit{Section 1.1}, or other sample entries;
    \item contain all required citations and references;
    \item contain no unresolved cross-references.
\end{itemize}

If the Faculty requires submission of the \LaTeX{} source, the complete
project shall include all \verb|.tex| files, the bibliography database,
figures, and any other files required to reproduce the submitted PDF.

\vspace{0.4cm}

\noindent
\textbf{Important:}
Before submitting the final internship report, students shall remove or
comment out the following command in the main document:

\begin{verbatim}
\input{Regulations.tex}
\end{verbatim}

For example:

\begin{verbatim}
% \input{Regulations.tex}
\end{verbatim}

This removes the guidance pages from the final submitted report while
preserving the regulations file in the \LaTeX{} project for reference.

} % End 13pt content

\clearpage

\end{verbatim}

For example:

\begin{verbatim}
% 
% =================================================================
% REGULATIONS FOR PREPARING THE GRADUATION INTERNSHIP REPORT
% USING LATEX
% =================================================================
%
% This file is intended to be included from the main document with:
%
%     \input{Regulations.tex}
%
% IMPORTANT:
% This section is provided for guidance only. Students shall remove
% or comment out the \input{Regulations.tex} command before submitting
% the final internship report.
%
% =================================================================

\clearpage
\thispagestyle{empty}

\begin{center}
    {\fontsize{16pt}{20pt}\selectfont\bfseries
    SPECIFIC REGULATIONS FOR PREPARING\\
    GRADUATION INTERNSHIP REPORTS USING \LaTeX}
\end{center}

\vspace{0.5cm}

{\fontsize{13pt}{17pt}\selectfont
\setstretch{1.3}
\justifying

\noindent
\textbf{Note:}
This section is provided for guidance only and is included in the main
document using the command \verb|\input{Regulations.tex}|.
Students shall remove or comment out this command before submitting the
final internship report.

\vspace{0.4cm}

% ================================================================
% I. GENERAL REQUIREMENTS
% ================================================================

\noindent
\textbf{I.\quad GENERAL REQUIREMENTS}

\vspace{0.15cm}

The Graduation Internship Report shall be prepared using the official
\LaTeX{} template provided by the Faculty. Students shall preserve the
predefined page layout, typography, line spacing, heading styles,
numbering system, figure and table formatting, and bibliography settings
unless otherwise instructed by the Faculty or the supervisor.

The report shall be written in English.

Students are responsible for the academic and technical content of the
report. The official \LaTeX{} template is responsible for standardized
document formatting.

\vspace{0.3cm}

% ================================================================
% II. REPORT STRUCTURE
% ================================================================

\noindent
\textbf{II.\quad REPORT STRUCTURE}

\vspace{0.15cm}

The report shall normally be arranged in the following order:

\begin{enumerate}[label=\arabic*., leftmargin=1.2cm, itemsep=2pt]
    \item Outer Cover;
    \item Comments of the Internship Host Organization;
    \item Supervisor's Evaluation and Graduation Internship Defense Results;
    \item Table of Contents;
    \item List of Symbols and Abbreviations, if applicable;
    \item List of Figures;
    \item List of Tables;
    \item Abstract;
    \item Main Chapters of the Internship Report;
    \item References;
    \item Appendices, if applicable.
\end{enumerate}

The preliminary pages shall not be numbered as chapters. The main report
content shall be organized into consecutively numbered chapters beginning
with Chapter 1.

\vspace{0.3cm}

% ================================================================
% III. PAGE SIZE AND MARGINS
% ================================================================

\noindent
\textbf{III.\quad PAGE SIZE AND MARGINS}

\vspace{0.15cm}

The report shall be formatted on A4 paper with the following margins:

\begin{center}
\begin{tabular}{ll}
    Left margin:   & 3.5 cm \\
    Right margin:  & 2.0 cm \\
    Top margin:    & 2.5 cm \\
    Bottom margin: & 2.5 cm
\end{tabular}
\end{center}

These settings are predefined in the official template and shall not be
modified by students.

\vspace{0.3cm}

% ================================================================
% IV. FONT, FONT SIZE, AND LINE SPACING
% ================================================================

\noindent
\textbf{IV.\quad FONT, FONT SIZE, AND LINE SPACING}

\vspace{0.15cm}

The main text shall use a 13 pt font size with 1.3-line spacing.
A Times-family typeface shall be used in accordance with the official
template.

The text shall be fully justified.

Students shall not use repeated spaces, repeated manual line breaks, or
multiple blank lines to position text. Paragraph and heading spacing shall
be controlled by the template.

\vspace{0.3cm}

% ================================================================
% V. OUTER COVER
% ================================================================

\noindent
\textbf{V.\quad OUTER COVER}

\vspace{0.15cm}

The outer cover shall follow the official format provided in the template
and shall contain the following information:

\begin{itemize}[leftmargin=0.8cm, itemsep=2pt]
    \item Thai Nguyen University of Technology;
    \item Faculty of International Training;
    \item Engineering and Technology Department;
    \item University logo and Faculty logo;
    \item the title \textbf{INTERNSHIP REPORT};
    \item Internship Organization;
    \item Author;
    \item Class;
    \item Supervisor;
    \item place and year of completion.
\end{itemize}

Students shall only replace the required information fields and shall not
modify the cover-page layout unless instructed otherwise.

\vspace{0.3cm}

% ================================================================
% VI. COMMENTS OF THE INTERNSHIP HOST ORGANIZATION
% ================================================================

\noindent
\textbf{VI.\quad COMMENTS OF THE INTERNSHIP HOST ORGANIZATION}

\vspace{0.15cm}

The page entitled \textbf{COMMENTS OF THE INTERNSHIP HOST ORGANIZATION}
is reserved for the internship host organization.

It includes:

\begin{itemize}[leftmargin=0.8cm, itemsep=2pt]
    \item the student's full name;
    \item class;
    \item comments from the host organization;
    \item score;
    \item date;
    \item signature and official seal.
\end{itemize}

Students shall not alter or complete official evaluation fields unless
specifically authorized.

\vspace{0.3cm}

% ================================================================
% VII. SUPERVISOR'S EVALUATION AND DEFENSE RESULTS
% ================================================================

\noindent
\textbf{VII.\quad SUPERVISOR'S EVALUATION AND DEFENSE RESULTS}

\vspace{0.15cm}

The section entitled
\textbf{SUPERVISOR'S EVALUATION AND GRADUATION INTERNSHIP DEFENSE RESULTS}
is reserved for the supervisor and defense committee.

The evaluation may include:

\begin{itemize}[leftmargin=0.8cm, itemsep=2pt]
    \item Internship Progress and Student's Attitude;
    \item Report Content;
    \item Report Presentation;
    \item Other Comments;
    \item Supervisor's Opinion;
    \item Defense Result;
    \item Score.
\end{itemize}

Students shall not enter evaluation scores or defense results before they
are officially determined.

\vspace{0.3cm}

% ================================================================
% VIII. CHAPTERS, SECTIONS, AND SUBSECTIONS
% ================================================================

\noindent
\textbf{VIII.\quad CHAPTERS, SECTIONS, AND SUBSECTIONS}

\vspace{0.15cm}

The report structure shall be generated using the standard \LaTeX{}
commands:

\begin{verbatim}
\chapter{...}
\section{...}
\subsection{...}
\end{verbatim}

Chapter, section, and subsection numbers shall be generated automatically
by \LaTeX{}.

Students shall not manually type chapter or section numbers and shall not
reproduce heading styles using commands such as \verb|\textbf|,
\verb|\fontsize|, or \verb|\centering| instead of the appropriate
structural command.

Excessive subdivision should be avoided. In most cases, the hierarchy
Chapter--Section--Subsection is sufficient.

\vspace{0.3cm}

% ================================================================
% IX. ORGANIZATION OF THE MAIN CONTENT
% ================================================================

\noindent
\textbf{IX.\quad ORGANIZATION OF THE MAIN CONTENT}

\vspace{0.15cm}

The specific content of the report shall depend on the internship
organization and the assigned tasks. A typical structure may include:

\begin{itemize}[leftmargin=0.8cm, itemsep=2pt]
    \item \textbf{Chapter 1. Overview of the Internship Organization};
    \item \textbf{Chapter 2. Internship Activities and Assigned Tasks};
    \item \textbf{Chapter 3. Results, Experience, and Evaluation}.
\end{itemize}

These chapter titles are recommendations only and may be adapted to the
actual internship content.

The report shall reflect work actually performed, observed, analyzed, or
learned during the internship. Unrelated theoretical material should not
be included merely to increase report length.

\vspace{0.3cm}

% ================================================================
% X. TABLE OF CONTENTS
% ================================================================

\noindent
\textbf{X.\quad TABLE OF CONTENTS}

\vspace{0.15cm}

The Table of Contents shall be generated automatically using:

\begin{verbatim}
\tableofcontents
\end{verbatim}

Students shall not create or manually edit page numbers in the Table of
Contents.

The document shall be compiled sufficiently to ensure that all page
numbers and entries are correctly updated.

\vspace{0.3cm}

% ================================================================
% XI. LIST OF SYMBOLS AND ABBREVIATIONS
% ================================================================

\noindent
\textbf{XI.\quad LIST OF SYMBOLS AND ABBREVIATIONS}

\vspace{0.15cm}

Symbols and abbreviations that are used repeatedly in the report shall be
listed in the \textbf{List of Symbols and Abbreviations}.

The list may contain the following columns:

\begin{center}
\begin{tabular}{llll}
    No. & Symbol / Abbreviation & Unit & Definition
\end{tabular}
\end{center}

Only symbols and abbreviations actually used in the final report shall be
retained. Example entries contained in the template shall be removed if
they are not relevant to the student's report.

Each abbreviation should normally be defined when it first appears in
the main text.

\vspace{0.3cm}

% ================================================================
% XII. FIGURES, DIAGRAMS, AND PHOTOGRAPHS
% ================================================================

\noindent
\textbf{XII.\quad FIGURES, DIAGRAMS, AND PHOTOGRAPHS}

\vspace{0.15cm}

Figures shall be clear, legible, and directly relevant to the report
content.

Figures shall be inserted using the standard \verb|figure| environment.

Example:

\begin{verbatim}
\begin{figure}[H]
    \centering
    \includegraphics[width=0.8\textwidth]
                    {figures/company_structure.pdf}
    \caption{Organizational structure of the company}
    \label{fig:company-structure}
\end{figure}
\end{verbatim}

Figure captions shall be placed below the figure.

Each referenced figure shall have a unique \verb|\label| and shall be
referenced in the text using, for example:

\begin{verbatim}
Figure~\ref{fig:company-structure}
\end{verbatim}

Students shall not manually type figure numbers.

Figures reproduced or adapted from external sources shall include an
appropriate citation.

Photographs taken during the internship shall be relevant and sufficiently
clear. Students shall not include confidential, restricted, or proprietary
information without permission from the internship organization.

Vector graphics or high-resolution images are recommended whenever
possible.

\vspace{0.3cm}

% ================================================================
% XIII. TABLES
% ================================================================

\noindent
\textbf{XIII.\quad TABLES}

\vspace{0.15cm}

Tables shall be created using suitable \LaTeX{} environments such as
\verb|table|, \verb|tabular|, \verb|tabularx|, or \verb|longtable|.

Table captions shall be placed above the table.

Example:

\begin{verbatim}
\begin{table}[H]
    \centering
    \caption{Internship schedule}
    \label{tab:internship-schedule}

    \begin{tabular}{lll}
        \toprule
        Week & Task & Result \\
        \midrule
        1 & Orientation & Completed \\
        2 & System study & Completed \\
        \bottomrule
    \end{tabular}
\end{table}
\end{verbatim}

Tables shall be referenced using:

\begin{verbatim}
Table~\ref{tab:internship-schedule}
\end{verbatim}

Students shall not manually type table numbers.

Tables should be created directly in \LaTeX{} rather than inserted as
screenshots whenever practical.

\vspace{0.3cm}

% ================================================================
% XIV. EQUATIONS AND TECHNICAL QUANTITIES
% ================================================================

\noindent
\textbf{XIV.\quad EQUATIONS AND TECHNICAL QUANTITIES}

\vspace{0.15cm}

Mathematical expressions shall be typeset directly using \LaTeX{}.
Images or screenshots of equations shall not be used.

Short expressions may be written in inline mathematical mode, for example:

\begin{verbatim}
$R = 0.7\,\Omega$
\end{verbatim}

Important equations shall use a numbered equation environment.

Example:

\begin{verbatim}
\begin{equation}
    P = UI
    \label{eq:power}
\end{equation}
\end{verbatim}

Equations shall be referenced using:

\begin{verbatim}
Equation~\eqref{eq:power}
\end{verbatim}

Equation numbers shall not be manually entered.

SI units shall be used wherever applicable. The \verb|siunitx| package
should be used where appropriate.

Examples:

\begin{verbatim}
\SI{220}{\volt}
\SI{50}{\hertz}
\SI{2.5}{\ampere}
\end{verbatim}

Notation and units shall be consistent throughout the report.

\vspace{0.3cm}

% ================================================================
% XV. SOURCE CODE
% ================================================================

\noindent
\textbf{XV.\quad SOURCE CODE}

\vspace{0.15cm}

If the internship involves programming, PLCs, microcontrollers,
simulation, or software development, relevant source code may be included
in the report.

Source code shall be presented using the \verb|listings| environment or
another environment defined by the official template.

Example:

\begin{verbatim}
\begin{lstlisting}[language=st]
IF Start = TRUE THEN
    Motor := TRUE;
ELSE
    Motor := FALSE;
END_IF;
\end{lstlisting}
\end{verbatim}

Only code necessary for explaining the work should appear in the main
report. Long source-code listings should be moved to an appendix or
submitted as supplementary files if required.

Source code should not normally be inserted as screenshots.

\vspace{0.3cm}

% ================================================================
% XVI. ABSTRACT
% ================================================================

\noindent
\textbf{XVI.\quad ABSTRACT}

\vspace{0.15cm}

The Abstract shall briefly summarize:

\begin{itemize}[leftmargin=0.8cm, itemsep=2pt]
    \item the internship organization;
    \item the objectives of the internship;
    \item the main tasks performed;
    \item the main results, knowledge, and skills obtained.
\end{itemize}

The Abstract should be concise and should not reproduce the full content
of the report chapter by chapter.

\vspace{0.3cm}

% ================================================================
% XVII. REFERENCES AND CITATIONS
% ================================================================

\noindent
\textbf{XVII.\quad REFERENCES AND CITATIONS}

\vspace{0.15cm}

All books, journal papers, conference papers, standards, websites,
datasheets, manuals, technical documents, figures, tables, or other
external sources used in preparing the report shall be properly cited.

References shall be managed through the bibliography system provided with
the official template.

Bibliographic entries shall normally be stored in:

\begin{verbatim}
References.bib
\end{verbatim}

References shall be cited in the report using:

\begin{verbatim}
\cite{reference-key}
\end{verbatim}

Students shall not manually type reference numbers such as ``[1]'',
``[2]'', etc.

Titles of foreign-language publications shall remain in their original
language in the bibliography.

All borrowed ideas, data, figures, tables, equations, methods, or results
shall be properly acknowledged.

\vspace{0.3cm}

% ================================================================
% XVIII. CROSS-REFERENCING
% ================================================================

\noindent
\textbf{XVIII.\quad CROSS-REFERENCING}

\vspace{0.15cm}

Numbered elements shall be cross-referenced using the standard \LaTeX{}
mechanism.

Use the following commands where appropriate:

\begin{verbatim}
\label{...}
\ref{...}
\eqref{...}
\end{verbatim}

For example:

\begin{verbatim}
As shown in Figure~\ref{fig:system}, ...
\end{verbatim}

Students shall not manually type numbers such as ``Figure 2.3'' when
those numbers can be generated automatically.

\vspace{0.3cm}

% ================================================================
% XIX. PAGE NUMBERING
% ================================================================

\noindent
\textbf{XIX.\quad PAGE NUMBERING}

\vspace{0.15cm}

Cover pages and official evaluation forms at the beginning of the report
shall not display page numbers.

The main body of the report shall use Arabic page numbers beginning from
Chapter 1.

Page numbering shall be generated automatically by the official template.
Students shall not manually enter page numbers.

\vspace{0.3cm}

% ================================================================
% XX. PROJECT FILE ORGANIZATION
% ================================================================

\noindent
\textbf{XX.\quad PROJECT FILE ORGANIZATION}

\vspace{0.15cm}

Students should maintain a clear and organized \LaTeX{} project
structure. A recommended structure is:

\begin{verbatim}
InternshipReport/
|
|-- main.tex
|-- References.bib
|
|-- figures/
|   |-- company.pdf
|   |-- system.png
|   `-- ...
|
`-- chapters/
    |-- chapter1.tex
    |-- chapter2.tex
    `-- chapter3.tex
\end{verbatim}

Figures should normally be stored in the \verb|figures| directory
specified by the official template.

File and directory names should preferably contain Latin letters,
numbers, hyphens, or underscores.

Recommended examples include:

\begin{verbatim}
company_structure.pdf
motor_control.png
chapter_01.tex
\end{verbatim}

Spaces and unnecessary special characters in file names should be
avoided where possible.

\vspace{0.3cm}

% ================================================================
% XXI. DOCUMENT COMPILATION AND FINAL CHECK
% ================================================================

\noindent
\textbf{XXI.\quad DOCUMENT COMPILATION AND FINAL CHECK}

\vspace{0.15cm}

Before submission, the complete \LaTeX{} project shall compile
successfully and produce the final PDF report.

The final document shall not contain unresolved references such as:

\begin{verbatim}
??
[?]
\end{verbatim}

Before generating the final PDF, students shall check:

\begin{itemize}[leftmargin=0.8cm, itemsep=2pt]
    \item the Table of Contents;
    \item page numbering;
    \item chapter, section, and subsection numbering;
    \item figure numbering;
    \item table numbering;
    \item equation numbering;
    \item cross-references;
    \item citations;
    \item bibliography;
    \item font display;
    \item mathematical symbols;
    \item image quality;
    \item completeness of all required forms.
\end{itemize}

The document may be compiled using TeX Live, MiKTeX, or another
compatible \LaTeX{} environment. Regardless of the compilation
environment, the resulting PDF shall conform to the official Faculty
template.

\vspace{0.3cm}

% ================================================================
% XXII. CONTENT THAT STUDENTS MAY EDIT
% ================================================================

\noindent
\textbf{XXII.\quad CONTENT THAT STUDENTS MAY EDIT}

\vspace{0.15cm}

Students may edit or replace:

\begin{itemize}[leftmargin=0.8cm, itemsep=2pt]
    \item student name;
    \item class;
    \item internship organization;
    \item supervisor name;
    \item Abstract;
    \item List of Symbols and Abbreviations;
    \item chapter titles and chapter content;
    \item figures and photographs;
    \item tables;
    \item equations;
    \item citations and references.
\end{itemize}

Students shall not modify the main document formatting settings in the
preamble, including the document class, page geometry, line spacing,
heading formatting, table-of-contents formatting, font configuration,
page layout, or cover-page macros unless specifically instructed.

\vspace{0.3cm}

% ================================================================
% XXIII. FINAL SUBMISSION REQUIREMENTS
% ================================================================

\noindent
\textbf{XXIII.\quad FINAL SUBMISSION REQUIREMENTS}

\vspace{0.15cm}

Before submission, the final report shall:

\begin{itemize}[leftmargin=0.8cm, itemsep=2pt]
    \item compile without errors that affect the report content or presentation;
    \item contain all required forms and report sections;
    \item contain correct student, class, internship organization, and supervisor information;
    \item contain an updated Table of Contents, List of Figures, and List of Tables;
    \item contain no unused example content from the template;
    \item contain no placeholder text such as \textit{Title of Chapter 1},
    \textit{Section 1.1}, or other sample entries;
    \item contain all required citations and references;
    \item contain no unresolved cross-references.
\end{itemize}

If the Faculty requires submission of the \LaTeX{} source, the complete
project shall include all \verb|.tex| files, the bibliography database,
figures, and any other files required to reproduce the submitted PDF.

\vspace{0.4cm}

\noindent
\textbf{Important:}
Before submitting the final internship report, students shall remove or
comment out the following command in the main document:

\begin{verbatim}
\input{Regulations.tex}
\end{verbatim}

For example:

\begin{verbatim}
% \input{Regulations.tex}
\end{verbatim}

This removes the guidance pages from the final submitted report while
preserving the regulations file in the \LaTeX{} project for reference.

} % End 13pt content

\clearpage

\end{verbatim}

This removes the guidance pages from the final submitted report while
preserving the regulations file in the \LaTeX{} project for reference.

} % End 13pt content

\clearpage
|.
Students shall remove or comment out this command before submitting the
final internship report.

\vspace{0.4cm}

% ================================================================
% I. GENERAL REQUIREMENTS
% ================================================================

\noindent
\textbf{I.\quad GENERAL REQUIREMENTS}

\vspace{0.15cm}

The Graduation Internship Report shall be prepared using the official
\LaTeX{} template provided by the Faculty. Students shall preserve the
predefined page layout, typography, line spacing, heading styles,
numbering system, figure and table formatting, and bibliography settings
unless otherwise instructed by the Faculty or the supervisor.

The report shall be written in English.

Students are responsible for the academic and technical content of the
report. The official \LaTeX{} template is responsible for standardized
document formatting.

\vspace{0.3cm}

% ================================================================
% II. REPORT STRUCTURE
% ================================================================

\noindent
\textbf{II.\quad REPORT STRUCTURE}

\vspace{0.15cm}

The report shall normally be arranged in the following order:

\begin{enumerate}[label=\arabic*., leftmargin=1.2cm, itemsep=2pt]
    \item Outer Cover;
    \item Comments of the Internship Host Organization;
    \item Supervisor's Evaluation and Graduation Internship Defense Results;
    \item Table of Contents;
    \item List of Symbols and Abbreviations, if applicable;
    \item List of Figures;
    \item List of Tables;
    \item Abstract;
    \item Main Chapters of the Internship Report;
    \item References;
    \item Appendices, if applicable.
\end{enumerate}

The preliminary pages shall not be numbered as chapters. The main report
content shall be organized into consecutively numbered chapters beginning
with Chapter 1.

\vspace{0.3cm}

% ================================================================
% III. PAGE SIZE AND MARGINS
% ================================================================

\noindent
\textbf{III.\quad PAGE SIZE AND MARGINS}

\vspace{0.15cm}

The report shall be formatted on A4 paper with the following margins:

\begin{center}
\begin{tabular}{ll}
    Left margin:   & 3.5 cm \\
    Right margin:  & 2.0 cm \\
    Top margin:    & 2.5 cm \\
    Bottom margin: & 2.5 cm
\end{tabular}
\end{center}

These settings are predefined in the official template and shall not be
modified by students.

\vspace{0.3cm}

% ================================================================
% IV. FONT, FONT SIZE, AND LINE SPACING
% ================================================================

\noindent
\textbf{IV.\quad FONT, FONT SIZE, AND LINE SPACING}

\vspace{0.15cm}

The main text shall use a 13 pt font size with 1.3-line spacing.
A Times-family typeface shall be used in accordance with the official
template.

The text shall be fully justified.

Students shall not use repeated spaces, repeated manual line breaks, or
multiple blank lines to position text. Paragraph and heading spacing shall
be controlled by the template.

\vspace{0.3cm}

% ================================================================
% V. OUTER COVER
% ================================================================

\noindent
\textbf{V.\quad OUTER COVER}

\vspace{0.15cm}

The outer cover shall follow the official format provided in the template
and shall contain the following information:

\begin{itemize}[leftmargin=0.8cm, itemsep=2pt]
    \item Thai Nguyen University of Technology;
    \item Faculty of International Training;
    \item Engineering and Technology Department;
    \item University logo and Faculty logo;
    \item the title \textbf{INTERNSHIP REPORT};
    \item Internship Organization;
    \item Author;
    \item Class;
    \item Supervisor;
    \item place and year of completion.
\end{itemize}

Students shall only replace the required information fields and shall not
modify the cover-page layout unless instructed otherwise.

\vspace{0.3cm}

% ================================================================
% VI. COMMENTS OF THE INTERNSHIP HOST ORGANIZATION
% ================================================================

\noindent
\textbf{VI.\quad COMMENTS OF THE INTERNSHIP HOST ORGANIZATION}

\vspace{0.15cm}

The page entitled \textbf{COMMENTS OF THE INTERNSHIP HOST ORGANIZATION}
is reserved for the internship host organization.

It includes:

\begin{itemize}[leftmargin=0.8cm, itemsep=2pt]
    \item the student's full name;
    \item class;
    \item comments from the host organization;
    \item score;
    \item date;
    \item signature and official seal.
\end{itemize}

Students shall not alter or complete official evaluation fields unless
specifically authorized.

\vspace{0.3cm}

% ================================================================
% VII. SUPERVISOR'S EVALUATION AND DEFENSE RESULTS
% ================================================================

\noindent
\textbf{VII.\quad SUPERVISOR'S EVALUATION AND DEFENSE RESULTS}

\vspace{0.15cm}

The section entitled
\textbf{SUPERVISOR'S EVALUATION AND GRADUATION INTERNSHIP DEFENSE RESULTS}
is reserved for the supervisor and defense committee.

The evaluation may include:

\begin{itemize}[leftmargin=0.8cm, itemsep=2pt]
    \item Internship Progress and Student's Attitude;
    \item Report Content;
    \item Report Presentation;
    \item Other Comments;
    \item Supervisor's Opinion;
    \item Defense Result;
    \item Score.
\end{itemize}

Students shall not enter evaluation scores or defense results before they
are officially determined.

\vspace{0.3cm}

% ================================================================
% VIII. CHAPTERS, SECTIONS, AND SUBSECTIONS
% ================================================================

\noindent
\textbf{VIII.\quad CHAPTERS, SECTIONS, AND SUBSECTIONS}

\vspace{0.15cm}

The report structure shall be generated using the standard \LaTeX{}
commands:

\begin{verbatim}
\chapter{...}
\section{...}
\subsection{...}
\end{verbatim}

Chapter, section, and subsection numbers shall be generated automatically
by \LaTeX{}.

Students shall not manually type chapter or section numbers and shall not
reproduce heading styles using commands such as \verb|\textbf|,
\verb|\fontsize|, or \verb|\centering| instead of the appropriate
structural command.

Excessive subdivision should be avoided. In most cases, the hierarchy
Chapter--Section--Subsection is sufficient.

\vspace{0.3cm}

% ================================================================
% IX. ORGANIZATION OF THE MAIN CONTENT
% ================================================================

\noindent
\textbf{IX.\quad ORGANIZATION OF THE MAIN CONTENT}

\vspace{0.15cm}

The specific content of the report shall depend on the internship
organization and the assigned tasks. A typical structure may include:

\begin{itemize}[leftmargin=0.8cm, itemsep=2pt]
    \item \textbf{Chapter 1. Overview of the Internship Organization};
    \item \textbf{Chapter 2. Internship Activities and Assigned Tasks};
    \item \textbf{Chapter 3. Results, Experience, and Evaluation}.
\end{itemize}

These chapter titles are recommendations only and may be adapted to the
actual internship content.

The report shall reflect work actually performed, observed, analyzed, or
learned during the internship. Unrelated theoretical material should not
be included merely to increase report length.

\vspace{0.3cm}

% ================================================================
% X. TABLE OF CONTENTS
% ================================================================

\noindent
\textbf{X.\quad TABLE OF CONTENTS}

\vspace{0.15cm}

The Table of Contents shall be generated automatically using:

\begin{verbatim}
\tableofcontents
\end{verbatim}

Students shall not create or manually edit page numbers in the Table of
Contents.

The document shall be compiled sufficiently to ensure that all page
numbers and entries are correctly updated.

\vspace{0.3cm}

% ================================================================
% XI. LIST OF SYMBOLS AND ABBREVIATIONS
% ================================================================

\noindent
\textbf{XI.\quad LIST OF SYMBOLS AND ABBREVIATIONS}

\vspace{0.15cm}

Symbols and abbreviations that are used repeatedly in the report shall be
listed in the \textbf{List of Symbols and Abbreviations}.

The list may contain the following columns:

\begin{center}
\begin{tabular}{llll}
    No. & Symbol / Abbreviation & Unit & Definition
\end{tabular}
\end{center}

Only symbols and abbreviations actually used in the final report shall be
retained. Example entries contained in the template shall be removed if
they are not relevant to the student's report.

Each abbreviation should normally be defined when it first appears in
the main text.

\vspace{0.3cm}

% ================================================================
% XII. FIGURES, DIAGRAMS, AND PHOTOGRAPHS
% ================================================================

\noindent
\textbf{XII.\quad FIGURES, DIAGRAMS, AND PHOTOGRAPHS}

\vspace{0.15cm}

Figures shall be clear, legible, and directly relevant to the report
content.

Figures shall be inserted using the standard \verb|figure| environment.

Example:

\begin{verbatim}
\begin{figure}[H]
    \centering
    \includegraphics[width=0.8\textwidth]
                    {figures/company_structure.pdf}
    \caption{Organizational structure of the company}
    \label{fig:company-structure}
\end{figure}
\end{verbatim}

Figure captions shall be placed below the figure.

Each referenced figure shall have a unique \verb|\label| and shall be
referenced in the text using, for example:

\begin{verbatim}
Figure~\ref{fig:company-structure}
\end{verbatim}

Students shall not manually type figure numbers.

Figures reproduced or adapted from external sources shall include an
appropriate citation.

Photographs taken during the internship shall be relevant and sufficiently
clear. Students shall not include confidential, restricted, or proprietary
information without permission from the internship organization.

Vector graphics or high-resolution images are recommended whenever
possible.

\vspace{0.3cm}

% ================================================================
% XIII. TABLES
% ================================================================

\noindent
\textbf{XIII.\quad TABLES}

\vspace{0.15cm}

Tables shall be created using suitable \LaTeX{} environments such as
\verb|table|, \verb|tabular|, \verb|tabularx|, or \verb|longtable|.

Table captions shall be placed above the table.

Example:

\begin{verbatim}
\begin{table}[H]
    \centering
    \caption{Internship schedule}
    \label{tab:internship-schedule}

    \begin{tabular}{lll}
        \toprule
        Week & Task & Result \\
        \midrule
        1 & Orientation & Completed \\
        2 & System study & Completed \\
        \bottomrule
    \end{tabular}
\end{table}
\end{verbatim}

Tables shall be referenced using:

\begin{verbatim}
Table~\ref{tab:internship-schedule}
\end{verbatim}

Students shall not manually type table numbers.

Tables should be created directly in \LaTeX{} rather than inserted as
screenshots whenever practical.

\vspace{0.3cm}

% ================================================================
% XIV. EQUATIONS AND TECHNICAL QUANTITIES
% ================================================================

\noindent
\textbf{XIV.\quad EQUATIONS AND TECHNICAL QUANTITIES}

\vspace{0.15cm}

Mathematical expressions shall be typeset directly using \LaTeX{}.
Images or screenshots of equations shall not be used.

Short expressions may be written in inline mathematical mode, for example:

\begin{verbatim}
$R = 0.7\,\Omega$
\end{verbatim}

Important equations shall use a numbered equation environment.

Example:

\begin{verbatim}
\begin{equation}
    P = UI
    \label{eq:power}
\end{equation}
\end{verbatim}

Equations shall be referenced using:

\begin{verbatim}
Equation~\eqref{eq:power}
\end{verbatim}

Equation numbers shall not be manually entered.

SI units shall be used wherever applicable. The \verb|siunitx| package
should be used where appropriate.

Examples:

\begin{verbatim}
\SI{220}{\volt}
\SI{50}{\hertz}
\SI{2.5}{\ampere}
\end{verbatim}

Notation and units shall be consistent throughout the report.

\vspace{0.3cm}

% ================================================================
% XV. SOURCE CODE
% ================================================================

\noindent
\textbf{XV.\quad SOURCE CODE}

\vspace{0.15cm}

If the internship involves programming, PLCs, microcontrollers,
simulation, or software development, relevant source code may be included
in the report.

Source code shall be presented using the \verb|listings| environment or
another environment defined by the official template.

Example:

\begin{verbatim}
\begin{lstlisting}[language=st]
IF Start = TRUE THEN
    Motor := TRUE;
ELSE
    Motor := FALSE;
END_IF;
\end{lstlisting}
\end{verbatim}

Only code necessary for explaining the work should appear in the main
report. Long source-code listings should be moved to an appendix or
submitted as supplementary files if required.

Source code should not normally be inserted as screenshots.

\vspace{0.3cm}

% ================================================================
% XVI. ABSTRACT
% ================================================================

\noindent
\textbf{XVI.\quad ABSTRACT}

\vspace{0.15cm}

The Abstract shall briefly summarize:

\begin{itemize}[leftmargin=0.8cm, itemsep=2pt]
    \item the internship organization;
    \item the objectives of the internship;
    \item the main tasks performed;
    \item the main results, knowledge, and skills obtained.
\end{itemize}

The Abstract should be concise and should not reproduce the full content
of the report chapter by chapter.

\vspace{0.3cm}

% ================================================================
% XVII. REFERENCES AND CITATIONS
% ================================================================

\noindent
\textbf{XVII.\quad REFERENCES AND CITATIONS}

\vspace{0.15cm}

All books, journal papers, conference papers, standards, websites,
datasheets, manuals, technical documents, figures, tables, or other
external sources used in preparing the report shall be properly cited.

References shall be managed through the bibliography system provided with
the official template.

Bibliographic entries shall normally be stored in:

\begin{verbatim}
References.bib
\end{verbatim}

References shall be cited in the report using:

\begin{verbatim}
\cite{reference-key}
\end{verbatim}

Students shall not manually type reference numbers such as ``[1]'',
``[2]'', etc.

Titles of foreign-language publications shall remain in their original
language in the bibliography.

All borrowed ideas, data, figures, tables, equations, methods, or results
shall be properly acknowledged.

\vspace{0.3cm}

% ================================================================
% XVIII. CROSS-REFERENCING
% ================================================================

\noindent
\textbf{XVIII.\quad CROSS-REFERENCING}

\vspace{0.15cm}

Numbered elements shall be cross-referenced using the standard \LaTeX{}
mechanism.

Use the following commands where appropriate:

\begin{verbatim}
\label{...}
\ref{...}
\eqref{...}
\end{verbatim}

For example:

\begin{verbatim}
As shown in Figure~\ref{fig:system}, ...
\end{verbatim}

Students shall not manually type numbers such as ``Figure 2.3'' when
those numbers can be generated automatically.

\vspace{0.3cm}

% ================================================================
% XIX. PAGE NUMBERING
% ================================================================

\noindent
\textbf{XIX.\quad PAGE NUMBERING}

\vspace{0.15cm}

Cover pages and official evaluation forms at the beginning of the report
shall not display page numbers.

The main body of the report shall use Arabic page numbers beginning from
Chapter 1.

Page numbering shall be generated automatically by the official template.
Students shall not manually enter page numbers.

\vspace{0.3cm}

% ================================================================
% XX. PROJECT FILE ORGANIZATION
% ================================================================

\noindent
\textbf{XX.\quad PROJECT FILE ORGANIZATION}

\vspace{0.15cm}

Students should maintain a clear and organized \LaTeX{} project
structure. A recommended structure is:

\begin{verbatim}
InternshipReport/
|
|-- main.tex
|-- References.bib
|
|-- figures/
|   |-- company.pdf
|   |-- system.png
|   `-- ...
|
`-- chapters/
    |-- chapter1.tex
    |-- chapter2.tex
    `-- chapter3.tex
\end{verbatim}

Figures should normally be stored in the \verb|figures| directory
specified by the official template.

File and directory names should preferably contain Latin letters,
numbers, hyphens, or underscores.

Recommended examples include:

\begin{verbatim}
company_structure.pdf
motor_control.png
chapter_01.tex
\end{verbatim}

Spaces and unnecessary special characters in file names should be
avoided where possible.

\vspace{0.3cm}

% ================================================================
% XXI. DOCUMENT COMPILATION AND FINAL CHECK
% ================================================================

\noindent
\textbf{XXI.\quad DOCUMENT COMPILATION AND FINAL CHECK}

\vspace{0.15cm}

Before submission, the complete \LaTeX{} project shall compile
successfully and produce the final PDF report.

The final document shall not contain unresolved references such as:

\begin{verbatim}
??
[?]
\end{verbatim}

Before generating the final PDF, students shall check:

\begin{itemize}[leftmargin=0.8cm, itemsep=2pt]
    \item the Table of Contents;
    \item page numbering;
    \item chapter, section, and subsection numbering;
    \item figure numbering;
    \item table numbering;
    \item equation numbering;
    \item cross-references;
    \item citations;
    \item bibliography;
    \item font display;
    \item mathematical symbols;
    \item image quality;
    \item completeness of all required forms.
\end{itemize}

The document may be compiled using TeX Live, MiKTeX, or another
compatible \LaTeX{} environment. Regardless of the compilation
environment, the resulting PDF shall conform to the official Faculty
template.

\vspace{0.3cm}

% ================================================================
% XXII. CONTENT THAT STUDENTS MAY EDIT
% ================================================================

\noindent
\textbf{XXII.\quad CONTENT THAT STUDENTS MAY EDIT}

\vspace{0.15cm}

Students may edit or replace:

\begin{itemize}[leftmargin=0.8cm, itemsep=2pt]
    \item student name;
    \item class;
    \item internship organization;
    \item supervisor name;
    \item Abstract;
    \item List of Symbols and Abbreviations;
    \item chapter titles and chapter content;
    \item figures and photographs;
    \item tables;
    \item equations;
    \item citations and references.
\end{itemize}

Students shall not modify the main document formatting settings in the
preamble, including the document class, page geometry, line spacing,
heading formatting, table-of-contents formatting, font configuration,
page layout, or cover-page macros unless specifically instructed.

\vspace{0.3cm}

% ================================================================
% XXIII. FINAL SUBMISSION REQUIREMENTS
% ================================================================

\noindent
\textbf{XXIII.\quad FINAL SUBMISSION REQUIREMENTS}

\vspace{0.15cm}

Before submission, the final report shall:

\begin{itemize}[leftmargin=0.8cm, itemsep=2pt]
    \item compile without errors that affect the report content or presentation;
    \item contain all required forms and report sections;
    \item contain correct student, class, internship organization, and supervisor information;
    \item contain an updated Table of Contents, List of Figures, and List of Tables;
    \item contain no unused example content from the template;
    \item contain no placeholder text such as \textit{Title of Chapter 1},
    \textit{Section 1.1}, or other sample entries;
    \item contain all required citations and references;
    \item contain no unresolved cross-references.
\end{itemize}

If the Faculty requires submission of the \LaTeX{} source, the complete
project shall include all \verb|.tex| files, the bibliography database,
figures, and any other files required to reproduce the submitted PDF.

\vspace{0.4cm}

\noindent
\textbf{Important:}
Before submitting the final internship report, students shall remove or
comment out the following command in the main document:

\begin{verbatim}

% =================================================================
% REGULATIONS FOR PREPARING THE GRADUATION INTERNSHIP REPORT
% USING LATEX
% =================================================================
%
% This file is intended to be included from the main document with:
%
%     
% =================================================================
% REGULATIONS FOR PREPARING THE GRADUATION INTERNSHIP REPORT
% USING LATEX
% =================================================================
%
% This file is intended to be included from the main document with:
%
%     \input{Regulations.tex}
%
% IMPORTANT:
% This section is provided for guidance only. Students shall remove
% or comment out the \input{Regulations.tex} command before submitting
% the final internship report.
%
% =================================================================

\clearpage
\thispagestyle{empty}

\begin{center}
    {\fontsize{16pt}{20pt}\selectfont\bfseries
    SPECIFIC REGULATIONS FOR PREPARING\\
    GRADUATION INTERNSHIP REPORTS USING \LaTeX}
\end{center}

\vspace{0.5cm}

{\fontsize{13pt}{17pt}\selectfont
\setstretch{1.3}
\justifying

\noindent
\textbf{Note:}
This section is provided for guidance only and is included in the main
document using the command \verb|\input{Regulations.tex}|.
Students shall remove or comment out this command before submitting the
final internship report.

\vspace{0.4cm}

% ================================================================
% I. GENERAL REQUIREMENTS
% ================================================================

\noindent
\textbf{I.\quad GENERAL REQUIREMENTS}

\vspace{0.15cm}

The Graduation Internship Report shall be prepared using the official
\LaTeX{} template provided by the Faculty. Students shall preserve the
predefined page layout, typography, line spacing, heading styles,
numbering system, figure and table formatting, and bibliography settings
unless otherwise instructed by the Faculty or the supervisor.

The report shall be written in English.

Students are responsible for the academic and technical content of the
report. The official \LaTeX{} template is responsible for standardized
document formatting.

\vspace{0.3cm}

% ================================================================
% II. REPORT STRUCTURE
% ================================================================

\noindent
\textbf{II.\quad REPORT STRUCTURE}

\vspace{0.15cm}

The report shall normally be arranged in the following order:

\begin{enumerate}[label=\arabic*., leftmargin=1.2cm, itemsep=2pt]
    \item Outer Cover;
    \item Comments of the Internship Host Organization;
    \item Supervisor's Evaluation and Graduation Internship Defense Results;
    \item Table of Contents;
    \item List of Symbols and Abbreviations, if applicable;
    \item List of Figures;
    \item List of Tables;
    \item Abstract;
    \item Main Chapters of the Internship Report;
    \item References;
    \item Appendices, if applicable.
\end{enumerate}

The preliminary pages shall not be numbered as chapters. The main report
content shall be organized into consecutively numbered chapters beginning
with Chapter 1.

\vspace{0.3cm}

% ================================================================
% III. PAGE SIZE AND MARGINS
% ================================================================

\noindent
\textbf{III.\quad PAGE SIZE AND MARGINS}

\vspace{0.15cm}

The report shall be formatted on A4 paper with the following margins:

\begin{center}
\begin{tabular}{ll}
    Left margin:   & 3.5 cm \\
    Right margin:  & 2.0 cm \\
    Top margin:    & 2.5 cm \\
    Bottom margin: & 2.5 cm
\end{tabular}
\end{center}

These settings are predefined in the official template and shall not be
modified by students.

\vspace{0.3cm}

% ================================================================
% IV. FONT, FONT SIZE, AND LINE SPACING
% ================================================================

\noindent
\textbf{IV.\quad FONT, FONT SIZE, AND LINE SPACING}

\vspace{0.15cm}

The main text shall use a 13 pt font size with 1.3-line spacing.
A Times-family typeface shall be used in accordance with the official
template.

The text shall be fully justified.

Students shall not use repeated spaces, repeated manual line breaks, or
multiple blank lines to position text. Paragraph and heading spacing shall
be controlled by the template.

\vspace{0.3cm}

% ================================================================
% V. OUTER COVER
% ================================================================

\noindent
\textbf{V.\quad OUTER COVER}

\vspace{0.15cm}

The outer cover shall follow the official format provided in the template
and shall contain the following information:

\begin{itemize}[leftmargin=0.8cm, itemsep=2pt]
    \item Thai Nguyen University of Technology;
    \item Faculty of International Training;
    \item Engineering and Technology Department;
    \item University logo and Faculty logo;
    \item the title \textbf{INTERNSHIP REPORT};
    \item Internship Organization;
    \item Author;
    \item Class;
    \item Supervisor;
    \item place and year of completion.
\end{itemize}

Students shall only replace the required information fields and shall not
modify the cover-page layout unless instructed otherwise.

\vspace{0.3cm}

% ================================================================
% VI. COMMENTS OF THE INTERNSHIP HOST ORGANIZATION
% ================================================================

\noindent
\textbf{VI.\quad COMMENTS OF THE INTERNSHIP HOST ORGANIZATION}

\vspace{0.15cm}

The page entitled \textbf{COMMENTS OF THE INTERNSHIP HOST ORGANIZATION}
is reserved for the internship host organization.

It includes:

\begin{itemize}[leftmargin=0.8cm, itemsep=2pt]
    \item the student's full name;
    \item class;
    \item comments from the host organization;
    \item score;
    \item date;
    \item signature and official seal.
\end{itemize}

Students shall not alter or complete official evaluation fields unless
specifically authorized.

\vspace{0.3cm}

% ================================================================
% VII. SUPERVISOR'S EVALUATION AND DEFENSE RESULTS
% ================================================================

\noindent
\textbf{VII.\quad SUPERVISOR'S EVALUATION AND DEFENSE RESULTS}

\vspace{0.15cm}

The section entitled
\textbf{SUPERVISOR'S EVALUATION AND GRADUATION INTERNSHIP DEFENSE RESULTS}
is reserved for the supervisor and defense committee.

The evaluation may include:

\begin{itemize}[leftmargin=0.8cm, itemsep=2pt]
    \item Internship Progress and Student's Attitude;
    \item Report Content;
    \item Report Presentation;
    \item Other Comments;
    \item Supervisor's Opinion;
    \item Defense Result;
    \item Score.
\end{itemize}

Students shall not enter evaluation scores or defense results before they
are officially determined.

\vspace{0.3cm}

% ================================================================
% VIII. CHAPTERS, SECTIONS, AND SUBSECTIONS
% ================================================================

\noindent
\textbf{VIII.\quad CHAPTERS, SECTIONS, AND SUBSECTIONS}

\vspace{0.15cm}

The report structure shall be generated using the standard \LaTeX{}
commands:

\begin{verbatim}
\chapter{...}
\section{...}
\subsection{...}
\end{verbatim}

Chapter, section, and subsection numbers shall be generated automatically
by \LaTeX{}.

Students shall not manually type chapter or section numbers and shall not
reproduce heading styles using commands such as \verb|\textbf|,
\verb|\fontsize|, or \verb|\centering| instead of the appropriate
structural command.

Excessive subdivision should be avoided. In most cases, the hierarchy
Chapter--Section--Subsection is sufficient.

\vspace{0.3cm}

% ================================================================
% IX. ORGANIZATION OF THE MAIN CONTENT
% ================================================================

\noindent
\textbf{IX.\quad ORGANIZATION OF THE MAIN CONTENT}

\vspace{0.15cm}

The specific content of the report shall depend on the internship
organization and the assigned tasks. A typical structure may include:

\begin{itemize}[leftmargin=0.8cm, itemsep=2pt]
    \item \textbf{Chapter 1. Overview of the Internship Organization};
    \item \textbf{Chapter 2. Internship Activities and Assigned Tasks};
    \item \textbf{Chapter 3. Results, Experience, and Evaluation}.
\end{itemize}

These chapter titles are recommendations only and may be adapted to the
actual internship content.

The report shall reflect work actually performed, observed, analyzed, or
learned during the internship. Unrelated theoretical material should not
be included merely to increase report length.

\vspace{0.3cm}

% ================================================================
% X. TABLE OF CONTENTS
% ================================================================

\noindent
\textbf{X.\quad TABLE OF CONTENTS}

\vspace{0.15cm}

The Table of Contents shall be generated automatically using:

\begin{verbatim}
\tableofcontents
\end{verbatim}

Students shall not create or manually edit page numbers in the Table of
Contents.

The document shall be compiled sufficiently to ensure that all page
numbers and entries are correctly updated.

\vspace{0.3cm}

% ================================================================
% XI. LIST OF SYMBOLS AND ABBREVIATIONS
% ================================================================

\noindent
\textbf{XI.\quad LIST OF SYMBOLS AND ABBREVIATIONS}

\vspace{0.15cm}

Symbols and abbreviations that are used repeatedly in the report shall be
listed in the \textbf{List of Symbols and Abbreviations}.

The list may contain the following columns:

\begin{center}
\begin{tabular}{llll}
    No. & Symbol / Abbreviation & Unit & Definition
\end{tabular}
\end{center}

Only symbols and abbreviations actually used in the final report shall be
retained. Example entries contained in the template shall be removed if
they are not relevant to the student's report.

Each abbreviation should normally be defined when it first appears in
the main text.

\vspace{0.3cm}

% ================================================================
% XII. FIGURES, DIAGRAMS, AND PHOTOGRAPHS
% ================================================================

\noindent
\textbf{XII.\quad FIGURES, DIAGRAMS, AND PHOTOGRAPHS}

\vspace{0.15cm}

Figures shall be clear, legible, and directly relevant to the report
content.

Figures shall be inserted using the standard \verb|figure| environment.

Example:

\begin{verbatim}
\begin{figure}[H]
    \centering
    \includegraphics[width=0.8\textwidth]
                    {figures/company_structure.pdf}
    \caption{Organizational structure of the company}
    \label{fig:company-structure}
\end{figure}
\end{verbatim}

Figure captions shall be placed below the figure.

Each referenced figure shall have a unique \verb|\label| and shall be
referenced in the text using, for example:

\begin{verbatim}
Figure~\ref{fig:company-structure}
\end{verbatim}

Students shall not manually type figure numbers.

Figures reproduced or adapted from external sources shall include an
appropriate citation.

Photographs taken during the internship shall be relevant and sufficiently
clear. Students shall not include confidential, restricted, or proprietary
information without permission from the internship organization.

Vector graphics or high-resolution images are recommended whenever
possible.

\vspace{0.3cm}

% ================================================================
% XIII. TABLES
% ================================================================

\noindent
\textbf{XIII.\quad TABLES}

\vspace{0.15cm}

Tables shall be created using suitable \LaTeX{} environments such as
\verb|table|, \verb|tabular|, \verb|tabularx|, or \verb|longtable|.

Table captions shall be placed above the table.

Example:

\begin{verbatim}
\begin{table}[H]
    \centering
    \caption{Internship schedule}
    \label{tab:internship-schedule}

    \begin{tabular}{lll}
        \toprule
        Week & Task & Result \\
        \midrule
        1 & Orientation & Completed \\
        2 & System study & Completed \\
        \bottomrule
    \end{tabular}
\end{table}
\end{verbatim}

Tables shall be referenced using:

\begin{verbatim}
Table~\ref{tab:internship-schedule}
\end{verbatim}

Students shall not manually type table numbers.

Tables should be created directly in \LaTeX{} rather than inserted as
screenshots whenever practical.

\vspace{0.3cm}

% ================================================================
% XIV. EQUATIONS AND TECHNICAL QUANTITIES
% ================================================================

\noindent
\textbf{XIV.\quad EQUATIONS AND TECHNICAL QUANTITIES}

\vspace{0.15cm}

Mathematical expressions shall be typeset directly using \LaTeX{}.
Images or screenshots of equations shall not be used.

Short expressions may be written in inline mathematical mode, for example:

\begin{verbatim}
$R = 0.7\,\Omega$
\end{verbatim}

Important equations shall use a numbered equation environment.

Example:

\begin{verbatim}
\begin{equation}
    P = UI
    \label{eq:power}
\end{equation}
\end{verbatim}

Equations shall be referenced using:

\begin{verbatim}
Equation~\eqref{eq:power}
\end{verbatim}

Equation numbers shall not be manually entered.

SI units shall be used wherever applicable. The \verb|siunitx| package
should be used where appropriate.

Examples:

\begin{verbatim}
\SI{220}{\volt}
\SI{50}{\hertz}
\SI{2.5}{\ampere}
\end{verbatim}

Notation and units shall be consistent throughout the report.

\vspace{0.3cm}

% ================================================================
% XV. SOURCE CODE
% ================================================================

\noindent
\textbf{XV.\quad SOURCE CODE}

\vspace{0.15cm}

If the internship involves programming, PLCs, microcontrollers,
simulation, or software development, relevant source code may be included
in the report.

Source code shall be presented using the \verb|listings| environment or
another environment defined by the official template.

Example:

\begin{verbatim}
\begin{lstlisting}[language=st]
IF Start = TRUE THEN
    Motor := TRUE;
ELSE
    Motor := FALSE;
END_IF;
\end{lstlisting}
\end{verbatim}

Only code necessary for explaining the work should appear in the main
report. Long source-code listings should be moved to an appendix or
submitted as supplementary files if required.

Source code should not normally be inserted as screenshots.

\vspace{0.3cm}

% ================================================================
% XVI. ABSTRACT
% ================================================================

\noindent
\textbf{XVI.\quad ABSTRACT}

\vspace{0.15cm}

The Abstract shall briefly summarize:

\begin{itemize}[leftmargin=0.8cm, itemsep=2pt]
    \item the internship organization;
    \item the objectives of the internship;
    \item the main tasks performed;
    \item the main results, knowledge, and skills obtained.
\end{itemize}

The Abstract should be concise and should not reproduce the full content
of the report chapter by chapter.

\vspace{0.3cm}

% ================================================================
% XVII. REFERENCES AND CITATIONS
% ================================================================

\noindent
\textbf{XVII.\quad REFERENCES AND CITATIONS}

\vspace{0.15cm}

All books, journal papers, conference papers, standards, websites,
datasheets, manuals, technical documents, figures, tables, or other
external sources used in preparing the report shall be properly cited.

References shall be managed through the bibliography system provided with
the official template.

Bibliographic entries shall normally be stored in:

\begin{verbatim}
References.bib
\end{verbatim}

References shall be cited in the report using:

\begin{verbatim}
\cite{reference-key}
\end{verbatim}

Students shall not manually type reference numbers such as ``[1]'',
``[2]'', etc.

Titles of foreign-language publications shall remain in their original
language in the bibliography.

All borrowed ideas, data, figures, tables, equations, methods, or results
shall be properly acknowledged.

\vspace{0.3cm}

% ================================================================
% XVIII. CROSS-REFERENCING
% ================================================================

\noindent
\textbf{XVIII.\quad CROSS-REFERENCING}

\vspace{0.15cm}

Numbered elements shall be cross-referenced using the standard \LaTeX{}
mechanism.

Use the following commands where appropriate:

\begin{verbatim}
\label{...}
\ref{...}
\eqref{...}
\end{verbatim}

For example:

\begin{verbatim}
As shown in Figure~\ref{fig:system}, ...
\end{verbatim}

Students shall not manually type numbers such as ``Figure 2.3'' when
those numbers can be generated automatically.

\vspace{0.3cm}

% ================================================================
% XIX. PAGE NUMBERING
% ================================================================

\noindent
\textbf{XIX.\quad PAGE NUMBERING}

\vspace{0.15cm}

Cover pages and official evaluation forms at the beginning of the report
shall not display page numbers.

The main body of the report shall use Arabic page numbers beginning from
Chapter 1.

Page numbering shall be generated automatically by the official template.
Students shall not manually enter page numbers.

\vspace{0.3cm}

% ================================================================
% XX. PROJECT FILE ORGANIZATION
% ================================================================

\noindent
\textbf{XX.\quad PROJECT FILE ORGANIZATION}

\vspace{0.15cm}

Students should maintain a clear and organized \LaTeX{} project
structure. A recommended structure is:

\begin{verbatim}
InternshipReport/
|
|-- main.tex
|-- References.bib
|
|-- figures/
|   |-- company.pdf
|   |-- system.png
|   `-- ...
|
`-- chapters/
    |-- chapter1.tex
    |-- chapter2.tex
    `-- chapter3.tex
\end{verbatim}

Figures should normally be stored in the \verb|figures| directory
specified by the official template.

File and directory names should preferably contain Latin letters,
numbers, hyphens, or underscores.

Recommended examples include:

\begin{verbatim}
company_structure.pdf
motor_control.png
chapter_01.tex
\end{verbatim}

Spaces and unnecessary special characters in file names should be
avoided where possible.

\vspace{0.3cm}

% ================================================================
% XXI. DOCUMENT COMPILATION AND FINAL CHECK
% ================================================================

\noindent
\textbf{XXI.\quad DOCUMENT COMPILATION AND FINAL CHECK}

\vspace{0.15cm}

Before submission, the complete \LaTeX{} project shall compile
successfully and produce the final PDF report.

The final document shall not contain unresolved references such as:

\begin{verbatim}
??
[?]
\end{verbatim}

Before generating the final PDF, students shall check:

\begin{itemize}[leftmargin=0.8cm, itemsep=2pt]
    \item the Table of Contents;
    \item page numbering;
    \item chapter, section, and subsection numbering;
    \item figure numbering;
    \item table numbering;
    \item equation numbering;
    \item cross-references;
    \item citations;
    \item bibliography;
    \item font display;
    \item mathematical symbols;
    \item image quality;
    \item completeness of all required forms.
\end{itemize}

The document may be compiled using TeX Live, MiKTeX, or another
compatible \LaTeX{} environment. Regardless of the compilation
environment, the resulting PDF shall conform to the official Faculty
template.

\vspace{0.3cm}

% ================================================================
% XXII. CONTENT THAT STUDENTS MAY EDIT
% ================================================================

\noindent
\textbf{XXII.\quad CONTENT THAT STUDENTS MAY EDIT}

\vspace{0.15cm}

Students may edit or replace:

\begin{itemize}[leftmargin=0.8cm, itemsep=2pt]
    \item student name;
    \item class;
    \item internship organization;
    \item supervisor name;
    \item Abstract;
    \item List of Symbols and Abbreviations;
    \item chapter titles and chapter content;
    \item figures and photographs;
    \item tables;
    \item equations;
    \item citations and references.
\end{itemize}

Students shall not modify the main document formatting settings in the
preamble, including the document class, page geometry, line spacing,
heading formatting, table-of-contents formatting, font configuration,
page layout, or cover-page macros unless specifically instructed.

\vspace{0.3cm}

% ================================================================
% XXIII. FINAL SUBMISSION REQUIREMENTS
% ================================================================

\noindent
\textbf{XXIII.\quad FINAL SUBMISSION REQUIREMENTS}

\vspace{0.15cm}

Before submission, the final report shall:

\begin{itemize}[leftmargin=0.8cm, itemsep=2pt]
    \item compile without errors that affect the report content or presentation;
    \item contain all required forms and report sections;
    \item contain correct student, class, internship organization, and supervisor information;
    \item contain an updated Table of Contents, List of Figures, and List of Tables;
    \item contain no unused example content from the template;
    \item contain no placeholder text such as \textit{Title of Chapter 1},
    \textit{Section 1.1}, or other sample entries;
    \item contain all required citations and references;
    \item contain no unresolved cross-references.
\end{itemize}

If the Faculty requires submission of the \LaTeX{} source, the complete
project shall include all \verb|.tex| files, the bibliography database,
figures, and any other files required to reproduce the submitted PDF.

\vspace{0.4cm}

\noindent
\textbf{Important:}
Before submitting the final internship report, students shall remove or
comment out the following command in the main document:

\begin{verbatim}
\input{Regulations.tex}
\end{verbatim}

For example:

\begin{verbatim}
% \input{Regulations.tex}
\end{verbatim}

This removes the guidance pages from the final submitted report while
preserving the regulations file in the \LaTeX{} project for reference.

} % End 13pt content

\clearpage

%
% IMPORTANT:
% This section is provided for guidance only. Students shall remove
% or comment out the 
% =================================================================
% REGULATIONS FOR PREPARING THE GRADUATION INTERNSHIP REPORT
% USING LATEX
% =================================================================
%
% This file is intended to be included from the main document with:
%
%     \input{Regulations.tex}
%
% IMPORTANT:
% This section is provided for guidance only. Students shall remove
% or comment out the \input{Regulations.tex} command before submitting
% the final internship report.
%
% =================================================================

\clearpage
\thispagestyle{empty}

\begin{center}
    {\fontsize{16pt}{20pt}\selectfont\bfseries
    SPECIFIC REGULATIONS FOR PREPARING\\
    GRADUATION INTERNSHIP REPORTS USING \LaTeX}
\end{center}

\vspace{0.5cm}

{\fontsize{13pt}{17pt}\selectfont
\setstretch{1.3}
\justifying

\noindent
\textbf{Note:}
This section is provided for guidance only and is included in the main
document using the command \verb|\input{Regulations.tex}|.
Students shall remove or comment out this command before submitting the
final internship report.

\vspace{0.4cm}

% ================================================================
% I. GENERAL REQUIREMENTS
% ================================================================

\noindent
\textbf{I.\quad GENERAL REQUIREMENTS}

\vspace{0.15cm}

The Graduation Internship Report shall be prepared using the official
\LaTeX{} template provided by the Faculty. Students shall preserve the
predefined page layout, typography, line spacing, heading styles,
numbering system, figure and table formatting, and bibliography settings
unless otherwise instructed by the Faculty or the supervisor.

The report shall be written in English.

Students are responsible for the academic and technical content of the
report. The official \LaTeX{} template is responsible for standardized
document formatting.

\vspace{0.3cm}

% ================================================================
% II. REPORT STRUCTURE
% ================================================================

\noindent
\textbf{II.\quad REPORT STRUCTURE}

\vspace{0.15cm}

The report shall normally be arranged in the following order:

\begin{enumerate}[label=\arabic*., leftmargin=1.2cm, itemsep=2pt]
    \item Outer Cover;
    \item Comments of the Internship Host Organization;
    \item Supervisor's Evaluation and Graduation Internship Defense Results;
    \item Table of Contents;
    \item List of Symbols and Abbreviations, if applicable;
    \item List of Figures;
    \item List of Tables;
    \item Abstract;
    \item Main Chapters of the Internship Report;
    \item References;
    \item Appendices, if applicable.
\end{enumerate}

The preliminary pages shall not be numbered as chapters. The main report
content shall be organized into consecutively numbered chapters beginning
with Chapter 1.

\vspace{0.3cm}

% ================================================================
% III. PAGE SIZE AND MARGINS
% ================================================================

\noindent
\textbf{III.\quad PAGE SIZE AND MARGINS}

\vspace{0.15cm}

The report shall be formatted on A4 paper with the following margins:

\begin{center}
\begin{tabular}{ll}
    Left margin:   & 3.5 cm \\
    Right margin:  & 2.0 cm \\
    Top margin:    & 2.5 cm \\
    Bottom margin: & 2.5 cm
\end{tabular}
\end{center}

These settings are predefined in the official template and shall not be
modified by students.

\vspace{0.3cm}

% ================================================================
% IV. FONT, FONT SIZE, AND LINE SPACING
% ================================================================

\noindent
\textbf{IV.\quad FONT, FONT SIZE, AND LINE SPACING}

\vspace{0.15cm}

The main text shall use a 13 pt font size with 1.3-line spacing.
A Times-family typeface shall be used in accordance with the official
template.

The text shall be fully justified.

Students shall not use repeated spaces, repeated manual line breaks, or
multiple blank lines to position text. Paragraph and heading spacing shall
be controlled by the template.

\vspace{0.3cm}

% ================================================================
% V. OUTER COVER
% ================================================================

\noindent
\textbf{V.\quad OUTER COVER}

\vspace{0.15cm}

The outer cover shall follow the official format provided in the template
and shall contain the following information:

\begin{itemize}[leftmargin=0.8cm, itemsep=2pt]
    \item Thai Nguyen University of Technology;
    \item Faculty of International Training;
    \item Engineering and Technology Department;
    \item University logo and Faculty logo;
    \item the title \textbf{INTERNSHIP REPORT};
    \item Internship Organization;
    \item Author;
    \item Class;
    \item Supervisor;
    \item place and year of completion.
\end{itemize}

Students shall only replace the required information fields and shall not
modify the cover-page layout unless instructed otherwise.

\vspace{0.3cm}

% ================================================================
% VI. COMMENTS OF THE INTERNSHIP HOST ORGANIZATION
% ================================================================

\noindent
\textbf{VI.\quad COMMENTS OF THE INTERNSHIP HOST ORGANIZATION}

\vspace{0.15cm}

The page entitled \textbf{COMMENTS OF THE INTERNSHIP HOST ORGANIZATION}
is reserved for the internship host organization.

It includes:

\begin{itemize}[leftmargin=0.8cm, itemsep=2pt]
    \item the student's full name;
    \item class;
    \item comments from the host organization;
    \item score;
    \item date;
    \item signature and official seal.
\end{itemize}

Students shall not alter or complete official evaluation fields unless
specifically authorized.

\vspace{0.3cm}

% ================================================================
% VII. SUPERVISOR'S EVALUATION AND DEFENSE RESULTS
% ================================================================

\noindent
\textbf{VII.\quad SUPERVISOR'S EVALUATION AND DEFENSE RESULTS}

\vspace{0.15cm}

The section entitled
\textbf{SUPERVISOR'S EVALUATION AND GRADUATION INTERNSHIP DEFENSE RESULTS}
is reserved for the supervisor and defense committee.

The evaluation may include:

\begin{itemize}[leftmargin=0.8cm, itemsep=2pt]
    \item Internship Progress and Student's Attitude;
    \item Report Content;
    \item Report Presentation;
    \item Other Comments;
    \item Supervisor's Opinion;
    \item Defense Result;
    \item Score.
\end{itemize}

Students shall not enter evaluation scores or defense results before they
are officially determined.

\vspace{0.3cm}

% ================================================================
% VIII. CHAPTERS, SECTIONS, AND SUBSECTIONS
% ================================================================

\noindent
\textbf{VIII.\quad CHAPTERS, SECTIONS, AND SUBSECTIONS}

\vspace{0.15cm}

The report structure shall be generated using the standard \LaTeX{}
commands:

\begin{verbatim}
\chapter{...}
\section{...}
\subsection{...}
\end{verbatim}

Chapter, section, and subsection numbers shall be generated automatically
by \LaTeX{}.

Students shall not manually type chapter or section numbers and shall not
reproduce heading styles using commands such as \verb|\textbf|,
\verb|\fontsize|, or \verb|\centering| instead of the appropriate
structural command.

Excessive subdivision should be avoided. In most cases, the hierarchy
Chapter--Section--Subsection is sufficient.

\vspace{0.3cm}

% ================================================================
% IX. ORGANIZATION OF THE MAIN CONTENT
% ================================================================

\noindent
\textbf{IX.\quad ORGANIZATION OF THE MAIN CONTENT}

\vspace{0.15cm}

The specific content of the report shall depend on the internship
organization and the assigned tasks. A typical structure may include:

\begin{itemize}[leftmargin=0.8cm, itemsep=2pt]
    \item \textbf{Chapter 1. Overview of the Internship Organization};
    \item \textbf{Chapter 2. Internship Activities and Assigned Tasks};
    \item \textbf{Chapter 3. Results, Experience, and Evaluation}.
\end{itemize}

These chapter titles are recommendations only and may be adapted to the
actual internship content.

The report shall reflect work actually performed, observed, analyzed, or
learned during the internship. Unrelated theoretical material should not
be included merely to increase report length.

\vspace{0.3cm}

% ================================================================
% X. TABLE OF CONTENTS
% ================================================================

\noindent
\textbf{X.\quad TABLE OF CONTENTS}

\vspace{0.15cm}

The Table of Contents shall be generated automatically using:

\begin{verbatim}
\tableofcontents
\end{verbatim}

Students shall not create or manually edit page numbers in the Table of
Contents.

The document shall be compiled sufficiently to ensure that all page
numbers and entries are correctly updated.

\vspace{0.3cm}

% ================================================================
% XI. LIST OF SYMBOLS AND ABBREVIATIONS
% ================================================================

\noindent
\textbf{XI.\quad LIST OF SYMBOLS AND ABBREVIATIONS}

\vspace{0.15cm}

Symbols and abbreviations that are used repeatedly in the report shall be
listed in the \textbf{List of Symbols and Abbreviations}.

The list may contain the following columns:

\begin{center}
\begin{tabular}{llll}
    No. & Symbol / Abbreviation & Unit & Definition
\end{tabular}
\end{center}

Only symbols and abbreviations actually used in the final report shall be
retained. Example entries contained in the template shall be removed if
they are not relevant to the student's report.

Each abbreviation should normally be defined when it first appears in
the main text.

\vspace{0.3cm}

% ================================================================
% XII. FIGURES, DIAGRAMS, AND PHOTOGRAPHS
% ================================================================

\noindent
\textbf{XII.\quad FIGURES, DIAGRAMS, AND PHOTOGRAPHS}

\vspace{0.15cm}

Figures shall be clear, legible, and directly relevant to the report
content.

Figures shall be inserted using the standard \verb|figure| environment.

Example:

\begin{verbatim}
\begin{figure}[H]
    \centering
    \includegraphics[width=0.8\textwidth]
                    {figures/company_structure.pdf}
    \caption{Organizational structure of the company}
    \label{fig:company-structure}
\end{figure}
\end{verbatim}

Figure captions shall be placed below the figure.

Each referenced figure shall have a unique \verb|\label| and shall be
referenced in the text using, for example:

\begin{verbatim}
Figure~\ref{fig:company-structure}
\end{verbatim}

Students shall not manually type figure numbers.

Figures reproduced or adapted from external sources shall include an
appropriate citation.

Photographs taken during the internship shall be relevant and sufficiently
clear. Students shall not include confidential, restricted, or proprietary
information without permission from the internship organization.

Vector graphics or high-resolution images are recommended whenever
possible.

\vspace{0.3cm}

% ================================================================
% XIII. TABLES
% ================================================================

\noindent
\textbf{XIII.\quad TABLES}

\vspace{0.15cm}

Tables shall be created using suitable \LaTeX{} environments such as
\verb|table|, \verb|tabular|, \verb|tabularx|, or \verb|longtable|.

Table captions shall be placed above the table.

Example:

\begin{verbatim}
\begin{table}[H]
    \centering
    \caption{Internship schedule}
    \label{tab:internship-schedule}

    \begin{tabular}{lll}
        \toprule
        Week & Task & Result \\
        \midrule
        1 & Orientation & Completed \\
        2 & System study & Completed \\
        \bottomrule
    \end{tabular}
\end{table}
\end{verbatim}

Tables shall be referenced using:

\begin{verbatim}
Table~\ref{tab:internship-schedule}
\end{verbatim}

Students shall not manually type table numbers.

Tables should be created directly in \LaTeX{} rather than inserted as
screenshots whenever practical.

\vspace{0.3cm}

% ================================================================
% XIV. EQUATIONS AND TECHNICAL QUANTITIES
% ================================================================

\noindent
\textbf{XIV.\quad EQUATIONS AND TECHNICAL QUANTITIES}

\vspace{0.15cm}

Mathematical expressions shall be typeset directly using \LaTeX{}.
Images or screenshots of equations shall not be used.

Short expressions may be written in inline mathematical mode, for example:

\begin{verbatim}
$R = 0.7\,\Omega$
\end{verbatim}

Important equations shall use a numbered equation environment.

Example:

\begin{verbatim}
\begin{equation}
    P = UI
    \label{eq:power}
\end{equation}
\end{verbatim}

Equations shall be referenced using:

\begin{verbatim}
Equation~\eqref{eq:power}
\end{verbatim}

Equation numbers shall not be manually entered.

SI units shall be used wherever applicable. The \verb|siunitx| package
should be used where appropriate.

Examples:

\begin{verbatim}
\SI{220}{\volt}
\SI{50}{\hertz}
\SI{2.5}{\ampere}
\end{verbatim}

Notation and units shall be consistent throughout the report.

\vspace{0.3cm}

% ================================================================
% XV. SOURCE CODE
% ================================================================

\noindent
\textbf{XV.\quad SOURCE CODE}

\vspace{0.15cm}

If the internship involves programming, PLCs, microcontrollers,
simulation, or software development, relevant source code may be included
in the report.

Source code shall be presented using the \verb|listings| environment or
another environment defined by the official template.

Example:

\begin{verbatim}
\begin{lstlisting}[language=st]
IF Start = TRUE THEN
    Motor := TRUE;
ELSE
    Motor := FALSE;
END_IF;
\end{lstlisting}
\end{verbatim}

Only code necessary for explaining the work should appear in the main
report. Long source-code listings should be moved to an appendix or
submitted as supplementary files if required.

Source code should not normally be inserted as screenshots.

\vspace{0.3cm}

% ================================================================
% XVI. ABSTRACT
% ================================================================

\noindent
\textbf{XVI.\quad ABSTRACT}

\vspace{0.15cm}

The Abstract shall briefly summarize:

\begin{itemize}[leftmargin=0.8cm, itemsep=2pt]
    \item the internship organization;
    \item the objectives of the internship;
    \item the main tasks performed;
    \item the main results, knowledge, and skills obtained.
\end{itemize}

The Abstract should be concise and should not reproduce the full content
of the report chapter by chapter.

\vspace{0.3cm}

% ================================================================
% XVII. REFERENCES AND CITATIONS
% ================================================================

\noindent
\textbf{XVII.\quad REFERENCES AND CITATIONS}

\vspace{0.15cm}

All books, journal papers, conference papers, standards, websites,
datasheets, manuals, technical documents, figures, tables, or other
external sources used in preparing the report shall be properly cited.

References shall be managed through the bibliography system provided with
the official template.

Bibliographic entries shall normally be stored in:

\begin{verbatim}
References.bib
\end{verbatim}

References shall be cited in the report using:

\begin{verbatim}
\cite{reference-key}
\end{verbatim}

Students shall not manually type reference numbers such as ``[1]'',
``[2]'', etc.

Titles of foreign-language publications shall remain in their original
language in the bibliography.

All borrowed ideas, data, figures, tables, equations, methods, or results
shall be properly acknowledged.

\vspace{0.3cm}

% ================================================================
% XVIII. CROSS-REFERENCING
% ================================================================

\noindent
\textbf{XVIII.\quad CROSS-REFERENCING}

\vspace{0.15cm}

Numbered elements shall be cross-referenced using the standard \LaTeX{}
mechanism.

Use the following commands where appropriate:

\begin{verbatim}
\label{...}
\ref{...}
\eqref{...}
\end{verbatim}

For example:

\begin{verbatim}
As shown in Figure~\ref{fig:system}, ...
\end{verbatim}

Students shall not manually type numbers such as ``Figure 2.3'' when
those numbers can be generated automatically.

\vspace{0.3cm}

% ================================================================
% XIX. PAGE NUMBERING
% ================================================================

\noindent
\textbf{XIX.\quad PAGE NUMBERING}

\vspace{0.15cm}

Cover pages and official evaluation forms at the beginning of the report
shall not display page numbers.

The main body of the report shall use Arabic page numbers beginning from
Chapter 1.

Page numbering shall be generated automatically by the official template.
Students shall not manually enter page numbers.

\vspace{0.3cm}

% ================================================================
% XX. PROJECT FILE ORGANIZATION
% ================================================================

\noindent
\textbf{XX.\quad PROJECT FILE ORGANIZATION}

\vspace{0.15cm}

Students should maintain a clear and organized \LaTeX{} project
structure. A recommended structure is:

\begin{verbatim}
InternshipReport/
|
|-- main.tex
|-- References.bib
|
|-- figures/
|   |-- company.pdf
|   |-- system.png
|   `-- ...
|
`-- chapters/
    |-- chapter1.tex
    |-- chapter2.tex
    `-- chapter3.tex
\end{verbatim}

Figures should normally be stored in the \verb|figures| directory
specified by the official template.

File and directory names should preferably contain Latin letters,
numbers, hyphens, or underscores.

Recommended examples include:

\begin{verbatim}
company_structure.pdf
motor_control.png
chapter_01.tex
\end{verbatim}

Spaces and unnecessary special characters in file names should be
avoided where possible.

\vspace{0.3cm}

% ================================================================
% XXI. DOCUMENT COMPILATION AND FINAL CHECK
% ================================================================

\noindent
\textbf{XXI.\quad DOCUMENT COMPILATION AND FINAL CHECK}

\vspace{0.15cm}

Before submission, the complete \LaTeX{} project shall compile
successfully and produce the final PDF report.

The final document shall not contain unresolved references such as:

\begin{verbatim}
??
[?]
\end{verbatim}

Before generating the final PDF, students shall check:

\begin{itemize}[leftmargin=0.8cm, itemsep=2pt]
    \item the Table of Contents;
    \item page numbering;
    \item chapter, section, and subsection numbering;
    \item figure numbering;
    \item table numbering;
    \item equation numbering;
    \item cross-references;
    \item citations;
    \item bibliography;
    \item font display;
    \item mathematical symbols;
    \item image quality;
    \item completeness of all required forms.
\end{itemize}

The document may be compiled using TeX Live, MiKTeX, or another
compatible \LaTeX{} environment. Regardless of the compilation
environment, the resulting PDF shall conform to the official Faculty
template.

\vspace{0.3cm}

% ================================================================
% XXII. CONTENT THAT STUDENTS MAY EDIT
% ================================================================

\noindent
\textbf{XXII.\quad CONTENT THAT STUDENTS MAY EDIT}

\vspace{0.15cm}

Students may edit or replace:

\begin{itemize}[leftmargin=0.8cm, itemsep=2pt]
    \item student name;
    \item class;
    \item internship organization;
    \item supervisor name;
    \item Abstract;
    \item List of Symbols and Abbreviations;
    \item chapter titles and chapter content;
    \item figures and photographs;
    \item tables;
    \item equations;
    \item citations and references.
\end{itemize}

Students shall not modify the main document formatting settings in the
preamble, including the document class, page geometry, line spacing,
heading formatting, table-of-contents formatting, font configuration,
page layout, or cover-page macros unless specifically instructed.

\vspace{0.3cm}

% ================================================================
% XXIII. FINAL SUBMISSION REQUIREMENTS
% ================================================================

\noindent
\textbf{XXIII.\quad FINAL SUBMISSION REQUIREMENTS}

\vspace{0.15cm}

Before submission, the final report shall:

\begin{itemize}[leftmargin=0.8cm, itemsep=2pt]
    \item compile without errors that affect the report content or presentation;
    \item contain all required forms and report sections;
    \item contain correct student, class, internship organization, and supervisor information;
    \item contain an updated Table of Contents, List of Figures, and List of Tables;
    \item contain no unused example content from the template;
    \item contain no placeholder text such as \textit{Title of Chapter 1},
    \textit{Section 1.1}, or other sample entries;
    \item contain all required citations and references;
    \item contain no unresolved cross-references.
\end{itemize}

If the Faculty requires submission of the \LaTeX{} source, the complete
project shall include all \verb|.tex| files, the bibliography database,
figures, and any other files required to reproduce the submitted PDF.

\vspace{0.4cm}

\noindent
\textbf{Important:}
Before submitting the final internship report, students shall remove or
comment out the following command in the main document:

\begin{verbatim}
\input{Regulations.tex}
\end{verbatim}

For example:

\begin{verbatim}
% \input{Regulations.tex}
\end{verbatim}

This removes the guidance pages from the final submitted report while
preserving the regulations file in the \LaTeX{} project for reference.

} % End 13pt content

\clearpage
 command before submitting
% the final internship report.
%
% =================================================================

\clearpage
\thispagestyle{empty}

\begin{center}
    {\fontsize{16pt}{20pt}\selectfont\bfseries
    SPECIFIC REGULATIONS FOR PREPARING\\
    GRADUATION INTERNSHIP REPORTS USING \LaTeX}
\end{center}

\vspace{0.5cm}

{\fontsize{13pt}{17pt}\selectfont
\setstretch{1.3}
\justifying

\noindent
\textbf{Note:}
This section is provided for guidance only and is included in the main
document using the command \verb|
% =================================================================
% REGULATIONS FOR PREPARING THE GRADUATION INTERNSHIP REPORT
% USING LATEX
% =================================================================
%
% This file is intended to be included from the main document with:
%
%     \input{Regulations.tex}
%
% IMPORTANT:
% This section is provided for guidance only. Students shall remove
% or comment out the \input{Regulations.tex} command before submitting
% the final internship report.
%
% =================================================================

\clearpage
\thispagestyle{empty}

\begin{center}
    {\fontsize{16pt}{20pt}\selectfont\bfseries
    SPECIFIC REGULATIONS FOR PREPARING\\
    GRADUATION INTERNSHIP REPORTS USING \LaTeX}
\end{center}

\vspace{0.5cm}

{\fontsize{13pt}{17pt}\selectfont
\setstretch{1.3}
\justifying

\noindent
\textbf{Note:}
This section is provided for guidance only and is included in the main
document using the command \verb|\input{Regulations.tex}|.
Students shall remove or comment out this command before submitting the
final internship report.

\vspace{0.4cm}

% ================================================================
% I. GENERAL REQUIREMENTS
% ================================================================

\noindent
\textbf{I.\quad GENERAL REQUIREMENTS}

\vspace{0.15cm}

The Graduation Internship Report shall be prepared using the official
\LaTeX{} template provided by the Faculty. Students shall preserve the
predefined page layout, typography, line spacing, heading styles,
numbering system, figure and table formatting, and bibliography settings
unless otherwise instructed by the Faculty or the supervisor.

The report shall be written in English.

Students are responsible for the academic and technical content of the
report. The official \LaTeX{} template is responsible for standardized
document formatting.

\vspace{0.3cm}

% ================================================================
% II. REPORT STRUCTURE
% ================================================================

\noindent
\textbf{II.\quad REPORT STRUCTURE}

\vspace{0.15cm}

The report shall normally be arranged in the following order:

\begin{enumerate}[label=\arabic*., leftmargin=1.2cm, itemsep=2pt]
    \item Outer Cover;
    \item Comments of the Internship Host Organization;
    \item Supervisor's Evaluation and Graduation Internship Defense Results;
    \item Table of Contents;
    \item List of Symbols and Abbreviations, if applicable;
    \item List of Figures;
    \item List of Tables;
    \item Abstract;
    \item Main Chapters of the Internship Report;
    \item References;
    \item Appendices, if applicable.
\end{enumerate}

The preliminary pages shall not be numbered as chapters. The main report
content shall be organized into consecutively numbered chapters beginning
with Chapter 1.

\vspace{0.3cm}

% ================================================================
% III. PAGE SIZE AND MARGINS
% ================================================================

\noindent
\textbf{III.\quad PAGE SIZE AND MARGINS}

\vspace{0.15cm}

The report shall be formatted on A4 paper with the following margins:

\begin{center}
\begin{tabular}{ll}
    Left margin:   & 3.5 cm \\
    Right margin:  & 2.0 cm \\
    Top margin:    & 2.5 cm \\
    Bottom margin: & 2.5 cm
\end{tabular}
\end{center}

These settings are predefined in the official template and shall not be
modified by students.

\vspace{0.3cm}

% ================================================================
% IV. FONT, FONT SIZE, AND LINE SPACING
% ================================================================

\noindent
\textbf{IV.\quad FONT, FONT SIZE, AND LINE SPACING}

\vspace{0.15cm}

The main text shall use a 13 pt font size with 1.3-line spacing.
A Times-family typeface shall be used in accordance with the official
template.

The text shall be fully justified.

Students shall not use repeated spaces, repeated manual line breaks, or
multiple blank lines to position text. Paragraph and heading spacing shall
be controlled by the template.

\vspace{0.3cm}

% ================================================================
% V. OUTER COVER
% ================================================================

\noindent
\textbf{V.\quad OUTER COVER}

\vspace{0.15cm}

The outer cover shall follow the official format provided in the template
and shall contain the following information:

\begin{itemize}[leftmargin=0.8cm, itemsep=2pt]
    \item Thai Nguyen University of Technology;
    \item Faculty of International Training;
    \item Engineering and Technology Department;
    \item University logo and Faculty logo;
    \item the title \textbf{INTERNSHIP REPORT};
    \item Internship Organization;
    \item Author;
    \item Class;
    \item Supervisor;
    \item place and year of completion.
\end{itemize}

Students shall only replace the required information fields and shall not
modify the cover-page layout unless instructed otherwise.

\vspace{0.3cm}

% ================================================================
% VI. COMMENTS OF THE INTERNSHIP HOST ORGANIZATION
% ================================================================

\noindent
\textbf{VI.\quad COMMENTS OF THE INTERNSHIP HOST ORGANIZATION}

\vspace{0.15cm}

The page entitled \textbf{COMMENTS OF THE INTERNSHIP HOST ORGANIZATION}
is reserved for the internship host organization.

It includes:

\begin{itemize}[leftmargin=0.8cm, itemsep=2pt]
    \item the student's full name;
    \item class;
    \item comments from the host organization;
    \item score;
    \item date;
    \item signature and official seal.
\end{itemize}

Students shall not alter or complete official evaluation fields unless
specifically authorized.

\vspace{0.3cm}

% ================================================================
% VII. SUPERVISOR'S EVALUATION AND DEFENSE RESULTS
% ================================================================

\noindent
\textbf{VII.\quad SUPERVISOR'S EVALUATION AND DEFENSE RESULTS}

\vspace{0.15cm}

The section entitled
\textbf{SUPERVISOR'S EVALUATION AND GRADUATION INTERNSHIP DEFENSE RESULTS}
is reserved for the supervisor and defense committee.

The evaluation may include:

\begin{itemize}[leftmargin=0.8cm, itemsep=2pt]
    \item Internship Progress and Student's Attitude;
    \item Report Content;
    \item Report Presentation;
    \item Other Comments;
    \item Supervisor's Opinion;
    \item Defense Result;
    \item Score.
\end{itemize}

Students shall not enter evaluation scores or defense results before they
are officially determined.

\vspace{0.3cm}

% ================================================================
% VIII. CHAPTERS, SECTIONS, AND SUBSECTIONS
% ================================================================

\noindent
\textbf{VIII.\quad CHAPTERS, SECTIONS, AND SUBSECTIONS}

\vspace{0.15cm}

The report structure shall be generated using the standard \LaTeX{}
commands:

\begin{verbatim}
\chapter{...}
\section{...}
\subsection{...}
\end{verbatim}

Chapter, section, and subsection numbers shall be generated automatically
by \LaTeX{}.

Students shall not manually type chapter or section numbers and shall not
reproduce heading styles using commands such as \verb|\textbf|,
\verb|\fontsize|, or \verb|\centering| instead of the appropriate
structural command.

Excessive subdivision should be avoided. In most cases, the hierarchy
Chapter--Section--Subsection is sufficient.

\vspace{0.3cm}

% ================================================================
% IX. ORGANIZATION OF THE MAIN CONTENT
% ================================================================

\noindent
\textbf{IX.\quad ORGANIZATION OF THE MAIN CONTENT}

\vspace{0.15cm}

The specific content of the report shall depend on the internship
organization and the assigned tasks. A typical structure may include:

\begin{itemize}[leftmargin=0.8cm, itemsep=2pt]
    \item \textbf{Chapter 1. Overview of the Internship Organization};
    \item \textbf{Chapter 2. Internship Activities and Assigned Tasks};
    \item \textbf{Chapter 3. Results, Experience, and Evaluation}.
\end{itemize}

These chapter titles are recommendations only and may be adapted to the
actual internship content.

The report shall reflect work actually performed, observed, analyzed, or
learned during the internship. Unrelated theoretical material should not
be included merely to increase report length.

\vspace{0.3cm}

% ================================================================
% X. TABLE OF CONTENTS
% ================================================================

\noindent
\textbf{X.\quad TABLE OF CONTENTS}

\vspace{0.15cm}

The Table of Contents shall be generated automatically using:

\begin{verbatim}
\tableofcontents
\end{verbatim}

Students shall not create or manually edit page numbers in the Table of
Contents.

The document shall be compiled sufficiently to ensure that all page
numbers and entries are correctly updated.

\vspace{0.3cm}

% ================================================================
% XI. LIST OF SYMBOLS AND ABBREVIATIONS
% ================================================================

\noindent
\textbf{XI.\quad LIST OF SYMBOLS AND ABBREVIATIONS}

\vspace{0.15cm}

Symbols and abbreviations that are used repeatedly in the report shall be
listed in the \textbf{List of Symbols and Abbreviations}.

The list may contain the following columns:

\begin{center}
\begin{tabular}{llll}
    No. & Symbol / Abbreviation & Unit & Definition
\end{tabular}
\end{center}

Only symbols and abbreviations actually used in the final report shall be
retained. Example entries contained in the template shall be removed if
they are not relevant to the student's report.

Each abbreviation should normally be defined when it first appears in
the main text.

\vspace{0.3cm}

% ================================================================
% XII. FIGURES, DIAGRAMS, AND PHOTOGRAPHS
% ================================================================

\noindent
\textbf{XII.\quad FIGURES, DIAGRAMS, AND PHOTOGRAPHS}

\vspace{0.15cm}

Figures shall be clear, legible, and directly relevant to the report
content.

Figures shall be inserted using the standard \verb|figure| environment.

Example:

\begin{verbatim}
\begin{figure}[H]
    \centering
    \includegraphics[width=0.8\textwidth]
                    {figures/company_structure.pdf}
    \caption{Organizational structure of the company}
    \label{fig:company-structure}
\end{figure}
\end{verbatim}

Figure captions shall be placed below the figure.

Each referenced figure shall have a unique \verb|\label| and shall be
referenced in the text using, for example:

\begin{verbatim}
Figure~\ref{fig:company-structure}
\end{verbatim}

Students shall not manually type figure numbers.

Figures reproduced or adapted from external sources shall include an
appropriate citation.

Photographs taken during the internship shall be relevant and sufficiently
clear. Students shall not include confidential, restricted, or proprietary
information without permission from the internship organization.

Vector graphics or high-resolution images are recommended whenever
possible.

\vspace{0.3cm}

% ================================================================
% XIII. TABLES
% ================================================================

\noindent
\textbf{XIII.\quad TABLES}

\vspace{0.15cm}

Tables shall be created using suitable \LaTeX{} environments such as
\verb|table|, \verb|tabular|, \verb|tabularx|, or \verb|longtable|.

Table captions shall be placed above the table.

Example:

\begin{verbatim}
\begin{table}[H]
    \centering
    \caption{Internship schedule}
    \label{tab:internship-schedule}

    \begin{tabular}{lll}
        \toprule
        Week & Task & Result \\
        \midrule
        1 & Orientation & Completed \\
        2 & System study & Completed \\
        \bottomrule
    \end{tabular}
\end{table}
\end{verbatim}

Tables shall be referenced using:

\begin{verbatim}
Table~\ref{tab:internship-schedule}
\end{verbatim}

Students shall not manually type table numbers.

Tables should be created directly in \LaTeX{} rather than inserted as
screenshots whenever practical.

\vspace{0.3cm}

% ================================================================
% XIV. EQUATIONS AND TECHNICAL QUANTITIES
% ================================================================

\noindent
\textbf{XIV.\quad EQUATIONS AND TECHNICAL QUANTITIES}

\vspace{0.15cm}

Mathematical expressions shall be typeset directly using \LaTeX{}.
Images or screenshots of equations shall not be used.

Short expressions may be written in inline mathematical mode, for example:

\begin{verbatim}
$R = 0.7\,\Omega$
\end{verbatim}

Important equations shall use a numbered equation environment.

Example:

\begin{verbatim}
\begin{equation}
    P = UI
    \label{eq:power}
\end{equation}
\end{verbatim}

Equations shall be referenced using:

\begin{verbatim}
Equation~\eqref{eq:power}
\end{verbatim}

Equation numbers shall not be manually entered.

SI units shall be used wherever applicable. The \verb|siunitx| package
should be used where appropriate.

Examples:

\begin{verbatim}
\SI{220}{\volt}
\SI{50}{\hertz}
\SI{2.5}{\ampere}
\end{verbatim}

Notation and units shall be consistent throughout the report.

\vspace{0.3cm}

% ================================================================
% XV. SOURCE CODE
% ================================================================

\noindent
\textbf{XV.\quad SOURCE CODE}

\vspace{0.15cm}

If the internship involves programming, PLCs, microcontrollers,
simulation, or software development, relevant source code may be included
in the report.

Source code shall be presented using the \verb|listings| environment or
another environment defined by the official template.

Example:

\begin{verbatim}
\begin{lstlisting}[language=st]
IF Start = TRUE THEN
    Motor := TRUE;
ELSE
    Motor := FALSE;
END_IF;
\end{lstlisting}
\end{verbatim}

Only code necessary for explaining the work should appear in the main
report. Long source-code listings should be moved to an appendix or
submitted as supplementary files if required.

Source code should not normally be inserted as screenshots.

\vspace{0.3cm}

% ================================================================
% XVI. ABSTRACT
% ================================================================

\noindent
\textbf{XVI.\quad ABSTRACT}

\vspace{0.15cm}

The Abstract shall briefly summarize:

\begin{itemize}[leftmargin=0.8cm, itemsep=2pt]
    \item the internship organization;
    \item the objectives of the internship;
    \item the main tasks performed;
    \item the main results, knowledge, and skills obtained.
\end{itemize}

The Abstract should be concise and should not reproduce the full content
of the report chapter by chapter.

\vspace{0.3cm}

% ================================================================
% XVII. REFERENCES AND CITATIONS
% ================================================================

\noindent
\textbf{XVII.\quad REFERENCES AND CITATIONS}

\vspace{0.15cm}

All books, journal papers, conference papers, standards, websites,
datasheets, manuals, technical documents, figures, tables, or other
external sources used in preparing the report shall be properly cited.

References shall be managed through the bibliography system provided with
the official template.

Bibliographic entries shall normally be stored in:

\begin{verbatim}
References.bib
\end{verbatim}

References shall be cited in the report using:

\begin{verbatim}
\cite{reference-key}
\end{verbatim}

Students shall not manually type reference numbers such as ``[1]'',
``[2]'', etc.

Titles of foreign-language publications shall remain in their original
language in the bibliography.

All borrowed ideas, data, figures, tables, equations, methods, or results
shall be properly acknowledged.

\vspace{0.3cm}

% ================================================================
% XVIII. CROSS-REFERENCING
% ================================================================

\noindent
\textbf{XVIII.\quad CROSS-REFERENCING}

\vspace{0.15cm}

Numbered elements shall be cross-referenced using the standard \LaTeX{}
mechanism.

Use the following commands where appropriate:

\begin{verbatim}
\label{...}
\ref{...}
\eqref{...}
\end{verbatim}

For example:

\begin{verbatim}
As shown in Figure~\ref{fig:system}, ...
\end{verbatim}

Students shall not manually type numbers such as ``Figure 2.3'' when
those numbers can be generated automatically.

\vspace{0.3cm}

% ================================================================
% XIX. PAGE NUMBERING
% ================================================================

\noindent
\textbf{XIX.\quad PAGE NUMBERING}

\vspace{0.15cm}

Cover pages and official evaluation forms at the beginning of the report
shall not display page numbers.

The main body of the report shall use Arabic page numbers beginning from
Chapter 1.

Page numbering shall be generated automatically by the official template.
Students shall not manually enter page numbers.

\vspace{0.3cm}

% ================================================================
% XX. PROJECT FILE ORGANIZATION
% ================================================================

\noindent
\textbf{XX.\quad PROJECT FILE ORGANIZATION}

\vspace{0.15cm}

Students should maintain a clear and organized \LaTeX{} project
structure. A recommended structure is:

\begin{verbatim}
InternshipReport/
|
|-- main.tex
|-- References.bib
|
|-- figures/
|   |-- company.pdf
|   |-- system.png
|   `-- ...
|
`-- chapters/
    |-- chapter1.tex
    |-- chapter2.tex
    `-- chapter3.tex
\end{verbatim}

Figures should normally be stored in the \verb|figures| directory
specified by the official template.

File and directory names should preferably contain Latin letters,
numbers, hyphens, or underscores.

Recommended examples include:

\begin{verbatim}
company_structure.pdf
motor_control.png
chapter_01.tex
\end{verbatim}

Spaces and unnecessary special characters in file names should be
avoided where possible.

\vspace{0.3cm}

% ================================================================
% XXI. DOCUMENT COMPILATION AND FINAL CHECK
% ================================================================

\noindent
\textbf{XXI.\quad DOCUMENT COMPILATION AND FINAL CHECK}

\vspace{0.15cm}

Before submission, the complete \LaTeX{} project shall compile
successfully and produce the final PDF report.

The final document shall not contain unresolved references such as:

\begin{verbatim}
??
[?]
\end{verbatim}

Before generating the final PDF, students shall check:

\begin{itemize}[leftmargin=0.8cm, itemsep=2pt]
    \item the Table of Contents;
    \item page numbering;
    \item chapter, section, and subsection numbering;
    \item figure numbering;
    \item table numbering;
    \item equation numbering;
    \item cross-references;
    \item citations;
    \item bibliography;
    \item font display;
    \item mathematical symbols;
    \item image quality;
    \item completeness of all required forms.
\end{itemize}

The document may be compiled using TeX Live, MiKTeX, or another
compatible \LaTeX{} environment. Regardless of the compilation
environment, the resulting PDF shall conform to the official Faculty
template.

\vspace{0.3cm}

% ================================================================
% XXII. CONTENT THAT STUDENTS MAY EDIT
% ================================================================

\noindent
\textbf{XXII.\quad CONTENT THAT STUDENTS MAY EDIT}

\vspace{0.15cm}

Students may edit or replace:

\begin{itemize}[leftmargin=0.8cm, itemsep=2pt]
    \item student name;
    \item class;
    \item internship organization;
    \item supervisor name;
    \item Abstract;
    \item List of Symbols and Abbreviations;
    \item chapter titles and chapter content;
    \item figures and photographs;
    \item tables;
    \item equations;
    \item citations and references.
\end{itemize}

Students shall not modify the main document formatting settings in the
preamble, including the document class, page geometry, line spacing,
heading formatting, table-of-contents formatting, font configuration,
page layout, or cover-page macros unless specifically instructed.

\vspace{0.3cm}

% ================================================================
% XXIII. FINAL SUBMISSION REQUIREMENTS
% ================================================================

\noindent
\textbf{XXIII.\quad FINAL SUBMISSION REQUIREMENTS}

\vspace{0.15cm}

Before submission, the final report shall:

\begin{itemize}[leftmargin=0.8cm, itemsep=2pt]
    \item compile without errors that affect the report content or presentation;
    \item contain all required forms and report sections;
    \item contain correct student, class, internship organization, and supervisor information;
    \item contain an updated Table of Contents, List of Figures, and List of Tables;
    \item contain no unused example content from the template;
    \item contain no placeholder text such as \textit{Title of Chapter 1},
    \textit{Section 1.1}, or other sample entries;
    \item contain all required citations and references;
    \item contain no unresolved cross-references.
\end{itemize}

If the Faculty requires submission of the \LaTeX{} source, the complete
project shall include all \verb|.tex| files, the bibliography database,
figures, and any other files required to reproduce the submitted PDF.

\vspace{0.4cm}

\noindent
\textbf{Important:}
Before submitting the final internship report, students shall remove or
comment out the following command in the main document:

\begin{verbatim}
\input{Regulations.tex}
\end{verbatim}

For example:

\begin{verbatim}
% \input{Regulations.tex}
\end{verbatim}

This removes the guidance pages from the final submitted report while
preserving the regulations file in the \LaTeX{} project for reference.

} % End 13pt content

\clearpage
|.
Students shall remove or comment out this command before submitting the
final internship report.

\vspace{0.4cm}

% ================================================================
% I. GENERAL REQUIREMENTS
% ================================================================

\noindent
\textbf{I.\quad GENERAL REQUIREMENTS}

\vspace{0.15cm}

The Graduation Internship Report shall be prepared using the official
\LaTeX{} template provided by the Faculty. Students shall preserve the
predefined page layout, typography, line spacing, heading styles,
numbering system, figure and table formatting, and bibliography settings
unless otherwise instructed by the Faculty or the supervisor.

The report shall be written in English.

Students are responsible for the academic and technical content of the
report. The official \LaTeX{} template is responsible for standardized
document formatting.

\vspace{0.3cm}

% ================================================================
% II. REPORT STRUCTURE
% ================================================================

\noindent
\textbf{II.\quad REPORT STRUCTURE}

\vspace{0.15cm}

The report shall normally be arranged in the following order:

\begin{enumerate}[label=\arabic*., leftmargin=1.2cm, itemsep=2pt]
    \item Outer Cover;
    \item Comments of the Internship Host Organization;
    \item Supervisor's Evaluation and Graduation Internship Defense Results;
    \item Table of Contents;
    \item List of Symbols and Abbreviations, if applicable;
    \item List of Figures;
    \item List of Tables;
    \item Abstract;
    \item Main Chapters of the Internship Report;
    \item References;
    \item Appendices, if applicable.
\end{enumerate}

The preliminary pages shall not be numbered as chapters. The main report
content shall be organized into consecutively numbered chapters beginning
with Chapter 1.

\vspace{0.3cm}

% ================================================================
% III. PAGE SIZE AND MARGINS
% ================================================================

\noindent
\textbf{III.\quad PAGE SIZE AND MARGINS}

\vspace{0.15cm}

The report shall be formatted on A4 paper with the following margins:

\begin{center}
\begin{tabular}{ll}
    Left margin:   & 3.5 cm \\
    Right margin:  & 2.0 cm \\
    Top margin:    & 2.5 cm \\
    Bottom margin: & 2.5 cm
\end{tabular}
\end{center}

These settings are predefined in the official template and shall not be
modified by students.

\vspace{0.3cm}

% ================================================================
% IV. FONT, FONT SIZE, AND LINE SPACING
% ================================================================

\noindent
\textbf{IV.\quad FONT, FONT SIZE, AND LINE SPACING}

\vspace{0.15cm}

The main text shall use a 13 pt font size with 1.3-line spacing.
A Times-family typeface shall be used in accordance with the official
template.

The text shall be fully justified.

Students shall not use repeated spaces, repeated manual line breaks, or
multiple blank lines to position text. Paragraph and heading spacing shall
be controlled by the template.

\vspace{0.3cm}

% ================================================================
% V. OUTER COVER
% ================================================================

\noindent
\textbf{V.\quad OUTER COVER}

\vspace{0.15cm}

The outer cover shall follow the official format provided in the template
and shall contain the following information:

\begin{itemize}[leftmargin=0.8cm, itemsep=2pt]
    \item Thai Nguyen University of Technology;
    \item Faculty of International Training;
    \item Engineering and Technology Department;
    \item University logo and Faculty logo;
    \item the title \textbf{INTERNSHIP REPORT};
    \item Internship Organization;
    \item Author;
    \item Class;
    \item Supervisor;
    \item place and year of completion.
\end{itemize}

Students shall only replace the required information fields and shall not
modify the cover-page layout unless instructed otherwise.

\vspace{0.3cm}

% ================================================================
% VI. COMMENTS OF THE INTERNSHIP HOST ORGANIZATION
% ================================================================

\noindent
\textbf{VI.\quad COMMENTS OF THE INTERNSHIP HOST ORGANIZATION}

\vspace{0.15cm}

The page entitled \textbf{COMMENTS OF THE INTERNSHIP HOST ORGANIZATION}
is reserved for the internship host organization.

It includes:

\begin{itemize}[leftmargin=0.8cm, itemsep=2pt]
    \item the student's full name;
    \item class;
    \item comments from the host organization;
    \item score;
    \item date;
    \item signature and official seal.
\end{itemize}

Students shall not alter or complete official evaluation fields unless
specifically authorized.

\vspace{0.3cm}

% ================================================================
% VII. SUPERVISOR'S EVALUATION AND DEFENSE RESULTS
% ================================================================

\noindent
\textbf{VII.\quad SUPERVISOR'S EVALUATION AND DEFENSE RESULTS}

\vspace{0.15cm}

The section entitled
\textbf{SUPERVISOR'S EVALUATION AND GRADUATION INTERNSHIP DEFENSE RESULTS}
is reserved for the supervisor and defense committee.

The evaluation may include:

\begin{itemize}[leftmargin=0.8cm, itemsep=2pt]
    \item Internship Progress and Student's Attitude;
    \item Report Content;
    \item Report Presentation;
    \item Other Comments;
    \item Supervisor's Opinion;
    \item Defense Result;
    \item Score.
\end{itemize}

Students shall not enter evaluation scores or defense results before they
are officially determined.

\vspace{0.3cm}

% ================================================================
% VIII. CHAPTERS, SECTIONS, AND SUBSECTIONS
% ================================================================

\noindent
\textbf{VIII.\quad CHAPTERS, SECTIONS, AND SUBSECTIONS}

\vspace{0.15cm}

The report structure shall be generated using the standard \LaTeX{}
commands:

\begin{verbatim}
\chapter{...}
\section{...}
\subsection{...}
\end{verbatim}

Chapter, section, and subsection numbers shall be generated automatically
by \LaTeX{}.

Students shall not manually type chapter or section numbers and shall not
reproduce heading styles using commands such as \verb|\textbf|,
\verb|\fontsize|, or \verb|\centering| instead of the appropriate
structural command.

Excessive subdivision should be avoided. In most cases, the hierarchy
Chapter--Section--Subsection is sufficient.

\vspace{0.3cm}

% ================================================================
% IX. ORGANIZATION OF THE MAIN CONTENT
% ================================================================

\noindent
\textbf{IX.\quad ORGANIZATION OF THE MAIN CONTENT}

\vspace{0.15cm}

The specific content of the report shall depend on the internship
organization and the assigned tasks. A typical structure may include:

\begin{itemize}[leftmargin=0.8cm, itemsep=2pt]
    \item \textbf{Chapter 1. Overview of the Internship Organization};
    \item \textbf{Chapter 2. Internship Activities and Assigned Tasks};
    \item \textbf{Chapter 3. Results, Experience, and Evaluation}.
\end{itemize}

These chapter titles are recommendations only and may be adapted to the
actual internship content.

The report shall reflect work actually performed, observed, analyzed, or
learned during the internship. Unrelated theoretical material should not
be included merely to increase report length.

\vspace{0.3cm}

% ================================================================
% X. TABLE OF CONTENTS
% ================================================================

\noindent
\textbf{X.\quad TABLE OF CONTENTS}

\vspace{0.15cm}

The Table of Contents shall be generated automatically using:

\begin{verbatim}
\tableofcontents
\end{verbatim}

Students shall not create or manually edit page numbers in the Table of
Contents.

The document shall be compiled sufficiently to ensure that all page
numbers and entries are correctly updated.

\vspace{0.3cm}

% ================================================================
% XI. LIST OF SYMBOLS AND ABBREVIATIONS
% ================================================================

\noindent
\textbf{XI.\quad LIST OF SYMBOLS AND ABBREVIATIONS}

\vspace{0.15cm}

Symbols and abbreviations that are used repeatedly in the report shall be
listed in the \textbf{List of Symbols and Abbreviations}.

The list may contain the following columns:

\begin{center}
\begin{tabular}{llll}
    No. & Symbol / Abbreviation & Unit & Definition
\end{tabular}
\end{center}

Only symbols and abbreviations actually used in the final report shall be
retained. Example entries contained in the template shall be removed if
they are not relevant to the student's report.

Each abbreviation should normally be defined when it first appears in
the main text.

\vspace{0.3cm}

% ================================================================
% XII. FIGURES, DIAGRAMS, AND PHOTOGRAPHS
% ================================================================

\noindent
\textbf{XII.\quad FIGURES, DIAGRAMS, AND PHOTOGRAPHS}

\vspace{0.15cm}

Figures shall be clear, legible, and directly relevant to the report
content.

Figures shall be inserted using the standard \verb|figure| environment.

Example:

\begin{verbatim}
\begin{figure}[H]
    \centering
    \includegraphics[width=0.8\textwidth]
                    {figures/company_structure.pdf}
    \caption{Organizational structure of the company}
    \label{fig:company-structure}
\end{figure}
\end{verbatim}

Figure captions shall be placed below the figure.

Each referenced figure shall have a unique \verb|\label| and shall be
referenced in the text using, for example:

\begin{verbatim}
Figure~\ref{fig:company-structure}
\end{verbatim}

Students shall not manually type figure numbers.

Figures reproduced or adapted from external sources shall include an
appropriate citation.

Photographs taken during the internship shall be relevant and sufficiently
clear. Students shall not include confidential, restricted, or proprietary
information without permission from the internship organization.

Vector graphics or high-resolution images are recommended whenever
possible.

\vspace{0.3cm}

% ================================================================
% XIII. TABLES
% ================================================================

\noindent
\textbf{XIII.\quad TABLES}

\vspace{0.15cm}

Tables shall be created using suitable \LaTeX{} environments such as
\verb|table|, \verb|tabular|, \verb|tabularx|, or \verb|longtable|.

Table captions shall be placed above the table.

Example:

\begin{verbatim}
\begin{table}[H]
    \centering
    \caption{Internship schedule}
    \label{tab:internship-schedule}

    \begin{tabular}{lll}
        \toprule
        Week & Task & Result \\
        \midrule
        1 & Orientation & Completed \\
        2 & System study & Completed \\
        \bottomrule
    \end{tabular}
\end{table}
\end{verbatim}

Tables shall be referenced using:

\begin{verbatim}
Table~\ref{tab:internship-schedule}
\end{verbatim}

Students shall not manually type table numbers.

Tables should be created directly in \LaTeX{} rather than inserted as
screenshots whenever practical.

\vspace{0.3cm}

% ================================================================
% XIV. EQUATIONS AND TECHNICAL QUANTITIES
% ================================================================

\noindent
\textbf{XIV.\quad EQUATIONS AND TECHNICAL QUANTITIES}

\vspace{0.15cm}

Mathematical expressions shall be typeset directly using \LaTeX{}.
Images or screenshots of equations shall not be used.

Short expressions may be written in inline mathematical mode, for example:

\begin{verbatim}
$R = 0.7\,\Omega$
\end{verbatim}

Important equations shall use a numbered equation environment.

Example:

\begin{verbatim}
\begin{equation}
    P = UI
    \label{eq:power}
\end{equation}
\end{verbatim}

Equations shall be referenced using:

\begin{verbatim}
Equation~\eqref{eq:power}
\end{verbatim}

Equation numbers shall not be manually entered.

SI units shall be used wherever applicable. The \verb|siunitx| package
should be used where appropriate.

Examples:

\begin{verbatim}
\SI{220}{\volt}
\SI{50}{\hertz}
\SI{2.5}{\ampere}
\end{verbatim}

Notation and units shall be consistent throughout the report.

\vspace{0.3cm}

% ================================================================
% XV. SOURCE CODE
% ================================================================

\noindent
\textbf{XV.\quad SOURCE CODE}

\vspace{0.15cm}

If the internship involves programming, PLCs, microcontrollers,
simulation, or software development, relevant source code may be included
in the report.

Source code shall be presented using the \verb|listings| environment or
another environment defined by the official template.

Example:

\begin{verbatim}
\begin{lstlisting}[language=st]
IF Start = TRUE THEN
    Motor := TRUE;
ELSE
    Motor := FALSE;
END_IF;
\end{lstlisting}
\end{verbatim}

Only code necessary for explaining the work should appear in the main
report. Long source-code listings should be moved to an appendix or
submitted as supplementary files if required.

Source code should not normally be inserted as screenshots.

\vspace{0.3cm}

% ================================================================
% XVI. ABSTRACT
% ================================================================

\noindent
\textbf{XVI.\quad ABSTRACT}

\vspace{0.15cm}

The Abstract shall briefly summarize:

\begin{itemize}[leftmargin=0.8cm, itemsep=2pt]
    \item the internship organization;
    \item the objectives of the internship;
    \item the main tasks performed;
    \item the main results, knowledge, and skills obtained.
\end{itemize}

The Abstract should be concise and should not reproduce the full content
of the report chapter by chapter.

\vspace{0.3cm}

% ================================================================
% XVII. REFERENCES AND CITATIONS
% ================================================================

\noindent
\textbf{XVII.\quad REFERENCES AND CITATIONS}

\vspace{0.15cm}

All books, journal papers, conference papers, standards, websites,
datasheets, manuals, technical documents, figures, tables, or other
external sources used in preparing the report shall be properly cited.

References shall be managed through the bibliography system provided with
the official template.

Bibliographic entries shall normally be stored in:

\begin{verbatim}
References.bib
\end{verbatim}

References shall be cited in the report using:

\begin{verbatim}
\cite{reference-key}
\end{verbatim}

Students shall not manually type reference numbers such as ``[1]'',
``[2]'', etc.

Titles of foreign-language publications shall remain in their original
language in the bibliography.

All borrowed ideas, data, figures, tables, equations, methods, or results
shall be properly acknowledged.

\vspace{0.3cm}

% ================================================================
% XVIII. CROSS-REFERENCING
% ================================================================

\noindent
\textbf{XVIII.\quad CROSS-REFERENCING}

\vspace{0.15cm}

Numbered elements shall be cross-referenced using the standard \LaTeX{}
mechanism.

Use the following commands where appropriate:

\begin{verbatim}
\label{...}
\ref{...}
\eqref{...}
\end{verbatim}

For example:

\begin{verbatim}
As shown in Figure~\ref{fig:system}, ...
\end{verbatim}

Students shall not manually type numbers such as ``Figure 2.3'' when
those numbers can be generated automatically.

\vspace{0.3cm}

% ================================================================
% XIX. PAGE NUMBERING
% ================================================================

\noindent
\textbf{XIX.\quad PAGE NUMBERING}

\vspace{0.15cm}

Cover pages and official evaluation forms at the beginning of the report
shall not display page numbers.

The main body of the report shall use Arabic page numbers beginning from
Chapter 1.

Page numbering shall be generated automatically by the official template.
Students shall not manually enter page numbers.

\vspace{0.3cm}

% ================================================================
% XX. PROJECT FILE ORGANIZATION
% ================================================================

\noindent
\textbf{XX.\quad PROJECT FILE ORGANIZATION}

\vspace{0.15cm}

Students should maintain a clear and organized \LaTeX{} project
structure. A recommended structure is:

\begin{verbatim}
InternshipReport/
|
|-- main.tex
|-- References.bib
|
|-- figures/
|   |-- company.pdf
|   |-- system.png
|   `-- ...
|
`-- chapters/
    |-- chapter1.tex
    |-- chapter2.tex
    `-- chapter3.tex
\end{verbatim}

Figures should normally be stored in the \verb|figures| directory
specified by the official template.

File and directory names should preferably contain Latin letters,
numbers, hyphens, or underscores.

Recommended examples include:

\begin{verbatim}
company_structure.pdf
motor_control.png
chapter_01.tex
\end{verbatim}

Spaces and unnecessary special characters in file names should be
avoided where possible.

\vspace{0.3cm}

% ================================================================
% XXI. DOCUMENT COMPILATION AND FINAL CHECK
% ================================================================

\noindent
\textbf{XXI.\quad DOCUMENT COMPILATION AND FINAL CHECK}

\vspace{0.15cm}

Before submission, the complete \LaTeX{} project shall compile
successfully and produce the final PDF report.

The final document shall not contain unresolved references such as:

\begin{verbatim}
??
[?]
\end{verbatim}

Before generating the final PDF, students shall check:

\begin{itemize}[leftmargin=0.8cm, itemsep=2pt]
    \item the Table of Contents;
    \item page numbering;
    \item chapter, section, and subsection numbering;
    \item figure numbering;
    \item table numbering;
    \item equation numbering;
    \item cross-references;
    \item citations;
    \item bibliography;
    \item font display;
    \item mathematical symbols;
    \item image quality;
    \item completeness of all required forms.
\end{itemize}

The document may be compiled using TeX Live, MiKTeX, or another
compatible \LaTeX{} environment. Regardless of the compilation
environment, the resulting PDF shall conform to the official Faculty
template.

\vspace{0.3cm}

% ================================================================
% XXII. CONTENT THAT STUDENTS MAY EDIT
% ================================================================

\noindent
\textbf{XXII.\quad CONTENT THAT STUDENTS MAY EDIT}

\vspace{0.15cm}

Students may edit or replace:

\begin{itemize}[leftmargin=0.8cm, itemsep=2pt]
    \item student name;
    \item class;
    \item internship organization;
    \item supervisor name;
    \item Abstract;
    \item List of Symbols and Abbreviations;
    \item chapter titles and chapter content;
    \item figures and photographs;
    \item tables;
    \item equations;
    \item citations and references.
\end{itemize}

Students shall not modify the main document formatting settings in the
preamble, including the document class, page geometry, line spacing,
heading formatting, table-of-contents formatting, font configuration,
page layout, or cover-page macros unless specifically instructed.

\vspace{0.3cm}

% ================================================================
% XXIII. FINAL SUBMISSION REQUIREMENTS
% ================================================================

\noindent
\textbf{XXIII.\quad FINAL SUBMISSION REQUIREMENTS}

\vspace{0.15cm}

Before submission, the final report shall:

\begin{itemize}[leftmargin=0.8cm, itemsep=2pt]
    \item compile without errors that affect the report content or presentation;
    \item contain all required forms and report sections;
    \item contain correct student, class, internship organization, and supervisor information;
    \item contain an updated Table of Contents, List of Figures, and List of Tables;
    \item contain no unused example content from the template;
    \item contain no placeholder text such as \textit{Title of Chapter 1},
    \textit{Section 1.1}, or other sample entries;
    \item contain all required citations and references;
    \item contain no unresolved cross-references.
\end{itemize}

If the Faculty requires submission of the \LaTeX{} source, the complete
project shall include all \verb|.tex| files, the bibliography database,
figures, and any other files required to reproduce the submitted PDF.

\vspace{0.4cm}

\noindent
\textbf{Important:}
Before submitting the final internship report, students shall remove or
comment out the following command in the main document:

\begin{verbatim}

% =================================================================
% REGULATIONS FOR PREPARING THE GRADUATION INTERNSHIP REPORT
% USING LATEX
% =================================================================
%
% This file is intended to be included from the main document with:
%
%     \input{Regulations.tex}
%
% IMPORTANT:
% This section is provided for guidance only. Students shall remove
% or comment out the \input{Regulations.tex} command before submitting
% the final internship report.
%
% =================================================================

\clearpage
\thispagestyle{empty}

\begin{center}
    {\fontsize{16pt}{20pt}\selectfont\bfseries
    SPECIFIC REGULATIONS FOR PREPARING\\
    GRADUATION INTERNSHIP REPORTS USING \LaTeX}
\end{center}

\vspace{0.5cm}

{\fontsize{13pt}{17pt}\selectfont
\setstretch{1.3}
\justifying

\noindent
\textbf{Note:}
This section is provided for guidance only and is included in the main
document using the command \verb|\input{Regulations.tex}|.
Students shall remove or comment out this command before submitting the
final internship report.

\vspace{0.4cm}

% ================================================================
% I. GENERAL REQUIREMENTS
% ================================================================

\noindent
\textbf{I.\quad GENERAL REQUIREMENTS}

\vspace{0.15cm}

The Graduation Internship Report shall be prepared using the official
\LaTeX{} template provided by the Faculty. Students shall preserve the
predefined page layout, typography, line spacing, heading styles,
numbering system, figure and table formatting, and bibliography settings
unless otherwise instructed by the Faculty or the supervisor.

The report shall be written in English.

Students are responsible for the academic and technical content of the
report. The official \LaTeX{} template is responsible for standardized
document formatting.

\vspace{0.3cm}

% ================================================================
% II. REPORT STRUCTURE
% ================================================================

\noindent
\textbf{II.\quad REPORT STRUCTURE}

\vspace{0.15cm}

The report shall normally be arranged in the following order:

\begin{enumerate}[label=\arabic*., leftmargin=1.2cm, itemsep=2pt]
    \item Outer Cover;
    \item Comments of the Internship Host Organization;
    \item Supervisor's Evaluation and Graduation Internship Defense Results;
    \item Table of Contents;
    \item List of Symbols and Abbreviations, if applicable;
    \item List of Figures;
    \item List of Tables;
    \item Abstract;
    \item Main Chapters of the Internship Report;
    \item References;
    \item Appendices, if applicable.
\end{enumerate}

The preliminary pages shall not be numbered as chapters. The main report
content shall be organized into consecutively numbered chapters beginning
with Chapter 1.

\vspace{0.3cm}

% ================================================================
% III. PAGE SIZE AND MARGINS
% ================================================================

\noindent
\textbf{III.\quad PAGE SIZE AND MARGINS}

\vspace{0.15cm}

The report shall be formatted on A4 paper with the following margins:

\begin{center}
\begin{tabular}{ll}
    Left margin:   & 3.5 cm \\
    Right margin:  & 2.0 cm \\
    Top margin:    & 2.5 cm \\
    Bottom margin: & 2.5 cm
\end{tabular}
\end{center}

These settings are predefined in the official template and shall not be
modified by students.

\vspace{0.3cm}

% ================================================================
% IV. FONT, FONT SIZE, AND LINE SPACING
% ================================================================

\noindent
\textbf{IV.\quad FONT, FONT SIZE, AND LINE SPACING}

\vspace{0.15cm}

The main text shall use a 13 pt font size with 1.3-line spacing.
A Times-family typeface shall be used in accordance with the official
template.

The text shall be fully justified.

Students shall not use repeated spaces, repeated manual line breaks, or
multiple blank lines to position text. Paragraph and heading spacing shall
be controlled by the template.

\vspace{0.3cm}

% ================================================================
% V. OUTER COVER
% ================================================================

\noindent
\textbf{V.\quad OUTER COVER}

\vspace{0.15cm}

The outer cover shall follow the official format provided in the template
and shall contain the following information:

\begin{itemize}[leftmargin=0.8cm, itemsep=2pt]
    \item Thai Nguyen University of Technology;
    \item Faculty of International Training;
    \item Engineering and Technology Department;
    \item University logo and Faculty logo;
    \item the title \textbf{INTERNSHIP REPORT};
    \item Internship Organization;
    \item Author;
    \item Class;
    \item Supervisor;
    \item place and year of completion.
\end{itemize}

Students shall only replace the required information fields and shall not
modify the cover-page layout unless instructed otherwise.

\vspace{0.3cm}

% ================================================================
% VI. COMMENTS OF THE INTERNSHIP HOST ORGANIZATION
% ================================================================

\noindent
\textbf{VI.\quad COMMENTS OF THE INTERNSHIP HOST ORGANIZATION}

\vspace{0.15cm}

The page entitled \textbf{COMMENTS OF THE INTERNSHIP HOST ORGANIZATION}
is reserved for the internship host organization.

It includes:

\begin{itemize}[leftmargin=0.8cm, itemsep=2pt]
    \item the student's full name;
    \item class;
    \item comments from the host organization;
    \item score;
    \item date;
    \item signature and official seal.
\end{itemize}

Students shall not alter or complete official evaluation fields unless
specifically authorized.

\vspace{0.3cm}

% ================================================================
% VII. SUPERVISOR'S EVALUATION AND DEFENSE RESULTS
% ================================================================

\noindent
\textbf{VII.\quad SUPERVISOR'S EVALUATION AND DEFENSE RESULTS}

\vspace{0.15cm}

The section entitled
\textbf{SUPERVISOR'S EVALUATION AND GRADUATION INTERNSHIP DEFENSE RESULTS}
is reserved for the supervisor and defense committee.

The evaluation may include:

\begin{itemize}[leftmargin=0.8cm, itemsep=2pt]
    \item Internship Progress and Student's Attitude;
    \item Report Content;
    \item Report Presentation;
    \item Other Comments;
    \item Supervisor's Opinion;
    \item Defense Result;
    \item Score.
\end{itemize}

Students shall not enter evaluation scores or defense results before they
are officially determined.

\vspace{0.3cm}

% ================================================================
% VIII. CHAPTERS, SECTIONS, AND SUBSECTIONS
% ================================================================

\noindent
\textbf{VIII.\quad CHAPTERS, SECTIONS, AND SUBSECTIONS}

\vspace{0.15cm}

The report structure shall be generated using the standard \LaTeX{}
commands:

\begin{verbatim}
\chapter{...}
\section{...}
\subsection{...}
\end{verbatim}

Chapter, section, and subsection numbers shall be generated automatically
by \LaTeX{}.

Students shall not manually type chapter or section numbers and shall not
reproduce heading styles using commands such as \verb|\textbf|,
\verb|\fontsize|, or \verb|\centering| instead of the appropriate
structural command.

Excessive subdivision should be avoided. In most cases, the hierarchy
Chapter--Section--Subsection is sufficient.

\vspace{0.3cm}

% ================================================================
% IX. ORGANIZATION OF THE MAIN CONTENT
% ================================================================

\noindent
\textbf{IX.\quad ORGANIZATION OF THE MAIN CONTENT}

\vspace{0.15cm}

The specific content of the report shall depend on the internship
organization and the assigned tasks. A typical structure may include:

\begin{itemize}[leftmargin=0.8cm, itemsep=2pt]
    \item \textbf{Chapter 1. Overview of the Internship Organization};
    \item \textbf{Chapter 2. Internship Activities and Assigned Tasks};
    \item \textbf{Chapter 3. Results, Experience, and Evaluation}.
\end{itemize}

These chapter titles are recommendations only and may be adapted to the
actual internship content.

The report shall reflect work actually performed, observed, analyzed, or
learned during the internship. Unrelated theoretical material should not
be included merely to increase report length.

\vspace{0.3cm}

% ================================================================
% X. TABLE OF CONTENTS
% ================================================================

\noindent
\textbf{X.\quad TABLE OF CONTENTS}

\vspace{0.15cm}

The Table of Contents shall be generated automatically using:

\begin{verbatim}
\tableofcontents
\end{verbatim}

Students shall not create or manually edit page numbers in the Table of
Contents.

The document shall be compiled sufficiently to ensure that all page
numbers and entries are correctly updated.

\vspace{0.3cm}

% ================================================================
% XI. LIST OF SYMBOLS AND ABBREVIATIONS
% ================================================================

\noindent
\textbf{XI.\quad LIST OF SYMBOLS AND ABBREVIATIONS}

\vspace{0.15cm}

Symbols and abbreviations that are used repeatedly in the report shall be
listed in the \textbf{List of Symbols and Abbreviations}.

The list may contain the following columns:

\begin{center}
\begin{tabular}{llll}
    No. & Symbol / Abbreviation & Unit & Definition
\end{tabular}
\end{center}

Only symbols and abbreviations actually used in the final report shall be
retained. Example entries contained in the template shall be removed if
they are not relevant to the student's report.

Each abbreviation should normally be defined when it first appears in
the main text.

\vspace{0.3cm}

% ================================================================
% XII. FIGURES, DIAGRAMS, AND PHOTOGRAPHS
% ================================================================

\noindent
\textbf{XII.\quad FIGURES, DIAGRAMS, AND PHOTOGRAPHS}

\vspace{0.15cm}

Figures shall be clear, legible, and directly relevant to the report
content.

Figures shall be inserted using the standard \verb|figure| environment.

Example:

\begin{verbatim}
\begin{figure}[H]
    \centering
    \includegraphics[width=0.8\textwidth]
                    {figures/company_structure.pdf}
    \caption{Organizational structure of the company}
    \label{fig:company-structure}
\end{figure}
\end{verbatim}

Figure captions shall be placed below the figure.

Each referenced figure shall have a unique \verb|\label| and shall be
referenced in the text using, for example:

\begin{verbatim}
Figure~\ref{fig:company-structure}
\end{verbatim}

Students shall not manually type figure numbers.

Figures reproduced or adapted from external sources shall include an
appropriate citation.

Photographs taken during the internship shall be relevant and sufficiently
clear. Students shall not include confidential, restricted, or proprietary
information without permission from the internship organization.

Vector graphics or high-resolution images are recommended whenever
possible.

\vspace{0.3cm}

% ================================================================
% XIII. TABLES
% ================================================================

\noindent
\textbf{XIII.\quad TABLES}

\vspace{0.15cm}

Tables shall be created using suitable \LaTeX{} environments such as
\verb|table|, \verb|tabular|, \verb|tabularx|, or \verb|longtable|.

Table captions shall be placed above the table.

Example:

\begin{verbatim}
\begin{table}[H]
    \centering
    \caption{Internship schedule}
    \label{tab:internship-schedule}

    \begin{tabular}{lll}
        \toprule
        Week & Task & Result \\
        \midrule
        1 & Orientation & Completed \\
        2 & System study & Completed \\
        \bottomrule
    \end{tabular}
\end{table}
\end{verbatim}

Tables shall be referenced using:

\begin{verbatim}
Table~\ref{tab:internship-schedule}
\end{verbatim}

Students shall not manually type table numbers.

Tables should be created directly in \LaTeX{} rather than inserted as
screenshots whenever practical.

\vspace{0.3cm}

% ================================================================
% XIV. EQUATIONS AND TECHNICAL QUANTITIES
% ================================================================

\noindent
\textbf{XIV.\quad EQUATIONS AND TECHNICAL QUANTITIES}

\vspace{0.15cm}

Mathematical expressions shall be typeset directly using \LaTeX{}.
Images or screenshots of equations shall not be used.

Short expressions may be written in inline mathematical mode, for example:

\begin{verbatim}
$R = 0.7\,\Omega$
\end{verbatim}

Important equations shall use a numbered equation environment.

Example:

\begin{verbatim}
\begin{equation}
    P = UI
    \label{eq:power}
\end{equation}
\end{verbatim}

Equations shall be referenced using:

\begin{verbatim}
Equation~\eqref{eq:power}
\end{verbatim}

Equation numbers shall not be manually entered.

SI units shall be used wherever applicable. The \verb|siunitx| package
should be used where appropriate.

Examples:

\begin{verbatim}
\SI{220}{\volt}
\SI{50}{\hertz}
\SI{2.5}{\ampere}
\end{verbatim}

Notation and units shall be consistent throughout the report.

\vspace{0.3cm}

% ================================================================
% XV. SOURCE CODE
% ================================================================

\noindent
\textbf{XV.\quad SOURCE CODE}

\vspace{0.15cm}

If the internship involves programming, PLCs, microcontrollers,
simulation, or software development, relevant source code may be included
in the report.

Source code shall be presented using the \verb|listings| environment or
another environment defined by the official template.

Example:

\begin{verbatim}
\begin{lstlisting}[language=st]
IF Start = TRUE THEN
    Motor := TRUE;
ELSE
    Motor := FALSE;
END_IF;
\end{lstlisting}
\end{verbatim}

Only code necessary for explaining the work should appear in the main
report. Long source-code listings should be moved to an appendix or
submitted as supplementary files if required.

Source code should not normally be inserted as screenshots.

\vspace{0.3cm}

% ================================================================
% XVI. ABSTRACT
% ================================================================

\noindent
\textbf{XVI.\quad ABSTRACT}

\vspace{0.15cm}

The Abstract shall briefly summarize:

\begin{itemize}[leftmargin=0.8cm, itemsep=2pt]
    \item the internship organization;
    \item the objectives of the internship;
    \item the main tasks performed;
    \item the main results, knowledge, and skills obtained.
\end{itemize}

The Abstract should be concise and should not reproduce the full content
of the report chapter by chapter.

\vspace{0.3cm}

% ================================================================
% XVII. REFERENCES AND CITATIONS
% ================================================================

\noindent
\textbf{XVII.\quad REFERENCES AND CITATIONS}

\vspace{0.15cm}

All books, journal papers, conference papers, standards, websites,
datasheets, manuals, technical documents, figures, tables, or other
external sources used in preparing the report shall be properly cited.

References shall be managed through the bibliography system provided with
the official template.

Bibliographic entries shall normally be stored in:

\begin{verbatim}
References.bib
\end{verbatim}

References shall be cited in the report using:

\begin{verbatim}
\cite{reference-key}
\end{verbatim}

Students shall not manually type reference numbers such as ``[1]'',
``[2]'', etc.

Titles of foreign-language publications shall remain in their original
language in the bibliography.

All borrowed ideas, data, figures, tables, equations, methods, or results
shall be properly acknowledged.

\vspace{0.3cm}

% ================================================================
% XVIII. CROSS-REFERENCING
% ================================================================

\noindent
\textbf{XVIII.\quad CROSS-REFERENCING}

\vspace{0.15cm}

Numbered elements shall be cross-referenced using the standard \LaTeX{}
mechanism.

Use the following commands where appropriate:

\begin{verbatim}
\label{...}
\ref{...}
\eqref{...}
\end{verbatim}

For example:

\begin{verbatim}
As shown in Figure~\ref{fig:system}, ...
\end{verbatim}

Students shall not manually type numbers such as ``Figure 2.3'' when
those numbers can be generated automatically.

\vspace{0.3cm}

% ================================================================
% XIX. PAGE NUMBERING
% ================================================================

\noindent
\textbf{XIX.\quad PAGE NUMBERING}

\vspace{0.15cm}

Cover pages and official evaluation forms at the beginning of the report
shall not display page numbers.

The main body of the report shall use Arabic page numbers beginning from
Chapter 1.

Page numbering shall be generated automatically by the official template.
Students shall not manually enter page numbers.

\vspace{0.3cm}

% ================================================================
% XX. PROJECT FILE ORGANIZATION
% ================================================================

\noindent
\textbf{XX.\quad PROJECT FILE ORGANIZATION}

\vspace{0.15cm}

Students should maintain a clear and organized \LaTeX{} project
structure. A recommended structure is:

\begin{verbatim}
InternshipReport/
|
|-- main.tex
|-- References.bib
|
|-- figures/
|   |-- company.pdf
|   |-- system.png
|   `-- ...
|
`-- chapters/
    |-- chapter1.tex
    |-- chapter2.tex
    `-- chapter3.tex
\end{verbatim}

Figures should normally be stored in the \verb|figures| directory
specified by the official template.

File and directory names should preferably contain Latin letters,
numbers, hyphens, or underscores.

Recommended examples include:

\begin{verbatim}
company_structure.pdf
motor_control.png
chapter_01.tex
\end{verbatim}

Spaces and unnecessary special characters in file names should be
avoided where possible.

\vspace{0.3cm}

% ================================================================
% XXI. DOCUMENT COMPILATION AND FINAL CHECK
% ================================================================

\noindent
\textbf{XXI.\quad DOCUMENT COMPILATION AND FINAL CHECK}

\vspace{0.15cm}

Before submission, the complete \LaTeX{} project shall compile
successfully and produce the final PDF report.

The final document shall not contain unresolved references such as:

\begin{verbatim}
??
[?]
\end{verbatim}

Before generating the final PDF, students shall check:

\begin{itemize}[leftmargin=0.8cm, itemsep=2pt]
    \item the Table of Contents;
    \item page numbering;
    \item chapter, section, and subsection numbering;
    \item figure numbering;
    \item table numbering;
    \item equation numbering;
    \item cross-references;
    \item citations;
    \item bibliography;
    \item font display;
    \item mathematical symbols;
    \item image quality;
    \item completeness of all required forms.
\end{itemize}

The document may be compiled using TeX Live, MiKTeX, or another
compatible \LaTeX{} environment. Regardless of the compilation
environment, the resulting PDF shall conform to the official Faculty
template.

\vspace{0.3cm}

% ================================================================
% XXII. CONTENT THAT STUDENTS MAY EDIT
% ================================================================

\noindent
\textbf{XXII.\quad CONTENT THAT STUDENTS MAY EDIT}

\vspace{0.15cm}

Students may edit or replace:

\begin{itemize}[leftmargin=0.8cm, itemsep=2pt]
    \item student name;
    \item class;
    \item internship organization;
    \item supervisor name;
    \item Abstract;
    \item List of Symbols and Abbreviations;
    \item chapter titles and chapter content;
    \item figures and photographs;
    \item tables;
    \item equations;
    \item citations and references.
\end{itemize}

Students shall not modify the main document formatting settings in the
preamble, including the document class, page geometry, line spacing,
heading formatting, table-of-contents formatting, font configuration,
page layout, or cover-page macros unless specifically instructed.

\vspace{0.3cm}

% ================================================================
% XXIII. FINAL SUBMISSION REQUIREMENTS
% ================================================================

\noindent
\textbf{XXIII.\quad FINAL SUBMISSION REQUIREMENTS}

\vspace{0.15cm}

Before submission, the final report shall:

\begin{itemize}[leftmargin=0.8cm, itemsep=2pt]
    \item compile without errors that affect the report content or presentation;
    \item contain all required forms and report sections;
    \item contain correct student, class, internship organization, and supervisor information;
    \item contain an updated Table of Contents, List of Figures, and List of Tables;
    \item contain no unused example content from the template;
    \item contain no placeholder text such as \textit{Title of Chapter 1},
    \textit{Section 1.1}, or other sample entries;
    \item contain all required citations and references;
    \item contain no unresolved cross-references.
\end{itemize}

If the Faculty requires submission of the \LaTeX{} source, the complete
project shall include all \verb|.tex| files, the bibliography database,
figures, and any other files required to reproduce the submitted PDF.

\vspace{0.4cm}

\noindent
\textbf{Important:}
Before submitting the final internship report, students shall remove or
comment out the following command in the main document:

\begin{verbatim}
\input{Regulations.tex}
\end{verbatim}

For example:

\begin{verbatim}
% \input{Regulations.tex}
\end{verbatim}

This removes the guidance pages from the final submitted report while
preserving the regulations file in the \LaTeX{} project for reference.

} % End 13pt content

\clearpage

\end{verbatim}

For example:

\begin{verbatim}
% 
% =================================================================
% REGULATIONS FOR PREPARING THE GRADUATION INTERNSHIP REPORT
% USING LATEX
% =================================================================
%
% This file is intended to be included from the main document with:
%
%     \input{Regulations.tex}
%
% IMPORTANT:
% This section is provided for guidance only. Students shall remove
% or comment out the \input{Regulations.tex} command before submitting
% the final internship report.
%
% =================================================================

\clearpage
\thispagestyle{empty}

\begin{center}
    {\fontsize{16pt}{20pt}\selectfont\bfseries
    SPECIFIC REGULATIONS FOR PREPARING\\
    GRADUATION INTERNSHIP REPORTS USING \LaTeX}
\end{center}

\vspace{0.5cm}

{\fontsize{13pt}{17pt}\selectfont
\setstretch{1.3}
\justifying

\noindent
\textbf{Note:}
This section is provided for guidance only and is included in the main
document using the command \verb|\input{Regulations.tex}|.
Students shall remove or comment out this command before submitting the
final internship report.

\vspace{0.4cm}

% ================================================================
% I. GENERAL REQUIREMENTS
% ================================================================

\noindent
\textbf{I.\quad GENERAL REQUIREMENTS}

\vspace{0.15cm}

The Graduation Internship Report shall be prepared using the official
\LaTeX{} template provided by the Faculty. Students shall preserve the
predefined page layout, typography, line spacing, heading styles,
numbering system, figure and table formatting, and bibliography settings
unless otherwise instructed by the Faculty or the supervisor.

The report shall be written in English.

Students are responsible for the academic and technical content of the
report. The official \LaTeX{} template is responsible for standardized
document formatting.

\vspace{0.3cm}

% ================================================================
% II. REPORT STRUCTURE
% ================================================================

\noindent
\textbf{II.\quad REPORT STRUCTURE}

\vspace{0.15cm}

The report shall normally be arranged in the following order:

\begin{enumerate}[label=\arabic*., leftmargin=1.2cm, itemsep=2pt]
    \item Outer Cover;
    \item Comments of the Internship Host Organization;
    \item Supervisor's Evaluation and Graduation Internship Defense Results;
    \item Table of Contents;
    \item List of Symbols and Abbreviations, if applicable;
    \item List of Figures;
    \item List of Tables;
    \item Abstract;
    \item Main Chapters of the Internship Report;
    \item References;
    \item Appendices, if applicable.
\end{enumerate}

The preliminary pages shall not be numbered as chapters. The main report
content shall be organized into consecutively numbered chapters beginning
with Chapter 1.

\vspace{0.3cm}

% ================================================================
% III. PAGE SIZE AND MARGINS
% ================================================================

\noindent
\textbf{III.\quad PAGE SIZE AND MARGINS}

\vspace{0.15cm}

The report shall be formatted on A4 paper with the following margins:

\begin{center}
\begin{tabular}{ll}
    Left margin:   & 3.5 cm \\
    Right margin:  & 2.0 cm \\
    Top margin:    & 2.5 cm \\
    Bottom margin: & 2.5 cm
\end{tabular}
\end{center}

These settings are predefined in the official template and shall not be
modified by students.

\vspace{0.3cm}

% ================================================================
% IV. FONT, FONT SIZE, AND LINE SPACING
% ================================================================

\noindent
\textbf{IV.\quad FONT, FONT SIZE, AND LINE SPACING}

\vspace{0.15cm}

The main text shall use a 13 pt font size with 1.3-line spacing.
A Times-family typeface shall be used in accordance with the official
template.

The text shall be fully justified.

Students shall not use repeated spaces, repeated manual line breaks, or
multiple blank lines to position text. Paragraph and heading spacing shall
be controlled by the template.

\vspace{0.3cm}

% ================================================================
% V. OUTER COVER
% ================================================================

\noindent
\textbf{V.\quad OUTER COVER}

\vspace{0.15cm}

The outer cover shall follow the official format provided in the template
and shall contain the following information:

\begin{itemize}[leftmargin=0.8cm, itemsep=2pt]
    \item Thai Nguyen University of Technology;
    \item Faculty of International Training;
    \item Engineering and Technology Department;
    \item University logo and Faculty logo;
    \item the title \textbf{INTERNSHIP REPORT};
    \item Internship Organization;
    \item Author;
    \item Class;
    \item Supervisor;
    \item place and year of completion.
\end{itemize}

Students shall only replace the required information fields and shall not
modify the cover-page layout unless instructed otherwise.

\vspace{0.3cm}

% ================================================================
% VI. COMMENTS OF THE INTERNSHIP HOST ORGANIZATION
% ================================================================

\noindent
\textbf{VI.\quad COMMENTS OF THE INTERNSHIP HOST ORGANIZATION}

\vspace{0.15cm}

The page entitled \textbf{COMMENTS OF THE INTERNSHIP HOST ORGANIZATION}
is reserved for the internship host organization.

It includes:

\begin{itemize}[leftmargin=0.8cm, itemsep=2pt]
    \item the student's full name;
    \item class;
    \item comments from the host organization;
    \item score;
    \item date;
    \item signature and official seal.
\end{itemize}

Students shall not alter or complete official evaluation fields unless
specifically authorized.

\vspace{0.3cm}

% ================================================================
% VII. SUPERVISOR'S EVALUATION AND DEFENSE RESULTS
% ================================================================

\noindent
\textbf{VII.\quad SUPERVISOR'S EVALUATION AND DEFENSE RESULTS}

\vspace{0.15cm}

The section entitled
\textbf{SUPERVISOR'S EVALUATION AND GRADUATION INTERNSHIP DEFENSE RESULTS}
is reserved for the supervisor and defense committee.

The evaluation may include:

\begin{itemize}[leftmargin=0.8cm, itemsep=2pt]
    \item Internship Progress and Student's Attitude;
    \item Report Content;
    \item Report Presentation;
    \item Other Comments;
    \item Supervisor's Opinion;
    \item Defense Result;
    \item Score.
\end{itemize}

Students shall not enter evaluation scores or defense results before they
are officially determined.

\vspace{0.3cm}

% ================================================================
% VIII. CHAPTERS, SECTIONS, AND SUBSECTIONS
% ================================================================

\noindent
\textbf{VIII.\quad CHAPTERS, SECTIONS, AND SUBSECTIONS}

\vspace{0.15cm}

The report structure shall be generated using the standard \LaTeX{}
commands:

\begin{verbatim}
\chapter{...}
\section{...}
\subsection{...}
\end{verbatim}

Chapter, section, and subsection numbers shall be generated automatically
by \LaTeX{}.

Students shall not manually type chapter or section numbers and shall not
reproduce heading styles using commands such as \verb|\textbf|,
\verb|\fontsize|, or \verb|\centering| instead of the appropriate
structural command.

Excessive subdivision should be avoided. In most cases, the hierarchy
Chapter--Section--Subsection is sufficient.

\vspace{0.3cm}

% ================================================================
% IX. ORGANIZATION OF THE MAIN CONTENT
% ================================================================

\noindent
\textbf{IX.\quad ORGANIZATION OF THE MAIN CONTENT}

\vspace{0.15cm}

The specific content of the report shall depend on the internship
organization and the assigned tasks. A typical structure may include:

\begin{itemize}[leftmargin=0.8cm, itemsep=2pt]
    \item \textbf{Chapter 1. Overview of the Internship Organization};
    \item \textbf{Chapter 2. Internship Activities and Assigned Tasks};
    \item \textbf{Chapter 3. Results, Experience, and Evaluation}.
\end{itemize}

These chapter titles are recommendations only and may be adapted to the
actual internship content.

The report shall reflect work actually performed, observed, analyzed, or
learned during the internship. Unrelated theoretical material should not
be included merely to increase report length.

\vspace{0.3cm}

% ================================================================
% X. TABLE OF CONTENTS
% ================================================================

\noindent
\textbf{X.\quad TABLE OF CONTENTS}

\vspace{0.15cm}

The Table of Contents shall be generated automatically using:

\begin{verbatim}
\tableofcontents
\end{verbatim}

Students shall not create or manually edit page numbers in the Table of
Contents.

The document shall be compiled sufficiently to ensure that all page
numbers and entries are correctly updated.

\vspace{0.3cm}

% ================================================================
% XI. LIST OF SYMBOLS AND ABBREVIATIONS
% ================================================================

\noindent
\textbf{XI.\quad LIST OF SYMBOLS AND ABBREVIATIONS}

\vspace{0.15cm}

Symbols and abbreviations that are used repeatedly in the report shall be
listed in the \textbf{List of Symbols and Abbreviations}.

The list may contain the following columns:

\begin{center}
\begin{tabular}{llll}
    No. & Symbol / Abbreviation & Unit & Definition
\end{tabular}
\end{center}

Only symbols and abbreviations actually used in the final report shall be
retained. Example entries contained in the template shall be removed if
they are not relevant to the student's report.

Each abbreviation should normally be defined when it first appears in
the main text.

\vspace{0.3cm}

% ================================================================
% XII. FIGURES, DIAGRAMS, AND PHOTOGRAPHS
% ================================================================

\noindent
\textbf{XII.\quad FIGURES, DIAGRAMS, AND PHOTOGRAPHS}

\vspace{0.15cm}

Figures shall be clear, legible, and directly relevant to the report
content.

Figures shall be inserted using the standard \verb|figure| environment.

Example:

\begin{verbatim}
\begin{figure}[H]
    \centering
    \includegraphics[width=0.8\textwidth]
                    {figures/company_structure.pdf}
    \caption{Organizational structure of the company}
    \label{fig:company-structure}
\end{figure}
\end{verbatim}

Figure captions shall be placed below the figure.

Each referenced figure shall have a unique \verb|\label| and shall be
referenced in the text using, for example:

\begin{verbatim}
Figure~\ref{fig:company-structure}
\end{verbatim}

Students shall not manually type figure numbers.

Figures reproduced or adapted from external sources shall include an
appropriate citation.

Photographs taken during the internship shall be relevant and sufficiently
clear. Students shall not include confidential, restricted, or proprietary
information without permission from the internship organization.

Vector graphics or high-resolution images are recommended whenever
possible.

\vspace{0.3cm}

% ================================================================
% XIII. TABLES
% ================================================================

\noindent
\textbf{XIII.\quad TABLES}

\vspace{0.15cm}

Tables shall be created using suitable \LaTeX{} environments such as
\verb|table|, \verb|tabular|, \verb|tabularx|, or \verb|longtable|.

Table captions shall be placed above the table.

Example:

\begin{verbatim}
\begin{table}[H]
    \centering
    \caption{Internship schedule}
    \label{tab:internship-schedule}

    \begin{tabular}{lll}
        \toprule
        Week & Task & Result \\
        \midrule
        1 & Orientation & Completed \\
        2 & System study & Completed \\
        \bottomrule
    \end{tabular}
\end{table}
\end{verbatim}

Tables shall be referenced using:

\begin{verbatim}
Table~\ref{tab:internship-schedule}
\end{verbatim}

Students shall not manually type table numbers.

Tables should be created directly in \LaTeX{} rather than inserted as
screenshots whenever practical.

\vspace{0.3cm}

% ================================================================
% XIV. EQUATIONS AND TECHNICAL QUANTITIES
% ================================================================

\noindent
\textbf{XIV.\quad EQUATIONS AND TECHNICAL QUANTITIES}

\vspace{0.15cm}

Mathematical expressions shall be typeset directly using \LaTeX{}.
Images or screenshots of equations shall not be used.

Short expressions may be written in inline mathematical mode, for example:

\begin{verbatim}
$R = 0.7\,\Omega$
\end{verbatim}

Important equations shall use a numbered equation environment.

Example:

\begin{verbatim}
\begin{equation}
    P = UI
    \label{eq:power}
\end{equation}
\end{verbatim}

Equations shall be referenced using:

\begin{verbatim}
Equation~\eqref{eq:power}
\end{verbatim}

Equation numbers shall not be manually entered.

SI units shall be used wherever applicable. The \verb|siunitx| package
should be used where appropriate.

Examples:

\begin{verbatim}
\SI{220}{\volt}
\SI{50}{\hertz}
\SI{2.5}{\ampere}
\end{verbatim}

Notation and units shall be consistent throughout the report.

\vspace{0.3cm}

% ================================================================
% XV. SOURCE CODE
% ================================================================

\noindent
\textbf{XV.\quad SOURCE CODE}

\vspace{0.15cm}

If the internship involves programming, PLCs, microcontrollers,
simulation, or software development, relevant source code may be included
in the report.

Source code shall be presented using the \verb|listings| environment or
another environment defined by the official template.

Example:

\begin{verbatim}
\begin{lstlisting}[language=st]
IF Start = TRUE THEN
    Motor := TRUE;
ELSE
    Motor := FALSE;
END_IF;
\end{lstlisting}
\end{verbatim}

Only code necessary for explaining the work should appear in the main
report. Long source-code listings should be moved to an appendix or
submitted as supplementary files if required.

Source code should not normally be inserted as screenshots.

\vspace{0.3cm}

% ================================================================
% XVI. ABSTRACT
% ================================================================

\noindent
\textbf{XVI.\quad ABSTRACT}

\vspace{0.15cm}

The Abstract shall briefly summarize:

\begin{itemize}[leftmargin=0.8cm, itemsep=2pt]
    \item the internship organization;
    \item the objectives of the internship;
    \item the main tasks performed;
    \item the main results, knowledge, and skills obtained.
\end{itemize}

The Abstract should be concise and should not reproduce the full content
of the report chapter by chapter.

\vspace{0.3cm}

% ================================================================
% XVII. REFERENCES AND CITATIONS
% ================================================================

\noindent
\textbf{XVII.\quad REFERENCES AND CITATIONS}

\vspace{0.15cm}

All books, journal papers, conference papers, standards, websites,
datasheets, manuals, technical documents, figures, tables, or other
external sources used in preparing the report shall be properly cited.

References shall be managed through the bibliography system provided with
the official template.

Bibliographic entries shall normally be stored in:

\begin{verbatim}
References.bib
\end{verbatim}

References shall be cited in the report using:

\begin{verbatim}
\cite{reference-key}
\end{verbatim}

Students shall not manually type reference numbers such as ``[1]'',
``[2]'', etc.

Titles of foreign-language publications shall remain in their original
language in the bibliography.

All borrowed ideas, data, figures, tables, equations, methods, or results
shall be properly acknowledged.

\vspace{0.3cm}

% ================================================================
% XVIII. CROSS-REFERENCING
% ================================================================

\noindent
\textbf{XVIII.\quad CROSS-REFERENCING}

\vspace{0.15cm}

Numbered elements shall be cross-referenced using the standard \LaTeX{}
mechanism.

Use the following commands where appropriate:

\begin{verbatim}
\label{...}
\ref{...}
\eqref{...}
\end{verbatim}

For example:

\begin{verbatim}
As shown in Figure~\ref{fig:system}, ...
\end{verbatim}

Students shall not manually type numbers such as ``Figure 2.3'' when
those numbers can be generated automatically.

\vspace{0.3cm}

% ================================================================
% XIX. PAGE NUMBERING
% ================================================================

\noindent
\textbf{XIX.\quad PAGE NUMBERING}

\vspace{0.15cm}

Cover pages and official evaluation forms at the beginning of the report
shall not display page numbers.

The main body of the report shall use Arabic page numbers beginning from
Chapter 1.

Page numbering shall be generated automatically by the official template.
Students shall not manually enter page numbers.

\vspace{0.3cm}

% ================================================================
% XX. PROJECT FILE ORGANIZATION
% ================================================================

\noindent
\textbf{XX.\quad PROJECT FILE ORGANIZATION}

\vspace{0.15cm}

Students should maintain a clear and organized \LaTeX{} project
structure. A recommended structure is:

\begin{verbatim}
InternshipReport/
|
|-- main.tex
|-- References.bib
|
|-- figures/
|   |-- company.pdf
|   |-- system.png
|   `-- ...
|
`-- chapters/
    |-- chapter1.tex
    |-- chapter2.tex
    `-- chapter3.tex
\end{verbatim}

Figures should normally be stored in the \verb|figures| directory
specified by the official template.

File and directory names should preferably contain Latin letters,
numbers, hyphens, or underscores.

Recommended examples include:

\begin{verbatim}
company_structure.pdf
motor_control.png
chapter_01.tex
\end{verbatim}

Spaces and unnecessary special characters in file names should be
avoided where possible.

\vspace{0.3cm}

% ================================================================
% XXI. DOCUMENT COMPILATION AND FINAL CHECK
% ================================================================

\noindent
\textbf{XXI.\quad DOCUMENT COMPILATION AND FINAL CHECK}

\vspace{0.15cm}

Before submission, the complete \LaTeX{} project shall compile
successfully and produce the final PDF report.

The final document shall not contain unresolved references such as:

\begin{verbatim}
??
[?]
\end{verbatim}

Before generating the final PDF, students shall check:

\begin{itemize}[leftmargin=0.8cm, itemsep=2pt]
    \item the Table of Contents;
    \item page numbering;
    \item chapter, section, and subsection numbering;
    \item figure numbering;
    \item table numbering;
    \item equation numbering;
    \item cross-references;
    \item citations;
    \item bibliography;
    \item font display;
    \item mathematical symbols;
    \item image quality;
    \item completeness of all required forms.
\end{itemize}

The document may be compiled using TeX Live, MiKTeX, or another
compatible \LaTeX{} environment. Regardless of the compilation
environment, the resulting PDF shall conform to the official Faculty
template.

\vspace{0.3cm}

% ================================================================
% XXII. CONTENT THAT STUDENTS MAY EDIT
% ================================================================

\noindent
\textbf{XXII.\quad CONTENT THAT STUDENTS MAY EDIT}

\vspace{0.15cm}

Students may edit or replace:

\begin{itemize}[leftmargin=0.8cm, itemsep=2pt]
    \item student name;
    \item class;
    \item internship organization;
    \item supervisor name;
    \item Abstract;
    \item List of Symbols and Abbreviations;
    \item chapter titles and chapter content;
    \item figures and photographs;
    \item tables;
    \item equations;
    \item citations and references.
\end{itemize}

Students shall not modify the main document formatting settings in the
preamble, including the document class, page geometry, line spacing,
heading formatting, table-of-contents formatting, font configuration,
page layout, or cover-page macros unless specifically instructed.

\vspace{0.3cm}

% ================================================================
% XXIII. FINAL SUBMISSION REQUIREMENTS
% ================================================================

\noindent
\textbf{XXIII.\quad FINAL SUBMISSION REQUIREMENTS}

\vspace{0.15cm}

Before submission, the final report shall:

\begin{itemize}[leftmargin=0.8cm, itemsep=2pt]
    \item compile without errors that affect the report content or presentation;
    \item contain all required forms and report sections;
    \item contain correct student, class, internship organization, and supervisor information;
    \item contain an updated Table of Contents, List of Figures, and List of Tables;
    \item contain no unused example content from the template;
    \item contain no placeholder text such as \textit{Title of Chapter 1},
    \textit{Section 1.1}, or other sample entries;
    \item contain all required citations and references;
    \item contain no unresolved cross-references.
\end{itemize}

If the Faculty requires submission of the \LaTeX{} source, the complete
project shall include all \verb|.tex| files, the bibliography database,
figures, and any other files required to reproduce the submitted PDF.

\vspace{0.4cm}

\noindent
\textbf{Important:}
Before submitting the final internship report, students shall remove or
comment out the following command in the main document:

\begin{verbatim}
\input{Regulations.tex}
\end{verbatim}

For example:

\begin{verbatim}
% \input{Regulations.tex}
\end{verbatim}

This removes the guidance pages from the final submitted report while
preserving the regulations file in the \LaTeX{} project for reference.

} % End 13pt content

\clearpage

\end{verbatim}

This removes the guidance pages from the final submitted report while
preserving the regulations file in the \LaTeX{} project for reference.

} % End 13pt content

\clearpage

\end{verbatim}

For example:

\begin{verbatim}
% 
% =================================================================
% REGULATIONS FOR PREPARING THE GRADUATION INTERNSHIP REPORT
% USING LATEX
% =================================================================
%
% This file is intended to be included from the main document with:
%
%     
% =================================================================
% REGULATIONS FOR PREPARING THE GRADUATION INTERNSHIP REPORT
% USING LATEX
% =================================================================
%
% This file is intended to be included from the main document with:
%
%     \input{Regulations.tex}
%
% IMPORTANT:
% This section is provided for guidance only. Students shall remove
% or comment out the \input{Regulations.tex} command before submitting
% the final internship report.
%
% =================================================================

\clearpage
\thispagestyle{empty}

\begin{center}
    {\fontsize{16pt}{20pt}\selectfont\bfseries
    SPECIFIC REGULATIONS FOR PREPARING\\
    GRADUATION INTERNSHIP REPORTS USING \LaTeX}
\end{center}

\vspace{0.5cm}

{\fontsize{13pt}{17pt}\selectfont
\setstretch{1.3}
\justifying

\noindent
\textbf{Note:}
This section is provided for guidance only and is included in the main
document using the command \verb|\input{Regulations.tex}|.
Students shall remove or comment out this command before submitting the
final internship report.

\vspace{0.4cm}

% ================================================================
% I. GENERAL REQUIREMENTS
% ================================================================

\noindent
\textbf{I.\quad GENERAL REQUIREMENTS}

\vspace{0.15cm}

The Graduation Internship Report shall be prepared using the official
\LaTeX{} template provided by the Faculty. Students shall preserve the
predefined page layout, typography, line spacing, heading styles,
numbering system, figure and table formatting, and bibliography settings
unless otherwise instructed by the Faculty or the supervisor.

The report shall be written in English.

Students are responsible for the academic and technical content of the
report. The official \LaTeX{} template is responsible for standardized
document formatting.

\vspace{0.3cm}

% ================================================================
% II. REPORT STRUCTURE
% ================================================================

\noindent
\textbf{II.\quad REPORT STRUCTURE}

\vspace{0.15cm}

The report shall normally be arranged in the following order:

\begin{enumerate}[label=\arabic*., leftmargin=1.2cm, itemsep=2pt]
    \item Outer Cover;
    \item Comments of the Internship Host Organization;
    \item Supervisor's Evaluation and Graduation Internship Defense Results;
    \item Table of Contents;
    \item List of Symbols and Abbreviations, if applicable;
    \item List of Figures;
    \item List of Tables;
    \item Abstract;
    \item Main Chapters of the Internship Report;
    \item References;
    \item Appendices, if applicable.
\end{enumerate}

The preliminary pages shall not be numbered as chapters. The main report
content shall be organized into consecutively numbered chapters beginning
with Chapter 1.

\vspace{0.3cm}

% ================================================================
% III. PAGE SIZE AND MARGINS
% ================================================================

\noindent
\textbf{III.\quad PAGE SIZE AND MARGINS}

\vspace{0.15cm}

The report shall be formatted on A4 paper with the following margins:

\begin{center}
\begin{tabular}{ll}
    Left margin:   & 3.5 cm \\
    Right margin:  & 2.0 cm \\
    Top margin:    & 2.5 cm \\
    Bottom margin: & 2.5 cm
\end{tabular}
\end{center}

These settings are predefined in the official template and shall not be
modified by students.

\vspace{0.3cm}

% ================================================================
% IV. FONT, FONT SIZE, AND LINE SPACING
% ================================================================

\noindent
\textbf{IV.\quad FONT, FONT SIZE, AND LINE SPACING}

\vspace{0.15cm}

The main text shall use a 13 pt font size with 1.3-line spacing.
A Times-family typeface shall be used in accordance with the official
template.

The text shall be fully justified.

Students shall not use repeated spaces, repeated manual line breaks, or
multiple blank lines to position text. Paragraph and heading spacing shall
be controlled by the template.

\vspace{0.3cm}

% ================================================================
% V. OUTER COVER
% ================================================================

\noindent
\textbf{V.\quad OUTER COVER}

\vspace{0.15cm}

The outer cover shall follow the official format provided in the template
and shall contain the following information:

\begin{itemize}[leftmargin=0.8cm, itemsep=2pt]
    \item Thai Nguyen University of Technology;
    \item Faculty of International Training;
    \item Engineering and Technology Department;
    \item University logo and Faculty logo;
    \item the title \textbf{INTERNSHIP REPORT};
    \item Internship Organization;
    \item Author;
    \item Class;
    \item Supervisor;
    \item place and year of completion.
\end{itemize}

Students shall only replace the required information fields and shall not
modify the cover-page layout unless instructed otherwise.

\vspace{0.3cm}

% ================================================================
% VI. COMMENTS OF THE INTERNSHIP HOST ORGANIZATION
% ================================================================

\noindent
\textbf{VI.\quad COMMENTS OF THE INTERNSHIP HOST ORGANIZATION}

\vspace{0.15cm}

The page entitled \textbf{COMMENTS OF THE INTERNSHIP HOST ORGANIZATION}
is reserved for the internship host organization.

It includes:

\begin{itemize}[leftmargin=0.8cm, itemsep=2pt]
    \item the student's full name;
    \item class;
    \item comments from the host organization;
    \item score;
    \item date;
    \item signature and official seal.
\end{itemize}

Students shall not alter or complete official evaluation fields unless
specifically authorized.

\vspace{0.3cm}

% ================================================================
% VII. SUPERVISOR'S EVALUATION AND DEFENSE RESULTS
% ================================================================

\noindent
\textbf{VII.\quad SUPERVISOR'S EVALUATION AND DEFENSE RESULTS}

\vspace{0.15cm}

The section entitled
\textbf{SUPERVISOR'S EVALUATION AND GRADUATION INTERNSHIP DEFENSE RESULTS}
is reserved for the supervisor and defense committee.

The evaluation may include:

\begin{itemize}[leftmargin=0.8cm, itemsep=2pt]
    \item Internship Progress and Student's Attitude;
    \item Report Content;
    \item Report Presentation;
    \item Other Comments;
    \item Supervisor's Opinion;
    \item Defense Result;
    \item Score.
\end{itemize}

Students shall not enter evaluation scores or defense results before they
are officially determined.

\vspace{0.3cm}

% ================================================================
% VIII. CHAPTERS, SECTIONS, AND SUBSECTIONS
% ================================================================

\noindent
\textbf{VIII.\quad CHAPTERS, SECTIONS, AND SUBSECTIONS}

\vspace{0.15cm}

The report structure shall be generated using the standard \LaTeX{}
commands:

\begin{verbatim}
\chapter{...}
\section{...}
\subsection{...}
\end{verbatim}

Chapter, section, and subsection numbers shall be generated automatically
by \LaTeX{}.

Students shall not manually type chapter or section numbers and shall not
reproduce heading styles using commands such as \verb|\textbf|,
\verb|\fontsize|, or \verb|\centering| instead of the appropriate
structural command.

Excessive subdivision should be avoided. In most cases, the hierarchy
Chapter--Section--Subsection is sufficient.

\vspace{0.3cm}

% ================================================================
% IX. ORGANIZATION OF THE MAIN CONTENT
% ================================================================

\noindent
\textbf{IX.\quad ORGANIZATION OF THE MAIN CONTENT}

\vspace{0.15cm}

The specific content of the report shall depend on the internship
organization and the assigned tasks. A typical structure may include:

\begin{itemize}[leftmargin=0.8cm, itemsep=2pt]
    \item \textbf{Chapter 1. Overview of the Internship Organization};
    \item \textbf{Chapter 2. Internship Activities and Assigned Tasks};
    \item \textbf{Chapter 3. Results, Experience, and Evaluation}.
\end{itemize}

These chapter titles are recommendations only and may be adapted to the
actual internship content.

The report shall reflect work actually performed, observed, analyzed, or
learned during the internship. Unrelated theoretical material should not
be included merely to increase report length.

\vspace{0.3cm}

% ================================================================
% X. TABLE OF CONTENTS
% ================================================================

\noindent
\textbf{X.\quad TABLE OF CONTENTS}

\vspace{0.15cm}

The Table of Contents shall be generated automatically using:

\begin{verbatim}
\tableofcontents
\end{verbatim}

Students shall not create or manually edit page numbers in the Table of
Contents.

The document shall be compiled sufficiently to ensure that all page
numbers and entries are correctly updated.

\vspace{0.3cm}

% ================================================================
% XI. LIST OF SYMBOLS AND ABBREVIATIONS
% ================================================================

\noindent
\textbf{XI.\quad LIST OF SYMBOLS AND ABBREVIATIONS}

\vspace{0.15cm}

Symbols and abbreviations that are used repeatedly in the report shall be
listed in the \textbf{List of Symbols and Abbreviations}.

The list may contain the following columns:

\begin{center}
\begin{tabular}{llll}
    No. & Symbol / Abbreviation & Unit & Definition
\end{tabular}
\end{center}

Only symbols and abbreviations actually used in the final report shall be
retained. Example entries contained in the template shall be removed if
they are not relevant to the student's report.

Each abbreviation should normally be defined when it first appears in
the main text.

\vspace{0.3cm}

% ================================================================
% XII. FIGURES, DIAGRAMS, AND PHOTOGRAPHS
% ================================================================

\noindent
\textbf{XII.\quad FIGURES, DIAGRAMS, AND PHOTOGRAPHS}

\vspace{0.15cm}

Figures shall be clear, legible, and directly relevant to the report
content.

Figures shall be inserted using the standard \verb|figure| environment.

Example:

\begin{verbatim}
\begin{figure}[H]
    \centering
    \includegraphics[width=0.8\textwidth]
                    {figures/company_structure.pdf}
    \caption{Organizational structure of the company}
    \label{fig:company-structure}
\end{figure}
\end{verbatim}

Figure captions shall be placed below the figure.

Each referenced figure shall have a unique \verb|\label| and shall be
referenced in the text using, for example:

\begin{verbatim}
Figure~\ref{fig:company-structure}
\end{verbatim}

Students shall not manually type figure numbers.

Figures reproduced or adapted from external sources shall include an
appropriate citation.

Photographs taken during the internship shall be relevant and sufficiently
clear. Students shall not include confidential, restricted, or proprietary
information without permission from the internship organization.

Vector graphics or high-resolution images are recommended whenever
possible.

\vspace{0.3cm}

% ================================================================
% XIII. TABLES
% ================================================================

\noindent
\textbf{XIII.\quad TABLES}

\vspace{0.15cm}

Tables shall be created using suitable \LaTeX{} environments such as
\verb|table|, \verb|tabular|, \verb|tabularx|, or \verb|longtable|.

Table captions shall be placed above the table.

Example:

\begin{verbatim}
\begin{table}[H]
    \centering
    \caption{Internship schedule}
    \label{tab:internship-schedule}

    \begin{tabular}{lll}
        \toprule
        Week & Task & Result \\
        \midrule
        1 & Orientation & Completed \\
        2 & System study & Completed \\
        \bottomrule
    \end{tabular}
\end{table}
\end{verbatim}

Tables shall be referenced using:

\begin{verbatim}
Table~\ref{tab:internship-schedule}
\end{verbatim}

Students shall not manually type table numbers.

Tables should be created directly in \LaTeX{} rather than inserted as
screenshots whenever practical.

\vspace{0.3cm}

% ================================================================
% XIV. EQUATIONS AND TECHNICAL QUANTITIES
% ================================================================

\noindent
\textbf{XIV.\quad EQUATIONS AND TECHNICAL QUANTITIES}

\vspace{0.15cm}

Mathematical expressions shall be typeset directly using \LaTeX{}.
Images or screenshots of equations shall not be used.

Short expressions may be written in inline mathematical mode, for example:

\begin{verbatim}
$R = 0.7\,\Omega$
\end{verbatim}

Important equations shall use a numbered equation environment.

Example:

\begin{verbatim}
\begin{equation}
    P = UI
    \label{eq:power}
\end{equation}
\end{verbatim}

Equations shall be referenced using:

\begin{verbatim}
Equation~\eqref{eq:power}
\end{verbatim}

Equation numbers shall not be manually entered.

SI units shall be used wherever applicable. The \verb|siunitx| package
should be used where appropriate.

Examples:

\begin{verbatim}
\SI{220}{\volt}
\SI{50}{\hertz}
\SI{2.5}{\ampere}
\end{verbatim}

Notation and units shall be consistent throughout the report.

\vspace{0.3cm}

% ================================================================
% XV. SOURCE CODE
% ================================================================

\noindent
\textbf{XV.\quad SOURCE CODE}

\vspace{0.15cm}

If the internship involves programming, PLCs, microcontrollers,
simulation, or software development, relevant source code may be included
in the report.

Source code shall be presented using the \verb|listings| environment or
another environment defined by the official template.

Example:

\begin{verbatim}
\begin{lstlisting}[language=st]
IF Start = TRUE THEN
    Motor := TRUE;
ELSE
    Motor := FALSE;
END_IF;
\end{lstlisting}
\end{verbatim}

Only code necessary for explaining the work should appear in the main
report. Long source-code listings should be moved to an appendix or
submitted as supplementary files if required.

Source code should not normally be inserted as screenshots.

\vspace{0.3cm}

% ================================================================
% XVI. ABSTRACT
% ================================================================

\noindent
\textbf{XVI.\quad ABSTRACT}

\vspace{0.15cm}

The Abstract shall briefly summarize:

\begin{itemize}[leftmargin=0.8cm, itemsep=2pt]
    \item the internship organization;
    \item the objectives of the internship;
    \item the main tasks performed;
    \item the main results, knowledge, and skills obtained.
\end{itemize}

The Abstract should be concise and should not reproduce the full content
of the report chapter by chapter.

\vspace{0.3cm}

% ================================================================
% XVII. REFERENCES AND CITATIONS
% ================================================================

\noindent
\textbf{XVII.\quad REFERENCES AND CITATIONS}

\vspace{0.15cm}

All books, journal papers, conference papers, standards, websites,
datasheets, manuals, technical documents, figures, tables, or other
external sources used in preparing the report shall be properly cited.

References shall be managed through the bibliography system provided with
the official template.

Bibliographic entries shall normally be stored in:

\begin{verbatim}
References.bib
\end{verbatim}

References shall be cited in the report using:

\begin{verbatim}
\cite{reference-key}
\end{verbatim}

Students shall not manually type reference numbers such as ``[1]'',
``[2]'', etc.

Titles of foreign-language publications shall remain in their original
language in the bibliography.

All borrowed ideas, data, figures, tables, equations, methods, or results
shall be properly acknowledged.

\vspace{0.3cm}

% ================================================================
% XVIII. CROSS-REFERENCING
% ================================================================

\noindent
\textbf{XVIII.\quad CROSS-REFERENCING}

\vspace{0.15cm}

Numbered elements shall be cross-referenced using the standard \LaTeX{}
mechanism.

Use the following commands where appropriate:

\begin{verbatim}
\label{...}
\ref{...}
\eqref{...}
\end{verbatim}

For example:

\begin{verbatim}
As shown in Figure~\ref{fig:system}, ...
\end{verbatim}

Students shall not manually type numbers such as ``Figure 2.3'' when
those numbers can be generated automatically.

\vspace{0.3cm}

% ================================================================
% XIX. PAGE NUMBERING
% ================================================================

\noindent
\textbf{XIX.\quad PAGE NUMBERING}

\vspace{0.15cm}

Cover pages and official evaluation forms at the beginning of the report
shall not display page numbers.

The main body of the report shall use Arabic page numbers beginning from
Chapter 1.

Page numbering shall be generated automatically by the official template.
Students shall not manually enter page numbers.

\vspace{0.3cm}

% ================================================================
% XX. PROJECT FILE ORGANIZATION
% ================================================================

\noindent
\textbf{XX.\quad PROJECT FILE ORGANIZATION}

\vspace{0.15cm}

Students should maintain a clear and organized \LaTeX{} project
structure. A recommended structure is:

\begin{verbatim}
InternshipReport/
|
|-- main.tex
|-- References.bib
|
|-- figures/
|   |-- company.pdf
|   |-- system.png
|   `-- ...
|
`-- chapters/
    |-- chapter1.tex
    |-- chapter2.tex
    `-- chapter3.tex
\end{verbatim}

Figures should normally be stored in the \verb|figures| directory
specified by the official template.

File and directory names should preferably contain Latin letters,
numbers, hyphens, or underscores.

Recommended examples include:

\begin{verbatim}
company_structure.pdf
motor_control.png
chapter_01.tex
\end{verbatim}

Spaces and unnecessary special characters in file names should be
avoided where possible.

\vspace{0.3cm}

% ================================================================
% XXI. DOCUMENT COMPILATION AND FINAL CHECK
% ================================================================

\noindent
\textbf{XXI.\quad DOCUMENT COMPILATION AND FINAL CHECK}

\vspace{0.15cm}

Before submission, the complete \LaTeX{} project shall compile
successfully and produce the final PDF report.

The final document shall not contain unresolved references such as:

\begin{verbatim}
??
[?]
\end{verbatim}

Before generating the final PDF, students shall check:

\begin{itemize}[leftmargin=0.8cm, itemsep=2pt]
    \item the Table of Contents;
    \item page numbering;
    \item chapter, section, and subsection numbering;
    \item figure numbering;
    \item table numbering;
    \item equation numbering;
    \item cross-references;
    \item citations;
    \item bibliography;
    \item font display;
    \item mathematical symbols;
    \item image quality;
    \item completeness of all required forms.
\end{itemize}

The document may be compiled using TeX Live, MiKTeX, or another
compatible \LaTeX{} environment. Regardless of the compilation
environment, the resulting PDF shall conform to the official Faculty
template.

\vspace{0.3cm}

% ================================================================
% XXII. CONTENT THAT STUDENTS MAY EDIT
% ================================================================

\noindent
\textbf{XXII.\quad CONTENT THAT STUDENTS MAY EDIT}

\vspace{0.15cm}

Students may edit or replace:

\begin{itemize}[leftmargin=0.8cm, itemsep=2pt]
    \item student name;
    \item class;
    \item internship organization;
    \item supervisor name;
    \item Abstract;
    \item List of Symbols and Abbreviations;
    \item chapter titles and chapter content;
    \item figures and photographs;
    \item tables;
    \item equations;
    \item citations and references.
\end{itemize}

Students shall not modify the main document formatting settings in the
preamble, including the document class, page geometry, line spacing,
heading formatting, table-of-contents formatting, font configuration,
page layout, or cover-page macros unless specifically instructed.

\vspace{0.3cm}

% ================================================================
% XXIII. FINAL SUBMISSION REQUIREMENTS
% ================================================================

\noindent
\textbf{XXIII.\quad FINAL SUBMISSION REQUIREMENTS}

\vspace{0.15cm}

Before submission, the final report shall:

\begin{itemize}[leftmargin=0.8cm, itemsep=2pt]
    \item compile without errors that affect the report content or presentation;
    \item contain all required forms and report sections;
    \item contain correct student, class, internship organization, and supervisor information;
    \item contain an updated Table of Contents, List of Figures, and List of Tables;
    \item contain no unused example content from the template;
    \item contain no placeholder text such as \textit{Title of Chapter 1},
    \textit{Section 1.1}, or other sample entries;
    \item contain all required citations and references;
    \item contain no unresolved cross-references.
\end{itemize}

If the Faculty requires submission of the \LaTeX{} source, the complete
project shall include all \verb|.tex| files, the bibliography database,
figures, and any other files required to reproduce the submitted PDF.

\vspace{0.4cm}

\noindent
\textbf{Important:}
Before submitting the final internship report, students shall remove or
comment out the following command in the main document:

\begin{verbatim}
\input{Regulations.tex}
\end{verbatim}

For example:

\begin{verbatim}
% \input{Regulations.tex}
\end{verbatim}

This removes the guidance pages from the final submitted report while
preserving the regulations file in the \LaTeX{} project for reference.

} % End 13pt content

\clearpage

%
% IMPORTANT:
% This section is provided for guidance only. Students shall remove
% or comment out the 
% =================================================================
% REGULATIONS FOR PREPARING THE GRADUATION INTERNSHIP REPORT
% USING LATEX
% =================================================================
%
% This file is intended to be included from the main document with:
%
%     \input{Regulations.tex}
%
% IMPORTANT:
% This section is provided for guidance only. Students shall remove
% or comment out the \input{Regulations.tex} command before submitting
% the final internship report.
%
% =================================================================

\clearpage
\thispagestyle{empty}

\begin{center}
    {\fontsize{16pt}{20pt}\selectfont\bfseries
    SPECIFIC REGULATIONS FOR PREPARING\\
    GRADUATION INTERNSHIP REPORTS USING \LaTeX}
\end{center}

\vspace{0.5cm}

{\fontsize{13pt}{17pt}\selectfont
\setstretch{1.3}
\justifying

\noindent
\textbf{Note:}
This section is provided for guidance only and is included in the main
document using the command \verb|\input{Regulations.tex}|.
Students shall remove or comment out this command before submitting the
final internship report.

\vspace{0.4cm}

% ================================================================
% I. GENERAL REQUIREMENTS
% ================================================================

\noindent
\textbf{I.\quad GENERAL REQUIREMENTS}

\vspace{0.15cm}

The Graduation Internship Report shall be prepared using the official
\LaTeX{} template provided by the Faculty. Students shall preserve the
predefined page layout, typography, line spacing, heading styles,
numbering system, figure and table formatting, and bibliography settings
unless otherwise instructed by the Faculty or the supervisor.

The report shall be written in English.

Students are responsible for the academic and technical content of the
report. The official \LaTeX{} template is responsible for standardized
document formatting.

\vspace{0.3cm}

% ================================================================
% II. REPORT STRUCTURE
% ================================================================

\noindent
\textbf{II.\quad REPORT STRUCTURE}

\vspace{0.15cm}

The report shall normally be arranged in the following order:

\begin{enumerate}[label=\arabic*., leftmargin=1.2cm, itemsep=2pt]
    \item Outer Cover;
    \item Comments of the Internship Host Organization;
    \item Supervisor's Evaluation and Graduation Internship Defense Results;
    \item Table of Contents;
    \item List of Symbols and Abbreviations, if applicable;
    \item List of Figures;
    \item List of Tables;
    \item Abstract;
    \item Main Chapters of the Internship Report;
    \item References;
    \item Appendices, if applicable.
\end{enumerate}

The preliminary pages shall not be numbered as chapters. The main report
content shall be organized into consecutively numbered chapters beginning
with Chapter 1.

\vspace{0.3cm}

% ================================================================
% III. PAGE SIZE AND MARGINS
% ================================================================

\noindent
\textbf{III.\quad PAGE SIZE AND MARGINS}

\vspace{0.15cm}

The report shall be formatted on A4 paper with the following margins:

\begin{center}
\begin{tabular}{ll}
    Left margin:   & 3.5 cm \\
    Right margin:  & 2.0 cm \\
    Top margin:    & 2.5 cm \\
    Bottom margin: & 2.5 cm
\end{tabular}
\end{center}

These settings are predefined in the official template and shall not be
modified by students.

\vspace{0.3cm}

% ================================================================
% IV. FONT, FONT SIZE, AND LINE SPACING
% ================================================================

\noindent
\textbf{IV.\quad FONT, FONT SIZE, AND LINE SPACING}

\vspace{0.15cm}

The main text shall use a 13 pt font size with 1.3-line spacing.
A Times-family typeface shall be used in accordance with the official
template.

The text shall be fully justified.

Students shall not use repeated spaces, repeated manual line breaks, or
multiple blank lines to position text. Paragraph and heading spacing shall
be controlled by the template.

\vspace{0.3cm}

% ================================================================
% V. OUTER COVER
% ================================================================

\noindent
\textbf{V.\quad OUTER COVER}

\vspace{0.15cm}

The outer cover shall follow the official format provided in the template
and shall contain the following information:

\begin{itemize}[leftmargin=0.8cm, itemsep=2pt]
    \item Thai Nguyen University of Technology;
    \item Faculty of International Training;
    \item Engineering and Technology Department;
    \item University logo and Faculty logo;
    \item the title \textbf{INTERNSHIP REPORT};
    \item Internship Organization;
    \item Author;
    \item Class;
    \item Supervisor;
    \item place and year of completion.
\end{itemize}

Students shall only replace the required information fields and shall not
modify the cover-page layout unless instructed otherwise.

\vspace{0.3cm}

% ================================================================
% VI. COMMENTS OF THE INTERNSHIP HOST ORGANIZATION
% ================================================================

\noindent
\textbf{VI.\quad COMMENTS OF THE INTERNSHIP HOST ORGANIZATION}

\vspace{0.15cm}

The page entitled \textbf{COMMENTS OF THE INTERNSHIP HOST ORGANIZATION}
is reserved for the internship host organization.

It includes:

\begin{itemize}[leftmargin=0.8cm, itemsep=2pt]
    \item the student's full name;
    \item class;
    \item comments from the host organization;
    \item score;
    \item date;
    \item signature and official seal.
\end{itemize}

Students shall not alter or complete official evaluation fields unless
specifically authorized.

\vspace{0.3cm}

% ================================================================
% VII. SUPERVISOR'S EVALUATION AND DEFENSE RESULTS
% ================================================================

\noindent
\textbf{VII.\quad SUPERVISOR'S EVALUATION AND DEFENSE RESULTS}

\vspace{0.15cm}

The section entitled
\textbf{SUPERVISOR'S EVALUATION AND GRADUATION INTERNSHIP DEFENSE RESULTS}
is reserved for the supervisor and defense committee.

The evaluation may include:

\begin{itemize}[leftmargin=0.8cm, itemsep=2pt]
    \item Internship Progress and Student's Attitude;
    \item Report Content;
    \item Report Presentation;
    \item Other Comments;
    \item Supervisor's Opinion;
    \item Defense Result;
    \item Score.
\end{itemize}

Students shall not enter evaluation scores or defense results before they
are officially determined.

\vspace{0.3cm}

% ================================================================
% VIII. CHAPTERS, SECTIONS, AND SUBSECTIONS
% ================================================================

\noindent
\textbf{VIII.\quad CHAPTERS, SECTIONS, AND SUBSECTIONS}

\vspace{0.15cm}

The report structure shall be generated using the standard \LaTeX{}
commands:

\begin{verbatim}
\chapter{...}
\section{...}
\subsection{...}
\end{verbatim}

Chapter, section, and subsection numbers shall be generated automatically
by \LaTeX{}.

Students shall not manually type chapter or section numbers and shall not
reproduce heading styles using commands such as \verb|\textbf|,
\verb|\fontsize|, or \verb|\centering| instead of the appropriate
structural command.

Excessive subdivision should be avoided. In most cases, the hierarchy
Chapter--Section--Subsection is sufficient.

\vspace{0.3cm}

% ================================================================
% IX. ORGANIZATION OF THE MAIN CONTENT
% ================================================================

\noindent
\textbf{IX.\quad ORGANIZATION OF THE MAIN CONTENT}

\vspace{0.15cm}

The specific content of the report shall depend on the internship
organization and the assigned tasks. A typical structure may include:

\begin{itemize}[leftmargin=0.8cm, itemsep=2pt]
    \item \textbf{Chapter 1. Overview of the Internship Organization};
    \item \textbf{Chapter 2. Internship Activities and Assigned Tasks};
    \item \textbf{Chapter 3. Results, Experience, and Evaluation}.
\end{itemize}

These chapter titles are recommendations only and may be adapted to the
actual internship content.

The report shall reflect work actually performed, observed, analyzed, or
learned during the internship. Unrelated theoretical material should not
be included merely to increase report length.

\vspace{0.3cm}

% ================================================================
% X. TABLE OF CONTENTS
% ================================================================

\noindent
\textbf{X.\quad TABLE OF CONTENTS}

\vspace{0.15cm}

The Table of Contents shall be generated automatically using:

\begin{verbatim}
\tableofcontents
\end{verbatim}

Students shall not create or manually edit page numbers in the Table of
Contents.

The document shall be compiled sufficiently to ensure that all page
numbers and entries are correctly updated.

\vspace{0.3cm}

% ================================================================
% XI. LIST OF SYMBOLS AND ABBREVIATIONS
% ================================================================

\noindent
\textbf{XI.\quad LIST OF SYMBOLS AND ABBREVIATIONS}

\vspace{0.15cm}

Symbols and abbreviations that are used repeatedly in the report shall be
listed in the \textbf{List of Symbols and Abbreviations}.

The list may contain the following columns:

\begin{center}
\begin{tabular}{llll}
    No. & Symbol / Abbreviation & Unit & Definition
\end{tabular}
\end{center}

Only symbols and abbreviations actually used in the final report shall be
retained. Example entries contained in the template shall be removed if
they are not relevant to the student's report.

Each abbreviation should normally be defined when it first appears in
the main text.

\vspace{0.3cm}

% ================================================================
% XII. FIGURES, DIAGRAMS, AND PHOTOGRAPHS
% ================================================================

\noindent
\textbf{XII.\quad FIGURES, DIAGRAMS, AND PHOTOGRAPHS}

\vspace{0.15cm}

Figures shall be clear, legible, and directly relevant to the report
content.

Figures shall be inserted using the standard \verb|figure| environment.

Example:

\begin{verbatim}
\begin{figure}[H]
    \centering
    \includegraphics[width=0.8\textwidth]
                    {figures/company_structure.pdf}
    \caption{Organizational structure of the company}
    \label{fig:company-structure}
\end{figure}
\end{verbatim}

Figure captions shall be placed below the figure.

Each referenced figure shall have a unique \verb|\label| and shall be
referenced in the text using, for example:

\begin{verbatim}
Figure~\ref{fig:company-structure}
\end{verbatim}

Students shall not manually type figure numbers.

Figures reproduced or adapted from external sources shall include an
appropriate citation.

Photographs taken during the internship shall be relevant and sufficiently
clear. Students shall not include confidential, restricted, or proprietary
information without permission from the internship organization.

Vector graphics or high-resolution images are recommended whenever
possible.

\vspace{0.3cm}

% ================================================================
% XIII. TABLES
% ================================================================

\noindent
\textbf{XIII.\quad TABLES}

\vspace{0.15cm}

Tables shall be created using suitable \LaTeX{} environments such as
\verb|table|, \verb|tabular|, \verb|tabularx|, or \verb|longtable|.

Table captions shall be placed above the table.

Example:

\begin{verbatim}
\begin{table}[H]
    \centering
    \caption{Internship schedule}
    \label{tab:internship-schedule}

    \begin{tabular}{lll}
        \toprule
        Week & Task & Result \\
        \midrule
        1 & Orientation & Completed \\
        2 & System study & Completed \\
        \bottomrule
    \end{tabular}
\end{table}
\end{verbatim}

Tables shall be referenced using:

\begin{verbatim}
Table~\ref{tab:internship-schedule}
\end{verbatim}

Students shall not manually type table numbers.

Tables should be created directly in \LaTeX{} rather than inserted as
screenshots whenever practical.

\vspace{0.3cm}

% ================================================================
% XIV. EQUATIONS AND TECHNICAL QUANTITIES
% ================================================================

\noindent
\textbf{XIV.\quad EQUATIONS AND TECHNICAL QUANTITIES}

\vspace{0.15cm}

Mathematical expressions shall be typeset directly using \LaTeX{}.
Images or screenshots of equations shall not be used.

Short expressions may be written in inline mathematical mode, for example:

\begin{verbatim}
$R = 0.7\,\Omega$
\end{verbatim}

Important equations shall use a numbered equation environment.

Example:

\begin{verbatim}
\begin{equation}
    P = UI
    \label{eq:power}
\end{equation}
\end{verbatim}

Equations shall be referenced using:

\begin{verbatim}
Equation~\eqref{eq:power}
\end{verbatim}

Equation numbers shall not be manually entered.

SI units shall be used wherever applicable. The \verb|siunitx| package
should be used where appropriate.

Examples:

\begin{verbatim}
\SI{220}{\volt}
\SI{50}{\hertz}
\SI{2.5}{\ampere}
\end{verbatim}

Notation and units shall be consistent throughout the report.

\vspace{0.3cm}

% ================================================================
% XV. SOURCE CODE
% ================================================================

\noindent
\textbf{XV.\quad SOURCE CODE}

\vspace{0.15cm}

If the internship involves programming, PLCs, microcontrollers,
simulation, or software development, relevant source code may be included
in the report.

Source code shall be presented using the \verb|listings| environment or
another environment defined by the official template.

Example:

\begin{verbatim}
\begin{lstlisting}[language=st]
IF Start = TRUE THEN
    Motor := TRUE;
ELSE
    Motor := FALSE;
END_IF;
\end{lstlisting}
\end{verbatim}

Only code necessary for explaining the work should appear in the main
report. Long source-code listings should be moved to an appendix or
submitted as supplementary files if required.

Source code should not normally be inserted as screenshots.

\vspace{0.3cm}

% ================================================================
% XVI. ABSTRACT
% ================================================================

\noindent
\textbf{XVI.\quad ABSTRACT}

\vspace{0.15cm}

The Abstract shall briefly summarize:

\begin{itemize}[leftmargin=0.8cm, itemsep=2pt]
    \item the internship organization;
    \item the objectives of the internship;
    \item the main tasks performed;
    \item the main results, knowledge, and skills obtained.
\end{itemize}

The Abstract should be concise and should not reproduce the full content
of the report chapter by chapter.

\vspace{0.3cm}

% ================================================================
% XVII. REFERENCES AND CITATIONS
% ================================================================

\noindent
\textbf{XVII.\quad REFERENCES AND CITATIONS}

\vspace{0.15cm}

All books, journal papers, conference papers, standards, websites,
datasheets, manuals, technical documents, figures, tables, or other
external sources used in preparing the report shall be properly cited.

References shall be managed through the bibliography system provided with
the official template.

Bibliographic entries shall normally be stored in:

\begin{verbatim}
References.bib
\end{verbatim}

References shall be cited in the report using:

\begin{verbatim}
\cite{reference-key}
\end{verbatim}

Students shall not manually type reference numbers such as ``[1]'',
``[2]'', etc.

Titles of foreign-language publications shall remain in their original
language in the bibliography.

All borrowed ideas, data, figures, tables, equations, methods, or results
shall be properly acknowledged.

\vspace{0.3cm}

% ================================================================
% XVIII. CROSS-REFERENCING
% ================================================================

\noindent
\textbf{XVIII.\quad CROSS-REFERENCING}

\vspace{0.15cm}

Numbered elements shall be cross-referenced using the standard \LaTeX{}
mechanism.

Use the following commands where appropriate:

\begin{verbatim}
\label{...}
\ref{...}
\eqref{...}
\end{verbatim}

For example:

\begin{verbatim}
As shown in Figure~\ref{fig:system}, ...
\end{verbatim}

Students shall not manually type numbers such as ``Figure 2.3'' when
those numbers can be generated automatically.

\vspace{0.3cm}

% ================================================================
% XIX. PAGE NUMBERING
% ================================================================

\noindent
\textbf{XIX.\quad PAGE NUMBERING}

\vspace{0.15cm}

Cover pages and official evaluation forms at the beginning of the report
shall not display page numbers.

The main body of the report shall use Arabic page numbers beginning from
Chapter 1.

Page numbering shall be generated automatically by the official template.
Students shall not manually enter page numbers.

\vspace{0.3cm}

% ================================================================
% XX. PROJECT FILE ORGANIZATION
% ================================================================

\noindent
\textbf{XX.\quad PROJECT FILE ORGANIZATION}

\vspace{0.15cm}

Students should maintain a clear and organized \LaTeX{} project
structure. A recommended structure is:

\begin{verbatim}
InternshipReport/
|
|-- main.tex
|-- References.bib
|
|-- figures/
|   |-- company.pdf
|   |-- system.png
|   `-- ...
|
`-- chapters/
    |-- chapter1.tex
    |-- chapter2.tex
    `-- chapter3.tex
\end{verbatim}

Figures should normally be stored in the \verb|figures| directory
specified by the official template.

File and directory names should preferably contain Latin letters,
numbers, hyphens, or underscores.

Recommended examples include:

\begin{verbatim}
company_structure.pdf
motor_control.png
chapter_01.tex
\end{verbatim}

Spaces and unnecessary special characters in file names should be
avoided where possible.

\vspace{0.3cm}

% ================================================================
% XXI. DOCUMENT COMPILATION AND FINAL CHECK
% ================================================================

\noindent
\textbf{XXI.\quad DOCUMENT COMPILATION AND FINAL CHECK}

\vspace{0.15cm}

Before submission, the complete \LaTeX{} project shall compile
successfully and produce the final PDF report.

The final document shall not contain unresolved references such as:

\begin{verbatim}
??
[?]
\end{verbatim}

Before generating the final PDF, students shall check:

\begin{itemize}[leftmargin=0.8cm, itemsep=2pt]
    \item the Table of Contents;
    \item page numbering;
    \item chapter, section, and subsection numbering;
    \item figure numbering;
    \item table numbering;
    \item equation numbering;
    \item cross-references;
    \item citations;
    \item bibliography;
    \item font display;
    \item mathematical symbols;
    \item image quality;
    \item completeness of all required forms.
\end{itemize}

The document may be compiled using TeX Live, MiKTeX, or another
compatible \LaTeX{} environment. Regardless of the compilation
environment, the resulting PDF shall conform to the official Faculty
template.

\vspace{0.3cm}

% ================================================================
% XXII. CONTENT THAT STUDENTS MAY EDIT
% ================================================================

\noindent
\textbf{XXII.\quad CONTENT THAT STUDENTS MAY EDIT}

\vspace{0.15cm}

Students may edit or replace:

\begin{itemize}[leftmargin=0.8cm, itemsep=2pt]
    \item student name;
    \item class;
    \item internship organization;
    \item supervisor name;
    \item Abstract;
    \item List of Symbols and Abbreviations;
    \item chapter titles and chapter content;
    \item figures and photographs;
    \item tables;
    \item equations;
    \item citations and references.
\end{itemize}

Students shall not modify the main document formatting settings in the
preamble, including the document class, page geometry, line spacing,
heading formatting, table-of-contents formatting, font configuration,
page layout, or cover-page macros unless specifically instructed.

\vspace{0.3cm}

% ================================================================
% XXIII. FINAL SUBMISSION REQUIREMENTS
% ================================================================

\noindent
\textbf{XXIII.\quad FINAL SUBMISSION REQUIREMENTS}

\vspace{0.15cm}

Before submission, the final report shall:

\begin{itemize}[leftmargin=0.8cm, itemsep=2pt]
    \item compile without errors that affect the report content or presentation;
    \item contain all required forms and report sections;
    \item contain correct student, class, internship organization, and supervisor information;
    \item contain an updated Table of Contents, List of Figures, and List of Tables;
    \item contain no unused example content from the template;
    \item contain no placeholder text such as \textit{Title of Chapter 1},
    \textit{Section 1.1}, or other sample entries;
    \item contain all required citations and references;
    \item contain no unresolved cross-references.
\end{itemize}

If the Faculty requires submission of the \LaTeX{} source, the complete
project shall include all \verb|.tex| files, the bibliography database,
figures, and any other files required to reproduce the submitted PDF.

\vspace{0.4cm}

\noindent
\textbf{Important:}
Before submitting the final internship report, students shall remove or
comment out the following command in the main document:

\begin{verbatim}
\input{Regulations.tex}
\end{verbatim}

For example:

\begin{verbatim}
% \input{Regulations.tex}
\end{verbatim}

This removes the guidance pages from the final submitted report while
preserving the regulations file in the \LaTeX{} project for reference.

} % End 13pt content

\clearpage
 command before submitting
% the final internship report.
%
% =================================================================

\clearpage
\thispagestyle{empty}

\begin{center}
    {\fontsize{16pt}{20pt}\selectfont\bfseries
    SPECIFIC REGULATIONS FOR PREPARING\\
    GRADUATION INTERNSHIP REPORTS USING \LaTeX}
\end{center}

\vspace{0.5cm}

{\fontsize{13pt}{17pt}\selectfont
\setstretch{1.3}
\justifying

\noindent
\textbf{Note:}
This section is provided for guidance only and is included in the main
document using the command \verb|
% =================================================================
% REGULATIONS FOR PREPARING THE GRADUATION INTERNSHIP REPORT
% USING LATEX
% =================================================================
%
% This file is intended to be included from the main document with:
%
%     \input{Regulations.tex}
%
% IMPORTANT:
% This section is provided for guidance only. Students shall remove
% or comment out the \input{Regulations.tex} command before submitting
% the final internship report.
%
% =================================================================

\clearpage
\thispagestyle{empty}

\begin{center}
    {\fontsize{16pt}{20pt}\selectfont\bfseries
    SPECIFIC REGULATIONS FOR PREPARING\\
    GRADUATION INTERNSHIP REPORTS USING \LaTeX}
\end{center}

\vspace{0.5cm}

{\fontsize{13pt}{17pt}\selectfont
\setstretch{1.3}
\justifying

\noindent
\textbf{Note:}
This section is provided for guidance only and is included in the main
document using the command \verb|\input{Regulations.tex}|.
Students shall remove or comment out this command before submitting the
final internship report.

\vspace{0.4cm}

% ================================================================
% I. GENERAL REQUIREMENTS
% ================================================================

\noindent
\textbf{I.\quad GENERAL REQUIREMENTS}

\vspace{0.15cm}

The Graduation Internship Report shall be prepared using the official
\LaTeX{} template provided by the Faculty. Students shall preserve the
predefined page layout, typography, line spacing, heading styles,
numbering system, figure and table formatting, and bibliography settings
unless otherwise instructed by the Faculty or the supervisor.

The report shall be written in English.

Students are responsible for the academic and technical content of the
report. The official \LaTeX{} template is responsible for standardized
document formatting.

\vspace{0.3cm}

% ================================================================
% II. REPORT STRUCTURE
% ================================================================

\noindent
\textbf{II.\quad REPORT STRUCTURE}

\vspace{0.15cm}

The report shall normally be arranged in the following order:

\begin{enumerate}[label=\arabic*., leftmargin=1.2cm, itemsep=2pt]
    \item Outer Cover;
    \item Comments of the Internship Host Organization;
    \item Supervisor's Evaluation and Graduation Internship Defense Results;
    \item Table of Contents;
    \item List of Symbols and Abbreviations, if applicable;
    \item List of Figures;
    \item List of Tables;
    \item Abstract;
    \item Main Chapters of the Internship Report;
    \item References;
    \item Appendices, if applicable.
\end{enumerate}

The preliminary pages shall not be numbered as chapters. The main report
content shall be organized into consecutively numbered chapters beginning
with Chapter 1.

\vspace{0.3cm}

% ================================================================
% III. PAGE SIZE AND MARGINS
% ================================================================

\noindent
\textbf{III.\quad PAGE SIZE AND MARGINS}

\vspace{0.15cm}

The report shall be formatted on A4 paper with the following margins:

\begin{center}
\begin{tabular}{ll}
    Left margin:   & 3.5 cm \\
    Right margin:  & 2.0 cm \\
    Top margin:    & 2.5 cm \\
    Bottom margin: & 2.5 cm
\end{tabular}
\end{center}

These settings are predefined in the official template and shall not be
modified by students.

\vspace{0.3cm}

% ================================================================
% IV. FONT, FONT SIZE, AND LINE SPACING
% ================================================================

\noindent
\textbf{IV.\quad FONT, FONT SIZE, AND LINE SPACING}

\vspace{0.15cm}

The main text shall use a 13 pt font size with 1.3-line spacing.
A Times-family typeface shall be used in accordance with the official
template.

The text shall be fully justified.

Students shall not use repeated spaces, repeated manual line breaks, or
multiple blank lines to position text. Paragraph and heading spacing shall
be controlled by the template.

\vspace{0.3cm}

% ================================================================
% V. OUTER COVER
% ================================================================

\noindent
\textbf{V.\quad OUTER COVER}

\vspace{0.15cm}

The outer cover shall follow the official format provided in the template
and shall contain the following information:

\begin{itemize}[leftmargin=0.8cm, itemsep=2pt]
    \item Thai Nguyen University of Technology;
    \item Faculty of International Training;
    \item Engineering and Technology Department;
    \item University logo and Faculty logo;
    \item the title \textbf{INTERNSHIP REPORT};
    \item Internship Organization;
    \item Author;
    \item Class;
    \item Supervisor;
    \item place and year of completion.
\end{itemize}

Students shall only replace the required information fields and shall not
modify the cover-page layout unless instructed otherwise.

\vspace{0.3cm}

% ================================================================
% VI. COMMENTS OF THE INTERNSHIP HOST ORGANIZATION
% ================================================================

\noindent
\textbf{VI.\quad COMMENTS OF THE INTERNSHIP HOST ORGANIZATION}

\vspace{0.15cm}

The page entitled \textbf{COMMENTS OF THE INTERNSHIP HOST ORGANIZATION}
is reserved for the internship host organization.

It includes:

\begin{itemize}[leftmargin=0.8cm, itemsep=2pt]
    \item the student's full name;
    \item class;
    \item comments from the host organization;
    \item score;
    \item date;
    \item signature and official seal.
\end{itemize}

Students shall not alter or complete official evaluation fields unless
specifically authorized.

\vspace{0.3cm}

% ================================================================
% VII. SUPERVISOR'S EVALUATION AND DEFENSE RESULTS
% ================================================================

\noindent
\textbf{VII.\quad SUPERVISOR'S EVALUATION AND DEFENSE RESULTS}

\vspace{0.15cm}

The section entitled
\textbf{SUPERVISOR'S EVALUATION AND GRADUATION INTERNSHIP DEFENSE RESULTS}
is reserved for the supervisor and defense committee.

The evaluation may include:

\begin{itemize}[leftmargin=0.8cm, itemsep=2pt]
    \item Internship Progress and Student's Attitude;
    \item Report Content;
    \item Report Presentation;
    \item Other Comments;
    \item Supervisor's Opinion;
    \item Defense Result;
    \item Score.
\end{itemize}

Students shall not enter evaluation scores or defense results before they
are officially determined.

\vspace{0.3cm}

% ================================================================
% VIII. CHAPTERS, SECTIONS, AND SUBSECTIONS
% ================================================================

\noindent
\textbf{VIII.\quad CHAPTERS, SECTIONS, AND SUBSECTIONS}

\vspace{0.15cm}

The report structure shall be generated using the standard \LaTeX{}
commands:

\begin{verbatim}
\chapter{...}
\section{...}
\subsection{...}
\end{verbatim}

Chapter, section, and subsection numbers shall be generated automatically
by \LaTeX{}.

Students shall not manually type chapter or section numbers and shall not
reproduce heading styles using commands such as \verb|\textbf|,
\verb|\fontsize|, or \verb|\centering| instead of the appropriate
structural command.

Excessive subdivision should be avoided. In most cases, the hierarchy
Chapter--Section--Subsection is sufficient.

\vspace{0.3cm}

% ================================================================
% IX. ORGANIZATION OF THE MAIN CONTENT
% ================================================================

\noindent
\textbf{IX.\quad ORGANIZATION OF THE MAIN CONTENT}

\vspace{0.15cm}

The specific content of the report shall depend on the internship
organization and the assigned tasks. A typical structure may include:

\begin{itemize}[leftmargin=0.8cm, itemsep=2pt]
    \item \textbf{Chapter 1. Overview of the Internship Organization};
    \item \textbf{Chapter 2. Internship Activities and Assigned Tasks};
    \item \textbf{Chapter 3. Results, Experience, and Evaluation}.
\end{itemize}

These chapter titles are recommendations only and may be adapted to the
actual internship content.

The report shall reflect work actually performed, observed, analyzed, or
learned during the internship. Unrelated theoretical material should not
be included merely to increase report length.

\vspace{0.3cm}

% ================================================================
% X. TABLE OF CONTENTS
% ================================================================

\noindent
\textbf{X.\quad TABLE OF CONTENTS}

\vspace{0.15cm}

The Table of Contents shall be generated automatically using:

\begin{verbatim}
\tableofcontents
\end{verbatim}

Students shall not create or manually edit page numbers in the Table of
Contents.

The document shall be compiled sufficiently to ensure that all page
numbers and entries are correctly updated.

\vspace{0.3cm}

% ================================================================
% XI. LIST OF SYMBOLS AND ABBREVIATIONS
% ================================================================

\noindent
\textbf{XI.\quad LIST OF SYMBOLS AND ABBREVIATIONS}

\vspace{0.15cm}

Symbols and abbreviations that are used repeatedly in the report shall be
listed in the \textbf{List of Symbols and Abbreviations}.

The list may contain the following columns:

\begin{center}
\begin{tabular}{llll}
    No. & Symbol / Abbreviation & Unit & Definition
\end{tabular}
\end{center}

Only symbols and abbreviations actually used in the final report shall be
retained. Example entries contained in the template shall be removed if
they are not relevant to the student's report.

Each abbreviation should normally be defined when it first appears in
the main text.

\vspace{0.3cm}

% ================================================================
% XII. FIGURES, DIAGRAMS, AND PHOTOGRAPHS
% ================================================================

\noindent
\textbf{XII.\quad FIGURES, DIAGRAMS, AND PHOTOGRAPHS}

\vspace{0.15cm}

Figures shall be clear, legible, and directly relevant to the report
content.

Figures shall be inserted using the standard \verb|figure| environment.

Example:

\begin{verbatim}
\begin{figure}[H]
    \centering
    \includegraphics[width=0.8\textwidth]
                    {figures/company_structure.pdf}
    \caption{Organizational structure of the company}
    \label{fig:company-structure}
\end{figure}
\end{verbatim}

Figure captions shall be placed below the figure.

Each referenced figure shall have a unique \verb|\label| and shall be
referenced in the text using, for example:

\begin{verbatim}
Figure~\ref{fig:company-structure}
\end{verbatim}

Students shall not manually type figure numbers.

Figures reproduced or adapted from external sources shall include an
appropriate citation.

Photographs taken during the internship shall be relevant and sufficiently
clear. Students shall not include confidential, restricted, or proprietary
information without permission from the internship organization.

Vector graphics or high-resolution images are recommended whenever
possible.

\vspace{0.3cm}

% ================================================================
% XIII. TABLES
% ================================================================

\noindent
\textbf{XIII.\quad TABLES}

\vspace{0.15cm}

Tables shall be created using suitable \LaTeX{} environments such as
\verb|table|, \verb|tabular|, \verb|tabularx|, or \verb|longtable|.

Table captions shall be placed above the table.

Example:

\begin{verbatim}
\begin{table}[H]
    \centering
    \caption{Internship schedule}
    \label{tab:internship-schedule}

    \begin{tabular}{lll}
        \toprule
        Week & Task & Result \\
        \midrule
        1 & Orientation & Completed \\
        2 & System study & Completed \\
        \bottomrule
    \end{tabular}
\end{table}
\end{verbatim}

Tables shall be referenced using:

\begin{verbatim}
Table~\ref{tab:internship-schedule}
\end{verbatim}

Students shall not manually type table numbers.

Tables should be created directly in \LaTeX{} rather than inserted as
screenshots whenever practical.

\vspace{0.3cm}

% ================================================================
% XIV. EQUATIONS AND TECHNICAL QUANTITIES
% ================================================================

\noindent
\textbf{XIV.\quad EQUATIONS AND TECHNICAL QUANTITIES}

\vspace{0.15cm}

Mathematical expressions shall be typeset directly using \LaTeX{}.
Images or screenshots of equations shall not be used.

Short expressions may be written in inline mathematical mode, for example:

\begin{verbatim}
$R = 0.7\,\Omega$
\end{verbatim}

Important equations shall use a numbered equation environment.

Example:

\begin{verbatim}
\begin{equation}
    P = UI
    \label{eq:power}
\end{equation}
\end{verbatim}

Equations shall be referenced using:

\begin{verbatim}
Equation~\eqref{eq:power}
\end{verbatim}

Equation numbers shall not be manually entered.

SI units shall be used wherever applicable. The \verb|siunitx| package
should be used where appropriate.

Examples:

\begin{verbatim}
\SI{220}{\volt}
\SI{50}{\hertz}
\SI{2.5}{\ampere}
\end{verbatim}

Notation and units shall be consistent throughout the report.

\vspace{0.3cm}

% ================================================================
% XV. SOURCE CODE
% ================================================================

\noindent
\textbf{XV.\quad SOURCE CODE}

\vspace{0.15cm}

If the internship involves programming, PLCs, microcontrollers,
simulation, or software development, relevant source code may be included
in the report.

Source code shall be presented using the \verb|listings| environment or
another environment defined by the official template.

Example:

\begin{verbatim}
\begin{lstlisting}[language=st]
IF Start = TRUE THEN
    Motor := TRUE;
ELSE
    Motor := FALSE;
END_IF;
\end{lstlisting}
\end{verbatim}

Only code necessary for explaining the work should appear in the main
report. Long source-code listings should be moved to an appendix or
submitted as supplementary files if required.

Source code should not normally be inserted as screenshots.

\vspace{0.3cm}

% ================================================================
% XVI. ABSTRACT
% ================================================================

\noindent
\textbf{XVI.\quad ABSTRACT}

\vspace{0.15cm}

The Abstract shall briefly summarize:

\begin{itemize}[leftmargin=0.8cm, itemsep=2pt]
    \item the internship organization;
    \item the objectives of the internship;
    \item the main tasks performed;
    \item the main results, knowledge, and skills obtained.
\end{itemize}

The Abstract should be concise and should not reproduce the full content
of the report chapter by chapter.

\vspace{0.3cm}

% ================================================================
% XVII. REFERENCES AND CITATIONS
% ================================================================

\noindent
\textbf{XVII.\quad REFERENCES AND CITATIONS}

\vspace{0.15cm}

All books, journal papers, conference papers, standards, websites,
datasheets, manuals, technical documents, figures, tables, or other
external sources used in preparing the report shall be properly cited.

References shall be managed through the bibliography system provided with
the official template.

Bibliographic entries shall normally be stored in:

\begin{verbatim}
References.bib
\end{verbatim}

References shall be cited in the report using:

\begin{verbatim}
\cite{reference-key}
\end{verbatim}

Students shall not manually type reference numbers such as ``[1]'',
``[2]'', etc.

Titles of foreign-language publications shall remain in their original
language in the bibliography.

All borrowed ideas, data, figures, tables, equations, methods, or results
shall be properly acknowledged.

\vspace{0.3cm}

% ================================================================
% XVIII. CROSS-REFERENCING
% ================================================================

\noindent
\textbf{XVIII.\quad CROSS-REFERENCING}

\vspace{0.15cm}

Numbered elements shall be cross-referenced using the standard \LaTeX{}
mechanism.

Use the following commands where appropriate:

\begin{verbatim}
\label{...}
\ref{...}
\eqref{...}
\end{verbatim}

For example:

\begin{verbatim}
As shown in Figure~\ref{fig:system}, ...
\end{verbatim}

Students shall not manually type numbers such as ``Figure 2.3'' when
those numbers can be generated automatically.

\vspace{0.3cm}

% ================================================================
% XIX. PAGE NUMBERING
% ================================================================

\noindent
\textbf{XIX.\quad PAGE NUMBERING}

\vspace{0.15cm}

Cover pages and official evaluation forms at the beginning of the report
shall not display page numbers.

The main body of the report shall use Arabic page numbers beginning from
Chapter 1.

Page numbering shall be generated automatically by the official template.
Students shall not manually enter page numbers.

\vspace{0.3cm}

% ================================================================
% XX. PROJECT FILE ORGANIZATION
% ================================================================

\noindent
\textbf{XX.\quad PROJECT FILE ORGANIZATION}

\vspace{0.15cm}

Students should maintain a clear and organized \LaTeX{} project
structure. A recommended structure is:

\begin{verbatim}
InternshipReport/
|
|-- main.tex
|-- References.bib
|
|-- figures/
|   |-- company.pdf
|   |-- system.png
|   `-- ...
|
`-- chapters/
    |-- chapter1.tex
    |-- chapter2.tex
    `-- chapter3.tex
\end{verbatim}

Figures should normally be stored in the \verb|figures| directory
specified by the official template.

File and directory names should preferably contain Latin letters,
numbers, hyphens, or underscores.

Recommended examples include:

\begin{verbatim}
company_structure.pdf
motor_control.png
chapter_01.tex
\end{verbatim}

Spaces and unnecessary special characters in file names should be
avoided where possible.

\vspace{0.3cm}

% ================================================================
% XXI. DOCUMENT COMPILATION AND FINAL CHECK
% ================================================================

\noindent
\textbf{XXI.\quad DOCUMENT COMPILATION AND FINAL CHECK}

\vspace{0.15cm}

Before submission, the complete \LaTeX{} project shall compile
successfully and produce the final PDF report.

The final document shall not contain unresolved references such as:

\begin{verbatim}
??
[?]
\end{verbatim}

Before generating the final PDF, students shall check:

\begin{itemize}[leftmargin=0.8cm, itemsep=2pt]
    \item the Table of Contents;
    \item page numbering;
    \item chapter, section, and subsection numbering;
    \item figure numbering;
    \item table numbering;
    \item equation numbering;
    \item cross-references;
    \item citations;
    \item bibliography;
    \item font display;
    \item mathematical symbols;
    \item image quality;
    \item completeness of all required forms.
\end{itemize}

The document may be compiled using TeX Live, MiKTeX, or another
compatible \LaTeX{} environment. Regardless of the compilation
environment, the resulting PDF shall conform to the official Faculty
template.

\vspace{0.3cm}

% ================================================================
% XXII. CONTENT THAT STUDENTS MAY EDIT
% ================================================================

\noindent
\textbf{XXII.\quad CONTENT THAT STUDENTS MAY EDIT}

\vspace{0.15cm}

Students may edit or replace:

\begin{itemize}[leftmargin=0.8cm, itemsep=2pt]
    \item student name;
    \item class;
    \item internship organization;
    \item supervisor name;
    \item Abstract;
    \item List of Symbols and Abbreviations;
    \item chapter titles and chapter content;
    \item figures and photographs;
    \item tables;
    \item equations;
    \item citations and references.
\end{itemize}

Students shall not modify the main document formatting settings in the
preamble, including the document class, page geometry, line spacing,
heading formatting, table-of-contents formatting, font configuration,
page layout, or cover-page macros unless specifically instructed.

\vspace{0.3cm}

% ================================================================
% XXIII. FINAL SUBMISSION REQUIREMENTS
% ================================================================

\noindent
\textbf{XXIII.\quad FINAL SUBMISSION REQUIREMENTS}

\vspace{0.15cm}

Before submission, the final report shall:

\begin{itemize}[leftmargin=0.8cm, itemsep=2pt]
    \item compile without errors that affect the report content or presentation;
    \item contain all required forms and report sections;
    \item contain correct student, class, internship organization, and supervisor information;
    \item contain an updated Table of Contents, List of Figures, and List of Tables;
    \item contain no unused example content from the template;
    \item contain no placeholder text such as \textit{Title of Chapter 1},
    \textit{Section 1.1}, or other sample entries;
    \item contain all required citations and references;
    \item contain no unresolved cross-references.
\end{itemize}

If the Faculty requires submission of the \LaTeX{} source, the complete
project shall include all \verb|.tex| files, the bibliography database,
figures, and any other files required to reproduce the submitted PDF.

\vspace{0.4cm}

\noindent
\textbf{Important:}
Before submitting the final internship report, students shall remove or
comment out the following command in the main document:

\begin{verbatim}
\input{Regulations.tex}
\end{verbatim}

For example:

\begin{verbatim}
% \input{Regulations.tex}
\end{verbatim}

This removes the guidance pages from the final submitted report while
preserving the regulations file in the \LaTeX{} project for reference.

} % End 13pt content

\clearpage
|.
Students shall remove or comment out this command before submitting the
final internship report.

\vspace{0.4cm}

% ================================================================
% I. GENERAL REQUIREMENTS
% ================================================================

\noindent
\textbf{I.\quad GENERAL REQUIREMENTS}

\vspace{0.15cm}

The Graduation Internship Report shall be prepared using the official
\LaTeX{} template provided by the Faculty. Students shall preserve the
predefined page layout, typography, line spacing, heading styles,
numbering system, figure and table formatting, and bibliography settings
unless otherwise instructed by the Faculty or the supervisor.

The report shall be written in English.

Students are responsible for the academic and technical content of the
report. The official \LaTeX{} template is responsible for standardized
document formatting.

\vspace{0.3cm}

% ================================================================
% II. REPORT STRUCTURE
% ================================================================

\noindent
\textbf{II.\quad REPORT STRUCTURE}

\vspace{0.15cm}

The report shall normally be arranged in the following order:

\begin{enumerate}[label=\arabic*., leftmargin=1.2cm, itemsep=2pt]
    \item Outer Cover;
    \item Comments of the Internship Host Organization;
    \item Supervisor's Evaluation and Graduation Internship Defense Results;
    \item Table of Contents;
    \item List of Symbols and Abbreviations, if applicable;
    \item List of Figures;
    \item List of Tables;
    \item Abstract;
    \item Main Chapters of the Internship Report;
    \item References;
    \item Appendices, if applicable.
\end{enumerate}

The preliminary pages shall not be numbered as chapters. The main report
content shall be organized into consecutively numbered chapters beginning
with Chapter 1.

\vspace{0.3cm}

% ================================================================
% III. PAGE SIZE AND MARGINS
% ================================================================

\noindent
\textbf{III.\quad PAGE SIZE AND MARGINS}

\vspace{0.15cm}

The report shall be formatted on A4 paper with the following margins:

\begin{center}
\begin{tabular}{ll}
    Left margin:   & 3.5 cm \\
    Right margin:  & 2.0 cm \\
    Top margin:    & 2.5 cm \\
    Bottom margin: & 2.5 cm
\end{tabular}
\end{center}

These settings are predefined in the official template and shall not be
modified by students.

\vspace{0.3cm}

% ================================================================
% IV. FONT, FONT SIZE, AND LINE SPACING
% ================================================================

\noindent
\textbf{IV.\quad FONT, FONT SIZE, AND LINE SPACING}

\vspace{0.15cm}

The main text shall use a 13 pt font size with 1.3-line spacing.
A Times-family typeface shall be used in accordance with the official
template.

The text shall be fully justified.

Students shall not use repeated spaces, repeated manual line breaks, or
multiple blank lines to position text. Paragraph and heading spacing shall
be controlled by the template.

\vspace{0.3cm}

% ================================================================
% V. OUTER COVER
% ================================================================

\noindent
\textbf{V.\quad OUTER COVER}

\vspace{0.15cm}

The outer cover shall follow the official format provided in the template
and shall contain the following information:

\begin{itemize}[leftmargin=0.8cm, itemsep=2pt]
    \item Thai Nguyen University of Technology;
    \item Faculty of International Training;
    \item Engineering and Technology Department;
    \item University logo and Faculty logo;
    \item the title \textbf{INTERNSHIP REPORT};
    \item Internship Organization;
    \item Author;
    \item Class;
    \item Supervisor;
    \item place and year of completion.
\end{itemize}

Students shall only replace the required information fields and shall not
modify the cover-page layout unless instructed otherwise.

\vspace{0.3cm}

% ================================================================
% VI. COMMENTS OF THE INTERNSHIP HOST ORGANIZATION
% ================================================================

\noindent
\textbf{VI.\quad COMMENTS OF THE INTERNSHIP HOST ORGANIZATION}

\vspace{0.15cm}

The page entitled \textbf{COMMENTS OF THE INTERNSHIP HOST ORGANIZATION}
is reserved for the internship host organization.

It includes:

\begin{itemize}[leftmargin=0.8cm, itemsep=2pt]
    \item the student's full name;
    \item class;
    \item comments from the host organization;
    \item score;
    \item date;
    \item signature and official seal.
\end{itemize}

Students shall not alter or complete official evaluation fields unless
specifically authorized.

\vspace{0.3cm}

% ================================================================
% VII. SUPERVISOR'S EVALUATION AND DEFENSE RESULTS
% ================================================================

\noindent
\textbf{VII.\quad SUPERVISOR'S EVALUATION AND DEFENSE RESULTS}

\vspace{0.15cm}

The section entitled
\textbf{SUPERVISOR'S EVALUATION AND GRADUATION INTERNSHIP DEFENSE RESULTS}
is reserved for the supervisor and defense committee.

The evaluation may include:

\begin{itemize}[leftmargin=0.8cm, itemsep=2pt]
    \item Internship Progress and Student's Attitude;
    \item Report Content;
    \item Report Presentation;
    \item Other Comments;
    \item Supervisor's Opinion;
    \item Defense Result;
    \item Score.
\end{itemize}

Students shall not enter evaluation scores or defense results before they
are officially determined.

\vspace{0.3cm}

% ================================================================
% VIII. CHAPTERS, SECTIONS, AND SUBSECTIONS
% ================================================================

\noindent
\textbf{VIII.\quad CHAPTERS, SECTIONS, AND SUBSECTIONS}

\vspace{0.15cm}

The report structure shall be generated using the standard \LaTeX{}
commands:

\begin{verbatim}
\chapter{...}
\section{...}
\subsection{...}
\end{verbatim}

Chapter, section, and subsection numbers shall be generated automatically
by \LaTeX{}.

Students shall not manually type chapter or section numbers and shall not
reproduce heading styles using commands such as \verb|\textbf|,
\verb|\fontsize|, or \verb|\centering| instead of the appropriate
structural command.

Excessive subdivision should be avoided. In most cases, the hierarchy
Chapter--Section--Subsection is sufficient.

\vspace{0.3cm}

% ================================================================
% IX. ORGANIZATION OF THE MAIN CONTENT
% ================================================================

\noindent
\textbf{IX.\quad ORGANIZATION OF THE MAIN CONTENT}

\vspace{0.15cm}

The specific content of the report shall depend on the internship
organization and the assigned tasks. A typical structure may include:

\begin{itemize}[leftmargin=0.8cm, itemsep=2pt]
    \item \textbf{Chapter 1. Overview of the Internship Organization};
    \item \textbf{Chapter 2. Internship Activities and Assigned Tasks};
    \item \textbf{Chapter 3. Results, Experience, and Evaluation}.
\end{itemize}

These chapter titles are recommendations only and may be adapted to the
actual internship content.

The report shall reflect work actually performed, observed, analyzed, or
learned during the internship. Unrelated theoretical material should not
be included merely to increase report length.

\vspace{0.3cm}

% ================================================================
% X. TABLE OF CONTENTS
% ================================================================

\noindent
\textbf{X.\quad TABLE OF CONTENTS}

\vspace{0.15cm}

The Table of Contents shall be generated automatically using:

\begin{verbatim}
\tableofcontents
\end{verbatim}

Students shall not create or manually edit page numbers in the Table of
Contents.

The document shall be compiled sufficiently to ensure that all page
numbers and entries are correctly updated.

\vspace{0.3cm}

% ================================================================
% XI. LIST OF SYMBOLS AND ABBREVIATIONS
% ================================================================

\noindent
\textbf{XI.\quad LIST OF SYMBOLS AND ABBREVIATIONS}

\vspace{0.15cm}

Symbols and abbreviations that are used repeatedly in the report shall be
listed in the \textbf{List of Symbols and Abbreviations}.

The list may contain the following columns:

\begin{center}
\begin{tabular}{llll}
    No. & Symbol / Abbreviation & Unit & Definition
\end{tabular}
\end{center}

Only symbols and abbreviations actually used in the final report shall be
retained. Example entries contained in the template shall be removed if
they are not relevant to the student's report.

Each abbreviation should normally be defined when it first appears in
the main text.

\vspace{0.3cm}

% ================================================================
% XII. FIGURES, DIAGRAMS, AND PHOTOGRAPHS
% ================================================================

\noindent
\textbf{XII.\quad FIGURES, DIAGRAMS, AND PHOTOGRAPHS}

\vspace{0.15cm}

Figures shall be clear, legible, and directly relevant to the report
content.

Figures shall be inserted using the standard \verb|figure| environment.

Example:

\begin{verbatim}
\begin{figure}[H]
    \centering
    \includegraphics[width=0.8\textwidth]
                    {figures/company_structure.pdf}
    \caption{Organizational structure of the company}
    \label{fig:company-structure}
\end{figure}
\end{verbatim}

Figure captions shall be placed below the figure.

Each referenced figure shall have a unique \verb|\label| and shall be
referenced in the text using, for example:

\begin{verbatim}
Figure~\ref{fig:company-structure}
\end{verbatim}

Students shall not manually type figure numbers.

Figures reproduced or adapted from external sources shall include an
appropriate citation.

Photographs taken during the internship shall be relevant and sufficiently
clear. Students shall not include confidential, restricted, or proprietary
information without permission from the internship organization.

Vector graphics or high-resolution images are recommended whenever
possible.

\vspace{0.3cm}

% ================================================================
% XIII. TABLES
% ================================================================

\noindent
\textbf{XIII.\quad TABLES}

\vspace{0.15cm}

Tables shall be created using suitable \LaTeX{} environments such as
\verb|table|, \verb|tabular|, \verb|tabularx|, or \verb|longtable|.

Table captions shall be placed above the table.

Example:

\begin{verbatim}
\begin{table}[H]
    \centering
    \caption{Internship schedule}
    \label{tab:internship-schedule}

    \begin{tabular}{lll}
        \toprule
        Week & Task & Result \\
        \midrule
        1 & Orientation & Completed \\
        2 & System study & Completed \\
        \bottomrule
    \end{tabular}
\end{table}
\end{verbatim}

Tables shall be referenced using:

\begin{verbatim}
Table~\ref{tab:internship-schedule}
\end{verbatim}

Students shall not manually type table numbers.

Tables should be created directly in \LaTeX{} rather than inserted as
screenshots whenever practical.

\vspace{0.3cm}

% ================================================================
% XIV. EQUATIONS AND TECHNICAL QUANTITIES
% ================================================================

\noindent
\textbf{XIV.\quad EQUATIONS AND TECHNICAL QUANTITIES}

\vspace{0.15cm}

Mathematical expressions shall be typeset directly using \LaTeX{}.
Images or screenshots of equations shall not be used.

Short expressions may be written in inline mathematical mode, for example:

\begin{verbatim}
$R = 0.7\,\Omega$
\end{verbatim}

Important equations shall use a numbered equation environment.

Example:

\begin{verbatim}
\begin{equation}
    P = UI
    \label{eq:power}
\end{equation}
\end{verbatim}

Equations shall be referenced using:

\begin{verbatim}
Equation~\eqref{eq:power}
\end{verbatim}

Equation numbers shall not be manually entered.

SI units shall be used wherever applicable. The \verb|siunitx| package
should be used where appropriate.

Examples:

\begin{verbatim}
\SI{220}{\volt}
\SI{50}{\hertz}
\SI{2.5}{\ampere}
\end{verbatim}

Notation and units shall be consistent throughout the report.

\vspace{0.3cm}

% ================================================================
% XV. SOURCE CODE
% ================================================================

\noindent
\textbf{XV.\quad SOURCE CODE}

\vspace{0.15cm}

If the internship involves programming, PLCs, microcontrollers,
simulation, or software development, relevant source code may be included
in the report.

Source code shall be presented using the \verb|listings| environment or
another environment defined by the official template.

Example:

\begin{verbatim}
\begin{lstlisting}[language=st]
IF Start = TRUE THEN
    Motor := TRUE;
ELSE
    Motor := FALSE;
END_IF;
\end{lstlisting}
\end{verbatim}

Only code necessary for explaining the work should appear in the main
report. Long source-code listings should be moved to an appendix or
submitted as supplementary files if required.

Source code should not normally be inserted as screenshots.

\vspace{0.3cm}

% ================================================================
% XVI. ABSTRACT
% ================================================================

\noindent
\textbf{XVI.\quad ABSTRACT}

\vspace{0.15cm}

The Abstract shall briefly summarize:

\begin{itemize}[leftmargin=0.8cm, itemsep=2pt]
    \item the internship organization;
    \item the objectives of the internship;
    \item the main tasks performed;
    \item the main results, knowledge, and skills obtained.
\end{itemize}

The Abstract should be concise and should not reproduce the full content
of the report chapter by chapter.

\vspace{0.3cm}

% ================================================================
% XVII. REFERENCES AND CITATIONS
% ================================================================

\noindent
\textbf{XVII.\quad REFERENCES AND CITATIONS}

\vspace{0.15cm}

All books, journal papers, conference papers, standards, websites,
datasheets, manuals, technical documents, figures, tables, or other
external sources used in preparing the report shall be properly cited.

References shall be managed through the bibliography system provided with
the official template.

Bibliographic entries shall normally be stored in:

\begin{verbatim}
References.bib
\end{verbatim}

References shall be cited in the report using:

\begin{verbatim}
\cite{reference-key}
\end{verbatim}

Students shall not manually type reference numbers such as ``[1]'',
``[2]'', etc.

Titles of foreign-language publications shall remain in their original
language in the bibliography.

All borrowed ideas, data, figures, tables, equations, methods, or results
shall be properly acknowledged.

\vspace{0.3cm}

% ================================================================
% XVIII. CROSS-REFERENCING
% ================================================================

\noindent
\textbf{XVIII.\quad CROSS-REFERENCING}

\vspace{0.15cm}

Numbered elements shall be cross-referenced using the standard \LaTeX{}
mechanism.

Use the following commands where appropriate:

\begin{verbatim}
\label{...}
\ref{...}
\eqref{...}
\end{verbatim}

For example:

\begin{verbatim}
As shown in Figure~\ref{fig:system}, ...
\end{verbatim}

Students shall not manually type numbers such as ``Figure 2.3'' when
those numbers can be generated automatically.

\vspace{0.3cm}

% ================================================================
% XIX. PAGE NUMBERING
% ================================================================

\noindent
\textbf{XIX.\quad PAGE NUMBERING}

\vspace{0.15cm}

Cover pages and official evaluation forms at the beginning of the report
shall not display page numbers.

The main body of the report shall use Arabic page numbers beginning from
Chapter 1.

Page numbering shall be generated automatically by the official template.
Students shall not manually enter page numbers.

\vspace{0.3cm}

% ================================================================
% XX. PROJECT FILE ORGANIZATION
% ================================================================

\noindent
\textbf{XX.\quad PROJECT FILE ORGANIZATION}

\vspace{0.15cm}

Students should maintain a clear and organized \LaTeX{} project
structure. A recommended structure is:

\begin{verbatim}
InternshipReport/
|
|-- main.tex
|-- References.bib
|
|-- figures/
|   |-- company.pdf
|   |-- system.png
|   `-- ...
|
`-- chapters/
    |-- chapter1.tex
    |-- chapter2.tex
    `-- chapter3.tex
\end{verbatim}

Figures should normally be stored in the \verb|figures| directory
specified by the official template.

File and directory names should preferably contain Latin letters,
numbers, hyphens, or underscores.

Recommended examples include:

\begin{verbatim}
company_structure.pdf
motor_control.png
chapter_01.tex
\end{verbatim}

Spaces and unnecessary special characters in file names should be
avoided where possible.

\vspace{0.3cm}

% ================================================================
% XXI. DOCUMENT COMPILATION AND FINAL CHECK
% ================================================================

\noindent
\textbf{XXI.\quad DOCUMENT COMPILATION AND FINAL CHECK}

\vspace{0.15cm}

Before submission, the complete \LaTeX{} project shall compile
successfully and produce the final PDF report.

The final document shall not contain unresolved references such as:

\begin{verbatim}
??
[?]
\end{verbatim}

Before generating the final PDF, students shall check:

\begin{itemize}[leftmargin=0.8cm, itemsep=2pt]
    \item the Table of Contents;
    \item page numbering;
    \item chapter, section, and subsection numbering;
    \item figure numbering;
    \item table numbering;
    \item equation numbering;
    \item cross-references;
    \item citations;
    \item bibliography;
    \item font display;
    \item mathematical symbols;
    \item image quality;
    \item completeness of all required forms.
\end{itemize}

The document may be compiled using TeX Live, MiKTeX, or another
compatible \LaTeX{} environment. Regardless of the compilation
environment, the resulting PDF shall conform to the official Faculty
template.

\vspace{0.3cm}

% ================================================================
% XXII. CONTENT THAT STUDENTS MAY EDIT
% ================================================================

\noindent
\textbf{XXII.\quad CONTENT THAT STUDENTS MAY EDIT}

\vspace{0.15cm}

Students may edit or replace:

\begin{itemize}[leftmargin=0.8cm, itemsep=2pt]
    \item student name;
    \item class;
    \item internship organization;
    \item supervisor name;
    \item Abstract;
    \item List of Symbols and Abbreviations;
    \item chapter titles and chapter content;
    \item figures and photographs;
    \item tables;
    \item equations;
    \item citations and references.
\end{itemize}

Students shall not modify the main document formatting settings in the
preamble, including the document class, page geometry, line spacing,
heading formatting, table-of-contents formatting, font configuration,
page layout, or cover-page macros unless specifically instructed.

\vspace{0.3cm}

% ================================================================
% XXIII. FINAL SUBMISSION REQUIREMENTS
% ================================================================

\noindent
\textbf{XXIII.\quad FINAL SUBMISSION REQUIREMENTS}

\vspace{0.15cm}

Before submission, the final report shall:

\begin{itemize}[leftmargin=0.8cm, itemsep=2pt]
    \item compile without errors that affect the report content or presentation;
    \item contain all required forms and report sections;
    \item contain correct student, class, internship organization, and supervisor information;
    \item contain an updated Table of Contents, List of Figures, and List of Tables;
    \item contain no unused example content from the template;
    \item contain no placeholder text such as \textit{Title of Chapter 1},
    \textit{Section 1.1}, or other sample entries;
    \item contain all required citations and references;
    \item contain no unresolved cross-references.
\end{itemize}

If the Faculty requires submission of the \LaTeX{} source, the complete
project shall include all \verb|.tex| files, the bibliography database,
figures, and any other files required to reproduce the submitted PDF.

\vspace{0.4cm}

\noindent
\textbf{Important:}
Before submitting the final internship report, students shall remove or
comment out the following command in the main document:

\begin{verbatim}

% =================================================================
% REGULATIONS FOR PREPARING THE GRADUATION INTERNSHIP REPORT
% USING LATEX
% =================================================================
%
% This file is intended to be included from the main document with:
%
%     \input{Regulations.tex}
%
% IMPORTANT:
% This section is provided for guidance only. Students shall remove
% or comment out the \input{Regulations.tex} command before submitting
% the final internship report.
%
% =================================================================

\clearpage
\thispagestyle{empty}

\begin{center}
    {\fontsize{16pt}{20pt}\selectfont\bfseries
    SPECIFIC REGULATIONS FOR PREPARING\\
    GRADUATION INTERNSHIP REPORTS USING \LaTeX}
\end{center}

\vspace{0.5cm}

{\fontsize{13pt}{17pt}\selectfont
\setstretch{1.3}
\justifying

\noindent
\textbf{Note:}
This section is provided for guidance only and is included in the main
document using the command \verb|\input{Regulations.tex}|.
Students shall remove or comment out this command before submitting the
final internship report.

\vspace{0.4cm}

% ================================================================
% I. GENERAL REQUIREMENTS
% ================================================================

\noindent
\textbf{I.\quad GENERAL REQUIREMENTS}

\vspace{0.15cm}

The Graduation Internship Report shall be prepared using the official
\LaTeX{} template provided by the Faculty. Students shall preserve the
predefined page layout, typography, line spacing, heading styles,
numbering system, figure and table formatting, and bibliography settings
unless otherwise instructed by the Faculty or the supervisor.

The report shall be written in English.

Students are responsible for the academic and technical content of the
report. The official \LaTeX{} template is responsible for standardized
document formatting.

\vspace{0.3cm}

% ================================================================
% II. REPORT STRUCTURE
% ================================================================

\noindent
\textbf{II.\quad REPORT STRUCTURE}

\vspace{0.15cm}

The report shall normally be arranged in the following order:

\begin{enumerate}[label=\arabic*., leftmargin=1.2cm, itemsep=2pt]
    \item Outer Cover;
    \item Comments of the Internship Host Organization;
    \item Supervisor's Evaluation and Graduation Internship Defense Results;
    \item Table of Contents;
    \item List of Symbols and Abbreviations, if applicable;
    \item List of Figures;
    \item List of Tables;
    \item Abstract;
    \item Main Chapters of the Internship Report;
    \item References;
    \item Appendices, if applicable.
\end{enumerate}

The preliminary pages shall not be numbered as chapters. The main report
content shall be organized into consecutively numbered chapters beginning
with Chapter 1.

\vspace{0.3cm}

% ================================================================
% III. PAGE SIZE AND MARGINS
% ================================================================

\noindent
\textbf{III.\quad PAGE SIZE AND MARGINS}

\vspace{0.15cm}

The report shall be formatted on A4 paper with the following margins:

\begin{center}
\begin{tabular}{ll}
    Left margin:   & 3.5 cm \\
    Right margin:  & 2.0 cm \\
    Top margin:    & 2.5 cm \\
    Bottom margin: & 2.5 cm
\end{tabular}
\end{center}

These settings are predefined in the official template and shall not be
modified by students.

\vspace{0.3cm}

% ================================================================
% IV. FONT, FONT SIZE, AND LINE SPACING
% ================================================================

\noindent
\textbf{IV.\quad FONT, FONT SIZE, AND LINE SPACING}

\vspace{0.15cm}

The main text shall use a 13 pt font size with 1.3-line spacing.
A Times-family typeface shall be used in accordance with the official
template.

The text shall be fully justified.

Students shall not use repeated spaces, repeated manual line breaks, or
multiple blank lines to position text. Paragraph and heading spacing shall
be controlled by the template.

\vspace{0.3cm}

% ================================================================
% V. OUTER COVER
% ================================================================

\noindent
\textbf{V.\quad OUTER COVER}

\vspace{0.15cm}

The outer cover shall follow the official format provided in the template
and shall contain the following information:

\begin{itemize}[leftmargin=0.8cm, itemsep=2pt]
    \item Thai Nguyen University of Technology;
    \item Faculty of International Training;
    \item Engineering and Technology Department;
    \item University logo and Faculty logo;
    \item the title \textbf{INTERNSHIP REPORT};
    \item Internship Organization;
    \item Author;
    \item Class;
    \item Supervisor;
    \item place and year of completion.
\end{itemize}

Students shall only replace the required information fields and shall not
modify the cover-page layout unless instructed otherwise.

\vspace{0.3cm}

% ================================================================
% VI. COMMENTS OF THE INTERNSHIP HOST ORGANIZATION
% ================================================================

\noindent
\textbf{VI.\quad COMMENTS OF THE INTERNSHIP HOST ORGANIZATION}

\vspace{0.15cm}

The page entitled \textbf{COMMENTS OF THE INTERNSHIP HOST ORGANIZATION}
is reserved for the internship host organization.

It includes:

\begin{itemize}[leftmargin=0.8cm, itemsep=2pt]
    \item the student's full name;
    \item class;
    \item comments from the host organization;
    \item score;
    \item date;
    \item signature and official seal.
\end{itemize}

Students shall not alter or complete official evaluation fields unless
specifically authorized.

\vspace{0.3cm}

% ================================================================
% VII. SUPERVISOR'S EVALUATION AND DEFENSE RESULTS
% ================================================================

\noindent
\textbf{VII.\quad SUPERVISOR'S EVALUATION AND DEFENSE RESULTS}

\vspace{0.15cm}

The section entitled
\textbf{SUPERVISOR'S EVALUATION AND GRADUATION INTERNSHIP DEFENSE RESULTS}
is reserved for the supervisor and defense committee.

The evaluation may include:

\begin{itemize}[leftmargin=0.8cm, itemsep=2pt]
    \item Internship Progress and Student's Attitude;
    \item Report Content;
    \item Report Presentation;
    \item Other Comments;
    \item Supervisor's Opinion;
    \item Defense Result;
    \item Score.
\end{itemize}

Students shall not enter evaluation scores or defense results before they
are officially determined.

\vspace{0.3cm}

% ================================================================
% VIII. CHAPTERS, SECTIONS, AND SUBSECTIONS
% ================================================================

\noindent
\textbf{VIII.\quad CHAPTERS, SECTIONS, AND SUBSECTIONS}

\vspace{0.15cm}

The report structure shall be generated using the standard \LaTeX{}
commands:

\begin{verbatim}
\chapter{...}
\section{...}
\subsection{...}
\end{verbatim}

Chapter, section, and subsection numbers shall be generated automatically
by \LaTeX{}.

Students shall not manually type chapter or section numbers and shall not
reproduce heading styles using commands such as \verb|\textbf|,
\verb|\fontsize|, or \verb|\centering| instead of the appropriate
structural command.

Excessive subdivision should be avoided. In most cases, the hierarchy
Chapter--Section--Subsection is sufficient.

\vspace{0.3cm}

% ================================================================
% IX. ORGANIZATION OF THE MAIN CONTENT
% ================================================================

\noindent
\textbf{IX.\quad ORGANIZATION OF THE MAIN CONTENT}

\vspace{0.15cm}

The specific content of the report shall depend on the internship
organization and the assigned tasks. A typical structure may include:

\begin{itemize}[leftmargin=0.8cm, itemsep=2pt]
    \item \textbf{Chapter 1. Overview of the Internship Organization};
    \item \textbf{Chapter 2. Internship Activities and Assigned Tasks};
    \item \textbf{Chapter 3. Results, Experience, and Evaluation}.
\end{itemize}

These chapter titles are recommendations only and may be adapted to the
actual internship content.

The report shall reflect work actually performed, observed, analyzed, or
learned during the internship. Unrelated theoretical material should not
be included merely to increase report length.

\vspace{0.3cm}

% ================================================================
% X. TABLE OF CONTENTS
% ================================================================

\noindent
\textbf{X.\quad TABLE OF CONTENTS}

\vspace{0.15cm}

The Table of Contents shall be generated automatically using:

\begin{verbatim}
\tableofcontents
\end{verbatim}

Students shall not create or manually edit page numbers in the Table of
Contents.

The document shall be compiled sufficiently to ensure that all page
numbers and entries are correctly updated.

\vspace{0.3cm}

% ================================================================
% XI. LIST OF SYMBOLS AND ABBREVIATIONS
% ================================================================

\noindent
\textbf{XI.\quad LIST OF SYMBOLS AND ABBREVIATIONS}

\vspace{0.15cm}

Symbols and abbreviations that are used repeatedly in the report shall be
listed in the \textbf{List of Symbols and Abbreviations}.

The list may contain the following columns:

\begin{center}
\begin{tabular}{llll}
    No. & Symbol / Abbreviation & Unit & Definition
\end{tabular}
\end{center}

Only symbols and abbreviations actually used in the final report shall be
retained. Example entries contained in the template shall be removed if
they are not relevant to the student's report.

Each abbreviation should normally be defined when it first appears in
the main text.

\vspace{0.3cm}

% ================================================================
% XII. FIGURES, DIAGRAMS, AND PHOTOGRAPHS
% ================================================================

\noindent
\textbf{XII.\quad FIGURES, DIAGRAMS, AND PHOTOGRAPHS}

\vspace{0.15cm}

Figures shall be clear, legible, and directly relevant to the report
content.

Figures shall be inserted using the standard \verb|figure| environment.

Example:

\begin{verbatim}
\begin{figure}[H]
    \centering
    \includegraphics[width=0.8\textwidth]
                    {figures/company_structure.pdf}
    \caption{Organizational structure of the company}
    \label{fig:company-structure}
\end{figure}
\end{verbatim}

Figure captions shall be placed below the figure.

Each referenced figure shall have a unique \verb|\label| and shall be
referenced in the text using, for example:

\begin{verbatim}
Figure~\ref{fig:company-structure}
\end{verbatim}

Students shall not manually type figure numbers.

Figures reproduced or adapted from external sources shall include an
appropriate citation.

Photographs taken during the internship shall be relevant and sufficiently
clear. Students shall not include confidential, restricted, or proprietary
information without permission from the internship organization.

Vector graphics or high-resolution images are recommended whenever
possible.

\vspace{0.3cm}

% ================================================================
% XIII. TABLES
% ================================================================

\noindent
\textbf{XIII.\quad TABLES}

\vspace{0.15cm}

Tables shall be created using suitable \LaTeX{} environments such as
\verb|table|, \verb|tabular|, \verb|tabularx|, or \verb|longtable|.

Table captions shall be placed above the table.

Example:

\begin{verbatim}
\begin{table}[H]
    \centering
    \caption{Internship schedule}
    \label{tab:internship-schedule}

    \begin{tabular}{lll}
        \toprule
        Week & Task & Result \\
        \midrule
        1 & Orientation & Completed \\
        2 & System study & Completed \\
        \bottomrule
    \end{tabular}
\end{table}
\end{verbatim}

Tables shall be referenced using:

\begin{verbatim}
Table~\ref{tab:internship-schedule}
\end{verbatim}

Students shall not manually type table numbers.

Tables should be created directly in \LaTeX{} rather than inserted as
screenshots whenever practical.

\vspace{0.3cm}

% ================================================================
% XIV. EQUATIONS AND TECHNICAL QUANTITIES
% ================================================================

\noindent
\textbf{XIV.\quad EQUATIONS AND TECHNICAL QUANTITIES}

\vspace{0.15cm}

Mathematical expressions shall be typeset directly using \LaTeX{}.
Images or screenshots of equations shall not be used.

Short expressions may be written in inline mathematical mode, for example:

\begin{verbatim}
$R = 0.7\,\Omega$
\end{verbatim}

Important equations shall use a numbered equation environment.

Example:

\begin{verbatim}
\begin{equation}
    P = UI
    \label{eq:power}
\end{equation}
\end{verbatim}

Equations shall be referenced using:

\begin{verbatim}
Equation~\eqref{eq:power}
\end{verbatim}

Equation numbers shall not be manually entered.

SI units shall be used wherever applicable. The \verb|siunitx| package
should be used where appropriate.

Examples:

\begin{verbatim}
\SI{220}{\volt}
\SI{50}{\hertz}
\SI{2.5}{\ampere}
\end{verbatim}

Notation and units shall be consistent throughout the report.

\vspace{0.3cm}

% ================================================================
% XV. SOURCE CODE
% ================================================================

\noindent
\textbf{XV.\quad SOURCE CODE}

\vspace{0.15cm}

If the internship involves programming, PLCs, microcontrollers,
simulation, or software development, relevant source code may be included
in the report.

Source code shall be presented using the \verb|listings| environment or
another environment defined by the official template.

Example:

\begin{verbatim}
\begin{lstlisting}[language=st]
IF Start = TRUE THEN
    Motor := TRUE;
ELSE
    Motor := FALSE;
END_IF;
\end{lstlisting}
\end{verbatim}

Only code necessary for explaining the work should appear in the main
report. Long source-code listings should be moved to an appendix or
submitted as supplementary files if required.

Source code should not normally be inserted as screenshots.

\vspace{0.3cm}

% ================================================================
% XVI. ABSTRACT
% ================================================================

\noindent
\textbf{XVI.\quad ABSTRACT}

\vspace{0.15cm}

The Abstract shall briefly summarize:

\begin{itemize}[leftmargin=0.8cm, itemsep=2pt]
    \item the internship organization;
    \item the objectives of the internship;
    \item the main tasks performed;
    \item the main results, knowledge, and skills obtained.
\end{itemize}

The Abstract should be concise and should not reproduce the full content
of the report chapter by chapter.

\vspace{0.3cm}

% ================================================================
% XVII. REFERENCES AND CITATIONS
% ================================================================

\noindent
\textbf{XVII.\quad REFERENCES AND CITATIONS}

\vspace{0.15cm}

All books, journal papers, conference papers, standards, websites,
datasheets, manuals, technical documents, figures, tables, or other
external sources used in preparing the report shall be properly cited.

References shall be managed through the bibliography system provided with
the official template.

Bibliographic entries shall normally be stored in:

\begin{verbatim}
References.bib
\end{verbatim}

References shall be cited in the report using:

\begin{verbatim}
\cite{reference-key}
\end{verbatim}

Students shall not manually type reference numbers such as ``[1]'',
``[2]'', etc.

Titles of foreign-language publications shall remain in their original
language in the bibliography.

All borrowed ideas, data, figures, tables, equations, methods, or results
shall be properly acknowledged.

\vspace{0.3cm}

% ================================================================
% XVIII. CROSS-REFERENCING
% ================================================================

\noindent
\textbf{XVIII.\quad CROSS-REFERENCING}

\vspace{0.15cm}

Numbered elements shall be cross-referenced using the standard \LaTeX{}
mechanism.

Use the following commands where appropriate:

\begin{verbatim}
\label{...}
\ref{...}
\eqref{...}
\end{verbatim}

For example:

\begin{verbatim}
As shown in Figure~\ref{fig:system}, ...
\end{verbatim}

Students shall not manually type numbers such as ``Figure 2.3'' when
those numbers can be generated automatically.

\vspace{0.3cm}

% ================================================================
% XIX. PAGE NUMBERING
% ================================================================

\noindent
\textbf{XIX.\quad PAGE NUMBERING}

\vspace{0.15cm}

Cover pages and official evaluation forms at the beginning of the report
shall not display page numbers.

The main body of the report shall use Arabic page numbers beginning from
Chapter 1.

Page numbering shall be generated automatically by the official template.
Students shall not manually enter page numbers.

\vspace{0.3cm}

% ================================================================
% XX. PROJECT FILE ORGANIZATION
% ================================================================

\noindent
\textbf{XX.\quad PROJECT FILE ORGANIZATION}

\vspace{0.15cm}

Students should maintain a clear and organized \LaTeX{} project
structure. A recommended structure is:

\begin{verbatim}
InternshipReport/
|
|-- main.tex
|-- References.bib
|
|-- figures/
|   |-- company.pdf
|   |-- system.png
|   `-- ...
|
`-- chapters/
    |-- chapter1.tex
    |-- chapter2.tex
    `-- chapter3.tex
\end{verbatim}

Figures should normally be stored in the \verb|figures| directory
specified by the official template.

File and directory names should preferably contain Latin letters,
numbers, hyphens, or underscores.

Recommended examples include:

\begin{verbatim}
company_structure.pdf
motor_control.png
chapter_01.tex
\end{verbatim}

Spaces and unnecessary special characters in file names should be
avoided where possible.

\vspace{0.3cm}

% ================================================================
% XXI. DOCUMENT COMPILATION AND FINAL CHECK
% ================================================================

\noindent
\textbf{XXI.\quad DOCUMENT COMPILATION AND FINAL CHECK}

\vspace{0.15cm}

Before submission, the complete \LaTeX{} project shall compile
successfully and produce the final PDF report.

The final document shall not contain unresolved references such as:

\begin{verbatim}
??
[?]
\end{verbatim}

Before generating the final PDF, students shall check:

\begin{itemize}[leftmargin=0.8cm, itemsep=2pt]
    \item the Table of Contents;
    \item page numbering;
    \item chapter, section, and subsection numbering;
    \item figure numbering;
    \item table numbering;
    \item equation numbering;
    \item cross-references;
    \item citations;
    \item bibliography;
    \item font display;
    \item mathematical symbols;
    \item image quality;
    \item completeness of all required forms.
\end{itemize}

The document may be compiled using TeX Live, MiKTeX, or another
compatible \LaTeX{} environment. Regardless of the compilation
environment, the resulting PDF shall conform to the official Faculty
template.

\vspace{0.3cm}

% ================================================================
% XXII. CONTENT THAT STUDENTS MAY EDIT
% ================================================================

\noindent
\textbf{XXII.\quad CONTENT THAT STUDENTS MAY EDIT}

\vspace{0.15cm}

Students may edit or replace:

\begin{itemize}[leftmargin=0.8cm, itemsep=2pt]
    \item student name;
    \item class;
    \item internship organization;
    \item supervisor name;
    \item Abstract;
    \item List of Symbols and Abbreviations;
    \item chapter titles and chapter content;
    \item figures and photographs;
    \item tables;
    \item equations;
    \item citations and references.
\end{itemize}

Students shall not modify the main document formatting settings in the
preamble, including the document class, page geometry, line spacing,
heading formatting, table-of-contents formatting, font configuration,
page layout, or cover-page macros unless specifically instructed.

\vspace{0.3cm}

% ================================================================
% XXIII. FINAL SUBMISSION REQUIREMENTS
% ================================================================

\noindent
\textbf{XXIII.\quad FINAL SUBMISSION REQUIREMENTS}

\vspace{0.15cm}

Before submission, the final report shall:

\begin{itemize}[leftmargin=0.8cm, itemsep=2pt]
    \item compile without errors that affect the report content or presentation;
    \item contain all required forms and report sections;
    \item contain correct student, class, internship organization, and supervisor information;
    \item contain an updated Table of Contents, List of Figures, and List of Tables;
    \item contain no unused example content from the template;
    \item contain no placeholder text such as \textit{Title of Chapter 1},
    \textit{Section 1.1}, or other sample entries;
    \item contain all required citations and references;
    \item contain no unresolved cross-references.
\end{itemize}

If the Faculty requires submission of the \LaTeX{} source, the complete
project shall include all \verb|.tex| files, the bibliography database,
figures, and any other files required to reproduce the submitted PDF.

\vspace{0.4cm}

\noindent
\textbf{Important:}
Before submitting the final internship report, students shall remove or
comment out the following command in the main document:

\begin{verbatim}
\input{Regulations.tex}
\end{verbatim}

For example:

\begin{verbatim}
% \input{Regulations.tex}
\end{verbatim}

This removes the guidance pages from the final submitted report while
preserving the regulations file in the \LaTeX{} project for reference.

} % End 13pt content

\clearpage

\end{verbatim}

For example:

\begin{verbatim}
% 
% =================================================================
% REGULATIONS FOR PREPARING THE GRADUATION INTERNSHIP REPORT
% USING LATEX
% =================================================================
%
% This file is intended to be included from the main document with:
%
%     \input{Regulations.tex}
%
% IMPORTANT:
% This section is provided for guidance only. Students shall remove
% or comment out the \input{Regulations.tex} command before submitting
% the final internship report.
%
% =================================================================

\clearpage
\thispagestyle{empty}

\begin{center}
    {\fontsize{16pt}{20pt}\selectfont\bfseries
    SPECIFIC REGULATIONS FOR PREPARING\\
    GRADUATION INTERNSHIP REPORTS USING \LaTeX}
\end{center}

\vspace{0.5cm}

{\fontsize{13pt}{17pt}\selectfont
\setstretch{1.3}
\justifying

\noindent
\textbf{Note:}
This section is provided for guidance only and is included in the main
document using the command \verb|\input{Regulations.tex}|.
Students shall remove or comment out this command before submitting the
final internship report.

\vspace{0.4cm}

% ================================================================
% I. GENERAL REQUIREMENTS
% ================================================================

\noindent
\textbf{I.\quad GENERAL REQUIREMENTS}

\vspace{0.15cm}

The Graduation Internship Report shall be prepared using the official
\LaTeX{} template provided by the Faculty. Students shall preserve the
predefined page layout, typography, line spacing, heading styles,
numbering system, figure and table formatting, and bibliography settings
unless otherwise instructed by the Faculty or the supervisor.

The report shall be written in English.

Students are responsible for the academic and technical content of the
report. The official \LaTeX{} template is responsible for standardized
document formatting.

\vspace{0.3cm}

% ================================================================
% II. REPORT STRUCTURE
% ================================================================

\noindent
\textbf{II.\quad REPORT STRUCTURE}

\vspace{0.15cm}

The report shall normally be arranged in the following order:

\begin{enumerate}[label=\arabic*., leftmargin=1.2cm, itemsep=2pt]
    \item Outer Cover;
    \item Comments of the Internship Host Organization;
    \item Supervisor's Evaluation and Graduation Internship Defense Results;
    \item Table of Contents;
    \item List of Symbols and Abbreviations, if applicable;
    \item List of Figures;
    \item List of Tables;
    \item Abstract;
    \item Main Chapters of the Internship Report;
    \item References;
    \item Appendices, if applicable.
\end{enumerate}

The preliminary pages shall not be numbered as chapters. The main report
content shall be organized into consecutively numbered chapters beginning
with Chapter 1.

\vspace{0.3cm}

% ================================================================
% III. PAGE SIZE AND MARGINS
% ================================================================

\noindent
\textbf{III.\quad PAGE SIZE AND MARGINS}

\vspace{0.15cm}

The report shall be formatted on A4 paper with the following margins:

\begin{center}
\begin{tabular}{ll}
    Left margin:   & 3.5 cm \\
    Right margin:  & 2.0 cm \\
    Top margin:    & 2.5 cm \\
    Bottom margin: & 2.5 cm
\end{tabular}
\end{center}

These settings are predefined in the official template and shall not be
modified by students.

\vspace{0.3cm}

% ================================================================
% IV. FONT, FONT SIZE, AND LINE SPACING
% ================================================================

\noindent
\textbf{IV.\quad FONT, FONT SIZE, AND LINE SPACING}

\vspace{0.15cm}

The main text shall use a 13 pt font size with 1.3-line spacing.
A Times-family typeface shall be used in accordance with the official
template.

The text shall be fully justified.

Students shall not use repeated spaces, repeated manual line breaks, or
multiple blank lines to position text. Paragraph and heading spacing shall
be controlled by the template.

\vspace{0.3cm}

% ================================================================
% V. OUTER COVER
% ================================================================

\noindent
\textbf{V.\quad OUTER COVER}

\vspace{0.15cm}

The outer cover shall follow the official format provided in the template
and shall contain the following information:

\begin{itemize}[leftmargin=0.8cm, itemsep=2pt]
    \item Thai Nguyen University of Technology;
    \item Faculty of International Training;
    \item Engineering and Technology Department;
    \item University logo and Faculty logo;
    \item the title \textbf{INTERNSHIP REPORT};
    \item Internship Organization;
    \item Author;
    \item Class;
    \item Supervisor;
    \item place and year of completion.
\end{itemize}

Students shall only replace the required information fields and shall not
modify the cover-page layout unless instructed otherwise.

\vspace{0.3cm}

% ================================================================
% VI. COMMENTS OF THE INTERNSHIP HOST ORGANIZATION
% ================================================================

\noindent
\textbf{VI.\quad COMMENTS OF THE INTERNSHIP HOST ORGANIZATION}

\vspace{0.15cm}

The page entitled \textbf{COMMENTS OF THE INTERNSHIP HOST ORGANIZATION}
is reserved for the internship host organization.

It includes:

\begin{itemize}[leftmargin=0.8cm, itemsep=2pt]
    \item the student's full name;
    \item class;
    \item comments from the host organization;
    \item score;
    \item date;
    \item signature and official seal.
\end{itemize}

Students shall not alter or complete official evaluation fields unless
specifically authorized.

\vspace{0.3cm}

% ================================================================
% VII. SUPERVISOR'S EVALUATION AND DEFENSE RESULTS
% ================================================================

\noindent
\textbf{VII.\quad SUPERVISOR'S EVALUATION AND DEFENSE RESULTS}

\vspace{0.15cm}

The section entitled
\textbf{SUPERVISOR'S EVALUATION AND GRADUATION INTERNSHIP DEFENSE RESULTS}
is reserved for the supervisor and defense committee.

The evaluation may include:

\begin{itemize}[leftmargin=0.8cm, itemsep=2pt]
    \item Internship Progress and Student's Attitude;
    \item Report Content;
    \item Report Presentation;
    \item Other Comments;
    \item Supervisor's Opinion;
    \item Defense Result;
    \item Score.
\end{itemize}

Students shall not enter evaluation scores or defense results before they
are officially determined.

\vspace{0.3cm}

% ================================================================
% VIII. CHAPTERS, SECTIONS, AND SUBSECTIONS
% ================================================================

\noindent
\textbf{VIII.\quad CHAPTERS, SECTIONS, AND SUBSECTIONS}

\vspace{0.15cm}

The report structure shall be generated using the standard \LaTeX{}
commands:

\begin{verbatim}
\chapter{...}
\section{...}
\subsection{...}
\end{verbatim}

Chapter, section, and subsection numbers shall be generated automatically
by \LaTeX{}.

Students shall not manually type chapter or section numbers and shall not
reproduce heading styles using commands such as \verb|\textbf|,
\verb|\fontsize|, or \verb|\centering| instead of the appropriate
structural command.

Excessive subdivision should be avoided. In most cases, the hierarchy
Chapter--Section--Subsection is sufficient.

\vspace{0.3cm}

% ================================================================
% IX. ORGANIZATION OF THE MAIN CONTENT
% ================================================================

\noindent
\textbf{IX.\quad ORGANIZATION OF THE MAIN CONTENT}

\vspace{0.15cm}

The specific content of the report shall depend on the internship
organization and the assigned tasks. A typical structure may include:

\begin{itemize}[leftmargin=0.8cm, itemsep=2pt]
    \item \textbf{Chapter 1. Overview of the Internship Organization};
    \item \textbf{Chapter 2. Internship Activities and Assigned Tasks};
    \item \textbf{Chapter 3. Results, Experience, and Evaluation}.
\end{itemize}

These chapter titles are recommendations only and may be adapted to the
actual internship content.

The report shall reflect work actually performed, observed, analyzed, or
learned during the internship. Unrelated theoretical material should not
be included merely to increase report length.

\vspace{0.3cm}

% ================================================================
% X. TABLE OF CONTENTS
% ================================================================

\noindent
\textbf{X.\quad TABLE OF CONTENTS}

\vspace{0.15cm}

The Table of Contents shall be generated automatically using:

\begin{verbatim}
\tableofcontents
\end{verbatim}

Students shall not create or manually edit page numbers in the Table of
Contents.

The document shall be compiled sufficiently to ensure that all page
numbers and entries are correctly updated.

\vspace{0.3cm}

% ================================================================
% XI. LIST OF SYMBOLS AND ABBREVIATIONS
% ================================================================

\noindent
\textbf{XI.\quad LIST OF SYMBOLS AND ABBREVIATIONS}

\vspace{0.15cm}

Symbols and abbreviations that are used repeatedly in the report shall be
listed in the \textbf{List of Symbols and Abbreviations}.

The list may contain the following columns:

\begin{center}
\begin{tabular}{llll}
    No. & Symbol / Abbreviation & Unit & Definition
\end{tabular}
\end{center}

Only symbols and abbreviations actually used in the final report shall be
retained. Example entries contained in the template shall be removed if
they are not relevant to the student's report.

Each abbreviation should normally be defined when it first appears in
the main text.

\vspace{0.3cm}

% ================================================================
% XII. FIGURES, DIAGRAMS, AND PHOTOGRAPHS
% ================================================================

\noindent
\textbf{XII.\quad FIGURES, DIAGRAMS, AND PHOTOGRAPHS}

\vspace{0.15cm}

Figures shall be clear, legible, and directly relevant to the report
content.

Figures shall be inserted using the standard \verb|figure| environment.

Example:

\begin{verbatim}
\begin{figure}[H]
    \centering
    \includegraphics[width=0.8\textwidth]
                    {figures/company_structure.pdf}
    \caption{Organizational structure of the company}
    \label{fig:company-structure}
\end{figure}
\end{verbatim}

Figure captions shall be placed below the figure.

Each referenced figure shall have a unique \verb|\label| and shall be
referenced in the text using, for example:

\begin{verbatim}
Figure~\ref{fig:company-structure}
\end{verbatim}

Students shall not manually type figure numbers.

Figures reproduced or adapted from external sources shall include an
appropriate citation.

Photographs taken during the internship shall be relevant and sufficiently
clear. Students shall not include confidential, restricted, or proprietary
information without permission from the internship organization.

Vector graphics or high-resolution images are recommended whenever
possible.

\vspace{0.3cm}

% ================================================================
% XIII. TABLES
% ================================================================

\noindent
\textbf{XIII.\quad TABLES}

\vspace{0.15cm}

Tables shall be created using suitable \LaTeX{} environments such as
\verb|table|, \verb|tabular|, \verb|tabularx|, or \verb|longtable|.

Table captions shall be placed above the table.

Example:

\begin{verbatim}
\begin{table}[H]
    \centering
    \caption{Internship schedule}
    \label{tab:internship-schedule}

    \begin{tabular}{lll}
        \toprule
        Week & Task & Result \\
        \midrule
        1 & Orientation & Completed \\
        2 & System study & Completed \\
        \bottomrule
    \end{tabular}
\end{table}
\end{verbatim}

Tables shall be referenced using:

\begin{verbatim}
Table~\ref{tab:internship-schedule}
\end{verbatim}

Students shall not manually type table numbers.

Tables should be created directly in \LaTeX{} rather than inserted as
screenshots whenever practical.

\vspace{0.3cm}

% ================================================================
% XIV. EQUATIONS AND TECHNICAL QUANTITIES
% ================================================================

\noindent
\textbf{XIV.\quad EQUATIONS AND TECHNICAL QUANTITIES}

\vspace{0.15cm}

Mathematical expressions shall be typeset directly using \LaTeX{}.
Images or screenshots of equations shall not be used.

Short expressions may be written in inline mathematical mode, for example:

\begin{verbatim}
$R = 0.7\,\Omega$
\end{verbatim}

Important equations shall use a numbered equation environment.

Example:

\begin{verbatim}
\begin{equation}
    P = UI
    \label{eq:power}
\end{equation}
\end{verbatim}

Equations shall be referenced using:

\begin{verbatim}
Equation~\eqref{eq:power}
\end{verbatim}

Equation numbers shall not be manually entered.

SI units shall be used wherever applicable. The \verb|siunitx| package
should be used where appropriate.

Examples:

\begin{verbatim}
\SI{220}{\volt}
\SI{50}{\hertz}
\SI{2.5}{\ampere}
\end{verbatim}

Notation and units shall be consistent throughout the report.

\vspace{0.3cm}

% ================================================================
% XV. SOURCE CODE
% ================================================================

\noindent
\textbf{XV.\quad SOURCE CODE}

\vspace{0.15cm}

If the internship involves programming, PLCs, microcontrollers,
simulation, or software development, relevant source code may be included
in the report.

Source code shall be presented using the \verb|listings| environment or
another environment defined by the official template.

Example:

\begin{verbatim}
\begin{lstlisting}[language=st]
IF Start = TRUE THEN
    Motor := TRUE;
ELSE
    Motor := FALSE;
END_IF;
\end{lstlisting}
\end{verbatim}

Only code necessary for explaining the work should appear in the main
report. Long source-code listings should be moved to an appendix or
submitted as supplementary files if required.

Source code should not normally be inserted as screenshots.

\vspace{0.3cm}

% ================================================================
% XVI. ABSTRACT
% ================================================================

\noindent
\textbf{XVI.\quad ABSTRACT}

\vspace{0.15cm}

The Abstract shall briefly summarize:

\begin{itemize}[leftmargin=0.8cm, itemsep=2pt]
    \item the internship organization;
    \item the objectives of the internship;
    \item the main tasks performed;
    \item the main results, knowledge, and skills obtained.
\end{itemize}

The Abstract should be concise and should not reproduce the full content
of the report chapter by chapter.

\vspace{0.3cm}

% ================================================================
% XVII. REFERENCES AND CITATIONS
% ================================================================

\noindent
\textbf{XVII.\quad REFERENCES AND CITATIONS}

\vspace{0.15cm}

All books, journal papers, conference papers, standards, websites,
datasheets, manuals, technical documents, figures, tables, or other
external sources used in preparing the report shall be properly cited.

References shall be managed through the bibliography system provided with
the official template.

Bibliographic entries shall normally be stored in:

\begin{verbatim}
References.bib
\end{verbatim}

References shall be cited in the report using:

\begin{verbatim}
\cite{reference-key}
\end{verbatim}

Students shall not manually type reference numbers such as ``[1]'',
``[2]'', etc.

Titles of foreign-language publications shall remain in their original
language in the bibliography.

All borrowed ideas, data, figures, tables, equations, methods, or results
shall be properly acknowledged.

\vspace{0.3cm}

% ================================================================
% XVIII. CROSS-REFERENCING
% ================================================================

\noindent
\textbf{XVIII.\quad CROSS-REFERENCING}

\vspace{0.15cm}

Numbered elements shall be cross-referenced using the standard \LaTeX{}
mechanism.

Use the following commands where appropriate:

\begin{verbatim}
\label{...}
\ref{...}
\eqref{...}
\end{verbatim}

For example:

\begin{verbatim}
As shown in Figure~\ref{fig:system}, ...
\end{verbatim}

Students shall not manually type numbers such as ``Figure 2.3'' when
those numbers can be generated automatically.

\vspace{0.3cm}

% ================================================================
% XIX. PAGE NUMBERING
% ================================================================

\noindent
\textbf{XIX.\quad PAGE NUMBERING}

\vspace{0.15cm}

Cover pages and official evaluation forms at the beginning of the report
shall not display page numbers.

The main body of the report shall use Arabic page numbers beginning from
Chapter 1.

Page numbering shall be generated automatically by the official template.
Students shall not manually enter page numbers.

\vspace{0.3cm}

% ================================================================
% XX. PROJECT FILE ORGANIZATION
% ================================================================

\noindent
\textbf{XX.\quad PROJECT FILE ORGANIZATION}

\vspace{0.15cm}

Students should maintain a clear and organized \LaTeX{} project
structure. A recommended structure is:

\begin{verbatim}
InternshipReport/
|
|-- main.tex
|-- References.bib
|
|-- figures/
|   |-- company.pdf
|   |-- system.png
|   `-- ...
|
`-- chapters/
    |-- chapter1.tex
    |-- chapter2.tex
    `-- chapter3.tex
\end{verbatim}

Figures should normally be stored in the \verb|figures| directory
specified by the official template.

File and directory names should preferably contain Latin letters,
numbers, hyphens, or underscores.

Recommended examples include:

\begin{verbatim}
company_structure.pdf
motor_control.png
chapter_01.tex
\end{verbatim}

Spaces and unnecessary special characters in file names should be
avoided where possible.

\vspace{0.3cm}

% ================================================================
% XXI. DOCUMENT COMPILATION AND FINAL CHECK
% ================================================================

\noindent
\textbf{XXI.\quad DOCUMENT COMPILATION AND FINAL CHECK}

\vspace{0.15cm}

Before submission, the complete \LaTeX{} project shall compile
successfully and produce the final PDF report.

The final document shall not contain unresolved references such as:

\begin{verbatim}
??
[?]
\end{verbatim}

Before generating the final PDF, students shall check:

\begin{itemize}[leftmargin=0.8cm, itemsep=2pt]
    \item the Table of Contents;
    \item page numbering;
    \item chapter, section, and subsection numbering;
    \item figure numbering;
    \item table numbering;
    \item equation numbering;
    \item cross-references;
    \item citations;
    \item bibliography;
    \item font display;
    \item mathematical symbols;
    \item image quality;
    \item completeness of all required forms.
\end{itemize}

The document may be compiled using TeX Live, MiKTeX, or another
compatible \LaTeX{} environment. Regardless of the compilation
environment, the resulting PDF shall conform to the official Faculty
template.

\vspace{0.3cm}

% ================================================================
% XXII. CONTENT THAT STUDENTS MAY EDIT
% ================================================================

\noindent
\textbf{XXII.\quad CONTENT THAT STUDENTS MAY EDIT}

\vspace{0.15cm}

Students may edit or replace:

\begin{itemize}[leftmargin=0.8cm, itemsep=2pt]
    \item student name;
    \item class;
    \item internship organization;
    \item supervisor name;
    \item Abstract;
    \item List of Symbols and Abbreviations;
    \item chapter titles and chapter content;
    \item figures and photographs;
    \item tables;
    \item equations;
    \item citations and references.
\end{itemize}

Students shall not modify the main document formatting settings in the
preamble, including the document class, page geometry, line spacing,
heading formatting, table-of-contents formatting, font configuration,
page layout, or cover-page macros unless specifically instructed.

\vspace{0.3cm}

% ================================================================
% XXIII. FINAL SUBMISSION REQUIREMENTS
% ================================================================

\noindent
\textbf{XXIII.\quad FINAL SUBMISSION REQUIREMENTS}

\vspace{0.15cm}

Before submission, the final report shall:

\begin{itemize}[leftmargin=0.8cm, itemsep=2pt]
    \item compile without errors that affect the report content or presentation;
    \item contain all required forms and report sections;
    \item contain correct student, class, internship organization, and supervisor information;
    \item contain an updated Table of Contents, List of Figures, and List of Tables;
    \item contain no unused example content from the template;
    \item contain no placeholder text such as \textit{Title of Chapter 1},
    \textit{Section 1.1}, or other sample entries;
    \item contain all required citations and references;
    \item contain no unresolved cross-references.
\end{itemize}

If the Faculty requires submission of the \LaTeX{} source, the complete
project shall include all \verb|.tex| files, the bibliography database,
figures, and any other files required to reproduce the submitted PDF.

\vspace{0.4cm}

\noindent
\textbf{Important:}
Before submitting the final internship report, students shall remove or
comment out the following command in the main document:

\begin{verbatim}
\input{Regulations.tex}
\end{verbatim}

For example:

\begin{verbatim}
% \input{Regulations.tex}
\end{verbatim}

This removes the guidance pages from the final submitted report while
preserving the regulations file in the \LaTeX{} project for reference.

} % End 13pt content

\clearpage

\end{verbatim}

This removes the guidance pages from the final submitted report while
preserving the regulations file in the \LaTeX{} project for reference.

} % End 13pt content

\clearpage

\end{verbatim}

This removes the guidance pages from the final submitted report while
preserving the regulations file in the \LaTeX{} project for reference.

} % End 13pt content

\clearpage
 command before submitting
% the final internship report.
%
% =================================================================

\clearpage
\thispagestyle{empty}

\begin{center}
    {\fontsize{16pt}{20pt}\selectfont\bfseries
    SPECIFIC REGULATIONS FOR PREPARING\\
    GRADUATION INTERNSHIP REPORTS USING \LaTeX}
\end{center}

\vspace{0.5cm}

{\fontsize{13pt}{17pt}\selectfont
\setstretch{1.3}
\justifying

\noindent
\textbf{Note:}
This section is provided for guidance only and is included in the main
document using the command \verb|
% =================================================================
% REGULATIONS FOR PREPARING THE GRADUATION INTERNSHIP REPORT
% USING LATEX
% =================================================================
%
% This file is intended to be included from the main document with:
%
%     
% =================================================================
% REGULATIONS FOR PREPARING THE GRADUATION INTERNSHIP REPORT
% USING LATEX
% =================================================================
%
% This file is intended to be included from the main document with:
%
%     
% =================================================================
% REGULATIONS FOR PREPARING THE GRADUATION INTERNSHIP REPORT
% USING LATEX
% =================================================================
%
% This file is intended to be included from the main document with:
%
%     \input{Regulations.tex}
%
% IMPORTANT:
% This section is provided for guidance only. Students shall remove
% or comment out the \input{Regulations.tex} command before submitting
% the final internship report.
%
% =================================================================

\clearpage
\thispagestyle{empty}

\begin{center}
    {\fontsize{16pt}{20pt}\selectfont\bfseries
    SPECIFIC REGULATIONS FOR PREPARING\\
    GRADUATION INTERNSHIP REPORTS USING \LaTeX}
\end{center}

\vspace{0.5cm}

{\fontsize{13pt}{17pt}\selectfont
\setstretch{1.3}
\justifying

\noindent
\textbf{Note:}
This section is provided for guidance only and is included in the main
document using the command \verb|\input{Regulations.tex}|.
Students shall remove or comment out this command before submitting the
final internship report.

\vspace{0.4cm}

% ================================================================
% I. GENERAL REQUIREMENTS
% ================================================================

\noindent
\textbf{I.\quad GENERAL REQUIREMENTS}

\vspace{0.15cm}

The Graduation Internship Report shall be prepared using the official
\LaTeX{} template provided by the Faculty. Students shall preserve the
predefined page layout, typography, line spacing, heading styles,
numbering system, figure and table formatting, and bibliography settings
unless otherwise instructed by the Faculty or the supervisor.

The report shall be written in English.

Students are responsible for the academic and technical content of the
report. The official \LaTeX{} template is responsible for standardized
document formatting.

\vspace{0.3cm}

% ================================================================
% II. REPORT STRUCTURE
% ================================================================

\noindent
\textbf{II.\quad REPORT STRUCTURE}

\vspace{0.15cm}

The report shall normally be arranged in the following order:

\begin{enumerate}[label=\arabic*., leftmargin=1.2cm, itemsep=2pt]
    \item Outer Cover;
    \item Comments of the Internship Host Organization;
    \item Supervisor's Evaluation and Graduation Internship Defense Results;
    \item Table of Contents;
    \item List of Symbols and Abbreviations, if applicable;
    \item List of Figures;
    \item List of Tables;
    \item Abstract;
    \item Main Chapters of the Internship Report;
    \item References;
    \item Appendices, if applicable.
\end{enumerate}

The preliminary pages shall not be numbered as chapters. The main report
content shall be organized into consecutively numbered chapters beginning
with Chapter 1.

\vspace{0.3cm}

% ================================================================
% III. PAGE SIZE AND MARGINS
% ================================================================

\noindent
\textbf{III.\quad PAGE SIZE AND MARGINS}

\vspace{0.15cm}

The report shall be formatted on A4 paper with the following margins:

\begin{center}
\begin{tabular}{ll}
    Left margin:   & 3.5 cm \\
    Right margin:  & 2.0 cm \\
    Top margin:    & 2.5 cm \\
    Bottom margin: & 2.5 cm
\end{tabular}
\end{center}

These settings are predefined in the official template and shall not be
modified by students.

\vspace{0.3cm}

% ================================================================
% IV. FONT, FONT SIZE, AND LINE SPACING
% ================================================================

\noindent
\textbf{IV.\quad FONT, FONT SIZE, AND LINE SPACING}

\vspace{0.15cm}

The main text shall use a 13 pt font size with 1.3-line spacing.
A Times-family typeface shall be used in accordance with the official
template.

The text shall be fully justified.

Students shall not use repeated spaces, repeated manual line breaks, or
multiple blank lines to position text. Paragraph and heading spacing shall
be controlled by the template.

\vspace{0.3cm}

% ================================================================
% V. OUTER COVER
% ================================================================

\noindent
\textbf{V.\quad OUTER COVER}

\vspace{0.15cm}

The outer cover shall follow the official format provided in the template
and shall contain the following information:

\begin{itemize}[leftmargin=0.8cm, itemsep=2pt]
    \item Thai Nguyen University of Technology;
    \item Faculty of International Training;
    \item Engineering and Technology Department;
    \item University logo and Faculty logo;
    \item the title \textbf{INTERNSHIP REPORT};
    \item Internship Organization;
    \item Author;
    \item Class;
    \item Supervisor;
    \item place and year of completion.
\end{itemize}

Students shall only replace the required information fields and shall not
modify the cover-page layout unless instructed otherwise.

\vspace{0.3cm}

% ================================================================
% VI. COMMENTS OF THE INTERNSHIP HOST ORGANIZATION
% ================================================================

\noindent
\textbf{VI.\quad COMMENTS OF THE INTERNSHIP HOST ORGANIZATION}

\vspace{0.15cm}

The page entitled \textbf{COMMENTS OF THE INTERNSHIP HOST ORGANIZATION}
is reserved for the internship host organization.

It includes:

\begin{itemize}[leftmargin=0.8cm, itemsep=2pt]
    \item the student's full name;
    \item class;
    \item comments from the host organization;
    \item score;
    \item date;
    \item signature and official seal.
\end{itemize}

Students shall not alter or complete official evaluation fields unless
specifically authorized.

\vspace{0.3cm}

% ================================================================
% VII. SUPERVISOR'S EVALUATION AND DEFENSE RESULTS
% ================================================================

\noindent
\textbf{VII.\quad SUPERVISOR'S EVALUATION AND DEFENSE RESULTS}

\vspace{0.15cm}

The section entitled
\textbf{SUPERVISOR'S EVALUATION AND GRADUATION INTERNSHIP DEFENSE RESULTS}
is reserved for the supervisor and defense committee.

The evaluation may include:

\begin{itemize}[leftmargin=0.8cm, itemsep=2pt]
    \item Internship Progress and Student's Attitude;
    \item Report Content;
    \item Report Presentation;
    \item Other Comments;
    \item Supervisor's Opinion;
    \item Defense Result;
    \item Score.
\end{itemize}

Students shall not enter evaluation scores or defense results before they
are officially determined.

\vspace{0.3cm}

% ================================================================
% VIII. CHAPTERS, SECTIONS, AND SUBSECTIONS
% ================================================================

\noindent
\textbf{VIII.\quad CHAPTERS, SECTIONS, AND SUBSECTIONS}

\vspace{0.15cm}

The report structure shall be generated using the standard \LaTeX{}
commands:

\begin{verbatim}
\chapter{...}
\section{...}
\subsection{...}
\end{verbatim}

Chapter, section, and subsection numbers shall be generated automatically
by \LaTeX{}.

Students shall not manually type chapter or section numbers and shall not
reproduce heading styles using commands such as \verb|\textbf|,
\verb|\fontsize|, or \verb|\centering| instead of the appropriate
structural command.

Excessive subdivision should be avoided. In most cases, the hierarchy
Chapter--Section--Subsection is sufficient.

\vspace{0.3cm}

% ================================================================
% IX. ORGANIZATION OF THE MAIN CONTENT
% ================================================================

\noindent
\textbf{IX.\quad ORGANIZATION OF THE MAIN CONTENT}

\vspace{0.15cm}

The specific content of the report shall depend on the internship
organization and the assigned tasks. A typical structure may include:

\begin{itemize}[leftmargin=0.8cm, itemsep=2pt]
    \item \textbf{Chapter 1. Overview of the Internship Organization};
    \item \textbf{Chapter 2. Internship Activities and Assigned Tasks};
    \item \textbf{Chapter 3. Results, Experience, and Evaluation}.
\end{itemize}

These chapter titles are recommendations only and may be adapted to the
actual internship content.

The report shall reflect work actually performed, observed, analyzed, or
learned during the internship. Unrelated theoretical material should not
be included merely to increase report length.

\vspace{0.3cm}

% ================================================================
% X. TABLE OF CONTENTS
% ================================================================

\noindent
\textbf{X.\quad TABLE OF CONTENTS}

\vspace{0.15cm}

The Table of Contents shall be generated automatically using:

\begin{verbatim}
\tableofcontents
\end{verbatim}

Students shall not create or manually edit page numbers in the Table of
Contents.

The document shall be compiled sufficiently to ensure that all page
numbers and entries are correctly updated.

\vspace{0.3cm}

% ================================================================
% XI. LIST OF SYMBOLS AND ABBREVIATIONS
% ================================================================

\noindent
\textbf{XI.\quad LIST OF SYMBOLS AND ABBREVIATIONS}

\vspace{0.15cm}

Symbols and abbreviations that are used repeatedly in the report shall be
listed in the \textbf{List of Symbols and Abbreviations}.

The list may contain the following columns:

\begin{center}
\begin{tabular}{llll}
    No. & Symbol / Abbreviation & Unit & Definition
\end{tabular}
\end{center}

Only symbols and abbreviations actually used in the final report shall be
retained. Example entries contained in the template shall be removed if
they are not relevant to the student's report.

Each abbreviation should normally be defined when it first appears in
the main text.

\vspace{0.3cm}

% ================================================================
% XII. FIGURES, DIAGRAMS, AND PHOTOGRAPHS
% ================================================================

\noindent
\textbf{XII.\quad FIGURES, DIAGRAMS, AND PHOTOGRAPHS}

\vspace{0.15cm}

Figures shall be clear, legible, and directly relevant to the report
content.

Figures shall be inserted using the standard \verb|figure| environment.

Example:

\begin{verbatim}
\begin{figure}[H]
    \centering
    \includegraphics[width=0.8\textwidth]
                    {figures/company_structure.pdf}
    \caption{Organizational structure of the company}
    \label{fig:company-structure}
\end{figure}
\end{verbatim}

Figure captions shall be placed below the figure.

Each referenced figure shall have a unique \verb|\label| and shall be
referenced in the text using, for example:

\begin{verbatim}
Figure~\ref{fig:company-structure}
\end{verbatim}

Students shall not manually type figure numbers.

Figures reproduced or adapted from external sources shall include an
appropriate citation.

Photographs taken during the internship shall be relevant and sufficiently
clear. Students shall not include confidential, restricted, or proprietary
information without permission from the internship organization.

Vector graphics or high-resolution images are recommended whenever
possible.

\vspace{0.3cm}

% ================================================================
% XIII. TABLES
% ================================================================

\noindent
\textbf{XIII.\quad TABLES}

\vspace{0.15cm}

Tables shall be created using suitable \LaTeX{} environments such as
\verb|table|, \verb|tabular|, \verb|tabularx|, or \verb|longtable|.

Table captions shall be placed above the table.

Example:

\begin{verbatim}
\begin{table}[H]
    \centering
    \caption{Internship schedule}
    \label{tab:internship-schedule}

    \begin{tabular}{lll}
        \toprule
        Week & Task & Result \\
        \midrule
        1 & Orientation & Completed \\
        2 & System study & Completed \\
        \bottomrule
    \end{tabular}
\end{table}
\end{verbatim}

Tables shall be referenced using:

\begin{verbatim}
Table~\ref{tab:internship-schedule}
\end{verbatim}

Students shall not manually type table numbers.

Tables should be created directly in \LaTeX{} rather than inserted as
screenshots whenever practical.

\vspace{0.3cm}

% ================================================================
% XIV. EQUATIONS AND TECHNICAL QUANTITIES
% ================================================================

\noindent
\textbf{XIV.\quad EQUATIONS AND TECHNICAL QUANTITIES}

\vspace{0.15cm}

Mathematical expressions shall be typeset directly using \LaTeX{}.
Images or screenshots of equations shall not be used.

Short expressions may be written in inline mathematical mode, for example:

\begin{verbatim}
$R = 0.7\,\Omega$
\end{verbatim}

Important equations shall use a numbered equation environment.

Example:

\begin{verbatim}
\begin{equation}
    P = UI
    \label{eq:power}
\end{equation}
\end{verbatim}

Equations shall be referenced using:

\begin{verbatim}
Equation~\eqref{eq:power}
\end{verbatim}

Equation numbers shall not be manually entered.

SI units shall be used wherever applicable. The \verb|siunitx| package
should be used where appropriate.

Examples:

\begin{verbatim}
\SI{220}{\volt}
\SI{50}{\hertz}
\SI{2.5}{\ampere}
\end{verbatim}

Notation and units shall be consistent throughout the report.

\vspace{0.3cm}

% ================================================================
% XV. SOURCE CODE
% ================================================================

\noindent
\textbf{XV.\quad SOURCE CODE}

\vspace{0.15cm}

If the internship involves programming, PLCs, microcontrollers,
simulation, or software development, relevant source code may be included
in the report.

Source code shall be presented using the \verb|listings| environment or
another environment defined by the official template.

Example:

\begin{verbatim}
\begin{lstlisting}[language=st]
IF Start = TRUE THEN
    Motor := TRUE;
ELSE
    Motor := FALSE;
END_IF;
\end{lstlisting}
\end{verbatim}

Only code necessary for explaining the work should appear in the main
report. Long source-code listings should be moved to an appendix or
submitted as supplementary files if required.

Source code should not normally be inserted as screenshots.

\vspace{0.3cm}

% ================================================================
% XVI. ABSTRACT
% ================================================================

\noindent
\textbf{XVI.\quad ABSTRACT}

\vspace{0.15cm}

The Abstract shall briefly summarize:

\begin{itemize}[leftmargin=0.8cm, itemsep=2pt]
    \item the internship organization;
    \item the objectives of the internship;
    \item the main tasks performed;
    \item the main results, knowledge, and skills obtained.
\end{itemize}

The Abstract should be concise and should not reproduce the full content
of the report chapter by chapter.

\vspace{0.3cm}

% ================================================================
% XVII. REFERENCES AND CITATIONS
% ================================================================

\noindent
\textbf{XVII.\quad REFERENCES AND CITATIONS}

\vspace{0.15cm}

All books, journal papers, conference papers, standards, websites,
datasheets, manuals, technical documents, figures, tables, or other
external sources used in preparing the report shall be properly cited.

References shall be managed through the bibliography system provided with
the official template.

Bibliographic entries shall normally be stored in:

\begin{verbatim}
References.bib
\end{verbatim}

References shall be cited in the report using:

\begin{verbatim}
\cite{reference-key}
\end{verbatim}

Students shall not manually type reference numbers such as ``[1]'',
``[2]'', etc.

Titles of foreign-language publications shall remain in their original
language in the bibliography.

All borrowed ideas, data, figures, tables, equations, methods, or results
shall be properly acknowledged.

\vspace{0.3cm}

% ================================================================
% XVIII. CROSS-REFERENCING
% ================================================================

\noindent
\textbf{XVIII.\quad CROSS-REFERENCING}

\vspace{0.15cm}

Numbered elements shall be cross-referenced using the standard \LaTeX{}
mechanism.

Use the following commands where appropriate:

\begin{verbatim}
\label{...}
\ref{...}
\eqref{...}
\end{verbatim}

For example:

\begin{verbatim}
As shown in Figure~\ref{fig:system}, ...
\end{verbatim}

Students shall not manually type numbers such as ``Figure 2.3'' when
those numbers can be generated automatically.

\vspace{0.3cm}

% ================================================================
% XIX. PAGE NUMBERING
% ================================================================

\noindent
\textbf{XIX.\quad PAGE NUMBERING}

\vspace{0.15cm}

Cover pages and official evaluation forms at the beginning of the report
shall not display page numbers.

The main body of the report shall use Arabic page numbers beginning from
Chapter 1.

Page numbering shall be generated automatically by the official template.
Students shall not manually enter page numbers.

\vspace{0.3cm}

% ================================================================
% XX. PROJECT FILE ORGANIZATION
% ================================================================

\noindent
\textbf{XX.\quad PROJECT FILE ORGANIZATION}

\vspace{0.15cm}

Students should maintain a clear and organized \LaTeX{} project
structure. A recommended structure is:

\begin{verbatim}
InternshipReport/
|
|-- main.tex
|-- References.bib
|
|-- figures/
|   |-- company.pdf
|   |-- system.png
|   `-- ...
|
`-- chapters/
    |-- chapter1.tex
    |-- chapter2.tex
    `-- chapter3.tex
\end{verbatim}

Figures should normally be stored in the \verb|figures| directory
specified by the official template.

File and directory names should preferably contain Latin letters,
numbers, hyphens, or underscores.

Recommended examples include:

\begin{verbatim}
company_structure.pdf
motor_control.png
chapter_01.tex
\end{verbatim}

Spaces and unnecessary special characters in file names should be
avoided where possible.

\vspace{0.3cm}

% ================================================================
% XXI. DOCUMENT COMPILATION AND FINAL CHECK
% ================================================================

\noindent
\textbf{XXI.\quad DOCUMENT COMPILATION AND FINAL CHECK}

\vspace{0.15cm}

Before submission, the complete \LaTeX{} project shall compile
successfully and produce the final PDF report.

The final document shall not contain unresolved references such as:

\begin{verbatim}
??
[?]
\end{verbatim}

Before generating the final PDF, students shall check:

\begin{itemize}[leftmargin=0.8cm, itemsep=2pt]
    \item the Table of Contents;
    \item page numbering;
    \item chapter, section, and subsection numbering;
    \item figure numbering;
    \item table numbering;
    \item equation numbering;
    \item cross-references;
    \item citations;
    \item bibliography;
    \item font display;
    \item mathematical symbols;
    \item image quality;
    \item completeness of all required forms.
\end{itemize}

The document may be compiled using TeX Live, MiKTeX, or another
compatible \LaTeX{} environment. Regardless of the compilation
environment, the resulting PDF shall conform to the official Faculty
template.

\vspace{0.3cm}

% ================================================================
% XXII. CONTENT THAT STUDENTS MAY EDIT
% ================================================================

\noindent
\textbf{XXII.\quad CONTENT THAT STUDENTS MAY EDIT}

\vspace{0.15cm}

Students may edit or replace:

\begin{itemize}[leftmargin=0.8cm, itemsep=2pt]
    \item student name;
    \item class;
    \item internship organization;
    \item supervisor name;
    \item Abstract;
    \item List of Symbols and Abbreviations;
    \item chapter titles and chapter content;
    \item figures and photographs;
    \item tables;
    \item equations;
    \item citations and references.
\end{itemize}

Students shall not modify the main document formatting settings in the
preamble, including the document class, page geometry, line spacing,
heading formatting, table-of-contents formatting, font configuration,
page layout, or cover-page macros unless specifically instructed.

\vspace{0.3cm}

% ================================================================
% XXIII. FINAL SUBMISSION REQUIREMENTS
% ================================================================

\noindent
\textbf{XXIII.\quad FINAL SUBMISSION REQUIREMENTS}

\vspace{0.15cm}

Before submission, the final report shall:

\begin{itemize}[leftmargin=0.8cm, itemsep=2pt]
    \item compile without errors that affect the report content or presentation;
    \item contain all required forms and report sections;
    \item contain correct student, class, internship organization, and supervisor information;
    \item contain an updated Table of Contents, List of Figures, and List of Tables;
    \item contain no unused example content from the template;
    \item contain no placeholder text such as \textit{Title of Chapter 1},
    \textit{Section 1.1}, or other sample entries;
    \item contain all required citations and references;
    \item contain no unresolved cross-references.
\end{itemize}

If the Faculty requires submission of the \LaTeX{} source, the complete
project shall include all \verb|.tex| files, the bibliography database,
figures, and any other files required to reproduce the submitted PDF.

\vspace{0.4cm}

\noindent
\textbf{Important:}
Before submitting the final internship report, students shall remove or
comment out the following command in the main document:

\begin{verbatim}
\input{Regulations.tex}
\end{verbatim}

For example:

\begin{verbatim}
% \input{Regulations.tex}
\end{verbatim}

This removes the guidance pages from the final submitted report while
preserving the regulations file in the \LaTeX{} project for reference.

} % End 13pt content

\clearpage

%
% IMPORTANT:
% This section is provided for guidance only. Students shall remove
% or comment out the 
% =================================================================
% REGULATIONS FOR PREPARING THE GRADUATION INTERNSHIP REPORT
% USING LATEX
% =================================================================
%
% This file is intended to be included from the main document with:
%
%     \input{Regulations.tex}
%
% IMPORTANT:
% This section is provided for guidance only. Students shall remove
% or comment out the \input{Regulations.tex} command before submitting
% the final internship report.
%
% =================================================================

\clearpage
\thispagestyle{empty}

\begin{center}
    {\fontsize{16pt}{20pt}\selectfont\bfseries
    SPECIFIC REGULATIONS FOR PREPARING\\
    GRADUATION INTERNSHIP REPORTS USING \LaTeX}
\end{center}

\vspace{0.5cm}

{\fontsize{13pt}{17pt}\selectfont
\setstretch{1.3}
\justifying

\noindent
\textbf{Note:}
This section is provided for guidance only and is included in the main
document using the command \verb|\input{Regulations.tex}|.
Students shall remove or comment out this command before submitting the
final internship report.

\vspace{0.4cm}

% ================================================================
% I. GENERAL REQUIREMENTS
% ================================================================

\noindent
\textbf{I.\quad GENERAL REQUIREMENTS}

\vspace{0.15cm}

The Graduation Internship Report shall be prepared using the official
\LaTeX{} template provided by the Faculty. Students shall preserve the
predefined page layout, typography, line spacing, heading styles,
numbering system, figure and table formatting, and bibliography settings
unless otherwise instructed by the Faculty or the supervisor.

The report shall be written in English.

Students are responsible for the academic and technical content of the
report. The official \LaTeX{} template is responsible for standardized
document formatting.

\vspace{0.3cm}

% ================================================================
% II. REPORT STRUCTURE
% ================================================================

\noindent
\textbf{II.\quad REPORT STRUCTURE}

\vspace{0.15cm}

The report shall normally be arranged in the following order:

\begin{enumerate}[label=\arabic*., leftmargin=1.2cm, itemsep=2pt]
    \item Outer Cover;
    \item Comments of the Internship Host Organization;
    \item Supervisor's Evaluation and Graduation Internship Defense Results;
    \item Table of Contents;
    \item List of Symbols and Abbreviations, if applicable;
    \item List of Figures;
    \item List of Tables;
    \item Abstract;
    \item Main Chapters of the Internship Report;
    \item References;
    \item Appendices, if applicable.
\end{enumerate}

The preliminary pages shall not be numbered as chapters. The main report
content shall be organized into consecutively numbered chapters beginning
with Chapter 1.

\vspace{0.3cm}

% ================================================================
% III. PAGE SIZE AND MARGINS
% ================================================================

\noindent
\textbf{III.\quad PAGE SIZE AND MARGINS}

\vspace{0.15cm}

The report shall be formatted on A4 paper with the following margins:

\begin{center}
\begin{tabular}{ll}
    Left margin:   & 3.5 cm \\
    Right margin:  & 2.0 cm \\
    Top margin:    & 2.5 cm \\
    Bottom margin: & 2.5 cm
\end{tabular}
\end{center}

These settings are predefined in the official template and shall not be
modified by students.

\vspace{0.3cm}

% ================================================================
% IV. FONT, FONT SIZE, AND LINE SPACING
% ================================================================

\noindent
\textbf{IV.\quad FONT, FONT SIZE, AND LINE SPACING}

\vspace{0.15cm}

The main text shall use a 13 pt font size with 1.3-line spacing.
A Times-family typeface shall be used in accordance with the official
template.

The text shall be fully justified.

Students shall not use repeated spaces, repeated manual line breaks, or
multiple blank lines to position text. Paragraph and heading spacing shall
be controlled by the template.

\vspace{0.3cm}

% ================================================================
% V. OUTER COVER
% ================================================================

\noindent
\textbf{V.\quad OUTER COVER}

\vspace{0.15cm}

The outer cover shall follow the official format provided in the template
and shall contain the following information:

\begin{itemize}[leftmargin=0.8cm, itemsep=2pt]
    \item Thai Nguyen University of Technology;
    \item Faculty of International Training;
    \item Engineering and Technology Department;
    \item University logo and Faculty logo;
    \item the title \textbf{INTERNSHIP REPORT};
    \item Internship Organization;
    \item Author;
    \item Class;
    \item Supervisor;
    \item place and year of completion.
\end{itemize}

Students shall only replace the required information fields and shall not
modify the cover-page layout unless instructed otherwise.

\vspace{0.3cm}

% ================================================================
% VI. COMMENTS OF THE INTERNSHIP HOST ORGANIZATION
% ================================================================

\noindent
\textbf{VI.\quad COMMENTS OF THE INTERNSHIP HOST ORGANIZATION}

\vspace{0.15cm}

The page entitled \textbf{COMMENTS OF THE INTERNSHIP HOST ORGANIZATION}
is reserved for the internship host organization.

It includes:

\begin{itemize}[leftmargin=0.8cm, itemsep=2pt]
    \item the student's full name;
    \item class;
    \item comments from the host organization;
    \item score;
    \item date;
    \item signature and official seal.
\end{itemize}

Students shall not alter or complete official evaluation fields unless
specifically authorized.

\vspace{0.3cm}

% ================================================================
% VII. SUPERVISOR'S EVALUATION AND DEFENSE RESULTS
% ================================================================

\noindent
\textbf{VII.\quad SUPERVISOR'S EVALUATION AND DEFENSE RESULTS}

\vspace{0.15cm}

The section entitled
\textbf{SUPERVISOR'S EVALUATION AND GRADUATION INTERNSHIP DEFENSE RESULTS}
is reserved for the supervisor and defense committee.

The evaluation may include:

\begin{itemize}[leftmargin=0.8cm, itemsep=2pt]
    \item Internship Progress and Student's Attitude;
    \item Report Content;
    \item Report Presentation;
    \item Other Comments;
    \item Supervisor's Opinion;
    \item Defense Result;
    \item Score.
\end{itemize}

Students shall not enter evaluation scores or defense results before they
are officially determined.

\vspace{0.3cm}

% ================================================================
% VIII. CHAPTERS, SECTIONS, AND SUBSECTIONS
% ================================================================

\noindent
\textbf{VIII.\quad CHAPTERS, SECTIONS, AND SUBSECTIONS}

\vspace{0.15cm}

The report structure shall be generated using the standard \LaTeX{}
commands:

\begin{verbatim}
\chapter{...}
\section{...}
\subsection{...}
\end{verbatim}

Chapter, section, and subsection numbers shall be generated automatically
by \LaTeX{}.

Students shall not manually type chapter or section numbers and shall not
reproduce heading styles using commands such as \verb|\textbf|,
\verb|\fontsize|, or \verb|\centering| instead of the appropriate
structural command.

Excessive subdivision should be avoided. In most cases, the hierarchy
Chapter--Section--Subsection is sufficient.

\vspace{0.3cm}

% ================================================================
% IX. ORGANIZATION OF THE MAIN CONTENT
% ================================================================

\noindent
\textbf{IX.\quad ORGANIZATION OF THE MAIN CONTENT}

\vspace{0.15cm}

The specific content of the report shall depend on the internship
organization and the assigned tasks. A typical structure may include:

\begin{itemize}[leftmargin=0.8cm, itemsep=2pt]
    \item \textbf{Chapter 1. Overview of the Internship Organization};
    \item \textbf{Chapter 2. Internship Activities and Assigned Tasks};
    \item \textbf{Chapter 3. Results, Experience, and Evaluation}.
\end{itemize}

These chapter titles are recommendations only and may be adapted to the
actual internship content.

The report shall reflect work actually performed, observed, analyzed, or
learned during the internship. Unrelated theoretical material should not
be included merely to increase report length.

\vspace{0.3cm}

% ================================================================
% X. TABLE OF CONTENTS
% ================================================================

\noindent
\textbf{X.\quad TABLE OF CONTENTS}

\vspace{0.15cm}

The Table of Contents shall be generated automatically using:

\begin{verbatim}
\tableofcontents
\end{verbatim}

Students shall not create or manually edit page numbers in the Table of
Contents.

The document shall be compiled sufficiently to ensure that all page
numbers and entries are correctly updated.

\vspace{0.3cm}

% ================================================================
% XI. LIST OF SYMBOLS AND ABBREVIATIONS
% ================================================================

\noindent
\textbf{XI.\quad LIST OF SYMBOLS AND ABBREVIATIONS}

\vspace{0.15cm}

Symbols and abbreviations that are used repeatedly in the report shall be
listed in the \textbf{List of Symbols and Abbreviations}.

The list may contain the following columns:

\begin{center}
\begin{tabular}{llll}
    No. & Symbol / Abbreviation & Unit & Definition
\end{tabular}
\end{center}

Only symbols and abbreviations actually used in the final report shall be
retained. Example entries contained in the template shall be removed if
they are not relevant to the student's report.

Each abbreviation should normally be defined when it first appears in
the main text.

\vspace{0.3cm}

% ================================================================
% XII. FIGURES, DIAGRAMS, AND PHOTOGRAPHS
% ================================================================

\noindent
\textbf{XII.\quad FIGURES, DIAGRAMS, AND PHOTOGRAPHS}

\vspace{0.15cm}

Figures shall be clear, legible, and directly relevant to the report
content.

Figures shall be inserted using the standard \verb|figure| environment.

Example:

\begin{verbatim}
\begin{figure}[H]
    \centering
    \includegraphics[width=0.8\textwidth]
                    {figures/company_structure.pdf}
    \caption{Organizational structure of the company}
    \label{fig:company-structure}
\end{figure}
\end{verbatim}

Figure captions shall be placed below the figure.

Each referenced figure shall have a unique \verb|\label| and shall be
referenced in the text using, for example:

\begin{verbatim}
Figure~\ref{fig:company-structure}
\end{verbatim}

Students shall not manually type figure numbers.

Figures reproduced or adapted from external sources shall include an
appropriate citation.

Photographs taken during the internship shall be relevant and sufficiently
clear. Students shall not include confidential, restricted, or proprietary
information without permission from the internship organization.

Vector graphics or high-resolution images are recommended whenever
possible.

\vspace{0.3cm}

% ================================================================
% XIII. TABLES
% ================================================================

\noindent
\textbf{XIII.\quad TABLES}

\vspace{0.15cm}

Tables shall be created using suitable \LaTeX{} environments such as
\verb|table|, \verb|tabular|, \verb|tabularx|, or \verb|longtable|.

Table captions shall be placed above the table.

Example:

\begin{verbatim}
\begin{table}[H]
    \centering
    \caption{Internship schedule}
    \label{tab:internship-schedule}

    \begin{tabular}{lll}
        \toprule
        Week & Task & Result \\
        \midrule
        1 & Orientation & Completed \\
        2 & System study & Completed \\
        \bottomrule
    \end{tabular}
\end{table}
\end{verbatim}

Tables shall be referenced using:

\begin{verbatim}
Table~\ref{tab:internship-schedule}
\end{verbatim}

Students shall not manually type table numbers.

Tables should be created directly in \LaTeX{} rather than inserted as
screenshots whenever practical.

\vspace{0.3cm}

% ================================================================
% XIV. EQUATIONS AND TECHNICAL QUANTITIES
% ================================================================

\noindent
\textbf{XIV.\quad EQUATIONS AND TECHNICAL QUANTITIES}

\vspace{0.15cm}

Mathematical expressions shall be typeset directly using \LaTeX{}.
Images or screenshots of equations shall not be used.

Short expressions may be written in inline mathematical mode, for example:

\begin{verbatim}
$R = 0.7\,\Omega$
\end{verbatim}

Important equations shall use a numbered equation environment.

Example:

\begin{verbatim}
\begin{equation}
    P = UI
    \label{eq:power}
\end{equation}
\end{verbatim}

Equations shall be referenced using:

\begin{verbatim}
Equation~\eqref{eq:power}
\end{verbatim}

Equation numbers shall not be manually entered.

SI units shall be used wherever applicable. The \verb|siunitx| package
should be used where appropriate.

Examples:

\begin{verbatim}
\SI{220}{\volt}
\SI{50}{\hertz}
\SI{2.5}{\ampere}
\end{verbatim}

Notation and units shall be consistent throughout the report.

\vspace{0.3cm}

% ================================================================
% XV. SOURCE CODE
% ================================================================

\noindent
\textbf{XV.\quad SOURCE CODE}

\vspace{0.15cm}

If the internship involves programming, PLCs, microcontrollers,
simulation, or software development, relevant source code may be included
in the report.

Source code shall be presented using the \verb|listings| environment or
another environment defined by the official template.

Example:

\begin{verbatim}
\begin{lstlisting}[language=st]
IF Start = TRUE THEN
    Motor := TRUE;
ELSE
    Motor := FALSE;
END_IF;
\end{lstlisting}
\end{verbatim}

Only code necessary for explaining the work should appear in the main
report. Long source-code listings should be moved to an appendix or
submitted as supplementary files if required.

Source code should not normally be inserted as screenshots.

\vspace{0.3cm}

% ================================================================
% XVI. ABSTRACT
% ================================================================

\noindent
\textbf{XVI.\quad ABSTRACT}

\vspace{0.15cm}

The Abstract shall briefly summarize:

\begin{itemize}[leftmargin=0.8cm, itemsep=2pt]
    \item the internship organization;
    \item the objectives of the internship;
    \item the main tasks performed;
    \item the main results, knowledge, and skills obtained.
\end{itemize}

The Abstract should be concise and should not reproduce the full content
of the report chapter by chapter.

\vspace{0.3cm}

% ================================================================
% XVII. REFERENCES AND CITATIONS
% ================================================================

\noindent
\textbf{XVII.\quad REFERENCES AND CITATIONS}

\vspace{0.15cm}

All books, journal papers, conference papers, standards, websites,
datasheets, manuals, technical documents, figures, tables, or other
external sources used in preparing the report shall be properly cited.

References shall be managed through the bibliography system provided with
the official template.

Bibliographic entries shall normally be stored in:

\begin{verbatim}
References.bib
\end{verbatim}

References shall be cited in the report using:

\begin{verbatim}
\cite{reference-key}
\end{verbatim}

Students shall not manually type reference numbers such as ``[1]'',
``[2]'', etc.

Titles of foreign-language publications shall remain in their original
language in the bibliography.

All borrowed ideas, data, figures, tables, equations, methods, or results
shall be properly acknowledged.

\vspace{0.3cm}

% ================================================================
% XVIII. CROSS-REFERENCING
% ================================================================

\noindent
\textbf{XVIII.\quad CROSS-REFERENCING}

\vspace{0.15cm}

Numbered elements shall be cross-referenced using the standard \LaTeX{}
mechanism.

Use the following commands where appropriate:

\begin{verbatim}
\label{...}
\ref{...}
\eqref{...}
\end{verbatim}

For example:

\begin{verbatim}
As shown in Figure~\ref{fig:system}, ...
\end{verbatim}

Students shall not manually type numbers such as ``Figure 2.3'' when
those numbers can be generated automatically.

\vspace{0.3cm}

% ================================================================
% XIX. PAGE NUMBERING
% ================================================================

\noindent
\textbf{XIX.\quad PAGE NUMBERING}

\vspace{0.15cm}

Cover pages and official evaluation forms at the beginning of the report
shall not display page numbers.

The main body of the report shall use Arabic page numbers beginning from
Chapter 1.

Page numbering shall be generated automatically by the official template.
Students shall not manually enter page numbers.

\vspace{0.3cm}

% ================================================================
% XX. PROJECT FILE ORGANIZATION
% ================================================================

\noindent
\textbf{XX.\quad PROJECT FILE ORGANIZATION}

\vspace{0.15cm}

Students should maintain a clear and organized \LaTeX{} project
structure. A recommended structure is:

\begin{verbatim}
InternshipReport/
|
|-- main.tex
|-- References.bib
|
|-- figures/
|   |-- company.pdf
|   |-- system.png
|   `-- ...
|
`-- chapters/
    |-- chapter1.tex
    |-- chapter2.tex
    `-- chapter3.tex
\end{verbatim}

Figures should normally be stored in the \verb|figures| directory
specified by the official template.

File and directory names should preferably contain Latin letters,
numbers, hyphens, or underscores.

Recommended examples include:

\begin{verbatim}
company_structure.pdf
motor_control.png
chapter_01.tex
\end{verbatim}

Spaces and unnecessary special characters in file names should be
avoided where possible.

\vspace{0.3cm}

% ================================================================
% XXI. DOCUMENT COMPILATION AND FINAL CHECK
% ================================================================

\noindent
\textbf{XXI.\quad DOCUMENT COMPILATION AND FINAL CHECK}

\vspace{0.15cm}

Before submission, the complete \LaTeX{} project shall compile
successfully and produce the final PDF report.

The final document shall not contain unresolved references such as:

\begin{verbatim}
??
[?]
\end{verbatim}

Before generating the final PDF, students shall check:

\begin{itemize}[leftmargin=0.8cm, itemsep=2pt]
    \item the Table of Contents;
    \item page numbering;
    \item chapter, section, and subsection numbering;
    \item figure numbering;
    \item table numbering;
    \item equation numbering;
    \item cross-references;
    \item citations;
    \item bibliography;
    \item font display;
    \item mathematical symbols;
    \item image quality;
    \item completeness of all required forms.
\end{itemize}

The document may be compiled using TeX Live, MiKTeX, or another
compatible \LaTeX{} environment. Regardless of the compilation
environment, the resulting PDF shall conform to the official Faculty
template.

\vspace{0.3cm}

% ================================================================
% XXII. CONTENT THAT STUDENTS MAY EDIT
% ================================================================

\noindent
\textbf{XXII.\quad CONTENT THAT STUDENTS MAY EDIT}

\vspace{0.15cm}

Students may edit or replace:

\begin{itemize}[leftmargin=0.8cm, itemsep=2pt]
    \item student name;
    \item class;
    \item internship organization;
    \item supervisor name;
    \item Abstract;
    \item List of Symbols and Abbreviations;
    \item chapter titles and chapter content;
    \item figures and photographs;
    \item tables;
    \item equations;
    \item citations and references.
\end{itemize}

Students shall not modify the main document formatting settings in the
preamble, including the document class, page geometry, line spacing,
heading formatting, table-of-contents formatting, font configuration,
page layout, or cover-page macros unless specifically instructed.

\vspace{0.3cm}

% ================================================================
% XXIII. FINAL SUBMISSION REQUIREMENTS
% ================================================================

\noindent
\textbf{XXIII.\quad FINAL SUBMISSION REQUIREMENTS}

\vspace{0.15cm}

Before submission, the final report shall:

\begin{itemize}[leftmargin=0.8cm, itemsep=2pt]
    \item compile without errors that affect the report content or presentation;
    \item contain all required forms and report sections;
    \item contain correct student, class, internship organization, and supervisor information;
    \item contain an updated Table of Contents, List of Figures, and List of Tables;
    \item contain no unused example content from the template;
    \item contain no placeholder text such as \textit{Title of Chapter 1},
    \textit{Section 1.1}, or other sample entries;
    \item contain all required citations and references;
    \item contain no unresolved cross-references.
\end{itemize}

If the Faculty requires submission of the \LaTeX{} source, the complete
project shall include all \verb|.tex| files, the bibliography database,
figures, and any other files required to reproduce the submitted PDF.

\vspace{0.4cm}

\noindent
\textbf{Important:}
Before submitting the final internship report, students shall remove or
comment out the following command in the main document:

\begin{verbatim}
\input{Regulations.tex}
\end{verbatim}

For example:

\begin{verbatim}
% \input{Regulations.tex}
\end{verbatim}

This removes the guidance pages from the final submitted report while
preserving the regulations file in the \LaTeX{} project for reference.

} % End 13pt content

\clearpage
 command before submitting
% the final internship report.
%
% =================================================================

\clearpage
\thispagestyle{empty}

\begin{center}
    {\fontsize{16pt}{20pt}\selectfont\bfseries
    SPECIFIC REGULATIONS FOR PREPARING\\
    GRADUATION INTERNSHIP REPORTS USING \LaTeX}
\end{center}

\vspace{0.5cm}

{\fontsize{13pt}{17pt}\selectfont
\setstretch{1.3}
\justifying

\noindent
\textbf{Note:}
This section is provided for guidance only and is included in the main
document using the command \verb|
% =================================================================
% REGULATIONS FOR PREPARING THE GRADUATION INTERNSHIP REPORT
% USING LATEX
% =================================================================
%
% This file is intended to be included from the main document with:
%
%     \input{Regulations.tex}
%
% IMPORTANT:
% This section is provided for guidance only. Students shall remove
% or comment out the \input{Regulations.tex} command before submitting
% the final internship report.
%
% =================================================================

\clearpage
\thispagestyle{empty}

\begin{center}
    {\fontsize{16pt}{20pt}\selectfont\bfseries
    SPECIFIC REGULATIONS FOR PREPARING\\
    GRADUATION INTERNSHIP REPORTS USING \LaTeX}
\end{center}

\vspace{0.5cm}

{\fontsize{13pt}{17pt}\selectfont
\setstretch{1.3}
\justifying

\noindent
\textbf{Note:}
This section is provided for guidance only and is included in the main
document using the command \verb|\input{Regulations.tex}|.
Students shall remove or comment out this command before submitting the
final internship report.

\vspace{0.4cm}

% ================================================================
% I. GENERAL REQUIREMENTS
% ================================================================

\noindent
\textbf{I.\quad GENERAL REQUIREMENTS}

\vspace{0.15cm}

The Graduation Internship Report shall be prepared using the official
\LaTeX{} template provided by the Faculty. Students shall preserve the
predefined page layout, typography, line spacing, heading styles,
numbering system, figure and table formatting, and bibliography settings
unless otherwise instructed by the Faculty or the supervisor.

The report shall be written in English.

Students are responsible for the academic and technical content of the
report. The official \LaTeX{} template is responsible for standardized
document formatting.

\vspace{0.3cm}

% ================================================================
% II. REPORT STRUCTURE
% ================================================================

\noindent
\textbf{II.\quad REPORT STRUCTURE}

\vspace{0.15cm}

The report shall normally be arranged in the following order:

\begin{enumerate}[label=\arabic*., leftmargin=1.2cm, itemsep=2pt]
    \item Outer Cover;
    \item Comments of the Internship Host Organization;
    \item Supervisor's Evaluation and Graduation Internship Defense Results;
    \item Table of Contents;
    \item List of Symbols and Abbreviations, if applicable;
    \item List of Figures;
    \item List of Tables;
    \item Abstract;
    \item Main Chapters of the Internship Report;
    \item References;
    \item Appendices, if applicable.
\end{enumerate}

The preliminary pages shall not be numbered as chapters. The main report
content shall be organized into consecutively numbered chapters beginning
with Chapter 1.

\vspace{0.3cm}

% ================================================================
% III. PAGE SIZE AND MARGINS
% ================================================================

\noindent
\textbf{III.\quad PAGE SIZE AND MARGINS}

\vspace{0.15cm}

The report shall be formatted on A4 paper with the following margins:

\begin{center}
\begin{tabular}{ll}
    Left margin:   & 3.5 cm \\
    Right margin:  & 2.0 cm \\
    Top margin:    & 2.5 cm \\
    Bottom margin: & 2.5 cm
\end{tabular}
\end{center}

These settings are predefined in the official template and shall not be
modified by students.

\vspace{0.3cm}

% ================================================================
% IV. FONT, FONT SIZE, AND LINE SPACING
% ================================================================

\noindent
\textbf{IV.\quad FONT, FONT SIZE, AND LINE SPACING}

\vspace{0.15cm}

The main text shall use a 13 pt font size with 1.3-line spacing.
A Times-family typeface shall be used in accordance with the official
template.

The text shall be fully justified.

Students shall not use repeated spaces, repeated manual line breaks, or
multiple blank lines to position text. Paragraph and heading spacing shall
be controlled by the template.

\vspace{0.3cm}

% ================================================================
% V. OUTER COVER
% ================================================================

\noindent
\textbf{V.\quad OUTER COVER}

\vspace{0.15cm}

The outer cover shall follow the official format provided in the template
and shall contain the following information:

\begin{itemize}[leftmargin=0.8cm, itemsep=2pt]
    \item Thai Nguyen University of Technology;
    \item Faculty of International Training;
    \item Engineering and Technology Department;
    \item University logo and Faculty logo;
    \item the title \textbf{INTERNSHIP REPORT};
    \item Internship Organization;
    \item Author;
    \item Class;
    \item Supervisor;
    \item place and year of completion.
\end{itemize}

Students shall only replace the required information fields and shall not
modify the cover-page layout unless instructed otherwise.

\vspace{0.3cm}

% ================================================================
% VI. COMMENTS OF THE INTERNSHIP HOST ORGANIZATION
% ================================================================

\noindent
\textbf{VI.\quad COMMENTS OF THE INTERNSHIP HOST ORGANIZATION}

\vspace{0.15cm}

The page entitled \textbf{COMMENTS OF THE INTERNSHIP HOST ORGANIZATION}
is reserved for the internship host organization.

It includes:

\begin{itemize}[leftmargin=0.8cm, itemsep=2pt]
    \item the student's full name;
    \item class;
    \item comments from the host organization;
    \item score;
    \item date;
    \item signature and official seal.
\end{itemize}

Students shall not alter or complete official evaluation fields unless
specifically authorized.

\vspace{0.3cm}

% ================================================================
% VII. SUPERVISOR'S EVALUATION AND DEFENSE RESULTS
% ================================================================

\noindent
\textbf{VII.\quad SUPERVISOR'S EVALUATION AND DEFENSE RESULTS}

\vspace{0.15cm}

The section entitled
\textbf{SUPERVISOR'S EVALUATION AND GRADUATION INTERNSHIP DEFENSE RESULTS}
is reserved for the supervisor and defense committee.

The evaluation may include:

\begin{itemize}[leftmargin=0.8cm, itemsep=2pt]
    \item Internship Progress and Student's Attitude;
    \item Report Content;
    \item Report Presentation;
    \item Other Comments;
    \item Supervisor's Opinion;
    \item Defense Result;
    \item Score.
\end{itemize}

Students shall not enter evaluation scores or defense results before they
are officially determined.

\vspace{0.3cm}

% ================================================================
% VIII. CHAPTERS, SECTIONS, AND SUBSECTIONS
% ================================================================

\noindent
\textbf{VIII.\quad CHAPTERS, SECTIONS, AND SUBSECTIONS}

\vspace{0.15cm}

The report structure shall be generated using the standard \LaTeX{}
commands:

\begin{verbatim}
\chapter{...}
\section{...}
\subsection{...}
\end{verbatim}

Chapter, section, and subsection numbers shall be generated automatically
by \LaTeX{}.

Students shall not manually type chapter or section numbers and shall not
reproduce heading styles using commands such as \verb|\textbf|,
\verb|\fontsize|, or \verb|\centering| instead of the appropriate
structural command.

Excessive subdivision should be avoided. In most cases, the hierarchy
Chapter--Section--Subsection is sufficient.

\vspace{0.3cm}

% ================================================================
% IX. ORGANIZATION OF THE MAIN CONTENT
% ================================================================

\noindent
\textbf{IX.\quad ORGANIZATION OF THE MAIN CONTENT}

\vspace{0.15cm}

The specific content of the report shall depend on the internship
organization and the assigned tasks. A typical structure may include:

\begin{itemize}[leftmargin=0.8cm, itemsep=2pt]
    \item \textbf{Chapter 1. Overview of the Internship Organization};
    \item \textbf{Chapter 2. Internship Activities and Assigned Tasks};
    \item \textbf{Chapter 3. Results, Experience, and Evaluation}.
\end{itemize}

These chapter titles are recommendations only and may be adapted to the
actual internship content.

The report shall reflect work actually performed, observed, analyzed, or
learned during the internship. Unrelated theoretical material should not
be included merely to increase report length.

\vspace{0.3cm}

% ================================================================
% X. TABLE OF CONTENTS
% ================================================================

\noindent
\textbf{X.\quad TABLE OF CONTENTS}

\vspace{0.15cm}

The Table of Contents shall be generated automatically using:

\begin{verbatim}
\tableofcontents
\end{verbatim}

Students shall not create or manually edit page numbers in the Table of
Contents.

The document shall be compiled sufficiently to ensure that all page
numbers and entries are correctly updated.

\vspace{0.3cm}

% ================================================================
% XI. LIST OF SYMBOLS AND ABBREVIATIONS
% ================================================================

\noindent
\textbf{XI.\quad LIST OF SYMBOLS AND ABBREVIATIONS}

\vspace{0.15cm}

Symbols and abbreviations that are used repeatedly in the report shall be
listed in the \textbf{List of Symbols and Abbreviations}.

The list may contain the following columns:

\begin{center}
\begin{tabular}{llll}
    No. & Symbol / Abbreviation & Unit & Definition
\end{tabular}
\end{center}

Only symbols and abbreviations actually used in the final report shall be
retained. Example entries contained in the template shall be removed if
they are not relevant to the student's report.

Each abbreviation should normally be defined when it first appears in
the main text.

\vspace{0.3cm}

% ================================================================
% XII. FIGURES, DIAGRAMS, AND PHOTOGRAPHS
% ================================================================

\noindent
\textbf{XII.\quad FIGURES, DIAGRAMS, AND PHOTOGRAPHS}

\vspace{0.15cm}

Figures shall be clear, legible, and directly relevant to the report
content.

Figures shall be inserted using the standard \verb|figure| environment.

Example:

\begin{verbatim}
\begin{figure}[H]
    \centering
    \includegraphics[width=0.8\textwidth]
                    {figures/company_structure.pdf}
    \caption{Organizational structure of the company}
    \label{fig:company-structure}
\end{figure}
\end{verbatim}

Figure captions shall be placed below the figure.

Each referenced figure shall have a unique \verb|\label| and shall be
referenced in the text using, for example:

\begin{verbatim}
Figure~\ref{fig:company-structure}
\end{verbatim}

Students shall not manually type figure numbers.

Figures reproduced or adapted from external sources shall include an
appropriate citation.

Photographs taken during the internship shall be relevant and sufficiently
clear. Students shall not include confidential, restricted, or proprietary
information without permission from the internship organization.

Vector graphics or high-resolution images are recommended whenever
possible.

\vspace{0.3cm}

% ================================================================
% XIII. TABLES
% ================================================================

\noindent
\textbf{XIII.\quad TABLES}

\vspace{0.15cm}

Tables shall be created using suitable \LaTeX{} environments such as
\verb|table|, \verb|tabular|, \verb|tabularx|, or \verb|longtable|.

Table captions shall be placed above the table.

Example:

\begin{verbatim}
\begin{table}[H]
    \centering
    \caption{Internship schedule}
    \label{tab:internship-schedule}

    \begin{tabular}{lll}
        \toprule
        Week & Task & Result \\
        \midrule
        1 & Orientation & Completed \\
        2 & System study & Completed \\
        \bottomrule
    \end{tabular}
\end{table}
\end{verbatim}

Tables shall be referenced using:

\begin{verbatim}
Table~\ref{tab:internship-schedule}
\end{verbatim}

Students shall not manually type table numbers.

Tables should be created directly in \LaTeX{} rather than inserted as
screenshots whenever practical.

\vspace{0.3cm}

% ================================================================
% XIV. EQUATIONS AND TECHNICAL QUANTITIES
% ================================================================

\noindent
\textbf{XIV.\quad EQUATIONS AND TECHNICAL QUANTITIES}

\vspace{0.15cm}

Mathematical expressions shall be typeset directly using \LaTeX{}.
Images or screenshots of equations shall not be used.

Short expressions may be written in inline mathematical mode, for example:

\begin{verbatim}
$R = 0.7\,\Omega$
\end{verbatim}

Important equations shall use a numbered equation environment.

Example:

\begin{verbatim}
\begin{equation}
    P = UI
    \label{eq:power}
\end{equation}
\end{verbatim}

Equations shall be referenced using:

\begin{verbatim}
Equation~\eqref{eq:power}
\end{verbatim}

Equation numbers shall not be manually entered.

SI units shall be used wherever applicable. The \verb|siunitx| package
should be used where appropriate.

Examples:

\begin{verbatim}
\SI{220}{\volt}
\SI{50}{\hertz}
\SI{2.5}{\ampere}
\end{verbatim}

Notation and units shall be consistent throughout the report.

\vspace{0.3cm}

% ================================================================
% XV. SOURCE CODE
% ================================================================

\noindent
\textbf{XV.\quad SOURCE CODE}

\vspace{0.15cm}

If the internship involves programming, PLCs, microcontrollers,
simulation, or software development, relevant source code may be included
in the report.

Source code shall be presented using the \verb|listings| environment or
another environment defined by the official template.

Example:

\begin{verbatim}
\begin{lstlisting}[language=st]
IF Start = TRUE THEN
    Motor := TRUE;
ELSE
    Motor := FALSE;
END_IF;
\end{lstlisting}
\end{verbatim}

Only code necessary for explaining the work should appear in the main
report. Long source-code listings should be moved to an appendix or
submitted as supplementary files if required.

Source code should not normally be inserted as screenshots.

\vspace{0.3cm}

% ================================================================
% XVI. ABSTRACT
% ================================================================

\noindent
\textbf{XVI.\quad ABSTRACT}

\vspace{0.15cm}

The Abstract shall briefly summarize:

\begin{itemize}[leftmargin=0.8cm, itemsep=2pt]
    \item the internship organization;
    \item the objectives of the internship;
    \item the main tasks performed;
    \item the main results, knowledge, and skills obtained.
\end{itemize}

The Abstract should be concise and should not reproduce the full content
of the report chapter by chapter.

\vspace{0.3cm}

% ================================================================
% XVII. REFERENCES AND CITATIONS
% ================================================================

\noindent
\textbf{XVII.\quad REFERENCES AND CITATIONS}

\vspace{0.15cm}

All books, journal papers, conference papers, standards, websites,
datasheets, manuals, technical documents, figures, tables, or other
external sources used in preparing the report shall be properly cited.

References shall be managed through the bibliography system provided with
the official template.

Bibliographic entries shall normally be stored in:

\begin{verbatim}
References.bib
\end{verbatim}

References shall be cited in the report using:

\begin{verbatim}
\cite{reference-key}
\end{verbatim}

Students shall not manually type reference numbers such as ``[1]'',
``[2]'', etc.

Titles of foreign-language publications shall remain in their original
language in the bibliography.

All borrowed ideas, data, figures, tables, equations, methods, or results
shall be properly acknowledged.

\vspace{0.3cm}

% ================================================================
% XVIII. CROSS-REFERENCING
% ================================================================

\noindent
\textbf{XVIII.\quad CROSS-REFERENCING}

\vspace{0.15cm}

Numbered elements shall be cross-referenced using the standard \LaTeX{}
mechanism.

Use the following commands where appropriate:

\begin{verbatim}
\label{...}
\ref{...}
\eqref{...}
\end{verbatim}

For example:

\begin{verbatim}
As shown in Figure~\ref{fig:system}, ...
\end{verbatim}

Students shall not manually type numbers such as ``Figure 2.3'' when
those numbers can be generated automatically.

\vspace{0.3cm}

% ================================================================
% XIX. PAGE NUMBERING
% ================================================================

\noindent
\textbf{XIX.\quad PAGE NUMBERING}

\vspace{0.15cm}

Cover pages and official evaluation forms at the beginning of the report
shall not display page numbers.

The main body of the report shall use Arabic page numbers beginning from
Chapter 1.

Page numbering shall be generated automatically by the official template.
Students shall not manually enter page numbers.

\vspace{0.3cm}

% ================================================================
% XX. PROJECT FILE ORGANIZATION
% ================================================================

\noindent
\textbf{XX.\quad PROJECT FILE ORGANIZATION}

\vspace{0.15cm}

Students should maintain a clear and organized \LaTeX{} project
structure. A recommended structure is:

\begin{verbatim}
InternshipReport/
|
|-- main.tex
|-- References.bib
|
|-- figures/
|   |-- company.pdf
|   |-- system.png
|   `-- ...
|
`-- chapters/
    |-- chapter1.tex
    |-- chapter2.tex
    `-- chapter3.tex
\end{verbatim}

Figures should normally be stored in the \verb|figures| directory
specified by the official template.

File and directory names should preferably contain Latin letters,
numbers, hyphens, or underscores.

Recommended examples include:

\begin{verbatim}
company_structure.pdf
motor_control.png
chapter_01.tex
\end{verbatim}

Spaces and unnecessary special characters in file names should be
avoided where possible.

\vspace{0.3cm}

% ================================================================
% XXI. DOCUMENT COMPILATION AND FINAL CHECK
% ================================================================

\noindent
\textbf{XXI.\quad DOCUMENT COMPILATION AND FINAL CHECK}

\vspace{0.15cm}

Before submission, the complete \LaTeX{} project shall compile
successfully and produce the final PDF report.

The final document shall not contain unresolved references such as:

\begin{verbatim}
??
[?]
\end{verbatim}

Before generating the final PDF, students shall check:

\begin{itemize}[leftmargin=0.8cm, itemsep=2pt]
    \item the Table of Contents;
    \item page numbering;
    \item chapter, section, and subsection numbering;
    \item figure numbering;
    \item table numbering;
    \item equation numbering;
    \item cross-references;
    \item citations;
    \item bibliography;
    \item font display;
    \item mathematical symbols;
    \item image quality;
    \item completeness of all required forms.
\end{itemize}

The document may be compiled using TeX Live, MiKTeX, or another
compatible \LaTeX{} environment. Regardless of the compilation
environment, the resulting PDF shall conform to the official Faculty
template.

\vspace{0.3cm}

% ================================================================
% XXII. CONTENT THAT STUDENTS MAY EDIT
% ================================================================

\noindent
\textbf{XXII.\quad CONTENT THAT STUDENTS MAY EDIT}

\vspace{0.15cm}

Students may edit or replace:

\begin{itemize}[leftmargin=0.8cm, itemsep=2pt]
    \item student name;
    \item class;
    \item internship organization;
    \item supervisor name;
    \item Abstract;
    \item List of Symbols and Abbreviations;
    \item chapter titles and chapter content;
    \item figures and photographs;
    \item tables;
    \item equations;
    \item citations and references.
\end{itemize}

Students shall not modify the main document formatting settings in the
preamble, including the document class, page geometry, line spacing,
heading formatting, table-of-contents formatting, font configuration,
page layout, or cover-page macros unless specifically instructed.

\vspace{0.3cm}

% ================================================================
% XXIII. FINAL SUBMISSION REQUIREMENTS
% ================================================================

\noindent
\textbf{XXIII.\quad FINAL SUBMISSION REQUIREMENTS}

\vspace{0.15cm}

Before submission, the final report shall:

\begin{itemize}[leftmargin=0.8cm, itemsep=2pt]
    \item compile without errors that affect the report content or presentation;
    \item contain all required forms and report sections;
    \item contain correct student, class, internship organization, and supervisor information;
    \item contain an updated Table of Contents, List of Figures, and List of Tables;
    \item contain no unused example content from the template;
    \item contain no placeholder text such as \textit{Title of Chapter 1},
    \textit{Section 1.1}, or other sample entries;
    \item contain all required citations and references;
    \item contain no unresolved cross-references.
\end{itemize}

If the Faculty requires submission of the \LaTeX{} source, the complete
project shall include all \verb|.tex| files, the bibliography database,
figures, and any other files required to reproduce the submitted PDF.

\vspace{0.4cm}

\noindent
\textbf{Important:}
Before submitting the final internship report, students shall remove or
comment out the following command in the main document:

\begin{verbatim}
\input{Regulations.tex}
\end{verbatim}

For example:

\begin{verbatim}
% \input{Regulations.tex}
\end{verbatim}

This removes the guidance pages from the final submitted report while
preserving the regulations file in the \LaTeX{} project for reference.

} % End 13pt content

\clearpage
|.
Students shall remove or comment out this command before submitting the
final internship report.

\vspace{0.4cm}

% ================================================================
% I. GENERAL REQUIREMENTS
% ================================================================

\noindent
\textbf{I.\quad GENERAL REQUIREMENTS}

\vspace{0.15cm}

The Graduation Internship Report shall be prepared using the official
\LaTeX{} template provided by the Faculty. Students shall preserve the
predefined page layout, typography, line spacing, heading styles,
numbering system, figure and table formatting, and bibliography settings
unless otherwise instructed by the Faculty or the supervisor.

The report shall be written in English.

Students are responsible for the academic and technical content of the
report. The official \LaTeX{} template is responsible for standardized
document formatting.

\vspace{0.3cm}

% ================================================================
% II. REPORT STRUCTURE
% ================================================================

\noindent
\textbf{II.\quad REPORT STRUCTURE}

\vspace{0.15cm}

The report shall normally be arranged in the following order:

\begin{enumerate}[label=\arabic*., leftmargin=1.2cm, itemsep=2pt]
    \item Outer Cover;
    \item Comments of the Internship Host Organization;
    \item Supervisor's Evaluation and Graduation Internship Defense Results;
    \item Table of Contents;
    \item List of Symbols and Abbreviations, if applicable;
    \item List of Figures;
    \item List of Tables;
    \item Abstract;
    \item Main Chapters of the Internship Report;
    \item References;
    \item Appendices, if applicable.
\end{enumerate}

The preliminary pages shall not be numbered as chapters. The main report
content shall be organized into consecutively numbered chapters beginning
with Chapter 1.

\vspace{0.3cm}

% ================================================================
% III. PAGE SIZE AND MARGINS
% ================================================================

\noindent
\textbf{III.\quad PAGE SIZE AND MARGINS}

\vspace{0.15cm}

The report shall be formatted on A4 paper with the following margins:

\begin{center}
\begin{tabular}{ll}
    Left margin:   & 3.5 cm \\
    Right margin:  & 2.0 cm \\
    Top margin:    & 2.5 cm \\
    Bottom margin: & 2.5 cm
\end{tabular}
\end{center}

These settings are predefined in the official template and shall not be
modified by students.

\vspace{0.3cm}

% ================================================================
% IV. FONT, FONT SIZE, AND LINE SPACING
% ================================================================

\noindent
\textbf{IV.\quad FONT, FONT SIZE, AND LINE SPACING}

\vspace{0.15cm}

The main text shall use a 13 pt font size with 1.3-line spacing.
A Times-family typeface shall be used in accordance with the official
template.

The text shall be fully justified.

Students shall not use repeated spaces, repeated manual line breaks, or
multiple blank lines to position text. Paragraph and heading spacing shall
be controlled by the template.

\vspace{0.3cm}

% ================================================================
% V. OUTER COVER
% ================================================================

\noindent
\textbf{V.\quad OUTER COVER}

\vspace{0.15cm}

The outer cover shall follow the official format provided in the template
and shall contain the following information:

\begin{itemize}[leftmargin=0.8cm, itemsep=2pt]
    \item Thai Nguyen University of Technology;
    \item Faculty of International Training;
    \item Engineering and Technology Department;
    \item University logo and Faculty logo;
    \item the title \textbf{INTERNSHIP REPORT};
    \item Internship Organization;
    \item Author;
    \item Class;
    \item Supervisor;
    \item place and year of completion.
\end{itemize}

Students shall only replace the required information fields and shall not
modify the cover-page layout unless instructed otherwise.

\vspace{0.3cm}

% ================================================================
% VI. COMMENTS OF THE INTERNSHIP HOST ORGANIZATION
% ================================================================

\noindent
\textbf{VI.\quad COMMENTS OF THE INTERNSHIP HOST ORGANIZATION}

\vspace{0.15cm}

The page entitled \textbf{COMMENTS OF THE INTERNSHIP HOST ORGANIZATION}
is reserved for the internship host organization.

It includes:

\begin{itemize}[leftmargin=0.8cm, itemsep=2pt]
    \item the student's full name;
    \item class;
    \item comments from the host organization;
    \item score;
    \item date;
    \item signature and official seal.
\end{itemize}

Students shall not alter or complete official evaluation fields unless
specifically authorized.

\vspace{0.3cm}

% ================================================================
% VII. SUPERVISOR'S EVALUATION AND DEFENSE RESULTS
% ================================================================

\noindent
\textbf{VII.\quad SUPERVISOR'S EVALUATION AND DEFENSE RESULTS}

\vspace{0.15cm}

The section entitled
\textbf{SUPERVISOR'S EVALUATION AND GRADUATION INTERNSHIP DEFENSE RESULTS}
is reserved for the supervisor and defense committee.

The evaluation may include:

\begin{itemize}[leftmargin=0.8cm, itemsep=2pt]
    \item Internship Progress and Student's Attitude;
    \item Report Content;
    \item Report Presentation;
    \item Other Comments;
    \item Supervisor's Opinion;
    \item Defense Result;
    \item Score.
\end{itemize}

Students shall not enter evaluation scores or defense results before they
are officially determined.

\vspace{0.3cm}

% ================================================================
% VIII. CHAPTERS, SECTIONS, AND SUBSECTIONS
% ================================================================

\noindent
\textbf{VIII.\quad CHAPTERS, SECTIONS, AND SUBSECTIONS}

\vspace{0.15cm}

The report structure shall be generated using the standard \LaTeX{}
commands:

\begin{verbatim}
\chapter{...}
\section{...}
\subsection{...}
\end{verbatim}

Chapter, section, and subsection numbers shall be generated automatically
by \LaTeX{}.

Students shall not manually type chapter or section numbers and shall not
reproduce heading styles using commands such as \verb|\textbf|,
\verb|\fontsize|, or \verb|\centering| instead of the appropriate
structural command.

Excessive subdivision should be avoided. In most cases, the hierarchy
Chapter--Section--Subsection is sufficient.

\vspace{0.3cm}

% ================================================================
% IX. ORGANIZATION OF THE MAIN CONTENT
% ================================================================

\noindent
\textbf{IX.\quad ORGANIZATION OF THE MAIN CONTENT}

\vspace{0.15cm}

The specific content of the report shall depend on the internship
organization and the assigned tasks. A typical structure may include:

\begin{itemize}[leftmargin=0.8cm, itemsep=2pt]
    \item \textbf{Chapter 1. Overview of the Internship Organization};
    \item \textbf{Chapter 2. Internship Activities and Assigned Tasks};
    \item \textbf{Chapter 3. Results, Experience, and Evaluation}.
\end{itemize}

These chapter titles are recommendations only and may be adapted to the
actual internship content.

The report shall reflect work actually performed, observed, analyzed, or
learned during the internship. Unrelated theoretical material should not
be included merely to increase report length.

\vspace{0.3cm}

% ================================================================
% X. TABLE OF CONTENTS
% ================================================================

\noindent
\textbf{X.\quad TABLE OF CONTENTS}

\vspace{0.15cm}

The Table of Contents shall be generated automatically using:

\begin{verbatim}
\tableofcontents
\end{verbatim}

Students shall not create or manually edit page numbers in the Table of
Contents.

The document shall be compiled sufficiently to ensure that all page
numbers and entries are correctly updated.

\vspace{0.3cm}

% ================================================================
% XI. LIST OF SYMBOLS AND ABBREVIATIONS
% ================================================================

\noindent
\textbf{XI.\quad LIST OF SYMBOLS AND ABBREVIATIONS}

\vspace{0.15cm}

Symbols and abbreviations that are used repeatedly in the report shall be
listed in the \textbf{List of Symbols and Abbreviations}.

The list may contain the following columns:

\begin{center}
\begin{tabular}{llll}
    No. & Symbol / Abbreviation & Unit & Definition
\end{tabular}
\end{center}

Only symbols and abbreviations actually used in the final report shall be
retained. Example entries contained in the template shall be removed if
they are not relevant to the student's report.

Each abbreviation should normally be defined when it first appears in
the main text.

\vspace{0.3cm}

% ================================================================
% XII. FIGURES, DIAGRAMS, AND PHOTOGRAPHS
% ================================================================

\noindent
\textbf{XII.\quad FIGURES, DIAGRAMS, AND PHOTOGRAPHS}

\vspace{0.15cm}

Figures shall be clear, legible, and directly relevant to the report
content.

Figures shall be inserted using the standard \verb|figure| environment.

Example:

\begin{verbatim}
\begin{figure}[H]
    \centering
    \includegraphics[width=0.8\textwidth]
                    {figures/company_structure.pdf}
    \caption{Organizational structure of the company}
    \label{fig:company-structure}
\end{figure}
\end{verbatim}

Figure captions shall be placed below the figure.

Each referenced figure shall have a unique \verb|\label| and shall be
referenced in the text using, for example:

\begin{verbatim}
Figure~\ref{fig:company-structure}
\end{verbatim}

Students shall not manually type figure numbers.

Figures reproduced or adapted from external sources shall include an
appropriate citation.

Photographs taken during the internship shall be relevant and sufficiently
clear. Students shall not include confidential, restricted, or proprietary
information without permission from the internship organization.

Vector graphics or high-resolution images are recommended whenever
possible.

\vspace{0.3cm}

% ================================================================
% XIII. TABLES
% ================================================================

\noindent
\textbf{XIII.\quad TABLES}

\vspace{0.15cm}

Tables shall be created using suitable \LaTeX{} environments such as
\verb|table|, \verb|tabular|, \verb|tabularx|, or \verb|longtable|.

Table captions shall be placed above the table.

Example:

\begin{verbatim}
\begin{table}[H]
    \centering
    \caption{Internship schedule}
    \label{tab:internship-schedule}

    \begin{tabular}{lll}
        \toprule
        Week & Task & Result \\
        \midrule
        1 & Orientation & Completed \\
        2 & System study & Completed \\
        \bottomrule
    \end{tabular}
\end{table}
\end{verbatim}

Tables shall be referenced using:

\begin{verbatim}
Table~\ref{tab:internship-schedule}
\end{verbatim}

Students shall not manually type table numbers.

Tables should be created directly in \LaTeX{} rather than inserted as
screenshots whenever practical.

\vspace{0.3cm}

% ================================================================
% XIV. EQUATIONS AND TECHNICAL QUANTITIES
% ================================================================

\noindent
\textbf{XIV.\quad EQUATIONS AND TECHNICAL QUANTITIES}

\vspace{0.15cm}

Mathematical expressions shall be typeset directly using \LaTeX{}.
Images or screenshots of equations shall not be used.

Short expressions may be written in inline mathematical mode, for example:

\begin{verbatim}
$R = 0.7\,\Omega$
\end{verbatim}

Important equations shall use a numbered equation environment.

Example:

\begin{verbatim}
\begin{equation}
    P = UI
    \label{eq:power}
\end{equation}
\end{verbatim}

Equations shall be referenced using:

\begin{verbatim}
Equation~\eqref{eq:power}
\end{verbatim}

Equation numbers shall not be manually entered.

SI units shall be used wherever applicable. The \verb|siunitx| package
should be used where appropriate.

Examples:

\begin{verbatim}
\SI{220}{\volt}
\SI{50}{\hertz}
\SI{2.5}{\ampere}
\end{verbatim}

Notation and units shall be consistent throughout the report.

\vspace{0.3cm}

% ================================================================
% XV. SOURCE CODE
% ================================================================

\noindent
\textbf{XV.\quad SOURCE CODE}

\vspace{0.15cm}

If the internship involves programming, PLCs, microcontrollers,
simulation, or software development, relevant source code may be included
in the report.

Source code shall be presented using the \verb|listings| environment or
another environment defined by the official template.

Example:

\begin{verbatim}
\begin{lstlisting}[language=st]
IF Start = TRUE THEN
    Motor := TRUE;
ELSE
    Motor := FALSE;
END_IF;
\end{lstlisting}
\end{verbatim}

Only code necessary for explaining the work should appear in the main
report. Long source-code listings should be moved to an appendix or
submitted as supplementary files if required.

Source code should not normally be inserted as screenshots.

\vspace{0.3cm}

% ================================================================
% XVI. ABSTRACT
% ================================================================

\noindent
\textbf{XVI.\quad ABSTRACT}

\vspace{0.15cm}

The Abstract shall briefly summarize:

\begin{itemize}[leftmargin=0.8cm, itemsep=2pt]
    \item the internship organization;
    \item the objectives of the internship;
    \item the main tasks performed;
    \item the main results, knowledge, and skills obtained.
\end{itemize}

The Abstract should be concise and should not reproduce the full content
of the report chapter by chapter.

\vspace{0.3cm}

% ================================================================
% XVII. REFERENCES AND CITATIONS
% ================================================================

\noindent
\textbf{XVII.\quad REFERENCES AND CITATIONS}

\vspace{0.15cm}

All books, journal papers, conference papers, standards, websites,
datasheets, manuals, technical documents, figures, tables, or other
external sources used in preparing the report shall be properly cited.

References shall be managed through the bibliography system provided with
the official template.

Bibliographic entries shall normally be stored in:

\begin{verbatim}
References.bib
\end{verbatim}

References shall be cited in the report using:

\begin{verbatim}
\cite{reference-key}
\end{verbatim}

Students shall not manually type reference numbers such as ``[1]'',
``[2]'', etc.

Titles of foreign-language publications shall remain in their original
language in the bibliography.

All borrowed ideas, data, figures, tables, equations, methods, or results
shall be properly acknowledged.

\vspace{0.3cm}

% ================================================================
% XVIII. CROSS-REFERENCING
% ================================================================

\noindent
\textbf{XVIII.\quad CROSS-REFERENCING}

\vspace{0.15cm}

Numbered elements shall be cross-referenced using the standard \LaTeX{}
mechanism.

Use the following commands where appropriate:

\begin{verbatim}
\label{...}
\ref{...}
\eqref{...}
\end{verbatim}

For example:

\begin{verbatim}
As shown in Figure~\ref{fig:system}, ...
\end{verbatim}

Students shall not manually type numbers such as ``Figure 2.3'' when
those numbers can be generated automatically.

\vspace{0.3cm}

% ================================================================
% XIX. PAGE NUMBERING
% ================================================================

\noindent
\textbf{XIX.\quad PAGE NUMBERING}

\vspace{0.15cm}

Cover pages and official evaluation forms at the beginning of the report
shall not display page numbers.

The main body of the report shall use Arabic page numbers beginning from
Chapter 1.

Page numbering shall be generated automatically by the official template.
Students shall not manually enter page numbers.

\vspace{0.3cm}

% ================================================================
% XX. PROJECT FILE ORGANIZATION
% ================================================================

\noindent
\textbf{XX.\quad PROJECT FILE ORGANIZATION}

\vspace{0.15cm}

Students should maintain a clear and organized \LaTeX{} project
structure. A recommended structure is:

\begin{verbatim}
InternshipReport/
|
|-- main.tex
|-- References.bib
|
|-- figures/
|   |-- company.pdf
|   |-- system.png
|   `-- ...
|
`-- chapters/
    |-- chapter1.tex
    |-- chapter2.tex
    `-- chapter3.tex
\end{verbatim}

Figures should normally be stored in the \verb|figures| directory
specified by the official template.

File and directory names should preferably contain Latin letters,
numbers, hyphens, or underscores.

Recommended examples include:

\begin{verbatim}
company_structure.pdf
motor_control.png
chapter_01.tex
\end{verbatim}

Spaces and unnecessary special characters in file names should be
avoided where possible.

\vspace{0.3cm}

% ================================================================
% XXI. DOCUMENT COMPILATION AND FINAL CHECK
% ================================================================

\noindent
\textbf{XXI.\quad DOCUMENT COMPILATION AND FINAL CHECK}

\vspace{0.15cm}

Before submission, the complete \LaTeX{} project shall compile
successfully and produce the final PDF report.

The final document shall not contain unresolved references such as:

\begin{verbatim}
??
[?]
\end{verbatim}

Before generating the final PDF, students shall check:

\begin{itemize}[leftmargin=0.8cm, itemsep=2pt]
    \item the Table of Contents;
    \item page numbering;
    \item chapter, section, and subsection numbering;
    \item figure numbering;
    \item table numbering;
    \item equation numbering;
    \item cross-references;
    \item citations;
    \item bibliography;
    \item font display;
    \item mathematical symbols;
    \item image quality;
    \item completeness of all required forms.
\end{itemize}

The document may be compiled using TeX Live, MiKTeX, or another
compatible \LaTeX{} environment. Regardless of the compilation
environment, the resulting PDF shall conform to the official Faculty
template.

\vspace{0.3cm}

% ================================================================
% XXII. CONTENT THAT STUDENTS MAY EDIT
% ================================================================

\noindent
\textbf{XXII.\quad CONTENT THAT STUDENTS MAY EDIT}

\vspace{0.15cm}

Students may edit or replace:

\begin{itemize}[leftmargin=0.8cm, itemsep=2pt]
    \item student name;
    \item class;
    \item internship organization;
    \item supervisor name;
    \item Abstract;
    \item List of Symbols and Abbreviations;
    \item chapter titles and chapter content;
    \item figures and photographs;
    \item tables;
    \item equations;
    \item citations and references.
\end{itemize}

Students shall not modify the main document formatting settings in the
preamble, including the document class, page geometry, line spacing,
heading formatting, table-of-contents formatting, font configuration,
page layout, or cover-page macros unless specifically instructed.

\vspace{0.3cm}

% ================================================================
% XXIII. FINAL SUBMISSION REQUIREMENTS
% ================================================================

\noindent
\textbf{XXIII.\quad FINAL SUBMISSION REQUIREMENTS}

\vspace{0.15cm}

Before submission, the final report shall:

\begin{itemize}[leftmargin=0.8cm, itemsep=2pt]
    \item compile without errors that affect the report content or presentation;
    \item contain all required forms and report sections;
    \item contain correct student, class, internship organization, and supervisor information;
    \item contain an updated Table of Contents, List of Figures, and List of Tables;
    \item contain no unused example content from the template;
    \item contain no placeholder text such as \textit{Title of Chapter 1},
    \textit{Section 1.1}, or other sample entries;
    \item contain all required citations and references;
    \item contain no unresolved cross-references.
\end{itemize}

If the Faculty requires submission of the \LaTeX{} source, the complete
project shall include all \verb|.tex| files, the bibliography database,
figures, and any other files required to reproduce the submitted PDF.

\vspace{0.4cm}

\noindent
\textbf{Important:}
Before submitting the final internship report, students shall remove or
comment out the following command in the main document:

\begin{verbatim}

% =================================================================
% REGULATIONS FOR PREPARING THE GRADUATION INTERNSHIP REPORT
% USING LATEX
% =================================================================
%
% This file is intended to be included from the main document with:
%
%     \input{Regulations.tex}
%
% IMPORTANT:
% This section is provided for guidance only. Students shall remove
% or comment out the \input{Regulations.tex} command before submitting
% the final internship report.
%
% =================================================================

\clearpage
\thispagestyle{empty}

\begin{center}
    {\fontsize{16pt}{20pt}\selectfont\bfseries
    SPECIFIC REGULATIONS FOR PREPARING\\
    GRADUATION INTERNSHIP REPORTS USING \LaTeX}
\end{center}

\vspace{0.5cm}

{\fontsize{13pt}{17pt}\selectfont
\setstretch{1.3}
\justifying

\noindent
\textbf{Note:}
This section is provided for guidance only and is included in the main
document using the command \verb|\input{Regulations.tex}|.
Students shall remove or comment out this command before submitting the
final internship report.

\vspace{0.4cm}

% ================================================================
% I. GENERAL REQUIREMENTS
% ================================================================

\noindent
\textbf{I.\quad GENERAL REQUIREMENTS}

\vspace{0.15cm}

The Graduation Internship Report shall be prepared using the official
\LaTeX{} template provided by the Faculty. Students shall preserve the
predefined page layout, typography, line spacing, heading styles,
numbering system, figure and table formatting, and bibliography settings
unless otherwise instructed by the Faculty or the supervisor.

The report shall be written in English.

Students are responsible for the academic and technical content of the
report. The official \LaTeX{} template is responsible for standardized
document formatting.

\vspace{0.3cm}

% ================================================================
% II. REPORT STRUCTURE
% ================================================================

\noindent
\textbf{II.\quad REPORT STRUCTURE}

\vspace{0.15cm}

The report shall normally be arranged in the following order:

\begin{enumerate}[label=\arabic*., leftmargin=1.2cm, itemsep=2pt]
    \item Outer Cover;
    \item Comments of the Internship Host Organization;
    \item Supervisor's Evaluation and Graduation Internship Defense Results;
    \item Table of Contents;
    \item List of Symbols and Abbreviations, if applicable;
    \item List of Figures;
    \item List of Tables;
    \item Abstract;
    \item Main Chapters of the Internship Report;
    \item References;
    \item Appendices, if applicable.
\end{enumerate}

The preliminary pages shall not be numbered as chapters. The main report
content shall be organized into consecutively numbered chapters beginning
with Chapter 1.

\vspace{0.3cm}

% ================================================================
% III. PAGE SIZE AND MARGINS
% ================================================================

\noindent
\textbf{III.\quad PAGE SIZE AND MARGINS}

\vspace{0.15cm}

The report shall be formatted on A4 paper with the following margins:

\begin{center}
\begin{tabular}{ll}
    Left margin:   & 3.5 cm \\
    Right margin:  & 2.0 cm \\
    Top margin:    & 2.5 cm \\
    Bottom margin: & 2.5 cm
\end{tabular}
\end{center}

These settings are predefined in the official template and shall not be
modified by students.

\vspace{0.3cm}

% ================================================================
% IV. FONT, FONT SIZE, AND LINE SPACING
% ================================================================

\noindent
\textbf{IV.\quad FONT, FONT SIZE, AND LINE SPACING}

\vspace{0.15cm}

The main text shall use a 13 pt font size with 1.3-line spacing.
A Times-family typeface shall be used in accordance with the official
template.

The text shall be fully justified.

Students shall not use repeated spaces, repeated manual line breaks, or
multiple blank lines to position text. Paragraph and heading spacing shall
be controlled by the template.

\vspace{0.3cm}

% ================================================================
% V. OUTER COVER
% ================================================================

\noindent
\textbf{V.\quad OUTER COVER}

\vspace{0.15cm}

The outer cover shall follow the official format provided in the template
and shall contain the following information:

\begin{itemize}[leftmargin=0.8cm, itemsep=2pt]
    \item Thai Nguyen University of Technology;
    \item Faculty of International Training;
    \item Engineering and Technology Department;
    \item University logo and Faculty logo;
    \item the title \textbf{INTERNSHIP REPORT};
    \item Internship Organization;
    \item Author;
    \item Class;
    \item Supervisor;
    \item place and year of completion.
\end{itemize}

Students shall only replace the required information fields and shall not
modify the cover-page layout unless instructed otherwise.

\vspace{0.3cm}

% ================================================================
% VI. COMMENTS OF THE INTERNSHIP HOST ORGANIZATION
% ================================================================

\noindent
\textbf{VI.\quad COMMENTS OF THE INTERNSHIP HOST ORGANIZATION}

\vspace{0.15cm}

The page entitled \textbf{COMMENTS OF THE INTERNSHIP HOST ORGANIZATION}
is reserved for the internship host organization.

It includes:

\begin{itemize}[leftmargin=0.8cm, itemsep=2pt]
    \item the student's full name;
    \item class;
    \item comments from the host organization;
    \item score;
    \item date;
    \item signature and official seal.
\end{itemize}

Students shall not alter or complete official evaluation fields unless
specifically authorized.

\vspace{0.3cm}

% ================================================================
% VII. SUPERVISOR'S EVALUATION AND DEFENSE RESULTS
% ================================================================

\noindent
\textbf{VII.\quad SUPERVISOR'S EVALUATION AND DEFENSE RESULTS}

\vspace{0.15cm}

The section entitled
\textbf{SUPERVISOR'S EVALUATION AND GRADUATION INTERNSHIP DEFENSE RESULTS}
is reserved for the supervisor and defense committee.

The evaluation may include:

\begin{itemize}[leftmargin=0.8cm, itemsep=2pt]
    \item Internship Progress and Student's Attitude;
    \item Report Content;
    \item Report Presentation;
    \item Other Comments;
    \item Supervisor's Opinion;
    \item Defense Result;
    \item Score.
\end{itemize}

Students shall not enter evaluation scores or defense results before they
are officially determined.

\vspace{0.3cm}

% ================================================================
% VIII. CHAPTERS, SECTIONS, AND SUBSECTIONS
% ================================================================

\noindent
\textbf{VIII.\quad CHAPTERS, SECTIONS, AND SUBSECTIONS}

\vspace{0.15cm}

The report structure shall be generated using the standard \LaTeX{}
commands:

\begin{verbatim}
\chapter{...}
\section{...}
\subsection{...}
\end{verbatim}

Chapter, section, and subsection numbers shall be generated automatically
by \LaTeX{}.

Students shall not manually type chapter or section numbers and shall not
reproduce heading styles using commands such as \verb|\textbf|,
\verb|\fontsize|, or \verb|\centering| instead of the appropriate
structural command.

Excessive subdivision should be avoided. In most cases, the hierarchy
Chapter--Section--Subsection is sufficient.

\vspace{0.3cm}

% ================================================================
% IX. ORGANIZATION OF THE MAIN CONTENT
% ================================================================

\noindent
\textbf{IX.\quad ORGANIZATION OF THE MAIN CONTENT}

\vspace{0.15cm}

The specific content of the report shall depend on the internship
organization and the assigned tasks. A typical structure may include:

\begin{itemize}[leftmargin=0.8cm, itemsep=2pt]
    \item \textbf{Chapter 1. Overview of the Internship Organization};
    \item \textbf{Chapter 2. Internship Activities and Assigned Tasks};
    \item \textbf{Chapter 3. Results, Experience, and Evaluation}.
\end{itemize}

These chapter titles are recommendations only and may be adapted to the
actual internship content.

The report shall reflect work actually performed, observed, analyzed, or
learned during the internship. Unrelated theoretical material should not
be included merely to increase report length.

\vspace{0.3cm}

% ================================================================
% X. TABLE OF CONTENTS
% ================================================================

\noindent
\textbf{X.\quad TABLE OF CONTENTS}

\vspace{0.15cm}

The Table of Contents shall be generated automatically using:

\begin{verbatim}
\tableofcontents
\end{verbatim}

Students shall not create or manually edit page numbers in the Table of
Contents.

The document shall be compiled sufficiently to ensure that all page
numbers and entries are correctly updated.

\vspace{0.3cm}

% ================================================================
% XI. LIST OF SYMBOLS AND ABBREVIATIONS
% ================================================================

\noindent
\textbf{XI.\quad LIST OF SYMBOLS AND ABBREVIATIONS}

\vspace{0.15cm}

Symbols and abbreviations that are used repeatedly in the report shall be
listed in the \textbf{List of Symbols and Abbreviations}.

The list may contain the following columns:

\begin{center}
\begin{tabular}{llll}
    No. & Symbol / Abbreviation & Unit & Definition
\end{tabular}
\end{center}

Only symbols and abbreviations actually used in the final report shall be
retained. Example entries contained in the template shall be removed if
they are not relevant to the student's report.

Each abbreviation should normally be defined when it first appears in
the main text.

\vspace{0.3cm}

% ================================================================
% XII. FIGURES, DIAGRAMS, AND PHOTOGRAPHS
% ================================================================

\noindent
\textbf{XII.\quad FIGURES, DIAGRAMS, AND PHOTOGRAPHS}

\vspace{0.15cm}

Figures shall be clear, legible, and directly relevant to the report
content.

Figures shall be inserted using the standard \verb|figure| environment.

Example:

\begin{verbatim}
\begin{figure}[H]
    \centering
    \includegraphics[width=0.8\textwidth]
                    {figures/company_structure.pdf}
    \caption{Organizational structure of the company}
    \label{fig:company-structure}
\end{figure}
\end{verbatim}

Figure captions shall be placed below the figure.

Each referenced figure shall have a unique \verb|\label| and shall be
referenced in the text using, for example:

\begin{verbatim}
Figure~\ref{fig:company-structure}
\end{verbatim}

Students shall not manually type figure numbers.

Figures reproduced or adapted from external sources shall include an
appropriate citation.

Photographs taken during the internship shall be relevant and sufficiently
clear. Students shall not include confidential, restricted, or proprietary
information without permission from the internship organization.

Vector graphics or high-resolution images are recommended whenever
possible.

\vspace{0.3cm}

% ================================================================
% XIII. TABLES
% ================================================================

\noindent
\textbf{XIII.\quad TABLES}

\vspace{0.15cm}

Tables shall be created using suitable \LaTeX{} environments such as
\verb|table|, \verb|tabular|, \verb|tabularx|, or \verb|longtable|.

Table captions shall be placed above the table.

Example:

\begin{verbatim}
\begin{table}[H]
    \centering
    \caption{Internship schedule}
    \label{tab:internship-schedule}

    \begin{tabular}{lll}
        \toprule
        Week & Task & Result \\
        \midrule
        1 & Orientation & Completed \\
        2 & System study & Completed \\
        \bottomrule
    \end{tabular}
\end{table}
\end{verbatim}

Tables shall be referenced using:

\begin{verbatim}
Table~\ref{tab:internship-schedule}
\end{verbatim}

Students shall not manually type table numbers.

Tables should be created directly in \LaTeX{} rather than inserted as
screenshots whenever practical.

\vspace{0.3cm}

% ================================================================
% XIV. EQUATIONS AND TECHNICAL QUANTITIES
% ================================================================

\noindent
\textbf{XIV.\quad EQUATIONS AND TECHNICAL QUANTITIES}

\vspace{0.15cm}

Mathematical expressions shall be typeset directly using \LaTeX{}.
Images or screenshots of equations shall not be used.

Short expressions may be written in inline mathematical mode, for example:

\begin{verbatim}
$R = 0.7\,\Omega$
\end{verbatim}

Important equations shall use a numbered equation environment.

Example:

\begin{verbatim}
\begin{equation}
    P = UI
    \label{eq:power}
\end{equation}
\end{verbatim}

Equations shall be referenced using:

\begin{verbatim}
Equation~\eqref{eq:power}
\end{verbatim}

Equation numbers shall not be manually entered.

SI units shall be used wherever applicable. The \verb|siunitx| package
should be used where appropriate.

Examples:

\begin{verbatim}
\SI{220}{\volt}
\SI{50}{\hertz}
\SI{2.5}{\ampere}
\end{verbatim}

Notation and units shall be consistent throughout the report.

\vspace{0.3cm}

% ================================================================
% XV. SOURCE CODE
% ================================================================

\noindent
\textbf{XV.\quad SOURCE CODE}

\vspace{0.15cm}

If the internship involves programming, PLCs, microcontrollers,
simulation, or software development, relevant source code may be included
in the report.

Source code shall be presented using the \verb|listings| environment or
another environment defined by the official template.

Example:

\begin{verbatim}
\begin{lstlisting}[language=st]
IF Start = TRUE THEN
    Motor := TRUE;
ELSE
    Motor := FALSE;
END_IF;
\end{lstlisting}
\end{verbatim}

Only code necessary for explaining the work should appear in the main
report. Long source-code listings should be moved to an appendix or
submitted as supplementary files if required.

Source code should not normally be inserted as screenshots.

\vspace{0.3cm}

% ================================================================
% XVI. ABSTRACT
% ================================================================

\noindent
\textbf{XVI.\quad ABSTRACT}

\vspace{0.15cm}

The Abstract shall briefly summarize:

\begin{itemize}[leftmargin=0.8cm, itemsep=2pt]
    \item the internship organization;
    \item the objectives of the internship;
    \item the main tasks performed;
    \item the main results, knowledge, and skills obtained.
\end{itemize}

The Abstract should be concise and should not reproduce the full content
of the report chapter by chapter.

\vspace{0.3cm}

% ================================================================
% XVII. REFERENCES AND CITATIONS
% ================================================================

\noindent
\textbf{XVII.\quad REFERENCES AND CITATIONS}

\vspace{0.15cm}

All books, journal papers, conference papers, standards, websites,
datasheets, manuals, technical documents, figures, tables, or other
external sources used in preparing the report shall be properly cited.

References shall be managed through the bibliography system provided with
the official template.

Bibliographic entries shall normally be stored in:

\begin{verbatim}
References.bib
\end{verbatim}

References shall be cited in the report using:

\begin{verbatim}
\cite{reference-key}
\end{verbatim}

Students shall not manually type reference numbers such as ``[1]'',
``[2]'', etc.

Titles of foreign-language publications shall remain in their original
language in the bibliography.

All borrowed ideas, data, figures, tables, equations, methods, or results
shall be properly acknowledged.

\vspace{0.3cm}

% ================================================================
% XVIII. CROSS-REFERENCING
% ================================================================

\noindent
\textbf{XVIII.\quad CROSS-REFERENCING}

\vspace{0.15cm}

Numbered elements shall be cross-referenced using the standard \LaTeX{}
mechanism.

Use the following commands where appropriate:

\begin{verbatim}
\label{...}
\ref{...}
\eqref{...}
\end{verbatim}

For example:

\begin{verbatim}
As shown in Figure~\ref{fig:system}, ...
\end{verbatim}

Students shall not manually type numbers such as ``Figure 2.3'' when
those numbers can be generated automatically.

\vspace{0.3cm}

% ================================================================
% XIX. PAGE NUMBERING
% ================================================================

\noindent
\textbf{XIX.\quad PAGE NUMBERING}

\vspace{0.15cm}

Cover pages and official evaluation forms at the beginning of the report
shall not display page numbers.

The main body of the report shall use Arabic page numbers beginning from
Chapter 1.

Page numbering shall be generated automatically by the official template.
Students shall not manually enter page numbers.

\vspace{0.3cm}

% ================================================================
% XX. PROJECT FILE ORGANIZATION
% ================================================================

\noindent
\textbf{XX.\quad PROJECT FILE ORGANIZATION}

\vspace{0.15cm}

Students should maintain a clear and organized \LaTeX{} project
structure. A recommended structure is:

\begin{verbatim}
InternshipReport/
|
|-- main.tex
|-- References.bib
|
|-- figures/
|   |-- company.pdf
|   |-- system.png
|   `-- ...
|
`-- chapters/
    |-- chapter1.tex
    |-- chapter2.tex
    `-- chapter3.tex
\end{verbatim}

Figures should normally be stored in the \verb|figures| directory
specified by the official template.

File and directory names should preferably contain Latin letters,
numbers, hyphens, or underscores.

Recommended examples include:

\begin{verbatim}
company_structure.pdf
motor_control.png
chapter_01.tex
\end{verbatim}

Spaces and unnecessary special characters in file names should be
avoided where possible.

\vspace{0.3cm}

% ================================================================
% XXI. DOCUMENT COMPILATION AND FINAL CHECK
% ================================================================

\noindent
\textbf{XXI.\quad DOCUMENT COMPILATION AND FINAL CHECK}

\vspace{0.15cm}

Before submission, the complete \LaTeX{} project shall compile
successfully and produce the final PDF report.

The final document shall not contain unresolved references such as:

\begin{verbatim}
??
[?]
\end{verbatim}

Before generating the final PDF, students shall check:

\begin{itemize}[leftmargin=0.8cm, itemsep=2pt]
    \item the Table of Contents;
    \item page numbering;
    \item chapter, section, and subsection numbering;
    \item figure numbering;
    \item table numbering;
    \item equation numbering;
    \item cross-references;
    \item citations;
    \item bibliography;
    \item font display;
    \item mathematical symbols;
    \item image quality;
    \item completeness of all required forms.
\end{itemize}

The document may be compiled using TeX Live, MiKTeX, or another
compatible \LaTeX{} environment. Regardless of the compilation
environment, the resulting PDF shall conform to the official Faculty
template.

\vspace{0.3cm}

% ================================================================
% XXII. CONTENT THAT STUDENTS MAY EDIT
% ================================================================

\noindent
\textbf{XXII.\quad CONTENT THAT STUDENTS MAY EDIT}

\vspace{0.15cm}

Students may edit or replace:

\begin{itemize}[leftmargin=0.8cm, itemsep=2pt]
    \item student name;
    \item class;
    \item internship organization;
    \item supervisor name;
    \item Abstract;
    \item List of Symbols and Abbreviations;
    \item chapter titles and chapter content;
    \item figures and photographs;
    \item tables;
    \item equations;
    \item citations and references.
\end{itemize}

Students shall not modify the main document formatting settings in the
preamble, including the document class, page geometry, line spacing,
heading formatting, table-of-contents formatting, font configuration,
page layout, or cover-page macros unless specifically instructed.

\vspace{0.3cm}

% ================================================================
% XXIII. FINAL SUBMISSION REQUIREMENTS
% ================================================================

\noindent
\textbf{XXIII.\quad FINAL SUBMISSION REQUIREMENTS}

\vspace{0.15cm}

Before submission, the final report shall:

\begin{itemize}[leftmargin=0.8cm, itemsep=2pt]
    \item compile without errors that affect the report content or presentation;
    \item contain all required forms and report sections;
    \item contain correct student, class, internship organization, and supervisor information;
    \item contain an updated Table of Contents, List of Figures, and List of Tables;
    \item contain no unused example content from the template;
    \item contain no placeholder text such as \textit{Title of Chapter 1},
    \textit{Section 1.1}, or other sample entries;
    \item contain all required citations and references;
    \item contain no unresolved cross-references.
\end{itemize}

If the Faculty requires submission of the \LaTeX{} source, the complete
project shall include all \verb|.tex| files, the bibliography database,
figures, and any other files required to reproduce the submitted PDF.

\vspace{0.4cm}

\noindent
\textbf{Important:}
Before submitting the final internship report, students shall remove or
comment out the following command in the main document:

\begin{verbatim}
\input{Regulations.tex}
\end{verbatim}

For example:

\begin{verbatim}
% \input{Regulations.tex}
\end{verbatim}

This removes the guidance pages from the final submitted report while
preserving the regulations file in the \LaTeX{} project for reference.

} % End 13pt content

\clearpage

\end{verbatim}

For example:

\begin{verbatim}
% 
% =================================================================
% REGULATIONS FOR PREPARING THE GRADUATION INTERNSHIP REPORT
% USING LATEX
% =================================================================
%
% This file is intended to be included from the main document with:
%
%     \input{Regulations.tex}
%
% IMPORTANT:
% This section is provided for guidance only. Students shall remove
% or comment out the \input{Regulations.tex} command before submitting
% the final internship report.
%
% =================================================================

\clearpage
\thispagestyle{empty}

\begin{center}
    {\fontsize{16pt}{20pt}\selectfont\bfseries
    SPECIFIC REGULATIONS FOR PREPARING\\
    GRADUATION INTERNSHIP REPORTS USING \LaTeX}
\end{center}

\vspace{0.5cm}

{\fontsize{13pt}{17pt}\selectfont
\setstretch{1.3}
\justifying

\noindent
\textbf{Note:}
This section is provided for guidance only and is included in the main
document using the command \verb|\input{Regulations.tex}|.
Students shall remove or comment out this command before submitting the
final internship report.

\vspace{0.4cm}

% ================================================================
% I. GENERAL REQUIREMENTS
% ================================================================

\noindent
\textbf{I.\quad GENERAL REQUIREMENTS}

\vspace{0.15cm}

The Graduation Internship Report shall be prepared using the official
\LaTeX{} template provided by the Faculty. Students shall preserve the
predefined page layout, typography, line spacing, heading styles,
numbering system, figure and table formatting, and bibliography settings
unless otherwise instructed by the Faculty or the supervisor.

The report shall be written in English.

Students are responsible for the academic and technical content of the
report. The official \LaTeX{} template is responsible for standardized
document formatting.

\vspace{0.3cm}

% ================================================================
% II. REPORT STRUCTURE
% ================================================================

\noindent
\textbf{II.\quad REPORT STRUCTURE}

\vspace{0.15cm}

The report shall normally be arranged in the following order:

\begin{enumerate}[label=\arabic*., leftmargin=1.2cm, itemsep=2pt]
    \item Outer Cover;
    \item Comments of the Internship Host Organization;
    \item Supervisor's Evaluation and Graduation Internship Defense Results;
    \item Table of Contents;
    \item List of Symbols and Abbreviations, if applicable;
    \item List of Figures;
    \item List of Tables;
    \item Abstract;
    \item Main Chapters of the Internship Report;
    \item References;
    \item Appendices, if applicable.
\end{enumerate}

The preliminary pages shall not be numbered as chapters. The main report
content shall be organized into consecutively numbered chapters beginning
with Chapter 1.

\vspace{0.3cm}

% ================================================================
% III. PAGE SIZE AND MARGINS
% ================================================================

\noindent
\textbf{III.\quad PAGE SIZE AND MARGINS}

\vspace{0.15cm}

The report shall be formatted on A4 paper with the following margins:

\begin{center}
\begin{tabular}{ll}
    Left margin:   & 3.5 cm \\
    Right margin:  & 2.0 cm \\
    Top margin:    & 2.5 cm \\
    Bottom margin: & 2.5 cm
\end{tabular}
\end{center}

These settings are predefined in the official template and shall not be
modified by students.

\vspace{0.3cm}

% ================================================================
% IV. FONT, FONT SIZE, AND LINE SPACING
% ================================================================

\noindent
\textbf{IV.\quad FONT, FONT SIZE, AND LINE SPACING}

\vspace{0.15cm}

The main text shall use a 13 pt font size with 1.3-line spacing.
A Times-family typeface shall be used in accordance with the official
template.

The text shall be fully justified.

Students shall not use repeated spaces, repeated manual line breaks, or
multiple blank lines to position text. Paragraph and heading spacing shall
be controlled by the template.

\vspace{0.3cm}

% ================================================================
% V. OUTER COVER
% ================================================================

\noindent
\textbf{V.\quad OUTER COVER}

\vspace{0.15cm}

The outer cover shall follow the official format provided in the template
and shall contain the following information:

\begin{itemize}[leftmargin=0.8cm, itemsep=2pt]
    \item Thai Nguyen University of Technology;
    \item Faculty of International Training;
    \item Engineering and Technology Department;
    \item University logo and Faculty logo;
    \item the title \textbf{INTERNSHIP REPORT};
    \item Internship Organization;
    \item Author;
    \item Class;
    \item Supervisor;
    \item place and year of completion.
\end{itemize}

Students shall only replace the required information fields and shall not
modify the cover-page layout unless instructed otherwise.

\vspace{0.3cm}

% ================================================================
% VI. COMMENTS OF THE INTERNSHIP HOST ORGANIZATION
% ================================================================

\noindent
\textbf{VI.\quad COMMENTS OF THE INTERNSHIP HOST ORGANIZATION}

\vspace{0.15cm}

The page entitled \textbf{COMMENTS OF THE INTERNSHIP HOST ORGANIZATION}
is reserved for the internship host organization.

It includes:

\begin{itemize}[leftmargin=0.8cm, itemsep=2pt]
    \item the student's full name;
    \item class;
    \item comments from the host organization;
    \item score;
    \item date;
    \item signature and official seal.
\end{itemize}

Students shall not alter or complete official evaluation fields unless
specifically authorized.

\vspace{0.3cm}

% ================================================================
% VII. SUPERVISOR'S EVALUATION AND DEFENSE RESULTS
% ================================================================

\noindent
\textbf{VII.\quad SUPERVISOR'S EVALUATION AND DEFENSE RESULTS}

\vspace{0.15cm}

The section entitled
\textbf{SUPERVISOR'S EVALUATION AND GRADUATION INTERNSHIP DEFENSE RESULTS}
is reserved for the supervisor and defense committee.

The evaluation may include:

\begin{itemize}[leftmargin=0.8cm, itemsep=2pt]
    \item Internship Progress and Student's Attitude;
    \item Report Content;
    \item Report Presentation;
    \item Other Comments;
    \item Supervisor's Opinion;
    \item Defense Result;
    \item Score.
\end{itemize}

Students shall not enter evaluation scores or defense results before they
are officially determined.

\vspace{0.3cm}

% ================================================================
% VIII. CHAPTERS, SECTIONS, AND SUBSECTIONS
% ================================================================

\noindent
\textbf{VIII.\quad CHAPTERS, SECTIONS, AND SUBSECTIONS}

\vspace{0.15cm}

The report structure shall be generated using the standard \LaTeX{}
commands:

\begin{verbatim}
\chapter{...}
\section{...}
\subsection{...}
\end{verbatim}

Chapter, section, and subsection numbers shall be generated automatically
by \LaTeX{}.

Students shall not manually type chapter or section numbers and shall not
reproduce heading styles using commands such as \verb|\textbf|,
\verb|\fontsize|, or \verb|\centering| instead of the appropriate
structural command.

Excessive subdivision should be avoided. In most cases, the hierarchy
Chapter--Section--Subsection is sufficient.

\vspace{0.3cm}

% ================================================================
% IX. ORGANIZATION OF THE MAIN CONTENT
% ================================================================

\noindent
\textbf{IX.\quad ORGANIZATION OF THE MAIN CONTENT}

\vspace{0.15cm}

The specific content of the report shall depend on the internship
organization and the assigned tasks. A typical structure may include:

\begin{itemize}[leftmargin=0.8cm, itemsep=2pt]
    \item \textbf{Chapter 1. Overview of the Internship Organization};
    \item \textbf{Chapter 2. Internship Activities and Assigned Tasks};
    \item \textbf{Chapter 3. Results, Experience, and Evaluation}.
\end{itemize}

These chapter titles are recommendations only and may be adapted to the
actual internship content.

The report shall reflect work actually performed, observed, analyzed, or
learned during the internship. Unrelated theoretical material should not
be included merely to increase report length.

\vspace{0.3cm}

% ================================================================
% X. TABLE OF CONTENTS
% ================================================================

\noindent
\textbf{X.\quad TABLE OF CONTENTS}

\vspace{0.15cm}

The Table of Contents shall be generated automatically using:

\begin{verbatim}
\tableofcontents
\end{verbatim}

Students shall not create or manually edit page numbers in the Table of
Contents.

The document shall be compiled sufficiently to ensure that all page
numbers and entries are correctly updated.

\vspace{0.3cm}

% ================================================================
% XI. LIST OF SYMBOLS AND ABBREVIATIONS
% ================================================================

\noindent
\textbf{XI.\quad LIST OF SYMBOLS AND ABBREVIATIONS}

\vspace{0.15cm}

Symbols and abbreviations that are used repeatedly in the report shall be
listed in the \textbf{List of Symbols and Abbreviations}.

The list may contain the following columns:

\begin{center}
\begin{tabular}{llll}
    No. & Symbol / Abbreviation & Unit & Definition
\end{tabular}
\end{center}

Only symbols and abbreviations actually used in the final report shall be
retained. Example entries contained in the template shall be removed if
they are not relevant to the student's report.

Each abbreviation should normally be defined when it first appears in
the main text.

\vspace{0.3cm}

% ================================================================
% XII. FIGURES, DIAGRAMS, AND PHOTOGRAPHS
% ================================================================

\noindent
\textbf{XII.\quad FIGURES, DIAGRAMS, AND PHOTOGRAPHS}

\vspace{0.15cm}

Figures shall be clear, legible, and directly relevant to the report
content.

Figures shall be inserted using the standard \verb|figure| environment.

Example:

\begin{verbatim}
\begin{figure}[H]
    \centering
    \includegraphics[width=0.8\textwidth]
                    {figures/company_structure.pdf}
    \caption{Organizational structure of the company}
    \label{fig:company-structure}
\end{figure}
\end{verbatim}

Figure captions shall be placed below the figure.

Each referenced figure shall have a unique \verb|\label| and shall be
referenced in the text using, for example:

\begin{verbatim}
Figure~\ref{fig:company-structure}
\end{verbatim}

Students shall not manually type figure numbers.

Figures reproduced or adapted from external sources shall include an
appropriate citation.

Photographs taken during the internship shall be relevant and sufficiently
clear. Students shall not include confidential, restricted, or proprietary
information without permission from the internship organization.

Vector graphics or high-resolution images are recommended whenever
possible.

\vspace{0.3cm}

% ================================================================
% XIII. TABLES
% ================================================================

\noindent
\textbf{XIII.\quad TABLES}

\vspace{0.15cm}

Tables shall be created using suitable \LaTeX{} environments such as
\verb|table|, \verb|tabular|, \verb|tabularx|, or \verb|longtable|.

Table captions shall be placed above the table.

Example:

\begin{verbatim}
\begin{table}[H]
    \centering
    \caption{Internship schedule}
    \label{tab:internship-schedule}

    \begin{tabular}{lll}
        \toprule
        Week & Task & Result \\
        \midrule
        1 & Orientation & Completed \\
        2 & System study & Completed \\
        \bottomrule
    \end{tabular}
\end{table}
\end{verbatim}

Tables shall be referenced using:

\begin{verbatim}
Table~\ref{tab:internship-schedule}
\end{verbatim}

Students shall not manually type table numbers.

Tables should be created directly in \LaTeX{} rather than inserted as
screenshots whenever practical.

\vspace{0.3cm}

% ================================================================
% XIV. EQUATIONS AND TECHNICAL QUANTITIES
% ================================================================

\noindent
\textbf{XIV.\quad EQUATIONS AND TECHNICAL QUANTITIES}

\vspace{0.15cm}

Mathematical expressions shall be typeset directly using \LaTeX{}.
Images or screenshots of equations shall not be used.

Short expressions may be written in inline mathematical mode, for example:

\begin{verbatim}
$R = 0.7\,\Omega$
\end{verbatim}

Important equations shall use a numbered equation environment.

Example:

\begin{verbatim}
\begin{equation}
    P = UI
    \label{eq:power}
\end{equation}
\end{verbatim}

Equations shall be referenced using:

\begin{verbatim}
Equation~\eqref{eq:power}
\end{verbatim}

Equation numbers shall not be manually entered.

SI units shall be used wherever applicable. The \verb|siunitx| package
should be used where appropriate.

Examples:

\begin{verbatim}
\SI{220}{\volt}
\SI{50}{\hertz}
\SI{2.5}{\ampere}
\end{verbatim}

Notation and units shall be consistent throughout the report.

\vspace{0.3cm}

% ================================================================
% XV. SOURCE CODE
% ================================================================

\noindent
\textbf{XV.\quad SOURCE CODE}

\vspace{0.15cm}

If the internship involves programming, PLCs, microcontrollers,
simulation, or software development, relevant source code may be included
in the report.

Source code shall be presented using the \verb|listings| environment or
another environment defined by the official template.

Example:

\begin{verbatim}
\begin{lstlisting}[language=st]
IF Start = TRUE THEN
    Motor := TRUE;
ELSE
    Motor := FALSE;
END_IF;
\end{lstlisting}
\end{verbatim}

Only code necessary for explaining the work should appear in the main
report. Long source-code listings should be moved to an appendix or
submitted as supplementary files if required.

Source code should not normally be inserted as screenshots.

\vspace{0.3cm}

% ================================================================
% XVI. ABSTRACT
% ================================================================

\noindent
\textbf{XVI.\quad ABSTRACT}

\vspace{0.15cm}

The Abstract shall briefly summarize:

\begin{itemize}[leftmargin=0.8cm, itemsep=2pt]
    \item the internship organization;
    \item the objectives of the internship;
    \item the main tasks performed;
    \item the main results, knowledge, and skills obtained.
\end{itemize}

The Abstract should be concise and should not reproduce the full content
of the report chapter by chapter.

\vspace{0.3cm}

% ================================================================
% XVII. REFERENCES AND CITATIONS
% ================================================================

\noindent
\textbf{XVII.\quad REFERENCES AND CITATIONS}

\vspace{0.15cm}

All books, journal papers, conference papers, standards, websites,
datasheets, manuals, technical documents, figures, tables, or other
external sources used in preparing the report shall be properly cited.

References shall be managed through the bibliography system provided with
the official template.

Bibliographic entries shall normally be stored in:

\begin{verbatim}
References.bib
\end{verbatim}

References shall be cited in the report using:

\begin{verbatim}
\cite{reference-key}
\end{verbatim}

Students shall not manually type reference numbers such as ``[1]'',
``[2]'', etc.

Titles of foreign-language publications shall remain in their original
language in the bibliography.

All borrowed ideas, data, figures, tables, equations, methods, or results
shall be properly acknowledged.

\vspace{0.3cm}

% ================================================================
% XVIII. CROSS-REFERENCING
% ================================================================

\noindent
\textbf{XVIII.\quad CROSS-REFERENCING}

\vspace{0.15cm}

Numbered elements shall be cross-referenced using the standard \LaTeX{}
mechanism.

Use the following commands where appropriate:

\begin{verbatim}
\label{...}
\ref{...}
\eqref{...}
\end{verbatim}

For example:

\begin{verbatim}
As shown in Figure~\ref{fig:system}, ...
\end{verbatim}

Students shall not manually type numbers such as ``Figure 2.3'' when
those numbers can be generated automatically.

\vspace{0.3cm}

% ================================================================
% XIX. PAGE NUMBERING
% ================================================================

\noindent
\textbf{XIX.\quad PAGE NUMBERING}

\vspace{0.15cm}

Cover pages and official evaluation forms at the beginning of the report
shall not display page numbers.

The main body of the report shall use Arabic page numbers beginning from
Chapter 1.

Page numbering shall be generated automatically by the official template.
Students shall not manually enter page numbers.

\vspace{0.3cm}

% ================================================================
% XX. PROJECT FILE ORGANIZATION
% ================================================================

\noindent
\textbf{XX.\quad PROJECT FILE ORGANIZATION}

\vspace{0.15cm}

Students should maintain a clear and organized \LaTeX{} project
structure. A recommended structure is:

\begin{verbatim}
InternshipReport/
|
|-- main.tex
|-- References.bib
|
|-- figures/
|   |-- company.pdf
|   |-- system.png
|   `-- ...
|
`-- chapters/
    |-- chapter1.tex
    |-- chapter2.tex
    `-- chapter3.tex
\end{verbatim}

Figures should normally be stored in the \verb|figures| directory
specified by the official template.

File and directory names should preferably contain Latin letters,
numbers, hyphens, or underscores.

Recommended examples include:

\begin{verbatim}
company_structure.pdf
motor_control.png
chapter_01.tex
\end{verbatim}

Spaces and unnecessary special characters in file names should be
avoided where possible.

\vspace{0.3cm}

% ================================================================
% XXI. DOCUMENT COMPILATION AND FINAL CHECK
% ================================================================

\noindent
\textbf{XXI.\quad DOCUMENT COMPILATION AND FINAL CHECK}

\vspace{0.15cm}

Before submission, the complete \LaTeX{} project shall compile
successfully and produce the final PDF report.

The final document shall not contain unresolved references such as:

\begin{verbatim}
??
[?]
\end{verbatim}

Before generating the final PDF, students shall check:

\begin{itemize}[leftmargin=0.8cm, itemsep=2pt]
    \item the Table of Contents;
    \item page numbering;
    \item chapter, section, and subsection numbering;
    \item figure numbering;
    \item table numbering;
    \item equation numbering;
    \item cross-references;
    \item citations;
    \item bibliography;
    \item font display;
    \item mathematical symbols;
    \item image quality;
    \item completeness of all required forms.
\end{itemize}

The document may be compiled using TeX Live, MiKTeX, or another
compatible \LaTeX{} environment. Regardless of the compilation
environment, the resulting PDF shall conform to the official Faculty
template.

\vspace{0.3cm}

% ================================================================
% XXII. CONTENT THAT STUDENTS MAY EDIT
% ================================================================

\noindent
\textbf{XXII.\quad CONTENT THAT STUDENTS MAY EDIT}

\vspace{0.15cm}

Students may edit or replace:

\begin{itemize}[leftmargin=0.8cm, itemsep=2pt]
    \item student name;
    \item class;
    \item internship organization;
    \item supervisor name;
    \item Abstract;
    \item List of Symbols and Abbreviations;
    \item chapter titles and chapter content;
    \item figures and photographs;
    \item tables;
    \item equations;
    \item citations and references.
\end{itemize}

Students shall not modify the main document formatting settings in the
preamble, including the document class, page geometry, line spacing,
heading formatting, table-of-contents formatting, font configuration,
page layout, or cover-page macros unless specifically instructed.

\vspace{0.3cm}

% ================================================================
% XXIII. FINAL SUBMISSION REQUIREMENTS
% ================================================================

\noindent
\textbf{XXIII.\quad FINAL SUBMISSION REQUIREMENTS}

\vspace{0.15cm}

Before submission, the final report shall:

\begin{itemize}[leftmargin=0.8cm, itemsep=2pt]
    \item compile without errors that affect the report content or presentation;
    \item contain all required forms and report sections;
    \item contain correct student, class, internship organization, and supervisor information;
    \item contain an updated Table of Contents, List of Figures, and List of Tables;
    \item contain no unused example content from the template;
    \item contain no placeholder text such as \textit{Title of Chapter 1},
    \textit{Section 1.1}, or other sample entries;
    \item contain all required citations and references;
    \item contain no unresolved cross-references.
\end{itemize}

If the Faculty requires submission of the \LaTeX{} source, the complete
project shall include all \verb|.tex| files, the bibliography database,
figures, and any other files required to reproduce the submitted PDF.

\vspace{0.4cm}

\noindent
\textbf{Important:}
Before submitting the final internship report, students shall remove or
comment out the following command in the main document:

\begin{verbatim}
\input{Regulations.tex}
\end{verbatim}

For example:

\begin{verbatim}
% \input{Regulations.tex}
\end{verbatim}

This removes the guidance pages from the final submitted report while
preserving the regulations file in the \LaTeX{} project for reference.

} % End 13pt content

\clearpage

\end{verbatim}

This removes the guidance pages from the final submitted report while
preserving the regulations file in the \LaTeX{} project for reference.

} % End 13pt content

\clearpage

%
% IMPORTANT:
% This section is provided for guidance only. Students shall remove
% or comment out the 
% =================================================================
% REGULATIONS FOR PREPARING THE GRADUATION INTERNSHIP REPORT
% USING LATEX
% =================================================================
%
% This file is intended to be included from the main document with:
%
%     
% =================================================================
% REGULATIONS FOR PREPARING THE GRADUATION INTERNSHIP REPORT
% USING LATEX
% =================================================================
%
% This file is intended to be included from the main document with:
%
%     \input{Regulations.tex}
%
% IMPORTANT:
% This section is provided for guidance only. Students shall remove
% or comment out the \input{Regulations.tex} command before submitting
% the final internship report.
%
% =================================================================

\clearpage
\thispagestyle{empty}

\begin{center}
    {\fontsize{16pt}{20pt}\selectfont\bfseries
    SPECIFIC REGULATIONS FOR PREPARING\\
    GRADUATION INTERNSHIP REPORTS USING \LaTeX}
\end{center}

\vspace{0.5cm}

{\fontsize{13pt}{17pt}\selectfont
\setstretch{1.3}
\justifying

\noindent
\textbf{Note:}
This section is provided for guidance only and is included in the main
document using the command \verb|\input{Regulations.tex}|.
Students shall remove or comment out this command before submitting the
final internship report.

\vspace{0.4cm}

% ================================================================
% I. GENERAL REQUIREMENTS
% ================================================================

\noindent
\textbf{I.\quad GENERAL REQUIREMENTS}

\vspace{0.15cm}

The Graduation Internship Report shall be prepared using the official
\LaTeX{} template provided by the Faculty. Students shall preserve the
predefined page layout, typography, line spacing, heading styles,
numbering system, figure and table formatting, and bibliography settings
unless otherwise instructed by the Faculty or the supervisor.

The report shall be written in English.

Students are responsible for the academic and technical content of the
report. The official \LaTeX{} template is responsible for standardized
document formatting.

\vspace{0.3cm}

% ================================================================
% II. REPORT STRUCTURE
% ================================================================

\noindent
\textbf{II.\quad REPORT STRUCTURE}

\vspace{0.15cm}

The report shall normally be arranged in the following order:

\begin{enumerate}[label=\arabic*., leftmargin=1.2cm, itemsep=2pt]
    \item Outer Cover;
    \item Comments of the Internship Host Organization;
    \item Supervisor's Evaluation and Graduation Internship Defense Results;
    \item Table of Contents;
    \item List of Symbols and Abbreviations, if applicable;
    \item List of Figures;
    \item List of Tables;
    \item Abstract;
    \item Main Chapters of the Internship Report;
    \item References;
    \item Appendices, if applicable.
\end{enumerate}

The preliminary pages shall not be numbered as chapters. The main report
content shall be organized into consecutively numbered chapters beginning
with Chapter 1.

\vspace{0.3cm}

% ================================================================
% III. PAGE SIZE AND MARGINS
% ================================================================

\noindent
\textbf{III.\quad PAGE SIZE AND MARGINS}

\vspace{0.15cm}

The report shall be formatted on A4 paper with the following margins:

\begin{center}
\begin{tabular}{ll}
    Left margin:   & 3.5 cm \\
    Right margin:  & 2.0 cm \\
    Top margin:    & 2.5 cm \\
    Bottom margin: & 2.5 cm
\end{tabular}
\end{center}

These settings are predefined in the official template and shall not be
modified by students.

\vspace{0.3cm}

% ================================================================
% IV. FONT, FONT SIZE, AND LINE SPACING
% ================================================================

\noindent
\textbf{IV.\quad FONT, FONT SIZE, AND LINE SPACING}

\vspace{0.15cm}

The main text shall use a 13 pt font size with 1.3-line spacing.
A Times-family typeface shall be used in accordance with the official
template.

The text shall be fully justified.

Students shall not use repeated spaces, repeated manual line breaks, or
multiple blank lines to position text. Paragraph and heading spacing shall
be controlled by the template.

\vspace{0.3cm}

% ================================================================
% V. OUTER COVER
% ================================================================

\noindent
\textbf{V.\quad OUTER COVER}

\vspace{0.15cm}

The outer cover shall follow the official format provided in the template
and shall contain the following information:

\begin{itemize}[leftmargin=0.8cm, itemsep=2pt]
    \item Thai Nguyen University of Technology;
    \item Faculty of International Training;
    \item Engineering and Technology Department;
    \item University logo and Faculty logo;
    \item the title \textbf{INTERNSHIP REPORT};
    \item Internship Organization;
    \item Author;
    \item Class;
    \item Supervisor;
    \item place and year of completion.
\end{itemize}

Students shall only replace the required information fields and shall not
modify the cover-page layout unless instructed otherwise.

\vspace{0.3cm}

% ================================================================
% VI. COMMENTS OF THE INTERNSHIP HOST ORGANIZATION
% ================================================================

\noindent
\textbf{VI.\quad COMMENTS OF THE INTERNSHIP HOST ORGANIZATION}

\vspace{0.15cm}

The page entitled \textbf{COMMENTS OF THE INTERNSHIP HOST ORGANIZATION}
is reserved for the internship host organization.

It includes:

\begin{itemize}[leftmargin=0.8cm, itemsep=2pt]
    \item the student's full name;
    \item class;
    \item comments from the host organization;
    \item score;
    \item date;
    \item signature and official seal.
\end{itemize}

Students shall not alter or complete official evaluation fields unless
specifically authorized.

\vspace{0.3cm}

% ================================================================
% VII. SUPERVISOR'S EVALUATION AND DEFENSE RESULTS
% ================================================================

\noindent
\textbf{VII.\quad SUPERVISOR'S EVALUATION AND DEFENSE RESULTS}

\vspace{0.15cm}

The section entitled
\textbf{SUPERVISOR'S EVALUATION AND GRADUATION INTERNSHIP DEFENSE RESULTS}
is reserved for the supervisor and defense committee.

The evaluation may include:

\begin{itemize}[leftmargin=0.8cm, itemsep=2pt]
    \item Internship Progress and Student's Attitude;
    \item Report Content;
    \item Report Presentation;
    \item Other Comments;
    \item Supervisor's Opinion;
    \item Defense Result;
    \item Score.
\end{itemize}

Students shall not enter evaluation scores or defense results before they
are officially determined.

\vspace{0.3cm}

% ================================================================
% VIII. CHAPTERS, SECTIONS, AND SUBSECTIONS
% ================================================================

\noindent
\textbf{VIII.\quad CHAPTERS, SECTIONS, AND SUBSECTIONS}

\vspace{0.15cm}

The report structure shall be generated using the standard \LaTeX{}
commands:

\begin{verbatim}
\chapter{...}
\section{...}
\subsection{...}
\end{verbatim}

Chapter, section, and subsection numbers shall be generated automatically
by \LaTeX{}.

Students shall not manually type chapter or section numbers and shall not
reproduce heading styles using commands such as \verb|\textbf|,
\verb|\fontsize|, or \verb|\centering| instead of the appropriate
structural command.

Excessive subdivision should be avoided. In most cases, the hierarchy
Chapter--Section--Subsection is sufficient.

\vspace{0.3cm}

% ================================================================
% IX. ORGANIZATION OF THE MAIN CONTENT
% ================================================================

\noindent
\textbf{IX.\quad ORGANIZATION OF THE MAIN CONTENT}

\vspace{0.15cm}

The specific content of the report shall depend on the internship
organization and the assigned tasks. A typical structure may include:

\begin{itemize}[leftmargin=0.8cm, itemsep=2pt]
    \item \textbf{Chapter 1. Overview of the Internship Organization};
    \item \textbf{Chapter 2. Internship Activities and Assigned Tasks};
    \item \textbf{Chapter 3. Results, Experience, and Evaluation}.
\end{itemize}

These chapter titles are recommendations only and may be adapted to the
actual internship content.

The report shall reflect work actually performed, observed, analyzed, or
learned during the internship. Unrelated theoretical material should not
be included merely to increase report length.

\vspace{0.3cm}

% ================================================================
% X. TABLE OF CONTENTS
% ================================================================

\noindent
\textbf{X.\quad TABLE OF CONTENTS}

\vspace{0.15cm}

The Table of Contents shall be generated automatically using:

\begin{verbatim}
\tableofcontents
\end{verbatim}

Students shall not create or manually edit page numbers in the Table of
Contents.

The document shall be compiled sufficiently to ensure that all page
numbers and entries are correctly updated.

\vspace{0.3cm}

% ================================================================
% XI. LIST OF SYMBOLS AND ABBREVIATIONS
% ================================================================

\noindent
\textbf{XI.\quad LIST OF SYMBOLS AND ABBREVIATIONS}

\vspace{0.15cm}

Symbols and abbreviations that are used repeatedly in the report shall be
listed in the \textbf{List of Symbols and Abbreviations}.

The list may contain the following columns:

\begin{center}
\begin{tabular}{llll}
    No. & Symbol / Abbreviation & Unit & Definition
\end{tabular}
\end{center}

Only symbols and abbreviations actually used in the final report shall be
retained. Example entries contained in the template shall be removed if
they are not relevant to the student's report.

Each abbreviation should normally be defined when it first appears in
the main text.

\vspace{0.3cm}

% ================================================================
% XII. FIGURES, DIAGRAMS, AND PHOTOGRAPHS
% ================================================================

\noindent
\textbf{XII.\quad FIGURES, DIAGRAMS, AND PHOTOGRAPHS}

\vspace{0.15cm}

Figures shall be clear, legible, and directly relevant to the report
content.

Figures shall be inserted using the standard \verb|figure| environment.

Example:

\begin{verbatim}
\begin{figure}[H]
    \centering
    \includegraphics[width=0.8\textwidth]
                    {figures/company_structure.pdf}
    \caption{Organizational structure of the company}
    \label{fig:company-structure}
\end{figure}
\end{verbatim}

Figure captions shall be placed below the figure.

Each referenced figure shall have a unique \verb|\label| and shall be
referenced in the text using, for example:

\begin{verbatim}
Figure~\ref{fig:company-structure}
\end{verbatim}

Students shall not manually type figure numbers.

Figures reproduced or adapted from external sources shall include an
appropriate citation.

Photographs taken during the internship shall be relevant and sufficiently
clear. Students shall not include confidential, restricted, or proprietary
information without permission from the internship organization.

Vector graphics or high-resolution images are recommended whenever
possible.

\vspace{0.3cm}

% ================================================================
% XIII. TABLES
% ================================================================

\noindent
\textbf{XIII.\quad TABLES}

\vspace{0.15cm}

Tables shall be created using suitable \LaTeX{} environments such as
\verb|table|, \verb|tabular|, \verb|tabularx|, or \verb|longtable|.

Table captions shall be placed above the table.

Example:

\begin{verbatim}
\begin{table}[H]
    \centering
    \caption{Internship schedule}
    \label{tab:internship-schedule}

    \begin{tabular}{lll}
        \toprule
        Week & Task & Result \\
        \midrule
        1 & Orientation & Completed \\
        2 & System study & Completed \\
        \bottomrule
    \end{tabular}
\end{table}
\end{verbatim}

Tables shall be referenced using:

\begin{verbatim}
Table~\ref{tab:internship-schedule}
\end{verbatim}

Students shall not manually type table numbers.

Tables should be created directly in \LaTeX{} rather than inserted as
screenshots whenever practical.

\vspace{0.3cm}

% ================================================================
% XIV. EQUATIONS AND TECHNICAL QUANTITIES
% ================================================================

\noindent
\textbf{XIV.\quad EQUATIONS AND TECHNICAL QUANTITIES}

\vspace{0.15cm}

Mathematical expressions shall be typeset directly using \LaTeX{}.
Images or screenshots of equations shall not be used.

Short expressions may be written in inline mathematical mode, for example:

\begin{verbatim}
$R = 0.7\,\Omega$
\end{verbatim}

Important equations shall use a numbered equation environment.

Example:

\begin{verbatim}
\begin{equation}
    P = UI
    \label{eq:power}
\end{equation}
\end{verbatim}

Equations shall be referenced using:

\begin{verbatim}
Equation~\eqref{eq:power}
\end{verbatim}

Equation numbers shall not be manually entered.

SI units shall be used wherever applicable. The \verb|siunitx| package
should be used where appropriate.

Examples:

\begin{verbatim}
\SI{220}{\volt}
\SI{50}{\hertz}
\SI{2.5}{\ampere}
\end{verbatim}

Notation and units shall be consistent throughout the report.

\vspace{0.3cm}

% ================================================================
% XV. SOURCE CODE
% ================================================================

\noindent
\textbf{XV.\quad SOURCE CODE}

\vspace{0.15cm}

If the internship involves programming, PLCs, microcontrollers,
simulation, or software development, relevant source code may be included
in the report.

Source code shall be presented using the \verb|listings| environment or
another environment defined by the official template.

Example:

\begin{verbatim}
\begin{lstlisting}[language=st]
IF Start = TRUE THEN
    Motor := TRUE;
ELSE
    Motor := FALSE;
END_IF;
\end{lstlisting}
\end{verbatim}

Only code necessary for explaining the work should appear in the main
report. Long source-code listings should be moved to an appendix or
submitted as supplementary files if required.

Source code should not normally be inserted as screenshots.

\vspace{0.3cm}

% ================================================================
% XVI. ABSTRACT
% ================================================================

\noindent
\textbf{XVI.\quad ABSTRACT}

\vspace{0.15cm}

The Abstract shall briefly summarize:

\begin{itemize}[leftmargin=0.8cm, itemsep=2pt]
    \item the internship organization;
    \item the objectives of the internship;
    \item the main tasks performed;
    \item the main results, knowledge, and skills obtained.
\end{itemize}

The Abstract should be concise and should not reproduce the full content
of the report chapter by chapter.

\vspace{0.3cm}

% ================================================================
% XVII. REFERENCES AND CITATIONS
% ================================================================

\noindent
\textbf{XVII.\quad REFERENCES AND CITATIONS}

\vspace{0.15cm}

All books, journal papers, conference papers, standards, websites,
datasheets, manuals, technical documents, figures, tables, or other
external sources used in preparing the report shall be properly cited.

References shall be managed through the bibliography system provided with
the official template.

Bibliographic entries shall normally be stored in:

\begin{verbatim}
References.bib
\end{verbatim}

References shall be cited in the report using:

\begin{verbatim}
\cite{reference-key}
\end{verbatim}

Students shall not manually type reference numbers such as ``[1]'',
``[2]'', etc.

Titles of foreign-language publications shall remain in their original
language in the bibliography.

All borrowed ideas, data, figures, tables, equations, methods, or results
shall be properly acknowledged.

\vspace{0.3cm}

% ================================================================
% XVIII. CROSS-REFERENCING
% ================================================================

\noindent
\textbf{XVIII.\quad CROSS-REFERENCING}

\vspace{0.15cm}

Numbered elements shall be cross-referenced using the standard \LaTeX{}
mechanism.

Use the following commands where appropriate:

\begin{verbatim}
\label{...}
\ref{...}
\eqref{...}
\end{verbatim}

For example:

\begin{verbatim}
As shown in Figure~\ref{fig:system}, ...
\end{verbatim}

Students shall not manually type numbers such as ``Figure 2.3'' when
those numbers can be generated automatically.

\vspace{0.3cm}

% ================================================================
% XIX. PAGE NUMBERING
% ================================================================

\noindent
\textbf{XIX.\quad PAGE NUMBERING}

\vspace{0.15cm}

Cover pages and official evaluation forms at the beginning of the report
shall not display page numbers.

The main body of the report shall use Arabic page numbers beginning from
Chapter 1.

Page numbering shall be generated automatically by the official template.
Students shall not manually enter page numbers.

\vspace{0.3cm}

% ================================================================
% XX. PROJECT FILE ORGANIZATION
% ================================================================

\noindent
\textbf{XX.\quad PROJECT FILE ORGANIZATION}

\vspace{0.15cm}

Students should maintain a clear and organized \LaTeX{} project
structure. A recommended structure is:

\begin{verbatim}
InternshipReport/
|
|-- main.tex
|-- References.bib
|
|-- figures/
|   |-- company.pdf
|   |-- system.png
|   `-- ...
|
`-- chapters/
    |-- chapter1.tex
    |-- chapter2.tex
    `-- chapter3.tex
\end{verbatim}

Figures should normally be stored in the \verb|figures| directory
specified by the official template.

File and directory names should preferably contain Latin letters,
numbers, hyphens, or underscores.

Recommended examples include:

\begin{verbatim}
company_structure.pdf
motor_control.png
chapter_01.tex
\end{verbatim}

Spaces and unnecessary special characters in file names should be
avoided where possible.

\vspace{0.3cm}

% ================================================================
% XXI. DOCUMENT COMPILATION AND FINAL CHECK
% ================================================================

\noindent
\textbf{XXI.\quad DOCUMENT COMPILATION AND FINAL CHECK}

\vspace{0.15cm}

Before submission, the complete \LaTeX{} project shall compile
successfully and produce the final PDF report.

The final document shall not contain unresolved references such as:

\begin{verbatim}
??
[?]
\end{verbatim}

Before generating the final PDF, students shall check:

\begin{itemize}[leftmargin=0.8cm, itemsep=2pt]
    \item the Table of Contents;
    \item page numbering;
    \item chapter, section, and subsection numbering;
    \item figure numbering;
    \item table numbering;
    \item equation numbering;
    \item cross-references;
    \item citations;
    \item bibliography;
    \item font display;
    \item mathematical symbols;
    \item image quality;
    \item completeness of all required forms.
\end{itemize}

The document may be compiled using TeX Live, MiKTeX, or another
compatible \LaTeX{} environment. Regardless of the compilation
environment, the resulting PDF shall conform to the official Faculty
template.

\vspace{0.3cm}

% ================================================================
% XXII. CONTENT THAT STUDENTS MAY EDIT
% ================================================================

\noindent
\textbf{XXII.\quad CONTENT THAT STUDENTS MAY EDIT}

\vspace{0.15cm}

Students may edit or replace:

\begin{itemize}[leftmargin=0.8cm, itemsep=2pt]
    \item student name;
    \item class;
    \item internship organization;
    \item supervisor name;
    \item Abstract;
    \item List of Symbols and Abbreviations;
    \item chapter titles and chapter content;
    \item figures and photographs;
    \item tables;
    \item equations;
    \item citations and references.
\end{itemize}

Students shall not modify the main document formatting settings in the
preamble, including the document class, page geometry, line spacing,
heading formatting, table-of-contents formatting, font configuration,
page layout, or cover-page macros unless specifically instructed.

\vspace{0.3cm}

% ================================================================
% XXIII. FINAL SUBMISSION REQUIREMENTS
% ================================================================

\noindent
\textbf{XXIII.\quad FINAL SUBMISSION REQUIREMENTS}

\vspace{0.15cm}

Before submission, the final report shall:

\begin{itemize}[leftmargin=0.8cm, itemsep=2pt]
    \item compile without errors that affect the report content or presentation;
    \item contain all required forms and report sections;
    \item contain correct student, class, internship organization, and supervisor information;
    \item contain an updated Table of Contents, List of Figures, and List of Tables;
    \item contain no unused example content from the template;
    \item contain no placeholder text such as \textit{Title of Chapter 1},
    \textit{Section 1.1}, or other sample entries;
    \item contain all required citations and references;
    \item contain no unresolved cross-references.
\end{itemize}

If the Faculty requires submission of the \LaTeX{} source, the complete
project shall include all \verb|.tex| files, the bibliography database,
figures, and any other files required to reproduce the submitted PDF.

\vspace{0.4cm}

\noindent
\textbf{Important:}
Before submitting the final internship report, students shall remove or
comment out the following command in the main document:

\begin{verbatim}
\input{Regulations.tex}
\end{verbatim}

For example:

\begin{verbatim}
% \input{Regulations.tex}
\end{verbatim}

This removes the guidance pages from the final submitted report while
preserving the regulations file in the \LaTeX{} project for reference.

} % End 13pt content

\clearpage

%
% IMPORTANT:
% This section is provided for guidance only. Students shall remove
% or comment out the 
% =================================================================
% REGULATIONS FOR PREPARING THE GRADUATION INTERNSHIP REPORT
% USING LATEX
% =================================================================
%
% This file is intended to be included from the main document with:
%
%     \input{Regulations.tex}
%
% IMPORTANT:
% This section is provided for guidance only. Students shall remove
% or comment out the \input{Regulations.tex} command before submitting
% the final internship report.
%
% =================================================================

\clearpage
\thispagestyle{empty}

\begin{center}
    {\fontsize{16pt}{20pt}\selectfont\bfseries
    SPECIFIC REGULATIONS FOR PREPARING\\
    GRADUATION INTERNSHIP REPORTS USING \LaTeX}
\end{center}

\vspace{0.5cm}

{\fontsize{13pt}{17pt}\selectfont
\setstretch{1.3}
\justifying

\noindent
\textbf{Note:}
This section is provided for guidance only and is included in the main
document using the command \verb|\input{Regulations.tex}|.
Students shall remove or comment out this command before submitting the
final internship report.

\vspace{0.4cm}

% ================================================================
% I. GENERAL REQUIREMENTS
% ================================================================

\noindent
\textbf{I.\quad GENERAL REQUIREMENTS}

\vspace{0.15cm}

The Graduation Internship Report shall be prepared using the official
\LaTeX{} template provided by the Faculty. Students shall preserve the
predefined page layout, typography, line spacing, heading styles,
numbering system, figure and table formatting, and bibliography settings
unless otherwise instructed by the Faculty or the supervisor.

The report shall be written in English.

Students are responsible for the academic and technical content of the
report. The official \LaTeX{} template is responsible for standardized
document formatting.

\vspace{0.3cm}

% ================================================================
% II. REPORT STRUCTURE
% ================================================================

\noindent
\textbf{II.\quad REPORT STRUCTURE}

\vspace{0.15cm}

The report shall normally be arranged in the following order:

\begin{enumerate}[label=\arabic*., leftmargin=1.2cm, itemsep=2pt]
    \item Outer Cover;
    \item Comments of the Internship Host Organization;
    \item Supervisor's Evaluation and Graduation Internship Defense Results;
    \item Table of Contents;
    \item List of Symbols and Abbreviations, if applicable;
    \item List of Figures;
    \item List of Tables;
    \item Abstract;
    \item Main Chapters of the Internship Report;
    \item References;
    \item Appendices, if applicable.
\end{enumerate}

The preliminary pages shall not be numbered as chapters. The main report
content shall be organized into consecutively numbered chapters beginning
with Chapter 1.

\vspace{0.3cm}

% ================================================================
% III. PAGE SIZE AND MARGINS
% ================================================================

\noindent
\textbf{III.\quad PAGE SIZE AND MARGINS}

\vspace{0.15cm}

The report shall be formatted on A4 paper with the following margins:

\begin{center}
\begin{tabular}{ll}
    Left margin:   & 3.5 cm \\
    Right margin:  & 2.0 cm \\
    Top margin:    & 2.5 cm \\
    Bottom margin: & 2.5 cm
\end{tabular}
\end{center}

These settings are predefined in the official template and shall not be
modified by students.

\vspace{0.3cm}

% ================================================================
% IV. FONT, FONT SIZE, AND LINE SPACING
% ================================================================

\noindent
\textbf{IV.\quad FONT, FONT SIZE, AND LINE SPACING}

\vspace{0.15cm}

The main text shall use a 13 pt font size with 1.3-line spacing.
A Times-family typeface shall be used in accordance with the official
template.

The text shall be fully justified.

Students shall not use repeated spaces, repeated manual line breaks, or
multiple blank lines to position text. Paragraph and heading spacing shall
be controlled by the template.

\vspace{0.3cm}

% ================================================================
% V. OUTER COVER
% ================================================================

\noindent
\textbf{V.\quad OUTER COVER}

\vspace{0.15cm}

The outer cover shall follow the official format provided in the template
and shall contain the following information:

\begin{itemize}[leftmargin=0.8cm, itemsep=2pt]
    \item Thai Nguyen University of Technology;
    \item Faculty of International Training;
    \item Engineering and Technology Department;
    \item University logo and Faculty logo;
    \item the title \textbf{INTERNSHIP REPORT};
    \item Internship Organization;
    \item Author;
    \item Class;
    \item Supervisor;
    \item place and year of completion.
\end{itemize}

Students shall only replace the required information fields and shall not
modify the cover-page layout unless instructed otherwise.

\vspace{0.3cm}

% ================================================================
% VI. COMMENTS OF THE INTERNSHIP HOST ORGANIZATION
% ================================================================

\noindent
\textbf{VI.\quad COMMENTS OF THE INTERNSHIP HOST ORGANIZATION}

\vspace{0.15cm}

The page entitled \textbf{COMMENTS OF THE INTERNSHIP HOST ORGANIZATION}
is reserved for the internship host organization.

It includes:

\begin{itemize}[leftmargin=0.8cm, itemsep=2pt]
    \item the student's full name;
    \item class;
    \item comments from the host organization;
    \item score;
    \item date;
    \item signature and official seal.
\end{itemize}

Students shall not alter or complete official evaluation fields unless
specifically authorized.

\vspace{0.3cm}

% ================================================================
% VII. SUPERVISOR'S EVALUATION AND DEFENSE RESULTS
% ================================================================

\noindent
\textbf{VII.\quad SUPERVISOR'S EVALUATION AND DEFENSE RESULTS}

\vspace{0.15cm}

The section entitled
\textbf{SUPERVISOR'S EVALUATION AND GRADUATION INTERNSHIP DEFENSE RESULTS}
is reserved for the supervisor and defense committee.

The evaluation may include:

\begin{itemize}[leftmargin=0.8cm, itemsep=2pt]
    \item Internship Progress and Student's Attitude;
    \item Report Content;
    \item Report Presentation;
    \item Other Comments;
    \item Supervisor's Opinion;
    \item Defense Result;
    \item Score.
\end{itemize}

Students shall not enter evaluation scores or defense results before they
are officially determined.

\vspace{0.3cm}

% ================================================================
% VIII. CHAPTERS, SECTIONS, AND SUBSECTIONS
% ================================================================

\noindent
\textbf{VIII.\quad CHAPTERS, SECTIONS, AND SUBSECTIONS}

\vspace{0.15cm}

The report structure shall be generated using the standard \LaTeX{}
commands:

\begin{verbatim}
\chapter{...}
\section{...}
\subsection{...}
\end{verbatim}

Chapter, section, and subsection numbers shall be generated automatically
by \LaTeX{}.

Students shall not manually type chapter or section numbers and shall not
reproduce heading styles using commands such as \verb|\textbf|,
\verb|\fontsize|, or \verb|\centering| instead of the appropriate
structural command.

Excessive subdivision should be avoided. In most cases, the hierarchy
Chapter--Section--Subsection is sufficient.

\vspace{0.3cm}

% ================================================================
% IX. ORGANIZATION OF THE MAIN CONTENT
% ================================================================

\noindent
\textbf{IX.\quad ORGANIZATION OF THE MAIN CONTENT}

\vspace{0.15cm}

The specific content of the report shall depend on the internship
organization and the assigned tasks. A typical structure may include:

\begin{itemize}[leftmargin=0.8cm, itemsep=2pt]
    \item \textbf{Chapter 1. Overview of the Internship Organization};
    \item \textbf{Chapter 2. Internship Activities and Assigned Tasks};
    \item \textbf{Chapter 3. Results, Experience, and Evaluation}.
\end{itemize}

These chapter titles are recommendations only and may be adapted to the
actual internship content.

The report shall reflect work actually performed, observed, analyzed, or
learned during the internship. Unrelated theoretical material should not
be included merely to increase report length.

\vspace{0.3cm}

% ================================================================
% X. TABLE OF CONTENTS
% ================================================================

\noindent
\textbf{X.\quad TABLE OF CONTENTS}

\vspace{0.15cm}

The Table of Contents shall be generated automatically using:

\begin{verbatim}
\tableofcontents
\end{verbatim}

Students shall not create or manually edit page numbers in the Table of
Contents.

The document shall be compiled sufficiently to ensure that all page
numbers and entries are correctly updated.

\vspace{0.3cm}

% ================================================================
% XI. LIST OF SYMBOLS AND ABBREVIATIONS
% ================================================================

\noindent
\textbf{XI.\quad LIST OF SYMBOLS AND ABBREVIATIONS}

\vspace{0.15cm}

Symbols and abbreviations that are used repeatedly in the report shall be
listed in the \textbf{List of Symbols and Abbreviations}.

The list may contain the following columns:

\begin{center}
\begin{tabular}{llll}
    No. & Symbol / Abbreviation & Unit & Definition
\end{tabular}
\end{center}

Only symbols and abbreviations actually used in the final report shall be
retained. Example entries contained in the template shall be removed if
they are not relevant to the student's report.

Each abbreviation should normally be defined when it first appears in
the main text.

\vspace{0.3cm}

% ================================================================
% XII. FIGURES, DIAGRAMS, AND PHOTOGRAPHS
% ================================================================

\noindent
\textbf{XII.\quad FIGURES, DIAGRAMS, AND PHOTOGRAPHS}

\vspace{0.15cm}

Figures shall be clear, legible, and directly relevant to the report
content.

Figures shall be inserted using the standard \verb|figure| environment.

Example:

\begin{verbatim}
\begin{figure}[H]
    \centering
    \includegraphics[width=0.8\textwidth]
                    {figures/company_structure.pdf}
    \caption{Organizational structure of the company}
    \label{fig:company-structure}
\end{figure}
\end{verbatim}

Figure captions shall be placed below the figure.

Each referenced figure shall have a unique \verb|\label| and shall be
referenced in the text using, for example:

\begin{verbatim}
Figure~\ref{fig:company-structure}
\end{verbatim}

Students shall not manually type figure numbers.

Figures reproduced or adapted from external sources shall include an
appropriate citation.

Photographs taken during the internship shall be relevant and sufficiently
clear. Students shall not include confidential, restricted, or proprietary
information without permission from the internship organization.

Vector graphics or high-resolution images are recommended whenever
possible.

\vspace{0.3cm}

% ================================================================
% XIII. TABLES
% ================================================================

\noindent
\textbf{XIII.\quad TABLES}

\vspace{0.15cm}

Tables shall be created using suitable \LaTeX{} environments such as
\verb|table|, \verb|tabular|, \verb|tabularx|, or \verb|longtable|.

Table captions shall be placed above the table.

Example:

\begin{verbatim}
\begin{table}[H]
    \centering
    \caption{Internship schedule}
    \label{tab:internship-schedule}

    \begin{tabular}{lll}
        \toprule
        Week & Task & Result \\
        \midrule
        1 & Orientation & Completed \\
        2 & System study & Completed \\
        \bottomrule
    \end{tabular}
\end{table}
\end{verbatim}

Tables shall be referenced using:

\begin{verbatim}
Table~\ref{tab:internship-schedule}
\end{verbatim}

Students shall not manually type table numbers.

Tables should be created directly in \LaTeX{} rather than inserted as
screenshots whenever practical.

\vspace{0.3cm}

% ================================================================
% XIV. EQUATIONS AND TECHNICAL QUANTITIES
% ================================================================

\noindent
\textbf{XIV.\quad EQUATIONS AND TECHNICAL QUANTITIES}

\vspace{0.15cm}

Mathematical expressions shall be typeset directly using \LaTeX{}.
Images or screenshots of equations shall not be used.

Short expressions may be written in inline mathematical mode, for example:

\begin{verbatim}
$R = 0.7\,\Omega$
\end{verbatim}

Important equations shall use a numbered equation environment.

Example:

\begin{verbatim}
\begin{equation}
    P = UI
    \label{eq:power}
\end{equation}
\end{verbatim}

Equations shall be referenced using:

\begin{verbatim}
Equation~\eqref{eq:power}
\end{verbatim}

Equation numbers shall not be manually entered.

SI units shall be used wherever applicable. The \verb|siunitx| package
should be used where appropriate.

Examples:

\begin{verbatim}
\SI{220}{\volt}
\SI{50}{\hertz}
\SI{2.5}{\ampere}
\end{verbatim}

Notation and units shall be consistent throughout the report.

\vspace{0.3cm}

% ================================================================
% XV. SOURCE CODE
% ================================================================

\noindent
\textbf{XV.\quad SOURCE CODE}

\vspace{0.15cm}

If the internship involves programming, PLCs, microcontrollers,
simulation, or software development, relevant source code may be included
in the report.

Source code shall be presented using the \verb|listings| environment or
another environment defined by the official template.

Example:

\begin{verbatim}
\begin{lstlisting}[language=st]
IF Start = TRUE THEN
    Motor := TRUE;
ELSE
    Motor := FALSE;
END_IF;
\end{lstlisting}
\end{verbatim}

Only code necessary for explaining the work should appear in the main
report. Long source-code listings should be moved to an appendix or
submitted as supplementary files if required.

Source code should not normally be inserted as screenshots.

\vspace{0.3cm}

% ================================================================
% XVI. ABSTRACT
% ================================================================

\noindent
\textbf{XVI.\quad ABSTRACT}

\vspace{0.15cm}

The Abstract shall briefly summarize:

\begin{itemize}[leftmargin=0.8cm, itemsep=2pt]
    \item the internship organization;
    \item the objectives of the internship;
    \item the main tasks performed;
    \item the main results, knowledge, and skills obtained.
\end{itemize}

The Abstract should be concise and should not reproduce the full content
of the report chapter by chapter.

\vspace{0.3cm}

% ================================================================
% XVII. REFERENCES AND CITATIONS
% ================================================================

\noindent
\textbf{XVII.\quad REFERENCES AND CITATIONS}

\vspace{0.15cm}

All books, journal papers, conference papers, standards, websites,
datasheets, manuals, technical documents, figures, tables, or other
external sources used in preparing the report shall be properly cited.

References shall be managed through the bibliography system provided with
the official template.

Bibliographic entries shall normally be stored in:

\begin{verbatim}
References.bib
\end{verbatim}

References shall be cited in the report using:

\begin{verbatim}
\cite{reference-key}
\end{verbatim}

Students shall not manually type reference numbers such as ``[1]'',
``[2]'', etc.

Titles of foreign-language publications shall remain in their original
language in the bibliography.

All borrowed ideas, data, figures, tables, equations, methods, or results
shall be properly acknowledged.

\vspace{0.3cm}

% ================================================================
% XVIII. CROSS-REFERENCING
% ================================================================

\noindent
\textbf{XVIII.\quad CROSS-REFERENCING}

\vspace{0.15cm}

Numbered elements shall be cross-referenced using the standard \LaTeX{}
mechanism.

Use the following commands where appropriate:

\begin{verbatim}
\label{...}
\ref{...}
\eqref{...}
\end{verbatim}

For example:

\begin{verbatim}
As shown in Figure~\ref{fig:system}, ...
\end{verbatim}

Students shall not manually type numbers such as ``Figure 2.3'' when
those numbers can be generated automatically.

\vspace{0.3cm}

% ================================================================
% XIX. PAGE NUMBERING
% ================================================================

\noindent
\textbf{XIX.\quad PAGE NUMBERING}

\vspace{0.15cm}

Cover pages and official evaluation forms at the beginning of the report
shall not display page numbers.

The main body of the report shall use Arabic page numbers beginning from
Chapter 1.

Page numbering shall be generated automatically by the official template.
Students shall not manually enter page numbers.

\vspace{0.3cm}

% ================================================================
% XX. PROJECT FILE ORGANIZATION
% ================================================================

\noindent
\textbf{XX.\quad PROJECT FILE ORGANIZATION}

\vspace{0.15cm}

Students should maintain a clear and organized \LaTeX{} project
structure. A recommended structure is:

\begin{verbatim}
InternshipReport/
|
|-- main.tex
|-- References.bib
|
|-- figures/
|   |-- company.pdf
|   |-- system.png
|   `-- ...
|
`-- chapters/
    |-- chapter1.tex
    |-- chapter2.tex
    `-- chapter3.tex
\end{verbatim}

Figures should normally be stored in the \verb|figures| directory
specified by the official template.

File and directory names should preferably contain Latin letters,
numbers, hyphens, or underscores.

Recommended examples include:

\begin{verbatim}
company_structure.pdf
motor_control.png
chapter_01.tex
\end{verbatim}

Spaces and unnecessary special characters in file names should be
avoided where possible.

\vspace{0.3cm}

% ================================================================
% XXI. DOCUMENT COMPILATION AND FINAL CHECK
% ================================================================

\noindent
\textbf{XXI.\quad DOCUMENT COMPILATION AND FINAL CHECK}

\vspace{0.15cm}

Before submission, the complete \LaTeX{} project shall compile
successfully and produce the final PDF report.

The final document shall not contain unresolved references such as:

\begin{verbatim}
??
[?]
\end{verbatim}

Before generating the final PDF, students shall check:

\begin{itemize}[leftmargin=0.8cm, itemsep=2pt]
    \item the Table of Contents;
    \item page numbering;
    \item chapter, section, and subsection numbering;
    \item figure numbering;
    \item table numbering;
    \item equation numbering;
    \item cross-references;
    \item citations;
    \item bibliography;
    \item font display;
    \item mathematical symbols;
    \item image quality;
    \item completeness of all required forms.
\end{itemize}

The document may be compiled using TeX Live, MiKTeX, or another
compatible \LaTeX{} environment. Regardless of the compilation
environment, the resulting PDF shall conform to the official Faculty
template.

\vspace{0.3cm}

% ================================================================
% XXII. CONTENT THAT STUDENTS MAY EDIT
% ================================================================

\noindent
\textbf{XXII.\quad CONTENT THAT STUDENTS MAY EDIT}

\vspace{0.15cm}

Students may edit or replace:

\begin{itemize}[leftmargin=0.8cm, itemsep=2pt]
    \item student name;
    \item class;
    \item internship organization;
    \item supervisor name;
    \item Abstract;
    \item List of Symbols and Abbreviations;
    \item chapter titles and chapter content;
    \item figures and photographs;
    \item tables;
    \item equations;
    \item citations and references.
\end{itemize}

Students shall not modify the main document formatting settings in the
preamble, including the document class, page geometry, line spacing,
heading formatting, table-of-contents formatting, font configuration,
page layout, or cover-page macros unless specifically instructed.

\vspace{0.3cm}

% ================================================================
% XXIII. FINAL SUBMISSION REQUIREMENTS
% ================================================================

\noindent
\textbf{XXIII.\quad FINAL SUBMISSION REQUIREMENTS}

\vspace{0.15cm}

Before submission, the final report shall:

\begin{itemize}[leftmargin=0.8cm, itemsep=2pt]
    \item compile without errors that affect the report content or presentation;
    \item contain all required forms and report sections;
    \item contain correct student, class, internship organization, and supervisor information;
    \item contain an updated Table of Contents, List of Figures, and List of Tables;
    \item contain no unused example content from the template;
    \item contain no placeholder text such as \textit{Title of Chapter 1},
    \textit{Section 1.1}, or other sample entries;
    \item contain all required citations and references;
    \item contain no unresolved cross-references.
\end{itemize}

If the Faculty requires submission of the \LaTeX{} source, the complete
project shall include all \verb|.tex| files, the bibliography database,
figures, and any other files required to reproduce the submitted PDF.

\vspace{0.4cm}

\noindent
\textbf{Important:}
Before submitting the final internship report, students shall remove or
comment out the following command in the main document:

\begin{verbatim}
\input{Regulations.tex}
\end{verbatim}

For example:

\begin{verbatim}
% \input{Regulations.tex}
\end{verbatim}

This removes the guidance pages from the final submitted report while
preserving the regulations file in the \LaTeX{} project for reference.

} % End 13pt content

\clearpage
 command before submitting
% the final internship report.
%
% =================================================================

\clearpage
\thispagestyle{empty}

\begin{center}
    {\fontsize{16pt}{20pt}\selectfont\bfseries
    SPECIFIC REGULATIONS FOR PREPARING\\
    GRADUATION INTERNSHIP REPORTS USING \LaTeX}
\end{center}

\vspace{0.5cm}

{\fontsize{13pt}{17pt}\selectfont
\setstretch{1.3}
\justifying

\noindent
\textbf{Note:}
This section is provided for guidance only and is included in the main
document using the command \verb|
% =================================================================
% REGULATIONS FOR PREPARING THE GRADUATION INTERNSHIP REPORT
% USING LATEX
% =================================================================
%
% This file is intended to be included from the main document with:
%
%     \input{Regulations.tex}
%
% IMPORTANT:
% This section is provided for guidance only. Students shall remove
% or comment out the \input{Regulations.tex} command before submitting
% the final internship report.
%
% =================================================================

\clearpage
\thispagestyle{empty}

\begin{center}
    {\fontsize{16pt}{20pt}\selectfont\bfseries
    SPECIFIC REGULATIONS FOR PREPARING\\
    GRADUATION INTERNSHIP REPORTS USING \LaTeX}
\end{center}

\vspace{0.5cm}

{\fontsize{13pt}{17pt}\selectfont
\setstretch{1.3}
\justifying

\noindent
\textbf{Note:}
This section is provided for guidance only and is included in the main
document using the command \verb|\input{Regulations.tex}|.
Students shall remove or comment out this command before submitting the
final internship report.

\vspace{0.4cm}

% ================================================================
% I. GENERAL REQUIREMENTS
% ================================================================

\noindent
\textbf{I.\quad GENERAL REQUIREMENTS}

\vspace{0.15cm}

The Graduation Internship Report shall be prepared using the official
\LaTeX{} template provided by the Faculty. Students shall preserve the
predefined page layout, typography, line spacing, heading styles,
numbering system, figure and table formatting, and bibliography settings
unless otherwise instructed by the Faculty or the supervisor.

The report shall be written in English.

Students are responsible for the academic and technical content of the
report. The official \LaTeX{} template is responsible for standardized
document formatting.

\vspace{0.3cm}

% ================================================================
% II. REPORT STRUCTURE
% ================================================================

\noindent
\textbf{II.\quad REPORT STRUCTURE}

\vspace{0.15cm}

The report shall normally be arranged in the following order:

\begin{enumerate}[label=\arabic*., leftmargin=1.2cm, itemsep=2pt]
    \item Outer Cover;
    \item Comments of the Internship Host Organization;
    \item Supervisor's Evaluation and Graduation Internship Defense Results;
    \item Table of Contents;
    \item List of Symbols and Abbreviations, if applicable;
    \item List of Figures;
    \item List of Tables;
    \item Abstract;
    \item Main Chapters of the Internship Report;
    \item References;
    \item Appendices, if applicable.
\end{enumerate}

The preliminary pages shall not be numbered as chapters. The main report
content shall be organized into consecutively numbered chapters beginning
with Chapter 1.

\vspace{0.3cm}

% ================================================================
% III. PAGE SIZE AND MARGINS
% ================================================================

\noindent
\textbf{III.\quad PAGE SIZE AND MARGINS}

\vspace{0.15cm}

The report shall be formatted on A4 paper with the following margins:

\begin{center}
\begin{tabular}{ll}
    Left margin:   & 3.5 cm \\
    Right margin:  & 2.0 cm \\
    Top margin:    & 2.5 cm \\
    Bottom margin: & 2.5 cm
\end{tabular}
\end{center}

These settings are predefined in the official template and shall not be
modified by students.

\vspace{0.3cm}

% ================================================================
% IV. FONT, FONT SIZE, AND LINE SPACING
% ================================================================

\noindent
\textbf{IV.\quad FONT, FONT SIZE, AND LINE SPACING}

\vspace{0.15cm}

The main text shall use a 13 pt font size with 1.3-line spacing.
A Times-family typeface shall be used in accordance with the official
template.

The text shall be fully justified.

Students shall not use repeated spaces, repeated manual line breaks, or
multiple blank lines to position text. Paragraph and heading spacing shall
be controlled by the template.

\vspace{0.3cm}

% ================================================================
% V. OUTER COVER
% ================================================================

\noindent
\textbf{V.\quad OUTER COVER}

\vspace{0.15cm}

The outer cover shall follow the official format provided in the template
and shall contain the following information:

\begin{itemize}[leftmargin=0.8cm, itemsep=2pt]
    \item Thai Nguyen University of Technology;
    \item Faculty of International Training;
    \item Engineering and Technology Department;
    \item University logo and Faculty logo;
    \item the title \textbf{INTERNSHIP REPORT};
    \item Internship Organization;
    \item Author;
    \item Class;
    \item Supervisor;
    \item place and year of completion.
\end{itemize}

Students shall only replace the required information fields and shall not
modify the cover-page layout unless instructed otherwise.

\vspace{0.3cm}

% ================================================================
% VI. COMMENTS OF THE INTERNSHIP HOST ORGANIZATION
% ================================================================

\noindent
\textbf{VI.\quad COMMENTS OF THE INTERNSHIP HOST ORGANIZATION}

\vspace{0.15cm}

The page entitled \textbf{COMMENTS OF THE INTERNSHIP HOST ORGANIZATION}
is reserved for the internship host organization.

It includes:

\begin{itemize}[leftmargin=0.8cm, itemsep=2pt]
    \item the student's full name;
    \item class;
    \item comments from the host organization;
    \item score;
    \item date;
    \item signature and official seal.
\end{itemize}

Students shall not alter or complete official evaluation fields unless
specifically authorized.

\vspace{0.3cm}

% ================================================================
% VII. SUPERVISOR'S EVALUATION AND DEFENSE RESULTS
% ================================================================

\noindent
\textbf{VII.\quad SUPERVISOR'S EVALUATION AND DEFENSE RESULTS}

\vspace{0.15cm}

The section entitled
\textbf{SUPERVISOR'S EVALUATION AND GRADUATION INTERNSHIP DEFENSE RESULTS}
is reserved for the supervisor and defense committee.

The evaluation may include:

\begin{itemize}[leftmargin=0.8cm, itemsep=2pt]
    \item Internship Progress and Student's Attitude;
    \item Report Content;
    \item Report Presentation;
    \item Other Comments;
    \item Supervisor's Opinion;
    \item Defense Result;
    \item Score.
\end{itemize}

Students shall not enter evaluation scores or defense results before they
are officially determined.

\vspace{0.3cm}

% ================================================================
% VIII. CHAPTERS, SECTIONS, AND SUBSECTIONS
% ================================================================

\noindent
\textbf{VIII.\quad CHAPTERS, SECTIONS, AND SUBSECTIONS}

\vspace{0.15cm}

The report structure shall be generated using the standard \LaTeX{}
commands:

\begin{verbatim}
\chapter{...}
\section{...}
\subsection{...}
\end{verbatim}

Chapter, section, and subsection numbers shall be generated automatically
by \LaTeX{}.

Students shall not manually type chapter or section numbers and shall not
reproduce heading styles using commands such as \verb|\textbf|,
\verb|\fontsize|, or \verb|\centering| instead of the appropriate
structural command.

Excessive subdivision should be avoided. In most cases, the hierarchy
Chapter--Section--Subsection is sufficient.

\vspace{0.3cm}

% ================================================================
% IX. ORGANIZATION OF THE MAIN CONTENT
% ================================================================

\noindent
\textbf{IX.\quad ORGANIZATION OF THE MAIN CONTENT}

\vspace{0.15cm}

The specific content of the report shall depend on the internship
organization and the assigned tasks. A typical structure may include:

\begin{itemize}[leftmargin=0.8cm, itemsep=2pt]
    \item \textbf{Chapter 1. Overview of the Internship Organization};
    \item \textbf{Chapter 2. Internship Activities and Assigned Tasks};
    \item \textbf{Chapter 3. Results, Experience, and Evaluation}.
\end{itemize}

These chapter titles are recommendations only and may be adapted to the
actual internship content.

The report shall reflect work actually performed, observed, analyzed, or
learned during the internship. Unrelated theoretical material should not
be included merely to increase report length.

\vspace{0.3cm}

% ================================================================
% X. TABLE OF CONTENTS
% ================================================================

\noindent
\textbf{X.\quad TABLE OF CONTENTS}

\vspace{0.15cm}

The Table of Contents shall be generated automatically using:

\begin{verbatim}
\tableofcontents
\end{verbatim}

Students shall not create or manually edit page numbers in the Table of
Contents.

The document shall be compiled sufficiently to ensure that all page
numbers and entries are correctly updated.

\vspace{0.3cm}

% ================================================================
% XI. LIST OF SYMBOLS AND ABBREVIATIONS
% ================================================================

\noindent
\textbf{XI.\quad LIST OF SYMBOLS AND ABBREVIATIONS}

\vspace{0.15cm}

Symbols and abbreviations that are used repeatedly in the report shall be
listed in the \textbf{List of Symbols and Abbreviations}.

The list may contain the following columns:

\begin{center}
\begin{tabular}{llll}
    No. & Symbol / Abbreviation & Unit & Definition
\end{tabular}
\end{center}

Only symbols and abbreviations actually used in the final report shall be
retained. Example entries contained in the template shall be removed if
they are not relevant to the student's report.

Each abbreviation should normally be defined when it first appears in
the main text.

\vspace{0.3cm}

% ================================================================
% XII. FIGURES, DIAGRAMS, AND PHOTOGRAPHS
% ================================================================

\noindent
\textbf{XII.\quad FIGURES, DIAGRAMS, AND PHOTOGRAPHS}

\vspace{0.15cm}

Figures shall be clear, legible, and directly relevant to the report
content.

Figures shall be inserted using the standard \verb|figure| environment.

Example:

\begin{verbatim}
\begin{figure}[H]
    \centering
    \includegraphics[width=0.8\textwidth]
                    {figures/company_structure.pdf}
    \caption{Organizational structure of the company}
    \label{fig:company-structure}
\end{figure}
\end{verbatim}

Figure captions shall be placed below the figure.

Each referenced figure shall have a unique \verb|\label| and shall be
referenced in the text using, for example:

\begin{verbatim}
Figure~\ref{fig:company-structure}
\end{verbatim}

Students shall not manually type figure numbers.

Figures reproduced or adapted from external sources shall include an
appropriate citation.

Photographs taken during the internship shall be relevant and sufficiently
clear. Students shall not include confidential, restricted, or proprietary
information without permission from the internship organization.

Vector graphics or high-resolution images are recommended whenever
possible.

\vspace{0.3cm}

% ================================================================
% XIII. TABLES
% ================================================================

\noindent
\textbf{XIII.\quad TABLES}

\vspace{0.15cm}

Tables shall be created using suitable \LaTeX{} environments such as
\verb|table|, \verb|tabular|, \verb|tabularx|, or \verb|longtable|.

Table captions shall be placed above the table.

Example:

\begin{verbatim}
\begin{table}[H]
    \centering
    \caption{Internship schedule}
    \label{tab:internship-schedule}

    \begin{tabular}{lll}
        \toprule
        Week & Task & Result \\
        \midrule
        1 & Orientation & Completed \\
        2 & System study & Completed \\
        \bottomrule
    \end{tabular}
\end{table}
\end{verbatim}

Tables shall be referenced using:

\begin{verbatim}
Table~\ref{tab:internship-schedule}
\end{verbatim}

Students shall not manually type table numbers.

Tables should be created directly in \LaTeX{} rather than inserted as
screenshots whenever practical.

\vspace{0.3cm}

% ================================================================
% XIV. EQUATIONS AND TECHNICAL QUANTITIES
% ================================================================

\noindent
\textbf{XIV.\quad EQUATIONS AND TECHNICAL QUANTITIES}

\vspace{0.15cm}

Mathematical expressions shall be typeset directly using \LaTeX{}.
Images or screenshots of equations shall not be used.

Short expressions may be written in inline mathematical mode, for example:

\begin{verbatim}
$R = 0.7\,\Omega$
\end{verbatim}

Important equations shall use a numbered equation environment.

Example:

\begin{verbatim}
\begin{equation}
    P = UI
    \label{eq:power}
\end{equation}
\end{verbatim}

Equations shall be referenced using:

\begin{verbatim}
Equation~\eqref{eq:power}
\end{verbatim}

Equation numbers shall not be manually entered.

SI units shall be used wherever applicable. The \verb|siunitx| package
should be used where appropriate.

Examples:

\begin{verbatim}
\SI{220}{\volt}
\SI{50}{\hertz}
\SI{2.5}{\ampere}
\end{verbatim}

Notation and units shall be consistent throughout the report.

\vspace{0.3cm}

% ================================================================
% XV. SOURCE CODE
% ================================================================

\noindent
\textbf{XV.\quad SOURCE CODE}

\vspace{0.15cm}

If the internship involves programming, PLCs, microcontrollers,
simulation, or software development, relevant source code may be included
in the report.

Source code shall be presented using the \verb|listings| environment or
another environment defined by the official template.

Example:

\begin{verbatim}
\begin{lstlisting}[language=st]
IF Start = TRUE THEN
    Motor := TRUE;
ELSE
    Motor := FALSE;
END_IF;
\end{lstlisting}
\end{verbatim}

Only code necessary for explaining the work should appear in the main
report. Long source-code listings should be moved to an appendix or
submitted as supplementary files if required.

Source code should not normally be inserted as screenshots.

\vspace{0.3cm}

% ================================================================
% XVI. ABSTRACT
% ================================================================

\noindent
\textbf{XVI.\quad ABSTRACT}

\vspace{0.15cm}

The Abstract shall briefly summarize:

\begin{itemize}[leftmargin=0.8cm, itemsep=2pt]
    \item the internship organization;
    \item the objectives of the internship;
    \item the main tasks performed;
    \item the main results, knowledge, and skills obtained.
\end{itemize}

The Abstract should be concise and should not reproduce the full content
of the report chapter by chapter.

\vspace{0.3cm}

% ================================================================
% XVII. REFERENCES AND CITATIONS
% ================================================================

\noindent
\textbf{XVII.\quad REFERENCES AND CITATIONS}

\vspace{0.15cm}

All books, journal papers, conference papers, standards, websites,
datasheets, manuals, technical documents, figures, tables, or other
external sources used in preparing the report shall be properly cited.

References shall be managed through the bibliography system provided with
the official template.

Bibliographic entries shall normally be stored in:

\begin{verbatim}
References.bib
\end{verbatim}

References shall be cited in the report using:

\begin{verbatim}
\cite{reference-key}
\end{verbatim}

Students shall not manually type reference numbers such as ``[1]'',
``[2]'', etc.

Titles of foreign-language publications shall remain in their original
language in the bibliography.

All borrowed ideas, data, figures, tables, equations, methods, or results
shall be properly acknowledged.

\vspace{0.3cm}

% ================================================================
% XVIII. CROSS-REFERENCING
% ================================================================

\noindent
\textbf{XVIII.\quad CROSS-REFERENCING}

\vspace{0.15cm}

Numbered elements shall be cross-referenced using the standard \LaTeX{}
mechanism.

Use the following commands where appropriate:

\begin{verbatim}
\label{...}
\ref{...}
\eqref{...}
\end{verbatim}

For example:

\begin{verbatim}
As shown in Figure~\ref{fig:system}, ...
\end{verbatim}

Students shall not manually type numbers such as ``Figure 2.3'' when
those numbers can be generated automatically.

\vspace{0.3cm}

% ================================================================
% XIX. PAGE NUMBERING
% ================================================================

\noindent
\textbf{XIX.\quad PAGE NUMBERING}

\vspace{0.15cm}

Cover pages and official evaluation forms at the beginning of the report
shall not display page numbers.

The main body of the report shall use Arabic page numbers beginning from
Chapter 1.

Page numbering shall be generated automatically by the official template.
Students shall not manually enter page numbers.

\vspace{0.3cm}

% ================================================================
% XX. PROJECT FILE ORGANIZATION
% ================================================================

\noindent
\textbf{XX.\quad PROJECT FILE ORGANIZATION}

\vspace{0.15cm}

Students should maintain a clear and organized \LaTeX{} project
structure. A recommended structure is:

\begin{verbatim}
InternshipReport/
|
|-- main.tex
|-- References.bib
|
|-- figures/
|   |-- company.pdf
|   |-- system.png
|   `-- ...
|
`-- chapters/
    |-- chapter1.tex
    |-- chapter2.tex
    `-- chapter3.tex
\end{verbatim}

Figures should normally be stored in the \verb|figures| directory
specified by the official template.

File and directory names should preferably contain Latin letters,
numbers, hyphens, or underscores.

Recommended examples include:

\begin{verbatim}
company_structure.pdf
motor_control.png
chapter_01.tex
\end{verbatim}

Spaces and unnecessary special characters in file names should be
avoided where possible.

\vspace{0.3cm}

% ================================================================
% XXI. DOCUMENT COMPILATION AND FINAL CHECK
% ================================================================

\noindent
\textbf{XXI.\quad DOCUMENT COMPILATION AND FINAL CHECK}

\vspace{0.15cm}

Before submission, the complete \LaTeX{} project shall compile
successfully and produce the final PDF report.

The final document shall not contain unresolved references such as:

\begin{verbatim}
??
[?]
\end{verbatim}

Before generating the final PDF, students shall check:

\begin{itemize}[leftmargin=0.8cm, itemsep=2pt]
    \item the Table of Contents;
    \item page numbering;
    \item chapter, section, and subsection numbering;
    \item figure numbering;
    \item table numbering;
    \item equation numbering;
    \item cross-references;
    \item citations;
    \item bibliography;
    \item font display;
    \item mathematical symbols;
    \item image quality;
    \item completeness of all required forms.
\end{itemize}

The document may be compiled using TeX Live, MiKTeX, or another
compatible \LaTeX{} environment. Regardless of the compilation
environment, the resulting PDF shall conform to the official Faculty
template.

\vspace{0.3cm}

% ================================================================
% XXII. CONTENT THAT STUDENTS MAY EDIT
% ================================================================

\noindent
\textbf{XXII.\quad CONTENT THAT STUDENTS MAY EDIT}

\vspace{0.15cm}

Students may edit or replace:

\begin{itemize}[leftmargin=0.8cm, itemsep=2pt]
    \item student name;
    \item class;
    \item internship organization;
    \item supervisor name;
    \item Abstract;
    \item List of Symbols and Abbreviations;
    \item chapter titles and chapter content;
    \item figures and photographs;
    \item tables;
    \item equations;
    \item citations and references.
\end{itemize}

Students shall not modify the main document formatting settings in the
preamble, including the document class, page geometry, line spacing,
heading formatting, table-of-contents formatting, font configuration,
page layout, or cover-page macros unless specifically instructed.

\vspace{0.3cm}

% ================================================================
% XXIII. FINAL SUBMISSION REQUIREMENTS
% ================================================================

\noindent
\textbf{XXIII.\quad FINAL SUBMISSION REQUIREMENTS}

\vspace{0.15cm}

Before submission, the final report shall:

\begin{itemize}[leftmargin=0.8cm, itemsep=2pt]
    \item compile without errors that affect the report content or presentation;
    \item contain all required forms and report sections;
    \item contain correct student, class, internship organization, and supervisor information;
    \item contain an updated Table of Contents, List of Figures, and List of Tables;
    \item contain no unused example content from the template;
    \item contain no placeholder text such as \textit{Title of Chapter 1},
    \textit{Section 1.1}, or other sample entries;
    \item contain all required citations and references;
    \item contain no unresolved cross-references.
\end{itemize}

If the Faculty requires submission of the \LaTeX{} source, the complete
project shall include all \verb|.tex| files, the bibliography database,
figures, and any other files required to reproduce the submitted PDF.

\vspace{0.4cm}

\noindent
\textbf{Important:}
Before submitting the final internship report, students shall remove or
comment out the following command in the main document:

\begin{verbatim}
\input{Regulations.tex}
\end{verbatim}

For example:

\begin{verbatim}
% \input{Regulations.tex}
\end{verbatim}

This removes the guidance pages from the final submitted report while
preserving the regulations file in the \LaTeX{} project for reference.

} % End 13pt content

\clearpage
|.
Students shall remove or comment out this command before submitting the
final internship report.

\vspace{0.4cm}

% ================================================================
% I. GENERAL REQUIREMENTS
% ================================================================

\noindent
\textbf{I.\quad GENERAL REQUIREMENTS}

\vspace{0.15cm}

The Graduation Internship Report shall be prepared using the official
\LaTeX{} template provided by the Faculty. Students shall preserve the
predefined page layout, typography, line spacing, heading styles,
numbering system, figure and table formatting, and bibliography settings
unless otherwise instructed by the Faculty or the supervisor.

The report shall be written in English.

Students are responsible for the academic and technical content of the
report. The official \LaTeX{} template is responsible for standardized
document formatting.

\vspace{0.3cm}

% ================================================================
% II. REPORT STRUCTURE
% ================================================================

\noindent
\textbf{II.\quad REPORT STRUCTURE}

\vspace{0.15cm}

The report shall normally be arranged in the following order:

\begin{enumerate}[label=\arabic*., leftmargin=1.2cm, itemsep=2pt]
    \item Outer Cover;
    \item Comments of the Internship Host Organization;
    \item Supervisor's Evaluation and Graduation Internship Defense Results;
    \item Table of Contents;
    \item List of Symbols and Abbreviations, if applicable;
    \item List of Figures;
    \item List of Tables;
    \item Abstract;
    \item Main Chapters of the Internship Report;
    \item References;
    \item Appendices, if applicable.
\end{enumerate}

The preliminary pages shall not be numbered as chapters. The main report
content shall be organized into consecutively numbered chapters beginning
with Chapter 1.

\vspace{0.3cm}

% ================================================================
% III. PAGE SIZE AND MARGINS
% ================================================================

\noindent
\textbf{III.\quad PAGE SIZE AND MARGINS}

\vspace{0.15cm}

The report shall be formatted on A4 paper with the following margins:

\begin{center}
\begin{tabular}{ll}
    Left margin:   & 3.5 cm \\
    Right margin:  & 2.0 cm \\
    Top margin:    & 2.5 cm \\
    Bottom margin: & 2.5 cm
\end{tabular}
\end{center}

These settings are predefined in the official template and shall not be
modified by students.

\vspace{0.3cm}

% ================================================================
% IV. FONT, FONT SIZE, AND LINE SPACING
% ================================================================

\noindent
\textbf{IV.\quad FONT, FONT SIZE, AND LINE SPACING}

\vspace{0.15cm}

The main text shall use a 13 pt font size with 1.3-line spacing.
A Times-family typeface shall be used in accordance with the official
template.

The text shall be fully justified.

Students shall not use repeated spaces, repeated manual line breaks, or
multiple blank lines to position text. Paragraph and heading spacing shall
be controlled by the template.

\vspace{0.3cm}

% ================================================================
% V. OUTER COVER
% ================================================================

\noindent
\textbf{V.\quad OUTER COVER}

\vspace{0.15cm}

The outer cover shall follow the official format provided in the template
and shall contain the following information:

\begin{itemize}[leftmargin=0.8cm, itemsep=2pt]
    \item Thai Nguyen University of Technology;
    \item Faculty of International Training;
    \item Engineering and Technology Department;
    \item University logo and Faculty logo;
    \item the title \textbf{INTERNSHIP REPORT};
    \item Internship Organization;
    \item Author;
    \item Class;
    \item Supervisor;
    \item place and year of completion.
\end{itemize}

Students shall only replace the required information fields and shall not
modify the cover-page layout unless instructed otherwise.

\vspace{0.3cm}

% ================================================================
% VI. COMMENTS OF THE INTERNSHIP HOST ORGANIZATION
% ================================================================

\noindent
\textbf{VI.\quad COMMENTS OF THE INTERNSHIP HOST ORGANIZATION}

\vspace{0.15cm}

The page entitled \textbf{COMMENTS OF THE INTERNSHIP HOST ORGANIZATION}
is reserved for the internship host organization.

It includes:

\begin{itemize}[leftmargin=0.8cm, itemsep=2pt]
    \item the student's full name;
    \item class;
    \item comments from the host organization;
    \item score;
    \item date;
    \item signature and official seal.
\end{itemize}

Students shall not alter or complete official evaluation fields unless
specifically authorized.

\vspace{0.3cm}

% ================================================================
% VII. SUPERVISOR'S EVALUATION AND DEFENSE RESULTS
% ================================================================

\noindent
\textbf{VII.\quad SUPERVISOR'S EVALUATION AND DEFENSE RESULTS}

\vspace{0.15cm}

The section entitled
\textbf{SUPERVISOR'S EVALUATION AND GRADUATION INTERNSHIP DEFENSE RESULTS}
is reserved for the supervisor and defense committee.

The evaluation may include:

\begin{itemize}[leftmargin=0.8cm, itemsep=2pt]
    \item Internship Progress and Student's Attitude;
    \item Report Content;
    \item Report Presentation;
    \item Other Comments;
    \item Supervisor's Opinion;
    \item Defense Result;
    \item Score.
\end{itemize}

Students shall not enter evaluation scores or defense results before they
are officially determined.

\vspace{0.3cm}

% ================================================================
% VIII. CHAPTERS, SECTIONS, AND SUBSECTIONS
% ================================================================

\noindent
\textbf{VIII.\quad CHAPTERS, SECTIONS, AND SUBSECTIONS}

\vspace{0.15cm}

The report structure shall be generated using the standard \LaTeX{}
commands:

\begin{verbatim}
\chapter{...}
\section{...}
\subsection{...}
\end{verbatim}

Chapter, section, and subsection numbers shall be generated automatically
by \LaTeX{}.

Students shall not manually type chapter or section numbers and shall not
reproduce heading styles using commands such as \verb|\textbf|,
\verb|\fontsize|, or \verb|\centering| instead of the appropriate
structural command.

Excessive subdivision should be avoided. In most cases, the hierarchy
Chapter--Section--Subsection is sufficient.

\vspace{0.3cm}

% ================================================================
% IX. ORGANIZATION OF THE MAIN CONTENT
% ================================================================

\noindent
\textbf{IX.\quad ORGANIZATION OF THE MAIN CONTENT}

\vspace{0.15cm}

The specific content of the report shall depend on the internship
organization and the assigned tasks. A typical structure may include:

\begin{itemize}[leftmargin=0.8cm, itemsep=2pt]
    \item \textbf{Chapter 1. Overview of the Internship Organization};
    \item \textbf{Chapter 2. Internship Activities and Assigned Tasks};
    \item \textbf{Chapter 3. Results, Experience, and Evaluation}.
\end{itemize}

These chapter titles are recommendations only and may be adapted to the
actual internship content.

The report shall reflect work actually performed, observed, analyzed, or
learned during the internship. Unrelated theoretical material should not
be included merely to increase report length.

\vspace{0.3cm}

% ================================================================
% X. TABLE OF CONTENTS
% ================================================================

\noindent
\textbf{X.\quad TABLE OF CONTENTS}

\vspace{0.15cm}

The Table of Contents shall be generated automatically using:

\begin{verbatim}
\tableofcontents
\end{verbatim}

Students shall not create or manually edit page numbers in the Table of
Contents.

The document shall be compiled sufficiently to ensure that all page
numbers and entries are correctly updated.

\vspace{0.3cm}

% ================================================================
% XI. LIST OF SYMBOLS AND ABBREVIATIONS
% ================================================================

\noindent
\textbf{XI.\quad LIST OF SYMBOLS AND ABBREVIATIONS}

\vspace{0.15cm}

Symbols and abbreviations that are used repeatedly in the report shall be
listed in the \textbf{List of Symbols and Abbreviations}.

The list may contain the following columns:

\begin{center}
\begin{tabular}{llll}
    No. & Symbol / Abbreviation & Unit & Definition
\end{tabular}
\end{center}

Only symbols and abbreviations actually used in the final report shall be
retained. Example entries contained in the template shall be removed if
they are not relevant to the student's report.

Each abbreviation should normally be defined when it first appears in
the main text.

\vspace{0.3cm}

% ================================================================
% XII. FIGURES, DIAGRAMS, AND PHOTOGRAPHS
% ================================================================

\noindent
\textbf{XII.\quad FIGURES, DIAGRAMS, AND PHOTOGRAPHS}

\vspace{0.15cm}

Figures shall be clear, legible, and directly relevant to the report
content.

Figures shall be inserted using the standard \verb|figure| environment.

Example:

\begin{verbatim}
\begin{figure}[H]
    \centering
    \includegraphics[width=0.8\textwidth]
                    {figures/company_structure.pdf}
    \caption{Organizational structure of the company}
    \label{fig:company-structure}
\end{figure}
\end{verbatim}

Figure captions shall be placed below the figure.

Each referenced figure shall have a unique \verb|\label| and shall be
referenced in the text using, for example:

\begin{verbatim}
Figure~\ref{fig:company-structure}
\end{verbatim}

Students shall not manually type figure numbers.

Figures reproduced or adapted from external sources shall include an
appropriate citation.

Photographs taken during the internship shall be relevant and sufficiently
clear. Students shall not include confidential, restricted, or proprietary
information without permission from the internship organization.

Vector graphics or high-resolution images are recommended whenever
possible.

\vspace{0.3cm}

% ================================================================
% XIII. TABLES
% ================================================================

\noindent
\textbf{XIII.\quad TABLES}

\vspace{0.15cm}

Tables shall be created using suitable \LaTeX{} environments such as
\verb|table|, \verb|tabular|, \verb|tabularx|, or \verb|longtable|.

Table captions shall be placed above the table.

Example:

\begin{verbatim}
\begin{table}[H]
    \centering
    \caption{Internship schedule}
    \label{tab:internship-schedule}

    \begin{tabular}{lll}
        \toprule
        Week & Task & Result \\
        \midrule
        1 & Orientation & Completed \\
        2 & System study & Completed \\
        \bottomrule
    \end{tabular}
\end{table}
\end{verbatim}

Tables shall be referenced using:

\begin{verbatim}
Table~\ref{tab:internship-schedule}
\end{verbatim}

Students shall not manually type table numbers.

Tables should be created directly in \LaTeX{} rather than inserted as
screenshots whenever practical.

\vspace{0.3cm}

% ================================================================
% XIV. EQUATIONS AND TECHNICAL QUANTITIES
% ================================================================

\noindent
\textbf{XIV.\quad EQUATIONS AND TECHNICAL QUANTITIES}

\vspace{0.15cm}

Mathematical expressions shall be typeset directly using \LaTeX{}.
Images or screenshots of equations shall not be used.

Short expressions may be written in inline mathematical mode, for example:

\begin{verbatim}
$R = 0.7\,\Omega$
\end{verbatim}

Important equations shall use a numbered equation environment.

Example:

\begin{verbatim}
\begin{equation}
    P = UI
    \label{eq:power}
\end{equation}
\end{verbatim}

Equations shall be referenced using:

\begin{verbatim}
Equation~\eqref{eq:power}
\end{verbatim}

Equation numbers shall not be manually entered.

SI units shall be used wherever applicable. The \verb|siunitx| package
should be used where appropriate.

Examples:

\begin{verbatim}
\SI{220}{\volt}
\SI{50}{\hertz}
\SI{2.5}{\ampere}
\end{verbatim}

Notation and units shall be consistent throughout the report.

\vspace{0.3cm}

% ================================================================
% XV. SOURCE CODE
% ================================================================

\noindent
\textbf{XV.\quad SOURCE CODE}

\vspace{0.15cm}

If the internship involves programming, PLCs, microcontrollers,
simulation, or software development, relevant source code may be included
in the report.

Source code shall be presented using the \verb|listings| environment or
another environment defined by the official template.

Example:

\begin{verbatim}
\begin{lstlisting}[language=st]
IF Start = TRUE THEN
    Motor := TRUE;
ELSE
    Motor := FALSE;
END_IF;
\end{lstlisting}
\end{verbatim}

Only code necessary for explaining the work should appear in the main
report. Long source-code listings should be moved to an appendix or
submitted as supplementary files if required.

Source code should not normally be inserted as screenshots.

\vspace{0.3cm}

% ================================================================
% XVI. ABSTRACT
% ================================================================

\noindent
\textbf{XVI.\quad ABSTRACT}

\vspace{0.15cm}

The Abstract shall briefly summarize:

\begin{itemize}[leftmargin=0.8cm, itemsep=2pt]
    \item the internship organization;
    \item the objectives of the internship;
    \item the main tasks performed;
    \item the main results, knowledge, and skills obtained.
\end{itemize}

The Abstract should be concise and should not reproduce the full content
of the report chapter by chapter.

\vspace{0.3cm}

% ================================================================
% XVII. REFERENCES AND CITATIONS
% ================================================================

\noindent
\textbf{XVII.\quad REFERENCES AND CITATIONS}

\vspace{0.15cm}

All books, journal papers, conference papers, standards, websites,
datasheets, manuals, technical documents, figures, tables, or other
external sources used in preparing the report shall be properly cited.

References shall be managed through the bibliography system provided with
the official template.

Bibliographic entries shall normally be stored in:

\begin{verbatim}
References.bib
\end{verbatim}

References shall be cited in the report using:

\begin{verbatim}
\cite{reference-key}
\end{verbatim}

Students shall not manually type reference numbers such as ``[1]'',
``[2]'', etc.

Titles of foreign-language publications shall remain in their original
language in the bibliography.

All borrowed ideas, data, figures, tables, equations, methods, or results
shall be properly acknowledged.

\vspace{0.3cm}

% ================================================================
% XVIII. CROSS-REFERENCING
% ================================================================

\noindent
\textbf{XVIII.\quad CROSS-REFERENCING}

\vspace{0.15cm}

Numbered elements shall be cross-referenced using the standard \LaTeX{}
mechanism.

Use the following commands where appropriate:

\begin{verbatim}
\label{...}
\ref{...}
\eqref{...}
\end{verbatim}

For example:

\begin{verbatim}
As shown in Figure~\ref{fig:system}, ...
\end{verbatim}

Students shall not manually type numbers such as ``Figure 2.3'' when
those numbers can be generated automatically.

\vspace{0.3cm}

% ================================================================
% XIX. PAGE NUMBERING
% ================================================================

\noindent
\textbf{XIX.\quad PAGE NUMBERING}

\vspace{0.15cm}

Cover pages and official evaluation forms at the beginning of the report
shall not display page numbers.

The main body of the report shall use Arabic page numbers beginning from
Chapter 1.

Page numbering shall be generated automatically by the official template.
Students shall not manually enter page numbers.

\vspace{0.3cm}

% ================================================================
% XX. PROJECT FILE ORGANIZATION
% ================================================================

\noindent
\textbf{XX.\quad PROJECT FILE ORGANIZATION}

\vspace{0.15cm}

Students should maintain a clear and organized \LaTeX{} project
structure. A recommended structure is:

\begin{verbatim}
InternshipReport/
|
|-- main.tex
|-- References.bib
|
|-- figures/
|   |-- company.pdf
|   |-- system.png
|   `-- ...
|
`-- chapters/
    |-- chapter1.tex
    |-- chapter2.tex
    `-- chapter3.tex
\end{verbatim}

Figures should normally be stored in the \verb|figures| directory
specified by the official template.

File and directory names should preferably contain Latin letters,
numbers, hyphens, or underscores.

Recommended examples include:

\begin{verbatim}
company_structure.pdf
motor_control.png
chapter_01.tex
\end{verbatim}

Spaces and unnecessary special characters in file names should be
avoided where possible.

\vspace{0.3cm}

% ================================================================
% XXI. DOCUMENT COMPILATION AND FINAL CHECK
% ================================================================

\noindent
\textbf{XXI.\quad DOCUMENT COMPILATION AND FINAL CHECK}

\vspace{0.15cm}

Before submission, the complete \LaTeX{} project shall compile
successfully and produce the final PDF report.

The final document shall not contain unresolved references such as:

\begin{verbatim}
??
[?]
\end{verbatim}

Before generating the final PDF, students shall check:

\begin{itemize}[leftmargin=0.8cm, itemsep=2pt]
    \item the Table of Contents;
    \item page numbering;
    \item chapter, section, and subsection numbering;
    \item figure numbering;
    \item table numbering;
    \item equation numbering;
    \item cross-references;
    \item citations;
    \item bibliography;
    \item font display;
    \item mathematical symbols;
    \item image quality;
    \item completeness of all required forms.
\end{itemize}

The document may be compiled using TeX Live, MiKTeX, or another
compatible \LaTeX{} environment. Regardless of the compilation
environment, the resulting PDF shall conform to the official Faculty
template.

\vspace{0.3cm}

% ================================================================
% XXII. CONTENT THAT STUDENTS MAY EDIT
% ================================================================

\noindent
\textbf{XXII.\quad CONTENT THAT STUDENTS MAY EDIT}

\vspace{0.15cm}

Students may edit or replace:

\begin{itemize}[leftmargin=0.8cm, itemsep=2pt]
    \item student name;
    \item class;
    \item internship organization;
    \item supervisor name;
    \item Abstract;
    \item List of Symbols and Abbreviations;
    \item chapter titles and chapter content;
    \item figures and photographs;
    \item tables;
    \item equations;
    \item citations and references.
\end{itemize}

Students shall not modify the main document formatting settings in the
preamble, including the document class, page geometry, line spacing,
heading formatting, table-of-contents formatting, font configuration,
page layout, or cover-page macros unless specifically instructed.

\vspace{0.3cm}

% ================================================================
% XXIII. FINAL SUBMISSION REQUIREMENTS
% ================================================================

\noindent
\textbf{XXIII.\quad FINAL SUBMISSION REQUIREMENTS}

\vspace{0.15cm}

Before submission, the final report shall:

\begin{itemize}[leftmargin=0.8cm, itemsep=2pt]
    \item compile without errors that affect the report content or presentation;
    \item contain all required forms and report sections;
    \item contain correct student, class, internship organization, and supervisor information;
    \item contain an updated Table of Contents, List of Figures, and List of Tables;
    \item contain no unused example content from the template;
    \item contain no placeholder text such as \textit{Title of Chapter 1},
    \textit{Section 1.1}, or other sample entries;
    \item contain all required citations and references;
    \item contain no unresolved cross-references.
\end{itemize}

If the Faculty requires submission of the \LaTeX{} source, the complete
project shall include all \verb|.tex| files, the bibliography database,
figures, and any other files required to reproduce the submitted PDF.

\vspace{0.4cm}

\noindent
\textbf{Important:}
Before submitting the final internship report, students shall remove or
comment out the following command in the main document:

\begin{verbatim}

% =================================================================
% REGULATIONS FOR PREPARING THE GRADUATION INTERNSHIP REPORT
% USING LATEX
% =================================================================
%
% This file is intended to be included from the main document with:
%
%     \input{Regulations.tex}
%
% IMPORTANT:
% This section is provided for guidance only. Students shall remove
% or comment out the \input{Regulations.tex} command before submitting
% the final internship report.
%
% =================================================================

\clearpage
\thispagestyle{empty}

\begin{center}
    {\fontsize{16pt}{20pt}\selectfont\bfseries
    SPECIFIC REGULATIONS FOR PREPARING\\
    GRADUATION INTERNSHIP REPORTS USING \LaTeX}
\end{center}

\vspace{0.5cm}

{\fontsize{13pt}{17pt}\selectfont
\setstretch{1.3}
\justifying

\noindent
\textbf{Note:}
This section is provided for guidance only and is included in the main
document using the command \verb|\input{Regulations.tex}|.
Students shall remove or comment out this command before submitting the
final internship report.

\vspace{0.4cm}

% ================================================================
% I. GENERAL REQUIREMENTS
% ================================================================

\noindent
\textbf{I.\quad GENERAL REQUIREMENTS}

\vspace{0.15cm}

The Graduation Internship Report shall be prepared using the official
\LaTeX{} template provided by the Faculty. Students shall preserve the
predefined page layout, typography, line spacing, heading styles,
numbering system, figure and table formatting, and bibliography settings
unless otherwise instructed by the Faculty or the supervisor.

The report shall be written in English.

Students are responsible for the academic and technical content of the
report. The official \LaTeX{} template is responsible for standardized
document formatting.

\vspace{0.3cm}

% ================================================================
% II. REPORT STRUCTURE
% ================================================================

\noindent
\textbf{II.\quad REPORT STRUCTURE}

\vspace{0.15cm}

The report shall normally be arranged in the following order:

\begin{enumerate}[label=\arabic*., leftmargin=1.2cm, itemsep=2pt]
    \item Outer Cover;
    \item Comments of the Internship Host Organization;
    \item Supervisor's Evaluation and Graduation Internship Defense Results;
    \item Table of Contents;
    \item List of Symbols and Abbreviations, if applicable;
    \item List of Figures;
    \item List of Tables;
    \item Abstract;
    \item Main Chapters of the Internship Report;
    \item References;
    \item Appendices, if applicable.
\end{enumerate}

The preliminary pages shall not be numbered as chapters. The main report
content shall be organized into consecutively numbered chapters beginning
with Chapter 1.

\vspace{0.3cm}

% ================================================================
% III. PAGE SIZE AND MARGINS
% ================================================================

\noindent
\textbf{III.\quad PAGE SIZE AND MARGINS}

\vspace{0.15cm}

The report shall be formatted on A4 paper with the following margins:

\begin{center}
\begin{tabular}{ll}
    Left margin:   & 3.5 cm \\
    Right margin:  & 2.0 cm \\
    Top margin:    & 2.5 cm \\
    Bottom margin: & 2.5 cm
\end{tabular}
\end{center}

These settings are predefined in the official template and shall not be
modified by students.

\vspace{0.3cm}

% ================================================================
% IV. FONT, FONT SIZE, AND LINE SPACING
% ================================================================

\noindent
\textbf{IV.\quad FONT, FONT SIZE, AND LINE SPACING}

\vspace{0.15cm}

The main text shall use a 13 pt font size with 1.3-line spacing.
A Times-family typeface shall be used in accordance with the official
template.

The text shall be fully justified.

Students shall not use repeated spaces, repeated manual line breaks, or
multiple blank lines to position text. Paragraph and heading spacing shall
be controlled by the template.

\vspace{0.3cm}

% ================================================================
% V. OUTER COVER
% ================================================================

\noindent
\textbf{V.\quad OUTER COVER}

\vspace{0.15cm}

The outer cover shall follow the official format provided in the template
and shall contain the following information:

\begin{itemize}[leftmargin=0.8cm, itemsep=2pt]
    \item Thai Nguyen University of Technology;
    \item Faculty of International Training;
    \item Engineering and Technology Department;
    \item University logo and Faculty logo;
    \item the title \textbf{INTERNSHIP REPORT};
    \item Internship Organization;
    \item Author;
    \item Class;
    \item Supervisor;
    \item place and year of completion.
\end{itemize}

Students shall only replace the required information fields and shall not
modify the cover-page layout unless instructed otherwise.

\vspace{0.3cm}

% ================================================================
% VI. COMMENTS OF THE INTERNSHIP HOST ORGANIZATION
% ================================================================

\noindent
\textbf{VI.\quad COMMENTS OF THE INTERNSHIP HOST ORGANIZATION}

\vspace{0.15cm}

The page entitled \textbf{COMMENTS OF THE INTERNSHIP HOST ORGANIZATION}
is reserved for the internship host organization.

It includes:

\begin{itemize}[leftmargin=0.8cm, itemsep=2pt]
    \item the student's full name;
    \item class;
    \item comments from the host organization;
    \item score;
    \item date;
    \item signature and official seal.
\end{itemize}

Students shall not alter or complete official evaluation fields unless
specifically authorized.

\vspace{0.3cm}

% ================================================================
% VII. SUPERVISOR'S EVALUATION AND DEFENSE RESULTS
% ================================================================

\noindent
\textbf{VII.\quad SUPERVISOR'S EVALUATION AND DEFENSE RESULTS}

\vspace{0.15cm}

The section entitled
\textbf{SUPERVISOR'S EVALUATION AND GRADUATION INTERNSHIP DEFENSE RESULTS}
is reserved for the supervisor and defense committee.

The evaluation may include:

\begin{itemize}[leftmargin=0.8cm, itemsep=2pt]
    \item Internship Progress and Student's Attitude;
    \item Report Content;
    \item Report Presentation;
    \item Other Comments;
    \item Supervisor's Opinion;
    \item Defense Result;
    \item Score.
\end{itemize}

Students shall not enter evaluation scores or defense results before they
are officially determined.

\vspace{0.3cm}

% ================================================================
% VIII. CHAPTERS, SECTIONS, AND SUBSECTIONS
% ================================================================

\noindent
\textbf{VIII.\quad CHAPTERS, SECTIONS, AND SUBSECTIONS}

\vspace{0.15cm}

The report structure shall be generated using the standard \LaTeX{}
commands:

\begin{verbatim}
\chapter{...}
\section{...}
\subsection{...}
\end{verbatim}

Chapter, section, and subsection numbers shall be generated automatically
by \LaTeX{}.

Students shall not manually type chapter or section numbers and shall not
reproduce heading styles using commands such as \verb|\textbf|,
\verb|\fontsize|, or \verb|\centering| instead of the appropriate
structural command.

Excessive subdivision should be avoided. In most cases, the hierarchy
Chapter--Section--Subsection is sufficient.

\vspace{0.3cm}

% ================================================================
% IX. ORGANIZATION OF THE MAIN CONTENT
% ================================================================

\noindent
\textbf{IX.\quad ORGANIZATION OF THE MAIN CONTENT}

\vspace{0.15cm}

The specific content of the report shall depend on the internship
organization and the assigned tasks. A typical structure may include:

\begin{itemize}[leftmargin=0.8cm, itemsep=2pt]
    \item \textbf{Chapter 1. Overview of the Internship Organization};
    \item \textbf{Chapter 2. Internship Activities and Assigned Tasks};
    \item \textbf{Chapter 3. Results, Experience, and Evaluation}.
\end{itemize}

These chapter titles are recommendations only and may be adapted to the
actual internship content.

The report shall reflect work actually performed, observed, analyzed, or
learned during the internship. Unrelated theoretical material should not
be included merely to increase report length.

\vspace{0.3cm}

% ================================================================
% X. TABLE OF CONTENTS
% ================================================================

\noindent
\textbf{X.\quad TABLE OF CONTENTS}

\vspace{0.15cm}

The Table of Contents shall be generated automatically using:

\begin{verbatim}
\tableofcontents
\end{verbatim}

Students shall not create or manually edit page numbers in the Table of
Contents.

The document shall be compiled sufficiently to ensure that all page
numbers and entries are correctly updated.

\vspace{0.3cm}

% ================================================================
% XI. LIST OF SYMBOLS AND ABBREVIATIONS
% ================================================================

\noindent
\textbf{XI.\quad LIST OF SYMBOLS AND ABBREVIATIONS}

\vspace{0.15cm}

Symbols and abbreviations that are used repeatedly in the report shall be
listed in the \textbf{List of Symbols and Abbreviations}.

The list may contain the following columns:

\begin{center}
\begin{tabular}{llll}
    No. & Symbol / Abbreviation & Unit & Definition
\end{tabular}
\end{center}

Only symbols and abbreviations actually used in the final report shall be
retained. Example entries contained in the template shall be removed if
they are not relevant to the student's report.

Each abbreviation should normally be defined when it first appears in
the main text.

\vspace{0.3cm}

% ================================================================
% XII. FIGURES, DIAGRAMS, AND PHOTOGRAPHS
% ================================================================

\noindent
\textbf{XII.\quad FIGURES, DIAGRAMS, AND PHOTOGRAPHS}

\vspace{0.15cm}

Figures shall be clear, legible, and directly relevant to the report
content.

Figures shall be inserted using the standard \verb|figure| environment.

Example:

\begin{verbatim}
\begin{figure}[H]
    \centering
    \includegraphics[width=0.8\textwidth]
                    {figures/company_structure.pdf}
    \caption{Organizational structure of the company}
    \label{fig:company-structure}
\end{figure}
\end{verbatim}

Figure captions shall be placed below the figure.

Each referenced figure shall have a unique \verb|\label| and shall be
referenced in the text using, for example:

\begin{verbatim}
Figure~\ref{fig:company-structure}
\end{verbatim}

Students shall not manually type figure numbers.

Figures reproduced or adapted from external sources shall include an
appropriate citation.

Photographs taken during the internship shall be relevant and sufficiently
clear. Students shall not include confidential, restricted, or proprietary
information without permission from the internship organization.

Vector graphics or high-resolution images are recommended whenever
possible.

\vspace{0.3cm}

% ================================================================
% XIII. TABLES
% ================================================================

\noindent
\textbf{XIII.\quad TABLES}

\vspace{0.15cm}

Tables shall be created using suitable \LaTeX{} environments such as
\verb|table|, \verb|tabular|, \verb|tabularx|, or \verb|longtable|.

Table captions shall be placed above the table.

Example:

\begin{verbatim}
\begin{table}[H]
    \centering
    \caption{Internship schedule}
    \label{tab:internship-schedule}

    \begin{tabular}{lll}
        \toprule
        Week & Task & Result \\
        \midrule
        1 & Orientation & Completed \\
        2 & System study & Completed \\
        \bottomrule
    \end{tabular}
\end{table}
\end{verbatim}

Tables shall be referenced using:

\begin{verbatim}
Table~\ref{tab:internship-schedule}
\end{verbatim}

Students shall not manually type table numbers.

Tables should be created directly in \LaTeX{} rather than inserted as
screenshots whenever practical.

\vspace{0.3cm}

% ================================================================
% XIV. EQUATIONS AND TECHNICAL QUANTITIES
% ================================================================

\noindent
\textbf{XIV.\quad EQUATIONS AND TECHNICAL QUANTITIES}

\vspace{0.15cm}

Mathematical expressions shall be typeset directly using \LaTeX{}.
Images or screenshots of equations shall not be used.

Short expressions may be written in inline mathematical mode, for example:

\begin{verbatim}
$R = 0.7\,\Omega$
\end{verbatim}

Important equations shall use a numbered equation environment.

Example:

\begin{verbatim}
\begin{equation}
    P = UI
    \label{eq:power}
\end{equation}
\end{verbatim}

Equations shall be referenced using:

\begin{verbatim}
Equation~\eqref{eq:power}
\end{verbatim}

Equation numbers shall not be manually entered.

SI units shall be used wherever applicable. The \verb|siunitx| package
should be used where appropriate.

Examples:

\begin{verbatim}
\SI{220}{\volt}
\SI{50}{\hertz}
\SI{2.5}{\ampere}
\end{verbatim}

Notation and units shall be consistent throughout the report.

\vspace{0.3cm}

% ================================================================
% XV. SOURCE CODE
% ================================================================

\noindent
\textbf{XV.\quad SOURCE CODE}

\vspace{0.15cm}

If the internship involves programming, PLCs, microcontrollers,
simulation, or software development, relevant source code may be included
in the report.

Source code shall be presented using the \verb|listings| environment or
another environment defined by the official template.

Example:

\begin{verbatim}
\begin{lstlisting}[language=st]
IF Start = TRUE THEN
    Motor := TRUE;
ELSE
    Motor := FALSE;
END_IF;
\end{lstlisting}
\end{verbatim}

Only code necessary for explaining the work should appear in the main
report. Long source-code listings should be moved to an appendix or
submitted as supplementary files if required.

Source code should not normally be inserted as screenshots.

\vspace{0.3cm}

% ================================================================
% XVI. ABSTRACT
% ================================================================

\noindent
\textbf{XVI.\quad ABSTRACT}

\vspace{0.15cm}

The Abstract shall briefly summarize:

\begin{itemize}[leftmargin=0.8cm, itemsep=2pt]
    \item the internship organization;
    \item the objectives of the internship;
    \item the main tasks performed;
    \item the main results, knowledge, and skills obtained.
\end{itemize}

The Abstract should be concise and should not reproduce the full content
of the report chapter by chapter.

\vspace{0.3cm}

% ================================================================
% XVII. REFERENCES AND CITATIONS
% ================================================================

\noindent
\textbf{XVII.\quad REFERENCES AND CITATIONS}

\vspace{0.15cm}

All books, journal papers, conference papers, standards, websites,
datasheets, manuals, technical documents, figures, tables, or other
external sources used in preparing the report shall be properly cited.

References shall be managed through the bibliography system provided with
the official template.

Bibliographic entries shall normally be stored in:

\begin{verbatim}
References.bib
\end{verbatim}

References shall be cited in the report using:

\begin{verbatim}
\cite{reference-key}
\end{verbatim}

Students shall not manually type reference numbers such as ``[1]'',
``[2]'', etc.

Titles of foreign-language publications shall remain in their original
language in the bibliography.

All borrowed ideas, data, figures, tables, equations, methods, or results
shall be properly acknowledged.

\vspace{0.3cm}

% ================================================================
% XVIII. CROSS-REFERENCING
% ================================================================

\noindent
\textbf{XVIII.\quad CROSS-REFERENCING}

\vspace{0.15cm}

Numbered elements shall be cross-referenced using the standard \LaTeX{}
mechanism.

Use the following commands where appropriate:

\begin{verbatim}
\label{...}
\ref{...}
\eqref{...}
\end{verbatim}

For example:

\begin{verbatim}
As shown in Figure~\ref{fig:system}, ...
\end{verbatim}

Students shall not manually type numbers such as ``Figure 2.3'' when
those numbers can be generated automatically.

\vspace{0.3cm}

% ================================================================
% XIX. PAGE NUMBERING
% ================================================================

\noindent
\textbf{XIX.\quad PAGE NUMBERING}

\vspace{0.15cm}

Cover pages and official evaluation forms at the beginning of the report
shall not display page numbers.

The main body of the report shall use Arabic page numbers beginning from
Chapter 1.

Page numbering shall be generated automatically by the official template.
Students shall not manually enter page numbers.

\vspace{0.3cm}

% ================================================================
% XX. PROJECT FILE ORGANIZATION
% ================================================================

\noindent
\textbf{XX.\quad PROJECT FILE ORGANIZATION}

\vspace{0.15cm}

Students should maintain a clear and organized \LaTeX{} project
structure. A recommended structure is:

\begin{verbatim}
InternshipReport/
|
|-- main.tex
|-- References.bib
|
|-- figures/
|   |-- company.pdf
|   |-- system.png
|   `-- ...
|
`-- chapters/
    |-- chapter1.tex
    |-- chapter2.tex
    `-- chapter3.tex
\end{verbatim}

Figures should normally be stored in the \verb|figures| directory
specified by the official template.

File and directory names should preferably contain Latin letters,
numbers, hyphens, or underscores.

Recommended examples include:

\begin{verbatim}
company_structure.pdf
motor_control.png
chapter_01.tex
\end{verbatim}

Spaces and unnecessary special characters in file names should be
avoided where possible.

\vspace{0.3cm}

% ================================================================
% XXI. DOCUMENT COMPILATION AND FINAL CHECK
% ================================================================

\noindent
\textbf{XXI.\quad DOCUMENT COMPILATION AND FINAL CHECK}

\vspace{0.15cm}

Before submission, the complete \LaTeX{} project shall compile
successfully and produce the final PDF report.

The final document shall not contain unresolved references such as:

\begin{verbatim}
??
[?]
\end{verbatim}

Before generating the final PDF, students shall check:

\begin{itemize}[leftmargin=0.8cm, itemsep=2pt]
    \item the Table of Contents;
    \item page numbering;
    \item chapter, section, and subsection numbering;
    \item figure numbering;
    \item table numbering;
    \item equation numbering;
    \item cross-references;
    \item citations;
    \item bibliography;
    \item font display;
    \item mathematical symbols;
    \item image quality;
    \item completeness of all required forms.
\end{itemize}

The document may be compiled using TeX Live, MiKTeX, or another
compatible \LaTeX{} environment. Regardless of the compilation
environment, the resulting PDF shall conform to the official Faculty
template.

\vspace{0.3cm}

% ================================================================
% XXII. CONTENT THAT STUDENTS MAY EDIT
% ================================================================

\noindent
\textbf{XXII.\quad CONTENT THAT STUDENTS MAY EDIT}

\vspace{0.15cm}

Students may edit or replace:

\begin{itemize}[leftmargin=0.8cm, itemsep=2pt]
    \item student name;
    \item class;
    \item internship organization;
    \item supervisor name;
    \item Abstract;
    \item List of Symbols and Abbreviations;
    \item chapter titles and chapter content;
    \item figures and photographs;
    \item tables;
    \item equations;
    \item citations and references.
\end{itemize}

Students shall not modify the main document formatting settings in the
preamble, including the document class, page geometry, line spacing,
heading formatting, table-of-contents formatting, font configuration,
page layout, or cover-page macros unless specifically instructed.

\vspace{0.3cm}

% ================================================================
% XXIII. FINAL SUBMISSION REQUIREMENTS
% ================================================================

\noindent
\textbf{XXIII.\quad FINAL SUBMISSION REQUIREMENTS}

\vspace{0.15cm}

Before submission, the final report shall:

\begin{itemize}[leftmargin=0.8cm, itemsep=2pt]
    \item compile without errors that affect the report content or presentation;
    \item contain all required forms and report sections;
    \item contain correct student, class, internship organization, and supervisor information;
    \item contain an updated Table of Contents, List of Figures, and List of Tables;
    \item contain no unused example content from the template;
    \item contain no placeholder text such as \textit{Title of Chapter 1},
    \textit{Section 1.1}, or other sample entries;
    \item contain all required citations and references;
    \item contain no unresolved cross-references.
\end{itemize}

If the Faculty requires submission of the \LaTeX{} source, the complete
project shall include all \verb|.tex| files, the bibliography database,
figures, and any other files required to reproduce the submitted PDF.

\vspace{0.4cm}

\noindent
\textbf{Important:}
Before submitting the final internship report, students shall remove or
comment out the following command in the main document:

\begin{verbatim}
\input{Regulations.tex}
\end{verbatim}

For example:

\begin{verbatim}
% \input{Regulations.tex}
\end{verbatim}

This removes the guidance pages from the final submitted report while
preserving the regulations file in the \LaTeX{} project for reference.

} % End 13pt content

\clearpage

\end{verbatim}

For example:

\begin{verbatim}
% 
% =================================================================
% REGULATIONS FOR PREPARING THE GRADUATION INTERNSHIP REPORT
% USING LATEX
% =================================================================
%
% This file is intended to be included from the main document with:
%
%     \input{Regulations.tex}
%
% IMPORTANT:
% This section is provided for guidance only. Students shall remove
% or comment out the \input{Regulations.tex} command before submitting
% the final internship report.
%
% =================================================================

\clearpage
\thispagestyle{empty}

\begin{center}
    {\fontsize{16pt}{20pt}\selectfont\bfseries
    SPECIFIC REGULATIONS FOR PREPARING\\
    GRADUATION INTERNSHIP REPORTS USING \LaTeX}
\end{center}

\vspace{0.5cm}

{\fontsize{13pt}{17pt}\selectfont
\setstretch{1.3}
\justifying

\noindent
\textbf{Note:}
This section is provided for guidance only and is included in the main
document using the command \verb|\input{Regulations.tex}|.
Students shall remove or comment out this command before submitting the
final internship report.

\vspace{0.4cm}

% ================================================================
% I. GENERAL REQUIREMENTS
% ================================================================

\noindent
\textbf{I.\quad GENERAL REQUIREMENTS}

\vspace{0.15cm}

The Graduation Internship Report shall be prepared using the official
\LaTeX{} template provided by the Faculty. Students shall preserve the
predefined page layout, typography, line spacing, heading styles,
numbering system, figure and table formatting, and bibliography settings
unless otherwise instructed by the Faculty or the supervisor.

The report shall be written in English.

Students are responsible for the academic and technical content of the
report. The official \LaTeX{} template is responsible for standardized
document formatting.

\vspace{0.3cm}

% ================================================================
% II. REPORT STRUCTURE
% ================================================================

\noindent
\textbf{II.\quad REPORT STRUCTURE}

\vspace{0.15cm}

The report shall normally be arranged in the following order:

\begin{enumerate}[label=\arabic*., leftmargin=1.2cm, itemsep=2pt]
    \item Outer Cover;
    \item Comments of the Internship Host Organization;
    \item Supervisor's Evaluation and Graduation Internship Defense Results;
    \item Table of Contents;
    \item List of Symbols and Abbreviations, if applicable;
    \item List of Figures;
    \item List of Tables;
    \item Abstract;
    \item Main Chapters of the Internship Report;
    \item References;
    \item Appendices, if applicable.
\end{enumerate}

The preliminary pages shall not be numbered as chapters. The main report
content shall be organized into consecutively numbered chapters beginning
with Chapter 1.

\vspace{0.3cm}

% ================================================================
% III. PAGE SIZE AND MARGINS
% ================================================================

\noindent
\textbf{III.\quad PAGE SIZE AND MARGINS}

\vspace{0.15cm}

The report shall be formatted on A4 paper with the following margins:

\begin{center}
\begin{tabular}{ll}
    Left margin:   & 3.5 cm \\
    Right margin:  & 2.0 cm \\
    Top margin:    & 2.5 cm \\
    Bottom margin: & 2.5 cm
\end{tabular}
\end{center}

These settings are predefined in the official template and shall not be
modified by students.

\vspace{0.3cm}

% ================================================================
% IV. FONT, FONT SIZE, AND LINE SPACING
% ================================================================

\noindent
\textbf{IV.\quad FONT, FONT SIZE, AND LINE SPACING}

\vspace{0.15cm}

The main text shall use a 13 pt font size with 1.3-line spacing.
A Times-family typeface shall be used in accordance with the official
template.

The text shall be fully justified.

Students shall not use repeated spaces, repeated manual line breaks, or
multiple blank lines to position text. Paragraph and heading spacing shall
be controlled by the template.

\vspace{0.3cm}

% ================================================================
% V. OUTER COVER
% ================================================================

\noindent
\textbf{V.\quad OUTER COVER}

\vspace{0.15cm}

The outer cover shall follow the official format provided in the template
and shall contain the following information:

\begin{itemize}[leftmargin=0.8cm, itemsep=2pt]
    \item Thai Nguyen University of Technology;
    \item Faculty of International Training;
    \item Engineering and Technology Department;
    \item University logo and Faculty logo;
    \item the title \textbf{INTERNSHIP REPORT};
    \item Internship Organization;
    \item Author;
    \item Class;
    \item Supervisor;
    \item place and year of completion.
\end{itemize}

Students shall only replace the required information fields and shall not
modify the cover-page layout unless instructed otherwise.

\vspace{0.3cm}

% ================================================================
% VI. COMMENTS OF THE INTERNSHIP HOST ORGANIZATION
% ================================================================

\noindent
\textbf{VI.\quad COMMENTS OF THE INTERNSHIP HOST ORGANIZATION}

\vspace{0.15cm}

The page entitled \textbf{COMMENTS OF THE INTERNSHIP HOST ORGANIZATION}
is reserved for the internship host organization.

It includes:

\begin{itemize}[leftmargin=0.8cm, itemsep=2pt]
    \item the student's full name;
    \item class;
    \item comments from the host organization;
    \item score;
    \item date;
    \item signature and official seal.
\end{itemize}

Students shall not alter or complete official evaluation fields unless
specifically authorized.

\vspace{0.3cm}

% ================================================================
% VII. SUPERVISOR'S EVALUATION AND DEFENSE RESULTS
% ================================================================

\noindent
\textbf{VII.\quad SUPERVISOR'S EVALUATION AND DEFENSE RESULTS}

\vspace{0.15cm}

The section entitled
\textbf{SUPERVISOR'S EVALUATION AND GRADUATION INTERNSHIP DEFENSE RESULTS}
is reserved for the supervisor and defense committee.

The evaluation may include:

\begin{itemize}[leftmargin=0.8cm, itemsep=2pt]
    \item Internship Progress and Student's Attitude;
    \item Report Content;
    \item Report Presentation;
    \item Other Comments;
    \item Supervisor's Opinion;
    \item Defense Result;
    \item Score.
\end{itemize}

Students shall not enter evaluation scores or defense results before they
are officially determined.

\vspace{0.3cm}

% ================================================================
% VIII. CHAPTERS, SECTIONS, AND SUBSECTIONS
% ================================================================

\noindent
\textbf{VIII.\quad CHAPTERS, SECTIONS, AND SUBSECTIONS}

\vspace{0.15cm}

The report structure shall be generated using the standard \LaTeX{}
commands:

\begin{verbatim}
\chapter{...}
\section{...}
\subsection{...}
\end{verbatim}

Chapter, section, and subsection numbers shall be generated automatically
by \LaTeX{}.

Students shall not manually type chapter or section numbers and shall not
reproduce heading styles using commands such as \verb|\textbf|,
\verb|\fontsize|, or \verb|\centering| instead of the appropriate
structural command.

Excessive subdivision should be avoided. In most cases, the hierarchy
Chapter--Section--Subsection is sufficient.

\vspace{0.3cm}

% ================================================================
% IX. ORGANIZATION OF THE MAIN CONTENT
% ================================================================

\noindent
\textbf{IX.\quad ORGANIZATION OF THE MAIN CONTENT}

\vspace{0.15cm}

The specific content of the report shall depend on the internship
organization and the assigned tasks. A typical structure may include:

\begin{itemize}[leftmargin=0.8cm, itemsep=2pt]
    \item \textbf{Chapter 1. Overview of the Internship Organization};
    \item \textbf{Chapter 2. Internship Activities and Assigned Tasks};
    \item \textbf{Chapter 3. Results, Experience, and Evaluation}.
\end{itemize}

These chapter titles are recommendations only and may be adapted to the
actual internship content.

The report shall reflect work actually performed, observed, analyzed, or
learned during the internship. Unrelated theoretical material should not
be included merely to increase report length.

\vspace{0.3cm}

% ================================================================
% X. TABLE OF CONTENTS
% ================================================================

\noindent
\textbf{X.\quad TABLE OF CONTENTS}

\vspace{0.15cm}

The Table of Contents shall be generated automatically using:

\begin{verbatim}
\tableofcontents
\end{verbatim}

Students shall not create or manually edit page numbers in the Table of
Contents.

The document shall be compiled sufficiently to ensure that all page
numbers and entries are correctly updated.

\vspace{0.3cm}

% ================================================================
% XI. LIST OF SYMBOLS AND ABBREVIATIONS
% ================================================================

\noindent
\textbf{XI.\quad LIST OF SYMBOLS AND ABBREVIATIONS}

\vspace{0.15cm}

Symbols and abbreviations that are used repeatedly in the report shall be
listed in the \textbf{List of Symbols and Abbreviations}.

The list may contain the following columns:

\begin{center}
\begin{tabular}{llll}
    No. & Symbol / Abbreviation & Unit & Definition
\end{tabular}
\end{center}

Only symbols and abbreviations actually used in the final report shall be
retained. Example entries contained in the template shall be removed if
they are not relevant to the student's report.

Each abbreviation should normally be defined when it first appears in
the main text.

\vspace{0.3cm}

% ================================================================
% XII. FIGURES, DIAGRAMS, AND PHOTOGRAPHS
% ================================================================

\noindent
\textbf{XII.\quad FIGURES, DIAGRAMS, AND PHOTOGRAPHS}

\vspace{0.15cm}

Figures shall be clear, legible, and directly relevant to the report
content.

Figures shall be inserted using the standard \verb|figure| environment.

Example:

\begin{verbatim}
\begin{figure}[H]
    \centering
    \includegraphics[width=0.8\textwidth]
                    {figures/company_structure.pdf}
    \caption{Organizational structure of the company}
    \label{fig:company-structure}
\end{figure}
\end{verbatim}

Figure captions shall be placed below the figure.

Each referenced figure shall have a unique \verb|\label| and shall be
referenced in the text using, for example:

\begin{verbatim}
Figure~\ref{fig:company-structure}
\end{verbatim}

Students shall not manually type figure numbers.

Figures reproduced or adapted from external sources shall include an
appropriate citation.

Photographs taken during the internship shall be relevant and sufficiently
clear. Students shall not include confidential, restricted, or proprietary
information without permission from the internship organization.

Vector graphics or high-resolution images are recommended whenever
possible.

\vspace{0.3cm}

% ================================================================
% XIII. TABLES
% ================================================================

\noindent
\textbf{XIII.\quad TABLES}

\vspace{0.15cm}

Tables shall be created using suitable \LaTeX{} environments such as
\verb|table|, \verb|tabular|, \verb|tabularx|, or \verb|longtable|.

Table captions shall be placed above the table.

Example:

\begin{verbatim}
\begin{table}[H]
    \centering
    \caption{Internship schedule}
    \label{tab:internship-schedule}

    \begin{tabular}{lll}
        \toprule
        Week & Task & Result \\
        \midrule
        1 & Orientation & Completed \\
        2 & System study & Completed \\
        \bottomrule
    \end{tabular}
\end{table}
\end{verbatim}

Tables shall be referenced using:

\begin{verbatim}
Table~\ref{tab:internship-schedule}
\end{verbatim}

Students shall not manually type table numbers.

Tables should be created directly in \LaTeX{} rather than inserted as
screenshots whenever practical.

\vspace{0.3cm}

% ================================================================
% XIV. EQUATIONS AND TECHNICAL QUANTITIES
% ================================================================

\noindent
\textbf{XIV.\quad EQUATIONS AND TECHNICAL QUANTITIES}

\vspace{0.15cm}

Mathematical expressions shall be typeset directly using \LaTeX{}.
Images or screenshots of equations shall not be used.

Short expressions may be written in inline mathematical mode, for example:

\begin{verbatim}
$R = 0.7\,\Omega$
\end{verbatim}

Important equations shall use a numbered equation environment.

Example:

\begin{verbatim}
\begin{equation}
    P = UI
    \label{eq:power}
\end{equation}
\end{verbatim}

Equations shall be referenced using:

\begin{verbatim}
Equation~\eqref{eq:power}
\end{verbatim}

Equation numbers shall not be manually entered.

SI units shall be used wherever applicable. The \verb|siunitx| package
should be used where appropriate.

Examples:

\begin{verbatim}
\SI{220}{\volt}
\SI{50}{\hertz}
\SI{2.5}{\ampere}
\end{verbatim}

Notation and units shall be consistent throughout the report.

\vspace{0.3cm}

% ================================================================
% XV. SOURCE CODE
% ================================================================

\noindent
\textbf{XV.\quad SOURCE CODE}

\vspace{0.15cm}

If the internship involves programming, PLCs, microcontrollers,
simulation, or software development, relevant source code may be included
in the report.

Source code shall be presented using the \verb|listings| environment or
another environment defined by the official template.

Example:

\begin{verbatim}
\begin{lstlisting}[language=st]
IF Start = TRUE THEN
    Motor := TRUE;
ELSE
    Motor := FALSE;
END_IF;
\end{lstlisting}
\end{verbatim}

Only code necessary for explaining the work should appear in the main
report. Long source-code listings should be moved to an appendix or
submitted as supplementary files if required.

Source code should not normally be inserted as screenshots.

\vspace{0.3cm}

% ================================================================
% XVI. ABSTRACT
% ================================================================

\noindent
\textbf{XVI.\quad ABSTRACT}

\vspace{0.15cm}

The Abstract shall briefly summarize:

\begin{itemize}[leftmargin=0.8cm, itemsep=2pt]
    \item the internship organization;
    \item the objectives of the internship;
    \item the main tasks performed;
    \item the main results, knowledge, and skills obtained.
\end{itemize}

The Abstract should be concise and should not reproduce the full content
of the report chapter by chapter.

\vspace{0.3cm}

% ================================================================
% XVII. REFERENCES AND CITATIONS
% ================================================================

\noindent
\textbf{XVII.\quad REFERENCES AND CITATIONS}

\vspace{0.15cm}

All books, journal papers, conference papers, standards, websites,
datasheets, manuals, technical documents, figures, tables, or other
external sources used in preparing the report shall be properly cited.

References shall be managed through the bibliography system provided with
the official template.

Bibliographic entries shall normally be stored in:

\begin{verbatim}
References.bib
\end{verbatim}

References shall be cited in the report using:

\begin{verbatim}
\cite{reference-key}
\end{verbatim}

Students shall not manually type reference numbers such as ``[1]'',
``[2]'', etc.

Titles of foreign-language publications shall remain in their original
language in the bibliography.

All borrowed ideas, data, figures, tables, equations, methods, or results
shall be properly acknowledged.

\vspace{0.3cm}

% ================================================================
% XVIII. CROSS-REFERENCING
% ================================================================

\noindent
\textbf{XVIII.\quad CROSS-REFERENCING}

\vspace{0.15cm}

Numbered elements shall be cross-referenced using the standard \LaTeX{}
mechanism.

Use the following commands where appropriate:

\begin{verbatim}
\label{...}
\ref{...}
\eqref{...}
\end{verbatim}

For example:

\begin{verbatim}
As shown in Figure~\ref{fig:system}, ...
\end{verbatim}

Students shall not manually type numbers such as ``Figure 2.3'' when
those numbers can be generated automatically.

\vspace{0.3cm}

% ================================================================
% XIX. PAGE NUMBERING
% ================================================================

\noindent
\textbf{XIX.\quad PAGE NUMBERING}

\vspace{0.15cm}

Cover pages and official evaluation forms at the beginning of the report
shall not display page numbers.

The main body of the report shall use Arabic page numbers beginning from
Chapter 1.

Page numbering shall be generated automatically by the official template.
Students shall not manually enter page numbers.

\vspace{0.3cm}

% ================================================================
% XX. PROJECT FILE ORGANIZATION
% ================================================================

\noindent
\textbf{XX.\quad PROJECT FILE ORGANIZATION}

\vspace{0.15cm}

Students should maintain a clear and organized \LaTeX{} project
structure. A recommended structure is:

\begin{verbatim}
InternshipReport/
|
|-- main.tex
|-- References.bib
|
|-- figures/
|   |-- company.pdf
|   |-- system.png
|   `-- ...
|
`-- chapters/
    |-- chapter1.tex
    |-- chapter2.tex
    `-- chapter3.tex
\end{verbatim}

Figures should normally be stored in the \verb|figures| directory
specified by the official template.

File and directory names should preferably contain Latin letters,
numbers, hyphens, or underscores.

Recommended examples include:

\begin{verbatim}
company_structure.pdf
motor_control.png
chapter_01.tex
\end{verbatim}

Spaces and unnecessary special characters in file names should be
avoided where possible.

\vspace{0.3cm}

% ================================================================
% XXI. DOCUMENT COMPILATION AND FINAL CHECK
% ================================================================

\noindent
\textbf{XXI.\quad DOCUMENT COMPILATION AND FINAL CHECK}

\vspace{0.15cm}

Before submission, the complete \LaTeX{} project shall compile
successfully and produce the final PDF report.

The final document shall not contain unresolved references such as:

\begin{verbatim}
??
[?]
\end{verbatim}

Before generating the final PDF, students shall check:

\begin{itemize}[leftmargin=0.8cm, itemsep=2pt]
    \item the Table of Contents;
    \item page numbering;
    \item chapter, section, and subsection numbering;
    \item figure numbering;
    \item table numbering;
    \item equation numbering;
    \item cross-references;
    \item citations;
    \item bibliography;
    \item font display;
    \item mathematical symbols;
    \item image quality;
    \item completeness of all required forms.
\end{itemize}

The document may be compiled using TeX Live, MiKTeX, or another
compatible \LaTeX{} environment. Regardless of the compilation
environment, the resulting PDF shall conform to the official Faculty
template.

\vspace{0.3cm}

% ================================================================
% XXII. CONTENT THAT STUDENTS MAY EDIT
% ================================================================

\noindent
\textbf{XXII.\quad CONTENT THAT STUDENTS MAY EDIT}

\vspace{0.15cm}

Students may edit or replace:

\begin{itemize}[leftmargin=0.8cm, itemsep=2pt]
    \item student name;
    \item class;
    \item internship organization;
    \item supervisor name;
    \item Abstract;
    \item List of Symbols and Abbreviations;
    \item chapter titles and chapter content;
    \item figures and photographs;
    \item tables;
    \item equations;
    \item citations and references.
\end{itemize}

Students shall not modify the main document formatting settings in the
preamble, including the document class, page geometry, line spacing,
heading formatting, table-of-contents formatting, font configuration,
page layout, or cover-page macros unless specifically instructed.

\vspace{0.3cm}

% ================================================================
% XXIII. FINAL SUBMISSION REQUIREMENTS
% ================================================================

\noindent
\textbf{XXIII.\quad FINAL SUBMISSION REQUIREMENTS}

\vspace{0.15cm}

Before submission, the final report shall:

\begin{itemize}[leftmargin=0.8cm, itemsep=2pt]
    \item compile without errors that affect the report content or presentation;
    \item contain all required forms and report sections;
    \item contain correct student, class, internship organization, and supervisor information;
    \item contain an updated Table of Contents, List of Figures, and List of Tables;
    \item contain no unused example content from the template;
    \item contain no placeholder text such as \textit{Title of Chapter 1},
    \textit{Section 1.1}, or other sample entries;
    \item contain all required citations and references;
    \item contain no unresolved cross-references.
\end{itemize}

If the Faculty requires submission of the \LaTeX{} source, the complete
project shall include all \verb|.tex| files, the bibliography database,
figures, and any other files required to reproduce the submitted PDF.

\vspace{0.4cm}

\noindent
\textbf{Important:}
Before submitting the final internship report, students shall remove or
comment out the following command in the main document:

\begin{verbatim}
\input{Regulations.tex}
\end{verbatim}

For example:

\begin{verbatim}
% \input{Regulations.tex}
\end{verbatim}

This removes the guidance pages from the final submitted report while
preserving the regulations file in the \LaTeX{} project for reference.

} % End 13pt content

\clearpage

\end{verbatim}

This removes the guidance pages from the final submitted report while
preserving the regulations file in the \LaTeX{} project for reference.

} % End 13pt content

\clearpage
 command before submitting
% the final internship report.
%
% =================================================================

\clearpage
\thispagestyle{empty}

\begin{center}
    {\fontsize{16pt}{20pt}\selectfont\bfseries
    SPECIFIC REGULATIONS FOR PREPARING\\
    GRADUATION INTERNSHIP REPORTS USING \LaTeX}
\end{center}

\vspace{0.5cm}

{\fontsize{13pt}{17pt}\selectfont
\setstretch{1.3}
\justifying

\noindent
\textbf{Note:}
This section is provided for guidance only and is included in the main
document using the command \verb|
% =================================================================
% REGULATIONS FOR PREPARING THE GRADUATION INTERNSHIP REPORT
% USING LATEX
% =================================================================
%
% This file is intended to be included from the main document with:
%
%     
% =================================================================
% REGULATIONS FOR PREPARING THE GRADUATION INTERNSHIP REPORT
% USING LATEX
% =================================================================
%
% This file is intended to be included from the main document with:
%
%     \input{Regulations.tex}
%
% IMPORTANT:
% This section is provided for guidance only. Students shall remove
% or comment out the \input{Regulations.tex} command before submitting
% the final internship report.
%
% =================================================================

\clearpage
\thispagestyle{empty}

\begin{center}
    {\fontsize{16pt}{20pt}\selectfont\bfseries
    SPECIFIC REGULATIONS FOR PREPARING\\
    GRADUATION INTERNSHIP REPORTS USING \LaTeX}
\end{center}

\vspace{0.5cm}

{\fontsize{13pt}{17pt}\selectfont
\setstretch{1.3}
\justifying

\noindent
\textbf{Note:}
This section is provided for guidance only and is included in the main
document using the command \verb|\input{Regulations.tex}|.
Students shall remove or comment out this command before submitting the
final internship report.

\vspace{0.4cm}

% ================================================================
% I. GENERAL REQUIREMENTS
% ================================================================

\noindent
\textbf{I.\quad GENERAL REQUIREMENTS}

\vspace{0.15cm}

The Graduation Internship Report shall be prepared using the official
\LaTeX{} template provided by the Faculty. Students shall preserve the
predefined page layout, typography, line spacing, heading styles,
numbering system, figure and table formatting, and bibliography settings
unless otherwise instructed by the Faculty or the supervisor.

The report shall be written in English.

Students are responsible for the academic and technical content of the
report. The official \LaTeX{} template is responsible for standardized
document formatting.

\vspace{0.3cm}

% ================================================================
% II. REPORT STRUCTURE
% ================================================================

\noindent
\textbf{II.\quad REPORT STRUCTURE}

\vspace{0.15cm}

The report shall normally be arranged in the following order:

\begin{enumerate}[label=\arabic*., leftmargin=1.2cm, itemsep=2pt]
    \item Outer Cover;
    \item Comments of the Internship Host Organization;
    \item Supervisor's Evaluation and Graduation Internship Defense Results;
    \item Table of Contents;
    \item List of Symbols and Abbreviations, if applicable;
    \item List of Figures;
    \item List of Tables;
    \item Abstract;
    \item Main Chapters of the Internship Report;
    \item References;
    \item Appendices, if applicable.
\end{enumerate}

The preliminary pages shall not be numbered as chapters. The main report
content shall be organized into consecutively numbered chapters beginning
with Chapter 1.

\vspace{0.3cm}

% ================================================================
% III. PAGE SIZE AND MARGINS
% ================================================================

\noindent
\textbf{III.\quad PAGE SIZE AND MARGINS}

\vspace{0.15cm}

The report shall be formatted on A4 paper with the following margins:

\begin{center}
\begin{tabular}{ll}
    Left margin:   & 3.5 cm \\
    Right margin:  & 2.0 cm \\
    Top margin:    & 2.5 cm \\
    Bottom margin: & 2.5 cm
\end{tabular}
\end{center}

These settings are predefined in the official template and shall not be
modified by students.

\vspace{0.3cm}

% ================================================================
% IV. FONT, FONT SIZE, AND LINE SPACING
% ================================================================

\noindent
\textbf{IV.\quad FONT, FONT SIZE, AND LINE SPACING}

\vspace{0.15cm}

The main text shall use a 13 pt font size with 1.3-line spacing.
A Times-family typeface shall be used in accordance with the official
template.

The text shall be fully justified.

Students shall not use repeated spaces, repeated manual line breaks, or
multiple blank lines to position text. Paragraph and heading spacing shall
be controlled by the template.

\vspace{0.3cm}

% ================================================================
% V. OUTER COVER
% ================================================================

\noindent
\textbf{V.\quad OUTER COVER}

\vspace{0.15cm}

The outer cover shall follow the official format provided in the template
and shall contain the following information:

\begin{itemize}[leftmargin=0.8cm, itemsep=2pt]
    \item Thai Nguyen University of Technology;
    \item Faculty of International Training;
    \item Engineering and Technology Department;
    \item University logo and Faculty logo;
    \item the title \textbf{INTERNSHIP REPORT};
    \item Internship Organization;
    \item Author;
    \item Class;
    \item Supervisor;
    \item place and year of completion.
\end{itemize}

Students shall only replace the required information fields and shall not
modify the cover-page layout unless instructed otherwise.

\vspace{0.3cm}

% ================================================================
% VI. COMMENTS OF THE INTERNSHIP HOST ORGANIZATION
% ================================================================

\noindent
\textbf{VI.\quad COMMENTS OF THE INTERNSHIP HOST ORGANIZATION}

\vspace{0.15cm}

The page entitled \textbf{COMMENTS OF THE INTERNSHIP HOST ORGANIZATION}
is reserved for the internship host organization.

It includes:

\begin{itemize}[leftmargin=0.8cm, itemsep=2pt]
    \item the student's full name;
    \item class;
    \item comments from the host organization;
    \item score;
    \item date;
    \item signature and official seal.
\end{itemize}

Students shall not alter or complete official evaluation fields unless
specifically authorized.

\vspace{0.3cm}

% ================================================================
% VII. SUPERVISOR'S EVALUATION AND DEFENSE RESULTS
% ================================================================

\noindent
\textbf{VII.\quad SUPERVISOR'S EVALUATION AND DEFENSE RESULTS}

\vspace{0.15cm}

The section entitled
\textbf{SUPERVISOR'S EVALUATION AND GRADUATION INTERNSHIP DEFENSE RESULTS}
is reserved for the supervisor and defense committee.

The evaluation may include:

\begin{itemize}[leftmargin=0.8cm, itemsep=2pt]
    \item Internship Progress and Student's Attitude;
    \item Report Content;
    \item Report Presentation;
    \item Other Comments;
    \item Supervisor's Opinion;
    \item Defense Result;
    \item Score.
\end{itemize}

Students shall not enter evaluation scores or defense results before they
are officially determined.

\vspace{0.3cm}

% ================================================================
% VIII. CHAPTERS, SECTIONS, AND SUBSECTIONS
% ================================================================

\noindent
\textbf{VIII.\quad CHAPTERS, SECTIONS, AND SUBSECTIONS}

\vspace{0.15cm}

The report structure shall be generated using the standard \LaTeX{}
commands:

\begin{verbatim}
\chapter{...}
\section{...}
\subsection{...}
\end{verbatim}

Chapter, section, and subsection numbers shall be generated automatically
by \LaTeX{}.

Students shall not manually type chapter or section numbers and shall not
reproduce heading styles using commands such as \verb|\textbf|,
\verb|\fontsize|, or \verb|\centering| instead of the appropriate
structural command.

Excessive subdivision should be avoided. In most cases, the hierarchy
Chapter--Section--Subsection is sufficient.

\vspace{0.3cm}

% ================================================================
% IX. ORGANIZATION OF THE MAIN CONTENT
% ================================================================

\noindent
\textbf{IX.\quad ORGANIZATION OF THE MAIN CONTENT}

\vspace{0.15cm}

The specific content of the report shall depend on the internship
organization and the assigned tasks. A typical structure may include:

\begin{itemize}[leftmargin=0.8cm, itemsep=2pt]
    \item \textbf{Chapter 1. Overview of the Internship Organization};
    \item \textbf{Chapter 2. Internship Activities and Assigned Tasks};
    \item \textbf{Chapter 3. Results, Experience, and Evaluation}.
\end{itemize}

These chapter titles are recommendations only and may be adapted to the
actual internship content.

The report shall reflect work actually performed, observed, analyzed, or
learned during the internship. Unrelated theoretical material should not
be included merely to increase report length.

\vspace{0.3cm}

% ================================================================
% X. TABLE OF CONTENTS
% ================================================================

\noindent
\textbf{X.\quad TABLE OF CONTENTS}

\vspace{0.15cm}

The Table of Contents shall be generated automatically using:

\begin{verbatim}
\tableofcontents
\end{verbatim}

Students shall not create or manually edit page numbers in the Table of
Contents.

The document shall be compiled sufficiently to ensure that all page
numbers and entries are correctly updated.

\vspace{0.3cm}

% ================================================================
% XI. LIST OF SYMBOLS AND ABBREVIATIONS
% ================================================================

\noindent
\textbf{XI.\quad LIST OF SYMBOLS AND ABBREVIATIONS}

\vspace{0.15cm}

Symbols and abbreviations that are used repeatedly in the report shall be
listed in the \textbf{List of Symbols and Abbreviations}.

The list may contain the following columns:

\begin{center}
\begin{tabular}{llll}
    No. & Symbol / Abbreviation & Unit & Definition
\end{tabular}
\end{center}

Only symbols and abbreviations actually used in the final report shall be
retained. Example entries contained in the template shall be removed if
they are not relevant to the student's report.

Each abbreviation should normally be defined when it first appears in
the main text.

\vspace{0.3cm}

% ================================================================
% XII. FIGURES, DIAGRAMS, AND PHOTOGRAPHS
% ================================================================

\noindent
\textbf{XII.\quad FIGURES, DIAGRAMS, AND PHOTOGRAPHS}

\vspace{0.15cm}

Figures shall be clear, legible, and directly relevant to the report
content.

Figures shall be inserted using the standard \verb|figure| environment.

Example:

\begin{verbatim}
\begin{figure}[H]
    \centering
    \includegraphics[width=0.8\textwidth]
                    {figures/company_structure.pdf}
    \caption{Organizational structure of the company}
    \label{fig:company-structure}
\end{figure}
\end{verbatim}

Figure captions shall be placed below the figure.

Each referenced figure shall have a unique \verb|\label| and shall be
referenced in the text using, for example:

\begin{verbatim}
Figure~\ref{fig:company-structure}
\end{verbatim}

Students shall not manually type figure numbers.

Figures reproduced or adapted from external sources shall include an
appropriate citation.

Photographs taken during the internship shall be relevant and sufficiently
clear. Students shall not include confidential, restricted, or proprietary
information without permission from the internship organization.

Vector graphics or high-resolution images are recommended whenever
possible.

\vspace{0.3cm}

% ================================================================
% XIII. TABLES
% ================================================================

\noindent
\textbf{XIII.\quad TABLES}

\vspace{0.15cm}

Tables shall be created using suitable \LaTeX{} environments such as
\verb|table|, \verb|tabular|, \verb|tabularx|, or \verb|longtable|.

Table captions shall be placed above the table.

Example:

\begin{verbatim}
\begin{table}[H]
    \centering
    \caption{Internship schedule}
    \label{tab:internship-schedule}

    \begin{tabular}{lll}
        \toprule
        Week & Task & Result \\
        \midrule
        1 & Orientation & Completed \\
        2 & System study & Completed \\
        \bottomrule
    \end{tabular}
\end{table}
\end{verbatim}

Tables shall be referenced using:

\begin{verbatim}
Table~\ref{tab:internship-schedule}
\end{verbatim}

Students shall not manually type table numbers.

Tables should be created directly in \LaTeX{} rather than inserted as
screenshots whenever practical.

\vspace{0.3cm}

% ================================================================
% XIV. EQUATIONS AND TECHNICAL QUANTITIES
% ================================================================

\noindent
\textbf{XIV.\quad EQUATIONS AND TECHNICAL QUANTITIES}

\vspace{0.15cm}

Mathematical expressions shall be typeset directly using \LaTeX{}.
Images or screenshots of equations shall not be used.

Short expressions may be written in inline mathematical mode, for example:

\begin{verbatim}
$R = 0.7\,\Omega$
\end{verbatim}

Important equations shall use a numbered equation environment.

Example:

\begin{verbatim}
\begin{equation}
    P = UI
    \label{eq:power}
\end{equation}
\end{verbatim}

Equations shall be referenced using:

\begin{verbatim}
Equation~\eqref{eq:power}
\end{verbatim}

Equation numbers shall not be manually entered.

SI units shall be used wherever applicable. The \verb|siunitx| package
should be used where appropriate.

Examples:

\begin{verbatim}
\SI{220}{\volt}
\SI{50}{\hertz}
\SI{2.5}{\ampere}
\end{verbatim}

Notation and units shall be consistent throughout the report.

\vspace{0.3cm}

% ================================================================
% XV. SOURCE CODE
% ================================================================

\noindent
\textbf{XV.\quad SOURCE CODE}

\vspace{0.15cm}

If the internship involves programming, PLCs, microcontrollers,
simulation, or software development, relevant source code may be included
in the report.

Source code shall be presented using the \verb|listings| environment or
another environment defined by the official template.

Example:

\begin{verbatim}
\begin{lstlisting}[language=st]
IF Start = TRUE THEN
    Motor := TRUE;
ELSE
    Motor := FALSE;
END_IF;
\end{lstlisting}
\end{verbatim}

Only code necessary for explaining the work should appear in the main
report. Long source-code listings should be moved to an appendix or
submitted as supplementary files if required.

Source code should not normally be inserted as screenshots.

\vspace{0.3cm}

% ================================================================
% XVI. ABSTRACT
% ================================================================

\noindent
\textbf{XVI.\quad ABSTRACT}

\vspace{0.15cm}

The Abstract shall briefly summarize:

\begin{itemize}[leftmargin=0.8cm, itemsep=2pt]
    \item the internship organization;
    \item the objectives of the internship;
    \item the main tasks performed;
    \item the main results, knowledge, and skills obtained.
\end{itemize}

The Abstract should be concise and should not reproduce the full content
of the report chapter by chapter.

\vspace{0.3cm}

% ================================================================
% XVII. REFERENCES AND CITATIONS
% ================================================================

\noindent
\textbf{XVII.\quad REFERENCES AND CITATIONS}

\vspace{0.15cm}

All books, journal papers, conference papers, standards, websites,
datasheets, manuals, technical documents, figures, tables, or other
external sources used in preparing the report shall be properly cited.

References shall be managed through the bibliography system provided with
the official template.

Bibliographic entries shall normally be stored in:

\begin{verbatim}
References.bib
\end{verbatim}

References shall be cited in the report using:

\begin{verbatim}
\cite{reference-key}
\end{verbatim}

Students shall not manually type reference numbers such as ``[1]'',
``[2]'', etc.

Titles of foreign-language publications shall remain in their original
language in the bibliography.

All borrowed ideas, data, figures, tables, equations, methods, or results
shall be properly acknowledged.

\vspace{0.3cm}

% ================================================================
% XVIII. CROSS-REFERENCING
% ================================================================

\noindent
\textbf{XVIII.\quad CROSS-REFERENCING}

\vspace{0.15cm}

Numbered elements shall be cross-referenced using the standard \LaTeX{}
mechanism.

Use the following commands where appropriate:

\begin{verbatim}
\label{...}
\ref{...}
\eqref{...}
\end{verbatim}

For example:

\begin{verbatim}
As shown in Figure~\ref{fig:system}, ...
\end{verbatim}

Students shall not manually type numbers such as ``Figure 2.3'' when
those numbers can be generated automatically.

\vspace{0.3cm}

% ================================================================
% XIX. PAGE NUMBERING
% ================================================================

\noindent
\textbf{XIX.\quad PAGE NUMBERING}

\vspace{0.15cm}

Cover pages and official evaluation forms at the beginning of the report
shall not display page numbers.

The main body of the report shall use Arabic page numbers beginning from
Chapter 1.

Page numbering shall be generated automatically by the official template.
Students shall not manually enter page numbers.

\vspace{0.3cm}

% ================================================================
% XX. PROJECT FILE ORGANIZATION
% ================================================================

\noindent
\textbf{XX.\quad PROJECT FILE ORGANIZATION}

\vspace{0.15cm}

Students should maintain a clear and organized \LaTeX{} project
structure. A recommended structure is:

\begin{verbatim}
InternshipReport/
|
|-- main.tex
|-- References.bib
|
|-- figures/
|   |-- company.pdf
|   |-- system.png
|   `-- ...
|
`-- chapters/
    |-- chapter1.tex
    |-- chapter2.tex
    `-- chapter3.tex
\end{verbatim}

Figures should normally be stored in the \verb|figures| directory
specified by the official template.

File and directory names should preferably contain Latin letters,
numbers, hyphens, or underscores.

Recommended examples include:

\begin{verbatim}
company_structure.pdf
motor_control.png
chapter_01.tex
\end{verbatim}

Spaces and unnecessary special characters in file names should be
avoided where possible.

\vspace{0.3cm}

% ================================================================
% XXI. DOCUMENT COMPILATION AND FINAL CHECK
% ================================================================

\noindent
\textbf{XXI.\quad DOCUMENT COMPILATION AND FINAL CHECK}

\vspace{0.15cm}

Before submission, the complete \LaTeX{} project shall compile
successfully and produce the final PDF report.

The final document shall not contain unresolved references such as:

\begin{verbatim}
??
[?]
\end{verbatim}

Before generating the final PDF, students shall check:

\begin{itemize}[leftmargin=0.8cm, itemsep=2pt]
    \item the Table of Contents;
    \item page numbering;
    \item chapter, section, and subsection numbering;
    \item figure numbering;
    \item table numbering;
    \item equation numbering;
    \item cross-references;
    \item citations;
    \item bibliography;
    \item font display;
    \item mathematical symbols;
    \item image quality;
    \item completeness of all required forms.
\end{itemize}

The document may be compiled using TeX Live, MiKTeX, or another
compatible \LaTeX{} environment. Regardless of the compilation
environment, the resulting PDF shall conform to the official Faculty
template.

\vspace{0.3cm}

% ================================================================
% XXII. CONTENT THAT STUDENTS MAY EDIT
% ================================================================

\noindent
\textbf{XXII.\quad CONTENT THAT STUDENTS MAY EDIT}

\vspace{0.15cm}

Students may edit or replace:

\begin{itemize}[leftmargin=0.8cm, itemsep=2pt]
    \item student name;
    \item class;
    \item internship organization;
    \item supervisor name;
    \item Abstract;
    \item List of Symbols and Abbreviations;
    \item chapter titles and chapter content;
    \item figures and photographs;
    \item tables;
    \item equations;
    \item citations and references.
\end{itemize}

Students shall not modify the main document formatting settings in the
preamble, including the document class, page geometry, line spacing,
heading formatting, table-of-contents formatting, font configuration,
page layout, or cover-page macros unless specifically instructed.

\vspace{0.3cm}

% ================================================================
% XXIII. FINAL SUBMISSION REQUIREMENTS
% ================================================================

\noindent
\textbf{XXIII.\quad FINAL SUBMISSION REQUIREMENTS}

\vspace{0.15cm}

Before submission, the final report shall:

\begin{itemize}[leftmargin=0.8cm, itemsep=2pt]
    \item compile without errors that affect the report content or presentation;
    \item contain all required forms and report sections;
    \item contain correct student, class, internship organization, and supervisor information;
    \item contain an updated Table of Contents, List of Figures, and List of Tables;
    \item contain no unused example content from the template;
    \item contain no placeholder text such as \textit{Title of Chapter 1},
    \textit{Section 1.1}, or other sample entries;
    \item contain all required citations and references;
    \item contain no unresolved cross-references.
\end{itemize}

If the Faculty requires submission of the \LaTeX{} source, the complete
project shall include all \verb|.tex| files, the bibliography database,
figures, and any other files required to reproduce the submitted PDF.

\vspace{0.4cm}

\noindent
\textbf{Important:}
Before submitting the final internship report, students shall remove or
comment out the following command in the main document:

\begin{verbatim}
\input{Regulations.tex}
\end{verbatim}

For example:

\begin{verbatim}
% \input{Regulations.tex}
\end{verbatim}

This removes the guidance pages from the final submitted report while
preserving the regulations file in the \LaTeX{} project for reference.

} % End 13pt content

\clearpage

%
% IMPORTANT:
% This section is provided for guidance only. Students shall remove
% or comment out the 
% =================================================================
% REGULATIONS FOR PREPARING THE GRADUATION INTERNSHIP REPORT
% USING LATEX
% =================================================================
%
% This file is intended to be included from the main document with:
%
%     \input{Regulations.tex}
%
% IMPORTANT:
% This section is provided for guidance only. Students shall remove
% or comment out the \input{Regulations.tex} command before submitting
% the final internship report.
%
% =================================================================

\clearpage
\thispagestyle{empty}

\begin{center}
    {\fontsize{16pt}{20pt}\selectfont\bfseries
    SPECIFIC REGULATIONS FOR PREPARING\\
    GRADUATION INTERNSHIP REPORTS USING \LaTeX}
\end{center}

\vspace{0.5cm}

{\fontsize{13pt}{17pt}\selectfont
\setstretch{1.3}
\justifying

\noindent
\textbf{Note:}
This section is provided for guidance only and is included in the main
document using the command \verb|\input{Regulations.tex}|.
Students shall remove or comment out this command before submitting the
final internship report.

\vspace{0.4cm}

% ================================================================
% I. GENERAL REQUIREMENTS
% ================================================================

\noindent
\textbf{I.\quad GENERAL REQUIREMENTS}

\vspace{0.15cm}

The Graduation Internship Report shall be prepared using the official
\LaTeX{} template provided by the Faculty. Students shall preserve the
predefined page layout, typography, line spacing, heading styles,
numbering system, figure and table formatting, and bibliography settings
unless otherwise instructed by the Faculty or the supervisor.

The report shall be written in English.

Students are responsible for the academic and technical content of the
report. The official \LaTeX{} template is responsible for standardized
document formatting.

\vspace{0.3cm}

% ================================================================
% II. REPORT STRUCTURE
% ================================================================

\noindent
\textbf{II.\quad REPORT STRUCTURE}

\vspace{0.15cm}

The report shall normally be arranged in the following order:

\begin{enumerate}[label=\arabic*., leftmargin=1.2cm, itemsep=2pt]
    \item Outer Cover;
    \item Comments of the Internship Host Organization;
    \item Supervisor's Evaluation and Graduation Internship Defense Results;
    \item Table of Contents;
    \item List of Symbols and Abbreviations, if applicable;
    \item List of Figures;
    \item List of Tables;
    \item Abstract;
    \item Main Chapters of the Internship Report;
    \item References;
    \item Appendices, if applicable.
\end{enumerate}

The preliminary pages shall not be numbered as chapters. The main report
content shall be organized into consecutively numbered chapters beginning
with Chapter 1.

\vspace{0.3cm}

% ================================================================
% III. PAGE SIZE AND MARGINS
% ================================================================

\noindent
\textbf{III.\quad PAGE SIZE AND MARGINS}

\vspace{0.15cm}

The report shall be formatted on A4 paper with the following margins:

\begin{center}
\begin{tabular}{ll}
    Left margin:   & 3.5 cm \\
    Right margin:  & 2.0 cm \\
    Top margin:    & 2.5 cm \\
    Bottom margin: & 2.5 cm
\end{tabular}
\end{center}

These settings are predefined in the official template and shall not be
modified by students.

\vspace{0.3cm}

% ================================================================
% IV. FONT, FONT SIZE, AND LINE SPACING
% ================================================================

\noindent
\textbf{IV.\quad FONT, FONT SIZE, AND LINE SPACING}

\vspace{0.15cm}

The main text shall use a 13 pt font size with 1.3-line spacing.
A Times-family typeface shall be used in accordance with the official
template.

The text shall be fully justified.

Students shall not use repeated spaces, repeated manual line breaks, or
multiple blank lines to position text. Paragraph and heading spacing shall
be controlled by the template.

\vspace{0.3cm}

% ================================================================
% V. OUTER COVER
% ================================================================

\noindent
\textbf{V.\quad OUTER COVER}

\vspace{0.15cm}

The outer cover shall follow the official format provided in the template
and shall contain the following information:

\begin{itemize}[leftmargin=0.8cm, itemsep=2pt]
    \item Thai Nguyen University of Technology;
    \item Faculty of International Training;
    \item Engineering and Technology Department;
    \item University logo and Faculty logo;
    \item the title \textbf{INTERNSHIP REPORT};
    \item Internship Organization;
    \item Author;
    \item Class;
    \item Supervisor;
    \item place and year of completion.
\end{itemize}

Students shall only replace the required information fields and shall not
modify the cover-page layout unless instructed otherwise.

\vspace{0.3cm}

% ================================================================
% VI. COMMENTS OF THE INTERNSHIP HOST ORGANIZATION
% ================================================================

\noindent
\textbf{VI.\quad COMMENTS OF THE INTERNSHIP HOST ORGANIZATION}

\vspace{0.15cm}

The page entitled \textbf{COMMENTS OF THE INTERNSHIP HOST ORGANIZATION}
is reserved for the internship host organization.

It includes:

\begin{itemize}[leftmargin=0.8cm, itemsep=2pt]
    \item the student's full name;
    \item class;
    \item comments from the host organization;
    \item score;
    \item date;
    \item signature and official seal.
\end{itemize}

Students shall not alter or complete official evaluation fields unless
specifically authorized.

\vspace{0.3cm}

% ================================================================
% VII. SUPERVISOR'S EVALUATION AND DEFENSE RESULTS
% ================================================================

\noindent
\textbf{VII.\quad SUPERVISOR'S EVALUATION AND DEFENSE RESULTS}

\vspace{0.15cm}

The section entitled
\textbf{SUPERVISOR'S EVALUATION AND GRADUATION INTERNSHIP DEFENSE RESULTS}
is reserved for the supervisor and defense committee.

The evaluation may include:

\begin{itemize}[leftmargin=0.8cm, itemsep=2pt]
    \item Internship Progress and Student's Attitude;
    \item Report Content;
    \item Report Presentation;
    \item Other Comments;
    \item Supervisor's Opinion;
    \item Defense Result;
    \item Score.
\end{itemize}

Students shall not enter evaluation scores or defense results before they
are officially determined.

\vspace{0.3cm}

% ================================================================
% VIII. CHAPTERS, SECTIONS, AND SUBSECTIONS
% ================================================================

\noindent
\textbf{VIII.\quad CHAPTERS, SECTIONS, AND SUBSECTIONS}

\vspace{0.15cm}

The report structure shall be generated using the standard \LaTeX{}
commands:

\begin{verbatim}
\chapter{...}
\section{...}
\subsection{...}
\end{verbatim}

Chapter, section, and subsection numbers shall be generated automatically
by \LaTeX{}.

Students shall not manually type chapter or section numbers and shall not
reproduce heading styles using commands such as \verb|\textbf|,
\verb|\fontsize|, or \verb|\centering| instead of the appropriate
structural command.

Excessive subdivision should be avoided. In most cases, the hierarchy
Chapter--Section--Subsection is sufficient.

\vspace{0.3cm}

% ================================================================
% IX. ORGANIZATION OF THE MAIN CONTENT
% ================================================================

\noindent
\textbf{IX.\quad ORGANIZATION OF THE MAIN CONTENT}

\vspace{0.15cm}

The specific content of the report shall depend on the internship
organization and the assigned tasks. A typical structure may include:

\begin{itemize}[leftmargin=0.8cm, itemsep=2pt]
    \item \textbf{Chapter 1. Overview of the Internship Organization};
    \item \textbf{Chapter 2. Internship Activities and Assigned Tasks};
    \item \textbf{Chapter 3. Results, Experience, and Evaluation}.
\end{itemize}

These chapter titles are recommendations only and may be adapted to the
actual internship content.

The report shall reflect work actually performed, observed, analyzed, or
learned during the internship. Unrelated theoretical material should not
be included merely to increase report length.

\vspace{0.3cm}

% ================================================================
% X. TABLE OF CONTENTS
% ================================================================

\noindent
\textbf{X.\quad TABLE OF CONTENTS}

\vspace{0.15cm}

The Table of Contents shall be generated automatically using:

\begin{verbatim}
\tableofcontents
\end{verbatim}

Students shall not create or manually edit page numbers in the Table of
Contents.

The document shall be compiled sufficiently to ensure that all page
numbers and entries are correctly updated.

\vspace{0.3cm}

% ================================================================
% XI. LIST OF SYMBOLS AND ABBREVIATIONS
% ================================================================

\noindent
\textbf{XI.\quad LIST OF SYMBOLS AND ABBREVIATIONS}

\vspace{0.15cm}

Symbols and abbreviations that are used repeatedly in the report shall be
listed in the \textbf{List of Symbols and Abbreviations}.

The list may contain the following columns:

\begin{center}
\begin{tabular}{llll}
    No. & Symbol / Abbreviation & Unit & Definition
\end{tabular}
\end{center}

Only symbols and abbreviations actually used in the final report shall be
retained. Example entries contained in the template shall be removed if
they are not relevant to the student's report.

Each abbreviation should normally be defined when it first appears in
the main text.

\vspace{0.3cm}

% ================================================================
% XII. FIGURES, DIAGRAMS, AND PHOTOGRAPHS
% ================================================================

\noindent
\textbf{XII.\quad FIGURES, DIAGRAMS, AND PHOTOGRAPHS}

\vspace{0.15cm}

Figures shall be clear, legible, and directly relevant to the report
content.

Figures shall be inserted using the standard \verb|figure| environment.

Example:

\begin{verbatim}
\begin{figure}[H]
    \centering
    \includegraphics[width=0.8\textwidth]
                    {figures/company_structure.pdf}
    \caption{Organizational structure of the company}
    \label{fig:company-structure}
\end{figure}
\end{verbatim}

Figure captions shall be placed below the figure.

Each referenced figure shall have a unique \verb|\label| and shall be
referenced in the text using, for example:

\begin{verbatim}
Figure~\ref{fig:company-structure}
\end{verbatim}

Students shall not manually type figure numbers.

Figures reproduced or adapted from external sources shall include an
appropriate citation.

Photographs taken during the internship shall be relevant and sufficiently
clear. Students shall not include confidential, restricted, or proprietary
information without permission from the internship organization.

Vector graphics or high-resolution images are recommended whenever
possible.

\vspace{0.3cm}

% ================================================================
% XIII. TABLES
% ================================================================

\noindent
\textbf{XIII.\quad TABLES}

\vspace{0.15cm}

Tables shall be created using suitable \LaTeX{} environments such as
\verb|table|, \verb|tabular|, \verb|tabularx|, or \verb|longtable|.

Table captions shall be placed above the table.

Example:

\begin{verbatim}
\begin{table}[H]
    \centering
    \caption{Internship schedule}
    \label{tab:internship-schedule}

    \begin{tabular}{lll}
        \toprule
        Week & Task & Result \\
        \midrule
        1 & Orientation & Completed \\
        2 & System study & Completed \\
        \bottomrule
    \end{tabular}
\end{table}
\end{verbatim}

Tables shall be referenced using:

\begin{verbatim}
Table~\ref{tab:internship-schedule}
\end{verbatim}

Students shall not manually type table numbers.

Tables should be created directly in \LaTeX{} rather than inserted as
screenshots whenever practical.

\vspace{0.3cm}

% ================================================================
% XIV. EQUATIONS AND TECHNICAL QUANTITIES
% ================================================================

\noindent
\textbf{XIV.\quad EQUATIONS AND TECHNICAL QUANTITIES}

\vspace{0.15cm}

Mathematical expressions shall be typeset directly using \LaTeX{}.
Images or screenshots of equations shall not be used.

Short expressions may be written in inline mathematical mode, for example:

\begin{verbatim}
$R = 0.7\,\Omega$
\end{verbatim}

Important equations shall use a numbered equation environment.

Example:

\begin{verbatim}
\begin{equation}
    P = UI
    \label{eq:power}
\end{equation}
\end{verbatim}

Equations shall be referenced using:

\begin{verbatim}
Equation~\eqref{eq:power}
\end{verbatim}

Equation numbers shall not be manually entered.

SI units shall be used wherever applicable. The \verb|siunitx| package
should be used where appropriate.

Examples:

\begin{verbatim}
\SI{220}{\volt}
\SI{50}{\hertz}
\SI{2.5}{\ampere}
\end{verbatim}

Notation and units shall be consistent throughout the report.

\vspace{0.3cm}

% ================================================================
% XV. SOURCE CODE
% ================================================================

\noindent
\textbf{XV.\quad SOURCE CODE}

\vspace{0.15cm}

If the internship involves programming, PLCs, microcontrollers,
simulation, or software development, relevant source code may be included
in the report.

Source code shall be presented using the \verb|listings| environment or
another environment defined by the official template.

Example:

\begin{verbatim}
\begin{lstlisting}[language=st]
IF Start = TRUE THEN
    Motor := TRUE;
ELSE
    Motor := FALSE;
END_IF;
\end{lstlisting}
\end{verbatim}

Only code necessary for explaining the work should appear in the main
report. Long source-code listings should be moved to an appendix or
submitted as supplementary files if required.

Source code should not normally be inserted as screenshots.

\vspace{0.3cm}

% ================================================================
% XVI. ABSTRACT
% ================================================================

\noindent
\textbf{XVI.\quad ABSTRACT}

\vspace{0.15cm}

The Abstract shall briefly summarize:

\begin{itemize}[leftmargin=0.8cm, itemsep=2pt]
    \item the internship organization;
    \item the objectives of the internship;
    \item the main tasks performed;
    \item the main results, knowledge, and skills obtained.
\end{itemize}

The Abstract should be concise and should not reproduce the full content
of the report chapter by chapter.

\vspace{0.3cm}

% ================================================================
% XVII. REFERENCES AND CITATIONS
% ================================================================

\noindent
\textbf{XVII.\quad REFERENCES AND CITATIONS}

\vspace{0.15cm}

All books, journal papers, conference papers, standards, websites,
datasheets, manuals, technical documents, figures, tables, or other
external sources used in preparing the report shall be properly cited.

References shall be managed through the bibliography system provided with
the official template.

Bibliographic entries shall normally be stored in:

\begin{verbatim}
References.bib
\end{verbatim}

References shall be cited in the report using:

\begin{verbatim}
\cite{reference-key}
\end{verbatim}

Students shall not manually type reference numbers such as ``[1]'',
``[2]'', etc.

Titles of foreign-language publications shall remain in their original
language in the bibliography.

All borrowed ideas, data, figures, tables, equations, methods, or results
shall be properly acknowledged.

\vspace{0.3cm}

% ================================================================
% XVIII. CROSS-REFERENCING
% ================================================================

\noindent
\textbf{XVIII.\quad CROSS-REFERENCING}

\vspace{0.15cm}

Numbered elements shall be cross-referenced using the standard \LaTeX{}
mechanism.

Use the following commands where appropriate:

\begin{verbatim}
\label{...}
\ref{...}
\eqref{...}
\end{verbatim}

For example:

\begin{verbatim}
As shown in Figure~\ref{fig:system}, ...
\end{verbatim}

Students shall not manually type numbers such as ``Figure 2.3'' when
those numbers can be generated automatically.

\vspace{0.3cm}

% ================================================================
% XIX. PAGE NUMBERING
% ================================================================

\noindent
\textbf{XIX.\quad PAGE NUMBERING}

\vspace{0.15cm}

Cover pages and official evaluation forms at the beginning of the report
shall not display page numbers.

The main body of the report shall use Arabic page numbers beginning from
Chapter 1.

Page numbering shall be generated automatically by the official template.
Students shall not manually enter page numbers.

\vspace{0.3cm}

% ================================================================
% XX. PROJECT FILE ORGANIZATION
% ================================================================

\noindent
\textbf{XX.\quad PROJECT FILE ORGANIZATION}

\vspace{0.15cm}

Students should maintain a clear and organized \LaTeX{} project
structure. A recommended structure is:

\begin{verbatim}
InternshipReport/
|
|-- main.tex
|-- References.bib
|
|-- figures/
|   |-- company.pdf
|   |-- system.png
|   `-- ...
|
`-- chapters/
    |-- chapter1.tex
    |-- chapter2.tex
    `-- chapter3.tex
\end{verbatim}

Figures should normally be stored in the \verb|figures| directory
specified by the official template.

File and directory names should preferably contain Latin letters,
numbers, hyphens, or underscores.

Recommended examples include:

\begin{verbatim}
company_structure.pdf
motor_control.png
chapter_01.tex
\end{verbatim}

Spaces and unnecessary special characters in file names should be
avoided where possible.

\vspace{0.3cm}

% ================================================================
% XXI. DOCUMENT COMPILATION AND FINAL CHECK
% ================================================================

\noindent
\textbf{XXI.\quad DOCUMENT COMPILATION AND FINAL CHECK}

\vspace{0.15cm}

Before submission, the complete \LaTeX{} project shall compile
successfully and produce the final PDF report.

The final document shall not contain unresolved references such as:

\begin{verbatim}
??
[?]
\end{verbatim}

Before generating the final PDF, students shall check:

\begin{itemize}[leftmargin=0.8cm, itemsep=2pt]
    \item the Table of Contents;
    \item page numbering;
    \item chapter, section, and subsection numbering;
    \item figure numbering;
    \item table numbering;
    \item equation numbering;
    \item cross-references;
    \item citations;
    \item bibliography;
    \item font display;
    \item mathematical symbols;
    \item image quality;
    \item completeness of all required forms.
\end{itemize}

The document may be compiled using TeX Live, MiKTeX, or another
compatible \LaTeX{} environment. Regardless of the compilation
environment, the resulting PDF shall conform to the official Faculty
template.

\vspace{0.3cm}

% ================================================================
% XXII. CONTENT THAT STUDENTS MAY EDIT
% ================================================================

\noindent
\textbf{XXII.\quad CONTENT THAT STUDENTS MAY EDIT}

\vspace{0.15cm}

Students may edit or replace:

\begin{itemize}[leftmargin=0.8cm, itemsep=2pt]
    \item student name;
    \item class;
    \item internship organization;
    \item supervisor name;
    \item Abstract;
    \item List of Symbols and Abbreviations;
    \item chapter titles and chapter content;
    \item figures and photographs;
    \item tables;
    \item equations;
    \item citations and references.
\end{itemize}

Students shall not modify the main document formatting settings in the
preamble, including the document class, page geometry, line spacing,
heading formatting, table-of-contents formatting, font configuration,
page layout, or cover-page macros unless specifically instructed.

\vspace{0.3cm}

% ================================================================
% XXIII. FINAL SUBMISSION REQUIREMENTS
% ================================================================

\noindent
\textbf{XXIII.\quad FINAL SUBMISSION REQUIREMENTS}

\vspace{0.15cm}

Before submission, the final report shall:

\begin{itemize}[leftmargin=0.8cm, itemsep=2pt]
    \item compile without errors that affect the report content or presentation;
    \item contain all required forms and report sections;
    \item contain correct student, class, internship organization, and supervisor information;
    \item contain an updated Table of Contents, List of Figures, and List of Tables;
    \item contain no unused example content from the template;
    \item contain no placeholder text such as \textit{Title of Chapter 1},
    \textit{Section 1.1}, or other sample entries;
    \item contain all required citations and references;
    \item contain no unresolved cross-references.
\end{itemize}

If the Faculty requires submission of the \LaTeX{} source, the complete
project shall include all \verb|.tex| files, the bibliography database,
figures, and any other files required to reproduce the submitted PDF.

\vspace{0.4cm}

\noindent
\textbf{Important:}
Before submitting the final internship report, students shall remove or
comment out the following command in the main document:

\begin{verbatim}
\input{Regulations.tex}
\end{verbatim}

For example:

\begin{verbatim}
% \input{Regulations.tex}
\end{verbatim}

This removes the guidance pages from the final submitted report while
preserving the regulations file in the \LaTeX{} project for reference.

} % End 13pt content

\clearpage
 command before submitting
% the final internship report.
%
% =================================================================

\clearpage
\thispagestyle{empty}

\begin{center}
    {\fontsize{16pt}{20pt}\selectfont\bfseries
    SPECIFIC REGULATIONS FOR PREPARING\\
    GRADUATION INTERNSHIP REPORTS USING \LaTeX}
\end{center}

\vspace{0.5cm}

{\fontsize{13pt}{17pt}\selectfont
\setstretch{1.3}
\justifying

\noindent
\textbf{Note:}
This section is provided for guidance only and is included in the main
document using the command \verb|
% =================================================================
% REGULATIONS FOR PREPARING THE GRADUATION INTERNSHIP REPORT
% USING LATEX
% =================================================================
%
% This file is intended to be included from the main document with:
%
%     \input{Regulations.tex}
%
% IMPORTANT:
% This section is provided for guidance only. Students shall remove
% or comment out the \input{Regulations.tex} command before submitting
% the final internship report.
%
% =================================================================

\clearpage
\thispagestyle{empty}

\begin{center}
    {\fontsize{16pt}{20pt}\selectfont\bfseries
    SPECIFIC REGULATIONS FOR PREPARING\\
    GRADUATION INTERNSHIP REPORTS USING \LaTeX}
\end{center}

\vspace{0.5cm}

{\fontsize{13pt}{17pt}\selectfont
\setstretch{1.3}
\justifying

\noindent
\textbf{Note:}
This section is provided for guidance only and is included in the main
document using the command \verb|\input{Regulations.tex}|.
Students shall remove or comment out this command before submitting the
final internship report.

\vspace{0.4cm}

% ================================================================
% I. GENERAL REQUIREMENTS
% ================================================================

\noindent
\textbf{I.\quad GENERAL REQUIREMENTS}

\vspace{0.15cm}

The Graduation Internship Report shall be prepared using the official
\LaTeX{} template provided by the Faculty. Students shall preserve the
predefined page layout, typography, line spacing, heading styles,
numbering system, figure and table formatting, and bibliography settings
unless otherwise instructed by the Faculty or the supervisor.

The report shall be written in English.

Students are responsible for the academic and technical content of the
report. The official \LaTeX{} template is responsible for standardized
document formatting.

\vspace{0.3cm}

% ================================================================
% II. REPORT STRUCTURE
% ================================================================

\noindent
\textbf{II.\quad REPORT STRUCTURE}

\vspace{0.15cm}

The report shall normally be arranged in the following order:

\begin{enumerate}[label=\arabic*., leftmargin=1.2cm, itemsep=2pt]
    \item Outer Cover;
    \item Comments of the Internship Host Organization;
    \item Supervisor's Evaluation and Graduation Internship Defense Results;
    \item Table of Contents;
    \item List of Symbols and Abbreviations, if applicable;
    \item List of Figures;
    \item List of Tables;
    \item Abstract;
    \item Main Chapters of the Internship Report;
    \item References;
    \item Appendices, if applicable.
\end{enumerate}

The preliminary pages shall not be numbered as chapters. The main report
content shall be organized into consecutively numbered chapters beginning
with Chapter 1.

\vspace{0.3cm}

% ================================================================
% III. PAGE SIZE AND MARGINS
% ================================================================

\noindent
\textbf{III.\quad PAGE SIZE AND MARGINS}

\vspace{0.15cm}

The report shall be formatted on A4 paper with the following margins:

\begin{center}
\begin{tabular}{ll}
    Left margin:   & 3.5 cm \\
    Right margin:  & 2.0 cm \\
    Top margin:    & 2.5 cm \\
    Bottom margin: & 2.5 cm
\end{tabular}
\end{center}

These settings are predefined in the official template and shall not be
modified by students.

\vspace{0.3cm}

% ================================================================
% IV. FONT, FONT SIZE, AND LINE SPACING
% ================================================================

\noindent
\textbf{IV.\quad FONT, FONT SIZE, AND LINE SPACING}

\vspace{0.15cm}

The main text shall use a 13 pt font size with 1.3-line spacing.
A Times-family typeface shall be used in accordance with the official
template.

The text shall be fully justified.

Students shall not use repeated spaces, repeated manual line breaks, or
multiple blank lines to position text. Paragraph and heading spacing shall
be controlled by the template.

\vspace{0.3cm}

% ================================================================
% V. OUTER COVER
% ================================================================

\noindent
\textbf{V.\quad OUTER COVER}

\vspace{0.15cm}

The outer cover shall follow the official format provided in the template
and shall contain the following information:

\begin{itemize}[leftmargin=0.8cm, itemsep=2pt]
    \item Thai Nguyen University of Technology;
    \item Faculty of International Training;
    \item Engineering and Technology Department;
    \item University logo and Faculty logo;
    \item the title \textbf{INTERNSHIP REPORT};
    \item Internship Organization;
    \item Author;
    \item Class;
    \item Supervisor;
    \item place and year of completion.
\end{itemize}

Students shall only replace the required information fields and shall not
modify the cover-page layout unless instructed otherwise.

\vspace{0.3cm}

% ================================================================
% VI. COMMENTS OF THE INTERNSHIP HOST ORGANIZATION
% ================================================================

\noindent
\textbf{VI.\quad COMMENTS OF THE INTERNSHIP HOST ORGANIZATION}

\vspace{0.15cm}

The page entitled \textbf{COMMENTS OF THE INTERNSHIP HOST ORGANIZATION}
is reserved for the internship host organization.

It includes:

\begin{itemize}[leftmargin=0.8cm, itemsep=2pt]
    \item the student's full name;
    \item class;
    \item comments from the host organization;
    \item score;
    \item date;
    \item signature and official seal.
\end{itemize}

Students shall not alter or complete official evaluation fields unless
specifically authorized.

\vspace{0.3cm}

% ================================================================
% VII. SUPERVISOR'S EVALUATION AND DEFENSE RESULTS
% ================================================================

\noindent
\textbf{VII.\quad SUPERVISOR'S EVALUATION AND DEFENSE RESULTS}

\vspace{0.15cm}

The section entitled
\textbf{SUPERVISOR'S EVALUATION AND GRADUATION INTERNSHIP DEFENSE RESULTS}
is reserved for the supervisor and defense committee.

The evaluation may include:

\begin{itemize}[leftmargin=0.8cm, itemsep=2pt]
    \item Internship Progress and Student's Attitude;
    \item Report Content;
    \item Report Presentation;
    \item Other Comments;
    \item Supervisor's Opinion;
    \item Defense Result;
    \item Score.
\end{itemize}

Students shall not enter evaluation scores or defense results before they
are officially determined.

\vspace{0.3cm}

% ================================================================
% VIII. CHAPTERS, SECTIONS, AND SUBSECTIONS
% ================================================================

\noindent
\textbf{VIII.\quad CHAPTERS, SECTIONS, AND SUBSECTIONS}

\vspace{0.15cm}

The report structure shall be generated using the standard \LaTeX{}
commands:

\begin{verbatim}
\chapter{...}
\section{...}
\subsection{...}
\end{verbatim}

Chapter, section, and subsection numbers shall be generated automatically
by \LaTeX{}.

Students shall not manually type chapter or section numbers and shall not
reproduce heading styles using commands such as \verb|\textbf|,
\verb|\fontsize|, or \verb|\centering| instead of the appropriate
structural command.

Excessive subdivision should be avoided. In most cases, the hierarchy
Chapter--Section--Subsection is sufficient.

\vspace{0.3cm}

% ================================================================
% IX. ORGANIZATION OF THE MAIN CONTENT
% ================================================================

\noindent
\textbf{IX.\quad ORGANIZATION OF THE MAIN CONTENT}

\vspace{0.15cm}

The specific content of the report shall depend on the internship
organization and the assigned tasks. A typical structure may include:

\begin{itemize}[leftmargin=0.8cm, itemsep=2pt]
    \item \textbf{Chapter 1. Overview of the Internship Organization};
    \item \textbf{Chapter 2. Internship Activities and Assigned Tasks};
    \item \textbf{Chapter 3. Results, Experience, and Evaluation}.
\end{itemize}

These chapter titles are recommendations only and may be adapted to the
actual internship content.

The report shall reflect work actually performed, observed, analyzed, or
learned during the internship. Unrelated theoretical material should not
be included merely to increase report length.

\vspace{0.3cm}

% ================================================================
% X. TABLE OF CONTENTS
% ================================================================

\noindent
\textbf{X.\quad TABLE OF CONTENTS}

\vspace{0.15cm}

The Table of Contents shall be generated automatically using:

\begin{verbatim}
\tableofcontents
\end{verbatim}

Students shall not create or manually edit page numbers in the Table of
Contents.

The document shall be compiled sufficiently to ensure that all page
numbers and entries are correctly updated.

\vspace{0.3cm}

% ================================================================
% XI. LIST OF SYMBOLS AND ABBREVIATIONS
% ================================================================

\noindent
\textbf{XI.\quad LIST OF SYMBOLS AND ABBREVIATIONS}

\vspace{0.15cm}

Symbols and abbreviations that are used repeatedly in the report shall be
listed in the \textbf{List of Symbols and Abbreviations}.

The list may contain the following columns:

\begin{center}
\begin{tabular}{llll}
    No. & Symbol / Abbreviation & Unit & Definition
\end{tabular}
\end{center}

Only symbols and abbreviations actually used in the final report shall be
retained. Example entries contained in the template shall be removed if
they are not relevant to the student's report.

Each abbreviation should normally be defined when it first appears in
the main text.

\vspace{0.3cm}

% ================================================================
% XII. FIGURES, DIAGRAMS, AND PHOTOGRAPHS
% ================================================================

\noindent
\textbf{XII.\quad FIGURES, DIAGRAMS, AND PHOTOGRAPHS}

\vspace{0.15cm}

Figures shall be clear, legible, and directly relevant to the report
content.

Figures shall be inserted using the standard \verb|figure| environment.

Example:

\begin{verbatim}
\begin{figure}[H]
    \centering
    \includegraphics[width=0.8\textwidth]
                    {figures/company_structure.pdf}
    \caption{Organizational structure of the company}
    \label{fig:company-structure}
\end{figure}
\end{verbatim}

Figure captions shall be placed below the figure.

Each referenced figure shall have a unique \verb|\label| and shall be
referenced in the text using, for example:

\begin{verbatim}
Figure~\ref{fig:company-structure}
\end{verbatim}

Students shall not manually type figure numbers.

Figures reproduced or adapted from external sources shall include an
appropriate citation.

Photographs taken during the internship shall be relevant and sufficiently
clear. Students shall not include confidential, restricted, or proprietary
information without permission from the internship organization.

Vector graphics or high-resolution images are recommended whenever
possible.

\vspace{0.3cm}

% ================================================================
% XIII. TABLES
% ================================================================

\noindent
\textbf{XIII.\quad TABLES}

\vspace{0.15cm}

Tables shall be created using suitable \LaTeX{} environments such as
\verb|table|, \verb|tabular|, \verb|tabularx|, or \verb|longtable|.

Table captions shall be placed above the table.

Example:

\begin{verbatim}
\begin{table}[H]
    \centering
    \caption{Internship schedule}
    \label{tab:internship-schedule}

    \begin{tabular}{lll}
        \toprule
        Week & Task & Result \\
        \midrule
        1 & Orientation & Completed \\
        2 & System study & Completed \\
        \bottomrule
    \end{tabular}
\end{table}
\end{verbatim}

Tables shall be referenced using:

\begin{verbatim}
Table~\ref{tab:internship-schedule}
\end{verbatim}

Students shall not manually type table numbers.

Tables should be created directly in \LaTeX{} rather than inserted as
screenshots whenever practical.

\vspace{0.3cm}

% ================================================================
% XIV. EQUATIONS AND TECHNICAL QUANTITIES
% ================================================================

\noindent
\textbf{XIV.\quad EQUATIONS AND TECHNICAL QUANTITIES}

\vspace{0.15cm}

Mathematical expressions shall be typeset directly using \LaTeX{}.
Images or screenshots of equations shall not be used.

Short expressions may be written in inline mathematical mode, for example:

\begin{verbatim}
$R = 0.7\,\Omega$
\end{verbatim}

Important equations shall use a numbered equation environment.

Example:

\begin{verbatim}
\begin{equation}
    P = UI
    \label{eq:power}
\end{equation}
\end{verbatim}

Equations shall be referenced using:

\begin{verbatim}
Equation~\eqref{eq:power}
\end{verbatim}

Equation numbers shall not be manually entered.

SI units shall be used wherever applicable. The \verb|siunitx| package
should be used where appropriate.

Examples:

\begin{verbatim}
\SI{220}{\volt}
\SI{50}{\hertz}
\SI{2.5}{\ampere}
\end{verbatim}

Notation and units shall be consistent throughout the report.

\vspace{0.3cm}

% ================================================================
% XV. SOURCE CODE
% ================================================================

\noindent
\textbf{XV.\quad SOURCE CODE}

\vspace{0.15cm}

If the internship involves programming, PLCs, microcontrollers,
simulation, or software development, relevant source code may be included
in the report.

Source code shall be presented using the \verb|listings| environment or
another environment defined by the official template.

Example:

\begin{verbatim}
\begin{lstlisting}[language=st]
IF Start = TRUE THEN
    Motor := TRUE;
ELSE
    Motor := FALSE;
END_IF;
\end{lstlisting}
\end{verbatim}

Only code necessary for explaining the work should appear in the main
report. Long source-code listings should be moved to an appendix or
submitted as supplementary files if required.

Source code should not normally be inserted as screenshots.

\vspace{0.3cm}

% ================================================================
% XVI. ABSTRACT
% ================================================================

\noindent
\textbf{XVI.\quad ABSTRACT}

\vspace{0.15cm}

The Abstract shall briefly summarize:

\begin{itemize}[leftmargin=0.8cm, itemsep=2pt]
    \item the internship organization;
    \item the objectives of the internship;
    \item the main tasks performed;
    \item the main results, knowledge, and skills obtained.
\end{itemize}

The Abstract should be concise and should not reproduce the full content
of the report chapter by chapter.

\vspace{0.3cm}

% ================================================================
% XVII. REFERENCES AND CITATIONS
% ================================================================

\noindent
\textbf{XVII.\quad REFERENCES AND CITATIONS}

\vspace{0.15cm}

All books, journal papers, conference papers, standards, websites,
datasheets, manuals, technical documents, figures, tables, or other
external sources used in preparing the report shall be properly cited.

References shall be managed through the bibliography system provided with
the official template.

Bibliographic entries shall normally be stored in:

\begin{verbatim}
References.bib
\end{verbatim}

References shall be cited in the report using:

\begin{verbatim}
\cite{reference-key}
\end{verbatim}

Students shall not manually type reference numbers such as ``[1]'',
``[2]'', etc.

Titles of foreign-language publications shall remain in their original
language in the bibliography.

All borrowed ideas, data, figures, tables, equations, methods, or results
shall be properly acknowledged.

\vspace{0.3cm}

% ================================================================
% XVIII. CROSS-REFERENCING
% ================================================================

\noindent
\textbf{XVIII.\quad CROSS-REFERENCING}

\vspace{0.15cm}

Numbered elements shall be cross-referenced using the standard \LaTeX{}
mechanism.

Use the following commands where appropriate:

\begin{verbatim}
\label{...}
\ref{...}
\eqref{...}
\end{verbatim}

For example:

\begin{verbatim}
As shown in Figure~\ref{fig:system}, ...
\end{verbatim}

Students shall not manually type numbers such as ``Figure 2.3'' when
those numbers can be generated automatically.

\vspace{0.3cm}

% ================================================================
% XIX. PAGE NUMBERING
% ================================================================

\noindent
\textbf{XIX.\quad PAGE NUMBERING}

\vspace{0.15cm}

Cover pages and official evaluation forms at the beginning of the report
shall not display page numbers.

The main body of the report shall use Arabic page numbers beginning from
Chapter 1.

Page numbering shall be generated automatically by the official template.
Students shall not manually enter page numbers.

\vspace{0.3cm}

% ================================================================
% XX. PROJECT FILE ORGANIZATION
% ================================================================

\noindent
\textbf{XX.\quad PROJECT FILE ORGANIZATION}

\vspace{0.15cm}

Students should maintain a clear and organized \LaTeX{} project
structure. A recommended structure is:

\begin{verbatim}
InternshipReport/
|
|-- main.tex
|-- References.bib
|
|-- figures/
|   |-- company.pdf
|   |-- system.png
|   `-- ...
|
`-- chapters/
    |-- chapter1.tex
    |-- chapter2.tex
    `-- chapter3.tex
\end{verbatim}

Figures should normally be stored in the \verb|figures| directory
specified by the official template.

File and directory names should preferably contain Latin letters,
numbers, hyphens, or underscores.

Recommended examples include:

\begin{verbatim}
company_structure.pdf
motor_control.png
chapter_01.tex
\end{verbatim}

Spaces and unnecessary special characters in file names should be
avoided where possible.

\vspace{0.3cm}

% ================================================================
% XXI. DOCUMENT COMPILATION AND FINAL CHECK
% ================================================================

\noindent
\textbf{XXI.\quad DOCUMENT COMPILATION AND FINAL CHECK}

\vspace{0.15cm}

Before submission, the complete \LaTeX{} project shall compile
successfully and produce the final PDF report.

The final document shall not contain unresolved references such as:

\begin{verbatim}
??
[?]
\end{verbatim}

Before generating the final PDF, students shall check:

\begin{itemize}[leftmargin=0.8cm, itemsep=2pt]
    \item the Table of Contents;
    \item page numbering;
    \item chapter, section, and subsection numbering;
    \item figure numbering;
    \item table numbering;
    \item equation numbering;
    \item cross-references;
    \item citations;
    \item bibliography;
    \item font display;
    \item mathematical symbols;
    \item image quality;
    \item completeness of all required forms.
\end{itemize}

The document may be compiled using TeX Live, MiKTeX, or another
compatible \LaTeX{} environment. Regardless of the compilation
environment, the resulting PDF shall conform to the official Faculty
template.

\vspace{0.3cm}

% ================================================================
% XXII. CONTENT THAT STUDENTS MAY EDIT
% ================================================================

\noindent
\textbf{XXII.\quad CONTENT THAT STUDENTS MAY EDIT}

\vspace{0.15cm}

Students may edit or replace:

\begin{itemize}[leftmargin=0.8cm, itemsep=2pt]
    \item student name;
    \item class;
    \item internship organization;
    \item supervisor name;
    \item Abstract;
    \item List of Symbols and Abbreviations;
    \item chapter titles and chapter content;
    \item figures and photographs;
    \item tables;
    \item equations;
    \item citations and references.
\end{itemize}

Students shall not modify the main document formatting settings in the
preamble, including the document class, page geometry, line spacing,
heading formatting, table-of-contents formatting, font configuration,
page layout, or cover-page macros unless specifically instructed.

\vspace{0.3cm}

% ================================================================
% XXIII. FINAL SUBMISSION REQUIREMENTS
% ================================================================

\noindent
\textbf{XXIII.\quad FINAL SUBMISSION REQUIREMENTS}

\vspace{0.15cm}

Before submission, the final report shall:

\begin{itemize}[leftmargin=0.8cm, itemsep=2pt]
    \item compile without errors that affect the report content or presentation;
    \item contain all required forms and report sections;
    \item contain correct student, class, internship organization, and supervisor information;
    \item contain an updated Table of Contents, List of Figures, and List of Tables;
    \item contain no unused example content from the template;
    \item contain no placeholder text such as \textit{Title of Chapter 1},
    \textit{Section 1.1}, or other sample entries;
    \item contain all required citations and references;
    \item contain no unresolved cross-references.
\end{itemize}

If the Faculty requires submission of the \LaTeX{} source, the complete
project shall include all \verb|.tex| files, the bibliography database,
figures, and any other files required to reproduce the submitted PDF.

\vspace{0.4cm}

\noindent
\textbf{Important:}
Before submitting the final internship report, students shall remove or
comment out the following command in the main document:

\begin{verbatim}
\input{Regulations.tex}
\end{verbatim}

For example:

\begin{verbatim}
% \input{Regulations.tex}
\end{verbatim}

This removes the guidance pages from the final submitted report while
preserving the regulations file in the \LaTeX{} project for reference.

} % End 13pt content

\clearpage
|.
Students shall remove or comment out this command before submitting the
final internship report.

\vspace{0.4cm}

% ================================================================
% I. GENERAL REQUIREMENTS
% ================================================================

\noindent
\textbf{I.\quad GENERAL REQUIREMENTS}

\vspace{0.15cm}

The Graduation Internship Report shall be prepared using the official
\LaTeX{} template provided by the Faculty. Students shall preserve the
predefined page layout, typography, line spacing, heading styles,
numbering system, figure and table formatting, and bibliography settings
unless otherwise instructed by the Faculty or the supervisor.

The report shall be written in English.

Students are responsible for the academic and technical content of the
report. The official \LaTeX{} template is responsible for standardized
document formatting.

\vspace{0.3cm}

% ================================================================
% II. REPORT STRUCTURE
% ================================================================

\noindent
\textbf{II.\quad REPORT STRUCTURE}

\vspace{0.15cm}

The report shall normally be arranged in the following order:

\begin{enumerate}[label=\arabic*., leftmargin=1.2cm, itemsep=2pt]
    \item Outer Cover;
    \item Comments of the Internship Host Organization;
    \item Supervisor's Evaluation and Graduation Internship Defense Results;
    \item Table of Contents;
    \item List of Symbols and Abbreviations, if applicable;
    \item List of Figures;
    \item List of Tables;
    \item Abstract;
    \item Main Chapters of the Internship Report;
    \item References;
    \item Appendices, if applicable.
\end{enumerate}

The preliminary pages shall not be numbered as chapters. The main report
content shall be organized into consecutively numbered chapters beginning
with Chapter 1.

\vspace{0.3cm}

% ================================================================
% III. PAGE SIZE AND MARGINS
% ================================================================

\noindent
\textbf{III.\quad PAGE SIZE AND MARGINS}

\vspace{0.15cm}

The report shall be formatted on A4 paper with the following margins:

\begin{center}
\begin{tabular}{ll}
    Left margin:   & 3.5 cm \\
    Right margin:  & 2.0 cm \\
    Top margin:    & 2.5 cm \\
    Bottom margin: & 2.5 cm
\end{tabular}
\end{center}

These settings are predefined in the official template and shall not be
modified by students.

\vspace{0.3cm}

% ================================================================
% IV. FONT, FONT SIZE, AND LINE SPACING
% ================================================================

\noindent
\textbf{IV.\quad FONT, FONT SIZE, AND LINE SPACING}

\vspace{0.15cm}

The main text shall use a 13 pt font size with 1.3-line spacing.
A Times-family typeface shall be used in accordance with the official
template.

The text shall be fully justified.

Students shall not use repeated spaces, repeated manual line breaks, or
multiple blank lines to position text. Paragraph and heading spacing shall
be controlled by the template.

\vspace{0.3cm}

% ================================================================
% V. OUTER COVER
% ================================================================

\noindent
\textbf{V.\quad OUTER COVER}

\vspace{0.15cm}

The outer cover shall follow the official format provided in the template
and shall contain the following information:

\begin{itemize}[leftmargin=0.8cm, itemsep=2pt]
    \item Thai Nguyen University of Technology;
    \item Faculty of International Training;
    \item Engineering and Technology Department;
    \item University logo and Faculty logo;
    \item the title \textbf{INTERNSHIP REPORT};
    \item Internship Organization;
    \item Author;
    \item Class;
    \item Supervisor;
    \item place and year of completion.
\end{itemize}

Students shall only replace the required information fields and shall not
modify the cover-page layout unless instructed otherwise.

\vspace{0.3cm}

% ================================================================
% VI. COMMENTS OF THE INTERNSHIP HOST ORGANIZATION
% ================================================================

\noindent
\textbf{VI.\quad COMMENTS OF THE INTERNSHIP HOST ORGANIZATION}

\vspace{0.15cm}

The page entitled \textbf{COMMENTS OF THE INTERNSHIP HOST ORGANIZATION}
is reserved for the internship host organization.

It includes:

\begin{itemize}[leftmargin=0.8cm, itemsep=2pt]
    \item the student's full name;
    \item class;
    \item comments from the host organization;
    \item score;
    \item date;
    \item signature and official seal.
\end{itemize}

Students shall not alter or complete official evaluation fields unless
specifically authorized.

\vspace{0.3cm}

% ================================================================
% VII. SUPERVISOR'S EVALUATION AND DEFENSE RESULTS
% ================================================================

\noindent
\textbf{VII.\quad SUPERVISOR'S EVALUATION AND DEFENSE RESULTS}

\vspace{0.15cm}

The section entitled
\textbf{SUPERVISOR'S EVALUATION AND GRADUATION INTERNSHIP DEFENSE RESULTS}
is reserved for the supervisor and defense committee.

The evaluation may include:

\begin{itemize}[leftmargin=0.8cm, itemsep=2pt]
    \item Internship Progress and Student's Attitude;
    \item Report Content;
    \item Report Presentation;
    \item Other Comments;
    \item Supervisor's Opinion;
    \item Defense Result;
    \item Score.
\end{itemize}

Students shall not enter evaluation scores or defense results before they
are officially determined.

\vspace{0.3cm}

% ================================================================
% VIII. CHAPTERS, SECTIONS, AND SUBSECTIONS
% ================================================================

\noindent
\textbf{VIII.\quad CHAPTERS, SECTIONS, AND SUBSECTIONS}

\vspace{0.15cm}

The report structure shall be generated using the standard \LaTeX{}
commands:

\begin{verbatim}
\chapter{...}
\section{...}
\subsection{...}
\end{verbatim}

Chapter, section, and subsection numbers shall be generated automatically
by \LaTeX{}.

Students shall not manually type chapter or section numbers and shall not
reproduce heading styles using commands such as \verb|\textbf|,
\verb|\fontsize|, or \verb|\centering| instead of the appropriate
structural command.

Excessive subdivision should be avoided. In most cases, the hierarchy
Chapter--Section--Subsection is sufficient.

\vspace{0.3cm}

% ================================================================
% IX. ORGANIZATION OF THE MAIN CONTENT
% ================================================================

\noindent
\textbf{IX.\quad ORGANIZATION OF THE MAIN CONTENT}

\vspace{0.15cm}

The specific content of the report shall depend on the internship
organization and the assigned tasks. A typical structure may include:

\begin{itemize}[leftmargin=0.8cm, itemsep=2pt]
    \item \textbf{Chapter 1. Overview of the Internship Organization};
    \item \textbf{Chapter 2. Internship Activities and Assigned Tasks};
    \item \textbf{Chapter 3. Results, Experience, and Evaluation}.
\end{itemize}

These chapter titles are recommendations only and may be adapted to the
actual internship content.

The report shall reflect work actually performed, observed, analyzed, or
learned during the internship. Unrelated theoretical material should not
be included merely to increase report length.

\vspace{0.3cm}

% ================================================================
% X. TABLE OF CONTENTS
% ================================================================

\noindent
\textbf{X.\quad TABLE OF CONTENTS}

\vspace{0.15cm}

The Table of Contents shall be generated automatically using:

\begin{verbatim}
\tableofcontents
\end{verbatim}

Students shall not create or manually edit page numbers in the Table of
Contents.

The document shall be compiled sufficiently to ensure that all page
numbers and entries are correctly updated.

\vspace{0.3cm}

% ================================================================
% XI. LIST OF SYMBOLS AND ABBREVIATIONS
% ================================================================

\noindent
\textbf{XI.\quad LIST OF SYMBOLS AND ABBREVIATIONS}

\vspace{0.15cm}

Symbols and abbreviations that are used repeatedly in the report shall be
listed in the \textbf{List of Symbols and Abbreviations}.

The list may contain the following columns:

\begin{center}
\begin{tabular}{llll}
    No. & Symbol / Abbreviation & Unit & Definition
\end{tabular}
\end{center}

Only symbols and abbreviations actually used in the final report shall be
retained. Example entries contained in the template shall be removed if
they are not relevant to the student's report.

Each abbreviation should normally be defined when it first appears in
the main text.

\vspace{0.3cm}

% ================================================================
% XII. FIGURES, DIAGRAMS, AND PHOTOGRAPHS
% ================================================================

\noindent
\textbf{XII.\quad FIGURES, DIAGRAMS, AND PHOTOGRAPHS}

\vspace{0.15cm}

Figures shall be clear, legible, and directly relevant to the report
content.

Figures shall be inserted using the standard \verb|figure| environment.

Example:

\begin{verbatim}
\begin{figure}[H]
    \centering
    \includegraphics[width=0.8\textwidth]
                    {figures/company_structure.pdf}
    \caption{Organizational structure of the company}
    \label{fig:company-structure}
\end{figure}
\end{verbatim}

Figure captions shall be placed below the figure.

Each referenced figure shall have a unique \verb|\label| and shall be
referenced in the text using, for example:

\begin{verbatim}
Figure~\ref{fig:company-structure}
\end{verbatim}

Students shall not manually type figure numbers.

Figures reproduced or adapted from external sources shall include an
appropriate citation.

Photographs taken during the internship shall be relevant and sufficiently
clear. Students shall not include confidential, restricted, or proprietary
information without permission from the internship organization.

Vector graphics or high-resolution images are recommended whenever
possible.

\vspace{0.3cm}

% ================================================================
% XIII. TABLES
% ================================================================

\noindent
\textbf{XIII.\quad TABLES}

\vspace{0.15cm}

Tables shall be created using suitable \LaTeX{} environments such as
\verb|table|, \verb|tabular|, \verb|tabularx|, or \verb|longtable|.

Table captions shall be placed above the table.

Example:

\begin{verbatim}
\begin{table}[H]
    \centering
    \caption{Internship schedule}
    \label{tab:internship-schedule}

    \begin{tabular}{lll}
        \toprule
        Week & Task & Result \\
        \midrule
        1 & Orientation & Completed \\
        2 & System study & Completed \\
        \bottomrule
    \end{tabular}
\end{table}
\end{verbatim}

Tables shall be referenced using:

\begin{verbatim}
Table~\ref{tab:internship-schedule}
\end{verbatim}

Students shall not manually type table numbers.

Tables should be created directly in \LaTeX{} rather than inserted as
screenshots whenever practical.

\vspace{0.3cm}

% ================================================================
% XIV. EQUATIONS AND TECHNICAL QUANTITIES
% ================================================================

\noindent
\textbf{XIV.\quad EQUATIONS AND TECHNICAL QUANTITIES}

\vspace{0.15cm}

Mathematical expressions shall be typeset directly using \LaTeX{}.
Images or screenshots of equations shall not be used.

Short expressions may be written in inline mathematical mode, for example:

\begin{verbatim}
$R = 0.7\,\Omega$
\end{verbatim}

Important equations shall use a numbered equation environment.

Example:

\begin{verbatim}
\begin{equation}
    P = UI
    \label{eq:power}
\end{equation}
\end{verbatim}

Equations shall be referenced using:

\begin{verbatim}
Equation~\eqref{eq:power}
\end{verbatim}

Equation numbers shall not be manually entered.

SI units shall be used wherever applicable. The \verb|siunitx| package
should be used where appropriate.

Examples:

\begin{verbatim}
\SI{220}{\volt}
\SI{50}{\hertz}
\SI{2.5}{\ampere}
\end{verbatim}

Notation and units shall be consistent throughout the report.

\vspace{0.3cm}

% ================================================================
% XV. SOURCE CODE
% ================================================================

\noindent
\textbf{XV.\quad SOURCE CODE}

\vspace{0.15cm}

If the internship involves programming, PLCs, microcontrollers,
simulation, or software development, relevant source code may be included
in the report.

Source code shall be presented using the \verb|listings| environment or
another environment defined by the official template.

Example:

\begin{verbatim}
\begin{lstlisting}[language=st]
IF Start = TRUE THEN
    Motor := TRUE;
ELSE
    Motor := FALSE;
END_IF;
\end{lstlisting}
\end{verbatim}

Only code necessary for explaining the work should appear in the main
report. Long source-code listings should be moved to an appendix or
submitted as supplementary files if required.

Source code should not normally be inserted as screenshots.

\vspace{0.3cm}

% ================================================================
% XVI. ABSTRACT
% ================================================================

\noindent
\textbf{XVI.\quad ABSTRACT}

\vspace{0.15cm}

The Abstract shall briefly summarize:

\begin{itemize}[leftmargin=0.8cm, itemsep=2pt]
    \item the internship organization;
    \item the objectives of the internship;
    \item the main tasks performed;
    \item the main results, knowledge, and skills obtained.
\end{itemize}

The Abstract should be concise and should not reproduce the full content
of the report chapter by chapter.

\vspace{0.3cm}

% ================================================================
% XVII. REFERENCES AND CITATIONS
% ================================================================

\noindent
\textbf{XVII.\quad REFERENCES AND CITATIONS}

\vspace{0.15cm}

All books, journal papers, conference papers, standards, websites,
datasheets, manuals, technical documents, figures, tables, or other
external sources used in preparing the report shall be properly cited.

References shall be managed through the bibliography system provided with
the official template.

Bibliographic entries shall normally be stored in:

\begin{verbatim}
References.bib
\end{verbatim}

References shall be cited in the report using:

\begin{verbatim}
\cite{reference-key}
\end{verbatim}

Students shall not manually type reference numbers such as ``[1]'',
``[2]'', etc.

Titles of foreign-language publications shall remain in their original
language in the bibliography.

All borrowed ideas, data, figures, tables, equations, methods, or results
shall be properly acknowledged.

\vspace{0.3cm}

% ================================================================
% XVIII. CROSS-REFERENCING
% ================================================================

\noindent
\textbf{XVIII.\quad CROSS-REFERENCING}

\vspace{0.15cm}

Numbered elements shall be cross-referenced using the standard \LaTeX{}
mechanism.

Use the following commands where appropriate:

\begin{verbatim}
\label{...}
\ref{...}
\eqref{...}
\end{verbatim}

For example:

\begin{verbatim}
As shown in Figure~\ref{fig:system}, ...
\end{verbatim}

Students shall not manually type numbers such as ``Figure 2.3'' when
those numbers can be generated automatically.

\vspace{0.3cm}

% ================================================================
% XIX. PAGE NUMBERING
% ================================================================

\noindent
\textbf{XIX.\quad PAGE NUMBERING}

\vspace{0.15cm}

Cover pages and official evaluation forms at the beginning of the report
shall not display page numbers.

The main body of the report shall use Arabic page numbers beginning from
Chapter 1.

Page numbering shall be generated automatically by the official template.
Students shall not manually enter page numbers.

\vspace{0.3cm}

% ================================================================
% XX. PROJECT FILE ORGANIZATION
% ================================================================

\noindent
\textbf{XX.\quad PROJECT FILE ORGANIZATION}

\vspace{0.15cm}

Students should maintain a clear and organized \LaTeX{} project
structure. A recommended structure is:

\begin{verbatim}
InternshipReport/
|
|-- main.tex
|-- References.bib
|
|-- figures/
|   |-- company.pdf
|   |-- system.png
|   `-- ...
|
`-- chapters/
    |-- chapter1.tex
    |-- chapter2.tex
    `-- chapter3.tex
\end{verbatim}

Figures should normally be stored in the \verb|figures| directory
specified by the official template.

File and directory names should preferably contain Latin letters,
numbers, hyphens, or underscores.

Recommended examples include:

\begin{verbatim}
company_structure.pdf
motor_control.png
chapter_01.tex
\end{verbatim}

Spaces and unnecessary special characters in file names should be
avoided where possible.

\vspace{0.3cm}

% ================================================================
% XXI. DOCUMENT COMPILATION AND FINAL CHECK
% ================================================================

\noindent
\textbf{XXI.\quad DOCUMENT COMPILATION AND FINAL CHECK}

\vspace{0.15cm}

Before submission, the complete \LaTeX{} project shall compile
successfully and produce the final PDF report.

The final document shall not contain unresolved references such as:

\begin{verbatim}
??
[?]
\end{verbatim}

Before generating the final PDF, students shall check:

\begin{itemize}[leftmargin=0.8cm, itemsep=2pt]
    \item the Table of Contents;
    \item page numbering;
    \item chapter, section, and subsection numbering;
    \item figure numbering;
    \item table numbering;
    \item equation numbering;
    \item cross-references;
    \item citations;
    \item bibliography;
    \item font display;
    \item mathematical symbols;
    \item image quality;
    \item completeness of all required forms.
\end{itemize}

The document may be compiled using TeX Live, MiKTeX, or another
compatible \LaTeX{} environment. Regardless of the compilation
environment, the resulting PDF shall conform to the official Faculty
template.

\vspace{0.3cm}

% ================================================================
% XXII. CONTENT THAT STUDENTS MAY EDIT
% ================================================================

\noindent
\textbf{XXII.\quad CONTENT THAT STUDENTS MAY EDIT}

\vspace{0.15cm}

Students may edit or replace:

\begin{itemize}[leftmargin=0.8cm, itemsep=2pt]
    \item student name;
    \item class;
    \item internship organization;
    \item supervisor name;
    \item Abstract;
    \item List of Symbols and Abbreviations;
    \item chapter titles and chapter content;
    \item figures and photographs;
    \item tables;
    \item equations;
    \item citations and references.
\end{itemize}

Students shall not modify the main document formatting settings in the
preamble, including the document class, page geometry, line spacing,
heading formatting, table-of-contents formatting, font configuration,
page layout, or cover-page macros unless specifically instructed.

\vspace{0.3cm}

% ================================================================
% XXIII. FINAL SUBMISSION REQUIREMENTS
% ================================================================

\noindent
\textbf{XXIII.\quad FINAL SUBMISSION REQUIREMENTS}

\vspace{0.15cm}

Before submission, the final report shall:

\begin{itemize}[leftmargin=0.8cm, itemsep=2pt]
    \item compile without errors that affect the report content or presentation;
    \item contain all required forms and report sections;
    \item contain correct student, class, internship organization, and supervisor information;
    \item contain an updated Table of Contents, List of Figures, and List of Tables;
    \item contain no unused example content from the template;
    \item contain no placeholder text such as \textit{Title of Chapter 1},
    \textit{Section 1.1}, or other sample entries;
    \item contain all required citations and references;
    \item contain no unresolved cross-references.
\end{itemize}

If the Faculty requires submission of the \LaTeX{} source, the complete
project shall include all \verb|.tex| files, the bibliography database,
figures, and any other files required to reproduce the submitted PDF.

\vspace{0.4cm}

\noindent
\textbf{Important:}
Before submitting the final internship report, students shall remove or
comment out the following command in the main document:

\begin{verbatim}

% =================================================================
% REGULATIONS FOR PREPARING THE GRADUATION INTERNSHIP REPORT
% USING LATEX
% =================================================================
%
% This file is intended to be included from the main document with:
%
%     \input{Regulations.tex}
%
% IMPORTANT:
% This section is provided for guidance only. Students shall remove
% or comment out the \input{Regulations.tex} command before submitting
% the final internship report.
%
% =================================================================

\clearpage
\thispagestyle{empty}

\begin{center}
    {\fontsize{16pt}{20pt}\selectfont\bfseries
    SPECIFIC REGULATIONS FOR PREPARING\\
    GRADUATION INTERNSHIP REPORTS USING \LaTeX}
\end{center}

\vspace{0.5cm}

{\fontsize{13pt}{17pt}\selectfont
\setstretch{1.3}
\justifying

\noindent
\textbf{Note:}
This section is provided for guidance only and is included in the main
document using the command \verb|\input{Regulations.tex}|.
Students shall remove or comment out this command before submitting the
final internship report.

\vspace{0.4cm}

% ================================================================
% I. GENERAL REQUIREMENTS
% ================================================================

\noindent
\textbf{I.\quad GENERAL REQUIREMENTS}

\vspace{0.15cm}

The Graduation Internship Report shall be prepared using the official
\LaTeX{} template provided by the Faculty. Students shall preserve the
predefined page layout, typography, line spacing, heading styles,
numbering system, figure and table formatting, and bibliography settings
unless otherwise instructed by the Faculty or the supervisor.

The report shall be written in English.

Students are responsible for the academic and technical content of the
report. The official \LaTeX{} template is responsible for standardized
document formatting.

\vspace{0.3cm}

% ================================================================
% II. REPORT STRUCTURE
% ================================================================

\noindent
\textbf{II.\quad REPORT STRUCTURE}

\vspace{0.15cm}

The report shall normally be arranged in the following order:

\begin{enumerate}[label=\arabic*., leftmargin=1.2cm, itemsep=2pt]
    \item Outer Cover;
    \item Comments of the Internship Host Organization;
    \item Supervisor's Evaluation and Graduation Internship Defense Results;
    \item Table of Contents;
    \item List of Symbols and Abbreviations, if applicable;
    \item List of Figures;
    \item List of Tables;
    \item Abstract;
    \item Main Chapters of the Internship Report;
    \item References;
    \item Appendices, if applicable.
\end{enumerate}

The preliminary pages shall not be numbered as chapters. The main report
content shall be organized into consecutively numbered chapters beginning
with Chapter 1.

\vspace{0.3cm}

% ================================================================
% III. PAGE SIZE AND MARGINS
% ================================================================

\noindent
\textbf{III.\quad PAGE SIZE AND MARGINS}

\vspace{0.15cm}

The report shall be formatted on A4 paper with the following margins:

\begin{center}
\begin{tabular}{ll}
    Left margin:   & 3.5 cm \\
    Right margin:  & 2.0 cm \\
    Top margin:    & 2.5 cm \\
    Bottom margin: & 2.5 cm
\end{tabular}
\end{center}

These settings are predefined in the official template and shall not be
modified by students.

\vspace{0.3cm}

% ================================================================
% IV. FONT, FONT SIZE, AND LINE SPACING
% ================================================================

\noindent
\textbf{IV.\quad FONT, FONT SIZE, AND LINE SPACING}

\vspace{0.15cm}

The main text shall use a 13 pt font size with 1.3-line spacing.
A Times-family typeface shall be used in accordance with the official
template.

The text shall be fully justified.

Students shall not use repeated spaces, repeated manual line breaks, or
multiple blank lines to position text. Paragraph and heading spacing shall
be controlled by the template.

\vspace{0.3cm}

% ================================================================
% V. OUTER COVER
% ================================================================

\noindent
\textbf{V.\quad OUTER COVER}

\vspace{0.15cm}

The outer cover shall follow the official format provided in the template
and shall contain the following information:

\begin{itemize}[leftmargin=0.8cm, itemsep=2pt]
    \item Thai Nguyen University of Technology;
    \item Faculty of International Training;
    \item Engineering and Technology Department;
    \item University logo and Faculty logo;
    \item the title \textbf{INTERNSHIP REPORT};
    \item Internship Organization;
    \item Author;
    \item Class;
    \item Supervisor;
    \item place and year of completion.
\end{itemize}

Students shall only replace the required information fields and shall not
modify the cover-page layout unless instructed otherwise.

\vspace{0.3cm}

% ================================================================
% VI. COMMENTS OF THE INTERNSHIP HOST ORGANIZATION
% ================================================================

\noindent
\textbf{VI.\quad COMMENTS OF THE INTERNSHIP HOST ORGANIZATION}

\vspace{0.15cm}

The page entitled \textbf{COMMENTS OF THE INTERNSHIP HOST ORGANIZATION}
is reserved for the internship host organization.

It includes:

\begin{itemize}[leftmargin=0.8cm, itemsep=2pt]
    \item the student's full name;
    \item class;
    \item comments from the host organization;
    \item score;
    \item date;
    \item signature and official seal.
\end{itemize}

Students shall not alter or complete official evaluation fields unless
specifically authorized.

\vspace{0.3cm}

% ================================================================
% VII. SUPERVISOR'S EVALUATION AND DEFENSE RESULTS
% ================================================================

\noindent
\textbf{VII.\quad SUPERVISOR'S EVALUATION AND DEFENSE RESULTS}

\vspace{0.15cm}

The section entitled
\textbf{SUPERVISOR'S EVALUATION AND GRADUATION INTERNSHIP DEFENSE RESULTS}
is reserved for the supervisor and defense committee.

The evaluation may include:

\begin{itemize}[leftmargin=0.8cm, itemsep=2pt]
    \item Internship Progress and Student's Attitude;
    \item Report Content;
    \item Report Presentation;
    \item Other Comments;
    \item Supervisor's Opinion;
    \item Defense Result;
    \item Score.
\end{itemize}

Students shall not enter evaluation scores or defense results before they
are officially determined.

\vspace{0.3cm}

% ================================================================
% VIII. CHAPTERS, SECTIONS, AND SUBSECTIONS
% ================================================================

\noindent
\textbf{VIII.\quad CHAPTERS, SECTIONS, AND SUBSECTIONS}

\vspace{0.15cm}

The report structure shall be generated using the standard \LaTeX{}
commands:

\begin{verbatim}
\chapter{...}
\section{...}
\subsection{...}
\end{verbatim}

Chapter, section, and subsection numbers shall be generated automatically
by \LaTeX{}.

Students shall not manually type chapter or section numbers and shall not
reproduce heading styles using commands such as \verb|\textbf|,
\verb|\fontsize|, or \verb|\centering| instead of the appropriate
structural command.

Excessive subdivision should be avoided. In most cases, the hierarchy
Chapter--Section--Subsection is sufficient.

\vspace{0.3cm}

% ================================================================
% IX. ORGANIZATION OF THE MAIN CONTENT
% ================================================================

\noindent
\textbf{IX.\quad ORGANIZATION OF THE MAIN CONTENT}

\vspace{0.15cm}

The specific content of the report shall depend on the internship
organization and the assigned tasks. A typical structure may include:

\begin{itemize}[leftmargin=0.8cm, itemsep=2pt]
    \item \textbf{Chapter 1. Overview of the Internship Organization};
    \item \textbf{Chapter 2. Internship Activities and Assigned Tasks};
    \item \textbf{Chapter 3. Results, Experience, and Evaluation}.
\end{itemize}

These chapter titles are recommendations only and may be adapted to the
actual internship content.

The report shall reflect work actually performed, observed, analyzed, or
learned during the internship. Unrelated theoretical material should not
be included merely to increase report length.

\vspace{0.3cm}

% ================================================================
% X. TABLE OF CONTENTS
% ================================================================

\noindent
\textbf{X.\quad TABLE OF CONTENTS}

\vspace{0.15cm}

The Table of Contents shall be generated automatically using:

\begin{verbatim}
\tableofcontents
\end{verbatim}

Students shall not create or manually edit page numbers in the Table of
Contents.

The document shall be compiled sufficiently to ensure that all page
numbers and entries are correctly updated.

\vspace{0.3cm}

% ================================================================
% XI. LIST OF SYMBOLS AND ABBREVIATIONS
% ================================================================

\noindent
\textbf{XI.\quad LIST OF SYMBOLS AND ABBREVIATIONS}

\vspace{0.15cm}

Symbols and abbreviations that are used repeatedly in the report shall be
listed in the \textbf{List of Symbols and Abbreviations}.

The list may contain the following columns:

\begin{center}
\begin{tabular}{llll}
    No. & Symbol / Abbreviation & Unit & Definition
\end{tabular}
\end{center}

Only symbols and abbreviations actually used in the final report shall be
retained. Example entries contained in the template shall be removed if
they are not relevant to the student's report.

Each abbreviation should normally be defined when it first appears in
the main text.

\vspace{0.3cm}

% ================================================================
% XII. FIGURES, DIAGRAMS, AND PHOTOGRAPHS
% ================================================================

\noindent
\textbf{XII.\quad FIGURES, DIAGRAMS, AND PHOTOGRAPHS}

\vspace{0.15cm}

Figures shall be clear, legible, and directly relevant to the report
content.

Figures shall be inserted using the standard \verb|figure| environment.

Example:

\begin{verbatim}
\begin{figure}[H]
    \centering
    \includegraphics[width=0.8\textwidth]
                    {figures/company_structure.pdf}
    \caption{Organizational structure of the company}
    \label{fig:company-structure}
\end{figure}
\end{verbatim}

Figure captions shall be placed below the figure.

Each referenced figure shall have a unique \verb|\label| and shall be
referenced in the text using, for example:

\begin{verbatim}
Figure~\ref{fig:company-structure}
\end{verbatim}

Students shall not manually type figure numbers.

Figures reproduced or adapted from external sources shall include an
appropriate citation.

Photographs taken during the internship shall be relevant and sufficiently
clear. Students shall not include confidential, restricted, or proprietary
information without permission from the internship organization.

Vector graphics or high-resolution images are recommended whenever
possible.

\vspace{0.3cm}

% ================================================================
% XIII. TABLES
% ================================================================

\noindent
\textbf{XIII.\quad TABLES}

\vspace{0.15cm}

Tables shall be created using suitable \LaTeX{} environments such as
\verb|table|, \verb|tabular|, \verb|tabularx|, or \verb|longtable|.

Table captions shall be placed above the table.

Example:

\begin{verbatim}
\begin{table}[H]
    \centering
    \caption{Internship schedule}
    \label{tab:internship-schedule}

    \begin{tabular}{lll}
        \toprule
        Week & Task & Result \\
        \midrule
        1 & Orientation & Completed \\
        2 & System study & Completed \\
        \bottomrule
    \end{tabular}
\end{table}
\end{verbatim}

Tables shall be referenced using:

\begin{verbatim}
Table~\ref{tab:internship-schedule}
\end{verbatim}

Students shall not manually type table numbers.

Tables should be created directly in \LaTeX{} rather than inserted as
screenshots whenever practical.

\vspace{0.3cm}

% ================================================================
% XIV. EQUATIONS AND TECHNICAL QUANTITIES
% ================================================================

\noindent
\textbf{XIV.\quad EQUATIONS AND TECHNICAL QUANTITIES}

\vspace{0.15cm}

Mathematical expressions shall be typeset directly using \LaTeX{}.
Images or screenshots of equations shall not be used.

Short expressions may be written in inline mathematical mode, for example:

\begin{verbatim}
$R = 0.7\,\Omega$
\end{verbatim}

Important equations shall use a numbered equation environment.

Example:

\begin{verbatim}
\begin{equation}
    P = UI
    \label{eq:power}
\end{equation}
\end{verbatim}

Equations shall be referenced using:

\begin{verbatim}
Equation~\eqref{eq:power}
\end{verbatim}

Equation numbers shall not be manually entered.

SI units shall be used wherever applicable. The \verb|siunitx| package
should be used where appropriate.

Examples:

\begin{verbatim}
\SI{220}{\volt}
\SI{50}{\hertz}
\SI{2.5}{\ampere}
\end{verbatim}

Notation and units shall be consistent throughout the report.

\vspace{0.3cm}

% ================================================================
% XV. SOURCE CODE
% ================================================================

\noindent
\textbf{XV.\quad SOURCE CODE}

\vspace{0.15cm}

If the internship involves programming, PLCs, microcontrollers,
simulation, or software development, relevant source code may be included
in the report.

Source code shall be presented using the \verb|listings| environment or
another environment defined by the official template.

Example:

\begin{verbatim}
\begin{lstlisting}[language=st]
IF Start = TRUE THEN
    Motor := TRUE;
ELSE
    Motor := FALSE;
END_IF;
\end{lstlisting}
\end{verbatim}

Only code necessary for explaining the work should appear in the main
report. Long source-code listings should be moved to an appendix or
submitted as supplementary files if required.

Source code should not normally be inserted as screenshots.

\vspace{0.3cm}

% ================================================================
% XVI. ABSTRACT
% ================================================================

\noindent
\textbf{XVI.\quad ABSTRACT}

\vspace{0.15cm}

The Abstract shall briefly summarize:

\begin{itemize}[leftmargin=0.8cm, itemsep=2pt]
    \item the internship organization;
    \item the objectives of the internship;
    \item the main tasks performed;
    \item the main results, knowledge, and skills obtained.
\end{itemize}

The Abstract should be concise and should not reproduce the full content
of the report chapter by chapter.

\vspace{0.3cm}

% ================================================================
% XVII. REFERENCES AND CITATIONS
% ================================================================

\noindent
\textbf{XVII.\quad REFERENCES AND CITATIONS}

\vspace{0.15cm}

All books, journal papers, conference papers, standards, websites,
datasheets, manuals, technical documents, figures, tables, or other
external sources used in preparing the report shall be properly cited.

References shall be managed through the bibliography system provided with
the official template.

Bibliographic entries shall normally be stored in:

\begin{verbatim}
References.bib
\end{verbatim}

References shall be cited in the report using:

\begin{verbatim}
\cite{reference-key}
\end{verbatim}

Students shall not manually type reference numbers such as ``[1]'',
``[2]'', etc.

Titles of foreign-language publications shall remain in their original
language in the bibliography.

All borrowed ideas, data, figures, tables, equations, methods, or results
shall be properly acknowledged.

\vspace{0.3cm}

% ================================================================
% XVIII. CROSS-REFERENCING
% ================================================================

\noindent
\textbf{XVIII.\quad CROSS-REFERENCING}

\vspace{0.15cm}

Numbered elements shall be cross-referenced using the standard \LaTeX{}
mechanism.

Use the following commands where appropriate:

\begin{verbatim}
\label{...}
\ref{...}
\eqref{...}
\end{verbatim}

For example:

\begin{verbatim}
As shown in Figure~\ref{fig:system}, ...
\end{verbatim}

Students shall not manually type numbers such as ``Figure 2.3'' when
those numbers can be generated automatically.

\vspace{0.3cm}

% ================================================================
% XIX. PAGE NUMBERING
% ================================================================

\noindent
\textbf{XIX.\quad PAGE NUMBERING}

\vspace{0.15cm}

Cover pages and official evaluation forms at the beginning of the report
shall not display page numbers.

The main body of the report shall use Arabic page numbers beginning from
Chapter 1.

Page numbering shall be generated automatically by the official template.
Students shall not manually enter page numbers.

\vspace{0.3cm}

% ================================================================
% XX. PROJECT FILE ORGANIZATION
% ================================================================

\noindent
\textbf{XX.\quad PROJECT FILE ORGANIZATION}

\vspace{0.15cm}

Students should maintain a clear and organized \LaTeX{} project
structure. A recommended structure is:

\begin{verbatim}
InternshipReport/
|
|-- main.tex
|-- References.bib
|
|-- figures/
|   |-- company.pdf
|   |-- system.png
|   `-- ...
|
`-- chapters/
    |-- chapter1.tex
    |-- chapter2.tex
    `-- chapter3.tex
\end{verbatim}

Figures should normally be stored in the \verb|figures| directory
specified by the official template.

File and directory names should preferably contain Latin letters,
numbers, hyphens, or underscores.

Recommended examples include:

\begin{verbatim}
company_structure.pdf
motor_control.png
chapter_01.tex
\end{verbatim}

Spaces and unnecessary special characters in file names should be
avoided where possible.

\vspace{0.3cm}

% ================================================================
% XXI. DOCUMENT COMPILATION AND FINAL CHECK
% ================================================================

\noindent
\textbf{XXI.\quad DOCUMENT COMPILATION AND FINAL CHECK}

\vspace{0.15cm}

Before submission, the complete \LaTeX{} project shall compile
successfully and produce the final PDF report.

The final document shall not contain unresolved references such as:

\begin{verbatim}
??
[?]
\end{verbatim}

Before generating the final PDF, students shall check:

\begin{itemize}[leftmargin=0.8cm, itemsep=2pt]
    \item the Table of Contents;
    \item page numbering;
    \item chapter, section, and subsection numbering;
    \item figure numbering;
    \item table numbering;
    \item equation numbering;
    \item cross-references;
    \item citations;
    \item bibliography;
    \item font display;
    \item mathematical symbols;
    \item image quality;
    \item completeness of all required forms.
\end{itemize}

The document may be compiled using TeX Live, MiKTeX, or another
compatible \LaTeX{} environment. Regardless of the compilation
environment, the resulting PDF shall conform to the official Faculty
template.

\vspace{0.3cm}

% ================================================================
% XXII. CONTENT THAT STUDENTS MAY EDIT
% ================================================================

\noindent
\textbf{XXII.\quad CONTENT THAT STUDENTS MAY EDIT}

\vspace{0.15cm}

Students may edit or replace:

\begin{itemize}[leftmargin=0.8cm, itemsep=2pt]
    \item student name;
    \item class;
    \item internship organization;
    \item supervisor name;
    \item Abstract;
    \item List of Symbols and Abbreviations;
    \item chapter titles and chapter content;
    \item figures and photographs;
    \item tables;
    \item equations;
    \item citations and references.
\end{itemize}

Students shall not modify the main document formatting settings in the
preamble, including the document class, page geometry, line spacing,
heading formatting, table-of-contents formatting, font configuration,
page layout, or cover-page macros unless specifically instructed.

\vspace{0.3cm}

% ================================================================
% XXIII. FINAL SUBMISSION REQUIREMENTS
% ================================================================

\noindent
\textbf{XXIII.\quad FINAL SUBMISSION REQUIREMENTS}

\vspace{0.15cm}

Before submission, the final report shall:

\begin{itemize}[leftmargin=0.8cm, itemsep=2pt]
    \item compile without errors that affect the report content or presentation;
    \item contain all required forms and report sections;
    \item contain correct student, class, internship organization, and supervisor information;
    \item contain an updated Table of Contents, List of Figures, and List of Tables;
    \item contain no unused example content from the template;
    \item contain no placeholder text such as \textit{Title of Chapter 1},
    \textit{Section 1.1}, or other sample entries;
    \item contain all required citations and references;
    \item contain no unresolved cross-references.
\end{itemize}

If the Faculty requires submission of the \LaTeX{} source, the complete
project shall include all \verb|.tex| files, the bibliography database,
figures, and any other files required to reproduce the submitted PDF.

\vspace{0.4cm}

\noindent
\textbf{Important:}
Before submitting the final internship report, students shall remove or
comment out the following command in the main document:

\begin{verbatim}
\input{Regulations.tex}
\end{verbatim}

For example:

\begin{verbatim}
% \input{Regulations.tex}
\end{verbatim}

This removes the guidance pages from the final submitted report while
preserving the regulations file in the \LaTeX{} project for reference.

} % End 13pt content

\clearpage

\end{verbatim}

For example:

\begin{verbatim}
% 
% =================================================================
% REGULATIONS FOR PREPARING THE GRADUATION INTERNSHIP REPORT
% USING LATEX
% =================================================================
%
% This file is intended to be included from the main document with:
%
%     \input{Regulations.tex}
%
% IMPORTANT:
% This section is provided for guidance only. Students shall remove
% or comment out the \input{Regulations.tex} command before submitting
% the final internship report.
%
% =================================================================

\clearpage
\thispagestyle{empty}

\begin{center}
    {\fontsize{16pt}{20pt}\selectfont\bfseries
    SPECIFIC REGULATIONS FOR PREPARING\\
    GRADUATION INTERNSHIP REPORTS USING \LaTeX}
\end{center}

\vspace{0.5cm}

{\fontsize{13pt}{17pt}\selectfont
\setstretch{1.3}
\justifying

\noindent
\textbf{Note:}
This section is provided for guidance only and is included in the main
document using the command \verb|\input{Regulations.tex}|.
Students shall remove or comment out this command before submitting the
final internship report.

\vspace{0.4cm}

% ================================================================
% I. GENERAL REQUIREMENTS
% ================================================================

\noindent
\textbf{I.\quad GENERAL REQUIREMENTS}

\vspace{0.15cm}

The Graduation Internship Report shall be prepared using the official
\LaTeX{} template provided by the Faculty. Students shall preserve the
predefined page layout, typography, line spacing, heading styles,
numbering system, figure and table formatting, and bibliography settings
unless otherwise instructed by the Faculty or the supervisor.

The report shall be written in English.

Students are responsible for the academic and technical content of the
report. The official \LaTeX{} template is responsible for standardized
document formatting.

\vspace{0.3cm}

% ================================================================
% II. REPORT STRUCTURE
% ================================================================

\noindent
\textbf{II.\quad REPORT STRUCTURE}

\vspace{0.15cm}

The report shall normally be arranged in the following order:

\begin{enumerate}[label=\arabic*., leftmargin=1.2cm, itemsep=2pt]
    \item Outer Cover;
    \item Comments of the Internship Host Organization;
    \item Supervisor's Evaluation and Graduation Internship Defense Results;
    \item Table of Contents;
    \item List of Symbols and Abbreviations, if applicable;
    \item List of Figures;
    \item List of Tables;
    \item Abstract;
    \item Main Chapters of the Internship Report;
    \item References;
    \item Appendices, if applicable.
\end{enumerate}

The preliminary pages shall not be numbered as chapters. The main report
content shall be organized into consecutively numbered chapters beginning
with Chapter 1.

\vspace{0.3cm}

% ================================================================
% III. PAGE SIZE AND MARGINS
% ================================================================

\noindent
\textbf{III.\quad PAGE SIZE AND MARGINS}

\vspace{0.15cm}

The report shall be formatted on A4 paper with the following margins:

\begin{center}
\begin{tabular}{ll}
    Left margin:   & 3.5 cm \\
    Right margin:  & 2.0 cm \\
    Top margin:    & 2.5 cm \\
    Bottom margin: & 2.5 cm
\end{tabular}
\end{center}

These settings are predefined in the official template and shall not be
modified by students.

\vspace{0.3cm}

% ================================================================
% IV. FONT, FONT SIZE, AND LINE SPACING
% ================================================================

\noindent
\textbf{IV.\quad FONT, FONT SIZE, AND LINE SPACING}

\vspace{0.15cm}

The main text shall use a 13 pt font size with 1.3-line spacing.
A Times-family typeface shall be used in accordance with the official
template.

The text shall be fully justified.

Students shall not use repeated spaces, repeated manual line breaks, or
multiple blank lines to position text. Paragraph and heading spacing shall
be controlled by the template.

\vspace{0.3cm}

% ================================================================
% V. OUTER COVER
% ================================================================

\noindent
\textbf{V.\quad OUTER COVER}

\vspace{0.15cm}

The outer cover shall follow the official format provided in the template
and shall contain the following information:

\begin{itemize}[leftmargin=0.8cm, itemsep=2pt]
    \item Thai Nguyen University of Technology;
    \item Faculty of International Training;
    \item Engineering and Technology Department;
    \item University logo and Faculty logo;
    \item the title \textbf{INTERNSHIP REPORT};
    \item Internship Organization;
    \item Author;
    \item Class;
    \item Supervisor;
    \item place and year of completion.
\end{itemize}

Students shall only replace the required information fields and shall not
modify the cover-page layout unless instructed otherwise.

\vspace{0.3cm}

% ================================================================
% VI. COMMENTS OF THE INTERNSHIP HOST ORGANIZATION
% ================================================================

\noindent
\textbf{VI.\quad COMMENTS OF THE INTERNSHIP HOST ORGANIZATION}

\vspace{0.15cm}

The page entitled \textbf{COMMENTS OF THE INTERNSHIP HOST ORGANIZATION}
is reserved for the internship host organization.

It includes:

\begin{itemize}[leftmargin=0.8cm, itemsep=2pt]
    \item the student's full name;
    \item class;
    \item comments from the host organization;
    \item score;
    \item date;
    \item signature and official seal.
\end{itemize}

Students shall not alter or complete official evaluation fields unless
specifically authorized.

\vspace{0.3cm}

% ================================================================
% VII. SUPERVISOR'S EVALUATION AND DEFENSE RESULTS
% ================================================================

\noindent
\textbf{VII.\quad SUPERVISOR'S EVALUATION AND DEFENSE RESULTS}

\vspace{0.15cm}

The section entitled
\textbf{SUPERVISOR'S EVALUATION AND GRADUATION INTERNSHIP DEFENSE RESULTS}
is reserved for the supervisor and defense committee.

The evaluation may include:

\begin{itemize}[leftmargin=0.8cm, itemsep=2pt]
    \item Internship Progress and Student's Attitude;
    \item Report Content;
    \item Report Presentation;
    \item Other Comments;
    \item Supervisor's Opinion;
    \item Defense Result;
    \item Score.
\end{itemize}

Students shall not enter evaluation scores or defense results before they
are officially determined.

\vspace{0.3cm}

% ================================================================
% VIII. CHAPTERS, SECTIONS, AND SUBSECTIONS
% ================================================================

\noindent
\textbf{VIII.\quad CHAPTERS, SECTIONS, AND SUBSECTIONS}

\vspace{0.15cm}

The report structure shall be generated using the standard \LaTeX{}
commands:

\begin{verbatim}
\chapter{...}
\section{...}
\subsection{...}
\end{verbatim}

Chapter, section, and subsection numbers shall be generated automatically
by \LaTeX{}.

Students shall not manually type chapter or section numbers and shall not
reproduce heading styles using commands such as \verb|\textbf|,
\verb|\fontsize|, or \verb|\centering| instead of the appropriate
structural command.

Excessive subdivision should be avoided. In most cases, the hierarchy
Chapter--Section--Subsection is sufficient.

\vspace{0.3cm}

% ================================================================
% IX. ORGANIZATION OF THE MAIN CONTENT
% ================================================================

\noindent
\textbf{IX.\quad ORGANIZATION OF THE MAIN CONTENT}

\vspace{0.15cm}

The specific content of the report shall depend on the internship
organization and the assigned tasks. A typical structure may include:

\begin{itemize}[leftmargin=0.8cm, itemsep=2pt]
    \item \textbf{Chapter 1. Overview of the Internship Organization};
    \item \textbf{Chapter 2. Internship Activities and Assigned Tasks};
    \item \textbf{Chapter 3. Results, Experience, and Evaluation}.
\end{itemize}

These chapter titles are recommendations only and may be adapted to the
actual internship content.

The report shall reflect work actually performed, observed, analyzed, or
learned during the internship. Unrelated theoretical material should not
be included merely to increase report length.

\vspace{0.3cm}

% ================================================================
% X. TABLE OF CONTENTS
% ================================================================

\noindent
\textbf{X.\quad TABLE OF CONTENTS}

\vspace{0.15cm}

The Table of Contents shall be generated automatically using:

\begin{verbatim}
\tableofcontents
\end{verbatim}

Students shall not create or manually edit page numbers in the Table of
Contents.

The document shall be compiled sufficiently to ensure that all page
numbers and entries are correctly updated.

\vspace{0.3cm}

% ================================================================
% XI. LIST OF SYMBOLS AND ABBREVIATIONS
% ================================================================

\noindent
\textbf{XI.\quad LIST OF SYMBOLS AND ABBREVIATIONS}

\vspace{0.15cm}

Symbols and abbreviations that are used repeatedly in the report shall be
listed in the \textbf{List of Symbols and Abbreviations}.

The list may contain the following columns:

\begin{center}
\begin{tabular}{llll}
    No. & Symbol / Abbreviation & Unit & Definition
\end{tabular}
\end{center}

Only symbols and abbreviations actually used in the final report shall be
retained. Example entries contained in the template shall be removed if
they are not relevant to the student's report.

Each abbreviation should normally be defined when it first appears in
the main text.

\vspace{0.3cm}

% ================================================================
% XII. FIGURES, DIAGRAMS, AND PHOTOGRAPHS
% ================================================================

\noindent
\textbf{XII.\quad FIGURES, DIAGRAMS, AND PHOTOGRAPHS}

\vspace{0.15cm}

Figures shall be clear, legible, and directly relevant to the report
content.

Figures shall be inserted using the standard \verb|figure| environment.

Example:

\begin{verbatim}
\begin{figure}[H]
    \centering
    \includegraphics[width=0.8\textwidth]
                    {figures/company_structure.pdf}
    \caption{Organizational structure of the company}
    \label{fig:company-structure}
\end{figure}
\end{verbatim}

Figure captions shall be placed below the figure.

Each referenced figure shall have a unique \verb|\label| and shall be
referenced in the text using, for example:

\begin{verbatim}
Figure~\ref{fig:company-structure}
\end{verbatim}

Students shall not manually type figure numbers.

Figures reproduced or adapted from external sources shall include an
appropriate citation.

Photographs taken during the internship shall be relevant and sufficiently
clear. Students shall not include confidential, restricted, or proprietary
information without permission from the internship organization.

Vector graphics or high-resolution images are recommended whenever
possible.

\vspace{0.3cm}

% ================================================================
% XIII. TABLES
% ================================================================

\noindent
\textbf{XIII.\quad TABLES}

\vspace{0.15cm}

Tables shall be created using suitable \LaTeX{} environments such as
\verb|table|, \verb|tabular|, \verb|tabularx|, or \verb|longtable|.

Table captions shall be placed above the table.

Example:

\begin{verbatim}
\begin{table}[H]
    \centering
    \caption{Internship schedule}
    \label{tab:internship-schedule}

    \begin{tabular}{lll}
        \toprule
        Week & Task & Result \\
        \midrule
        1 & Orientation & Completed \\
        2 & System study & Completed \\
        \bottomrule
    \end{tabular}
\end{table}
\end{verbatim}

Tables shall be referenced using:

\begin{verbatim}
Table~\ref{tab:internship-schedule}
\end{verbatim}

Students shall not manually type table numbers.

Tables should be created directly in \LaTeX{} rather than inserted as
screenshots whenever practical.

\vspace{0.3cm}

% ================================================================
% XIV. EQUATIONS AND TECHNICAL QUANTITIES
% ================================================================

\noindent
\textbf{XIV.\quad EQUATIONS AND TECHNICAL QUANTITIES}

\vspace{0.15cm}

Mathematical expressions shall be typeset directly using \LaTeX{}.
Images or screenshots of equations shall not be used.

Short expressions may be written in inline mathematical mode, for example:

\begin{verbatim}
$R = 0.7\,\Omega$
\end{verbatim}

Important equations shall use a numbered equation environment.

Example:

\begin{verbatim}
\begin{equation}
    P = UI
    \label{eq:power}
\end{equation}
\end{verbatim}

Equations shall be referenced using:

\begin{verbatim}
Equation~\eqref{eq:power}
\end{verbatim}

Equation numbers shall not be manually entered.

SI units shall be used wherever applicable. The \verb|siunitx| package
should be used where appropriate.

Examples:

\begin{verbatim}
\SI{220}{\volt}
\SI{50}{\hertz}
\SI{2.5}{\ampere}
\end{verbatim}

Notation and units shall be consistent throughout the report.

\vspace{0.3cm}

% ================================================================
% XV. SOURCE CODE
% ================================================================

\noindent
\textbf{XV.\quad SOURCE CODE}

\vspace{0.15cm}

If the internship involves programming, PLCs, microcontrollers,
simulation, or software development, relevant source code may be included
in the report.

Source code shall be presented using the \verb|listings| environment or
another environment defined by the official template.

Example:

\begin{verbatim}
\begin{lstlisting}[language=st]
IF Start = TRUE THEN
    Motor := TRUE;
ELSE
    Motor := FALSE;
END_IF;
\end{lstlisting}
\end{verbatim}

Only code necessary for explaining the work should appear in the main
report. Long source-code listings should be moved to an appendix or
submitted as supplementary files if required.

Source code should not normally be inserted as screenshots.

\vspace{0.3cm}

% ================================================================
% XVI. ABSTRACT
% ================================================================

\noindent
\textbf{XVI.\quad ABSTRACT}

\vspace{0.15cm}

The Abstract shall briefly summarize:

\begin{itemize}[leftmargin=0.8cm, itemsep=2pt]
    \item the internship organization;
    \item the objectives of the internship;
    \item the main tasks performed;
    \item the main results, knowledge, and skills obtained.
\end{itemize}

The Abstract should be concise and should not reproduce the full content
of the report chapter by chapter.

\vspace{0.3cm}

% ================================================================
% XVII. REFERENCES AND CITATIONS
% ================================================================

\noindent
\textbf{XVII.\quad REFERENCES AND CITATIONS}

\vspace{0.15cm}

All books, journal papers, conference papers, standards, websites,
datasheets, manuals, technical documents, figures, tables, or other
external sources used in preparing the report shall be properly cited.

References shall be managed through the bibliography system provided with
the official template.

Bibliographic entries shall normally be stored in:

\begin{verbatim}
References.bib
\end{verbatim}

References shall be cited in the report using:

\begin{verbatim}
\cite{reference-key}
\end{verbatim}

Students shall not manually type reference numbers such as ``[1]'',
``[2]'', etc.

Titles of foreign-language publications shall remain in their original
language in the bibliography.

All borrowed ideas, data, figures, tables, equations, methods, or results
shall be properly acknowledged.

\vspace{0.3cm}

% ================================================================
% XVIII. CROSS-REFERENCING
% ================================================================

\noindent
\textbf{XVIII.\quad CROSS-REFERENCING}

\vspace{0.15cm}

Numbered elements shall be cross-referenced using the standard \LaTeX{}
mechanism.

Use the following commands where appropriate:

\begin{verbatim}
\label{...}
\ref{...}
\eqref{...}
\end{verbatim}

For example:

\begin{verbatim}
As shown in Figure~\ref{fig:system}, ...
\end{verbatim}

Students shall not manually type numbers such as ``Figure 2.3'' when
those numbers can be generated automatically.

\vspace{0.3cm}

% ================================================================
% XIX. PAGE NUMBERING
% ================================================================

\noindent
\textbf{XIX.\quad PAGE NUMBERING}

\vspace{0.15cm}

Cover pages and official evaluation forms at the beginning of the report
shall not display page numbers.

The main body of the report shall use Arabic page numbers beginning from
Chapter 1.

Page numbering shall be generated automatically by the official template.
Students shall not manually enter page numbers.

\vspace{0.3cm}

% ================================================================
% XX. PROJECT FILE ORGANIZATION
% ================================================================

\noindent
\textbf{XX.\quad PROJECT FILE ORGANIZATION}

\vspace{0.15cm}

Students should maintain a clear and organized \LaTeX{} project
structure. A recommended structure is:

\begin{verbatim}
InternshipReport/
|
|-- main.tex
|-- References.bib
|
|-- figures/
|   |-- company.pdf
|   |-- system.png
|   `-- ...
|
`-- chapters/
    |-- chapter1.tex
    |-- chapter2.tex
    `-- chapter3.tex
\end{verbatim}

Figures should normally be stored in the \verb|figures| directory
specified by the official template.

File and directory names should preferably contain Latin letters,
numbers, hyphens, or underscores.

Recommended examples include:

\begin{verbatim}
company_structure.pdf
motor_control.png
chapter_01.tex
\end{verbatim}

Spaces and unnecessary special characters in file names should be
avoided where possible.

\vspace{0.3cm}

% ================================================================
% XXI. DOCUMENT COMPILATION AND FINAL CHECK
% ================================================================

\noindent
\textbf{XXI.\quad DOCUMENT COMPILATION AND FINAL CHECK}

\vspace{0.15cm}

Before submission, the complete \LaTeX{} project shall compile
successfully and produce the final PDF report.

The final document shall not contain unresolved references such as:

\begin{verbatim}
??
[?]
\end{verbatim}

Before generating the final PDF, students shall check:

\begin{itemize}[leftmargin=0.8cm, itemsep=2pt]
    \item the Table of Contents;
    \item page numbering;
    \item chapter, section, and subsection numbering;
    \item figure numbering;
    \item table numbering;
    \item equation numbering;
    \item cross-references;
    \item citations;
    \item bibliography;
    \item font display;
    \item mathematical symbols;
    \item image quality;
    \item completeness of all required forms.
\end{itemize}

The document may be compiled using TeX Live, MiKTeX, or another
compatible \LaTeX{} environment. Regardless of the compilation
environment, the resulting PDF shall conform to the official Faculty
template.

\vspace{0.3cm}

% ================================================================
% XXII. CONTENT THAT STUDENTS MAY EDIT
% ================================================================

\noindent
\textbf{XXII.\quad CONTENT THAT STUDENTS MAY EDIT}

\vspace{0.15cm}

Students may edit or replace:

\begin{itemize}[leftmargin=0.8cm, itemsep=2pt]
    \item student name;
    \item class;
    \item internship organization;
    \item supervisor name;
    \item Abstract;
    \item List of Symbols and Abbreviations;
    \item chapter titles and chapter content;
    \item figures and photographs;
    \item tables;
    \item equations;
    \item citations and references.
\end{itemize}

Students shall not modify the main document formatting settings in the
preamble, including the document class, page geometry, line spacing,
heading formatting, table-of-contents formatting, font configuration,
page layout, or cover-page macros unless specifically instructed.

\vspace{0.3cm}

% ================================================================
% XXIII. FINAL SUBMISSION REQUIREMENTS
% ================================================================

\noindent
\textbf{XXIII.\quad FINAL SUBMISSION REQUIREMENTS}

\vspace{0.15cm}

Before submission, the final report shall:

\begin{itemize}[leftmargin=0.8cm, itemsep=2pt]
    \item compile without errors that affect the report content or presentation;
    \item contain all required forms and report sections;
    \item contain correct student, class, internship organization, and supervisor information;
    \item contain an updated Table of Contents, List of Figures, and List of Tables;
    \item contain no unused example content from the template;
    \item contain no placeholder text such as \textit{Title of Chapter 1},
    \textit{Section 1.1}, or other sample entries;
    \item contain all required citations and references;
    \item contain no unresolved cross-references.
\end{itemize}

If the Faculty requires submission of the \LaTeX{} source, the complete
project shall include all \verb|.tex| files, the bibliography database,
figures, and any other files required to reproduce the submitted PDF.

\vspace{0.4cm}

\noindent
\textbf{Important:}
Before submitting the final internship report, students shall remove or
comment out the following command in the main document:

\begin{verbatim}
\input{Regulations.tex}
\end{verbatim}

For example:

\begin{verbatim}
% \input{Regulations.tex}
\end{verbatim}

This removes the guidance pages from the final submitted report while
preserving the regulations file in the \LaTeX{} project for reference.

} % End 13pt content

\clearpage

\end{verbatim}

This removes the guidance pages from the final submitted report while
preserving the regulations file in the \LaTeX{} project for reference.

} % End 13pt content

\clearpage
|.
Students shall remove or comment out this command before submitting the
final internship report.

\vspace{0.4cm}

% ================================================================
% I. GENERAL REQUIREMENTS
% ================================================================

\noindent
\textbf{I.\quad GENERAL REQUIREMENTS}

\vspace{0.15cm}

The Graduation Internship Report shall be prepared using the official
\LaTeX{} template provided by the Faculty. Students shall preserve the
predefined page layout, typography, line spacing, heading styles,
numbering system, figure and table formatting, and bibliography settings
unless otherwise instructed by the Faculty or the supervisor.

The report shall be written in English.

Students are responsible for the academic and technical content of the
report. The official \LaTeX{} template is responsible for standardized
document formatting.

\vspace{0.3cm}

% ================================================================
% II. REPORT STRUCTURE
% ================================================================

\noindent
\textbf{II.\quad REPORT STRUCTURE}

\vspace{0.15cm}

The report shall normally be arranged in the following order:

\begin{enumerate}[label=\arabic*., leftmargin=1.2cm, itemsep=2pt]
    \item Outer Cover;
    \item Comments of the Internship Host Organization;
    \item Supervisor's Evaluation and Graduation Internship Defense Results;
    \item Table of Contents;
    \item List of Symbols and Abbreviations, if applicable;
    \item List of Figures;
    \item List of Tables;
    \item Abstract;
    \item Main Chapters of the Internship Report;
    \item References;
    \item Appendices, if applicable.
\end{enumerate}

The preliminary pages shall not be numbered as chapters. The main report
content shall be organized into consecutively numbered chapters beginning
with Chapter 1.

\vspace{0.3cm}

% ================================================================
% III. PAGE SIZE AND MARGINS
% ================================================================

\noindent
\textbf{III.\quad PAGE SIZE AND MARGINS}

\vspace{0.15cm}

The report shall be formatted on A4 paper with the following margins:

\begin{center}
\begin{tabular}{ll}
    Left margin:   & 3.5 cm \\
    Right margin:  & 2.0 cm \\
    Top margin:    & 2.5 cm \\
    Bottom margin: & 2.5 cm
\end{tabular}
\end{center}

These settings are predefined in the official template and shall not be
modified by students.

\vspace{0.3cm}

% ================================================================
% IV. FONT, FONT SIZE, AND LINE SPACING
% ================================================================

\noindent
\textbf{IV.\quad FONT, FONT SIZE, AND LINE SPACING}

\vspace{0.15cm}

The main text shall use a 13 pt font size with 1.3-line spacing.
A Times-family typeface shall be used in accordance with the official
template.

The text shall be fully justified.

Students shall not use repeated spaces, repeated manual line breaks, or
multiple blank lines to position text. Paragraph and heading spacing shall
be controlled by the template.

\vspace{0.3cm}

% ================================================================
% V. OUTER COVER
% ================================================================

\noindent
\textbf{V.\quad OUTER COVER}

\vspace{0.15cm}

The outer cover shall follow the official format provided in the template
and shall contain the following information:

\begin{itemize}[leftmargin=0.8cm, itemsep=2pt]
    \item Thai Nguyen University of Technology;
    \item Faculty of International Training;
    \item Engineering and Technology Department;
    \item University logo and Faculty logo;
    \item the title \textbf{INTERNSHIP REPORT};
    \item Internship Organization;
    \item Author;
    \item Class;
    \item Supervisor;
    \item place and year of completion.
\end{itemize}

Students shall only replace the required information fields and shall not
modify the cover-page layout unless instructed otherwise.

\vspace{0.3cm}

% ================================================================
% VI. COMMENTS OF THE INTERNSHIP HOST ORGANIZATION
% ================================================================

\noindent
\textbf{VI.\quad COMMENTS OF THE INTERNSHIP HOST ORGANIZATION}

\vspace{0.15cm}

The page entitled \textbf{COMMENTS OF THE INTERNSHIP HOST ORGANIZATION}
is reserved for the internship host organization.

It includes:

\begin{itemize}[leftmargin=0.8cm, itemsep=2pt]
    \item the student's full name;
    \item class;
    \item comments from the host organization;
    \item score;
    \item date;
    \item signature and official seal.
\end{itemize}

Students shall not alter or complete official evaluation fields unless
specifically authorized.

\vspace{0.3cm}

% ================================================================
% VII. SUPERVISOR'S EVALUATION AND DEFENSE RESULTS
% ================================================================

\noindent
\textbf{VII.\quad SUPERVISOR'S EVALUATION AND DEFENSE RESULTS}

\vspace{0.15cm}

The section entitled
\textbf{SUPERVISOR'S EVALUATION AND GRADUATION INTERNSHIP DEFENSE RESULTS}
is reserved for the supervisor and defense committee.

The evaluation may include:

\begin{itemize}[leftmargin=0.8cm, itemsep=2pt]
    \item Internship Progress and Student's Attitude;
    \item Report Content;
    \item Report Presentation;
    \item Other Comments;
    \item Supervisor's Opinion;
    \item Defense Result;
    \item Score.
\end{itemize}

Students shall not enter evaluation scores or defense results before they
are officially determined.

\vspace{0.3cm}

% ================================================================
% VIII. CHAPTERS, SECTIONS, AND SUBSECTIONS
% ================================================================

\noindent
\textbf{VIII.\quad CHAPTERS, SECTIONS, AND SUBSECTIONS}

\vspace{0.15cm}

The report structure shall be generated using the standard \LaTeX{}
commands:

\begin{verbatim}
\chapter{...}
\section{...}
\subsection{...}
\end{verbatim}

Chapter, section, and subsection numbers shall be generated automatically
by \LaTeX{}.

Students shall not manually type chapter or section numbers and shall not
reproduce heading styles using commands such as \verb|\textbf|,
\verb|\fontsize|, or \verb|\centering| instead of the appropriate
structural command.

Excessive subdivision should be avoided. In most cases, the hierarchy
Chapter--Section--Subsection is sufficient.

\vspace{0.3cm}

% ================================================================
% IX. ORGANIZATION OF THE MAIN CONTENT
% ================================================================

\noindent
\textbf{IX.\quad ORGANIZATION OF THE MAIN CONTENT}

\vspace{0.15cm}

The specific content of the report shall depend on the internship
organization and the assigned tasks. A typical structure may include:

\begin{itemize}[leftmargin=0.8cm, itemsep=2pt]
    \item \textbf{Chapter 1. Overview of the Internship Organization};
    \item \textbf{Chapter 2. Internship Activities and Assigned Tasks};
    \item \textbf{Chapter 3. Results, Experience, and Evaluation}.
\end{itemize}

These chapter titles are recommendations only and may be adapted to the
actual internship content.

The report shall reflect work actually performed, observed, analyzed, or
learned during the internship. Unrelated theoretical material should not
be included merely to increase report length.

\vspace{0.3cm}

% ================================================================
% X. TABLE OF CONTENTS
% ================================================================

\noindent
\textbf{X.\quad TABLE OF CONTENTS}

\vspace{0.15cm}

The Table of Contents shall be generated automatically using:

\begin{verbatim}
\tableofcontents
\end{verbatim}

Students shall not create or manually edit page numbers in the Table of
Contents.

The document shall be compiled sufficiently to ensure that all page
numbers and entries are correctly updated.

\vspace{0.3cm}

% ================================================================
% XI. LIST OF SYMBOLS AND ABBREVIATIONS
% ================================================================

\noindent
\textbf{XI.\quad LIST OF SYMBOLS AND ABBREVIATIONS}

\vspace{0.15cm}

Symbols and abbreviations that are used repeatedly in the report shall be
listed in the \textbf{List of Symbols and Abbreviations}.

The list may contain the following columns:

\begin{center}
\begin{tabular}{llll}
    No. & Symbol / Abbreviation & Unit & Definition
\end{tabular}
\end{center}

Only symbols and abbreviations actually used in the final report shall be
retained. Example entries contained in the template shall be removed if
they are not relevant to the student's report.

Each abbreviation should normally be defined when it first appears in
the main text.

\vspace{0.3cm}

% ================================================================
% XII. FIGURES, DIAGRAMS, AND PHOTOGRAPHS
% ================================================================

\noindent
\textbf{XII.\quad FIGURES, DIAGRAMS, AND PHOTOGRAPHS}

\vspace{0.15cm}

Figures shall be clear, legible, and directly relevant to the report
content.

Figures shall be inserted using the standard \verb|figure| environment.

Example:

\begin{verbatim}
\begin{figure}[H]
    \centering
    \includegraphics[width=0.8\textwidth]
                    {figures/company_structure.pdf}
    \caption{Organizational structure of the company}
    \label{fig:company-structure}
\end{figure}
\end{verbatim}

Figure captions shall be placed below the figure.

Each referenced figure shall have a unique \verb|\label| and shall be
referenced in the text using, for example:

\begin{verbatim}
Figure~\ref{fig:company-structure}
\end{verbatim}

Students shall not manually type figure numbers.

Figures reproduced or adapted from external sources shall include an
appropriate citation.

Photographs taken during the internship shall be relevant and sufficiently
clear. Students shall not include confidential, restricted, or proprietary
information without permission from the internship organization.

Vector graphics or high-resolution images are recommended whenever
possible.

\vspace{0.3cm}

% ================================================================
% XIII. TABLES
% ================================================================

\noindent
\textbf{XIII.\quad TABLES}

\vspace{0.15cm}

Tables shall be created using suitable \LaTeX{} environments such as
\verb|table|, \verb|tabular|, \verb|tabularx|, or \verb|longtable|.

Table captions shall be placed above the table.

Example:

\begin{verbatim}
\begin{table}[H]
    \centering
    \caption{Internship schedule}
    \label{tab:internship-schedule}

    \begin{tabular}{lll}
        \toprule
        Week & Task & Result \\
        \midrule
        1 & Orientation & Completed \\
        2 & System study & Completed \\
        \bottomrule
    \end{tabular}
\end{table}
\end{verbatim}

Tables shall be referenced using:

\begin{verbatim}
Table~\ref{tab:internship-schedule}
\end{verbatim}

Students shall not manually type table numbers.

Tables should be created directly in \LaTeX{} rather than inserted as
screenshots whenever practical.

\vspace{0.3cm}

% ================================================================
% XIV. EQUATIONS AND TECHNICAL QUANTITIES
% ================================================================

\noindent
\textbf{XIV.\quad EQUATIONS AND TECHNICAL QUANTITIES}

\vspace{0.15cm}

Mathematical expressions shall be typeset directly using \LaTeX{}.
Images or screenshots of equations shall not be used.

Short expressions may be written in inline mathematical mode, for example:

\begin{verbatim}
$R = 0.7\,\Omega$
\end{verbatim}

Important equations shall use a numbered equation environment.

Example:

\begin{verbatim}
\begin{equation}
    P = UI
    \label{eq:power}
\end{equation}
\end{verbatim}

Equations shall be referenced using:

\begin{verbatim}
Equation~\eqref{eq:power}
\end{verbatim}

Equation numbers shall not be manually entered.

SI units shall be used wherever applicable. The \verb|siunitx| package
should be used where appropriate.

Examples:

\begin{verbatim}
\SI{220}{\volt}
\SI{50}{\hertz}
\SI{2.5}{\ampere}
\end{verbatim}

Notation and units shall be consistent throughout the report.

\vspace{0.3cm}

% ================================================================
% XV. SOURCE CODE
% ================================================================

\noindent
\textbf{XV.\quad SOURCE CODE}

\vspace{0.15cm}

If the internship involves programming, PLCs, microcontrollers,
simulation, or software development, relevant source code may be included
in the report.

Source code shall be presented using the \verb|listings| environment or
another environment defined by the official template.

Example:

\begin{verbatim}
\begin{lstlisting}[language=st]
IF Start = TRUE THEN
    Motor := TRUE;
ELSE
    Motor := FALSE;
END_IF;
\end{lstlisting}
\end{verbatim}

Only code necessary for explaining the work should appear in the main
report. Long source-code listings should be moved to an appendix or
submitted as supplementary files if required.

Source code should not normally be inserted as screenshots.

\vspace{0.3cm}

% ================================================================
% XVI. ABSTRACT
% ================================================================

\noindent
\textbf{XVI.\quad ABSTRACT}

\vspace{0.15cm}

The Abstract shall briefly summarize:

\begin{itemize}[leftmargin=0.8cm, itemsep=2pt]
    \item the internship organization;
    \item the objectives of the internship;
    \item the main tasks performed;
    \item the main results, knowledge, and skills obtained.
\end{itemize}

The Abstract should be concise and should not reproduce the full content
of the report chapter by chapter.

\vspace{0.3cm}

% ================================================================
% XVII. REFERENCES AND CITATIONS
% ================================================================

\noindent
\textbf{XVII.\quad REFERENCES AND CITATIONS}

\vspace{0.15cm}

All books, journal papers, conference papers, standards, websites,
datasheets, manuals, technical documents, figures, tables, or other
external sources used in preparing the report shall be properly cited.

References shall be managed through the bibliography system provided with
the official template.

Bibliographic entries shall normally be stored in:

\begin{verbatim}
References.bib
\end{verbatim}

References shall be cited in the report using:

\begin{verbatim}
\cite{reference-key}
\end{verbatim}

Students shall not manually type reference numbers such as ``[1]'',
``[2]'', etc.

Titles of foreign-language publications shall remain in their original
language in the bibliography.

All borrowed ideas, data, figures, tables, equations, methods, or results
shall be properly acknowledged.

\vspace{0.3cm}

% ================================================================
% XVIII. CROSS-REFERENCING
% ================================================================

\noindent
\textbf{XVIII.\quad CROSS-REFERENCING}

\vspace{0.15cm}

Numbered elements shall be cross-referenced using the standard \LaTeX{}
mechanism.

Use the following commands where appropriate:

\begin{verbatim}
\label{...}
\ref{...}
\eqref{...}
\end{verbatim}

For example:

\begin{verbatim}
As shown in Figure~\ref{fig:system}, ...
\end{verbatim}

Students shall not manually type numbers such as ``Figure 2.3'' when
those numbers can be generated automatically.

\vspace{0.3cm}

% ================================================================
% XIX. PAGE NUMBERING
% ================================================================

\noindent
\textbf{XIX.\quad PAGE NUMBERING}

\vspace{0.15cm}

Cover pages and official evaluation forms at the beginning of the report
shall not display page numbers.

The main body of the report shall use Arabic page numbers beginning from
Chapter 1.

Page numbering shall be generated automatically by the official template.
Students shall not manually enter page numbers.

\vspace{0.3cm}

% ================================================================
% XX. PROJECT FILE ORGANIZATION
% ================================================================

\noindent
\textbf{XX.\quad PROJECT FILE ORGANIZATION}

\vspace{0.15cm}

Students should maintain a clear and organized \LaTeX{} project
structure. A recommended structure is:

\begin{verbatim}
InternshipReport/
|
|-- main.tex
|-- References.bib
|
|-- figures/
|   |-- company.pdf
|   |-- system.png
|   `-- ...
|
`-- chapters/
    |-- chapter1.tex
    |-- chapter2.tex
    `-- chapter3.tex
\end{verbatim}

Figures should normally be stored in the \verb|figures| directory
specified by the official template.

File and directory names should preferably contain Latin letters,
numbers, hyphens, or underscores.

Recommended examples include:

\begin{verbatim}
company_structure.pdf
motor_control.png
chapter_01.tex
\end{verbatim}

Spaces and unnecessary special characters in file names should be
avoided where possible.

\vspace{0.3cm}

% ================================================================
% XXI. DOCUMENT COMPILATION AND FINAL CHECK
% ================================================================

\noindent
\textbf{XXI.\quad DOCUMENT COMPILATION AND FINAL CHECK}

\vspace{0.15cm}

Before submission, the complete \LaTeX{} project shall compile
successfully and produce the final PDF report.

The final document shall not contain unresolved references such as:

\begin{verbatim}
??
[?]
\end{verbatim}

Before generating the final PDF, students shall check:

\begin{itemize}[leftmargin=0.8cm, itemsep=2pt]
    \item the Table of Contents;
    \item page numbering;
    \item chapter, section, and subsection numbering;
    \item figure numbering;
    \item table numbering;
    \item equation numbering;
    \item cross-references;
    \item citations;
    \item bibliography;
    \item font display;
    \item mathematical symbols;
    \item image quality;
    \item completeness of all required forms.
\end{itemize}

The document may be compiled using TeX Live, MiKTeX, or another
compatible \LaTeX{} environment. Regardless of the compilation
environment, the resulting PDF shall conform to the official Faculty
template.

\vspace{0.3cm}

% ================================================================
% XXII. CONTENT THAT STUDENTS MAY EDIT
% ================================================================

\noindent
\textbf{XXII.\quad CONTENT THAT STUDENTS MAY EDIT}

\vspace{0.15cm}

Students may edit or replace:

\begin{itemize}[leftmargin=0.8cm, itemsep=2pt]
    \item student name;
    \item class;
    \item internship organization;
    \item supervisor name;
    \item Abstract;
    \item List of Symbols and Abbreviations;
    \item chapter titles and chapter content;
    \item figures and photographs;
    \item tables;
    \item equations;
    \item citations and references.
\end{itemize}

Students shall not modify the main document formatting settings in the
preamble, including the document class, page geometry, line spacing,
heading formatting, table-of-contents formatting, font configuration,
page layout, or cover-page macros unless specifically instructed.

\vspace{0.3cm}

% ================================================================
% XXIII. FINAL SUBMISSION REQUIREMENTS
% ================================================================

\noindent
\textbf{XXIII.\quad FINAL SUBMISSION REQUIREMENTS}

\vspace{0.15cm}

Before submission, the final report shall:

\begin{itemize}[leftmargin=0.8cm, itemsep=2pt]
    \item compile without errors that affect the report content or presentation;
    \item contain all required forms and report sections;
    \item contain correct student, class, internship organization, and supervisor information;
    \item contain an updated Table of Contents, List of Figures, and List of Tables;
    \item contain no unused example content from the template;
    \item contain no placeholder text such as \textit{Title of Chapter 1},
    \textit{Section 1.1}, or other sample entries;
    \item contain all required citations and references;
    \item contain no unresolved cross-references.
\end{itemize}

If the Faculty requires submission of the \LaTeX{} source, the complete
project shall include all \verb|.tex| files, the bibliography database,
figures, and any other files required to reproduce the submitted PDF.

\vspace{0.4cm}

\noindent
\textbf{Important:}
Before submitting the final internship report, students shall remove or
comment out the following command in the main document:

\begin{verbatim}

% =================================================================
% REGULATIONS FOR PREPARING THE GRADUATION INTERNSHIP REPORT
% USING LATEX
% =================================================================
%
% This file is intended to be included from the main document with:
%
%     
% =================================================================
% REGULATIONS FOR PREPARING THE GRADUATION INTERNSHIP REPORT
% USING LATEX
% =================================================================
%
% This file is intended to be included from the main document with:
%
%     \input{Regulations.tex}
%
% IMPORTANT:
% This section is provided for guidance only. Students shall remove
% or comment out the \input{Regulations.tex} command before submitting
% the final internship report.
%
% =================================================================

\clearpage
\thispagestyle{empty}

\begin{center}
    {\fontsize{16pt}{20pt}\selectfont\bfseries
    SPECIFIC REGULATIONS FOR PREPARING\\
    GRADUATION INTERNSHIP REPORTS USING \LaTeX}
\end{center}

\vspace{0.5cm}

{\fontsize{13pt}{17pt}\selectfont
\setstretch{1.3}
\justifying

\noindent
\textbf{Note:}
This section is provided for guidance only and is included in the main
document using the command \verb|\input{Regulations.tex}|.
Students shall remove or comment out this command before submitting the
final internship report.

\vspace{0.4cm}

% ================================================================
% I. GENERAL REQUIREMENTS
% ================================================================

\noindent
\textbf{I.\quad GENERAL REQUIREMENTS}

\vspace{0.15cm}

The Graduation Internship Report shall be prepared using the official
\LaTeX{} template provided by the Faculty. Students shall preserve the
predefined page layout, typography, line spacing, heading styles,
numbering system, figure and table formatting, and bibliography settings
unless otherwise instructed by the Faculty or the supervisor.

The report shall be written in English.

Students are responsible for the academic and technical content of the
report. The official \LaTeX{} template is responsible for standardized
document formatting.

\vspace{0.3cm}

% ================================================================
% II. REPORT STRUCTURE
% ================================================================

\noindent
\textbf{II.\quad REPORT STRUCTURE}

\vspace{0.15cm}

The report shall normally be arranged in the following order:

\begin{enumerate}[label=\arabic*., leftmargin=1.2cm, itemsep=2pt]
    \item Outer Cover;
    \item Comments of the Internship Host Organization;
    \item Supervisor's Evaluation and Graduation Internship Defense Results;
    \item Table of Contents;
    \item List of Symbols and Abbreviations, if applicable;
    \item List of Figures;
    \item List of Tables;
    \item Abstract;
    \item Main Chapters of the Internship Report;
    \item References;
    \item Appendices, if applicable.
\end{enumerate}

The preliminary pages shall not be numbered as chapters. The main report
content shall be organized into consecutively numbered chapters beginning
with Chapter 1.

\vspace{0.3cm}

% ================================================================
% III. PAGE SIZE AND MARGINS
% ================================================================

\noindent
\textbf{III.\quad PAGE SIZE AND MARGINS}

\vspace{0.15cm}

The report shall be formatted on A4 paper with the following margins:

\begin{center}
\begin{tabular}{ll}
    Left margin:   & 3.5 cm \\
    Right margin:  & 2.0 cm \\
    Top margin:    & 2.5 cm \\
    Bottom margin: & 2.5 cm
\end{tabular}
\end{center}

These settings are predefined in the official template and shall not be
modified by students.

\vspace{0.3cm}

% ================================================================
% IV. FONT, FONT SIZE, AND LINE SPACING
% ================================================================

\noindent
\textbf{IV.\quad FONT, FONT SIZE, AND LINE SPACING}

\vspace{0.15cm}

The main text shall use a 13 pt font size with 1.3-line spacing.
A Times-family typeface shall be used in accordance with the official
template.

The text shall be fully justified.

Students shall not use repeated spaces, repeated manual line breaks, or
multiple blank lines to position text. Paragraph and heading spacing shall
be controlled by the template.

\vspace{0.3cm}

% ================================================================
% V. OUTER COVER
% ================================================================

\noindent
\textbf{V.\quad OUTER COVER}

\vspace{0.15cm}

The outer cover shall follow the official format provided in the template
and shall contain the following information:

\begin{itemize}[leftmargin=0.8cm, itemsep=2pt]
    \item Thai Nguyen University of Technology;
    \item Faculty of International Training;
    \item Engineering and Technology Department;
    \item University logo and Faculty logo;
    \item the title \textbf{INTERNSHIP REPORT};
    \item Internship Organization;
    \item Author;
    \item Class;
    \item Supervisor;
    \item place and year of completion.
\end{itemize}

Students shall only replace the required information fields and shall not
modify the cover-page layout unless instructed otherwise.

\vspace{0.3cm}

% ================================================================
% VI. COMMENTS OF THE INTERNSHIP HOST ORGANIZATION
% ================================================================

\noindent
\textbf{VI.\quad COMMENTS OF THE INTERNSHIP HOST ORGANIZATION}

\vspace{0.15cm}

The page entitled \textbf{COMMENTS OF THE INTERNSHIP HOST ORGANIZATION}
is reserved for the internship host organization.

It includes:

\begin{itemize}[leftmargin=0.8cm, itemsep=2pt]
    \item the student's full name;
    \item class;
    \item comments from the host organization;
    \item score;
    \item date;
    \item signature and official seal.
\end{itemize}

Students shall not alter or complete official evaluation fields unless
specifically authorized.

\vspace{0.3cm}

% ================================================================
% VII. SUPERVISOR'S EVALUATION AND DEFENSE RESULTS
% ================================================================

\noindent
\textbf{VII.\quad SUPERVISOR'S EVALUATION AND DEFENSE RESULTS}

\vspace{0.15cm}

The section entitled
\textbf{SUPERVISOR'S EVALUATION AND GRADUATION INTERNSHIP DEFENSE RESULTS}
is reserved for the supervisor and defense committee.

The evaluation may include:

\begin{itemize}[leftmargin=0.8cm, itemsep=2pt]
    \item Internship Progress and Student's Attitude;
    \item Report Content;
    \item Report Presentation;
    \item Other Comments;
    \item Supervisor's Opinion;
    \item Defense Result;
    \item Score.
\end{itemize}

Students shall not enter evaluation scores or defense results before they
are officially determined.

\vspace{0.3cm}

% ================================================================
% VIII. CHAPTERS, SECTIONS, AND SUBSECTIONS
% ================================================================

\noindent
\textbf{VIII.\quad CHAPTERS, SECTIONS, AND SUBSECTIONS}

\vspace{0.15cm}

The report structure shall be generated using the standard \LaTeX{}
commands:

\begin{verbatim}
\chapter{...}
\section{...}
\subsection{...}
\end{verbatim}

Chapter, section, and subsection numbers shall be generated automatically
by \LaTeX{}.

Students shall not manually type chapter or section numbers and shall not
reproduce heading styles using commands such as \verb|\textbf|,
\verb|\fontsize|, or \verb|\centering| instead of the appropriate
structural command.

Excessive subdivision should be avoided. In most cases, the hierarchy
Chapter--Section--Subsection is sufficient.

\vspace{0.3cm}

% ================================================================
% IX. ORGANIZATION OF THE MAIN CONTENT
% ================================================================

\noindent
\textbf{IX.\quad ORGANIZATION OF THE MAIN CONTENT}

\vspace{0.15cm}

The specific content of the report shall depend on the internship
organization and the assigned tasks. A typical structure may include:

\begin{itemize}[leftmargin=0.8cm, itemsep=2pt]
    \item \textbf{Chapter 1. Overview of the Internship Organization};
    \item \textbf{Chapter 2. Internship Activities and Assigned Tasks};
    \item \textbf{Chapter 3. Results, Experience, and Evaluation}.
\end{itemize}

These chapter titles are recommendations only and may be adapted to the
actual internship content.

The report shall reflect work actually performed, observed, analyzed, or
learned during the internship. Unrelated theoretical material should not
be included merely to increase report length.

\vspace{0.3cm}

% ================================================================
% X. TABLE OF CONTENTS
% ================================================================

\noindent
\textbf{X.\quad TABLE OF CONTENTS}

\vspace{0.15cm}

The Table of Contents shall be generated automatically using:

\begin{verbatim}
\tableofcontents
\end{verbatim}

Students shall not create or manually edit page numbers in the Table of
Contents.

The document shall be compiled sufficiently to ensure that all page
numbers and entries are correctly updated.

\vspace{0.3cm}

% ================================================================
% XI. LIST OF SYMBOLS AND ABBREVIATIONS
% ================================================================

\noindent
\textbf{XI.\quad LIST OF SYMBOLS AND ABBREVIATIONS}

\vspace{0.15cm}

Symbols and abbreviations that are used repeatedly in the report shall be
listed in the \textbf{List of Symbols and Abbreviations}.

The list may contain the following columns:

\begin{center}
\begin{tabular}{llll}
    No. & Symbol / Abbreviation & Unit & Definition
\end{tabular}
\end{center}

Only symbols and abbreviations actually used in the final report shall be
retained. Example entries contained in the template shall be removed if
they are not relevant to the student's report.

Each abbreviation should normally be defined when it first appears in
the main text.

\vspace{0.3cm}

% ================================================================
% XII. FIGURES, DIAGRAMS, AND PHOTOGRAPHS
% ================================================================

\noindent
\textbf{XII.\quad FIGURES, DIAGRAMS, AND PHOTOGRAPHS}

\vspace{0.15cm}

Figures shall be clear, legible, and directly relevant to the report
content.

Figures shall be inserted using the standard \verb|figure| environment.

Example:

\begin{verbatim}
\begin{figure}[H]
    \centering
    \includegraphics[width=0.8\textwidth]
                    {figures/company_structure.pdf}
    \caption{Organizational structure of the company}
    \label{fig:company-structure}
\end{figure}
\end{verbatim}

Figure captions shall be placed below the figure.

Each referenced figure shall have a unique \verb|\label| and shall be
referenced in the text using, for example:

\begin{verbatim}
Figure~\ref{fig:company-structure}
\end{verbatim}

Students shall not manually type figure numbers.

Figures reproduced or adapted from external sources shall include an
appropriate citation.

Photographs taken during the internship shall be relevant and sufficiently
clear. Students shall not include confidential, restricted, or proprietary
information without permission from the internship organization.

Vector graphics or high-resolution images are recommended whenever
possible.

\vspace{0.3cm}

% ================================================================
% XIII. TABLES
% ================================================================

\noindent
\textbf{XIII.\quad TABLES}

\vspace{0.15cm}

Tables shall be created using suitable \LaTeX{} environments such as
\verb|table|, \verb|tabular|, \verb|tabularx|, or \verb|longtable|.

Table captions shall be placed above the table.

Example:

\begin{verbatim}
\begin{table}[H]
    \centering
    \caption{Internship schedule}
    \label{tab:internship-schedule}

    \begin{tabular}{lll}
        \toprule
        Week & Task & Result \\
        \midrule
        1 & Orientation & Completed \\
        2 & System study & Completed \\
        \bottomrule
    \end{tabular}
\end{table}
\end{verbatim}

Tables shall be referenced using:

\begin{verbatim}
Table~\ref{tab:internship-schedule}
\end{verbatim}

Students shall not manually type table numbers.

Tables should be created directly in \LaTeX{} rather than inserted as
screenshots whenever practical.

\vspace{0.3cm}

% ================================================================
% XIV. EQUATIONS AND TECHNICAL QUANTITIES
% ================================================================

\noindent
\textbf{XIV.\quad EQUATIONS AND TECHNICAL QUANTITIES}

\vspace{0.15cm}

Mathematical expressions shall be typeset directly using \LaTeX{}.
Images or screenshots of equations shall not be used.

Short expressions may be written in inline mathematical mode, for example:

\begin{verbatim}
$R = 0.7\,\Omega$
\end{verbatim}

Important equations shall use a numbered equation environment.

Example:

\begin{verbatim}
\begin{equation}
    P = UI
    \label{eq:power}
\end{equation}
\end{verbatim}

Equations shall be referenced using:

\begin{verbatim}
Equation~\eqref{eq:power}
\end{verbatim}

Equation numbers shall not be manually entered.

SI units shall be used wherever applicable. The \verb|siunitx| package
should be used where appropriate.

Examples:

\begin{verbatim}
\SI{220}{\volt}
\SI{50}{\hertz}
\SI{2.5}{\ampere}
\end{verbatim}

Notation and units shall be consistent throughout the report.

\vspace{0.3cm}

% ================================================================
% XV. SOURCE CODE
% ================================================================

\noindent
\textbf{XV.\quad SOURCE CODE}

\vspace{0.15cm}

If the internship involves programming, PLCs, microcontrollers,
simulation, or software development, relevant source code may be included
in the report.

Source code shall be presented using the \verb|listings| environment or
another environment defined by the official template.

Example:

\begin{verbatim}
\begin{lstlisting}[language=st]
IF Start = TRUE THEN
    Motor := TRUE;
ELSE
    Motor := FALSE;
END_IF;
\end{lstlisting}
\end{verbatim}

Only code necessary for explaining the work should appear in the main
report. Long source-code listings should be moved to an appendix or
submitted as supplementary files if required.

Source code should not normally be inserted as screenshots.

\vspace{0.3cm}

% ================================================================
% XVI. ABSTRACT
% ================================================================

\noindent
\textbf{XVI.\quad ABSTRACT}

\vspace{0.15cm}

The Abstract shall briefly summarize:

\begin{itemize}[leftmargin=0.8cm, itemsep=2pt]
    \item the internship organization;
    \item the objectives of the internship;
    \item the main tasks performed;
    \item the main results, knowledge, and skills obtained.
\end{itemize}

The Abstract should be concise and should not reproduce the full content
of the report chapter by chapter.

\vspace{0.3cm}

% ================================================================
% XVII. REFERENCES AND CITATIONS
% ================================================================

\noindent
\textbf{XVII.\quad REFERENCES AND CITATIONS}

\vspace{0.15cm}

All books, journal papers, conference papers, standards, websites,
datasheets, manuals, technical documents, figures, tables, or other
external sources used in preparing the report shall be properly cited.

References shall be managed through the bibliography system provided with
the official template.

Bibliographic entries shall normally be stored in:

\begin{verbatim}
References.bib
\end{verbatim}

References shall be cited in the report using:

\begin{verbatim}
\cite{reference-key}
\end{verbatim}

Students shall not manually type reference numbers such as ``[1]'',
``[2]'', etc.

Titles of foreign-language publications shall remain in their original
language in the bibliography.

All borrowed ideas, data, figures, tables, equations, methods, or results
shall be properly acknowledged.

\vspace{0.3cm}

% ================================================================
% XVIII. CROSS-REFERENCING
% ================================================================

\noindent
\textbf{XVIII.\quad CROSS-REFERENCING}

\vspace{0.15cm}

Numbered elements shall be cross-referenced using the standard \LaTeX{}
mechanism.

Use the following commands where appropriate:

\begin{verbatim}
\label{...}
\ref{...}
\eqref{...}
\end{verbatim}

For example:

\begin{verbatim}
As shown in Figure~\ref{fig:system}, ...
\end{verbatim}

Students shall not manually type numbers such as ``Figure 2.3'' when
those numbers can be generated automatically.

\vspace{0.3cm}

% ================================================================
% XIX. PAGE NUMBERING
% ================================================================

\noindent
\textbf{XIX.\quad PAGE NUMBERING}

\vspace{0.15cm}

Cover pages and official evaluation forms at the beginning of the report
shall not display page numbers.

The main body of the report shall use Arabic page numbers beginning from
Chapter 1.

Page numbering shall be generated automatically by the official template.
Students shall not manually enter page numbers.

\vspace{0.3cm}

% ================================================================
% XX. PROJECT FILE ORGANIZATION
% ================================================================

\noindent
\textbf{XX.\quad PROJECT FILE ORGANIZATION}

\vspace{0.15cm}

Students should maintain a clear and organized \LaTeX{} project
structure. A recommended structure is:

\begin{verbatim}
InternshipReport/
|
|-- main.tex
|-- References.bib
|
|-- figures/
|   |-- company.pdf
|   |-- system.png
|   `-- ...
|
`-- chapters/
    |-- chapter1.tex
    |-- chapter2.tex
    `-- chapter3.tex
\end{verbatim}

Figures should normally be stored in the \verb|figures| directory
specified by the official template.

File and directory names should preferably contain Latin letters,
numbers, hyphens, or underscores.

Recommended examples include:

\begin{verbatim}
company_structure.pdf
motor_control.png
chapter_01.tex
\end{verbatim}

Spaces and unnecessary special characters in file names should be
avoided where possible.

\vspace{0.3cm}

% ================================================================
% XXI. DOCUMENT COMPILATION AND FINAL CHECK
% ================================================================

\noindent
\textbf{XXI.\quad DOCUMENT COMPILATION AND FINAL CHECK}

\vspace{0.15cm}

Before submission, the complete \LaTeX{} project shall compile
successfully and produce the final PDF report.

The final document shall not contain unresolved references such as:

\begin{verbatim}
??
[?]
\end{verbatim}

Before generating the final PDF, students shall check:

\begin{itemize}[leftmargin=0.8cm, itemsep=2pt]
    \item the Table of Contents;
    \item page numbering;
    \item chapter, section, and subsection numbering;
    \item figure numbering;
    \item table numbering;
    \item equation numbering;
    \item cross-references;
    \item citations;
    \item bibliography;
    \item font display;
    \item mathematical symbols;
    \item image quality;
    \item completeness of all required forms.
\end{itemize}

The document may be compiled using TeX Live, MiKTeX, or another
compatible \LaTeX{} environment. Regardless of the compilation
environment, the resulting PDF shall conform to the official Faculty
template.

\vspace{0.3cm}

% ================================================================
% XXII. CONTENT THAT STUDENTS MAY EDIT
% ================================================================

\noindent
\textbf{XXII.\quad CONTENT THAT STUDENTS MAY EDIT}

\vspace{0.15cm}

Students may edit or replace:

\begin{itemize}[leftmargin=0.8cm, itemsep=2pt]
    \item student name;
    \item class;
    \item internship organization;
    \item supervisor name;
    \item Abstract;
    \item List of Symbols and Abbreviations;
    \item chapter titles and chapter content;
    \item figures and photographs;
    \item tables;
    \item equations;
    \item citations and references.
\end{itemize}

Students shall not modify the main document formatting settings in the
preamble, including the document class, page geometry, line spacing,
heading formatting, table-of-contents formatting, font configuration,
page layout, or cover-page macros unless specifically instructed.

\vspace{0.3cm}

% ================================================================
% XXIII. FINAL SUBMISSION REQUIREMENTS
% ================================================================

\noindent
\textbf{XXIII.\quad FINAL SUBMISSION REQUIREMENTS}

\vspace{0.15cm}

Before submission, the final report shall:

\begin{itemize}[leftmargin=0.8cm, itemsep=2pt]
    \item compile without errors that affect the report content or presentation;
    \item contain all required forms and report sections;
    \item contain correct student, class, internship organization, and supervisor information;
    \item contain an updated Table of Contents, List of Figures, and List of Tables;
    \item contain no unused example content from the template;
    \item contain no placeholder text such as \textit{Title of Chapter 1},
    \textit{Section 1.1}, or other sample entries;
    \item contain all required citations and references;
    \item contain no unresolved cross-references.
\end{itemize}

If the Faculty requires submission of the \LaTeX{} source, the complete
project shall include all \verb|.tex| files, the bibliography database,
figures, and any other files required to reproduce the submitted PDF.

\vspace{0.4cm}

\noindent
\textbf{Important:}
Before submitting the final internship report, students shall remove or
comment out the following command in the main document:

\begin{verbatim}
\input{Regulations.tex}
\end{verbatim}

For example:

\begin{verbatim}
% \input{Regulations.tex}
\end{verbatim}

This removes the guidance pages from the final submitted report while
preserving the regulations file in the \LaTeX{} project for reference.

} % End 13pt content

\clearpage

%
% IMPORTANT:
% This section is provided for guidance only. Students shall remove
% or comment out the 
% =================================================================
% REGULATIONS FOR PREPARING THE GRADUATION INTERNSHIP REPORT
% USING LATEX
% =================================================================
%
% This file is intended to be included from the main document with:
%
%     \input{Regulations.tex}
%
% IMPORTANT:
% This section is provided for guidance only. Students shall remove
% or comment out the \input{Regulations.tex} command before submitting
% the final internship report.
%
% =================================================================

\clearpage
\thispagestyle{empty}

\begin{center}
    {\fontsize{16pt}{20pt}\selectfont\bfseries
    SPECIFIC REGULATIONS FOR PREPARING\\
    GRADUATION INTERNSHIP REPORTS USING \LaTeX}
\end{center}

\vspace{0.5cm}

{\fontsize{13pt}{17pt}\selectfont
\setstretch{1.3}
\justifying

\noindent
\textbf{Note:}
This section is provided for guidance only and is included in the main
document using the command \verb|\input{Regulations.tex}|.
Students shall remove or comment out this command before submitting the
final internship report.

\vspace{0.4cm}

% ================================================================
% I. GENERAL REQUIREMENTS
% ================================================================

\noindent
\textbf{I.\quad GENERAL REQUIREMENTS}

\vspace{0.15cm}

The Graduation Internship Report shall be prepared using the official
\LaTeX{} template provided by the Faculty. Students shall preserve the
predefined page layout, typography, line spacing, heading styles,
numbering system, figure and table formatting, and bibliography settings
unless otherwise instructed by the Faculty or the supervisor.

The report shall be written in English.

Students are responsible for the academic and technical content of the
report. The official \LaTeX{} template is responsible for standardized
document formatting.

\vspace{0.3cm}

% ================================================================
% II. REPORT STRUCTURE
% ================================================================

\noindent
\textbf{II.\quad REPORT STRUCTURE}

\vspace{0.15cm}

The report shall normally be arranged in the following order:

\begin{enumerate}[label=\arabic*., leftmargin=1.2cm, itemsep=2pt]
    \item Outer Cover;
    \item Comments of the Internship Host Organization;
    \item Supervisor's Evaluation and Graduation Internship Defense Results;
    \item Table of Contents;
    \item List of Symbols and Abbreviations, if applicable;
    \item List of Figures;
    \item List of Tables;
    \item Abstract;
    \item Main Chapters of the Internship Report;
    \item References;
    \item Appendices, if applicable.
\end{enumerate}

The preliminary pages shall not be numbered as chapters. The main report
content shall be organized into consecutively numbered chapters beginning
with Chapter 1.

\vspace{0.3cm}

% ================================================================
% III. PAGE SIZE AND MARGINS
% ================================================================

\noindent
\textbf{III.\quad PAGE SIZE AND MARGINS}

\vspace{0.15cm}

The report shall be formatted on A4 paper with the following margins:

\begin{center}
\begin{tabular}{ll}
    Left margin:   & 3.5 cm \\
    Right margin:  & 2.0 cm \\
    Top margin:    & 2.5 cm \\
    Bottom margin: & 2.5 cm
\end{tabular}
\end{center}

These settings are predefined in the official template and shall not be
modified by students.

\vspace{0.3cm}

% ================================================================
% IV. FONT, FONT SIZE, AND LINE SPACING
% ================================================================

\noindent
\textbf{IV.\quad FONT, FONT SIZE, AND LINE SPACING}

\vspace{0.15cm}

The main text shall use a 13 pt font size with 1.3-line spacing.
A Times-family typeface shall be used in accordance with the official
template.

The text shall be fully justified.

Students shall not use repeated spaces, repeated manual line breaks, or
multiple blank lines to position text. Paragraph and heading spacing shall
be controlled by the template.

\vspace{0.3cm}

% ================================================================
% V. OUTER COVER
% ================================================================

\noindent
\textbf{V.\quad OUTER COVER}

\vspace{0.15cm}

The outer cover shall follow the official format provided in the template
and shall contain the following information:

\begin{itemize}[leftmargin=0.8cm, itemsep=2pt]
    \item Thai Nguyen University of Technology;
    \item Faculty of International Training;
    \item Engineering and Technology Department;
    \item University logo and Faculty logo;
    \item the title \textbf{INTERNSHIP REPORT};
    \item Internship Organization;
    \item Author;
    \item Class;
    \item Supervisor;
    \item place and year of completion.
\end{itemize}

Students shall only replace the required information fields and shall not
modify the cover-page layout unless instructed otherwise.

\vspace{0.3cm}

% ================================================================
% VI. COMMENTS OF THE INTERNSHIP HOST ORGANIZATION
% ================================================================

\noindent
\textbf{VI.\quad COMMENTS OF THE INTERNSHIP HOST ORGANIZATION}

\vspace{0.15cm}

The page entitled \textbf{COMMENTS OF THE INTERNSHIP HOST ORGANIZATION}
is reserved for the internship host organization.

It includes:

\begin{itemize}[leftmargin=0.8cm, itemsep=2pt]
    \item the student's full name;
    \item class;
    \item comments from the host organization;
    \item score;
    \item date;
    \item signature and official seal.
\end{itemize}

Students shall not alter or complete official evaluation fields unless
specifically authorized.

\vspace{0.3cm}

% ================================================================
% VII. SUPERVISOR'S EVALUATION AND DEFENSE RESULTS
% ================================================================

\noindent
\textbf{VII.\quad SUPERVISOR'S EVALUATION AND DEFENSE RESULTS}

\vspace{0.15cm}

The section entitled
\textbf{SUPERVISOR'S EVALUATION AND GRADUATION INTERNSHIP DEFENSE RESULTS}
is reserved for the supervisor and defense committee.

The evaluation may include:

\begin{itemize}[leftmargin=0.8cm, itemsep=2pt]
    \item Internship Progress and Student's Attitude;
    \item Report Content;
    \item Report Presentation;
    \item Other Comments;
    \item Supervisor's Opinion;
    \item Defense Result;
    \item Score.
\end{itemize}

Students shall not enter evaluation scores or defense results before they
are officially determined.

\vspace{0.3cm}

% ================================================================
% VIII. CHAPTERS, SECTIONS, AND SUBSECTIONS
% ================================================================

\noindent
\textbf{VIII.\quad CHAPTERS, SECTIONS, AND SUBSECTIONS}

\vspace{0.15cm}

The report structure shall be generated using the standard \LaTeX{}
commands:

\begin{verbatim}
\chapter{...}
\section{...}
\subsection{...}
\end{verbatim}

Chapter, section, and subsection numbers shall be generated automatically
by \LaTeX{}.

Students shall not manually type chapter or section numbers and shall not
reproduce heading styles using commands such as \verb|\textbf|,
\verb|\fontsize|, or \verb|\centering| instead of the appropriate
structural command.

Excessive subdivision should be avoided. In most cases, the hierarchy
Chapter--Section--Subsection is sufficient.

\vspace{0.3cm}

% ================================================================
% IX. ORGANIZATION OF THE MAIN CONTENT
% ================================================================

\noindent
\textbf{IX.\quad ORGANIZATION OF THE MAIN CONTENT}

\vspace{0.15cm}

The specific content of the report shall depend on the internship
organization and the assigned tasks. A typical structure may include:

\begin{itemize}[leftmargin=0.8cm, itemsep=2pt]
    \item \textbf{Chapter 1. Overview of the Internship Organization};
    \item \textbf{Chapter 2. Internship Activities and Assigned Tasks};
    \item \textbf{Chapter 3. Results, Experience, and Evaluation}.
\end{itemize}

These chapter titles are recommendations only and may be adapted to the
actual internship content.

The report shall reflect work actually performed, observed, analyzed, or
learned during the internship. Unrelated theoretical material should not
be included merely to increase report length.

\vspace{0.3cm}

% ================================================================
% X. TABLE OF CONTENTS
% ================================================================

\noindent
\textbf{X.\quad TABLE OF CONTENTS}

\vspace{0.15cm}

The Table of Contents shall be generated automatically using:

\begin{verbatim}
\tableofcontents
\end{verbatim}

Students shall not create or manually edit page numbers in the Table of
Contents.

The document shall be compiled sufficiently to ensure that all page
numbers and entries are correctly updated.

\vspace{0.3cm}

% ================================================================
% XI. LIST OF SYMBOLS AND ABBREVIATIONS
% ================================================================

\noindent
\textbf{XI.\quad LIST OF SYMBOLS AND ABBREVIATIONS}

\vspace{0.15cm}

Symbols and abbreviations that are used repeatedly in the report shall be
listed in the \textbf{List of Symbols and Abbreviations}.

The list may contain the following columns:

\begin{center}
\begin{tabular}{llll}
    No. & Symbol / Abbreviation & Unit & Definition
\end{tabular}
\end{center}

Only symbols and abbreviations actually used in the final report shall be
retained. Example entries contained in the template shall be removed if
they are not relevant to the student's report.

Each abbreviation should normally be defined when it first appears in
the main text.

\vspace{0.3cm}

% ================================================================
% XII. FIGURES, DIAGRAMS, AND PHOTOGRAPHS
% ================================================================

\noindent
\textbf{XII.\quad FIGURES, DIAGRAMS, AND PHOTOGRAPHS}

\vspace{0.15cm}

Figures shall be clear, legible, and directly relevant to the report
content.

Figures shall be inserted using the standard \verb|figure| environment.

Example:

\begin{verbatim}
\begin{figure}[H]
    \centering
    \includegraphics[width=0.8\textwidth]
                    {figures/company_structure.pdf}
    \caption{Organizational structure of the company}
    \label{fig:company-structure}
\end{figure}
\end{verbatim}

Figure captions shall be placed below the figure.

Each referenced figure shall have a unique \verb|\label| and shall be
referenced in the text using, for example:

\begin{verbatim}
Figure~\ref{fig:company-structure}
\end{verbatim}

Students shall not manually type figure numbers.

Figures reproduced or adapted from external sources shall include an
appropriate citation.

Photographs taken during the internship shall be relevant and sufficiently
clear. Students shall not include confidential, restricted, or proprietary
information without permission from the internship organization.

Vector graphics or high-resolution images are recommended whenever
possible.

\vspace{0.3cm}

% ================================================================
% XIII. TABLES
% ================================================================

\noindent
\textbf{XIII.\quad TABLES}

\vspace{0.15cm}

Tables shall be created using suitable \LaTeX{} environments such as
\verb|table|, \verb|tabular|, \verb|tabularx|, or \verb|longtable|.

Table captions shall be placed above the table.

Example:

\begin{verbatim}
\begin{table}[H]
    \centering
    \caption{Internship schedule}
    \label{tab:internship-schedule}

    \begin{tabular}{lll}
        \toprule
        Week & Task & Result \\
        \midrule
        1 & Orientation & Completed \\
        2 & System study & Completed \\
        \bottomrule
    \end{tabular}
\end{table}
\end{verbatim}

Tables shall be referenced using:

\begin{verbatim}
Table~\ref{tab:internship-schedule}
\end{verbatim}

Students shall not manually type table numbers.

Tables should be created directly in \LaTeX{} rather than inserted as
screenshots whenever practical.

\vspace{0.3cm}

% ================================================================
% XIV. EQUATIONS AND TECHNICAL QUANTITIES
% ================================================================

\noindent
\textbf{XIV.\quad EQUATIONS AND TECHNICAL QUANTITIES}

\vspace{0.15cm}

Mathematical expressions shall be typeset directly using \LaTeX{}.
Images or screenshots of equations shall not be used.

Short expressions may be written in inline mathematical mode, for example:

\begin{verbatim}
$R = 0.7\,\Omega$
\end{verbatim}

Important equations shall use a numbered equation environment.

Example:

\begin{verbatim}
\begin{equation}
    P = UI
    \label{eq:power}
\end{equation}
\end{verbatim}

Equations shall be referenced using:

\begin{verbatim}
Equation~\eqref{eq:power}
\end{verbatim}

Equation numbers shall not be manually entered.

SI units shall be used wherever applicable. The \verb|siunitx| package
should be used where appropriate.

Examples:

\begin{verbatim}
\SI{220}{\volt}
\SI{50}{\hertz}
\SI{2.5}{\ampere}
\end{verbatim}

Notation and units shall be consistent throughout the report.

\vspace{0.3cm}

% ================================================================
% XV. SOURCE CODE
% ================================================================

\noindent
\textbf{XV.\quad SOURCE CODE}

\vspace{0.15cm}

If the internship involves programming, PLCs, microcontrollers,
simulation, or software development, relevant source code may be included
in the report.

Source code shall be presented using the \verb|listings| environment or
another environment defined by the official template.

Example:

\begin{verbatim}
\begin{lstlisting}[language=st]
IF Start = TRUE THEN
    Motor := TRUE;
ELSE
    Motor := FALSE;
END_IF;
\end{lstlisting}
\end{verbatim}

Only code necessary for explaining the work should appear in the main
report. Long source-code listings should be moved to an appendix or
submitted as supplementary files if required.

Source code should not normally be inserted as screenshots.

\vspace{0.3cm}

% ================================================================
% XVI. ABSTRACT
% ================================================================

\noindent
\textbf{XVI.\quad ABSTRACT}

\vspace{0.15cm}

The Abstract shall briefly summarize:

\begin{itemize}[leftmargin=0.8cm, itemsep=2pt]
    \item the internship organization;
    \item the objectives of the internship;
    \item the main tasks performed;
    \item the main results, knowledge, and skills obtained.
\end{itemize}

The Abstract should be concise and should not reproduce the full content
of the report chapter by chapter.

\vspace{0.3cm}

% ================================================================
% XVII. REFERENCES AND CITATIONS
% ================================================================

\noindent
\textbf{XVII.\quad REFERENCES AND CITATIONS}

\vspace{0.15cm}

All books, journal papers, conference papers, standards, websites,
datasheets, manuals, technical documents, figures, tables, or other
external sources used in preparing the report shall be properly cited.

References shall be managed through the bibliography system provided with
the official template.

Bibliographic entries shall normally be stored in:

\begin{verbatim}
References.bib
\end{verbatim}

References shall be cited in the report using:

\begin{verbatim}
\cite{reference-key}
\end{verbatim}

Students shall not manually type reference numbers such as ``[1]'',
``[2]'', etc.

Titles of foreign-language publications shall remain in their original
language in the bibliography.

All borrowed ideas, data, figures, tables, equations, methods, or results
shall be properly acknowledged.

\vspace{0.3cm}

% ================================================================
% XVIII. CROSS-REFERENCING
% ================================================================

\noindent
\textbf{XVIII.\quad CROSS-REFERENCING}

\vspace{0.15cm}

Numbered elements shall be cross-referenced using the standard \LaTeX{}
mechanism.

Use the following commands where appropriate:

\begin{verbatim}
\label{...}
\ref{...}
\eqref{...}
\end{verbatim}

For example:

\begin{verbatim}
As shown in Figure~\ref{fig:system}, ...
\end{verbatim}

Students shall not manually type numbers such as ``Figure 2.3'' when
those numbers can be generated automatically.

\vspace{0.3cm}

% ================================================================
% XIX. PAGE NUMBERING
% ================================================================

\noindent
\textbf{XIX.\quad PAGE NUMBERING}

\vspace{0.15cm}

Cover pages and official evaluation forms at the beginning of the report
shall not display page numbers.

The main body of the report shall use Arabic page numbers beginning from
Chapter 1.

Page numbering shall be generated automatically by the official template.
Students shall not manually enter page numbers.

\vspace{0.3cm}

% ================================================================
% XX. PROJECT FILE ORGANIZATION
% ================================================================

\noindent
\textbf{XX.\quad PROJECT FILE ORGANIZATION}

\vspace{0.15cm}

Students should maintain a clear and organized \LaTeX{} project
structure. A recommended structure is:

\begin{verbatim}
InternshipReport/
|
|-- main.tex
|-- References.bib
|
|-- figures/
|   |-- company.pdf
|   |-- system.png
|   `-- ...
|
`-- chapters/
    |-- chapter1.tex
    |-- chapter2.tex
    `-- chapter3.tex
\end{verbatim}

Figures should normally be stored in the \verb|figures| directory
specified by the official template.

File and directory names should preferably contain Latin letters,
numbers, hyphens, or underscores.

Recommended examples include:

\begin{verbatim}
company_structure.pdf
motor_control.png
chapter_01.tex
\end{verbatim}

Spaces and unnecessary special characters in file names should be
avoided where possible.

\vspace{0.3cm}

% ================================================================
% XXI. DOCUMENT COMPILATION AND FINAL CHECK
% ================================================================

\noindent
\textbf{XXI.\quad DOCUMENT COMPILATION AND FINAL CHECK}

\vspace{0.15cm}

Before submission, the complete \LaTeX{} project shall compile
successfully and produce the final PDF report.

The final document shall not contain unresolved references such as:

\begin{verbatim}
??
[?]
\end{verbatim}

Before generating the final PDF, students shall check:

\begin{itemize}[leftmargin=0.8cm, itemsep=2pt]
    \item the Table of Contents;
    \item page numbering;
    \item chapter, section, and subsection numbering;
    \item figure numbering;
    \item table numbering;
    \item equation numbering;
    \item cross-references;
    \item citations;
    \item bibliography;
    \item font display;
    \item mathematical symbols;
    \item image quality;
    \item completeness of all required forms.
\end{itemize}

The document may be compiled using TeX Live, MiKTeX, or another
compatible \LaTeX{} environment. Regardless of the compilation
environment, the resulting PDF shall conform to the official Faculty
template.

\vspace{0.3cm}

% ================================================================
% XXII. CONTENT THAT STUDENTS MAY EDIT
% ================================================================

\noindent
\textbf{XXII.\quad CONTENT THAT STUDENTS MAY EDIT}

\vspace{0.15cm}

Students may edit or replace:

\begin{itemize}[leftmargin=0.8cm, itemsep=2pt]
    \item student name;
    \item class;
    \item internship organization;
    \item supervisor name;
    \item Abstract;
    \item List of Symbols and Abbreviations;
    \item chapter titles and chapter content;
    \item figures and photographs;
    \item tables;
    \item equations;
    \item citations and references.
\end{itemize}

Students shall not modify the main document formatting settings in the
preamble, including the document class, page geometry, line spacing,
heading formatting, table-of-contents formatting, font configuration,
page layout, or cover-page macros unless specifically instructed.

\vspace{0.3cm}

% ================================================================
% XXIII. FINAL SUBMISSION REQUIREMENTS
% ================================================================

\noindent
\textbf{XXIII.\quad FINAL SUBMISSION REQUIREMENTS}

\vspace{0.15cm}

Before submission, the final report shall:

\begin{itemize}[leftmargin=0.8cm, itemsep=2pt]
    \item compile without errors that affect the report content or presentation;
    \item contain all required forms and report sections;
    \item contain correct student, class, internship organization, and supervisor information;
    \item contain an updated Table of Contents, List of Figures, and List of Tables;
    \item contain no unused example content from the template;
    \item contain no placeholder text such as \textit{Title of Chapter 1},
    \textit{Section 1.1}, or other sample entries;
    \item contain all required citations and references;
    \item contain no unresolved cross-references.
\end{itemize}

If the Faculty requires submission of the \LaTeX{} source, the complete
project shall include all \verb|.tex| files, the bibliography database,
figures, and any other files required to reproduce the submitted PDF.

\vspace{0.4cm}

\noindent
\textbf{Important:}
Before submitting the final internship report, students shall remove or
comment out the following command in the main document:

\begin{verbatim}
\input{Regulations.tex}
\end{verbatim}

For example:

\begin{verbatim}
% \input{Regulations.tex}
\end{verbatim}

This removes the guidance pages from the final submitted report while
preserving the regulations file in the \LaTeX{} project for reference.

} % End 13pt content

\clearpage
 command before submitting
% the final internship report.
%
% =================================================================

\clearpage
\thispagestyle{empty}

\begin{center}
    {\fontsize{16pt}{20pt}\selectfont\bfseries
    SPECIFIC REGULATIONS FOR PREPARING\\
    GRADUATION INTERNSHIP REPORTS USING \LaTeX}
\end{center}

\vspace{0.5cm}

{\fontsize{13pt}{17pt}\selectfont
\setstretch{1.3}
\justifying

\noindent
\textbf{Note:}
This section is provided for guidance only and is included in the main
document using the command \verb|
% =================================================================
% REGULATIONS FOR PREPARING THE GRADUATION INTERNSHIP REPORT
% USING LATEX
% =================================================================
%
% This file is intended to be included from the main document with:
%
%     \input{Regulations.tex}
%
% IMPORTANT:
% This section is provided for guidance only. Students shall remove
% or comment out the \input{Regulations.tex} command before submitting
% the final internship report.
%
% =================================================================

\clearpage
\thispagestyle{empty}

\begin{center}
    {\fontsize{16pt}{20pt}\selectfont\bfseries
    SPECIFIC REGULATIONS FOR PREPARING\\
    GRADUATION INTERNSHIP REPORTS USING \LaTeX}
\end{center}

\vspace{0.5cm}

{\fontsize{13pt}{17pt}\selectfont
\setstretch{1.3}
\justifying

\noindent
\textbf{Note:}
This section is provided for guidance only and is included in the main
document using the command \verb|\input{Regulations.tex}|.
Students shall remove or comment out this command before submitting the
final internship report.

\vspace{0.4cm}

% ================================================================
% I. GENERAL REQUIREMENTS
% ================================================================

\noindent
\textbf{I.\quad GENERAL REQUIREMENTS}

\vspace{0.15cm}

The Graduation Internship Report shall be prepared using the official
\LaTeX{} template provided by the Faculty. Students shall preserve the
predefined page layout, typography, line spacing, heading styles,
numbering system, figure and table formatting, and bibliography settings
unless otherwise instructed by the Faculty or the supervisor.

The report shall be written in English.

Students are responsible for the academic and technical content of the
report. The official \LaTeX{} template is responsible for standardized
document formatting.

\vspace{0.3cm}

% ================================================================
% II. REPORT STRUCTURE
% ================================================================

\noindent
\textbf{II.\quad REPORT STRUCTURE}

\vspace{0.15cm}

The report shall normally be arranged in the following order:

\begin{enumerate}[label=\arabic*., leftmargin=1.2cm, itemsep=2pt]
    \item Outer Cover;
    \item Comments of the Internship Host Organization;
    \item Supervisor's Evaluation and Graduation Internship Defense Results;
    \item Table of Contents;
    \item List of Symbols and Abbreviations, if applicable;
    \item List of Figures;
    \item List of Tables;
    \item Abstract;
    \item Main Chapters of the Internship Report;
    \item References;
    \item Appendices, if applicable.
\end{enumerate}

The preliminary pages shall not be numbered as chapters. The main report
content shall be organized into consecutively numbered chapters beginning
with Chapter 1.

\vspace{0.3cm}

% ================================================================
% III. PAGE SIZE AND MARGINS
% ================================================================

\noindent
\textbf{III.\quad PAGE SIZE AND MARGINS}

\vspace{0.15cm}

The report shall be formatted on A4 paper with the following margins:

\begin{center}
\begin{tabular}{ll}
    Left margin:   & 3.5 cm \\
    Right margin:  & 2.0 cm \\
    Top margin:    & 2.5 cm \\
    Bottom margin: & 2.5 cm
\end{tabular}
\end{center}

These settings are predefined in the official template and shall not be
modified by students.

\vspace{0.3cm}

% ================================================================
% IV. FONT, FONT SIZE, AND LINE SPACING
% ================================================================

\noindent
\textbf{IV.\quad FONT, FONT SIZE, AND LINE SPACING}

\vspace{0.15cm}

The main text shall use a 13 pt font size with 1.3-line spacing.
A Times-family typeface shall be used in accordance with the official
template.

The text shall be fully justified.

Students shall not use repeated spaces, repeated manual line breaks, or
multiple blank lines to position text. Paragraph and heading spacing shall
be controlled by the template.

\vspace{0.3cm}

% ================================================================
% V. OUTER COVER
% ================================================================

\noindent
\textbf{V.\quad OUTER COVER}

\vspace{0.15cm}

The outer cover shall follow the official format provided in the template
and shall contain the following information:

\begin{itemize}[leftmargin=0.8cm, itemsep=2pt]
    \item Thai Nguyen University of Technology;
    \item Faculty of International Training;
    \item Engineering and Technology Department;
    \item University logo and Faculty logo;
    \item the title \textbf{INTERNSHIP REPORT};
    \item Internship Organization;
    \item Author;
    \item Class;
    \item Supervisor;
    \item place and year of completion.
\end{itemize}

Students shall only replace the required information fields and shall not
modify the cover-page layout unless instructed otherwise.

\vspace{0.3cm}

% ================================================================
% VI. COMMENTS OF THE INTERNSHIP HOST ORGANIZATION
% ================================================================

\noindent
\textbf{VI.\quad COMMENTS OF THE INTERNSHIP HOST ORGANIZATION}

\vspace{0.15cm}

The page entitled \textbf{COMMENTS OF THE INTERNSHIP HOST ORGANIZATION}
is reserved for the internship host organization.

It includes:

\begin{itemize}[leftmargin=0.8cm, itemsep=2pt]
    \item the student's full name;
    \item class;
    \item comments from the host organization;
    \item score;
    \item date;
    \item signature and official seal.
\end{itemize}

Students shall not alter or complete official evaluation fields unless
specifically authorized.

\vspace{0.3cm}

% ================================================================
% VII. SUPERVISOR'S EVALUATION AND DEFENSE RESULTS
% ================================================================

\noindent
\textbf{VII.\quad SUPERVISOR'S EVALUATION AND DEFENSE RESULTS}

\vspace{0.15cm}

The section entitled
\textbf{SUPERVISOR'S EVALUATION AND GRADUATION INTERNSHIP DEFENSE RESULTS}
is reserved for the supervisor and defense committee.

The evaluation may include:

\begin{itemize}[leftmargin=0.8cm, itemsep=2pt]
    \item Internship Progress and Student's Attitude;
    \item Report Content;
    \item Report Presentation;
    \item Other Comments;
    \item Supervisor's Opinion;
    \item Defense Result;
    \item Score.
\end{itemize}

Students shall not enter evaluation scores or defense results before they
are officially determined.

\vspace{0.3cm}

% ================================================================
% VIII. CHAPTERS, SECTIONS, AND SUBSECTIONS
% ================================================================

\noindent
\textbf{VIII.\quad CHAPTERS, SECTIONS, AND SUBSECTIONS}

\vspace{0.15cm}

The report structure shall be generated using the standard \LaTeX{}
commands:

\begin{verbatim}
\chapter{...}
\section{...}
\subsection{...}
\end{verbatim}

Chapter, section, and subsection numbers shall be generated automatically
by \LaTeX{}.

Students shall not manually type chapter or section numbers and shall not
reproduce heading styles using commands such as \verb|\textbf|,
\verb|\fontsize|, or \verb|\centering| instead of the appropriate
structural command.

Excessive subdivision should be avoided. In most cases, the hierarchy
Chapter--Section--Subsection is sufficient.

\vspace{0.3cm}

% ================================================================
% IX. ORGANIZATION OF THE MAIN CONTENT
% ================================================================

\noindent
\textbf{IX.\quad ORGANIZATION OF THE MAIN CONTENT}

\vspace{0.15cm}

The specific content of the report shall depend on the internship
organization and the assigned tasks. A typical structure may include:

\begin{itemize}[leftmargin=0.8cm, itemsep=2pt]
    \item \textbf{Chapter 1. Overview of the Internship Organization};
    \item \textbf{Chapter 2. Internship Activities and Assigned Tasks};
    \item \textbf{Chapter 3. Results, Experience, and Evaluation}.
\end{itemize}

These chapter titles are recommendations only and may be adapted to the
actual internship content.

The report shall reflect work actually performed, observed, analyzed, or
learned during the internship. Unrelated theoretical material should not
be included merely to increase report length.

\vspace{0.3cm}

% ================================================================
% X. TABLE OF CONTENTS
% ================================================================

\noindent
\textbf{X.\quad TABLE OF CONTENTS}

\vspace{0.15cm}

The Table of Contents shall be generated automatically using:

\begin{verbatim}
\tableofcontents
\end{verbatim}

Students shall not create or manually edit page numbers in the Table of
Contents.

The document shall be compiled sufficiently to ensure that all page
numbers and entries are correctly updated.

\vspace{0.3cm}

% ================================================================
% XI. LIST OF SYMBOLS AND ABBREVIATIONS
% ================================================================

\noindent
\textbf{XI.\quad LIST OF SYMBOLS AND ABBREVIATIONS}

\vspace{0.15cm}

Symbols and abbreviations that are used repeatedly in the report shall be
listed in the \textbf{List of Symbols and Abbreviations}.

The list may contain the following columns:

\begin{center}
\begin{tabular}{llll}
    No. & Symbol / Abbreviation & Unit & Definition
\end{tabular}
\end{center}

Only symbols and abbreviations actually used in the final report shall be
retained. Example entries contained in the template shall be removed if
they are not relevant to the student's report.

Each abbreviation should normally be defined when it first appears in
the main text.

\vspace{0.3cm}

% ================================================================
% XII. FIGURES, DIAGRAMS, AND PHOTOGRAPHS
% ================================================================

\noindent
\textbf{XII.\quad FIGURES, DIAGRAMS, AND PHOTOGRAPHS}

\vspace{0.15cm}

Figures shall be clear, legible, and directly relevant to the report
content.

Figures shall be inserted using the standard \verb|figure| environment.

Example:

\begin{verbatim}
\begin{figure}[H]
    \centering
    \includegraphics[width=0.8\textwidth]
                    {figures/company_structure.pdf}
    \caption{Organizational structure of the company}
    \label{fig:company-structure}
\end{figure}
\end{verbatim}

Figure captions shall be placed below the figure.

Each referenced figure shall have a unique \verb|\label| and shall be
referenced in the text using, for example:

\begin{verbatim}
Figure~\ref{fig:company-structure}
\end{verbatim}

Students shall not manually type figure numbers.

Figures reproduced or adapted from external sources shall include an
appropriate citation.

Photographs taken during the internship shall be relevant and sufficiently
clear. Students shall not include confidential, restricted, or proprietary
information without permission from the internship organization.

Vector graphics or high-resolution images are recommended whenever
possible.

\vspace{0.3cm}

% ================================================================
% XIII. TABLES
% ================================================================

\noindent
\textbf{XIII.\quad TABLES}

\vspace{0.15cm}

Tables shall be created using suitable \LaTeX{} environments such as
\verb|table|, \verb|tabular|, \verb|tabularx|, or \verb|longtable|.

Table captions shall be placed above the table.

Example:

\begin{verbatim}
\begin{table}[H]
    \centering
    \caption{Internship schedule}
    \label{tab:internship-schedule}

    \begin{tabular}{lll}
        \toprule
        Week & Task & Result \\
        \midrule
        1 & Orientation & Completed \\
        2 & System study & Completed \\
        \bottomrule
    \end{tabular}
\end{table}
\end{verbatim}

Tables shall be referenced using:

\begin{verbatim}
Table~\ref{tab:internship-schedule}
\end{verbatim}

Students shall not manually type table numbers.

Tables should be created directly in \LaTeX{} rather than inserted as
screenshots whenever practical.

\vspace{0.3cm}

% ================================================================
% XIV. EQUATIONS AND TECHNICAL QUANTITIES
% ================================================================

\noindent
\textbf{XIV.\quad EQUATIONS AND TECHNICAL QUANTITIES}

\vspace{0.15cm}

Mathematical expressions shall be typeset directly using \LaTeX{}.
Images or screenshots of equations shall not be used.

Short expressions may be written in inline mathematical mode, for example:

\begin{verbatim}
$R = 0.7\,\Omega$
\end{verbatim}

Important equations shall use a numbered equation environment.

Example:

\begin{verbatim}
\begin{equation}
    P = UI
    \label{eq:power}
\end{equation}
\end{verbatim}

Equations shall be referenced using:

\begin{verbatim}
Equation~\eqref{eq:power}
\end{verbatim}

Equation numbers shall not be manually entered.

SI units shall be used wherever applicable. The \verb|siunitx| package
should be used where appropriate.

Examples:

\begin{verbatim}
\SI{220}{\volt}
\SI{50}{\hertz}
\SI{2.5}{\ampere}
\end{verbatim}

Notation and units shall be consistent throughout the report.

\vspace{0.3cm}

% ================================================================
% XV. SOURCE CODE
% ================================================================

\noindent
\textbf{XV.\quad SOURCE CODE}

\vspace{0.15cm}

If the internship involves programming, PLCs, microcontrollers,
simulation, or software development, relevant source code may be included
in the report.

Source code shall be presented using the \verb|listings| environment or
another environment defined by the official template.

Example:

\begin{verbatim}
\begin{lstlisting}[language=st]
IF Start = TRUE THEN
    Motor := TRUE;
ELSE
    Motor := FALSE;
END_IF;
\end{lstlisting}
\end{verbatim}

Only code necessary for explaining the work should appear in the main
report. Long source-code listings should be moved to an appendix or
submitted as supplementary files if required.

Source code should not normally be inserted as screenshots.

\vspace{0.3cm}

% ================================================================
% XVI. ABSTRACT
% ================================================================

\noindent
\textbf{XVI.\quad ABSTRACT}

\vspace{0.15cm}

The Abstract shall briefly summarize:

\begin{itemize}[leftmargin=0.8cm, itemsep=2pt]
    \item the internship organization;
    \item the objectives of the internship;
    \item the main tasks performed;
    \item the main results, knowledge, and skills obtained.
\end{itemize}

The Abstract should be concise and should not reproduce the full content
of the report chapter by chapter.

\vspace{0.3cm}

% ================================================================
% XVII. REFERENCES AND CITATIONS
% ================================================================

\noindent
\textbf{XVII.\quad REFERENCES AND CITATIONS}

\vspace{0.15cm}

All books, journal papers, conference papers, standards, websites,
datasheets, manuals, technical documents, figures, tables, or other
external sources used in preparing the report shall be properly cited.

References shall be managed through the bibliography system provided with
the official template.

Bibliographic entries shall normally be stored in:

\begin{verbatim}
References.bib
\end{verbatim}

References shall be cited in the report using:

\begin{verbatim}
\cite{reference-key}
\end{verbatim}

Students shall not manually type reference numbers such as ``[1]'',
``[2]'', etc.

Titles of foreign-language publications shall remain in their original
language in the bibliography.

All borrowed ideas, data, figures, tables, equations, methods, or results
shall be properly acknowledged.

\vspace{0.3cm}

% ================================================================
% XVIII. CROSS-REFERENCING
% ================================================================

\noindent
\textbf{XVIII.\quad CROSS-REFERENCING}

\vspace{0.15cm}

Numbered elements shall be cross-referenced using the standard \LaTeX{}
mechanism.

Use the following commands where appropriate:

\begin{verbatim}
\label{...}
\ref{...}
\eqref{...}
\end{verbatim}

For example:

\begin{verbatim}
As shown in Figure~\ref{fig:system}, ...
\end{verbatim}

Students shall not manually type numbers such as ``Figure 2.3'' when
those numbers can be generated automatically.

\vspace{0.3cm}

% ================================================================
% XIX. PAGE NUMBERING
% ================================================================

\noindent
\textbf{XIX.\quad PAGE NUMBERING}

\vspace{0.15cm}

Cover pages and official evaluation forms at the beginning of the report
shall not display page numbers.

The main body of the report shall use Arabic page numbers beginning from
Chapter 1.

Page numbering shall be generated automatically by the official template.
Students shall not manually enter page numbers.

\vspace{0.3cm}

% ================================================================
% XX. PROJECT FILE ORGANIZATION
% ================================================================

\noindent
\textbf{XX.\quad PROJECT FILE ORGANIZATION}

\vspace{0.15cm}

Students should maintain a clear and organized \LaTeX{} project
structure. A recommended structure is:

\begin{verbatim}
InternshipReport/
|
|-- main.tex
|-- References.bib
|
|-- figures/
|   |-- company.pdf
|   |-- system.png
|   `-- ...
|
`-- chapters/
    |-- chapter1.tex
    |-- chapter2.tex
    `-- chapter3.tex
\end{verbatim}

Figures should normally be stored in the \verb|figures| directory
specified by the official template.

File and directory names should preferably contain Latin letters,
numbers, hyphens, or underscores.

Recommended examples include:

\begin{verbatim}
company_structure.pdf
motor_control.png
chapter_01.tex
\end{verbatim}

Spaces and unnecessary special characters in file names should be
avoided where possible.

\vspace{0.3cm}

% ================================================================
% XXI. DOCUMENT COMPILATION AND FINAL CHECK
% ================================================================

\noindent
\textbf{XXI.\quad DOCUMENT COMPILATION AND FINAL CHECK}

\vspace{0.15cm}

Before submission, the complete \LaTeX{} project shall compile
successfully and produce the final PDF report.

The final document shall not contain unresolved references such as:

\begin{verbatim}
??
[?]
\end{verbatim}

Before generating the final PDF, students shall check:

\begin{itemize}[leftmargin=0.8cm, itemsep=2pt]
    \item the Table of Contents;
    \item page numbering;
    \item chapter, section, and subsection numbering;
    \item figure numbering;
    \item table numbering;
    \item equation numbering;
    \item cross-references;
    \item citations;
    \item bibliography;
    \item font display;
    \item mathematical symbols;
    \item image quality;
    \item completeness of all required forms.
\end{itemize}

The document may be compiled using TeX Live, MiKTeX, or another
compatible \LaTeX{} environment. Regardless of the compilation
environment, the resulting PDF shall conform to the official Faculty
template.

\vspace{0.3cm}

% ================================================================
% XXII. CONTENT THAT STUDENTS MAY EDIT
% ================================================================

\noindent
\textbf{XXII.\quad CONTENT THAT STUDENTS MAY EDIT}

\vspace{0.15cm}

Students may edit or replace:

\begin{itemize}[leftmargin=0.8cm, itemsep=2pt]
    \item student name;
    \item class;
    \item internship organization;
    \item supervisor name;
    \item Abstract;
    \item List of Symbols and Abbreviations;
    \item chapter titles and chapter content;
    \item figures and photographs;
    \item tables;
    \item equations;
    \item citations and references.
\end{itemize}

Students shall not modify the main document formatting settings in the
preamble, including the document class, page geometry, line spacing,
heading formatting, table-of-contents formatting, font configuration,
page layout, or cover-page macros unless specifically instructed.

\vspace{0.3cm}

% ================================================================
% XXIII. FINAL SUBMISSION REQUIREMENTS
% ================================================================

\noindent
\textbf{XXIII.\quad FINAL SUBMISSION REQUIREMENTS}

\vspace{0.15cm}

Before submission, the final report shall:

\begin{itemize}[leftmargin=0.8cm, itemsep=2pt]
    \item compile without errors that affect the report content or presentation;
    \item contain all required forms and report sections;
    \item contain correct student, class, internship organization, and supervisor information;
    \item contain an updated Table of Contents, List of Figures, and List of Tables;
    \item contain no unused example content from the template;
    \item contain no placeholder text such as \textit{Title of Chapter 1},
    \textit{Section 1.1}, or other sample entries;
    \item contain all required citations and references;
    \item contain no unresolved cross-references.
\end{itemize}

If the Faculty requires submission of the \LaTeX{} source, the complete
project shall include all \verb|.tex| files, the bibliography database,
figures, and any other files required to reproduce the submitted PDF.

\vspace{0.4cm}

\noindent
\textbf{Important:}
Before submitting the final internship report, students shall remove or
comment out the following command in the main document:

\begin{verbatim}
\input{Regulations.tex}
\end{verbatim}

For example:

\begin{verbatim}
% \input{Regulations.tex}
\end{verbatim}

This removes the guidance pages from the final submitted report while
preserving the regulations file in the \LaTeX{} project for reference.

} % End 13pt content

\clearpage
|.
Students shall remove or comment out this command before submitting the
final internship report.

\vspace{0.4cm}

% ================================================================
% I. GENERAL REQUIREMENTS
% ================================================================

\noindent
\textbf{I.\quad GENERAL REQUIREMENTS}

\vspace{0.15cm}

The Graduation Internship Report shall be prepared using the official
\LaTeX{} template provided by the Faculty. Students shall preserve the
predefined page layout, typography, line spacing, heading styles,
numbering system, figure and table formatting, and bibliography settings
unless otherwise instructed by the Faculty or the supervisor.

The report shall be written in English.

Students are responsible for the academic and technical content of the
report. The official \LaTeX{} template is responsible for standardized
document formatting.

\vspace{0.3cm}

% ================================================================
% II. REPORT STRUCTURE
% ================================================================

\noindent
\textbf{II.\quad REPORT STRUCTURE}

\vspace{0.15cm}

The report shall normally be arranged in the following order:

\begin{enumerate}[label=\arabic*., leftmargin=1.2cm, itemsep=2pt]
    \item Outer Cover;
    \item Comments of the Internship Host Organization;
    \item Supervisor's Evaluation and Graduation Internship Defense Results;
    \item Table of Contents;
    \item List of Symbols and Abbreviations, if applicable;
    \item List of Figures;
    \item List of Tables;
    \item Abstract;
    \item Main Chapters of the Internship Report;
    \item References;
    \item Appendices, if applicable.
\end{enumerate}

The preliminary pages shall not be numbered as chapters. The main report
content shall be organized into consecutively numbered chapters beginning
with Chapter 1.

\vspace{0.3cm}

% ================================================================
% III. PAGE SIZE AND MARGINS
% ================================================================

\noindent
\textbf{III.\quad PAGE SIZE AND MARGINS}

\vspace{0.15cm}

The report shall be formatted on A4 paper with the following margins:

\begin{center}
\begin{tabular}{ll}
    Left margin:   & 3.5 cm \\
    Right margin:  & 2.0 cm \\
    Top margin:    & 2.5 cm \\
    Bottom margin: & 2.5 cm
\end{tabular}
\end{center}

These settings are predefined in the official template and shall not be
modified by students.

\vspace{0.3cm}

% ================================================================
% IV. FONT, FONT SIZE, AND LINE SPACING
% ================================================================

\noindent
\textbf{IV.\quad FONT, FONT SIZE, AND LINE SPACING}

\vspace{0.15cm}

The main text shall use a 13 pt font size with 1.3-line spacing.
A Times-family typeface shall be used in accordance with the official
template.

The text shall be fully justified.

Students shall not use repeated spaces, repeated manual line breaks, or
multiple blank lines to position text. Paragraph and heading spacing shall
be controlled by the template.

\vspace{0.3cm}

% ================================================================
% V. OUTER COVER
% ================================================================

\noindent
\textbf{V.\quad OUTER COVER}

\vspace{0.15cm}

The outer cover shall follow the official format provided in the template
and shall contain the following information:

\begin{itemize}[leftmargin=0.8cm, itemsep=2pt]
    \item Thai Nguyen University of Technology;
    \item Faculty of International Training;
    \item Engineering and Technology Department;
    \item University logo and Faculty logo;
    \item the title \textbf{INTERNSHIP REPORT};
    \item Internship Organization;
    \item Author;
    \item Class;
    \item Supervisor;
    \item place and year of completion.
\end{itemize}

Students shall only replace the required information fields and shall not
modify the cover-page layout unless instructed otherwise.

\vspace{0.3cm}

% ================================================================
% VI. COMMENTS OF THE INTERNSHIP HOST ORGANIZATION
% ================================================================

\noindent
\textbf{VI.\quad COMMENTS OF THE INTERNSHIP HOST ORGANIZATION}

\vspace{0.15cm}

The page entitled \textbf{COMMENTS OF THE INTERNSHIP HOST ORGANIZATION}
is reserved for the internship host organization.

It includes:

\begin{itemize}[leftmargin=0.8cm, itemsep=2pt]
    \item the student's full name;
    \item class;
    \item comments from the host organization;
    \item score;
    \item date;
    \item signature and official seal.
\end{itemize}

Students shall not alter or complete official evaluation fields unless
specifically authorized.

\vspace{0.3cm}

% ================================================================
% VII. SUPERVISOR'S EVALUATION AND DEFENSE RESULTS
% ================================================================

\noindent
\textbf{VII.\quad SUPERVISOR'S EVALUATION AND DEFENSE RESULTS}

\vspace{0.15cm}

The section entitled
\textbf{SUPERVISOR'S EVALUATION AND GRADUATION INTERNSHIP DEFENSE RESULTS}
is reserved for the supervisor and defense committee.

The evaluation may include:

\begin{itemize}[leftmargin=0.8cm, itemsep=2pt]
    \item Internship Progress and Student's Attitude;
    \item Report Content;
    \item Report Presentation;
    \item Other Comments;
    \item Supervisor's Opinion;
    \item Defense Result;
    \item Score.
\end{itemize}

Students shall not enter evaluation scores or defense results before they
are officially determined.

\vspace{0.3cm}

% ================================================================
% VIII. CHAPTERS, SECTIONS, AND SUBSECTIONS
% ================================================================

\noindent
\textbf{VIII.\quad CHAPTERS, SECTIONS, AND SUBSECTIONS}

\vspace{0.15cm}

The report structure shall be generated using the standard \LaTeX{}
commands:

\begin{verbatim}
\chapter{...}
\section{...}
\subsection{...}
\end{verbatim}

Chapter, section, and subsection numbers shall be generated automatically
by \LaTeX{}.

Students shall not manually type chapter or section numbers and shall not
reproduce heading styles using commands such as \verb|\textbf|,
\verb|\fontsize|, or \verb|\centering| instead of the appropriate
structural command.

Excessive subdivision should be avoided. In most cases, the hierarchy
Chapter--Section--Subsection is sufficient.

\vspace{0.3cm}

% ================================================================
% IX. ORGANIZATION OF THE MAIN CONTENT
% ================================================================

\noindent
\textbf{IX.\quad ORGANIZATION OF THE MAIN CONTENT}

\vspace{0.15cm}

The specific content of the report shall depend on the internship
organization and the assigned tasks. A typical structure may include:

\begin{itemize}[leftmargin=0.8cm, itemsep=2pt]
    \item \textbf{Chapter 1. Overview of the Internship Organization};
    \item \textbf{Chapter 2. Internship Activities and Assigned Tasks};
    \item \textbf{Chapter 3. Results, Experience, and Evaluation}.
\end{itemize}

These chapter titles are recommendations only and may be adapted to the
actual internship content.

The report shall reflect work actually performed, observed, analyzed, or
learned during the internship. Unrelated theoretical material should not
be included merely to increase report length.

\vspace{0.3cm}

% ================================================================
% X. TABLE OF CONTENTS
% ================================================================

\noindent
\textbf{X.\quad TABLE OF CONTENTS}

\vspace{0.15cm}

The Table of Contents shall be generated automatically using:

\begin{verbatim}
\tableofcontents
\end{verbatim}

Students shall not create or manually edit page numbers in the Table of
Contents.

The document shall be compiled sufficiently to ensure that all page
numbers and entries are correctly updated.

\vspace{0.3cm}

% ================================================================
% XI. LIST OF SYMBOLS AND ABBREVIATIONS
% ================================================================

\noindent
\textbf{XI.\quad LIST OF SYMBOLS AND ABBREVIATIONS}

\vspace{0.15cm}

Symbols and abbreviations that are used repeatedly in the report shall be
listed in the \textbf{List of Symbols and Abbreviations}.

The list may contain the following columns:

\begin{center}
\begin{tabular}{llll}
    No. & Symbol / Abbreviation & Unit & Definition
\end{tabular}
\end{center}

Only symbols and abbreviations actually used in the final report shall be
retained. Example entries contained in the template shall be removed if
they are not relevant to the student's report.

Each abbreviation should normally be defined when it first appears in
the main text.

\vspace{0.3cm}

% ================================================================
% XII. FIGURES, DIAGRAMS, AND PHOTOGRAPHS
% ================================================================

\noindent
\textbf{XII.\quad FIGURES, DIAGRAMS, AND PHOTOGRAPHS}

\vspace{0.15cm}

Figures shall be clear, legible, and directly relevant to the report
content.

Figures shall be inserted using the standard \verb|figure| environment.

Example:

\begin{verbatim}
\begin{figure}[H]
    \centering
    \includegraphics[width=0.8\textwidth]
                    {figures/company_structure.pdf}
    \caption{Organizational structure of the company}
    \label{fig:company-structure}
\end{figure}
\end{verbatim}

Figure captions shall be placed below the figure.

Each referenced figure shall have a unique \verb|\label| and shall be
referenced in the text using, for example:

\begin{verbatim}
Figure~\ref{fig:company-structure}
\end{verbatim}

Students shall not manually type figure numbers.

Figures reproduced or adapted from external sources shall include an
appropriate citation.

Photographs taken during the internship shall be relevant and sufficiently
clear. Students shall not include confidential, restricted, or proprietary
information without permission from the internship organization.

Vector graphics or high-resolution images are recommended whenever
possible.

\vspace{0.3cm}

% ================================================================
% XIII. TABLES
% ================================================================

\noindent
\textbf{XIII.\quad TABLES}

\vspace{0.15cm}

Tables shall be created using suitable \LaTeX{} environments such as
\verb|table|, \verb|tabular|, \verb|tabularx|, or \verb|longtable|.

Table captions shall be placed above the table.

Example:

\begin{verbatim}
\begin{table}[H]
    \centering
    \caption{Internship schedule}
    \label{tab:internship-schedule}

    \begin{tabular}{lll}
        \toprule
        Week & Task & Result \\
        \midrule
        1 & Orientation & Completed \\
        2 & System study & Completed \\
        \bottomrule
    \end{tabular}
\end{table}
\end{verbatim}

Tables shall be referenced using:

\begin{verbatim}
Table~\ref{tab:internship-schedule}
\end{verbatim}

Students shall not manually type table numbers.

Tables should be created directly in \LaTeX{} rather than inserted as
screenshots whenever practical.

\vspace{0.3cm}

% ================================================================
% XIV. EQUATIONS AND TECHNICAL QUANTITIES
% ================================================================

\noindent
\textbf{XIV.\quad EQUATIONS AND TECHNICAL QUANTITIES}

\vspace{0.15cm}

Mathematical expressions shall be typeset directly using \LaTeX{}.
Images or screenshots of equations shall not be used.

Short expressions may be written in inline mathematical mode, for example:

\begin{verbatim}
$R = 0.7\,\Omega$
\end{verbatim}

Important equations shall use a numbered equation environment.

Example:

\begin{verbatim}
\begin{equation}
    P = UI
    \label{eq:power}
\end{equation}
\end{verbatim}

Equations shall be referenced using:

\begin{verbatim}
Equation~\eqref{eq:power}
\end{verbatim}

Equation numbers shall not be manually entered.

SI units shall be used wherever applicable. The \verb|siunitx| package
should be used where appropriate.

Examples:

\begin{verbatim}
\SI{220}{\volt}
\SI{50}{\hertz}
\SI{2.5}{\ampere}
\end{verbatim}

Notation and units shall be consistent throughout the report.

\vspace{0.3cm}

% ================================================================
% XV. SOURCE CODE
% ================================================================

\noindent
\textbf{XV.\quad SOURCE CODE}

\vspace{0.15cm}

If the internship involves programming, PLCs, microcontrollers,
simulation, or software development, relevant source code may be included
in the report.

Source code shall be presented using the \verb|listings| environment or
another environment defined by the official template.

Example:

\begin{verbatim}
\begin{lstlisting}[language=st]
IF Start = TRUE THEN
    Motor := TRUE;
ELSE
    Motor := FALSE;
END_IF;
\end{lstlisting}
\end{verbatim}

Only code necessary for explaining the work should appear in the main
report. Long source-code listings should be moved to an appendix or
submitted as supplementary files if required.

Source code should not normally be inserted as screenshots.

\vspace{0.3cm}

% ================================================================
% XVI. ABSTRACT
% ================================================================

\noindent
\textbf{XVI.\quad ABSTRACT}

\vspace{0.15cm}

The Abstract shall briefly summarize:

\begin{itemize}[leftmargin=0.8cm, itemsep=2pt]
    \item the internship organization;
    \item the objectives of the internship;
    \item the main tasks performed;
    \item the main results, knowledge, and skills obtained.
\end{itemize}

The Abstract should be concise and should not reproduce the full content
of the report chapter by chapter.

\vspace{0.3cm}

% ================================================================
% XVII. REFERENCES AND CITATIONS
% ================================================================

\noindent
\textbf{XVII.\quad REFERENCES AND CITATIONS}

\vspace{0.15cm}

All books, journal papers, conference papers, standards, websites,
datasheets, manuals, technical documents, figures, tables, or other
external sources used in preparing the report shall be properly cited.

References shall be managed through the bibliography system provided with
the official template.

Bibliographic entries shall normally be stored in:

\begin{verbatim}
References.bib
\end{verbatim}

References shall be cited in the report using:

\begin{verbatim}
\cite{reference-key}
\end{verbatim}

Students shall not manually type reference numbers such as ``[1]'',
``[2]'', etc.

Titles of foreign-language publications shall remain in their original
language in the bibliography.

All borrowed ideas, data, figures, tables, equations, methods, or results
shall be properly acknowledged.

\vspace{0.3cm}

% ================================================================
% XVIII. CROSS-REFERENCING
% ================================================================

\noindent
\textbf{XVIII.\quad CROSS-REFERENCING}

\vspace{0.15cm}

Numbered elements shall be cross-referenced using the standard \LaTeX{}
mechanism.

Use the following commands where appropriate:

\begin{verbatim}
\label{...}
\ref{...}
\eqref{...}
\end{verbatim}

For example:

\begin{verbatim}
As shown in Figure~\ref{fig:system}, ...
\end{verbatim}

Students shall not manually type numbers such as ``Figure 2.3'' when
those numbers can be generated automatically.

\vspace{0.3cm}

% ================================================================
% XIX. PAGE NUMBERING
% ================================================================

\noindent
\textbf{XIX.\quad PAGE NUMBERING}

\vspace{0.15cm}

Cover pages and official evaluation forms at the beginning of the report
shall not display page numbers.

The main body of the report shall use Arabic page numbers beginning from
Chapter 1.

Page numbering shall be generated automatically by the official template.
Students shall not manually enter page numbers.

\vspace{0.3cm}

% ================================================================
% XX. PROJECT FILE ORGANIZATION
% ================================================================

\noindent
\textbf{XX.\quad PROJECT FILE ORGANIZATION}

\vspace{0.15cm}

Students should maintain a clear and organized \LaTeX{} project
structure. A recommended structure is:

\begin{verbatim}
InternshipReport/
|
|-- main.tex
|-- References.bib
|
|-- figures/
|   |-- company.pdf
|   |-- system.png
|   `-- ...
|
`-- chapters/
    |-- chapter1.tex
    |-- chapter2.tex
    `-- chapter3.tex
\end{verbatim}

Figures should normally be stored in the \verb|figures| directory
specified by the official template.

File and directory names should preferably contain Latin letters,
numbers, hyphens, or underscores.

Recommended examples include:

\begin{verbatim}
company_structure.pdf
motor_control.png
chapter_01.tex
\end{verbatim}

Spaces and unnecessary special characters in file names should be
avoided where possible.

\vspace{0.3cm}

% ================================================================
% XXI. DOCUMENT COMPILATION AND FINAL CHECK
% ================================================================

\noindent
\textbf{XXI.\quad DOCUMENT COMPILATION AND FINAL CHECK}

\vspace{0.15cm}

Before submission, the complete \LaTeX{} project shall compile
successfully and produce the final PDF report.

The final document shall not contain unresolved references such as:

\begin{verbatim}
??
[?]
\end{verbatim}

Before generating the final PDF, students shall check:

\begin{itemize}[leftmargin=0.8cm, itemsep=2pt]
    \item the Table of Contents;
    \item page numbering;
    \item chapter, section, and subsection numbering;
    \item figure numbering;
    \item table numbering;
    \item equation numbering;
    \item cross-references;
    \item citations;
    \item bibliography;
    \item font display;
    \item mathematical symbols;
    \item image quality;
    \item completeness of all required forms.
\end{itemize}

The document may be compiled using TeX Live, MiKTeX, or another
compatible \LaTeX{} environment. Regardless of the compilation
environment, the resulting PDF shall conform to the official Faculty
template.

\vspace{0.3cm}

% ================================================================
% XXII. CONTENT THAT STUDENTS MAY EDIT
% ================================================================

\noindent
\textbf{XXII.\quad CONTENT THAT STUDENTS MAY EDIT}

\vspace{0.15cm}

Students may edit or replace:

\begin{itemize}[leftmargin=0.8cm, itemsep=2pt]
    \item student name;
    \item class;
    \item internship organization;
    \item supervisor name;
    \item Abstract;
    \item List of Symbols and Abbreviations;
    \item chapter titles and chapter content;
    \item figures and photographs;
    \item tables;
    \item equations;
    \item citations and references.
\end{itemize}

Students shall not modify the main document formatting settings in the
preamble, including the document class, page geometry, line spacing,
heading formatting, table-of-contents formatting, font configuration,
page layout, or cover-page macros unless specifically instructed.

\vspace{0.3cm}

% ================================================================
% XXIII. FINAL SUBMISSION REQUIREMENTS
% ================================================================

\noindent
\textbf{XXIII.\quad FINAL SUBMISSION REQUIREMENTS}

\vspace{0.15cm}

Before submission, the final report shall:

\begin{itemize}[leftmargin=0.8cm, itemsep=2pt]
    \item compile without errors that affect the report content or presentation;
    \item contain all required forms and report sections;
    \item contain correct student, class, internship organization, and supervisor information;
    \item contain an updated Table of Contents, List of Figures, and List of Tables;
    \item contain no unused example content from the template;
    \item contain no placeholder text such as \textit{Title of Chapter 1},
    \textit{Section 1.1}, or other sample entries;
    \item contain all required citations and references;
    \item contain no unresolved cross-references.
\end{itemize}

If the Faculty requires submission of the \LaTeX{} source, the complete
project shall include all \verb|.tex| files, the bibliography database,
figures, and any other files required to reproduce the submitted PDF.

\vspace{0.4cm}

\noindent
\textbf{Important:}
Before submitting the final internship report, students shall remove or
comment out the following command in the main document:

\begin{verbatim}

% =================================================================
% REGULATIONS FOR PREPARING THE GRADUATION INTERNSHIP REPORT
% USING LATEX
% =================================================================
%
% This file is intended to be included from the main document with:
%
%     \input{Regulations.tex}
%
% IMPORTANT:
% This section is provided for guidance only. Students shall remove
% or comment out the \input{Regulations.tex} command before submitting
% the final internship report.
%
% =================================================================

\clearpage
\thispagestyle{empty}

\begin{center}
    {\fontsize{16pt}{20pt}\selectfont\bfseries
    SPECIFIC REGULATIONS FOR PREPARING\\
    GRADUATION INTERNSHIP REPORTS USING \LaTeX}
\end{center}

\vspace{0.5cm}

{\fontsize{13pt}{17pt}\selectfont
\setstretch{1.3}
\justifying

\noindent
\textbf{Note:}
This section is provided for guidance only and is included in the main
document using the command \verb|\input{Regulations.tex}|.
Students shall remove or comment out this command before submitting the
final internship report.

\vspace{0.4cm}

% ================================================================
% I. GENERAL REQUIREMENTS
% ================================================================

\noindent
\textbf{I.\quad GENERAL REQUIREMENTS}

\vspace{0.15cm}

The Graduation Internship Report shall be prepared using the official
\LaTeX{} template provided by the Faculty. Students shall preserve the
predefined page layout, typography, line spacing, heading styles,
numbering system, figure and table formatting, and bibliography settings
unless otherwise instructed by the Faculty or the supervisor.

The report shall be written in English.

Students are responsible for the academic and technical content of the
report. The official \LaTeX{} template is responsible for standardized
document formatting.

\vspace{0.3cm}

% ================================================================
% II. REPORT STRUCTURE
% ================================================================

\noindent
\textbf{II.\quad REPORT STRUCTURE}

\vspace{0.15cm}

The report shall normally be arranged in the following order:

\begin{enumerate}[label=\arabic*., leftmargin=1.2cm, itemsep=2pt]
    \item Outer Cover;
    \item Comments of the Internship Host Organization;
    \item Supervisor's Evaluation and Graduation Internship Defense Results;
    \item Table of Contents;
    \item List of Symbols and Abbreviations, if applicable;
    \item List of Figures;
    \item List of Tables;
    \item Abstract;
    \item Main Chapters of the Internship Report;
    \item References;
    \item Appendices, if applicable.
\end{enumerate}

The preliminary pages shall not be numbered as chapters. The main report
content shall be organized into consecutively numbered chapters beginning
with Chapter 1.

\vspace{0.3cm}

% ================================================================
% III. PAGE SIZE AND MARGINS
% ================================================================

\noindent
\textbf{III.\quad PAGE SIZE AND MARGINS}

\vspace{0.15cm}

The report shall be formatted on A4 paper with the following margins:

\begin{center}
\begin{tabular}{ll}
    Left margin:   & 3.5 cm \\
    Right margin:  & 2.0 cm \\
    Top margin:    & 2.5 cm \\
    Bottom margin: & 2.5 cm
\end{tabular}
\end{center}

These settings are predefined in the official template and shall not be
modified by students.

\vspace{0.3cm}

% ================================================================
% IV. FONT, FONT SIZE, AND LINE SPACING
% ================================================================

\noindent
\textbf{IV.\quad FONT, FONT SIZE, AND LINE SPACING}

\vspace{0.15cm}

The main text shall use a 13 pt font size with 1.3-line spacing.
A Times-family typeface shall be used in accordance with the official
template.

The text shall be fully justified.

Students shall not use repeated spaces, repeated manual line breaks, or
multiple blank lines to position text. Paragraph and heading spacing shall
be controlled by the template.

\vspace{0.3cm}

% ================================================================
% V. OUTER COVER
% ================================================================

\noindent
\textbf{V.\quad OUTER COVER}

\vspace{0.15cm}

The outer cover shall follow the official format provided in the template
and shall contain the following information:

\begin{itemize}[leftmargin=0.8cm, itemsep=2pt]
    \item Thai Nguyen University of Technology;
    \item Faculty of International Training;
    \item Engineering and Technology Department;
    \item University logo and Faculty logo;
    \item the title \textbf{INTERNSHIP REPORT};
    \item Internship Organization;
    \item Author;
    \item Class;
    \item Supervisor;
    \item place and year of completion.
\end{itemize}

Students shall only replace the required information fields and shall not
modify the cover-page layout unless instructed otherwise.

\vspace{0.3cm}

% ================================================================
% VI. COMMENTS OF THE INTERNSHIP HOST ORGANIZATION
% ================================================================

\noindent
\textbf{VI.\quad COMMENTS OF THE INTERNSHIP HOST ORGANIZATION}

\vspace{0.15cm}

The page entitled \textbf{COMMENTS OF THE INTERNSHIP HOST ORGANIZATION}
is reserved for the internship host organization.

It includes:

\begin{itemize}[leftmargin=0.8cm, itemsep=2pt]
    \item the student's full name;
    \item class;
    \item comments from the host organization;
    \item score;
    \item date;
    \item signature and official seal.
\end{itemize}

Students shall not alter or complete official evaluation fields unless
specifically authorized.

\vspace{0.3cm}

% ================================================================
% VII. SUPERVISOR'S EVALUATION AND DEFENSE RESULTS
% ================================================================

\noindent
\textbf{VII.\quad SUPERVISOR'S EVALUATION AND DEFENSE RESULTS}

\vspace{0.15cm}

The section entitled
\textbf{SUPERVISOR'S EVALUATION AND GRADUATION INTERNSHIP DEFENSE RESULTS}
is reserved for the supervisor and defense committee.

The evaluation may include:

\begin{itemize}[leftmargin=0.8cm, itemsep=2pt]
    \item Internship Progress and Student's Attitude;
    \item Report Content;
    \item Report Presentation;
    \item Other Comments;
    \item Supervisor's Opinion;
    \item Defense Result;
    \item Score.
\end{itemize}

Students shall not enter evaluation scores or defense results before they
are officially determined.

\vspace{0.3cm}

% ================================================================
% VIII. CHAPTERS, SECTIONS, AND SUBSECTIONS
% ================================================================

\noindent
\textbf{VIII.\quad CHAPTERS, SECTIONS, AND SUBSECTIONS}

\vspace{0.15cm}

The report structure shall be generated using the standard \LaTeX{}
commands:

\begin{verbatim}
\chapter{...}
\section{...}
\subsection{...}
\end{verbatim}

Chapter, section, and subsection numbers shall be generated automatically
by \LaTeX{}.

Students shall not manually type chapter or section numbers and shall not
reproduce heading styles using commands such as \verb|\textbf|,
\verb|\fontsize|, or \verb|\centering| instead of the appropriate
structural command.

Excessive subdivision should be avoided. In most cases, the hierarchy
Chapter--Section--Subsection is sufficient.

\vspace{0.3cm}

% ================================================================
% IX. ORGANIZATION OF THE MAIN CONTENT
% ================================================================

\noindent
\textbf{IX.\quad ORGANIZATION OF THE MAIN CONTENT}

\vspace{0.15cm}

The specific content of the report shall depend on the internship
organization and the assigned tasks. A typical structure may include:

\begin{itemize}[leftmargin=0.8cm, itemsep=2pt]
    \item \textbf{Chapter 1. Overview of the Internship Organization};
    \item \textbf{Chapter 2. Internship Activities and Assigned Tasks};
    \item \textbf{Chapter 3. Results, Experience, and Evaluation}.
\end{itemize}

These chapter titles are recommendations only and may be adapted to the
actual internship content.

The report shall reflect work actually performed, observed, analyzed, or
learned during the internship. Unrelated theoretical material should not
be included merely to increase report length.

\vspace{0.3cm}

% ================================================================
% X. TABLE OF CONTENTS
% ================================================================

\noindent
\textbf{X.\quad TABLE OF CONTENTS}

\vspace{0.15cm}

The Table of Contents shall be generated automatically using:

\begin{verbatim}
\tableofcontents
\end{verbatim}

Students shall not create or manually edit page numbers in the Table of
Contents.

The document shall be compiled sufficiently to ensure that all page
numbers and entries are correctly updated.

\vspace{0.3cm}

% ================================================================
% XI. LIST OF SYMBOLS AND ABBREVIATIONS
% ================================================================

\noindent
\textbf{XI.\quad LIST OF SYMBOLS AND ABBREVIATIONS}

\vspace{0.15cm}

Symbols and abbreviations that are used repeatedly in the report shall be
listed in the \textbf{List of Symbols and Abbreviations}.

The list may contain the following columns:

\begin{center}
\begin{tabular}{llll}
    No. & Symbol / Abbreviation & Unit & Definition
\end{tabular}
\end{center}

Only symbols and abbreviations actually used in the final report shall be
retained. Example entries contained in the template shall be removed if
they are not relevant to the student's report.

Each abbreviation should normally be defined when it first appears in
the main text.

\vspace{0.3cm}

% ================================================================
% XII. FIGURES, DIAGRAMS, AND PHOTOGRAPHS
% ================================================================

\noindent
\textbf{XII.\quad FIGURES, DIAGRAMS, AND PHOTOGRAPHS}

\vspace{0.15cm}

Figures shall be clear, legible, and directly relevant to the report
content.

Figures shall be inserted using the standard \verb|figure| environment.

Example:

\begin{verbatim}
\begin{figure}[H]
    \centering
    \includegraphics[width=0.8\textwidth]
                    {figures/company_structure.pdf}
    \caption{Organizational structure of the company}
    \label{fig:company-structure}
\end{figure}
\end{verbatim}

Figure captions shall be placed below the figure.

Each referenced figure shall have a unique \verb|\label| and shall be
referenced in the text using, for example:

\begin{verbatim}
Figure~\ref{fig:company-structure}
\end{verbatim}

Students shall not manually type figure numbers.

Figures reproduced or adapted from external sources shall include an
appropriate citation.

Photographs taken during the internship shall be relevant and sufficiently
clear. Students shall not include confidential, restricted, or proprietary
information without permission from the internship organization.

Vector graphics or high-resolution images are recommended whenever
possible.

\vspace{0.3cm}

% ================================================================
% XIII. TABLES
% ================================================================

\noindent
\textbf{XIII.\quad TABLES}

\vspace{0.15cm}

Tables shall be created using suitable \LaTeX{} environments such as
\verb|table|, \verb|tabular|, \verb|tabularx|, or \verb|longtable|.

Table captions shall be placed above the table.

Example:

\begin{verbatim}
\begin{table}[H]
    \centering
    \caption{Internship schedule}
    \label{tab:internship-schedule}

    \begin{tabular}{lll}
        \toprule
        Week & Task & Result \\
        \midrule
        1 & Orientation & Completed \\
        2 & System study & Completed \\
        \bottomrule
    \end{tabular}
\end{table}
\end{verbatim}

Tables shall be referenced using:

\begin{verbatim}
Table~\ref{tab:internship-schedule}
\end{verbatim}

Students shall not manually type table numbers.

Tables should be created directly in \LaTeX{} rather than inserted as
screenshots whenever practical.

\vspace{0.3cm}

% ================================================================
% XIV. EQUATIONS AND TECHNICAL QUANTITIES
% ================================================================

\noindent
\textbf{XIV.\quad EQUATIONS AND TECHNICAL QUANTITIES}

\vspace{0.15cm}

Mathematical expressions shall be typeset directly using \LaTeX{}.
Images or screenshots of equations shall not be used.

Short expressions may be written in inline mathematical mode, for example:

\begin{verbatim}
$R = 0.7\,\Omega$
\end{verbatim}

Important equations shall use a numbered equation environment.

Example:

\begin{verbatim}
\begin{equation}
    P = UI
    \label{eq:power}
\end{equation}
\end{verbatim}

Equations shall be referenced using:

\begin{verbatim}
Equation~\eqref{eq:power}
\end{verbatim}

Equation numbers shall not be manually entered.

SI units shall be used wherever applicable. The \verb|siunitx| package
should be used where appropriate.

Examples:

\begin{verbatim}
\SI{220}{\volt}
\SI{50}{\hertz}
\SI{2.5}{\ampere}
\end{verbatim}

Notation and units shall be consistent throughout the report.

\vspace{0.3cm}

% ================================================================
% XV. SOURCE CODE
% ================================================================

\noindent
\textbf{XV.\quad SOURCE CODE}

\vspace{0.15cm}

If the internship involves programming, PLCs, microcontrollers,
simulation, or software development, relevant source code may be included
in the report.

Source code shall be presented using the \verb|listings| environment or
another environment defined by the official template.

Example:

\begin{verbatim}
\begin{lstlisting}[language=st]
IF Start = TRUE THEN
    Motor := TRUE;
ELSE
    Motor := FALSE;
END_IF;
\end{lstlisting}
\end{verbatim}

Only code necessary for explaining the work should appear in the main
report. Long source-code listings should be moved to an appendix or
submitted as supplementary files if required.

Source code should not normally be inserted as screenshots.

\vspace{0.3cm}

% ================================================================
% XVI. ABSTRACT
% ================================================================

\noindent
\textbf{XVI.\quad ABSTRACT}

\vspace{0.15cm}

The Abstract shall briefly summarize:

\begin{itemize}[leftmargin=0.8cm, itemsep=2pt]
    \item the internship organization;
    \item the objectives of the internship;
    \item the main tasks performed;
    \item the main results, knowledge, and skills obtained.
\end{itemize}

The Abstract should be concise and should not reproduce the full content
of the report chapter by chapter.

\vspace{0.3cm}

% ================================================================
% XVII. REFERENCES AND CITATIONS
% ================================================================

\noindent
\textbf{XVII.\quad REFERENCES AND CITATIONS}

\vspace{0.15cm}

All books, journal papers, conference papers, standards, websites,
datasheets, manuals, technical documents, figures, tables, or other
external sources used in preparing the report shall be properly cited.

References shall be managed through the bibliography system provided with
the official template.

Bibliographic entries shall normally be stored in:

\begin{verbatim}
References.bib
\end{verbatim}

References shall be cited in the report using:

\begin{verbatim}
\cite{reference-key}
\end{verbatim}

Students shall not manually type reference numbers such as ``[1]'',
``[2]'', etc.

Titles of foreign-language publications shall remain in their original
language in the bibliography.

All borrowed ideas, data, figures, tables, equations, methods, or results
shall be properly acknowledged.

\vspace{0.3cm}

% ================================================================
% XVIII. CROSS-REFERENCING
% ================================================================

\noindent
\textbf{XVIII.\quad CROSS-REFERENCING}

\vspace{0.15cm}

Numbered elements shall be cross-referenced using the standard \LaTeX{}
mechanism.

Use the following commands where appropriate:

\begin{verbatim}
\label{...}
\ref{...}
\eqref{...}
\end{verbatim}

For example:

\begin{verbatim}
As shown in Figure~\ref{fig:system}, ...
\end{verbatim}

Students shall not manually type numbers such as ``Figure 2.3'' when
those numbers can be generated automatically.

\vspace{0.3cm}

% ================================================================
% XIX. PAGE NUMBERING
% ================================================================

\noindent
\textbf{XIX.\quad PAGE NUMBERING}

\vspace{0.15cm}

Cover pages and official evaluation forms at the beginning of the report
shall not display page numbers.

The main body of the report shall use Arabic page numbers beginning from
Chapter 1.

Page numbering shall be generated automatically by the official template.
Students shall not manually enter page numbers.

\vspace{0.3cm}

% ================================================================
% XX. PROJECT FILE ORGANIZATION
% ================================================================

\noindent
\textbf{XX.\quad PROJECT FILE ORGANIZATION}

\vspace{0.15cm}

Students should maintain a clear and organized \LaTeX{} project
structure. A recommended structure is:

\begin{verbatim}
InternshipReport/
|
|-- main.tex
|-- References.bib
|
|-- figures/
|   |-- company.pdf
|   |-- system.png
|   `-- ...
|
`-- chapters/
    |-- chapter1.tex
    |-- chapter2.tex
    `-- chapter3.tex
\end{verbatim}

Figures should normally be stored in the \verb|figures| directory
specified by the official template.

File and directory names should preferably contain Latin letters,
numbers, hyphens, or underscores.

Recommended examples include:

\begin{verbatim}
company_structure.pdf
motor_control.png
chapter_01.tex
\end{verbatim}

Spaces and unnecessary special characters in file names should be
avoided where possible.

\vspace{0.3cm}

% ================================================================
% XXI. DOCUMENT COMPILATION AND FINAL CHECK
% ================================================================

\noindent
\textbf{XXI.\quad DOCUMENT COMPILATION AND FINAL CHECK}

\vspace{0.15cm}

Before submission, the complete \LaTeX{} project shall compile
successfully and produce the final PDF report.

The final document shall not contain unresolved references such as:

\begin{verbatim}
??
[?]
\end{verbatim}

Before generating the final PDF, students shall check:

\begin{itemize}[leftmargin=0.8cm, itemsep=2pt]
    \item the Table of Contents;
    \item page numbering;
    \item chapter, section, and subsection numbering;
    \item figure numbering;
    \item table numbering;
    \item equation numbering;
    \item cross-references;
    \item citations;
    \item bibliography;
    \item font display;
    \item mathematical symbols;
    \item image quality;
    \item completeness of all required forms.
\end{itemize}

The document may be compiled using TeX Live, MiKTeX, or another
compatible \LaTeX{} environment. Regardless of the compilation
environment, the resulting PDF shall conform to the official Faculty
template.

\vspace{0.3cm}

% ================================================================
% XXII. CONTENT THAT STUDENTS MAY EDIT
% ================================================================

\noindent
\textbf{XXII.\quad CONTENT THAT STUDENTS MAY EDIT}

\vspace{0.15cm}

Students may edit or replace:

\begin{itemize}[leftmargin=0.8cm, itemsep=2pt]
    \item student name;
    \item class;
    \item internship organization;
    \item supervisor name;
    \item Abstract;
    \item List of Symbols and Abbreviations;
    \item chapter titles and chapter content;
    \item figures and photographs;
    \item tables;
    \item equations;
    \item citations and references.
\end{itemize}

Students shall not modify the main document formatting settings in the
preamble, including the document class, page geometry, line spacing,
heading formatting, table-of-contents formatting, font configuration,
page layout, or cover-page macros unless specifically instructed.

\vspace{0.3cm}

% ================================================================
% XXIII. FINAL SUBMISSION REQUIREMENTS
% ================================================================

\noindent
\textbf{XXIII.\quad FINAL SUBMISSION REQUIREMENTS}

\vspace{0.15cm}

Before submission, the final report shall:

\begin{itemize}[leftmargin=0.8cm, itemsep=2pt]
    \item compile without errors that affect the report content or presentation;
    \item contain all required forms and report sections;
    \item contain correct student, class, internship organization, and supervisor information;
    \item contain an updated Table of Contents, List of Figures, and List of Tables;
    \item contain no unused example content from the template;
    \item contain no placeholder text such as \textit{Title of Chapter 1},
    \textit{Section 1.1}, or other sample entries;
    \item contain all required citations and references;
    \item contain no unresolved cross-references.
\end{itemize}

If the Faculty requires submission of the \LaTeX{} source, the complete
project shall include all \verb|.tex| files, the bibliography database,
figures, and any other files required to reproduce the submitted PDF.

\vspace{0.4cm}

\noindent
\textbf{Important:}
Before submitting the final internship report, students shall remove or
comment out the following command in the main document:

\begin{verbatim}
\input{Regulations.tex}
\end{verbatim}

For example:

\begin{verbatim}
% \input{Regulations.tex}
\end{verbatim}

This removes the guidance pages from the final submitted report while
preserving the regulations file in the \LaTeX{} project for reference.

} % End 13pt content

\clearpage

\end{verbatim}

For example:

\begin{verbatim}
% 
% =================================================================
% REGULATIONS FOR PREPARING THE GRADUATION INTERNSHIP REPORT
% USING LATEX
% =================================================================
%
% This file is intended to be included from the main document with:
%
%     \input{Regulations.tex}
%
% IMPORTANT:
% This section is provided for guidance only. Students shall remove
% or comment out the \input{Regulations.tex} command before submitting
% the final internship report.
%
% =================================================================

\clearpage
\thispagestyle{empty}

\begin{center}
    {\fontsize{16pt}{20pt}\selectfont\bfseries
    SPECIFIC REGULATIONS FOR PREPARING\\
    GRADUATION INTERNSHIP REPORTS USING \LaTeX}
\end{center}

\vspace{0.5cm}

{\fontsize{13pt}{17pt}\selectfont
\setstretch{1.3}
\justifying

\noindent
\textbf{Note:}
This section is provided for guidance only and is included in the main
document using the command \verb|\input{Regulations.tex}|.
Students shall remove or comment out this command before submitting the
final internship report.

\vspace{0.4cm}

% ================================================================
% I. GENERAL REQUIREMENTS
% ================================================================

\noindent
\textbf{I.\quad GENERAL REQUIREMENTS}

\vspace{0.15cm}

The Graduation Internship Report shall be prepared using the official
\LaTeX{} template provided by the Faculty. Students shall preserve the
predefined page layout, typography, line spacing, heading styles,
numbering system, figure and table formatting, and bibliography settings
unless otherwise instructed by the Faculty or the supervisor.

The report shall be written in English.

Students are responsible for the academic and technical content of the
report. The official \LaTeX{} template is responsible for standardized
document formatting.

\vspace{0.3cm}

% ================================================================
% II. REPORT STRUCTURE
% ================================================================

\noindent
\textbf{II.\quad REPORT STRUCTURE}

\vspace{0.15cm}

The report shall normally be arranged in the following order:

\begin{enumerate}[label=\arabic*., leftmargin=1.2cm, itemsep=2pt]
    \item Outer Cover;
    \item Comments of the Internship Host Organization;
    \item Supervisor's Evaluation and Graduation Internship Defense Results;
    \item Table of Contents;
    \item List of Symbols and Abbreviations, if applicable;
    \item List of Figures;
    \item List of Tables;
    \item Abstract;
    \item Main Chapters of the Internship Report;
    \item References;
    \item Appendices, if applicable.
\end{enumerate}

The preliminary pages shall not be numbered as chapters. The main report
content shall be organized into consecutively numbered chapters beginning
with Chapter 1.

\vspace{0.3cm}

% ================================================================
% III. PAGE SIZE AND MARGINS
% ================================================================

\noindent
\textbf{III.\quad PAGE SIZE AND MARGINS}

\vspace{0.15cm}

The report shall be formatted on A4 paper with the following margins:

\begin{center}
\begin{tabular}{ll}
    Left margin:   & 3.5 cm \\
    Right margin:  & 2.0 cm \\
    Top margin:    & 2.5 cm \\
    Bottom margin: & 2.5 cm
\end{tabular}
\end{center}

These settings are predefined in the official template and shall not be
modified by students.

\vspace{0.3cm}

% ================================================================
% IV. FONT, FONT SIZE, AND LINE SPACING
% ================================================================

\noindent
\textbf{IV.\quad FONT, FONT SIZE, AND LINE SPACING}

\vspace{0.15cm}

The main text shall use a 13 pt font size with 1.3-line spacing.
A Times-family typeface shall be used in accordance with the official
template.

The text shall be fully justified.

Students shall not use repeated spaces, repeated manual line breaks, or
multiple blank lines to position text. Paragraph and heading spacing shall
be controlled by the template.

\vspace{0.3cm}

% ================================================================
% V. OUTER COVER
% ================================================================

\noindent
\textbf{V.\quad OUTER COVER}

\vspace{0.15cm}

The outer cover shall follow the official format provided in the template
and shall contain the following information:

\begin{itemize}[leftmargin=0.8cm, itemsep=2pt]
    \item Thai Nguyen University of Technology;
    \item Faculty of International Training;
    \item Engineering and Technology Department;
    \item University logo and Faculty logo;
    \item the title \textbf{INTERNSHIP REPORT};
    \item Internship Organization;
    \item Author;
    \item Class;
    \item Supervisor;
    \item place and year of completion.
\end{itemize}

Students shall only replace the required information fields and shall not
modify the cover-page layout unless instructed otherwise.

\vspace{0.3cm}

% ================================================================
% VI. COMMENTS OF THE INTERNSHIP HOST ORGANIZATION
% ================================================================

\noindent
\textbf{VI.\quad COMMENTS OF THE INTERNSHIP HOST ORGANIZATION}

\vspace{0.15cm}

The page entitled \textbf{COMMENTS OF THE INTERNSHIP HOST ORGANIZATION}
is reserved for the internship host organization.

It includes:

\begin{itemize}[leftmargin=0.8cm, itemsep=2pt]
    \item the student's full name;
    \item class;
    \item comments from the host organization;
    \item score;
    \item date;
    \item signature and official seal.
\end{itemize}

Students shall not alter or complete official evaluation fields unless
specifically authorized.

\vspace{0.3cm}

% ================================================================
% VII. SUPERVISOR'S EVALUATION AND DEFENSE RESULTS
% ================================================================

\noindent
\textbf{VII.\quad SUPERVISOR'S EVALUATION AND DEFENSE RESULTS}

\vspace{0.15cm}

The section entitled
\textbf{SUPERVISOR'S EVALUATION AND GRADUATION INTERNSHIP DEFENSE RESULTS}
is reserved for the supervisor and defense committee.

The evaluation may include:

\begin{itemize}[leftmargin=0.8cm, itemsep=2pt]
    \item Internship Progress and Student's Attitude;
    \item Report Content;
    \item Report Presentation;
    \item Other Comments;
    \item Supervisor's Opinion;
    \item Defense Result;
    \item Score.
\end{itemize}

Students shall not enter evaluation scores or defense results before they
are officially determined.

\vspace{0.3cm}

% ================================================================
% VIII. CHAPTERS, SECTIONS, AND SUBSECTIONS
% ================================================================

\noindent
\textbf{VIII.\quad CHAPTERS, SECTIONS, AND SUBSECTIONS}

\vspace{0.15cm}

The report structure shall be generated using the standard \LaTeX{}
commands:

\begin{verbatim}
\chapter{...}
\section{...}
\subsection{...}
\end{verbatim}

Chapter, section, and subsection numbers shall be generated automatically
by \LaTeX{}.

Students shall not manually type chapter or section numbers and shall not
reproduce heading styles using commands such as \verb|\textbf|,
\verb|\fontsize|, or \verb|\centering| instead of the appropriate
structural command.

Excessive subdivision should be avoided. In most cases, the hierarchy
Chapter--Section--Subsection is sufficient.

\vspace{0.3cm}

% ================================================================
% IX. ORGANIZATION OF THE MAIN CONTENT
% ================================================================

\noindent
\textbf{IX.\quad ORGANIZATION OF THE MAIN CONTENT}

\vspace{0.15cm}

The specific content of the report shall depend on the internship
organization and the assigned tasks. A typical structure may include:

\begin{itemize}[leftmargin=0.8cm, itemsep=2pt]
    \item \textbf{Chapter 1. Overview of the Internship Organization};
    \item \textbf{Chapter 2. Internship Activities and Assigned Tasks};
    \item \textbf{Chapter 3. Results, Experience, and Evaluation}.
\end{itemize}

These chapter titles are recommendations only and may be adapted to the
actual internship content.

The report shall reflect work actually performed, observed, analyzed, or
learned during the internship. Unrelated theoretical material should not
be included merely to increase report length.

\vspace{0.3cm}

% ================================================================
% X. TABLE OF CONTENTS
% ================================================================

\noindent
\textbf{X.\quad TABLE OF CONTENTS}

\vspace{0.15cm}

The Table of Contents shall be generated automatically using:

\begin{verbatim}
\tableofcontents
\end{verbatim}

Students shall not create or manually edit page numbers in the Table of
Contents.

The document shall be compiled sufficiently to ensure that all page
numbers and entries are correctly updated.

\vspace{0.3cm}

% ================================================================
% XI. LIST OF SYMBOLS AND ABBREVIATIONS
% ================================================================

\noindent
\textbf{XI.\quad LIST OF SYMBOLS AND ABBREVIATIONS}

\vspace{0.15cm}

Symbols and abbreviations that are used repeatedly in the report shall be
listed in the \textbf{List of Symbols and Abbreviations}.

The list may contain the following columns:

\begin{center}
\begin{tabular}{llll}
    No. & Symbol / Abbreviation & Unit & Definition
\end{tabular}
\end{center}

Only symbols and abbreviations actually used in the final report shall be
retained. Example entries contained in the template shall be removed if
they are not relevant to the student's report.

Each abbreviation should normally be defined when it first appears in
the main text.

\vspace{0.3cm}

% ================================================================
% XII. FIGURES, DIAGRAMS, AND PHOTOGRAPHS
% ================================================================

\noindent
\textbf{XII.\quad FIGURES, DIAGRAMS, AND PHOTOGRAPHS}

\vspace{0.15cm}

Figures shall be clear, legible, and directly relevant to the report
content.

Figures shall be inserted using the standard \verb|figure| environment.

Example:

\begin{verbatim}
\begin{figure}[H]
    \centering
    \includegraphics[width=0.8\textwidth]
                    {figures/company_structure.pdf}
    \caption{Organizational structure of the company}
    \label{fig:company-structure}
\end{figure}
\end{verbatim}

Figure captions shall be placed below the figure.

Each referenced figure shall have a unique \verb|\label| and shall be
referenced in the text using, for example:

\begin{verbatim}
Figure~\ref{fig:company-structure}
\end{verbatim}

Students shall not manually type figure numbers.

Figures reproduced or adapted from external sources shall include an
appropriate citation.

Photographs taken during the internship shall be relevant and sufficiently
clear. Students shall not include confidential, restricted, or proprietary
information without permission from the internship organization.

Vector graphics or high-resolution images are recommended whenever
possible.

\vspace{0.3cm}

% ================================================================
% XIII. TABLES
% ================================================================

\noindent
\textbf{XIII.\quad TABLES}

\vspace{0.15cm}

Tables shall be created using suitable \LaTeX{} environments such as
\verb|table|, \verb|tabular|, \verb|tabularx|, or \verb|longtable|.

Table captions shall be placed above the table.

Example:

\begin{verbatim}
\begin{table}[H]
    \centering
    \caption{Internship schedule}
    \label{tab:internship-schedule}

    \begin{tabular}{lll}
        \toprule
        Week & Task & Result \\
        \midrule
        1 & Orientation & Completed \\
        2 & System study & Completed \\
        \bottomrule
    \end{tabular}
\end{table}
\end{verbatim}

Tables shall be referenced using:

\begin{verbatim}
Table~\ref{tab:internship-schedule}
\end{verbatim}

Students shall not manually type table numbers.

Tables should be created directly in \LaTeX{} rather than inserted as
screenshots whenever practical.

\vspace{0.3cm}

% ================================================================
% XIV. EQUATIONS AND TECHNICAL QUANTITIES
% ================================================================

\noindent
\textbf{XIV.\quad EQUATIONS AND TECHNICAL QUANTITIES}

\vspace{0.15cm}

Mathematical expressions shall be typeset directly using \LaTeX{}.
Images or screenshots of equations shall not be used.

Short expressions may be written in inline mathematical mode, for example:

\begin{verbatim}
$R = 0.7\,\Omega$
\end{verbatim}

Important equations shall use a numbered equation environment.

Example:

\begin{verbatim}
\begin{equation}
    P = UI
    \label{eq:power}
\end{equation}
\end{verbatim}

Equations shall be referenced using:

\begin{verbatim}
Equation~\eqref{eq:power}
\end{verbatim}

Equation numbers shall not be manually entered.

SI units shall be used wherever applicable. The \verb|siunitx| package
should be used where appropriate.

Examples:

\begin{verbatim}
\SI{220}{\volt}
\SI{50}{\hertz}
\SI{2.5}{\ampere}
\end{verbatim}

Notation and units shall be consistent throughout the report.

\vspace{0.3cm}

% ================================================================
% XV. SOURCE CODE
% ================================================================

\noindent
\textbf{XV.\quad SOURCE CODE}

\vspace{0.15cm}

If the internship involves programming, PLCs, microcontrollers,
simulation, or software development, relevant source code may be included
in the report.

Source code shall be presented using the \verb|listings| environment or
another environment defined by the official template.

Example:

\begin{verbatim}
\begin{lstlisting}[language=st]
IF Start = TRUE THEN
    Motor := TRUE;
ELSE
    Motor := FALSE;
END_IF;
\end{lstlisting}
\end{verbatim}

Only code necessary for explaining the work should appear in the main
report. Long source-code listings should be moved to an appendix or
submitted as supplementary files if required.

Source code should not normally be inserted as screenshots.

\vspace{0.3cm}

% ================================================================
% XVI. ABSTRACT
% ================================================================

\noindent
\textbf{XVI.\quad ABSTRACT}

\vspace{0.15cm}

The Abstract shall briefly summarize:

\begin{itemize}[leftmargin=0.8cm, itemsep=2pt]
    \item the internship organization;
    \item the objectives of the internship;
    \item the main tasks performed;
    \item the main results, knowledge, and skills obtained.
\end{itemize}

The Abstract should be concise and should not reproduce the full content
of the report chapter by chapter.

\vspace{0.3cm}

% ================================================================
% XVII. REFERENCES AND CITATIONS
% ================================================================

\noindent
\textbf{XVII.\quad REFERENCES AND CITATIONS}

\vspace{0.15cm}

All books, journal papers, conference papers, standards, websites,
datasheets, manuals, technical documents, figures, tables, or other
external sources used in preparing the report shall be properly cited.

References shall be managed through the bibliography system provided with
the official template.

Bibliographic entries shall normally be stored in:

\begin{verbatim}
References.bib
\end{verbatim}

References shall be cited in the report using:

\begin{verbatim}
\cite{reference-key}
\end{verbatim}

Students shall not manually type reference numbers such as ``[1]'',
``[2]'', etc.

Titles of foreign-language publications shall remain in their original
language in the bibliography.

All borrowed ideas, data, figures, tables, equations, methods, or results
shall be properly acknowledged.

\vspace{0.3cm}

% ================================================================
% XVIII. CROSS-REFERENCING
% ================================================================

\noindent
\textbf{XVIII.\quad CROSS-REFERENCING}

\vspace{0.15cm}

Numbered elements shall be cross-referenced using the standard \LaTeX{}
mechanism.

Use the following commands where appropriate:

\begin{verbatim}
\label{...}
\ref{...}
\eqref{...}
\end{verbatim}

For example:

\begin{verbatim}
As shown in Figure~\ref{fig:system}, ...
\end{verbatim}

Students shall not manually type numbers such as ``Figure 2.3'' when
those numbers can be generated automatically.

\vspace{0.3cm}

% ================================================================
% XIX. PAGE NUMBERING
% ================================================================

\noindent
\textbf{XIX.\quad PAGE NUMBERING}

\vspace{0.15cm}

Cover pages and official evaluation forms at the beginning of the report
shall not display page numbers.

The main body of the report shall use Arabic page numbers beginning from
Chapter 1.

Page numbering shall be generated automatically by the official template.
Students shall not manually enter page numbers.

\vspace{0.3cm}

% ================================================================
% XX. PROJECT FILE ORGANIZATION
% ================================================================

\noindent
\textbf{XX.\quad PROJECT FILE ORGANIZATION}

\vspace{0.15cm}

Students should maintain a clear and organized \LaTeX{} project
structure. A recommended structure is:

\begin{verbatim}
InternshipReport/
|
|-- main.tex
|-- References.bib
|
|-- figures/
|   |-- company.pdf
|   |-- system.png
|   `-- ...
|
`-- chapters/
    |-- chapter1.tex
    |-- chapter2.tex
    `-- chapter3.tex
\end{verbatim}

Figures should normally be stored in the \verb|figures| directory
specified by the official template.

File and directory names should preferably contain Latin letters,
numbers, hyphens, or underscores.

Recommended examples include:

\begin{verbatim}
company_structure.pdf
motor_control.png
chapter_01.tex
\end{verbatim}

Spaces and unnecessary special characters in file names should be
avoided where possible.

\vspace{0.3cm}

% ================================================================
% XXI. DOCUMENT COMPILATION AND FINAL CHECK
% ================================================================

\noindent
\textbf{XXI.\quad DOCUMENT COMPILATION AND FINAL CHECK}

\vspace{0.15cm}

Before submission, the complete \LaTeX{} project shall compile
successfully and produce the final PDF report.

The final document shall not contain unresolved references such as:

\begin{verbatim}
??
[?]
\end{verbatim}

Before generating the final PDF, students shall check:

\begin{itemize}[leftmargin=0.8cm, itemsep=2pt]
    \item the Table of Contents;
    \item page numbering;
    \item chapter, section, and subsection numbering;
    \item figure numbering;
    \item table numbering;
    \item equation numbering;
    \item cross-references;
    \item citations;
    \item bibliography;
    \item font display;
    \item mathematical symbols;
    \item image quality;
    \item completeness of all required forms.
\end{itemize}

The document may be compiled using TeX Live, MiKTeX, or another
compatible \LaTeX{} environment. Regardless of the compilation
environment, the resulting PDF shall conform to the official Faculty
template.

\vspace{0.3cm}

% ================================================================
% XXII. CONTENT THAT STUDENTS MAY EDIT
% ================================================================

\noindent
\textbf{XXII.\quad CONTENT THAT STUDENTS MAY EDIT}

\vspace{0.15cm}

Students may edit or replace:

\begin{itemize}[leftmargin=0.8cm, itemsep=2pt]
    \item student name;
    \item class;
    \item internship organization;
    \item supervisor name;
    \item Abstract;
    \item List of Symbols and Abbreviations;
    \item chapter titles and chapter content;
    \item figures and photographs;
    \item tables;
    \item equations;
    \item citations and references.
\end{itemize}

Students shall not modify the main document formatting settings in the
preamble, including the document class, page geometry, line spacing,
heading formatting, table-of-contents formatting, font configuration,
page layout, or cover-page macros unless specifically instructed.

\vspace{0.3cm}

% ================================================================
% XXIII. FINAL SUBMISSION REQUIREMENTS
% ================================================================

\noindent
\textbf{XXIII.\quad FINAL SUBMISSION REQUIREMENTS}

\vspace{0.15cm}

Before submission, the final report shall:

\begin{itemize}[leftmargin=0.8cm, itemsep=2pt]
    \item compile without errors that affect the report content or presentation;
    \item contain all required forms and report sections;
    \item contain correct student, class, internship organization, and supervisor information;
    \item contain an updated Table of Contents, List of Figures, and List of Tables;
    \item contain no unused example content from the template;
    \item contain no placeholder text such as \textit{Title of Chapter 1},
    \textit{Section 1.1}, or other sample entries;
    \item contain all required citations and references;
    \item contain no unresolved cross-references.
\end{itemize}

If the Faculty requires submission of the \LaTeX{} source, the complete
project shall include all \verb|.tex| files, the bibliography database,
figures, and any other files required to reproduce the submitted PDF.

\vspace{0.4cm}

\noindent
\textbf{Important:}
Before submitting the final internship report, students shall remove or
comment out the following command in the main document:

\begin{verbatim}
\input{Regulations.tex}
\end{verbatim}

For example:

\begin{verbatim}
% \input{Regulations.tex}
\end{verbatim}

This removes the guidance pages from the final submitted report while
preserving the regulations file in the \LaTeX{} project for reference.

} % End 13pt content

\clearpage

\end{verbatim}

This removes the guidance pages from the final submitted report while
preserving the regulations file in the \LaTeX{} project for reference.

} % End 13pt content

\clearpage

\end{verbatim}

For example:

\begin{verbatim}
% 
% =================================================================
% REGULATIONS FOR PREPARING THE GRADUATION INTERNSHIP REPORT
% USING LATEX
% =================================================================
%
% This file is intended to be included from the main document with:
%
%     
% =================================================================
% REGULATIONS FOR PREPARING THE GRADUATION INTERNSHIP REPORT
% USING LATEX
% =================================================================
%
% This file is intended to be included from the main document with:
%
%     \input{Regulations.tex}
%
% IMPORTANT:
% This section is provided for guidance only. Students shall remove
% or comment out the \input{Regulations.tex} command before submitting
% the final internship report.
%
% =================================================================

\clearpage
\thispagestyle{empty}

\begin{center}
    {\fontsize{16pt}{20pt}\selectfont\bfseries
    SPECIFIC REGULATIONS FOR PREPARING\\
    GRADUATION INTERNSHIP REPORTS USING \LaTeX}
\end{center}

\vspace{0.5cm}

{\fontsize{13pt}{17pt}\selectfont
\setstretch{1.3}
\justifying

\noindent
\textbf{Note:}
This section is provided for guidance only and is included in the main
document using the command \verb|\input{Regulations.tex}|.
Students shall remove or comment out this command before submitting the
final internship report.

\vspace{0.4cm}

% ================================================================
% I. GENERAL REQUIREMENTS
% ================================================================

\noindent
\textbf{I.\quad GENERAL REQUIREMENTS}

\vspace{0.15cm}

The Graduation Internship Report shall be prepared using the official
\LaTeX{} template provided by the Faculty. Students shall preserve the
predefined page layout, typography, line spacing, heading styles,
numbering system, figure and table formatting, and bibliography settings
unless otherwise instructed by the Faculty or the supervisor.

The report shall be written in English.

Students are responsible for the academic and technical content of the
report. The official \LaTeX{} template is responsible for standardized
document formatting.

\vspace{0.3cm}

% ================================================================
% II. REPORT STRUCTURE
% ================================================================

\noindent
\textbf{II.\quad REPORT STRUCTURE}

\vspace{0.15cm}

The report shall normally be arranged in the following order:

\begin{enumerate}[label=\arabic*., leftmargin=1.2cm, itemsep=2pt]
    \item Outer Cover;
    \item Comments of the Internship Host Organization;
    \item Supervisor's Evaluation and Graduation Internship Defense Results;
    \item Table of Contents;
    \item List of Symbols and Abbreviations, if applicable;
    \item List of Figures;
    \item List of Tables;
    \item Abstract;
    \item Main Chapters of the Internship Report;
    \item References;
    \item Appendices, if applicable.
\end{enumerate}

The preliminary pages shall not be numbered as chapters. The main report
content shall be organized into consecutively numbered chapters beginning
with Chapter 1.

\vspace{0.3cm}

% ================================================================
% III. PAGE SIZE AND MARGINS
% ================================================================

\noindent
\textbf{III.\quad PAGE SIZE AND MARGINS}

\vspace{0.15cm}

The report shall be formatted on A4 paper with the following margins:

\begin{center}
\begin{tabular}{ll}
    Left margin:   & 3.5 cm \\
    Right margin:  & 2.0 cm \\
    Top margin:    & 2.5 cm \\
    Bottom margin: & 2.5 cm
\end{tabular}
\end{center}

These settings are predefined in the official template and shall not be
modified by students.

\vspace{0.3cm}

% ================================================================
% IV. FONT, FONT SIZE, AND LINE SPACING
% ================================================================

\noindent
\textbf{IV.\quad FONT, FONT SIZE, AND LINE SPACING}

\vspace{0.15cm}

The main text shall use a 13 pt font size with 1.3-line spacing.
A Times-family typeface shall be used in accordance with the official
template.

The text shall be fully justified.

Students shall not use repeated spaces, repeated manual line breaks, or
multiple blank lines to position text. Paragraph and heading spacing shall
be controlled by the template.

\vspace{0.3cm}

% ================================================================
% V. OUTER COVER
% ================================================================

\noindent
\textbf{V.\quad OUTER COVER}

\vspace{0.15cm}

The outer cover shall follow the official format provided in the template
and shall contain the following information:

\begin{itemize}[leftmargin=0.8cm, itemsep=2pt]
    \item Thai Nguyen University of Technology;
    \item Faculty of International Training;
    \item Engineering and Technology Department;
    \item University logo and Faculty logo;
    \item the title \textbf{INTERNSHIP REPORT};
    \item Internship Organization;
    \item Author;
    \item Class;
    \item Supervisor;
    \item place and year of completion.
\end{itemize}

Students shall only replace the required information fields and shall not
modify the cover-page layout unless instructed otherwise.

\vspace{0.3cm}

% ================================================================
% VI. COMMENTS OF THE INTERNSHIP HOST ORGANIZATION
% ================================================================

\noindent
\textbf{VI.\quad COMMENTS OF THE INTERNSHIP HOST ORGANIZATION}

\vspace{0.15cm}

The page entitled \textbf{COMMENTS OF THE INTERNSHIP HOST ORGANIZATION}
is reserved for the internship host organization.

It includes:

\begin{itemize}[leftmargin=0.8cm, itemsep=2pt]
    \item the student's full name;
    \item class;
    \item comments from the host organization;
    \item score;
    \item date;
    \item signature and official seal.
\end{itemize}

Students shall not alter or complete official evaluation fields unless
specifically authorized.

\vspace{0.3cm}

% ================================================================
% VII. SUPERVISOR'S EVALUATION AND DEFENSE RESULTS
% ================================================================

\noindent
\textbf{VII.\quad SUPERVISOR'S EVALUATION AND DEFENSE RESULTS}

\vspace{0.15cm}

The section entitled
\textbf{SUPERVISOR'S EVALUATION AND GRADUATION INTERNSHIP DEFENSE RESULTS}
is reserved for the supervisor and defense committee.

The evaluation may include:

\begin{itemize}[leftmargin=0.8cm, itemsep=2pt]
    \item Internship Progress and Student's Attitude;
    \item Report Content;
    \item Report Presentation;
    \item Other Comments;
    \item Supervisor's Opinion;
    \item Defense Result;
    \item Score.
\end{itemize}

Students shall not enter evaluation scores or defense results before they
are officially determined.

\vspace{0.3cm}

% ================================================================
% VIII. CHAPTERS, SECTIONS, AND SUBSECTIONS
% ================================================================

\noindent
\textbf{VIII.\quad CHAPTERS, SECTIONS, AND SUBSECTIONS}

\vspace{0.15cm}

The report structure shall be generated using the standard \LaTeX{}
commands:

\begin{verbatim}
\chapter{...}
\section{...}
\subsection{...}
\end{verbatim}

Chapter, section, and subsection numbers shall be generated automatically
by \LaTeX{}.

Students shall not manually type chapter or section numbers and shall not
reproduce heading styles using commands such as \verb|\textbf|,
\verb|\fontsize|, or \verb|\centering| instead of the appropriate
structural command.

Excessive subdivision should be avoided. In most cases, the hierarchy
Chapter--Section--Subsection is sufficient.

\vspace{0.3cm}

% ================================================================
% IX. ORGANIZATION OF THE MAIN CONTENT
% ================================================================

\noindent
\textbf{IX.\quad ORGANIZATION OF THE MAIN CONTENT}

\vspace{0.15cm}

The specific content of the report shall depend on the internship
organization and the assigned tasks. A typical structure may include:

\begin{itemize}[leftmargin=0.8cm, itemsep=2pt]
    \item \textbf{Chapter 1. Overview of the Internship Organization};
    \item \textbf{Chapter 2. Internship Activities and Assigned Tasks};
    \item \textbf{Chapter 3. Results, Experience, and Evaluation}.
\end{itemize}

These chapter titles are recommendations only and may be adapted to the
actual internship content.

The report shall reflect work actually performed, observed, analyzed, or
learned during the internship. Unrelated theoretical material should not
be included merely to increase report length.

\vspace{0.3cm}

% ================================================================
% X. TABLE OF CONTENTS
% ================================================================

\noindent
\textbf{X.\quad TABLE OF CONTENTS}

\vspace{0.15cm}

The Table of Contents shall be generated automatically using:

\begin{verbatim}
\tableofcontents
\end{verbatim}

Students shall not create or manually edit page numbers in the Table of
Contents.

The document shall be compiled sufficiently to ensure that all page
numbers and entries are correctly updated.

\vspace{0.3cm}

% ================================================================
% XI. LIST OF SYMBOLS AND ABBREVIATIONS
% ================================================================

\noindent
\textbf{XI.\quad LIST OF SYMBOLS AND ABBREVIATIONS}

\vspace{0.15cm}

Symbols and abbreviations that are used repeatedly in the report shall be
listed in the \textbf{List of Symbols and Abbreviations}.

The list may contain the following columns:

\begin{center}
\begin{tabular}{llll}
    No. & Symbol / Abbreviation & Unit & Definition
\end{tabular}
\end{center}

Only symbols and abbreviations actually used in the final report shall be
retained. Example entries contained in the template shall be removed if
they are not relevant to the student's report.

Each abbreviation should normally be defined when it first appears in
the main text.

\vspace{0.3cm}

% ================================================================
% XII. FIGURES, DIAGRAMS, AND PHOTOGRAPHS
% ================================================================

\noindent
\textbf{XII.\quad FIGURES, DIAGRAMS, AND PHOTOGRAPHS}

\vspace{0.15cm}

Figures shall be clear, legible, and directly relevant to the report
content.

Figures shall be inserted using the standard \verb|figure| environment.

Example:

\begin{verbatim}
\begin{figure}[H]
    \centering
    \includegraphics[width=0.8\textwidth]
                    {figures/company_structure.pdf}
    \caption{Organizational structure of the company}
    \label{fig:company-structure}
\end{figure}
\end{verbatim}

Figure captions shall be placed below the figure.

Each referenced figure shall have a unique \verb|\label| and shall be
referenced in the text using, for example:

\begin{verbatim}
Figure~\ref{fig:company-structure}
\end{verbatim}

Students shall not manually type figure numbers.

Figures reproduced or adapted from external sources shall include an
appropriate citation.

Photographs taken during the internship shall be relevant and sufficiently
clear. Students shall not include confidential, restricted, or proprietary
information without permission from the internship organization.

Vector graphics or high-resolution images are recommended whenever
possible.

\vspace{0.3cm}

% ================================================================
% XIII. TABLES
% ================================================================

\noindent
\textbf{XIII.\quad TABLES}

\vspace{0.15cm}

Tables shall be created using suitable \LaTeX{} environments such as
\verb|table|, \verb|tabular|, \verb|tabularx|, or \verb|longtable|.

Table captions shall be placed above the table.

Example:

\begin{verbatim}
\begin{table}[H]
    \centering
    \caption{Internship schedule}
    \label{tab:internship-schedule}

    \begin{tabular}{lll}
        \toprule
        Week & Task & Result \\
        \midrule
        1 & Orientation & Completed \\
        2 & System study & Completed \\
        \bottomrule
    \end{tabular}
\end{table}
\end{verbatim}

Tables shall be referenced using:

\begin{verbatim}
Table~\ref{tab:internship-schedule}
\end{verbatim}

Students shall not manually type table numbers.

Tables should be created directly in \LaTeX{} rather than inserted as
screenshots whenever practical.

\vspace{0.3cm}

% ================================================================
% XIV. EQUATIONS AND TECHNICAL QUANTITIES
% ================================================================

\noindent
\textbf{XIV.\quad EQUATIONS AND TECHNICAL QUANTITIES}

\vspace{0.15cm}

Mathematical expressions shall be typeset directly using \LaTeX{}.
Images or screenshots of equations shall not be used.

Short expressions may be written in inline mathematical mode, for example:

\begin{verbatim}
$R = 0.7\,\Omega$
\end{verbatim}

Important equations shall use a numbered equation environment.

Example:

\begin{verbatim}
\begin{equation}
    P = UI
    \label{eq:power}
\end{equation}
\end{verbatim}

Equations shall be referenced using:

\begin{verbatim}
Equation~\eqref{eq:power}
\end{verbatim}

Equation numbers shall not be manually entered.

SI units shall be used wherever applicable. The \verb|siunitx| package
should be used where appropriate.

Examples:

\begin{verbatim}
\SI{220}{\volt}
\SI{50}{\hertz}
\SI{2.5}{\ampere}
\end{verbatim}

Notation and units shall be consistent throughout the report.

\vspace{0.3cm}

% ================================================================
% XV. SOURCE CODE
% ================================================================

\noindent
\textbf{XV.\quad SOURCE CODE}

\vspace{0.15cm}

If the internship involves programming, PLCs, microcontrollers,
simulation, or software development, relevant source code may be included
in the report.

Source code shall be presented using the \verb|listings| environment or
another environment defined by the official template.

Example:

\begin{verbatim}
\begin{lstlisting}[language=st]
IF Start = TRUE THEN
    Motor := TRUE;
ELSE
    Motor := FALSE;
END_IF;
\end{lstlisting}
\end{verbatim}

Only code necessary for explaining the work should appear in the main
report. Long source-code listings should be moved to an appendix or
submitted as supplementary files if required.

Source code should not normally be inserted as screenshots.

\vspace{0.3cm}

% ================================================================
% XVI. ABSTRACT
% ================================================================

\noindent
\textbf{XVI.\quad ABSTRACT}

\vspace{0.15cm}

The Abstract shall briefly summarize:

\begin{itemize}[leftmargin=0.8cm, itemsep=2pt]
    \item the internship organization;
    \item the objectives of the internship;
    \item the main tasks performed;
    \item the main results, knowledge, and skills obtained.
\end{itemize}

The Abstract should be concise and should not reproduce the full content
of the report chapter by chapter.

\vspace{0.3cm}

% ================================================================
% XVII. REFERENCES AND CITATIONS
% ================================================================

\noindent
\textbf{XVII.\quad REFERENCES AND CITATIONS}

\vspace{0.15cm}

All books, journal papers, conference papers, standards, websites,
datasheets, manuals, technical documents, figures, tables, or other
external sources used in preparing the report shall be properly cited.

References shall be managed through the bibliography system provided with
the official template.

Bibliographic entries shall normally be stored in:

\begin{verbatim}
References.bib
\end{verbatim}

References shall be cited in the report using:

\begin{verbatim}
\cite{reference-key}
\end{verbatim}

Students shall not manually type reference numbers such as ``[1]'',
``[2]'', etc.

Titles of foreign-language publications shall remain in their original
language in the bibliography.

All borrowed ideas, data, figures, tables, equations, methods, or results
shall be properly acknowledged.

\vspace{0.3cm}

% ================================================================
% XVIII. CROSS-REFERENCING
% ================================================================

\noindent
\textbf{XVIII.\quad CROSS-REFERENCING}

\vspace{0.15cm}

Numbered elements shall be cross-referenced using the standard \LaTeX{}
mechanism.

Use the following commands where appropriate:

\begin{verbatim}
\label{...}
\ref{...}
\eqref{...}
\end{verbatim}

For example:

\begin{verbatim}
As shown in Figure~\ref{fig:system}, ...
\end{verbatim}

Students shall not manually type numbers such as ``Figure 2.3'' when
those numbers can be generated automatically.

\vspace{0.3cm}

% ================================================================
% XIX. PAGE NUMBERING
% ================================================================

\noindent
\textbf{XIX.\quad PAGE NUMBERING}

\vspace{0.15cm}

Cover pages and official evaluation forms at the beginning of the report
shall not display page numbers.

The main body of the report shall use Arabic page numbers beginning from
Chapter 1.

Page numbering shall be generated automatically by the official template.
Students shall not manually enter page numbers.

\vspace{0.3cm}

% ================================================================
% XX. PROJECT FILE ORGANIZATION
% ================================================================

\noindent
\textbf{XX.\quad PROJECT FILE ORGANIZATION}

\vspace{0.15cm}

Students should maintain a clear and organized \LaTeX{} project
structure. A recommended structure is:

\begin{verbatim}
InternshipReport/
|
|-- main.tex
|-- References.bib
|
|-- figures/
|   |-- company.pdf
|   |-- system.png
|   `-- ...
|
`-- chapters/
    |-- chapter1.tex
    |-- chapter2.tex
    `-- chapter3.tex
\end{verbatim}

Figures should normally be stored in the \verb|figures| directory
specified by the official template.

File and directory names should preferably contain Latin letters,
numbers, hyphens, or underscores.

Recommended examples include:

\begin{verbatim}
company_structure.pdf
motor_control.png
chapter_01.tex
\end{verbatim}

Spaces and unnecessary special characters in file names should be
avoided where possible.

\vspace{0.3cm}

% ================================================================
% XXI. DOCUMENT COMPILATION AND FINAL CHECK
% ================================================================

\noindent
\textbf{XXI.\quad DOCUMENT COMPILATION AND FINAL CHECK}

\vspace{0.15cm}

Before submission, the complete \LaTeX{} project shall compile
successfully and produce the final PDF report.

The final document shall not contain unresolved references such as:

\begin{verbatim}
??
[?]
\end{verbatim}

Before generating the final PDF, students shall check:

\begin{itemize}[leftmargin=0.8cm, itemsep=2pt]
    \item the Table of Contents;
    \item page numbering;
    \item chapter, section, and subsection numbering;
    \item figure numbering;
    \item table numbering;
    \item equation numbering;
    \item cross-references;
    \item citations;
    \item bibliography;
    \item font display;
    \item mathematical symbols;
    \item image quality;
    \item completeness of all required forms.
\end{itemize}

The document may be compiled using TeX Live, MiKTeX, or another
compatible \LaTeX{} environment. Regardless of the compilation
environment, the resulting PDF shall conform to the official Faculty
template.

\vspace{0.3cm}

% ================================================================
% XXII. CONTENT THAT STUDENTS MAY EDIT
% ================================================================

\noindent
\textbf{XXII.\quad CONTENT THAT STUDENTS MAY EDIT}

\vspace{0.15cm}

Students may edit or replace:

\begin{itemize}[leftmargin=0.8cm, itemsep=2pt]
    \item student name;
    \item class;
    \item internship organization;
    \item supervisor name;
    \item Abstract;
    \item List of Symbols and Abbreviations;
    \item chapter titles and chapter content;
    \item figures and photographs;
    \item tables;
    \item equations;
    \item citations and references.
\end{itemize}

Students shall not modify the main document formatting settings in the
preamble, including the document class, page geometry, line spacing,
heading formatting, table-of-contents formatting, font configuration,
page layout, or cover-page macros unless specifically instructed.

\vspace{0.3cm}

% ================================================================
% XXIII. FINAL SUBMISSION REQUIREMENTS
% ================================================================

\noindent
\textbf{XXIII.\quad FINAL SUBMISSION REQUIREMENTS}

\vspace{0.15cm}

Before submission, the final report shall:

\begin{itemize}[leftmargin=0.8cm, itemsep=2pt]
    \item compile without errors that affect the report content or presentation;
    \item contain all required forms and report sections;
    \item contain correct student, class, internship organization, and supervisor information;
    \item contain an updated Table of Contents, List of Figures, and List of Tables;
    \item contain no unused example content from the template;
    \item contain no placeholder text such as \textit{Title of Chapter 1},
    \textit{Section 1.1}, or other sample entries;
    \item contain all required citations and references;
    \item contain no unresolved cross-references.
\end{itemize}

If the Faculty requires submission of the \LaTeX{} source, the complete
project shall include all \verb|.tex| files, the bibliography database,
figures, and any other files required to reproduce the submitted PDF.

\vspace{0.4cm}

\noindent
\textbf{Important:}
Before submitting the final internship report, students shall remove or
comment out the following command in the main document:

\begin{verbatim}
\input{Regulations.tex}
\end{verbatim}

For example:

\begin{verbatim}
% \input{Regulations.tex}
\end{verbatim}

This removes the guidance pages from the final submitted report while
preserving the regulations file in the \LaTeX{} project for reference.

} % End 13pt content

\clearpage

%
% IMPORTANT:
% This section is provided for guidance only. Students shall remove
% or comment out the 
% =================================================================
% REGULATIONS FOR PREPARING THE GRADUATION INTERNSHIP REPORT
% USING LATEX
% =================================================================
%
% This file is intended to be included from the main document with:
%
%     \input{Regulations.tex}
%
% IMPORTANT:
% This section is provided for guidance only. Students shall remove
% or comment out the \input{Regulations.tex} command before submitting
% the final internship report.
%
% =================================================================

\clearpage
\thispagestyle{empty}

\begin{center}
    {\fontsize{16pt}{20pt}\selectfont\bfseries
    SPECIFIC REGULATIONS FOR PREPARING\\
    GRADUATION INTERNSHIP REPORTS USING \LaTeX}
\end{center}

\vspace{0.5cm}

{\fontsize{13pt}{17pt}\selectfont
\setstretch{1.3}
\justifying

\noindent
\textbf{Note:}
This section is provided for guidance only and is included in the main
document using the command \verb|\input{Regulations.tex}|.
Students shall remove or comment out this command before submitting the
final internship report.

\vspace{0.4cm}

% ================================================================
% I. GENERAL REQUIREMENTS
% ================================================================

\noindent
\textbf{I.\quad GENERAL REQUIREMENTS}

\vspace{0.15cm}

The Graduation Internship Report shall be prepared using the official
\LaTeX{} template provided by the Faculty. Students shall preserve the
predefined page layout, typography, line spacing, heading styles,
numbering system, figure and table formatting, and bibliography settings
unless otherwise instructed by the Faculty or the supervisor.

The report shall be written in English.

Students are responsible for the academic and technical content of the
report. The official \LaTeX{} template is responsible for standardized
document formatting.

\vspace{0.3cm}

% ================================================================
% II. REPORT STRUCTURE
% ================================================================

\noindent
\textbf{II.\quad REPORT STRUCTURE}

\vspace{0.15cm}

The report shall normally be arranged in the following order:

\begin{enumerate}[label=\arabic*., leftmargin=1.2cm, itemsep=2pt]
    \item Outer Cover;
    \item Comments of the Internship Host Organization;
    \item Supervisor's Evaluation and Graduation Internship Defense Results;
    \item Table of Contents;
    \item List of Symbols and Abbreviations, if applicable;
    \item List of Figures;
    \item List of Tables;
    \item Abstract;
    \item Main Chapters of the Internship Report;
    \item References;
    \item Appendices, if applicable.
\end{enumerate}

The preliminary pages shall not be numbered as chapters. The main report
content shall be organized into consecutively numbered chapters beginning
with Chapter 1.

\vspace{0.3cm}

% ================================================================
% III. PAGE SIZE AND MARGINS
% ================================================================

\noindent
\textbf{III.\quad PAGE SIZE AND MARGINS}

\vspace{0.15cm}

The report shall be formatted on A4 paper with the following margins:

\begin{center}
\begin{tabular}{ll}
    Left margin:   & 3.5 cm \\
    Right margin:  & 2.0 cm \\
    Top margin:    & 2.5 cm \\
    Bottom margin: & 2.5 cm
\end{tabular}
\end{center}

These settings are predefined in the official template and shall not be
modified by students.

\vspace{0.3cm}

% ================================================================
% IV. FONT, FONT SIZE, AND LINE SPACING
% ================================================================

\noindent
\textbf{IV.\quad FONT, FONT SIZE, AND LINE SPACING}

\vspace{0.15cm}

The main text shall use a 13 pt font size with 1.3-line spacing.
A Times-family typeface shall be used in accordance with the official
template.

The text shall be fully justified.

Students shall not use repeated spaces, repeated manual line breaks, or
multiple blank lines to position text. Paragraph and heading spacing shall
be controlled by the template.

\vspace{0.3cm}

% ================================================================
% V. OUTER COVER
% ================================================================

\noindent
\textbf{V.\quad OUTER COVER}

\vspace{0.15cm}

The outer cover shall follow the official format provided in the template
and shall contain the following information:

\begin{itemize}[leftmargin=0.8cm, itemsep=2pt]
    \item Thai Nguyen University of Technology;
    \item Faculty of International Training;
    \item Engineering and Technology Department;
    \item University logo and Faculty logo;
    \item the title \textbf{INTERNSHIP REPORT};
    \item Internship Organization;
    \item Author;
    \item Class;
    \item Supervisor;
    \item place and year of completion.
\end{itemize}

Students shall only replace the required information fields and shall not
modify the cover-page layout unless instructed otherwise.

\vspace{0.3cm}

% ================================================================
% VI. COMMENTS OF THE INTERNSHIP HOST ORGANIZATION
% ================================================================

\noindent
\textbf{VI.\quad COMMENTS OF THE INTERNSHIP HOST ORGANIZATION}

\vspace{0.15cm}

The page entitled \textbf{COMMENTS OF THE INTERNSHIP HOST ORGANIZATION}
is reserved for the internship host organization.

It includes:

\begin{itemize}[leftmargin=0.8cm, itemsep=2pt]
    \item the student's full name;
    \item class;
    \item comments from the host organization;
    \item score;
    \item date;
    \item signature and official seal.
\end{itemize}

Students shall not alter or complete official evaluation fields unless
specifically authorized.

\vspace{0.3cm}

% ================================================================
% VII. SUPERVISOR'S EVALUATION AND DEFENSE RESULTS
% ================================================================

\noindent
\textbf{VII.\quad SUPERVISOR'S EVALUATION AND DEFENSE RESULTS}

\vspace{0.15cm}

The section entitled
\textbf{SUPERVISOR'S EVALUATION AND GRADUATION INTERNSHIP DEFENSE RESULTS}
is reserved for the supervisor and defense committee.

The evaluation may include:

\begin{itemize}[leftmargin=0.8cm, itemsep=2pt]
    \item Internship Progress and Student's Attitude;
    \item Report Content;
    \item Report Presentation;
    \item Other Comments;
    \item Supervisor's Opinion;
    \item Defense Result;
    \item Score.
\end{itemize}

Students shall not enter evaluation scores or defense results before they
are officially determined.

\vspace{0.3cm}

% ================================================================
% VIII. CHAPTERS, SECTIONS, AND SUBSECTIONS
% ================================================================

\noindent
\textbf{VIII.\quad CHAPTERS, SECTIONS, AND SUBSECTIONS}

\vspace{0.15cm}

The report structure shall be generated using the standard \LaTeX{}
commands:

\begin{verbatim}
\chapter{...}
\section{...}
\subsection{...}
\end{verbatim}

Chapter, section, and subsection numbers shall be generated automatically
by \LaTeX{}.

Students shall not manually type chapter or section numbers and shall not
reproduce heading styles using commands such as \verb|\textbf|,
\verb|\fontsize|, or \verb|\centering| instead of the appropriate
structural command.

Excessive subdivision should be avoided. In most cases, the hierarchy
Chapter--Section--Subsection is sufficient.

\vspace{0.3cm}

% ================================================================
% IX. ORGANIZATION OF THE MAIN CONTENT
% ================================================================

\noindent
\textbf{IX.\quad ORGANIZATION OF THE MAIN CONTENT}

\vspace{0.15cm}

The specific content of the report shall depend on the internship
organization and the assigned tasks. A typical structure may include:

\begin{itemize}[leftmargin=0.8cm, itemsep=2pt]
    \item \textbf{Chapter 1. Overview of the Internship Organization};
    \item \textbf{Chapter 2. Internship Activities and Assigned Tasks};
    \item \textbf{Chapter 3. Results, Experience, and Evaluation}.
\end{itemize}

These chapter titles are recommendations only and may be adapted to the
actual internship content.

The report shall reflect work actually performed, observed, analyzed, or
learned during the internship. Unrelated theoretical material should not
be included merely to increase report length.

\vspace{0.3cm}

% ================================================================
% X. TABLE OF CONTENTS
% ================================================================

\noindent
\textbf{X.\quad TABLE OF CONTENTS}

\vspace{0.15cm}

The Table of Contents shall be generated automatically using:

\begin{verbatim}
\tableofcontents
\end{verbatim}

Students shall not create or manually edit page numbers in the Table of
Contents.

The document shall be compiled sufficiently to ensure that all page
numbers and entries are correctly updated.

\vspace{0.3cm}

% ================================================================
% XI. LIST OF SYMBOLS AND ABBREVIATIONS
% ================================================================

\noindent
\textbf{XI.\quad LIST OF SYMBOLS AND ABBREVIATIONS}

\vspace{0.15cm}

Symbols and abbreviations that are used repeatedly in the report shall be
listed in the \textbf{List of Symbols and Abbreviations}.

The list may contain the following columns:

\begin{center}
\begin{tabular}{llll}
    No. & Symbol / Abbreviation & Unit & Definition
\end{tabular}
\end{center}

Only symbols and abbreviations actually used in the final report shall be
retained. Example entries contained in the template shall be removed if
they are not relevant to the student's report.

Each abbreviation should normally be defined when it first appears in
the main text.

\vspace{0.3cm}

% ================================================================
% XII. FIGURES, DIAGRAMS, AND PHOTOGRAPHS
% ================================================================

\noindent
\textbf{XII.\quad FIGURES, DIAGRAMS, AND PHOTOGRAPHS}

\vspace{0.15cm}

Figures shall be clear, legible, and directly relevant to the report
content.

Figures shall be inserted using the standard \verb|figure| environment.

Example:

\begin{verbatim}
\begin{figure}[H]
    \centering
    \includegraphics[width=0.8\textwidth]
                    {figures/company_structure.pdf}
    \caption{Organizational structure of the company}
    \label{fig:company-structure}
\end{figure}
\end{verbatim}

Figure captions shall be placed below the figure.

Each referenced figure shall have a unique \verb|\label| and shall be
referenced in the text using, for example:

\begin{verbatim}
Figure~\ref{fig:company-structure}
\end{verbatim}

Students shall not manually type figure numbers.

Figures reproduced or adapted from external sources shall include an
appropriate citation.

Photographs taken during the internship shall be relevant and sufficiently
clear. Students shall not include confidential, restricted, or proprietary
information without permission from the internship organization.

Vector graphics or high-resolution images are recommended whenever
possible.

\vspace{0.3cm}

% ================================================================
% XIII. TABLES
% ================================================================

\noindent
\textbf{XIII.\quad TABLES}

\vspace{0.15cm}

Tables shall be created using suitable \LaTeX{} environments such as
\verb|table|, \verb|tabular|, \verb|tabularx|, or \verb|longtable|.

Table captions shall be placed above the table.

Example:

\begin{verbatim}
\begin{table}[H]
    \centering
    \caption{Internship schedule}
    \label{tab:internship-schedule}

    \begin{tabular}{lll}
        \toprule
        Week & Task & Result \\
        \midrule
        1 & Orientation & Completed \\
        2 & System study & Completed \\
        \bottomrule
    \end{tabular}
\end{table}
\end{verbatim}

Tables shall be referenced using:

\begin{verbatim}
Table~\ref{tab:internship-schedule}
\end{verbatim}

Students shall not manually type table numbers.

Tables should be created directly in \LaTeX{} rather than inserted as
screenshots whenever practical.

\vspace{0.3cm}

% ================================================================
% XIV. EQUATIONS AND TECHNICAL QUANTITIES
% ================================================================

\noindent
\textbf{XIV.\quad EQUATIONS AND TECHNICAL QUANTITIES}

\vspace{0.15cm}

Mathematical expressions shall be typeset directly using \LaTeX{}.
Images or screenshots of equations shall not be used.

Short expressions may be written in inline mathematical mode, for example:

\begin{verbatim}
$R = 0.7\,\Omega$
\end{verbatim}

Important equations shall use a numbered equation environment.

Example:

\begin{verbatim}
\begin{equation}
    P = UI
    \label{eq:power}
\end{equation}
\end{verbatim}

Equations shall be referenced using:

\begin{verbatim}
Equation~\eqref{eq:power}
\end{verbatim}

Equation numbers shall not be manually entered.

SI units shall be used wherever applicable. The \verb|siunitx| package
should be used where appropriate.

Examples:

\begin{verbatim}
\SI{220}{\volt}
\SI{50}{\hertz}
\SI{2.5}{\ampere}
\end{verbatim}

Notation and units shall be consistent throughout the report.

\vspace{0.3cm}

% ================================================================
% XV. SOURCE CODE
% ================================================================

\noindent
\textbf{XV.\quad SOURCE CODE}

\vspace{0.15cm}

If the internship involves programming, PLCs, microcontrollers,
simulation, or software development, relevant source code may be included
in the report.

Source code shall be presented using the \verb|listings| environment or
another environment defined by the official template.

Example:

\begin{verbatim}
\begin{lstlisting}[language=st]
IF Start = TRUE THEN
    Motor := TRUE;
ELSE
    Motor := FALSE;
END_IF;
\end{lstlisting}
\end{verbatim}

Only code necessary for explaining the work should appear in the main
report. Long source-code listings should be moved to an appendix or
submitted as supplementary files if required.

Source code should not normally be inserted as screenshots.

\vspace{0.3cm}

% ================================================================
% XVI. ABSTRACT
% ================================================================

\noindent
\textbf{XVI.\quad ABSTRACT}

\vspace{0.15cm}

The Abstract shall briefly summarize:

\begin{itemize}[leftmargin=0.8cm, itemsep=2pt]
    \item the internship organization;
    \item the objectives of the internship;
    \item the main tasks performed;
    \item the main results, knowledge, and skills obtained.
\end{itemize}

The Abstract should be concise and should not reproduce the full content
of the report chapter by chapter.

\vspace{0.3cm}

% ================================================================
% XVII. REFERENCES AND CITATIONS
% ================================================================

\noindent
\textbf{XVII.\quad REFERENCES AND CITATIONS}

\vspace{0.15cm}

All books, journal papers, conference papers, standards, websites,
datasheets, manuals, technical documents, figures, tables, or other
external sources used in preparing the report shall be properly cited.

References shall be managed through the bibliography system provided with
the official template.

Bibliographic entries shall normally be stored in:

\begin{verbatim}
References.bib
\end{verbatim}

References shall be cited in the report using:

\begin{verbatim}
\cite{reference-key}
\end{verbatim}

Students shall not manually type reference numbers such as ``[1]'',
``[2]'', etc.

Titles of foreign-language publications shall remain in their original
language in the bibliography.

All borrowed ideas, data, figures, tables, equations, methods, or results
shall be properly acknowledged.

\vspace{0.3cm}

% ================================================================
% XVIII. CROSS-REFERENCING
% ================================================================

\noindent
\textbf{XVIII.\quad CROSS-REFERENCING}

\vspace{0.15cm}

Numbered elements shall be cross-referenced using the standard \LaTeX{}
mechanism.

Use the following commands where appropriate:

\begin{verbatim}
\label{...}
\ref{...}
\eqref{...}
\end{verbatim}

For example:

\begin{verbatim}
As shown in Figure~\ref{fig:system}, ...
\end{verbatim}

Students shall not manually type numbers such as ``Figure 2.3'' when
those numbers can be generated automatically.

\vspace{0.3cm}

% ================================================================
% XIX. PAGE NUMBERING
% ================================================================

\noindent
\textbf{XIX.\quad PAGE NUMBERING}

\vspace{0.15cm}

Cover pages and official evaluation forms at the beginning of the report
shall not display page numbers.

The main body of the report shall use Arabic page numbers beginning from
Chapter 1.

Page numbering shall be generated automatically by the official template.
Students shall not manually enter page numbers.

\vspace{0.3cm}

% ================================================================
% XX. PROJECT FILE ORGANIZATION
% ================================================================

\noindent
\textbf{XX.\quad PROJECT FILE ORGANIZATION}

\vspace{0.15cm}

Students should maintain a clear and organized \LaTeX{} project
structure. A recommended structure is:

\begin{verbatim}
InternshipReport/
|
|-- main.tex
|-- References.bib
|
|-- figures/
|   |-- company.pdf
|   |-- system.png
|   `-- ...
|
`-- chapters/
    |-- chapter1.tex
    |-- chapter2.tex
    `-- chapter3.tex
\end{verbatim}

Figures should normally be stored in the \verb|figures| directory
specified by the official template.

File and directory names should preferably contain Latin letters,
numbers, hyphens, or underscores.

Recommended examples include:

\begin{verbatim}
company_structure.pdf
motor_control.png
chapter_01.tex
\end{verbatim}

Spaces and unnecessary special characters in file names should be
avoided where possible.

\vspace{0.3cm}

% ================================================================
% XXI. DOCUMENT COMPILATION AND FINAL CHECK
% ================================================================

\noindent
\textbf{XXI.\quad DOCUMENT COMPILATION AND FINAL CHECK}

\vspace{0.15cm}

Before submission, the complete \LaTeX{} project shall compile
successfully and produce the final PDF report.

The final document shall not contain unresolved references such as:

\begin{verbatim}
??
[?]
\end{verbatim}

Before generating the final PDF, students shall check:

\begin{itemize}[leftmargin=0.8cm, itemsep=2pt]
    \item the Table of Contents;
    \item page numbering;
    \item chapter, section, and subsection numbering;
    \item figure numbering;
    \item table numbering;
    \item equation numbering;
    \item cross-references;
    \item citations;
    \item bibliography;
    \item font display;
    \item mathematical symbols;
    \item image quality;
    \item completeness of all required forms.
\end{itemize}

The document may be compiled using TeX Live, MiKTeX, or another
compatible \LaTeX{} environment. Regardless of the compilation
environment, the resulting PDF shall conform to the official Faculty
template.

\vspace{0.3cm}

% ================================================================
% XXII. CONTENT THAT STUDENTS MAY EDIT
% ================================================================

\noindent
\textbf{XXII.\quad CONTENT THAT STUDENTS MAY EDIT}

\vspace{0.15cm}

Students may edit or replace:

\begin{itemize}[leftmargin=0.8cm, itemsep=2pt]
    \item student name;
    \item class;
    \item internship organization;
    \item supervisor name;
    \item Abstract;
    \item List of Symbols and Abbreviations;
    \item chapter titles and chapter content;
    \item figures and photographs;
    \item tables;
    \item equations;
    \item citations and references.
\end{itemize}

Students shall not modify the main document formatting settings in the
preamble, including the document class, page geometry, line spacing,
heading formatting, table-of-contents formatting, font configuration,
page layout, or cover-page macros unless specifically instructed.

\vspace{0.3cm}

% ================================================================
% XXIII. FINAL SUBMISSION REQUIREMENTS
% ================================================================

\noindent
\textbf{XXIII.\quad FINAL SUBMISSION REQUIREMENTS}

\vspace{0.15cm}

Before submission, the final report shall:

\begin{itemize}[leftmargin=0.8cm, itemsep=2pt]
    \item compile without errors that affect the report content or presentation;
    \item contain all required forms and report sections;
    \item contain correct student, class, internship organization, and supervisor information;
    \item contain an updated Table of Contents, List of Figures, and List of Tables;
    \item contain no unused example content from the template;
    \item contain no placeholder text such as \textit{Title of Chapter 1},
    \textit{Section 1.1}, or other sample entries;
    \item contain all required citations and references;
    \item contain no unresolved cross-references.
\end{itemize}

If the Faculty requires submission of the \LaTeX{} source, the complete
project shall include all \verb|.tex| files, the bibliography database,
figures, and any other files required to reproduce the submitted PDF.

\vspace{0.4cm}

\noindent
\textbf{Important:}
Before submitting the final internship report, students shall remove or
comment out the following command in the main document:

\begin{verbatim}
\input{Regulations.tex}
\end{verbatim}

For example:

\begin{verbatim}
% \input{Regulations.tex}
\end{verbatim}

This removes the guidance pages from the final submitted report while
preserving the regulations file in the \LaTeX{} project for reference.

} % End 13pt content

\clearpage
 command before submitting
% the final internship report.
%
% =================================================================

\clearpage
\thispagestyle{empty}

\begin{center}
    {\fontsize{16pt}{20pt}\selectfont\bfseries
    SPECIFIC REGULATIONS FOR PREPARING\\
    GRADUATION INTERNSHIP REPORTS USING \LaTeX}
\end{center}

\vspace{0.5cm}

{\fontsize{13pt}{17pt}\selectfont
\setstretch{1.3}
\justifying

\noindent
\textbf{Note:}
This section is provided for guidance only and is included in the main
document using the command \verb|
% =================================================================
% REGULATIONS FOR PREPARING THE GRADUATION INTERNSHIP REPORT
% USING LATEX
% =================================================================
%
% This file is intended to be included from the main document with:
%
%     \input{Regulations.tex}
%
% IMPORTANT:
% This section is provided for guidance only. Students shall remove
% or comment out the \input{Regulations.tex} command before submitting
% the final internship report.
%
% =================================================================

\clearpage
\thispagestyle{empty}

\begin{center}
    {\fontsize{16pt}{20pt}\selectfont\bfseries
    SPECIFIC REGULATIONS FOR PREPARING\\
    GRADUATION INTERNSHIP REPORTS USING \LaTeX}
\end{center}

\vspace{0.5cm}

{\fontsize{13pt}{17pt}\selectfont
\setstretch{1.3}
\justifying

\noindent
\textbf{Note:}
This section is provided for guidance only and is included in the main
document using the command \verb|\input{Regulations.tex}|.
Students shall remove or comment out this command before submitting the
final internship report.

\vspace{0.4cm}

% ================================================================
% I. GENERAL REQUIREMENTS
% ================================================================

\noindent
\textbf{I.\quad GENERAL REQUIREMENTS}

\vspace{0.15cm}

The Graduation Internship Report shall be prepared using the official
\LaTeX{} template provided by the Faculty. Students shall preserve the
predefined page layout, typography, line spacing, heading styles,
numbering system, figure and table formatting, and bibliography settings
unless otherwise instructed by the Faculty or the supervisor.

The report shall be written in English.

Students are responsible for the academic and technical content of the
report. The official \LaTeX{} template is responsible for standardized
document formatting.

\vspace{0.3cm}

% ================================================================
% II. REPORT STRUCTURE
% ================================================================

\noindent
\textbf{II.\quad REPORT STRUCTURE}

\vspace{0.15cm}

The report shall normally be arranged in the following order:

\begin{enumerate}[label=\arabic*., leftmargin=1.2cm, itemsep=2pt]
    \item Outer Cover;
    \item Comments of the Internship Host Organization;
    \item Supervisor's Evaluation and Graduation Internship Defense Results;
    \item Table of Contents;
    \item List of Symbols and Abbreviations, if applicable;
    \item List of Figures;
    \item List of Tables;
    \item Abstract;
    \item Main Chapters of the Internship Report;
    \item References;
    \item Appendices, if applicable.
\end{enumerate}

The preliminary pages shall not be numbered as chapters. The main report
content shall be organized into consecutively numbered chapters beginning
with Chapter 1.

\vspace{0.3cm}

% ================================================================
% III. PAGE SIZE AND MARGINS
% ================================================================

\noindent
\textbf{III.\quad PAGE SIZE AND MARGINS}

\vspace{0.15cm}

The report shall be formatted on A4 paper with the following margins:

\begin{center}
\begin{tabular}{ll}
    Left margin:   & 3.5 cm \\
    Right margin:  & 2.0 cm \\
    Top margin:    & 2.5 cm \\
    Bottom margin: & 2.5 cm
\end{tabular}
\end{center}

These settings are predefined in the official template and shall not be
modified by students.

\vspace{0.3cm}

% ================================================================
% IV. FONT, FONT SIZE, AND LINE SPACING
% ================================================================

\noindent
\textbf{IV.\quad FONT, FONT SIZE, AND LINE SPACING}

\vspace{0.15cm}

The main text shall use a 13 pt font size with 1.3-line spacing.
A Times-family typeface shall be used in accordance with the official
template.

The text shall be fully justified.

Students shall not use repeated spaces, repeated manual line breaks, or
multiple blank lines to position text. Paragraph and heading spacing shall
be controlled by the template.

\vspace{0.3cm}

% ================================================================
% V. OUTER COVER
% ================================================================

\noindent
\textbf{V.\quad OUTER COVER}

\vspace{0.15cm}

The outer cover shall follow the official format provided in the template
and shall contain the following information:

\begin{itemize}[leftmargin=0.8cm, itemsep=2pt]
    \item Thai Nguyen University of Technology;
    \item Faculty of International Training;
    \item Engineering and Technology Department;
    \item University logo and Faculty logo;
    \item the title \textbf{INTERNSHIP REPORT};
    \item Internship Organization;
    \item Author;
    \item Class;
    \item Supervisor;
    \item place and year of completion.
\end{itemize}

Students shall only replace the required information fields and shall not
modify the cover-page layout unless instructed otherwise.

\vspace{0.3cm}

% ================================================================
% VI. COMMENTS OF THE INTERNSHIP HOST ORGANIZATION
% ================================================================

\noindent
\textbf{VI.\quad COMMENTS OF THE INTERNSHIP HOST ORGANIZATION}

\vspace{0.15cm}

The page entitled \textbf{COMMENTS OF THE INTERNSHIP HOST ORGANIZATION}
is reserved for the internship host organization.

It includes:

\begin{itemize}[leftmargin=0.8cm, itemsep=2pt]
    \item the student's full name;
    \item class;
    \item comments from the host organization;
    \item score;
    \item date;
    \item signature and official seal.
\end{itemize}

Students shall not alter or complete official evaluation fields unless
specifically authorized.

\vspace{0.3cm}

% ================================================================
% VII. SUPERVISOR'S EVALUATION AND DEFENSE RESULTS
% ================================================================

\noindent
\textbf{VII.\quad SUPERVISOR'S EVALUATION AND DEFENSE RESULTS}

\vspace{0.15cm}

The section entitled
\textbf{SUPERVISOR'S EVALUATION AND GRADUATION INTERNSHIP DEFENSE RESULTS}
is reserved for the supervisor and defense committee.

The evaluation may include:

\begin{itemize}[leftmargin=0.8cm, itemsep=2pt]
    \item Internship Progress and Student's Attitude;
    \item Report Content;
    \item Report Presentation;
    \item Other Comments;
    \item Supervisor's Opinion;
    \item Defense Result;
    \item Score.
\end{itemize}

Students shall not enter evaluation scores or defense results before they
are officially determined.

\vspace{0.3cm}

% ================================================================
% VIII. CHAPTERS, SECTIONS, AND SUBSECTIONS
% ================================================================

\noindent
\textbf{VIII.\quad CHAPTERS, SECTIONS, AND SUBSECTIONS}

\vspace{0.15cm}

The report structure shall be generated using the standard \LaTeX{}
commands:

\begin{verbatim}
\chapter{...}
\section{...}
\subsection{...}
\end{verbatim}

Chapter, section, and subsection numbers shall be generated automatically
by \LaTeX{}.

Students shall not manually type chapter or section numbers and shall not
reproduce heading styles using commands such as \verb|\textbf|,
\verb|\fontsize|, or \verb|\centering| instead of the appropriate
structural command.

Excessive subdivision should be avoided. In most cases, the hierarchy
Chapter--Section--Subsection is sufficient.

\vspace{0.3cm}

% ================================================================
% IX. ORGANIZATION OF THE MAIN CONTENT
% ================================================================

\noindent
\textbf{IX.\quad ORGANIZATION OF THE MAIN CONTENT}

\vspace{0.15cm}

The specific content of the report shall depend on the internship
organization and the assigned tasks. A typical structure may include:

\begin{itemize}[leftmargin=0.8cm, itemsep=2pt]
    \item \textbf{Chapter 1. Overview of the Internship Organization};
    \item \textbf{Chapter 2. Internship Activities and Assigned Tasks};
    \item \textbf{Chapter 3. Results, Experience, and Evaluation}.
\end{itemize}

These chapter titles are recommendations only and may be adapted to the
actual internship content.

The report shall reflect work actually performed, observed, analyzed, or
learned during the internship. Unrelated theoretical material should not
be included merely to increase report length.

\vspace{0.3cm}

% ================================================================
% X. TABLE OF CONTENTS
% ================================================================

\noindent
\textbf{X.\quad TABLE OF CONTENTS}

\vspace{0.15cm}

The Table of Contents shall be generated automatically using:

\begin{verbatim}
\tableofcontents
\end{verbatim}

Students shall not create or manually edit page numbers in the Table of
Contents.

The document shall be compiled sufficiently to ensure that all page
numbers and entries are correctly updated.

\vspace{0.3cm}

% ================================================================
% XI. LIST OF SYMBOLS AND ABBREVIATIONS
% ================================================================

\noindent
\textbf{XI.\quad LIST OF SYMBOLS AND ABBREVIATIONS}

\vspace{0.15cm}

Symbols and abbreviations that are used repeatedly in the report shall be
listed in the \textbf{List of Symbols and Abbreviations}.

The list may contain the following columns:

\begin{center}
\begin{tabular}{llll}
    No. & Symbol / Abbreviation & Unit & Definition
\end{tabular}
\end{center}

Only symbols and abbreviations actually used in the final report shall be
retained. Example entries contained in the template shall be removed if
they are not relevant to the student's report.

Each abbreviation should normally be defined when it first appears in
the main text.

\vspace{0.3cm}

% ================================================================
% XII. FIGURES, DIAGRAMS, AND PHOTOGRAPHS
% ================================================================

\noindent
\textbf{XII.\quad FIGURES, DIAGRAMS, AND PHOTOGRAPHS}

\vspace{0.15cm}

Figures shall be clear, legible, and directly relevant to the report
content.

Figures shall be inserted using the standard \verb|figure| environment.

Example:

\begin{verbatim}
\begin{figure}[H]
    \centering
    \includegraphics[width=0.8\textwidth]
                    {figures/company_structure.pdf}
    \caption{Organizational structure of the company}
    \label{fig:company-structure}
\end{figure}
\end{verbatim}

Figure captions shall be placed below the figure.

Each referenced figure shall have a unique \verb|\label| and shall be
referenced in the text using, for example:

\begin{verbatim}
Figure~\ref{fig:company-structure}
\end{verbatim}

Students shall not manually type figure numbers.

Figures reproduced or adapted from external sources shall include an
appropriate citation.

Photographs taken during the internship shall be relevant and sufficiently
clear. Students shall not include confidential, restricted, or proprietary
information without permission from the internship organization.

Vector graphics or high-resolution images are recommended whenever
possible.

\vspace{0.3cm}

% ================================================================
% XIII. TABLES
% ================================================================

\noindent
\textbf{XIII.\quad TABLES}

\vspace{0.15cm}

Tables shall be created using suitable \LaTeX{} environments such as
\verb|table|, \verb|tabular|, \verb|tabularx|, or \verb|longtable|.

Table captions shall be placed above the table.

Example:

\begin{verbatim}
\begin{table}[H]
    \centering
    \caption{Internship schedule}
    \label{tab:internship-schedule}

    \begin{tabular}{lll}
        \toprule
        Week & Task & Result \\
        \midrule
        1 & Orientation & Completed \\
        2 & System study & Completed \\
        \bottomrule
    \end{tabular}
\end{table}
\end{verbatim}

Tables shall be referenced using:

\begin{verbatim}
Table~\ref{tab:internship-schedule}
\end{verbatim}

Students shall not manually type table numbers.

Tables should be created directly in \LaTeX{} rather than inserted as
screenshots whenever practical.

\vspace{0.3cm}

% ================================================================
% XIV. EQUATIONS AND TECHNICAL QUANTITIES
% ================================================================

\noindent
\textbf{XIV.\quad EQUATIONS AND TECHNICAL QUANTITIES}

\vspace{0.15cm}

Mathematical expressions shall be typeset directly using \LaTeX{}.
Images or screenshots of equations shall not be used.

Short expressions may be written in inline mathematical mode, for example:

\begin{verbatim}
$R = 0.7\,\Omega$
\end{verbatim}

Important equations shall use a numbered equation environment.

Example:

\begin{verbatim}
\begin{equation}
    P = UI
    \label{eq:power}
\end{equation}
\end{verbatim}

Equations shall be referenced using:

\begin{verbatim}
Equation~\eqref{eq:power}
\end{verbatim}

Equation numbers shall not be manually entered.

SI units shall be used wherever applicable. The \verb|siunitx| package
should be used where appropriate.

Examples:

\begin{verbatim}
\SI{220}{\volt}
\SI{50}{\hertz}
\SI{2.5}{\ampere}
\end{verbatim}

Notation and units shall be consistent throughout the report.

\vspace{0.3cm}

% ================================================================
% XV. SOURCE CODE
% ================================================================

\noindent
\textbf{XV.\quad SOURCE CODE}

\vspace{0.15cm}

If the internship involves programming, PLCs, microcontrollers,
simulation, or software development, relevant source code may be included
in the report.

Source code shall be presented using the \verb|listings| environment or
another environment defined by the official template.

Example:

\begin{verbatim}
\begin{lstlisting}[language=st]
IF Start = TRUE THEN
    Motor := TRUE;
ELSE
    Motor := FALSE;
END_IF;
\end{lstlisting}
\end{verbatim}

Only code necessary for explaining the work should appear in the main
report. Long source-code listings should be moved to an appendix or
submitted as supplementary files if required.

Source code should not normally be inserted as screenshots.

\vspace{0.3cm}

% ================================================================
% XVI. ABSTRACT
% ================================================================

\noindent
\textbf{XVI.\quad ABSTRACT}

\vspace{0.15cm}

The Abstract shall briefly summarize:

\begin{itemize}[leftmargin=0.8cm, itemsep=2pt]
    \item the internship organization;
    \item the objectives of the internship;
    \item the main tasks performed;
    \item the main results, knowledge, and skills obtained.
\end{itemize}

The Abstract should be concise and should not reproduce the full content
of the report chapter by chapter.

\vspace{0.3cm}

% ================================================================
% XVII. REFERENCES AND CITATIONS
% ================================================================

\noindent
\textbf{XVII.\quad REFERENCES AND CITATIONS}

\vspace{0.15cm}

All books, journal papers, conference papers, standards, websites,
datasheets, manuals, technical documents, figures, tables, or other
external sources used in preparing the report shall be properly cited.

References shall be managed through the bibliography system provided with
the official template.

Bibliographic entries shall normally be stored in:

\begin{verbatim}
References.bib
\end{verbatim}

References shall be cited in the report using:

\begin{verbatim}
\cite{reference-key}
\end{verbatim}

Students shall not manually type reference numbers such as ``[1]'',
``[2]'', etc.

Titles of foreign-language publications shall remain in their original
language in the bibliography.

All borrowed ideas, data, figures, tables, equations, methods, or results
shall be properly acknowledged.

\vspace{0.3cm}

% ================================================================
% XVIII. CROSS-REFERENCING
% ================================================================

\noindent
\textbf{XVIII.\quad CROSS-REFERENCING}

\vspace{0.15cm}

Numbered elements shall be cross-referenced using the standard \LaTeX{}
mechanism.

Use the following commands where appropriate:

\begin{verbatim}
\label{...}
\ref{...}
\eqref{...}
\end{verbatim}

For example:

\begin{verbatim}
As shown in Figure~\ref{fig:system}, ...
\end{verbatim}

Students shall not manually type numbers such as ``Figure 2.3'' when
those numbers can be generated automatically.

\vspace{0.3cm}

% ================================================================
% XIX. PAGE NUMBERING
% ================================================================

\noindent
\textbf{XIX.\quad PAGE NUMBERING}

\vspace{0.15cm}

Cover pages and official evaluation forms at the beginning of the report
shall not display page numbers.

The main body of the report shall use Arabic page numbers beginning from
Chapter 1.

Page numbering shall be generated automatically by the official template.
Students shall not manually enter page numbers.

\vspace{0.3cm}

% ================================================================
% XX. PROJECT FILE ORGANIZATION
% ================================================================

\noindent
\textbf{XX.\quad PROJECT FILE ORGANIZATION}

\vspace{0.15cm}

Students should maintain a clear and organized \LaTeX{} project
structure. A recommended structure is:

\begin{verbatim}
InternshipReport/
|
|-- main.tex
|-- References.bib
|
|-- figures/
|   |-- company.pdf
|   |-- system.png
|   `-- ...
|
`-- chapters/
    |-- chapter1.tex
    |-- chapter2.tex
    `-- chapter3.tex
\end{verbatim}

Figures should normally be stored in the \verb|figures| directory
specified by the official template.

File and directory names should preferably contain Latin letters,
numbers, hyphens, or underscores.

Recommended examples include:

\begin{verbatim}
company_structure.pdf
motor_control.png
chapter_01.tex
\end{verbatim}

Spaces and unnecessary special characters in file names should be
avoided where possible.

\vspace{0.3cm}

% ================================================================
% XXI. DOCUMENT COMPILATION AND FINAL CHECK
% ================================================================

\noindent
\textbf{XXI.\quad DOCUMENT COMPILATION AND FINAL CHECK}

\vspace{0.15cm}

Before submission, the complete \LaTeX{} project shall compile
successfully and produce the final PDF report.

The final document shall not contain unresolved references such as:

\begin{verbatim}
??
[?]
\end{verbatim}

Before generating the final PDF, students shall check:

\begin{itemize}[leftmargin=0.8cm, itemsep=2pt]
    \item the Table of Contents;
    \item page numbering;
    \item chapter, section, and subsection numbering;
    \item figure numbering;
    \item table numbering;
    \item equation numbering;
    \item cross-references;
    \item citations;
    \item bibliography;
    \item font display;
    \item mathematical symbols;
    \item image quality;
    \item completeness of all required forms.
\end{itemize}

The document may be compiled using TeX Live, MiKTeX, or another
compatible \LaTeX{} environment. Regardless of the compilation
environment, the resulting PDF shall conform to the official Faculty
template.

\vspace{0.3cm}

% ================================================================
% XXII. CONTENT THAT STUDENTS MAY EDIT
% ================================================================

\noindent
\textbf{XXII.\quad CONTENT THAT STUDENTS MAY EDIT}

\vspace{0.15cm}

Students may edit or replace:

\begin{itemize}[leftmargin=0.8cm, itemsep=2pt]
    \item student name;
    \item class;
    \item internship organization;
    \item supervisor name;
    \item Abstract;
    \item List of Symbols and Abbreviations;
    \item chapter titles and chapter content;
    \item figures and photographs;
    \item tables;
    \item equations;
    \item citations and references.
\end{itemize}

Students shall not modify the main document formatting settings in the
preamble, including the document class, page geometry, line spacing,
heading formatting, table-of-contents formatting, font configuration,
page layout, or cover-page macros unless specifically instructed.

\vspace{0.3cm}

% ================================================================
% XXIII. FINAL SUBMISSION REQUIREMENTS
% ================================================================

\noindent
\textbf{XXIII.\quad FINAL SUBMISSION REQUIREMENTS}

\vspace{0.15cm}

Before submission, the final report shall:

\begin{itemize}[leftmargin=0.8cm, itemsep=2pt]
    \item compile without errors that affect the report content or presentation;
    \item contain all required forms and report sections;
    \item contain correct student, class, internship organization, and supervisor information;
    \item contain an updated Table of Contents, List of Figures, and List of Tables;
    \item contain no unused example content from the template;
    \item contain no placeholder text such as \textit{Title of Chapter 1},
    \textit{Section 1.1}, or other sample entries;
    \item contain all required citations and references;
    \item contain no unresolved cross-references.
\end{itemize}

If the Faculty requires submission of the \LaTeX{} source, the complete
project shall include all \verb|.tex| files, the bibliography database,
figures, and any other files required to reproduce the submitted PDF.

\vspace{0.4cm}

\noindent
\textbf{Important:}
Before submitting the final internship report, students shall remove or
comment out the following command in the main document:

\begin{verbatim}
\input{Regulations.tex}
\end{verbatim}

For example:

\begin{verbatim}
% \input{Regulations.tex}
\end{verbatim}

This removes the guidance pages from the final submitted report while
preserving the regulations file in the \LaTeX{} project for reference.

} % End 13pt content

\clearpage
|.
Students shall remove or comment out this command before submitting the
final internship report.

\vspace{0.4cm}

% ================================================================
% I. GENERAL REQUIREMENTS
% ================================================================

\noindent
\textbf{I.\quad GENERAL REQUIREMENTS}

\vspace{0.15cm}

The Graduation Internship Report shall be prepared using the official
\LaTeX{} template provided by the Faculty. Students shall preserve the
predefined page layout, typography, line spacing, heading styles,
numbering system, figure and table formatting, and bibliography settings
unless otherwise instructed by the Faculty or the supervisor.

The report shall be written in English.

Students are responsible for the academic and technical content of the
report. The official \LaTeX{} template is responsible for standardized
document formatting.

\vspace{0.3cm}

% ================================================================
% II. REPORT STRUCTURE
% ================================================================

\noindent
\textbf{II.\quad REPORT STRUCTURE}

\vspace{0.15cm}

The report shall normally be arranged in the following order:

\begin{enumerate}[label=\arabic*., leftmargin=1.2cm, itemsep=2pt]
    \item Outer Cover;
    \item Comments of the Internship Host Organization;
    \item Supervisor's Evaluation and Graduation Internship Defense Results;
    \item Table of Contents;
    \item List of Symbols and Abbreviations, if applicable;
    \item List of Figures;
    \item List of Tables;
    \item Abstract;
    \item Main Chapters of the Internship Report;
    \item References;
    \item Appendices, if applicable.
\end{enumerate}

The preliminary pages shall not be numbered as chapters. The main report
content shall be organized into consecutively numbered chapters beginning
with Chapter 1.

\vspace{0.3cm}

% ================================================================
% III. PAGE SIZE AND MARGINS
% ================================================================

\noindent
\textbf{III.\quad PAGE SIZE AND MARGINS}

\vspace{0.15cm}

The report shall be formatted on A4 paper with the following margins:

\begin{center}
\begin{tabular}{ll}
    Left margin:   & 3.5 cm \\
    Right margin:  & 2.0 cm \\
    Top margin:    & 2.5 cm \\
    Bottom margin: & 2.5 cm
\end{tabular}
\end{center}

These settings are predefined in the official template and shall not be
modified by students.

\vspace{0.3cm}

% ================================================================
% IV. FONT, FONT SIZE, AND LINE SPACING
% ================================================================

\noindent
\textbf{IV.\quad FONT, FONT SIZE, AND LINE SPACING}

\vspace{0.15cm}

The main text shall use a 13 pt font size with 1.3-line spacing.
A Times-family typeface shall be used in accordance with the official
template.

The text shall be fully justified.

Students shall not use repeated spaces, repeated manual line breaks, or
multiple blank lines to position text. Paragraph and heading spacing shall
be controlled by the template.

\vspace{0.3cm}

% ================================================================
% V. OUTER COVER
% ================================================================

\noindent
\textbf{V.\quad OUTER COVER}

\vspace{0.15cm}

The outer cover shall follow the official format provided in the template
and shall contain the following information:

\begin{itemize}[leftmargin=0.8cm, itemsep=2pt]
    \item Thai Nguyen University of Technology;
    \item Faculty of International Training;
    \item Engineering and Technology Department;
    \item University logo and Faculty logo;
    \item the title \textbf{INTERNSHIP REPORT};
    \item Internship Organization;
    \item Author;
    \item Class;
    \item Supervisor;
    \item place and year of completion.
\end{itemize}

Students shall only replace the required information fields and shall not
modify the cover-page layout unless instructed otherwise.

\vspace{0.3cm}

% ================================================================
% VI. COMMENTS OF THE INTERNSHIP HOST ORGANIZATION
% ================================================================

\noindent
\textbf{VI.\quad COMMENTS OF THE INTERNSHIP HOST ORGANIZATION}

\vspace{0.15cm}

The page entitled \textbf{COMMENTS OF THE INTERNSHIP HOST ORGANIZATION}
is reserved for the internship host organization.

It includes:

\begin{itemize}[leftmargin=0.8cm, itemsep=2pt]
    \item the student's full name;
    \item class;
    \item comments from the host organization;
    \item score;
    \item date;
    \item signature and official seal.
\end{itemize}

Students shall not alter or complete official evaluation fields unless
specifically authorized.

\vspace{0.3cm}

% ================================================================
% VII. SUPERVISOR'S EVALUATION AND DEFENSE RESULTS
% ================================================================

\noindent
\textbf{VII.\quad SUPERVISOR'S EVALUATION AND DEFENSE RESULTS}

\vspace{0.15cm}

The section entitled
\textbf{SUPERVISOR'S EVALUATION AND GRADUATION INTERNSHIP DEFENSE RESULTS}
is reserved for the supervisor and defense committee.

The evaluation may include:

\begin{itemize}[leftmargin=0.8cm, itemsep=2pt]
    \item Internship Progress and Student's Attitude;
    \item Report Content;
    \item Report Presentation;
    \item Other Comments;
    \item Supervisor's Opinion;
    \item Defense Result;
    \item Score.
\end{itemize}

Students shall not enter evaluation scores or defense results before they
are officially determined.

\vspace{0.3cm}

% ================================================================
% VIII. CHAPTERS, SECTIONS, AND SUBSECTIONS
% ================================================================

\noindent
\textbf{VIII.\quad CHAPTERS, SECTIONS, AND SUBSECTIONS}

\vspace{0.15cm}

The report structure shall be generated using the standard \LaTeX{}
commands:

\begin{verbatim}
\chapter{...}
\section{...}
\subsection{...}
\end{verbatim}

Chapter, section, and subsection numbers shall be generated automatically
by \LaTeX{}.

Students shall not manually type chapter or section numbers and shall not
reproduce heading styles using commands such as \verb|\textbf|,
\verb|\fontsize|, or \verb|\centering| instead of the appropriate
structural command.

Excessive subdivision should be avoided. In most cases, the hierarchy
Chapter--Section--Subsection is sufficient.

\vspace{0.3cm}

% ================================================================
% IX. ORGANIZATION OF THE MAIN CONTENT
% ================================================================

\noindent
\textbf{IX.\quad ORGANIZATION OF THE MAIN CONTENT}

\vspace{0.15cm}

The specific content of the report shall depend on the internship
organization and the assigned tasks. A typical structure may include:

\begin{itemize}[leftmargin=0.8cm, itemsep=2pt]
    \item \textbf{Chapter 1. Overview of the Internship Organization};
    \item \textbf{Chapter 2. Internship Activities and Assigned Tasks};
    \item \textbf{Chapter 3. Results, Experience, and Evaluation}.
\end{itemize}

These chapter titles are recommendations only and may be adapted to the
actual internship content.

The report shall reflect work actually performed, observed, analyzed, or
learned during the internship. Unrelated theoretical material should not
be included merely to increase report length.

\vspace{0.3cm}

% ================================================================
% X. TABLE OF CONTENTS
% ================================================================

\noindent
\textbf{X.\quad TABLE OF CONTENTS}

\vspace{0.15cm}

The Table of Contents shall be generated automatically using:

\begin{verbatim}
\tableofcontents
\end{verbatim}

Students shall not create or manually edit page numbers in the Table of
Contents.

The document shall be compiled sufficiently to ensure that all page
numbers and entries are correctly updated.

\vspace{0.3cm}

% ================================================================
% XI. LIST OF SYMBOLS AND ABBREVIATIONS
% ================================================================

\noindent
\textbf{XI.\quad LIST OF SYMBOLS AND ABBREVIATIONS}

\vspace{0.15cm}

Symbols and abbreviations that are used repeatedly in the report shall be
listed in the \textbf{List of Symbols and Abbreviations}.

The list may contain the following columns:

\begin{center}
\begin{tabular}{llll}
    No. & Symbol / Abbreviation & Unit & Definition
\end{tabular}
\end{center}

Only symbols and abbreviations actually used in the final report shall be
retained. Example entries contained in the template shall be removed if
they are not relevant to the student's report.

Each abbreviation should normally be defined when it first appears in
the main text.

\vspace{0.3cm}

% ================================================================
% XII. FIGURES, DIAGRAMS, AND PHOTOGRAPHS
% ================================================================

\noindent
\textbf{XII.\quad FIGURES, DIAGRAMS, AND PHOTOGRAPHS}

\vspace{0.15cm}

Figures shall be clear, legible, and directly relevant to the report
content.

Figures shall be inserted using the standard \verb|figure| environment.

Example:

\begin{verbatim}
\begin{figure}[H]
    \centering
    \includegraphics[width=0.8\textwidth]
                    {figures/company_structure.pdf}
    \caption{Organizational structure of the company}
    \label{fig:company-structure}
\end{figure}
\end{verbatim}

Figure captions shall be placed below the figure.

Each referenced figure shall have a unique \verb|\label| and shall be
referenced in the text using, for example:

\begin{verbatim}
Figure~\ref{fig:company-structure}
\end{verbatim}

Students shall not manually type figure numbers.

Figures reproduced or adapted from external sources shall include an
appropriate citation.

Photographs taken during the internship shall be relevant and sufficiently
clear. Students shall not include confidential, restricted, or proprietary
information without permission from the internship organization.

Vector graphics or high-resolution images are recommended whenever
possible.

\vspace{0.3cm}

% ================================================================
% XIII. TABLES
% ================================================================

\noindent
\textbf{XIII.\quad TABLES}

\vspace{0.15cm}

Tables shall be created using suitable \LaTeX{} environments such as
\verb|table|, \verb|tabular|, \verb|tabularx|, or \verb|longtable|.

Table captions shall be placed above the table.

Example:

\begin{verbatim}
\begin{table}[H]
    \centering
    \caption{Internship schedule}
    \label{tab:internship-schedule}

    \begin{tabular}{lll}
        \toprule
        Week & Task & Result \\
        \midrule
        1 & Orientation & Completed \\
        2 & System study & Completed \\
        \bottomrule
    \end{tabular}
\end{table}
\end{verbatim}

Tables shall be referenced using:

\begin{verbatim}
Table~\ref{tab:internship-schedule}
\end{verbatim}

Students shall not manually type table numbers.

Tables should be created directly in \LaTeX{} rather than inserted as
screenshots whenever practical.

\vspace{0.3cm}

% ================================================================
% XIV. EQUATIONS AND TECHNICAL QUANTITIES
% ================================================================

\noindent
\textbf{XIV.\quad EQUATIONS AND TECHNICAL QUANTITIES}

\vspace{0.15cm}

Mathematical expressions shall be typeset directly using \LaTeX{}.
Images or screenshots of equations shall not be used.

Short expressions may be written in inline mathematical mode, for example:

\begin{verbatim}
$R = 0.7\,\Omega$
\end{verbatim}

Important equations shall use a numbered equation environment.

Example:

\begin{verbatim}
\begin{equation}
    P = UI
    \label{eq:power}
\end{equation}
\end{verbatim}

Equations shall be referenced using:

\begin{verbatim}
Equation~\eqref{eq:power}
\end{verbatim}

Equation numbers shall not be manually entered.

SI units shall be used wherever applicable. The \verb|siunitx| package
should be used where appropriate.

Examples:

\begin{verbatim}
\SI{220}{\volt}
\SI{50}{\hertz}
\SI{2.5}{\ampere}
\end{verbatim}

Notation and units shall be consistent throughout the report.

\vspace{0.3cm}

% ================================================================
% XV. SOURCE CODE
% ================================================================

\noindent
\textbf{XV.\quad SOURCE CODE}

\vspace{0.15cm}

If the internship involves programming, PLCs, microcontrollers,
simulation, or software development, relevant source code may be included
in the report.

Source code shall be presented using the \verb|listings| environment or
another environment defined by the official template.

Example:

\begin{verbatim}
\begin{lstlisting}[language=st]
IF Start = TRUE THEN
    Motor := TRUE;
ELSE
    Motor := FALSE;
END_IF;
\end{lstlisting}
\end{verbatim}

Only code necessary for explaining the work should appear in the main
report. Long source-code listings should be moved to an appendix or
submitted as supplementary files if required.

Source code should not normally be inserted as screenshots.

\vspace{0.3cm}

% ================================================================
% XVI. ABSTRACT
% ================================================================

\noindent
\textbf{XVI.\quad ABSTRACT}

\vspace{0.15cm}

The Abstract shall briefly summarize:

\begin{itemize}[leftmargin=0.8cm, itemsep=2pt]
    \item the internship organization;
    \item the objectives of the internship;
    \item the main tasks performed;
    \item the main results, knowledge, and skills obtained.
\end{itemize}

The Abstract should be concise and should not reproduce the full content
of the report chapter by chapter.

\vspace{0.3cm}

% ================================================================
% XVII. REFERENCES AND CITATIONS
% ================================================================

\noindent
\textbf{XVII.\quad REFERENCES AND CITATIONS}

\vspace{0.15cm}

All books, journal papers, conference papers, standards, websites,
datasheets, manuals, technical documents, figures, tables, or other
external sources used in preparing the report shall be properly cited.

References shall be managed through the bibliography system provided with
the official template.

Bibliographic entries shall normally be stored in:

\begin{verbatim}
References.bib
\end{verbatim}

References shall be cited in the report using:

\begin{verbatim}
\cite{reference-key}
\end{verbatim}

Students shall not manually type reference numbers such as ``[1]'',
``[2]'', etc.

Titles of foreign-language publications shall remain in their original
language in the bibliography.

All borrowed ideas, data, figures, tables, equations, methods, or results
shall be properly acknowledged.

\vspace{0.3cm}

% ================================================================
% XVIII. CROSS-REFERENCING
% ================================================================

\noindent
\textbf{XVIII.\quad CROSS-REFERENCING}

\vspace{0.15cm}

Numbered elements shall be cross-referenced using the standard \LaTeX{}
mechanism.

Use the following commands where appropriate:

\begin{verbatim}
\label{...}
\ref{...}
\eqref{...}
\end{verbatim}

For example:

\begin{verbatim}
As shown in Figure~\ref{fig:system}, ...
\end{verbatim}

Students shall not manually type numbers such as ``Figure 2.3'' when
those numbers can be generated automatically.

\vspace{0.3cm}

% ================================================================
% XIX. PAGE NUMBERING
% ================================================================

\noindent
\textbf{XIX.\quad PAGE NUMBERING}

\vspace{0.15cm}

Cover pages and official evaluation forms at the beginning of the report
shall not display page numbers.

The main body of the report shall use Arabic page numbers beginning from
Chapter 1.

Page numbering shall be generated automatically by the official template.
Students shall not manually enter page numbers.

\vspace{0.3cm}

% ================================================================
% XX. PROJECT FILE ORGANIZATION
% ================================================================

\noindent
\textbf{XX.\quad PROJECT FILE ORGANIZATION}

\vspace{0.15cm}

Students should maintain a clear and organized \LaTeX{} project
structure. A recommended structure is:

\begin{verbatim}
InternshipReport/
|
|-- main.tex
|-- References.bib
|
|-- figures/
|   |-- company.pdf
|   |-- system.png
|   `-- ...
|
`-- chapters/
    |-- chapter1.tex
    |-- chapter2.tex
    `-- chapter3.tex
\end{verbatim}

Figures should normally be stored in the \verb|figures| directory
specified by the official template.

File and directory names should preferably contain Latin letters,
numbers, hyphens, or underscores.

Recommended examples include:

\begin{verbatim}
company_structure.pdf
motor_control.png
chapter_01.tex
\end{verbatim}

Spaces and unnecessary special characters in file names should be
avoided where possible.

\vspace{0.3cm}

% ================================================================
% XXI. DOCUMENT COMPILATION AND FINAL CHECK
% ================================================================

\noindent
\textbf{XXI.\quad DOCUMENT COMPILATION AND FINAL CHECK}

\vspace{0.15cm}

Before submission, the complete \LaTeX{} project shall compile
successfully and produce the final PDF report.

The final document shall not contain unresolved references such as:

\begin{verbatim}
??
[?]
\end{verbatim}

Before generating the final PDF, students shall check:

\begin{itemize}[leftmargin=0.8cm, itemsep=2pt]
    \item the Table of Contents;
    \item page numbering;
    \item chapter, section, and subsection numbering;
    \item figure numbering;
    \item table numbering;
    \item equation numbering;
    \item cross-references;
    \item citations;
    \item bibliography;
    \item font display;
    \item mathematical symbols;
    \item image quality;
    \item completeness of all required forms.
\end{itemize}

The document may be compiled using TeX Live, MiKTeX, or another
compatible \LaTeX{} environment. Regardless of the compilation
environment, the resulting PDF shall conform to the official Faculty
template.

\vspace{0.3cm}

% ================================================================
% XXII. CONTENT THAT STUDENTS MAY EDIT
% ================================================================

\noindent
\textbf{XXII.\quad CONTENT THAT STUDENTS MAY EDIT}

\vspace{0.15cm}

Students may edit or replace:

\begin{itemize}[leftmargin=0.8cm, itemsep=2pt]
    \item student name;
    \item class;
    \item internship organization;
    \item supervisor name;
    \item Abstract;
    \item List of Symbols and Abbreviations;
    \item chapter titles and chapter content;
    \item figures and photographs;
    \item tables;
    \item equations;
    \item citations and references.
\end{itemize}

Students shall not modify the main document formatting settings in the
preamble, including the document class, page geometry, line spacing,
heading formatting, table-of-contents formatting, font configuration,
page layout, or cover-page macros unless specifically instructed.

\vspace{0.3cm}

% ================================================================
% XXIII. FINAL SUBMISSION REQUIREMENTS
% ================================================================

\noindent
\textbf{XXIII.\quad FINAL SUBMISSION REQUIREMENTS}

\vspace{0.15cm}

Before submission, the final report shall:

\begin{itemize}[leftmargin=0.8cm, itemsep=2pt]
    \item compile without errors that affect the report content or presentation;
    \item contain all required forms and report sections;
    \item contain correct student, class, internship organization, and supervisor information;
    \item contain an updated Table of Contents, List of Figures, and List of Tables;
    \item contain no unused example content from the template;
    \item contain no placeholder text such as \textit{Title of Chapter 1},
    \textit{Section 1.1}, or other sample entries;
    \item contain all required citations and references;
    \item contain no unresolved cross-references.
\end{itemize}

If the Faculty requires submission of the \LaTeX{} source, the complete
project shall include all \verb|.tex| files, the bibliography database,
figures, and any other files required to reproduce the submitted PDF.

\vspace{0.4cm}

\noindent
\textbf{Important:}
Before submitting the final internship report, students shall remove or
comment out the following command in the main document:

\begin{verbatim}

% =================================================================
% REGULATIONS FOR PREPARING THE GRADUATION INTERNSHIP REPORT
% USING LATEX
% =================================================================
%
% This file is intended to be included from the main document with:
%
%     \input{Regulations.tex}
%
% IMPORTANT:
% This section is provided for guidance only. Students shall remove
% or comment out the \input{Regulations.tex} command before submitting
% the final internship report.
%
% =================================================================

\clearpage
\thispagestyle{empty}

\begin{center}
    {\fontsize{16pt}{20pt}\selectfont\bfseries
    SPECIFIC REGULATIONS FOR PREPARING\\
    GRADUATION INTERNSHIP REPORTS USING \LaTeX}
\end{center}

\vspace{0.5cm}

{\fontsize{13pt}{17pt}\selectfont
\setstretch{1.3}
\justifying

\noindent
\textbf{Note:}
This section is provided for guidance only and is included in the main
document using the command \verb|\input{Regulations.tex}|.
Students shall remove or comment out this command before submitting the
final internship report.

\vspace{0.4cm}

% ================================================================
% I. GENERAL REQUIREMENTS
% ================================================================

\noindent
\textbf{I.\quad GENERAL REQUIREMENTS}

\vspace{0.15cm}

The Graduation Internship Report shall be prepared using the official
\LaTeX{} template provided by the Faculty. Students shall preserve the
predefined page layout, typography, line spacing, heading styles,
numbering system, figure and table formatting, and bibliography settings
unless otherwise instructed by the Faculty or the supervisor.

The report shall be written in English.

Students are responsible for the academic and technical content of the
report. The official \LaTeX{} template is responsible for standardized
document formatting.

\vspace{0.3cm}

% ================================================================
% II. REPORT STRUCTURE
% ================================================================

\noindent
\textbf{II.\quad REPORT STRUCTURE}

\vspace{0.15cm}

The report shall normally be arranged in the following order:

\begin{enumerate}[label=\arabic*., leftmargin=1.2cm, itemsep=2pt]
    \item Outer Cover;
    \item Comments of the Internship Host Organization;
    \item Supervisor's Evaluation and Graduation Internship Defense Results;
    \item Table of Contents;
    \item List of Symbols and Abbreviations, if applicable;
    \item List of Figures;
    \item List of Tables;
    \item Abstract;
    \item Main Chapters of the Internship Report;
    \item References;
    \item Appendices, if applicable.
\end{enumerate}

The preliminary pages shall not be numbered as chapters. The main report
content shall be organized into consecutively numbered chapters beginning
with Chapter 1.

\vspace{0.3cm}

% ================================================================
% III. PAGE SIZE AND MARGINS
% ================================================================

\noindent
\textbf{III.\quad PAGE SIZE AND MARGINS}

\vspace{0.15cm}

The report shall be formatted on A4 paper with the following margins:

\begin{center}
\begin{tabular}{ll}
    Left margin:   & 3.5 cm \\
    Right margin:  & 2.0 cm \\
    Top margin:    & 2.5 cm \\
    Bottom margin: & 2.5 cm
\end{tabular}
\end{center}

These settings are predefined in the official template and shall not be
modified by students.

\vspace{0.3cm}

% ================================================================
% IV. FONT, FONT SIZE, AND LINE SPACING
% ================================================================

\noindent
\textbf{IV.\quad FONT, FONT SIZE, AND LINE SPACING}

\vspace{0.15cm}

The main text shall use a 13 pt font size with 1.3-line spacing.
A Times-family typeface shall be used in accordance with the official
template.

The text shall be fully justified.

Students shall not use repeated spaces, repeated manual line breaks, or
multiple blank lines to position text. Paragraph and heading spacing shall
be controlled by the template.

\vspace{0.3cm}

% ================================================================
% V. OUTER COVER
% ================================================================

\noindent
\textbf{V.\quad OUTER COVER}

\vspace{0.15cm}

The outer cover shall follow the official format provided in the template
and shall contain the following information:

\begin{itemize}[leftmargin=0.8cm, itemsep=2pt]
    \item Thai Nguyen University of Technology;
    \item Faculty of International Training;
    \item Engineering and Technology Department;
    \item University logo and Faculty logo;
    \item the title \textbf{INTERNSHIP REPORT};
    \item Internship Organization;
    \item Author;
    \item Class;
    \item Supervisor;
    \item place and year of completion.
\end{itemize}

Students shall only replace the required information fields and shall not
modify the cover-page layout unless instructed otherwise.

\vspace{0.3cm}

% ================================================================
% VI. COMMENTS OF THE INTERNSHIP HOST ORGANIZATION
% ================================================================

\noindent
\textbf{VI.\quad COMMENTS OF THE INTERNSHIP HOST ORGANIZATION}

\vspace{0.15cm}

The page entitled \textbf{COMMENTS OF THE INTERNSHIP HOST ORGANIZATION}
is reserved for the internship host organization.

It includes:

\begin{itemize}[leftmargin=0.8cm, itemsep=2pt]
    \item the student's full name;
    \item class;
    \item comments from the host organization;
    \item score;
    \item date;
    \item signature and official seal.
\end{itemize}

Students shall not alter or complete official evaluation fields unless
specifically authorized.

\vspace{0.3cm}

% ================================================================
% VII. SUPERVISOR'S EVALUATION AND DEFENSE RESULTS
% ================================================================

\noindent
\textbf{VII.\quad SUPERVISOR'S EVALUATION AND DEFENSE RESULTS}

\vspace{0.15cm}

The section entitled
\textbf{SUPERVISOR'S EVALUATION AND GRADUATION INTERNSHIP DEFENSE RESULTS}
is reserved for the supervisor and defense committee.

The evaluation may include:

\begin{itemize}[leftmargin=0.8cm, itemsep=2pt]
    \item Internship Progress and Student's Attitude;
    \item Report Content;
    \item Report Presentation;
    \item Other Comments;
    \item Supervisor's Opinion;
    \item Defense Result;
    \item Score.
\end{itemize}

Students shall not enter evaluation scores or defense results before they
are officially determined.

\vspace{0.3cm}

% ================================================================
% VIII. CHAPTERS, SECTIONS, AND SUBSECTIONS
% ================================================================

\noindent
\textbf{VIII.\quad CHAPTERS, SECTIONS, AND SUBSECTIONS}

\vspace{0.15cm}

The report structure shall be generated using the standard \LaTeX{}
commands:

\begin{verbatim}
\chapter{...}
\section{...}
\subsection{...}
\end{verbatim}

Chapter, section, and subsection numbers shall be generated automatically
by \LaTeX{}.

Students shall not manually type chapter or section numbers and shall not
reproduce heading styles using commands such as \verb|\textbf|,
\verb|\fontsize|, or \verb|\centering| instead of the appropriate
structural command.

Excessive subdivision should be avoided. In most cases, the hierarchy
Chapter--Section--Subsection is sufficient.

\vspace{0.3cm}

% ================================================================
% IX. ORGANIZATION OF THE MAIN CONTENT
% ================================================================

\noindent
\textbf{IX.\quad ORGANIZATION OF THE MAIN CONTENT}

\vspace{0.15cm}

The specific content of the report shall depend on the internship
organization and the assigned tasks. A typical structure may include:

\begin{itemize}[leftmargin=0.8cm, itemsep=2pt]
    \item \textbf{Chapter 1. Overview of the Internship Organization};
    \item \textbf{Chapter 2. Internship Activities and Assigned Tasks};
    \item \textbf{Chapter 3. Results, Experience, and Evaluation}.
\end{itemize}

These chapter titles are recommendations only and may be adapted to the
actual internship content.

The report shall reflect work actually performed, observed, analyzed, or
learned during the internship. Unrelated theoretical material should not
be included merely to increase report length.

\vspace{0.3cm}

% ================================================================
% X. TABLE OF CONTENTS
% ================================================================

\noindent
\textbf{X.\quad TABLE OF CONTENTS}

\vspace{0.15cm}

The Table of Contents shall be generated automatically using:

\begin{verbatim}
\tableofcontents
\end{verbatim}

Students shall not create or manually edit page numbers in the Table of
Contents.

The document shall be compiled sufficiently to ensure that all page
numbers and entries are correctly updated.

\vspace{0.3cm}

% ================================================================
% XI. LIST OF SYMBOLS AND ABBREVIATIONS
% ================================================================

\noindent
\textbf{XI.\quad LIST OF SYMBOLS AND ABBREVIATIONS}

\vspace{0.15cm}

Symbols and abbreviations that are used repeatedly in the report shall be
listed in the \textbf{List of Symbols and Abbreviations}.

The list may contain the following columns:

\begin{center}
\begin{tabular}{llll}
    No. & Symbol / Abbreviation & Unit & Definition
\end{tabular}
\end{center}

Only symbols and abbreviations actually used in the final report shall be
retained. Example entries contained in the template shall be removed if
they are not relevant to the student's report.

Each abbreviation should normally be defined when it first appears in
the main text.

\vspace{0.3cm}

% ================================================================
% XII. FIGURES, DIAGRAMS, AND PHOTOGRAPHS
% ================================================================

\noindent
\textbf{XII.\quad FIGURES, DIAGRAMS, AND PHOTOGRAPHS}

\vspace{0.15cm}

Figures shall be clear, legible, and directly relevant to the report
content.

Figures shall be inserted using the standard \verb|figure| environment.

Example:

\begin{verbatim}
\begin{figure}[H]
    \centering
    \includegraphics[width=0.8\textwidth]
                    {figures/company_structure.pdf}
    \caption{Organizational structure of the company}
    \label{fig:company-structure}
\end{figure}
\end{verbatim}

Figure captions shall be placed below the figure.

Each referenced figure shall have a unique \verb|\label| and shall be
referenced in the text using, for example:

\begin{verbatim}
Figure~\ref{fig:company-structure}
\end{verbatim}

Students shall not manually type figure numbers.

Figures reproduced or adapted from external sources shall include an
appropriate citation.

Photographs taken during the internship shall be relevant and sufficiently
clear. Students shall not include confidential, restricted, or proprietary
information without permission from the internship organization.

Vector graphics or high-resolution images are recommended whenever
possible.

\vspace{0.3cm}

% ================================================================
% XIII. TABLES
% ================================================================

\noindent
\textbf{XIII.\quad TABLES}

\vspace{0.15cm}

Tables shall be created using suitable \LaTeX{} environments such as
\verb|table|, \verb|tabular|, \verb|tabularx|, or \verb|longtable|.

Table captions shall be placed above the table.

Example:

\begin{verbatim}
\begin{table}[H]
    \centering
    \caption{Internship schedule}
    \label{tab:internship-schedule}

    \begin{tabular}{lll}
        \toprule
        Week & Task & Result \\
        \midrule
        1 & Orientation & Completed \\
        2 & System study & Completed \\
        \bottomrule
    \end{tabular}
\end{table}
\end{verbatim}

Tables shall be referenced using:

\begin{verbatim}
Table~\ref{tab:internship-schedule}
\end{verbatim}

Students shall not manually type table numbers.

Tables should be created directly in \LaTeX{} rather than inserted as
screenshots whenever practical.

\vspace{0.3cm}

% ================================================================
% XIV. EQUATIONS AND TECHNICAL QUANTITIES
% ================================================================

\noindent
\textbf{XIV.\quad EQUATIONS AND TECHNICAL QUANTITIES}

\vspace{0.15cm}

Mathematical expressions shall be typeset directly using \LaTeX{}.
Images or screenshots of equations shall not be used.

Short expressions may be written in inline mathematical mode, for example:

\begin{verbatim}
$R = 0.7\,\Omega$
\end{verbatim}

Important equations shall use a numbered equation environment.

Example:

\begin{verbatim}
\begin{equation}
    P = UI
    \label{eq:power}
\end{equation}
\end{verbatim}

Equations shall be referenced using:

\begin{verbatim}
Equation~\eqref{eq:power}
\end{verbatim}

Equation numbers shall not be manually entered.

SI units shall be used wherever applicable. The \verb|siunitx| package
should be used where appropriate.

Examples:

\begin{verbatim}
\SI{220}{\volt}
\SI{50}{\hertz}
\SI{2.5}{\ampere}
\end{verbatim}

Notation and units shall be consistent throughout the report.

\vspace{0.3cm}

% ================================================================
% XV. SOURCE CODE
% ================================================================

\noindent
\textbf{XV.\quad SOURCE CODE}

\vspace{0.15cm}

If the internship involves programming, PLCs, microcontrollers,
simulation, or software development, relevant source code may be included
in the report.

Source code shall be presented using the \verb|listings| environment or
another environment defined by the official template.

Example:

\begin{verbatim}
\begin{lstlisting}[language=st]
IF Start = TRUE THEN
    Motor := TRUE;
ELSE
    Motor := FALSE;
END_IF;
\end{lstlisting}
\end{verbatim}

Only code necessary for explaining the work should appear in the main
report. Long source-code listings should be moved to an appendix or
submitted as supplementary files if required.

Source code should not normally be inserted as screenshots.

\vspace{0.3cm}

% ================================================================
% XVI. ABSTRACT
% ================================================================

\noindent
\textbf{XVI.\quad ABSTRACT}

\vspace{0.15cm}

The Abstract shall briefly summarize:

\begin{itemize}[leftmargin=0.8cm, itemsep=2pt]
    \item the internship organization;
    \item the objectives of the internship;
    \item the main tasks performed;
    \item the main results, knowledge, and skills obtained.
\end{itemize}

The Abstract should be concise and should not reproduce the full content
of the report chapter by chapter.

\vspace{0.3cm}

% ================================================================
% XVII. REFERENCES AND CITATIONS
% ================================================================

\noindent
\textbf{XVII.\quad REFERENCES AND CITATIONS}

\vspace{0.15cm}

All books, journal papers, conference papers, standards, websites,
datasheets, manuals, technical documents, figures, tables, or other
external sources used in preparing the report shall be properly cited.

References shall be managed through the bibliography system provided with
the official template.

Bibliographic entries shall normally be stored in:

\begin{verbatim}
References.bib
\end{verbatim}

References shall be cited in the report using:

\begin{verbatim}
\cite{reference-key}
\end{verbatim}

Students shall not manually type reference numbers such as ``[1]'',
``[2]'', etc.

Titles of foreign-language publications shall remain in their original
language in the bibliography.

All borrowed ideas, data, figures, tables, equations, methods, or results
shall be properly acknowledged.

\vspace{0.3cm}

% ================================================================
% XVIII. CROSS-REFERENCING
% ================================================================

\noindent
\textbf{XVIII.\quad CROSS-REFERENCING}

\vspace{0.15cm}

Numbered elements shall be cross-referenced using the standard \LaTeX{}
mechanism.

Use the following commands where appropriate:

\begin{verbatim}
\label{...}
\ref{...}
\eqref{...}
\end{verbatim}

For example:

\begin{verbatim}
As shown in Figure~\ref{fig:system}, ...
\end{verbatim}

Students shall not manually type numbers such as ``Figure 2.3'' when
those numbers can be generated automatically.

\vspace{0.3cm}

% ================================================================
% XIX. PAGE NUMBERING
% ================================================================

\noindent
\textbf{XIX.\quad PAGE NUMBERING}

\vspace{0.15cm}

Cover pages and official evaluation forms at the beginning of the report
shall not display page numbers.

The main body of the report shall use Arabic page numbers beginning from
Chapter 1.

Page numbering shall be generated automatically by the official template.
Students shall not manually enter page numbers.

\vspace{0.3cm}

% ================================================================
% XX. PROJECT FILE ORGANIZATION
% ================================================================

\noindent
\textbf{XX.\quad PROJECT FILE ORGANIZATION}

\vspace{0.15cm}

Students should maintain a clear and organized \LaTeX{} project
structure. A recommended structure is:

\begin{verbatim}
InternshipReport/
|
|-- main.tex
|-- References.bib
|
|-- figures/
|   |-- company.pdf
|   |-- system.png
|   `-- ...
|
`-- chapters/
    |-- chapter1.tex
    |-- chapter2.tex
    `-- chapter3.tex
\end{verbatim}

Figures should normally be stored in the \verb|figures| directory
specified by the official template.

File and directory names should preferably contain Latin letters,
numbers, hyphens, or underscores.

Recommended examples include:

\begin{verbatim}
company_structure.pdf
motor_control.png
chapter_01.tex
\end{verbatim}

Spaces and unnecessary special characters in file names should be
avoided where possible.

\vspace{0.3cm}

% ================================================================
% XXI. DOCUMENT COMPILATION AND FINAL CHECK
% ================================================================

\noindent
\textbf{XXI.\quad DOCUMENT COMPILATION AND FINAL CHECK}

\vspace{0.15cm}

Before submission, the complete \LaTeX{} project shall compile
successfully and produce the final PDF report.

The final document shall not contain unresolved references such as:

\begin{verbatim}
??
[?]
\end{verbatim}

Before generating the final PDF, students shall check:

\begin{itemize}[leftmargin=0.8cm, itemsep=2pt]
    \item the Table of Contents;
    \item page numbering;
    \item chapter, section, and subsection numbering;
    \item figure numbering;
    \item table numbering;
    \item equation numbering;
    \item cross-references;
    \item citations;
    \item bibliography;
    \item font display;
    \item mathematical symbols;
    \item image quality;
    \item completeness of all required forms.
\end{itemize}

The document may be compiled using TeX Live, MiKTeX, or another
compatible \LaTeX{} environment. Regardless of the compilation
environment, the resulting PDF shall conform to the official Faculty
template.

\vspace{0.3cm}

% ================================================================
% XXII. CONTENT THAT STUDENTS MAY EDIT
% ================================================================

\noindent
\textbf{XXII.\quad CONTENT THAT STUDENTS MAY EDIT}

\vspace{0.15cm}

Students may edit or replace:

\begin{itemize}[leftmargin=0.8cm, itemsep=2pt]
    \item student name;
    \item class;
    \item internship organization;
    \item supervisor name;
    \item Abstract;
    \item List of Symbols and Abbreviations;
    \item chapter titles and chapter content;
    \item figures and photographs;
    \item tables;
    \item equations;
    \item citations and references.
\end{itemize}

Students shall not modify the main document formatting settings in the
preamble, including the document class, page geometry, line spacing,
heading formatting, table-of-contents formatting, font configuration,
page layout, or cover-page macros unless specifically instructed.

\vspace{0.3cm}

% ================================================================
% XXIII. FINAL SUBMISSION REQUIREMENTS
% ================================================================

\noindent
\textbf{XXIII.\quad FINAL SUBMISSION REQUIREMENTS}

\vspace{0.15cm}

Before submission, the final report shall:

\begin{itemize}[leftmargin=0.8cm, itemsep=2pt]
    \item compile without errors that affect the report content or presentation;
    \item contain all required forms and report sections;
    \item contain correct student, class, internship organization, and supervisor information;
    \item contain an updated Table of Contents, List of Figures, and List of Tables;
    \item contain no unused example content from the template;
    \item contain no placeholder text such as \textit{Title of Chapter 1},
    \textit{Section 1.1}, or other sample entries;
    \item contain all required citations and references;
    \item contain no unresolved cross-references.
\end{itemize}

If the Faculty requires submission of the \LaTeX{} source, the complete
project shall include all \verb|.tex| files, the bibliography database,
figures, and any other files required to reproduce the submitted PDF.

\vspace{0.4cm}

\noindent
\textbf{Important:}
Before submitting the final internship report, students shall remove or
comment out the following command in the main document:

\begin{verbatim}
\input{Regulations.tex}
\end{verbatim}

For example:

\begin{verbatim}
% \input{Regulations.tex}
\end{verbatim}

This removes the guidance pages from the final submitted report while
preserving the regulations file in the \LaTeX{} project for reference.

} % End 13pt content

\clearpage

\end{verbatim}

For example:

\begin{verbatim}
% 
% =================================================================
% REGULATIONS FOR PREPARING THE GRADUATION INTERNSHIP REPORT
% USING LATEX
% =================================================================
%
% This file is intended to be included from the main document with:
%
%     \input{Regulations.tex}
%
% IMPORTANT:
% This section is provided for guidance only. Students shall remove
% or comment out the \input{Regulations.tex} command before submitting
% the final internship report.
%
% =================================================================

\clearpage
\thispagestyle{empty}

\begin{center}
    {\fontsize{16pt}{20pt}\selectfont\bfseries
    SPECIFIC REGULATIONS FOR PREPARING\\
    GRADUATION INTERNSHIP REPORTS USING \LaTeX}
\end{center}

\vspace{0.5cm}

{\fontsize{13pt}{17pt}\selectfont
\setstretch{1.3}
\justifying

\noindent
\textbf{Note:}
This section is provided for guidance only and is included in the main
document using the command \verb|\input{Regulations.tex}|.
Students shall remove or comment out this command before submitting the
final internship report.

\vspace{0.4cm}

% ================================================================
% I. GENERAL REQUIREMENTS
% ================================================================

\noindent
\textbf{I.\quad GENERAL REQUIREMENTS}

\vspace{0.15cm}

The Graduation Internship Report shall be prepared using the official
\LaTeX{} template provided by the Faculty. Students shall preserve the
predefined page layout, typography, line spacing, heading styles,
numbering system, figure and table formatting, and bibliography settings
unless otherwise instructed by the Faculty or the supervisor.

The report shall be written in English.

Students are responsible for the academic and technical content of the
report. The official \LaTeX{} template is responsible for standardized
document formatting.

\vspace{0.3cm}

% ================================================================
% II. REPORT STRUCTURE
% ================================================================

\noindent
\textbf{II.\quad REPORT STRUCTURE}

\vspace{0.15cm}

The report shall normally be arranged in the following order:

\begin{enumerate}[label=\arabic*., leftmargin=1.2cm, itemsep=2pt]
    \item Outer Cover;
    \item Comments of the Internship Host Organization;
    \item Supervisor's Evaluation and Graduation Internship Defense Results;
    \item Table of Contents;
    \item List of Symbols and Abbreviations, if applicable;
    \item List of Figures;
    \item List of Tables;
    \item Abstract;
    \item Main Chapters of the Internship Report;
    \item References;
    \item Appendices, if applicable.
\end{enumerate}

The preliminary pages shall not be numbered as chapters. The main report
content shall be organized into consecutively numbered chapters beginning
with Chapter 1.

\vspace{0.3cm}

% ================================================================
% III. PAGE SIZE AND MARGINS
% ================================================================

\noindent
\textbf{III.\quad PAGE SIZE AND MARGINS}

\vspace{0.15cm}

The report shall be formatted on A4 paper with the following margins:

\begin{center}
\begin{tabular}{ll}
    Left margin:   & 3.5 cm \\
    Right margin:  & 2.0 cm \\
    Top margin:    & 2.5 cm \\
    Bottom margin: & 2.5 cm
\end{tabular}
\end{center}

These settings are predefined in the official template and shall not be
modified by students.

\vspace{0.3cm}

% ================================================================
% IV. FONT, FONT SIZE, AND LINE SPACING
% ================================================================

\noindent
\textbf{IV.\quad FONT, FONT SIZE, AND LINE SPACING}

\vspace{0.15cm}

The main text shall use a 13 pt font size with 1.3-line spacing.
A Times-family typeface shall be used in accordance with the official
template.

The text shall be fully justified.

Students shall not use repeated spaces, repeated manual line breaks, or
multiple blank lines to position text. Paragraph and heading spacing shall
be controlled by the template.

\vspace{0.3cm}

% ================================================================
% V. OUTER COVER
% ================================================================

\noindent
\textbf{V.\quad OUTER COVER}

\vspace{0.15cm}

The outer cover shall follow the official format provided in the template
and shall contain the following information:

\begin{itemize}[leftmargin=0.8cm, itemsep=2pt]
    \item Thai Nguyen University of Technology;
    \item Faculty of International Training;
    \item Engineering and Technology Department;
    \item University logo and Faculty logo;
    \item the title \textbf{INTERNSHIP REPORT};
    \item Internship Organization;
    \item Author;
    \item Class;
    \item Supervisor;
    \item place and year of completion.
\end{itemize}

Students shall only replace the required information fields and shall not
modify the cover-page layout unless instructed otherwise.

\vspace{0.3cm}

% ================================================================
% VI. COMMENTS OF THE INTERNSHIP HOST ORGANIZATION
% ================================================================

\noindent
\textbf{VI.\quad COMMENTS OF THE INTERNSHIP HOST ORGANIZATION}

\vspace{0.15cm}

The page entitled \textbf{COMMENTS OF THE INTERNSHIP HOST ORGANIZATION}
is reserved for the internship host organization.

It includes:

\begin{itemize}[leftmargin=0.8cm, itemsep=2pt]
    \item the student's full name;
    \item class;
    \item comments from the host organization;
    \item score;
    \item date;
    \item signature and official seal.
\end{itemize}

Students shall not alter or complete official evaluation fields unless
specifically authorized.

\vspace{0.3cm}

% ================================================================
% VII. SUPERVISOR'S EVALUATION AND DEFENSE RESULTS
% ================================================================

\noindent
\textbf{VII.\quad SUPERVISOR'S EVALUATION AND DEFENSE RESULTS}

\vspace{0.15cm}

The section entitled
\textbf{SUPERVISOR'S EVALUATION AND GRADUATION INTERNSHIP DEFENSE RESULTS}
is reserved for the supervisor and defense committee.

The evaluation may include:

\begin{itemize}[leftmargin=0.8cm, itemsep=2pt]
    \item Internship Progress and Student's Attitude;
    \item Report Content;
    \item Report Presentation;
    \item Other Comments;
    \item Supervisor's Opinion;
    \item Defense Result;
    \item Score.
\end{itemize}

Students shall not enter evaluation scores or defense results before they
are officially determined.

\vspace{0.3cm}

% ================================================================
% VIII. CHAPTERS, SECTIONS, AND SUBSECTIONS
% ================================================================

\noindent
\textbf{VIII.\quad CHAPTERS, SECTIONS, AND SUBSECTIONS}

\vspace{0.15cm}

The report structure shall be generated using the standard \LaTeX{}
commands:

\begin{verbatim}
\chapter{...}
\section{...}
\subsection{...}
\end{verbatim}

Chapter, section, and subsection numbers shall be generated automatically
by \LaTeX{}.

Students shall not manually type chapter or section numbers and shall not
reproduce heading styles using commands such as \verb|\textbf|,
\verb|\fontsize|, or \verb|\centering| instead of the appropriate
structural command.

Excessive subdivision should be avoided. In most cases, the hierarchy
Chapter--Section--Subsection is sufficient.

\vspace{0.3cm}

% ================================================================
% IX. ORGANIZATION OF THE MAIN CONTENT
% ================================================================

\noindent
\textbf{IX.\quad ORGANIZATION OF THE MAIN CONTENT}

\vspace{0.15cm}

The specific content of the report shall depend on the internship
organization and the assigned tasks. A typical structure may include:

\begin{itemize}[leftmargin=0.8cm, itemsep=2pt]
    \item \textbf{Chapter 1. Overview of the Internship Organization};
    \item \textbf{Chapter 2. Internship Activities and Assigned Tasks};
    \item \textbf{Chapter 3. Results, Experience, and Evaluation}.
\end{itemize}

These chapter titles are recommendations only and may be adapted to the
actual internship content.

The report shall reflect work actually performed, observed, analyzed, or
learned during the internship. Unrelated theoretical material should not
be included merely to increase report length.

\vspace{0.3cm}

% ================================================================
% X. TABLE OF CONTENTS
% ================================================================

\noindent
\textbf{X.\quad TABLE OF CONTENTS}

\vspace{0.15cm}

The Table of Contents shall be generated automatically using:

\begin{verbatim}
\tableofcontents
\end{verbatim}

Students shall not create or manually edit page numbers in the Table of
Contents.

The document shall be compiled sufficiently to ensure that all page
numbers and entries are correctly updated.

\vspace{0.3cm}

% ================================================================
% XI. LIST OF SYMBOLS AND ABBREVIATIONS
% ================================================================

\noindent
\textbf{XI.\quad LIST OF SYMBOLS AND ABBREVIATIONS}

\vspace{0.15cm}

Symbols and abbreviations that are used repeatedly in the report shall be
listed in the \textbf{List of Symbols and Abbreviations}.

The list may contain the following columns:

\begin{center}
\begin{tabular}{llll}
    No. & Symbol / Abbreviation & Unit & Definition
\end{tabular}
\end{center}

Only symbols and abbreviations actually used in the final report shall be
retained. Example entries contained in the template shall be removed if
they are not relevant to the student's report.

Each abbreviation should normally be defined when it first appears in
the main text.

\vspace{0.3cm}

% ================================================================
% XII. FIGURES, DIAGRAMS, AND PHOTOGRAPHS
% ================================================================

\noindent
\textbf{XII.\quad FIGURES, DIAGRAMS, AND PHOTOGRAPHS}

\vspace{0.15cm}

Figures shall be clear, legible, and directly relevant to the report
content.

Figures shall be inserted using the standard \verb|figure| environment.

Example:

\begin{verbatim}
\begin{figure}[H]
    \centering
    \includegraphics[width=0.8\textwidth]
                    {figures/company_structure.pdf}
    \caption{Organizational structure of the company}
    \label{fig:company-structure}
\end{figure}
\end{verbatim}

Figure captions shall be placed below the figure.

Each referenced figure shall have a unique \verb|\label| and shall be
referenced in the text using, for example:

\begin{verbatim}
Figure~\ref{fig:company-structure}
\end{verbatim}

Students shall not manually type figure numbers.

Figures reproduced or adapted from external sources shall include an
appropriate citation.

Photographs taken during the internship shall be relevant and sufficiently
clear. Students shall not include confidential, restricted, or proprietary
information without permission from the internship organization.

Vector graphics or high-resolution images are recommended whenever
possible.

\vspace{0.3cm}

% ================================================================
% XIII. TABLES
% ================================================================

\noindent
\textbf{XIII.\quad TABLES}

\vspace{0.15cm}

Tables shall be created using suitable \LaTeX{} environments such as
\verb|table|, \verb|tabular|, \verb|tabularx|, or \verb|longtable|.

Table captions shall be placed above the table.

Example:

\begin{verbatim}
\begin{table}[H]
    \centering
    \caption{Internship schedule}
    \label{tab:internship-schedule}

    \begin{tabular}{lll}
        \toprule
        Week & Task & Result \\
        \midrule
        1 & Orientation & Completed \\
        2 & System study & Completed \\
        \bottomrule
    \end{tabular}
\end{table}
\end{verbatim}

Tables shall be referenced using:

\begin{verbatim}
Table~\ref{tab:internship-schedule}
\end{verbatim}

Students shall not manually type table numbers.

Tables should be created directly in \LaTeX{} rather than inserted as
screenshots whenever practical.

\vspace{0.3cm}

% ================================================================
% XIV. EQUATIONS AND TECHNICAL QUANTITIES
% ================================================================

\noindent
\textbf{XIV.\quad EQUATIONS AND TECHNICAL QUANTITIES}

\vspace{0.15cm}

Mathematical expressions shall be typeset directly using \LaTeX{}.
Images or screenshots of equations shall not be used.

Short expressions may be written in inline mathematical mode, for example:

\begin{verbatim}
$R = 0.7\,\Omega$
\end{verbatim}

Important equations shall use a numbered equation environment.

Example:

\begin{verbatim}
\begin{equation}
    P = UI
    \label{eq:power}
\end{equation}
\end{verbatim}

Equations shall be referenced using:

\begin{verbatim}
Equation~\eqref{eq:power}
\end{verbatim}

Equation numbers shall not be manually entered.

SI units shall be used wherever applicable. The \verb|siunitx| package
should be used where appropriate.

Examples:

\begin{verbatim}
\SI{220}{\volt}
\SI{50}{\hertz}
\SI{2.5}{\ampere}
\end{verbatim}

Notation and units shall be consistent throughout the report.

\vspace{0.3cm}

% ================================================================
% XV. SOURCE CODE
% ================================================================

\noindent
\textbf{XV.\quad SOURCE CODE}

\vspace{0.15cm}

If the internship involves programming, PLCs, microcontrollers,
simulation, or software development, relevant source code may be included
in the report.

Source code shall be presented using the \verb|listings| environment or
another environment defined by the official template.

Example:

\begin{verbatim}
\begin{lstlisting}[language=st]
IF Start = TRUE THEN
    Motor := TRUE;
ELSE
    Motor := FALSE;
END_IF;
\end{lstlisting}
\end{verbatim}

Only code necessary for explaining the work should appear in the main
report. Long source-code listings should be moved to an appendix or
submitted as supplementary files if required.

Source code should not normally be inserted as screenshots.

\vspace{0.3cm}

% ================================================================
% XVI. ABSTRACT
% ================================================================

\noindent
\textbf{XVI.\quad ABSTRACT}

\vspace{0.15cm}

The Abstract shall briefly summarize:

\begin{itemize}[leftmargin=0.8cm, itemsep=2pt]
    \item the internship organization;
    \item the objectives of the internship;
    \item the main tasks performed;
    \item the main results, knowledge, and skills obtained.
\end{itemize}

The Abstract should be concise and should not reproduce the full content
of the report chapter by chapter.

\vspace{0.3cm}

% ================================================================
% XVII. REFERENCES AND CITATIONS
% ================================================================

\noindent
\textbf{XVII.\quad REFERENCES AND CITATIONS}

\vspace{0.15cm}

All books, journal papers, conference papers, standards, websites,
datasheets, manuals, technical documents, figures, tables, or other
external sources used in preparing the report shall be properly cited.

References shall be managed through the bibliography system provided with
the official template.

Bibliographic entries shall normally be stored in:

\begin{verbatim}
References.bib
\end{verbatim}

References shall be cited in the report using:

\begin{verbatim}
\cite{reference-key}
\end{verbatim}

Students shall not manually type reference numbers such as ``[1]'',
``[2]'', etc.

Titles of foreign-language publications shall remain in their original
language in the bibliography.

All borrowed ideas, data, figures, tables, equations, methods, or results
shall be properly acknowledged.

\vspace{0.3cm}

% ================================================================
% XVIII. CROSS-REFERENCING
% ================================================================

\noindent
\textbf{XVIII.\quad CROSS-REFERENCING}

\vspace{0.15cm}

Numbered elements shall be cross-referenced using the standard \LaTeX{}
mechanism.

Use the following commands where appropriate:

\begin{verbatim}
\label{...}
\ref{...}
\eqref{...}
\end{verbatim}

For example:

\begin{verbatim}
As shown in Figure~\ref{fig:system}, ...
\end{verbatim}

Students shall not manually type numbers such as ``Figure 2.3'' when
those numbers can be generated automatically.

\vspace{0.3cm}

% ================================================================
% XIX. PAGE NUMBERING
% ================================================================

\noindent
\textbf{XIX.\quad PAGE NUMBERING}

\vspace{0.15cm}

Cover pages and official evaluation forms at the beginning of the report
shall not display page numbers.

The main body of the report shall use Arabic page numbers beginning from
Chapter 1.

Page numbering shall be generated automatically by the official template.
Students shall not manually enter page numbers.

\vspace{0.3cm}

% ================================================================
% XX. PROJECT FILE ORGANIZATION
% ================================================================

\noindent
\textbf{XX.\quad PROJECT FILE ORGANIZATION}

\vspace{0.15cm}

Students should maintain a clear and organized \LaTeX{} project
structure. A recommended structure is:

\begin{verbatim}
InternshipReport/
|
|-- main.tex
|-- References.bib
|
|-- figures/
|   |-- company.pdf
|   |-- system.png
|   `-- ...
|
`-- chapters/
    |-- chapter1.tex
    |-- chapter2.tex
    `-- chapter3.tex
\end{verbatim}

Figures should normally be stored in the \verb|figures| directory
specified by the official template.

File and directory names should preferably contain Latin letters,
numbers, hyphens, or underscores.

Recommended examples include:

\begin{verbatim}
company_structure.pdf
motor_control.png
chapter_01.tex
\end{verbatim}

Spaces and unnecessary special characters in file names should be
avoided where possible.

\vspace{0.3cm}

% ================================================================
% XXI. DOCUMENT COMPILATION AND FINAL CHECK
% ================================================================

\noindent
\textbf{XXI.\quad DOCUMENT COMPILATION AND FINAL CHECK}

\vspace{0.15cm}

Before submission, the complete \LaTeX{} project shall compile
successfully and produce the final PDF report.

The final document shall not contain unresolved references such as:

\begin{verbatim}
??
[?]
\end{verbatim}

Before generating the final PDF, students shall check:

\begin{itemize}[leftmargin=0.8cm, itemsep=2pt]
    \item the Table of Contents;
    \item page numbering;
    \item chapter, section, and subsection numbering;
    \item figure numbering;
    \item table numbering;
    \item equation numbering;
    \item cross-references;
    \item citations;
    \item bibliography;
    \item font display;
    \item mathematical symbols;
    \item image quality;
    \item completeness of all required forms.
\end{itemize}

The document may be compiled using TeX Live, MiKTeX, or another
compatible \LaTeX{} environment. Regardless of the compilation
environment, the resulting PDF shall conform to the official Faculty
template.

\vspace{0.3cm}

% ================================================================
% XXII. CONTENT THAT STUDENTS MAY EDIT
% ================================================================

\noindent
\textbf{XXII.\quad CONTENT THAT STUDENTS MAY EDIT}

\vspace{0.15cm}

Students may edit or replace:

\begin{itemize}[leftmargin=0.8cm, itemsep=2pt]
    \item student name;
    \item class;
    \item internship organization;
    \item supervisor name;
    \item Abstract;
    \item List of Symbols and Abbreviations;
    \item chapter titles and chapter content;
    \item figures and photographs;
    \item tables;
    \item equations;
    \item citations and references.
\end{itemize}

Students shall not modify the main document formatting settings in the
preamble, including the document class, page geometry, line spacing,
heading formatting, table-of-contents formatting, font configuration,
page layout, or cover-page macros unless specifically instructed.

\vspace{0.3cm}

% ================================================================
% XXIII. FINAL SUBMISSION REQUIREMENTS
% ================================================================

\noindent
\textbf{XXIII.\quad FINAL SUBMISSION REQUIREMENTS}

\vspace{0.15cm}

Before submission, the final report shall:

\begin{itemize}[leftmargin=0.8cm, itemsep=2pt]
    \item compile without errors that affect the report content or presentation;
    \item contain all required forms and report sections;
    \item contain correct student, class, internship organization, and supervisor information;
    \item contain an updated Table of Contents, List of Figures, and List of Tables;
    \item contain no unused example content from the template;
    \item contain no placeholder text such as \textit{Title of Chapter 1},
    \textit{Section 1.1}, or other sample entries;
    \item contain all required citations and references;
    \item contain no unresolved cross-references.
\end{itemize}

If the Faculty requires submission of the \LaTeX{} source, the complete
project shall include all \verb|.tex| files, the bibliography database,
figures, and any other files required to reproduce the submitted PDF.

\vspace{0.4cm}

\noindent
\textbf{Important:}
Before submitting the final internship report, students shall remove or
comment out the following command in the main document:

\begin{verbatim}
\input{Regulations.tex}
\end{verbatim}

For example:

\begin{verbatim}
% \input{Regulations.tex}
\end{verbatim}

This removes the guidance pages from the final submitted report while
preserving the regulations file in the \LaTeX{} project for reference.

} % End 13pt content

\clearpage

\end{verbatim}

This removes the guidance pages from the final submitted report while
preserving the regulations file in the \LaTeX{} project for reference.

} % End 13pt content

\clearpage

\end{verbatim}

This removes the guidance pages from the final submitted report while
preserving the regulations file in the \LaTeX{} project for reference.

} % End 13pt content

\clearpage
|.
Students shall remove or comment out this command before submitting the
final internship report.

\vspace{0.4cm}

% ================================================================
% I. GENERAL REQUIREMENTS
% ================================================================

\noindent
\textbf{I.\quad GENERAL REQUIREMENTS}

\vspace{0.15cm}

The Graduation Internship Report shall be prepared using the official
\LaTeX{} template provided by the Faculty. Students shall preserve the
predefined page layout, typography, line spacing, heading styles,
numbering system, figure and table formatting, and bibliography settings
unless otherwise instructed by the Faculty or the supervisor.

The report shall be written in English.

Students are responsible for the academic and technical content of the
report. The official \LaTeX{} template is responsible for standardized
document formatting.

\vspace{0.3cm}

% ================================================================
% II. REPORT STRUCTURE
% ================================================================

\noindent
\textbf{II.\quad REPORT STRUCTURE}

\vspace{0.15cm}

The report shall normally be arranged in the following order:

\begin{enumerate}[label=\arabic*., leftmargin=1.2cm, itemsep=2pt]
    \item Outer Cover;
    \item Comments of the Internship Host Organization;
    \item Supervisor's Evaluation and Graduation Internship Defense Results;
    \item Table of Contents;
    \item List of Symbols and Abbreviations, if applicable;
    \item List of Figures;
    \item List of Tables;
    \item Abstract;
    \item Main Chapters of the Internship Report;
    \item References;
    \item Appendices, if applicable.
\end{enumerate}

The preliminary pages shall not be numbered as chapters. The main report
content shall be organized into consecutively numbered chapters beginning
with Chapter 1.

\vspace{0.3cm}

% ================================================================
% III. PAGE SIZE AND MARGINS
% ================================================================

\noindent
\textbf{III.\quad PAGE SIZE AND MARGINS}

\vspace{0.15cm}

The report shall be formatted on A4 paper with the following margins:

\begin{center}
\begin{tabular}{ll}
    Left margin:   & 3.5 cm \\
    Right margin:  & 2.0 cm \\
    Top margin:    & 2.5 cm \\
    Bottom margin: & 2.5 cm
\end{tabular}
\end{center}

These settings are predefined in the official template and shall not be
modified by students.

\vspace{0.3cm}

% ================================================================
% IV. FONT, FONT SIZE, AND LINE SPACING
% ================================================================

\noindent
\textbf{IV.\quad FONT, FONT SIZE, AND LINE SPACING}

\vspace{0.15cm}

The main text shall use a 13 pt font size with 1.3-line spacing.
A Times-family typeface shall be used in accordance with the official
template.

The text shall be fully justified.

Students shall not use repeated spaces, repeated manual line breaks, or
multiple blank lines to position text. Paragraph and heading spacing shall
be controlled by the template.

\vspace{0.3cm}

% ================================================================
% V. OUTER COVER
% ================================================================

\noindent
\textbf{V.\quad OUTER COVER}

\vspace{0.15cm}

The outer cover shall follow the official format provided in the template
and shall contain the following information:

\begin{itemize}[leftmargin=0.8cm, itemsep=2pt]
    \item Thai Nguyen University of Technology;
    \item Faculty of International Training;
    \item Engineering and Technology Department;
    \item University logo and Faculty logo;
    \item the title \textbf{INTERNSHIP REPORT};
    \item Internship Organization;
    \item Author;
    \item Class;
    \item Supervisor;
    \item place and year of completion.
\end{itemize}

Students shall only replace the required information fields and shall not
modify the cover-page layout unless instructed otherwise.

\vspace{0.3cm}

% ================================================================
% VI. COMMENTS OF THE INTERNSHIP HOST ORGANIZATION
% ================================================================

\noindent
\textbf{VI.\quad COMMENTS OF THE INTERNSHIP HOST ORGANIZATION}

\vspace{0.15cm}

The page entitled \textbf{COMMENTS OF THE INTERNSHIP HOST ORGANIZATION}
is reserved for the internship host organization.

It includes:

\begin{itemize}[leftmargin=0.8cm, itemsep=2pt]
    \item the student's full name;
    \item class;
    \item comments from the host organization;
    \item score;
    \item date;
    \item signature and official seal.
\end{itemize}

Students shall not alter or complete official evaluation fields unless
specifically authorized.

\vspace{0.3cm}

% ================================================================
% VII. SUPERVISOR'S EVALUATION AND DEFENSE RESULTS
% ================================================================

\noindent
\textbf{VII.\quad SUPERVISOR'S EVALUATION AND DEFENSE RESULTS}

\vspace{0.15cm}

The section entitled
\textbf{SUPERVISOR'S EVALUATION AND GRADUATION INTERNSHIP DEFENSE RESULTS}
is reserved for the supervisor and defense committee.

The evaluation may include:

\begin{itemize}[leftmargin=0.8cm, itemsep=2pt]
    \item Internship Progress and Student's Attitude;
    \item Report Content;
    \item Report Presentation;
    \item Other Comments;
    \item Supervisor's Opinion;
    \item Defense Result;
    \item Score.
\end{itemize}

Students shall not enter evaluation scores or defense results before they
are officially determined.

\vspace{0.3cm}

% ================================================================
% VIII. CHAPTERS, SECTIONS, AND SUBSECTIONS
% ================================================================

\noindent
\textbf{VIII.\quad CHAPTERS, SECTIONS, AND SUBSECTIONS}

\vspace{0.15cm}

The report structure shall be generated using the standard \LaTeX{}
commands:

\begin{verbatim}
\chapter{...}
\section{...}
\subsection{...}
\end{verbatim}

Chapter, section, and subsection numbers shall be generated automatically
by \LaTeX{}.

Students shall not manually type chapter or section numbers and shall not
reproduce heading styles using commands such as \verb|\textbf|,
\verb|\fontsize|, or \verb|\centering| instead of the appropriate
structural command.

Excessive subdivision should be avoided. In most cases, the hierarchy
Chapter--Section--Subsection is sufficient.

\vspace{0.3cm}

% ================================================================
% IX. ORGANIZATION OF THE MAIN CONTENT
% ================================================================

\noindent
\textbf{IX.\quad ORGANIZATION OF THE MAIN CONTENT}

\vspace{0.15cm}

The specific content of the report shall depend on the internship
organization and the assigned tasks. A typical structure may include:

\begin{itemize}[leftmargin=0.8cm, itemsep=2pt]
    \item \textbf{Chapter 1. Overview of the Internship Organization};
    \item \textbf{Chapter 2. Internship Activities and Assigned Tasks};
    \item \textbf{Chapter 3. Results, Experience, and Evaluation}.
\end{itemize}

These chapter titles are recommendations only and may be adapted to the
actual internship content.

The report shall reflect work actually performed, observed, analyzed, or
learned during the internship. Unrelated theoretical material should not
be included merely to increase report length.

\vspace{0.3cm}

% ================================================================
% X. TABLE OF CONTENTS
% ================================================================

\noindent
\textbf{X.\quad TABLE OF CONTENTS}

\vspace{0.15cm}

The Table of Contents shall be generated automatically using:

\begin{verbatim}
\tableofcontents
\end{verbatim}

Students shall not create or manually edit page numbers in the Table of
Contents.

The document shall be compiled sufficiently to ensure that all page
numbers and entries are correctly updated.

\vspace{0.3cm}

% ================================================================
% XI. LIST OF SYMBOLS AND ABBREVIATIONS
% ================================================================

\noindent
\textbf{XI.\quad LIST OF SYMBOLS AND ABBREVIATIONS}

\vspace{0.15cm}

Symbols and abbreviations that are used repeatedly in the report shall be
listed in the \textbf{List of Symbols and Abbreviations}.

The list may contain the following columns:

\begin{center}
\begin{tabular}{llll}
    No. & Symbol / Abbreviation & Unit & Definition
\end{tabular}
\end{center}

Only symbols and abbreviations actually used in the final report shall be
retained. Example entries contained in the template shall be removed if
they are not relevant to the student's report.

Each abbreviation should normally be defined when it first appears in
the main text.

\vspace{0.3cm}

% ================================================================
% XII. FIGURES, DIAGRAMS, AND PHOTOGRAPHS
% ================================================================

\noindent
\textbf{XII.\quad FIGURES, DIAGRAMS, AND PHOTOGRAPHS}

\vspace{0.15cm}

Figures shall be clear, legible, and directly relevant to the report
content.

Figures shall be inserted using the standard \verb|figure| environment.

Example:

\begin{verbatim}
\begin{figure}[H]
    \centering
    \includegraphics[width=0.8\textwidth]
                    {figures/company_structure.pdf}
    \caption{Organizational structure of the company}
    \label{fig:company-structure}
\end{figure}
\end{verbatim}

Figure captions shall be placed below the figure.

Each referenced figure shall have a unique \verb|\label| and shall be
referenced in the text using, for example:

\begin{verbatim}
Figure~\ref{fig:company-structure}
\end{verbatim}

Students shall not manually type figure numbers.

Figures reproduced or adapted from external sources shall include an
appropriate citation.

Photographs taken during the internship shall be relevant and sufficiently
clear. Students shall not include confidential, restricted, or proprietary
information without permission from the internship organization.

Vector graphics or high-resolution images are recommended whenever
possible.

\vspace{0.3cm}

% ================================================================
% XIII. TABLES
% ================================================================

\noindent
\textbf{XIII.\quad TABLES}

\vspace{0.15cm}

Tables shall be created using suitable \LaTeX{} environments such as
\verb|table|, \verb|tabular|, \verb|tabularx|, or \verb|longtable|.

Table captions shall be placed above the table.

Example:

\begin{verbatim}
\begin{table}[H]
    \centering
    \caption{Internship schedule}
    \label{tab:internship-schedule}

    \begin{tabular}{lll}
        \toprule
        Week & Task & Result \\
        \midrule
        1 & Orientation & Completed \\
        2 & System study & Completed \\
        \bottomrule
    \end{tabular}
\end{table}
\end{verbatim}

Tables shall be referenced using:

\begin{verbatim}
Table~\ref{tab:internship-schedule}
\end{verbatim}

Students shall not manually type table numbers.

Tables should be created directly in \LaTeX{} rather than inserted as
screenshots whenever practical.

\vspace{0.3cm}

% ================================================================
% XIV. EQUATIONS AND TECHNICAL QUANTITIES
% ================================================================

\noindent
\textbf{XIV.\quad EQUATIONS AND TECHNICAL QUANTITIES}

\vspace{0.15cm}

Mathematical expressions shall be typeset directly using \LaTeX{}.
Images or screenshots of equations shall not be used.

Short expressions may be written in inline mathematical mode, for example:

\begin{verbatim}
$R = 0.7\,\Omega$
\end{verbatim}

Important equations shall use a numbered equation environment.

Example:

\begin{verbatim}
\begin{equation}
    P = UI
    \label{eq:power}
\end{equation}
\end{verbatim}

Equations shall be referenced using:

\begin{verbatim}
Equation~\eqref{eq:power}
\end{verbatim}

Equation numbers shall not be manually entered.

SI units shall be used wherever applicable. The \verb|siunitx| package
should be used where appropriate.

Examples:

\begin{verbatim}
\SI{220}{\volt}
\SI{50}{\hertz}
\SI{2.5}{\ampere}
\end{verbatim}

Notation and units shall be consistent throughout the report.

\vspace{0.3cm}

% ================================================================
% XV. SOURCE CODE
% ================================================================

\noindent
\textbf{XV.\quad SOURCE CODE}

\vspace{0.15cm}

If the internship involves programming, PLCs, microcontrollers,
simulation, or software development, relevant source code may be included
in the report.

Source code shall be presented using the \verb|listings| environment or
another environment defined by the official template.

Example:

\begin{verbatim}
\begin{lstlisting}[language=st]
IF Start = TRUE THEN
    Motor := TRUE;
ELSE
    Motor := FALSE;
END_IF;
\end{lstlisting}
\end{verbatim}

Only code necessary for explaining the work should appear in the main
report. Long source-code listings should be moved to an appendix or
submitted as supplementary files if required.

Source code should not normally be inserted as screenshots.

\vspace{0.3cm}

% ================================================================
% XVI. ABSTRACT
% ================================================================

\noindent
\textbf{XVI.\quad ABSTRACT}

\vspace{0.15cm}

The Abstract shall briefly summarize:

\begin{itemize}[leftmargin=0.8cm, itemsep=2pt]
    \item the internship organization;
    \item the objectives of the internship;
    \item the main tasks performed;
    \item the main results, knowledge, and skills obtained.
\end{itemize}

The Abstract should be concise and should not reproduce the full content
of the report chapter by chapter.

\vspace{0.3cm}

% ================================================================
% XVII. REFERENCES AND CITATIONS
% ================================================================

\noindent
\textbf{XVII.\quad REFERENCES AND CITATIONS}

\vspace{0.15cm}

All books, journal papers, conference papers, standards, websites,
datasheets, manuals, technical documents, figures, tables, or other
external sources used in preparing the report shall be properly cited.

References shall be managed through the bibliography system provided with
the official template.

Bibliographic entries shall normally be stored in:

\begin{verbatim}
References.bib
\end{verbatim}

References shall be cited in the report using:

\begin{verbatim}
\cite{reference-key}
\end{verbatim}

Students shall not manually type reference numbers such as ``[1]'',
``[2]'', etc.

Titles of foreign-language publications shall remain in their original
language in the bibliography.

All borrowed ideas, data, figures, tables, equations, methods, or results
shall be properly acknowledged.

\vspace{0.3cm}

% ================================================================
% XVIII. CROSS-REFERENCING
% ================================================================

\noindent
\textbf{XVIII.\quad CROSS-REFERENCING}

\vspace{0.15cm}

Numbered elements shall be cross-referenced using the standard \LaTeX{}
mechanism.

Use the following commands where appropriate:

\begin{verbatim}
\label{...}
\ref{...}
\eqref{...}
\end{verbatim}

For example:

\begin{verbatim}
As shown in Figure~\ref{fig:system}, ...
\end{verbatim}

Students shall not manually type numbers such as ``Figure 2.3'' when
those numbers can be generated automatically.

\vspace{0.3cm}

% ================================================================
% XIX. PAGE NUMBERING
% ================================================================

\noindent
\textbf{XIX.\quad PAGE NUMBERING}

\vspace{0.15cm}

Cover pages and official evaluation forms at the beginning of the report
shall not display page numbers.

The main body of the report shall use Arabic page numbers beginning from
Chapter 1.

Page numbering shall be generated automatically by the official template.
Students shall not manually enter page numbers.

\vspace{0.3cm}

% ================================================================
% XX. PROJECT FILE ORGANIZATION
% ================================================================

\noindent
\textbf{XX.\quad PROJECT FILE ORGANIZATION}

\vspace{0.15cm}

Students should maintain a clear and organized \LaTeX{} project
structure. A recommended structure is:

\begin{verbatim}
InternshipReport/
|
|-- main.tex
|-- References.bib
|
|-- figures/
|   |-- company.pdf
|   |-- system.png
|   `-- ...
|
`-- chapters/
    |-- chapter1.tex
    |-- chapter2.tex
    `-- chapter3.tex
\end{verbatim}

Figures should normally be stored in the \verb|figures| directory
specified by the official template.

File and directory names should preferably contain Latin letters,
numbers, hyphens, or underscores.

Recommended examples include:

\begin{verbatim}
company_structure.pdf
motor_control.png
chapter_01.tex
\end{verbatim}

Spaces and unnecessary special characters in file names should be
avoided where possible.

\vspace{0.3cm}

% ================================================================
% XXI. DOCUMENT COMPILATION AND FINAL CHECK
% ================================================================

\noindent
\textbf{XXI.\quad DOCUMENT COMPILATION AND FINAL CHECK}

\vspace{0.15cm}

Before submission, the complete \LaTeX{} project shall compile
successfully and produce the final PDF report.

The final document shall not contain unresolved references such as:

\begin{verbatim}
??
[?]
\end{verbatim}

Before generating the final PDF, students shall check:

\begin{itemize}[leftmargin=0.8cm, itemsep=2pt]
    \item the Table of Contents;
    \item page numbering;
    \item chapter, section, and subsection numbering;
    \item figure numbering;
    \item table numbering;
    \item equation numbering;
    \item cross-references;
    \item citations;
    \item bibliography;
    \item font display;
    \item mathematical symbols;
    \item image quality;
    \item completeness of all required forms.
\end{itemize}

The document may be compiled using TeX Live, MiKTeX, or another
compatible \LaTeX{} environment. Regardless of the compilation
environment, the resulting PDF shall conform to the official Faculty
template.

\vspace{0.3cm}

% ================================================================
% XXII. CONTENT THAT STUDENTS MAY EDIT
% ================================================================

\noindent
\textbf{XXII.\quad CONTENT THAT STUDENTS MAY EDIT}

\vspace{0.15cm}

Students may edit or replace:

\begin{itemize}[leftmargin=0.8cm, itemsep=2pt]
    \item student name;
    \item class;
    \item internship organization;
    \item supervisor name;
    \item Abstract;
    \item List of Symbols and Abbreviations;
    \item chapter titles and chapter content;
    \item figures and photographs;
    \item tables;
    \item equations;
    \item citations and references.
\end{itemize}

Students shall not modify the main document formatting settings in the
preamble, including the document class, page geometry, line spacing,
heading formatting, table-of-contents formatting, font configuration,
page layout, or cover-page macros unless specifically instructed.

\vspace{0.3cm}

% ================================================================
% XXIII. FINAL SUBMISSION REQUIREMENTS
% ================================================================

\noindent
\textbf{XXIII.\quad FINAL SUBMISSION REQUIREMENTS}

\vspace{0.15cm}

Before submission, the final report shall:

\begin{itemize}[leftmargin=0.8cm, itemsep=2pt]
    \item compile without errors that affect the report content or presentation;
    \item contain all required forms and report sections;
    \item contain correct student, class, internship organization, and supervisor information;
    \item contain an updated Table of Contents, List of Figures, and List of Tables;
    \item contain no unused example content from the template;
    \item contain no placeholder text such as \textit{Title of Chapter 1},
    \textit{Section 1.1}, or other sample entries;
    \item contain all required citations and references;
    \item contain no unresolved cross-references.
\end{itemize}

If the Faculty requires submission of the \LaTeX{} source, the complete
project shall include all \verb|.tex| files, the bibliography database,
figures, and any other files required to reproduce the submitted PDF.

\vspace{0.4cm}

\noindent
\textbf{Important:}
Before submitting the final internship report, students shall remove or
comment out the following command in the main document:

\begin{verbatim}

% =================================================================
% REGULATIONS FOR PREPARING THE GRADUATION INTERNSHIP REPORT
% USING LATEX
% =================================================================
%
% This file is intended to be included from the main document with:
%
%     
% =================================================================
% REGULATIONS FOR PREPARING THE GRADUATION INTERNSHIP REPORT
% USING LATEX
% =================================================================
%
% This file is intended to be included from the main document with:
%
%     
% =================================================================
% REGULATIONS FOR PREPARING THE GRADUATION INTERNSHIP REPORT
% USING LATEX
% =================================================================
%
% This file is intended to be included from the main document with:
%
%     \input{Regulations.tex}
%
% IMPORTANT:
% This section is provided for guidance only. Students shall remove
% or comment out the \input{Regulations.tex} command before submitting
% the final internship report.
%
% =================================================================

\clearpage
\thispagestyle{empty}

\begin{center}
    {\fontsize{16pt}{20pt}\selectfont\bfseries
    SPECIFIC REGULATIONS FOR PREPARING\\
    GRADUATION INTERNSHIP REPORTS USING \LaTeX}
\end{center}

\vspace{0.5cm}

{\fontsize{13pt}{17pt}\selectfont
\setstretch{1.3}
\justifying

\noindent
\textbf{Note:}
This section is provided for guidance only and is included in the main
document using the command \verb|\input{Regulations.tex}|.
Students shall remove or comment out this command before submitting the
final internship report.

\vspace{0.4cm}

% ================================================================
% I. GENERAL REQUIREMENTS
% ================================================================

\noindent
\textbf{I.\quad GENERAL REQUIREMENTS}

\vspace{0.15cm}

The Graduation Internship Report shall be prepared using the official
\LaTeX{} template provided by the Faculty. Students shall preserve the
predefined page layout, typography, line spacing, heading styles,
numbering system, figure and table formatting, and bibliography settings
unless otherwise instructed by the Faculty or the supervisor.

The report shall be written in English.

Students are responsible for the academic and technical content of the
report. The official \LaTeX{} template is responsible for standardized
document formatting.

\vspace{0.3cm}

% ================================================================
% II. REPORT STRUCTURE
% ================================================================

\noindent
\textbf{II.\quad REPORT STRUCTURE}

\vspace{0.15cm}

The report shall normally be arranged in the following order:

\begin{enumerate}[label=\arabic*., leftmargin=1.2cm, itemsep=2pt]
    \item Outer Cover;
    \item Comments of the Internship Host Organization;
    \item Supervisor's Evaluation and Graduation Internship Defense Results;
    \item Table of Contents;
    \item List of Symbols and Abbreviations, if applicable;
    \item List of Figures;
    \item List of Tables;
    \item Abstract;
    \item Main Chapters of the Internship Report;
    \item References;
    \item Appendices, if applicable.
\end{enumerate}

The preliminary pages shall not be numbered as chapters. The main report
content shall be organized into consecutively numbered chapters beginning
with Chapter 1.

\vspace{0.3cm}

% ================================================================
% III. PAGE SIZE AND MARGINS
% ================================================================

\noindent
\textbf{III.\quad PAGE SIZE AND MARGINS}

\vspace{0.15cm}

The report shall be formatted on A4 paper with the following margins:

\begin{center}
\begin{tabular}{ll}
    Left margin:   & 3.5 cm \\
    Right margin:  & 2.0 cm \\
    Top margin:    & 2.5 cm \\
    Bottom margin: & 2.5 cm
\end{tabular}
\end{center}

These settings are predefined in the official template and shall not be
modified by students.

\vspace{0.3cm}

% ================================================================
% IV. FONT, FONT SIZE, AND LINE SPACING
% ================================================================

\noindent
\textbf{IV.\quad FONT, FONT SIZE, AND LINE SPACING}

\vspace{0.15cm}

The main text shall use a 13 pt font size with 1.3-line spacing.
A Times-family typeface shall be used in accordance with the official
template.

The text shall be fully justified.

Students shall not use repeated spaces, repeated manual line breaks, or
multiple blank lines to position text. Paragraph and heading spacing shall
be controlled by the template.

\vspace{0.3cm}

% ================================================================
% V. OUTER COVER
% ================================================================

\noindent
\textbf{V.\quad OUTER COVER}

\vspace{0.15cm}

The outer cover shall follow the official format provided in the template
and shall contain the following information:

\begin{itemize}[leftmargin=0.8cm, itemsep=2pt]
    \item Thai Nguyen University of Technology;
    \item Faculty of International Training;
    \item Engineering and Technology Department;
    \item University logo and Faculty logo;
    \item the title \textbf{INTERNSHIP REPORT};
    \item Internship Organization;
    \item Author;
    \item Class;
    \item Supervisor;
    \item place and year of completion.
\end{itemize}

Students shall only replace the required information fields and shall not
modify the cover-page layout unless instructed otherwise.

\vspace{0.3cm}

% ================================================================
% VI. COMMENTS OF THE INTERNSHIP HOST ORGANIZATION
% ================================================================

\noindent
\textbf{VI.\quad COMMENTS OF THE INTERNSHIP HOST ORGANIZATION}

\vspace{0.15cm}

The page entitled \textbf{COMMENTS OF THE INTERNSHIP HOST ORGANIZATION}
is reserved for the internship host organization.

It includes:

\begin{itemize}[leftmargin=0.8cm, itemsep=2pt]
    \item the student's full name;
    \item class;
    \item comments from the host organization;
    \item score;
    \item date;
    \item signature and official seal.
\end{itemize}

Students shall not alter or complete official evaluation fields unless
specifically authorized.

\vspace{0.3cm}

% ================================================================
% VII. SUPERVISOR'S EVALUATION AND DEFENSE RESULTS
% ================================================================

\noindent
\textbf{VII.\quad SUPERVISOR'S EVALUATION AND DEFENSE RESULTS}

\vspace{0.15cm}

The section entitled
\textbf{SUPERVISOR'S EVALUATION AND GRADUATION INTERNSHIP DEFENSE RESULTS}
is reserved for the supervisor and defense committee.

The evaluation may include:

\begin{itemize}[leftmargin=0.8cm, itemsep=2pt]
    \item Internship Progress and Student's Attitude;
    \item Report Content;
    \item Report Presentation;
    \item Other Comments;
    \item Supervisor's Opinion;
    \item Defense Result;
    \item Score.
\end{itemize}

Students shall not enter evaluation scores or defense results before they
are officially determined.

\vspace{0.3cm}

% ================================================================
% VIII. CHAPTERS, SECTIONS, AND SUBSECTIONS
% ================================================================

\noindent
\textbf{VIII.\quad CHAPTERS, SECTIONS, AND SUBSECTIONS}

\vspace{0.15cm}

The report structure shall be generated using the standard \LaTeX{}
commands:

\begin{verbatim}
\chapter{...}
\section{...}
\subsection{...}
\end{verbatim}

Chapter, section, and subsection numbers shall be generated automatically
by \LaTeX{}.

Students shall not manually type chapter or section numbers and shall not
reproduce heading styles using commands such as \verb|\textbf|,
\verb|\fontsize|, or \verb|\centering| instead of the appropriate
structural command.

Excessive subdivision should be avoided. In most cases, the hierarchy
Chapter--Section--Subsection is sufficient.

\vspace{0.3cm}

% ================================================================
% IX. ORGANIZATION OF THE MAIN CONTENT
% ================================================================

\noindent
\textbf{IX.\quad ORGANIZATION OF THE MAIN CONTENT}

\vspace{0.15cm}

The specific content of the report shall depend on the internship
organization and the assigned tasks. A typical structure may include:

\begin{itemize}[leftmargin=0.8cm, itemsep=2pt]
    \item \textbf{Chapter 1. Overview of the Internship Organization};
    \item \textbf{Chapter 2. Internship Activities and Assigned Tasks};
    \item \textbf{Chapter 3. Results, Experience, and Evaluation}.
\end{itemize}

These chapter titles are recommendations only and may be adapted to the
actual internship content.

The report shall reflect work actually performed, observed, analyzed, or
learned during the internship. Unrelated theoretical material should not
be included merely to increase report length.

\vspace{0.3cm}

% ================================================================
% X. TABLE OF CONTENTS
% ================================================================

\noindent
\textbf{X.\quad TABLE OF CONTENTS}

\vspace{0.15cm}

The Table of Contents shall be generated automatically using:

\begin{verbatim}
\tableofcontents
\end{verbatim}

Students shall not create or manually edit page numbers in the Table of
Contents.

The document shall be compiled sufficiently to ensure that all page
numbers and entries are correctly updated.

\vspace{0.3cm}

% ================================================================
% XI. LIST OF SYMBOLS AND ABBREVIATIONS
% ================================================================

\noindent
\textbf{XI.\quad LIST OF SYMBOLS AND ABBREVIATIONS}

\vspace{0.15cm}

Symbols and abbreviations that are used repeatedly in the report shall be
listed in the \textbf{List of Symbols and Abbreviations}.

The list may contain the following columns:

\begin{center}
\begin{tabular}{llll}
    No. & Symbol / Abbreviation & Unit & Definition
\end{tabular}
\end{center}

Only symbols and abbreviations actually used in the final report shall be
retained. Example entries contained in the template shall be removed if
they are not relevant to the student's report.

Each abbreviation should normally be defined when it first appears in
the main text.

\vspace{0.3cm}

% ================================================================
% XII. FIGURES, DIAGRAMS, AND PHOTOGRAPHS
% ================================================================

\noindent
\textbf{XII.\quad FIGURES, DIAGRAMS, AND PHOTOGRAPHS}

\vspace{0.15cm}

Figures shall be clear, legible, and directly relevant to the report
content.

Figures shall be inserted using the standard \verb|figure| environment.

Example:

\begin{verbatim}
\begin{figure}[H]
    \centering
    \includegraphics[width=0.8\textwidth]
                    {figures/company_structure.pdf}
    \caption{Organizational structure of the company}
    \label{fig:company-structure}
\end{figure}
\end{verbatim}

Figure captions shall be placed below the figure.

Each referenced figure shall have a unique \verb|\label| and shall be
referenced in the text using, for example:

\begin{verbatim}
Figure~\ref{fig:company-structure}
\end{verbatim}

Students shall not manually type figure numbers.

Figures reproduced or adapted from external sources shall include an
appropriate citation.

Photographs taken during the internship shall be relevant and sufficiently
clear. Students shall not include confidential, restricted, or proprietary
information without permission from the internship organization.

Vector graphics or high-resolution images are recommended whenever
possible.

\vspace{0.3cm}

% ================================================================
% XIII. TABLES
% ================================================================

\noindent
\textbf{XIII.\quad TABLES}

\vspace{0.15cm}

Tables shall be created using suitable \LaTeX{} environments such as
\verb|table|, \verb|tabular|, \verb|tabularx|, or \verb|longtable|.

Table captions shall be placed above the table.

Example:

\begin{verbatim}
\begin{table}[H]
    \centering
    \caption{Internship schedule}
    \label{tab:internship-schedule}

    \begin{tabular}{lll}
        \toprule
        Week & Task & Result \\
        \midrule
        1 & Orientation & Completed \\
        2 & System study & Completed \\
        \bottomrule
    \end{tabular}
\end{table}
\end{verbatim}

Tables shall be referenced using:

\begin{verbatim}
Table~\ref{tab:internship-schedule}
\end{verbatim}

Students shall not manually type table numbers.

Tables should be created directly in \LaTeX{} rather than inserted as
screenshots whenever practical.

\vspace{0.3cm}

% ================================================================
% XIV. EQUATIONS AND TECHNICAL QUANTITIES
% ================================================================

\noindent
\textbf{XIV.\quad EQUATIONS AND TECHNICAL QUANTITIES}

\vspace{0.15cm}

Mathematical expressions shall be typeset directly using \LaTeX{}.
Images or screenshots of equations shall not be used.

Short expressions may be written in inline mathematical mode, for example:

\begin{verbatim}
$R = 0.7\,\Omega$
\end{verbatim}

Important equations shall use a numbered equation environment.

Example:

\begin{verbatim}
\begin{equation}
    P = UI
    \label{eq:power}
\end{equation}
\end{verbatim}

Equations shall be referenced using:

\begin{verbatim}
Equation~\eqref{eq:power}
\end{verbatim}

Equation numbers shall not be manually entered.

SI units shall be used wherever applicable. The \verb|siunitx| package
should be used where appropriate.

Examples:

\begin{verbatim}
\SI{220}{\volt}
\SI{50}{\hertz}
\SI{2.5}{\ampere}
\end{verbatim}

Notation and units shall be consistent throughout the report.

\vspace{0.3cm}

% ================================================================
% XV. SOURCE CODE
% ================================================================

\noindent
\textbf{XV.\quad SOURCE CODE}

\vspace{0.15cm}

If the internship involves programming, PLCs, microcontrollers,
simulation, or software development, relevant source code may be included
in the report.

Source code shall be presented using the \verb|listings| environment or
another environment defined by the official template.

Example:

\begin{verbatim}
\begin{lstlisting}[language=st]
IF Start = TRUE THEN
    Motor := TRUE;
ELSE
    Motor := FALSE;
END_IF;
\end{lstlisting}
\end{verbatim}

Only code necessary for explaining the work should appear in the main
report. Long source-code listings should be moved to an appendix or
submitted as supplementary files if required.

Source code should not normally be inserted as screenshots.

\vspace{0.3cm}

% ================================================================
% XVI. ABSTRACT
% ================================================================

\noindent
\textbf{XVI.\quad ABSTRACT}

\vspace{0.15cm}

The Abstract shall briefly summarize:

\begin{itemize}[leftmargin=0.8cm, itemsep=2pt]
    \item the internship organization;
    \item the objectives of the internship;
    \item the main tasks performed;
    \item the main results, knowledge, and skills obtained.
\end{itemize}

The Abstract should be concise and should not reproduce the full content
of the report chapter by chapter.

\vspace{0.3cm}

% ================================================================
% XVII. REFERENCES AND CITATIONS
% ================================================================

\noindent
\textbf{XVII.\quad REFERENCES AND CITATIONS}

\vspace{0.15cm}

All books, journal papers, conference papers, standards, websites,
datasheets, manuals, technical documents, figures, tables, or other
external sources used in preparing the report shall be properly cited.

References shall be managed through the bibliography system provided with
the official template.

Bibliographic entries shall normally be stored in:

\begin{verbatim}
References.bib
\end{verbatim}

References shall be cited in the report using:

\begin{verbatim}
\cite{reference-key}
\end{verbatim}

Students shall not manually type reference numbers such as ``[1]'',
``[2]'', etc.

Titles of foreign-language publications shall remain in their original
language in the bibliography.

All borrowed ideas, data, figures, tables, equations, methods, or results
shall be properly acknowledged.

\vspace{0.3cm}

% ================================================================
% XVIII. CROSS-REFERENCING
% ================================================================

\noindent
\textbf{XVIII.\quad CROSS-REFERENCING}

\vspace{0.15cm}

Numbered elements shall be cross-referenced using the standard \LaTeX{}
mechanism.

Use the following commands where appropriate:

\begin{verbatim}
\label{...}
\ref{...}
\eqref{...}
\end{verbatim}

For example:

\begin{verbatim}
As shown in Figure~\ref{fig:system}, ...
\end{verbatim}

Students shall not manually type numbers such as ``Figure 2.3'' when
those numbers can be generated automatically.

\vspace{0.3cm}

% ================================================================
% XIX. PAGE NUMBERING
% ================================================================

\noindent
\textbf{XIX.\quad PAGE NUMBERING}

\vspace{0.15cm}

Cover pages and official evaluation forms at the beginning of the report
shall not display page numbers.

The main body of the report shall use Arabic page numbers beginning from
Chapter 1.

Page numbering shall be generated automatically by the official template.
Students shall not manually enter page numbers.

\vspace{0.3cm}

% ================================================================
% XX. PROJECT FILE ORGANIZATION
% ================================================================

\noindent
\textbf{XX.\quad PROJECT FILE ORGANIZATION}

\vspace{0.15cm}

Students should maintain a clear and organized \LaTeX{} project
structure. A recommended structure is:

\begin{verbatim}
InternshipReport/
|
|-- main.tex
|-- References.bib
|
|-- figures/
|   |-- company.pdf
|   |-- system.png
|   `-- ...
|
`-- chapters/
    |-- chapter1.tex
    |-- chapter2.tex
    `-- chapter3.tex
\end{verbatim}

Figures should normally be stored in the \verb|figures| directory
specified by the official template.

File and directory names should preferably contain Latin letters,
numbers, hyphens, or underscores.

Recommended examples include:

\begin{verbatim}
company_structure.pdf
motor_control.png
chapter_01.tex
\end{verbatim}

Spaces and unnecessary special characters in file names should be
avoided where possible.

\vspace{0.3cm}

% ================================================================
% XXI. DOCUMENT COMPILATION AND FINAL CHECK
% ================================================================

\noindent
\textbf{XXI.\quad DOCUMENT COMPILATION AND FINAL CHECK}

\vspace{0.15cm}

Before submission, the complete \LaTeX{} project shall compile
successfully and produce the final PDF report.

The final document shall not contain unresolved references such as:

\begin{verbatim}
??
[?]
\end{verbatim}

Before generating the final PDF, students shall check:

\begin{itemize}[leftmargin=0.8cm, itemsep=2pt]
    \item the Table of Contents;
    \item page numbering;
    \item chapter, section, and subsection numbering;
    \item figure numbering;
    \item table numbering;
    \item equation numbering;
    \item cross-references;
    \item citations;
    \item bibliography;
    \item font display;
    \item mathematical symbols;
    \item image quality;
    \item completeness of all required forms.
\end{itemize}

The document may be compiled using TeX Live, MiKTeX, or another
compatible \LaTeX{} environment. Regardless of the compilation
environment, the resulting PDF shall conform to the official Faculty
template.

\vspace{0.3cm}

% ================================================================
% XXII. CONTENT THAT STUDENTS MAY EDIT
% ================================================================

\noindent
\textbf{XXII.\quad CONTENT THAT STUDENTS MAY EDIT}

\vspace{0.15cm}

Students may edit or replace:

\begin{itemize}[leftmargin=0.8cm, itemsep=2pt]
    \item student name;
    \item class;
    \item internship organization;
    \item supervisor name;
    \item Abstract;
    \item List of Symbols and Abbreviations;
    \item chapter titles and chapter content;
    \item figures and photographs;
    \item tables;
    \item equations;
    \item citations and references.
\end{itemize}

Students shall not modify the main document formatting settings in the
preamble, including the document class, page geometry, line spacing,
heading formatting, table-of-contents formatting, font configuration,
page layout, or cover-page macros unless specifically instructed.

\vspace{0.3cm}

% ================================================================
% XXIII. FINAL SUBMISSION REQUIREMENTS
% ================================================================

\noindent
\textbf{XXIII.\quad FINAL SUBMISSION REQUIREMENTS}

\vspace{0.15cm}

Before submission, the final report shall:

\begin{itemize}[leftmargin=0.8cm, itemsep=2pt]
    \item compile without errors that affect the report content or presentation;
    \item contain all required forms and report sections;
    \item contain correct student, class, internship organization, and supervisor information;
    \item contain an updated Table of Contents, List of Figures, and List of Tables;
    \item contain no unused example content from the template;
    \item contain no placeholder text such as \textit{Title of Chapter 1},
    \textit{Section 1.1}, or other sample entries;
    \item contain all required citations and references;
    \item contain no unresolved cross-references.
\end{itemize}

If the Faculty requires submission of the \LaTeX{} source, the complete
project shall include all \verb|.tex| files, the bibliography database,
figures, and any other files required to reproduce the submitted PDF.

\vspace{0.4cm}

\noindent
\textbf{Important:}
Before submitting the final internship report, students shall remove or
comment out the following command in the main document:

\begin{verbatim}
\input{Regulations.tex}
\end{verbatim}

For example:

\begin{verbatim}
% \input{Regulations.tex}
\end{verbatim}

This removes the guidance pages from the final submitted report while
preserving the regulations file in the \LaTeX{} project for reference.

} % End 13pt content

\clearpage

%
% IMPORTANT:
% This section is provided for guidance only. Students shall remove
% or comment out the 
% =================================================================
% REGULATIONS FOR PREPARING THE GRADUATION INTERNSHIP REPORT
% USING LATEX
% =================================================================
%
% This file is intended to be included from the main document with:
%
%     \input{Regulations.tex}
%
% IMPORTANT:
% This section is provided for guidance only. Students shall remove
% or comment out the \input{Regulations.tex} command before submitting
% the final internship report.
%
% =================================================================

\clearpage
\thispagestyle{empty}

\begin{center}
    {\fontsize{16pt}{20pt}\selectfont\bfseries
    SPECIFIC REGULATIONS FOR PREPARING\\
    GRADUATION INTERNSHIP REPORTS USING \LaTeX}
\end{center}

\vspace{0.5cm}

{\fontsize{13pt}{17pt}\selectfont
\setstretch{1.3}
\justifying

\noindent
\textbf{Note:}
This section is provided for guidance only and is included in the main
document using the command \verb|\input{Regulations.tex}|.
Students shall remove or comment out this command before submitting the
final internship report.

\vspace{0.4cm}

% ================================================================
% I. GENERAL REQUIREMENTS
% ================================================================

\noindent
\textbf{I.\quad GENERAL REQUIREMENTS}

\vspace{0.15cm}

The Graduation Internship Report shall be prepared using the official
\LaTeX{} template provided by the Faculty. Students shall preserve the
predefined page layout, typography, line spacing, heading styles,
numbering system, figure and table formatting, and bibliography settings
unless otherwise instructed by the Faculty or the supervisor.

The report shall be written in English.

Students are responsible for the academic and technical content of the
report. The official \LaTeX{} template is responsible for standardized
document formatting.

\vspace{0.3cm}

% ================================================================
% II. REPORT STRUCTURE
% ================================================================

\noindent
\textbf{II.\quad REPORT STRUCTURE}

\vspace{0.15cm}

The report shall normally be arranged in the following order:

\begin{enumerate}[label=\arabic*., leftmargin=1.2cm, itemsep=2pt]
    \item Outer Cover;
    \item Comments of the Internship Host Organization;
    \item Supervisor's Evaluation and Graduation Internship Defense Results;
    \item Table of Contents;
    \item List of Symbols and Abbreviations, if applicable;
    \item List of Figures;
    \item List of Tables;
    \item Abstract;
    \item Main Chapters of the Internship Report;
    \item References;
    \item Appendices, if applicable.
\end{enumerate}

The preliminary pages shall not be numbered as chapters. The main report
content shall be organized into consecutively numbered chapters beginning
with Chapter 1.

\vspace{0.3cm}

% ================================================================
% III. PAGE SIZE AND MARGINS
% ================================================================

\noindent
\textbf{III.\quad PAGE SIZE AND MARGINS}

\vspace{0.15cm}

The report shall be formatted on A4 paper with the following margins:

\begin{center}
\begin{tabular}{ll}
    Left margin:   & 3.5 cm \\
    Right margin:  & 2.0 cm \\
    Top margin:    & 2.5 cm \\
    Bottom margin: & 2.5 cm
\end{tabular}
\end{center}

These settings are predefined in the official template and shall not be
modified by students.

\vspace{0.3cm}

% ================================================================
% IV. FONT, FONT SIZE, AND LINE SPACING
% ================================================================

\noindent
\textbf{IV.\quad FONT, FONT SIZE, AND LINE SPACING}

\vspace{0.15cm}

The main text shall use a 13 pt font size with 1.3-line spacing.
A Times-family typeface shall be used in accordance with the official
template.

The text shall be fully justified.

Students shall not use repeated spaces, repeated manual line breaks, or
multiple blank lines to position text. Paragraph and heading spacing shall
be controlled by the template.

\vspace{0.3cm}

% ================================================================
% V. OUTER COVER
% ================================================================

\noindent
\textbf{V.\quad OUTER COVER}

\vspace{0.15cm}

The outer cover shall follow the official format provided in the template
and shall contain the following information:

\begin{itemize}[leftmargin=0.8cm, itemsep=2pt]
    \item Thai Nguyen University of Technology;
    \item Faculty of International Training;
    \item Engineering and Technology Department;
    \item University logo and Faculty logo;
    \item the title \textbf{INTERNSHIP REPORT};
    \item Internship Organization;
    \item Author;
    \item Class;
    \item Supervisor;
    \item place and year of completion.
\end{itemize}

Students shall only replace the required information fields and shall not
modify the cover-page layout unless instructed otherwise.

\vspace{0.3cm}

% ================================================================
% VI. COMMENTS OF THE INTERNSHIP HOST ORGANIZATION
% ================================================================

\noindent
\textbf{VI.\quad COMMENTS OF THE INTERNSHIP HOST ORGANIZATION}

\vspace{0.15cm}

The page entitled \textbf{COMMENTS OF THE INTERNSHIP HOST ORGANIZATION}
is reserved for the internship host organization.

It includes:

\begin{itemize}[leftmargin=0.8cm, itemsep=2pt]
    \item the student's full name;
    \item class;
    \item comments from the host organization;
    \item score;
    \item date;
    \item signature and official seal.
\end{itemize}

Students shall not alter or complete official evaluation fields unless
specifically authorized.

\vspace{0.3cm}

% ================================================================
% VII. SUPERVISOR'S EVALUATION AND DEFENSE RESULTS
% ================================================================

\noindent
\textbf{VII.\quad SUPERVISOR'S EVALUATION AND DEFENSE RESULTS}

\vspace{0.15cm}

The section entitled
\textbf{SUPERVISOR'S EVALUATION AND GRADUATION INTERNSHIP DEFENSE RESULTS}
is reserved for the supervisor and defense committee.

The evaluation may include:

\begin{itemize}[leftmargin=0.8cm, itemsep=2pt]
    \item Internship Progress and Student's Attitude;
    \item Report Content;
    \item Report Presentation;
    \item Other Comments;
    \item Supervisor's Opinion;
    \item Defense Result;
    \item Score.
\end{itemize}

Students shall not enter evaluation scores or defense results before they
are officially determined.

\vspace{0.3cm}

% ================================================================
% VIII. CHAPTERS, SECTIONS, AND SUBSECTIONS
% ================================================================

\noindent
\textbf{VIII.\quad CHAPTERS, SECTIONS, AND SUBSECTIONS}

\vspace{0.15cm}

The report structure shall be generated using the standard \LaTeX{}
commands:

\begin{verbatim}
\chapter{...}
\section{...}
\subsection{...}
\end{verbatim}

Chapter, section, and subsection numbers shall be generated automatically
by \LaTeX{}.

Students shall not manually type chapter or section numbers and shall not
reproduce heading styles using commands such as \verb|\textbf|,
\verb|\fontsize|, or \verb|\centering| instead of the appropriate
structural command.

Excessive subdivision should be avoided. In most cases, the hierarchy
Chapter--Section--Subsection is sufficient.

\vspace{0.3cm}

% ================================================================
% IX. ORGANIZATION OF THE MAIN CONTENT
% ================================================================

\noindent
\textbf{IX.\quad ORGANIZATION OF THE MAIN CONTENT}

\vspace{0.15cm}

The specific content of the report shall depend on the internship
organization and the assigned tasks. A typical structure may include:

\begin{itemize}[leftmargin=0.8cm, itemsep=2pt]
    \item \textbf{Chapter 1. Overview of the Internship Organization};
    \item \textbf{Chapter 2. Internship Activities and Assigned Tasks};
    \item \textbf{Chapter 3. Results, Experience, and Evaluation}.
\end{itemize}

These chapter titles are recommendations only and may be adapted to the
actual internship content.

The report shall reflect work actually performed, observed, analyzed, or
learned during the internship. Unrelated theoretical material should not
be included merely to increase report length.

\vspace{0.3cm}

% ================================================================
% X. TABLE OF CONTENTS
% ================================================================

\noindent
\textbf{X.\quad TABLE OF CONTENTS}

\vspace{0.15cm}

The Table of Contents shall be generated automatically using:

\begin{verbatim}
\tableofcontents
\end{verbatim}

Students shall not create or manually edit page numbers in the Table of
Contents.

The document shall be compiled sufficiently to ensure that all page
numbers and entries are correctly updated.

\vspace{0.3cm}

% ================================================================
% XI. LIST OF SYMBOLS AND ABBREVIATIONS
% ================================================================

\noindent
\textbf{XI.\quad LIST OF SYMBOLS AND ABBREVIATIONS}

\vspace{0.15cm}

Symbols and abbreviations that are used repeatedly in the report shall be
listed in the \textbf{List of Symbols and Abbreviations}.

The list may contain the following columns:

\begin{center}
\begin{tabular}{llll}
    No. & Symbol / Abbreviation & Unit & Definition
\end{tabular}
\end{center}

Only symbols and abbreviations actually used in the final report shall be
retained. Example entries contained in the template shall be removed if
they are not relevant to the student's report.

Each abbreviation should normally be defined when it first appears in
the main text.

\vspace{0.3cm}

% ================================================================
% XII. FIGURES, DIAGRAMS, AND PHOTOGRAPHS
% ================================================================

\noindent
\textbf{XII.\quad FIGURES, DIAGRAMS, AND PHOTOGRAPHS}

\vspace{0.15cm}

Figures shall be clear, legible, and directly relevant to the report
content.

Figures shall be inserted using the standard \verb|figure| environment.

Example:

\begin{verbatim}
\begin{figure}[H]
    \centering
    \includegraphics[width=0.8\textwidth]
                    {figures/company_structure.pdf}
    \caption{Organizational structure of the company}
    \label{fig:company-structure}
\end{figure}
\end{verbatim}

Figure captions shall be placed below the figure.

Each referenced figure shall have a unique \verb|\label| and shall be
referenced in the text using, for example:

\begin{verbatim}
Figure~\ref{fig:company-structure}
\end{verbatim}

Students shall not manually type figure numbers.

Figures reproduced or adapted from external sources shall include an
appropriate citation.

Photographs taken during the internship shall be relevant and sufficiently
clear. Students shall not include confidential, restricted, or proprietary
information without permission from the internship organization.

Vector graphics or high-resolution images are recommended whenever
possible.

\vspace{0.3cm}

% ================================================================
% XIII. TABLES
% ================================================================

\noindent
\textbf{XIII.\quad TABLES}

\vspace{0.15cm}

Tables shall be created using suitable \LaTeX{} environments such as
\verb|table|, \verb|tabular|, \verb|tabularx|, or \verb|longtable|.

Table captions shall be placed above the table.

Example:

\begin{verbatim}
\begin{table}[H]
    \centering
    \caption{Internship schedule}
    \label{tab:internship-schedule}

    \begin{tabular}{lll}
        \toprule
        Week & Task & Result \\
        \midrule
        1 & Orientation & Completed \\
        2 & System study & Completed \\
        \bottomrule
    \end{tabular}
\end{table}
\end{verbatim}

Tables shall be referenced using:

\begin{verbatim}
Table~\ref{tab:internship-schedule}
\end{verbatim}

Students shall not manually type table numbers.

Tables should be created directly in \LaTeX{} rather than inserted as
screenshots whenever practical.

\vspace{0.3cm}

% ================================================================
% XIV. EQUATIONS AND TECHNICAL QUANTITIES
% ================================================================

\noindent
\textbf{XIV.\quad EQUATIONS AND TECHNICAL QUANTITIES}

\vspace{0.15cm}

Mathematical expressions shall be typeset directly using \LaTeX{}.
Images or screenshots of equations shall not be used.

Short expressions may be written in inline mathematical mode, for example:

\begin{verbatim}
$R = 0.7\,\Omega$
\end{verbatim}

Important equations shall use a numbered equation environment.

Example:

\begin{verbatim}
\begin{equation}
    P = UI
    \label{eq:power}
\end{equation}
\end{verbatim}

Equations shall be referenced using:

\begin{verbatim}
Equation~\eqref{eq:power}
\end{verbatim}

Equation numbers shall not be manually entered.

SI units shall be used wherever applicable. The \verb|siunitx| package
should be used where appropriate.

Examples:

\begin{verbatim}
\SI{220}{\volt}
\SI{50}{\hertz}
\SI{2.5}{\ampere}
\end{verbatim}

Notation and units shall be consistent throughout the report.

\vspace{0.3cm}

% ================================================================
% XV. SOURCE CODE
% ================================================================

\noindent
\textbf{XV.\quad SOURCE CODE}

\vspace{0.15cm}

If the internship involves programming, PLCs, microcontrollers,
simulation, or software development, relevant source code may be included
in the report.

Source code shall be presented using the \verb|listings| environment or
another environment defined by the official template.

Example:

\begin{verbatim}
\begin{lstlisting}[language=st]
IF Start = TRUE THEN
    Motor := TRUE;
ELSE
    Motor := FALSE;
END_IF;
\end{lstlisting}
\end{verbatim}

Only code necessary for explaining the work should appear in the main
report. Long source-code listings should be moved to an appendix or
submitted as supplementary files if required.

Source code should not normally be inserted as screenshots.

\vspace{0.3cm}

% ================================================================
% XVI. ABSTRACT
% ================================================================

\noindent
\textbf{XVI.\quad ABSTRACT}

\vspace{0.15cm}

The Abstract shall briefly summarize:

\begin{itemize}[leftmargin=0.8cm, itemsep=2pt]
    \item the internship organization;
    \item the objectives of the internship;
    \item the main tasks performed;
    \item the main results, knowledge, and skills obtained.
\end{itemize}

The Abstract should be concise and should not reproduce the full content
of the report chapter by chapter.

\vspace{0.3cm}

% ================================================================
% XVII. REFERENCES AND CITATIONS
% ================================================================

\noindent
\textbf{XVII.\quad REFERENCES AND CITATIONS}

\vspace{0.15cm}

All books, journal papers, conference papers, standards, websites,
datasheets, manuals, technical documents, figures, tables, or other
external sources used in preparing the report shall be properly cited.

References shall be managed through the bibliography system provided with
the official template.

Bibliographic entries shall normally be stored in:

\begin{verbatim}
References.bib
\end{verbatim}

References shall be cited in the report using:

\begin{verbatim}
\cite{reference-key}
\end{verbatim}

Students shall not manually type reference numbers such as ``[1]'',
``[2]'', etc.

Titles of foreign-language publications shall remain in their original
language in the bibliography.

All borrowed ideas, data, figures, tables, equations, methods, or results
shall be properly acknowledged.

\vspace{0.3cm}

% ================================================================
% XVIII. CROSS-REFERENCING
% ================================================================

\noindent
\textbf{XVIII.\quad CROSS-REFERENCING}

\vspace{0.15cm}

Numbered elements shall be cross-referenced using the standard \LaTeX{}
mechanism.

Use the following commands where appropriate:

\begin{verbatim}
\label{...}
\ref{...}
\eqref{...}
\end{verbatim}

For example:

\begin{verbatim}
As shown in Figure~\ref{fig:system}, ...
\end{verbatim}

Students shall not manually type numbers such as ``Figure 2.3'' when
those numbers can be generated automatically.

\vspace{0.3cm}

% ================================================================
% XIX. PAGE NUMBERING
% ================================================================

\noindent
\textbf{XIX.\quad PAGE NUMBERING}

\vspace{0.15cm}

Cover pages and official evaluation forms at the beginning of the report
shall not display page numbers.

The main body of the report shall use Arabic page numbers beginning from
Chapter 1.

Page numbering shall be generated automatically by the official template.
Students shall not manually enter page numbers.

\vspace{0.3cm}

% ================================================================
% XX. PROJECT FILE ORGANIZATION
% ================================================================

\noindent
\textbf{XX.\quad PROJECT FILE ORGANIZATION}

\vspace{0.15cm}

Students should maintain a clear and organized \LaTeX{} project
structure. A recommended structure is:

\begin{verbatim}
InternshipReport/
|
|-- main.tex
|-- References.bib
|
|-- figures/
|   |-- company.pdf
|   |-- system.png
|   `-- ...
|
`-- chapters/
    |-- chapter1.tex
    |-- chapter2.tex
    `-- chapter3.tex
\end{verbatim}

Figures should normally be stored in the \verb|figures| directory
specified by the official template.

File and directory names should preferably contain Latin letters,
numbers, hyphens, or underscores.

Recommended examples include:

\begin{verbatim}
company_structure.pdf
motor_control.png
chapter_01.tex
\end{verbatim}

Spaces and unnecessary special characters in file names should be
avoided where possible.

\vspace{0.3cm}

% ================================================================
% XXI. DOCUMENT COMPILATION AND FINAL CHECK
% ================================================================

\noindent
\textbf{XXI.\quad DOCUMENT COMPILATION AND FINAL CHECK}

\vspace{0.15cm}

Before submission, the complete \LaTeX{} project shall compile
successfully and produce the final PDF report.

The final document shall not contain unresolved references such as:

\begin{verbatim}
??
[?]
\end{verbatim}

Before generating the final PDF, students shall check:

\begin{itemize}[leftmargin=0.8cm, itemsep=2pt]
    \item the Table of Contents;
    \item page numbering;
    \item chapter, section, and subsection numbering;
    \item figure numbering;
    \item table numbering;
    \item equation numbering;
    \item cross-references;
    \item citations;
    \item bibliography;
    \item font display;
    \item mathematical symbols;
    \item image quality;
    \item completeness of all required forms.
\end{itemize}

The document may be compiled using TeX Live, MiKTeX, or another
compatible \LaTeX{} environment. Regardless of the compilation
environment, the resulting PDF shall conform to the official Faculty
template.

\vspace{0.3cm}

% ================================================================
% XXII. CONTENT THAT STUDENTS MAY EDIT
% ================================================================

\noindent
\textbf{XXII.\quad CONTENT THAT STUDENTS MAY EDIT}

\vspace{0.15cm}

Students may edit or replace:

\begin{itemize}[leftmargin=0.8cm, itemsep=2pt]
    \item student name;
    \item class;
    \item internship organization;
    \item supervisor name;
    \item Abstract;
    \item List of Symbols and Abbreviations;
    \item chapter titles and chapter content;
    \item figures and photographs;
    \item tables;
    \item equations;
    \item citations and references.
\end{itemize}

Students shall not modify the main document formatting settings in the
preamble, including the document class, page geometry, line spacing,
heading formatting, table-of-contents formatting, font configuration,
page layout, or cover-page macros unless specifically instructed.

\vspace{0.3cm}

% ================================================================
% XXIII. FINAL SUBMISSION REQUIREMENTS
% ================================================================

\noindent
\textbf{XXIII.\quad FINAL SUBMISSION REQUIREMENTS}

\vspace{0.15cm}

Before submission, the final report shall:

\begin{itemize}[leftmargin=0.8cm, itemsep=2pt]
    \item compile without errors that affect the report content or presentation;
    \item contain all required forms and report sections;
    \item contain correct student, class, internship organization, and supervisor information;
    \item contain an updated Table of Contents, List of Figures, and List of Tables;
    \item contain no unused example content from the template;
    \item contain no placeholder text such as \textit{Title of Chapter 1},
    \textit{Section 1.1}, or other sample entries;
    \item contain all required citations and references;
    \item contain no unresolved cross-references.
\end{itemize}

If the Faculty requires submission of the \LaTeX{} source, the complete
project shall include all \verb|.tex| files, the bibliography database,
figures, and any other files required to reproduce the submitted PDF.

\vspace{0.4cm}

\noindent
\textbf{Important:}
Before submitting the final internship report, students shall remove or
comment out the following command in the main document:

\begin{verbatim}
\input{Regulations.tex}
\end{verbatim}

For example:

\begin{verbatim}
% \input{Regulations.tex}
\end{verbatim}

This removes the guidance pages from the final submitted report while
preserving the regulations file in the \LaTeX{} project for reference.

} % End 13pt content

\clearpage
 command before submitting
% the final internship report.
%
% =================================================================

\clearpage
\thispagestyle{empty}

\begin{center}
    {\fontsize{16pt}{20pt}\selectfont\bfseries
    SPECIFIC REGULATIONS FOR PREPARING\\
    GRADUATION INTERNSHIP REPORTS USING \LaTeX}
\end{center}

\vspace{0.5cm}

{\fontsize{13pt}{17pt}\selectfont
\setstretch{1.3}
\justifying

\noindent
\textbf{Note:}
This section is provided for guidance only and is included in the main
document using the command \verb|
% =================================================================
% REGULATIONS FOR PREPARING THE GRADUATION INTERNSHIP REPORT
% USING LATEX
% =================================================================
%
% This file is intended to be included from the main document with:
%
%     \input{Regulations.tex}
%
% IMPORTANT:
% This section is provided for guidance only. Students shall remove
% or comment out the \input{Regulations.tex} command before submitting
% the final internship report.
%
% =================================================================

\clearpage
\thispagestyle{empty}

\begin{center}
    {\fontsize{16pt}{20pt}\selectfont\bfseries
    SPECIFIC REGULATIONS FOR PREPARING\\
    GRADUATION INTERNSHIP REPORTS USING \LaTeX}
\end{center}

\vspace{0.5cm}

{\fontsize{13pt}{17pt}\selectfont
\setstretch{1.3}
\justifying

\noindent
\textbf{Note:}
This section is provided for guidance only and is included in the main
document using the command \verb|\input{Regulations.tex}|.
Students shall remove or comment out this command before submitting the
final internship report.

\vspace{0.4cm}

% ================================================================
% I. GENERAL REQUIREMENTS
% ================================================================

\noindent
\textbf{I.\quad GENERAL REQUIREMENTS}

\vspace{0.15cm}

The Graduation Internship Report shall be prepared using the official
\LaTeX{} template provided by the Faculty. Students shall preserve the
predefined page layout, typography, line spacing, heading styles,
numbering system, figure and table formatting, and bibliography settings
unless otherwise instructed by the Faculty or the supervisor.

The report shall be written in English.

Students are responsible for the academic and technical content of the
report. The official \LaTeX{} template is responsible for standardized
document formatting.

\vspace{0.3cm}

% ================================================================
% II. REPORT STRUCTURE
% ================================================================

\noindent
\textbf{II.\quad REPORT STRUCTURE}

\vspace{0.15cm}

The report shall normally be arranged in the following order:

\begin{enumerate}[label=\arabic*., leftmargin=1.2cm, itemsep=2pt]
    \item Outer Cover;
    \item Comments of the Internship Host Organization;
    \item Supervisor's Evaluation and Graduation Internship Defense Results;
    \item Table of Contents;
    \item List of Symbols and Abbreviations, if applicable;
    \item List of Figures;
    \item List of Tables;
    \item Abstract;
    \item Main Chapters of the Internship Report;
    \item References;
    \item Appendices, if applicable.
\end{enumerate}

The preliminary pages shall not be numbered as chapters. The main report
content shall be organized into consecutively numbered chapters beginning
with Chapter 1.

\vspace{0.3cm}

% ================================================================
% III. PAGE SIZE AND MARGINS
% ================================================================

\noindent
\textbf{III.\quad PAGE SIZE AND MARGINS}

\vspace{0.15cm}

The report shall be formatted on A4 paper with the following margins:

\begin{center}
\begin{tabular}{ll}
    Left margin:   & 3.5 cm \\
    Right margin:  & 2.0 cm \\
    Top margin:    & 2.5 cm \\
    Bottom margin: & 2.5 cm
\end{tabular}
\end{center}

These settings are predefined in the official template and shall not be
modified by students.

\vspace{0.3cm}

% ================================================================
% IV. FONT, FONT SIZE, AND LINE SPACING
% ================================================================

\noindent
\textbf{IV.\quad FONT, FONT SIZE, AND LINE SPACING}

\vspace{0.15cm}

The main text shall use a 13 pt font size with 1.3-line spacing.
A Times-family typeface shall be used in accordance with the official
template.

The text shall be fully justified.

Students shall not use repeated spaces, repeated manual line breaks, or
multiple blank lines to position text. Paragraph and heading spacing shall
be controlled by the template.

\vspace{0.3cm}

% ================================================================
% V. OUTER COVER
% ================================================================

\noindent
\textbf{V.\quad OUTER COVER}

\vspace{0.15cm}

The outer cover shall follow the official format provided in the template
and shall contain the following information:

\begin{itemize}[leftmargin=0.8cm, itemsep=2pt]
    \item Thai Nguyen University of Technology;
    \item Faculty of International Training;
    \item Engineering and Technology Department;
    \item University logo and Faculty logo;
    \item the title \textbf{INTERNSHIP REPORT};
    \item Internship Organization;
    \item Author;
    \item Class;
    \item Supervisor;
    \item place and year of completion.
\end{itemize}

Students shall only replace the required information fields and shall not
modify the cover-page layout unless instructed otherwise.

\vspace{0.3cm}

% ================================================================
% VI. COMMENTS OF THE INTERNSHIP HOST ORGANIZATION
% ================================================================

\noindent
\textbf{VI.\quad COMMENTS OF THE INTERNSHIP HOST ORGANIZATION}

\vspace{0.15cm}

The page entitled \textbf{COMMENTS OF THE INTERNSHIP HOST ORGANIZATION}
is reserved for the internship host organization.

It includes:

\begin{itemize}[leftmargin=0.8cm, itemsep=2pt]
    \item the student's full name;
    \item class;
    \item comments from the host organization;
    \item score;
    \item date;
    \item signature and official seal.
\end{itemize}

Students shall not alter or complete official evaluation fields unless
specifically authorized.

\vspace{0.3cm}

% ================================================================
% VII. SUPERVISOR'S EVALUATION AND DEFENSE RESULTS
% ================================================================

\noindent
\textbf{VII.\quad SUPERVISOR'S EVALUATION AND DEFENSE RESULTS}

\vspace{0.15cm}

The section entitled
\textbf{SUPERVISOR'S EVALUATION AND GRADUATION INTERNSHIP DEFENSE RESULTS}
is reserved for the supervisor and defense committee.

The evaluation may include:

\begin{itemize}[leftmargin=0.8cm, itemsep=2pt]
    \item Internship Progress and Student's Attitude;
    \item Report Content;
    \item Report Presentation;
    \item Other Comments;
    \item Supervisor's Opinion;
    \item Defense Result;
    \item Score.
\end{itemize}

Students shall not enter evaluation scores or defense results before they
are officially determined.

\vspace{0.3cm}

% ================================================================
% VIII. CHAPTERS, SECTIONS, AND SUBSECTIONS
% ================================================================

\noindent
\textbf{VIII.\quad CHAPTERS, SECTIONS, AND SUBSECTIONS}

\vspace{0.15cm}

The report structure shall be generated using the standard \LaTeX{}
commands:

\begin{verbatim}
\chapter{...}
\section{...}
\subsection{...}
\end{verbatim}

Chapter, section, and subsection numbers shall be generated automatically
by \LaTeX{}.

Students shall not manually type chapter or section numbers and shall not
reproduce heading styles using commands such as \verb|\textbf|,
\verb|\fontsize|, or \verb|\centering| instead of the appropriate
structural command.

Excessive subdivision should be avoided. In most cases, the hierarchy
Chapter--Section--Subsection is sufficient.

\vspace{0.3cm}

% ================================================================
% IX. ORGANIZATION OF THE MAIN CONTENT
% ================================================================

\noindent
\textbf{IX.\quad ORGANIZATION OF THE MAIN CONTENT}

\vspace{0.15cm}

The specific content of the report shall depend on the internship
organization and the assigned tasks. A typical structure may include:

\begin{itemize}[leftmargin=0.8cm, itemsep=2pt]
    \item \textbf{Chapter 1. Overview of the Internship Organization};
    \item \textbf{Chapter 2. Internship Activities and Assigned Tasks};
    \item \textbf{Chapter 3. Results, Experience, and Evaluation}.
\end{itemize}

These chapter titles are recommendations only and may be adapted to the
actual internship content.

The report shall reflect work actually performed, observed, analyzed, or
learned during the internship. Unrelated theoretical material should not
be included merely to increase report length.

\vspace{0.3cm}

% ================================================================
% X. TABLE OF CONTENTS
% ================================================================

\noindent
\textbf{X.\quad TABLE OF CONTENTS}

\vspace{0.15cm}

The Table of Contents shall be generated automatically using:

\begin{verbatim}
\tableofcontents
\end{verbatim}

Students shall not create or manually edit page numbers in the Table of
Contents.

The document shall be compiled sufficiently to ensure that all page
numbers and entries are correctly updated.

\vspace{0.3cm}

% ================================================================
% XI. LIST OF SYMBOLS AND ABBREVIATIONS
% ================================================================

\noindent
\textbf{XI.\quad LIST OF SYMBOLS AND ABBREVIATIONS}

\vspace{0.15cm}

Symbols and abbreviations that are used repeatedly in the report shall be
listed in the \textbf{List of Symbols and Abbreviations}.

The list may contain the following columns:

\begin{center}
\begin{tabular}{llll}
    No. & Symbol / Abbreviation & Unit & Definition
\end{tabular}
\end{center}

Only symbols and abbreviations actually used in the final report shall be
retained. Example entries contained in the template shall be removed if
they are not relevant to the student's report.

Each abbreviation should normally be defined when it first appears in
the main text.

\vspace{0.3cm}

% ================================================================
% XII. FIGURES, DIAGRAMS, AND PHOTOGRAPHS
% ================================================================

\noindent
\textbf{XII.\quad FIGURES, DIAGRAMS, AND PHOTOGRAPHS}

\vspace{0.15cm}

Figures shall be clear, legible, and directly relevant to the report
content.

Figures shall be inserted using the standard \verb|figure| environment.

Example:

\begin{verbatim}
\begin{figure}[H]
    \centering
    \includegraphics[width=0.8\textwidth]
                    {figures/company_structure.pdf}
    \caption{Organizational structure of the company}
    \label{fig:company-structure}
\end{figure}
\end{verbatim}

Figure captions shall be placed below the figure.

Each referenced figure shall have a unique \verb|\label| and shall be
referenced in the text using, for example:

\begin{verbatim}
Figure~\ref{fig:company-structure}
\end{verbatim}

Students shall not manually type figure numbers.

Figures reproduced or adapted from external sources shall include an
appropriate citation.

Photographs taken during the internship shall be relevant and sufficiently
clear. Students shall not include confidential, restricted, or proprietary
information without permission from the internship organization.

Vector graphics or high-resolution images are recommended whenever
possible.

\vspace{0.3cm}

% ================================================================
% XIII. TABLES
% ================================================================

\noindent
\textbf{XIII.\quad TABLES}

\vspace{0.15cm}

Tables shall be created using suitable \LaTeX{} environments such as
\verb|table|, \verb|tabular|, \verb|tabularx|, or \verb|longtable|.

Table captions shall be placed above the table.

Example:

\begin{verbatim}
\begin{table}[H]
    \centering
    \caption{Internship schedule}
    \label{tab:internship-schedule}

    \begin{tabular}{lll}
        \toprule
        Week & Task & Result \\
        \midrule
        1 & Orientation & Completed \\
        2 & System study & Completed \\
        \bottomrule
    \end{tabular}
\end{table}
\end{verbatim}

Tables shall be referenced using:

\begin{verbatim}
Table~\ref{tab:internship-schedule}
\end{verbatim}

Students shall not manually type table numbers.

Tables should be created directly in \LaTeX{} rather than inserted as
screenshots whenever practical.

\vspace{0.3cm}

% ================================================================
% XIV. EQUATIONS AND TECHNICAL QUANTITIES
% ================================================================

\noindent
\textbf{XIV.\quad EQUATIONS AND TECHNICAL QUANTITIES}

\vspace{0.15cm}

Mathematical expressions shall be typeset directly using \LaTeX{}.
Images or screenshots of equations shall not be used.

Short expressions may be written in inline mathematical mode, for example:

\begin{verbatim}
$R = 0.7\,\Omega$
\end{verbatim}

Important equations shall use a numbered equation environment.

Example:

\begin{verbatim}
\begin{equation}
    P = UI
    \label{eq:power}
\end{equation}
\end{verbatim}

Equations shall be referenced using:

\begin{verbatim}
Equation~\eqref{eq:power}
\end{verbatim}

Equation numbers shall not be manually entered.

SI units shall be used wherever applicable. The \verb|siunitx| package
should be used where appropriate.

Examples:

\begin{verbatim}
\SI{220}{\volt}
\SI{50}{\hertz}
\SI{2.5}{\ampere}
\end{verbatim}

Notation and units shall be consistent throughout the report.

\vspace{0.3cm}

% ================================================================
% XV. SOURCE CODE
% ================================================================

\noindent
\textbf{XV.\quad SOURCE CODE}

\vspace{0.15cm}

If the internship involves programming, PLCs, microcontrollers,
simulation, or software development, relevant source code may be included
in the report.

Source code shall be presented using the \verb|listings| environment or
another environment defined by the official template.

Example:

\begin{verbatim}
\begin{lstlisting}[language=st]
IF Start = TRUE THEN
    Motor := TRUE;
ELSE
    Motor := FALSE;
END_IF;
\end{lstlisting}
\end{verbatim}

Only code necessary for explaining the work should appear in the main
report. Long source-code listings should be moved to an appendix or
submitted as supplementary files if required.

Source code should not normally be inserted as screenshots.

\vspace{0.3cm}

% ================================================================
% XVI. ABSTRACT
% ================================================================

\noindent
\textbf{XVI.\quad ABSTRACT}

\vspace{0.15cm}

The Abstract shall briefly summarize:

\begin{itemize}[leftmargin=0.8cm, itemsep=2pt]
    \item the internship organization;
    \item the objectives of the internship;
    \item the main tasks performed;
    \item the main results, knowledge, and skills obtained.
\end{itemize}

The Abstract should be concise and should not reproduce the full content
of the report chapter by chapter.

\vspace{0.3cm}

% ================================================================
% XVII. REFERENCES AND CITATIONS
% ================================================================

\noindent
\textbf{XVII.\quad REFERENCES AND CITATIONS}

\vspace{0.15cm}

All books, journal papers, conference papers, standards, websites,
datasheets, manuals, technical documents, figures, tables, or other
external sources used in preparing the report shall be properly cited.

References shall be managed through the bibliography system provided with
the official template.

Bibliographic entries shall normally be stored in:

\begin{verbatim}
References.bib
\end{verbatim}

References shall be cited in the report using:

\begin{verbatim}
\cite{reference-key}
\end{verbatim}

Students shall not manually type reference numbers such as ``[1]'',
``[2]'', etc.

Titles of foreign-language publications shall remain in their original
language in the bibliography.

All borrowed ideas, data, figures, tables, equations, methods, or results
shall be properly acknowledged.

\vspace{0.3cm}

% ================================================================
% XVIII. CROSS-REFERENCING
% ================================================================

\noindent
\textbf{XVIII.\quad CROSS-REFERENCING}

\vspace{0.15cm}

Numbered elements shall be cross-referenced using the standard \LaTeX{}
mechanism.

Use the following commands where appropriate:

\begin{verbatim}
\label{...}
\ref{...}
\eqref{...}
\end{verbatim}

For example:

\begin{verbatim}
As shown in Figure~\ref{fig:system}, ...
\end{verbatim}

Students shall not manually type numbers such as ``Figure 2.3'' when
those numbers can be generated automatically.

\vspace{0.3cm}

% ================================================================
% XIX. PAGE NUMBERING
% ================================================================

\noindent
\textbf{XIX.\quad PAGE NUMBERING}

\vspace{0.15cm}

Cover pages and official evaluation forms at the beginning of the report
shall not display page numbers.

The main body of the report shall use Arabic page numbers beginning from
Chapter 1.

Page numbering shall be generated automatically by the official template.
Students shall not manually enter page numbers.

\vspace{0.3cm}

% ================================================================
% XX. PROJECT FILE ORGANIZATION
% ================================================================

\noindent
\textbf{XX.\quad PROJECT FILE ORGANIZATION}

\vspace{0.15cm}

Students should maintain a clear and organized \LaTeX{} project
structure. A recommended structure is:

\begin{verbatim}
InternshipReport/
|
|-- main.tex
|-- References.bib
|
|-- figures/
|   |-- company.pdf
|   |-- system.png
|   `-- ...
|
`-- chapters/
    |-- chapter1.tex
    |-- chapter2.tex
    `-- chapter3.tex
\end{verbatim}

Figures should normally be stored in the \verb|figures| directory
specified by the official template.

File and directory names should preferably contain Latin letters,
numbers, hyphens, or underscores.

Recommended examples include:

\begin{verbatim}
company_structure.pdf
motor_control.png
chapter_01.tex
\end{verbatim}

Spaces and unnecessary special characters in file names should be
avoided where possible.

\vspace{0.3cm}

% ================================================================
% XXI. DOCUMENT COMPILATION AND FINAL CHECK
% ================================================================

\noindent
\textbf{XXI.\quad DOCUMENT COMPILATION AND FINAL CHECK}

\vspace{0.15cm}

Before submission, the complete \LaTeX{} project shall compile
successfully and produce the final PDF report.

The final document shall not contain unresolved references such as:

\begin{verbatim}
??
[?]
\end{verbatim}

Before generating the final PDF, students shall check:

\begin{itemize}[leftmargin=0.8cm, itemsep=2pt]
    \item the Table of Contents;
    \item page numbering;
    \item chapter, section, and subsection numbering;
    \item figure numbering;
    \item table numbering;
    \item equation numbering;
    \item cross-references;
    \item citations;
    \item bibliography;
    \item font display;
    \item mathematical symbols;
    \item image quality;
    \item completeness of all required forms.
\end{itemize}

The document may be compiled using TeX Live, MiKTeX, or another
compatible \LaTeX{} environment. Regardless of the compilation
environment, the resulting PDF shall conform to the official Faculty
template.

\vspace{0.3cm}

% ================================================================
% XXII. CONTENT THAT STUDENTS MAY EDIT
% ================================================================

\noindent
\textbf{XXII.\quad CONTENT THAT STUDENTS MAY EDIT}

\vspace{0.15cm}

Students may edit or replace:

\begin{itemize}[leftmargin=0.8cm, itemsep=2pt]
    \item student name;
    \item class;
    \item internship organization;
    \item supervisor name;
    \item Abstract;
    \item List of Symbols and Abbreviations;
    \item chapter titles and chapter content;
    \item figures and photographs;
    \item tables;
    \item equations;
    \item citations and references.
\end{itemize}

Students shall not modify the main document formatting settings in the
preamble, including the document class, page geometry, line spacing,
heading formatting, table-of-contents formatting, font configuration,
page layout, or cover-page macros unless specifically instructed.

\vspace{0.3cm}

% ================================================================
% XXIII. FINAL SUBMISSION REQUIREMENTS
% ================================================================

\noindent
\textbf{XXIII.\quad FINAL SUBMISSION REQUIREMENTS}

\vspace{0.15cm}

Before submission, the final report shall:

\begin{itemize}[leftmargin=0.8cm, itemsep=2pt]
    \item compile without errors that affect the report content or presentation;
    \item contain all required forms and report sections;
    \item contain correct student, class, internship organization, and supervisor information;
    \item contain an updated Table of Contents, List of Figures, and List of Tables;
    \item contain no unused example content from the template;
    \item contain no placeholder text such as \textit{Title of Chapter 1},
    \textit{Section 1.1}, or other sample entries;
    \item contain all required citations and references;
    \item contain no unresolved cross-references.
\end{itemize}

If the Faculty requires submission of the \LaTeX{} source, the complete
project shall include all \verb|.tex| files, the bibliography database,
figures, and any other files required to reproduce the submitted PDF.

\vspace{0.4cm}

\noindent
\textbf{Important:}
Before submitting the final internship report, students shall remove or
comment out the following command in the main document:

\begin{verbatim}
\input{Regulations.tex}
\end{verbatim}

For example:

\begin{verbatim}
% \input{Regulations.tex}
\end{verbatim}

This removes the guidance pages from the final submitted report while
preserving the regulations file in the \LaTeX{} project for reference.

} % End 13pt content

\clearpage
|.
Students shall remove or comment out this command before submitting the
final internship report.

\vspace{0.4cm}

% ================================================================
% I. GENERAL REQUIREMENTS
% ================================================================

\noindent
\textbf{I.\quad GENERAL REQUIREMENTS}

\vspace{0.15cm}

The Graduation Internship Report shall be prepared using the official
\LaTeX{} template provided by the Faculty. Students shall preserve the
predefined page layout, typography, line spacing, heading styles,
numbering system, figure and table formatting, and bibliography settings
unless otherwise instructed by the Faculty or the supervisor.

The report shall be written in English.

Students are responsible for the academic and technical content of the
report. The official \LaTeX{} template is responsible for standardized
document formatting.

\vspace{0.3cm}

% ================================================================
% II. REPORT STRUCTURE
% ================================================================

\noindent
\textbf{II.\quad REPORT STRUCTURE}

\vspace{0.15cm}

The report shall normally be arranged in the following order:

\begin{enumerate}[label=\arabic*., leftmargin=1.2cm, itemsep=2pt]
    \item Outer Cover;
    \item Comments of the Internship Host Organization;
    \item Supervisor's Evaluation and Graduation Internship Defense Results;
    \item Table of Contents;
    \item List of Symbols and Abbreviations, if applicable;
    \item List of Figures;
    \item List of Tables;
    \item Abstract;
    \item Main Chapters of the Internship Report;
    \item References;
    \item Appendices, if applicable.
\end{enumerate}

The preliminary pages shall not be numbered as chapters. The main report
content shall be organized into consecutively numbered chapters beginning
with Chapter 1.

\vspace{0.3cm}

% ================================================================
% III. PAGE SIZE AND MARGINS
% ================================================================

\noindent
\textbf{III.\quad PAGE SIZE AND MARGINS}

\vspace{0.15cm}

The report shall be formatted on A4 paper with the following margins:

\begin{center}
\begin{tabular}{ll}
    Left margin:   & 3.5 cm \\
    Right margin:  & 2.0 cm \\
    Top margin:    & 2.5 cm \\
    Bottom margin: & 2.5 cm
\end{tabular}
\end{center}

These settings are predefined in the official template and shall not be
modified by students.

\vspace{0.3cm}

% ================================================================
% IV. FONT, FONT SIZE, AND LINE SPACING
% ================================================================

\noindent
\textbf{IV.\quad FONT, FONT SIZE, AND LINE SPACING}

\vspace{0.15cm}

The main text shall use a 13 pt font size with 1.3-line spacing.
A Times-family typeface shall be used in accordance with the official
template.

The text shall be fully justified.

Students shall not use repeated spaces, repeated manual line breaks, or
multiple blank lines to position text. Paragraph and heading spacing shall
be controlled by the template.

\vspace{0.3cm}

% ================================================================
% V. OUTER COVER
% ================================================================

\noindent
\textbf{V.\quad OUTER COVER}

\vspace{0.15cm}

The outer cover shall follow the official format provided in the template
and shall contain the following information:

\begin{itemize}[leftmargin=0.8cm, itemsep=2pt]
    \item Thai Nguyen University of Technology;
    \item Faculty of International Training;
    \item Engineering and Technology Department;
    \item University logo and Faculty logo;
    \item the title \textbf{INTERNSHIP REPORT};
    \item Internship Organization;
    \item Author;
    \item Class;
    \item Supervisor;
    \item place and year of completion.
\end{itemize}

Students shall only replace the required information fields and shall not
modify the cover-page layout unless instructed otherwise.

\vspace{0.3cm}

% ================================================================
% VI. COMMENTS OF THE INTERNSHIP HOST ORGANIZATION
% ================================================================

\noindent
\textbf{VI.\quad COMMENTS OF THE INTERNSHIP HOST ORGANIZATION}

\vspace{0.15cm}

The page entitled \textbf{COMMENTS OF THE INTERNSHIP HOST ORGANIZATION}
is reserved for the internship host organization.

It includes:

\begin{itemize}[leftmargin=0.8cm, itemsep=2pt]
    \item the student's full name;
    \item class;
    \item comments from the host organization;
    \item score;
    \item date;
    \item signature and official seal.
\end{itemize}

Students shall not alter or complete official evaluation fields unless
specifically authorized.

\vspace{0.3cm}

% ================================================================
% VII. SUPERVISOR'S EVALUATION AND DEFENSE RESULTS
% ================================================================

\noindent
\textbf{VII.\quad SUPERVISOR'S EVALUATION AND DEFENSE RESULTS}

\vspace{0.15cm}

The section entitled
\textbf{SUPERVISOR'S EVALUATION AND GRADUATION INTERNSHIP DEFENSE RESULTS}
is reserved for the supervisor and defense committee.

The evaluation may include:

\begin{itemize}[leftmargin=0.8cm, itemsep=2pt]
    \item Internship Progress and Student's Attitude;
    \item Report Content;
    \item Report Presentation;
    \item Other Comments;
    \item Supervisor's Opinion;
    \item Defense Result;
    \item Score.
\end{itemize}

Students shall not enter evaluation scores or defense results before they
are officially determined.

\vspace{0.3cm}

% ================================================================
% VIII. CHAPTERS, SECTIONS, AND SUBSECTIONS
% ================================================================

\noindent
\textbf{VIII.\quad CHAPTERS, SECTIONS, AND SUBSECTIONS}

\vspace{0.15cm}

The report structure shall be generated using the standard \LaTeX{}
commands:

\begin{verbatim}
\chapter{...}
\section{...}
\subsection{...}
\end{verbatim}

Chapter, section, and subsection numbers shall be generated automatically
by \LaTeX{}.

Students shall not manually type chapter or section numbers and shall not
reproduce heading styles using commands such as \verb|\textbf|,
\verb|\fontsize|, or \verb|\centering| instead of the appropriate
structural command.

Excessive subdivision should be avoided. In most cases, the hierarchy
Chapter--Section--Subsection is sufficient.

\vspace{0.3cm}

% ================================================================
% IX. ORGANIZATION OF THE MAIN CONTENT
% ================================================================

\noindent
\textbf{IX.\quad ORGANIZATION OF THE MAIN CONTENT}

\vspace{0.15cm}

The specific content of the report shall depend on the internship
organization and the assigned tasks. A typical structure may include:

\begin{itemize}[leftmargin=0.8cm, itemsep=2pt]
    \item \textbf{Chapter 1. Overview of the Internship Organization};
    \item \textbf{Chapter 2. Internship Activities and Assigned Tasks};
    \item \textbf{Chapter 3. Results, Experience, and Evaluation}.
\end{itemize}

These chapter titles are recommendations only and may be adapted to the
actual internship content.

The report shall reflect work actually performed, observed, analyzed, or
learned during the internship. Unrelated theoretical material should not
be included merely to increase report length.

\vspace{0.3cm}

% ================================================================
% X. TABLE OF CONTENTS
% ================================================================

\noindent
\textbf{X.\quad TABLE OF CONTENTS}

\vspace{0.15cm}

The Table of Contents shall be generated automatically using:

\begin{verbatim}
\tableofcontents
\end{verbatim}

Students shall not create or manually edit page numbers in the Table of
Contents.

The document shall be compiled sufficiently to ensure that all page
numbers and entries are correctly updated.

\vspace{0.3cm}

% ================================================================
% XI. LIST OF SYMBOLS AND ABBREVIATIONS
% ================================================================

\noindent
\textbf{XI.\quad LIST OF SYMBOLS AND ABBREVIATIONS}

\vspace{0.15cm}

Symbols and abbreviations that are used repeatedly in the report shall be
listed in the \textbf{List of Symbols and Abbreviations}.

The list may contain the following columns:

\begin{center}
\begin{tabular}{llll}
    No. & Symbol / Abbreviation & Unit & Definition
\end{tabular}
\end{center}

Only symbols and abbreviations actually used in the final report shall be
retained. Example entries contained in the template shall be removed if
they are not relevant to the student's report.

Each abbreviation should normally be defined when it first appears in
the main text.

\vspace{0.3cm}

% ================================================================
% XII. FIGURES, DIAGRAMS, AND PHOTOGRAPHS
% ================================================================

\noindent
\textbf{XII.\quad FIGURES, DIAGRAMS, AND PHOTOGRAPHS}

\vspace{0.15cm}

Figures shall be clear, legible, and directly relevant to the report
content.

Figures shall be inserted using the standard \verb|figure| environment.

Example:

\begin{verbatim}
\begin{figure}[H]
    \centering
    \includegraphics[width=0.8\textwidth]
                    {figures/company_structure.pdf}
    \caption{Organizational structure of the company}
    \label{fig:company-structure}
\end{figure}
\end{verbatim}

Figure captions shall be placed below the figure.

Each referenced figure shall have a unique \verb|\label| and shall be
referenced in the text using, for example:

\begin{verbatim}
Figure~\ref{fig:company-structure}
\end{verbatim}

Students shall not manually type figure numbers.

Figures reproduced or adapted from external sources shall include an
appropriate citation.

Photographs taken during the internship shall be relevant and sufficiently
clear. Students shall not include confidential, restricted, or proprietary
information without permission from the internship organization.

Vector graphics or high-resolution images are recommended whenever
possible.

\vspace{0.3cm}

% ================================================================
% XIII. TABLES
% ================================================================

\noindent
\textbf{XIII.\quad TABLES}

\vspace{0.15cm}

Tables shall be created using suitable \LaTeX{} environments such as
\verb|table|, \verb|tabular|, \verb|tabularx|, or \verb|longtable|.

Table captions shall be placed above the table.

Example:

\begin{verbatim}
\begin{table}[H]
    \centering
    \caption{Internship schedule}
    \label{tab:internship-schedule}

    \begin{tabular}{lll}
        \toprule
        Week & Task & Result \\
        \midrule
        1 & Orientation & Completed \\
        2 & System study & Completed \\
        \bottomrule
    \end{tabular}
\end{table}
\end{verbatim}

Tables shall be referenced using:

\begin{verbatim}
Table~\ref{tab:internship-schedule}
\end{verbatim}

Students shall not manually type table numbers.

Tables should be created directly in \LaTeX{} rather than inserted as
screenshots whenever practical.

\vspace{0.3cm}

% ================================================================
% XIV. EQUATIONS AND TECHNICAL QUANTITIES
% ================================================================

\noindent
\textbf{XIV.\quad EQUATIONS AND TECHNICAL QUANTITIES}

\vspace{0.15cm}

Mathematical expressions shall be typeset directly using \LaTeX{}.
Images or screenshots of equations shall not be used.

Short expressions may be written in inline mathematical mode, for example:

\begin{verbatim}
$R = 0.7\,\Omega$
\end{verbatim}

Important equations shall use a numbered equation environment.

Example:

\begin{verbatim}
\begin{equation}
    P = UI
    \label{eq:power}
\end{equation}
\end{verbatim}

Equations shall be referenced using:

\begin{verbatim}
Equation~\eqref{eq:power}
\end{verbatim}

Equation numbers shall not be manually entered.

SI units shall be used wherever applicable. The \verb|siunitx| package
should be used where appropriate.

Examples:

\begin{verbatim}
\SI{220}{\volt}
\SI{50}{\hertz}
\SI{2.5}{\ampere}
\end{verbatim}

Notation and units shall be consistent throughout the report.

\vspace{0.3cm}

% ================================================================
% XV. SOURCE CODE
% ================================================================

\noindent
\textbf{XV.\quad SOURCE CODE}

\vspace{0.15cm}

If the internship involves programming, PLCs, microcontrollers,
simulation, or software development, relevant source code may be included
in the report.

Source code shall be presented using the \verb|listings| environment or
another environment defined by the official template.

Example:

\begin{verbatim}
\begin{lstlisting}[language=st]
IF Start = TRUE THEN
    Motor := TRUE;
ELSE
    Motor := FALSE;
END_IF;
\end{lstlisting}
\end{verbatim}

Only code necessary for explaining the work should appear in the main
report. Long source-code listings should be moved to an appendix or
submitted as supplementary files if required.

Source code should not normally be inserted as screenshots.

\vspace{0.3cm}

% ================================================================
% XVI. ABSTRACT
% ================================================================

\noindent
\textbf{XVI.\quad ABSTRACT}

\vspace{0.15cm}

The Abstract shall briefly summarize:

\begin{itemize}[leftmargin=0.8cm, itemsep=2pt]
    \item the internship organization;
    \item the objectives of the internship;
    \item the main tasks performed;
    \item the main results, knowledge, and skills obtained.
\end{itemize}

The Abstract should be concise and should not reproduce the full content
of the report chapter by chapter.

\vspace{0.3cm}

% ================================================================
% XVII. REFERENCES AND CITATIONS
% ================================================================

\noindent
\textbf{XVII.\quad REFERENCES AND CITATIONS}

\vspace{0.15cm}

All books, journal papers, conference papers, standards, websites,
datasheets, manuals, technical documents, figures, tables, or other
external sources used in preparing the report shall be properly cited.

References shall be managed through the bibliography system provided with
the official template.

Bibliographic entries shall normally be stored in:

\begin{verbatim}
References.bib
\end{verbatim}

References shall be cited in the report using:

\begin{verbatim}
\cite{reference-key}
\end{verbatim}

Students shall not manually type reference numbers such as ``[1]'',
``[2]'', etc.

Titles of foreign-language publications shall remain in their original
language in the bibliography.

All borrowed ideas, data, figures, tables, equations, methods, or results
shall be properly acknowledged.

\vspace{0.3cm}

% ================================================================
% XVIII. CROSS-REFERENCING
% ================================================================

\noindent
\textbf{XVIII.\quad CROSS-REFERENCING}

\vspace{0.15cm}

Numbered elements shall be cross-referenced using the standard \LaTeX{}
mechanism.

Use the following commands where appropriate:

\begin{verbatim}
\label{...}
\ref{...}
\eqref{...}
\end{verbatim}

For example:

\begin{verbatim}
As shown in Figure~\ref{fig:system}, ...
\end{verbatim}

Students shall not manually type numbers such as ``Figure 2.3'' when
those numbers can be generated automatically.

\vspace{0.3cm}

% ================================================================
% XIX. PAGE NUMBERING
% ================================================================

\noindent
\textbf{XIX.\quad PAGE NUMBERING}

\vspace{0.15cm}

Cover pages and official evaluation forms at the beginning of the report
shall not display page numbers.

The main body of the report shall use Arabic page numbers beginning from
Chapter 1.

Page numbering shall be generated automatically by the official template.
Students shall not manually enter page numbers.

\vspace{0.3cm}

% ================================================================
% XX. PROJECT FILE ORGANIZATION
% ================================================================

\noindent
\textbf{XX.\quad PROJECT FILE ORGANIZATION}

\vspace{0.15cm}

Students should maintain a clear and organized \LaTeX{} project
structure. A recommended structure is:

\begin{verbatim}
InternshipReport/
|
|-- main.tex
|-- References.bib
|
|-- figures/
|   |-- company.pdf
|   |-- system.png
|   `-- ...
|
`-- chapters/
    |-- chapter1.tex
    |-- chapter2.tex
    `-- chapter3.tex
\end{verbatim}

Figures should normally be stored in the \verb|figures| directory
specified by the official template.

File and directory names should preferably contain Latin letters,
numbers, hyphens, or underscores.

Recommended examples include:

\begin{verbatim}
company_structure.pdf
motor_control.png
chapter_01.tex
\end{verbatim}

Spaces and unnecessary special characters in file names should be
avoided where possible.

\vspace{0.3cm}

% ================================================================
% XXI. DOCUMENT COMPILATION AND FINAL CHECK
% ================================================================

\noindent
\textbf{XXI.\quad DOCUMENT COMPILATION AND FINAL CHECK}

\vspace{0.15cm}

Before submission, the complete \LaTeX{} project shall compile
successfully and produce the final PDF report.

The final document shall not contain unresolved references such as:

\begin{verbatim}
??
[?]
\end{verbatim}

Before generating the final PDF, students shall check:

\begin{itemize}[leftmargin=0.8cm, itemsep=2pt]
    \item the Table of Contents;
    \item page numbering;
    \item chapter, section, and subsection numbering;
    \item figure numbering;
    \item table numbering;
    \item equation numbering;
    \item cross-references;
    \item citations;
    \item bibliography;
    \item font display;
    \item mathematical symbols;
    \item image quality;
    \item completeness of all required forms.
\end{itemize}

The document may be compiled using TeX Live, MiKTeX, or another
compatible \LaTeX{} environment. Regardless of the compilation
environment, the resulting PDF shall conform to the official Faculty
template.

\vspace{0.3cm}

% ================================================================
% XXII. CONTENT THAT STUDENTS MAY EDIT
% ================================================================

\noindent
\textbf{XXII.\quad CONTENT THAT STUDENTS MAY EDIT}

\vspace{0.15cm}

Students may edit or replace:

\begin{itemize}[leftmargin=0.8cm, itemsep=2pt]
    \item student name;
    \item class;
    \item internship organization;
    \item supervisor name;
    \item Abstract;
    \item List of Symbols and Abbreviations;
    \item chapter titles and chapter content;
    \item figures and photographs;
    \item tables;
    \item equations;
    \item citations and references.
\end{itemize}

Students shall not modify the main document formatting settings in the
preamble, including the document class, page geometry, line spacing,
heading formatting, table-of-contents formatting, font configuration,
page layout, or cover-page macros unless specifically instructed.

\vspace{0.3cm}

% ================================================================
% XXIII. FINAL SUBMISSION REQUIREMENTS
% ================================================================

\noindent
\textbf{XXIII.\quad FINAL SUBMISSION REQUIREMENTS}

\vspace{0.15cm}

Before submission, the final report shall:

\begin{itemize}[leftmargin=0.8cm, itemsep=2pt]
    \item compile without errors that affect the report content or presentation;
    \item contain all required forms and report sections;
    \item contain correct student, class, internship organization, and supervisor information;
    \item contain an updated Table of Contents, List of Figures, and List of Tables;
    \item contain no unused example content from the template;
    \item contain no placeholder text such as \textit{Title of Chapter 1},
    \textit{Section 1.1}, or other sample entries;
    \item contain all required citations and references;
    \item contain no unresolved cross-references.
\end{itemize}

If the Faculty requires submission of the \LaTeX{} source, the complete
project shall include all \verb|.tex| files, the bibliography database,
figures, and any other files required to reproduce the submitted PDF.

\vspace{0.4cm}

\noindent
\textbf{Important:}
Before submitting the final internship report, students shall remove or
comment out the following command in the main document:

\begin{verbatim}

% =================================================================
% REGULATIONS FOR PREPARING THE GRADUATION INTERNSHIP REPORT
% USING LATEX
% =================================================================
%
% This file is intended to be included from the main document with:
%
%     \input{Regulations.tex}
%
% IMPORTANT:
% This section is provided for guidance only. Students shall remove
% or comment out the \input{Regulations.tex} command before submitting
% the final internship report.
%
% =================================================================

\clearpage
\thispagestyle{empty}

\begin{center}
    {\fontsize{16pt}{20pt}\selectfont\bfseries
    SPECIFIC REGULATIONS FOR PREPARING\\
    GRADUATION INTERNSHIP REPORTS USING \LaTeX}
\end{center}

\vspace{0.5cm}

{\fontsize{13pt}{17pt}\selectfont
\setstretch{1.3}
\justifying

\noindent
\textbf{Note:}
This section is provided for guidance only and is included in the main
document using the command \verb|\input{Regulations.tex}|.
Students shall remove or comment out this command before submitting the
final internship report.

\vspace{0.4cm}

% ================================================================
% I. GENERAL REQUIREMENTS
% ================================================================

\noindent
\textbf{I.\quad GENERAL REQUIREMENTS}

\vspace{0.15cm}

The Graduation Internship Report shall be prepared using the official
\LaTeX{} template provided by the Faculty. Students shall preserve the
predefined page layout, typography, line spacing, heading styles,
numbering system, figure and table formatting, and bibliography settings
unless otherwise instructed by the Faculty or the supervisor.

The report shall be written in English.

Students are responsible for the academic and technical content of the
report. The official \LaTeX{} template is responsible for standardized
document formatting.

\vspace{0.3cm}

% ================================================================
% II. REPORT STRUCTURE
% ================================================================

\noindent
\textbf{II.\quad REPORT STRUCTURE}

\vspace{0.15cm}

The report shall normally be arranged in the following order:

\begin{enumerate}[label=\arabic*., leftmargin=1.2cm, itemsep=2pt]
    \item Outer Cover;
    \item Comments of the Internship Host Organization;
    \item Supervisor's Evaluation and Graduation Internship Defense Results;
    \item Table of Contents;
    \item List of Symbols and Abbreviations, if applicable;
    \item List of Figures;
    \item List of Tables;
    \item Abstract;
    \item Main Chapters of the Internship Report;
    \item References;
    \item Appendices, if applicable.
\end{enumerate}

The preliminary pages shall not be numbered as chapters. The main report
content shall be organized into consecutively numbered chapters beginning
with Chapter 1.

\vspace{0.3cm}

% ================================================================
% III. PAGE SIZE AND MARGINS
% ================================================================

\noindent
\textbf{III.\quad PAGE SIZE AND MARGINS}

\vspace{0.15cm}

The report shall be formatted on A4 paper with the following margins:

\begin{center}
\begin{tabular}{ll}
    Left margin:   & 3.5 cm \\
    Right margin:  & 2.0 cm \\
    Top margin:    & 2.5 cm \\
    Bottom margin: & 2.5 cm
\end{tabular}
\end{center}

These settings are predefined in the official template and shall not be
modified by students.

\vspace{0.3cm}

% ================================================================
% IV. FONT, FONT SIZE, AND LINE SPACING
% ================================================================

\noindent
\textbf{IV.\quad FONT, FONT SIZE, AND LINE SPACING}

\vspace{0.15cm}

The main text shall use a 13 pt font size with 1.3-line spacing.
A Times-family typeface shall be used in accordance with the official
template.

The text shall be fully justified.

Students shall not use repeated spaces, repeated manual line breaks, or
multiple blank lines to position text. Paragraph and heading spacing shall
be controlled by the template.

\vspace{0.3cm}

% ================================================================
% V. OUTER COVER
% ================================================================

\noindent
\textbf{V.\quad OUTER COVER}

\vspace{0.15cm}

The outer cover shall follow the official format provided in the template
and shall contain the following information:

\begin{itemize}[leftmargin=0.8cm, itemsep=2pt]
    \item Thai Nguyen University of Technology;
    \item Faculty of International Training;
    \item Engineering and Technology Department;
    \item University logo and Faculty logo;
    \item the title \textbf{INTERNSHIP REPORT};
    \item Internship Organization;
    \item Author;
    \item Class;
    \item Supervisor;
    \item place and year of completion.
\end{itemize}

Students shall only replace the required information fields and shall not
modify the cover-page layout unless instructed otherwise.

\vspace{0.3cm}

% ================================================================
% VI. COMMENTS OF THE INTERNSHIP HOST ORGANIZATION
% ================================================================

\noindent
\textbf{VI.\quad COMMENTS OF THE INTERNSHIP HOST ORGANIZATION}

\vspace{0.15cm}

The page entitled \textbf{COMMENTS OF THE INTERNSHIP HOST ORGANIZATION}
is reserved for the internship host organization.

It includes:

\begin{itemize}[leftmargin=0.8cm, itemsep=2pt]
    \item the student's full name;
    \item class;
    \item comments from the host organization;
    \item score;
    \item date;
    \item signature and official seal.
\end{itemize}

Students shall not alter or complete official evaluation fields unless
specifically authorized.

\vspace{0.3cm}

% ================================================================
% VII. SUPERVISOR'S EVALUATION AND DEFENSE RESULTS
% ================================================================

\noindent
\textbf{VII.\quad SUPERVISOR'S EVALUATION AND DEFENSE RESULTS}

\vspace{0.15cm}

The section entitled
\textbf{SUPERVISOR'S EVALUATION AND GRADUATION INTERNSHIP DEFENSE RESULTS}
is reserved for the supervisor and defense committee.

The evaluation may include:

\begin{itemize}[leftmargin=0.8cm, itemsep=2pt]
    \item Internship Progress and Student's Attitude;
    \item Report Content;
    \item Report Presentation;
    \item Other Comments;
    \item Supervisor's Opinion;
    \item Defense Result;
    \item Score.
\end{itemize}

Students shall not enter evaluation scores or defense results before they
are officially determined.

\vspace{0.3cm}

% ================================================================
% VIII. CHAPTERS, SECTIONS, AND SUBSECTIONS
% ================================================================

\noindent
\textbf{VIII.\quad CHAPTERS, SECTIONS, AND SUBSECTIONS}

\vspace{0.15cm}

The report structure shall be generated using the standard \LaTeX{}
commands:

\begin{verbatim}
\chapter{...}
\section{...}
\subsection{...}
\end{verbatim}

Chapter, section, and subsection numbers shall be generated automatically
by \LaTeX{}.

Students shall not manually type chapter or section numbers and shall not
reproduce heading styles using commands such as \verb|\textbf|,
\verb|\fontsize|, or \verb|\centering| instead of the appropriate
structural command.

Excessive subdivision should be avoided. In most cases, the hierarchy
Chapter--Section--Subsection is sufficient.

\vspace{0.3cm}

% ================================================================
% IX. ORGANIZATION OF THE MAIN CONTENT
% ================================================================

\noindent
\textbf{IX.\quad ORGANIZATION OF THE MAIN CONTENT}

\vspace{0.15cm}

The specific content of the report shall depend on the internship
organization and the assigned tasks. A typical structure may include:

\begin{itemize}[leftmargin=0.8cm, itemsep=2pt]
    \item \textbf{Chapter 1. Overview of the Internship Organization};
    \item \textbf{Chapter 2. Internship Activities and Assigned Tasks};
    \item \textbf{Chapter 3. Results, Experience, and Evaluation}.
\end{itemize}

These chapter titles are recommendations only and may be adapted to the
actual internship content.

The report shall reflect work actually performed, observed, analyzed, or
learned during the internship. Unrelated theoretical material should not
be included merely to increase report length.

\vspace{0.3cm}

% ================================================================
% X. TABLE OF CONTENTS
% ================================================================

\noindent
\textbf{X.\quad TABLE OF CONTENTS}

\vspace{0.15cm}

The Table of Contents shall be generated automatically using:

\begin{verbatim}
\tableofcontents
\end{verbatim}

Students shall not create or manually edit page numbers in the Table of
Contents.

The document shall be compiled sufficiently to ensure that all page
numbers and entries are correctly updated.

\vspace{0.3cm}

% ================================================================
% XI. LIST OF SYMBOLS AND ABBREVIATIONS
% ================================================================

\noindent
\textbf{XI.\quad LIST OF SYMBOLS AND ABBREVIATIONS}

\vspace{0.15cm}

Symbols and abbreviations that are used repeatedly in the report shall be
listed in the \textbf{List of Symbols and Abbreviations}.

The list may contain the following columns:

\begin{center}
\begin{tabular}{llll}
    No. & Symbol / Abbreviation & Unit & Definition
\end{tabular}
\end{center}

Only symbols and abbreviations actually used in the final report shall be
retained. Example entries contained in the template shall be removed if
they are not relevant to the student's report.

Each abbreviation should normally be defined when it first appears in
the main text.

\vspace{0.3cm}

% ================================================================
% XII. FIGURES, DIAGRAMS, AND PHOTOGRAPHS
% ================================================================

\noindent
\textbf{XII.\quad FIGURES, DIAGRAMS, AND PHOTOGRAPHS}

\vspace{0.15cm}

Figures shall be clear, legible, and directly relevant to the report
content.

Figures shall be inserted using the standard \verb|figure| environment.

Example:

\begin{verbatim}
\begin{figure}[H]
    \centering
    \includegraphics[width=0.8\textwidth]
                    {figures/company_structure.pdf}
    \caption{Organizational structure of the company}
    \label{fig:company-structure}
\end{figure}
\end{verbatim}

Figure captions shall be placed below the figure.

Each referenced figure shall have a unique \verb|\label| and shall be
referenced in the text using, for example:

\begin{verbatim}
Figure~\ref{fig:company-structure}
\end{verbatim}

Students shall not manually type figure numbers.

Figures reproduced or adapted from external sources shall include an
appropriate citation.

Photographs taken during the internship shall be relevant and sufficiently
clear. Students shall not include confidential, restricted, or proprietary
information without permission from the internship organization.

Vector graphics or high-resolution images are recommended whenever
possible.

\vspace{0.3cm}

% ================================================================
% XIII. TABLES
% ================================================================

\noindent
\textbf{XIII.\quad TABLES}

\vspace{0.15cm}

Tables shall be created using suitable \LaTeX{} environments such as
\verb|table|, \verb|tabular|, \verb|tabularx|, or \verb|longtable|.

Table captions shall be placed above the table.

Example:

\begin{verbatim}
\begin{table}[H]
    \centering
    \caption{Internship schedule}
    \label{tab:internship-schedule}

    \begin{tabular}{lll}
        \toprule
        Week & Task & Result \\
        \midrule
        1 & Orientation & Completed \\
        2 & System study & Completed \\
        \bottomrule
    \end{tabular}
\end{table}
\end{verbatim}

Tables shall be referenced using:

\begin{verbatim}
Table~\ref{tab:internship-schedule}
\end{verbatim}

Students shall not manually type table numbers.

Tables should be created directly in \LaTeX{} rather than inserted as
screenshots whenever practical.

\vspace{0.3cm}

% ================================================================
% XIV. EQUATIONS AND TECHNICAL QUANTITIES
% ================================================================

\noindent
\textbf{XIV.\quad EQUATIONS AND TECHNICAL QUANTITIES}

\vspace{0.15cm}

Mathematical expressions shall be typeset directly using \LaTeX{}.
Images or screenshots of equations shall not be used.

Short expressions may be written in inline mathematical mode, for example:

\begin{verbatim}
$R = 0.7\,\Omega$
\end{verbatim}

Important equations shall use a numbered equation environment.

Example:

\begin{verbatim}
\begin{equation}
    P = UI
    \label{eq:power}
\end{equation}
\end{verbatim}

Equations shall be referenced using:

\begin{verbatim}
Equation~\eqref{eq:power}
\end{verbatim}

Equation numbers shall not be manually entered.

SI units shall be used wherever applicable. The \verb|siunitx| package
should be used where appropriate.

Examples:

\begin{verbatim}
\SI{220}{\volt}
\SI{50}{\hertz}
\SI{2.5}{\ampere}
\end{verbatim}

Notation and units shall be consistent throughout the report.

\vspace{0.3cm}

% ================================================================
% XV. SOURCE CODE
% ================================================================

\noindent
\textbf{XV.\quad SOURCE CODE}

\vspace{0.15cm}

If the internship involves programming, PLCs, microcontrollers,
simulation, or software development, relevant source code may be included
in the report.

Source code shall be presented using the \verb|listings| environment or
another environment defined by the official template.

Example:

\begin{verbatim}
\begin{lstlisting}[language=st]
IF Start = TRUE THEN
    Motor := TRUE;
ELSE
    Motor := FALSE;
END_IF;
\end{lstlisting}
\end{verbatim}

Only code necessary for explaining the work should appear in the main
report. Long source-code listings should be moved to an appendix or
submitted as supplementary files if required.

Source code should not normally be inserted as screenshots.

\vspace{0.3cm}

% ================================================================
% XVI. ABSTRACT
% ================================================================

\noindent
\textbf{XVI.\quad ABSTRACT}

\vspace{0.15cm}

The Abstract shall briefly summarize:

\begin{itemize}[leftmargin=0.8cm, itemsep=2pt]
    \item the internship organization;
    \item the objectives of the internship;
    \item the main tasks performed;
    \item the main results, knowledge, and skills obtained.
\end{itemize}

The Abstract should be concise and should not reproduce the full content
of the report chapter by chapter.

\vspace{0.3cm}

% ================================================================
% XVII. REFERENCES AND CITATIONS
% ================================================================

\noindent
\textbf{XVII.\quad REFERENCES AND CITATIONS}

\vspace{0.15cm}

All books, journal papers, conference papers, standards, websites,
datasheets, manuals, technical documents, figures, tables, or other
external sources used in preparing the report shall be properly cited.

References shall be managed through the bibliography system provided with
the official template.

Bibliographic entries shall normally be stored in:

\begin{verbatim}
References.bib
\end{verbatim}

References shall be cited in the report using:

\begin{verbatim}
\cite{reference-key}
\end{verbatim}

Students shall not manually type reference numbers such as ``[1]'',
``[2]'', etc.

Titles of foreign-language publications shall remain in their original
language in the bibliography.

All borrowed ideas, data, figures, tables, equations, methods, or results
shall be properly acknowledged.

\vspace{0.3cm}

% ================================================================
% XVIII. CROSS-REFERENCING
% ================================================================

\noindent
\textbf{XVIII.\quad CROSS-REFERENCING}

\vspace{0.15cm}

Numbered elements shall be cross-referenced using the standard \LaTeX{}
mechanism.

Use the following commands where appropriate:

\begin{verbatim}
\label{...}
\ref{...}
\eqref{...}
\end{verbatim}

For example:

\begin{verbatim}
As shown in Figure~\ref{fig:system}, ...
\end{verbatim}

Students shall not manually type numbers such as ``Figure 2.3'' when
those numbers can be generated automatically.

\vspace{0.3cm}

% ================================================================
% XIX. PAGE NUMBERING
% ================================================================

\noindent
\textbf{XIX.\quad PAGE NUMBERING}

\vspace{0.15cm}

Cover pages and official evaluation forms at the beginning of the report
shall not display page numbers.

The main body of the report shall use Arabic page numbers beginning from
Chapter 1.

Page numbering shall be generated automatically by the official template.
Students shall not manually enter page numbers.

\vspace{0.3cm}

% ================================================================
% XX. PROJECT FILE ORGANIZATION
% ================================================================

\noindent
\textbf{XX.\quad PROJECT FILE ORGANIZATION}

\vspace{0.15cm}

Students should maintain a clear and organized \LaTeX{} project
structure. A recommended structure is:

\begin{verbatim}
InternshipReport/
|
|-- main.tex
|-- References.bib
|
|-- figures/
|   |-- company.pdf
|   |-- system.png
|   `-- ...
|
`-- chapters/
    |-- chapter1.tex
    |-- chapter2.tex
    `-- chapter3.tex
\end{verbatim}

Figures should normally be stored in the \verb|figures| directory
specified by the official template.

File and directory names should preferably contain Latin letters,
numbers, hyphens, or underscores.

Recommended examples include:

\begin{verbatim}
company_structure.pdf
motor_control.png
chapter_01.tex
\end{verbatim}

Spaces and unnecessary special characters in file names should be
avoided where possible.

\vspace{0.3cm}

% ================================================================
% XXI. DOCUMENT COMPILATION AND FINAL CHECK
% ================================================================

\noindent
\textbf{XXI.\quad DOCUMENT COMPILATION AND FINAL CHECK}

\vspace{0.15cm}

Before submission, the complete \LaTeX{} project shall compile
successfully and produce the final PDF report.

The final document shall not contain unresolved references such as:

\begin{verbatim}
??
[?]
\end{verbatim}

Before generating the final PDF, students shall check:

\begin{itemize}[leftmargin=0.8cm, itemsep=2pt]
    \item the Table of Contents;
    \item page numbering;
    \item chapter, section, and subsection numbering;
    \item figure numbering;
    \item table numbering;
    \item equation numbering;
    \item cross-references;
    \item citations;
    \item bibliography;
    \item font display;
    \item mathematical symbols;
    \item image quality;
    \item completeness of all required forms.
\end{itemize}

The document may be compiled using TeX Live, MiKTeX, or another
compatible \LaTeX{} environment. Regardless of the compilation
environment, the resulting PDF shall conform to the official Faculty
template.

\vspace{0.3cm}

% ================================================================
% XXII. CONTENT THAT STUDENTS MAY EDIT
% ================================================================

\noindent
\textbf{XXII.\quad CONTENT THAT STUDENTS MAY EDIT}

\vspace{0.15cm}

Students may edit or replace:

\begin{itemize}[leftmargin=0.8cm, itemsep=2pt]
    \item student name;
    \item class;
    \item internship organization;
    \item supervisor name;
    \item Abstract;
    \item List of Symbols and Abbreviations;
    \item chapter titles and chapter content;
    \item figures and photographs;
    \item tables;
    \item equations;
    \item citations and references.
\end{itemize}

Students shall not modify the main document formatting settings in the
preamble, including the document class, page geometry, line spacing,
heading formatting, table-of-contents formatting, font configuration,
page layout, or cover-page macros unless specifically instructed.

\vspace{0.3cm}

% ================================================================
% XXIII. FINAL SUBMISSION REQUIREMENTS
% ================================================================

\noindent
\textbf{XXIII.\quad FINAL SUBMISSION REQUIREMENTS}

\vspace{0.15cm}

Before submission, the final report shall:

\begin{itemize}[leftmargin=0.8cm, itemsep=2pt]
    \item compile without errors that affect the report content or presentation;
    \item contain all required forms and report sections;
    \item contain correct student, class, internship organization, and supervisor information;
    \item contain an updated Table of Contents, List of Figures, and List of Tables;
    \item contain no unused example content from the template;
    \item contain no placeholder text such as \textit{Title of Chapter 1},
    \textit{Section 1.1}, or other sample entries;
    \item contain all required citations and references;
    \item contain no unresolved cross-references.
\end{itemize}

If the Faculty requires submission of the \LaTeX{} source, the complete
project shall include all \verb|.tex| files, the bibliography database,
figures, and any other files required to reproduce the submitted PDF.

\vspace{0.4cm}

\noindent
\textbf{Important:}
Before submitting the final internship report, students shall remove or
comment out the following command in the main document:

\begin{verbatim}
\input{Regulations.tex}
\end{verbatim}

For example:

\begin{verbatim}
% \input{Regulations.tex}
\end{verbatim}

This removes the guidance pages from the final submitted report while
preserving the regulations file in the \LaTeX{} project for reference.

} % End 13pt content

\clearpage

\end{verbatim}

For example:

\begin{verbatim}
% 
% =================================================================
% REGULATIONS FOR PREPARING THE GRADUATION INTERNSHIP REPORT
% USING LATEX
% =================================================================
%
% This file is intended to be included from the main document with:
%
%     \input{Regulations.tex}
%
% IMPORTANT:
% This section is provided for guidance only. Students shall remove
% or comment out the \input{Regulations.tex} command before submitting
% the final internship report.
%
% =================================================================

\clearpage
\thispagestyle{empty}

\begin{center}
    {\fontsize{16pt}{20pt}\selectfont\bfseries
    SPECIFIC REGULATIONS FOR PREPARING\\
    GRADUATION INTERNSHIP REPORTS USING \LaTeX}
\end{center}

\vspace{0.5cm}

{\fontsize{13pt}{17pt}\selectfont
\setstretch{1.3}
\justifying

\noindent
\textbf{Note:}
This section is provided for guidance only and is included in the main
document using the command \verb|\input{Regulations.tex}|.
Students shall remove or comment out this command before submitting the
final internship report.

\vspace{0.4cm}

% ================================================================
% I. GENERAL REQUIREMENTS
% ================================================================

\noindent
\textbf{I.\quad GENERAL REQUIREMENTS}

\vspace{0.15cm}

The Graduation Internship Report shall be prepared using the official
\LaTeX{} template provided by the Faculty. Students shall preserve the
predefined page layout, typography, line spacing, heading styles,
numbering system, figure and table formatting, and bibliography settings
unless otherwise instructed by the Faculty or the supervisor.

The report shall be written in English.

Students are responsible for the academic and technical content of the
report. The official \LaTeX{} template is responsible for standardized
document formatting.

\vspace{0.3cm}

% ================================================================
% II. REPORT STRUCTURE
% ================================================================

\noindent
\textbf{II.\quad REPORT STRUCTURE}

\vspace{0.15cm}

The report shall normally be arranged in the following order:

\begin{enumerate}[label=\arabic*., leftmargin=1.2cm, itemsep=2pt]
    \item Outer Cover;
    \item Comments of the Internship Host Organization;
    \item Supervisor's Evaluation and Graduation Internship Defense Results;
    \item Table of Contents;
    \item List of Symbols and Abbreviations, if applicable;
    \item List of Figures;
    \item List of Tables;
    \item Abstract;
    \item Main Chapters of the Internship Report;
    \item References;
    \item Appendices, if applicable.
\end{enumerate}

The preliminary pages shall not be numbered as chapters. The main report
content shall be organized into consecutively numbered chapters beginning
with Chapter 1.

\vspace{0.3cm}

% ================================================================
% III. PAGE SIZE AND MARGINS
% ================================================================

\noindent
\textbf{III.\quad PAGE SIZE AND MARGINS}

\vspace{0.15cm}

The report shall be formatted on A4 paper with the following margins:

\begin{center}
\begin{tabular}{ll}
    Left margin:   & 3.5 cm \\
    Right margin:  & 2.0 cm \\
    Top margin:    & 2.5 cm \\
    Bottom margin: & 2.5 cm
\end{tabular}
\end{center}

These settings are predefined in the official template and shall not be
modified by students.

\vspace{0.3cm}

% ================================================================
% IV. FONT, FONT SIZE, AND LINE SPACING
% ================================================================

\noindent
\textbf{IV.\quad FONT, FONT SIZE, AND LINE SPACING}

\vspace{0.15cm}

The main text shall use a 13 pt font size with 1.3-line spacing.
A Times-family typeface shall be used in accordance with the official
template.

The text shall be fully justified.

Students shall not use repeated spaces, repeated manual line breaks, or
multiple blank lines to position text. Paragraph and heading spacing shall
be controlled by the template.

\vspace{0.3cm}

% ================================================================
% V. OUTER COVER
% ================================================================

\noindent
\textbf{V.\quad OUTER COVER}

\vspace{0.15cm}

The outer cover shall follow the official format provided in the template
and shall contain the following information:

\begin{itemize}[leftmargin=0.8cm, itemsep=2pt]
    \item Thai Nguyen University of Technology;
    \item Faculty of International Training;
    \item Engineering and Technology Department;
    \item University logo and Faculty logo;
    \item the title \textbf{INTERNSHIP REPORT};
    \item Internship Organization;
    \item Author;
    \item Class;
    \item Supervisor;
    \item place and year of completion.
\end{itemize}

Students shall only replace the required information fields and shall not
modify the cover-page layout unless instructed otherwise.

\vspace{0.3cm}

% ================================================================
% VI. COMMENTS OF THE INTERNSHIP HOST ORGANIZATION
% ================================================================

\noindent
\textbf{VI.\quad COMMENTS OF THE INTERNSHIP HOST ORGANIZATION}

\vspace{0.15cm}

The page entitled \textbf{COMMENTS OF THE INTERNSHIP HOST ORGANIZATION}
is reserved for the internship host organization.

It includes:

\begin{itemize}[leftmargin=0.8cm, itemsep=2pt]
    \item the student's full name;
    \item class;
    \item comments from the host organization;
    \item score;
    \item date;
    \item signature and official seal.
\end{itemize}

Students shall not alter or complete official evaluation fields unless
specifically authorized.

\vspace{0.3cm}

% ================================================================
% VII. SUPERVISOR'S EVALUATION AND DEFENSE RESULTS
% ================================================================

\noindent
\textbf{VII.\quad SUPERVISOR'S EVALUATION AND DEFENSE RESULTS}

\vspace{0.15cm}

The section entitled
\textbf{SUPERVISOR'S EVALUATION AND GRADUATION INTERNSHIP DEFENSE RESULTS}
is reserved for the supervisor and defense committee.

The evaluation may include:

\begin{itemize}[leftmargin=0.8cm, itemsep=2pt]
    \item Internship Progress and Student's Attitude;
    \item Report Content;
    \item Report Presentation;
    \item Other Comments;
    \item Supervisor's Opinion;
    \item Defense Result;
    \item Score.
\end{itemize}

Students shall not enter evaluation scores or defense results before they
are officially determined.

\vspace{0.3cm}

% ================================================================
% VIII. CHAPTERS, SECTIONS, AND SUBSECTIONS
% ================================================================

\noindent
\textbf{VIII.\quad CHAPTERS, SECTIONS, AND SUBSECTIONS}

\vspace{0.15cm}

The report structure shall be generated using the standard \LaTeX{}
commands:

\begin{verbatim}
\chapter{...}
\section{...}
\subsection{...}
\end{verbatim}

Chapter, section, and subsection numbers shall be generated automatically
by \LaTeX{}.

Students shall not manually type chapter or section numbers and shall not
reproduce heading styles using commands such as \verb|\textbf|,
\verb|\fontsize|, or \verb|\centering| instead of the appropriate
structural command.

Excessive subdivision should be avoided. In most cases, the hierarchy
Chapter--Section--Subsection is sufficient.

\vspace{0.3cm}

% ================================================================
% IX. ORGANIZATION OF THE MAIN CONTENT
% ================================================================

\noindent
\textbf{IX.\quad ORGANIZATION OF THE MAIN CONTENT}

\vspace{0.15cm}

The specific content of the report shall depend on the internship
organization and the assigned tasks. A typical structure may include:

\begin{itemize}[leftmargin=0.8cm, itemsep=2pt]
    \item \textbf{Chapter 1. Overview of the Internship Organization};
    \item \textbf{Chapter 2. Internship Activities and Assigned Tasks};
    \item \textbf{Chapter 3. Results, Experience, and Evaluation}.
\end{itemize}

These chapter titles are recommendations only and may be adapted to the
actual internship content.

The report shall reflect work actually performed, observed, analyzed, or
learned during the internship. Unrelated theoretical material should not
be included merely to increase report length.

\vspace{0.3cm}

% ================================================================
% X. TABLE OF CONTENTS
% ================================================================

\noindent
\textbf{X.\quad TABLE OF CONTENTS}

\vspace{0.15cm}

The Table of Contents shall be generated automatically using:

\begin{verbatim}
\tableofcontents
\end{verbatim}

Students shall not create or manually edit page numbers in the Table of
Contents.

The document shall be compiled sufficiently to ensure that all page
numbers and entries are correctly updated.

\vspace{0.3cm}

% ================================================================
% XI. LIST OF SYMBOLS AND ABBREVIATIONS
% ================================================================

\noindent
\textbf{XI.\quad LIST OF SYMBOLS AND ABBREVIATIONS}

\vspace{0.15cm}

Symbols and abbreviations that are used repeatedly in the report shall be
listed in the \textbf{List of Symbols and Abbreviations}.

The list may contain the following columns:

\begin{center}
\begin{tabular}{llll}
    No. & Symbol / Abbreviation & Unit & Definition
\end{tabular}
\end{center}

Only symbols and abbreviations actually used in the final report shall be
retained. Example entries contained in the template shall be removed if
they are not relevant to the student's report.

Each abbreviation should normally be defined when it first appears in
the main text.

\vspace{0.3cm}

% ================================================================
% XII. FIGURES, DIAGRAMS, AND PHOTOGRAPHS
% ================================================================

\noindent
\textbf{XII.\quad FIGURES, DIAGRAMS, AND PHOTOGRAPHS}

\vspace{0.15cm}

Figures shall be clear, legible, and directly relevant to the report
content.

Figures shall be inserted using the standard \verb|figure| environment.

Example:

\begin{verbatim}
\begin{figure}[H]
    \centering
    \includegraphics[width=0.8\textwidth]
                    {figures/company_structure.pdf}
    \caption{Organizational structure of the company}
    \label{fig:company-structure}
\end{figure}
\end{verbatim}

Figure captions shall be placed below the figure.

Each referenced figure shall have a unique \verb|\label| and shall be
referenced in the text using, for example:

\begin{verbatim}
Figure~\ref{fig:company-structure}
\end{verbatim}

Students shall not manually type figure numbers.

Figures reproduced or adapted from external sources shall include an
appropriate citation.

Photographs taken during the internship shall be relevant and sufficiently
clear. Students shall not include confidential, restricted, or proprietary
information without permission from the internship organization.

Vector graphics or high-resolution images are recommended whenever
possible.

\vspace{0.3cm}

% ================================================================
% XIII. TABLES
% ================================================================

\noindent
\textbf{XIII.\quad TABLES}

\vspace{0.15cm}

Tables shall be created using suitable \LaTeX{} environments such as
\verb|table|, \verb|tabular|, \verb|tabularx|, or \verb|longtable|.

Table captions shall be placed above the table.

Example:

\begin{verbatim}
\begin{table}[H]
    \centering
    \caption{Internship schedule}
    \label{tab:internship-schedule}

    \begin{tabular}{lll}
        \toprule
        Week & Task & Result \\
        \midrule
        1 & Orientation & Completed \\
        2 & System study & Completed \\
        \bottomrule
    \end{tabular}
\end{table}
\end{verbatim}

Tables shall be referenced using:

\begin{verbatim}
Table~\ref{tab:internship-schedule}
\end{verbatim}

Students shall not manually type table numbers.

Tables should be created directly in \LaTeX{} rather than inserted as
screenshots whenever practical.

\vspace{0.3cm}

% ================================================================
% XIV. EQUATIONS AND TECHNICAL QUANTITIES
% ================================================================

\noindent
\textbf{XIV.\quad EQUATIONS AND TECHNICAL QUANTITIES}

\vspace{0.15cm}

Mathematical expressions shall be typeset directly using \LaTeX{}.
Images or screenshots of equations shall not be used.

Short expressions may be written in inline mathematical mode, for example:

\begin{verbatim}
$R = 0.7\,\Omega$
\end{verbatim}

Important equations shall use a numbered equation environment.

Example:

\begin{verbatim}
\begin{equation}
    P = UI
    \label{eq:power}
\end{equation}
\end{verbatim}

Equations shall be referenced using:

\begin{verbatim}
Equation~\eqref{eq:power}
\end{verbatim}

Equation numbers shall not be manually entered.

SI units shall be used wherever applicable. The \verb|siunitx| package
should be used where appropriate.

Examples:

\begin{verbatim}
\SI{220}{\volt}
\SI{50}{\hertz}
\SI{2.5}{\ampere}
\end{verbatim}

Notation and units shall be consistent throughout the report.

\vspace{0.3cm}

% ================================================================
% XV. SOURCE CODE
% ================================================================

\noindent
\textbf{XV.\quad SOURCE CODE}

\vspace{0.15cm}

If the internship involves programming, PLCs, microcontrollers,
simulation, or software development, relevant source code may be included
in the report.

Source code shall be presented using the \verb|listings| environment or
another environment defined by the official template.

Example:

\begin{verbatim}
\begin{lstlisting}[language=st]
IF Start = TRUE THEN
    Motor := TRUE;
ELSE
    Motor := FALSE;
END_IF;
\end{lstlisting}
\end{verbatim}

Only code necessary for explaining the work should appear in the main
report. Long source-code listings should be moved to an appendix or
submitted as supplementary files if required.

Source code should not normally be inserted as screenshots.

\vspace{0.3cm}

% ================================================================
% XVI. ABSTRACT
% ================================================================

\noindent
\textbf{XVI.\quad ABSTRACT}

\vspace{0.15cm}

The Abstract shall briefly summarize:

\begin{itemize}[leftmargin=0.8cm, itemsep=2pt]
    \item the internship organization;
    \item the objectives of the internship;
    \item the main tasks performed;
    \item the main results, knowledge, and skills obtained.
\end{itemize}

The Abstract should be concise and should not reproduce the full content
of the report chapter by chapter.

\vspace{0.3cm}

% ================================================================
% XVII. REFERENCES AND CITATIONS
% ================================================================

\noindent
\textbf{XVII.\quad REFERENCES AND CITATIONS}

\vspace{0.15cm}

All books, journal papers, conference papers, standards, websites,
datasheets, manuals, technical documents, figures, tables, or other
external sources used in preparing the report shall be properly cited.

References shall be managed through the bibliography system provided with
the official template.

Bibliographic entries shall normally be stored in:

\begin{verbatim}
References.bib
\end{verbatim}

References shall be cited in the report using:

\begin{verbatim}
\cite{reference-key}
\end{verbatim}

Students shall not manually type reference numbers such as ``[1]'',
``[2]'', etc.

Titles of foreign-language publications shall remain in their original
language in the bibliography.

All borrowed ideas, data, figures, tables, equations, methods, or results
shall be properly acknowledged.

\vspace{0.3cm}

% ================================================================
% XVIII. CROSS-REFERENCING
% ================================================================

\noindent
\textbf{XVIII.\quad CROSS-REFERENCING}

\vspace{0.15cm}

Numbered elements shall be cross-referenced using the standard \LaTeX{}
mechanism.

Use the following commands where appropriate:

\begin{verbatim}
\label{...}
\ref{...}
\eqref{...}
\end{verbatim}

For example:

\begin{verbatim}
As shown in Figure~\ref{fig:system}, ...
\end{verbatim}

Students shall not manually type numbers such as ``Figure 2.3'' when
those numbers can be generated automatically.

\vspace{0.3cm}

% ================================================================
% XIX. PAGE NUMBERING
% ================================================================

\noindent
\textbf{XIX.\quad PAGE NUMBERING}

\vspace{0.15cm}

Cover pages and official evaluation forms at the beginning of the report
shall not display page numbers.

The main body of the report shall use Arabic page numbers beginning from
Chapter 1.

Page numbering shall be generated automatically by the official template.
Students shall not manually enter page numbers.

\vspace{0.3cm}

% ================================================================
% XX. PROJECT FILE ORGANIZATION
% ================================================================

\noindent
\textbf{XX.\quad PROJECT FILE ORGANIZATION}

\vspace{0.15cm}

Students should maintain a clear and organized \LaTeX{} project
structure. A recommended structure is:

\begin{verbatim}
InternshipReport/
|
|-- main.tex
|-- References.bib
|
|-- figures/
|   |-- company.pdf
|   |-- system.png
|   `-- ...
|
`-- chapters/
    |-- chapter1.tex
    |-- chapter2.tex
    `-- chapter3.tex
\end{verbatim}

Figures should normally be stored in the \verb|figures| directory
specified by the official template.

File and directory names should preferably contain Latin letters,
numbers, hyphens, or underscores.

Recommended examples include:

\begin{verbatim}
company_structure.pdf
motor_control.png
chapter_01.tex
\end{verbatim}

Spaces and unnecessary special characters in file names should be
avoided where possible.

\vspace{0.3cm}

% ================================================================
% XXI. DOCUMENT COMPILATION AND FINAL CHECK
% ================================================================

\noindent
\textbf{XXI.\quad DOCUMENT COMPILATION AND FINAL CHECK}

\vspace{0.15cm}

Before submission, the complete \LaTeX{} project shall compile
successfully and produce the final PDF report.

The final document shall not contain unresolved references such as:

\begin{verbatim}
??
[?]
\end{verbatim}

Before generating the final PDF, students shall check:

\begin{itemize}[leftmargin=0.8cm, itemsep=2pt]
    \item the Table of Contents;
    \item page numbering;
    \item chapter, section, and subsection numbering;
    \item figure numbering;
    \item table numbering;
    \item equation numbering;
    \item cross-references;
    \item citations;
    \item bibliography;
    \item font display;
    \item mathematical symbols;
    \item image quality;
    \item completeness of all required forms.
\end{itemize}

The document may be compiled using TeX Live, MiKTeX, or another
compatible \LaTeX{} environment. Regardless of the compilation
environment, the resulting PDF shall conform to the official Faculty
template.

\vspace{0.3cm}

% ================================================================
% XXII. CONTENT THAT STUDENTS MAY EDIT
% ================================================================

\noindent
\textbf{XXII.\quad CONTENT THAT STUDENTS MAY EDIT}

\vspace{0.15cm}

Students may edit or replace:

\begin{itemize}[leftmargin=0.8cm, itemsep=2pt]
    \item student name;
    \item class;
    \item internship organization;
    \item supervisor name;
    \item Abstract;
    \item List of Symbols and Abbreviations;
    \item chapter titles and chapter content;
    \item figures and photographs;
    \item tables;
    \item equations;
    \item citations and references.
\end{itemize}

Students shall not modify the main document formatting settings in the
preamble, including the document class, page geometry, line spacing,
heading formatting, table-of-contents formatting, font configuration,
page layout, or cover-page macros unless specifically instructed.

\vspace{0.3cm}

% ================================================================
% XXIII. FINAL SUBMISSION REQUIREMENTS
% ================================================================

\noindent
\textbf{XXIII.\quad FINAL SUBMISSION REQUIREMENTS}

\vspace{0.15cm}

Before submission, the final report shall:

\begin{itemize}[leftmargin=0.8cm, itemsep=2pt]
    \item compile without errors that affect the report content or presentation;
    \item contain all required forms and report sections;
    \item contain correct student, class, internship organization, and supervisor information;
    \item contain an updated Table of Contents, List of Figures, and List of Tables;
    \item contain no unused example content from the template;
    \item contain no placeholder text such as \textit{Title of Chapter 1},
    \textit{Section 1.1}, or other sample entries;
    \item contain all required citations and references;
    \item contain no unresolved cross-references.
\end{itemize}

If the Faculty requires submission of the \LaTeX{} source, the complete
project shall include all \verb|.tex| files, the bibliography database,
figures, and any other files required to reproduce the submitted PDF.

\vspace{0.4cm}

\noindent
\textbf{Important:}
Before submitting the final internship report, students shall remove or
comment out the following command in the main document:

\begin{verbatim}
\input{Regulations.tex}
\end{verbatim}

For example:

\begin{verbatim}
% \input{Regulations.tex}
\end{verbatim}

This removes the guidance pages from the final submitted report while
preserving the regulations file in the \LaTeX{} project for reference.

} % End 13pt content

\clearpage

\end{verbatim}

This removes the guidance pages from the final submitted report while
preserving the regulations file in the \LaTeX{} project for reference.

} % End 13pt content

\clearpage

%
% IMPORTANT:
% This section is provided for guidance only. Students shall remove
% or comment out the 
% =================================================================
% REGULATIONS FOR PREPARING THE GRADUATION INTERNSHIP REPORT
% USING LATEX
% =================================================================
%
% This file is intended to be included from the main document with:
%
%     
% =================================================================
% REGULATIONS FOR PREPARING THE GRADUATION INTERNSHIP REPORT
% USING LATEX
% =================================================================
%
% This file is intended to be included from the main document with:
%
%     \input{Regulations.tex}
%
% IMPORTANT:
% This section is provided for guidance only. Students shall remove
% or comment out the \input{Regulations.tex} command before submitting
% the final internship report.
%
% =================================================================

\clearpage
\thispagestyle{empty}

\begin{center}
    {\fontsize{16pt}{20pt}\selectfont\bfseries
    SPECIFIC REGULATIONS FOR PREPARING\\
    GRADUATION INTERNSHIP REPORTS USING \LaTeX}
\end{center}

\vspace{0.5cm}

{\fontsize{13pt}{17pt}\selectfont
\setstretch{1.3}
\justifying

\noindent
\textbf{Note:}
This section is provided for guidance only and is included in the main
document using the command \verb|\input{Regulations.tex}|.
Students shall remove or comment out this command before submitting the
final internship report.

\vspace{0.4cm}

% ================================================================
% I. GENERAL REQUIREMENTS
% ================================================================

\noindent
\textbf{I.\quad GENERAL REQUIREMENTS}

\vspace{0.15cm}

The Graduation Internship Report shall be prepared using the official
\LaTeX{} template provided by the Faculty. Students shall preserve the
predefined page layout, typography, line spacing, heading styles,
numbering system, figure and table formatting, and bibliography settings
unless otherwise instructed by the Faculty or the supervisor.

The report shall be written in English.

Students are responsible for the academic and technical content of the
report. The official \LaTeX{} template is responsible for standardized
document formatting.

\vspace{0.3cm}

% ================================================================
% II. REPORT STRUCTURE
% ================================================================

\noindent
\textbf{II.\quad REPORT STRUCTURE}

\vspace{0.15cm}

The report shall normally be arranged in the following order:

\begin{enumerate}[label=\arabic*., leftmargin=1.2cm, itemsep=2pt]
    \item Outer Cover;
    \item Comments of the Internship Host Organization;
    \item Supervisor's Evaluation and Graduation Internship Defense Results;
    \item Table of Contents;
    \item List of Symbols and Abbreviations, if applicable;
    \item List of Figures;
    \item List of Tables;
    \item Abstract;
    \item Main Chapters of the Internship Report;
    \item References;
    \item Appendices, if applicable.
\end{enumerate}

The preliminary pages shall not be numbered as chapters. The main report
content shall be organized into consecutively numbered chapters beginning
with Chapter 1.

\vspace{0.3cm}

% ================================================================
% III. PAGE SIZE AND MARGINS
% ================================================================

\noindent
\textbf{III.\quad PAGE SIZE AND MARGINS}

\vspace{0.15cm}

The report shall be formatted on A4 paper with the following margins:

\begin{center}
\begin{tabular}{ll}
    Left margin:   & 3.5 cm \\
    Right margin:  & 2.0 cm \\
    Top margin:    & 2.5 cm \\
    Bottom margin: & 2.5 cm
\end{tabular}
\end{center}

These settings are predefined in the official template and shall not be
modified by students.

\vspace{0.3cm}

% ================================================================
% IV. FONT, FONT SIZE, AND LINE SPACING
% ================================================================

\noindent
\textbf{IV.\quad FONT, FONT SIZE, AND LINE SPACING}

\vspace{0.15cm}

The main text shall use a 13 pt font size with 1.3-line spacing.
A Times-family typeface shall be used in accordance with the official
template.

The text shall be fully justified.

Students shall not use repeated spaces, repeated manual line breaks, or
multiple blank lines to position text. Paragraph and heading spacing shall
be controlled by the template.

\vspace{0.3cm}

% ================================================================
% V. OUTER COVER
% ================================================================

\noindent
\textbf{V.\quad OUTER COVER}

\vspace{0.15cm}

The outer cover shall follow the official format provided in the template
and shall contain the following information:

\begin{itemize}[leftmargin=0.8cm, itemsep=2pt]
    \item Thai Nguyen University of Technology;
    \item Faculty of International Training;
    \item Engineering and Technology Department;
    \item University logo and Faculty logo;
    \item the title \textbf{INTERNSHIP REPORT};
    \item Internship Organization;
    \item Author;
    \item Class;
    \item Supervisor;
    \item place and year of completion.
\end{itemize}

Students shall only replace the required information fields and shall not
modify the cover-page layout unless instructed otherwise.

\vspace{0.3cm}

% ================================================================
% VI. COMMENTS OF THE INTERNSHIP HOST ORGANIZATION
% ================================================================

\noindent
\textbf{VI.\quad COMMENTS OF THE INTERNSHIP HOST ORGANIZATION}

\vspace{0.15cm}

The page entitled \textbf{COMMENTS OF THE INTERNSHIP HOST ORGANIZATION}
is reserved for the internship host organization.

It includes:

\begin{itemize}[leftmargin=0.8cm, itemsep=2pt]
    \item the student's full name;
    \item class;
    \item comments from the host organization;
    \item score;
    \item date;
    \item signature and official seal.
\end{itemize}

Students shall not alter or complete official evaluation fields unless
specifically authorized.

\vspace{0.3cm}

% ================================================================
% VII. SUPERVISOR'S EVALUATION AND DEFENSE RESULTS
% ================================================================

\noindent
\textbf{VII.\quad SUPERVISOR'S EVALUATION AND DEFENSE RESULTS}

\vspace{0.15cm}

The section entitled
\textbf{SUPERVISOR'S EVALUATION AND GRADUATION INTERNSHIP DEFENSE RESULTS}
is reserved for the supervisor and defense committee.

The evaluation may include:

\begin{itemize}[leftmargin=0.8cm, itemsep=2pt]
    \item Internship Progress and Student's Attitude;
    \item Report Content;
    \item Report Presentation;
    \item Other Comments;
    \item Supervisor's Opinion;
    \item Defense Result;
    \item Score.
\end{itemize}

Students shall not enter evaluation scores or defense results before they
are officially determined.

\vspace{0.3cm}

% ================================================================
% VIII. CHAPTERS, SECTIONS, AND SUBSECTIONS
% ================================================================

\noindent
\textbf{VIII.\quad CHAPTERS, SECTIONS, AND SUBSECTIONS}

\vspace{0.15cm}

The report structure shall be generated using the standard \LaTeX{}
commands:

\begin{verbatim}
\chapter{...}
\section{...}
\subsection{...}
\end{verbatim}

Chapter, section, and subsection numbers shall be generated automatically
by \LaTeX{}.

Students shall not manually type chapter or section numbers and shall not
reproduce heading styles using commands such as \verb|\textbf|,
\verb|\fontsize|, or \verb|\centering| instead of the appropriate
structural command.

Excessive subdivision should be avoided. In most cases, the hierarchy
Chapter--Section--Subsection is sufficient.

\vspace{0.3cm}

% ================================================================
% IX. ORGANIZATION OF THE MAIN CONTENT
% ================================================================

\noindent
\textbf{IX.\quad ORGANIZATION OF THE MAIN CONTENT}

\vspace{0.15cm}

The specific content of the report shall depend on the internship
organization and the assigned tasks. A typical structure may include:

\begin{itemize}[leftmargin=0.8cm, itemsep=2pt]
    \item \textbf{Chapter 1. Overview of the Internship Organization};
    \item \textbf{Chapter 2. Internship Activities and Assigned Tasks};
    \item \textbf{Chapter 3. Results, Experience, and Evaluation}.
\end{itemize}

These chapter titles are recommendations only and may be adapted to the
actual internship content.

The report shall reflect work actually performed, observed, analyzed, or
learned during the internship. Unrelated theoretical material should not
be included merely to increase report length.

\vspace{0.3cm}

% ================================================================
% X. TABLE OF CONTENTS
% ================================================================

\noindent
\textbf{X.\quad TABLE OF CONTENTS}

\vspace{0.15cm}

The Table of Contents shall be generated automatically using:

\begin{verbatim}
\tableofcontents
\end{verbatim}

Students shall not create or manually edit page numbers in the Table of
Contents.

The document shall be compiled sufficiently to ensure that all page
numbers and entries are correctly updated.

\vspace{0.3cm}

% ================================================================
% XI. LIST OF SYMBOLS AND ABBREVIATIONS
% ================================================================

\noindent
\textbf{XI.\quad LIST OF SYMBOLS AND ABBREVIATIONS}

\vspace{0.15cm}

Symbols and abbreviations that are used repeatedly in the report shall be
listed in the \textbf{List of Symbols and Abbreviations}.

The list may contain the following columns:

\begin{center}
\begin{tabular}{llll}
    No. & Symbol / Abbreviation & Unit & Definition
\end{tabular}
\end{center}

Only symbols and abbreviations actually used in the final report shall be
retained. Example entries contained in the template shall be removed if
they are not relevant to the student's report.

Each abbreviation should normally be defined when it first appears in
the main text.

\vspace{0.3cm}

% ================================================================
% XII. FIGURES, DIAGRAMS, AND PHOTOGRAPHS
% ================================================================

\noindent
\textbf{XII.\quad FIGURES, DIAGRAMS, AND PHOTOGRAPHS}

\vspace{0.15cm}

Figures shall be clear, legible, and directly relevant to the report
content.

Figures shall be inserted using the standard \verb|figure| environment.

Example:

\begin{verbatim}
\begin{figure}[H]
    \centering
    \includegraphics[width=0.8\textwidth]
                    {figures/company_structure.pdf}
    \caption{Organizational structure of the company}
    \label{fig:company-structure}
\end{figure}
\end{verbatim}

Figure captions shall be placed below the figure.

Each referenced figure shall have a unique \verb|\label| and shall be
referenced in the text using, for example:

\begin{verbatim}
Figure~\ref{fig:company-structure}
\end{verbatim}

Students shall not manually type figure numbers.

Figures reproduced or adapted from external sources shall include an
appropriate citation.

Photographs taken during the internship shall be relevant and sufficiently
clear. Students shall not include confidential, restricted, or proprietary
information without permission from the internship organization.

Vector graphics or high-resolution images are recommended whenever
possible.

\vspace{0.3cm}

% ================================================================
% XIII. TABLES
% ================================================================

\noindent
\textbf{XIII.\quad TABLES}

\vspace{0.15cm}

Tables shall be created using suitable \LaTeX{} environments such as
\verb|table|, \verb|tabular|, \verb|tabularx|, or \verb|longtable|.

Table captions shall be placed above the table.

Example:

\begin{verbatim}
\begin{table}[H]
    \centering
    \caption{Internship schedule}
    \label{tab:internship-schedule}

    \begin{tabular}{lll}
        \toprule
        Week & Task & Result \\
        \midrule
        1 & Orientation & Completed \\
        2 & System study & Completed \\
        \bottomrule
    \end{tabular}
\end{table}
\end{verbatim}

Tables shall be referenced using:

\begin{verbatim}
Table~\ref{tab:internship-schedule}
\end{verbatim}

Students shall not manually type table numbers.

Tables should be created directly in \LaTeX{} rather than inserted as
screenshots whenever practical.

\vspace{0.3cm}

% ================================================================
% XIV. EQUATIONS AND TECHNICAL QUANTITIES
% ================================================================

\noindent
\textbf{XIV.\quad EQUATIONS AND TECHNICAL QUANTITIES}

\vspace{0.15cm}

Mathematical expressions shall be typeset directly using \LaTeX{}.
Images or screenshots of equations shall not be used.

Short expressions may be written in inline mathematical mode, for example:

\begin{verbatim}
$R = 0.7\,\Omega$
\end{verbatim}

Important equations shall use a numbered equation environment.

Example:

\begin{verbatim}
\begin{equation}
    P = UI
    \label{eq:power}
\end{equation}
\end{verbatim}

Equations shall be referenced using:

\begin{verbatim}
Equation~\eqref{eq:power}
\end{verbatim}

Equation numbers shall not be manually entered.

SI units shall be used wherever applicable. The \verb|siunitx| package
should be used where appropriate.

Examples:

\begin{verbatim}
\SI{220}{\volt}
\SI{50}{\hertz}
\SI{2.5}{\ampere}
\end{verbatim}

Notation and units shall be consistent throughout the report.

\vspace{0.3cm}

% ================================================================
% XV. SOURCE CODE
% ================================================================

\noindent
\textbf{XV.\quad SOURCE CODE}

\vspace{0.15cm}

If the internship involves programming, PLCs, microcontrollers,
simulation, or software development, relevant source code may be included
in the report.

Source code shall be presented using the \verb|listings| environment or
another environment defined by the official template.

Example:

\begin{verbatim}
\begin{lstlisting}[language=st]
IF Start = TRUE THEN
    Motor := TRUE;
ELSE
    Motor := FALSE;
END_IF;
\end{lstlisting}
\end{verbatim}

Only code necessary for explaining the work should appear in the main
report. Long source-code listings should be moved to an appendix or
submitted as supplementary files if required.

Source code should not normally be inserted as screenshots.

\vspace{0.3cm}

% ================================================================
% XVI. ABSTRACT
% ================================================================

\noindent
\textbf{XVI.\quad ABSTRACT}

\vspace{0.15cm}

The Abstract shall briefly summarize:

\begin{itemize}[leftmargin=0.8cm, itemsep=2pt]
    \item the internship organization;
    \item the objectives of the internship;
    \item the main tasks performed;
    \item the main results, knowledge, and skills obtained.
\end{itemize}

The Abstract should be concise and should not reproduce the full content
of the report chapter by chapter.

\vspace{0.3cm}

% ================================================================
% XVII. REFERENCES AND CITATIONS
% ================================================================

\noindent
\textbf{XVII.\quad REFERENCES AND CITATIONS}

\vspace{0.15cm}

All books, journal papers, conference papers, standards, websites,
datasheets, manuals, technical documents, figures, tables, or other
external sources used in preparing the report shall be properly cited.

References shall be managed through the bibliography system provided with
the official template.

Bibliographic entries shall normally be stored in:

\begin{verbatim}
References.bib
\end{verbatim}

References shall be cited in the report using:

\begin{verbatim}
\cite{reference-key}
\end{verbatim}

Students shall not manually type reference numbers such as ``[1]'',
``[2]'', etc.

Titles of foreign-language publications shall remain in their original
language in the bibliography.

All borrowed ideas, data, figures, tables, equations, methods, or results
shall be properly acknowledged.

\vspace{0.3cm}

% ================================================================
% XVIII. CROSS-REFERENCING
% ================================================================

\noindent
\textbf{XVIII.\quad CROSS-REFERENCING}

\vspace{0.15cm}

Numbered elements shall be cross-referenced using the standard \LaTeX{}
mechanism.

Use the following commands where appropriate:

\begin{verbatim}
\label{...}
\ref{...}
\eqref{...}
\end{verbatim}

For example:

\begin{verbatim}
As shown in Figure~\ref{fig:system}, ...
\end{verbatim}

Students shall not manually type numbers such as ``Figure 2.3'' when
those numbers can be generated automatically.

\vspace{0.3cm}

% ================================================================
% XIX. PAGE NUMBERING
% ================================================================

\noindent
\textbf{XIX.\quad PAGE NUMBERING}

\vspace{0.15cm}

Cover pages and official evaluation forms at the beginning of the report
shall not display page numbers.

The main body of the report shall use Arabic page numbers beginning from
Chapter 1.

Page numbering shall be generated automatically by the official template.
Students shall not manually enter page numbers.

\vspace{0.3cm}

% ================================================================
% XX. PROJECT FILE ORGANIZATION
% ================================================================

\noindent
\textbf{XX.\quad PROJECT FILE ORGANIZATION}

\vspace{0.15cm}

Students should maintain a clear and organized \LaTeX{} project
structure. A recommended structure is:

\begin{verbatim}
InternshipReport/
|
|-- main.tex
|-- References.bib
|
|-- figures/
|   |-- company.pdf
|   |-- system.png
|   `-- ...
|
`-- chapters/
    |-- chapter1.tex
    |-- chapter2.tex
    `-- chapter3.tex
\end{verbatim}

Figures should normally be stored in the \verb|figures| directory
specified by the official template.

File and directory names should preferably contain Latin letters,
numbers, hyphens, or underscores.

Recommended examples include:

\begin{verbatim}
company_structure.pdf
motor_control.png
chapter_01.tex
\end{verbatim}

Spaces and unnecessary special characters in file names should be
avoided where possible.

\vspace{0.3cm}

% ================================================================
% XXI. DOCUMENT COMPILATION AND FINAL CHECK
% ================================================================

\noindent
\textbf{XXI.\quad DOCUMENT COMPILATION AND FINAL CHECK}

\vspace{0.15cm}

Before submission, the complete \LaTeX{} project shall compile
successfully and produce the final PDF report.

The final document shall not contain unresolved references such as:

\begin{verbatim}
??
[?]
\end{verbatim}

Before generating the final PDF, students shall check:

\begin{itemize}[leftmargin=0.8cm, itemsep=2pt]
    \item the Table of Contents;
    \item page numbering;
    \item chapter, section, and subsection numbering;
    \item figure numbering;
    \item table numbering;
    \item equation numbering;
    \item cross-references;
    \item citations;
    \item bibliography;
    \item font display;
    \item mathematical symbols;
    \item image quality;
    \item completeness of all required forms.
\end{itemize}

The document may be compiled using TeX Live, MiKTeX, or another
compatible \LaTeX{} environment. Regardless of the compilation
environment, the resulting PDF shall conform to the official Faculty
template.

\vspace{0.3cm}

% ================================================================
% XXII. CONTENT THAT STUDENTS MAY EDIT
% ================================================================

\noindent
\textbf{XXII.\quad CONTENT THAT STUDENTS MAY EDIT}

\vspace{0.15cm}

Students may edit or replace:

\begin{itemize}[leftmargin=0.8cm, itemsep=2pt]
    \item student name;
    \item class;
    \item internship organization;
    \item supervisor name;
    \item Abstract;
    \item List of Symbols and Abbreviations;
    \item chapter titles and chapter content;
    \item figures and photographs;
    \item tables;
    \item equations;
    \item citations and references.
\end{itemize}

Students shall not modify the main document formatting settings in the
preamble, including the document class, page geometry, line spacing,
heading formatting, table-of-contents formatting, font configuration,
page layout, or cover-page macros unless specifically instructed.

\vspace{0.3cm}

% ================================================================
% XXIII. FINAL SUBMISSION REQUIREMENTS
% ================================================================

\noindent
\textbf{XXIII.\quad FINAL SUBMISSION REQUIREMENTS}

\vspace{0.15cm}

Before submission, the final report shall:

\begin{itemize}[leftmargin=0.8cm, itemsep=2pt]
    \item compile without errors that affect the report content or presentation;
    \item contain all required forms and report sections;
    \item contain correct student, class, internship organization, and supervisor information;
    \item contain an updated Table of Contents, List of Figures, and List of Tables;
    \item contain no unused example content from the template;
    \item contain no placeholder text such as \textit{Title of Chapter 1},
    \textit{Section 1.1}, or other sample entries;
    \item contain all required citations and references;
    \item contain no unresolved cross-references.
\end{itemize}

If the Faculty requires submission of the \LaTeX{} source, the complete
project shall include all \verb|.tex| files, the bibliography database,
figures, and any other files required to reproduce the submitted PDF.

\vspace{0.4cm}

\noindent
\textbf{Important:}
Before submitting the final internship report, students shall remove or
comment out the following command in the main document:

\begin{verbatim}
\input{Regulations.tex}
\end{verbatim}

For example:

\begin{verbatim}
% \input{Regulations.tex}
\end{verbatim}

This removes the guidance pages from the final submitted report while
preserving the regulations file in the \LaTeX{} project for reference.

} % End 13pt content

\clearpage

%
% IMPORTANT:
% This section is provided for guidance only. Students shall remove
% or comment out the 
% =================================================================
% REGULATIONS FOR PREPARING THE GRADUATION INTERNSHIP REPORT
% USING LATEX
% =================================================================
%
% This file is intended to be included from the main document with:
%
%     \input{Regulations.tex}
%
% IMPORTANT:
% This section is provided for guidance only. Students shall remove
% or comment out the \input{Regulations.tex} command before submitting
% the final internship report.
%
% =================================================================

\clearpage
\thispagestyle{empty}

\begin{center}
    {\fontsize{16pt}{20pt}\selectfont\bfseries
    SPECIFIC REGULATIONS FOR PREPARING\\
    GRADUATION INTERNSHIP REPORTS USING \LaTeX}
\end{center}

\vspace{0.5cm}

{\fontsize{13pt}{17pt}\selectfont
\setstretch{1.3}
\justifying

\noindent
\textbf{Note:}
This section is provided for guidance only and is included in the main
document using the command \verb|\input{Regulations.tex}|.
Students shall remove or comment out this command before submitting the
final internship report.

\vspace{0.4cm}

% ================================================================
% I. GENERAL REQUIREMENTS
% ================================================================

\noindent
\textbf{I.\quad GENERAL REQUIREMENTS}

\vspace{0.15cm}

The Graduation Internship Report shall be prepared using the official
\LaTeX{} template provided by the Faculty. Students shall preserve the
predefined page layout, typography, line spacing, heading styles,
numbering system, figure and table formatting, and bibliography settings
unless otherwise instructed by the Faculty or the supervisor.

The report shall be written in English.

Students are responsible for the academic and technical content of the
report. The official \LaTeX{} template is responsible for standardized
document formatting.

\vspace{0.3cm}

% ================================================================
% II. REPORT STRUCTURE
% ================================================================

\noindent
\textbf{II.\quad REPORT STRUCTURE}

\vspace{0.15cm}

The report shall normally be arranged in the following order:

\begin{enumerate}[label=\arabic*., leftmargin=1.2cm, itemsep=2pt]
    \item Outer Cover;
    \item Comments of the Internship Host Organization;
    \item Supervisor's Evaluation and Graduation Internship Defense Results;
    \item Table of Contents;
    \item List of Symbols and Abbreviations, if applicable;
    \item List of Figures;
    \item List of Tables;
    \item Abstract;
    \item Main Chapters of the Internship Report;
    \item References;
    \item Appendices, if applicable.
\end{enumerate}

The preliminary pages shall not be numbered as chapters. The main report
content shall be organized into consecutively numbered chapters beginning
with Chapter 1.

\vspace{0.3cm}

% ================================================================
% III. PAGE SIZE AND MARGINS
% ================================================================

\noindent
\textbf{III.\quad PAGE SIZE AND MARGINS}

\vspace{0.15cm}

The report shall be formatted on A4 paper with the following margins:

\begin{center}
\begin{tabular}{ll}
    Left margin:   & 3.5 cm \\
    Right margin:  & 2.0 cm \\
    Top margin:    & 2.5 cm \\
    Bottom margin: & 2.5 cm
\end{tabular}
\end{center}

These settings are predefined in the official template and shall not be
modified by students.

\vspace{0.3cm}

% ================================================================
% IV. FONT, FONT SIZE, AND LINE SPACING
% ================================================================

\noindent
\textbf{IV.\quad FONT, FONT SIZE, AND LINE SPACING}

\vspace{0.15cm}

The main text shall use a 13 pt font size with 1.3-line spacing.
A Times-family typeface shall be used in accordance with the official
template.

The text shall be fully justified.

Students shall not use repeated spaces, repeated manual line breaks, or
multiple blank lines to position text. Paragraph and heading spacing shall
be controlled by the template.

\vspace{0.3cm}

% ================================================================
% V. OUTER COVER
% ================================================================

\noindent
\textbf{V.\quad OUTER COVER}

\vspace{0.15cm}

The outer cover shall follow the official format provided in the template
and shall contain the following information:

\begin{itemize}[leftmargin=0.8cm, itemsep=2pt]
    \item Thai Nguyen University of Technology;
    \item Faculty of International Training;
    \item Engineering and Technology Department;
    \item University logo and Faculty logo;
    \item the title \textbf{INTERNSHIP REPORT};
    \item Internship Organization;
    \item Author;
    \item Class;
    \item Supervisor;
    \item place and year of completion.
\end{itemize}

Students shall only replace the required information fields and shall not
modify the cover-page layout unless instructed otherwise.

\vspace{0.3cm}

% ================================================================
% VI. COMMENTS OF THE INTERNSHIP HOST ORGANIZATION
% ================================================================

\noindent
\textbf{VI.\quad COMMENTS OF THE INTERNSHIP HOST ORGANIZATION}

\vspace{0.15cm}

The page entitled \textbf{COMMENTS OF THE INTERNSHIP HOST ORGANIZATION}
is reserved for the internship host organization.

It includes:

\begin{itemize}[leftmargin=0.8cm, itemsep=2pt]
    \item the student's full name;
    \item class;
    \item comments from the host organization;
    \item score;
    \item date;
    \item signature and official seal.
\end{itemize}

Students shall not alter or complete official evaluation fields unless
specifically authorized.

\vspace{0.3cm}

% ================================================================
% VII. SUPERVISOR'S EVALUATION AND DEFENSE RESULTS
% ================================================================

\noindent
\textbf{VII.\quad SUPERVISOR'S EVALUATION AND DEFENSE RESULTS}

\vspace{0.15cm}

The section entitled
\textbf{SUPERVISOR'S EVALUATION AND GRADUATION INTERNSHIP DEFENSE RESULTS}
is reserved for the supervisor and defense committee.

The evaluation may include:

\begin{itemize}[leftmargin=0.8cm, itemsep=2pt]
    \item Internship Progress and Student's Attitude;
    \item Report Content;
    \item Report Presentation;
    \item Other Comments;
    \item Supervisor's Opinion;
    \item Defense Result;
    \item Score.
\end{itemize}

Students shall not enter evaluation scores or defense results before they
are officially determined.

\vspace{0.3cm}

% ================================================================
% VIII. CHAPTERS, SECTIONS, AND SUBSECTIONS
% ================================================================

\noindent
\textbf{VIII.\quad CHAPTERS, SECTIONS, AND SUBSECTIONS}

\vspace{0.15cm}

The report structure shall be generated using the standard \LaTeX{}
commands:

\begin{verbatim}
\chapter{...}
\section{...}
\subsection{...}
\end{verbatim}

Chapter, section, and subsection numbers shall be generated automatically
by \LaTeX{}.

Students shall not manually type chapter or section numbers and shall not
reproduce heading styles using commands such as \verb|\textbf|,
\verb|\fontsize|, or \verb|\centering| instead of the appropriate
structural command.

Excessive subdivision should be avoided. In most cases, the hierarchy
Chapter--Section--Subsection is sufficient.

\vspace{0.3cm}

% ================================================================
% IX. ORGANIZATION OF THE MAIN CONTENT
% ================================================================

\noindent
\textbf{IX.\quad ORGANIZATION OF THE MAIN CONTENT}

\vspace{0.15cm}

The specific content of the report shall depend on the internship
organization and the assigned tasks. A typical structure may include:

\begin{itemize}[leftmargin=0.8cm, itemsep=2pt]
    \item \textbf{Chapter 1. Overview of the Internship Organization};
    \item \textbf{Chapter 2. Internship Activities and Assigned Tasks};
    \item \textbf{Chapter 3. Results, Experience, and Evaluation}.
\end{itemize}

These chapter titles are recommendations only and may be adapted to the
actual internship content.

The report shall reflect work actually performed, observed, analyzed, or
learned during the internship. Unrelated theoretical material should not
be included merely to increase report length.

\vspace{0.3cm}

% ================================================================
% X. TABLE OF CONTENTS
% ================================================================

\noindent
\textbf{X.\quad TABLE OF CONTENTS}

\vspace{0.15cm}

The Table of Contents shall be generated automatically using:

\begin{verbatim}
\tableofcontents
\end{verbatim}

Students shall not create or manually edit page numbers in the Table of
Contents.

The document shall be compiled sufficiently to ensure that all page
numbers and entries are correctly updated.

\vspace{0.3cm}

% ================================================================
% XI. LIST OF SYMBOLS AND ABBREVIATIONS
% ================================================================

\noindent
\textbf{XI.\quad LIST OF SYMBOLS AND ABBREVIATIONS}

\vspace{0.15cm}

Symbols and abbreviations that are used repeatedly in the report shall be
listed in the \textbf{List of Symbols and Abbreviations}.

The list may contain the following columns:

\begin{center}
\begin{tabular}{llll}
    No. & Symbol / Abbreviation & Unit & Definition
\end{tabular}
\end{center}

Only symbols and abbreviations actually used in the final report shall be
retained. Example entries contained in the template shall be removed if
they are not relevant to the student's report.

Each abbreviation should normally be defined when it first appears in
the main text.

\vspace{0.3cm}

% ================================================================
% XII. FIGURES, DIAGRAMS, AND PHOTOGRAPHS
% ================================================================

\noindent
\textbf{XII.\quad FIGURES, DIAGRAMS, AND PHOTOGRAPHS}

\vspace{0.15cm}

Figures shall be clear, legible, and directly relevant to the report
content.

Figures shall be inserted using the standard \verb|figure| environment.

Example:

\begin{verbatim}
\begin{figure}[H]
    \centering
    \includegraphics[width=0.8\textwidth]
                    {figures/company_structure.pdf}
    \caption{Organizational structure of the company}
    \label{fig:company-structure}
\end{figure}
\end{verbatim}

Figure captions shall be placed below the figure.

Each referenced figure shall have a unique \verb|\label| and shall be
referenced in the text using, for example:

\begin{verbatim}
Figure~\ref{fig:company-structure}
\end{verbatim}

Students shall not manually type figure numbers.

Figures reproduced or adapted from external sources shall include an
appropriate citation.

Photographs taken during the internship shall be relevant and sufficiently
clear. Students shall not include confidential, restricted, or proprietary
information without permission from the internship organization.

Vector graphics or high-resolution images are recommended whenever
possible.

\vspace{0.3cm}

% ================================================================
% XIII. TABLES
% ================================================================

\noindent
\textbf{XIII.\quad TABLES}

\vspace{0.15cm}

Tables shall be created using suitable \LaTeX{} environments such as
\verb|table|, \verb|tabular|, \verb|tabularx|, or \verb|longtable|.

Table captions shall be placed above the table.

Example:

\begin{verbatim}
\begin{table}[H]
    \centering
    \caption{Internship schedule}
    \label{tab:internship-schedule}

    \begin{tabular}{lll}
        \toprule
        Week & Task & Result \\
        \midrule
        1 & Orientation & Completed \\
        2 & System study & Completed \\
        \bottomrule
    \end{tabular}
\end{table}
\end{verbatim}

Tables shall be referenced using:

\begin{verbatim}
Table~\ref{tab:internship-schedule}
\end{verbatim}

Students shall not manually type table numbers.

Tables should be created directly in \LaTeX{} rather than inserted as
screenshots whenever practical.

\vspace{0.3cm}

% ================================================================
% XIV. EQUATIONS AND TECHNICAL QUANTITIES
% ================================================================

\noindent
\textbf{XIV.\quad EQUATIONS AND TECHNICAL QUANTITIES}

\vspace{0.15cm}

Mathematical expressions shall be typeset directly using \LaTeX{}.
Images or screenshots of equations shall not be used.

Short expressions may be written in inline mathematical mode, for example:

\begin{verbatim}
$R = 0.7\,\Omega$
\end{verbatim}

Important equations shall use a numbered equation environment.

Example:

\begin{verbatim}
\begin{equation}
    P = UI
    \label{eq:power}
\end{equation}
\end{verbatim}

Equations shall be referenced using:

\begin{verbatim}
Equation~\eqref{eq:power}
\end{verbatim}

Equation numbers shall not be manually entered.

SI units shall be used wherever applicable. The \verb|siunitx| package
should be used where appropriate.

Examples:

\begin{verbatim}
\SI{220}{\volt}
\SI{50}{\hertz}
\SI{2.5}{\ampere}
\end{verbatim}

Notation and units shall be consistent throughout the report.

\vspace{0.3cm}

% ================================================================
% XV. SOURCE CODE
% ================================================================

\noindent
\textbf{XV.\quad SOURCE CODE}

\vspace{0.15cm}

If the internship involves programming, PLCs, microcontrollers,
simulation, or software development, relevant source code may be included
in the report.

Source code shall be presented using the \verb|listings| environment or
another environment defined by the official template.

Example:

\begin{verbatim}
\begin{lstlisting}[language=st]
IF Start = TRUE THEN
    Motor := TRUE;
ELSE
    Motor := FALSE;
END_IF;
\end{lstlisting}
\end{verbatim}

Only code necessary for explaining the work should appear in the main
report. Long source-code listings should be moved to an appendix or
submitted as supplementary files if required.

Source code should not normally be inserted as screenshots.

\vspace{0.3cm}

% ================================================================
% XVI. ABSTRACT
% ================================================================

\noindent
\textbf{XVI.\quad ABSTRACT}

\vspace{0.15cm}

The Abstract shall briefly summarize:

\begin{itemize}[leftmargin=0.8cm, itemsep=2pt]
    \item the internship organization;
    \item the objectives of the internship;
    \item the main tasks performed;
    \item the main results, knowledge, and skills obtained.
\end{itemize}

The Abstract should be concise and should not reproduce the full content
of the report chapter by chapter.

\vspace{0.3cm}

% ================================================================
% XVII. REFERENCES AND CITATIONS
% ================================================================

\noindent
\textbf{XVII.\quad REFERENCES AND CITATIONS}

\vspace{0.15cm}

All books, journal papers, conference papers, standards, websites,
datasheets, manuals, technical documents, figures, tables, or other
external sources used in preparing the report shall be properly cited.

References shall be managed through the bibliography system provided with
the official template.

Bibliographic entries shall normally be stored in:

\begin{verbatim}
References.bib
\end{verbatim}

References shall be cited in the report using:

\begin{verbatim}
\cite{reference-key}
\end{verbatim}

Students shall not manually type reference numbers such as ``[1]'',
``[2]'', etc.

Titles of foreign-language publications shall remain in their original
language in the bibliography.

All borrowed ideas, data, figures, tables, equations, methods, or results
shall be properly acknowledged.

\vspace{0.3cm}

% ================================================================
% XVIII. CROSS-REFERENCING
% ================================================================

\noindent
\textbf{XVIII.\quad CROSS-REFERENCING}

\vspace{0.15cm}

Numbered elements shall be cross-referenced using the standard \LaTeX{}
mechanism.

Use the following commands where appropriate:

\begin{verbatim}
\label{...}
\ref{...}
\eqref{...}
\end{verbatim}

For example:

\begin{verbatim}
As shown in Figure~\ref{fig:system}, ...
\end{verbatim}

Students shall not manually type numbers such as ``Figure 2.3'' when
those numbers can be generated automatically.

\vspace{0.3cm}

% ================================================================
% XIX. PAGE NUMBERING
% ================================================================

\noindent
\textbf{XIX.\quad PAGE NUMBERING}

\vspace{0.15cm}

Cover pages and official evaluation forms at the beginning of the report
shall not display page numbers.

The main body of the report shall use Arabic page numbers beginning from
Chapter 1.

Page numbering shall be generated automatically by the official template.
Students shall not manually enter page numbers.

\vspace{0.3cm}

% ================================================================
% XX. PROJECT FILE ORGANIZATION
% ================================================================

\noindent
\textbf{XX.\quad PROJECT FILE ORGANIZATION}

\vspace{0.15cm}

Students should maintain a clear and organized \LaTeX{} project
structure. A recommended structure is:

\begin{verbatim}
InternshipReport/
|
|-- main.tex
|-- References.bib
|
|-- figures/
|   |-- company.pdf
|   |-- system.png
|   `-- ...
|
`-- chapters/
    |-- chapter1.tex
    |-- chapter2.tex
    `-- chapter3.tex
\end{verbatim}

Figures should normally be stored in the \verb|figures| directory
specified by the official template.

File and directory names should preferably contain Latin letters,
numbers, hyphens, or underscores.

Recommended examples include:

\begin{verbatim}
company_structure.pdf
motor_control.png
chapter_01.tex
\end{verbatim}

Spaces and unnecessary special characters in file names should be
avoided where possible.

\vspace{0.3cm}

% ================================================================
% XXI. DOCUMENT COMPILATION AND FINAL CHECK
% ================================================================

\noindent
\textbf{XXI.\quad DOCUMENT COMPILATION AND FINAL CHECK}

\vspace{0.15cm}

Before submission, the complete \LaTeX{} project shall compile
successfully and produce the final PDF report.

The final document shall not contain unresolved references such as:

\begin{verbatim}
??
[?]
\end{verbatim}

Before generating the final PDF, students shall check:

\begin{itemize}[leftmargin=0.8cm, itemsep=2pt]
    \item the Table of Contents;
    \item page numbering;
    \item chapter, section, and subsection numbering;
    \item figure numbering;
    \item table numbering;
    \item equation numbering;
    \item cross-references;
    \item citations;
    \item bibliography;
    \item font display;
    \item mathematical symbols;
    \item image quality;
    \item completeness of all required forms.
\end{itemize}

The document may be compiled using TeX Live, MiKTeX, or another
compatible \LaTeX{} environment. Regardless of the compilation
environment, the resulting PDF shall conform to the official Faculty
template.

\vspace{0.3cm}

% ================================================================
% XXII. CONTENT THAT STUDENTS MAY EDIT
% ================================================================

\noindent
\textbf{XXII.\quad CONTENT THAT STUDENTS MAY EDIT}

\vspace{0.15cm}

Students may edit or replace:

\begin{itemize}[leftmargin=0.8cm, itemsep=2pt]
    \item student name;
    \item class;
    \item internship organization;
    \item supervisor name;
    \item Abstract;
    \item List of Symbols and Abbreviations;
    \item chapter titles and chapter content;
    \item figures and photographs;
    \item tables;
    \item equations;
    \item citations and references.
\end{itemize}

Students shall not modify the main document formatting settings in the
preamble, including the document class, page geometry, line spacing,
heading formatting, table-of-contents formatting, font configuration,
page layout, or cover-page macros unless specifically instructed.

\vspace{0.3cm}

% ================================================================
% XXIII. FINAL SUBMISSION REQUIREMENTS
% ================================================================

\noindent
\textbf{XXIII.\quad FINAL SUBMISSION REQUIREMENTS}

\vspace{0.15cm}

Before submission, the final report shall:

\begin{itemize}[leftmargin=0.8cm, itemsep=2pt]
    \item compile without errors that affect the report content or presentation;
    \item contain all required forms and report sections;
    \item contain correct student, class, internship organization, and supervisor information;
    \item contain an updated Table of Contents, List of Figures, and List of Tables;
    \item contain no unused example content from the template;
    \item contain no placeholder text such as \textit{Title of Chapter 1},
    \textit{Section 1.1}, or other sample entries;
    \item contain all required citations and references;
    \item contain no unresolved cross-references.
\end{itemize}

If the Faculty requires submission of the \LaTeX{} source, the complete
project shall include all \verb|.tex| files, the bibliography database,
figures, and any other files required to reproduce the submitted PDF.

\vspace{0.4cm}

\noindent
\textbf{Important:}
Before submitting the final internship report, students shall remove or
comment out the following command in the main document:

\begin{verbatim}
\input{Regulations.tex}
\end{verbatim}

For example:

\begin{verbatim}
% \input{Regulations.tex}
\end{verbatim}

This removes the guidance pages from the final submitted report while
preserving the regulations file in the \LaTeX{} project for reference.

} % End 13pt content

\clearpage
 command before submitting
% the final internship report.
%
% =================================================================

\clearpage
\thispagestyle{empty}

\begin{center}
    {\fontsize{16pt}{20pt}\selectfont\bfseries
    SPECIFIC REGULATIONS FOR PREPARING\\
    GRADUATION INTERNSHIP REPORTS USING \LaTeX}
\end{center}

\vspace{0.5cm}

{\fontsize{13pt}{17pt}\selectfont
\setstretch{1.3}
\justifying

\noindent
\textbf{Note:}
This section is provided for guidance only and is included in the main
document using the command \verb|
% =================================================================
% REGULATIONS FOR PREPARING THE GRADUATION INTERNSHIP REPORT
% USING LATEX
% =================================================================
%
% This file is intended to be included from the main document with:
%
%     \input{Regulations.tex}
%
% IMPORTANT:
% This section is provided for guidance only. Students shall remove
% or comment out the \input{Regulations.tex} command before submitting
% the final internship report.
%
% =================================================================

\clearpage
\thispagestyle{empty}

\begin{center}
    {\fontsize{16pt}{20pt}\selectfont\bfseries
    SPECIFIC REGULATIONS FOR PREPARING\\
    GRADUATION INTERNSHIP REPORTS USING \LaTeX}
\end{center}

\vspace{0.5cm}

{\fontsize{13pt}{17pt}\selectfont
\setstretch{1.3}
\justifying

\noindent
\textbf{Note:}
This section is provided for guidance only and is included in the main
document using the command \verb|\input{Regulations.tex}|.
Students shall remove or comment out this command before submitting the
final internship report.

\vspace{0.4cm}

% ================================================================
% I. GENERAL REQUIREMENTS
% ================================================================

\noindent
\textbf{I.\quad GENERAL REQUIREMENTS}

\vspace{0.15cm}

The Graduation Internship Report shall be prepared using the official
\LaTeX{} template provided by the Faculty. Students shall preserve the
predefined page layout, typography, line spacing, heading styles,
numbering system, figure and table formatting, and bibliography settings
unless otherwise instructed by the Faculty or the supervisor.

The report shall be written in English.

Students are responsible for the academic and technical content of the
report. The official \LaTeX{} template is responsible for standardized
document formatting.

\vspace{0.3cm}

% ================================================================
% II. REPORT STRUCTURE
% ================================================================

\noindent
\textbf{II.\quad REPORT STRUCTURE}

\vspace{0.15cm}

The report shall normally be arranged in the following order:

\begin{enumerate}[label=\arabic*., leftmargin=1.2cm, itemsep=2pt]
    \item Outer Cover;
    \item Comments of the Internship Host Organization;
    \item Supervisor's Evaluation and Graduation Internship Defense Results;
    \item Table of Contents;
    \item List of Symbols and Abbreviations, if applicable;
    \item List of Figures;
    \item List of Tables;
    \item Abstract;
    \item Main Chapters of the Internship Report;
    \item References;
    \item Appendices, if applicable.
\end{enumerate}

The preliminary pages shall not be numbered as chapters. The main report
content shall be organized into consecutively numbered chapters beginning
with Chapter 1.

\vspace{0.3cm}

% ================================================================
% III. PAGE SIZE AND MARGINS
% ================================================================

\noindent
\textbf{III.\quad PAGE SIZE AND MARGINS}

\vspace{0.15cm}

The report shall be formatted on A4 paper with the following margins:

\begin{center}
\begin{tabular}{ll}
    Left margin:   & 3.5 cm \\
    Right margin:  & 2.0 cm \\
    Top margin:    & 2.5 cm \\
    Bottom margin: & 2.5 cm
\end{tabular}
\end{center}

These settings are predefined in the official template and shall not be
modified by students.

\vspace{0.3cm}

% ================================================================
% IV. FONT, FONT SIZE, AND LINE SPACING
% ================================================================

\noindent
\textbf{IV.\quad FONT, FONT SIZE, AND LINE SPACING}

\vspace{0.15cm}

The main text shall use a 13 pt font size with 1.3-line spacing.
A Times-family typeface shall be used in accordance with the official
template.

The text shall be fully justified.

Students shall not use repeated spaces, repeated manual line breaks, or
multiple blank lines to position text. Paragraph and heading spacing shall
be controlled by the template.

\vspace{0.3cm}

% ================================================================
% V. OUTER COVER
% ================================================================

\noindent
\textbf{V.\quad OUTER COVER}

\vspace{0.15cm}

The outer cover shall follow the official format provided in the template
and shall contain the following information:

\begin{itemize}[leftmargin=0.8cm, itemsep=2pt]
    \item Thai Nguyen University of Technology;
    \item Faculty of International Training;
    \item Engineering and Technology Department;
    \item University logo and Faculty logo;
    \item the title \textbf{INTERNSHIP REPORT};
    \item Internship Organization;
    \item Author;
    \item Class;
    \item Supervisor;
    \item place and year of completion.
\end{itemize}

Students shall only replace the required information fields and shall not
modify the cover-page layout unless instructed otherwise.

\vspace{0.3cm}

% ================================================================
% VI. COMMENTS OF THE INTERNSHIP HOST ORGANIZATION
% ================================================================

\noindent
\textbf{VI.\quad COMMENTS OF THE INTERNSHIP HOST ORGANIZATION}

\vspace{0.15cm}

The page entitled \textbf{COMMENTS OF THE INTERNSHIP HOST ORGANIZATION}
is reserved for the internship host organization.

It includes:

\begin{itemize}[leftmargin=0.8cm, itemsep=2pt]
    \item the student's full name;
    \item class;
    \item comments from the host organization;
    \item score;
    \item date;
    \item signature and official seal.
\end{itemize}

Students shall not alter or complete official evaluation fields unless
specifically authorized.

\vspace{0.3cm}

% ================================================================
% VII. SUPERVISOR'S EVALUATION AND DEFENSE RESULTS
% ================================================================

\noindent
\textbf{VII.\quad SUPERVISOR'S EVALUATION AND DEFENSE RESULTS}

\vspace{0.15cm}

The section entitled
\textbf{SUPERVISOR'S EVALUATION AND GRADUATION INTERNSHIP DEFENSE RESULTS}
is reserved for the supervisor and defense committee.

The evaluation may include:

\begin{itemize}[leftmargin=0.8cm, itemsep=2pt]
    \item Internship Progress and Student's Attitude;
    \item Report Content;
    \item Report Presentation;
    \item Other Comments;
    \item Supervisor's Opinion;
    \item Defense Result;
    \item Score.
\end{itemize}

Students shall not enter evaluation scores or defense results before they
are officially determined.

\vspace{0.3cm}

% ================================================================
% VIII. CHAPTERS, SECTIONS, AND SUBSECTIONS
% ================================================================

\noindent
\textbf{VIII.\quad CHAPTERS, SECTIONS, AND SUBSECTIONS}

\vspace{0.15cm}

The report structure shall be generated using the standard \LaTeX{}
commands:

\begin{verbatim}
\chapter{...}
\section{...}
\subsection{...}
\end{verbatim}

Chapter, section, and subsection numbers shall be generated automatically
by \LaTeX{}.

Students shall not manually type chapter or section numbers and shall not
reproduce heading styles using commands such as \verb|\textbf|,
\verb|\fontsize|, or \verb|\centering| instead of the appropriate
structural command.

Excessive subdivision should be avoided. In most cases, the hierarchy
Chapter--Section--Subsection is sufficient.

\vspace{0.3cm}

% ================================================================
% IX. ORGANIZATION OF THE MAIN CONTENT
% ================================================================

\noindent
\textbf{IX.\quad ORGANIZATION OF THE MAIN CONTENT}

\vspace{0.15cm}

The specific content of the report shall depend on the internship
organization and the assigned tasks. A typical structure may include:

\begin{itemize}[leftmargin=0.8cm, itemsep=2pt]
    \item \textbf{Chapter 1. Overview of the Internship Organization};
    \item \textbf{Chapter 2. Internship Activities and Assigned Tasks};
    \item \textbf{Chapter 3. Results, Experience, and Evaluation}.
\end{itemize}

These chapter titles are recommendations only and may be adapted to the
actual internship content.

The report shall reflect work actually performed, observed, analyzed, or
learned during the internship. Unrelated theoretical material should not
be included merely to increase report length.

\vspace{0.3cm}

% ================================================================
% X. TABLE OF CONTENTS
% ================================================================

\noindent
\textbf{X.\quad TABLE OF CONTENTS}

\vspace{0.15cm}

The Table of Contents shall be generated automatically using:

\begin{verbatim}
\tableofcontents
\end{verbatim}

Students shall not create or manually edit page numbers in the Table of
Contents.

The document shall be compiled sufficiently to ensure that all page
numbers and entries are correctly updated.

\vspace{0.3cm}

% ================================================================
% XI. LIST OF SYMBOLS AND ABBREVIATIONS
% ================================================================

\noindent
\textbf{XI.\quad LIST OF SYMBOLS AND ABBREVIATIONS}

\vspace{0.15cm}

Symbols and abbreviations that are used repeatedly in the report shall be
listed in the \textbf{List of Symbols and Abbreviations}.

The list may contain the following columns:

\begin{center}
\begin{tabular}{llll}
    No. & Symbol / Abbreviation & Unit & Definition
\end{tabular}
\end{center}

Only symbols and abbreviations actually used in the final report shall be
retained. Example entries contained in the template shall be removed if
they are not relevant to the student's report.

Each abbreviation should normally be defined when it first appears in
the main text.

\vspace{0.3cm}

% ================================================================
% XII. FIGURES, DIAGRAMS, AND PHOTOGRAPHS
% ================================================================

\noindent
\textbf{XII.\quad FIGURES, DIAGRAMS, AND PHOTOGRAPHS}

\vspace{0.15cm}

Figures shall be clear, legible, and directly relevant to the report
content.

Figures shall be inserted using the standard \verb|figure| environment.

Example:

\begin{verbatim}
\begin{figure}[H]
    \centering
    \includegraphics[width=0.8\textwidth]
                    {figures/company_structure.pdf}
    \caption{Organizational structure of the company}
    \label{fig:company-structure}
\end{figure}
\end{verbatim}

Figure captions shall be placed below the figure.

Each referenced figure shall have a unique \verb|\label| and shall be
referenced in the text using, for example:

\begin{verbatim}
Figure~\ref{fig:company-structure}
\end{verbatim}

Students shall not manually type figure numbers.

Figures reproduced or adapted from external sources shall include an
appropriate citation.

Photographs taken during the internship shall be relevant and sufficiently
clear. Students shall not include confidential, restricted, or proprietary
information without permission from the internship organization.

Vector graphics or high-resolution images are recommended whenever
possible.

\vspace{0.3cm}

% ================================================================
% XIII. TABLES
% ================================================================

\noindent
\textbf{XIII.\quad TABLES}

\vspace{0.15cm}

Tables shall be created using suitable \LaTeX{} environments such as
\verb|table|, \verb|tabular|, \verb|tabularx|, or \verb|longtable|.

Table captions shall be placed above the table.

Example:

\begin{verbatim}
\begin{table}[H]
    \centering
    \caption{Internship schedule}
    \label{tab:internship-schedule}

    \begin{tabular}{lll}
        \toprule
        Week & Task & Result \\
        \midrule
        1 & Orientation & Completed \\
        2 & System study & Completed \\
        \bottomrule
    \end{tabular}
\end{table}
\end{verbatim}

Tables shall be referenced using:

\begin{verbatim}
Table~\ref{tab:internship-schedule}
\end{verbatim}

Students shall not manually type table numbers.

Tables should be created directly in \LaTeX{} rather than inserted as
screenshots whenever practical.

\vspace{0.3cm}

% ================================================================
% XIV. EQUATIONS AND TECHNICAL QUANTITIES
% ================================================================

\noindent
\textbf{XIV.\quad EQUATIONS AND TECHNICAL QUANTITIES}

\vspace{0.15cm}

Mathematical expressions shall be typeset directly using \LaTeX{}.
Images or screenshots of equations shall not be used.

Short expressions may be written in inline mathematical mode, for example:

\begin{verbatim}
$R = 0.7\,\Omega$
\end{verbatim}

Important equations shall use a numbered equation environment.

Example:

\begin{verbatim}
\begin{equation}
    P = UI
    \label{eq:power}
\end{equation}
\end{verbatim}

Equations shall be referenced using:

\begin{verbatim}
Equation~\eqref{eq:power}
\end{verbatim}

Equation numbers shall not be manually entered.

SI units shall be used wherever applicable. The \verb|siunitx| package
should be used where appropriate.

Examples:

\begin{verbatim}
\SI{220}{\volt}
\SI{50}{\hertz}
\SI{2.5}{\ampere}
\end{verbatim}

Notation and units shall be consistent throughout the report.

\vspace{0.3cm}

% ================================================================
% XV. SOURCE CODE
% ================================================================

\noindent
\textbf{XV.\quad SOURCE CODE}

\vspace{0.15cm}

If the internship involves programming, PLCs, microcontrollers,
simulation, or software development, relevant source code may be included
in the report.

Source code shall be presented using the \verb|listings| environment or
another environment defined by the official template.

Example:

\begin{verbatim}
\begin{lstlisting}[language=st]
IF Start = TRUE THEN
    Motor := TRUE;
ELSE
    Motor := FALSE;
END_IF;
\end{lstlisting}
\end{verbatim}

Only code necessary for explaining the work should appear in the main
report. Long source-code listings should be moved to an appendix or
submitted as supplementary files if required.

Source code should not normally be inserted as screenshots.

\vspace{0.3cm}

% ================================================================
% XVI. ABSTRACT
% ================================================================

\noindent
\textbf{XVI.\quad ABSTRACT}

\vspace{0.15cm}

The Abstract shall briefly summarize:

\begin{itemize}[leftmargin=0.8cm, itemsep=2pt]
    \item the internship organization;
    \item the objectives of the internship;
    \item the main tasks performed;
    \item the main results, knowledge, and skills obtained.
\end{itemize}

The Abstract should be concise and should not reproduce the full content
of the report chapter by chapter.

\vspace{0.3cm}

% ================================================================
% XVII. REFERENCES AND CITATIONS
% ================================================================

\noindent
\textbf{XVII.\quad REFERENCES AND CITATIONS}

\vspace{0.15cm}

All books, journal papers, conference papers, standards, websites,
datasheets, manuals, technical documents, figures, tables, or other
external sources used in preparing the report shall be properly cited.

References shall be managed through the bibliography system provided with
the official template.

Bibliographic entries shall normally be stored in:

\begin{verbatim}
References.bib
\end{verbatim}

References shall be cited in the report using:

\begin{verbatim}
\cite{reference-key}
\end{verbatim}

Students shall not manually type reference numbers such as ``[1]'',
``[2]'', etc.

Titles of foreign-language publications shall remain in their original
language in the bibliography.

All borrowed ideas, data, figures, tables, equations, methods, or results
shall be properly acknowledged.

\vspace{0.3cm}

% ================================================================
% XVIII. CROSS-REFERENCING
% ================================================================

\noindent
\textbf{XVIII.\quad CROSS-REFERENCING}

\vspace{0.15cm}

Numbered elements shall be cross-referenced using the standard \LaTeX{}
mechanism.

Use the following commands where appropriate:

\begin{verbatim}
\label{...}
\ref{...}
\eqref{...}
\end{verbatim}

For example:

\begin{verbatim}
As shown in Figure~\ref{fig:system}, ...
\end{verbatim}

Students shall not manually type numbers such as ``Figure 2.3'' when
those numbers can be generated automatically.

\vspace{0.3cm}

% ================================================================
% XIX. PAGE NUMBERING
% ================================================================

\noindent
\textbf{XIX.\quad PAGE NUMBERING}

\vspace{0.15cm}

Cover pages and official evaluation forms at the beginning of the report
shall not display page numbers.

The main body of the report shall use Arabic page numbers beginning from
Chapter 1.

Page numbering shall be generated automatically by the official template.
Students shall not manually enter page numbers.

\vspace{0.3cm}

% ================================================================
% XX. PROJECT FILE ORGANIZATION
% ================================================================

\noindent
\textbf{XX.\quad PROJECT FILE ORGANIZATION}

\vspace{0.15cm}

Students should maintain a clear and organized \LaTeX{} project
structure. A recommended structure is:

\begin{verbatim}
InternshipReport/
|
|-- main.tex
|-- References.bib
|
|-- figures/
|   |-- company.pdf
|   |-- system.png
|   `-- ...
|
`-- chapters/
    |-- chapter1.tex
    |-- chapter2.tex
    `-- chapter3.tex
\end{verbatim}

Figures should normally be stored in the \verb|figures| directory
specified by the official template.

File and directory names should preferably contain Latin letters,
numbers, hyphens, or underscores.

Recommended examples include:

\begin{verbatim}
company_structure.pdf
motor_control.png
chapter_01.tex
\end{verbatim}

Spaces and unnecessary special characters in file names should be
avoided where possible.

\vspace{0.3cm}

% ================================================================
% XXI. DOCUMENT COMPILATION AND FINAL CHECK
% ================================================================

\noindent
\textbf{XXI.\quad DOCUMENT COMPILATION AND FINAL CHECK}

\vspace{0.15cm}

Before submission, the complete \LaTeX{} project shall compile
successfully and produce the final PDF report.

The final document shall not contain unresolved references such as:

\begin{verbatim}
??
[?]
\end{verbatim}

Before generating the final PDF, students shall check:

\begin{itemize}[leftmargin=0.8cm, itemsep=2pt]
    \item the Table of Contents;
    \item page numbering;
    \item chapter, section, and subsection numbering;
    \item figure numbering;
    \item table numbering;
    \item equation numbering;
    \item cross-references;
    \item citations;
    \item bibliography;
    \item font display;
    \item mathematical symbols;
    \item image quality;
    \item completeness of all required forms.
\end{itemize}

The document may be compiled using TeX Live, MiKTeX, or another
compatible \LaTeX{} environment. Regardless of the compilation
environment, the resulting PDF shall conform to the official Faculty
template.

\vspace{0.3cm}

% ================================================================
% XXII. CONTENT THAT STUDENTS MAY EDIT
% ================================================================

\noindent
\textbf{XXII.\quad CONTENT THAT STUDENTS MAY EDIT}

\vspace{0.15cm}

Students may edit or replace:

\begin{itemize}[leftmargin=0.8cm, itemsep=2pt]
    \item student name;
    \item class;
    \item internship organization;
    \item supervisor name;
    \item Abstract;
    \item List of Symbols and Abbreviations;
    \item chapter titles and chapter content;
    \item figures and photographs;
    \item tables;
    \item equations;
    \item citations and references.
\end{itemize}

Students shall not modify the main document formatting settings in the
preamble, including the document class, page geometry, line spacing,
heading formatting, table-of-contents formatting, font configuration,
page layout, or cover-page macros unless specifically instructed.

\vspace{0.3cm}

% ================================================================
% XXIII. FINAL SUBMISSION REQUIREMENTS
% ================================================================

\noindent
\textbf{XXIII.\quad FINAL SUBMISSION REQUIREMENTS}

\vspace{0.15cm}

Before submission, the final report shall:

\begin{itemize}[leftmargin=0.8cm, itemsep=2pt]
    \item compile without errors that affect the report content or presentation;
    \item contain all required forms and report sections;
    \item contain correct student, class, internship organization, and supervisor information;
    \item contain an updated Table of Contents, List of Figures, and List of Tables;
    \item contain no unused example content from the template;
    \item contain no placeholder text such as \textit{Title of Chapter 1},
    \textit{Section 1.1}, or other sample entries;
    \item contain all required citations and references;
    \item contain no unresolved cross-references.
\end{itemize}

If the Faculty requires submission of the \LaTeX{} source, the complete
project shall include all \verb|.tex| files, the bibliography database,
figures, and any other files required to reproduce the submitted PDF.

\vspace{0.4cm}

\noindent
\textbf{Important:}
Before submitting the final internship report, students shall remove or
comment out the following command in the main document:

\begin{verbatim}
\input{Regulations.tex}
\end{verbatim}

For example:

\begin{verbatim}
% \input{Regulations.tex}
\end{verbatim}

This removes the guidance pages from the final submitted report while
preserving the regulations file in the \LaTeX{} project for reference.

} % End 13pt content

\clearpage
|.
Students shall remove or comment out this command before submitting the
final internship report.

\vspace{0.4cm}

% ================================================================
% I. GENERAL REQUIREMENTS
% ================================================================

\noindent
\textbf{I.\quad GENERAL REQUIREMENTS}

\vspace{0.15cm}

The Graduation Internship Report shall be prepared using the official
\LaTeX{} template provided by the Faculty. Students shall preserve the
predefined page layout, typography, line spacing, heading styles,
numbering system, figure and table formatting, and bibliography settings
unless otherwise instructed by the Faculty or the supervisor.

The report shall be written in English.

Students are responsible for the academic and technical content of the
report. The official \LaTeX{} template is responsible for standardized
document formatting.

\vspace{0.3cm}

% ================================================================
% II. REPORT STRUCTURE
% ================================================================

\noindent
\textbf{II.\quad REPORT STRUCTURE}

\vspace{0.15cm}

The report shall normally be arranged in the following order:

\begin{enumerate}[label=\arabic*., leftmargin=1.2cm, itemsep=2pt]
    \item Outer Cover;
    \item Comments of the Internship Host Organization;
    \item Supervisor's Evaluation and Graduation Internship Defense Results;
    \item Table of Contents;
    \item List of Symbols and Abbreviations, if applicable;
    \item List of Figures;
    \item List of Tables;
    \item Abstract;
    \item Main Chapters of the Internship Report;
    \item References;
    \item Appendices, if applicable.
\end{enumerate}

The preliminary pages shall not be numbered as chapters. The main report
content shall be organized into consecutively numbered chapters beginning
with Chapter 1.

\vspace{0.3cm}

% ================================================================
% III. PAGE SIZE AND MARGINS
% ================================================================

\noindent
\textbf{III.\quad PAGE SIZE AND MARGINS}

\vspace{0.15cm}

The report shall be formatted on A4 paper with the following margins:

\begin{center}
\begin{tabular}{ll}
    Left margin:   & 3.5 cm \\
    Right margin:  & 2.0 cm \\
    Top margin:    & 2.5 cm \\
    Bottom margin: & 2.5 cm
\end{tabular}
\end{center}

These settings are predefined in the official template and shall not be
modified by students.

\vspace{0.3cm}

% ================================================================
% IV. FONT, FONT SIZE, AND LINE SPACING
% ================================================================

\noindent
\textbf{IV.\quad FONT, FONT SIZE, AND LINE SPACING}

\vspace{0.15cm}

The main text shall use a 13 pt font size with 1.3-line spacing.
A Times-family typeface shall be used in accordance with the official
template.

The text shall be fully justified.

Students shall not use repeated spaces, repeated manual line breaks, or
multiple blank lines to position text. Paragraph and heading spacing shall
be controlled by the template.

\vspace{0.3cm}

% ================================================================
% V. OUTER COVER
% ================================================================

\noindent
\textbf{V.\quad OUTER COVER}

\vspace{0.15cm}

The outer cover shall follow the official format provided in the template
and shall contain the following information:

\begin{itemize}[leftmargin=0.8cm, itemsep=2pt]
    \item Thai Nguyen University of Technology;
    \item Faculty of International Training;
    \item Engineering and Technology Department;
    \item University logo and Faculty logo;
    \item the title \textbf{INTERNSHIP REPORT};
    \item Internship Organization;
    \item Author;
    \item Class;
    \item Supervisor;
    \item place and year of completion.
\end{itemize}

Students shall only replace the required information fields and shall not
modify the cover-page layout unless instructed otherwise.

\vspace{0.3cm}

% ================================================================
% VI. COMMENTS OF THE INTERNSHIP HOST ORGANIZATION
% ================================================================

\noindent
\textbf{VI.\quad COMMENTS OF THE INTERNSHIP HOST ORGANIZATION}

\vspace{0.15cm}

The page entitled \textbf{COMMENTS OF THE INTERNSHIP HOST ORGANIZATION}
is reserved for the internship host organization.

It includes:

\begin{itemize}[leftmargin=0.8cm, itemsep=2pt]
    \item the student's full name;
    \item class;
    \item comments from the host organization;
    \item score;
    \item date;
    \item signature and official seal.
\end{itemize}

Students shall not alter or complete official evaluation fields unless
specifically authorized.

\vspace{0.3cm}

% ================================================================
% VII. SUPERVISOR'S EVALUATION AND DEFENSE RESULTS
% ================================================================

\noindent
\textbf{VII.\quad SUPERVISOR'S EVALUATION AND DEFENSE RESULTS}

\vspace{0.15cm}

The section entitled
\textbf{SUPERVISOR'S EVALUATION AND GRADUATION INTERNSHIP DEFENSE RESULTS}
is reserved for the supervisor and defense committee.

The evaluation may include:

\begin{itemize}[leftmargin=0.8cm, itemsep=2pt]
    \item Internship Progress and Student's Attitude;
    \item Report Content;
    \item Report Presentation;
    \item Other Comments;
    \item Supervisor's Opinion;
    \item Defense Result;
    \item Score.
\end{itemize}

Students shall not enter evaluation scores or defense results before they
are officially determined.

\vspace{0.3cm}

% ================================================================
% VIII. CHAPTERS, SECTIONS, AND SUBSECTIONS
% ================================================================

\noindent
\textbf{VIII.\quad CHAPTERS, SECTIONS, AND SUBSECTIONS}

\vspace{0.15cm}

The report structure shall be generated using the standard \LaTeX{}
commands:

\begin{verbatim}
\chapter{...}
\section{...}
\subsection{...}
\end{verbatim}

Chapter, section, and subsection numbers shall be generated automatically
by \LaTeX{}.

Students shall not manually type chapter or section numbers and shall not
reproduce heading styles using commands such as \verb|\textbf|,
\verb|\fontsize|, or \verb|\centering| instead of the appropriate
structural command.

Excessive subdivision should be avoided. In most cases, the hierarchy
Chapter--Section--Subsection is sufficient.

\vspace{0.3cm}

% ================================================================
% IX. ORGANIZATION OF THE MAIN CONTENT
% ================================================================

\noindent
\textbf{IX.\quad ORGANIZATION OF THE MAIN CONTENT}

\vspace{0.15cm}

The specific content of the report shall depend on the internship
organization and the assigned tasks. A typical structure may include:

\begin{itemize}[leftmargin=0.8cm, itemsep=2pt]
    \item \textbf{Chapter 1. Overview of the Internship Organization};
    \item \textbf{Chapter 2. Internship Activities and Assigned Tasks};
    \item \textbf{Chapter 3. Results, Experience, and Evaluation}.
\end{itemize}

These chapter titles are recommendations only and may be adapted to the
actual internship content.

The report shall reflect work actually performed, observed, analyzed, or
learned during the internship. Unrelated theoretical material should not
be included merely to increase report length.

\vspace{0.3cm}

% ================================================================
% X. TABLE OF CONTENTS
% ================================================================

\noindent
\textbf{X.\quad TABLE OF CONTENTS}

\vspace{0.15cm}

The Table of Contents shall be generated automatically using:

\begin{verbatim}
\tableofcontents
\end{verbatim}

Students shall not create or manually edit page numbers in the Table of
Contents.

The document shall be compiled sufficiently to ensure that all page
numbers and entries are correctly updated.

\vspace{0.3cm}

% ================================================================
% XI. LIST OF SYMBOLS AND ABBREVIATIONS
% ================================================================

\noindent
\textbf{XI.\quad LIST OF SYMBOLS AND ABBREVIATIONS}

\vspace{0.15cm}

Symbols and abbreviations that are used repeatedly in the report shall be
listed in the \textbf{List of Symbols and Abbreviations}.

The list may contain the following columns:

\begin{center}
\begin{tabular}{llll}
    No. & Symbol / Abbreviation & Unit & Definition
\end{tabular}
\end{center}

Only symbols and abbreviations actually used in the final report shall be
retained. Example entries contained in the template shall be removed if
they are not relevant to the student's report.

Each abbreviation should normally be defined when it first appears in
the main text.

\vspace{0.3cm}

% ================================================================
% XII. FIGURES, DIAGRAMS, AND PHOTOGRAPHS
% ================================================================

\noindent
\textbf{XII.\quad FIGURES, DIAGRAMS, AND PHOTOGRAPHS}

\vspace{0.15cm}

Figures shall be clear, legible, and directly relevant to the report
content.

Figures shall be inserted using the standard \verb|figure| environment.

Example:

\begin{verbatim}
\begin{figure}[H]
    \centering
    \includegraphics[width=0.8\textwidth]
                    {figures/company_structure.pdf}
    \caption{Organizational structure of the company}
    \label{fig:company-structure}
\end{figure}
\end{verbatim}

Figure captions shall be placed below the figure.

Each referenced figure shall have a unique \verb|\label| and shall be
referenced in the text using, for example:

\begin{verbatim}
Figure~\ref{fig:company-structure}
\end{verbatim}

Students shall not manually type figure numbers.

Figures reproduced or adapted from external sources shall include an
appropriate citation.

Photographs taken during the internship shall be relevant and sufficiently
clear. Students shall not include confidential, restricted, or proprietary
information without permission from the internship organization.

Vector graphics or high-resolution images are recommended whenever
possible.

\vspace{0.3cm}

% ================================================================
% XIII. TABLES
% ================================================================

\noindent
\textbf{XIII.\quad TABLES}

\vspace{0.15cm}

Tables shall be created using suitable \LaTeX{} environments such as
\verb|table|, \verb|tabular|, \verb|tabularx|, or \verb|longtable|.

Table captions shall be placed above the table.

Example:

\begin{verbatim}
\begin{table}[H]
    \centering
    \caption{Internship schedule}
    \label{tab:internship-schedule}

    \begin{tabular}{lll}
        \toprule
        Week & Task & Result \\
        \midrule
        1 & Orientation & Completed \\
        2 & System study & Completed \\
        \bottomrule
    \end{tabular}
\end{table}
\end{verbatim}

Tables shall be referenced using:

\begin{verbatim}
Table~\ref{tab:internship-schedule}
\end{verbatim}

Students shall not manually type table numbers.

Tables should be created directly in \LaTeX{} rather than inserted as
screenshots whenever practical.

\vspace{0.3cm}

% ================================================================
% XIV. EQUATIONS AND TECHNICAL QUANTITIES
% ================================================================

\noindent
\textbf{XIV.\quad EQUATIONS AND TECHNICAL QUANTITIES}

\vspace{0.15cm}

Mathematical expressions shall be typeset directly using \LaTeX{}.
Images or screenshots of equations shall not be used.

Short expressions may be written in inline mathematical mode, for example:

\begin{verbatim}
$R = 0.7\,\Omega$
\end{verbatim}

Important equations shall use a numbered equation environment.

Example:

\begin{verbatim}
\begin{equation}
    P = UI
    \label{eq:power}
\end{equation}
\end{verbatim}

Equations shall be referenced using:

\begin{verbatim}
Equation~\eqref{eq:power}
\end{verbatim}

Equation numbers shall not be manually entered.

SI units shall be used wherever applicable. The \verb|siunitx| package
should be used where appropriate.

Examples:

\begin{verbatim}
\SI{220}{\volt}
\SI{50}{\hertz}
\SI{2.5}{\ampere}
\end{verbatim}

Notation and units shall be consistent throughout the report.

\vspace{0.3cm}

% ================================================================
% XV. SOURCE CODE
% ================================================================

\noindent
\textbf{XV.\quad SOURCE CODE}

\vspace{0.15cm}

If the internship involves programming, PLCs, microcontrollers,
simulation, or software development, relevant source code may be included
in the report.

Source code shall be presented using the \verb|listings| environment or
another environment defined by the official template.

Example:

\begin{verbatim}
\begin{lstlisting}[language=st]
IF Start = TRUE THEN
    Motor := TRUE;
ELSE
    Motor := FALSE;
END_IF;
\end{lstlisting}
\end{verbatim}

Only code necessary for explaining the work should appear in the main
report. Long source-code listings should be moved to an appendix or
submitted as supplementary files if required.

Source code should not normally be inserted as screenshots.

\vspace{0.3cm}

% ================================================================
% XVI. ABSTRACT
% ================================================================

\noindent
\textbf{XVI.\quad ABSTRACT}

\vspace{0.15cm}

The Abstract shall briefly summarize:

\begin{itemize}[leftmargin=0.8cm, itemsep=2pt]
    \item the internship organization;
    \item the objectives of the internship;
    \item the main tasks performed;
    \item the main results, knowledge, and skills obtained.
\end{itemize}

The Abstract should be concise and should not reproduce the full content
of the report chapter by chapter.

\vspace{0.3cm}

% ================================================================
% XVII. REFERENCES AND CITATIONS
% ================================================================

\noindent
\textbf{XVII.\quad REFERENCES AND CITATIONS}

\vspace{0.15cm}

All books, journal papers, conference papers, standards, websites,
datasheets, manuals, technical documents, figures, tables, or other
external sources used in preparing the report shall be properly cited.

References shall be managed through the bibliography system provided with
the official template.

Bibliographic entries shall normally be stored in:

\begin{verbatim}
References.bib
\end{verbatim}

References shall be cited in the report using:

\begin{verbatim}
\cite{reference-key}
\end{verbatim}

Students shall not manually type reference numbers such as ``[1]'',
``[2]'', etc.

Titles of foreign-language publications shall remain in their original
language in the bibliography.

All borrowed ideas, data, figures, tables, equations, methods, or results
shall be properly acknowledged.

\vspace{0.3cm}

% ================================================================
% XVIII. CROSS-REFERENCING
% ================================================================

\noindent
\textbf{XVIII.\quad CROSS-REFERENCING}

\vspace{0.15cm}

Numbered elements shall be cross-referenced using the standard \LaTeX{}
mechanism.

Use the following commands where appropriate:

\begin{verbatim}
\label{...}
\ref{...}
\eqref{...}
\end{verbatim}

For example:

\begin{verbatim}
As shown in Figure~\ref{fig:system}, ...
\end{verbatim}

Students shall not manually type numbers such as ``Figure 2.3'' when
those numbers can be generated automatically.

\vspace{0.3cm}

% ================================================================
% XIX. PAGE NUMBERING
% ================================================================

\noindent
\textbf{XIX.\quad PAGE NUMBERING}

\vspace{0.15cm}

Cover pages and official evaluation forms at the beginning of the report
shall not display page numbers.

The main body of the report shall use Arabic page numbers beginning from
Chapter 1.

Page numbering shall be generated automatically by the official template.
Students shall not manually enter page numbers.

\vspace{0.3cm}

% ================================================================
% XX. PROJECT FILE ORGANIZATION
% ================================================================

\noindent
\textbf{XX.\quad PROJECT FILE ORGANIZATION}

\vspace{0.15cm}

Students should maintain a clear and organized \LaTeX{} project
structure. A recommended structure is:

\begin{verbatim}
InternshipReport/
|
|-- main.tex
|-- References.bib
|
|-- figures/
|   |-- company.pdf
|   |-- system.png
|   `-- ...
|
`-- chapters/
    |-- chapter1.tex
    |-- chapter2.tex
    `-- chapter3.tex
\end{verbatim}

Figures should normally be stored in the \verb|figures| directory
specified by the official template.

File and directory names should preferably contain Latin letters,
numbers, hyphens, or underscores.

Recommended examples include:

\begin{verbatim}
company_structure.pdf
motor_control.png
chapter_01.tex
\end{verbatim}

Spaces and unnecessary special characters in file names should be
avoided where possible.

\vspace{0.3cm}

% ================================================================
% XXI. DOCUMENT COMPILATION AND FINAL CHECK
% ================================================================

\noindent
\textbf{XXI.\quad DOCUMENT COMPILATION AND FINAL CHECK}

\vspace{0.15cm}

Before submission, the complete \LaTeX{} project shall compile
successfully and produce the final PDF report.

The final document shall not contain unresolved references such as:

\begin{verbatim}
??
[?]
\end{verbatim}

Before generating the final PDF, students shall check:

\begin{itemize}[leftmargin=0.8cm, itemsep=2pt]
    \item the Table of Contents;
    \item page numbering;
    \item chapter, section, and subsection numbering;
    \item figure numbering;
    \item table numbering;
    \item equation numbering;
    \item cross-references;
    \item citations;
    \item bibliography;
    \item font display;
    \item mathematical symbols;
    \item image quality;
    \item completeness of all required forms.
\end{itemize}

The document may be compiled using TeX Live, MiKTeX, or another
compatible \LaTeX{} environment. Regardless of the compilation
environment, the resulting PDF shall conform to the official Faculty
template.

\vspace{0.3cm}

% ================================================================
% XXII. CONTENT THAT STUDENTS MAY EDIT
% ================================================================

\noindent
\textbf{XXII.\quad CONTENT THAT STUDENTS MAY EDIT}

\vspace{0.15cm}

Students may edit or replace:

\begin{itemize}[leftmargin=0.8cm, itemsep=2pt]
    \item student name;
    \item class;
    \item internship organization;
    \item supervisor name;
    \item Abstract;
    \item List of Symbols and Abbreviations;
    \item chapter titles and chapter content;
    \item figures and photographs;
    \item tables;
    \item equations;
    \item citations and references.
\end{itemize}

Students shall not modify the main document formatting settings in the
preamble, including the document class, page geometry, line spacing,
heading formatting, table-of-contents formatting, font configuration,
page layout, or cover-page macros unless specifically instructed.

\vspace{0.3cm}

% ================================================================
% XXIII. FINAL SUBMISSION REQUIREMENTS
% ================================================================

\noindent
\textbf{XXIII.\quad FINAL SUBMISSION REQUIREMENTS}

\vspace{0.15cm}

Before submission, the final report shall:

\begin{itemize}[leftmargin=0.8cm, itemsep=2pt]
    \item compile without errors that affect the report content or presentation;
    \item contain all required forms and report sections;
    \item contain correct student, class, internship organization, and supervisor information;
    \item contain an updated Table of Contents, List of Figures, and List of Tables;
    \item contain no unused example content from the template;
    \item contain no placeholder text such as \textit{Title of Chapter 1},
    \textit{Section 1.1}, or other sample entries;
    \item contain all required citations and references;
    \item contain no unresolved cross-references.
\end{itemize}

If the Faculty requires submission of the \LaTeX{} source, the complete
project shall include all \verb|.tex| files, the bibliography database,
figures, and any other files required to reproduce the submitted PDF.

\vspace{0.4cm}

\noindent
\textbf{Important:}
Before submitting the final internship report, students shall remove or
comment out the following command in the main document:

\begin{verbatim}

% =================================================================
% REGULATIONS FOR PREPARING THE GRADUATION INTERNSHIP REPORT
% USING LATEX
% =================================================================
%
% This file is intended to be included from the main document with:
%
%     \input{Regulations.tex}
%
% IMPORTANT:
% This section is provided for guidance only. Students shall remove
% or comment out the \input{Regulations.tex} command before submitting
% the final internship report.
%
% =================================================================

\clearpage
\thispagestyle{empty}

\begin{center}
    {\fontsize{16pt}{20pt}\selectfont\bfseries
    SPECIFIC REGULATIONS FOR PREPARING\\
    GRADUATION INTERNSHIP REPORTS USING \LaTeX}
\end{center}

\vspace{0.5cm}

{\fontsize{13pt}{17pt}\selectfont
\setstretch{1.3}
\justifying

\noindent
\textbf{Note:}
This section is provided for guidance only and is included in the main
document using the command \verb|\input{Regulations.tex}|.
Students shall remove or comment out this command before submitting the
final internship report.

\vspace{0.4cm}

% ================================================================
% I. GENERAL REQUIREMENTS
% ================================================================

\noindent
\textbf{I.\quad GENERAL REQUIREMENTS}

\vspace{0.15cm}

The Graduation Internship Report shall be prepared using the official
\LaTeX{} template provided by the Faculty. Students shall preserve the
predefined page layout, typography, line spacing, heading styles,
numbering system, figure and table formatting, and bibliography settings
unless otherwise instructed by the Faculty or the supervisor.

The report shall be written in English.

Students are responsible for the academic and technical content of the
report. The official \LaTeX{} template is responsible for standardized
document formatting.

\vspace{0.3cm}

% ================================================================
% II. REPORT STRUCTURE
% ================================================================

\noindent
\textbf{II.\quad REPORT STRUCTURE}

\vspace{0.15cm}

The report shall normally be arranged in the following order:

\begin{enumerate}[label=\arabic*., leftmargin=1.2cm, itemsep=2pt]
    \item Outer Cover;
    \item Comments of the Internship Host Organization;
    \item Supervisor's Evaluation and Graduation Internship Defense Results;
    \item Table of Contents;
    \item List of Symbols and Abbreviations, if applicable;
    \item List of Figures;
    \item List of Tables;
    \item Abstract;
    \item Main Chapters of the Internship Report;
    \item References;
    \item Appendices, if applicable.
\end{enumerate}

The preliminary pages shall not be numbered as chapters. The main report
content shall be organized into consecutively numbered chapters beginning
with Chapter 1.

\vspace{0.3cm}

% ================================================================
% III. PAGE SIZE AND MARGINS
% ================================================================

\noindent
\textbf{III.\quad PAGE SIZE AND MARGINS}

\vspace{0.15cm}

The report shall be formatted on A4 paper with the following margins:

\begin{center}
\begin{tabular}{ll}
    Left margin:   & 3.5 cm \\
    Right margin:  & 2.0 cm \\
    Top margin:    & 2.5 cm \\
    Bottom margin: & 2.5 cm
\end{tabular}
\end{center}

These settings are predefined in the official template and shall not be
modified by students.

\vspace{0.3cm}

% ================================================================
% IV. FONT, FONT SIZE, AND LINE SPACING
% ================================================================

\noindent
\textbf{IV.\quad FONT, FONT SIZE, AND LINE SPACING}

\vspace{0.15cm}

The main text shall use a 13 pt font size with 1.3-line spacing.
A Times-family typeface shall be used in accordance with the official
template.

The text shall be fully justified.

Students shall not use repeated spaces, repeated manual line breaks, or
multiple blank lines to position text. Paragraph and heading spacing shall
be controlled by the template.

\vspace{0.3cm}

% ================================================================
% V. OUTER COVER
% ================================================================

\noindent
\textbf{V.\quad OUTER COVER}

\vspace{0.15cm}

The outer cover shall follow the official format provided in the template
and shall contain the following information:

\begin{itemize}[leftmargin=0.8cm, itemsep=2pt]
    \item Thai Nguyen University of Technology;
    \item Faculty of International Training;
    \item Engineering and Technology Department;
    \item University logo and Faculty logo;
    \item the title \textbf{INTERNSHIP REPORT};
    \item Internship Organization;
    \item Author;
    \item Class;
    \item Supervisor;
    \item place and year of completion.
\end{itemize}

Students shall only replace the required information fields and shall not
modify the cover-page layout unless instructed otherwise.

\vspace{0.3cm}

% ================================================================
% VI. COMMENTS OF THE INTERNSHIP HOST ORGANIZATION
% ================================================================

\noindent
\textbf{VI.\quad COMMENTS OF THE INTERNSHIP HOST ORGANIZATION}

\vspace{0.15cm}

The page entitled \textbf{COMMENTS OF THE INTERNSHIP HOST ORGANIZATION}
is reserved for the internship host organization.

It includes:

\begin{itemize}[leftmargin=0.8cm, itemsep=2pt]
    \item the student's full name;
    \item class;
    \item comments from the host organization;
    \item score;
    \item date;
    \item signature and official seal.
\end{itemize}

Students shall not alter or complete official evaluation fields unless
specifically authorized.

\vspace{0.3cm}

% ================================================================
% VII. SUPERVISOR'S EVALUATION AND DEFENSE RESULTS
% ================================================================

\noindent
\textbf{VII.\quad SUPERVISOR'S EVALUATION AND DEFENSE RESULTS}

\vspace{0.15cm}

The section entitled
\textbf{SUPERVISOR'S EVALUATION AND GRADUATION INTERNSHIP DEFENSE RESULTS}
is reserved for the supervisor and defense committee.

The evaluation may include:

\begin{itemize}[leftmargin=0.8cm, itemsep=2pt]
    \item Internship Progress and Student's Attitude;
    \item Report Content;
    \item Report Presentation;
    \item Other Comments;
    \item Supervisor's Opinion;
    \item Defense Result;
    \item Score.
\end{itemize}

Students shall not enter evaluation scores or defense results before they
are officially determined.

\vspace{0.3cm}

% ================================================================
% VIII. CHAPTERS, SECTIONS, AND SUBSECTIONS
% ================================================================

\noindent
\textbf{VIII.\quad CHAPTERS, SECTIONS, AND SUBSECTIONS}

\vspace{0.15cm}

The report structure shall be generated using the standard \LaTeX{}
commands:

\begin{verbatim}
\chapter{...}
\section{...}
\subsection{...}
\end{verbatim}

Chapter, section, and subsection numbers shall be generated automatically
by \LaTeX{}.

Students shall not manually type chapter or section numbers and shall not
reproduce heading styles using commands such as \verb|\textbf|,
\verb|\fontsize|, or \verb|\centering| instead of the appropriate
structural command.

Excessive subdivision should be avoided. In most cases, the hierarchy
Chapter--Section--Subsection is sufficient.

\vspace{0.3cm}

% ================================================================
% IX. ORGANIZATION OF THE MAIN CONTENT
% ================================================================

\noindent
\textbf{IX.\quad ORGANIZATION OF THE MAIN CONTENT}

\vspace{0.15cm}

The specific content of the report shall depend on the internship
organization and the assigned tasks. A typical structure may include:

\begin{itemize}[leftmargin=0.8cm, itemsep=2pt]
    \item \textbf{Chapter 1. Overview of the Internship Organization};
    \item \textbf{Chapter 2. Internship Activities and Assigned Tasks};
    \item \textbf{Chapter 3. Results, Experience, and Evaluation}.
\end{itemize}

These chapter titles are recommendations only and may be adapted to the
actual internship content.

The report shall reflect work actually performed, observed, analyzed, or
learned during the internship. Unrelated theoretical material should not
be included merely to increase report length.

\vspace{0.3cm}

% ================================================================
% X. TABLE OF CONTENTS
% ================================================================

\noindent
\textbf{X.\quad TABLE OF CONTENTS}

\vspace{0.15cm}

The Table of Contents shall be generated automatically using:

\begin{verbatim}
\tableofcontents
\end{verbatim}

Students shall not create or manually edit page numbers in the Table of
Contents.

The document shall be compiled sufficiently to ensure that all page
numbers and entries are correctly updated.

\vspace{0.3cm}

% ================================================================
% XI. LIST OF SYMBOLS AND ABBREVIATIONS
% ================================================================

\noindent
\textbf{XI.\quad LIST OF SYMBOLS AND ABBREVIATIONS}

\vspace{0.15cm}

Symbols and abbreviations that are used repeatedly in the report shall be
listed in the \textbf{List of Symbols and Abbreviations}.

The list may contain the following columns:

\begin{center}
\begin{tabular}{llll}
    No. & Symbol / Abbreviation & Unit & Definition
\end{tabular}
\end{center}

Only symbols and abbreviations actually used in the final report shall be
retained. Example entries contained in the template shall be removed if
they are not relevant to the student's report.

Each abbreviation should normally be defined when it first appears in
the main text.

\vspace{0.3cm}

% ================================================================
% XII. FIGURES, DIAGRAMS, AND PHOTOGRAPHS
% ================================================================

\noindent
\textbf{XII.\quad FIGURES, DIAGRAMS, AND PHOTOGRAPHS}

\vspace{0.15cm}

Figures shall be clear, legible, and directly relevant to the report
content.

Figures shall be inserted using the standard \verb|figure| environment.

Example:

\begin{verbatim}
\begin{figure}[H]
    \centering
    \includegraphics[width=0.8\textwidth]
                    {figures/company_structure.pdf}
    \caption{Organizational structure of the company}
    \label{fig:company-structure}
\end{figure}
\end{verbatim}

Figure captions shall be placed below the figure.

Each referenced figure shall have a unique \verb|\label| and shall be
referenced in the text using, for example:

\begin{verbatim}
Figure~\ref{fig:company-structure}
\end{verbatim}

Students shall not manually type figure numbers.

Figures reproduced or adapted from external sources shall include an
appropriate citation.

Photographs taken during the internship shall be relevant and sufficiently
clear. Students shall not include confidential, restricted, or proprietary
information without permission from the internship organization.

Vector graphics or high-resolution images are recommended whenever
possible.

\vspace{0.3cm}

% ================================================================
% XIII. TABLES
% ================================================================

\noindent
\textbf{XIII.\quad TABLES}

\vspace{0.15cm}

Tables shall be created using suitable \LaTeX{} environments such as
\verb|table|, \verb|tabular|, \verb|tabularx|, or \verb|longtable|.

Table captions shall be placed above the table.

Example:

\begin{verbatim}
\begin{table}[H]
    \centering
    \caption{Internship schedule}
    \label{tab:internship-schedule}

    \begin{tabular}{lll}
        \toprule
        Week & Task & Result \\
        \midrule
        1 & Orientation & Completed \\
        2 & System study & Completed \\
        \bottomrule
    \end{tabular}
\end{table}
\end{verbatim}

Tables shall be referenced using:

\begin{verbatim}
Table~\ref{tab:internship-schedule}
\end{verbatim}

Students shall not manually type table numbers.

Tables should be created directly in \LaTeX{} rather than inserted as
screenshots whenever practical.

\vspace{0.3cm}

% ================================================================
% XIV. EQUATIONS AND TECHNICAL QUANTITIES
% ================================================================

\noindent
\textbf{XIV.\quad EQUATIONS AND TECHNICAL QUANTITIES}

\vspace{0.15cm}

Mathematical expressions shall be typeset directly using \LaTeX{}.
Images or screenshots of equations shall not be used.

Short expressions may be written in inline mathematical mode, for example:

\begin{verbatim}
$R = 0.7\,\Omega$
\end{verbatim}

Important equations shall use a numbered equation environment.

Example:

\begin{verbatim}
\begin{equation}
    P = UI
    \label{eq:power}
\end{equation}
\end{verbatim}

Equations shall be referenced using:

\begin{verbatim}
Equation~\eqref{eq:power}
\end{verbatim}

Equation numbers shall not be manually entered.

SI units shall be used wherever applicable. The \verb|siunitx| package
should be used where appropriate.

Examples:

\begin{verbatim}
\SI{220}{\volt}
\SI{50}{\hertz}
\SI{2.5}{\ampere}
\end{verbatim}

Notation and units shall be consistent throughout the report.

\vspace{0.3cm}

% ================================================================
% XV. SOURCE CODE
% ================================================================

\noindent
\textbf{XV.\quad SOURCE CODE}

\vspace{0.15cm}

If the internship involves programming, PLCs, microcontrollers,
simulation, or software development, relevant source code may be included
in the report.

Source code shall be presented using the \verb|listings| environment or
another environment defined by the official template.

Example:

\begin{verbatim}
\begin{lstlisting}[language=st]
IF Start = TRUE THEN
    Motor := TRUE;
ELSE
    Motor := FALSE;
END_IF;
\end{lstlisting}
\end{verbatim}

Only code necessary for explaining the work should appear in the main
report. Long source-code listings should be moved to an appendix or
submitted as supplementary files if required.

Source code should not normally be inserted as screenshots.

\vspace{0.3cm}

% ================================================================
% XVI. ABSTRACT
% ================================================================

\noindent
\textbf{XVI.\quad ABSTRACT}

\vspace{0.15cm}

The Abstract shall briefly summarize:

\begin{itemize}[leftmargin=0.8cm, itemsep=2pt]
    \item the internship organization;
    \item the objectives of the internship;
    \item the main tasks performed;
    \item the main results, knowledge, and skills obtained.
\end{itemize}

The Abstract should be concise and should not reproduce the full content
of the report chapter by chapter.

\vspace{0.3cm}

% ================================================================
% XVII. REFERENCES AND CITATIONS
% ================================================================

\noindent
\textbf{XVII.\quad REFERENCES AND CITATIONS}

\vspace{0.15cm}

All books, journal papers, conference papers, standards, websites,
datasheets, manuals, technical documents, figures, tables, or other
external sources used in preparing the report shall be properly cited.

References shall be managed through the bibliography system provided with
the official template.

Bibliographic entries shall normally be stored in:

\begin{verbatim}
References.bib
\end{verbatim}

References shall be cited in the report using:

\begin{verbatim}
\cite{reference-key}
\end{verbatim}

Students shall not manually type reference numbers such as ``[1]'',
``[2]'', etc.

Titles of foreign-language publications shall remain in their original
language in the bibliography.

All borrowed ideas, data, figures, tables, equations, methods, or results
shall be properly acknowledged.

\vspace{0.3cm}

% ================================================================
% XVIII. CROSS-REFERENCING
% ================================================================

\noindent
\textbf{XVIII.\quad CROSS-REFERENCING}

\vspace{0.15cm}

Numbered elements shall be cross-referenced using the standard \LaTeX{}
mechanism.

Use the following commands where appropriate:

\begin{verbatim}
\label{...}
\ref{...}
\eqref{...}
\end{verbatim}

For example:

\begin{verbatim}
As shown in Figure~\ref{fig:system}, ...
\end{verbatim}

Students shall not manually type numbers such as ``Figure 2.3'' when
those numbers can be generated automatically.

\vspace{0.3cm}

% ================================================================
% XIX. PAGE NUMBERING
% ================================================================

\noindent
\textbf{XIX.\quad PAGE NUMBERING}

\vspace{0.15cm}

Cover pages and official evaluation forms at the beginning of the report
shall not display page numbers.

The main body of the report shall use Arabic page numbers beginning from
Chapter 1.

Page numbering shall be generated automatically by the official template.
Students shall not manually enter page numbers.

\vspace{0.3cm}

% ================================================================
% XX. PROJECT FILE ORGANIZATION
% ================================================================

\noindent
\textbf{XX.\quad PROJECT FILE ORGANIZATION}

\vspace{0.15cm}

Students should maintain a clear and organized \LaTeX{} project
structure. A recommended structure is:

\begin{verbatim}
InternshipReport/
|
|-- main.tex
|-- References.bib
|
|-- figures/
|   |-- company.pdf
|   |-- system.png
|   `-- ...
|
`-- chapters/
    |-- chapter1.tex
    |-- chapter2.tex
    `-- chapter3.tex
\end{verbatim}

Figures should normally be stored in the \verb|figures| directory
specified by the official template.

File and directory names should preferably contain Latin letters,
numbers, hyphens, or underscores.

Recommended examples include:

\begin{verbatim}
company_structure.pdf
motor_control.png
chapter_01.tex
\end{verbatim}

Spaces and unnecessary special characters in file names should be
avoided where possible.

\vspace{0.3cm}

% ================================================================
% XXI. DOCUMENT COMPILATION AND FINAL CHECK
% ================================================================

\noindent
\textbf{XXI.\quad DOCUMENT COMPILATION AND FINAL CHECK}

\vspace{0.15cm}

Before submission, the complete \LaTeX{} project shall compile
successfully and produce the final PDF report.

The final document shall not contain unresolved references such as:

\begin{verbatim}
??
[?]
\end{verbatim}

Before generating the final PDF, students shall check:

\begin{itemize}[leftmargin=0.8cm, itemsep=2pt]
    \item the Table of Contents;
    \item page numbering;
    \item chapter, section, and subsection numbering;
    \item figure numbering;
    \item table numbering;
    \item equation numbering;
    \item cross-references;
    \item citations;
    \item bibliography;
    \item font display;
    \item mathematical symbols;
    \item image quality;
    \item completeness of all required forms.
\end{itemize}

The document may be compiled using TeX Live, MiKTeX, or another
compatible \LaTeX{} environment. Regardless of the compilation
environment, the resulting PDF shall conform to the official Faculty
template.

\vspace{0.3cm}

% ================================================================
% XXII. CONTENT THAT STUDENTS MAY EDIT
% ================================================================

\noindent
\textbf{XXII.\quad CONTENT THAT STUDENTS MAY EDIT}

\vspace{0.15cm}

Students may edit or replace:

\begin{itemize}[leftmargin=0.8cm, itemsep=2pt]
    \item student name;
    \item class;
    \item internship organization;
    \item supervisor name;
    \item Abstract;
    \item List of Symbols and Abbreviations;
    \item chapter titles and chapter content;
    \item figures and photographs;
    \item tables;
    \item equations;
    \item citations and references.
\end{itemize}

Students shall not modify the main document formatting settings in the
preamble, including the document class, page geometry, line spacing,
heading formatting, table-of-contents formatting, font configuration,
page layout, or cover-page macros unless specifically instructed.

\vspace{0.3cm}

% ================================================================
% XXIII. FINAL SUBMISSION REQUIREMENTS
% ================================================================

\noindent
\textbf{XXIII.\quad FINAL SUBMISSION REQUIREMENTS}

\vspace{0.15cm}

Before submission, the final report shall:

\begin{itemize}[leftmargin=0.8cm, itemsep=2pt]
    \item compile without errors that affect the report content or presentation;
    \item contain all required forms and report sections;
    \item contain correct student, class, internship organization, and supervisor information;
    \item contain an updated Table of Contents, List of Figures, and List of Tables;
    \item contain no unused example content from the template;
    \item contain no placeholder text such as \textit{Title of Chapter 1},
    \textit{Section 1.1}, or other sample entries;
    \item contain all required citations and references;
    \item contain no unresolved cross-references.
\end{itemize}

If the Faculty requires submission of the \LaTeX{} source, the complete
project shall include all \verb|.tex| files, the bibliography database,
figures, and any other files required to reproduce the submitted PDF.

\vspace{0.4cm}

\noindent
\textbf{Important:}
Before submitting the final internship report, students shall remove or
comment out the following command in the main document:

\begin{verbatim}
\input{Regulations.tex}
\end{verbatim}

For example:

\begin{verbatim}
% \input{Regulations.tex}
\end{verbatim}

This removes the guidance pages from the final submitted report while
preserving the regulations file in the \LaTeX{} project for reference.

} % End 13pt content

\clearpage

\end{verbatim}

For example:

\begin{verbatim}
% 
% =================================================================
% REGULATIONS FOR PREPARING THE GRADUATION INTERNSHIP REPORT
% USING LATEX
% =================================================================
%
% This file is intended to be included from the main document with:
%
%     \input{Regulations.tex}
%
% IMPORTANT:
% This section is provided for guidance only. Students shall remove
% or comment out the \input{Regulations.tex} command before submitting
% the final internship report.
%
% =================================================================

\clearpage
\thispagestyle{empty}

\begin{center}
    {\fontsize{16pt}{20pt}\selectfont\bfseries
    SPECIFIC REGULATIONS FOR PREPARING\\
    GRADUATION INTERNSHIP REPORTS USING \LaTeX}
\end{center}

\vspace{0.5cm}

{\fontsize{13pt}{17pt}\selectfont
\setstretch{1.3}
\justifying

\noindent
\textbf{Note:}
This section is provided for guidance only and is included in the main
document using the command \verb|\input{Regulations.tex}|.
Students shall remove or comment out this command before submitting the
final internship report.

\vspace{0.4cm}

% ================================================================
% I. GENERAL REQUIREMENTS
% ================================================================

\noindent
\textbf{I.\quad GENERAL REQUIREMENTS}

\vspace{0.15cm}

The Graduation Internship Report shall be prepared using the official
\LaTeX{} template provided by the Faculty. Students shall preserve the
predefined page layout, typography, line spacing, heading styles,
numbering system, figure and table formatting, and bibliography settings
unless otherwise instructed by the Faculty or the supervisor.

The report shall be written in English.

Students are responsible for the academic and technical content of the
report. The official \LaTeX{} template is responsible for standardized
document formatting.

\vspace{0.3cm}

% ================================================================
% II. REPORT STRUCTURE
% ================================================================

\noindent
\textbf{II.\quad REPORT STRUCTURE}

\vspace{0.15cm}

The report shall normally be arranged in the following order:

\begin{enumerate}[label=\arabic*., leftmargin=1.2cm, itemsep=2pt]
    \item Outer Cover;
    \item Comments of the Internship Host Organization;
    \item Supervisor's Evaluation and Graduation Internship Defense Results;
    \item Table of Contents;
    \item List of Symbols and Abbreviations, if applicable;
    \item List of Figures;
    \item List of Tables;
    \item Abstract;
    \item Main Chapters of the Internship Report;
    \item References;
    \item Appendices, if applicable.
\end{enumerate}

The preliminary pages shall not be numbered as chapters. The main report
content shall be organized into consecutively numbered chapters beginning
with Chapter 1.

\vspace{0.3cm}

% ================================================================
% III. PAGE SIZE AND MARGINS
% ================================================================

\noindent
\textbf{III.\quad PAGE SIZE AND MARGINS}

\vspace{0.15cm}

The report shall be formatted on A4 paper with the following margins:

\begin{center}
\begin{tabular}{ll}
    Left margin:   & 3.5 cm \\
    Right margin:  & 2.0 cm \\
    Top margin:    & 2.5 cm \\
    Bottom margin: & 2.5 cm
\end{tabular}
\end{center}

These settings are predefined in the official template and shall not be
modified by students.

\vspace{0.3cm}

% ================================================================
% IV. FONT, FONT SIZE, AND LINE SPACING
% ================================================================

\noindent
\textbf{IV.\quad FONT, FONT SIZE, AND LINE SPACING}

\vspace{0.15cm}

The main text shall use a 13 pt font size with 1.3-line spacing.
A Times-family typeface shall be used in accordance with the official
template.

The text shall be fully justified.

Students shall not use repeated spaces, repeated manual line breaks, or
multiple blank lines to position text. Paragraph and heading spacing shall
be controlled by the template.

\vspace{0.3cm}

% ================================================================
% V. OUTER COVER
% ================================================================

\noindent
\textbf{V.\quad OUTER COVER}

\vspace{0.15cm}

The outer cover shall follow the official format provided in the template
and shall contain the following information:

\begin{itemize}[leftmargin=0.8cm, itemsep=2pt]
    \item Thai Nguyen University of Technology;
    \item Faculty of International Training;
    \item Engineering and Technology Department;
    \item University logo and Faculty logo;
    \item the title \textbf{INTERNSHIP REPORT};
    \item Internship Organization;
    \item Author;
    \item Class;
    \item Supervisor;
    \item place and year of completion.
\end{itemize}

Students shall only replace the required information fields and shall not
modify the cover-page layout unless instructed otherwise.

\vspace{0.3cm}

% ================================================================
% VI. COMMENTS OF THE INTERNSHIP HOST ORGANIZATION
% ================================================================

\noindent
\textbf{VI.\quad COMMENTS OF THE INTERNSHIP HOST ORGANIZATION}

\vspace{0.15cm}

The page entitled \textbf{COMMENTS OF THE INTERNSHIP HOST ORGANIZATION}
is reserved for the internship host organization.

It includes:

\begin{itemize}[leftmargin=0.8cm, itemsep=2pt]
    \item the student's full name;
    \item class;
    \item comments from the host organization;
    \item score;
    \item date;
    \item signature and official seal.
\end{itemize}

Students shall not alter or complete official evaluation fields unless
specifically authorized.

\vspace{0.3cm}

% ================================================================
% VII. SUPERVISOR'S EVALUATION AND DEFENSE RESULTS
% ================================================================

\noindent
\textbf{VII.\quad SUPERVISOR'S EVALUATION AND DEFENSE RESULTS}

\vspace{0.15cm}

The section entitled
\textbf{SUPERVISOR'S EVALUATION AND GRADUATION INTERNSHIP DEFENSE RESULTS}
is reserved for the supervisor and defense committee.

The evaluation may include:

\begin{itemize}[leftmargin=0.8cm, itemsep=2pt]
    \item Internship Progress and Student's Attitude;
    \item Report Content;
    \item Report Presentation;
    \item Other Comments;
    \item Supervisor's Opinion;
    \item Defense Result;
    \item Score.
\end{itemize}

Students shall not enter evaluation scores or defense results before they
are officially determined.

\vspace{0.3cm}

% ================================================================
% VIII. CHAPTERS, SECTIONS, AND SUBSECTIONS
% ================================================================

\noindent
\textbf{VIII.\quad CHAPTERS, SECTIONS, AND SUBSECTIONS}

\vspace{0.15cm}

The report structure shall be generated using the standard \LaTeX{}
commands:

\begin{verbatim}
\chapter{...}
\section{...}
\subsection{...}
\end{verbatim}

Chapter, section, and subsection numbers shall be generated automatically
by \LaTeX{}.

Students shall not manually type chapter or section numbers and shall not
reproduce heading styles using commands such as \verb|\textbf|,
\verb|\fontsize|, or \verb|\centering| instead of the appropriate
structural command.

Excessive subdivision should be avoided. In most cases, the hierarchy
Chapter--Section--Subsection is sufficient.

\vspace{0.3cm}

% ================================================================
% IX. ORGANIZATION OF THE MAIN CONTENT
% ================================================================

\noindent
\textbf{IX.\quad ORGANIZATION OF THE MAIN CONTENT}

\vspace{0.15cm}

The specific content of the report shall depend on the internship
organization and the assigned tasks. A typical structure may include:

\begin{itemize}[leftmargin=0.8cm, itemsep=2pt]
    \item \textbf{Chapter 1. Overview of the Internship Organization};
    \item \textbf{Chapter 2. Internship Activities and Assigned Tasks};
    \item \textbf{Chapter 3. Results, Experience, and Evaluation}.
\end{itemize}

These chapter titles are recommendations only and may be adapted to the
actual internship content.

The report shall reflect work actually performed, observed, analyzed, or
learned during the internship. Unrelated theoretical material should not
be included merely to increase report length.

\vspace{0.3cm}

% ================================================================
% X. TABLE OF CONTENTS
% ================================================================

\noindent
\textbf{X.\quad TABLE OF CONTENTS}

\vspace{0.15cm}

The Table of Contents shall be generated automatically using:

\begin{verbatim}
\tableofcontents
\end{verbatim}

Students shall not create or manually edit page numbers in the Table of
Contents.

The document shall be compiled sufficiently to ensure that all page
numbers and entries are correctly updated.

\vspace{0.3cm}

% ================================================================
% XI. LIST OF SYMBOLS AND ABBREVIATIONS
% ================================================================

\noindent
\textbf{XI.\quad LIST OF SYMBOLS AND ABBREVIATIONS}

\vspace{0.15cm}

Symbols and abbreviations that are used repeatedly in the report shall be
listed in the \textbf{List of Symbols and Abbreviations}.

The list may contain the following columns:

\begin{center}
\begin{tabular}{llll}
    No. & Symbol / Abbreviation & Unit & Definition
\end{tabular}
\end{center}

Only symbols and abbreviations actually used in the final report shall be
retained. Example entries contained in the template shall be removed if
they are not relevant to the student's report.

Each abbreviation should normally be defined when it first appears in
the main text.

\vspace{0.3cm}

% ================================================================
% XII. FIGURES, DIAGRAMS, AND PHOTOGRAPHS
% ================================================================

\noindent
\textbf{XII.\quad FIGURES, DIAGRAMS, AND PHOTOGRAPHS}

\vspace{0.15cm}

Figures shall be clear, legible, and directly relevant to the report
content.

Figures shall be inserted using the standard \verb|figure| environment.

Example:

\begin{verbatim}
\begin{figure}[H]
    \centering
    \includegraphics[width=0.8\textwidth]
                    {figures/company_structure.pdf}
    \caption{Organizational structure of the company}
    \label{fig:company-structure}
\end{figure}
\end{verbatim}

Figure captions shall be placed below the figure.

Each referenced figure shall have a unique \verb|\label| and shall be
referenced in the text using, for example:

\begin{verbatim}
Figure~\ref{fig:company-structure}
\end{verbatim}

Students shall not manually type figure numbers.

Figures reproduced or adapted from external sources shall include an
appropriate citation.

Photographs taken during the internship shall be relevant and sufficiently
clear. Students shall not include confidential, restricted, or proprietary
information without permission from the internship organization.

Vector graphics or high-resolution images are recommended whenever
possible.

\vspace{0.3cm}

% ================================================================
% XIII. TABLES
% ================================================================

\noindent
\textbf{XIII.\quad TABLES}

\vspace{0.15cm}

Tables shall be created using suitable \LaTeX{} environments such as
\verb|table|, \verb|tabular|, \verb|tabularx|, or \verb|longtable|.

Table captions shall be placed above the table.

Example:

\begin{verbatim}
\begin{table}[H]
    \centering
    \caption{Internship schedule}
    \label{tab:internship-schedule}

    \begin{tabular}{lll}
        \toprule
        Week & Task & Result \\
        \midrule
        1 & Orientation & Completed \\
        2 & System study & Completed \\
        \bottomrule
    \end{tabular}
\end{table}
\end{verbatim}

Tables shall be referenced using:

\begin{verbatim}
Table~\ref{tab:internship-schedule}
\end{verbatim}

Students shall not manually type table numbers.

Tables should be created directly in \LaTeX{} rather than inserted as
screenshots whenever practical.

\vspace{0.3cm}

% ================================================================
% XIV. EQUATIONS AND TECHNICAL QUANTITIES
% ================================================================

\noindent
\textbf{XIV.\quad EQUATIONS AND TECHNICAL QUANTITIES}

\vspace{0.15cm}

Mathematical expressions shall be typeset directly using \LaTeX{}.
Images or screenshots of equations shall not be used.

Short expressions may be written in inline mathematical mode, for example:

\begin{verbatim}
$R = 0.7\,\Omega$
\end{verbatim}

Important equations shall use a numbered equation environment.

Example:

\begin{verbatim}
\begin{equation}
    P = UI
    \label{eq:power}
\end{equation}
\end{verbatim}

Equations shall be referenced using:

\begin{verbatim}
Equation~\eqref{eq:power}
\end{verbatim}

Equation numbers shall not be manually entered.

SI units shall be used wherever applicable. The \verb|siunitx| package
should be used where appropriate.

Examples:

\begin{verbatim}
\SI{220}{\volt}
\SI{50}{\hertz}
\SI{2.5}{\ampere}
\end{verbatim}

Notation and units shall be consistent throughout the report.

\vspace{0.3cm}

% ================================================================
% XV. SOURCE CODE
% ================================================================

\noindent
\textbf{XV.\quad SOURCE CODE}

\vspace{0.15cm}

If the internship involves programming, PLCs, microcontrollers,
simulation, or software development, relevant source code may be included
in the report.

Source code shall be presented using the \verb|listings| environment or
another environment defined by the official template.

Example:

\begin{verbatim}
\begin{lstlisting}[language=st]
IF Start = TRUE THEN
    Motor := TRUE;
ELSE
    Motor := FALSE;
END_IF;
\end{lstlisting}
\end{verbatim}

Only code necessary for explaining the work should appear in the main
report. Long source-code listings should be moved to an appendix or
submitted as supplementary files if required.

Source code should not normally be inserted as screenshots.

\vspace{0.3cm}

% ================================================================
% XVI. ABSTRACT
% ================================================================

\noindent
\textbf{XVI.\quad ABSTRACT}

\vspace{0.15cm}

The Abstract shall briefly summarize:

\begin{itemize}[leftmargin=0.8cm, itemsep=2pt]
    \item the internship organization;
    \item the objectives of the internship;
    \item the main tasks performed;
    \item the main results, knowledge, and skills obtained.
\end{itemize}

The Abstract should be concise and should not reproduce the full content
of the report chapter by chapter.

\vspace{0.3cm}

% ================================================================
% XVII. REFERENCES AND CITATIONS
% ================================================================

\noindent
\textbf{XVII.\quad REFERENCES AND CITATIONS}

\vspace{0.15cm}

All books, journal papers, conference papers, standards, websites,
datasheets, manuals, technical documents, figures, tables, or other
external sources used in preparing the report shall be properly cited.

References shall be managed through the bibliography system provided with
the official template.

Bibliographic entries shall normally be stored in:

\begin{verbatim}
References.bib
\end{verbatim}

References shall be cited in the report using:

\begin{verbatim}
\cite{reference-key}
\end{verbatim}

Students shall not manually type reference numbers such as ``[1]'',
``[2]'', etc.

Titles of foreign-language publications shall remain in their original
language in the bibliography.

All borrowed ideas, data, figures, tables, equations, methods, or results
shall be properly acknowledged.

\vspace{0.3cm}

% ================================================================
% XVIII. CROSS-REFERENCING
% ================================================================

\noindent
\textbf{XVIII.\quad CROSS-REFERENCING}

\vspace{0.15cm}

Numbered elements shall be cross-referenced using the standard \LaTeX{}
mechanism.

Use the following commands where appropriate:

\begin{verbatim}
\label{...}
\ref{...}
\eqref{...}
\end{verbatim}

For example:

\begin{verbatim}
As shown in Figure~\ref{fig:system}, ...
\end{verbatim}

Students shall not manually type numbers such as ``Figure 2.3'' when
those numbers can be generated automatically.

\vspace{0.3cm}

% ================================================================
% XIX. PAGE NUMBERING
% ================================================================

\noindent
\textbf{XIX.\quad PAGE NUMBERING}

\vspace{0.15cm}

Cover pages and official evaluation forms at the beginning of the report
shall not display page numbers.

The main body of the report shall use Arabic page numbers beginning from
Chapter 1.

Page numbering shall be generated automatically by the official template.
Students shall not manually enter page numbers.

\vspace{0.3cm}

% ================================================================
% XX. PROJECT FILE ORGANIZATION
% ================================================================

\noindent
\textbf{XX.\quad PROJECT FILE ORGANIZATION}

\vspace{0.15cm}

Students should maintain a clear and organized \LaTeX{} project
structure. A recommended structure is:

\begin{verbatim}
InternshipReport/
|
|-- main.tex
|-- References.bib
|
|-- figures/
|   |-- company.pdf
|   |-- system.png
|   `-- ...
|
`-- chapters/
    |-- chapter1.tex
    |-- chapter2.tex
    `-- chapter3.tex
\end{verbatim}

Figures should normally be stored in the \verb|figures| directory
specified by the official template.

File and directory names should preferably contain Latin letters,
numbers, hyphens, or underscores.

Recommended examples include:

\begin{verbatim}
company_structure.pdf
motor_control.png
chapter_01.tex
\end{verbatim}

Spaces and unnecessary special characters in file names should be
avoided where possible.

\vspace{0.3cm}

% ================================================================
% XXI. DOCUMENT COMPILATION AND FINAL CHECK
% ================================================================

\noindent
\textbf{XXI.\quad DOCUMENT COMPILATION AND FINAL CHECK}

\vspace{0.15cm}

Before submission, the complete \LaTeX{} project shall compile
successfully and produce the final PDF report.

The final document shall not contain unresolved references such as:

\begin{verbatim}
??
[?]
\end{verbatim}

Before generating the final PDF, students shall check:

\begin{itemize}[leftmargin=0.8cm, itemsep=2pt]
    \item the Table of Contents;
    \item page numbering;
    \item chapter, section, and subsection numbering;
    \item figure numbering;
    \item table numbering;
    \item equation numbering;
    \item cross-references;
    \item citations;
    \item bibliography;
    \item font display;
    \item mathematical symbols;
    \item image quality;
    \item completeness of all required forms.
\end{itemize}

The document may be compiled using TeX Live, MiKTeX, or another
compatible \LaTeX{} environment. Regardless of the compilation
environment, the resulting PDF shall conform to the official Faculty
template.

\vspace{0.3cm}

% ================================================================
% XXII. CONTENT THAT STUDENTS MAY EDIT
% ================================================================

\noindent
\textbf{XXII.\quad CONTENT THAT STUDENTS MAY EDIT}

\vspace{0.15cm}

Students may edit or replace:

\begin{itemize}[leftmargin=0.8cm, itemsep=2pt]
    \item student name;
    \item class;
    \item internship organization;
    \item supervisor name;
    \item Abstract;
    \item List of Symbols and Abbreviations;
    \item chapter titles and chapter content;
    \item figures and photographs;
    \item tables;
    \item equations;
    \item citations and references.
\end{itemize}

Students shall not modify the main document formatting settings in the
preamble, including the document class, page geometry, line spacing,
heading formatting, table-of-contents formatting, font configuration,
page layout, or cover-page macros unless specifically instructed.

\vspace{0.3cm}

% ================================================================
% XXIII. FINAL SUBMISSION REQUIREMENTS
% ================================================================

\noindent
\textbf{XXIII.\quad FINAL SUBMISSION REQUIREMENTS}

\vspace{0.15cm}

Before submission, the final report shall:

\begin{itemize}[leftmargin=0.8cm, itemsep=2pt]
    \item compile without errors that affect the report content or presentation;
    \item contain all required forms and report sections;
    \item contain correct student, class, internship organization, and supervisor information;
    \item contain an updated Table of Contents, List of Figures, and List of Tables;
    \item contain no unused example content from the template;
    \item contain no placeholder text such as \textit{Title of Chapter 1},
    \textit{Section 1.1}, or other sample entries;
    \item contain all required citations and references;
    \item contain no unresolved cross-references.
\end{itemize}

If the Faculty requires submission of the \LaTeX{} source, the complete
project shall include all \verb|.tex| files, the bibliography database,
figures, and any other files required to reproduce the submitted PDF.

\vspace{0.4cm}

\noindent
\textbf{Important:}
Before submitting the final internship report, students shall remove or
comment out the following command in the main document:

\begin{verbatim}
\input{Regulations.tex}
\end{verbatim}

For example:

\begin{verbatim}
% \input{Regulations.tex}
\end{verbatim}

This removes the guidance pages from the final submitted report while
preserving the regulations file in the \LaTeX{} project for reference.

} % End 13pt content

\clearpage

\end{verbatim}

This removes the guidance pages from the final submitted report while
preserving the regulations file in the \LaTeX{} project for reference.

} % End 13pt content

\clearpage
 command before submitting
% the final internship report.
%
% =================================================================

\clearpage
\thispagestyle{empty}

\begin{center}
    {\fontsize{16pt}{20pt}\selectfont\bfseries
    SPECIFIC REGULATIONS FOR PREPARING\\
    GRADUATION INTERNSHIP REPORTS USING \LaTeX}
\end{center}

\vspace{0.5cm}

{\fontsize{13pt}{17pt}\selectfont
\setstretch{1.3}
\justifying

\noindent
\textbf{Note:}
This section is provided for guidance only and is included in the main
document using the command \verb|
% =================================================================
% REGULATIONS FOR PREPARING THE GRADUATION INTERNSHIP REPORT
% USING LATEX
% =================================================================
%
% This file is intended to be included from the main document with:
%
%     
% =================================================================
% REGULATIONS FOR PREPARING THE GRADUATION INTERNSHIP REPORT
% USING LATEX
% =================================================================
%
% This file is intended to be included from the main document with:
%
%     \input{Regulations.tex}
%
% IMPORTANT:
% This section is provided for guidance only. Students shall remove
% or comment out the \input{Regulations.tex} command before submitting
% the final internship report.
%
% =================================================================

\clearpage
\thispagestyle{empty}

\begin{center}
    {\fontsize{16pt}{20pt}\selectfont\bfseries
    SPECIFIC REGULATIONS FOR PREPARING\\
    GRADUATION INTERNSHIP REPORTS USING \LaTeX}
\end{center}

\vspace{0.5cm}

{\fontsize{13pt}{17pt}\selectfont
\setstretch{1.3}
\justifying

\noindent
\textbf{Note:}
This section is provided for guidance only and is included in the main
document using the command \verb|\input{Regulations.tex}|.
Students shall remove or comment out this command before submitting the
final internship report.

\vspace{0.4cm}

% ================================================================
% I. GENERAL REQUIREMENTS
% ================================================================

\noindent
\textbf{I.\quad GENERAL REQUIREMENTS}

\vspace{0.15cm}

The Graduation Internship Report shall be prepared using the official
\LaTeX{} template provided by the Faculty. Students shall preserve the
predefined page layout, typography, line spacing, heading styles,
numbering system, figure and table formatting, and bibliography settings
unless otherwise instructed by the Faculty or the supervisor.

The report shall be written in English.

Students are responsible for the academic and technical content of the
report. The official \LaTeX{} template is responsible for standardized
document formatting.

\vspace{0.3cm}

% ================================================================
% II. REPORT STRUCTURE
% ================================================================

\noindent
\textbf{II.\quad REPORT STRUCTURE}

\vspace{0.15cm}

The report shall normally be arranged in the following order:

\begin{enumerate}[label=\arabic*., leftmargin=1.2cm, itemsep=2pt]
    \item Outer Cover;
    \item Comments of the Internship Host Organization;
    \item Supervisor's Evaluation and Graduation Internship Defense Results;
    \item Table of Contents;
    \item List of Symbols and Abbreviations, if applicable;
    \item List of Figures;
    \item List of Tables;
    \item Abstract;
    \item Main Chapters of the Internship Report;
    \item References;
    \item Appendices, if applicable.
\end{enumerate}

The preliminary pages shall not be numbered as chapters. The main report
content shall be organized into consecutively numbered chapters beginning
with Chapter 1.

\vspace{0.3cm}

% ================================================================
% III. PAGE SIZE AND MARGINS
% ================================================================

\noindent
\textbf{III.\quad PAGE SIZE AND MARGINS}

\vspace{0.15cm}

The report shall be formatted on A4 paper with the following margins:

\begin{center}
\begin{tabular}{ll}
    Left margin:   & 3.5 cm \\
    Right margin:  & 2.0 cm \\
    Top margin:    & 2.5 cm \\
    Bottom margin: & 2.5 cm
\end{tabular}
\end{center}

These settings are predefined in the official template and shall not be
modified by students.

\vspace{0.3cm}

% ================================================================
% IV. FONT, FONT SIZE, AND LINE SPACING
% ================================================================

\noindent
\textbf{IV.\quad FONT, FONT SIZE, AND LINE SPACING}

\vspace{0.15cm}

The main text shall use a 13 pt font size with 1.3-line spacing.
A Times-family typeface shall be used in accordance with the official
template.

The text shall be fully justified.

Students shall not use repeated spaces, repeated manual line breaks, or
multiple blank lines to position text. Paragraph and heading spacing shall
be controlled by the template.

\vspace{0.3cm}

% ================================================================
% V. OUTER COVER
% ================================================================

\noindent
\textbf{V.\quad OUTER COVER}

\vspace{0.15cm}

The outer cover shall follow the official format provided in the template
and shall contain the following information:

\begin{itemize}[leftmargin=0.8cm, itemsep=2pt]
    \item Thai Nguyen University of Technology;
    \item Faculty of International Training;
    \item Engineering and Technology Department;
    \item University logo and Faculty logo;
    \item the title \textbf{INTERNSHIP REPORT};
    \item Internship Organization;
    \item Author;
    \item Class;
    \item Supervisor;
    \item place and year of completion.
\end{itemize}

Students shall only replace the required information fields and shall not
modify the cover-page layout unless instructed otherwise.

\vspace{0.3cm}

% ================================================================
% VI. COMMENTS OF THE INTERNSHIP HOST ORGANIZATION
% ================================================================

\noindent
\textbf{VI.\quad COMMENTS OF THE INTERNSHIP HOST ORGANIZATION}

\vspace{0.15cm}

The page entitled \textbf{COMMENTS OF THE INTERNSHIP HOST ORGANIZATION}
is reserved for the internship host organization.

It includes:

\begin{itemize}[leftmargin=0.8cm, itemsep=2pt]
    \item the student's full name;
    \item class;
    \item comments from the host organization;
    \item score;
    \item date;
    \item signature and official seal.
\end{itemize}

Students shall not alter or complete official evaluation fields unless
specifically authorized.

\vspace{0.3cm}

% ================================================================
% VII. SUPERVISOR'S EVALUATION AND DEFENSE RESULTS
% ================================================================

\noindent
\textbf{VII.\quad SUPERVISOR'S EVALUATION AND DEFENSE RESULTS}

\vspace{0.15cm}

The section entitled
\textbf{SUPERVISOR'S EVALUATION AND GRADUATION INTERNSHIP DEFENSE RESULTS}
is reserved for the supervisor and defense committee.

The evaluation may include:

\begin{itemize}[leftmargin=0.8cm, itemsep=2pt]
    \item Internship Progress and Student's Attitude;
    \item Report Content;
    \item Report Presentation;
    \item Other Comments;
    \item Supervisor's Opinion;
    \item Defense Result;
    \item Score.
\end{itemize}

Students shall not enter evaluation scores or defense results before they
are officially determined.

\vspace{0.3cm}

% ================================================================
% VIII. CHAPTERS, SECTIONS, AND SUBSECTIONS
% ================================================================

\noindent
\textbf{VIII.\quad CHAPTERS, SECTIONS, AND SUBSECTIONS}

\vspace{0.15cm}

The report structure shall be generated using the standard \LaTeX{}
commands:

\begin{verbatim}
\chapter{...}
\section{...}
\subsection{...}
\end{verbatim}

Chapter, section, and subsection numbers shall be generated automatically
by \LaTeX{}.

Students shall not manually type chapter or section numbers and shall not
reproduce heading styles using commands such as \verb|\textbf|,
\verb|\fontsize|, or \verb|\centering| instead of the appropriate
structural command.

Excessive subdivision should be avoided. In most cases, the hierarchy
Chapter--Section--Subsection is sufficient.

\vspace{0.3cm}

% ================================================================
% IX. ORGANIZATION OF THE MAIN CONTENT
% ================================================================

\noindent
\textbf{IX.\quad ORGANIZATION OF THE MAIN CONTENT}

\vspace{0.15cm}

The specific content of the report shall depend on the internship
organization and the assigned tasks. A typical structure may include:

\begin{itemize}[leftmargin=0.8cm, itemsep=2pt]
    \item \textbf{Chapter 1. Overview of the Internship Organization};
    \item \textbf{Chapter 2. Internship Activities and Assigned Tasks};
    \item \textbf{Chapter 3. Results, Experience, and Evaluation}.
\end{itemize}

These chapter titles are recommendations only and may be adapted to the
actual internship content.

The report shall reflect work actually performed, observed, analyzed, or
learned during the internship. Unrelated theoretical material should not
be included merely to increase report length.

\vspace{0.3cm}

% ================================================================
% X. TABLE OF CONTENTS
% ================================================================

\noindent
\textbf{X.\quad TABLE OF CONTENTS}

\vspace{0.15cm}

The Table of Contents shall be generated automatically using:

\begin{verbatim}
\tableofcontents
\end{verbatim}

Students shall not create or manually edit page numbers in the Table of
Contents.

The document shall be compiled sufficiently to ensure that all page
numbers and entries are correctly updated.

\vspace{0.3cm}

% ================================================================
% XI. LIST OF SYMBOLS AND ABBREVIATIONS
% ================================================================

\noindent
\textbf{XI.\quad LIST OF SYMBOLS AND ABBREVIATIONS}

\vspace{0.15cm}

Symbols and abbreviations that are used repeatedly in the report shall be
listed in the \textbf{List of Symbols and Abbreviations}.

The list may contain the following columns:

\begin{center}
\begin{tabular}{llll}
    No. & Symbol / Abbreviation & Unit & Definition
\end{tabular}
\end{center}

Only symbols and abbreviations actually used in the final report shall be
retained. Example entries contained in the template shall be removed if
they are not relevant to the student's report.

Each abbreviation should normally be defined when it first appears in
the main text.

\vspace{0.3cm}

% ================================================================
% XII. FIGURES, DIAGRAMS, AND PHOTOGRAPHS
% ================================================================

\noindent
\textbf{XII.\quad FIGURES, DIAGRAMS, AND PHOTOGRAPHS}

\vspace{0.15cm}

Figures shall be clear, legible, and directly relevant to the report
content.

Figures shall be inserted using the standard \verb|figure| environment.

Example:

\begin{verbatim}
\begin{figure}[H]
    \centering
    \includegraphics[width=0.8\textwidth]
                    {figures/company_structure.pdf}
    \caption{Organizational structure of the company}
    \label{fig:company-structure}
\end{figure}
\end{verbatim}

Figure captions shall be placed below the figure.

Each referenced figure shall have a unique \verb|\label| and shall be
referenced in the text using, for example:

\begin{verbatim}
Figure~\ref{fig:company-structure}
\end{verbatim}

Students shall not manually type figure numbers.

Figures reproduced or adapted from external sources shall include an
appropriate citation.

Photographs taken during the internship shall be relevant and sufficiently
clear. Students shall not include confidential, restricted, or proprietary
information without permission from the internship organization.

Vector graphics or high-resolution images are recommended whenever
possible.

\vspace{0.3cm}

% ================================================================
% XIII. TABLES
% ================================================================

\noindent
\textbf{XIII.\quad TABLES}

\vspace{0.15cm}

Tables shall be created using suitable \LaTeX{} environments such as
\verb|table|, \verb|tabular|, \verb|tabularx|, or \verb|longtable|.

Table captions shall be placed above the table.

Example:

\begin{verbatim}
\begin{table}[H]
    \centering
    \caption{Internship schedule}
    \label{tab:internship-schedule}

    \begin{tabular}{lll}
        \toprule
        Week & Task & Result \\
        \midrule
        1 & Orientation & Completed \\
        2 & System study & Completed \\
        \bottomrule
    \end{tabular}
\end{table}
\end{verbatim}

Tables shall be referenced using:

\begin{verbatim}
Table~\ref{tab:internship-schedule}
\end{verbatim}

Students shall not manually type table numbers.

Tables should be created directly in \LaTeX{} rather than inserted as
screenshots whenever practical.

\vspace{0.3cm}

% ================================================================
% XIV. EQUATIONS AND TECHNICAL QUANTITIES
% ================================================================

\noindent
\textbf{XIV.\quad EQUATIONS AND TECHNICAL QUANTITIES}

\vspace{0.15cm}

Mathematical expressions shall be typeset directly using \LaTeX{}.
Images or screenshots of equations shall not be used.

Short expressions may be written in inline mathematical mode, for example:

\begin{verbatim}
$R = 0.7\,\Omega$
\end{verbatim}

Important equations shall use a numbered equation environment.

Example:

\begin{verbatim}
\begin{equation}
    P = UI
    \label{eq:power}
\end{equation}
\end{verbatim}

Equations shall be referenced using:

\begin{verbatim}
Equation~\eqref{eq:power}
\end{verbatim}

Equation numbers shall not be manually entered.

SI units shall be used wherever applicable. The \verb|siunitx| package
should be used where appropriate.

Examples:

\begin{verbatim}
\SI{220}{\volt}
\SI{50}{\hertz}
\SI{2.5}{\ampere}
\end{verbatim}

Notation and units shall be consistent throughout the report.

\vspace{0.3cm}

% ================================================================
% XV. SOURCE CODE
% ================================================================

\noindent
\textbf{XV.\quad SOURCE CODE}

\vspace{0.15cm}

If the internship involves programming, PLCs, microcontrollers,
simulation, or software development, relevant source code may be included
in the report.

Source code shall be presented using the \verb|listings| environment or
another environment defined by the official template.

Example:

\begin{verbatim}
\begin{lstlisting}[language=st]
IF Start = TRUE THEN
    Motor := TRUE;
ELSE
    Motor := FALSE;
END_IF;
\end{lstlisting}
\end{verbatim}

Only code necessary for explaining the work should appear in the main
report. Long source-code listings should be moved to an appendix or
submitted as supplementary files if required.

Source code should not normally be inserted as screenshots.

\vspace{0.3cm}

% ================================================================
% XVI. ABSTRACT
% ================================================================

\noindent
\textbf{XVI.\quad ABSTRACT}

\vspace{0.15cm}

The Abstract shall briefly summarize:

\begin{itemize}[leftmargin=0.8cm, itemsep=2pt]
    \item the internship organization;
    \item the objectives of the internship;
    \item the main tasks performed;
    \item the main results, knowledge, and skills obtained.
\end{itemize}

The Abstract should be concise and should not reproduce the full content
of the report chapter by chapter.

\vspace{0.3cm}

% ================================================================
% XVII. REFERENCES AND CITATIONS
% ================================================================

\noindent
\textbf{XVII.\quad REFERENCES AND CITATIONS}

\vspace{0.15cm}

All books, journal papers, conference papers, standards, websites,
datasheets, manuals, technical documents, figures, tables, or other
external sources used in preparing the report shall be properly cited.

References shall be managed through the bibliography system provided with
the official template.

Bibliographic entries shall normally be stored in:

\begin{verbatim}
References.bib
\end{verbatim}

References shall be cited in the report using:

\begin{verbatim}
\cite{reference-key}
\end{verbatim}

Students shall not manually type reference numbers such as ``[1]'',
``[2]'', etc.

Titles of foreign-language publications shall remain in their original
language in the bibliography.

All borrowed ideas, data, figures, tables, equations, methods, or results
shall be properly acknowledged.

\vspace{0.3cm}

% ================================================================
% XVIII. CROSS-REFERENCING
% ================================================================

\noindent
\textbf{XVIII.\quad CROSS-REFERENCING}

\vspace{0.15cm}

Numbered elements shall be cross-referenced using the standard \LaTeX{}
mechanism.

Use the following commands where appropriate:

\begin{verbatim}
\label{...}
\ref{...}
\eqref{...}
\end{verbatim}

For example:

\begin{verbatim}
As shown in Figure~\ref{fig:system}, ...
\end{verbatim}

Students shall not manually type numbers such as ``Figure 2.3'' when
those numbers can be generated automatically.

\vspace{0.3cm}

% ================================================================
% XIX. PAGE NUMBERING
% ================================================================

\noindent
\textbf{XIX.\quad PAGE NUMBERING}

\vspace{0.15cm}

Cover pages and official evaluation forms at the beginning of the report
shall not display page numbers.

The main body of the report shall use Arabic page numbers beginning from
Chapter 1.

Page numbering shall be generated automatically by the official template.
Students shall not manually enter page numbers.

\vspace{0.3cm}

% ================================================================
% XX. PROJECT FILE ORGANIZATION
% ================================================================

\noindent
\textbf{XX.\quad PROJECT FILE ORGANIZATION}

\vspace{0.15cm}

Students should maintain a clear and organized \LaTeX{} project
structure. A recommended structure is:

\begin{verbatim}
InternshipReport/
|
|-- main.tex
|-- References.bib
|
|-- figures/
|   |-- company.pdf
|   |-- system.png
|   `-- ...
|
`-- chapters/
    |-- chapter1.tex
    |-- chapter2.tex
    `-- chapter3.tex
\end{verbatim}

Figures should normally be stored in the \verb|figures| directory
specified by the official template.

File and directory names should preferably contain Latin letters,
numbers, hyphens, or underscores.

Recommended examples include:

\begin{verbatim}
company_structure.pdf
motor_control.png
chapter_01.tex
\end{verbatim}

Spaces and unnecessary special characters in file names should be
avoided where possible.

\vspace{0.3cm}

% ================================================================
% XXI. DOCUMENT COMPILATION AND FINAL CHECK
% ================================================================

\noindent
\textbf{XXI.\quad DOCUMENT COMPILATION AND FINAL CHECK}

\vspace{0.15cm}

Before submission, the complete \LaTeX{} project shall compile
successfully and produce the final PDF report.

The final document shall not contain unresolved references such as:

\begin{verbatim}
??
[?]
\end{verbatim}

Before generating the final PDF, students shall check:

\begin{itemize}[leftmargin=0.8cm, itemsep=2pt]
    \item the Table of Contents;
    \item page numbering;
    \item chapter, section, and subsection numbering;
    \item figure numbering;
    \item table numbering;
    \item equation numbering;
    \item cross-references;
    \item citations;
    \item bibliography;
    \item font display;
    \item mathematical symbols;
    \item image quality;
    \item completeness of all required forms.
\end{itemize}

The document may be compiled using TeX Live, MiKTeX, or another
compatible \LaTeX{} environment. Regardless of the compilation
environment, the resulting PDF shall conform to the official Faculty
template.

\vspace{0.3cm}

% ================================================================
% XXII. CONTENT THAT STUDENTS MAY EDIT
% ================================================================

\noindent
\textbf{XXII.\quad CONTENT THAT STUDENTS MAY EDIT}

\vspace{0.15cm}

Students may edit or replace:

\begin{itemize}[leftmargin=0.8cm, itemsep=2pt]
    \item student name;
    \item class;
    \item internship organization;
    \item supervisor name;
    \item Abstract;
    \item List of Symbols and Abbreviations;
    \item chapter titles and chapter content;
    \item figures and photographs;
    \item tables;
    \item equations;
    \item citations and references.
\end{itemize}

Students shall not modify the main document formatting settings in the
preamble, including the document class, page geometry, line spacing,
heading formatting, table-of-contents formatting, font configuration,
page layout, or cover-page macros unless specifically instructed.

\vspace{0.3cm}

% ================================================================
% XXIII. FINAL SUBMISSION REQUIREMENTS
% ================================================================

\noindent
\textbf{XXIII.\quad FINAL SUBMISSION REQUIREMENTS}

\vspace{0.15cm}

Before submission, the final report shall:

\begin{itemize}[leftmargin=0.8cm, itemsep=2pt]
    \item compile without errors that affect the report content or presentation;
    \item contain all required forms and report sections;
    \item contain correct student, class, internship organization, and supervisor information;
    \item contain an updated Table of Contents, List of Figures, and List of Tables;
    \item contain no unused example content from the template;
    \item contain no placeholder text such as \textit{Title of Chapter 1},
    \textit{Section 1.1}, or other sample entries;
    \item contain all required citations and references;
    \item contain no unresolved cross-references.
\end{itemize}

If the Faculty requires submission of the \LaTeX{} source, the complete
project shall include all \verb|.tex| files, the bibliography database,
figures, and any other files required to reproduce the submitted PDF.

\vspace{0.4cm}

\noindent
\textbf{Important:}
Before submitting the final internship report, students shall remove or
comment out the following command in the main document:

\begin{verbatim}
\input{Regulations.tex}
\end{verbatim}

For example:

\begin{verbatim}
% \input{Regulations.tex}
\end{verbatim}

This removes the guidance pages from the final submitted report while
preserving the regulations file in the \LaTeX{} project for reference.

} % End 13pt content

\clearpage

%
% IMPORTANT:
% This section is provided for guidance only. Students shall remove
% or comment out the 
% =================================================================
% REGULATIONS FOR PREPARING THE GRADUATION INTERNSHIP REPORT
% USING LATEX
% =================================================================
%
% This file is intended to be included from the main document with:
%
%     \input{Regulations.tex}
%
% IMPORTANT:
% This section is provided for guidance only. Students shall remove
% or comment out the \input{Regulations.tex} command before submitting
% the final internship report.
%
% =================================================================

\clearpage
\thispagestyle{empty}

\begin{center}
    {\fontsize{16pt}{20pt}\selectfont\bfseries
    SPECIFIC REGULATIONS FOR PREPARING\\
    GRADUATION INTERNSHIP REPORTS USING \LaTeX}
\end{center}

\vspace{0.5cm}

{\fontsize{13pt}{17pt}\selectfont
\setstretch{1.3}
\justifying

\noindent
\textbf{Note:}
This section is provided for guidance only and is included in the main
document using the command \verb|\input{Regulations.tex}|.
Students shall remove or comment out this command before submitting the
final internship report.

\vspace{0.4cm}

% ================================================================
% I. GENERAL REQUIREMENTS
% ================================================================

\noindent
\textbf{I.\quad GENERAL REQUIREMENTS}

\vspace{0.15cm}

The Graduation Internship Report shall be prepared using the official
\LaTeX{} template provided by the Faculty. Students shall preserve the
predefined page layout, typography, line spacing, heading styles,
numbering system, figure and table formatting, and bibliography settings
unless otherwise instructed by the Faculty or the supervisor.

The report shall be written in English.

Students are responsible for the academic and technical content of the
report. The official \LaTeX{} template is responsible for standardized
document formatting.

\vspace{0.3cm}

% ================================================================
% II. REPORT STRUCTURE
% ================================================================

\noindent
\textbf{II.\quad REPORT STRUCTURE}

\vspace{0.15cm}

The report shall normally be arranged in the following order:

\begin{enumerate}[label=\arabic*., leftmargin=1.2cm, itemsep=2pt]
    \item Outer Cover;
    \item Comments of the Internship Host Organization;
    \item Supervisor's Evaluation and Graduation Internship Defense Results;
    \item Table of Contents;
    \item List of Symbols and Abbreviations, if applicable;
    \item List of Figures;
    \item List of Tables;
    \item Abstract;
    \item Main Chapters of the Internship Report;
    \item References;
    \item Appendices, if applicable.
\end{enumerate}

The preliminary pages shall not be numbered as chapters. The main report
content shall be organized into consecutively numbered chapters beginning
with Chapter 1.

\vspace{0.3cm}

% ================================================================
% III. PAGE SIZE AND MARGINS
% ================================================================

\noindent
\textbf{III.\quad PAGE SIZE AND MARGINS}

\vspace{0.15cm}

The report shall be formatted on A4 paper with the following margins:

\begin{center}
\begin{tabular}{ll}
    Left margin:   & 3.5 cm \\
    Right margin:  & 2.0 cm \\
    Top margin:    & 2.5 cm \\
    Bottom margin: & 2.5 cm
\end{tabular}
\end{center}

These settings are predefined in the official template and shall not be
modified by students.

\vspace{0.3cm}

% ================================================================
% IV. FONT, FONT SIZE, AND LINE SPACING
% ================================================================

\noindent
\textbf{IV.\quad FONT, FONT SIZE, AND LINE SPACING}

\vspace{0.15cm}

The main text shall use a 13 pt font size with 1.3-line spacing.
A Times-family typeface shall be used in accordance with the official
template.

The text shall be fully justified.

Students shall not use repeated spaces, repeated manual line breaks, or
multiple blank lines to position text. Paragraph and heading spacing shall
be controlled by the template.

\vspace{0.3cm}

% ================================================================
% V. OUTER COVER
% ================================================================

\noindent
\textbf{V.\quad OUTER COVER}

\vspace{0.15cm}

The outer cover shall follow the official format provided in the template
and shall contain the following information:

\begin{itemize}[leftmargin=0.8cm, itemsep=2pt]
    \item Thai Nguyen University of Technology;
    \item Faculty of International Training;
    \item Engineering and Technology Department;
    \item University logo and Faculty logo;
    \item the title \textbf{INTERNSHIP REPORT};
    \item Internship Organization;
    \item Author;
    \item Class;
    \item Supervisor;
    \item place and year of completion.
\end{itemize}

Students shall only replace the required information fields and shall not
modify the cover-page layout unless instructed otherwise.

\vspace{0.3cm}

% ================================================================
% VI. COMMENTS OF THE INTERNSHIP HOST ORGANIZATION
% ================================================================

\noindent
\textbf{VI.\quad COMMENTS OF THE INTERNSHIP HOST ORGANIZATION}

\vspace{0.15cm}

The page entitled \textbf{COMMENTS OF THE INTERNSHIP HOST ORGANIZATION}
is reserved for the internship host organization.

It includes:

\begin{itemize}[leftmargin=0.8cm, itemsep=2pt]
    \item the student's full name;
    \item class;
    \item comments from the host organization;
    \item score;
    \item date;
    \item signature and official seal.
\end{itemize}

Students shall not alter or complete official evaluation fields unless
specifically authorized.

\vspace{0.3cm}

% ================================================================
% VII. SUPERVISOR'S EVALUATION AND DEFENSE RESULTS
% ================================================================

\noindent
\textbf{VII.\quad SUPERVISOR'S EVALUATION AND DEFENSE RESULTS}

\vspace{0.15cm}

The section entitled
\textbf{SUPERVISOR'S EVALUATION AND GRADUATION INTERNSHIP DEFENSE RESULTS}
is reserved for the supervisor and defense committee.

The evaluation may include:

\begin{itemize}[leftmargin=0.8cm, itemsep=2pt]
    \item Internship Progress and Student's Attitude;
    \item Report Content;
    \item Report Presentation;
    \item Other Comments;
    \item Supervisor's Opinion;
    \item Defense Result;
    \item Score.
\end{itemize}

Students shall not enter evaluation scores or defense results before they
are officially determined.

\vspace{0.3cm}

% ================================================================
% VIII. CHAPTERS, SECTIONS, AND SUBSECTIONS
% ================================================================

\noindent
\textbf{VIII.\quad CHAPTERS, SECTIONS, AND SUBSECTIONS}

\vspace{0.15cm}

The report structure shall be generated using the standard \LaTeX{}
commands:

\begin{verbatim}
\chapter{...}
\section{...}
\subsection{...}
\end{verbatim}

Chapter, section, and subsection numbers shall be generated automatically
by \LaTeX{}.

Students shall not manually type chapter or section numbers and shall not
reproduce heading styles using commands such as \verb|\textbf|,
\verb|\fontsize|, or \verb|\centering| instead of the appropriate
structural command.

Excessive subdivision should be avoided. In most cases, the hierarchy
Chapter--Section--Subsection is sufficient.

\vspace{0.3cm}

% ================================================================
% IX. ORGANIZATION OF THE MAIN CONTENT
% ================================================================

\noindent
\textbf{IX.\quad ORGANIZATION OF THE MAIN CONTENT}

\vspace{0.15cm}

The specific content of the report shall depend on the internship
organization and the assigned tasks. A typical structure may include:

\begin{itemize}[leftmargin=0.8cm, itemsep=2pt]
    \item \textbf{Chapter 1. Overview of the Internship Organization};
    \item \textbf{Chapter 2. Internship Activities and Assigned Tasks};
    \item \textbf{Chapter 3. Results, Experience, and Evaluation}.
\end{itemize}

These chapter titles are recommendations only and may be adapted to the
actual internship content.

The report shall reflect work actually performed, observed, analyzed, or
learned during the internship. Unrelated theoretical material should not
be included merely to increase report length.

\vspace{0.3cm}

% ================================================================
% X. TABLE OF CONTENTS
% ================================================================

\noindent
\textbf{X.\quad TABLE OF CONTENTS}

\vspace{0.15cm}

The Table of Contents shall be generated automatically using:

\begin{verbatim}
\tableofcontents
\end{verbatim}

Students shall not create or manually edit page numbers in the Table of
Contents.

The document shall be compiled sufficiently to ensure that all page
numbers and entries are correctly updated.

\vspace{0.3cm}

% ================================================================
% XI. LIST OF SYMBOLS AND ABBREVIATIONS
% ================================================================

\noindent
\textbf{XI.\quad LIST OF SYMBOLS AND ABBREVIATIONS}

\vspace{0.15cm}

Symbols and abbreviations that are used repeatedly in the report shall be
listed in the \textbf{List of Symbols and Abbreviations}.

The list may contain the following columns:

\begin{center}
\begin{tabular}{llll}
    No. & Symbol / Abbreviation & Unit & Definition
\end{tabular}
\end{center}

Only symbols and abbreviations actually used in the final report shall be
retained. Example entries contained in the template shall be removed if
they are not relevant to the student's report.

Each abbreviation should normally be defined when it first appears in
the main text.

\vspace{0.3cm}

% ================================================================
% XII. FIGURES, DIAGRAMS, AND PHOTOGRAPHS
% ================================================================

\noindent
\textbf{XII.\quad FIGURES, DIAGRAMS, AND PHOTOGRAPHS}

\vspace{0.15cm}

Figures shall be clear, legible, and directly relevant to the report
content.

Figures shall be inserted using the standard \verb|figure| environment.

Example:

\begin{verbatim}
\begin{figure}[H]
    \centering
    \includegraphics[width=0.8\textwidth]
                    {figures/company_structure.pdf}
    \caption{Organizational structure of the company}
    \label{fig:company-structure}
\end{figure}
\end{verbatim}

Figure captions shall be placed below the figure.

Each referenced figure shall have a unique \verb|\label| and shall be
referenced in the text using, for example:

\begin{verbatim}
Figure~\ref{fig:company-structure}
\end{verbatim}

Students shall not manually type figure numbers.

Figures reproduced or adapted from external sources shall include an
appropriate citation.

Photographs taken during the internship shall be relevant and sufficiently
clear. Students shall not include confidential, restricted, or proprietary
information without permission from the internship organization.

Vector graphics or high-resolution images are recommended whenever
possible.

\vspace{0.3cm}

% ================================================================
% XIII. TABLES
% ================================================================

\noindent
\textbf{XIII.\quad TABLES}

\vspace{0.15cm}

Tables shall be created using suitable \LaTeX{} environments such as
\verb|table|, \verb|tabular|, \verb|tabularx|, or \verb|longtable|.

Table captions shall be placed above the table.

Example:

\begin{verbatim}
\begin{table}[H]
    \centering
    \caption{Internship schedule}
    \label{tab:internship-schedule}

    \begin{tabular}{lll}
        \toprule
        Week & Task & Result \\
        \midrule
        1 & Orientation & Completed \\
        2 & System study & Completed \\
        \bottomrule
    \end{tabular}
\end{table}
\end{verbatim}

Tables shall be referenced using:

\begin{verbatim}
Table~\ref{tab:internship-schedule}
\end{verbatim}

Students shall not manually type table numbers.

Tables should be created directly in \LaTeX{} rather than inserted as
screenshots whenever practical.

\vspace{0.3cm}

% ================================================================
% XIV. EQUATIONS AND TECHNICAL QUANTITIES
% ================================================================

\noindent
\textbf{XIV.\quad EQUATIONS AND TECHNICAL QUANTITIES}

\vspace{0.15cm}

Mathematical expressions shall be typeset directly using \LaTeX{}.
Images or screenshots of equations shall not be used.

Short expressions may be written in inline mathematical mode, for example:

\begin{verbatim}
$R = 0.7\,\Omega$
\end{verbatim}

Important equations shall use a numbered equation environment.

Example:

\begin{verbatim}
\begin{equation}
    P = UI
    \label{eq:power}
\end{equation}
\end{verbatim}

Equations shall be referenced using:

\begin{verbatim}
Equation~\eqref{eq:power}
\end{verbatim}

Equation numbers shall not be manually entered.

SI units shall be used wherever applicable. The \verb|siunitx| package
should be used where appropriate.

Examples:

\begin{verbatim}
\SI{220}{\volt}
\SI{50}{\hertz}
\SI{2.5}{\ampere}
\end{verbatim}

Notation and units shall be consistent throughout the report.

\vspace{0.3cm}

% ================================================================
% XV. SOURCE CODE
% ================================================================

\noindent
\textbf{XV.\quad SOURCE CODE}

\vspace{0.15cm}

If the internship involves programming, PLCs, microcontrollers,
simulation, or software development, relevant source code may be included
in the report.

Source code shall be presented using the \verb|listings| environment or
another environment defined by the official template.

Example:

\begin{verbatim}
\begin{lstlisting}[language=st]
IF Start = TRUE THEN
    Motor := TRUE;
ELSE
    Motor := FALSE;
END_IF;
\end{lstlisting}
\end{verbatim}

Only code necessary for explaining the work should appear in the main
report. Long source-code listings should be moved to an appendix or
submitted as supplementary files if required.

Source code should not normally be inserted as screenshots.

\vspace{0.3cm}

% ================================================================
% XVI. ABSTRACT
% ================================================================

\noindent
\textbf{XVI.\quad ABSTRACT}

\vspace{0.15cm}

The Abstract shall briefly summarize:

\begin{itemize}[leftmargin=0.8cm, itemsep=2pt]
    \item the internship organization;
    \item the objectives of the internship;
    \item the main tasks performed;
    \item the main results, knowledge, and skills obtained.
\end{itemize}

The Abstract should be concise and should not reproduce the full content
of the report chapter by chapter.

\vspace{0.3cm}

% ================================================================
% XVII. REFERENCES AND CITATIONS
% ================================================================

\noindent
\textbf{XVII.\quad REFERENCES AND CITATIONS}

\vspace{0.15cm}

All books, journal papers, conference papers, standards, websites,
datasheets, manuals, technical documents, figures, tables, or other
external sources used in preparing the report shall be properly cited.

References shall be managed through the bibliography system provided with
the official template.

Bibliographic entries shall normally be stored in:

\begin{verbatim}
References.bib
\end{verbatim}

References shall be cited in the report using:

\begin{verbatim}
\cite{reference-key}
\end{verbatim}

Students shall not manually type reference numbers such as ``[1]'',
``[2]'', etc.

Titles of foreign-language publications shall remain in their original
language in the bibliography.

All borrowed ideas, data, figures, tables, equations, methods, or results
shall be properly acknowledged.

\vspace{0.3cm}

% ================================================================
% XVIII. CROSS-REFERENCING
% ================================================================

\noindent
\textbf{XVIII.\quad CROSS-REFERENCING}

\vspace{0.15cm}

Numbered elements shall be cross-referenced using the standard \LaTeX{}
mechanism.

Use the following commands where appropriate:

\begin{verbatim}
\label{...}
\ref{...}
\eqref{...}
\end{verbatim}

For example:

\begin{verbatim}
As shown in Figure~\ref{fig:system}, ...
\end{verbatim}

Students shall not manually type numbers such as ``Figure 2.3'' when
those numbers can be generated automatically.

\vspace{0.3cm}

% ================================================================
% XIX. PAGE NUMBERING
% ================================================================

\noindent
\textbf{XIX.\quad PAGE NUMBERING}

\vspace{0.15cm}

Cover pages and official evaluation forms at the beginning of the report
shall not display page numbers.

The main body of the report shall use Arabic page numbers beginning from
Chapter 1.

Page numbering shall be generated automatically by the official template.
Students shall not manually enter page numbers.

\vspace{0.3cm}

% ================================================================
% XX. PROJECT FILE ORGANIZATION
% ================================================================

\noindent
\textbf{XX.\quad PROJECT FILE ORGANIZATION}

\vspace{0.15cm}

Students should maintain a clear and organized \LaTeX{} project
structure. A recommended structure is:

\begin{verbatim}
InternshipReport/
|
|-- main.tex
|-- References.bib
|
|-- figures/
|   |-- company.pdf
|   |-- system.png
|   `-- ...
|
`-- chapters/
    |-- chapter1.tex
    |-- chapter2.tex
    `-- chapter3.tex
\end{verbatim}

Figures should normally be stored in the \verb|figures| directory
specified by the official template.

File and directory names should preferably contain Latin letters,
numbers, hyphens, or underscores.

Recommended examples include:

\begin{verbatim}
company_structure.pdf
motor_control.png
chapter_01.tex
\end{verbatim}

Spaces and unnecessary special characters in file names should be
avoided where possible.

\vspace{0.3cm}

% ================================================================
% XXI. DOCUMENT COMPILATION AND FINAL CHECK
% ================================================================

\noindent
\textbf{XXI.\quad DOCUMENT COMPILATION AND FINAL CHECK}

\vspace{0.15cm}

Before submission, the complete \LaTeX{} project shall compile
successfully and produce the final PDF report.

The final document shall not contain unresolved references such as:

\begin{verbatim}
??
[?]
\end{verbatim}

Before generating the final PDF, students shall check:

\begin{itemize}[leftmargin=0.8cm, itemsep=2pt]
    \item the Table of Contents;
    \item page numbering;
    \item chapter, section, and subsection numbering;
    \item figure numbering;
    \item table numbering;
    \item equation numbering;
    \item cross-references;
    \item citations;
    \item bibliography;
    \item font display;
    \item mathematical symbols;
    \item image quality;
    \item completeness of all required forms.
\end{itemize}

The document may be compiled using TeX Live, MiKTeX, or another
compatible \LaTeX{} environment. Regardless of the compilation
environment, the resulting PDF shall conform to the official Faculty
template.

\vspace{0.3cm}

% ================================================================
% XXII. CONTENT THAT STUDENTS MAY EDIT
% ================================================================

\noindent
\textbf{XXII.\quad CONTENT THAT STUDENTS MAY EDIT}

\vspace{0.15cm}

Students may edit or replace:

\begin{itemize}[leftmargin=0.8cm, itemsep=2pt]
    \item student name;
    \item class;
    \item internship organization;
    \item supervisor name;
    \item Abstract;
    \item List of Symbols and Abbreviations;
    \item chapter titles and chapter content;
    \item figures and photographs;
    \item tables;
    \item equations;
    \item citations and references.
\end{itemize}

Students shall not modify the main document formatting settings in the
preamble, including the document class, page geometry, line spacing,
heading formatting, table-of-contents formatting, font configuration,
page layout, or cover-page macros unless specifically instructed.

\vspace{0.3cm}

% ================================================================
% XXIII. FINAL SUBMISSION REQUIREMENTS
% ================================================================

\noindent
\textbf{XXIII.\quad FINAL SUBMISSION REQUIREMENTS}

\vspace{0.15cm}

Before submission, the final report shall:

\begin{itemize}[leftmargin=0.8cm, itemsep=2pt]
    \item compile without errors that affect the report content or presentation;
    \item contain all required forms and report sections;
    \item contain correct student, class, internship organization, and supervisor information;
    \item contain an updated Table of Contents, List of Figures, and List of Tables;
    \item contain no unused example content from the template;
    \item contain no placeholder text such as \textit{Title of Chapter 1},
    \textit{Section 1.1}, or other sample entries;
    \item contain all required citations and references;
    \item contain no unresolved cross-references.
\end{itemize}

If the Faculty requires submission of the \LaTeX{} source, the complete
project shall include all \verb|.tex| files, the bibliography database,
figures, and any other files required to reproduce the submitted PDF.

\vspace{0.4cm}

\noindent
\textbf{Important:}
Before submitting the final internship report, students shall remove or
comment out the following command in the main document:

\begin{verbatim}
\input{Regulations.tex}
\end{verbatim}

For example:

\begin{verbatim}
% \input{Regulations.tex}
\end{verbatim}

This removes the guidance pages from the final submitted report while
preserving the regulations file in the \LaTeX{} project for reference.

} % End 13pt content

\clearpage
 command before submitting
% the final internship report.
%
% =================================================================

\clearpage
\thispagestyle{empty}

\begin{center}
    {\fontsize{16pt}{20pt}\selectfont\bfseries
    SPECIFIC REGULATIONS FOR PREPARING\\
    GRADUATION INTERNSHIP REPORTS USING \LaTeX}
\end{center}

\vspace{0.5cm}

{\fontsize{13pt}{17pt}\selectfont
\setstretch{1.3}
\justifying

\noindent
\textbf{Note:}
This section is provided for guidance only and is included in the main
document using the command \verb|
% =================================================================
% REGULATIONS FOR PREPARING THE GRADUATION INTERNSHIP REPORT
% USING LATEX
% =================================================================
%
% This file is intended to be included from the main document with:
%
%     \input{Regulations.tex}
%
% IMPORTANT:
% This section is provided for guidance only. Students shall remove
% or comment out the \input{Regulations.tex} command before submitting
% the final internship report.
%
% =================================================================

\clearpage
\thispagestyle{empty}

\begin{center}
    {\fontsize{16pt}{20pt}\selectfont\bfseries
    SPECIFIC REGULATIONS FOR PREPARING\\
    GRADUATION INTERNSHIP REPORTS USING \LaTeX}
\end{center}

\vspace{0.5cm}

{\fontsize{13pt}{17pt}\selectfont
\setstretch{1.3}
\justifying

\noindent
\textbf{Note:}
This section is provided for guidance only and is included in the main
document using the command \verb|\input{Regulations.tex}|.
Students shall remove or comment out this command before submitting the
final internship report.

\vspace{0.4cm}

% ================================================================
% I. GENERAL REQUIREMENTS
% ================================================================

\noindent
\textbf{I.\quad GENERAL REQUIREMENTS}

\vspace{0.15cm}

The Graduation Internship Report shall be prepared using the official
\LaTeX{} template provided by the Faculty. Students shall preserve the
predefined page layout, typography, line spacing, heading styles,
numbering system, figure and table formatting, and bibliography settings
unless otherwise instructed by the Faculty or the supervisor.

The report shall be written in English.

Students are responsible for the academic and technical content of the
report. The official \LaTeX{} template is responsible for standardized
document formatting.

\vspace{0.3cm}

% ================================================================
% II. REPORT STRUCTURE
% ================================================================

\noindent
\textbf{II.\quad REPORT STRUCTURE}

\vspace{0.15cm}

The report shall normally be arranged in the following order:

\begin{enumerate}[label=\arabic*., leftmargin=1.2cm, itemsep=2pt]
    \item Outer Cover;
    \item Comments of the Internship Host Organization;
    \item Supervisor's Evaluation and Graduation Internship Defense Results;
    \item Table of Contents;
    \item List of Symbols and Abbreviations, if applicable;
    \item List of Figures;
    \item List of Tables;
    \item Abstract;
    \item Main Chapters of the Internship Report;
    \item References;
    \item Appendices, if applicable.
\end{enumerate}

The preliminary pages shall not be numbered as chapters. The main report
content shall be organized into consecutively numbered chapters beginning
with Chapter 1.

\vspace{0.3cm}

% ================================================================
% III. PAGE SIZE AND MARGINS
% ================================================================

\noindent
\textbf{III.\quad PAGE SIZE AND MARGINS}

\vspace{0.15cm}

The report shall be formatted on A4 paper with the following margins:

\begin{center}
\begin{tabular}{ll}
    Left margin:   & 3.5 cm \\
    Right margin:  & 2.0 cm \\
    Top margin:    & 2.5 cm \\
    Bottom margin: & 2.5 cm
\end{tabular}
\end{center}

These settings are predefined in the official template and shall not be
modified by students.

\vspace{0.3cm}

% ================================================================
% IV. FONT, FONT SIZE, AND LINE SPACING
% ================================================================

\noindent
\textbf{IV.\quad FONT, FONT SIZE, AND LINE SPACING}

\vspace{0.15cm}

The main text shall use a 13 pt font size with 1.3-line spacing.
A Times-family typeface shall be used in accordance with the official
template.

The text shall be fully justified.

Students shall not use repeated spaces, repeated manual line breaks, or
multiple blank lines to position text. Paragraph and heading spacing shall
be controlled by the template.

\vspace{0.3cm}

% ================================================================
% V. OUTER COVER
% ================================================================

\noindent
\textbf{V.\quad OUTER COVER}

\vspace{0.15cm}

The outer cover shall follow the official format provided in the template
and shall contain the following information:

\begin{itemize}[leftmargin=0.8cm, itemsep=2pt]
    \item Thai Nguyen University of Technology;
    \item Faculty of International Training;
    \item Engineering and Technology Department;
    \item University logo and Faculty logo;
    \item the title \textbf{INTERNSHIP REPORT};
    \item Internship Organization;
    \item Author;
    \item Class;
    \item Supervisor;
    \item place and year of completion.
\end{itemize}

Students shall only replace the required information fields and shall not
modify the cover-page layout unless instructed otherwise.

\vspace{0.3cm}

% ================================================================
% VI. COMMENTS OF THE INTERNSHIP HOST ORGANIZATION
% ================================================================

\noindent
\textbf{VI.\quad COMMENTS OF THE INTERNSHIP HOST ORGANIZATION}

\vspace{0.15cm}

The page entitled \textbf{COMMENTS OF THE INTERNSHIP HOST ORGANIZATION}
is reserved for the internship host organization.

It includes:

\begin{itemize}[leftmargin=0.8cm, itemsep=2pt]
    \item the student's full name;
    \item class;
    \item comments from the host organization;
    \item score;
    \item date;
    \item signature and official seal.
\end{itemize}

Students shall not alter or complete official evaluation fields unless
specifically authorized.

\vspace{0.3cm}

% ================================================================
% VII. SUPERVISOR'S EVALUATION AND DEFENSE RESULTS
% ================================================================

\noindent
\textbf{VII.\quad SUPERVISOR'S EVALUATION AND DEFENSE RESULTS}

\vspace{0.15cm}

The section entitled
\textbf{SUPERVISOR'S EVALUATION AND GRADUATION INTERNSHIP DEFENSE RESULTS}
is reserved for the supervisor and defense committee.

The evaluation may include:

\begin{itemize}[leftmargin=0.8cm, itemsep=2pt]
    \item Internship Progress and Student's Attitude;
    \item Report Content;
    \item Report Presentation;
    \item Other Comments;
    \item Supervisor's Opinion;
    \item Defense Result;
    \item Score.
\end{itemize}

Students shall not enter evaluation scores or defense results before they
are officially determined.

\vspace{0.3cm}

% ================================================================
% VIII. CHAPTERS, SECTIONS, AND SUBSECTIONS
% ================================================================

\noindent
\textbf{VIII.\quad CHAPTERS, SECTIONS, AND SUBSECTIONS}

\vspace{0.15cm}

The report structure shall be generated using the standard \LaTeX{}
commands:

\begin{verbatim}
\chapter{...}
\section{...}
\subsection{...}
\end{verbatim}

Chapter, section, and subsection numbers shall be generated automatically
by \LaTeX{}.

Students shall not manually type chapter or section numbers and shall not
reproduce heading styles using commands such as \verb|\textbf|,
\verb|\fontsize|, or \verb|\centering| instead of the appropriate
structural command.

Excessive subdivision should be avoided. In most cases, the hierarchy
Chapter--Section--Subsection is sufficient.

\vspace{0.3cm}

% ================================================================
% IX. ORGANIZATION OF THE MAIN CONTENT
% ================================================================

\noindent
\textbf{IX.\quad ORGANIZATION OF THE MAIN CONTENT}

\vspace{0.15cm}

The specific content of the report shall depend on the internship
organization and the assigned tasks. A typical structure may include:

\begin{itemize}[leftmargin=0.8cm, itemsep=2pt]
    \item \textbf{Chapter 1. Overview of the Internship Organization};
    \item \textbf{Chapter 2. Internship Activities and Assigned Tasks};
    \item \textbf{Chapter 3. Results, Experience, and Evaluation}.
\end{itemize}

These chapter titles are recommendations only and may be adapted to the
actual internship content.

The report shall reflect work actually performed, observed, analyzed, or
learned during the internship. Unrelated theoretical material should not
be included merely to increase report length.

\vspace{0.3cm}

% ================================================================
% X. TABLE OF CONTENTS
% ================================================================

\noindent
\textbf{X.\quad TABLE OF CONTENTS}

\vspace{0.15cm}

The Table of Contents shall be generated automatically using:

\begin{verbatim}
\tableofcontents
\end{verbatim}

Students shall not create or manually edit page numbers in the Table of
Contents.

The document shall be compiled sufficiently to ensure that all page
numbers and entries are correctly updated.

\vspace{0.3cm}

% ================================================================
% XI. LIST OF SYMBOLS AND ABBREVIATIONS
% ================================================================

\noindent
\textbf{XI.\quad LIST OF SYMBOLS AND ABBREVIATIONS}

\vspace{0.15cm}

Symbols and abbreviations that are used repeatedly in the report shall be
listed in the \textbf{List of Symbols and Abbreviations}.

The list may contain the following columns:

\begin{center}
\begin{tabular}{llll}
    No. & Symbol / Abbreviation & Unit & Definition
\end{tabular}
\end{center}

Only symbols and abbreviations actually used in the final report shall be
retained. Example entries contained in the template shall be removed if
they are not relevant to the student's report.

Each abbreviation should normally be defined when it first appears in
the main text.

\vspace{0.3cm}

% ================================================================
% XII. FIGURES, DIAGRAMS, AND PHOTOGRAPHS
% ================================================================

\noindent
\textbf{XII.\quad FIGURES, DIAGRAMS, AND PHOTOGRAPHS}

\vspace{0.15cm}

Figures shall be clear, legible, and directly relevant to the report
content.

Figures shall be inserted using the standard \verb|figure| environment.

Example:

\begin{verbatim}
\begin{figure}[H]
    \centering
    \includegraphics[width=0.8\textwidth]
                    {figures/company_structure.pdf}
    \caption{Organizational structure of the company}
    \label{fig:company-structure}
\end{figure}
\end{verbatim}

Figure captions shall be placed below the figure.

Each referenced figure shall have a unique \verb|\label| and shall be
referenced in the text using, for example:

\begin{verbatim}
Figure~\ref{fig:company-structure}
\end{verbatim}

Students shall not manually type figure numbers.

Figures reproduced or adapted from external sources shall include an
appropriate citation.

Photographs taken during the internship shall be relevant and sufficiently
clear. Students shall not include confidential, restricted, or proprietary
information without permission from the internship organization.

Vector graphics or high-resolution images are recommended whenever
possible.

\vspace{0.3cm}

% ================================================================
% XIII. TABLES
% ================================================================

\noindent
\textbf{XIII.\quad TABLES}

\vspace{0.15cm}

Tables shall be created using suitable \LaTeX{} environments such as
\verb|table|, \verb|tabular|, \verb|tabularx|, or \verb|longtable|.

Table captions shall be placed above the table.

Example:

\begin{verbatim}
\begin{table}[H]
    \centering
    \caption{Internship schedule}
    \label{tab:internship-schedule}

    \begin{tabular}{lll}
        \toprule
        Week & Task & Result \\
        \midrule
        1 & Orientation & Completed \\
        2 & System study & Completed \\
        \bottomrule
    \end{tabular}
\end{table}
\end{verbatim}

Tables shall be referenced using:

\begin{verbatim}
Table~\ref{tab:internship-schedule}
\end{verbatim}

Students shall not manually type table numbers.

Tables should be created directly in \LaTeX{} rather than inserted as
screenshots whenever practical.

\vspace{0.3cm}

% ================================================================
% XIV. EQUATIONS AND TECHNICAL QUANTITIES
% ================================================================

\noindent
\textbf{XIV.\quad EQUATIONS AND TECHNICAL QUANTITIES}

\vspace{0.15cm}

Mathematical expressions shall be typeset directly using \LaTeX{}.
Images or screenshots of equations shall not be used.

Short expressions may be written in inline mathematical mode, for example:

\begin{verbatim}
$R = 0.7\,\Omega$
\end{verbatim}

Important equations shall use a numbered equation environment.

Example:

\begin{verbatim}
\begin{equation}
    P = UI
    \label{eq:power}
\end{equation}
\end{verbatim}

Equations shall be referenced using:

\begin{verbatim}
Equation~\eqref{eq:power}
\end{verbatim}

Equation numbers shall not be manually entered.

SI units shall be used wherever applicable. The \verb|siunitx| package
should be used where appropriate.

Examples:

\begin{verbatim}
\SI{220}{\volt}
\SI{50}{\hertz}
\SI{2.5}{\ampere}
\end{verbatim}

Notation and units shall be consistent throughout the report.

\vspace{0.3cm}

% ================================================================
% XV. SOURCE CODE
% ================================================================

\noindent
\textbf{XV.\quad SOURCE CODE}

\vspace{0.15cm}

If the internship involves programming, PLCs, microcontrollers,
simulation, or software development, relevant source code may be included
in the report.

Source code shall be presented using the \verb|listings| environment or
another environment defined by the official template.

Example:

\begin{verbatim}
\begin{lstlisting}[language=st]
IF Start = TRUE THEN
    Motor := TRUE;
ELSE
    Motor := FALSE;
END_IF;
\end{lstlisting}
\end{verbatim}

Only code necessary for explaining the work should appear in the main
report. Long source-code listings should be moved to an appendix or
submitted as supplementary files if required.

Source code should not normally be inserted as screenshots.

\vspace{0.3cm}

% ================================================================
% XVI. ABSTRACT
% ================================================================

\noindent
\textbf{XVI.\quad ABSTRACT}

\vspace{0.15cm}

The Abstract shall briefly summarize:

\begin{itemize}[leftmargin=0.8cm, itemsep=2pt]
    \item the internship organization;
    \item the objectives of the internship;
    \item the main tasks performed;
    \item the main results, knowledge, and skills obtained.
\end{itemize}

The Abstract should be concise and should not reproduce the full content
of the report chapter by chapter.

\vspace{0.3cm}

% ================================================================
% XVII. REFERENCES AND CITATIONS
% ================================================================

\noindent
\textbf{XVII.\quad REFERENCES AND CITATIONS}

\vspace{0.15cm}

All books, journal papers, conference papers, standards, websites,
datasheets, manuals, technical documents, figures, tables, or other
external sources used in preparing the report shall be properly cited.

References shall be managed through the bibliography system provided with
the official template.

Bibliographic entries shall normally be stored in:

\begin{verbatim}
References.bib
\end{verbatim}

References shall be cited in the report using:

\begin{verbatim}
\cite{reference-key}
\end{verbatim}

Students shall not manually type reference numbers such as ``[1]'',
``[2]'', etc.

Titles of foreign-language publications shall remain in their original
language in the bibliography.

All borrowed ideas, data, figures, tables, equations, methods, or results
shall be properly acknowledged.

\vspace{0.3cm}

% ================================================================
% XVIII. CROSS-REFERENCING
% ================================================================

\noindent
\textbf{XVIII.\quad CROSS-REFERENCING}

\vspace{0.15cm}

Numbered elements shall be cross-referenced using the standard \LaTeX{}
mechanism.

Use the following commands where appropriate:

\begin{verbatim}
\label{...}
\ref{...}
\eqref{...}
\end{verbatim}

For example:

\begin{verbatim}
As shown in Figure~\ref{fig:system}, ...
\end{verbatim}

Students shall not manually type numbers such as ``Figure 2.3'' when
those numbers can be generated automatically.

\vspace{0.3cm}

% ================================================================
% XIX. PAGE NUMBERING
% ================================================================

\noindent
\textbf{XIX.\quad PAGE NUMBERING}

\vspace{0.15cm}

Cover pages and official evaluation forms at the beginning of the report
shall not display page numbers.

The main body of the report shall use Arabic page numbers beginning from
Chapter 1.

Page numbering shall be generated automatically by the official template.
Students shall not manually enter page numbers.

\vspace{0.3cm}

% ================================================================
% XX. PROJECT FILE ORGANIZATION
% ================================================================

\noindent
\textbf{XX.\quad PROJECT FILE ORGANIZATION}

\vspace{0.15cm}

Students should maintain a clear and organized \LaTeX{} project
structure. A recommended structure is:

\begin{verbatim}
InternshipReport/
|
|-- main.tex
|-- References.bib
|
|-- figures/
|   |-- company.pdf
|   |-- system.png
|   `-- ...
|
`-- chapters/
    |-- chapter1.tex
    |-- chapter2.tex
    `-- chapter3.tex
\end{verbatim}

Figures should normally be stored in the \verb|figures| directory
specified by the official template.

File and directory names should preferably contain Latin letters,
numbers, hyphens, or underscores.

Recommended examples include:

\begin{verbatim}
company_structure.pdf
motor_control.png
chapter_01.tex
\end{verbatim}

Spaces and unnecessary special characters in file names should be
avoided where possible.

\vspace{0.3cm}

% ================================================================
% XXI. DOCUMENT COMPILATION AND FINAL CHECK
% ================================================================

\noindent
\textbf{XXI.\quad DOCUMENT COMPILATION AND FINAL CHECK}

\vspace{0.15cm}

Before submission, the complete \LaTeX{} project shall compile
successfully and produce the final PDF report.

The final document shall not contain unresolved references such as:

\begin{verbatim}
??
[?]
\end{verbatim}

Before generating the final PDF, students shall check:

\begin{itemize}[leftmargin=0.8cm, itemsep=2pt]
    \item the Table of Contents;
    \item page numbering;
    \item chapter, section, and subsection numbering;
    \item figure numbering;
    \item table numbering;
    \item equation numbering;
    \item cross-references;
    \item citations;
    \item bibliography;
    \item font display;
    \item mathematical symbols;
    \item image quality;
    \item completeness of all required forms.
\end{itemize}

The document may be compiled using TeX Live, MiKTeX, or another
compatible \LaTeX{} environment. Regardless of the compilation
environment, the resulting PDF shall conform to the official Faculty
template.

\vspace{0.3cm}

% ================================================================
% XXII. CONTENT THAT STUDENTS MAY EDIT
% ================================================================

\noindent
\textbf{XXII.\quad CONTENT THAT STUDENTS MAY EDIT}

\vspace{0.15cm}

Students may edit or replace:

\begin{itemize}[leftmargin=0.8cm, itemsep=2pt]
    \item student name;
    \item class;
    \item internship organization;
    \item supervisor name;
    \item Abstract;
    \item List of Symbols and Abbreviations;
    \item chapter titles and chapter content;
    \item figures and photographs;
    \item tables;
    \item equations;
    \item citations and references.
\end{itemize}

Students shall not modify the main document formatting settings in the
preamble, including the document class, page geometry, line spacing,
heading formatting, table-of-contents formatting, font configuration,
page layout, or cover-page macros unless specifically instructed.

\vspace{0.3cm}

% ================================================================
% XXIII. FINAL SUBMISSION REQUIREMENTS
% ================================================================

\noindent
\textbf{XXIII.\quad FINAL SUBMISSION REQUIREMENTS}

\vspace{0.15cm}

Before submission, the final report shall:

\begin{itemize}[leftmargin=0.8cm, itemsep=2pt]
    \item compile without errors that affect the report content or presentation;
    \item contain all required forms and report sections;
    \item contain correct student, class, internship organization, and supervisor information;
    \item contain an updated Table of Contents, List of Figures, and List of Tables;
    \item contain no unused example content from the template;
    \item contain no placeholder text such as \textit{Title of Chapter 1},
    \textit{Section 1.1}, or other sample entries;
    \item contain all required citations and references;
    \item contain no unresolved cross-references.
\end{itemize}

If the Faculty requires submission of the \LaTeX{} source, the complete
project shall include all \verb|.tex| files, the bibliography database,
figures, and any other files required to reproduce the submitted PDF.

\vspace{0.4cm}

\noindent
\textbf{Important:}
Before submitting the final internship report, students shall remove or
comment out the following command in the main document:

\begin{verbatim}
\input{Regulations.tex}
\end{verbatim}

For example:

\begin{verbatim}
% \input{Regulations.tex}
\end{verbatim}

This removes the guidance pages from the final submitted report while
preserving the regulations file in the \LaTeX{} project for reference.

} % End 13pt content

\clearpage
|.
Students shall remove or comment out this command before submitting the
final internship report.

\vspace{0.4cm}

% ================================================================
% I. GENERAL REQUIREMENTS
% ================================================================

\noindent
\textbf{I.\quad GENERAL REQUIREMENTS}

\vspace{0.15cm}

The Graduation Internship Report shall be prepared using the official
\LaTeX{} template provided by the Faculty. Students shall preserve the
predefined page layout, typography, line spacing, heading styles,
numbering system, figure and table formatting, and bibliography settings
unless otherwise instructed by the Faculty or the supervisor.

The report shall be written in English.

Students are responsible for the academic and technical content of the
report. The official \LaTeX{} template is responsible for standardized
document formatting.

\vspace{0.3cm}

% ================================================================
% II. REPORT STRUCTURE
% ================================================================

\noindent
\textbf{II.\quad REPORT STRUCTURE}

\vspace{0.15cm}

The report shall normally be arranged in the following order:

\begin{enumerate}[label=\arabic*., leftmargin=1.2cm, itemsep=2pt]
    \item Outer Cover;
    \item Comments of the Internship Host Organization;
    \item Supervisor's Evaluation and Graduation Internship Defense Results;
    \item Table of Contents;
    \item List of Symbols and Abbreviations, if applicable;
    \item List of Figures;
    \item List of Tables;
    \item Abstract;
    \item Main Chapters of the Internship Report;
    \item References;
    \item Appendices, if applicable.
\end{enumerate}

The preliminary pages shall not be numbered as chapters. The main report
content shall be organized into consecutively numbered chapters beginning
with Chapter 1.

\vspace{0.3cm}

% ================================================================
% III. PAGE SIZE AND MARGINS
% ================================================================

\noindent
\textbf{III.\quad PAGE SIZE AND MARGINS}

\vspace{0.15cm}

The report shall be formatted on A4 paper with the following margins:

\begin{center}
\begin{tabular}{ll}
    Left margin:   & 3.5 cm \\
    Right margin:  & 2.0 cm \\
    Top margin:    & 2.5 cm \\
    Bottom margin: & 2.5 cm
\end{tabular}
\end{center}

These settings are predefined in the official template and shall not be
modified by students.

\vspace{0.3cm}

% ================================================================
% IV. FONT, FONT SIZE, AND LINE SPACING
% ================================================================

\noindent
\textbf{IV.\quad FONT, FONT SIZE, AND LINE SPACING}

\vspace{0.15cm}

The main text shall use a 13 pt font size with 1.3-line spacing.
A Times-family typeface shall be used in accordance with the official
template.

The text shall be fully justified.

Students shall not use repeated spaces, repeated manual line breaks, or
multiple blank lines to position text. Paragraph and heading spacing shall
be controlled by the template.

\vspace{0.3cm}

% ================================================================
% V. OUTER COVER
% ================================================================

\noindent
\textbf{V.\quad OUTER COVER}

\vspace{0.15cm}

The outer cover shall follow the official format provided in the template
and shall contain the following information:

\begin{itemize}[leftmargin=0.8cm, itemsep=2pt]
    \item Thai Nguyen University of Technology;
    \item Faculty of International Training;
    \item Engineering and Technology Department;
    \item University logo and Faculty logo;
    \item the title \textbf{INTERNSHIP REPORT};
    \item Internship Organization;
    \item Author;
    \item Class;
    \item Supervisor;
    \item place and year of completion.
\end{itemize}

Students shall only replace the required information fields and shall not
modify the cover-page layout unless instructed otherwise.

\vspace{0.3cm}

% ================================================================
% VI. COMMENTS OF THE INTERNSHIP HOST ORGANIZATION
% ================================================================

\noindent
\textbf{VI.\quad COMMENTS OF THE INTERNSHIP HOST ORGANIZATION}

\vspace{0.15cm}

The page entitled \textbf{COMMENTS OF THE INTERNSHIP HOST ORGANIZATION}
is reserved for the internship host organization.

It includes:

\begin{itemize}[leftmargin=0.8cm, itemsep=2pt]
    \item the student's full name;
    \item class;
    \item comments from the host organization;
    \item score;
    \item date;
    \item signature and official seal.
\end{itemize}

Students shall not alter or complete official evaluation fields unless
specifically authorized.

\vspace{0.3cm}

% ================================================================
% VII. SUPERVISOR'S EVALUATION AND DEFENSE RESULTS
% ================================================================

\noindent
\textbf{VII.\quad SUPERVISOR'S EVALUATION AND DEFENSE RESULTS}

\vspace{0.15cm}

The section entitled
\textbf{SUPERVISOR'S EVALUATION AND GRADUATION INTERNSHIP DEFENSE RESULTS}
is reserved for the supervisor and defense committee.

The evaluation may include:

\begin{itemize}[leftmargin=0.8cm, itemsep=2pt]
    \item Internship Progress and Student's Attitude;
    \item Report Content;
    \item Report Presentation;
    \item Other Comments;
    \item Supervisor's Opinion;
    \item Defense Result;
    \item Score.
\end{itemize}

Students shall not enter evaluation scores or defense results before they
are officially determined.

\vspace{0.3cm}

% ================================================================
% VIII. CHAPTERS, SECTIONS, AND SUBSECTIONS
% ================================================================

\noindent
\textbf{VIII.\quad CHAPTERS, SECTIONS, AND SUBSECTIONS}

\vspace{0.15cm}

The report structure shall be generated using the standard \LaTeX{}
commands:

\begin{verbatim}
\chapter{...}
\section{...}
\subsection{...}
\end{verbatim}

Chapter, section, and subsection numbers shall be generated automatically
by \LaTeX{}.

Students shall not manually type chapter or section numbers and shall not
reproduce heading styles using commands such as \verb|\textbf|,
\verb|\fontsize|, or \verb|\centering| instead of the appropriate
structural command.

Excessive subdivision should be avoided. In most cases, the hierarchy
Chapter--Section--Subsection is sufficient.

\vspace{0.3cm}

% ================================================================
% IX. ORGANIZATION OF THE MAIN CONTENT
% ================================================================

\noindent
\textbf{IX.\quad ORGANIZATION OF THE MAIN CONTENT}

\vspace{0.15cm}

The specific content of the report shall depend on the internship
organization and the assigned tasks. A typical structure may include:

\begin{itemize}[leftmargin=0.8cm, itemsep=2pt]
    \item \textbf{Chapter 1. Overview of the Internship Organization};
    \item \textbf{Chapter 2. Internship Activities and Assigned Tasks};
    \item \textbf{Chapter 3. Results, Experience, and Evaluation}.
\end{itemize}

These chapter titles are recommendations only and may be adapted to the
actual internship content.

The report shall reflect work actually performed, observed, analyzed, or
learned during the internship. Unrelated theoretical material should not
be included merely to increase report length.

\vspace{0.3cm}

% ================================================================
% X. TABLE OF CONTENTS
% ================================================================

\noindent
\textbf{X.\quad TABLE OF CONTENTS}

\vspace{0.15cm}

The Table of Contents shall be generated automatically using:

\begin{verbatim}
\tableofcontents
\end{verbatim}

Students shall not create or manually edit page numbers in the Table of
Contents.

The document shall be compiled sufficiently to ensure that all page
numbers and entries are correctly updated.

\vspace{0.3cm}

% ================================================================
% XI. LIST OF SYMBOLS AND ABBREVIATIONS
% ================================================================

\noindent
\textbf{XI.\quad LIST OF SYMBOLS AND ABBREVIATIONS}

\vspace{0.15cm}

Symbols and abbreviations that are used repeatedly in the report shall be
listed in the \textbf{List of Symbols and Abbreviations}.

The list may contain the following columns:

\begin{center}
\begin{tabular}{llll}
    No. & Symbol / Abbreviation & Unit & Definition
\end{tabular}
\end{center}

Only symbols and abbreviations actually used in the final report shall be
retained. Example entries contained in the template shall be removed if
they are not relevant to the student's report.

Each abbreviation should normally be defined when it first appears in
the main text.

\vspace{0.3cm}

% ================================================================
% XII. FIGURES, DIAGRAMS, AND PHOTOGRAPHS
% ================================================================

\noindent
\textbf{XII.\quad FIGURES, DIAGRAMS, AND PHOTOGRAPHS}

\vspace{0.15cm}

Figures shall be clear, legible, and directly relevant to the report
content.

Figures shall be inserted using the standard \verb|figure| environment.

Example:

\begin{verbatim}
\begin{figure}[H]
    \centering
    \includegraphics[width=0.8\textwidth]
                    {figures/company_structure.pdf}
    \caption{Organizational structure of the company}
    \label{fig:company-structure}
\end{figure}
\end{verbatim}

Figure captions shall be placed below the figure.

Each referenced figure shall have a unique \verb|\label| and shall be
referenced in the text using, for example:

\begin{verbatim}
Figure~\ref{fig:company-structure}
\end{verbatim}

Students shall not manually type figure numbers.

Figures reproduced or adapted from external sources shall include an
appropriate citation.

Photographs taken during the internship shall be relevant and sufficiently
clear. Students shall not include confidential, restricted, or proprietary
information without permission from the internship organization.

Vector graphics or high-resolution images are recommended whenever
possible.

\vspace{0.3cm}

% ================================================================
% XIII. TABLES
% ================================================================

\noindent
\textbf{XIII.\quad TABLES}

\vspace{0.15cm}

Tables shall be created using suitable \LaTeX{} environments such as
\verb|table|, \verb|tabular|, \verb|tabularx|, or \verb|longtable|.

Table captions shall be placed above the table.

Example:

\begin{verbatim}
\begin{table}[H]
    \centering
    \caption{Internship schedule}
    \label{tab:internship-schedule}

    \begin{tabular}{lll}
        \toprule
        Week & Task & Result \\
        \midrule
        1 & Orientation & Completed \\
        2 & System study & Completed \\
        \bottomrule
    \end{tabular}
\end{table}
\end{verbatim}

Tables shall be referenced using:

\begin{verbatim}
Table~\ref{tab:internship-schedule}
\end{verbatim}

Students shall not manually type table numbers.

Tables should be created directly in \LaTeX{} rather than inserted as
screenshots whenever practical.

\vspace{0.3cm}

% ================================================================
% XIV. EQUATIONS AND TECHNICAL QUANTITIES
% ================================================================

\noindent
\textbf{XIV.\quad EQUATIONS AND TECHNICAL QUANTITIES}

\vspace{0.15cm}

Mathematical expressions shall be typeset directly using \LaTeX{}.
Images or screenshots of equations shall not be used.

Short expressions may be written in inline mathematical mode, for example:

\begin{verbatim}
$R = 0.7\,\Omega$
\end{verbatim}

Important equations shall use a numbered equation environment.

Example:

\begin{verbatim}
\begin{equation}
    P = UI
    \label{eq:power}
\end{equation}
\end{verbatim}

Equations shall be referenced using:

\begin{verbatim}
Equation~\eqref{eq:power}
\end{verbatim}

Equation numbers shall not be manually entered.

SI units shall be used wherever applicable. The \verb|siunitx| package
should be used where appropriate.

Examples:

\begin{verbatim}
\SI{220}{\volt}
\SI{50}{\hertz}
\SI{2.5}{\ampere}
\end{verbatim}

Notation and units shall be consistent throughout the report.

\vspace{0.3cm}

% ================================================================
% XV. SOURCE CODE
% ================================================================

\noindent
\textbf{XV.\quad SOURCE CODE}

\vspace{0.15cm}

If the internship involves programming, PLCs, microcontrollers,
simulation, or software development, relevant source code may be included
in the report.

Source code shall be presented using the \verb|listings| environment or
another environment defined by the official template.

Example:

\begin{verbatim}
\begin{lstlisting}[language=st]
IF Start = TRUE THEN
    Motor := TRUE;
ELSE
    Motor := FALSE;
END_IF;
\end{lstlisting}
\end{verbatim}

Only code necessary for explaining the work should appear in the main
report. Long source-code listings should be moved to an appendix or
submitted as supplementary files if required.

Source code should not normally be inserted as screenshots.

\vspace{0.3cm}

% ================================================================
% XVI. ABSTRACT
% ================================================================

\noindent
\textbf{XVI.\quad ABSTRACT}

\vspace{0.15cm}

The Abstract shall briefly summarize:

\begin{itemize}[leftmargin=0.8cm, itemsep=2pt]
    \item the internship organization;
    \item the objectives of the internship;
    \item the main tasks performed;
    \item the main results, knowledge, and skills obtained.
\end{itemize}

The Abstract should be concise and should not reproduce the full content
of the report chapter by chapter.

\vspace{0.3cm}

% ================================================================
% XVII. REFERENCES AND CITATIONS
% ================================================================

\noindent
\textbf{XVII.\quad REFERENCES AND CITATIONS}

\vspace{0.15cm}

All books, journal papers, conference papers, standards, websites,
datasheets, manuals, technical documents, figures, tables, or other
external sources used in preparing the report shall be properly cited.

References shall be managed through the bibliography system provided with
the official template.

Bibliographic entries shall normally be stored in:

\begin{verbatim}
References.bib
\end{verbatim}

References shall be cited in the report using:

\begin{verbatim}
\cite{reference-key}
\end{verbatim}

Students shall not manually type reference numbers such as ``[1]'',
``[2]'', etc.

Titles of foreign-language publications shall remain in their original
language in the bibliography.

All borrowed ideas, data, figures, tables, equations, methods, or results
shall be properly acknowledged.

\vspace{0.3cm}

% ================================================================
% XVIII. CROSS-REFERENCING
% ================================================================

\noindent
\textbf{XVIII.\quad CROSS-REFERENCING}

\vspace{0.15cm}

Numbered elements shall be cross-referenced using the standard \LaTeX{}
mechanism.

Use the following commands where appropriate:

\begin{verbatim}
\label{...}
\ref{...}
\eqref{...}
\end{verbatim}

For example:

\begin{verbatim}
As shown in Figure~\ref{fig:system}, ...
\end{verbatim}

Students shall not manually type numbers such as ``Figure 2.3'' when
those numbers can be generated automatically.

\vspace{0.3cm}

% ================================================================
% XIX. PAGE NUMBERING
% ================================================================

\noindent
\textbf{XIX.\quad PAGE NUMBERING}

\vspace{0.15cm}

Cover pages and official evaluation forms at the beginning of the report
shall not display page numbers.

The main body of the report shall use Arabic page numbers beginning from
Chapter 1.

Page numbering shall be generated automatically by the official template.
Students shall not manually enter page numbers.

\vspace{0.3cm}

% ================================================================
% XX. PROJECT FILE ORGANIZATION
% ================================================================

\noindent
\textbf{XX.\quad PROJECT FILE ORGANIZATION}

\vspace{0.15cm}

Students should maintain a clear and organized \LaTeX{} project
structure. A recommended structure is:

\begin{verbatim}
InternshipReport/
|
|-- main.tex
|-- References.bib
|
|-- figures/
|   |-- company.pdf
|   |-- system.png
|   `-- ...
|
`-- chapters/
    |-- chapter1.tex
    |-- chapter2.tex
    `-- chapter3.tex
\end{verbatim}

Figures should normally be stored in the \verb|figures| directory
specified by the official template.

File and directory names should preferably contain Latin letters,
numbers, hyphens, or underscores.

Recommended examples include:

\begin{verbatim}
company_structure.pdf
motor_control.png
chapter_01.tex
\end{verbatim}

Spaces and unnecessary special characters in file names should be
avoided where possible.

\vspace{0.3cm}

% ================================================================
% XXI. DOCUMENT COMPILATION AND FINAL CHECK
% ================================================================

\noindent
\textbf{XXI.\quad DOCUMENT COMPILATION AND FINAL CHECK}

\vspace{0.15cm}

Before submission, the complete \LaTeX{} project shall compile
successfully and produce the final PDF report.

The final document shall not contain unresolved references such as:

\begin{verbatim}
??
[?]
\end{verbatim}

Before generating the final PDF, students shall check:

\begin{itemize}[leftmargin=0.8cm, itemsep=2pt]
    \item the Table of Contents;
    \item page numbering;
    \item chapter, section, and subsection numbering;
    \item figure numbering;
    \item table numbering;
    \item equation numbering;
    \item cross-references;
    \item citations;
    \item bibliography;
    \item font display;
    \item mathematical symbols;
    \item image quality;
    \item completeness of all required forms.
\end{itemize}

The document may be compiled using TeX Live, MiKTeX, or another
compatible \LaTeX{} environment. Regardless of the compilation
environment, the resulting PDF shall conform to the official Faculty
template.

\vspace{0.3cm}

% ================================================================
% XXII. CONTENT THAT STUDENTS MAY EDIT
% ================================================================

\noindent
\textbf{XXII.\quad CONTENT THAT STUDENTS MAY EDIT}

\vspace{0.15cm}

Students may edit or replace:

\begin{itemize}[leftmargin=0.8cm, itemsep=2pt]
    \item student name;
    \item class;
    \item internship organization;
    \item supervisor name;
    \item Abstract;
    \item List of Symbols and Abbreviations;
    \item chapter titles and chapter content;
    \item figures and photographs;
    \item tables;
    \item equations;
    \item citations and references.
\end{itemize}

Students shall not modify the main document formatting settings in the
preamble, including the document class, page geometry, line spacing,
heading formatting, table-of-contents formatting, font configuration,
page layout, or cover-page macros unless specifically instructed.

\vspace{0.3cm}

% ================================================================
% XXIII. FINAL SUBMISSION REQUIREMENTS
% ================================================================

\noindent
\textbf{XXIII.\quad FINAL SUBMISSION REQUIREMENTS}

\vspace{0.15cm}

Before submission, the final report shall:

\begin{itemize}[leftmargin=0.8cm, itemsep=2pt]
    \item compile without errors that affect the report content or presentation;
    \item contain all required forms and report sections;
    \item contain correct student, class, internship organization, and supervisor information;
    \item contain an updated Table of Contents, List of Figures, and List of Tables;
    \item contain no unused example content from the template;
    \item contain no placeholder text such as \textit{Title of Chapter 1},
    \textit{Section 1.1}, or other sample entries;
    \item contain all required citations and references;
    \item contain no unresolved cross-references.
\end{itemize}

If the Faculty requires submission of the \LaTeX{} source, the complete
project shall include all \verb|.tex| files, the bibliography database,
figures, and any other files required to reproduce the submitted PDF.

\vspace{0.4cm}

\noindent
\textbf{Important:}
Before submitting the final internship report, students shall remove or
comment out the following command in the main document:

\begin{verbatim}

% =================================================================
% REGULATIONS FOR PREPARING THE GRADUATION INTERNSHIP REPORT
% USING LATEX
% =================================================================
%
% This file is intended to be included from the main document with:
%
%     \input{Regulations.tex}
%
% IMPORTANT:
% This section is provided for guidance only. Students shall remove
% or comment out the \input{Regulations.tex} command before submitting
% the final internship report.
%
% =================================================================

\clearpage
\thispagestyle{empty}

\begin{center}
    {\fontsize{16pt}{20pt}\selectfont\bfseries
    SPECIFIC REGULATIONS FOR PREPARING\\
    GRADUATION INTERNSHIP REPORTS USING \LaTeX}
\end{center}

\vspace{0.5cm}

{\fontsize{13pt}{17pt}\selectfont
\setstretch{1.3}
\justifying

\noindent
\textbf{Note:}
This section is provided for guidance only and is included in the main
document using the command \verb|\input{Regulations.tex}|.
Students shall remove or comment out this command before submitting the
final internship report.

\vspace{0.4cm}

% ================================================================
% I. GENERAL REQUIREMENTS
% ================================================================

\noindent
\textbf{I.\quad GENERAL REQUIREMENTS}

\vspace{0.15cm}

The Graduation Internship Report shall be prepared using the official
\LaTeX{} template provided by the Faculty. Students shall preserve the
predefined page layout, typography, line spacing, heading styles,
numbering system, figure and table formatting, and bibliography settings
unless otherwise instructed by the Faculty or the supervisor.

The report shall be written in English.

Students are responsible for the academic and technical content of the
report. The official \LaTeX{} template is responsible for standardized
document formatting.

\vspace{0.3cm}

% ================================================================
% II. REPORT STRUCTURE
% ================================================================

\noindent
\textbf{II.\quad REPORT STRUCTURE}

\vspace{0.15cm}

The report shall normally be arranged in the following order:

\begin{enumerate}[label=\arabic*., leftmargin=1.2cm, itemsep=2pt]
    \item Outer Cover;
    \item Comments of the Internship Host Organization;
    \item Supervisor's Evaluation and Graduation Internship Defense Results;
    \item Table of Contents;
    \item List of Symbols and Abbreviations, if applicable;
    \item List of Figures;
    \item List of Tables;
    \item Abstract;
    \item Main Chapters of the Internship Report;
    \item References;
    \item Appendices, if applicable.
\end{enumerate}

The preliminary pages shall not be numbered as chapters. The main report
content shall be organized into consecutively numbered chapters beginning
with Chapter 1.

\vspace{0.3cm}

% ================================================================
% III. PAGE SIZE AND MARGINS
% ================================================================

\noindent
\textbf{III.\quad PAGE SIZE AND MARGINS}

\vspace{0.15cm}

The report shall be formatted on A4 paper with the following margins:

\begin{center}
\begin{tabular}{ll}
    Left margin:   & 3.5 cm \\
    Right margin:  & 2.0 cm \\
    Top margin:    & 2.5 cm \\
    Bottom margin: & 2.5 cm
\end{tabular}
\end{center}

These settings are predefined in the official template and shall not be
modified by students.

\vspace{0.3cm}

% ================================================================
% IV. FONT, FONT SIZE, AND LINE SPACING
% ================================================================

\noindent
\textbf{IV.\quad FONT, FONT SIZE, AND LINE SPACING}

\vspace{0.15cm}

The main text shall use a 13 pt font size with 1.3-line spacing.
A Times-family typeface shall be used in accordance with the official
template.

The text shall be fully justified.

Students shall not use repeated spaces, repeated manual line breaks, or
multiple blank lines to position text. Paragraph and heading spacing shall
be controlled by the template.

\vspace{0.3cm}

% ================================================================
% V. OUTER COVER
% ================================================================

\noindent
\textbf{V.\quad OUTER COVER}

\vspace{0.15cm}

The outer cover shall follow the official format provided in the template
and shall contain the following information:

\begin{itemize}[leftmargin=0.8cm, itemsep=2pt]
    \item Thai Nguyen University of Technology;
    \item Faculty of International Training;
    \item Engineering and Technology Department;
    \item University logo and Faculty logo;
    \item the title \textbf{INTERNSHIP REPORT};
    \item Internship Organization;
    \item Author;
    \item Class;
    \item Supervisor;
    \item place and year of completion.
\end{itemize}

Students shall only replace the required information fields and shall not
modify the cover-page layout unless instructed otherwise.

\vspace{0.3cm}

% ================================================================
% VI. COMMENTS OF THE INTERNSHIP HOST ORGANIZATION
% ================================================================

\noindent
\textbf{VI.\quad COMMENTS OF THE INTERNSHIP HOST ORGANIZATION}

\vspace{0.15cm}

The page entitled \textbf{COMMENTS OF THE INTERNSHIP HOST ORGANIZATION}
is reserved for the internship host organization.

It includes:

\begin{itemize}[leftmargin=0.8cm, itemsep=2pt]
    \item the student's full name;
    \item class;
    \item comments from the host organization;
    \item score;
    \item date;
    \item signature and official seal.
\end{itemize}

Students shall not alter or complete official evaluation fields unless
specifically authorized.

\vspace{0.3cm}

% ================================================================
% VII. SUPERVISOR'S EVALUATION AND DEFENSE RESULTS
% ================================================================

\noindent
\textbf{VII.\quad SUPERVISOR'S EVALUATION AND DEFENSE RESULTS}

\vspace{0.15cm}

The section entitled
\textbf{SUPERVISOR'S EVALUATION AND GRADUATION INTERNSHIP DEFENSE RESULTS}
is reserved for the supervisor and defense committee.

The evaluation may include:

\begin{itemize}[leftmargin=0.8cm, itemsep=2pt]
    \item Internship Progress and Student's Attitude;
    \item Report Content;
    \item Report Presentation;
    \item Other Comments;
    \item Supervisor's Opinion;
    \item Defense Result;
    \item Score.
\end{itemize}

Students shall not enter evaluation scores or defense results before they
are officially determined.

\vspace{0.3cm}

% ================================================================
% VIII. CHAPTERS, SECTIONS, AND SUBSECTIONS
% ================================================================

\noindent
\textbf{VIII.\quad CHAPTERS, SECTIONS, AND SUBSECTIONS}

\vspace{0.15cm}

The report structure shall be generated using the standard \LaTeX{}
commands:

\begin{verbatim}
\chapter{...}
\section{...}
\subsection{...}
\end{verbatim}

Chapter, section, and subsection numbers shall be generated automatically
by \LaTeX{}.

Students shall not manually type chapter or section numbers and shall not
reproduce heading styles using commands such as \verb|\textbf|,
\verb|\fontsize|, or \verb|\centering| instead of the appropriate
structural command.

Excessive subdivision should be avoided. In most cases, the hierarchy
Chapter--Section--Subsection is sufficient.

\vspace{0.3cm}

% ================================================================
% IX. ORGANIZATION OF THE MAIN CONTENT
% ================================================================

\noindent
\textbf{IX.\quad ORGANIZATION OF THE MAIN CONTENT}

\vspace{0.15cm}

The specific content of the report shall depend on the internship
organization and the assigned tasks. A typical structure may include:

\begin{itemize}[leftmargin=0.8cm, itemsep=2pt]
    \item \textbf{Chapter 1. Overview of the Internship Organization};
    \item \textbf{Chapter 2. Internship Activities and Assigned Tasks};
    \item \textbf{Chapter 3. Results, Experience, and Evaluation}.
\end{itemize}

These chapter titles are recommendations only and may be adapted to the
actual internship content.

The report shall reflect work actually performed, observed, analyzed, or
learned during the internship. Unrelated theoretical material should not
be included merely to increase report length.

\vspace{0.3cm}

% ================================================================
% X. TABLE OF CONTENTS
% ================================================================

\noindent
\textbf{X.\quad TABLE OF CONTENTS}

\vspace{0.15cm}

The Table of Contents shall be generated automatically using:

\begin{verbatim}
\tableofcontents
\end{verbatim}

Students shall not create or manually edit page numbers in the Table of
Contents.

The document shall be compiled sufficiently to ensure that all page
numbers and entries are correctly updated.

\vspace{0.3cm}

% ================================================================
% XI. LIST OF SYMBOLS AND ABBREVIATIONS
% ================================================================

\noindent
\textbf{XI.\quad LIST OF SYMBOLS AND ABBREVIATIONS}

\vspace{0.15cm}

Symbols and abbreviations that are used repeatedly in the report shall be
listed in the \textbf{List of Symbols and Abbreviations}.

The list may contain the following columns:

\begin{center}
\begin{tabular}{llll}
    No. & Symbol / Abbreviation & Unit & Definition
\end{tabular}
\end{center}

Only symbols and abbreviations actually used in the final report shall be
retained. Example entries contained in the template shall be removed if
they are not relevant to the student's report.

Each abbreviation should normally be defined when it first appears in
the main text.

\vspace{0.3cm}

% ================================================================
% XII. FIGURES, DIAGRAMS, AND PHOTOGRAPHS
% ================================================================

\noindent
\textbf{XII.\quad FIGURES, DIAGRAMS, AND PHOTOGRAPHS}

\vspace{0.15cm}

Figures shall be clear, legible, and directly relevant to the report
content.

Figures shall be inserted using the standard \verb|figure| environment.

Example:

\begin{verbatim}
\begin{figure}[H]
    \centering
    \includegraphics[width=0.8\textwidth]
                    {figures/company_structure.pdf}
    \caption{Organizational structure of the company}
    \label{fig:company-structure}
\end{figure}
\end{verbatim}

Figure captions shall be placed below the figure.

Each referenced figure shall have a unique \verb|\label| and shall be
referenced in the text using, for example:

\begin{verbatim}
Figure~\ref{fig:company-structure}
\end{verbatim}

Students shall not manually type figure numbers.

Figures reproduced or adapted from external sources shall include an
appropriate citation.

Photographs taken during the internship shall be relevant and sufficiently
clear. Students shall not include confidential, restricted, or proprietary
information without permission from the internship organization.

Vector graphics or high-resolution images are recommended whenever
possible.

\vspace{0.3cm}

% ================================================================
% XIII. TABLES
% ================================================================

\noindent
\textbf{XIII.\quad TABLES}

\vspace{0.15cm}

Tables shall be created using suitable \LaTeX{} environments such as
\verb|table|, \verb|tabular|, \verb|tabularx|, or \verb|longtable|.

Table captions shall be placed above the table.

Example:

\begin{verbatim}
\begin{table}[H]
    \centering
    \caption{Internship schedule}
    \label{tab:internship-schedule}

    \begin{tabular}{lll}
        \toprule
        Week & Task & Result \\
        \midrule
        1 & Orientation & Completed \\
        2 & System study & Completed \\
        \bottomrule
    \end{tabular}
\end{table}
\end{verbatim}

Tables shall be referenced using:

\begin{verbatim}
Table~\ref{tab:internship-schedule}
\end{verbatim}

Students shall not manually type table numbers.

Tables should be created directly in \LaTeX{} rather than inserted as
screenshots whenever practical.

\vspace{0.3cm}

% ================================================================
% XIV. EQUATIONS AND TECHNICAL QUANTITIES
% ================================================================

\noindent
\textbf{XIV.\quad EQUATIONS AND TECHNICAL QUANTITIES}

\vspace{0.15cm}

Mathematical expressions shall be typeset directly using \LaTeX{}.
Images or screenshots of equations shall not be used.

Short expressions may be written in inline mathematical mode, for example:

\begin{verbatim}
$R = 0.7\,\Omega$
\end{verbatim}

Important equations shall use a numbered equation environment.

Example:

\begin{verbatim}
\begin{equation}
    P = UI
    \label{eq:power}
\end{equation}
\end{verbatim}

Equations shall be referenced using:

\begin{verbatim}
Equation~\eqref{eq:power}
\end{verbatim}

Equation numbers shall not be manually entered.

SI units shall be used wherever applicable. The \verb|siunitx| package
should be used where appropriate.

Examples:

\begin{verbatim}
\SI{220}{\volt}
\SI{50}{\hertz}
\SI{2.5}{\ampere}
\end{verbatim}

Notation and units shall be consistent throughout the report.

\vspace{0.3cm}

% ================================================================
% XV. SOURCE CODE
% ================================================================

\noindent
\textbf{XV.\quad SOURCE CODE}

\vspace{0.15cm}

If the internship involves programming, PLCs, microcontrollers,
simulation, or software development, relevant source code may be included
in the report.

Source code shall be presented using the \verb|listings| environment or
another environment defined by the official template.

Example:

\begin{verbatim}
\begin{lstlisting}[language=st]
IF Start = TRUE THEN
    Motor := TRUE;
ELSE
    Motor := FALSE;
END_IF;
\end{lstlisting}
\end{verbatim}

Only code necessary for explaining the work should appear in the main
report. Long source-code listings should be moved to an appendix or
submitted as supplementary files if required.

Source code should not normally be inserted as screenshots.

\vspace{0.3cm}

% ================================================================
% XVI. ABSTRACT
% ================================================================

\noindent
\textbf{XVI.\quad ABSTRACT}

\vspace{0.15cm}

The Abstract shall briefly summarize:

\begin{itemize}[leftmargin=0.8cm, itemsep=2pt]
    \item the internship organization;
    \item the objectives of the internship;
    \item the main tasks performed;
    \item the main results, knowledge, and skills obtained.
\end{itemize}

The Abstract should be concise and should not reproduce the full content
of the report chapter by chapter.

\vspace{0.3cm}

% ================================================================
% XVII. REFERENCES AND CITATIONS
% ================================================================

\noindent
\textbf{XVII.\quad REFERENCES AND CITATIONS}

\vspace{0.15cm}

All books, journal papers, conference papers, standards, websites,
datasheets, manuals, technical documents, figures, tables, or other
external sources used in preparing the report shall be properly cited.

References shall be managed through the bibliography system provided with
the official template.

Bibliographic entries shall normally be stored in:

\begin{verbatim}
References.bib
\end{verbatim}

References shall be cited in the report using:

\begin{verbatim}
\cite{reference-key}
\end{verbatim}

Students shall not manually type reference numbers such as ``[1]'',
``[2]'', etc.

Titles of foreign-language publications shall remain in their original
language in the bibliography.

All borrowed ideas, data, figures, tables, equations, methods, or results
shall be properly acknowledged.

\vspace{0.3cm}

% ================================================================
% XVIII. CROSS-REFERENCING
% ================================================================

\noindent
\textbf{XVIII.\quad CROSS-REFERENCING}

\vspace{0.15cm}

Numbered elements shall be cross-referenced using the standard \LaTeX{}
mechanism.

Use the following commands where appropriate:

\begin{verbatim}
\label{...}
\ref{...}
\eqref{...}
\end{verbatim}

For example:

\begin{verbatim}
As shown in Figure~\ref{fig:system}, ...
\end{verbatim}

Students shall not manually type numbers such as ``Figure 2.3'' when
those numbers can be generated automatically.

\vspace{0.3cm}

% ================================================================
% XIX. PAGE NUMBERING
% ================================================================

\noindent
\textbf{XIX.\quad PAGE NUMBERING}

\vspace{0.15cm}

Cover pages and official evaluation forms at the beginning of the report
shall not display page numbers.

The main body of the report shall use Arabic page numbers beginning from
Chapter 1.

Page numbering shall be generated automatically by the official template.
Students shall not manually enter page numbers.

\vspace{0.3cm}

% ================================================================
% XX. PROJECT FILE ORGANIZATION
% ================================================================

\noindent
\textbf{XX.\quad PROJECT FILE ORGANIZATION}

\vspace{0.15cm}

Students should maintain a clear and organized \LaTeX{} project
structure. A recommended structure is:

\begin{verbatim}
InternshipReport/
|
|-- main.tex
|-- References.bib
|
|-- figures/
|   |-- company.pdf
|   |-- system.png
|   `-- ...
|
`-- chapters/
    |-- chapter1.tex
    |-- chapter2.tex
    `-- chapter3.tex
\end{verbatim}

Figures should normally be stored in the \verb|figures| directory
specified by the official template.

File and directory names should preferably contain Latin letters,
numbers, hyphens, or underscores.

Recommended examples include:

\begin{verbatim}
company_structure.pdf
motor_control.png
chapter_01.tex
\end{verbatim}

Spaces and unnecessary special characters in file names should be
avoided where possible.

\vspace{0.3cm}

% ================================================================
% XXI. DOCUMENT COMPILATION AND FINAL CHECK
% ================================================================

\noindent
\textbf{XXI.\quad DOCUMENT COMPILATION AND FINAL CHECK}

\vspace{0.15cm}

Before submission, the complete \LaTeX{} project shall compile
successfully and produce the final PDF report.

The final document shall not contain unresolved references such as:

\begin{verbatim}
??
[?]
\end{verbatim}

Before generating the final PDF, students shall check:

\begin{itemize}[leftmargin=0.8cm, itemsep=2pt]
    \item the Table of Contents;
    \item page numbering;
    \item chapter, section, and subsection numbering;
    \item figure numbering;
    \item table numbering;
    \item equation numbering;
    \item cross-references;
    \item citations;
    \item bibliography;
    \item font display;
    \item mathematical symbols;
    \item image quality;
    \item completeness of all required forms.
\end{itemize}

The document may be compiled using TeX Live, MiKTeX, or another
compatible \LaTeX{} environment. Regardless of the compilation
environment, the resulting PDF shall conform to the official Faculty
template.

\vspace{0.3cm}

% ================================================================
% XXII. CONTENT THAT STUDENTS MAY EDIT
% ================================================================

\noindent
\textbf{XXII.\quad CONTENT THAT STUDENTS MAY EDIT}

\vspace{0.15cm}

Students may edit or replace:

\begin{itemize}[leftmargin=0.8cm, itemsep=2pt]
    \item student name;
    \item class;
    \item internship organization;
    \item supervisor name;
    \item Abstract;
    \item List of Symbols and Abbreviations;
    \item chapter titles and chapter content;
    \item figures and photographs;
    \item tables;
    \item equations;
    \item citations and references.
\end{itemize}

Students shall not modify the main document formatting settings in the
preamble, including the document class, page geometry, line spacing,
heading formatting, table-of-contents formatting, font configuration,
page layout, or cover-page macros unless specifically instructed.

\vspace{0.3cm}

% ================================================================
% XXIII. FINAL SUBMISSION REQUIREMENTS
% ================================================================

\noindent
\textbf{XXIII.\quad FINAL SUBMISSION REQUIREMENTS}

\vspace{0.15cm}

Before submission, the final report shall:

\begin{itemize}[leftmargin=0.8cm, itemsep=2pt]
    \item compile without errors that affect the report content or presentation;
    \item contain all required forms and report sections;
    \item contain correct student, class, internship organization, and supervisor information;
    \item contain an updated Table of Contents, List of Figures, and List of Tables;
    \item contain no unused example content from the template;
    \item contain no placeholder text such as \textit{Title of Chapter 1},
    \textit{Section 1.1}, or other sample entries;
    \item contain all required citations and references;
    \item contain no unresolved cross-references.
\end{itemize}

If the Faculty requires submission of the \LaTeX{} source, the complete
project shall include all \verb|.tex| files, the bibliography database,
figures, and any other files required to reproduce the submitted PDF.

\vspace{0.4cm}

\noindent
\textbf{Important:}
Before submitting the final internship report, students shall remove or
comment out the following command in the main document:

\begin{verbatim}
\input{Regulations.tex}
\end{verbatim}

For example:

\begin{verbatim}
% \input{Regulations.tex}
\end{verbatim}

This removes the guidance pages from the final submitted report while
preserving the regulations file in the \LaTeX{} project for reference.

} % End 13pt content

\clearpage

\end{verbatim}

For example:

\begin{verbatim}
% 
% =================================================================
% REGULATIONS FOR PREPARING THE GRADUATION INTERNSHIP REPORT
% USING LATEX
% =================================================================
%
% This file is intended to be included from the main document with:
%
%     \input{Regulations.tex}
%
% IMPORTANT:
% This section is provided for guidance only. Students shall remove
% or comment out the \input{Regulations.tex} command before submitting
% the final internship report.
%
% =================================================================

\clearpage
\thispagestyle{empty}

\begin{center}
    {\fontsize{16pt}{20pt}\selectfont\bfseries
    SPECIFIC REGULATIONS FOR PREPARING\\
    GRADUATION INTERNSHIP REPORTS USING \LaTeX}
\end{center}

\vspace{0.5cm}

{\fontsize{13pt}{17pt}\selectfont
\setstretch{1.3}
\justifying

\noindent
\textbf{Note:}
This section is provided for guidance only and is included in the main
document using the command \verb|\input{Regulations.tex}|.
Students shall remove or comment out this command before submitting the
final internship report.

\vspace{0.4cm}

% ================================================================
% I. GENERAL REQUIREMENTS
% ================================================================

\noindent
\textbf{I.\quad GENERAL REQUIREMENTS}

\vspace{0.15cm}

The Graduation Internship Report shall be prepared using the official
\LaTeX{} template provided by the Faculty. Students shall preserve the
predefined page layout, typography, line spacing, heading styles,
numbering system, figure and table formatting, and bibliography settings
unless otherwise instructed by the Faculty or the supervisor.

The report shall be written in English.

Students are responsible for the academic and technical content of the
report. The official \LaTeX{} template is responsible for standardized
document formatting.

\vspace{0.3cm}

% ================================================================
% II. REPORT STRUCTURE
% ================================================================

\noindent
\textbf{II.\quad REPORT STRUCTURE}

\vspace{0.15cm}

The report shall normally be arranged in the following order:

\begin{enumerate}[label=\arabic*., leftmargin=1.2cm, itemsep=2pt]
    \item Outer Cover;
    \item Comments of the Internship Host Organization;
    \item Supervisor's Evaluation and Graduation Internship Defense Results;
    \item Table of Contents;
    \item List of Symbols and Abbreviations, if applicable;
    \item List of Figures;
    \item List of Tables;
    \item Abstract;
    \item Main Chapters of the Internship Report;
    \item References;
    \item Appendices, if applicable.
\end{enumerate}

The preliminary pages shall not be numbered as chapters. The main report
content shall be organized into consecutively numbered chapters beginning
with Chapter 1.

\vspace{0.3cm}

% ================================================================
% III. PAGE SIZE AND MARGINS
% ================================================================

\noindent
\textbf{III.\quad PAGE SIZE AND MARGINS}

\vspace{0.15cm}

The report shall be formatted on A4 paper with the following margins:

\begin{center}
\begin{tabular}{ll}
    Left margin:   & 3.5 cm \\
    Right margin:  & 2.0 cm \\
    Top margin:    & 2.5 cm \\
    Bottom margin: & 2.5 cm
\end{tabular}
\end{center}

These settings are predefined in the official template and shall not be
modified by students.

\vspace{0.3cm}

% ================================================================
% IV. FONT, FONT SIZE, AND LINE SPACING
% ================================================================

\noindent
\textbf{IV.\quad FONT, FONT SIZE, AND LINE SPACING}

\vspace{0.15cm}

The main text shall use a 13 pt font size with 1.3-line spacing.
A Times-family typeface shall be used in accordance with the official
template.

The text shall be fully justified.

Students shall not use repeated spaces, repeated manual line breaks, or
multiple blank lines to position text. Paragraph and heading spacing shall
be controlled by the template.

\vspace{0.3cm}

% ================================================================
% V. OUTER COVER
% ================================================================

\noindent
\textbf{V.\quad OUTER COVER}

\vspace{0.15cm}

The outer cover shall follow the official format provided in the template
and shall contain the following information:

\begin{itemize}[leftmargin=0.8cm, itemsep=2pt]
    \item Thai Nguyen University of Technology;
    \item Faculty of International Training;
    \item Engineering and Technology Department;
    \item University logo and Faculty logo;
    \item the title \textbf{INTERNSHIP REPORT};
    \item Internship Organization;
    \item Author;
    \item Class;
    \item Supervisor;
    \item place and year of completion.
\end{itemize}

Students shall only replace the required information fields and shall not
modify the cover-page layout unless instructed otherwise.

\vspace{0.3cm}

% ================================================================
% VI. COMMENTS OF THE INTERNSHIP HOST ORGANIZATION
% ================================================================

\noindent
\textbf{VI.\quad COMMENTS OF THE INTERNSHIP HOST ORGANIZATION}

\vspace{0.15cm}

The page entitled \textbf{COMMENTS OF THE INTERNSHIP HOST ORGANIZATION}
is reserved for the internship host organization.

It includes:

\begin{itemize}[leftmargin=0.8cm, itemsep=2pt]
    \item the student's full name;
    \item class;
    \item comments from the host organization;
    \item score;
    \item date;
    \item signature and official seal.
\end{itemize}

Students shall not alter or complete official evaluation fields unless
specifically authorized.

\vspace{0.3cm}

% ================================================================
% VII. SUPERVISOR'S EVALUATION AND DEFENSE RESULTS
% ================================================================

\noindent
\textbf{VII.\quad SUPERVISOR'S EVALUATION AND DEFENSE RESULTS}

\vspace{0.15cm}

The section entitled
\textbf{SUPERVISOR'S EVALUATION AND GRADUATION INTERNSHIP DEFENSE RESULTS}
is reserved for the supervisor and defense committee.

The evaluation may include:

\begin{itemize}[leftmargin=0.8cm, itemsep=2pt]
    \item Internship Progress and Student's Attitude;
    \item Report Content;
    \item Report Presentation;
    \item Other Comments;
    \item Supervisor's Opinion;
    \item Defense Result;
    \item Score.
\end{itemize}

Students shall not enter evaluation scores or defense results before they
are officially determined.

\vspace{0.3cm}

% ================================================================
% VIII. CHAPTERS, SECTIONS, AND SUBSECTIONS
% ================================================================

\noindent
\textbf{VIII.\quad CHAPTERS, SECTIONS, AND SUBSECTIONS}

\vspace{0.15cm}

The report structure shall be generated using the standard \LaTeX{}
commands:

\begin{verbatim}
\chapter{...}
\section{...}
\subsection{...}
\end{verbatim}

Chapter, section, and subsection numbers shall be generated automatically
by \LaTeX{}.

Students shall not manually type chapter or section numbers and shall not
reproduce heading styles using commands such as \verb|\textbf|,
\verb|\fontsize|, or \verb|\centering| instead of the appropriate
structural command.

Excessive subdivision should be avoided. In most cases, the hierarchy
Chapter--Section--Subsection is sufficient.

\vspace{0.3cm}

% ================================================================
% IX. ORGANIZATION OF THE MAIN CONTENT
% ================================================================

\noindent
\textbf{IX.\quad ORGANIZATION OF THE MAIN CONTENT}

\vspace{0.15cm}

The specific content of the report shall depend on the internship
organization and the assigned tasks. A typical structure may include:

\begin{itemize}[leftmargin=0.8cm, itemsep=2pt]
    \item \textbf{Chapter 1. Overview of the Internship Organization};
    \item \textbf{Chapter 2. Internship Activities and Assigned Tasks};
    \item \textbf{Chapter 3. Results, Experience, and Evaluation}.
\end{itemize}

These chapter titles are recommendations only and may be adapted to the
actual internship content.

The report shall reflect work actually performed, observed, analyzed, or
learned during the internship. Unrelated theoretical material should not
be included merely to increase report length.

\vspace{0.3cm}

% ================================================================
% X. TABLE OF CONTENTS
% ================================================================

\noindent
\textbf{X.\quad TABLE OF CONTENTS}

\vspace{0.15cm}

The Table of Contents shall be generated automatically using:

\begin{verbatim}
\tableofcontents
\end{verbatim}

Students shall not create or manually edit page numbers in the Table of
Contents.

The document shall be compiled sufficiently to ensure that all page
numbers and entries are correctly updated.

\vspace{0.3cm}

% ================================================================
% XI. LIST OF SYMBOLS AND ABBREVIATIONS
% ================================================================

\noindent
\textbf{XI.\quad LIST OF SYMBOLS AND ABBREVIATIONS}

\vspace{0.15cm}

Symbols and abbreviations that are used repeatedly in the report shall be
listed in the \textbf{List of Symbols and Abbreviations}.

The list may contain the following columns:

\begin{center}
\begin{tabular}{llll}
    No. & Symbol / Abbreviation & Unit & Definition
\end{tabular}
\end{center}

Only symbols and abbreviations actually used in the final report shall be
retained. Example entries contained in the template shall be removed if
they are not relevant to the student's report.

Each abbreviation should normally be defined when it first appears in
the main text.

\vspace{0.3cm}

% ================================================================
% XII. FIGURES, DIAGRAMS, AND PHOTOGRAPHS
% ================================================================

\noindent
\textbf{XII.\quad FIGURES, DIAGRAMS, AND PHOTOGRAPHS}

\vspace{0.15cm}

Figures shall be clear, legible, and directly relevant to the report
content.

Figures shall be inserted using the standard \verb|figure| environment.

Example:

\begin{verbatim}
\begin{figure}[H]
    \centering
    \includegraphics[width=0.8\textwidth]
                    {figures/company_structure.pdf}
    \caption{Organizational structure of the company}
    \label{fig:company-structure}
\end{figure}
\end{verbatim}

Figure captions shall be placed below the figure.

Each referenced figure shall have a unique \verb|\label| and shall be
referenced in the text using, for example:

\begin{verbatim}
Figure~\ref{fig:company-structure}
\end{verbatim}

Students shall not manually type figure numbers.

Figures reproduced or adapted from external sources shall include an
appropriate citation.

Photographs taken during the internship shall be relevant and sufficiently
clear. Students shall not include confidential, restricted, or proprietary
information without permission from the internship organization.

Vector graphics or high-resolution images are recommended whenever
possible.

\vspace{0.3cm}

% ================================================================
% XIII. TABLES
% ================================================================

\noindent
\textbf{XIII.\quad TABLES}

\vspace{0.15cm}

Tables shall be created using suitable \LaTeX{} environments such as
\verb|table|, \verb|tabular|, \verb|tabularx|, or \verb|longtable|.

Table captions shall be placed above the table.

Example:

\begin{verbatim}
\begin{table}[H]
    \centering
    \caption{Internship schedule}
    \label{tab:internship-schedule}

    \begin{tabular}{lll}
        \toprule
        Week & Task & Result \\
        \midrule
        1 & Orientation & Completed \\
        2 & System study & Completed \\
        \bottomrule
    \end{tabular}
\end{table}
\end{verbatim}

Tables shall be referenced using:

\begin{verbatim}
Table~\ref{tab:internship-schedule}
\end{verbatim}

Students shall not manually type table numbers.

Tables should be created directly in \LaTeX{} rather than inserted as
screenshots whenever practical.

\vspace{0.3cm}

% ================================================================
% XIV. EQUATIONS AND TECHNICAL QUANTITIES
% ================================================================

\noindent
\textbf{XIV.\quad EQUATIONS AND TECHNICAL QUANTITIES}

\vspace{0.15cm}

Mathematical expressions shall be typeset directly using \LaTeX{}.
Images or screenshots of equations shall not be used.

Short expressions may be written in inline mathematical mode, for example:

\begin{verbatim}
$R = 0.7\,\Omega$
\end{verbatim}

Important equations shall use a numbered equation environment.

Example:

\begin{verbatim}
\begin{equation}
    P = UI
    \label{eq:power}
\end{equation}
\end{verbatim}

Equations shall be referenced using:

\begin{verbatim}
Equation~\eqref{eq:power}
\end{verbatim}

Equation numbers shall not be manually entered.

SI units shall be used wherever applicable. The \verb|siunitx| package
should be used where appropriate.

Examples:

\begin{verbatim}
\SI{220}{\volt}
\SI{50}{\hertz}
\SI{2.5}{\ampere}
\end{verbatim}

Notation and units shall be consistent throughout the report.

\vspace{0.3cm}

% ================================================================
% XV. SOURCE CODE
% ================================================================

\noindent
\textbf{XV.\quad SOURCE CODE}

\vspace{0.15cm}

If the internship involves programming, PLCs, microcontrollers,
simulation, or software development, relevant source code may be included
in the report.

Source code shall be presented using the \verb|listings| environment or
another environment defined by the official template.

Example:

\begin{verbatim}
\begin{lstlisting}[language=st]
IF Start = TRUE THEN
    Motor := TRUE;
ELSE
    Motor := FALSE;
END_IF;
\end{lstlisting}
\end{verbatim}

Only code necessary for explaining the work should appear in the main
report. Long source-code listings should be moved to an appendix or
submitted as supplementary files if required.

Source code should not normally be inserted as screenshots.

\vspace{0.3cm}

% ================================================================
% XVI. ABSTRACT
% ================================================================

\noindent
\textbf{XVI.\quad ABSTRACT}

\vspace{0.15cm}

The Abstract shall briefly summarize:

\begin{itemize}[leftmargin=0.8cm, itemsep=2pt]
    \item the internship organization;
    \item the objectives of the internship;
    \item the main tasks performed;
    \item the main results, knowledge, and skills obtained.
\end{itemize}

The Abstract should be concise and should not reproduce the full content
of the report chapter by chapter.

\vspace{0.3cm}

% ================================================================
% XVII. REFERENCES AND CITATIONS
% ================================================================

\noindent
\textbf{XVII.\quad REFERENCES AND CITATIONS}

\vspace{0.15cm}

All books, journal papers, conference papers, standards, websites,
datasheets, manuals, technical documents, figures, tables, or other
external sources used in preparing the report shall be properly cited.

References shall be managed through the bibliography system provided with
the official template.

Bibliographic entries shall normally be stored in:

\begin{verbatim}
References.bib
\end{verbatim}

References shall be cited in the report using:

\begin{verbatim}
\cite{reference-key}
\end{verbatim}

Students shall not manually type reference numbers such as ``[1]'',
``[2]'', etc.

Titles of foreign-language publications shall remain in their original
language in the bibliography.

All borrowed ideas, data, figures, tables, equations, methods, or results
shall be properly acknowledged.

\vspace{0.3cm}

% ================================================================
% XVIII. CROSS-REFERENCING
% ================================================================

\noindent
\textbf{XVIII.\quad CROSS-REFERENCING}

\vspace{0.15cm}

Numbered elements shall be cross-referenced using the standard \LaTeX{}
mechanism.

Use the following commands where appropriate:

\begin{verbatim}
\label{...}
\ref{...}
\eqref{...}
\end{verbatim}

For example:

\begin{verbatim}
As shown in Figure~\ref{fig:system}, ...
\end{verbatim}

Students shall not manually type numbers such as ``Figure 2.3'' when
those numbers can be generated automatically.

\vspace{0.3cm}

% ================================================================
% XIX. PAGE NUMBERING
% ================================================================

\noindent
\textbf{XIX.\quad PAGE NUMBERING}

\vspace{0.15cm}

Cover pages and official evaluation forms at the beginning of the report
shall not display page numbers.

The main body of the report shall use Arabic page numbers beginning from
Chapter 1.

Page numbering shall be generated automatically by the official template.
Students shall not manually enter page numbers.

\vspace{0.3cm}

% ================================================================
% XX. PROJECT FILE ORGANIZATION
% ================================================================

\noindent
\textbf{XX.\quad PROJECT FILE ORGANIZATION}

\vspace{0.15cm}

Students should maintain a clear and organized \LaTeX{} project
structure. A recommended structure is:

\begin{verbatim}
InternshipReport/
|
|-- main.tex
|-- References.bib
|
|-- figures/
|   |-- company.pdf
|   |-- system.png
|   `-- ...
|
`-- chapters/
    |-- chapter1.tex
    |-- chapter2.tex
    `-- chapter3.tex
\end{verbatim}

Figures should normally be stored in the \verb|figures| directory
specified by the official template.

File and directory names should preferably contain Latin letters,
numbers, hyphens, or underscores.

Recommended examples include:

\begin{verbatim}
company_structure.pdf
motor_control.png
chapter_01.tex
\end{verbatim}

Spaces and unnecessary special characters in file names should be
avoided where possible.

\vspace{0.3cm}

% ================================================================
% XXI. DOCUMENT COMPILATION AND FINAL CHECK
% ================================================================

\noindent
\textbf{XXI.\quad DOCUMENT COMPILATION AND FINAL CHECK}

\vspace{0.15cm}

Before submission, the complete \LaTeX{} project shall compile
successfully and produce the final PDF report.

The final document shall not contain unresolved references such as:

\begin{verbatim}
??
[?]
\end{verbatim}

Before generating the final PDF, students shall check:

\begin{itemize}[leftmargin=0.8cm, itemsep=2pt]
    \item the Table of Contents;
    \item page numbering;
    \item chapter, section, and subsection numbering;
    \item figure numbering;
    \item table numbering;
    \item equation numbering;
    \item cross-references;
    \item citations;
    \item bibliography;
    \item font display;
    \item mathematical symbols;
    \item image quality;
    \item completeness of all required forms.
\end{itemize}

The document may be compiled using TeX Live, MiKTeX, or another
compatible \LaTeX{} environment. Regardless of the compilation
environment, the resulting PDF shall conform to the official Faculty
template.

\vspace{0.3cm}

% ================================================================
% XXII. CONTENT THAT STUDENTS MAY EDIT
% ================================================================

\noindent
\textbf{XXII.\quad CONTENT THAT STUDENTS MAY EDIT}

\vspace{0.15cm}

Students may edit or replace:

\begin{itemize}[leftmargin=0.8cm, itemsep=2pt]
    \item student name;
    \item class;
    \item internship organization;
    \item supervisor name;
    \item Abstract;
    \item List of Symbols and Abbreviations;
    \item chapter titles and chapter content;
    \item figures and photographs;
    \item tables;
    \item equations;
    \item citations and references.
\end{itemize}

Students shall not modify the main document formatting settings in the
preamble, including the document class, page geometry, line spacing,
heading formatting, table-of-contents formatting, font configuration,
page layout, or cover-page macros unless specifically instructed.

\vspace{0.3cm}

% ================================================================
% XXIII. FINAL SUBMISSION REQUIREMENTS
% ================================================================

\noindent
\textbf{XXIII.\quad FINAL SUBMISSION REQUIREMENTS}

\vspace{0.15cm}

Before submission, the final report shall:

\begin{itemize}[leftmargin=0.8cm, itemsep=2pt]
    \item compile without errors that affect the report content or presentation;
    \item contain all required forms and report sections;
    \item contain correct student, class, internship organization, and supervisor information;
    \item contain an updated Table of Contents, List of Figures, and List of Tables;
    \item contain no unused example content from the template;
    \item contain no placeholder text such as \textit{Title of Chapter 1},
    \textit{Section 1.1}, or other sample entries;
    \item contain all required citations and references;
    \item contain no unresolved cross-references.
\end{itemize}

If the Faculty requires submission of the \LaTeX{} source, the complete
project shall include all \verb|.tex| files, the bibliography database,
figures, and any other files required to reproduce the submitted PDF.

\vspace{0.4cm}

\noindent
\textbf{Important:}
Before submitting the final internship report, students shall remove or
comment out the following command in the main document:

\begin{verbatim}
\input{Regulations.tex}
\end{verbatim}

For example:

\begin{verbatim}
% \input{Regulations.tex}
\end{verbatim}

This removes the guidance pages from the final submitted report while
preserving the regulations file in the \LaTeX{} project for reference.

} % End 13pt content

\clearpage

\end{verbatim}

This removes the guidance pages from the final submitted report while
preserving the regulations file in the \LaTeX{} project for reference.

} % End 13pt content

\clearpage
|.
Students shall remove or comment out this command before submitting the
final internship report.

\vspace{0.4cm}

% ================================================================
% I. GENERAL REQUIREMENTS
% ================================================================

\noindent
\textbf{I.\quad GENERAL REQUIREMENTS}

\vspace{0.15cm}

The Graduation Internship Report shall be prepared using the official
\LaTeX{} template provided by the Faculty. Students shall preserve the
predefined page layout, typography, line spacing, heading styles,
numbering system, figure and table formatting, and bibliography settings
unless otherwise instructed by the Faculty or the supervisor.

The report shall be written in English.

Students are responsible for the academic and technical content of the
report. The official \LaTeX{} template is responsible for standardized
document formatting.

\vspace{0.3cm}

% ================================================================
% II. REPORT STRUCTURE
% ================================================================

\noindent
\textbf{II.\quad REPORT STRUCTURE}

\vspace{0.15cm}

The report shall normally be arranged in the following order:

\begin{enumerate}[label=\arabic*., leftmargin=1.2cm, itemsep=2pt]
    \item Outer Cover;
    \item Comments of the Internship Host Organization;
    \item Supervisor's Evaluation and Graduation Internship Defense Results;
    \item Table of Contents;
    \item List of Symbols and Abbreviations, if applicable;
    \item List of Figures;
    \item List of Tables;
    \item Abstract;
    \item Main Chapters of the Internship Report;
    \item References;
    \item Appendices, if applicable.
\end{enumerate}

The preliminary pages shall not be numbered as chapters. The main report
content shall be organized into consecutively numbered chapters beginning
with Chapter 1.

\vspace{0.3cm}

% ================================================================
% III. PAGE SIZE AND MARGINS
% ================================================================

\noindent
\textbf{III.\quad PAGE SIZE AND MARGINS}

\vspace{0.15cm}

The report shall be formatted on A4 paper with the following margins:

\begin{center}
\begin{tabular}{ll}
    Left margin:   & 3.5 cm \\
    Right margin:  & 2.0 cm \\
    Top margin:    & 2.5 cm \\
    Bottom margin: & 2.5 cm
\end{tabular}
\end{center}

These settings are predefined in the official template and shall not be
modified by students.

\vspace{0.3cm}

% ================================================================
% IV. FONT, FONT SIZE, AND LINE SPACING
% ================================================================

\noindent
\textbf{IV.\quad FONT, FONT SIZE, AND LINE SPACING}

\vspace{0.15cm}

The main text shall use a 13 pt font size with 1.3-line spacing.
A Times-family typeface shall be used in accordance with the official
template.

The text shall be fully justified.

Students shall not use repeated spaces, repeated manual line breaks, or
multiple blank lines to position text. Paragraph and heading spacing shall
be controlled by the template.

\vspace{0.3cm}

% ================================================================
% V. OUTER COVER
% ================================================================

\noindent
\textbf{V.\quad OUTER COVER}

\vspace{0.15cm}

The outer cover shall follow the official format provided in the template
and shall contain the following information:

\begin{itemize}[leftmargin=0.8cm, itemsep=2pt]
    \item Thai Nguyen University of Technology;
    \item Faculty of International Training;
    \item Engineering and Technology Department;
    \item University logo and Faculty logo;
    \item the title \textbf{INTERNSHIP REPORT};
    \item Internship Organization;
    \item Author;
    \item Class;
    \item Supervisor;
    \item place and year of completion.
\end{itemize}

Students shall only replace the required information fields and shall not
modify the cover-page layout unless instructed otherwise.

\vspace{0.3cm}

% ================================================================
% VI. COMMENTS OF THE INTERNSHIP HOST ORGANIZATION
% ================================================================

\noindent
\textbf{VI.\quad COMMENTS OF THE INTERNSHIP HOST ORGANIZATION}

\vspace{0.15cm}

The page entitled \textbf{COMMENTS OF THE INTERNSHIP HOST ORGANIZATION}
is reserved for the internship host organization.

It includes:

\begin{itemize}[leftmargin=0.8cm, itemsep=2pt]
    \item the student's full name;
    \item class;
    \item comments from the host organization;
    \item score;
    \item date;
    \item signature and official seal.
\end{itemize}

Students shall not alter or complete official evaluation fields unless
specifically authorized.

\vspace{0.3cm}

% ================================================================
% VII. SUPERVISOR'S EVALUATION AND DEFENSE RESULTS
% ================================================================

\noindent
\textbf{VII.\quad SUPERVISOR'S EVALUATION AND DEFENSE RESULTS}

\vspace{0.15cm}

The section entitled
\textbf{SUPERVISOR'S EVALUATION AND GRADUATION INTERNSHIP DEFENSE RESULTS}
is reserved for the supervisor and defense committee.

The evaluation may include:

\begin{itemize}[leftmargin=0.8cm, itemsep=2pt]
    \item Internship Progress and Student's Attitude;
    \item Report Content;
    \item Report Presentation;
    \item Other Comments;
    \item Supervisor's Opinion;
    \item Defense Result;
    \item Score.
\end{itemize}

Students shall not enter evaluation scores or defense results before they
are officially determined.

\vspace{0.3cm}

% ================================================================
% VIII. CHAPTERS, SECTIONS, AND SUBSECTIONS
% ================================================================

\noindent
\textbf{VIII.\quad CHAPTERS, SECTIONS, AND SUBSECTIONS}

\vspace{0.15cm}

The report structure shall be generated using the standard \LaTeX{}
commands:

\begin{verbatim}
\chapter{...}
\section{...}
\subsection{...}
\end{verbatim}

Chapter, section, and subsection numbers shall be generated automatically
by \LaTeX{}.

Students shall not manually type chapter or section numbers and shall not
reproduce heading styles using commands such as \verb|\textbf|,
\verb|\fontsize|, or \verb|\centering| instead of the appropriate
structural command.

Excessive subdivision should be avoided. In most cases, the hierarchy
Chapter--Section--Subsection is sufficient.

\vspace{0.3cm}

% ================================================================
% IX. ORGANIZATION OF THE MAIN CONTENT
% ================================================================

\noindent
\textbf{IX.\quad ORGANIZATION OF THE MAIN CONTENT}

\vspace{0.15cm}

The specific content of the report shall depend on the internship
organization and the assigned tasks. A typical structure may include:

\begin{itemize}[leftmargin=0.8cm, itemsep=2pt]
    \item \textbf{Chapter 1. Overview of the Internship Organization};
    \item \textbf{Chapter 2. Internship Activities and Assigned Tasks};
    \item \textbf{Chapter 3. Results, Experience, and Evaluation}.
\end{itemize}

These chapter titles are recommendations only and may be adapted to the
actual internship content.

The report shall reflect work actually performed, observed, analyzed, or
learned during the internship. Unrelated theoretical material should not
be included merely to increase report length.

\vspace{0.3cm}

% ================================================================
% X. TABLE OF CONTENTS
% ================================================================

\noindent
\textbf{X.\quad TABLE OF CONTENTS}

\vspace{0.15cm}

The Table of Contents shall be generated automatically using:

\begin{verbatim}
\tableofcontents
\end{verbatim}

Students shall not create or manually edit page numbers in the Table of
Contents.

The document shall be compiled sufficiently to ensure that all page
numbers and entries are correctly updated.

\vspace{0.3cm}

% ================================================================
% XI. LIST OF SYMBOLS AND ABBREVIATIONS
% ================================================================

\noindent
\textbf{XI.\quad LIST OF SYMBOLS AND ABBREVIATIONS}

\vspace{0.15cm}

Symbols and abbreviations that are used repeatedly in the report shall be
listed in the \textbf{List of Symbols and Abbreviations}.

The list may contain the following columns:

\begin{center}
\begin{tabular}{llll}
    No. & Symbol / Abbreviation & Unit & Definition
\end{tabular}
\end{center}

Only symbols and abbreviations actually used in the final report shall be
retained. Example entries contained in the template shall be removed if
they are not relevant to the student's report.

Each abbreviation should normally be defined when it first appears in
the main text.

\vspace{0.3cm}

% ================================================================
% XII. FIGURES, DIAGRAMS, AND PHOTOGRAPHS
% ================================================================

\noindent
\textbf{XII.\quad FIGURES, DIAGRAMS, AND PHOTOGRAPHS}

\vspace{0.15cm}

Figures shall be clear, legible, and directly relevant to the report
content.

Figures shall be inserted using the standard \verb|figure| environment.

Example:

\begin{verbatim}
\begin{figure}[H]
    \centering
    \includegraphics[width=0.8\textwidth]
                    {figures/company_structure.pdf}
    \caption{Organizational structure of the company}
    \label{fig:company-structure}
\end{figure}
\end{verbatim}

Figure captions shall be placed below the figure.

Each referenced figure shall have a unique \verb|\label| and shall be
referenced in the text using, for example:

\begin{verbatim}
Figure~\ref{fig:company-structure}
\end{verbatim}

Students shall not manually type figure numbers.

Figures reproduced or adapted from external sources shall include an
appropriate citation.

Photographs taken during the internship shall be relevant and sufficiently
clear. Students shall not include confidential, restricted, or proprietary
information without permission from the internship organization.

Vector graphics or high-resolution images are recommended whenever
possible.

\vspace{0.3cm}

% ================================================================
% XIII. TABLES
% ================================================================

\noindent
\textbf{XIII.\quad TABLES}

\vspace{0.15cm}

Tables shall be created using suitable \LaTeX{} environments such as
\verb|table|, \verb|tabular|, \verb|tabularx|, or \verb|longtable|.

Table captions shall be placed above the table.

Example:

\begin{verbatim}
\begin{table}[H]
    \centering
    \caption{Internship schedule}
    \label{tab:internship-schedule}

    \begin{tabular}{lll}
        \toprule
        Week & Task & Result \\
        \midrule
        1 & Orientation & Completed \\
        2 & System study & Completed \\
        \bottomrule
    \end{tabular}
\end{table}
\end{verbatim}

Tables shall be referenced using:

\begin{verbatim}
Table~\ref{tab:internship-schedule}
\end{verbatim}

Students shall not manually type table numbers.

Tables should be created directly in \LaTeX{} rather than inserted as
screenshots whenever practical.

\vspace{0.3cm}

% ================================================================
% XIV. EQUATIONS AND TECHNICAL QUANTITIES
% ================================================================

\noindent
\textbf{XIV.\quad EQUATIONS AND TECHNICAL QUANTITIES}

\vspace{0.15cm}

Mathematical expressions shall be typeset directly using \LaTeX{}.
Images or screenshots of equations shall not be used.

Short expressions may be written in inline mathematical mode, for example:

\begin{verbatim}
$R = 0.7\,\Omega$
\end{verbatim}

Important equations shall use a numbered equation environment.

Example:

\begin{verbatim}
\begin{equation}
    P = UI
    \label{eq:power}
\end{equation}
\end{verbatim}

Equations shall be referenced using:

\begin{verbatim}
Equation~\eqref{eq:power}
\end{verbatim}

Equation numbers shall not be manually entered.

SI units shall be used wherever applicable. The \verb|siunitx| package
should be used where appropriate.

Examples:

\begin{verbatim}
\SI{220}{\volt}
\SI{50}{\hertz}
\SI{2.5}{\ampere}
\end{verbatim}

Notation and units shall be consistent throughout the report.

\vspace{0.3cm}

% ================================================================
% XV. SOURCE CODE
% ================================================================

\noindent
\textbf{XV.\quad SOURCE CODE}

\vspace{0.15cm}

If the internship involves programming, PLCs, microcontrollers,
simulation, or software development, relevant source code may be included
in the report.

Source code shall be presented using the \verb|listings| environment or
another environment defined by the official template.

Example:

\begin{verbatim}
\begin{lstlisting}[language=st]
IF Start = TRUE THEN
    Motor := TRUE;
ELSE
    Motor := FALSE;
END_IF;
\end{lstlisting}
\end{verbatim}

Only code necessary for explaining the work should appear in the main
report. Long source-code listings should be moved to an appendix or
submitted as supplementary files if required.

Source code should not normally be inserted as screenshots.

\vspace{0.3cm}

% ================================================================
% XVI. ABSTRACT
% ================================================================

\noindent
\textbf{XVI.\quad ABSTRACT}

\vspace{0.15cm}

The Abstract shall briefly summarize:

\begin{itemize}[leftmargin=0.8cm, itemsep=2pt]
    \item the internship organization;
    \item the objectives of the internship;
    \item the main tasks performed;
    \item the main results, knowledge, and skills obtained.
\end{itemize}

The Abstract should be concise and should not reproduce the full content
of the report chapter by chapter.

\vspace{0.3cm}

% ================================================================
% XVII. REFERENCES AND CITATIONS
% ================================================================

\noindent
\textbf{XVII.\quad REFERENCES AND CITATIONS}

\vspace{0.15cm}

All books, journal papers, conference papers, standards, websites,
datasheets, manuals, technical documents, figures, tables, or other
external sources used in preparing the report shall be properly cited.

References shall be managed through the bibliography system provided with
the official template.

Bibliographic entries shall normally be stored in:

\begin{verbatim}
References.bib
\end{verbatim}

References shall be cited in the report using:

\begin{verbatim}
\cite{reference-key}
\end{verbatim}

Students shall not manually type reference numbers such as ``[1]'',
``[2]'', etc.

Titles of foreign-language publications shall remain in their original
language in the bibliography.

All borrowed ideas, data, figures, tables, equations, methods, or results
shall be properly acknowledged.

\vspace{0.3cm}

% ================================================================
% XVIII. CROSS-REFERENCING
% ================================================================

\noindent
\textbf{XVIII.\quad CROSS-REFERENCING}

\vspace{0.15cm}

Numbered elements shall be cross-referenced using the standard \LaTeX{}
mechanism.

Use the following commands where appropriate:

\begin{verbatim}
\label{...}
\ref{...}
\eqref{...}
\end{verbatim}

For example:

\begin{verbatim}
As shown in Figure~\ref{fig:system}, ...
\end{verbatim}

Students shall not manually type numbers such as ``Figure 2.3'' when
those numbers can be generated automatically.

\vspace{0.3cm}

% ================================================================
% XIX. PAGE NUMBERING
% ================================================================

\noindent
\textbf{XIX.\quad PAGE NUMBERING}

\vspace{0.15cm}

Cover pages and official evaluation forms at the beginning of the report
shall not display page numbers.

The main body of the report shall use Arabic page numbers beginning from
Chapter 1.

Page numbering shall be generated automatically by the official template.
Students shall not manually enter page numbers.

\vspace{0.3cm}

% ================================================================
% XX. PROJECT FILE ORGANIZATION
% ================================================================

\noindent
\textbf{XX.\quad PROJECT FILE ORGANIZATION}

\vspace{0.15cm}

Students should maintain a clear and organized \LaTeX{} project
structure. A recommended structure is:

\begin{verbatim}
InternshipReport/
|
|-- main.tex
|-- References.bib
|
|-- figures/
|   |-- company.pdf
|   |-- system.png
|   `-- ...
|
`-- chapters/
    |-- chapter1.tex
    |-- chapter2.tex
    `-- chapter3.tex
\end{verbatim}

Figures should normally be stored in the \verb|figures| directory
specified by the official template.

File and directory names should preferably contain Latin letters,
numbers, hyphens, or underscores.

Recommended examples include:

\begin{verbatim}
company_structure.pdf
motor_control.png
chapter_01.tex
\end{verbatim}

Spaces and unnecessary special characters in file names should be
avoided where possible.

\vspace{0.3cm}

% ================================================================
% XXI. DOCUMENT COMPILATION AND FINAL CHECK
% ================================================================

\noindent
\textbf{XXI.\quad DOCUMENT COMPILATION AND FINAL CHECK}

\vspace{0.15cm}

Before submission, the complete \LaTeX{} project shall compile
successfully and produce the final PDF report.

The final document shall not contain unresolved references such as:

\begin{verbatim}
??
[?]
\end{verbatim}

Before generating the final PDF, students shall check:

\begin{itemize}[leftmargin=0.8cm, itemsep=2pt]
    \item the Table of Contents;
    \item page numbering;
    \item chapter, section, and subsection numbering;
    \item figure numbering;
    \item table numbering;
    \item equation numbering;
    \item cross-references;
    \item citations;
    \item bibliography;
    \item font display;
    \item mathematical symbols;
    \item image quality;
    \item completeness of all required forms.
\end{itemize}

The document may be compiled using TeX Live, MiKTeX, or another
compatible \LaTeX{} environment. Regardless of the compilation
environment, the resulting PDF shall conform to the official Faculty
template.

\vspace{0.3cm}

% ================================================================
% XXII. CONTENT THAT STUDENTS MAY EDIT
% ================================================================

\noindent
\textbf{XXII.\quad CONTENT THAT STUDENTS MAY EDIT}

\vspace{0.15cm}

Students may edit or replace:

\begin{itemize}[leftmargin=0.8cm, itemsep=2pt]
    \item student name;
    \item class;
    \item internship organization;
    \item supervisor name;
    \item Abstract;
    \item List of Symbols and Abbreviations;
    \item chapter titles and chapter content;
    \item figures and photographs;
    \item tables;
    \item equations;
    \item citations and references.
\end{itemize}

Students shall not modify the main document formatting settings in the
preamble, including the document class, page geometry, line spacing,
heading formatting, table-of-contents formatting, font configuration,
page layout, or cover-page macros unless specifically instructed.

\vspace{0.3cm}

% ================================================================
% XXIII. FINAL SUBMISSION REQUIREMENTS
% ================================================================

\noindent
\textbf{XXIII.\quad FINAL SUBMISSION REQUIREMENTS}

\vspace{0.15cm}

Before submission, the final report shall:

\begin{itemize}[leftmargin=0.8cm, itemsep=2pt]
    \item compile without errors that affect the report content or presentation;
    \item contain all required forms and report sections;
    \item contain correct student, class, internship organization, and supervisor information;
    \item contain an updated Table of Contents, List of Figures, and List of Tables;
    \item contain no unused example content from the template;
    \item contain no placeholder text such as \textit{Title of Chapter 1},
    \textit{Section 1.1}, or other sample entries;
    \item contain all required citations and references;
    \item contain no unresolved cross-references.
\end{itemize}

If the Faculty requires submission of the \LaTeX{} source, the complete
project shall include all \verb|.tex| files, the bibliography database,
figures, and any other files required to reproduce the submitted PDF.

\vspace{0.4cm}

\noindent
\textbf{Important:}
Before submitting the final internship report, students shall remove or
comment out the following command in the main document:

\begin{verbatim}

% =================================================================
% REGULATIONS FOR PREPARING THE GRADUATION INTERNSHIP REPORT
% USING LATEX
% =================================================================
%
% This file is intended to be included from the main document with:
%
%     
% =================================================================
% REGULATIONS FOR PREPARING THE GRADUATION INTERNSHIP REPORT
% USING LATEX
% =================================================================
%
% This file is intended to be included from the main document with:
%
%     \input{Regulations.tex}
%
% IMPORTANT:
% This section is provided for guidance only. Students shall remove
% or comment out the \input{Regulations.tex} command before submitting
% the final internship report.
%
% =================================================================

\clearpage
\thispagestyle{empty}

\begin{center}
    {\fontsize{16pt}{20pt}\selectfont\bfseries
    SPECIFIC REGULATIONS FOR PREPARING\\
    GRADUATION INTERNSHIP REPORTS USING \LaTeX}
\end{center}

\vspace{0.5cm}

{\fontsize{13pt}{17pt}\selectfont
\setstretch{1.3}
\justifying

\noindent
\textbf{Note:}
This section is provided for guidance only and is included in the main
document using the command \verb|\input{Regulations.tex}|.
Students shall remove or comment out this command before submitting the
final internship report.

\vspace{0.4cm}

% ================================================================
% I. GENERAL REQUIREMENTS
% ================================================================

\noindent
\textbf{I.\quad GENERAL REQUIREMENTS}

\vspace{0.15cm}

The Graduation Internship Report shall be prepared using the official
\LaTeX{} template provided by the Faculty. Students shall preserve the
predefined page layout, typography, line spacing, heading styles,
numbering system, figure and table formatting, and bibliography settings
unless otherwise instructed by the Faculty or the supervisor.

The report shall be written in English.

Students are responsible for the academic and technical content of the
report. The official \LaTeX{} template is responsible for standardized
document formatting.

\vspace{0.3cm}

% ================================================================
% II. REPORT STRUCTURE
% ================================================================

\noindent
\textbf{II.\quad REPORT STRUCTURE}

\vspace{0.15cm}

The report shall normally be arranged in the following order:

\begin{enumerate}[label=\arabic*., leftmargin=1.2cm, itemsep=2pt]
    \item Outer Cover;
    \item Comments of the Internship Host Organization;
    \item Supervisor's Evaluation and Graduation Internship Defense Results;
    \item Table of Contents;
    \item List of Symbols and Abbreviations, if applicable;
    \item List of Figures;
    \item List of Tables;
    \item Abstract;
    \item Main Chapters of the Internship Report;
    \item References;
    \item Appendices, if applicable.
\end{enumerate}

The preliminary pages shall not be numbered as chapters. The main report
content shall be organized into consecutively numbered chapters beginning
with Chapter 1.

\vspace{0.3cm}

% ================================================================
% III. PAGE SIZE AND MARGINS
% ================================================================

\noindent
\textbf{III.\quad PAGE SIZE AND MARGINS}

\vspace{0.15cm}

The report shall be formatted on A4 paper with the following margins:

\begin{center}
\begin{tabular}{ll}
    Left margin:   & 3.5 cm \\
    Right margin:  & 2.0 cm \\
    Top margin:    & 2.5 cm \\
    Bottom margin: & 2.5 cm
\end{tabular}
\end{center}

These settings are predefined in the official template and shall not be
modified by students.

\vspace{0.3cm}

% ================================================================
% IV. FONT, FONT SIZE, AND LINE SPACING
% ================================================================

\noindent
\textbf{IV.\quad FONT, FONT SIZE, AND LINE SPACING}

\vspace{0.15cm}

The main text shall use a 13 pt font size with 1.3-line spacing.
A Times-family typeface shall be used in accordance with the official
template.

The text shall be fully justified.

Students shall not use repeated spaces, repeated manual line breaks, or
multiple blank lines to position text. Paragraph and heading spacing shall
be controlled by the template.

\vspace{0.3cm}

% ================================================================
% V. OUTER COVER
% ================================================================

\noindent
\textbf{V.\quad OUTER COVER}

\vspace{0.15cm}

The outer cover shall follow the official format provided in the template
and shall contain the following information:

\begin{itemize}[leftmargin=0.8cm, itemsep=2pt]
    \item Thai Nguyen University of Technology;
    \item Faculty of International Training;
    \item Engineering and Technology Department;
    \item University logo and Faculty logo;
    \item the title \textbf{INTERNSHIP REPORT};
    \item Internship Organization;
    \item Author;
    \item Class;
    \item Supervisor;
    \item place and year of completion.
\end{itemize}

Students shall only replace the required information fields and shall not
modify the cover-page layout unless instructed otherwise.

\vspace{0.3cm}

% ================================================================
% VI. COMMENTS OF THE INTERNSHIP HOST ORGANIZATION
% ================================================================

\noindent
\textbf{VI.\quad COMMENTS OF THE INTERNSHIP HOST ORGANIZATION}

\vspace{0.15cm}

The page entitled \textbf{COMMENTS OF THE INTERNSHIP HOST ORGANIZATION}
is reserved for the internship host organization.

It includes:

\begin{itemize}[leftmargin=0.8cm, itemsep=2pt]
    \item the student's full name;
    \item class;
    \item comments from the host organization;
    \item score;
    \item date;
    \item signature and official seal.
\end{itemize}

Students shall not alter or complete official evaluation fields unless
specifically authorized.

\vspace{0.3cm}

% ================================================================
% VII. SUPERVISOR'S EVALUATION AND DEFENSE RESULTS
% ================================================================

\noindent
\textbf{VII.\quad SUPERVISOR'S EVALUATION AND DEFENSE RESULTS}

\vspace{0.15cm}

The section entitled
\textbf{SUPERVISOR'S EVALUATION AND GRADUATION INTERNSHIP DEFENSE RESULTS}
is reserved for the supervisor and defense committee.

The evaluation may include:

\begin{itemize}[leftmargin=0.8cm, itemsep=2pt]
    \item Internship Progress and Student's Attitude;
    \item Report Content;
    \item Report Presentation;
    \item Other Comments;
    \item Supervisor's Opinion;
    \item Defense Result;
    \item Score.
\end{itemize}

Students shall not enter evaluation scores or defense results before they
are officially determined.

\vspace{0.3cm}

% ================================================================
% VIII. CHAPTERS, SECTIONS, AND SUBSECTIONS
% ================================================================

\noindent
\textbf{VIII.\quad CHAPTERS, SECTIONS, AND SUBSECTIONS}

\vspace{0.15cm}

The report structure shall be generated using the standard \LaTeX{}
commands:

\begin{verbatim}
\chapter{...}
\section{...}
\subsection{...}
\end{verbatim}

Chapter, section, and subsection numbers shall be generated automatically
by \LaTeX{}.

Students shall not manually type chapter or section numbers and shall not
reproduce heading styles using commands such as \verb|\textbf|,
\verb|\fontsize|, or \verb|\centering| instead of the appropriate
structural command.

Excessive subdivision should be avoided. In most cases, the hierarchy
Chapter--Section--Subsection is sufficient.

\vspace{0.3cm}

% ================================================================
% IX. ORGANIZATION OF THE MAIN CONTENT
% ================================================================

\noindent
\textbf{IX.\quad ORGANIZATION OF THE MAIN CONTENT}

\vspace{0.15cm}

The specific content of the report shall depend on the internship
organization and the assigned tasks. A typical structure may include:

\begin{itemize}[leftmargin=0.8cm, itemsep=2pt]
    \item \textbf{Chapter 1. Overview of the Internship Organization};
    \item \textbf{Chapter 2. Internship Activities and Assigned Tasks};
    \item \textbf{Chapter 3. Results, Experience, and Evaluation}.
\end{itemize}

These chapter titles are recommendations only and may be adapted to the
actual internship content.

The report shall reflect work actually performed, observed, analyzed, or
learned during the internship. Unrelated theoretical material should not
be included merely to increase report length.

\vspace{0.3cm}

% ================================================================
% X. TABLE OF CONTENTS
% ================================================================

\noindent
\textbf{X.\quad TABLE OF CONTENTS}

\vspace{0.15cm}

The Table of Contents shall be generated automatically using:

\begin{verbatim}
\tableofcontents
\end{verbatim}

Students shall not create or manually edit page numbers in the Table of
Contents.

The document shall be compiled sufficiently to ensure that all page
numbers and entries are correctly updated.

\vspace{0.3cm}

% ================================================================
% XI. LIST OF SYMBOLS AND ABBREVIATIONS
% ================================================================

\noindent
\textbf{XI.\quad LIST OF SYMBOLS AND ABBREVIATIONS}

\vspace{0.15cm}

Symbols and abbreviations that are used repeatedly in the report shall be
listed in the \textbf{List of Symbols and Abbreviations}.

The list may contain the following columns:

\begin{center}
\begin{tabular}{llll}
    No. & Symbol / Abbreviation & Unit & Definition
\end{tabular}
\end{center}

Only symbols and abbreviations actually used in the final report shall be
retained. Example entries contained in the template shall be removed if
they are not relevant to the student's report.

Each abbreviation should normally be defined when it first appears in
the main text.

\vspace{0.3cm}

% ================================================================
% XII. FIGURES, DIAGRAMS, AND PHOTOGRAPHS
% ================================================================

\noindent
\textbf{XII.\quad FIGURES, DIAGRAMS, AND PHOTOGRAPHS}

\vspace{0.15cm}

Figures shall be clear, legible, and directly relevant to the report
content.

Figures shall be inserted using the standard \verb|figure| environment.

Example:

\begin{verbatim}
\begin{figure}[H]
    \centering
    \includegraphics[width=0.8\textwidth]
                    {figures/company_structure.pdf}
    \caption{Organizational structure of the company}
    \label{fig:company-structure}
\end{figure}
\end{verbatim}

Figure captions shall be placed below the figure.

Each referenced figure shall have a unique \verb|\label| and shall be
referenced in the text using, for example:

\begin{verbatim}
Figure~\ref{fig:company-structure}
\end{verbatim}

Students shall not manually type figure numbers.

Figures reproduced or adapted from external sources shall include an
appropriate citation.

Photographs taken during the internship shall be relevant and sufficiently
clear. Students shall not include confidential, restricted, or proprietary
information without permission from the internship organization.

Vector graphics or high-resolution images are recommended whenever
possible.

\vspace{0.3cm}

% ================================================================
% XIII. TABLES
% ================================================================

\noindent
\textbf{XIII.\quad TABLES}

\vspace{0.15cm}

Tables shall be created using suitable \LaTeX{} environments such as
\verb|table|, \verb|tabular|, \verb|tabularx|, or \verb|longtable|.

Table captions shall be placed above the table.

Example:

\begin{verbatim}
\begin{table}[H]
    \centering
    \caption{Internship schedule}
    \label{tab:internship-schedule}

    \begin{tabular}{lll}
        \toprule
        Week & Task & Result \\
        \midrule
        1 & Orientation & Completed \\
        2 & System study & Completed \\
        \bottomrule
    \end{tabular}
\end{table}
\end{verbatim}

Tables shall be referenced using:

\begin{verbatim}
Table~\ref{tab:internship-schedule}
\end{verbatim}

Students shall not manually type table numbers.

Tables should be created directly in \LaTeX{} rather than inserted as
screenshots whenever practical.

\vspace{0.3cm}

% ================================================================
% XIV. EQUATIONS AND TECHNICAL QUANTITIES
% ================================================================

\noindent
\textbf{XIV.\quad EQUATIONS AND TECHNICAL QUANTITIES}

\vspace{0.15cm}

Mathematical expressions shall be typeset directly using \LaTeX{}.
Images or screenshots of equations shall not be used.

Short expressions may be written in inline mathematical mode, for example:

\begin{verbatim}
$R = 0.7\,\Omega$
\end{verbatim}

Important equations shall use a numbered equation environment.

Example:

\begin{verbatim}
\begin{equation}
    P = UI
    \label{eq:power}
\end{equation}
\end{verbatim}

Equations shall be referenced using:

\begin{verbatim}
Equation~\eqref{eq:power}
\end{verbatim}

Equation numbers shall not be manually entered.

SI units shall be used wherever applicable. The \verb|siunitx| package
should be used where appropriate.

Examples:

\begin{verbatim}
\SI{220}{\volt}
\SI{50}{\hertz}
\SI{2.5}{\ampere}
\end{verbatim}

Notation and units shall be consistent throughout the report.

\vspace{0.3cm}

% ================================================================
% XV. SOURCE CODE
% ================================================================

\noindent
\textbf{XV.\quad SOURCE CODE}

\vspace{0.15cm}

If the internship involves programming, PLCs, microcontrollers,
simulation, or software development, relevant source code may be included
in the report.

Source code shall be presented using the \verb|listings| environment or
another environment defined by the official template.

Example:

\begin{verbatim}
\begin{lstlisting}[language=st]
IF Start = TRUE THEN
    Motor := TRUE;
ELSE
    Motor := FALSE;
END_IF;
\end{lstlisting}
\end{verbatim}

Only code necessary for explaining the work should appear in the main
report. Long source-code listings should be moved to an appendix or
submitted as supplementary files if required.

Source code should not normally be inserted as screenshots.

\vspace{0.3cm}

% ================================================================
% XVI. ABSTRACT
% ================================================================

\noindent
\textbf{XVI.\quad ABSTRACT}

\vspace{0.15cm}

The Abstract shall briefly summarize:

\begin{itemize}[leftmargin=0.8cm, itemsep=2pt]
    \item the internship organization;
    \item the objectives of the internship;
    \item the main tasks performed;
    \item the main results, knowledge, and skills obtained.
\end{itemize}

The Abstract should be concise and should not reproduce the full content
of the report chapter by chapter.

\vspace{0.3cm}

% ================================================================
% XVII. REFERENCES AND CITATIONS
% ================================================================

\noindent
\textbf{XVII.\quad REFERENCES AND CITATIONS}

\vspace{0.15cm}

All books, journal papers, conference papers, standards, websites,
datasheets, manuals, technical documents, figures, tables, or other
external sources used in preparing the report shall be properly cited.

References shall be managed through the bibliography system provided with
the official template.

Bibliographic entries shall normally be stored in:

\begin{verbatim}
References.bib
\end{verbatim}

References shall be cited in the report using:

\begin{verbatim}
\cite{reference-key}
\end{verbatim}

Students shall not manually type reference numbers such as ``[1]'',
``[2]'', etc.

Titles of foreign-language publications shall remain in their original
language in the bibliography.

All borrowed ideas, data, figures, tables, equations, methods, or results
shall be properly acknowledged.

\vspace{0.3cm}

% ================================================================
% XVIII. CROSS-REFERENCING
% ================================================================

\noindent
\textbf{XVIII.\quad CROSS-REFERENCING}

\vspace{0.15cm}

Numbered elements shall be cross-referenced using the standard \LaTeX{}
mechanism.

Use the following commands where appropriate:

\begin{verbatim}
\label{...}
\ref{...}
\eqref{...}
\end{verbatim}

For example:

\begin{verbatim}
As shown in Figure~\ref{fig:system}, ...
\end{verbatim}

Students shall not manually type numbers such as ``Figure 2.3'' when
those numbers can be generated automatically.

\vspace{0.3cm}

% ================================================================
% XIX. PAGE NUMBERING
% ================================================================

\noindent
\textbf{XIX.\quad PAGE NUMBERING}

\vspace{0.15cm}

Cover pages and official evaluation forms at the beginning of the report
shall not display page numbers.

The main body of the report shall use Arabic page numbers beginning from
Chapter 1.

Page numbering shall be generated automatically by the official template.
Students shall not manually enter page numbers.

\vspace{0.3cm}

% ================================================================
% XX. PROJECT FILE ORGANIZATION
% ================================================================

\noindent
\textbf{XX.\quad PROJECT FILE ORGANIZATION}

\vspace{0.15cm}

Students should maintain a clear and organized \LaTeX{} project
structure. A recommended structure is:

\begin{verbatim}
InternshipReport/
|
|-- main.tex
|-- References.bib
|
|-- figures/
|   |-- company.pdf
|   |-- system.png
|   `-- ...
|
`-- chapters/
    |-- chapter1.tex
    |-- chapter2.tex
    `-- chapter3.tex
\end{verbatim}

Figures should normally be stored in the \verb|figures| directory
specified by the official template.

File and directory names should preferably contain Latin letters,
numbers, hyphens, or underscores.

Recommended examples include:

\begin{verbatim}
company_structure.pdf
motor_control.png
chapter_01.tex
\end{verbatim}

Spaces and unnecessary special characters in file names should be
avoided where possible.

\vspace{0.3cm}

% ================================================================
% XXI. DOCUMENT COMPILATION AND FINAL CHECK
% ================================================================

\noindent
\textbf{XXI.\quad DOCUMENT COMPILATION AND FINAL CHECK}

\vspace{0.15cm}

Before submission, the complete \LaTeX{} project shall compile
successfully and produce the final PDF report.

The final document shall not contain unresolved references such as:

\begin{verbatim}
??
[?]
\end{verbatim}

Before generating the final PDF, students shall check:

\begin{itemize}[leftmargin=0.8cm, itemsep=2pt]
    \item the Table of Contents;
    \item page numbering;
    \item chapter, section, and subsection numbering;
    \item figure numbering;
    \item table numbering;
    \item equation numbering;
    \item cross-references;
    \item citations;
    \item bibliography;
    \item font display;
    \item mathematical symbols;
    \item image quality;
    \item completeness of all required forms.
\end{itemize}

The document may be compiled using TeX Live, MiKTeX, or another
compatible \LaTeX{} environment. Regardless of the compilation
environment, the resulting PDF shall conform to the official Faculty
template.

\vspace{0.3cm}

% ================================================================
% XXII. CONTENT THAT STUDENTS MAY EDIT
% ================================================================

\noindent
\textbf{XXII.\quad CONTENT THAT STUDENTS MAY EDIT}

\vspace{0.15cm}

Students may edit or replace:

\begin{itemize}[leftmargin=0.8cm, itemsep=2pt]
    \item student name;
    \item class;
    \item internship organization;
    \item supervisor name;
    \item Abstract;
    \item List of Symbols and Abbreviations;
    \item chapter titles and chapter content;
    \item figures and photographs;
    \item tables;
    \item equations;
    \item citations and references.
\end{itemize}

Students shall not modify the main document formatting settings in the
preamble, including the document class, page geometry, line spacing,
heading formatting, table-of-contents formatting, font configuration,
page layout, or cover-page macros unless specifically instructed.

\vspace{0.3cm}

% ================================================================
% XXIII. FINAL SUBMISSION REQUIREMENTS
% ================================================================

\noindent
\textbf{XXIII.\quad FINAL SUBMISSION REQUIREMENTS}

\vspace{0.15cm}

Before submission, the final report shall:

\begin{itemize}[leftmargin=0.8cm, itemsep=2pt]
    \item compile without errors that affect the report content or presentation;
    \item contain all required forms and report sections;
    \item contain correct student, class, internship organization, and supervisor information;
    \item contain an updated Table of Contents, List of Figures, and List of Tables;
    \item contain no unused example content from the template;
    \item contain no placeholder text such as \textit{Title of Chapter 1},
    \textit{Section 1.1}, or other sample entries;
    \item contain all required citations and references;
    \item contain no unresolved cross-references.
\end{itemize}

If the Faculty requires submission of the \LaTeX{} source, the complete
project shall include all \verb|.tex| files, the bibliography database,
figures, and any other files required to reproduce the submitted PDF.

\vspace{0.4cm}

\noindent
\textbf{Important:}
Before submitting the final internship report, students shall remove or
comment out the following command in the main document:

\begin{verbatim}
\input{Regulations.tex}
\end{verbatim}

For example:

\begin{verbatim}
% \input{Regulations.tex}
\end{verbatim}

This removes the guidance pages from the final submitted report while
preserving the regulations file in the \LaTeX{} project for reference.

} % End 13pt content

\clearpage

%
% IMPORTANT:
% This section is provided for guidance only. Students shall remove
% or comment out the 
% =================================================================
% REGULATIONS FOR PREPARING THE GRADUATION INTERNSHIP REPORT
% USING LATEX
% =================================================================
%
% This file is intended to be included from the main document with:
%
%     \input{Regulations.tex}
%
% IMPORTANT:
% This section is provided for guidance only. Students shall remove
% or comment out the \input{Regulations.tex} command before submitting
% the final internship report.
%
% =================================================================

\clearpage
\thispagestyle{empty}

\begin{center}
    {\fontsize{16pt}{20pt}\selectfont\bfseries
    SPECIFIC REGULATIONS FOR PREPARING\\
    GRADUATION INTERNSHIP REPORTS USING \LaTeX}
\end{center}

\vspace{0.5cm}

{\fontsize{13pt}{17pt}\selectfont
\setstretch{1.3}
\justifying

\noindent
\textbf{Note:}
This section is provided for guidance only and is included in the main
document using the command \verb|\input{Regulations.tex}|.
Students shall remove or comment out this command before submitting the
final internship report.

\vspace{0.4cm}

% ================================================================
% I. GENERAL REQUIREMENTS
% ================================================================

\noindent
\textbf{I.\quad GENERAL REQUIREMENTS}

\vspace{0.15cm}

The Graduation Internship Report shall be prepared using the official
\LaTeX{} template provided by the Faculty. Students shall preserve the
predefined page layout, typography, line spacing, heading styles,
numbering system, figure and table formatting, and bibliography settings
unless otherwise instructed by the Faculty or the supervisor.

The report shall be written in English.

Students are responsible for the academic and technical content of the
report. The official \LaTeX{} template is responsible for standardized
document formatting.

\vspace{0.3cm}

% ================================================================
% II. REPORT STRUCTURE
% ================================================================

\noindent
\textbf{II.\quad REPORT STRUCTURE}

\vspace{0.15cm}

The report shall normally be arranged in the following order:

\begin{enumerate}[label=\arabic*., leftmargin=1.2cm, itemsep=2pt]
    \item Outer Cover;
    \item Comments of the Internship Host Organization;
    \item Supervisor's Evaluation and Graduation Internship Defense Results;
    \item Table of Contents;
    \item List of Symbols and Abbreviations, if applicable;
    \item List of Figures;
    \item List of Tables;
    \item Abstract;
    \item Main Chapters of the Internship Report;
    \item References;
    \item Appendices, if applicable.
\end{enumerate}

The preliminary pages shall not be numbered as chapters. The main report
content shall be organized into consecutively numbered chapters beginning
with Chapter 1.

\vspace{0.3cm}

% ================================================================
% III. PAGE SIZE AND MARGINS
% ================================================================

\noindent
\textbf{III.\quad PAGE SIZE AND MARGINS}

\vspace{0.15cm}

The report shall be formatted on A4 paper with the following margins:

\begin{center}
\begin{tabular}{ll}
    Left margin:   & 3.5 cm \\
    Right margin:  & 2.0 cm \\
    Top margin:    & 2.5 cm \\
    Bottom margin: & 2.5 cm
\end{tabular}
\end{center}

These settings are predefined in the official template and shall not be
modified by students.

\vspace{0.3cm}

% ================================================================
% IV. FONT, FONT SIZE, AND LINE SPACING
% ================================================================

\noindent
\textbf{IV.\quad FONT, FONT SIZE, AND LINE SPACING}

\vspace{0.15cm}

The main text shall use a 13 pt font size with 1.3-line spacing.
A Times-family typeface shall be used in accordance with the official
template.

The text shall be fully justified.

Students shall not use repeated spaces, repeated manual line breaks, or
multiple blank lines to position text. Paragraph and heading spacing shall
be controlled by the template.

\vspace{0.3cm}

% ================================================================
% V. OUTER COVER
% ================================================================

\noindent
\textbf{V.\quad OUTER COVER}

\vspace{0.15cm}

The outer cover shall follow the official format provided in the template
and shall contain the following information:

\begin{itemize}[leftmargin=0.8cm, itemsep=2pt]
    \item Thai Nguyen University of Technology;
    \item Faculty of International Training;
    \item Engineering and Technology Department;
    \item University logo and Faculty logo;
    \item the title \textbf{INTERNSHIP REPORT};
    \item Internship Organization;
    \item Author;
    \item Class;
    \item Supervisor;
    \item place and year of completion.
\end{itemize}

Students shall only replace the required information fields and shall not
modify the cover-page layout unless instructed otherwise.

\vspace{0.3cm}

% ================================================================
% VI. COMMENTS OF THE INTERNSHIP HOST ORGANIZATION
% ================================================================

\noindent
\textbf{VI.\quad COMMENTS OF THE INTERNSHIP HOST ORGANIZATION}

\vspace{0.15cm}

The page entitled \textbf{COMMENTS OF THE INTERNSHIP HOST ORGANIZATION}
is reserved for the internship host organization.

It includes:

\begin{itemize}[leftmargin=0.8cm, itemsep=2pt]
    \item the student's full name;
    \item class;
    \item comments from the host organization;
    \item score;
    \item date;
    \item signature and official seal.
\end{itemize}

Students shall not alter or complete official evaluation fields unless
specifically authorized.

\vspace{0.3cm}

% ================================================================
% VII. SUPERVISOR'S EVALUATION AND DEFENSE RESULTS
% ================================================================

\noindent
\textbf{VII.\quad SUPERVISOR'S EVALUATION AND DEFENSE RESULTS}

\vspace{0.15cm}

The section entitled
\textbf{SUPERVISOR'S EVALUATION AND GRADUATION INTERNSHIP DEFENSE RESULTS}
is reserved for the supervisor and defense committee.

The evaluation may include:

\begin{itemize}[leftmargin=0.8cm, itemsep=2pt]
    \item Internship Progress and Student's Attitude;
    \item Report Content;
    \item Report Presentation;
    \item Other Comments;
    \item Supervisor's Opinion;
    \item Defense Result;
    \item Score.
\end{itemize}

Students shall not enter evaluation scores or defense results before they
are officially determined.

\vspace{0.3cm}

% ================================================================
% VIII. CHAPTERS, SECTIONS, AND SUBSECTIONS
% ================================================================

\noindent
\textbf{VIII.\quad CHAPTERS, SECTIONS, AND SUBSECTIONS}

\vspace{0.15cm}

The report structure shall be generated using the standard \LaTeX{}
commands:

\begin{verbatim}
\chapter{...}
\section{...}
\subsection{...}
\end{verbatim}

Chapter, section, and subsection numbers shall be generated automatically
by \LaTeX{}.

Students shall not manually type chapter or section numbers and shall not
reproduce heading styles using commands such as \verb|\textbf|,
\verb|\fontsize|, or \verb|\centering| instead of the appropriate
structural command.

Excessive subdivision should be avoided. In most cases, the hierarchy
Chapter--Section--Subsection is sufficient.

\vspace{0.3cm}

% ================================================================
% IX. ORGANIZATION OF THE MAIN CONTENT
% ================================================================

\noindent
\textbf{IX.\quad ORGANIZATION OF THE MAIN CONTENT}

\vspace{0.15cm}

The specific content of the report shall depend on the internship
organization and the assigned tasks. A typical structure may include:

\begin{itemize}[leftmargin=0.8cm, itemsep=2pt]
    \item \textbf{Chapter 1. Overview of the Internship Organization};
    \item \textbf{Chapter 2. Internship Activities and Assigned Tasks};
    \item \textbf{Chapter 3. Results, Experience, and Evaluation}.
\end{itemize}

These chapter titles are recommendations only and may be adapted to the
actual internship content.

The report shall reflect work actually performed, observed, analyzed, or
learned during the internship. Unrelated theoretical material should not
be included merely to increase report length.

\vspace{0.3cm}

% ================================================================
% X. TABLE OF CONTENTS
% ================================================================

\noindent
\textbf{X.\quad TABLE OF CONTENTS}

\vspace{0.15cm}

The Table of Contents shall be generated automatically using:

\begin{verbatim}
\tableofcontents
\end{verbatim}

Students shall not create or manually edit page numbers in the Table of
Contents.

The document shall be compiled sufficiently to ensure that all page
numbers and entries are correctly updated.

\vspace{0.3cm}

% ================================================================
% XI. LIST OF SYMBOLS AND ABBREVIATIONS
% ================================================================

\noindent
\textbf{XI.\quad LIST OF SYMBOLS AND ABBREVIATIONS}

\vspace{0.15cm}

Symbols and abbreviations that are used repeatedly in the report shall be
listed in the \textbf{List of Symbols and Abbreviations}.

The list may contain the following columns:

\begin{center}
\begin{tabular}{llll}
    No. & Symbol / Abbreviation & Unit & Definition
\end{tabular}
\end{center}

Only symbols and abbreviations actually used in the final report shall be
retained. Example entries contained in the template shall be removed if
they are not relevant to the student's report.

Each abbreviation should normally be defined when it first appears in
the main text.

\vspace{0.3cm}

% ================================================================
% XII. FIGURES, DIAGRAMS, AND PHOTOGRAPHS
% ================================================================

\noindent
\textbf{XII.\quad FIGURES, DIAGRAMS, AND PHOTOGRAPHS}

\vspace{0.15cm}

Figures shall be clear, legible, and directly relevant to the report
content.

Figures shall be inserted using the standard \verb|figure| environment.

Example:

\begin{verbatim}
\begin{figure}[H]
    \centering
    \includegraphics[width=0.8\textwidth]
                    {figures/company_structure.pdf}
    \caption{Organizational structure of the company}
    \label{fig:company-structure}
\end{figure}
\end{verbatim}

Figure captions shall be placed below the figure.

Each referenced figure shall have a unique \verb|\label| and shall be
referenced in the text using, for example:

\begin{verbatim}
Figure~\ref{fig:company-structure}
\end{verbatim}

Students shall not manually type figure numbers.

Figures reproduced or adapted from external sources shall include an
appropriate citation.

Photographs taken during the internship shall be relevant and sufficiently
clear. Students shall not include confidential, restricted, or proprietary
information without permission from the internship organization.

Vector graphics or high-resolution images are recommended whenever
possible.

\vspace{0.3cm}

% ================================================================
% XIII. TABLES
% ================================================================

\noindent
\textbf{XIII.\quad TABLES}

\vspace{0.15cm}

Tables shall be created using suitable \LaTeX{} environments such as
\verb|table|, \verb|tabular|, \verb|tabularx|, or \verb|longtable|.

Table captions shall be placed above the table.

Example:

\begin{verbatim}
\begin{table}[H]
    \centering
    \caption{Internship schedule}
    \label{tab:internship-schedule}

    \begin{tabular}{lll}
        \toprule
        Week & Task & Result \\
        \midrule
        1 & Orientation & Completed \\
        2 & System study & Completed \\
        \bottomrule
    \end{tabular}
\end{table}
\end{verbatim}

Tables shall be referenced using:

\begin{verbatim}
Table~\ref{tab:internship-schedule}
\end{verbatim}

Students shall not manually type table numbers.

Tables should be created directly in \LaTeX{} rather than inserted as
screenshots whenever practical.

\vspace{0.3cm}

% ================================================================
% XIV. EQUATIONS AND TECHNICAL QUANTITIES
% ================================================================

\noindent
\textbf{XIV.\quad EQUATIONS AND TECHNICAL QUANTITIES}

\vspace{0.15cm}

Mathematical expressions shall be typeset directly using \LaTeX{}.
Images or screenshots of equations shall not be used.

Short expressions may be written in inline mathematical mode, for example:

\begin{verbatim}
$R = 0.7\,\Omega$
\end{verbatim}

Important equations shall use a numbered equation environment.

Example:

\begin{verbatim}
\begin{equation}
    P = UI
    \label{eq:power}
\end{equation}
\end{verbatim}

Equations shall be referenced using:

\begin{verbatim}
Equation~\eqref{eq:power}
\end{verbatim}

Equation numbers shall not be manually entered.

SI units shall be used wherever applicable. The \verb|siunitx| package
should be used where appropriate.

Examples:

\begin{verbatim}
\SI{220}{\volt}
\SI{50}{\hertz}
\SI{2.5}{\ampere}
\end{verbatim}

Notation and units shall be consistent throughout the report.

\vspace{0.3cm}

% ================================================================
% XV. SOURCE CODE
% ================================================================

\noindent
\textbf{XV.\quad SOURCE CODE}

\vspace{0.15cm}

If the internship involves programming, PLCs, microcontrollers,
simulation, or software development, relevant source code may be included
in the report.

Source code shall be presented using the \verb|listings| environment or
another environment defined by the official template.

Example:

\begin{verbatim}
\begin{lstlisting}[language=st]
IF Start = TRUE THEN
    Motor := TRUE;
ELSE
    Motor := FALSE;
END_IF;
\end{lstlisting}
\end{verbatim}

Only code necessary for explaining the work should appear in the main
report. Long source-code listings should be moved to an appendix or
submitted as supplementary files if required.

Source code should not normally be inserted as screenshots.

\vspace{0.3cm}

% ================================================================
% XVI. ABSTRACT
% ================================================================

\noindent
\textbf{XVI.\quad ABSTRACT}

\vspace{0.15cm}

The Abstract shall briefly summarize:

\begin{itemize}[leftmargin=0.8cm, itemsep=2pt]
    \item the internship organization;
    \item the objectives of the internship;
    \item the main tasks performed;
    \item the main results, knowledge, and skills obtained.
\end{itemize}

The Abstract should be concise and should not reproduce the full content
of the report chapter by chapter.

\vspace{0.3cm}

% ================================================================
% XVII. REFERENCES AND CITATIONS
% ================================================================

\noindent
\textbf{XVII.\quad REFERENCES AND CITATIONS}

\vspace{0.15cm}

All books, journal papers, conference papers, standards, websites,
datasheets, manuals, technical documents, figures, tables, or other
external sources used in preparing the report shall be properly cited.

References shall be managed through the bibliography system provided with
the official template.

Bibliographic entries shall normally be stored in:

\begin{verbatim}
References.bib
\end{verbatim}

References shall be cited in the report using:

\begin{verbatim}
\cite{reference-key}
\end{verbatim}

Students shall not manually type reference numbers such as ``[1]'',
``[2]'', etc.

Titles of foreign-language publications shall remain in their original
language in the bibliography.

All borrowed ideas, data, figures, tables, equations, methods, or results
shall be properly acknowledged.

\vspace{0.3cm}

% ================================================================
% XVIII. CROSS-REFERENCING
% ================================================================

\noindent
\textbf{XVIII.\quad CROSS-REFERENCING}

\vspace{0.15cm}

Numbered elements shall be cross-referenced using the standard \LaTeX{}
mechanism.

Use the following commands where appropriate:

\begin{verbatim}
\label{...}
\ref{...}
\eqref{...}
\end{verbatim}

For example:

\begin{verbatim}
As shown in Figure~\ref{fig:system}, ...
\end{verbatim}

Students shall not manually type numbers such as ``Figure 2.3'' when
those numbers can be generated automatically.

\vspace{0.3cm}

% ================================================================
% XIX. PAGE NUMBERING
% ================================================================

\noindent
\textbf{XIX.\quad PAGE NUMBERING}

\vspace{0.15cm}

Cover pages and official evaluation forms at the beginning of the report
shall not display page numbers.

The main body of the report shall use Arabic page numbers beginning from
Chapter 1.

Page numbering shall be generated automatically by the official template.
Students shall not manually enter page numbers.

\vspace{0.3cm}

% ================================================================
% XX. PROJECT FILE ORGANIZATION
% ================================================================

\noindent
\textbf{XX.\quad PROJECT FILE ORGANIZATION}

\vspace{0.15cm}

Students should maintain a clear and organized \LaTeX{} project
structure. A recommended structure is:

\begin{verbatim}
InternshipReport/
|
|-- main.tex
|-- References.bib
|
|-- figures/
|   |-- company.pdf
|   |-- system.png
|   `-- ...
|
`-- chapters/
    |-- chapter1.tex
    |-- chapter2.tex
    `-- chapter3.tex
\end{verbatim}

Figures should normally be stored in the \verb|figures| directory
specified by the official template.

File and directory names should preferably contain Latin letters,
numbers, hyphens, or underscores.

Recommended examples include:

\begin{verbatim}
company_structure.pdf
motor_control.png
chapter_01.tex
\end{verbatim}

Spaces and unnecessary special characters in file names should be
avoided where possible.

\vspace{0.3cm}

% ================================================================
% XXI. DOCUMENT COMPILATION AND FINAL CHECK
% ================================================================

\noindent
\textbf{XXI.\quad DOCUMENT COMPILATION AND FINAL CHECK}

\vspace{0.15cm}

Before submission, the complete \LaTeX{} project shall compile
successfully and produce the final PDF report.

The final document shall not contain unresolved references such as:

\begin{verbatim}
??
[?]
\end{verbatim}

Before generating the final PDF, students shall check:

\begin{itemize}[leftmargin=0.8cm, itemsep=2pt]
    \item the Table of Contents;
    \item page numbering;
    \item chapter, section, and subsection numbering;
    \item figure numbering;
    \item table numbering;
    \item equation numbering;
    \item cross-references;
    \item citations;
    \item bibliography;
    \item font display;
    \item mathematical symbols;
    \item image quality;
    \item completeness of all required forms.
\end{itemize}

The document may be compiled using TeX Live, MiKTeX, or another
compatible \LaTeX{} environment. Regardless of the compilation
environment, the resulting PDF shall conform to the official Faculty
template.

\vspace{0.3cm}

% ================================================================
% XXII. CONTENT THAT STUDENTS MAY EDIT
% ================================================================

\noindent
\textbf{XXII.\quad CONTENT THAT STUDENTS MAY EDIT}

\vspace{0.15cm}

Students may edit or replace:

\begin{itemize}[leftmargin=0.8cm, itemsep=2pt]
    \item student name;
    \item class;
    \item internship organization;
    \item supervisor name;
    \item Abstract;
    \item List of Symbols and Abbreviations;
    \item chapter titles and chapter content;
    \item figures and photographs;
    \item tables;
    \item equations;
    \item citations and references.
\end{itemize}

Students shall not modify the main document formatting settings in the
preamble, including the document class, page geometry, line spacing,
heading formatting, table-of-contents formatting, font configuration,
page layout, or cover-page macros unless specifically instructed.

\vspace{0.3cm}

% ================================================================
% XXIII. FINAL SUBMISSION REQUIREMENTS
% ================================================================

\noindent
\textbf{XXIII.\quad FINAL SUBMISSION REQUIREMENTS}

\vspace{0.15cm}

Before submission, the final report shall:

\begin{itemize}[leftmargin=0.8cm, itemsep=2pt]
    \item compile without errors that affect the report content or presentation;
    \item contain all required forms and report sections;
    \item contain correct student, class, internship organization, and supervisor information;
    \item contain an updated Table of Contents, List of Figures, and List of Tables;
    \item contain no unused example content from the template;
    \item contain no placeholder text such as \textit{Title of Chapter 1},
    \textit{Section 1.1}, or other sample entries;
    \item contain all required citations and references;
    \item contain no unresolved cross-references.
\end{itemize}

If the Faculty requires submission of the \LaTeX{} source, the complete
project shall include all \verb|.tex| files, the bibliography database,
figures, and any other files required to reproduce the submitted PDF.

\vspace{0.4cm}

\noindent
\textbf{Important:}
Before submitting the final internship report, students shall remove or
comment out the following command in the main document:

\begin{verbatim}
\input{Regulations.tex}
\end{verbatim}

For example:

\begin{verbatim}
% \input{Regulations.tex}
\end{verbatim}

This removes the guidance pages from the final submitted report while
preserving the regulations file in the \LaTeX{} project for reference.

} % End 13pt content

\clearpage
 command before submitting
% the final internship report.
%
% =================================================================

\clearpage
\thispagestyle{empty}

\begin{center}
    {\fontsize{16pt}{20pt}\selectfont\bfseries
    SPECIFIC REGULATIONS FOR PREPARING\\
    GRADUATION INTERNSHIP REPORTS USING \LaTeX}
\end{center}

\vspace{0.5cm}

{\fontsize{13pt}{17pt}\selectfont
\setstretch{1.3}
\justifying

\noindent
\textbf{Note:}
This section is provided for guidance only and is included in the main
document using the command \verb|
% =================================================================
% REGULATIONS FOR PREPARING THE GRADUATION INTERNSHIP REPORT
% USING LATEX
% =================================================================
%
% This file is intended to be included from the main document with:
%
%     \input{Regulations.tex}
%
% IMPORTANT:
% This section is provided for guidance only. Students shall remove
% or comment out the \input{Regulations.tex} command before submitting
% the final internship report.
%
% =================================================================

\clearpage
\thispagestyle{empty}

\begin{center}
    {\fontsize{16pt}{20pt}\selectfont\bfseries
    SPECIFIC REGULATIONS FOR PREPARING\\
    GRADUATION INTERNSHIP REPORTS USING \LaTeX}
\end{center}

\vspace{0.5cm}

{\fontsize{13pt}{17pt}\selectfont
\setstretch{1.3}
\justifying

\noindent
\textbf{Note:}
This section is provided for guidance only and is included in the main
document using the command \verb|\input{Regulations.tex}|.
Students shall remove or comment out this command before submitting the
final internship report.

\vspace{0.4cm}

% ================================================================
% I. GENERAL REQUIREMENTS
% ================================================================

\noindent
\textbf{I.\quad GENERAL REQUIREMENTS}

\vspace{0.15cm}

The Graduation Internship Report shall be prepared using the official
\LaTeX{} template provided by the Faculty. Students shall preserve the
predefined page layout, typography, line spacing, heading styles,
numbering system, figure and table formatting, and bibliography settings
unless otherwise instructed by the Faculty or the supervisor.

The report shall be written in English.

Students are responsible for the academic and technical content of the
report. The official \LaTeX{} template is responsible for standardized
document formatting.

\vspace{0.3cm}

% ================================================================
% II. REPORT STRUCTURE
% ================================================================

\noindent
\textbf{II.\quad REPORT STRUCTURE}

\vspace{0.15cm}

The report shall normally be arranged in the following order:

\begin{enumerate}[label=\arabic*., leftmargin=1.2cm, itemsep=2pt]
    \item Outer Cover;
    \item Comments of the Internship Host Organization;
    \item Supervisor's Evaluation and Graduation Internship Defense Results;
    \item Table of Contents;
    \item List of Symbols and Abbreviations, if applicable;
    \item List of Figures;
    \item List of Tables;
    \item Abstract;
    \item Main Chapters of the Internship Report;
    \item References;
    \item Appendices, if applicable.
\end{enumerate}

The preliminary pages shall not be numbered as chapters. The main report
content shall be organized into consecutively numbered chapters beginning
with Chapter 1.

\vspace{0.3cm}

% ================================================================
% III. PAGE SIZE AND MARGINS
% ================================================================

\noindent
\textbf{III.\quad PAGE SIZE AND MARGINS}

\vspace{0.15cm}

The report shall be formatted on A4 paper with the following margins:

\begin{center}
\begin{tabular}{ll}
    Left margin:   & 3.5 cm \\
    Right margin:  & 2.0 cm \\
    Top margin:    & 2.5 cm \\
    Bottom margin: & 2.5 cm
\end{tabular}
\end{center}

These settings are predefined in the official template and shall not be
modified by students.

\vspace{0.3cm}

% ================================================================
% IV. FONT, FONT SIZE, AND LINE SPACING
% ================================================================

\noindent
\textbf{IV.\quad FONT, FONT SIZE, AND LINE SPACING}

\vspace{0.15cm}

The main text shall use a 13 pt font size with 1.3-line spacing.
A Times-family typeface shall be used in accordance with the official
template.

The text shall be fully justified.

Students shall not use repeated spaces, repeated manual line breaks, or
multiple blank lines to position text. Paragraph and heading spacing shall
be controlled by the template.

\vspace{0.3cm}

% ================================================================
% V. OUTER COVER
% ================================================================

\noindent
\textbf{V.\quad OUTER COVER}

\vspace{0.15cm}

The outer cover shall follow the official format provided in the template
and shall contain the following information:

\begin{itemize}[leftmargin=0.8cm, itemsep=2pt]
    \item Thai Nguyen University of Technology;
    \item Faculty of International Training;
    \item Engineering and Technology Department;
    \item University logo and Faculty logo;
    \item the title \textbf{INTERNSHIP REPORT};
    \item Internship Organization;
    \item Author;
    \item Class;
    \item Supervisor;
    \item place and year of completion.
\end{itemize}

Students shall only replace the required information fields and shall not
modify the cover-page layout unless instructed otherwise.

\vspace{0.3cm}

% ================================================================
% VI. COMMENTS OF THE INTERNSHIP HOST ORGANIZATION
% ================================================================

\noindent
\textbf{VI.\quad COMMENTS OF THE INTERNSHIP HOST ORGANIZATION}

\vspace{0.15cm}

The page entitled \textbf{COMMENTS OF THE INTERNSHIP HOST ORGANIZATION}
is reserved for the internship host organization.

It includes:

\begin{itemize}[leftmargin=0.8cm, itemsep=2pt]
    \item the student's full name;
    \item class;
    \item comments from the host organization;
    \item score;
    \item date;
    \item signature and official seal.
\end{itemize}

Students shall not alter or complete official evaluation fields unless
specifically authorized.

\vspace{0.3cm}

% ================================================================
% VII. SUPERVISOR'S EVALUATION AND DEFENSE RESULTS
% ================================================================

\noindent
\textbf{VII.\quad SUPERVISOR'S EVALUATION AND DEFENSE RESULTS}

\vspace{0.15cm}

The section entitled
\textbf{SUPERVISOR'S EVALUATION AND GRADUATION INTERNSHIP DEFENSE RESULTS}
is reserved for the supervisor and defense committee.

The evaluation may include:

\begin{itemize}[leftmargin=0.8cm, itemsep=2pt]
    \item Internship Progress and Student's Attitude;
    \item Report Content;
    \item Report Presentation;
    \item Other Comments;
    \item Supervisor's Opinion;
    \item Defense Result;
    \item Score.
\end{itemize}

Students shall not enter evaluation scores or defense results before they
are officially determined.

\vspace{0.3cm}

% ================================================================
% VIII. CHAPTERS, SECTIONS, AND SUBSECTIONS
% ================================================================

\noindent
\textbf{VIII.\quad CHAPTERS, SECTIONS, AND SUBSECTIONS}

\vspace{0.15cm}

The report structure shall be generated using the standard \LaTeX{}
commands:

\begin{verbatim}
\chapter{...}
\section{...}
\subsection{...}
\end{verbatim}

Chapter, section, and subsection numbers shall be generated automatically
by \LaTeX{}.

Students shall not manually type chapter or section numbers and shall not
reproduce heading styles using commands such as \verb|\textbf|,
\verb|\fontsize|, or \verb|\centering| instead of the appropriate
structural command.

Excessive subdivision should be avoided. In most cases, the hierarchy
Chapter--Section--Subsection is sufficient.

\vspace{0.3cm}

% ================================================================
% IX. ORGANIZATION OF THE MAIN CONTENT
% ================================================================

\noindent
\textbf{IX.\quad ORGANIZATION OF THE MAIN CONTENT}

\vspace{0.15cm}

The specific content of the report shall depend on the internship
organization and the assigned tasks. A typical structure may include:

\begin{itemize}[leftmargin=0.8cm, itemsep=2pt]
    \item \textbf{Chapter 1. Overview of the Internship Organization};
    \item \textbf{Chapter 2. Internship Activities and Assigned Tasks};
    \item \textbf{Chapter 3. Results, Experience, and Evaluation}.
\end{itemize}

These chapter titles are recommendations only and may be adapted to the
actual internship content.

The report shall reflect work actually performed, observed, analyzed, or
learned during the internship. Unrelated theoretical material should not
be included merely to increase report length.

\vspace{0.3cm}

% ================================================================
% X. TABLE OF CONTENTS
% ================================================================

\noindent
\textbf{X.\quad TABLE OF CONTENTS}

\vspace{0.15cm}

The Table of Contents shall be generated automatically using:

\begin{verbatim}
\tableofcontents
\end{verbatim}

Students shall not create or manually edit page numbers in the Table of
Contents.

The document shall be compiled sufficiently to ensure that all page
numbers and entries are correctly updated.

\vspace{0.3cm}

% ================================================================
% XI. LIST OF SYMBOLS AND ABBREVIATIONS
% ================================================================

\noindent
\textbf{XI.\quad LIST OF SYMBOLS AND ABBREVIATIONS}

\vspace{0.15cm}

Symbols and abbreviations that are used repeatedly in the report shall be
listed in the \textbf{List of Symbols and Abbreviations}.

The list may contain the following columns:

\begin{center}
\begin{tabular}{llll}
    No. & Symbol / Abbreviation & Unit & Definition
\end{tabular}
\end{center}

Only symbols and abbreviations actually used in the final report shall be
retained. Example entries contained in the template shall be removed if
they are not relevant to the student's report.

Each abbreviation should normally be defined when it first appears in
the main text.

\vspace{0.3cm}

% ================================================================
% XII. FIGURES, DIAGRAMS, AND PHOTOGRAPHS
% ================================================================

\noindent
\textbf{XII.\quad FIGURES, DIAGRAMS, AND PHOTOGRAPHS}

\vspace{0.15cm}

Figures shall be clear, legible, and directly relevant to the report
content.

Figures shall be inserted using the standard \verb|figure| environment.

Example:

\begin{verbatim}
\begin{figure}[H]
    \centering
    \includegraphics[width=0.8\textwidth]
                    {figures/company_structure.pdf}
    \caption{Organizational structure of the company}
    \label{fig:company-structure}
\end{figure}
\end{verbatim}

Figure captions shall be placed below the figure.

Each referenced figure shall have a unique \verb|\label| and shall be
referenced in the text using, for example:

\begin{verbatim}
Figure~\ref{fig:company-structure}
\end{verbatim}

Students shall not manually type figure numbers.

Figures reproduced or adapted from external sources shall include an
appropriate citation.

Photographs taken during the internship shall be relevant and sufficiently
clear. Students shall not include confidential, restricted, or proprietary
information without permission from the internship organization.

Vector graphics or high-resolution images are recommended whenever
possible.

\vspace{0.3cm}

% ================================================================
% XIII. TABLES
% ================================================================

\noindent
\textbf{XIII.\quad TABLES}

\vspace{0.15cm}

Tables shall be created using suitable \LaTeX{} environments such as
\verb|table|, \verb|tabular|, \verb|tabularx|, or \verb|longtable|.

Table captions shall be placed above the table.

Example:

\begin{verbatim}
\begin{table}[H]
    \centering
    \caption{Internship schedule}
    \label{tab:internship-schedule}

    \begin{tabular}{lll}
        \toprule
        Week & Task & Result \\
        \midrule
        1 & Orientation & Completed \\
        2 & System study & Completed \\
        \bottomrule
    \end{tabular}
\end{table}
\end{verbatim}

Tables shall be referenced using:

\begin{verbatim}
Table~\ref{tab:internship-schedule}
\end{verbatim}

Students shall not manually type table numbers.

Tables should be created directly in \LaTeX{} rather than inserted as
screenshots whenever practical.

\vspace{0.3cm}

% ================================================================
% XIV. EQUATIONS AND TECHNICAL QUANTITIES
% ================================================================

\noindent
\textbf{XIV.\quad EQUATIONS AND TECHNICAL QUANTITIES}

\vspace{0.15cm}

Mathematical expressions shall be typeset directly using \LaTeX{}.
Images or screenshots of equations shall not be used.

Short expressions may be written in inline mathematical mode, for example:

\begin{verbatim}
$R = 0.7\,\Omega$
\end{verbatim}

Important equations shall use a numbered equation environment.

Example:

\begin{verbatim}
\begin{equation}
    P = UI
    \label{eq:power}
\end{equation}
\end{verbatim}

Equations shall be referenced using:

\begin{verbatim}
Equation~\eqref{eq:power}
\end{verbatim}

Equation numbers shall not be manually entered.

SI units shall be used wherever applicable. The \verb|siunitx| package
should be used where appropriate.

Examples:

\begin{verbatim}
\SI{220}{\volt}
\SI{50}{\hertz}
\SI{2.5}{\ampere}
\end{verbatim}

Notation and units shall be consistent throughout the report.

\vspace{0.3cm}

% ================================================================
% XV. SOURCE CODE
% ================================================================

\noindent
\textbf{XV.\quad SOURCE CODE}

\vspace{0.15cm}

If the internship involves programming, PLCs, microcontrollers,
simulation, or software development, relevant source code may be included
in the report.

Source code shall be presented using the \verb|listings| environment or
another environment defined by the official template.

Example:

\begin{verbatim}
\begin{lstlisting}[language=st]
IF Start = TRUE THEN
    Motor := TRUE;
ELSE
    Motor := FALSE;
END_IF;
\end{lstlisting}
\end{verbatim}

Only code necessary for explaining the work should appear in the main
report. Long source-code listings should be moved to an appendix or
submitted as supplementary files if required.

Source code should not normally be inserted as screenshots.

\vspace{0.3cm}

% ================================================================
% XVI. ABSTRACT
% ================================================================

\noindent
\textbf{XVI.\quad ABSTRACT}

\vspace{0.15cm}

The Abstract shall briefly summarize:

\begin{itemize}[leftmargin=0.8cm, itemsep=2pt]
    \item the internship organization;
    \item the objectives of the internship;
    \item the main tasks performed;
    \item the main results, knowledge, and skills obtained.
\end{itemize}

The Abstract should be concise and should not reproduce the full content
of the report chapter by chapter.

\vspace{0.3cm}

% ================================================================
% XVII. REFERENCES AND CITATIONS
% ================================================================

\noindent
\textbf{XVII.\quad REFERENCES AND CITATIONS}

\vspace{0.15cm}

All books, journal papers, conference papers, standards, websites,
datasheets, manuals, technical documents, figures, tables, or other
external sources used in preparing the report shall be properly cited.

References shall be managed through the bibliography system provided with
the official template.

Bibliographic entries shall normally be stored in:

\begin{verbatim}
References.bib
\end{verbatim}

References shall be cited in the report using:

\begin{verbatim}
\cite{reference-key}
\end{verbatim}

Students shall not manually type reference numbers such as ``[1]'',
``[2]'', etc.

Titles of foreign-language publications shall remain in their original
language in the bibliography.

All borrowed ideas, data, figures, tables, equations, methods, or results
shall be properly acknowledged.

\vspace{0.3cm}

% ================================================================
% XVIII. CROSS-REFERENCING
% ================================================================

\noindent
\textbf{XVIII.\quad CROSS-REFERENCING}

\vspace{0.15cm}

Numbered elements shall be cross-referenced using the standard \LaTeX{}
mechanism.

Use the following commands where appropriate:

\begin{verbatim}
\label{...}
\ref{...}
\eqref{...}
\end{verbatim}

For example:

\begin{verbatim}
As shown in Figure~\ref{fig:system}, ...
\end{verbatim}

Students shall not manually type numbers such as ``Figure 2.3'' when
those numbers can be generated automatically.

\vspace{0.3cm}

% ================================================================
% XIX. PAGE NUMBERING
% ================================================================

\noindent
\textbf{XIX.\quad PAGE NUMBERING}

\vspace{0.15cm}

Cover pages and official evaluation forms at the beginning of the report
shall not display page numbers.

The main body of the report shall use Arabic page numbers beginning from
Chapter 1.

Page numbering shall be generated automatically by the official template.
Students shall not manually enter page numbers.

\vspace{0.3cm}

% ================================================================
% XX. PROJECT FILE ORGANIZATION
% ================================================================

\noindent
\textbf{XX.\quad PROJECT FILE ORGANIZATION}

\vspace{0.15cm}

Students should maintain a clear and organized \LaTeX{} project
structure. A recommended structure is:

\begin{verbatim}
InternshipReport/
|
|-- main.tex
|-- References.bib
|
|-- figures/
|   |-- company.pdf
|   |-- system.png
|   `-- ...
|
`-- chapters/
    |-- chapter1.tex
    |-- chapter2.tex
    `-- chapter3.tex
\end{verbatim}

Figures should normally be stored in the \verb|figures| directory
specified by the official template.

File and directory names should preferably contain Latin letters,
numbers, hyphens, or underscores.

Recommended examples include:

\begin{verbatim}
company_structure.pdf
motor_control.png
chapter_01.tex
\end{verbatim}

Spaces and unnecessary special characters in file names should be
avoided where possible.

\vspace{0.3cm}

% ================================================================
% XXI. DOCUMENT COMPILATION AND FINAL CHECK
% ================================================================

\noindent
\textbf{XXI.\quad DOCUMENT COMPILATION AND FINAL CHECK}

\vspace{0.15cm}

Before submission, the complete \LaTeX{} project shall compile
successfully and produce the final PDF report.

The final document shall not contain unresolved references such as:

\begin{verbatim}
??
[?]
\end{verbatim}

Before generating the final PDF, students shall check:

\begin{itemize}[leftmargin=0.8cm, itemsep=2pt]
    \item the Table of Contents;
    \item page numbering;
    \item chapter, section, and subsection numbering;
    \item figure numbering;
    \item table numbering;
    \item equation numbering;
    \item cross-references;
    \item citations;
    \item bibliography;
    \item font display;
    \item mathematical symbols;
    \item image quality;
    \item completeness of all required forms.
\end{itemize}

The document may be compiled using TeX Live, MiKTeX, or another
compatible \LaTeX{} environment. Regardless of the compilation
environment, the resulting PDF shall conform to the official Faculty
template.

\vspace{0.3cm}

% ================================================================
% XXII. CONTENT THAT STUDENTS MAY EDIT
% ================================================================

\noindent
\textbf{XXII.\quad CONTENT THAT STUDENTS MAY EDIT}

\vspace{0.15cm}

Students may edit or replace:

\begin{itemize}[leftmargin=0.8cm, itemsep=2pt]
    \item student name;
    \item class;
    \item internship organization;
    \item supervisor name;
    \item Abstract;
    \item List of Symbols and Abbreviations;
    \item chapter titles and chapter content;
    \item figures and photographs;
    \item tables;
    \item equations;
    \item citations and references.
\end{itemize}

Students shall not modify the main document formatting settings in the
preamble, including the document class, page geometry, line spacing,
heading formatting, table-of-contents formatting, font configuration,
page layout, or cover-page macros unless specifically instructed.

\vspace{0.3cm}

% ================================================================
% XXIII. FINAL SUBMISSION REQUIREMENTS
% ================================================================

\noindent
\textbf{XXIII.\quad FINAL SUBMISSION REQUIREMENTS}

\vspace{0.15cm}

Before submission, the final report shall:

\begin{itemize}[leftmargin=0.8cm, itemsep=2pt]
    \item compile without errors that affect the report content or presentation;
    \item contain all required forms and report sections;
    \item contain correct student, class, internship organization, and supervisor information;
    \item contain an updated Table of Contents, List of Figures, and List of Tables;
    \item contain no unused example content from the template;
    \item contain no placeholder text such as \textit{Title of Chapter 1},
    \textit{Section 1.1}, or other sample entries;
    \item contain all required citations and references;
    \item contain no unresolved cross-references.
\end{itemize}

If the Faculty requires submission of the \LaTeX{} source, the complete
project shall include all \verb|.tex| files, the bibliography database,
figures, and any other files required to reproduce the submitted PDF.

\vspace{0.4cm}

\noindent
\textbf{Important:}
Before submitting the final internship report, students shall remove or
comment out the following command in the main document:

\begin{verbatim}
\input{Regulations.tex}
\end{verbatim}

For example:

\begin{verbatim}
% \input{Regulations.tex}
\end{verbatim}

This removes the guidance pages from the final submitted report while
preserving the regulations file in the \LaTeX{} project for reference.

} % End 13pt content

\clearpage
|.
Students shall remove or comment out this command before submitting the
final internship report.

\vspace{0.4cm}

% ================================================================
% I. GENERAL REQUIREMENTS
% ================================================================

\noindent
\textbf{I.\quad GENERAL REQUIREMENTS}

\vspace{0.15cm}

The Graduation Internship Report shall be prepared using the official
\LaTeX{} template provided by the Faculty. Students shall preserve the
predefined page layout, typography, line spacing, heading styles,
numbering system, figure and table formatting, and bibliography settings
unless otherwise instructed by the Faculty or the supervisor.

The report shall be written in English.

Students are responsible for the academic and technical content of the
report. The official \LaTeX{} template is responsible for standardized
document formatting.

\vspace{0.3cm}

% ================================================================
% II. REPORT STRUCTURE
% ================================================================

\noindent
\textbf{II.\quad REPORT STRUCTURE}

\vspace{0.15cm}

The report shall normally be arranged in the following order:

\begin{enumerate}[label=\arabic*., leftmargin=1.2cm, itemsep=2pt]
    \item Outer Cover;
    \item Comments of the Internship Host Organization;
    \item Supervisor's Evaluation and Graduation Internship Defense Results;
    \item Table of Contents;
    \item List of Symbols and Abbreviations, if applicable;
    \item List of Figures;
    \item List of Tables;
    \item Abstract;
    \item Main Chapters of the Internship Report;
    \item References;
    \item Appendices, if applicable.
\end{enumerate}

The preliminary pages shall not be numbered as chapters. The main report
content shall be organized into consecutively numbered chapters beginning
with Chapter 1.

\vspace{0.3cm}

% ================================================================
% III. PAGE SIZE AND MARGINS
% ================================================================

\noindent
\textbf{III.\quad PAGE SIZE AND MARGINS}

\vspace{0.15cm}

The report shall be formatted on A4 paper with the following margins:

\begin{center}
\begin{tabular}{ll}
    Left margin:   & 3.5 cm \\
    Right margin:  & 2.0 cm \\
    Top margin:    & 2.5 cm \\
    Bottom margin: & 2.5 cm
\end{tabular}
\end{center}

These settings are predefined in the official template and shall not be
modified by students.

\vspace{0.3cm}

% ================================================================
% IV. FONT, FONT SIZE, AND LINE SPACING
% ================================================================

\noindent
\textbf{IV.\quad FONT, FONT SIZE, AND LINE SPACING}

\vspace{0.15cm}

The main text shall use a 13 pt font size with 1.3-line spacing.
A Times-family typeface shall be used in accordance with the official
template.

The text shall be fully justified.

Students shall not use repeated spaces, repeated manual line breaks, or
multiple blank lines to position text. Paragraph and heading spacing shall
be controlled by the template.

\vspace{0.3cm}

% ================================================================
% V. OUTER COVER
% ================================================================

\noindent
\textbf{V.\quad OUTER COVER}

\vspace{0.15cm}

The outer cover shall follow the official format provided in the template
and shall contain the following information:

\begin{itemize}[leftmargin=0.8cm, itemsep=2pt]
    \item Thai Nguyen University of Technology;
    \item Faculty of International Training;
    \item Engineering and Technology Department;
    \item University logo and Faculty logo;
    \item the title \textbf{INTERNSHIP REPORT};
    \item Internship Organization;
    \item Author;
    \item Class;
    \item Supervisor;
    \item place and year of completion.
\end{itemize}

Students shall only replace the required information fields and shall not
modify the cover-page layout unless instructed otherwise.

\vspace{0.3cm}

% ================================================================
% VI. COMMENTS OF THE INTERNSHIP HOST ORGANIZATION
% ================================================================

\noindent
\textbf{VI.\quad COMMENTS OF THE INTERNSHIP HOST ORGANIZATION}

\vspace{0.15cm}

The page entitled \textbf{COMMENTS OF THE INTERNSHIP HOST ORGANIZATION}
is reserved for the internship host organization.

It includes:

\begin{itemize}[leftmargin=0.8cm, itemsep=2pt]
    \item the student's full name;
    \item class;
    \item comments from the host organization;
    \item score;
    \item date;
    \item signature and official seal.
\end{itemize}

Students shall not alter or complete official evaluation fields unless
specifically authorized.

\vspace{0.3cm}

% ================================================================
% VII. SUPERVISOR'S EVALUATION AND DEFENSE RESULTS
% ================================================================

\noindent
\textbf{VII.\quad SUPERVISOR'S EVALUATION AND DEFENSE RESULTS}

\vspace{0.15cm}

The section entitled
\textbf{SUPERVISOR'S EVALUATION AND GRADUATION INTERNSHIP DEFENSE RESULTS}
is reserved for the supervisor and defense committee.

The evaluation may include:

\begin{itemize}[leftmargin=0.8cm, itemsep=2pt]
    \item Internship Progress and Student's Attitude;
    \item Report Content;
    \item Report Presentation;
    \item Other Comments;
    \item Supervisor's Opinion;
    \item Defense Result;
    \item Score.
\end{itemize}

Students shall not enter evaluation scores or defense results before they
are officially determined.

\vspace{0.3cm}

% ================================================================
% VIII. CHAPTERS, SECTIONS, AND SUBSECTIONS
% ================================================================

\noindent
\textbf{VIII.\quad CHAPTERS, SECTIONS, AND SUBSECTIONS}

\vspace{0.15cm}

The report structure shall be generated using the standard \LaTeX{}
commands:

\begin{verbatim}
\chapter{...}
\section{...}
\subsection{...}
\end{verbatim}

Chapter, section, and subsection numbers shall be generated automatically
by \LaTeX{}.

Students shall not manually type chapter or section numbers and shall not
reproduce heading styles using commands such as \verb|\textbf|,
\verb|\fontsize|, or \verb|\centering| instead of the appropriate
structural command.

Excessive subdivision should be avoided. In most cases, the hierarchy
Chapter--Section--Subsection is sufficient.

\vspace{0.3cm}

% ================================================================
% IX. ORGANIZATION OF THE MAIN CONTENT
% ================================================================

\noindent
\textbf{IX.\quad ORGANIZATION OF THE MAIN CONTENT}

\vspace{0.15cm}

The specific content of the report shall depend on the internship
organization and the assigned tasks. A typical structure may include:

\begin{itemize}[leftmargin=0.8cm, itemsep=2pt]
    \item \textbf{Chapter 1. Overview of the Internship Organization};
    \item \textbf{Chapter 2. Internship Activities and Assigned Tasks};
    \item \textbf{Chapter 3. Results, Experience, and Evaluation}.
\end{itemize}

These chapter titles are recommendations only and may be adapted to the
actual internship content.

The report shall reflect work actually performed, observed, analyzed, or
learned during the internship. Unrelated theoretical material should not
be included merely to increase report length.

\vspace{0.3cm}

% ================================================================
% X. TABLE OF CONTENTS
% ================================================================

\noindent
\textbf{X.\quad TABLE OF CONTENTS}

\vspace{0.15cm}

The Table of Contents shall be generated automatically using:

\begin{verbatim}
\tableofcontents
\end{verbatim}

Students shall not create or manually edit page numbers in the Table of
Contents.

The document shall be compiled sufficiently to ensure that all page
numbers and entries are correctly updated.

\vspace{0.3cm}

% ================================================================
% XI. LIST OF SYMBOLS AND ABBREVIATIONS
% ================================================================

\noindent
\textbf{XI.\quad LIST OF SYMBOLS AND ABBREVIATIONS}

\vspace{0.15cm}

Symbols and abbreviations that are used repeatedly in the report shall be
listed in the \textbf{List of Symbols and Abbreviations}.

The list may contain the following columns:

\begin{center}
\begin{tabular}{llll}
    No. & Symbol / Abbreviation & Unit & Definition
\end{tabular}
\end{center}

Only symbols and abbreviations actually used in the final report shall be
retained. Example entries contained in the template shall be removed if
they are not relevant to the student's report.

Each abbreviation should normally be defined when it first appears in
the main text.

\vspace{0.3cm}

% ================================================================
% XII. FIGURES, DIAGRAMS, AND PHOTOGRAPHS
% ================================================================

\noindent
\textbf{XII.\quad FIGURES, DIAGRAMS, AND PHOTOGRAPHS}

\vspace{0.15cm}

Figures shall be clear, legible, and directly relevant to the report
content.

Figures shall be inserted using the standard \verb|figure| environment.

Example:

\begin{verbatim}
\begin{figure}[H]
    \centering
    \includegraphics[width=0.8\textwidth]
                    {figures/company_structure.pdf}
    \caption{Organizational structure of the company}
    \label{fig:company-structure}
\end{figure}
\end{verbatim}

Figure captions shall be placed below the figure.

Each referenced figure shall have a unique \verb|\label| and shall be
referenced in the text using, for example:

\begin{verbatim}
Figure~\ref{fig:company-structure}
\end{verbatim}

Students shall not manually type figure numbers.

Figures reproduced or adapted from external sources shall include an
appropriate citation.

Photographs taken during the internship shall be relevant and sufficiently
clear. Students shall not include confidential, restricted, or proprietary
information without permission from the internship organization.

Vector graphics or high-resolution images are recommended whenever
possible.

\vspace{0.3cm}

% ================================================================
% XIII. TABLES
% ================================================================

\noindent
\textbf{XIII.\quad TABLES}

\vspace{0.15cm}

Tables shall be created using suitable \LaTeX{} environments such as
\verb|table|, \verb|tabular|, \verb|tabularx|, or \verb|longtable|.

Table captions shall be placed above the table.

Example:

\begin{verbatim}
\begin{table}[H]
    \centering
    \caption{Internship schedule}
    \label{tab:internship-schedule}

    \begin{tabular}{lll}
        \toprule
        Week & Task & Result \\
        \midrule
        1 & Orientation & Completed \\
        2 & System study & Completed \\
        \bottomrule
    \end{tabular}
\end{table}
\end{verbatim}

Tables shall be referenced using:

\begin{verbatim}
Table~\ref{tab:internship-schedule}
\end{verbatim}

Students shall not manually type table numbers.

Tables should be created directly in \LaTeX{} rather than inserted as
screenshots whenever practical.

\vspace{0.3cm}

% ================================================================
% XIV. EQUATIONS AND TECHNICAL QUANTITIES
% ================================================================

\noindent
\textbf{XIV.\quad EQUATIONS AND TECHNICAL QUANTITIES}

\vspace{0.15cm}

Mathematical expressions shall be typeset directly using \LaTeX{}.
Images or screenshots of equations shall not be used.

Short expressions may be written in inline mathematical mode, for example:

\begin{verbatim}
$R = 0.7\,\Omega$
\end{verbatim}

Important equations shall use a numbered equation environment.

Example:

\begin{verbatim}
\begin{equation}
    P = UI
    \label{eq:power}
\end{equation}
\end{verbatim}

Equations shall be referenced using:

\begin{verbatim}
Equation~\eqref{eq:power}
\end{verbatim}

Equation numbers shall not be manually entered.

SI units shall be used wherever applicable. The \verb|siunitx| package
should be used where appropriate.

Examples:

\begin{verbatim}
\SI{220}{\volt}
\SI{50}{\hertz}
\SI{2.5}{\ampere}
\end{verbatim}

Notation and units shall be consistent throughout the report.

\vspace{0.3cm}

% ================================================================
% XV. SOURCE CODE
% ================================================================

\noindent
\textbf{XV.\quad SOURCE CODE}

\vspace{0.15cm}

If the internship involves programming, PLCs, microcontrollers,
simulation, or software development, relevant source code may be included
in the report.

Source code shall be presented using the \verb|listings| environment or
another environment defined by the official template.

Example:

\begin{verbatim}
\begin{lstlisting}[language=st]
IF Start = TRUE THEN
    Motor := TRUE;
ELSE
    Motor := FALSE;
END_IF;
\end{lstlisting}
\end{verbatim}

Only code necessary for explaining the work should appear in the main
report. Long source-code listings should be moved to an appendix or
submitted as supplementary files if required.

Source code should not normally be inserted as screenshots.

\vspace{0.3cm}

% ================================================================
% XVI. ABSTRACT
% ================================================================

\noindent
\textbf{XVI.\quad ABSTRACT}

\vspace{0.15cm}

The Abstract shall briefly summarize:

\begin{itemize}[leftmargin=0.8cm, itemsep=2pt]
    \item the internship organization;
    \item the objectives of the internship;
    \item the main tasks performed;
    \item the main results, knowledge, and skills obtained.
\end{itemize}

The Abstract should be concise and should not reproduce the full content
of the report chapter by chapter.

\vspace{0.3cm}

% ================================================================
% XVII. REFERENCES AND CITATIONS
% ================================================================

\noindent
\textbf{XVII.\quad REFERENCES AND CITATIONS}

\vspace{0.15cm}

All books, journal papers, conference papers, standards, websites,
datasheets, manuals, technical documents, figures, tables, or other
external sources used in preparing the report shall be properly cited.

References shall be managed through the bibliography system provided with
the official template.

Bibliographic entries shall normally be stored in:

\begin{verbatim}
References.bib
\end{verbatim}

References shall be cited in the report using:

\begin{verbatim}
\cite{reference-key}
\end{verbatim}

Students shall not manually type reference numbers such as ``[1]'',
``[2]'', etc.

Titles of foreign-language publications shall remain in their original
language in the bibliography.

All borrowed ideas, data, figures, tables, equations, methods, or results
shall be properly acknowledged.

\vspace{0.3cm}

% ================================================================
% XVIII. CROSS-REFERENCING
% ================================================================

\noindent
\textbf{XVIII.\quad CROSS-REFERENCING}

\vspace{0.15cm}

Numbered elements shall be cross-referenced using the standard \LaTeX{}
mechanism.

Use the following commands where appropriate:

\begin{verbatim}
\label{...}
\ref{...}
\eqref{...}
\end{verbatim}

For example:

\begin{verbatim}
As shown in Figure~\ref{fig:system}, ...
\end{verbatim}

Students shall not manually type numbers such as ``Figure 2.3'' when
those numbers can be generated automatically.

\vspace{0.3cm}

% ================================================================
% XIX. PAGE NUMBERING
% ================================================================

\noindent
\textbf{XIX.\quad PAGE NUMBERING}

\vspace{0.15cm}

Cover pages and official evaluation forms at the beginning of the report
shall not display page numbers.

The main body of the report shall use Arabic page numbers beginning from
Chapter 1.

Page numbering shall be generated automatically by the official template.
Students shall not manually enter page numbers.

\vspace{0.3cm}

% ================================================================
% XX. PROJECT FILE ORGANIZATION
% ================================================================

\noindent
\textbf{XX.\quad PROJECT FILE ORGANIZATION}

\vspace{0.15cm}

Students should maintain a clear and organized \LaTeX{} project
structure. A recommended structure is:

\begin{verbatim}
InternshipReport/
|
|-- main.tex
|-- References.bib
|
|-- figures/
|   |-- company.pdf
|   |-- system.png
|   `-- ...
|
`-- chapters/
    |-- chapter1.tex
    |-- chapter2.tex
    `-- chapter3.tex
\end{verbatim}

Figures should normally be stored in the \verb|figures| directory
specified by the official template.

File and directory names should preferably contain Latin letters,
numbers, hyphens, or underscores.

Recommended examples include:

\begin{verbatim}
company_structure.pdf
motor_control.png
chapter_01.tex
\end{verbatim}

Spaces and unnecessary special characters in file names should be
avoided where possible.

\vspace{0.3cm}

% ================================================================
% XXI. DOCUMENT COMPILATION AND FINAL CHECK
% ================================================================

\noindent
\textbf{XXI.\quad DOCUMENT COMPILATION AND FINAL CHECK}

\vspace{0.15cm}

Before submission, the complete \LaTeX{} project shall compile
successfully and produce the final PDF report.

The final document shall not contain unresolved references such as:

\begin{verbatim}
??
[?]
\end{verbatim}

Before generating the final PDF, students shall check:

\begin{itemize}[leftmargin=0.8cm, itemsep=2pt]
    \item the Table of Contents;
    \item page numbering;
    \item chapter, section, and subsection numbering;
    \item figure numbering;
    \item table numbering;
    \item equation numbering;
    \item cross-references;
    \item citations;
    \item bibliography;
    \item font display;
    \item mathematical symbols;
    \item image quality;
    \item completeness of all required forms.
\end{itemize}

The document may be compiled using TeX Live, MiKTeX, or another
compatible \LaTeX{} environment. Regardless of the compilation
environment, the resulting PDF shall conform to the official Faculty
template.

\vspace{0.3cm}

% ================================================================
% XXII. CONTENT THAT STUDENTS MAY EDIT
% ================================================================

\noindent
\textbf{XXII.\quad CONTENT THAT STUDENTS MAY EDIT}

\vspace{0.15cm}

Students may edit or replace:

\begin{itemize}[leftmargin=0.8cm, itemsep=2pt]
    \item student name;
    \item class;
    \item internship organization;
    \item supervisor name;
    \item Abstract;
    \item List of Symbols and Abbreviations;
    \item chapter titles and chapter content;
    \item figures and photographs;
    \item tables;
    \item equations;
    \item citations and references.
\end{itemize}

Students shall not modify the main document formatting settings in the
preamble, including the document class, page geometry, line spacing,
heading formatting, table-of-contents formatting, font configuration,
page layout, or cover-page macros unless specifically instructed.

\vspace{0.3cm}

% ================================================================
% XXIII. FINAL SUBMISSION REQUIREMENTS
% ================================================================

\noindent
\textbf{XXIII.\quad FINAL SUBMISSION REQUIREMENTS}

\vspace{0.15cm}

Before submission, the final report shall:

\begin{itemize}[leftmargin=0.8cm, itemsep=2pt]
    \item compile without errors that affect the report content or presentation;
    \item contain all required forms and report sections;
    \item contain correct student, class, internship organization, and supervisor information;
    \item contain an updated Table of Contents, List of Figures, and List of Tables;
    \item contain no unused example content from the template;
    \item contain no placeholder text such as \textit{Title of Chapter 1},
    \textit{Section 1.1}, or other sample entries;
    \item contain all required citations and references;
    \item contain no unresolved cross-references.
\end{itemize}

If the Faculty requires submission of the \LaTeX{} source, the complete
project shall include all \verb|.tex| files, the bibliography database,
figures, and any other files required to reproduce the submitted PDF.

\vspace{0.4cm}

\noindent
\textbf{Important:}
Before submitting the final internship report, students shall remove or
comment out the following command in the main document:

\begin{verbatim}

% =================================================================
% REGULATIONS FOR PREPARING THE GRADUATION INTERNSHIP REPORT
% USING LATEX
% =================================================================
%
% This file is intended to be included from the main document with:
%
%     \input{Regulations.tex}
%
% IMPORTANT:
% This section is provided for guidance only. Students shall remove
% or comment out the \input{Regulations.tex} command before submitting
% the final internship report.
%
% =================================================================

\clearpage
\thispagestyle{empty}

\begin{center}
    {\fontsize{16pt}{20pt}\selectfont\bfseries
    SPECIFIC REGULATIONS FOR PREPARING\\
    GRADUATION INTERNSHIP REPORTS USING \LaTeX}
\end{center}

\vspace{0.5cm}

{\fontsize{13pt}{17pt}\selectfont
\setstretch{1.3}
\justifying

\noindent
\textbf{Note:}
This section is provided for guidance only and is included in the main
document using the command \verb|\input{Regulations.tex}|.
Students shall remove or comment out this command before submitting the
final internship report.

\vspace{0.4cm}

% ================================================================
% I. GENERAL REQUIREMENTS
% ================================================================

\noindent
\textbf{I.\quad GENERAL REQUIREMENTS}

\vspace{0.15cm}

The Graduation Internship Report shall be prepared using the official
\LaTeX{} template provided by the Faculty. Students shall preserve the
predefined page layout, typography, line spacing, heading styles,
numbering system, figure and table formatting, and bibliography settings
unless otherwise instructed by the Faculty or the supervisor.

The report shall be written in English.

Students are responsible for the academic and technical content of the
report. The official \LaTeX{} template is responsible for standardized
document formatting.

\vspace{0.3cm}

% ================================================================
% II. REPORT STRUCTURE
% ================================================================

\noindent
\textbf{II.\quad REPORT STRUCTURE}

\vspace{0.15cm}

The report shall normally be arranged in the following order:

\begin{enumerate}[label=\arabic*., leftmargin=1.2cm, itemsep=2pt]
    \item Outer Cover;
    \item Comments of the Internship Host Organization;
    \item Supervisor's Evaluation and Graduation Internship Defense Results;
    \item Table of Contents;
    \item List of Symbols and Abbreviations, if applicable;
    \item List of Figures;
    \item List of Tables;
    \item Abstract;
    \item Main Chapters of the Internship Report;
    \item References;
    \item Appendices, if applicable.
\end{enumerate}

The preliminary pages shall not be numbered as chapters. The main report
content shall be organized into consecutively numbered chapters beginning
with Chapter 1.

\vspace{0.3cm}

% ================================================================
% III. PAGE SIZE AND MARGINS
% ================================================================

\noindent
\textbf{III.\quad PAGE SIZE AND MARGINS}

\vspace{0.15cm}

The report shall be formatted on A4 paper with the following margins:

\begin{center}
\begin{tabular}{ll}
    Left margin:   & 3.5 cm \\
    Right margin:  & 2.0 cm \\
    Top margin:    & 2.5 cm \\
    Bottom margin: & 2.5 cm
\end{tabular}
\end{center}

These settings are predefined in the official template and shall not be
modified by students.

\vspace{0.3cm}

% ================================================================
% IV. FONT, FONT SIZE, AND LINE SPACING
% ================================================================

\noindent
\textbf{IV.\quad FONT, FONT SIZE, AND LINE SPACING}

\vspace{0.15cm}

The main text shall use a 13 pt font size with 1.3-line spacing.
A Times-family typeface shall be used in accordance with the official
template.

The text shall be fully justified.

Students shall not use repeated spaces, repeated manual line breaks, or
multiple blank lines to position text. Paragraph and heading spacing shall
be controlled by the template.

\vspace{0.3cm}

% ================================================================
% V. OUTER COVER
% ================================================================

\noindent
\textbf{V.\quad OUTER COVER}

\vspace{0.15cm}

The outer cover shall follow the official format provided in the template
and shall contain the following information:

\begin{itemize}[leftmargin=0.8cm, itemsep=2pt]
    \item Thai Nguyen University of Technology;
    \item Faculty of International Training;
    \item Engineering and Technology Department;
    \item University logo and Faculty logo;
    \item the title \textbf{INTERNSHIP REPORT};
    \item Internship Organization;
    \item Author;
    \item Class;
    \item Supervisor;
    \item place and year of completion.
\end{itemize}

Students shall only replace the required information fields and shall not
modify the cover-page layout unless instructed otherwise.

\vspace{0.3cm}

% ================================================================
% VI. COMMENTS OF THE INTERNSHIP HOST ORGANIZATION
% ================================================================

\noindent
\textbf{VI.\quad COMMENTS OF THE INTERNSHIP HOST ORGANIZATION}

\vspace{0.15cm}

The page entitled \textbf{COMMENTS OF THE INTERNSHIP HOST ORGANIZATION}
is reserved for the internship host organization.

It includes:

\begin{itemize}[leftmargin=0.8cm, itemsep=2pt]
    \item the student's full name;
    \item class;
    \item comments from the host organization;
    \item score;
    \item date;
    \item signature and official seal.
\end{itemize}

Students shall not alter or complete official evaluation fields unless
specifically authorized.

\vspace{0.3cm}

% ================================================================
% VII. SUPERVISOR'S EVALUATION AND DEFENSE RESULTS
% ================================================================

\noindent
\textbf{VII.\quad SUPERVISOR'S EVALUATION AND DEFENSE RESULTS}

\vspace{0.15cm}

The section entitled
\textbf{SUPERVISOR'S EVALUATION AND GRADUATION INTERNSHIP DEFENSE RESULTS}
is reserved for the supervisor and defense committee.

The evaluation may include:

\begin{itemize}[leftmargin=0.8cm, itemsep=2pt]
    \item Internship Progress and Student's Attitude;
    \item Report Content;
    \item Report Presentation;
    \item Other Comments;
    \item Supervisor's Opinion;
    \item Defense Result;
    \item Score.
\end{itemize}

Students shall not enter evaluation scores or defense results before they
are officially determined.

\vspace{0.3cm}

% ================================================================
% VIII. CHAPTERS, SECTIONS, AND SUBSECTIONS
% ================================================================

\noindent
\textbf{VIII.\quad CHAPTERS, SECTIONS, AND SUBSECTIONS}

\vspace{0.15cm}

The report structure shall be generated using the standard \LaTeX{}
commands:

\begin{verbatim}
\chapter{...}
\section{...}
\subsection{...}
\end{verbatim}

Chapter, section, and subsection numbers shall be generated automatically
by \LaTeX{}.

Students shall not manually type chapter or section numbers and shall not
reproduce heading styles using commands such as \verb|\textbf|,
\verb|\fontsize|, or \verb|\centering| instead of the appropriate
structural command.

Excessive subdivision should be avoided. In most cases, the hierarchy
Chapter--Section--Subsection is sufficient.

\vspace{0.3cm}

% ================================================================
% IX. ORGANIZATION OF THE MAIN CONTENT
% ================================================================

\noindent
\textbf{IX.\quad ORGANIZATION OF THE MAIN CONTENT}

\vspace{0.15cm}

The specific content of the report shall depend on the internship
organization and the assigned tasks. A typical structure may include:

\begin{itemize}[leftmargin=0.8cm, itemsep=2pt]
    \item \textbf{Chapter 1. Overview of the Internship Organization};
    \item \textbf{Chapter 2. Internship Activities and Assigned Tasks};
    \item \textbf{Chapter 3. Results, Experience, and Evaluation}.
\end{itemize}

These chapter titles are recommendations only and may be adapted to the
actual internship content.

The report shall reflect work actually performed, observed, analyzed, or
learned during the internship. Unrelated theoretical material should not
be included merely to increase report length.

\vspace{0.3cm}

% ================================================================
% X. TABLE OF CONTENTS
% ================================================================

\noindent
\textbf{X.\quad TABLE OF CONTENTS}

\vspace{0.15cm}

The Table of Contents shall be generated automatically using:

\begin{verbatim}
\tableofcontents
\end{verbatim}

Students shall not create or manually edit page numbers in the Table of
Contents.

The document shall be compiled sufficiently to ensure that all page
numbers and entries are correctly updated.

\vspace{0.3cm}

% ================================================================
% XI. LIST OF SYMBOLS AND ABBREVIATIONS
% ================================================================

\noindent
\textbf{XI.\quad LIST OF SYMBOLS AND ABBREVIATIONS}

\vspace{0.15cm}

Symbols and abbreviations that are used repeatedly in the report shall be
listed in the \textbf{List of Symbols and Abbreviations}.

The list may contain the following columns:

\begin{center}
\begin{tabular}{llll}
    No. & Symbol / Abbreviation & Unit & Definition
\end{tabular}
\end{center}

Only symbols and abbreviations actually used in the final report shall be
retained. Example entries contained in the template shall be removed if
they are not relevant to the student's report.

Each abbreviation should normally be defined when it first appears in
the main text.

\vspace{0.3cm}

% ================================================================
% XII. FIGURES, DIAGRAMS, AND PHOTOGRAPHS
% ================================================================

\noindent
\textbf{XII.\quad FIGURES, DIAGRAMS, AND PHOTOGRAPHS}

\vspace{0.15cm}

Figures shall be clear, legible, and directly relevant to the report
content.

Figures shall be inserted using the standard \verb|figure| environment.

Example:

\begin{verbatim}
\begin{figure}[H]
    \centering
    \includegraphics[width=0.8\textwidth]
                    {figures/company_structure.pdf}
    \caption{Organizational structure of the company}
    \label{fig:company-structure}
\end{figure}
\end{verbatim}

Figure captions shall be placed below the figure.

Each referenced figure shall have a unique \verb|\label| and shall be
referenced in the text using, for example:

\begin{verbatim}
Figure~\ref{fig:company-structure}
\end{verbatim}

Students shall not manually type figure numbers.

Figures reproduced or adapted from external sources shall include an
appropriate citation.

Photographs taken during the internship shall be relevant and sufficiently
clear. Students shall not include confidential, restricted, or proprietary
information without permission from the internship organization.

Vector graphics or high-resolution images are recommended whenever
possible.

\vspace{0.3cm}

% ================================================================
% XIII. TABLES
% ================================================================

\noindent
\textbf{XIII.\quad TABLES}

\vspace{0.15cm}

Tables shall be created using suitable \LaTeX{} environments such as
\verb|table|, \verb|tabular|, \verb|tabularx|, or \verb|longtable|.

Table captions shall be placed above the table.

Example:

\begin{verbatim}
\begin{table}[H]
    \centering
    \caption{Internship schedule}
    \label{tab:internship-schedule}

    \begin{tabular}{lll}
        \toprule
        Week & Task & Result \\
        \midrule
        1 & Orientation & Completed \\
        2 & System study & Completed \\
        \bottomrule
    \end{tabular}
\end{table}
\end{verbatim}

Tables shall be referenced using:

\begin{verbatim}
Table~\ref{tab:internship-schedule}
\end{verbatim}

Students shall not manually type table numbers.

Tables should be created directly in \LaTeX{} rather than inserted as
screenshots whenever practical.

\vspace{0.3cm}

% ================================================================
% XIV. EQUATIONS AND TECHNICAL QUANTITIES
% ================================================================

\noindent
\textbf{XIV.\quad EQUATIONS AND TECHNICAL QUANTITIES}

\vspace{0.15cm}

Mathematical expressions shall be typeset directly using \LaTeX{}.
Images or screenshots of equations shall not be used.

Short expressions may be written in inline mathematical mode, for example:

\begin{verbatim}
$R = 0.7\,\Omega$
\end{verbatim}

Important equations shall use a numbered equation environment.

Example:

\begin{verbatim}
\begin{equation}
    P = UI
    \label{eq:power}
\end{equation}
\end{verbatim}

Equations shall be referenced using:

\begin{verbatim}
Equation~\eqref{eq:power}
\end{verbatim}

Equation numbers shall not be manually entered.

SI units shall be used wherever applicable. The \verb|siunitx| package
should be used where appropriate.

Examples:

\begin{verbatim}
\SI{220}{\volt}
\SI{50}{\hertz}
\SI{2.5}{\ampere}
\end{verbatim}

Notation and units shall be consistent throughout the report.

\vspace{0.3cm}

% ================================================================
% XV. SOURCE CODE
% ================================================================

\noindent
\textbf{XV.\quad SOURCE CODE}

\vspace{0.15cm}

If the internship involves programming, PLCs, microcontrollers,
simulation, or software development, relevant source code may be included
in the report.

Source code shall be presented using the \verb|listings| environment or
another environment defined by the official template.

Example:

\begin{verbatim}
\begin{lstlisting}[language=st]
IF Start = TRUE THEN
    Motor := TRUE;
ELSE
    Motor := FALSE;
END_IF;
\end{lstlisting}
\end{verbatim}

Only code necessary for explaining the work should appear in the main
report. Long source-code listings should be moved to an appendix or
submitted as supplementary files if required.

Source code should not normally be inserted as screenshots.

\vspace{0.3cm}

% ================================================================
% XVI. ABSTRACT
% ================================================================

\noindent
\textbf{XVI.\quad ABSTRACT}

\vspace{0.15cm}

The Abstract shall briefly summarize:

\begin{itemize}[leftmargin=0.8cm, itemsep=2pt]
    \item the internship organization;
    \item the objectives of the internship;
    \item the main tasks performed;
    \item the main results, knowledge, and skills obtained.
\end{itemize}

The Abstract should be concise and should not reproduce the full content
of the report chapter by chapter.

\vspace{0.3cm}

% ================================================================
% XVII. REFERENCES AND CITATIONS
% ================================================================

\noindent
\textbf{XVII.\quad REFERENCES AND CITATIONS}

\vspace{0.15cm}

All books, journal papers, conference papers, standards, websites,
datasheets, manuals, technical documents, figures, tables, or other
external sources used in preparing the report shall be properly cited.

References shall be managed through the bibliography system provided with
the official template.

Bibliographic entries shall normally be stored in:

\begin{verbatim}
References.bib
\end{verbatim}

References shall be cited in the report using:

\begin{verbatim}
\cite{reference-key}
\end{verbatim}

Students shall not manually type reference numbers such as ``[1]'',
``[2]'', etc.

Titles of foreign-language publications shall remain in their original
language in the bibliography.

All borrowed ideas, data, figures, tables, equations, methods, or results
shall be properly acknowledged.

\vspace{0.3cm}

% ================================================================
% XVIII. CROSS-REFERENCING
% ================================================================

\noindent
\textbf{XVIII.\quad CROSS-REFERENCING}

\vspace{0.15cm}

Numbered elements shall be cross-referenced using the standard \LaTeX{}
mechanism.

Use the following commands where appropriate:

\begin{verbatim}
\label{...}
\ref{...}
\eqref{...}
\end{verbatim}

For example:

\begin{verbatim}
As shown in Figure~\ref{fig:system}, ...
\end{verbatim}

Students shall not manually type numbers such as ``Figure 2.3'' when
those numbers can be generated automatically.

\vspace{0.3cm}

% ================================================================
% XIX. PAGE NUMBERING
% ================================================================

\noindent
\textbf{XIX.\quad PAGE NUMBERING}

\vspace{0.15cm}

Cover pages and official evaluation forms at the beginning of the report
shall not display page numbers.

The main body of the report shall use Arabic page numbers beginning from
Chapter 1.

Page numbering shall be generated automatically by the official template.
Students shall not manually enter page numbers.

\vspace{0.3cm}

% ================================================================
% XX. PROJECT FILE ORGANIZATION
% ================================================================

\noindent
\textbf{XX.\quad PROJECT FILE ORGANIZATION}

\vspace{0.15cm}

Students should maintain a clear and organized \LaTeX{} project
structure. A recommended structure is:

\begin{verbatim}
InternshipReport/
|
|-- main.tex
|-- References.bib
|
|-- figures/
|   |-- company.pdf
|   |-- system.png
|   `-- ...
|
`-- chapters/
    |-- chapter1.tex
    |-- chapter2.tex
    `-- chapter3.tex
\end{verbatim}

Figures should normally be stored in the \verb|figures| directory
specified by the official template.

File and directory names should preferably contain Latin letters,
numbers, hyphens, or underscores.

Recommended examples include:

\begin{verbatim}
company_structure.pdf
motor_control.png
chapter_01.tex
\end{verbatim}

Spaces and unnecessary special characters in file names should be
avoided where possible.

\vspace{0.3cm}

% ================================================================
% XXI. DOCUMENT COMPILATION AND FINAL CHECK
% ================================================================

\noindent
\textbf{XXI.\quad DOCUMENT COMPILATION AND FINAL CHECK}

\vspace{0.15cm}

Before submission, the complete \LaTeX{} project shall compile
successfully and produce the final PDF report.

The final document shall not contain unresolved references such as:

\begin{verbatim}
??
[?]
\end{verbatim}

Before generating the final PDF, students shall check:

\begin{itemize}[leftmargin=0.8cm, itemsep=2pt]
    \item the Table of Contents;
    \item page numbering;
    \item chapter, section, and subsection numbering;
    \item figure numbering;
    \item table numbering;
    \item equation numbering;
    \item cross-references;
    \item citations;
    \item bibliography;
    \item font display;
    \item mathematical symbols;
    \item image quality;
    \item completeness of all required forms.
\end{itemize}

The document may be compiled using TeX Live, MiKTeX, or another
compatible \LaTeX{} environment. Regardless of the compilation
environment, the resulting PDF shall conform to the official Faculty
template.

\vspace{0.3cm}

% ================================================================
% XXII. CONTENT THAT STUDENTS MAY EDIT
% ================================================================

\noindent
\textbf{XXII.\quad CONTENT THAT STUDENTS MAY EDIT}

\vspace{0.15cm}

Students may edit or replace:

\begin{itemize}[leftmargin=0.8cm, itemsep=2pt]
    \item student name;
    \item class;
    \item internship organization;
    \item supervisor name;
    \item Abstract;
    \item List of Symbols and Abbreviations;
    \item chapter titles and chapter content;
    \item figures and photographs;
    \item tables;
    \item equations;
    \item citations and references.
\end{itemize}

Students shall not modify the main document formatting settings in the
preamble, including the document class, page geometry, line spacing,
heading formatting, table-of-contents formatting, font configuration,
page layout, or cover-page macros unless specifically instructed.

\vspace{0.3cm}

% ================================================================
% XXIII. FINAL SUBMISSION REQUIREMENTS
% ================================================================

\noindent
\textbf{XXIII.\quad FINAL SUBMISSION REQUIREMENTS}

\vspace{0.15cm}

Before submission, the final report shall:

\begin{itemize}[leftmargin=0.8cm, itemsep=2pt]
    \item compile without errors that affect the report content or presentation;
    \item contain all required forms and report sections;
    \item contain correct student, class, internship organization, and supervisor information;
    \item contain an updated Table of Contents, List of Figures, and List of Tables;
    \item contain no unused example content from the template;
    \item contain no placeholder text such as \textit{Title of Chapter 1},
    \textit{Section 1.1}, or other sample entries;
    \item contain all required citations and references;
    \item contain no unresolved cross-references.
\end{itemize}

If the Faculty requires submission of the \LaTeX{} source, the complete
project shall include all \verb|.tex| files, the bibliography database,
figures, and any other files required to reproduce the submitted PDF.

\vspace{0.4cm}

\noindent
\textbf{Important:}
Before submitting the final internship report, students shall remove or
comment out the following command in the main document:

\begin{verbatim}
\input{Regulations.tex}
\end{verbatim}

For example:

\begin{verbatim}
% \input{Regulations.tex}
\end{verbatim}

This removes the guidance pages from the final submitted report while
preserving the regulations file in the \LaTeX{} project for reference.

} % End 13pt content

\clearpage

\end{verbatim}

For example:

\begin{verbatim}
% 
% =================================================================
% REGULATIONS FOR PREPARING THE GRADUATION INTERNSHIP REPORT
% USING LATEX
% =================================================================
%
% This file is intended to be included from the main document with:
%
%     \input{Regulations.tex}
%
% IMPORTANT:
% This section is provided for guidance only. Students shall remove
% or comment out the \input{Regulations.tex} command before submitting
% the final internship report.
%
% =================================================================

\clearpage
\thispagestyle{empty}

\begin{center}
    {\fontsize{16pt}{20pt}\selectfont\bfseries
    SPECIFIC REGULATIONS FOR PREPARING\\
    GRADUATION INTERNSHIP REPORTS USING \LaTeX}
\end{center}

\vspace{0.5cm}

{\fontsize{13pt}{17pt}\selectfont
\setstretch{1.3}
\justifying

\noindent
\textbf{Note:}
This section is provided for guidance only and is included in the main
document using the command \verb|\input{Regulations.tex}|.
Students shall remove or comment out this command before submitting the
final internship report.

\vspace{0.4cm}

% ================================================================
% I. GENERAL REQUIREMENTS
% ================================================================

\noindent
\textbf{I.\quad GENERAL REQUIREMENTS}

\vspace{0.15cm}

The Graduation Internship Report shall be prepared using the official
\LaTeX{} template provided by the Faculty. Students shall preserve the
predefined page layout, typography, line spacing, heading styles,
numbering system, figure and table formatting, and bibliography settings
unless otherwise instructed by the Faculty or the supervisor.

The report shall be written in English.

Students are responsible for the academic and technical content of the
report. The official \LaTeX{} template is responsible for standardized
document formatting.

\vspace{0.3cm}

% ================================================================
% II. REPORT STRUCTURE
% ================================================================

\noindent
\textbf{II.\quad REPORT STRUCTURE}

\vspace{0.15cm}

The report shall normally be arranged in the following order:

\begin{enumerate}[label=\arabic*., leftmargin=1.2cm, itemsep=2pt]
    \item Outer Cover;
    \item Comments of the Internship Host Organization;
    \item Supervisor's Evaluation and Graduation Internship Defense Results;
    \item Table of Contents;
    \item List of Symbols and Abbreviations, if applicable;
    \item List of Figures;
    \item List of Tables;
    \item Abstract;
    \item Main Chapters of the Internship Report;
    \item References;
    \item Appendices, if applicable.
\end{enumerate}

The preliminary pages shall not be numbered as chapters. The main report
content shall be organized into consecutively numbered chapters beginning
with Chapter 1.

\vspace{0.3cm}

% ================================================================
% III. PAGE SIZE AND MARGINS
% ================================================================

\noindent
\textbf{III.\quad PAGE SIZE AND MARGINS}

\vspace{0.15cm}

The report shall be formatted on A4 paper with the following margins:

\begin{center}
\begin{tabular}{ll}
    Left margin:   & 3.5 cm \\
    Right margin:  & 2.0 cm \\
    Top margin:    & 2.5 cm \\
    Bottom margin: & 2.5 cm
\end{tabular}
\end{center}

These settings are predefined in the official template and shall not be
modified by students.

\vspace{0.3cm}

% ================================================================
% IV. FONT, FONT SIZE, AND LINE SPACING
% ================================================================

\noindent
\textbf{IV.\quad FONT, FONT SIZE, AND LINE SPACING}

\vspace{0.15cm}

The main text shall use a 13 pt font size with 1.3-line spacing.
A Times-family typeface shall be used in accordance with the official
template.

The text shall be fully justified.

Students shall not use repeated spaces, repeated manual line breaks, or
multiple blank lines to position text. Paragraph and heading spacing shall
be controlled by the template.

\vspace{0.3cm}

% ================================================================
% V. OUTER COVER
% ================================================================

\noindent
\textbf{V.\quad OUTER COVER}

\vspace{0.15cm}

The outer cover shall follow the official format provided in the template
and shall contain the following information:

\begin{itemize}[leftmargin=0.8cm, itemsep=2pt]
    \item Thai Nguyen University of Technology;
    \item Faculty of International Training;
    \item Engineering and Technology Department;
    \item University logo and Faculty logo;
    \item the title \textbf{INTERNSHIP REPORT};
    \item Internship Organization;
    \item Author;
    \item Class;
    \item Supervisor;
    \item place and year of completion.
\end{itemize}

Students shall only replace the required information fields and shall not
modify the cover-page layout unless instructed otherwise.

\vspace{0.3cm}

% ================================================================
% VI. COMMENTS OF THE INTERNSHIP HOST ORGANIZATION
% ================================================================

\noindent
\textbf{VI.\quad COMMENTS OF THE INTERNSHIP HOST ORGANIZATION}

\vspace{0.15cm}

The page entitled \textbf{COMMENTS OF THE INTERNSHIP HOST ORGANIZATION}
is reserved for the internship host organization.

It includes:

\begin{itemize}[leftmargin=0.8cm, itemsep=2pt]
    \item the student's full name;
    \item class;
    \item comments from the host organization;
    \item score;
    \item date;
    \item signature and official seal.
\end{itemize}

Students shall not alter or complete official evaluation fields unless
specifically authorized.

\vspace{0.3cm}

% ================================================================
% VII. SUPERVISOR'S EVALUATION AND DEFENSE RESULTS
% ================================================================

\noindent
\textbf{VII.\quad SUPERVISOR'S EVALUATION AND DEFENSE RESULTS}

\vspace{0.15cm}

The section entitled
\textbf{SUPERVISOR'S EVALUATION AND GRADUATION INTERNSHIP DEFENSE RESULTS}
is reserved for the supervisor and defense committee.

The evaluation may include:

\begin{itemize}[leftmargin=0.8cm, itemsep=2pt]
    \item Internship Progress and Student's Attitude;
    \item Report Content;
    \item Report Presentation;
    \item Other Comments;
    \item Supervisor's Opinion;
    \item Defense Result;
    \item Score.
\end{itemize}

Students shall not enter evaluation scores or defense results before they
are officially determined.

\vspace{0.3cm}

% ================================================================
% VIII. CHAPTERS, SECTIONS, AND SUBSECTIONS
% ================================================================

\noindent
\textbf{VIII.\quad CHAPTERS, SECTIONS, AND SUBSECTIONS}

\vspace{0.15cm}

The report structure shall be generated using the standard \LaTeX{}
commands:

\begin{verbatim}
\chapter{...}
\section{...}
\subsection{...}
\end{verbatim}

Chapter, section, and subsection numbers shall be generated automatically
by \LaTeX{}.

Students shall not manually type chapter or section numbers and shall not
reproduce heading styles using commands such as \verb|\textbf|,
\verb|\fontsize|, or \verb|\centering| instead of the appropriate
structural command.

Excessive subdivision should be avoided. In most cases, the hierarchy
Chapter--Section--Subsection is sufficient.

\vspace{0.3cm}

% ================================================================
% IX. ORGANIZATION OF THE MAIN CONTENT
% ================================================================

\noindent
\textbf{IX.\quad ORGANIZATION OF THE MAIN CONTENT}

\vspace{0.15cm}

The specific content of the report shall depend on the internship
organization and the assigned tasks. A typical structure may include:

\begin{itemize}[leftmargin=0.8cm, itemsep=2pt]
    \item \textbf{Chapter 1. Overview of the Internship Organization};
    \item \textbf{Chapter 2. Internship Activities and Assigned Tasks};
    \item \textbf{Chapter 3. Results, Experience, and Evaluation}.
\end{itemize}

These chapter titles are recommendations only and may be adapted to the
actual internship content.

The report shall reflect work actually performed, observed, analyzed, or
learned during the internship. Unrelated theoretical material should not
be included merely to increase report length.

\vspace{0.3cm}

% ================================================================
% X. TABLE OF CONTENTS
% ================================================================

\noindent
\textbf{X.\quad TABLE OF CONTENTS}

\vspace{0.15cm}

The Table of Contents shall be generated automatically using:

\begin{verbatim}
\tableofcontents
\end{verbatim}

Students shall not create or manually edit page numbers in the Table of
Contents.

The document shall be compiled sufficiently to ensure that all page
numbers and entries are correctly updated.

\vspace{0.3cm}

% ================================================================
% XI. LIST OF SYMBOLS AND ABBREVIATIONS
% ================================================================

\noindent
\textbf{XI.\quad LIST OF SYMBOLS AND ABBREVIATIONS}

\vspace{0.15cm}

Symbols and abbreviations that are used repeatedly in the report shall be
listed in the \textbf{List of Symbols and Abbreviations}.

The list may contain the following columns:

\begin{center}
\begin{tabular}{llll}
    No. & Symbol / Abbreviation & Unit & Definition
\end{tabular}
\end{center}

Only symbols and abbreviations actually used in the final report shall be
retained. Example entries contained in the template shall be removed if
they are not relevant to the student's report.

Each abbreviation should normally be defined when it first appears in
the main text.

\vspace{0.3cm}

% ================================================================
% XII. FIGURES, DIAGRAMS, AND PHOTOGRAPHS
% ================================================================

\noindent
\textbf{XII.\quad FIGURES, DIAGRAMS, AND PHOTOGRAPHS}

\vspace{0.15cm}

Figures shall be clear, legible, and directly relevant to the report
content.

Figures shall be inserted using the standard \verb|figure| environment.

Example:

\begin{verbatim}
\begin{figure}[H]
    \centering
    \includegraphics[width=0.8\textwidth]
                    {figures/company_structure.pdf}
    \caption{Organizational structure of the company}
    \label{fig:company-structure}
\end{figure}
\end{verbatim}

Figure captions shall be placed below the figure.

Each referenced figure shall have a unique \verb|\label| and shall be
referenced in the text using, for example:

\begin{verbatim}
Figure~\ref{fig:company-structure}
\end{verbatim}

Students shall not manually type figure numbers.

Figures reproduced or adapted from external sources shall include an
appropriate citation.

Photographs taken during the internship shall be relevant and sufficiently
clear. Students shall not include confidential, restricted, or proprietary
information without permission from the internship organization.

Vector graphics or high-resolution images are recommended whenever
possible.

\vspace{0.3cm}

% ================================================================
% XIII. TABLES
% ================================================================

\noindent
\textbf{XIII.\quad TABLES}

\vspace{0.15cm}

Tables shall be created using suitable \LaTeX{} environments such as
\verb|table|, \verb|tabular|, \verb|tabularx|, or \verb|longtable|.

Table captions shall be placed above the table.

Example:

\begin{verbatim}
\begin{table}[H]
    \centering
    \caption{Internship schedule}
    \label{tab:internship-schedule}

    \begin{tabular}{lll}
        \toprule
        Week & Task & Result \\
        \midrule
        1 & Orientation & Completed \\
        2 & System study & Completed \\
        \bottomrule
    \end{tabular}
\end{table}
\end{verbatim}

Tables shall be referenced using:

\begin{verbatim}
Table~\ref{tab:internship-schedule}
\end{verbatim}

Students shall not manually type table numbers.

Tables should be created directly in \LaTeX{} rather than inserted as
screenshots whenever practical.

\vspace{0.3cm}

% ================================================================
% XIV. EQUATIONS AND TECHNICAL QUANTITIES
% ================================================================

\noindent
\textbf{XIV.\quad EQUATIONS AND TECHNICAL QUANTITIES}

\vspace{0.15cm}

Mathematical expressions shall be typeset directly using \LaTeX{}.
Images or screenshots of equations shall not be used.

Short expressions may be written in inline mathematical mode, for example:

\begin{verbatim}
$R = 0.7\,\Omega$
\end{verbatim}

Important equations shall use a numbered equation environment.

Example:

\begin{verbatim}
\begin{equation}
    P = UI
    \label{eq:power}
\end{equation}
\end{verbatim}

Equations shall be referenced using:

\begin{verbatim}
Equation~\eqref{eq:power}
\end{verbatim}

Equation numbers shall not be manually entered.

SI units shall be used wherever applicable. The \verb|siunitx| package
should be used where appropriate.

Examples:

\begin{verbatim}
\SI{220}{\volt}
\SI{50}{\hertz}
\SI{2.5}{\ampere}
\end{verbatim}

Notation and units shall be consistent throughout the report.

\vspace{0.3cm}

% ================================================================
% XV. SOURCE CODE
% ================================================================

\noindent
\textbf{XV.\quad SOURCE CODE}

\vspace{0.15cm}

If the internship involves programming, PLCs, microcontrollers,
simulation, or software development, relevant source code may be included
in the report.

Source code shall be presented using the \verb|listings| environment or
another environment defined by the official template.

Example:

\begin{verbatim}
\begin{lstlisting}[language=st]
IF Start = TRUE THEN
    Motor := TRUE;
ELSE
    Motor := FALSE;
END_IF;
\end{lstlisting}
\end{verbatim}

Only code necessary for explaining the work should appear in the main
report. Long source-code listings should be moved to an appendix or
submitted as supplementary files if required.

Source code should not normally be inserted as screenshots.

\vspace{0.3cm}

% ================================================================
% XVI. ABSTRACT
% ================================================================

\noindent
\textbf{XVI.\quad ABSTRACT}

\vspace{0.15cm}

The Abstract shall briefly summarize:

\begin{itemize}[leftmargin=0.8cm, itemsep=2pt]
    \item the internship organization;
    \item the objectives of the internship;
    \item the main tasks performed;
    \item the main results, knowledge, and skills obtained.
\end{itemize}

The Abstract should be concise and should not reproduce the full content
of the report chapter by chapter.

\vspace{0.3cm}

% ================================================================
% XVII. REFERENCES AND CITATIONS
% ================================================================

\noindent
\textbf{XVII.\quad REFERENCES AND CITATIONS}

\vspace{0.15cm}

All books, journal papers, conference papers, standards, websites,
datasheets, manuals, technical documents, figures, tables, or other
external sources used in preparing the report shall be properly cited.

References shall be managed through the bibliography system provided with
the official template.

Bibliographic entries shall normally be stored in:

\begin{verbatim}
References.bib
\end{verbatim}

References shall be cited in the report using:

\begin{verbatim}
\cite{reference-key}
\end{verbatim}

Students shall not manually type reference numbers such as ``[1]'',
``[2]'', etc.

Titles of foreign-language publications shall remain in their original
language in the bibliography.

All borrowed ideas, data, figures, tables, equations, methods, or results
shall be properly acknowledged.

\vspace{0.3cm}

% ================================================================
% XVIII. CROSS-REFERENCING
% ================================================================

\noindent
\textbf{XVIII.\quad CROSS-REFERENCING}

\vspace{0.15cm}

Numbered elements shall be cross-referenced using the standard \LaTeX{}
mechanism.

Use the following commands where appropriate:

\begin{verbatim}
\label{...}
\ref{...}
\eqref{...}
\end{verbatim}

For example:

\begin{verbatim}
As shown in Figure~\ref{fig:system}, ...
\end{verbatim}

Students shall not manually type numbers such as ``Figure 2.3'' when
those numbers can be generated automatically.

\vspace{0.3cm}

% ================================================================
% XIX. PAGE NUMBERING
% ================================================================

\noindent
\textbf{XIX.\quad PAGE NUMBERING}

\vspace{0.15cm}

Cover pages and official evaluation forms at the beginning of the report
shall not display page numbers.

The main body of the report shall use Arabic page numbers beginning from
Chapter 1.

Page numbering shall be generated automatically by the official template.
Students shall not manually enter page numbers.

\vspace{0.3cm}

% ================================================================
% XX. PROJECT FILE ORGANIZATION
% ================================================================

\noindent
\textbf{XX.\quad PROJECT FILE ORGANIZATION}

\vspace{0.15cm}

Students should maintain a clear and organized \LaTeX{} project
structure. A recommended structure is:

\begin{verbatim}
InternshipReport/
|
|-- main.tex
|-- References.bib
|
|-- figures/
|   |-- company.pdf
|   |-- system.png
|   `-- ...
|
`-- chapters/
    |-- chapter1.tex
    |-- chapter2.tex
    `-- chapter3.tex
\end{verbatim}

Figures should normally be stored in the \verb|figures| directory
specified by the official template.

File and directory names should preferably contain Latin letters,
numbers, hyphens, or underscores.

Recommended examples include:

\begin{verbatim}
company_structure.pdf
motor_control.png
chapter_01.tex
\end{verbatim}

Spaces and unnecessary special characters in file names should be
avoided where possible.

\vspace{0.3cm}

% ================================================================
% XXI. DOCUMENT COMPILATION AND FINAL CHECK
% ================================================================

\noindent
\textbf{XXI.\quad DOCUMENT COMPILATION AND FINAL CHECK}

\vspace{0.15cm}

Before submission, the complete \LaTeX{} project shall compile
successfully and produce the final PDF report.

The final document shall not contain unresolved references such as:

\begin{verbatim}
??
[?]
\end{verbatim}

Before generating the final PDF, students shall check:

\begin{itemize}[leftmargin=0.8cm, itemsep=2pt]
    \item the Table of Contents;
    \item page numbering;
    \item chapter, section, and subsection numbering;
    \item figure numbering;
    \item table numbering;
    \item equation numbering;
    \item cross-references;
    \item citations;
    \item bibliography;
    \item font display;
    \item mathematical symbols;
    \item image quality;
    \item completeness of all required forms.
\end{itemize}

The document may be compiled using TeX Live, MiKTeX, or another
compatible \LaTeX{} environment. Regardless of the compilation
environment, the resulting PDF shall conform to the official Faculty
template.

\vspace{0.3cm}

% ================================================================
% XXII. CONTENT THAT STUDENTS MAY EDIT
% ================================================================

\noindent
\textbf{XXII.\quad CONTENT THAT STUDENTS MAY EDIT}

\vspace{0.15cm}

Students may edit or replace:

\begin{itemize}[leftmargin=0.8cm, itemsep=2pt]
    \item student name;
    \item class;
    \item internship organization;
    \item supervisor name;
    \item Abstract;
    \item List of Symbols and Abbreviations;
    \item chapter titles and chapter content;
    \item figures and photographs;
    \item tables;
    \item equations;
    \item citations and references.
\end{itemize}

Students shall not modify the main document formatting settings in the
preamble, including the document class, page geometry, line spacing,
heading formatting, table-of-contents formatting, font configuration,
page layout, or cover-page macros unless specifically instructed.

\vspace{0.3cm}

% ================================================================
% XXIII. FINAL SUBMISSION REQUIREMENTS
% ================================================================

\noindent
\textbf{XXIII.\quad FINAL SUBMISSION REQUIREMENTS}

\vspace{0.15cm}

Before submission, the final report shall:

\begin{itemize}[leftmargin=0.8cm, itemsep=2pt]
    \item compile without errors that affect the report content or presentation;
    \item contain all required forms and report sections;
    \item contain correct student, class, internship organization, and supervisor information;
    \item contain an updated Table of Contents, List of Figures, and List of Tables;
    \item contain no unused example content from the template;
    \item contain no placeholder text such as \textit{Title of Chapter 1},
    \textit{Section 1.1}, or other sample entries;
    \item contain all required citations and references;
    \item contain no unresolved cross-references.
\end{itemize}

If the Faculty requires submission of the \LaTeX{} source, the complete
project shall include all \verb|.tex| files, the bibliography database,
figures, and any other files required to reproduce the submitted PDF.

\vspace{0.4cm}

\noindent
\textbf{Important:}
Before submitting the final internship report, students shall remove or
comment out the following command in the main document:

\begin{verbatim}
\input{Regulations.tex}
\end{verbatim}

For example:

\begin{verbatim}
% \input{Regulations.tex}
\end{verbatim}

This removes the guidance pages from the final submitted report while
preserving the regulations file in the \LaTeX{} project for reference.

} % End 13pt content

\clearpage

\end{verbatim}

This removes the guidance pages from the final submitted report while
preserving the regulations file in the \LaTeX{} project for reference.

} % End 13pt content

\clearpage

\end{verbatim}

For example:

\begin{verbatim}
% 
% =================================================================
% REGULATIONS FOR PREPARING THE GRADUATION INTERNSHIP REPORT
% USING LATEX
% =================================================================
%
% This file is intended to be included from the main document with:
%
%     
% =================================================================
% REGULATIONS FOR PREPARING THE GRADUATION INTERNSHIP REPORT
% USING LATEX
% =================================================================
%
% This file is intended to be included from the main document with:
%
%     \input{Regulations.tex}
%
% IMPORTANT:
% This section is provided for guidance only. Students shall remove
% or comment out the \input{Regulations.tex} command before submitting
% the final internship report.
%
% =================================================================

\clearpage
\thispagestyle{empty}

\begin{center}
    {\fontsize{16pt}{20pt}\selectfont\bfseries
    SPECIFIC REGULATIONS FOR PREPARING\\
    GRADUATION INTERNSHIP REPORTS USING \LaTeX}
\end{center}

\vspace{0.5cm}

{\fontsize{13pt}{17pt}\selectfont
\setstretch{1.3}
\justifying

\noindent
\textbf{Note:}
This section is provided for guidance only and is included in the main
document using the command \verb|\input{Regulations.tex}|.
Students shall remove or comment out this command before submitting the
final internship report.

\vspace{0.4cm}

% ================================================================
% I. GENERAL REQUIREMENTS
% ================================================================

\noindent
\textbf{I.\quad GENERAL REQUIREMENTS}

\vspace{0.15cm}

The Graduation Internship Report shall be prepared using the official
\LaTeX{} template provided by the Faculty. Students shall preserve the
predefined page layout, typography, line spacing, heading styles,
numbering system, figure and table formatting, and bibliography settings
unless otherwise instructed by the Faculty or the supervisor.

The report shall be written in English.

Students are responsible for the academic and technical content of the
report. The official \LaTeX{} template is responsible for standardized
document formatting.

\vspace{0.3cm}

% ================================================================
% II. REPORT STRUCTURE
% ================================================================

\noindent
\textbf{II.\quad REPORT STRUCTURE}

\vspace{0.15cm}

The report shall normally be arranged in the following order:

\begin{enumerate}[label=\arabic*., leftmargin=1.2cm, itemsep=2pt]
    \item Outer Cover;
    \item Comments of the Internship Host Organization;
    \item Supervisor's Evaluation and Graduation Internship Defense Results;
    \item Table of Contents;
    \item List of Symbols and Abbreviations, if applicable;
    \item List of Figures;
    \item List of Tables;
    \item Abstract;
    \item Main Chapters of the Internship Report;
    \item References;
    \item Appendices, if applicable.
\end{enumerate}

The preliminary pages shall not be numbered as chapters. The main report
content shall be organized into consecutively numbered chapters beginning
with Chapter 1.

\vspace{0.3cm}

% ================================================================
% III. PAGE SIZE AND MARGINS
% ================================================================

\noindent
\textbf{III.\quad PAGE SIZE AND MARGINS}

\vspace{0.15cm}

The report shall be formatted on A4 paper with the following margins:

\begin{center}
\begin{tabular}{ll}
    Left margin:   & 3.5 cm \\
    Right margin:  & 2.0 cm \\
    Top margin:    & 2.5 cm \\
    Bottom margin: & 2.5 cm
\end{tabular}
\end{center}

These settings are predefined in the official template and shall not be
modified by students.

\vspace{0.3cm}

% ================================================================
% IV. FONT, FONT SIZE, AND LINE SPACING
% ================================================================

\noindent
\textbf{IV.\quad FONT, FONT SIZE, AND LINE SPACING}

\vspace{0.15cm}

The main text shall use a 13 pt font size with 1.3-line spacing.
A Times-family typeface shall be used in accordance with the official
template.

The text shall be fully justified.

Students shall not use repeated spaces, repeated manual line breaks, or
multiple blank lines to position text. Paragraph and heading spacing shall
be controlled by the template.

\vspace{0.3cm}

% ================================================================
% V. OUTER COVER
% ================================================================

\noindent
\textbf{V.\quad OUTER COVER}

\vspace{0.15cm}

The outer cover shall follow the official format provided in the template
and shall contain the following information:

\begin{itemize}[leftmargin=0.8cm, itemsep=2pt]
    \item Thai Nguyen University of Technology;
    \item Faculty of International Training;
    \item Engineering and Technology Department;
    \item University logo and Faculty logo;
    \item the title \textbf{INTERNSHIP REPORT};
    \item Internship Organization;
    \item Author;
    \item Class;
    \item Supervisor;
    \item place and year of completion.
\end{itemize}

Students shall only replace the required information fields and shall not
modify the cover-page layout unless instructed otherwise.

\vspace{0.3cm}

% ================================================================
% VI. COMMENTS OF THE INTERNSHIP HOST ORGANIZATION
% ================================================================

\noindent
\textbf{VI.\quad COMMENTS OF THE INTERNSHIP HOST ORGANIZATION}

\vspace{0.15cm}

The page entitled \textbf{COMMENTS OF THE INTERNSHIP HOST ORGANIZATION}
is reserved for the internship host organization.

It includes:

\begin{itemize}[leftmargin=0.8cm, itemsep=2pt]
    \item the student's full name;
    \item class;
    \item comments from the host organization;
    \item score;
    \item date;
    \item signature and official seal.
\end{itemize}

Students shall not alter or complete official evaluation fields unless
specifically authorized.

\vspace{0.3cm}

% ================================================================
% VII. SUPERVISOR'S EVALUATION AND DEFENSE RESULTS
% ================================================================

\noindent
\textbf{VII.\quad SUPERVISOR'S EVALUATION AND DEFENSE RESULTS}

\vspace{0.15cm}

The section entitled
\textbf{SUPERVISOR'S EVALUATION AND GRADUATION INTERNSHIP DEFENSE RESULTS}
is reserved for the supervisor and defense committee.

The evaluation may include:

\begin{itemize}[leftmargin=0.8cm, itemsep=2pt]
    \item Internship Progress and Student's Attitude;
    \item Report Content;
    \item Report Presentation;
    \item Other Comments;
    \item Supervisor's Opinion;
    \item Defense Result;
    \item Score.
\end{itemize}

Students shall not enter evaluation scores or defense results before they
are officially determined.

\vspace{0.3cm}

% ================================================================
% VIII. CHAPTERS, SECTIONS, AND SUBSECTIONS
% ================================================================

\noindent
\textbf{VIII.\quad CHAPTERS, SECTIONS, AND SUBSECTIONS}

\vspace{0.15cm}

The report structure shall be generated using the standard \LaTeX{}
commands:

\begin{verbatim}
\chapter{...}
\section{...}
\subsection{...}
\end{verbatim}

Chapter, section, and subsection numbers shall be generated automatically
by \LaTeX{}.

Students shall not manually type chapter or section numbers and shall not
reproduce heading styles using commands such as \verb|\textbf|,
\verb|\fontsize|, or \verb|\centering| instead of the appropriate
structural command.

Excessive subdivision should be avoided. In most cases, the hierarchy
Chapter--Section--Subsection is sufficient.

\vspace{0.3cm}

% ================================================================
% IX. ORGANIZATION OF THE MAIN CONTENT
% ================================================================

\noindent
\textbf{IX.\quad ORGANIZATION OF THE MAIN CONTENT}

\vspace{0.15cm}

The specific content of the report shall depend on the internship
organization and the assigned tasks. A typical structure may include:

\begin{itemize}[leftmargin=0.8cm, itemsep=2pt]
    \item \textbf{Chapter 1. Overview of the Internship Organization};
    \item \textbf{Chapter 2. Internship Activities and Assigned Tasks};
    \item \textbf{Chapter 3. Results, Experience, and Evaluation}.
\end{itemize}

These chapter titles are recommendations only and may be adapted to the
actual internship content.

The report shall reflect work actually performed, observed, analyzed, or
learned during the internship. Unrelated theoretical material should not
be included merely to increase report length.

\vspace{0.3cm}

% ================================================================
% X. TABLE OF CONTENTS
% ================================================================

\noindent
\textbf{X.\quad TABLE OF CONTENTS}

\vspace{0.15cm}

The Table of Contents shall be generated automatically using:

\begin{verbatim}
\tableofcontents
\end{verbatim}

Students shall not create or manually edit page numbers in the Table of
Contents.

The document shall be compiled sufficiently to ensure that all page
numbers and entries are correctly updated.

\vspace{0.3cm}

% ================================================================
% XI. LIST OF SYMBOLS AND ABBREVIATIONS
% ================================================================

\noindent
\textbf{XI.\quad LIST OF SYMBOLS AND ABBREVIATIONS}

\vspace{0.15cm}

Symbols and abbreviations that are used repeatedly in the report shall be
listed in the \textbf{List of Symbols and Abbreviations}.

The list may contain the following columns:

\begin{center}
\begin{tabular}{llll}
    No. & Symbol / Abbreviation & Unit & Definition
\end{tabular}
\end{center}

Only symbols and abbreviations actually used in the final report shall be
retained. Example entries contained in the template shall be removed if
they are not relevant to the student's report.

Each abbreviation should normally be defined when it first appears in
the main text.

\vspace{0.3cm}

% ================================================================
% XII. FIGURES, DIAGRAMS, AND PHOTOGRAPHS
% ================================================================

\noindent
\textbf{XII.\quad FIGURES, DIAGRAMS, AND PHOTOGRAPHS}

\vspace{0.15cm}

Figures shall be clear, legible, and directly relevant to the report
content.

Figures shall be inserted using the standard \verb|figure| environment.

Example:

\begin{verbatim}
\begin{figure}[H]
    \centering
    \includegraphics[width=0.8\textwidth]
                    {figures/company_structure.pdf}
    \caption{Organizational structure of the company}
    \label{fig:company-structure}
\end{figure}
\end{verbatim}

Figure captions shall be placed below the figure.

Each referenced figure shall have a unique \verb|\label| and shall be
referenced in the text using, for example:

\begin{verbatim}
Figure~\ref{fig:company-structure}
\end{verbatim}

Students shall not manually type figure numbers.

Figures reproduced or adapted from external sources shall include an
appropriate citation.

Photographs taken during the internship shall be relevant and sufficiently
clear. Students shall not include confidential, restricted, or proprietary
information without permission from the internship organization.

Vector graphics or high-resolution images are recommended whenever
possible.

\vspace{0.3cm}

% ================================================================
% XIII. TABLES
% ================================================================

\noindent
\textbf{XIII.\quad TABLES}

\vspace{0.15cm}

Tables shall be created using suitable \LaTeX{} environments such as
\verb|table|, \verb|tabular|, \verb|tabularx|, or \verb|longtable|.

Table captions shall be placed above the table.

Example:

\begin{verbatim}
\begin{table}[H]
    \centering
    \caption{Internship schedule}
    \label{tab:internship-schedule}

    \begin{tabular}{lll}
        \toprule
        Week & Task & Result \\
        \midrule
        1 & Orientation & Completed \\
        2 & System study & Completed \\
        \bottomrule
    \end{tabular}
\end{table}
\end{verbatim}

Tables shall be referenced using:

\begin{verbatim}
Table~\ref{tab:internship-schedule}
\end{verbatim}

Students shall not manually type table numbers.

Tables should be created directly in \LaTeX{} rather than inserted as
screenshots whenever practical.

\vspace{0.3cm}

% ================================================================
% XIV. EQUATIONS AND TECHNICAL QUANTITIES
% ================================================================

\noindent
\textbf{XIV.\quad EQUATIONS AND TECHNICAL QUANTITIES}

\vspace{0.15cm}

Mathematical expressions shall be typeset directly using \LaTeX{}.
Images or screenshots of equations shall not be used.

Short expressions may be written in inline mathematical mode, for example:

\begin{verbatim}
$R = 0.7\,\Omega$
\end{verbatim}

Important equations shall use a numbered equation environment.

Example:

\begin{verbatim}
\begin{equation}
    P = UI
    \label{eq:power}
\end{equation}
\end{verbatim}

Equations shall be referenced using:

\begin{verbatim}
Equation~\eqref{eq:power}
\end{verbatim}

Equation numbers shall not be manually entered.

SI units shall be used wherever applicable. The \verb|siunitx| package
should be used where appropriate.

Examples:

\begin{verbatim}
\SI{220}{\volt}
\SI{50}{\hertz}
\SI{2.5}{\ampere}
\end{verbatim}

Notation and units shall be consistent throughout the report.

\vspace{0.3cm}

% ================================================================
% XV. SOURCE CODE
% ================================================================

\noindent
\textbf{XV.\quad SOURCE CODE}

\vspace{0.15cm}

If the internship involves programming, PLCs, microcontrollers,
simulation, or software development, relevant source code may be included
in the report.

Source code shall be presented using the \verb|listings| environment or
another environment defined by the official template.

Example:

\begin{verbatim}
\begin{lstlisting}[language=st]
IF Start = TRUE THEN
    Motor := TRUE;
ELSE
    Motor := FALSE;
END_IF;
\end{lstlisting}
\end{verbatim}

Only code necessary for explaining the work should appear in the main
report. Long source-code listings should be moved to an appendix or
submitted as supplementary files if required.

Source code should not normally be inserted as screenshots.

\vspace{0.3cm}

% ================================================================
% XVI. ABSTRACT
% ================================================================

\noindent
\textbf{XVI.\quad ABSTRACT}

\vspace{0.15cm}

The Abstract shall briefly summarize:

\begin{itemize}[leftmargin=0.8cm, itemsep=2pt]
    \item the internship organization;
    \item the objectives of the internship;
    \item the main tasks performed;
    \item the main results, knowledge, and skills obtained.
\end{itemize}

The Abstract should be concise and should not reproduce the full content
of the report chapter by chapter.

\vspace{0.3cm}

% ================================================================
% XVII. REFERENCES AND CITATIONS
% ================================================================

\noindent
\textbf{XVII.\quad REFERENCES AND CITATIONS}

\vspace{0.15cm}

All books, journal papers, conference papers, standards, websites,
datasheets, manuals, technical documents, figures, tables, or other
external sources used in preparing the report shall be properly cited.

References shall be managed through the bibliography system provided with
the official template.

Bibliographic entries shall normally be stored in:

\begin{verbatim}
References.bib
\end{verbatim}

References shall be cited in the report using:

\begin{verbatim}
\cite{reference-key}
\end{verbatim}

Students shall not manually type reference numbers such as ``[1]'',
``[2]'', etc.

Titles of foreign-language publications shall remain in their original
language in the bibliography.

All borrowed ideas, data, figures, tables, equations, methods, or results
shall be properly acknowledged.

\vspace{0.3cm}

% ================================================================
% XVIII. CROSS-REFERENCING
% ================================================================

\noindent
\textbf{XVIII.\quad CROSS-REFERENCING}

\vspace{0.15cm}

Numbered elements shall be cross-referenced using the standard \LaTeX{}
mechanism.

Use the following commands where appropriate:

\begin{verbatim}
\label{...}
\ref{...}
\eqref{...}
\end{verbatim}

For example:

\begin{verbatim}
As shown in Figure~\ref{fig:system}, ...
\end{verbatim}

Students shall not manually type numbers such as ``Figure 2.3'' when
those numbers can be generated automatically.

\vspace{0.3cm}

% ================================================================
% XIX. PAGE NUMBERING
% ================================================================

\noindent
\textbf{XIX.\quad PAGE NUMBERING}

\vspace{0.15cm}

Cover pages and official evaluation forms at the beginning of the report
shall not display page numbers.

The main body of the report shall use Arabic page numbers beginning from
Chapter 1.

Page numbering shall be generated automatically by the official template.
Students shall not manually enter page numbers.

\vspace{0.3cm}

% ================================================================
% XX. PROJECT FILE ORGANIZATION
% ================================================================

\noindent
\textbf{XX.\quad PROJECT FILE ORGANIZATION}

\vspace{0.15cm}

Students should maintain a clear and organized \LaTeX{} project
structure. A recommended structure is:

\begin{verbatim}
InternshipReport/
|
|-- main.tex
|-- References.bib
|
|-- figures/
|   |-- company.pdf
|   |-- system.png
|   `-- ...
|
`-- chapters/
    |-- chapter1.tex
    |-- chapter2.tex
    `-- chapter3.tex
\end{verbatim}

Figures should normally be stored in the \verb|figures| directory
specified by the official template.

File and directory names should preferably contain Latin letters,
numbers, hyphens, or underscores.

Recommended examples include:

\begin{verbatim}
company_structure.pdf
motor_control.png
chapter_01.tex
\end{verbatim}

Spaces and unnecessary special characters in file names should be
avoided where possible.

\vspace{0.3cm}

% ================================================================
% XXI. DOCUMENT COMPILATION AND FINAL CHECK
% ================================================================

\noindent
\textbf{XXI.\quad DOCUMENT COMPILATION AND FINAL CHECK}

\vspace{0.15cm}

Before submission, the complete \LaTeX{} project shall compile
successfully and produce the final PDF report.

The final document shall not contain unresolved references such as:

\begin{verbatim}
??
[?]
\end{verbatim}

Before generating the final PDF, students shall check:

\begin{itemize}[leftmargin=0.8cm, itemsep=2pt]
    \item the Table of Contents;
    \item page numbering;
    \item chapter, section, and subsection numbering;
    \item figure numbering;
    \item table numbering;
    \item equation numbering;
    \item cross-references;
    \item citations;
    \item bibliography;
    \item font display;
    \item mathematical symbols;
    \item image quality;
    \item completeness of all required forms.
\end{itemize}

The document may be compiled using TeX Live, MiKTeX, or another
compatible \LaTeX{} environment. Regardless of the compilation
environment, the resulting PDF shall conform to the official Faculty
template.

\vspace{0.3cm}

% ================================================================
% XXII. CONTENT THAT STUDENTS MAY EDIT
% ================================================================

\noindent
\textbf{XXII.\quad CONTENT THAT STUDENTS MAY EDIT}

\vspace{0.15cm}

Students may edit or replace:

\begin{itemize}[leftmargin=0.8cm, itemsep=2pt]
    \item student name;
    \item class;
    \item internship organization;
    \item supervisor name;
    \item Abstract;
    \item List of Symbols and Abbreviations;
    \item chapter titles and chapter content;
    \item figures and photographs;
    \item tables;
    \item equations;
    \item citations and references.
\end{itemize}

Students shall not modify the main document formatting settings in the
preamble, including the document class, page geometry, line spacing,
heading formatting, table-of-contents formatting, font configuration,
page layout, or cover-page macros unless specifically instructed.

\vspace{0.3cm}

% ================================================================
% XXIII. FINAL SUBMISSION REQUIREMENTS
% ================================================================

\noindent
\textbf{XXIII.\quad FINAL SUBMISSION REQUIREMENTS}

\vspace{0.15cm}

Before submission, the final report shall:

\begin{itemize}[leftmargin=0.8cm, itemsep=2pt]
    \item compile without errors that affect the report content or presentation;
    \item contain all required forms and report sections;
    \item contain correct student, class, internship organization, and supervisor information;
    \item contain an updated Table of Contents, List of Figures, and List of Tables;
    \item contain no unused example content from the template;
    \item contain no placeholder text such as \textit{Title of Chapter 1},
    \textit{Section 1.1}, or other sample entries;
    \item contain all required citations and references;
    \item contain no unresolved cross-references.
\end{itemize}

If the Faculty requires submission of the \LaTeX{} source, the complete
project shall include all \verb|.tex| files, the bibliography database,
figures, and any other files required to reproduce the submitted PDF.

\vspace{0.4cm}

\noindent
\textbf{Important:}
Before submitting the final internship report, students shall remove or
comment out the following command in the main document:

\begin{verbatim}
\input{Regulations.tex}
\end{verbatim}

For example:

\begin{verbatim}
% \input{Regulations.tex}
\end{verbatim}

This removes the guidance pages from the final submitted report while
preserving the regulations file in the \LaTeX{} project for reference.

} % End 13pt content

\clearpage

%
% IMPORTANT:
% This section is provided for guidance only. Students shall remove
% or comment out the 
% =================================================================
% REGULATIONS FOR PREPARING THE GRADUATION INTERNSHIP REPORT
% USING LATEX
% =================================================================
%
% This file is intended to be included from the main document with:
%
%     \input{Regulations.tex}
%
% IMPORTANT:
% This section is provided for guidance only. Students shall remove
% or comment out the \input{Regulations.tex} command before submitting
% the final internship report.
%
% =================================================================

\clearpage
\thispagestyle{empty}

\begin{center}
    {\fontsize{16pt}{20pt}\selectfont\bfseries
    SPECIFIC REGULATIONS FOR PREPARING\\
    GRADUATION INTERNSHIP REPORTS USING \LaTeX}
\end{center}

\vspace{0.5cm}

{\fontsize{13pt}{17pt}\selectfont
\setstretch{1.3}
\justifying

\noindent
\textbf{Note:}
This section is provided for guidance only and is included in the main
document using the command \verb|\input{Regulations.tex}|.
Students shall remove or comment out this command before submitting the
final internship report.

\vspace{0.4cm}

% ================================================================
% I. GENERAL REQUIREMENTS
% ================================================================

\noindent
\textbf{I.\quad GENERAL REQUIREMENTS}

\vspace{0.15cm}

The Graduation Internship Report shall be prepared using the official
\LaTeX{} template provided by the Faculty. Students shall preserve the
predefined page layout, typography, line spacing, heading styles,
numbering system, figure and table formatting, and bibliography settings
unless otherwise instructed by the Faculty or the supervisor.

The report shall be written in English.

Students are responsible for the academic and technical content of the
report. The official \LaTeX{} template is responsible for standardized
document formatting.

\vspace{0.3cm}

% ================================================================
% II. REPORT STRUCTURE
% ================================================================

\noindent
\textbf{II.\quad REPORT STRUCTURE}

\vspace{0.15cm}

The report shall normally be arranged in the following order:

\begin{enumerate}[label=\arabic*., leftmargin=1.2cm, itemsep=2pt]
    \item Outer Cover;
    \item Comments of the Internship Host Organization;
    \item Supervisor's Evaluation and Graduation Internship Defense Results;
    \item Table of Contents;
    \item List of Symbols and Abbreviations, if applicable;
    \item List of Figures;
    \item List of Tables;
    \item Abstract;
    \item Main Chapters of the Internship Report;
    \item References;
    \item Appendices, if applicable.
\end{enumerate}

The preliminary pages shall not be numbered as chapters. The main report
content shall be organized into consecutively numbered chapters beginning
with Chapter 1.

\vspace{0.3cm}

% ================================================================
% III. PAGE SIZE AND MARGINS
% ================================================================

\noindent
\textbf{III.\quad PAGE SIZE AND MARGINS}

\vspace{0.15cm}

The report shall be formatted on A4 paper with the following margins:

\begin{center}
\begin{tabular}{ll}
    Left margin:   & 3.5 cm \\
    Right margin:  & 2.0 cm \\
    Top margin:    & 2.5 cm \\
    Bottom margin: & 2.5 cm
\end{tabular}
\end{center}

These settings are predefined in the official template and shall not be
modified by students.

\vspace{0.3cm}

% ================================================================
% IV. FONT, FONT SIZE, AND LINE SPACING
% ================================================================

\noindent
\textbf{IV.\quad FONT, FONT SIZE, AND LINE SPACING}

\vspace{0.15cm}

The main text shall use a 13 pt font size with 1.3-line spacing.
A Times-family typeface shall be used in accordance with the official
template.

The text shall be fully justified.

Students shall not use repeated spaces, repeated manual line breaks, or
multiple blank lines to position text. Paragraph and heading spacing shall
be controlled by the template.

\vspace{0.3cm}

% ================================================================
% V. OUTER COVER
% ================================================================

\noindent
\textbf{V.\quad OUTER COVER}

\vspace{0.15cm}

The outer cover shall follow the official format provided in the template
and shall contain the following information:

\begin{itemize}[leftmargin=0.8cm, itemsep=2pt]
    \item Thai Nguyen University of Technology;
    \item Faculty of International Training;
    \item Engineering and Technology Department;
    \item University logo and Faculty logo;
    \item the title \textbf{INTERNSHIP REPORT};
    \item Internship Organization;
    \item Author;
    \item Class;
    \item Supervisor;
    \item place and year of completion.
\end{itemize}

Students shall only replace the required information fields and shall not
modify the cover-page layout unless instructed otherwise.

\vspace{0.3cm}

% ================================================================
% VI. COMMENTS OF THE INTERNSHIP HOST ORGANIZATION
% ================================================================

\noindent
\textbf{VI.\quad COMMENTS OF THE INTERNSHIP HOST ORGANIZATION}

\vspace{0.15cm}

The page entitled \textbf{COMMENTS OF THE INTERNSHIP HOST ORGANIZATION}
is reserved for the internship host organization.

It includes:

\begin{itemize}[leftmargin=0.8cm, itemsep=2pt]
    \item the student's full name;
    \item class;
    \item comments from the host organization;
    \item score;
    \item date;
    \item signature and official seal.
\end{itemize}

Students shall not alter or complete official evaluation fields unless
specifically authorized.

\vspace{0.3cm}

% ================================================================
% VII. SUPERVISOR'S EVALUATION AND DEFENSE RESULTS
% ================================================================

\noindent
\textbf{VII.\quad SUPERVISOR'S EVALUATION AND DEFENSE RESULTS}

\vspace{0.15cm}

The section entitled
\textbf{SUPERVISOR'S EVALUATION AND GRADUATION INTERNSHIP DEFENSE RESULTS}
is reserved for the supervisor and defense committee.

The evaluation may include:

\begin{itemize}[leftmargin=0.8cm, itemsep=2pt]
    \item Internship Progress and Student's Attitude;
    \item Report Content;
    \item Report Presentation;
    \item Other Comments;
    \item Supervisor's Opinion;
    \item Defense Result;
    \item Score.
\end{itemize}

Students shall not enter evaluation scores or defense results before they
are officially determined.

\vspace{0.3cm}

% ================================================================
% VIII. CHAPTERS, SECTIONS, AND SUBSECTIONS
% ================================================================

\noindent
\textbf{VIII.\quad CHAPTERS, SECTIONS, AND SUBSECTIONS}

\vspace{0.15cm}

The report structure shall be generated using the standard \LaTeX{}
commands:

\begin{verbatim}
\chapter{...}
\section{...}
\subsection{...}
\end{verbatim}

Chapter, section, and subsection numbers shall be generated automatically
by \LaTeX{}.

Students shall not manually type chapter or section numbers and shall not
reproduce heading styles using commands such as \verb|\textbf|,
\verb|\fontsize|, or \verb|\centering| instead of the appropriate
structural command.

Excessive subdivision should be avoided. In most cases, the hierarchy
Chapter--Section--Subsection is sufficient.

\vspace{0.3cm}

% ================================================================
% IX. ORGANIZATION OF THE MAIN CONTENT
% ================================================================

\noindent
\textbf{IX.\quad ORGANIZATION OF THE MAIN CONTENT}

\vspace{0.15cm}

The specific content of the report shall depend on the internship
organization and the assigned tasks. A typical structure may include:

\begin{itemize}[leftmargin=0.8cm, itemsep=2pt]
    \item \textbf{Chapter 1. Overview of the Internship Organization};
    \item \textbf{Chapter 2. Internship Activities and Assigned Tasks};
    \item \textbf{Chapter 3. Results, Experience, and Evaluation}.
\end{itemize}

These chapter titles are recommendations only and may be adapted to the
actual internship content.

The report shall reflect work actually performed, observed, analyzed, or
learned during the internship. Unrelated theoretical material should not
be included merely to increase report length.

\vspace{0.3cm}

% ================================================================
% X. TABLE OF CONTENTS
% ================================================================

\noindent
\textbf{X.\quad TABLE OF CONTENTS}

\vspace{0.15cm}

The Table of Contents shall be generated automatically using:

\begin{verbatim}
\tableofcontents
\end{verbatim}

Students shall not create or manually edit page numbers in the Table of
Contents.

The document shall be compiled sufficiently to ensure that all page
numbers and entries are correctly updated.

\vspace{0.3cm}

% ================================================================
% XI. LIST OF SYMBOLS AND ABBREVIATIONS
% ================================================================

\noindent
\textbf{XI.\quad LIST OF SYMBOLS AND ABBREVIATIONS}

\vspace{0.15cm}

Symbols and abbreviations that are used repeatedly in the report shall be
listed in the \textbf{List of Symbols and Abbreviations}.

The list may contain the following columns:

\begin{center}
\begin{tabular}{llll}
    No. & Symbol / Abbreviation & Unit & Definition
\end{tabular}
\end{center}

Only symbols and abbreviations actually used in the final report shall be
retained. Example entries contained in the template shall be removed if
they are not relevant to the student's report.

Each abbreviation should normally be defined when it first appears in
the main text.

\vspace{0.3cm}

% ================================================================
% XII. FIGURES, DIAGRAMS, AND PHOTOGRAPHS
% ================================================================

\noindent
\textbf{XII.\quad FIGURES, DIAGRAMS, AND PHOTOGRAPHS}

\vspace{0.15cm}

Figures shall be clear, legible, and directly relevant to the report
content.

Figures shall be inserted using the standard \verb|figure| environment.

Example:

\begin{verbatim}
\begin{figure}[H]
    \centering
    \includegraphics[width=0.8\textwidth]
                    {figures/company_structure.pdf}
    \caption{Organizational structure of the company}
    \label{fig:company-structure}
\end{figure}
\end{verbatim}

Figure captions shall be placed below the figure.

Each referenced figure shall have a unique \verb|\label| and shall be
referenced in the text using, for example:

\begin{verbatim}
Figure~\ref{fig:company-structure}
\end{verbatim}

Students shall not manually type figure numbers.

Figures reproduced or adapted from external sources shall include an
appropriate citation.

Photographs taken during the internship shall be relevant and sufficiently
clear. Students shall not include confidential, restricted, or proprietary
information without permission from the internship organization.

Vector graphics or high-resolution images are recommended whenever
possible.

\vspace{0.3cm}

% ================================================================
% XIII. TABLES
% ================================================================

\noindent
\textbf{XIII.\quad TABLES}

\vspace{0.15cm}

Tables shall be created using suitable \LaTeX{} environments such as
\verb|table|, \verb|tabular|, \verb|tabularx|, or \verb|longtable|.

Table captions shall be placed above the table.

Example:

\begin{verbatim}
\begin{table}[H]
    \centering
    \caption{Internship schedule}
    \label{tab:internship-schedule}

    \begin{tabular}{lll}
        \toprule
        Week & Task & Result \\
        \midrule
        1 & Orientation & Completed \\
        2 & System study & Completed \\
        \bottomrule
    \end{tabular}
\end{table}
\end{verbatim}

Tables shall be referenced using:

\begin{verbatim}
Table~\ref{tab:internship-schedule}
\end{verbatim}

Students shall not manually type table numbers.

Tables should be created directly in \LaTeX{} rather than inserted as
screenshots whenever practical.

\vspace{0.3cm}

% ================================================================
% XIV. EQUATIONS AND TECHNICAL QUANTITIES
% ================================================================

\noindent
\textbf{XIV.\quad EQUATIONS AND TECHNICAL QUANTITIES}

\vspace{0.15cm}

Mathematical expressions shall be typeset directly using \LaTeX{}.
Images or screenshots of equations shall not be used.

Short expressions may be written in inline mathematical mode, for example:

\begin{verbatim}
$R = 0.7\,\Omega$
\end{verbatim}

Important equations shall use a numbered equation environment.

Example:

\begin{verbatim}
\begin{equation}
    P = UI
    \label{eq:power}
\end{equation}
\end{verbatim}

Equations shall be referenced using:

\begin{verbatim}
Equation~\eqref{eq:power}
\end{verbatim}

Equation numbers shall not be manually entered.

SI units shall be used wherever applicable. The \verb|siunitx| package
should be used where appropriate.

Examples:

\begin{verbatim}
\SI{220}{\volt}
\SI{50}{\hertz}
\SI{2.5}{\ampere}
\end{verbatim}

Notation and units shall be consistent throughout the report.

\vspace{0.3cm}

% ================================================================
% XV. SOURCE CODE
% ================================================================

\noindent
\textbf{XV.\quad SOURCE CODE}

\vspace{0.15cm}

If the internship involves programming, PLCs, microcontrollers,
simulation, or software development, relevant source code may be included
in the report.

Source code shall be presented using the \verb|listings| environment or
another environment defined by the official template.

Example:

\begin{verbatim}
\begin{lstlisting}[language=st]
IF Start = TRUE THEN
    Motor := TRUE;
ELSE
    Motor := FALSE;
END_IF;
\end{lstlisting}
\end{verbatim}

Only code necessary for explaining the work should appear in the main
report. Long source-code listings should be moved to an appendix or
submitted as supplementary files if required.

Source code should not normally be inserted as screenshots.

\vspace{0.3cm}

% ================================================================
% XVI. ABSTRACT
% ================================================================

\noindent
\textbf{XVI.\quad ABSTRACT}

\vspace{0.15cm}

The Abstract shall briefly summarize:

\begin{itemize}[leftmargin=0.8cm, itemsep=2pt]
    \item the internship organization;
    \item the objectives of the internship;
    \item the main tasks performed;
    \item the main results, knowledge, and skills obtained.
\end{itemize}

The Abstract should be concise and should not reproduce the full content
of the report chapter by chapter.

\vspace{0.3cm}

% ================================================================
% XVII. REFERENCES AND CITATIONS
% ================================================================

\noindent
\textbf{XVII.\quad REFERENCES AND CITATIONS}

\vspace{0.15cm}

All books, journal papers, conference papers, standards, websites,
datasheets, manuals, technical documents, figures, tables, or other
external sources used in preparing the report shall be properly cited.

References shall be managed through the bibliography system provided with
the official template.

Bibliographic entries shall normally be stored in:

\begin{verbatim}
References.bib
\end{verbatim}

References shall be cited in the report using:

\begin{verbatim}
\cite{reference-key}
\end{verbatim}

Students shall not manually type reference numbers such as ``[1]'',
``[2]'', etc.

Titles of foreign-language publications shall remain in their original
language in the bibliography.

All borrowed ideas, data, figures, tables, equations, methods, or results
shall be properly acknowledged.

\vspace{0.3cm}

% ================================================================
% XVIII. CROSS-REFERENCING
% ================================================================

\noindent
\textbf{XVIII.\quad CROSS-REFERENCING}

\vspace{0.15cm}

Numbered elements shall be cross-referenced using the standard \LaTeX{}
mechanism.

Use the following commands where appropriate:

\begin{verbatim}
\label{...}
\ref{...}
\eqref{...}
\end{verbatim}

For example:

\begin{verbatim}
As shown in Figure~\ref{fig:system}, ...
\end{verbatim}

Students shall not manually type numbers such as ``Figure 2.3'' when
those numbers can be generated automatically.

\vspace{0.3cm}

% ================================================================
% XIX. PAGE NUMBERING
% ================================================================

\noindent
\textbf{XIX.\quad PAGE NUMBERING}

\vspace{0.15cm}

Cover pages and official evaluation forms at the beginning of the report
shall not display page numbers.

The main body of the report shall use Arabic page numbers beginning from
Chapter 1.

Page numbering shall be generated automatically by the official template.
Students shall not manually enter page numbers.

\vspace{0.3cm}

% ================================================================
% XX. PROJECT FILE ORGANIZATION
% ================================================================

\noindent
\textbf{XX.\quad PROJECT FILE ORGANIZATION}

\vspace{0.15cm}

Students should maintain a clear and organized \LaTeX{} project
structure. A recommended structure is:

\begin{verbatim}
InternshipReport/
|
|-- main.tex
|-- References.bib
|
|-- figures/
|   |-- company.pdf
|   |-- system.png
|   `-- ...
|
`-- chapters/
    |-- chapter1.tex
    |-- chapter2.tex
    `-- chapter3.tex
\end{verbatim}

Figures should normally be stored in the \verb|figures| directory
specified by the official template.

File and directory names should preferably contain Latin letters,
numbers, hyphens, or underscores.

Recommended examples include:

\begin{verbatim}
company_structure.pdf
motor_control.png
chapter_01.tex
\end{verbatim}

Spaces and unnecessary special characters in file names should be
avoided where possible.

\vspace{0.3cm}

% ================================================================
% XXI. DOCUMENT COMPILATION AND FINAL CHECK
% ================================================================

\noindent
\textbf{XXI.\quad DOCUMENT COMPILATION AND FINAL CHECK}

\vspace{0.15cm}

Before submission, the complete \LaTeX{} project shall compile
successfully and produce the final PDF report.

The final document shall not contain unresolved references such as:

\begin{verbatim}
??
[?]
\end{verbatim}

Before generating the final PDF, students shall check:

\begin{itemize}[leftmargin=0.8cm, itemsep=2pt]
    \item the Table of Contents;
    \item page numbering;
    \item chapter, section, and subsection numbering;
    \item figure numbering;
    \item table numbering;
    \item equation numbering;
    \item cross-references;
    \item citations;
    \item bibliography;
    \item font display;
    \item mathematical symbols;
    \item image quality;
    \item completeness of all required forms.
\end{itemize}

The document may be compiled using TeX Live, MiKTeX, or another
compatible \LaTeX{} environment. Regardless of the compilation
environment, the resulting PDF shall conform to the official Faculty
template.

\vspace{0.3cm}

% ================================================================
% XXII. CONTENT THAT STUDENTS MAY EDIT
% ================================================================

\noindent
\textbf{XXII.\quad CONTENT THAT STUDENTS MAY EDIT}

\vspace{0.15cm}

Students may edit or replace:

\begin{itemize}[leftmargin=0.8cm, itemsep=2pt]
    \item student name;
    \item class;
    \item internship organization;
    \item supervisor name;
    \item Abstract;
    \item List of Symbols and Abbreviations;
    \item chapter titles and chapter content;
    \item figures and photographs;
    \item tables;
    \item equations;
    \item citations and references.
\end{itemize}

Students shall not modify the main document formatting settings in the
preamble, including the document class, page geometry, line spacing,
heading formatting, table-of-contents formatting, font configuration,
page layout, or cover-page macros unless specifically instructed.

\vspace{0.3cm}

% ================================================================
% XXIII. FINAL SUBMISSION REQUIREMENTS
% ================================================================

\noindent
\textbf{XXIII.\quad FINAL SUBMISSION REQUIREMENTS}

\vspace{0.15cm}

Before submission, the final report shall:

\begin{itemize}[leftmargin=0.8cm, itemsep=2pt]
    \item compile without errors that affect the report content or presentation;
    \item contain all required forms and report sections;
    \item contain correct student, class, internship organization, and supervisor information;
    \item contain an updated Table of Contents, List of Figures, and List of Tables;
    \item contain no unused example content from the template;
    \item contain no placeholder text such as \textit{Title of Chapter 1},
    \textit{Section 1.1}, or other sample entries;
    \item contain all required citations and references;
    \item contain no unresolved cross-references.
\end{itemize}

If the Faculty requires submission of the \LaTeX{} source, the complete
project shall include all \verb|.tex| files, the bibliography database,
figures, and any other files required to reproduce the submitted PDF.

\vspace{0.4cm}

\noindent
\textbf{Important:}
Before submitting the final internship report, students shall remove or
comment out the following command in the main document:

\begin{verbatim}
\input{Regulations.tex}
\end{verbatim}

For example:

\begin{verbatim}
% \input{Regulations.tex}
\end{verbatim}

This removes the guidance pages from the final submitted report while
preserving the regulations file in the \LaTeX{} project for reference.

} % End 13pt content

\clearpage
 command before submitting
% the final internship report.
%
% =================================================================

\clearpage
\thispagestyle{empty}

\begin{center}
    {\fontsize{16pt}{20pt}\selectfont\bfseries
    SPECIFIC REGULATIONS FOR PREPARING\\
    GRADUATION INTERNSHIP REPORTS USING \LaTeX}
\end{center}

\vspace{0.5cm}

{\fontsize{13pt}{17pt}\selectfont
\setstretch{1.3}
\justifying

\noindent
\textbf{Note:}
This section is provided for guidance only and is included in the main
document using the command \verb|
% =================================================================
% REGULATIONS FOR PREPARING THE GRADUATION INTERNSHIP REPORT
% USING LATEX
% =================================================================
%
% This file is intended to be included from the main document with:
%
%     \input{Regulations.tex}
%
% IMPORTANT:
% This section is provided for guidance only. Students shall remove
% or comment out the \input{Regulations.tex} command before submitting
% the final internship report.
%
% =================================================================

\clearpage
\thispagestyle{empty}

\begin{center}
    {\fontsize{16pt}{20pt}\selectfont\bfseries
    SPECIFIC REGULATIONS FOR PREPARING\\
    GRADUATION INTERNSHIP REPORTS USING \LaTeX}
\end{center}

\vspace{0.5cm}

{\fontsize{13pt}{17pt}\selectfont
\setstretch{1.3}
\justifying

\noindent
\textbf{Note:}
This section is provided for guidance only and is included in the main
document using the command \verb|\input{Regulations.tex}|.
Students shall remove or comment out this command before submitting the
final internship report.

\vspace{0.4cm}

% ================================================================
% I. GENERAL REQUIREMENTS
% ================================================================

\noindent
\textbf{I.\quad GENERAL REQUIREMENTS}

\vspace{0.15cm}

The Graduation Internship Report shall be prepared using the official
\LaTeX{} template provided by the Faculty. Students shall preserve the
predefined page layout, typography, line spacing, heading styles,
numbering system, figure and table formatting, and bibliography settings
unless otherwise instructed by the Faculty or the supervisor.

The report shall be written in English.

Students are responsible for the academic and technical content of the
report. The official \LaTeX{} template is responsible for standardized
document formatting.

\vspace{0.3cm}

% ================================================================
% II. REPORT STRUCTURE
% ================================================================

\noindent
\textbf{II.\quad REPORT STRUCTURE}

\vspace{0.15cm}

The report shall normally be arranged in the following order:

\begin{enumerate}[label=\arabic*., leftmargin=1.2cm, itemsep=2pt]
    \item Outer Cover;
    \item Comments of the Internship Host Organization;
    \item Supervisor's Evaluation and Graduation Internship Defense Results;
    \item Table of Contents;
    \item List of Symbols and Abbreviations, if applicable;
    \item List of Figures;
    \item List of Tables;
    \item Abstract;
    \item Main Chapters of the Internship Report;
    \item References;
    \item Appendices, if applicable.
\end{enumerate}

The preliminary pages shall not be numbered as chapters. The main report
content shall be organized into consecutively numbered chapters beginning
with Chapter 1.

\vspace{0.3cm}

% ================================================================
% III. PAGE SIZE AND MARGINS
% ================================================================

\noindent
\textbf{III.\quad PAGE SIZE AND MARGINS}

\vspace{0.15cm}

The report shall be formatted on A4 paper with the following margins:

\begin{center}
\begin{tabular}{ll}
    Left margin:   & 3.5 cm \\
    Right margin:  & 2.0 cm \\
    Top margin:    & 2.5 cm \\
    Bottom margin: & 2.5 cm
\end{tabular}
\end{center}

These settings are predefined in the official template and shall not be
modified by students.

\vspace{0.3cm}

% ================================================================
% IV. FONT, FONT SIZE, AND LINE SPACING
% ================================================================

\noindent
\textbf{IV.\quad FONT, FONT SIZE, AND LINE SPACING}

\vspace{0.15cm}

The main text shall use a 13 pt font size with 1.3-line spacing.
A Times-family typeface shall be used in accordance with the official
template.

The text shall be fully justified.

Students shall not use repeated spaces, repeated manual line breaks, or
multiple blank lines to position text. Paragraph and heading spacing shall
be controlled by the template.

\vspace{0.3cm}

% ================================================================
% V. OUTER COVER
% ================================================================

\noindent
\textbf{V.\quad OUTER COVER}

\vspace{0.15cm}

The outer cover shall follow the official format provided in the template
and shall contain the following information:

\begin{itemize}[leftmargin=0.8cm, itemsep=2pt]
    \item Thai Nguyen University of Technology;
    \item Faculty of International Training;
    \item Engineering and Technology Department;
    \item University logo and Faculty logo;
    \item the title \textbf{INTERNSHIP REPORT};
    \item Internship Organization;
    \item Author;
    \item Class;
    \item Supervisor;
    \item place and year of completion.
\end{itemize}

Students shall only replace the required information fields and shall not
modify the cover-page layout unless instructed otherwise.

\vspace{0.3cm}

% ================================================================
% VI. COMMENTS OF THE INTERNSHIP HOST ORGANIZATION
% ================================================================

\noindent
\textbf{VI.\quad COMMENTS OF THE INTERNSHIP HOST ORGANIZATION}

\vspace{0.15cm}

The page entitled \textbf{COMMENTS OF THE INTERNSHIP HOST ORGANIZATION}
is reserved for the internship host organization.

It includes:

\begin{itemize}[leftmargin=0.8cm, itemsep=2pt]
    \item the student's full name;
    \item class;
    \item comments from the host organization;
    \item score;
    \item date;
    \item signature and official seal.
\end{itemize}

Students shall not alter or complete official evaluation fields unless
specifically authorized.

\vspace{0.3cm}

% ================================================================
% VII. SUPERVISOR'S EVALUATION AND DEFENSE RESULTS
% ================================================================

\noindent
\textbf{VII.\quad SUPERVISOR'S EVALUATION AND DEFENSE RESULTS}

\vspace{0.15cm}

The section entitled
\textbf{SUPERVISOR'S EVALUATION AND GRADUATION INTERNSHIP DEFENSE RESULTS}
is reserved for the supervisor and defense committee.

The evaluation may include:

\begin{itemize}[leftmargin=0.8cm, itemsep=2pt]
    \item Internship Progress and Student's Attitude;
    \item Report Content;
    \item Report Presentation;
    \item Other Comments;
    \item Supervisor's Opinion;
    \item Defense Result;
    \item Score.
\end{itemize}

Students shall not enter evaluation scores or defense results before they
are officially determined.

\vspace{0.3cm}

% ================================================================
% VIII. CHAPTERS, SECTIONS, AND SUBSECTIONS
% ================================================================

\noindent
\textbf{VIII.\quad CHAPTERS, SECTIONS, AND SUBSECTIONS}

\vspace{0.15cm}

The report structure shall be generated using the standard \LaTeX{}
commands:

\begin{verbatim}
\chapter{...}
\section{...}
\subsection{...}
\end{verbatim}

Chapter, section, and subsection numbers shall be generated automatically
by \LaTeX{}.

Students shall not manually type chapter or section numbers and shall not
reproduce heading styles using commands such as \verb|\textbf|,
\verb|\fontsize|, or \verb|\centering| instead of the appropriate
structural command.

Excessive subdivision should be avoided. In most cases, the hierarchy
Chapter--Section--Subsection is sufficient.

\vspace{0.3cm}

% ================================================================
% IX. ORGANIZATION OF THE MAIN CONTENT
% ================================================================

\noindent
\textbf{IX.\quad ORGANIZATION OF THE MAIN CONTENT}

\vspace{0.15cm}

The specific content of the report shall depend on the internship
organization and the assigned tasks. A typical structure may include:

\begin{itemize}[leftmargin=0.8cm, itemsep=2pt]
    \item \textbf{Chapter 1. Overview of the Internship Organization};
    \item \textbf{Chapter 2. Internship Activities and Assigned Tasks};
    \item \textbf{Chapter 3. Results, Experience, and Evaluation}.
\end{itemize}

These chapter titles are recommendations only and may be adapted to the
actual internship content.

The report shall reflect work actually performed, observed, analyzed, or
learned during the internship. Unrelated theoretical material should not
be included merely to increase report length.

\vspace{0.3cm}

% ================================================================
% X. TABLE OF CONTENTS
% ================================================================

\noindent
\textbf{X.\quad TABLE OF CONTENTS}

\vspace{0.15cm}

The Table of Contents shall be generated automatically using:

\begin{verbatim}
\tableofcontents
\end{verbatim}

Students shall not create or manually edit page numbers in the Table of
Contents.

The document shall be compiled sufficiently to ensure that all page
numbers and entries are correctly updated.

\vspace{0.3cm}

% ================================================================
% XI. LIST OF SYMBOLS AND ABBREVIATIONS
% ================================================================

\noindent
\textbf{XI.\quad LIST OF SYMBOLS AND ABBREVIATIONS}

\vspace{0.15cm}

Symbols and abbreviations that are used repeatedly in the report shall be
listed in the \textbf{List of Symbols and Abbreviations}.

The list may contain the following columns:

\begin{center}
\begin{tabular}{llll}
    No. & Symbol / Abbreviation & Unit & Definition
\end{tabular}
\end{center}

Only symbols and abbreviations actually used in the final report shall be
retained. Example entries contained in the template shall be removed if
they are not relevant to the student's report.

Each abbreviation should normally be defined when it first appears in
the main text.

\vspace{0.3cm}

% ================================================================
% XII. FIGURES, DIAGRAMS, AND PHOTOGRAPHS
% ================================================================

\noindent
\textbf{XII.\quad FIGURES, DIAGRAMS, AND PHOTOGRAPHS}

\vspace{0.15cm}

Figures shall be clear, legible, and directly relevant to the report
content.

Figures shall be inserted using the standard \verb|figure| environment.

Example:

\begin{verbatim}
\begin{figure}[H]
    \centering
    \includegraphics[width=0.8\textwidth]
                    {figures/company_structure.pdf}
    \caption{Organizational structure of the company}
    \label{fig:company-structure}
\end{figure}
\end{verbatim}

Figure captions shall be placed below the figure.

Each referenced figure shall have a unique \verb|\label| and shall be
referenced in the text using, for example:

\begin{verbatim}
Figure~\ref{fig:company-structure}
\end{verbatim}

Students shall not manually type figure numbers.

Figures reproduced or adapted from external sources shall include an
appropriate citation.

Photographs taken during the internship shall be relevant and sufficiently
clear. Students shall not include confidential, restricted, or proprietary
information without permission from the internship organization.

Vector graphics or high-resolution images are recommended whenever
possible.

\vspace{0.3cm}

% ================================================================
% XIII. TABLES
% ================================================================

\noindent
\textbf{XIII.\quad TABLES}

\vspace{0.15cm}

Tables shall be created using suitable \LaTeX{} environments such as
\verb|table|, \verb|tabular|, \verb|tabularx|, or \verb|longtable|.

Table captions shall be placed above the table.

Example:

\begin{verbatim}
\begin{table}[H]
    \centering
    \caption{Internship schedule}
    \label{tab:internship-schedule}

    \begin{tabular}{lll}
        \toprule
        Week & Task & Result \\
        \midrule
        1 & Orientation & Completed \\
        2 & System study & Completed \\
        \bottomrule
    \end{tabular}
\end{table}
\end{verbatim}

Tables shall be referenced using:

\begin{verbatim}
Table~\ref{tab:internship-schedule}
\end{verbatim}

Students shall not manually type table numbers.

Tables should be created directly in \LaTeX{} rather than inserted as
screenshots whenever practical.

\vspace{0.3cm}

% ================================================================
% XIV. EQUATIONS AND TECHNICAL QUANTITIES
% ================================================================

\noindent
\textbf{XIV.\quad EQUATIONS AND TECHNICAL QUANTITIES}

\vspace{0.15cm}

Mathematical expressions shall be typeset directly using \LaTeX{}.
Images or screenshots of equations shall not be used.

Short expressions may be written in inline mathematical mode, for example:

\begin{verbatim}
$R = 0.7\,\Omega$
\end{verbatim}

Important equations shall use a numbered equation environment.

Example:

\begin{verbatim}
\begin{equation}
    P = UI
    \label{eq:power}
\end{equation}
\end{verbatim}

Equations shall be referenced using:

\begin{verbatim}
Equation~\eqref{eq:power}
\end{verbatim}

Equation numbers shall not be manually entered.

SI units shall be used wherever applicable. The \verb|siunitx| package
should be used where appropriate.

Examples:

\begin{verbatim}
\SI{220}{\volt}
\SI{50}{\hertz}
\SI{2.5}{\ampere}
\end{verbatim}

Notation and units shall be consistent throughout the report.

\vspace{0.3cm}

% ================================================================
% XV. SOURCE CODE
% ================================================================

\noindent
\textbf{XV.\quad SOURCE CODE}

\vspace{0.15cm}

If the internship involves programming, PLCs, microcontrollers,
simulation, or software development, relevant source code may be included
in the report.

Source code shall be presented using the \verb|listings| environment or
another environment defined by the official template.

Example:

\begin{verbatim}
\begin{lstlisting}[language=st]
IF Start = TRUE THEN
    Motor := TRUE;
ELSE
    Motor := FALSE;
END_IF;
\end{lstlisting}
\end{verbatim}

Only code necessary for explaining the work should appear in the main
report. Long source-code listings should be moved to an appendix or
submitted as supplementary files if required.

Source code should not normally be inserted as screenshots.

\vspace{0.3cm}

% ================================================================
% XVI. ABSTRACT
% ================================================================

\noindent
\textbf{XVI.\quad ABSTRACT}

\vspace{0.15cm}

The Abstract shall briefly summarize:

\begin{itemize}[leftmargin=0.8cm, itemsep=2pt]
    \item the internship organization;
    \item the objectives of the internship;
    \item the main tasks performed;
    \item the main results, knowledge, and skills obtained.
\end{itemize}

The Abstract should be concise and should not reproduce the full content
of the report chapter by chapter.

\vspace{0.3cm}

% ================================================================
% XVII. REFERENCES AND CITATIONS
% ================================================================

\noindent
\textbf{XVII.\quad REFERENCES AND CITATIONS}

\vspace{0.15cm}

All books, journal papers, conference papers, standards, websites,
datasheets, manuals, technical documents, figures, tables, or other
external sources used in preparing the report shall be properly cited.

References shall be managed through the bibliography system provided with
the official template.

Bibliographic entries shall normally be stored in:

\begin{verbatim}
References.bib
\end{verbatim}

References shall be cited in the report using:

\begin{verbatim}
\cite{reference-key}
\end{verbatim}

Students shall not manually type reference numbers such as ``[1]'',
``[2]'', etc.

Titles of foreign-language publications shall remain in their original
language in the bibliography.

All borrowed ideas, data, figures, tables, equations, methods, or results
shall be properly acknowledged.

\vspace{0.3cm}

% ================================================================
% XVIII. CROSS-REFERENCING
% ================================================================

\noindent
\textbf{XVIII.\quad CROSS-REFERENCING}

\vspace{0.15cm}

Numbered elements shall be cross-referenced using the standard \LaTeX{}
mechanism.

Use the following commands where appropriate:

\begin{verbatim}
\label{...}
\ref{...}
\eqref{...}
\end{verbatim}

For example:

\begin{verbatim}
As shown in Figure~\ref{fig:system}, ...
\end{verbatim}

Students shall not manually type numbers such as ``Figure 2.3'' when
those numbers can be generated automatically.

\vspace{0.3cm}

% ================================================================
% XIX. PAGE NUMBERING
% ================================================================

\noindent
\textbf{XIX.\quad PAGE NUMBERING}

\vspace{0.15cm}

Cover pages and official evaluation forms at the beginning of the report
shall not display page numbers.

The main body of the report shall use Arabic page numbers beginning from
Chapter 1.

Page numbering shall be generated automatically by the official template.
Students shall not manually enter page numbers.

\vspace{0.3cm}

% ================================================================
% XX. PROJECT FILE ORGANIZATION
% ================================================================

\noindent
\textbf{XX.\quad PROJECT FILE ORGANIZATION}

\vspace{0.15cm}

Students should maintain a clear and organized \LaTeX{} project
structure. A recommended structure is:

\begin{verbatim}
InternshipReport/
|
|-- main.tex
|-- References.bib
|
|-- figures/
|   |-- company.pdf
|   |-- system.png
|   `-- ...
|
`-- chapters/
    |-- chapter1.tex
    |-- chapter2.tex
    `-- chapter3.tex
\end{verbatim}

Figures should normally be stored in the \verb|figures| directory
specified by the official template.

File and directory names should preferably contain Latin letters,
numbers, hyphens, or underscores.

Recommended examples include:

\begin{verbatim}
company_structure.pdf
motor_control.png
chapter_01.tex
\end{verbatim}

Spaces and unnecessary special characters in file names should be
avoided where possible.

\vspace{0.3cm}

% ================================================================
% XXI. DOCUMENT COMPILATION AND FINAL CHECK
% ================================================================

\noindent
\textbf{XXI.\quad DOCUMENT COMPILATION AND FINAL CHECK}

\vspace{0.15cm}

Before submission, the complete \LaTeX{} project shall compile
successfully and produce the final PDF report.

The final document shall not contain unresolved references such as:

\begin{verbatim}
??
[?]
\end{verbatim}

Before generating the final PDF, students shall check:

\begin{itemize}[leftmargin=0.8cm, itemsep=2pt]
    \item the Table of Contents;
    \item page numbering;
    \item chapter, section, and subsection numbering;
    \item figure numbering;
    \item table numbering;
    \item equation numbering;
    \item cross-references;
    \item citations;
    \item bibliography;
    \item font display;
    \item mathematical symbols;
    \item image quality;
    \item completeness of all required forms.
\end{itemize}

The document may be compiled using TeX Live, MiKTeX, or another
compatible \LaTeX{} environment. Regardless of the compilation
environment, the resulting PDF shall conform to the official Faculty
template.

\vspace{0.3cm}

% ================================================================
% XXII. CONTENT THAT STUDENTS MAY EDIT
% ================================================================

\noindent
\textbf{XXII.\quad CONTENT THAT STUDENTS MAY EDIT}

\vspace{0.15cm}

Students may edit or replace:

\begin{itemize}[leftmargin=0.8cm, itemsep=2pt]
    \item student name;
    \item class;
    \item internship organization;
    \item supervisor name;
    \item Abstract;
    \item List of Symbols and Abbreviations;
    \item chapter titles and chapter content;
    \item figures and photographs;
    \item tables;
    \item equations;
    \item citations and references.
\end{itemize}

Students shall not modify the main document formatting settings in the
preamble, including the document class, page geometry, line spacing,
heading formatting, table-of-contents formatting, font configuration,
page layout, or cover-page macros unless specifically instructed.

\vspace{0.3cm}

% ================================================================
% XXIII. FINAL SUBMISSION REQUIREMENTS
% ================================================================

\noindent
\textbf{XXIII.\quad FINAL SUBMISSION REQUIREMENTS}

\vspace{0.15cm}

Before submission, the final report shall:

\begin{itemize}[leftmargin=0.8cm, itemsep=2pt]
    \item compile without errors that affect the report content or presentation;
    \item contain all required forms and report sections;
    \item contain correct student, class, internship organization, and supervisor information;
    \item contain an updated Table of Contents, List of Figures, and List of Tables;
    \item contain no unused example content from the template;
    \item contain no placeholder text such as \textit{Title of Chapter 1},
    \textit{Section 1.1}, or other sample entries;
    \item contain all required citations and references;
    \item contain no unresolved cross-references.
\end{itemize}

If the Faculty requires submission of the \LaTeX{} source, the complete
project shall include all \verb|.tex| files, the bibliography database,
figures, and any other files required to reproduce the submitted PDF.

\vspace{0.4cm}

\noindent
\textbf{Important:}
Before submitting the final internship report, students shall remove or
comment out the following command in the main document:

\begin{verbatim}
\input{Regulations.tex}
\end{verbatim}

For example:

\begin{verbatim}
% \input{Regulations.tex}
\end{verbatim}

This removes the guidance pages from the final submitted report while
preserving the regulations file in the \LaTeX{} project for reference.

} % End 13pt content

\clearpage
|.
Students shall remove or comment out this command before submitting the
final internship report.

\vspace{0.4cm}

% ================================================================
% I. GENERAL REQUIREMENTS
% ================================================================

\noindent
\textbf{I.\quad GENERAL REQUIREMENTS}

\vspace{0.15cm}

The Graduation Internship Report shall be prepared using the official
\LaTeX{} template provided by the Faculty. Students shall preserve the
predefined page layout, typography, line spacing, heading styles,
numbering system, figure and table formatting, and bibliography settings
unless otherwise instructed by the Faculty or the supervisor.

The report shall be written in English.

Students are responsible for the academic and technical content of the
report. The official \LaTeX{} template is responsible for standardized
document formatting.

\vspace{0.3cm}

% ================================================================
% II. REPORT STRUCTURE
% ================================================================

\noindent
\textbf{II.\quad REPORT STRUCTURE}

\vspace{0.15cm}

The report shall normally be arranged in the following order:

\begin{enumerate}[label=\arabic*., leftmargin=1.2cm, itemsep=2pt]
    \item Outer Cover;
    \item Comments of the Internship Host Organization;
    \item Supervisor's Evaluation and Graduation Internship Defense Results;
    \item Table of Contents;
    \item List of Symbols and Abbreviations, if applicable;
    \item List of Figures;
    \item List of Tables;
    \item Abstract;
    \item Main Chapters of the Internship Report;
    \item References;
    \item Appendices, if applicable.
\end{enumerate}

The preliminary pages shall not be numbered as chapters. The main report
content shall be organized into consecutively numbered chapters beginning
with Chapter 1.

\vspace{0.3cm}

% ================================================================
% III. PAGE SIZE AND MARGINS
% ================================================================

\noindent
\textbf{III.\quad PAGE SIZE AND MARGINS}

\vspace{0.15cm}

The report shall be formatted on A4 paper with the following margins:

\begin{center}
\begin{tabular}{ll}
    Left margin:   & 3.5 cm \\
    Right margin:  & 2.0 cm \\
    Top margin:    & 2.5 cm \\
    Bottom margin: & 2.5 cm
\end{tabular}
\end{center}

These settings are predefined in the official template and shall not be
modified by students.

\vspace{0.3cm}

% ================================================================
% IV. FONT, FONT SIZE, AND LINE SPACING
% ================================================================

\noindent
\textbf{IV.\quad FONT, FONT SIZE, AND LINE SPACING}

\vspace{0.15cm}

The main text shall use a 13 pt font size with 1.3-line spacing.
A Times-family typeface shall be used in accordance with the official
template.

The text shall be fully justified.

Students shall not use repeated spaces, repeated manual line breaks, or
multiple blank lines to position text. Paragraph and heading spacing shall
be controlled by the template.

\vspace{0.3cm}

% ================================================================
% V. OUTER COVER
% ================================================================

\noindent
\textbf{V.\quad OUTER COVER}

\vspace{0.15cm}

The outer cover shall follow the official format provided in the template
and shall contain the following information:

\begin{itemize}[leftmargin=0.8cm, itemsep=2pt]
    \item Thai Nguyen University of Technology;
    \item Faculty of International Training;
    \item Engineering and Technology Department;
    \item University logo and Faculty logo;
    \item the title \textbf{INTERNSHIP REPORT};
    \item Internship Organization;
    \item Author;
    \item Class;
    \item Supervisor;
    \item place and year of completion.
\end{itemize}

Students shall only replace the required information fields and shall not
modify the cover-page layout unless instructed otherwise.

\vspace{0.3cm}

% ================================================================
% VI. COMMENTS OF THE INTERNSHIP HOST ORGANIZATION
% ================================================================

\noindent
\textbf{VI.\quad COMMENTS OF THE INTERNSHIP HOST ORGANIZATION}

\vspace{0.15cm}

The page entitled \textbf{COMMENTS OF THE INTERNSHIP HOST ORGANIZATION}
is reserved for the internship host organization.

It includes:

\begin{itemize}[leftmargin=0.8cm, itemsep=2pt]
    \item the student's full name;
    \item class;
    \item comments from the host organization;
    \item score;
    \item date;
    \item signature and official seal.
\end{itemize}

Students shall not alter or complete official evaluation fields unless
specifically authorized.

\vspace{0.3cm}

% ================================================================
% VII. SUPERVISOR'S EVALUATION AND DEFENSE RESULTS
% ================================================================

\noindent
\textbf{VII.\quad SUPERVISOR'S EVALUATION AND DEFENSE RESULTS}

\vspace{0.15cm}

The section entitled
\textbf{SUPERVISOR'S EVALUATION AND GRADUATION INTERNSHIP DEFENSE RESULTS}
is reserved for the supervisor and defense committee.

The evaluation may include:

\begin{itemize}[leftmargin=0.8cm, itemsep=2pt]
    \item Internship Progress and Student's Attitude;
    \item Report Content;
    \item Report Presentation;
    \item Other Comments;
    \item Supervisor's Opinion;
    \item Defense Result;
    \item Score.
\end{itemize}

Students shall not enter evaluation scores or defense results before they
are officially determined.

\vspace{0.3cm}

% ================================================================
% VIII. CHAPTERS, SECTIONS, AND SUBSECTIONS
% ================================================================

\noindent
\textbf{VIII.\quad CHAPTERS, SECTIONS, AND SUBSECTIONS}

\vspace{0.15cm}

The report structure shall be generated using the standard \LaTeX{}
commands:

\begin{verbatim}
\chapter{...}
\section{...}
\subsection{...}
\end{verbatim}

Chapter, section, and subsection numbers shall be generated automatically
by \LaTeX{}.

Students shall not manually type chapter or section numbers and shall not
reproduce heading styles using commands such as \verb|\textbf|,
\verb|\fontsize|, or \verb|\centering| instead of the appropriate
structural command.

Excessive subdivision should be avoided. In most cases, the hierarchy
Chapter--Section--Subsection is sufficient.

\vspace{0.3cm}

% ================================================================
% IX. ORGANIZATION OF THE MAIN CONTENT
% ================================================================

\noindent
\textbf{IX.\quad ORGANIZATION OF THE MAIN CONTENT}

\vspace{0.15cm}

The specific content of the report shall depend on the internship
organization and the assigned tasks. A typical structure may include:

\begin{itemize}[leftmargin=0.8cm, itemsep=2pt]
    \item \textbf{Chapter 1. Overview of the Internship Organization};
    \item \textbf{Chapter 2. Internship Activities and Assigned Tasks};
    \item \textbf{Chapter 3. Results, Experience, and Evaluation}.
\end{itemize}

These chapter titles are recommendations only and may be adapted to the
actual internship content.

The report shall reflect work actually performed, observed, analyzed, or
learned during the internship. Unrelated theoretical material should not
be included merely to increase report length.

\vspace{0.3cm}

% ================================================================
% X. TABLE OF CONTENTS
% ================================================================

\noindent
\textbf{X.\quad TABLE OF CONTENTS}

\vspace{0.15cm}

The Table of Contents shall be generated automatically using:

\begin{verbatim}
\tableofcontents
\end{verbatim}

Students shall not create or manually edit page numbers in the Table of
Contents.

The document shall be compiled sufficiently to ensure that all page
numbers and entries are correctly updated.

\vspace{0.3cm}

% ================================================================
% XI. LIST OF SYMBOLS AND ABBREVIATIONS
% ================================================================

\noindent
\textbf{XI.\quad LIST OF SYMBOLS AND ABBREVIATIONS}

\vspace{0.15cm}

Symbols and abbreviations that are used repeatedly in the report shall be
listed in the \textbf{List of Symbols and Abbreviations}.

The list may contain the following columns:

\begin{center}
\begin{tabular}{llll}
    No. & Symbol / Abbreviation & Unit & Definition
\end{tabular}
\end{center}

Only symbols and abbreviations actually used in the final report shall be
retained. Example entries contained in the template shall be removed if
they are not relevant to the student's report.

Each abbreviation should normally be defined when it first appears in
the main text.

\vspace{0.3cm}

% ================================================================
% XII. FIGURES, DIAGRAMS, AND PHOTOGRAPHS
% ================================================================

\noindent
\textbf{XII.\quad FIGURES, DIAGRAMS, AND PHOTOGRAPHS}

\vspace{0.15cm}

Figures shall be clear, legible, and directly relevant to the report
content.

Figures shall be inserted using the standard \verb|figure| environment.

Example:

\begin{verbatim}
\begin{figure}[H]
    \centering
    \includegraphics[width=0.8\textwidth]
                    {figures/company_structure.pdf}
    \caption{Organizational structure of the company}
    \label{fig:company-structure}
\end{figure}
\end{verbatim}

Figure captions shall be placed below the figure.

Each referenced figure shall have a unique \verb|\label| and shall be
referenced in the text using, for example:

\begin{verbatim}
Figure~\ref{fig:company-structure}
\end{verbatim}

Students shall not manually type figure numbers.

Figures reproduced or adapted from external sources shall include an
appropriate citation.

Photographs taken during the internship shall be relevant and sufficiently
clear. Students shall not include confidential, restricted, or proprietary
information without permission from the internship organization.

Vector graphics or high-resolution images are recommended whenever
possible.

\vspace{0.3cm}

% ================================================================
% XIII. TABLES
% ================================================================

\noindent
\textbf{XIII.\quad TABLES}

\vspace{0.15cm}

Tables shall be created using suitable \LaTeX{} environments such as
\verb|table|, \verb|tabular|, \verb|tabularx|, or \verb|longtable|.

Table captions shall be placed above the table.

Example:

\begin{verbatim}
\begin{table}[H]
    \centering
    \caption{Internship schedule}
    \label{tab:internship-schedule}

    \begin{tabular}{lll}
        \toprule
        Week & Task & Result \\
        \midrule
        1 & Orientation & Completed \\
        2 & System study & Completed \\
        \bottomrule
    \end{tabular}
\end{table}
\end{verbatim}

Tables shall be referenced using:

\begin{verbatim}
Table~\ref{tab:internship-schedule}
\end{verbatim}

Students shall not manually type table numbers.

Tables should be created directly in \LaTeX{} rather than inserted as
screenshots whenever practical.

\vspace{0.3cm}

% ================================================================
% XIV. EQUATIONS AND TECHNICAL QUANTITIES
% ================================================================

\noindent
\textbf{XIV.\quad EQUATIONS AND TECHNICAL QUANTITIES}

\vspace{0.15cm}

Mathematical expressions shall be typeset directly using \LaTeX{}.
Images or screenshots of equations shall not be used.

Short expressions may be written in inline mathematical mode, for example:

\begin{verbatim}
$R = 0.7\,\Omega$
\end{verbatim}

Important equations shall use a numbered equation environment.

Example:

\begin{verbatim}
\begin{equation}
    P = UI
    \label{eq:power}
\end{equation}
\end{verbatim}

Equations shall be referenced using:

\begin{verbatim}
Equation~\eqref{eq:power}
\end{verbatim}

Equation numbers shall not be manually entered.

SI units shall be used wherever applicable. The \verb|siunitx| package
should be used where appropriate.

Examples:

\begin{verbatim}
\SI{220}{\volt}
\SI{50}{\hertz}
\SI{2.5}{\ampere}
\end{verbatim}

Notation and units shall be consistent throughout the report.

\vspace{0.3cm}

% ================================================================
% XV. SOURCE CODE
% ================================================================

\noindent
\textbf{XV.\quad SOURCE CODE}

\vspace{0.15cm}

If the internship involves programming, PLCs, microcontrollers,
simulation, or software development, relevant source code may be included
in the report.

Source code shall be presented using the \verb|listings| environment or
another environment defined by the official template.

Example:

\begin{verbatim}
\begin{lstlisting}[language=st]
IF Start = TRUE THEN
    Motor := TRUE;
ELSE
    Motor := FALSE;
END_IF;
\end{lstlisting}
\end{verbatim}

Only code necessary for explaining the work should appear in the main
report. Long source-code listings should be moved to an appendix or
submitted as supplementary files if required.

Source code should not normally be inserted as screenshots.

\vspace{0.3cm}

% ================================================================
% XVI. ABSTRACT
% ================================================================

\noindent
\textbf{XVI.\quad ABSTRACT}

\vspace{0.15cm}

The Abstract shall briefly summarize:

\begin{itemize}[leftmargin=0.8cm, itemsep=2pt]
    \item the internship organization;
    \item the objectives of the internship;
    \item the main tasks performed;
    \item the main results, knowledge, and skills obtained.
\end{itemize}

The Abstract should be concise and should not reproduce the full content
of the report chapter by chapter.

\vspace{0.3cm}

% ================================================================
% XVII. REFERENCES AND CITATIONS
% ================================================================

\noindent
\textbf{XVII.\quad REFERENCES AND CITATIONS}

\vspace{0.15cm}

All books, journal papers, conference papers, standards, websites,
datasheets, manuals, technical documents, figures, tables, or other
external sources used in preparing the report shall be properly cited.

References shall be managed through the bibliography system provided with
the official template.

Bibliographic entries shall normally be stored in:

\begin{verbatim}
References.bib
\end{verbatim}

References shall be cited in the report using:

\begin{verbatim}
\cite{reference-key}
\end{verbatim}

Students shall not manually type reference numbers such as ``[1]'',
``[2]'', etc.

Titles of foreign-language publications shall remain in their original
language in the bibliography.

All borrowed ideas, data, figures, tables, equations, methods, or results
shall be properly acknowledged.

\vspace{0.3cm}

% ================================================================
% XVIII. CROSS-REFERENCING
% ================================================================

\noindent
\textbf{XVIII.\quad CROSS-REFERENCING}

\vspace{0.15cm}

Numbered elements shall be cross-referenced using the standard \LaTeX{}
mechanism.

Use the following commands where appropriate:

\begin{verbatim}
\label{...}
\ref{...}
\eqref{...}
\end{verbatim}

For example:

\begin{verbatim}
As shown in Figure~\ref{fig:system}, ...
\end{verbatim}

Students shall not manually type numbers such as ``Figure 2.3'' when
those numbers can be generated automatically.

\vspace{0.3cm}

% ================================================================
% XIX. PAGE NUMBERING
% ================================================================

\noindent
\textbf{XIX.\quad PAGE NUMBERING}

\vspace{0.15cm}

Cover pages and official evaluation forms at the beginning of the report
shall not display page numbers.

The main body of the report shall use Arabic page numbers beginning from
Chapter 1.

Page numbering shall be generated automatically by the official template.
Students shall not manually enter page numbers.

\vspace{0.3cm}

% ================================================================
% XX. PROJECT FILE ORGANIZATION
% ================================================================

\noindent
\textbf{XX.\quad PROJECT FILE ORGANIZATION}

\vspace{0.15cm}

Students should maintain a clear and organized \LaTeX{} project
structure. A recommended structure is:

\begin{verbatim}
InternshipReport/
|
|-- main.tex
|-- References.bib
|
|-- figures/
|   |-- company.pdf
|   |-- system.png
|   `-- ...
|
`-- chapters/
    |-- chapter1.tex
    |-- chapter2.tex
    `-- chapter3.tex
\end{verbatim}

Figures should normally be stored in the \verb|figures| directory
specified by the official template.

File and directory names should preferably contain Latin letters,
numbers, hyphens, or underscores.

Recommended examples include:

\begin{verbatim}
company_structure.pdf
motor_control.png
chapter_01.tex
\end{verbatim}

Spaces and unnecessary special characters in file names should be
avoided where possible.

\vspace{0.3cm}

% ================================================================
% XXI. DOCUMENT COMPILATION AND FINAL CHECK
% ================================================================

\noindent
\textbf{XXI.\quad DOCUMENT COMPILATION AND FINAL CHECK}

\vspace{0.15cm}

Before submission, the complete \LaTeX{} project shall compile
successfully and produce the final PDF report.

The final document shall not contain unresolved references such as:

\begin{verbatim}
??
[?]
\end{verbatim}

Before generating the final PDF, students shall check:

\begin{itemize}[leftmargin=0.8cm, itemsep=2pt]
    \item the Table of Contents;
    \item page numbering;
    \item chapter, section, and subsection numbering;
    \item figure numbering;
    \item table numbering;
    \item equation numbering;
    \item cross-references;
    \item citations;
    \item bibliography;
    \item font display;
    \item mathematical symbols;
    \item image quality;
    \item completeness of all required forms.
\end{itemize}

The document may be compiled using TeX Live, MiKTeX, or another
compatible \LaTeX{} environment. Regardless of the compilation
environment, the resulting PDF shall conform to the official Faculty
template.

\vspace{0.3cm}

% ================================================================
% XXII. CONTENT THAT STUDENTS MAY EDIT
% ================================================================

\noindent
\textbf{XXII.\quad CONTENT THAT STUDENTS MAY EDIT}

\vspace{0.15cm}

Students may edit or replace:

\begin{itemize}[leftmargin=0.8cm, itemsep=2pt]
    \item student name;
    \item class;
    \item internship organization;
    \item supervisor name;
    \item Abstract;
    \item List of Symbols and Abbreviations;
    \item chapter titles and chapter content;
    \item figures and photographs;
    \item tables;
    \item equations;
    \item citations and references.
\end{itemize}

Students shall not modify the main document formatting settings in the
preamble, including the document class, page geometry, line spacing,
heading formatting, table-of-contents formatting, font configuration,
page layout, or cover-page macros unless specifically instructed.

\vspace{0.3cm}

% ================================================================
% XXIII. FINAL SUBMISSION REQUIREMENTS
% ================================================================

\noindent
\textbf{XXIII.\quad FINAL SUBMISSION REQUIREMENTS}

\vspace{0.15cm}

Before submission, the final report shall:

\begin{itemize}[leftmargin=0.8cm, itemsep=2pt]
    \item compile without errors that affect the report content or presentation;
    \item contain all required forms and report sections;
    \item contain correct student, class, internship organization, and supervisor information;
    \item contain an updated Table of Contents, List of Figures, and List of Tables;
    \item contain no unused example content from the template;
    \item contain no placeholder text such as \textit{Title of Chapter 1},
    \textit{Section 1.1}, or other sample entries;
    \item contain all required citations and references;
    \item contain no unresolved cross-references.
\end{itemize}

If the Faculty requires submission of the \LaTeX{} source, the complete
project shall include all \verb|.tex| files, the bibliography database,
figures, and any other files required to reproduce the submitted PDF.

\vspace{0.4cm}

\noindent
\textbf{Important:}
Before submitting the final internship report, students shall remove or
comment out the following command in the main document:

\begin{verbatim}

% =================================================================
% REGULATIONS FOR PREPARING THE GRADUATION INTERNSHIP REPORT
% USING LATEX
% =================================================================
%
% This file is intended to be included from the main document with:
%
%     \input{Regulations.tex}
%
% IMPORTANT:
% This section is provided for guidance only. Students shall remove
% or comment out the \input{Regulations.tex} command before submitting
% the final internship report.
%
% =================================================================

\clearpage
\thispagestyle{empty}

\begin{center}
    {\fontsize{16pt}{20pt}\selectfont\bfseries
    SPECIFIC REGULATIONS FOR PREPARING\\
    GRADUATION INTERNSHIP REPORTS USING \LaTeX}
\end{center}

\vspace{0.5cm}

{\fontsize{13pt}{17pt}\selectfont
\setstretch{1.3}
\justifying

\noindent
\textbf{Note:}
This section is provided for guidance only and is included in the main
document using the command \verb|\input{Regulations.tex}|.
Students shall remove or comment out this command before submitting the
final internship report.

\vspace{0.4cm}

% ================================================================
% I. GENERAL REQUIREMENTS
% ================================================================

\noindent
\textbf{I.\quad GENERAL REQUIREMENTS}

\vspace{0.15cm}

The Graduation Internship Report shall be prepared using the official
\LaTeX{} template provided by the Faculty. Students shall preserve the
predefined page layout, typography, line spacing, heading styles,
numbering system, figure and table formatting, and bibliography settings
unless otherwise instructed by the Faculty or the supervisor.

The report shall be written in English.

Students are responsible for the academic and technical content of the
report. The official \LaTeX{} template is responsible for standardized
document formatting.

\vspace{0.3cm}

% ================================================================
% II. REPORT STRUCTURE
% ================================================================

\noindent
\textbf{II.\quad REPORT STRUCTURE}

\vspace{0.15cm}

The report shall normally be arranged in the following order:

\begin{enumerate}[label=\arabic*., leftmargin=1.2cm, itemsep=2pt]
    \item Outer Cover;
    \item Comments of the Internship Host Organization;
    \item Supervisor's Evaluation and Graduation Internship Defense Results;
    \item Table of Contents;
    \item List of Symbols and Abbreviations, if applicable;
    \item List of Figures;
    \item List of Tables;
    \item Abstract;
    \item Main Chapters of the Internship Report;
    \item References;
    \item Appendices, if applicable.
\end{enumerate}

The preliminary pages shall not be numbered as chapters. The main report
content shall be organized into consecutively numbered chapters beginning
with Chapter 1.

\vspace{0.3cm}

% ================================================================
% III. PAGE SIZE AND MARGINS
% ================================================================

\noindent
\textbf{III.\quad PAGE SIZE AND MARGINS}

\vspace{0.15cm}

The report shall be formatted on A4 paper with the following margins:

\begin{center}
\begin{tabular}{ll}
    Left margin:   & 3.5 cm \\
    Right margin:  & 2.0 cm \\
    Top margin:    & 2.5 cm \\
    Bottom margin: & 2.5 cm
\end{tabular}
\end{center}

These settings are predefined in the official template and shall not be
modified by students.

\vspace{0.3cm}

% ================================================================
% IV. FONT, FONT SIZE, AND LINE SPACING
% ================================================================

\noindent
\textbf{IV.\quad FONT, FONT SIZE, AND LINE SPACING}

\vspace{0.15cm}

The main text shall use a 13 pt font size with 1.3-line spacing.
A Times-family typeface shall be used in accordance with the official
template.

The text shall be fully justified.

Students shall not use repeated spaces, repeated manual line breaks, or
multiple blank lines to position text. Paragraph and heading spacing shall
be controlled by the template.

\vspace{0.3cm}

% ================================================================
% V. OUTER COVER
% ================================================================

\noindent
\textbf{V.\quad OUTER COVER}

\vspace{0.15cm}

The outer cover shall follow the official format provided in the template
and shall contain the following information:

\begin{itemize}[leftmargin=0.8cm, itemsep=2pt]
    \item Thai Nguyen University of Technology;
    \item Faculty of International Training;
    \item Engineering and Technology Department;
    \item University logo and Faculty logo;
    \item the title \textbf{INTERNSHIP REPORT};
    \item Internship Organization;
    \item Author;
    \item Class;
    \item Supervisor;
    \item place and year of completion.
\end{itemize}

Students shall only replace the required information fields and shall not
modify the cover-page layout unless instructed otherwise.

\vspace{0.3cm}

% ================================================================
% VI. COMMENTS OF THE INTERNSHIP HOST ORGANIZATION
% ================================================================

\noindent
\textbf{VI.\quad COMMENTS OF THE INTERNSHIP HOST ORGANIZATION}

\vspace{0.15cm}

The page entitled \textbf{COMMENTS OF THE INTERNSHIP HOST ORGANIZATION}
is reserved for the internship host organization.

It includes:

\begin{itemize}[leftmargin=0.8cm, itemsep=2pt]
    \item the student's full name;
    \item class;
    \item comments from the host organization;
    \item score;
    \item date;
    \item signature and official seal.
\end{itemize}

Students shall not alter or complete official evaluation fields unless
specifically authorized.

\vspace{0.3cm}

% ================================================================
% VII. SUPERVISOR'S EVALUATION AND DEFENSE RESULTS
% ================================================================

\noindent
\textbf{VII.\quad SUPERVISOR'S EVALUATION AND DEFENSE RESULTS}

\vspace{0.15cm}

The section entitled
\textbf{SUPERVISOR'S EVALUATION AND GRADUATION INTERNSHIP DEFENSE RESULTS}
is reserved for the supervisor and defense committee.

The evaluation may include:

\begin{itemize}[leftmargin=0.8cm, itemsep=2pt]
    \item Internship Progress and Student's Attitude;
    \item Report Content;
    \item Report Presentation;
    \item Other Comments;
    \item Supervisor's Opinion;
    \item Defense Result;
    \item Score.
\end{itemize}

Students shall not enter evaluation scores or defense results before they
are officially determined.

\vspace{0.3cm}

% ================================================================
% VIII. CHAPTERS, SECTIONS, AND SUBSECTIONS
% ================================================================

\noindent
\textbf{VIII.\quad CHAPTERS, SECTIONS, AND SUBSECTIONS}

\vspace{0.15cm}

The report structure shall be generated using the standard \LaTeX{}
commands:

\begin{verbatim}
\chapter{...}
\section{...}
\subsection{...}
\end{verbatim}

Chapter, section, and subsection numbers shall be generated automatically
by \LaTeX{}.

Students shall not manually type chapter or section numbers and shall not
reproduce heading styles using commands such as \verb|\textbf|,
\verb|\fontsize|, or \verb|\centering| instead of the appropriate
structural command.

Excessive subdivision should be avoided. In most cases, the hierarchy
Chapter--Section--Subsection is sufficient.

\vspace{0.3cm}

% ================================================================
% IX. ORGANIZATION OF THE MAIN CONTENT
% ================================================================

\noindent
\textbf{IX.\quad ORGANIZATION OF THE MAIN CONTENT}

\vspace{0.15cm}

The specific content of the report shall depend on the internship
organization and the assigned tasks. A typical structure may include:

\begin{itemize}[leftmargin=0.8cm, itemsep=2pt]
    \item \textbf{Chapter 1. Overview of the Internship Organization};
    \item \textbf{Chapter 2. Internship Activities and Assigned Tasks};
    \item \textbf{Chapter 3. Results, Experience, and Evaluation}.
\end{itemize}

These chapter titles are recommendations only and may be adapted to the
actual internship content.

The report shall reflect work actually performed, observed, analyzed, or
learned during the internship. Unrelated theoretical material should not
be included merely to increase report length.

\vspace{0.3cm}

% ================================================================
% X. TABLE OF CONTENTS
% ================================================================

\noindent
\textbf{X.\quad TABLE OF CONTENTS}

\vspace{0.15cm}

The Table of Contents shall be generated automatically using:

\begin{verbatim}
\tableofcontents
\end{verbatim}

Students shall not create or manually edit page numbers in the Table of
Contents.

The document shall be compiled sufficiently to ensure that all page
numbers and entries are correctly updated.

\vspace{0.3cm}

% ================================================================
% XI. LIST OF SYMBOLS AND ABBREVIATIONS
% ================================================================

\noindent
\textbf{XI.\quad LIST OF SYMBOLS AND ABBREVIATIONS}

\vspace{0.15cm}

Symbols and abbreviations that are used repeatedly in the report shall be
listed in the \textbf{List of Symbols and Abbreviations}.

The list may contain the following columns:

\begin{center}
\begin{tabular}{llll}
    No. & Symbol / Abbreviation & Unit & Definition
\end{tabular}
\end{center}

Only symbols and abbreviations actually used in the final report shall be
retained. Example entries contained in the template shall be removed if
they are not relevant to the student's report.

Each abbreviation should normally be defined when it first appears in
the main text.

\vspace{0.3cm}

% ================================================================
% XII. FIGURES, DIAGRAMS, AND PHOTOGRAPHS
% ================================================================

\noindent
\textbf{XII.\quad FIGURES, DIAGRAMS, AND PHOTOGRAPHS}

\vspace{0.15cm}

Figures shall be clear, legible, and directly relevant to the report
content.

Figures shall be inserted using the standard \verb|figure| environment.

Example:

\begin{verbatim}
\begin{figure}[H]
    \centering
    \includegraphics[width=0.8\textwidth]
                    {figures/company_structure.pdf}
    \caption{Organizational structure of the company}
    \label{fig:company-structure}
\end{figure}
\end{verbatim}

Figure captions shall be placed below the figure.

Each referenced figure shall have a unique \verb|\label| and shall be
referenced in the text using, for example:

\begin{verbatim}
Figure~\ref{fig:company-structure}
\end{verbatim}

Students shall not manually type figure numbers.

Figures reproduced or adapted from external sources shall include an
appropriate citation.

Photographs taken during the internship shall be relevant and sufficiently
clear. Students shall not include confidential, restricted, or proprietary
information without permission from the internship organization.

Vector graphics or high-resolution images are recommended whenever
possible.

\vspace{0.3cm}

% ================================================================
% XIII. TABLES
% ================================================================

\noindent
\textbf{XIII.\quad TABLES}

\vspace{0.15cm}

Tables shall be created using suitable \LaTeX{} environments such as
\verb|table|, \verb|tabular|, \verb|tabularx|, or \verb|longtable|.

Table captions shall be placed above the table.

Example:

\begin{verbatim}
\begin{table}[H]
    \centering
    \caption{Internship schedule}
    \label{tab:internship-schedule}

    \begin{tabular}{lll}
        \toprule
        Week & Task & Result \\
        \midrule
        1 & Orientation & Completed \\
        2 & System study & Completed \\
        \bottomrule
    \end{tabular}
\end{table}
\end{verbatim}

Tables shall be referenced using:

\begin{verbatim}
Table~\ref{tab:internship-schedule}
\end{verbatim}

Students shall not manually type table numbers.

Tables should be created directly in \LaTeX{} rather than inserted as
screenshots whenever practical.

\vspace{0.3cm}

% ================================================================
% XIV. EQUATIONS AND TECHNICAL QUANTITIES
% ================================================================

\noindent
\textbf{XIV.\quad EQUATIONS AND TECHNICAL QUANTITIES}

\vspace{0.15cm}

Mathematical expressions shall be typeset directly using \LaTeX{}.
Images or screenshots of equations shall not be used.

Short expressions may be written in inline mathematical mode, for example:

\begin{verbatim}
$R = 0.7\,\Omega$
\end{verbatim}

Important equations shall use a numbered equation environment.

Example:

\begin{verbatim}
\begin{equation}
    P = UI
    \label{eq:power}
\end{equation}
\end{verbatim}

Equations shall be referenced using:

\begin{verbatim}
Equation~\eqref{eq:power}
\end{verbatim}

Equation numbers shall not be manually entered.

SI units shall be used wherever applicable. The \verb|siunitx| package
should be used where appropriate.

Examples:

\begin{verbatim}
\SI{220}{\volt}
\SI{50}{\hertz}
\SI{2.5}{\ampere}
\end{verbatim}

Notation and units shall be consistent throughout the report.

\vspace{0.3cm}

% ================================================================
% XV. SOURCE CODE
% ================================================================

\noindent
\textbf{XV.\quad SOURCE CODE}

\vspace{0.15cm}

If the internship involves programming, PLCs, microcontrollers,
simulation, or software development, relevant source code may be included
in the report.

Source code shall be presented using the \verb|listings| environment or
another environment defined by the official template.

Example:

\begin{verbatim}
\begin{lstlisting}[language=st]
IF Start = TRUE THEN
    Motor := TRUE;
ELSE
    Motor := FALSE;
END_IF;
\end{lstlisting}
\end{verbatim}

Only code necessary for explaining the work should appear in the main
report. Long source-code listings should be moved to an appendix or
submitted as supplementary files if required.

Source code should not normally be inserted as screenshots.

\vspace{0.3cm}

% ================================================================
% XVI. ABSTRACT
% ================================================================

\noindent
\textbf{XVI.\quad ABSTRACT}

\vspace{0.15cm}

The Abstract shall briefly summarize:

\begin{itemize}[leftmargin=0.8cm, itemsep=2pt]
    \item the internship organization;
    \item the objectives of the internship;
    \item the main tasks performed;
    \item the main results, knowledge, and skills obtained.
\end{itemize}

The Abstract should be concise and should not reproduce the full content
of the report chapter by chapter.

\vspace{0.3cm}

% ================================================================
% XVII. REFERENCES AND CITATIONS
% ================================================================

\noindent
\textbf{XVII.\quad REFERENCES AND CITATIONS}

\vspace{0.15cm}

All books, journal papers, conference papers, standards, websites,
datasheets, manuals, technical documents, figures, tables, or other
external sources used in preparing the report shall be properly cited.

References shall be managed through the bibliography system provided with
the official template.

Bibliographic entries shall normally be stored in:

\begin{verbatim}
References.bib
\end{verbatim}

References shall be cited in the report using:

\begin{verbatim}
\cite{reference-key}
\end{verbatim}

Students shall not manually type reference numbers such as ``[1]'',
``[2]'', etc.

Titles of foreign-language publications shall remain in their original
language in the bibliography.

All borrowed ideas, data, figures, tables, equations, methods, or results
shall be properly acknowledged.

\vspace{0.3cm}

% ================================================================
% XVIII. CROSS-REFERENCING
% ================================================================

\noindent
\textbf{XVIII.\quad CROSS-REFERENCING}

\vspace{0.15cm}

Numbered elements shall be cross-referenced using the standard \LaTeX{}
mechanism.

Use the following commands where appropriate:

\begin{verbatim}
\label{...}
\ref{...}
\eqref{...}
\end{verbatim}

For example:

\begin{verbatim}
As shown in Figure~\ref{fig:system}, ...
\end{verbatim}

Students shall not manually type numbers such as ``Figure 2.3'' when
those numbers can be generated automatically.

\vspace{0.3cm}

% ================================================================
% XIX. PAGE NUMBERING
% ================================================================

\noindent
\textbf{XIX.\quad PAGE NUMBERING}

\vspace{0.15cm}

Cover pages and official evaluation forms at the beginning of the report
shall not display page numbers.

The main body of the report shall use Arabic page numbers beginning from
Chapter 1.

Page numbering shall be generated automatically by the official template.
Students shall not manually enter page numbers.

\vspace{0.3cm}

% ================================================================
% XX. PROJECT FILE ORGANIZATION
% ================================================================

\noindent
\textbf{XX.\quad PROJECT FILE ORGANIZATION}

\vspace{0.15cm}

Students should maintain a clear and organized \LaTeX{} project
structure. A recommended structure is:

\begin{verbatim}
InternshipReport/
|
|-- main.tex
|-- References.bib
|
|-- figures/
|   |-- company.pdf
|   |-- system.png
|   `-- ...
|
`-- chapters/
    |-- chapter1.tex
    |-- chapter2.tex
    `-- chapter3.tex
\end{verbatim}

Figures should normally be stored in the \verb|figures| directory
specified by the official template.

File and directory names should preferably contain Latin letters,
numbers, hyphens, or underscores.

Recommended examples include:

\begin{verbatim}
company_structure.pdf
motor_control.png
chapter_01.tex
\end{verbatim}

Spaces and unnecessary special characters in file names should be
avoided where possible.

\vspace{0.3cm}

% ================================================================
% XXI. DOCUMENT COMPILATION AND FINAL CHECK
% ================================================================

\noindent
\textbf{XXI.\quad DOCUMENT COMPILATION AND FINAL CHECK}

\vspace{0.15cm}

Before submission, the complete \LaTeX{} project shall compile
successfully and produce the final PDF report.

The final document shall not contain unresolved references such as:

\begin{verbatim}
??
[?]
\end{verbatim}

Before generating the final PDF, students shall check:

\begin{itemize}[leftmargin=0.8cm, itemsep=2pt]
    \item the Table of Contents;
    \item page numbering;
    \item chapter, section, and subsection numbering;
    \item figure numbering;
    \item table numbering;
    \item equation numbering;
    \item cross-references;
    \item citations;
    \item bibliography;
    \item font display;
    \item mathematical symbols;
    \item image quality;
    \item completeness of all required forms.
\end{itemize}

The document may be compiled using TeX Live, MiKTeX, or another
compatible \LaTeX{} environment. Regardless of the compilation
environment, the resulting PDF shall conform to the official Faculty
template.

\vspace{0.3cm}

% ================================================================
% XXII. CONTENT THAT STUDENTS MAY EDIT
% ================================================================

\noindent
\textbf{XXII.\quad CONTENT THAT STUDENTS MAY EDIT}

\vspace{0.15cm}

Students may edit or replace:

\begin{itemize}[leftmargin=0.8cm, itemsep=2pt]
    \item student name;
    \item class;
    \item internship organization;
    \item supervisor name;
    \item Abstract;
    \item List of Symbols and Abbreviations;
    \item chapter titles and chapter content;
    \item figures and photographs;
    \item tables;
    \item equations;
    \item citations and references.
\end{itemize}

Students shall not modify the main document formatting settings in the
preamble, including the document class, page geometry, line spacing,
heading formatting, table-of-contents formatting, font configuration,
page layout, or cover-page macros unless specifically instructed.

\vspace{0.3cm}

% ================================================================
% XXIII. FINAL SUBMISSION REQUIREMENTS
% ================================================================

\noindent
\textbf{XXIII.\quad FINAL SUBMISSION REQUIREMENTS}

\vspace{0.15cm}

Before submission, the final report shall:

\begin{itemize}[leftmargin=0.8cm, itemsep=2pt]
    \item compile without errors that affect the report content or presentation;
    \item contain all required forms and report sections;
    \item contain correct student, class, internship organization, and supervisor information;
    \item contain an updated Table of Contents, List of Figures, and List of Tables;
    \item contain no unused example content from the template;
    \item contain no placeholder text such as \textit{Title of Chapter 1},
    \textit{Section 1.1}, or other sample entries;
    \item contain all required citations and references;
    \item contain no unresolved cross-references.
\end{itemize}

If the Faculty requires submission of the \LaTeX{} source, the complete
project shall include all \verb|.tex| files, the bibliography database,
figures, and any other files required to reproduce the submitted PDF.

\vspace{0.4cm}

\noindent
\textbf{Important:}
Before submitting the final internship report, students shall remove or
comment out the following command in the main document:

\begin{verbatim}
\input{Regulations.tex}
\end{verbatim}

For example:

\begin{verbatim}
% \input{Regulations.tex}
\end{verbatim}

This removes the guidance pages from the final submitted report while
preserving the regulations file in the \LaTeX{} project for reference.

} % End 13pt content

\clearpage

\end{verbatim}

For example:

\begin{verbatim}
% 
% =================================================================
% REGULATIONS FOR PREPARING THE GRADUATION INTERNSHIP REPORT
% USING LATEX
% =================================================================
%
% This file is intended to be included from the main document with:
%
%     \input{Regulations.tex}
%
% IMPORTANT:
% This section is provided for guidance only. Students shall remove
% or comment out the \input{Regulations.tex} command before submitting
% the final internship report.
%
% =================================================================

\clearpage
\thispagestyle{empty}

\begin{center}
    {\fontsize{16pt}{20pt}\selectfont\bfseries
    SPECIFIC REGULATIONS FOR PREPARING\\
    GRADUATION INTERNSHIP REPORTS USING \LaTeX}
\end{center}

\vspace{0.5cm}

{\fontsize{13pt}{17pt}\selectfont
\setstretch{1.3}
\justifying

\noindent
\textbf{Note:}
This section is provided for guidance only and is included in the main
document using the command \verb|\input{Regulations.tex}|.
Students shall remove or comment out this command before submitting the
final internship report.

\vspace{0.4cm}

% ================================================================
% I. GENERAL REQUIREMENTS
% ================================================================

\noindent
\textbf{I.\quad GENERAL REQUIREMENTS}

\vspace{0.15cm}

The Graduation Internship Report shall be prepared using the official
\LaTeX{} template provided by the Faculty. Students shall preserve the
predefined page layout, typography, line spacing, heading styles,
numbering system, figure and table formatting, and bibliography settings
unless otherwise instructed by the Faculty or the supervisor.

The report shall be written in English.

Students are responsible for the academic and technical content of the
report. The official \LaTeX{} template is responsible for standardized
document formatting.

\vspace{0.3cm}

% ================================================================
% II. REPORT STRUCTURE
% ================================================================

\noindent
\textbf{II.\quad REPORT STRUCTURE}

\vspace{0.15cm}

The report shall normally be arranged in the following order:

\begin{enumerate}[label=\arabic*., leftmargin=1.2cm, itemsep=2pt]
    \item Outer Cover;
    \item Comments of the Internship Host Organization;
    \item Supervisor's Evaluation and Graduation Internship Defense Results;
    \item Table of Contents;
    \item List of Symbols and Abbreviations, if applicable;
    \item List of Figures;
    \item List of Tables;
    \item Abstract;
    \item Main Chapters of the Internship Report;
    \item References;
    \item Appendices, if applicable.
\end{enumerate}

The preliminary pages shall not be numbered as chapters. The main report
content shall be organized into consecutively numbered chapters beginning
with Chapter 1.

\vspace{0.3cm}

% ================================================================
% III. PAGE SIZE AND MARGINS
% ================================================================

\noindent
\textbf{III.\quad PAGE SIZE AND MARGINS}

\vspace{0.15cm}

The report shall be formatted on A4 paper with the following margins:

\begin{center}
\begin{tabular}{ll}
    Left margin:   & 3.5 cm \\
    Right margin:  & 2.0 cm \\
    Top margin:    & 2.5 cm \\
    Bottom margin: & 2.5 cm
\end{tabular}
\end{center}

These settings are predefined in the official template and shall not be
modified by students.

\vspace{0.3cm}

% ================================================================
% IV. FONT, FONT SIZE, AND LINE SPACING
% ================================================================

\noindent
\textbf{IV.\quad FONT, FONT SIZE, AND LINE SPACING}

\vspace{0.15cm}

The main text shall use a 13 pt font size with 1.3-line spacing.
A Times-family typeface shall be used in accordance with the official
template.

The text shall be fully justified.

Students shall not use repeated spaces, repeated manual line breaks, or
multiple blank lines to position text. Paragraph and heading spacing shall
be controlled by the template.

\vspace{0.3cm}

% ================================================================
% V. OUTER COVER
% ================================================================

\noindent
\textbf{V.\quad OUTER COVER}

\vspace{0.15cm}

The outer cover shall follow the official format provided in the template
and shall contain the following information:

\begin{itemize}[leftmargin=0.8cm, itemsep=2pt]
    \item Thai Nguyen University of Technology;
    \item Faculty of International Training;
    \item Engineering and Technology Department;
    \item University logo and Faculty logo;
    \item the title \textbf{INTERNSHIP REPORT};
    \item Internship Organization;
    \item Author;
    \item Class;
    \item Supervisor;
    \item place and year of completion.
\end{itemize}

Students shall only replace the required information fields and shall not
modify the cover-page layout unless instructed otherwise.

\vspace{0.3cm}

% ================================================================
% VI. COMMENTS OF THE INTERNSHIP HOST ORGANIZATION
% ================================================================

\noindent
\textbf{VI.\quad COMMENTS OF THE INTERNSHIP HOST ORGANIZATION}

\vspace{0.15cm}

The page entitled \textbf{COMMENTS OF THE INTERNSHIP HOST ORGANIZATION}
is reserved for the internship host organization.

It includes:

\begin{itemize}[leftmargin=0.8cm, itemsep=2pt]
    \item the student's full name;
    \item class;
    \item comments from the host organization;
    \item score;
    \item date;
    \item signature and official seal.
\end{itemize}

Students shall not alter or complete official evaluation fields unless
specifically authorized.

\vspace{0.3cm}

% ================================================================
% VII. SUPERVISOR'S EVALUATION AND DEFENSE RESULTS
% ================================================================

\noindent
\textbf{VII.\quad SUPERVISOR'S EVALUATION AND DEFENSE RESULTS}

\vspace{0.15cm}

The section entitled
\textbf{SUPERVISOR'S EVALUATION AND GRADUATION INTERNSHIP DEFENSE RESULTS}
is reserved for the supervisor and defense committee.

The evaluation may include:

\begin{itemize}[leftmargin=0.8cm, itemsep=2pt]
    \item Internship Progress and Student's Attitude;
    \item Report Content;
    \item Report Presentation;
    \item Other Comments;
    \item Supervisor's Opinion;
    \item Defense Result;
    \item Score.
\end{itemize}

Students shall not enter evaluation scores or defense results before they
are officially determined.

\vspace{0.3cm}

% ================================================================
% VIII. CHAPTERS, SECTIONS, AND SUBSECTIONS
% ================================================================

\noindent
\textbf{VIII.\quad CHAPTERS, SECTIONS, AND SUBSECTIONS}

\vspace{0.15cm}

The report structure shall be generated using the standard \LaTeX{}
commands:

\begin{verbatim}
\chapter{...}
\section{...}
\subsection{...}
\end{verbatim}

Chapter, section, and subsection numbers shall be generated automatically
by \LaTeX{}.

Students shall not manually type chapter or section numbers and shall not
reproduce heading styles using commands such as \verb|\textbf|,
\verb|\fontsize|, or \verb|\centering| instead of the appropriate
structural command.

Excessive subdivision should be avoided. In most cases, the hierarchy
Chapter--Section--Subsection is sufficient.

\vspace{0.3cm}

% ================================================================
% IX. ORGANIZATION OF THE MAIN CONTENT
% ================================================================

\noindent
\textbf{IX.\quad ORGANIZATION OF THE MAIN CONTENT}

\vspace{0.15cm}

The specific content of the report shall depend on the internship
organization and the assigned tasks. A typical structure may include:

\begin{itemize}[leftmargin=0.8cm, itemsep=2pt]
    \item \textbf{Chapter 1. Overview of the Internship Organization};
    \item \textbf{Chapter 2. Internship Activities and Assigned Tasks};
    \item \textbf{Chapter 3. Results, Experience, and Evaluation}.
\end{itemize}

These chapter titles are recommendations only and may be adapted to the
actual internship content.

The report shall reflect work actually performed, observed, analyzed, or
learned during the internship. Unrelated theoretical material should not
be included merely to increase report length.

\vspace{0.3cm}

% ================================================================
% X. TABLE OF CONTENTS
% ================================================================

\noindent
\textbf{X.\quad TABLE OF CONTENTS}

\vspace{0.15cm}

The Table of Contents shall be generated automatically using:

\begin{verbatim}
\tableofcontents
\end{verbatim}

Students shall not create or manually edit page numbers in the Table of
Contents.

The document shall be compiled sufficiently to ensure that all page
numbers and entries are correctly updated.

\vspace{0.3cm}

% ================================================================
% XI. LIST OF SYMBOLS AND ABBREVIATIONS
% ================================================================

\noindent
\textbf{XI.\quad LIST OF SYMBOLS AND ABBREVIATIONS}

\vspace{0.15cm}

Symbols and abbreviations that are used repeatedly in the report shall be
listed in the \textbf{List of Symbols and Abbreviations}.

The list may contain the following columns:

\begin{center}
\begin{tabular}{llll}
    No. & Symbol / Abbreviation & Unit & Definition
\end{tabular}
\end{center}

Only symbols and abbreviations actually used in the final report shall be
retained. Example entries contained in the template shall be removed if
they are not relevant to the student's report.

Each abbreviation should normally be defined when it first appears in
the main text.

\vspace{0.3cm}

% ================================================================
% XII. FIGURES, DIAGRAMS, AND PHOTOGRAPHS
% ================================================================

\noindent
\textbf{XII.\quad FIGURES, DIAGRAMS, AND PHOTOGRAPHS}

\vspace{0.15cm}

Figures shall be clear, legible, and directly relevant to the report
content.

Figures shall be inserted using the standard \verb|figure| environment.

Example:

\begin{verbatim}
\begin{figure}[H]
    \centering
    \includegraphics[width=0.8\textwidth]
                    {figures/company_structure.pdf}
    \caption{Organizational structure of the company}
    \label{fig:company-structure}
\end{figure}
\end{verbatim}

Figure captions shall be placed below the figure.

Each referenced figure shall have a unique \verb|\label| and shall be
referenced in the text using, for example:

\begin{verbatim}
Figure~\ref{fig:company-structure}
\end{verbatim}

Students shall not manually type figure numbers.

Figures reproduced or adapted from external sources shall include an
appropriate citation.

Photographs taken during the internship shall be relevant and sufficiently
clear. Students shall not include confidential, restricted, or proprietary
information without permission from the internship organization.

Vector graphics or high-resolution images are recommended whenever
possible.

\vspace{0.3cm}

% ================================================================
% XIII. TABLES
% ================================================================

\noindent
\textbf{XIII.\quad TABLES}

\vspace{0.15cm}

Tables shall be created using suitable \LaTeX{} environments such as
\verb|table|, \verb|tabular|, \verb|tabularx|, or \verb|longtable|.

Table captions shall be placed above the table.

Example:

\begin{verbatim}
\begin{table}[H]
    \centering
    \caption{Internship schedule}
    \label{tab:internship-schedule}

    \begin{tabular}{lll}
        \toprule
        Week & Task & Result \\
        \midrule
        1 & Orientation & Completed \\
        2 & System study & Completed \\
        \bottomrule
    \end{tabular}
\end{table}
\end{verbatim}

Tables shall be referenced using:

\begin{verbatim}
Table~\ref{tab:internship-schedule}
\end{verbatim}

Students shall not manually type table numbers.

Tables should be created directly in \LaTeX{} rather than inserted as
screenshots whenever practical.

\vspace{0.3cm}

% ================================================================
% XIV. EQUATIONS AND TECHNICAL QUANTITIES
% ================================================================

\noindent
\textbf{XIV.\quad EQUATIONS AND TECHNICAL QUANTITIES}

\vspace{0.15cm}

Mathematical expressions shall be typeset directly using \LaTeX{}.
Images or screenshots of equations shall not be used.

Short expressions may be written in inline mathematical mode, for example:

\begin{verbatim}
$R = 0.7\,\Omega$
\end{verbatim}

Important equations shall use a numbered equation environment.

Example:

\begin{verbatim}
\begin{equation}
    P = UI
    \label{eq:power}
\end{equation}
\end{verbatim}

Equations shall be referenced using:

\begin{verbatim}
Equation~\eqref{eq:power}
\end{verbatim}

Equation numbers shall not be manually entered.

SI units shall be used wherever applicable. The \verb|siunitx| package
should be used where appropriate.

Examples:

\begin{verbatim}
\SI{220}{\volt}
\SI{50}{\hertz}
\SI{2.5}{\ampere}
\end{verbatim}

Notation and units shall be consistent throughout the report.

\vspace{0.3cm}

% ================================================================
% XV. SOURCE CODE
% ================================================================

\noindent
\textbf{XV.\quad SOURCE CODE}

\vspace{0.15cm}

If the internship involves programming, PLCs, microcontrollers,
simulation, or software development, relevant source code may be included
in the report.

Source code shall be presented using the \verb|listings| environment or
another environment defined by the official template.

Example:

\begin{verbatim}
\begin{lstlisting}[language=st]
IF Start = TRUE THEN
    Motor := TRUE;
ELSE
    Motor := FALSE;
END_IF;
\end{lstlisting}
\end{verbatim}

Only code necessary for explaining the work should appear in the main
report. Long source-code listings should be moved to an appendix or
submitted as supplementary files if required.

Source code should not normally be inserted as screenshots.

\vspace{0.3cm}

% ================================================================
% XVI. ABSTRACT
% ================================================================

\noindent
\textbf{XVI.\quad ABSTRACT}

\vspace{0.15cm}

The Abstract shall briefly summarize:

\begin{itemize}[leftmargin=0.8cm, itemsep=2pt]
    \item the internship organization;
    \item the objectives of the internship;
    \item the main tasks performed;
    \item the main results, knowledge, and skills obtained.
\end{itemize}

The Abstract should be concise and should not reproduce the full content
of the report chapter by chapter.

\vspace{0.3cm}

% ================================================================
% XVII. REFERENCES AND CITATIONS
% ================================================================

\noindent
\textbf{XVII.\quad REFERENCES AND CITATIONS}

\vspace{0.15cm}

All books, journal papers, conference papers, standards, websites,
datasheets, manuals, technical documents, figures, tables, or other
external sources used in preparing the report shall be properly cited.

References shall be managed through the bibliography system provided with
the official template.

Bibliographic entries shall normally be stored in:

\begin{verbatim}
References.bib
\end{verbatim}

References shall be cited in the report using:

\begin{verbatim}
\cite{reference-key}
\end{verbatim}

Students shall not manually type reference numbers such as ``[1]'',
``[2]'', etc.

Titles of foreign-language publications shall remain in their original
language in the bibliography.

All borrowed ideas, data, figures, tables, equations, methods, or results
shall be properly acknowledged.

\vspace{0.3cm}

% ================================================================
% XVIII. CROSS-REFERENCING
% ================================================================

\noindent
\textbf{XVIII.\quad CROSS-REFERENCING}

\vspace{0.15cm}

Numbered elements shall be cross-referenced using the standard \LaTeX{}
mechanism.

Use the following commands where appropriate:

\begin{verbatim}
\label{...}
\ref{...}
\eqref{...}
\end{verbatim}

For example:

\begin{verbatim}
As shown in Figure~\ref{fig:system}, ...
\end{verbatim}

Students shall not manually type numbers such as ``Figure 2.3'' when
those numbers can be generated automatically.

\vspace{0.3cm}

% ================================================================
% XIX. PAGE NUMBERING
% ================================================================

\noindent
\textbf{XIX.\quad PAGE NUMBERING}

\vspace{0.15cm}

Cover pages and official evaluation forms at the beginning of the report
shall not display page numbers.

The main body of the report shall use Arabic page numbers beginning from
Chapter 1.

Page numbering shall be generated automatically by the official template.
Students shall not manually enter page numbers.

\vspace{0.3cm}

% ================================================================
% XX. PROJECT FILE ORGANIZATION
% ================================================================

\noindent
\textbf{XX.\quad PROJECT FILE ORGANIZATION}

\vspace{0.15cm}

Students should maintain a clear and organized \LaTeX{} project
structure. A recommended structure is:

\begin{verbatim}
InternshipReport/
|
|-- main.tex
|-- References.bib
|
|-- figures/
|   |-- company.pdf
|   |-- system.png
|   `-- ...
|
`-- chapters/
    |-- chapter1.tex
    |-- chapter2.tex
    `-- chapter3.tex
\end{verbatim}

Figures should normally be stored in the \verb|figures| directory
specified by the official template.

File and directory names should preferably contain Latin letters,
numbers, hyphens, or underscores.

Recommended examples include:

\begin{verbatim}
company_structure.pdf
motor_control.png
chapter_01.tex
\end{verbatim}

Spaces and unnecessary special characters in file names should be
avoided where possible.

\vspace{0.3cm}

% ================================================================
% XXI. DOCUMENT COMPILATION AND FINAL CHECK
% ================================================================

\noindent
\textbf{XXI.\quad DOCUMENT COMPILATION AND FINAL CHECK}

\vspace{0.15cm}

Before submission, the complete \LaTeX{} project shall compile
successfully and produce the final PDF report.

The final document shall not contain unresolved references such as:

\begin{verbatim}
??
[?]
\end{verbatim}

Before generating the final PDF, students shall check:

\begin{itemize}[leftmargin=0.8cm, itemsep=2pt]
    \item the Table of Contents;
    \item page numbering;
    \item chapter, section, and subsection numbering;
    \item figure numbering;
    \item table numbering;
    \item equation numbering;
    \item cross-references;
    \item citations;
    \item bibliography;
    \item font display;
    \item mathematical symbols;
    \item image quality;
    \item completeness of all required forms.
\end{itemize}

The document may be compiled using TeX Live, MiKTeX, or another
compatible \LaTeX{} environment. Regardless of the compilation
environment, the resulting PDF shall conform to the official Faculty
template.

\vspace{0.3cm}

% ================================================================
% XXII. CONTENT THAT STUDENTS MAY EDIT
% ================================================================

\noindent
\textbf{XXII.\quad CONTENT THAT STUDENTS MAY EDIT}

\vspace{0.15cm}

Students may edit or replace:

\begin{itemize}[leftmargin=0.8cm, itemsep=2pt]
    \item student name;
    \item class;
    \item internship organization;
    \item supervisor name;
    \item Abstract;
    \item List of Symbols and Abbreviations;
    \item chapter titles and chapter content;
    \item figures and photographs;
    \item tables;
    \item equations;
    \item citations and references.
\end{itemize}

Students shall not modify the main document formatting settings in the
preamble, including the document class, page geometry, line spacing,
heading formatting, table-of-contents formatting, font configuration,
page layout, or cover-page macros unless specifically instructed.

\vspace{0.3cm}

% ================================================================
% XXIII. FINAL SUBMISSION REQUIREMENTS
% ================================================================

\noindent
\textbf{XXIII.\quad FINAL SUBMISSION REQUIREMENTS}

\vspace{0.15cm}

Before submission, the final report shall:

\begin{itemize}[leftmargin=0.8cm, itemsep=2pt]
    \item compile without errors that affect the report content or presentation;
    \item contain all required forms and report sections;
    \item contain correct student, class, internship organization, and supervisor information;
    \item contain an updated Table of Contents, List of Figures, and List of Tables;
    \item contain no unused example content from the template;
    \item contain no placeholder text such as \textit{Title of Chapter 1},
    \textit{Section 1.1}, or other sample entries;
    \item contain all required citations and references;
    \item contain no unresolved cross-references.
\end{itemize}

If the Faculty requires submission of the \LaTeX{} source, the complete
project shall include all \verb|.tex| files, the bibliography database,
figures, and any other files required to reproduce the submitted PDF.

\vspace{0.4cm}

\noindent
\textbf{Important:}
Before submitting the final internship report, students shall remove or
comment out the following command in the main document:

\begin{verbatim}
\input{Regulations.tex}
\end{verbatim}

For example:

\begin{verbatim}
% \input{Regulations.tex}
\end{verbatim}

This removes the guidance pages from the final submitted report while
preserving the regulations file in the \LaTeX{} project for reference.

} % End 13pt content

\clearpage

\end{verbatim}

This removes the guidance pages from the final submitted report while
preserving the regulations file in the \LaTeX{} project for reference.

} % End 13pt content

\clearpage

\end{verbatim}

This removes the guidance pages from the final submitted report while
preserving the regulations file in the \LaTeX{} project for reference.

} % End 13pt content

\clearpage

\end{verbatim}

For example:

\begin{verbatim}
% 
% =================================================================
% REGULATIONS FOR PREPARING THE GRADUATION INTERNSHIP REPORT
% USING LATEX
% =================================================================
%
% This file is intended to be included from the main document with:
%
%     
% =================================================================
% REGULATIONS FOR PREPARING THE GRADUATION INTERNSHIP REPORT
% USING LATEX
% =================================================================
%
% This file is intended to be included from the main document with:
%
%     
% =================================================================
% REGULATIONS FOR PREPARING THE GRADUATION INTERNSHIP REPORT
% USING LATEX
% =================================================================
%
% This file is intended to be included from the main document with:
%
%     \input{Regulations.tex}
%
% IMPORTANT:
% This section is provided for guidance only. Students shall remove
% or comment out the \input{Regulations.tex} command before submitting
% the final internship report.
%
% =================================================================

\clearpage
\thispagestyle{empty}

\begin{center}
    {\fontsize{16pt}{20pt}\selectfont\bfseries
    SPECIFIC REGULATIONS FOR PREPARING\\
    GRADUATION INTERNSHIP REPORTS USING \LaTeX}
\end{center}

\vspace{0.5cm}

{\fontsize{13pt}{17pt}\selectfont
\setstretch{1.3}
\justifying

\noindent
\textbf{Note:}
This section is provided for guidance only and is included in the main
document using the command \verb|\input{Regulations.tex}|.
Students shall remove or comment out this command before submitting the
final internship report.

\vspace{0.4cm}

% ================================================================
% I. GENERAL REQUIREMENTS
% ================================================================

\noindent
\textbf{I.\quad GENERAL REQUIREMENTS}

\vspace{0.15cm}

The Graduation Internship Report shall be prepared using the official
\LaTeX{} template provided by the Faculty. Students shall preserve the
predefined page layout, typography, line spacing, heading styles,
numbering system, figure and table formatting, and bibliography settings
unless otherwise instructed by the Faculty or the supervisor.

The report shall be written in English.

Students are responsible for the academic and technical content of the
report. The official \LaTeX{} template is responsible for standardized
document formatting.

\vspace{0.3cm}

% ================================================================
% II. REPORT STRUCTURE
% ================================================================

\noindent
\textbf{II.\quad REPORT STRUCTURE}

\vspace{0.15cm}

The report shall normally be arranged in the following order:

\begin{enumerate}[label=\arabic*., leftmargin=1.2cm, itemsep=2pt]
    \item Outer Cover;
    \item Comments of the Internship Host Organization;
    \item Supervisor's Evaluation and Graduation Internship Defense Results;
    \item Table of Contents;
    \item List of Symbols and Abbreviations, if applicable;
    \item List of Figures;
    \item List of Tables;
    \item Abstract;
    \item Main Chapters of the Internship Report;
    \item References;
    \item Appendices, if applicable.
\end{enumerate}

The preliminary pages shall not be numbered as chapters. The main report
content shall be organized into consecutively numbered chapters beginning
with Chapter 1.

\vspace{0.3cm}

% ================================================================
% III. PAGE SIZE AND MARGINS
% ================================================================

\noindent
\textbf{III.\quad PAGE SIZE AND MARGINS}

\vspace{0.15cm}

The report shall be formatted on A4 paper with the following margins:

\begin{center}
\begin{tabular}{ll}
    Left margin:   & 3.5 cm \\
    Right margin:  & 2.0 cm \\
    Top margin:    & 2.5 cm \\
    Bottom margin: & 2.5 cm
\end{tabular}
\end{center}

These settings are predefined in the official template and shall not be
modified by students.

\vspace{0.3cm}

% ================================================================
% IV. FONT, FONT SIZE, AND LINE SPACING
% ================================================================

\noindent
\textbf{IV.\quad FONT, FONT SIZE, AND LINE SPACING}

\vspace{0.15cm}

The main text shall use a 13 pt font size with 1.3-line spacing.
A Times-family typeface shall be used in accordance with the official
template.

The text shall be fully justified.

Students shall not use repeated spaces, repeated manual line breaks, or
multiple blank lines to position text. Paragraph and heading spacing shall
be controlled by the template.

\vspace{0.3cm}

% ================================================================
% V. OUTER COVER
% ================================================================

\noindent
\textbf{V.\quad OUTER COVER}

\vspace{0.15cm}

The outer cover shall follow the official format provided in the template
and shall contain the following information:

\begin{itemize}[leftmargin=0.8cm, itemsep=2pt]
    \item Thai Nguyen University of Technology;
    \item Faculty of International Training;
    \item Engineering and Technology Department;
    \item University logo and Faculty logo;
    \item the title \textbf{INTERNSHIP REPORT};
    \item Internship Organization;
    \item Author;
    \item Class;
    \item Supervisor;
    \item place and year of completion.
\end{itemize}

Students shall only replace the required information fields and shall not
modify the cover-page layout unless instructed otherwise.

\vspace{0.3cm}

% ================================================================
% VI. COMMENTS OF THE INTERNSHIP HOST ORGANIZATION
% ================================================================

\noindent
\textbf{VI.\quad COMMENTS OF THE INTERNSHIP HOST ORGANIZATION}

\vspace{0.15cm}

The page entitled \textbf{COMMENTS OF THE INTERNSHIP HOST ORGANIZATION}
is reserved for the internship host organization.

It includes:

\begin{itemize}[leftmargin=0.8cm, itemsep=2pt]
    \item the student's full name;
    \item class;
    \item comments from the host organization;
    \item score;
    \item date;
    \item signature and official seal.
\end{itemize}

Students shall not alter or complete official evaluation fields unless
specifically authorized.

\vspace{0.3cm}

% ================================================================
% VII. SUPERVISOR'S EVALUATION AND DEFENSE RESULTS
% ================================================================

\noindent
\textbf{VII.\quad SUPERVISOR'S EVALUATION AND DEFENSE RESULTS}

\vspace{0.15cm}

The section entitled
\textbf{SUPERVISOR'S EVALUATION AND GRADUATION INTERNSHIP DEFENSE RESULTS}
is reserved for the supervisor and defense committee.

The evaluation may include:

\begin{itemize}[leftmargin=0.8cm, itemsep=2pt]
    \item Internship Progress and Student's Attitude;
    \item Report Content;
    \item Report Presentation;
    \item Other Comments;
    \item Supervisor's Opinion;
    \item Defense Result;
    \item Score.
\end{itemize}

Students shall not enter evaluation scores or defense results before they
are officially determined.

\vspace{0.3cm}

% ================================================================
% VIII. CHAPTERS, SECTIONS, AND SUBSECTIONS
% ================================================================

\noindent
\textbf{VIII.\quad CHAPTERS, SECTIONS, AND SUBSECTIONS}

\vspace{0.15cm}

The report structure shall be generated using the standard \LaTeX{}
commands:

\begin{verbatim}
\chapter{...}
\section{...}
\subsection{...}
\end{verbatim}

Chapter, section, and subsection numbers shall be generated automatically
by \LaTeX{}.

Students shall not manually type chapter or section numbers and shall not
reproduce heading styles using commands such as \verb|\textbf|,
\verb|\fontsize|, or \verb|\centering| instead of the appropriate
structural command.

Excessive subdivision should be avoided. In most cases, the hierarchy
Chapter--Section--Subsection is sufficient.

\vspace{0.3cm}

% ================================================================
% IX. ORGANIZATION OF THE MAIN CONTENT
% ================================================================

\noindent
\textbf{IX.\quad ORGANIZATION OF THE MAIN CONTENT}

\vspace{0.15cm}

The specific content of the report shall depend on the internship
organization and the assigned tasks. A typical structure may include:

\begin{itemize}[leftmargin=0.8cm, itemsep=2pt]
    \item \textbf{Chapter 1. Overview of the Internship Organization};
    \item \textbf{Chapter 2. Internship Activities and Assigned Tasks};
    \item \textbf{Chapter 3. Results, Experience, and Evaluation}.
\end{itemize}

These chapter titles are recommendations only and may be adapted to the
actual internship content.

The report shall reflect work actually performed, observed, analyzed, or
learned during the internship. Unrelated theoretical material should not
be included merely to increase report length.

\vspace{0.3cm}

% ================================================================
% X. TABLE OF CONTENTS
% ================================================================

\noindent
\textbf{X.\quad TABLE OF CONTENTS}

\vspace{0.15cm}

The Table of Contents shall be generated automatically using:

\begin{verbatim}
\tableofcontents
\end{verbatim}

Students shall not create or manually edit page numbers in the Table of
Contents.

The document shall be compiled sufficiently to ensure that all page
numbers and entries are correctly updated.

\vspace{0.3cm}

% ================================================================
% XI. LIST OF SYMBOLS AND ABBREVIATIONS
% ================================================================

\noindent
\textbf{XI.\quad LIST OF SYMBOLS AND ABBREVIATIONS}

\vspace{0.15cm}

Symbols and abbreviations that are used repeatedly in the report shall be
listed in the \textbf{List of Symbols and Abbreviations}.

The list may contain the following columns:

\begin{center}
\begin{tabular}{llll}
    No. & Symbol / Abbreviation & Unit & Definition
\end{tabular}
\end{center}

Only symbols and abbreviations actually used in the final report shall be
retained. Example entries contained in the template shall be removed if
they are not relevant to the student's report.

Each abbreviation should normally be defined when it first appears in
the main text.

\vspace{0.3cm}

% ================================================================
% XII. FIGURES, DIAGRAMS, AND PHOTOGRAPHS
% ================================================================

\noindent
\textbf{XII.\quad FIGURES, DIAGRAMS, AND PHOTOGRAPHS}

\vspace{0.15cm}

Figures shall be clear, legible, and directly relevant to the report
content.

Figures shall be inserted using the standard \verb|figure| environment.

Example:

\begin{verbatim}
\begin{figure}[H]
    \centering
    \includegraphics[width=0.8\textwidth]
                    {figures/company_structure.pdf}
    \caption{Organizational structure of the company}
    \label{fig:company-structure}
\end{figure}
\end{verbatim}

Figure captions shall be placed below the figure.

Each referenced figure shall have a unique \verb|\label| and shall be
referenced in the text using, for example:

\begin{verbatim}
Figure~\ref{fig:company-structure}
\end{verbatim}

Students shall not manually type figure numbers.

Figures reproduced or adapted from external sources shall include an
appropriate citation.

Photographs taken during the internship shall be relevant and sufficiently
clear. Students shall not include confidential, restricted, or proprietary
information without permission from the internship organization.

Vector graphics or high-resolution images are recommended whenever
possible.

\vspace{0.3cm}

% ================================================================
% XIII. TABLES
% ================================================================

\noindent
\textbf{XIII.\quad TABLES}

\vspace{0.15cm}

Tables shall be created using suitable \LaTeX{} environments such as
\verb|table|, \verb|tabular|, \verb|tabularx|, or \verb|longtable|.

Table captions shall be placed above the table.

Example:

\begin{verbatim}
\begin{table}[H]
    \centering
    \caption{Internship schedule}
    \label{tab:internship-schedule}

    \begin{tabular}{lll}
        \toprule
        Week & Task & Result \\
        \midrule
        1 & Orientation & Completed \\
        2 & System study & Completed \\
        \bottomrule
    \end{tabular}
\end{table}
\end{verbatim}

Tables shall be referenced using:

\begin{verbatim}
Table~\ref{tab:internship-schedule}
\end{verbatim}

Students shall not manually type table numbers.

Tables should be created directly in \LaTeX{} rather than inserted as
screenshots whenever practical.

\vspace{0.3cm}

% ================================================================
% XIV. EQUATIONS AND TECHNICAL QUANTITIES
% ================================================================

\noindent
\textbf{XIV.\quad EQUATIONS AND TECHNICAL QUANTITIES}

\vspace{0.15cm}

Mathematical expressions shall be typeset directly using \LaTeX{}.
Images or screenshots of equations shall not be used.

Short expressions may be written in inline mathematical mode, for example:

\begin{verbatim}
$R = 0.7\,\Omega$
\end{verbatim}

Important equations shall use a numbered equation environment.

Example:

\begin{verbatim}
\begin{equation}
    P = UI
    \label{eq:power}
\end{equation}
\end{verbatim}

Equations shall be referenced using:

\begin{verbatim}
Equation~\eqref{eq:power}
\end{verbatim}

Equation numbers shall not be manually entered.

SI units shall be used wherever applicable. The \verb|siunitx| package
should be used where appropriate.

Examples:

\begin{verbatim}
\SI{220}{\volt}
\SI{50}{\hertz}
\SI{2.5}{\ampere}
\end{verbatim}

Notation and units shall be consistent throughout the report.

\vspace{0.3cm}

% ================================================================
% XV. SOURCE CODE
% ================================================================

\noindent
\textbf{XV.\quad SOURCE CODE}

\vspace{0.15cm}

If the internship involves programming, PLCs, microcontrollers,
simulation, or software development, relevant source code may be included
in the report.

Source code shall be presented using the \verb|listings| environment or
another environment defined by the official template.

Example:

\begin{verbatim}
\begin{lstlisting}[language=st]
IF Start = TRUE THEN
    Motor := TRUE;
ELSE
    Motor := FALSE;
END_IF;
\end{lstlisting}
\end{verbatim}

Only code necessary for explaining the work should appear in the main
report. Long source-code listings should be moved to an appendix or
submitted as supplementary files if required.

Source code should not normally be inserted as screenshots.

\vspace{0.3cm}

% ================================================================
% XVI. ABSTRACT
% ================================================================

\noindent
\textbf{XVI.\quad ABSTRACT}

\vspace{0.15cm}

The Abstract shall briefly summarize:

\begin{itemize}[leftmargin=0.8cm, itemsep=2pt]
    \item the internship organization;
    \item the objectives of the internship;
    \item the main tasks performed;
    \item the main results, knowledge, and skills obtained.
\end{itemize}

The Abstract should be concise and should not reproduce the full content
of the report chapter by chapter.

\vspace{0.3cm}

% ================================================================
% XVII. REFERENCES AND CITATIONS
% ================================================================

\noindent
\textbf{XVII.\quad REFERENCES AND CITATIONS}

\vspace{0.15cm}

All books, journal papers, conference papers, standards, websites,
datasheets, manuals, technical documents, figures, tables, or other
external sources used in preparing the report shall be properly cited.

References shall be managed through the bibliography system provided with
the official template.

Bibliographic entries shall normally be stored in:

\begin{verbatim}
References.bib
\end{verbatim}

References shall be cited in the report using:

\begin{verbatim}
\cite{reference-key}
\end{verbatim}

Students shall not manually type reference numbers such as ``[1]'',
``[2]'', etc.

Titles of foreign-language publications shall remain in their original
language in the bibliography.

All borrowed ideas, data, figures, tables, equations, methods, or results
shall be properly acknowledged.

\vspace{0.3cm}

% ================================================================
% XVIII. CROSS-REFERENCING
% ================================================================

\noindent
\textbf{XVIII.\quad CROSS-REFERENCING}

\vspace{0.15cm}

Numbered elements shall be cross-referenced using the standard \LaTeX{}
mechanism.

Use the following commands where appropriate:

\begin{verbatim}
\label{...}
\ref{...}
\eqref{...}
\end{verbatim}

For example:

\begin{verbatim}
As shown in Figure~\ref{fig:system}, ...
\end{verbatim}

Students shall not manually type numbers such as ``Figure 2.3'' when
those numbers can be generated automatically.

\vspace{0.3cm}

% ================================================================
% XIX. PAGE NUMBERING
% ================================================================

\noindent
\textbf{XIX.\quad PAGE NUMBERING}

\vspace{0.15cm}

Cover pages and official evaluation forms at the beginning of the report
shall not display page numbers.

The main body of the report shall use Arabic page numbers beginning from
Chapter 1.

Page numbering shall be generated automatically by the official template.
Students shall not manually enter page numbers.

\vspace{0.3cm}

% ================================================================
% XX. PROJECT FILE ORGANIZATION
% ================================================================

\noindent
\textbf{XX.\quad PROJECT FILE ORGANIZATION}

\vspace{0.15cm}

Students should maintain a clear and organized \LaTeX{} project
structure. A recommended structure is:

\begin{verbatim}
InternshipReport/
|
|-- main.tex
|-- References.bib
|
|-- figures/
|   |-- company.pdf
|   |-- system.png
|   `-- ...
|
`-- chapters/
    |-- chapter1.tex
    |-- chapter2.tex
    `-- chapter3.tex
\end{verbatim}

Figures should normally be stored in the \verb|figures| directory
specified by the official template.

File and directory names should preferably contain Latin letters,
numbers, hyphens, or underscores.

Recommended examples include:

\begin{verbatim}
company_structure.pdf
motor_control.png
chapter_01.tex
\end{verbatim}

Spaces and unnecessary special characters in file names should be
avoided where possible.

\vspace{0.3cm}

% ================================================================
% XXI. DOCUMENT COMPILATION AND FINAL CHECK
% ================================================================

\noindent
\textbf{XXI.\quad DOCUMENT COMPILATION AND FINAL CHECK}

\vspace{0.15cm}

Before submission, the complete \LaTeX{} project shall compile
successfully and produce the final PDF report.

The final document shall not contain unresolved references such as:

\begin{verbatim}
??
[?]
\end{verbatim}

Before generating the final PDF, students shall check:

\begin{itemize}[leftmargin=0.8cm, itemsep=2pt]
    \item the Table of Contents;
    \item page numbering;
    \item chapter, section, and subsection numbering;
    \item figure numbering;
    \item table numbering;
    \item equation numbering;
    \item cross-references;
    \item citations;
    \item bibliography;
    \item font display;
    \item mathematical symbols;
    \item image quality;
    \item completeness of all required forms.
\end{itemize}

The document may be compiled using TeX Live, MiKTeX, or another
compatible \LaTeX{} environment. Regardless of the compilation
environment, the resulting PDF shall conform to the official Faculty
template.

\vspace{0.3cm}

% ================================================================
% XXII. CONTENT THAT STUDENTS MAY EDIT
% ================================================================

\noindent
\textbf{XXII.\quad CONTENT THAT STUDENTS MAY EDIT}

\vspace{0.15cm}

Students may edit or replace:

\begin{itemize}[leftmargin=0.8cm, itemsep=2pt]
    \item student name;
    \item class;
    \item internship organization;
    \item supervisor name;
    \item Abstract;
    \item List of Symbols and Abbreviations;
    \item chapter titles and chapter content;
    \item figures and photographs;
    \item tables;
    \item equations;
    \item citations and references.
\end{itemize}

Students shall not modify the main document formatting settings in the
preamble, including the document class, page geometry, line spacing,
heading formatting, table-of-contents formatting, font configuration,
page layout, or cover-page macros unless specifically instructed.

\vspace{0.3cm}

% ================================================================
% XXIII. FINAL SUBMISSION REQUIREMENTS
% ================================================================

\noindent
\textbf{XXIII.\quad FINAL SUBMISSION REQUIREMENTS}

\vspace{0.15cm}

Before submission, the final report shall:

\begin{itemize}[leftmargin=0.8cm, itemsep=2pt]
    \item compile without errors that affect the report content or presentation;
    \item contain all required forms and report sections;
    \item contain correct student, class, internship organization, and supervisor information;
    \item contain an updated Table of Contents, List of Figures, and List of Tables;
    \item contain no unused example content from the template;
    \item contain no placeholder text such as \textit{Title of Chapter 1},
    \textit{Section 1.1}, or other sample entries;
    \item contain all required citations and references;
    \item contain no unresolved cross-references.
\end{itemize}

If the Faculty requires submission of the \LaTeX{} source, the complete
project shall include all \verb|.tex| files, the bibliography database,
figures, and any other files required to reproduce the submitted PDF.

\vspace{0.4cm}

\noindent
\textbf{Important:}
Before submitting the final internship report, students shall remove or
comment out the following command in the main document:

\begin{verbatim}
\input{Regulations.tex}
\end{verbatim}

For example:

\begin{verbatim}
% \input{Regulations.tex}
\end{verbatim}

This removes the guidance pages from the final submitted report while
preserving the regulations file in the \LaTeX{} project for reference.

} % End 13pt content

\clearpage

%
% IMPORTANT:
% This section is provided for guidance only. Students shall remove
% or comment out the 
% =================================================================
% REGULATIONS FOR PREPARING THE GRADUATION INTERNSHIP REPORT
% USING LATEX
% =================================================================
%
% This file is intended to be included from the main document with:
%
%     \input{Regulations.tex}
%
% IMPORTANT:
% This section is provided for guidance only. Students shall remove
% or comment out the \input{Regulations.tex} command before submitting
% the final internship report.
%
% =================================================================

\clearpage
\thispagestyle{empty}

\begin{center}
    {\fontsize{16pt}{20pt}\selectfont\bfseries
    SPECIFIC REGULATIONS FOR PREPARING\\
    GRADUATION INTERNSHIP REPORTS USING \LaTeX}
\end{center}

\vspace{0.5cm}

{\fontsize{13pt}{17pt}\selectfont
\setstretch{1.3}
\justifying

\noindent
\textbf{Note:}
This section is provided for guidance only and is included in the main
document using the command \verb|\input{Regulations.tex}|.
Students shall remove or comment out this command before submitting the
final internship report.

\vspace{0.4cm}

% ================================================================
% I. GENERAL REQUIREMENTS
% ================================================================

\noindent
\textbf{I.\quad GENERAL REQUIREMENTS}

\vspace{0.15cm}

The Graduation Internship Report shall be prepared using the official
\LaTeX{} template provided by the Faculty. Students shall preserve the
predefined page layout, typography, line spacing, heading styles,
numbering system, figure and table formatting, and bibliography settings
unless otherwise instructed by the Faculty or the supervisor.

The report shall be written in English.

Students are responsible for the academic and technical content of the
report. The official \LaTeX{} template is responsible for standardized
document formatting.

\vspace{0.3cm}

% ================================================================
% II. REPORT STRUCTURE
% ================================================================

\noindent
\textbf{II.\quad REPORT STRUCTURE}

\vspace{0.15cm}

The report shall normally be arranged in the following order:

\begin{enumerate}[label=\arabic*., leftmargin=1.2cm, itemsep=2pt]
    \item Outer Cover;
    \item Comments of the Internship Host Organization;
    \item Supervisor's Evaluation and Graduation Internship Defense Results;
    \item Table of Contents;
    \item List of Symbols and Abbreviations, if applicable;
    \item List of Figures;
    \item List of Tables;
    \item Abstract;
    \item Main Chapters of the Internship Report;
    \item References;
    \item Appendices, if applicable.
\end{enumerate}

The preliminary pages shall not be numbered as chapters. The main report
content shall be organized into consecutively numbered chapters beginning
with Chapter 1.

\vspace{0.3cm}

% ================================================================
% III. PAGE SIZE AND MARGINS
% ================================================================

\noindent
\textbf{III.\quad PAGE SIZE AND MARGINS}

\vspace{0.15cm}

The report shall be formatted on A4 paper with the following margins:

\begin{center}
\begin{tabular}{ll}
    Left margin:   & 3.5 cm \\
    Right margin:  & 2.0 cm \\
    Top margin:    & 2.5 cm \\
    Bottom margin: & 2.5 cm
\end{tabular}
\end{center}

These settings are predefined in the official template and shall not be
modified by students.

\vspace{0.3cm}

% ================================================================
% IV. FONT, FONT SIZE, AND LINE SPACING
% ================================================================

\noindent
\textbf{IV.\quad FONT, FONT SIZE, AND LINE SPACING}

\vspace{0.15cm}

The main text shall use a 13 pt font size with 1.3-line spacing.
A Times-family typeface shall be used in accordance with the official
template.

The text shall be fully justified.

Students shall not use repeated spaces, repeated manual line breaks, or
multiple blank lines to position text. Paragraph and heading spacing shall
be controlled by the template.

\vspace{0.3cm}

% ================================================================
% V. OUTER COVER
% ================================================================

\noindent
\textbf{V.\quad OUTER COVER}

\vspace{0.15cm}

The outer cover shall follow the official format provided in the template
and shall contain the following information:

\begin{itemize}[leftmargin=0.8cm, itemsep=2pt]
    \item Thai Nguyen University of Technology;
    \item Faculty of International Training;
    \item Engineering and Technology Department;
    \item University logo and Faculty logo;
    \item the title \textbf{INTERNSHIP REPORT};
    \item Internship Organization;
    \item Author;
    \item Class;
    \item Supervisor;
    \item place and year of completion.
\end{itemize}

Students shall only replace the required information fields and shall not
modify the cover-page layout unless instructed otherwise.

\vspace{0.3cm}

% ================================================================
% VI. COMMENTS OF THE INTERNSHIP HOST ORGANIZATION
% ================================================================

\noindent
\textbf{VI.\quad COMMENTS OF THE INTERNSHIP HOST ORGANIZATION}

\vspace{0.15cm}

The page entitled \textbf{COMMENTS OF THE INTERNSHIP HOST ORGANIZATION}
is reserved for the internship host organization.

It includes:

\begin{itemize}[leftmargin=0.8cm, itemsep=2pt]
    \item the student's full name;
    \item class;
    \item comments from the host organization;
    \item score;
    \item date;
    \item signature and official seal.
\end{itemize}

Students shall not alter or complete official evaluation fields unless
specifically authorized.

\vspace{0.3cm}

% ================================================================
% VII. SUPERVISOR'S EVALUATION AND DEFENSE RESULTS
% ================================================================

\noindent
\textbf{VII.\quad SUPERVISOR'S EVALUATION AND DEFENSE RESULTS}

\vspace{0.15cm}

The section entitled
\textbf{SUPERVISOR'S EVALUATION AND GRADUATION INTERNSHIP DEFENSE RESULTS}
is reserved for the supervisor and defense committee.

The evaluation may include:

\begin{itemize}[leftmargin=0.8cm, itemsep=2pt]
    \item Internship Progress and Student's Attitude;
    \item Report Content;
    \item Report Presentation;
    \item Other Comments;
    \item Supervisor's Opinion;
    \item Defense Result;
    \item Score.
\end{itemize}

Students shall not enter evaluation scores or defense results before they
are officially determined.

\vspace{0.3cm}

% ================================================================
% VIII. CHAPTERS, SECTIONS, AND SUBSECTIONS
% ================================================================

\noindent
\textbf{VIII.\quad CHAPTERS, SECTIONS, AND SUBSECTIONS}

\vspace{0.15cm}

The report structure shall be generated using the standard \LaTeX{}
commands:

\begin{verbatim}
\chapter{...}
\section{...}
\subsection{...}
\end{verbatim}

Chapter, section, and subsection numbers shall be generated automatically
by \LaTeX{}.

Students shall not manually type chapter or section numbers and shall not
reproduce heading styles using commands such as \verb|\textbf|,
\verb|\fontsize|, or \verb|\centering| instead of the appropriate
structural command.

Excessive subdivision should be avoided. In most cases, the hierarchy
Chapter--Section--Subsection is sufficient.

\vspace{0.3cm}

% ================================================================
% IX. ORGANIZATION OF THE MAIN CONTENT
% ================================================================

\noindent
\textbf{IX.\quad ORGANIZATION OF THE MAIN CONTENT}

\vspace{0.15cm}

The specific content of the report shall depend on the internship
organization and the assigned tasks. A typical structure may include:

\begin{itemize}[leftmargin=0.8cm, itemsep=2pt]
    \item \textbf{Chapter 1. Overview of the Internship Organization};
    \item \textbf{Chapter 2. Internship Activities and Assigned Tasks};
    \item \textbf{Chapter 3. Results, Experience, and Evaluation}.
\end{itemize}

These chapter titles are recommendations only and may be adapted to the
actual internship content.

The report shall reflect work actually performed, observed, analyzed, or
learned during the internship. Unrelated theoretical material should not
be included merely to increase report length.

\vspace{0.3cm}

% ================================================================
% X. TABLE OF CONTENTS
% ================================================================

\noindent
\textbf{X.\quad TABLE OF CONTENTS}

\vspace{0.15cm}

The Table of Contents shall be generated automatically using:

\begin{verbatim}
\tableofcontents
\end{verbatim}

Students shall not create or manually edit page numbers in the Table of
Contents.

The document shall be compiled sufficiently to ensure that all page
numbers and entries are correctly updated.

\vspace{0.3cm}

% ================================================================
% XI. LIST OF SYMBOLS AND ABBREVIATIONS
% ================================================================

\noindent
\textbf{XI.\quad LIST OF SYMBOLS AND ABBREVIATIONS}

\vspace{0.15cm}

Symbols and abbreviations that are used repeatedly in the report shall be
listed in the \textbf{List of Symbols and Abbreviations}.

The list may contain the following columns:

\begin{center}
\begin{tabular}{llll}
    No. & Symbol / Abbreviation & Unit & Definition
\end{tabular}
\end{center}

Only symbols and abbreviations actually used in the final report shall be
retained. Example entries contained in the template shall be removed if
they are not relevant to the student's report.

Each abbreviation should normally be defined when it first appears in
the main text.

\vspace{0.3cm}

% ================================================================
% XII. FIGURES, DIAGRAMS, AND PHOTOGRAPHS
% ================================================================

\noindent
\textbf{XII.\quad FIGURES, DIAGRAMS, AND PHOTOGRAPHS}

\vspace{0.15cm}

Figures shall be clear, legible, and directly relevant to the report
content.

Figures shall be inserted using the standard \verb|figure| environment.

Example:

\begin{verbatim}
\begin{figure}[H]
    \centering
    \includegraphics[width=0.8\textwidth]
                    {figures/company_structure.pdf}
    \caption{Organizational structure of the company}
    \label{fig:company-structure}
\end{figure}
\end{verbatim}

Figure captions shall be placed below the figure.

Each referenced figure shall have a unique \verb|\label| and shall be
referenced in the text using, for example:

\begin{verbatim}
Figure~\ref{fig:company-structure}
\end{verbatim}

Students shall not manually type figure numbers.

Figures reproduced or adapted from external sources shall include an
appropriate citation.

Photographs taken during the internship shall be relevant and sufficiently
clear. Students shall not include confidential, restricted, or proprietary
information without permission from the internship organization.

Vector graphics or high-resolution images are recommended whenever
possible.

\vspace{0.3cm}

% ================================================================
% XIII. TABLES
% ================================================================

\noindent
\textbf{XIII.\quad TABLES}

\vspace{0.15cm}

Tables shall be created using suitable \LaTeX{} environments such as
\verb|table|, \verb|tabular|, \verb|tabularx|, or \verb|longtable|.

Table captions shall be placed above the table.

Example:

\begin{verbatim}
\begin{table}[H]
    \centering
    \caption{Internship schedule}
    \label{tab:internship-schedule}

    \begin{tabular}{lll}
        \toprule
        Week & Task & Result \\
        \midrule
        1 & Orientation & Completed \\
        2 & System study & Completed \\
        \bottomrule
    \end{tabular}
\end{table}
\end{verbatim}

Tables shall be referenced using:

\begin{verbatim}
Table~\ref{tab:internship-schedule}
\end{verbatim}

Students shall not manually type table numbers.

Tables should be created directly in \LaTeX{} rather than inserted as
screenshots whenever practical.

\vspace{0.3cm}

% ================================================================
% XIV. EQUATIONS AND TECHNICAL QUANTITIES
% ================================================================

\noindent
\textbf{XIV.\quad EQUATIONS AND TECHNICAL QUANTITIES}

\vspace{0.15cm}

Mathematical expressions shall be typeset directly using \LaTeX{}.
Images or screenshots of equations shall not be used.

Short expressions may be written in inline mathematical mode, for example:

\begin{verbatim}
$R = 0.7\,\Omega$
\end{verbatim}

Important equations shall use a numbered equation environment.

Example:

\begin{verbatim}
\begin{equation}
    P = UI
    \label{eq:power}
\end{equation}
\end{verbatim}

Equations shall be referenced using:

\begin{verbatim}
Equation~\eqref{eq:power}
\end{verbatim}

Equation numbers shall not be manually entered.

SI units shall be used wherever applicable. The \verb|siunitx| package
should be used where appropriate.

Examples:

\begin{verbatim}
\SI{220}{\volt}
\SI{50}{\hertz}
\SI{2.5}{\ampere}
\end{verbatim}

Notation and units shall be consistent throughout the report.

\vspace{0.3cm}

% ================================================================
% XV. SOURCE CODE
% ================================================================

\noindent
\textbf{XV.\quad SOURCE CODE}

\vspace{0.15cm}

If the internship involves programming, PLCs, microcontrollers,
simulation, or software development, relevant source code may be included
in the report.

Source code shall be presented using the \verb|listings| environment or
another environment defined by the official template.

Example:

\begin{verbatim}
\begin{lstlisting}[language=st]
IF Start = TRUE THEN
    Motor := TRUE;
ELSE
    Motor := FALSE;
END_IF;
\end{lstlisting}
\end{verbatim}

Only code necessary for explaining the work should appear in the main
report. Long source-code listings should be moved to an appendix or
submitted as supplementary files if required.

Source code should not normally be inserted as screenshots.

\vspace{0.3cm}

% ================================================================
% XVI. ABSTRACT
% ================================================================

\noindent
\textbf{XVI.\quad ABSTRACT}

\vspace{0.15cm}

The Abstract shall briefly summarize:

\begin{itemize}[leftmargin=0.8cm, itemsep=2pt]
    \item the internship organization;
    \item the objectives of the internship;
    \item the main tasks performed;
    \item the main results, knowledge, and skills obtained.
\end{itemize}

The Abstract should be concise and should not reproduce the full content
of the report chapter by chapter.

\vspace{0.3cm}

% ================================================================
% XVII. REFERENCES AND CITATIONS
% ================================================================

\noindent
\textbf{XVII.\quad REFERENCES AND CITATIONS}

\vspace{0.15cm}

All books, journal papers, conference papers, standards, websites,
datasheets, manuals, technical documents, figures, tables, or other
external sources used in preparing the report shall be properly cited.

References shall be managed through the bibliography system provided with
the official template.

Bibliographic entries shall normally be stored in:

\begin{verbatim}
References.bib
\end{verbatim}

References shall be cited in the report using:

\begin{verbatim}
\cite{reference-key}
\end{verbatim}

Students shall not manually type reference numbers such as ``[1]'',
``[2]'', etc.

Titles of foreign-language publications shall remain in their original
language in the bibliography.

All borrowed ideas, data, figures, tables, equations, methods, or results
shall be properly acknowledged.

\vspace{0.3cm}

% ================================================================
% XVIII. CROSS-REFERENCING
% ================================================================

\noindent
\textbf{XVIII.\quad CROSS-REFERENCING}

\vspace{0.15cm}

Numbered elements shall be cross-referenced using the standard \LaTeX{}
mechanism.

Use the following commands where appropriate:

\begin{verbatim}
\label{...}
\ref{...}
\eqref{...}
\end{verbatim}

For example:

\begin{verbatim}
As shown in Figure~\ref{fig:system}, ...
\end{verbatim}

Students shall not manually type numbers such as ``Figure 2.3'' when
those numbers can be generated automatically.

\vspace{0.3cm}

% ================================================================
% XIX. PAGE NUMBERING
% ================================================================

\noindent
\textbf{XIX.\quad PAGE NUMBERING}

\vspace{0.15cm}

Cover pages and official evaluation forms at the beginning of the report
shall not display page numbers.

The main body of the report shall use Arabic page numbers beginning from
Chapter 1.

Page numbering shall be generated automatically by the official template.
Students shall not manually enter page numbers.

\vspace{0.3cm}

% ================================================================
% XX. PROJECT FILE ORGANIZATION
% ================================================================

\noindent
\textbf{XX.\quad PROJECT FILE ORGANIZATION}

\vspace{0.15cm}

Students should maintain a clear and organized \LaTeX{} project
structure. A recommended structure is:

\begin{verbatim}
InternshipReport/
|
|-- main.tex
|-- References.bib
|
|-- figures/
|   |-- company.pdf
|   |-- system.png
|   `-- ...
|
`-- chapters/
    |-- chapter1.tex
    |-- chapter2.tex
    `-- chapter3.tex
\end{verbatim}

Figures should normally be stored in the \verb|figures| directory
specified by the official template.

File and directory names should preferably contain Latin letters,
numbers, hyphens, or underscores.

Recommended examples include:

\begin{verbatim}
company_structure.pdf
motor_control.png
chapter_01.tex
\end{verbatim}

Spaces and unnecessary special characters in file names should be
avoided where possible.

\vspace{0.3cm}

% ================================================================
% XXI. DOCUMENT COMPILATION AND FINAL CHECK
% ================================================================

\noindent
\textbf{XXI.\quad DOCUMENT COMPILATION AND FINAL CHECK}

\vspace{0.15cm}

Before submission, the complete \LaTeX{} project shall compile
successfully and produce the final PDF report.

The final document shall not contain unresolved references such as:

\begin{verbatim}
??
[?]
\end{verbatim}

Before generating the final PDF, students shall check:

\begin{itemize}[leftmargin=0.8cm, itemsep=2pt]
    \item the Table of Contents;
    \item page numbering;
    \item chapter, section, and subsection numbering;
    \item figure numbering;
    \item table numbering;
    \item equation numbering;
    \item cross-references;
    \item citations;
    \item bibliography;
    \item font display;
    \item mathematical symbols;
    \item image quality;
    \item completeness of all required forms.
\end{itemize}

The document may be compiled using TeX Live, MiKTeX, or another
compatible \LaTeX{} environment. Regardless of the compilation
environment, the resulting PDF shall conform to the official Faculty
template.

\vspace{0.3cm}

% ================================================================
% XXII. CONTENT THAT STUDENTS MAY EDIT
% ================================================================

\noindent
\textbf{XXII.\quad CONTENT THAT STUDENTS MAY EDIT}

\vspace{0.15cm}

Students may edit or replace:

\begin{itemize}[leftmargin=0.8cm, itemsep=2pt]
    \item student name;
    \item class;
    \item internship organization;
    \item supervisor name;
    \item Abstract;
    \item List of Symbols and Abbreviations;
    \item chapter titles and chapter content;
    \item figures and photographs;
    \item tables;
    \item equations;
    \item citations and references.
\end{itemize}

Students shall not modify the main document formatting settings in the
preamble, including the document class, page geometry, line spacing,
heading formatting, table-of-contents formatting, font configuration,
page layout, or cover-page macros unless specifically instructed.

\vspace{0.3cm}

% ================================================================
% XXIII. FINAL SUBMISSION REQUIREMENTS
% ================================================================

\noindent
\textbf{XXIII.\quad FINAL SUBMISSION REQUIREMENTS}

\vspace{0.15cm}

Before submission, the final report shall:

\begin{itemize}[leftmargin=0.8cm, itemsep=2pt]
    \item compile without errors that affect the report content or presentation;
    \item contain all required forms and report sections;
    \item contain correct student, class, internship organization, and supervisor information;
    \item contain an updated Table of Contents, List of Figures, and List of Tables;
    \item contain no unused example content from the template;
    \item contain no placeholder text such as \textit{Title of Chapter 1},
    \textit{Section 1.1}, or other sample entries;
    \item contain all required citations and references;
    \item contain no unresolved cross-references.
\end{itemize}

If the Faculty requires submission of the \LaTeX{} source, the complete
project shall include all \verb|.tex| files, the bibliography database,
figures, and any other files required to reproduce the submitted PDF.

\vspace{0.4cm}

\noindent
\textbf{Important:}
Before submitting the final internship report, students shall remove or
comment out the following command in the main document:

\begin{verbatim}
\input{Regulations.tex}
\end{verbatim}

For example:

\begin{verbatim}
% \input{Regulations.tex}
\end{verbatim}

This removes the guidance pages from the final submitted report while
preserving the regulations file in the \LaTeX{} project for reference.

} % End 13pt content

\clearpage
 command before submitting
% the final internship report.
%
% =================================================================

\clearpage
\thispagestyle{empty}

\begin{center}
    {\fontsize{16pt}{20pt}\selectfont\bfseries
    SPECIFIC REGULATIONS FOR PREPARING\\
    GRADUATION INTERNSHIP REPORTS USING \LaTeX}
\end{center}

\vspace{0.5cm}

{\fontsize{13pt}{17pt}\selectfont
\setstretch{1.3}
\justifying

\noindent
\textbf{Note:}
This section is provided for guidance only and is included in the main
document using the command \verb|
% =================================================================
% REGULATIONS FOR PREPARING THE GRADUATION INTERNSHIP REPORT
% USING LATEX
% =================================================================
%
% This file is intended to be included from the main document with:
%
%     \input{Regulations.tex}
%
% IMPORTANT:
% This section is provided for guidance only. Students shall remove
% or comment out the \input{Regulations.tex} command before submitting
% the final internship report.
%
% =================================================================

\clearpage
\thispagestyle{empty}

\begin{center}
    {\fontsize{16pt}{20pt}\selectfont\bfseries
    SPECIFIC REGULATIONS FOR PREPARING\\
    GRADUATION INTERNSHIP REPORTS USING \LaTeX}
\end{center}

\vspace{0.5cm}

{\fontsize{13pt}{17pt}\selectfont
\setstretch{1.3}
\justifying

\noindent
\textbf{Note:}
This section is provided for guidance only and is included in the main
document using the command \verb|\input{Regulations.tex}|.
Students shall remove or comment out this command before submitting the
final internship report.

\vspace{0.4cm}

% ================================================================
% I. GENERAL REQUIREMENTS
% ================================================================

\noindent
\textbf{I.\quad GENERAL REQUIREMENTS}

\vspace{0.15cm}

The Graduation Internship Report shall be prepared using the official
\LaTeX{} template provided by the Faculty. Students shall preserve the
predefined page layout, typography, line spacing, heading styles,
numbering system, figure and table formatting, and bibliography settings
unless otherwise instructed by the Faculty or the supervisor.

The report shall be written in English.

Students are responsible for the academic and technical content of the
report. The official \LaTeX{} template is responsible for standardized
document formatting.

\vspace{0.3cm}

% ================================================================
% II. REPORT STRUCTURE
% ================================================================

\noindent
\textbf{II.\quad REPORT STRUCTURE}

\vspace{0.15cm}

The report shall normally be arranged in the following order:

\begin{enumerate}[label=\arabic*., leftmargin=1.2cm, itemsep=2pt]
    \item Outer Cover;
    \item Comments of the Internship Host Organization;
    \item Supervisor's Evaluation and Graduation Internship Defense Results;
    \item Table of Contents;
    \item List of Symbols and Abbreviations, if applicable;
    \item List of Figures;
    \item List of Tables;
    \item Abstract;
    \item Main Chapters of the Internship Report;
    \item References;
    \item Appendices, if applicable.
\end{enumerate}

The preliminary pages shall not be numbered as chapters. The main report
content shall be organized into consecutively numbered chapters beginning
with Chapter 1.

\vspace{0.3cm}

% ================================================================
% III. PAGE SIZE AND MARGINS
% ================================================================

\noindent
\textbf{III.\quad PAGE SIZE AND MARGINS}

\vspace{0.15cm}

The report shall be formatted on A4 paper with the following margins:

\begin{center}
\begin{tabular}{ll}
    Left margin:   & 3.5 cm \\
    Right margin:  & 2.0 cm \\
    Top margin:    & 2.5 cm \\
    Bottom margin: & 2.5 cm
\end{tabular}
\end{center}

These settings are predefined in the official template and shall not be
modified by students.

\vspace{0.3cm}

% ================================================================
% IV. FONT, FONT SIZE, AND LINE SPACING
% ================================================================

\noindent
\textbf{IV.\quad FONT, FONT SIZE, AND LINE SPACING}

\vspace{0.15cm}

The main text shall use a 13 pt font size with 1.3-line spacing.
A Times-family typeface shall be used in accordance with the official
template.

The text shall be fully justified.

Students shall not use repeated spaces, repeated manual line breaks, or
multiple blank lines to position text. Paragraph and heading spacing shall
be controlled by the template.

\vspace{0.3cm}

% ================================================================
% V. OUTER COVER
% ================================================================

\noindent
\textbf{V.\quad OUTER COVER}

\vspace{0.15cm}

The outer cover shall follow the official format provided in the template
and shall contain the following information:

\begin{itemize}[leftmargin=0.8cm, itemsep=2pt]
    \item Thai Nguyen University of Technology;
    \item Faculty of International Training;
    \item Engineering and Technology Department;
    \item University logo and Faculty logo;
    \item the title \textbf{INTERNSHIP REPORT};
    \item Internship Organization;
    \item Author;
    \item Class;
    \item Supervisor;
    \item place and year of completion.
\end{itemize}

Students shall only replace the required information fields and shall not
modify the cover-page layout unless instructed otherwise.

\vspace{0.3cm}

% ================================================================
% VI. COMMENTS OF THE INTERNSHIP HOST ORGANIZATION
% ================================================================

\noindent
\textbf{VI.\quad COMMENTS OF THE INTERNSHIP HOST ORGANIZATION}

\vspace{0.15cm}

The page entitled \textbf{COMMENTS OF THE INTERNSHIP HOST ORGANIZATION}
is reserved for the internship host organization.

It includes:

\begin{itemize}[leftmargin=0.8cm, itemsep=2pt]
    \item the student's full name;
    \item class;
    \item comments from the host organization;
    \item score;
    \item date;
    \item signature and official seal.
\end{itemize}

Students shall not alter or complete official evaluation fields unless
specifically authorized.

\vspace{0.3cm}

% ================================================================
% VII. SUPERVISOR'S EVALUATION AND DEFENSE RESULTS
% ================================================================

\noindent
\textbf{VII.\quad SUPERVISOR'S EVALUATION AND DEFENSE RESULTS}

\vspace{0.15cm}

The section entitled
\textbf{SUPERVISOR'S EVALUATION AND GRADUATION INTERNSHIP DEFENSE RESULTS}
is reserved for the supervisor and defense committee.

The evaluation may include:

\begin{itemize}[leftmargin=0.8cm, itemsep=2pt]
    \item Internship Progress and Student's Attitude;
    \item Report Content;
    \item Report Presentation;
    \item Other Comments;
    \item Supervisor's Opinion;
    \item Defense Result;
    \item Score.
\end{itemize}

Students shall not enter evaluation scores or defense results before they
are officially determined.

\vspace{0.3cm}

% ================================================================
% VIII. CHAPTERS, SECTIONS, AND SUBSECTIONS
% ================================================================

\noindent
\textbf{VIII.\quad CHAPTERS, SECTIONS, AND SUBSECTIONS}

\vspace{0.15cm}

The report structure shall be generated using the standard \LaTeX{}
commands:

\begin{verbatim}
\chapter{...}
\section{...}
\subsection{...}
\end{verbatim}

Chapter, section, and subsection numbers shall be generated automatically
by \LaTeX{}.

Students shall not manually type chapter or section numbers and shall not
reproduce heading styles using commands such as \verb|\textbf|,
\verb|\fontsize|, or \verb|\centering| instead of the appropriate
structural command.

Excessive subdivision should be avoided. In most cases, the hierarchy
Chapter--Section--Subsection is sufficient.

\vspace{0.3cm}

% ================================================================
% IX. ORGANIZATION OF THE MAIN CONTENT
% ================================================================

\noindent
\textbf{IX.\quad ORGANIZATION OF THE MAIN CONTENT}

\vspace{0.15cm}

The specific content of the report shall depend on the internship
organization and the assigned tasks. A typical structure may include:

\begin{itemize}[leftmargin=0.8cm, itemsep=2pt]
    \item \textbf{Chapter 1. Overview of the Internship Organization};
    \item \textbf{Chapter 2. Internship Activities and Assigned Tasks};
    \item \textbf{Chapter 3. Results, Experience, and Evaluation}.
\end{itemize}

These chapter titles are recommendations only and may be adapted to the
actual internship content.

The report shall reflect work actually performed, observed, analyzed, or
learned during the internship. Unrelated theoretical material should not
be included merely to increase report length.

\vspace{0.3cm}

% ================================================================
% X. TABLE OF CONTENTS
% ================================================================

\noindent
\textbf{X.\quad TABLE OF CONTENTS}

\vspace{0.15cm}

The Table of Contents shall be generated automatically using:

\begin{verbatim}
\tableofcontents
\end{verbatim}

Students shall not create or manually edit page numbers in the Table of
Contents.

The document shall be compiled sufficiently to ensure that all page
numbers and entries are correctly updated.

\vspace{0.3cm}

% ================================================================
% XI. LIST OF SYMBOLS AND ABBREVIATIONS
% ================================================================

\noindent
\textbf{XI.\quad LIST OF SYMBOLS AND ABBREVIATIONS}

\vspace{0.15cm}

Symbols and abbreviations that are used repeatedly in the report shall be
listed in the \textbf{List of Symbols and Abbreviations}.

The list may contain the following columns:

\begin{center}
\begin{tabular}{llll}
    No. & Symbol / Abbreviation & Unit & Definition
\end{tabular}
\end{center}

Only symbols and abbreviations actually used in the final report shall be
retained. Example entries contained in the template shall be removed if
they are not relevant to the student's report.

Each abbreviation should normally be defined when it first appears in
the main text.

\vspace{0.3cm}

% ================================================================
% XII. FIGURES, DIAGRAMS, AND PHOTOGRAPHS
% ================================================================

\noindent
\textbf{XII.\quad FIGURES, DIAGRAMS, AND PHOTOGRAPHS}

\vspace{0.15cm}

Figures shall be clear, legible, and directly relevant to the report
content.

Figures shall be inserted using the standard \verb|figure| environment.

Example:

\begin{verbatim}
\begin{figure}[H]
    \centering
    \includegraphics[width=0.8\textwidth]
                    {figures/company_structure.pdf}
    \caption{Organizational structure of the company}
    \label{fig:company-structure}
\end{figure}
\end{verbatim}

Figure captions shall be placed below the figure.

Each referenced figure shall have a unique \verb|\label| and shall be
referenced in the text using, for example:

\begin{verbatim}
Figure~\ref{fig:company-structure}
\end{verbatim}

Students shall not manually type figure numbers.

Figures reproduced or adapted from external sources shall include an
appropriate citation.

Photographs taken during the internship shall be relevant and sufficiently
clear. Students shall not include confidential, restricted, or proprietary
information without permission from the internship organization.

Vector graphics or high-resolution images are recommended whenever
possible.

\vspace{0.3cm}

% ================================================================
% XIII. TABLES
% ================================================================

\noindent
\textbf{XIII.\quad TABLES}

\vspace{0.15cm}

Tables shall be created using suitable \LaTeX{} environments such as
\verb|table|, \verb|tabular|, \verb|tabularx|, or \verb|longtable|.

Table captions shall be placed above the table.

Example:

\begin{verbatim}
\begin{table}[H]
    \centering
    \caption{Internship schedule}
    \label{tab:internship-schedule}

    \begin{tabular}{lll}
        \toprule
        Week & Task & Result \\
        \midrule
        1 & Orientation & Completed \\
        2 & System study & Completed \\
        \bottomrule
    \end{tabular}
\end{table}
\end{verbatim}

Tables shall be referenced using:

\begin{verbatim}
Table~\ref{tab:internship-schedule}
\end{verbatim}

Students shall not manually type table numbers.

Tables should be created directly in \LaTeX{} rather than inserted as
screenshots whenever practical.

\vspace{0.3cm}

% ================================================================
% XIV. EQUATIONS AND TECHNICAL QUANTITIES
% ================================================================

\noindent
\textbf{XIV.\quad EQUATIONS AND TECHNICAL QUANTITIES}

\vspace{0.15cm}

Mathematical expressions shall be typeset directly using \LaTeX{}.
Images or screenshots of equations shall not be used.

Short expressions may be written in inline mathematical mode, for example:

\begin{verbatim}
$R = 0.7\,\Omega$
\end{verbatim}

Important equations shall use a numbered equation environment.

Example:

\begin{verbatim}
\begin{equation}
    P = UI
    \label{eq:power}
\end{equation}
\end{verbatim}

Equations shall be referenced using:

\begin{verbatim}
Equation~\eqref{eq:power}
\end{verbatim}

Equation numbers shall not be manually entered.

SI units shall be used wherever applicable. The \verb|siunitx| package
should be used where appropriate.

Examples:

\begin{verbatim}
\SI{220}{\volt}
\SI{50}{\hertz}
\SI{2.5}{\ampere}
\end{verbatim}

Notation and units shall be consistent throughout the report.

\vspace{0.3cm}

% ================================================================
% XV. SOURCE CODE
% ================================================================

\noindent
\textbf{XV.\quad SOURCE CODE}

\vspace{0.15cm}

If the internship involves programming, PLCs, microcontrollers,
simulation, or software development, relevant source code may be included
in the report.

Source code shall be presented using the \verb|listings| environment or
another environment defined by the official template.

Example:

\begin{verbatim}
\begin{lstlisting}[language=st]
IF Start = TRUE THEN
    Motor := TRUE;
ELSE
    Motor := FALSE;
END_IF;
\end{lstlisting}
\end{verbatim}

Only code necessary for explaining the work should appear in the main
report. Long source-code listings should be moved to an appendix or
submitted as supplementary files if required.

Source code should not normally be inserted as screenshots.

\vspace{0.3cm}

% ================================================================
% XVI. ABSTRACT
% ================================================================

\noindent
\textbf{XVI.\quad ABSTRACT}

\vspace{0.15cm}

The Abstract shall briefly summarize:

\begin{itemize}[leftmargin=0.8cm, itemsep=2pt]
    \item the internship organization;
    \item the objectives of the internship;
    \item the main tasks performed;
    \item the main results, knowledge, and skills obtained.
\end{itemize}

The Abstract should be concise and should not reproduce the full content
of the report chapter by chapter.

\vspace{0.3cm}

% ================================================================
% XVII. REFERENCES AND CITATIONS
% ================================================================

\noindent
\textbf{XVII.\quad REFERENCES AND CITATIONS}

\vspace{0.15cm}

All books, journal papers, conference papers, standards, websites,
datasheets, manuals, technical documents, figures, tables, or other
external sources used in preparing the report shall be properly cited.

References shall be managed through the bibliography system provided with
the official template.

Bibliographic entries shall normally be stored in:

\begin{verbatim}
References.bib
\end{verbatim}

References shall be cited in the report using:

\begin{verbatim}
\cite{reference-key}
\end{verbatim}

Students shall not manually type reference numbers such as ``[1]'',
``[2]'', etc.

Titles of foreign-language publications shall remain in their original
language in the bibliography.

All borrowed ideas, data, figures, tables, equations, methods, or results
shall be properly acknowledged.

\vspace{0.3cm}

% ================================================================
% XVIII. CROSS-REFERENCING
% ================================================================

\noindent
\textbf{XVIII.\quad CROSS-REFERENCING}

\vspace{0.15cm}

Numbered elements shall be cross-referenced using the standard \LaTeX{}
mechanism.

Use the following commands where appropriate:

\begin{verbatim}
\label{...}
\ref{...}
\eqref{...}
\end{verbatim}

For example:

\begin{verbatim}
As shown in Figure~\ref{fig:system}, ...
\end{verbatim}

Students shall not manually type numbers such as ``Figure 2.3'' when
those numbers can be generated automatically.

\vspace{0.3cm}

% ================================================================
% XIX. PAGE NUMBERING
% ================================================================

\noindent
\textbf{XIX.\quad PAGE NUMBERING}

\vspace{0.15cm}

Cover pages and official evaluation forms at the beginning of the report
shall not display page numbers.

The main body of the report shall use Arabic page numbers beginning from
Chapter 1.

Page numbering shall be generated automatically by the official template.
Students shall not manually enter page numbers.

\vspace{0.3cm}

% ================================================================
% XX. PROJECT FILE ORGANIZATION
% ================================================================

\noindent
\textbf{XX.\quad PROJECT FILE ORGANIZATION}

\vspace{0.15cm}

Students should maintain a clear and organized \LaTeX{} project
structure. A recommended structure is:

\begin{verbatim}
InternshipReport/
|
|-- main.tex
|-- References.bib
|
|-- figures/
|   |-- company.pdf
|   |-- system.png
|   `-- ...
|
`-- chapters/
    |-- chapter1.tex
    |-- chapter2.tex
    `-- chapter3.tex
\end{verbatim}

Figures should normally be stored in the \verb|figures| directory
specified by the official template.

File and directory names should preferably contain Latin letters,
numbers, hyphens, or underscores.

Recommended examples include:

\begin{verbatim}
company_structure.pdf
motor_control.png
chapter_01.tex
\end{verbatim}

Spaces and unnecessary special characters in file names should be
avoided where possible.

\vspace{0.3cm}

% ================================================================
% XXI. DOCUMENT COMPILATION AND FINAL CHECK
% ================================================================

\noindent
\textbf{XXI.\quad DOCUMENT COMPILATION AND FINAL CHECK}

\vspace{0.15cm}

Before submission, the complete \LaTeX{} project shall compile
successfully and produce the final PDF report.

The final document shall not contain unresolved references such as:

\begin{verbatim}
??
[?]
\end{verbatim}

Before generating the final PDF, students shall check:

\begin{itemize}[leftmargin=0.8cm, itemsep=2pt]
    \item the Table of Contents;
    \item page numbering;
    \item chapter, section, and subsection numbering;
    \item figure numbering;
    \item table numbering;
    \item equation numbering;
    \item cross-references;
    \item citations;
    \item bibliography;
    \item font display;
    \item mathematical symbols;
    \item image quality;
    \item completeness of all required forms.
\end{itemize}

The document may be compiled using TeX Live, MiKTeX, or another
compatible \LaTeX{} environment. Regardless of the compilation
environment, the resulting PDF shall conform to the official Faculty
template.

\vspace{0.3cm}

% ================================================================
% XXII. CONTENT THAT STUDENTS MAY EDIT
% ================================================================

\noindent
\textbf{XXII.\quad CONTENT THAT STUDENTS MAY EDIT}

\vspace{0.15cm}

Students may edit or replace:

\begin{itemize}[leftmargin=0.8cm, itemsep=2pt]
    \item student name;
    \item class;
    \item internship organization;
    \item supervisor name;
    \item Abstract;
    \item List of Symbols and Abbreviations;
    \item chapter titles and chapter content;
    \item figures and photographs;
    \item tables;
    \item equations;
    \item citations and references.
\end{itemize}

Students shall not modify the main document formatting settings in the
preamble, including the document class, page geometry, line spacing,
heading formatting, table-of-contents formatting, font configuration,
page layout, or cover-page macros unless specifically instructed.

\vspace{0.3cm}

% ================================================================
% XXIII. FINAL SUBMISSION REQUIREMENTS
% ================================================================

\noindent
\textbf{XXIII.\quad FINAL SUBMISSION REQUIREMENTS}

\vspace{0.15cm}

Before submission, the final report shall:

\begin{itemize}[leftmargin=0.8cm, itemsep=2pt]
    \item compile without errors that affect the report content or presentation;
    \item contain all required forms and report sections;
    \item contain correct student, class, internship organization, and supervisor information;
    \item contain an updated Table of Contents, List of Figures, and List of Tables;
    \item contain no unused example content from the template;
    \item contain no placeholder text such as \textit{Title of Chapter 1},
    \textit{Section 1.1}, or other sample entries;
    \item contain all required citations and references;
    \item contain no unresolved cross-references.
\end{itemize}

If the Faculty requires submission of the \LaTeX{} source, the complete
project shall include all \verb|.tex| files, the bibliography database,
figures, and any other files required to reproduce the submitted PDF.

\vspace{0.4cm}

\noindent
\textbf{Important:}
Before submitting the final internship report, students shall remove or
comment out the following command in the main document:

\begin{verbatim}
\input{Regulations.tex}
\end{verbatim}

For example:

\begin{verbatim}
% \input{Regulations.tex}
\end{verbatim}

This removes the guidance pages from the final submitted report while
preserving the regulations file in the \LaTeX{} project for reference.

} % End 13pt content

\clearpage
|.
Students shall remove or comment out this command before submitting the
final internship report.

\vspace{0.4cm}

% ================================================================
% I. GENERAL REQUIREMENTS
% ================================================================

\noindent
\textbf{I.\quad GENERAL REQUIREMENTS}

\vspace{0.15cm}

The Graduation Internship Report shall be prepared using the official
\LaTeX{} template provided by the Faculty. Students shall preserve the
predefined page layout, typography, line spacing, heading styles,
numbering system, figure and table formatting, and bibliography settings
unless otherwise instructed by the Faculty or the supervisor.

The report shall be written in English.

Students are responsible for the academic and technical content of the
report. The official \LaTeX{} template is responsible for standardized
document formatting.

\vspace{0.3cm}

% ================================================================
% II. REPORT STRUCTURE
% ================================================================

\noindent
\textbf{II.\quad REPORT STRUCTURE}

\vspace{0.15cm}

The report shall normally be arranged in the following order:

\begin{enumerate}[label=\arabic*., leftmargin=1.2cm, itemsep=2pt]
    \item Outer Cover;
    \item Comments of the Internship Host Organization;
    \item Supervisor's Evaluation and Graduation Internship Defense Results;
    \item Table of Contents;
    \item List of Symbols and Abbreviations, if applicable;
    \item List of Figures;
    \item List of Tables;
    \item Abstract;
    \item Main Chapters of the Internship Report;
    \item References;
    \item Appendices, if applicable.
\end{enumerate}

The preliminary pages shall not be numbered as chapters. The main report
content shall be organized into consecutively numbered chapters beginning
with Chapter 1.

\vspace{0.3cm}

% ================================================================
% III. PAGE SIZE AND MARGINS
% ================================================================

\noindent
\textbf{III.\quad PAGE SIZE AND MARGINS}

\vspace{0.15cm}

The report shall be formatted on A4 paper with the following margins:

\begin{center}
\begin{tabular}{ll}
    Left margin:   & 3.5 cm \\
    Right margin:  & 2.0 cm \\
    Top margin:    & 2.5 cm \\
    Bottom margin: & 2.5 cm
\end{tabular}
\end{center}

These settings are predefined in the official template and shall not be
modified by students.

\vspace{0.3cm}

% ================================================================
% IV. FONT, FONT SIZE, AND LINE SPACING
% ================================================================

\noindent
\textbf{IV.\quad FONT, FONT SIZE, AND LINE SPACING}

\vspace{0.15cm}

The main text shall use a 13 pt font size with 1.3-line spacing.
A Times-family typeface shall be used in accordance with the official
template.

The text shall be fully justified.

Students shall not use repeated spaces, repeated manual line breaks, or
multiple blank lines to position text. Paragraph and heading spacing shall
be controlled by the template.

\vspace{0.3cm}

% ================================================================
% V. OUTER COVER
% ================================================================

\noindent
\textbf{V.\quad OUTER COVER}

\vspace{0.15cm}

The outer cover shall follow the official format provided in the template
and shall contain the following information:

\begin{itemize}[leftmargin=0.8cm, itemsep=2pt]
    \item Thai Nguyen University of Technology;
    \item Faculty of International Training;
    \item Engineering and Technology Department;
    \item University logo and Faculty logo;
    \item the title \textbf{INTERNSHIP REPORT};
    \item Internship Organization;
    \item Author;
    \item Class;
    \item Supervisor;
    \item place and year of completion.
\end{itemize}

Students shall only replace the required information fields and shall not
modify the cover-page layout unless instructed otherwise.

\vspace{0.3cm}

% ================================================================
% VI. COMMENTS OF THE INTERNSHIP HOST ORGANIZATION
% ================================================================

\noindent
\textbf{VI.\quad COMMENTS OF THE INTERNSHIP HOST ORGANIZATION}

\vspace{0.15cm}

The page entitled \textbf{COMMENTS OF THE INTERNSHIP HOST ORGANIZATION}
is reserved for the internship host organization.

It includes:

\begin{itemize}[leftmargin=0.8cm, itemsep=2pt]
    \item the student's full name;
    \item class;
    \item comments from the host organization;
    \item score;
    \item date;
    \item signature and official seal.
\end{itemize}

Students shall not alter or complete official evaluation fields unless
specifically authorized.

\vspace{0.3cm}

% ================================================================
% VII. SUPERVISOR'S EVALUATION AND DEFENSE RESULTS
% ================================================================

\noindent
\textbf{VII.\quad SUPERVISOR'S EVALUATION AND DEFENSE RESULTS}

\vspace{0.15cm}

The section entitled
\textbf{SUPERVISOR'S EVALUATION AND GRADUATION INTERNSHIP DEFENSE RESULTS}
is reserved for the supervisor and defense committee.

The evaluation may include:

\begin{itemize}[leftmargin=0.8cm, itemsep=2pt]
    \item Internship Progress and Student's Attitude;
    \item Report Content;
    \item Report Presentation;
    \item Other Comments;
    \item Supervisor's Opinion;
    \item Defense Result;
    \item Score.
\end{itemize}

Students shall not enter evaluation scores or defense results before they
are officially determined.

\vspace{0.3cm}

% ================================================================
% VIII. CHAPTERS, SECTIONS, AND SUBSECTIONS
% ================================================================

\noindent
\textbf{VIII.\quad CHAPTERS, SECTIONS, AND SUBSECTIONS}

\vspace{0.15cm}

The report structure shall be generated using the standard \LaTeX{}
commands:

\begin{verbatim}
\chapter{...}
\section{...}
\subsection{...}
\end{verbatim}

Chapter, section, and subsection numbers shall be generated automatically
by \LaTeX{}.

Students shall not manually type chapter or section numbers and shall not
reproduce heading styles using commands such as \verb|\textbf|,
\verb|\fontsize|, or \verb|\centering| instead of the appropriate
structural command.

Excessive subdivision should be avoided. In most cases, the hierarchy
Chapter--Section--Subsection is sufficient.

\vspace{0.3cm}

% ================================================================
% IX. ORGANIZATION OF THE MAIN CONTENT
% ================================================================

\noindent
\textbf{IX.\quad ORGANIZATION OF THE MAIN CONTENT}

\vspace{0.15cm}

The specific content of the report shall depend on the internship
organization and the assigned tasks. A typical structure may include:

\begin{itemize}[leftmargin=0.8cm, itemsep=2pt]
    \item \textbf{Chapter 1. Overview of the Internship Organization};
    \item \textbf{Chapter 2. Internship Activities and Assigned Tasks};
    \item \textbf{Chapter 3. Results, Experience, and Evaluation}.
\end{itemize}

These chapter titles are recommendations only and may be adapted to the
actual internship content.

The report shall reflect work actually performed, observed, analyzed, or
learned during the internship. Unrelated theoretical material should not
be included merely to increase report length.

\vspace{0.3cm}

% ================================================================
% X. TABLE OF CONTENTS
% ================================================================

\noindent
\textbf{X.\quad TABLE OF CONTENTS}

\vspace{0.15cm}

The Table of Contents shall be generated automatically using:

\begin{verbatim}
\tableofcontents
\end{verbatim}

Students shall not create or manually edit page numbers in the Table of
Contents.

The document shall be compiled sufficiently to ensure that all page
numbers and entries are correctly updated.

\vspace{0.3cm}

% ================================================================
% XI. LIST OF SYMBOLS AND ABBREVIATIONS
% ================================================================

\noindent
\textbf{XI.\quad LIST OF SYMBOLS AND ABBREVIATIONS}

\vspace{0.15cm}

Symbols and abbreviations that are used repeatedly in the report shall be
listed in the \textbf{List of Symbols and Abbreviations}.

The list may contain the following columns:

\begin{center}
\begin{tabular}{llll}
    No. & Symbol / Abbreviation & Unit & Definition
\end{tabular}
\end{center}

Only symbols and abbreviations actually used in the final report shall be
retained. Example entries contained in the template shall be removed if
they are not relevant to the student's report.

Each abbreviation should normally be defined when it first appears in
the main text.

\vspace{0.3cm}

% ================================================================
% XII. FIGURES, DIAGRAMS, AND PHOTOGRAPHS
% ================================================================

\noindent
\textbf{XII.\quad FIGURES, DIAGRAMS, AND PHOTOGRAPHS}

\vspace{0.15cm}

Figures shall be clear, legible, and directly relevant to the report
content.

Figures shall be inserted using the standard \verb|figure| environment.

Example:

\begin{verbatim}
\begin{figure}[H]
    \centering
    \includegraphics[width=0.8\textwidth]
                    {figures/company_structure.pdf}
    \caption{Organizational structure of the company}
    \label{fig:company-structure}
\end{figure}
\end{verbatim}

Figure captions shall be placed below the figure.

Each referenced figure shall have a unique \verb|\label| and shall be
referenced in the text using, for example:

\begin{verbatim}
Figure~\ref{fig:company-structure}
\end{verbatim}

Students shall not manually type figure numbers.

Figures reproduced or adapted from external sources shall include an
appropriate citation.

Photographs taken during the internship shall be relevant and sufficiently
clear. Students shall not include confidential, restricted, or proprietary
information without permission from the internship organization.

Vector graphics or high-resolution images are recommended whenever
possible.

\vspace{0.3cm}

% ================================================================
% XIII. TABLES
% ================================================================

\noindent
\textbf{XIII.\quad TABLES}

\vspace{0.15cm}

Tables shall be created using suitable \LaTeX{} environments such as
\verb|table|, \verb|tabular|, \verb|tabularx|, or \verb|longtable|.

Table captions shall be placed above the table.

Example:

\begin{verbatim}
\begin{table}[H]
    \centering
    \caption{Internship schedule}
    \label{tab:internship-schedule}

    \begin{tabular}{lll}
        \toprule
        Week & Task & Result \\
        \midrule
        1 & Orientation & Completed \\
        2 & System study & Completed \\
        \bottomrule
    \end{tabular}
\end{table}
\end{verbatim}

Tables shall be referenced using:

\begin{verbatim}
Table~\ref{tab:internship-schedule}
\end{verbatim}

Students shall not manually type table numbers.

Tables should be created directly in \LaTeX{} rather than inserted as
screenshots whenever practical.

\vspace{0.3cm}

% ================================================================
% XIV. EQUATIONS AND TECHNICAL QUANTITIES
% ================================================================

\noindent
\textbf{XIV.\quad EQUATIONS AND TECHNICAL QUANTITIES}

\vspace{0.15cm}

Mathematical expressions shall be typeset directly using \LaTeX{}.
Images or screenshots of equations shall not be used.

Short expressions may be written in inline mathematical mode, for example:

\begin{verbatim}
$R = 0.7\,\Omega$
\end{verbatim}

Important equations shall use a numbered equation environment.

Example:

\begin{verbatim}
\begin{equation}
    P = UI
    \label{eq:power}
\end{equation}
\end{verbatim}

Equations shall be referenced using:

\begin{verbatim}
Equation~\eqref{eq:power}
\end{verbatim}

Equation numbers shall not be manually entered.

SI units shall be used wherever applicable. The \verb|siunitx| package
should be used where appropriate.

Examples:

\begin{verbatim}
\SI{220}{\volt}
\SI{50}{\hertz}
\SI{2.5}{\ampere}
\end{verbatim}

Notation and units shall be consistent throughout the report.

\vspace{0.3cm}

% ================================================================
% XV. SOURCE CODE
% ================================================================

\noindent
\textbf{XV.\quad SOURCE CODE}

\vspace{0.15cm}

If the internship involves programming, PLCs, microcontrollers,
simulation, or software development, relevant source code may be included
in the report.

Source code shall be presented using the \verb|listings| environment or
another environment defined by the official template.

Example:

\begin{verbatim}
\begin{lstlisting}[language=st]
IF Start = TRUE THEN
    Motor := TRUE;
ELSE
    Motor := FALSE;
END_IF;
\end{lstlisting}
\end{verbatim}

Only code necessary for explaining the work should appear in the main
report. Long source-code listings should be moved to an appendix or
submitted as supplementary files if required.

Source code should not normally be inserted as screenshots.

\vspace{0.3cm}

% ================================================================
% XVI. ABSTRACT
% ================================================================

\noindent
\textbf{XVI.\quad ABSTRACT}

\vspace{0.15cm}

The Abstract shall briefly summarize:

\begin{itemize}[leftmargin=0.8cm, itemsep=2pt]
    \item the internship organization;
    \item the objectives of the internship;
    \item the main tasks performed;
    \item the main results, knowledge, and skills obtained.
\end{itemize}

The Abstract should be concise and should not reproduce the full content
of the report chapter by chapter.

\vspace{0.3cm}

% ================================================================
% XVII. REFERENCES AND CITATIONS
% ================================================================

\noindent
\textbf{XVII.\quad REFERENCES AND CITATIONS}

\vspace{0.15cm}

All books, journal papers, conference papers, standards, websites,
datasheets, manuals, technical documents, figures, tables, or other
external sources used in preparing the report shall be properly cited.

References shall be managed through the bibliography system provided with
the official template.

Bibliographic entries shall normally be stored in:

\begin{verbatim}
References.bib
\end{verbatim}

References shall be cited in the report using:

\begin{verbatim}
\cite{reference-key}
\end{verbatim}

Students shall not manually type reference numbers such as ``[1]'',
``[2]'', etc.

Titles of foreign-language publications shall remain in their original
language in the bibliography.

All borrowed ideas, data, figures, tables, equations, methods, or results
shall be properly acknowledged.

\vspace{0.3cm}

% ================================================================
% XVIII. CROSS-REFERENCING
% ================================================================

\noindent
\textbf{XVIII.\quad CROSS-REFERENCING}

\vspace{0.15cm}

Numbered elements shall be cross-referenced using the standard \LaTeX{}
mechanism.

Use the following commands where appropriate:

\begin{verbatim}
\label{...}
\ref{...}
\eqref{...}
\end{verbatim}

For example:

\begin{verbatim}
As shown in Figure~\ref{fig:system}, ...
\end{verbatim}

Students shall not manually type numbers such as ``Figure 2.3'' when
those numbers can be generated automatically.

\vspace{0.3cm}

% ================================================================
% XIX. PAGE NUMBERING
% ================================================================

\noindent
\textbf{XIX.\quad PAGE NUMBERING}

\vspace{0.15cm}

Cover pages and official evaluation forms at the beginning of the report
shall not display page numbers.

The main body of the report shall use Arabic page numbers beginning from
Chapter 1.

Page numbering shall be generated automatically by the official template.
Students shall not manually enter page numbers.

\vspace{0.3cm}

% ================================================================
% XX. PROJECT FILE ORGANIZATION
% ================================================================

\noindent
\textbf{XX.\quad PROJECT FILE ORGANIZATION}

\vspace{0.15cm}

Students should maintain a clear and organized \LaTeX{} project
structure. A recommended structure is:

\begin{verbatim}
InternshipReport/
|
|-- main.tex
|-- References.bib
|
|-- figures/
|   |-- company.pdf
|   |-- system.png
|   `-- ...
|
`-- chapters/
    |-- chapter1.tex
    |-- chapter2.tex
    `-- chapter3.tex
\end{verbatim}

Figures should normally be stored in the \verb|figures| directory
specified by the official template.

File and directory names should preferably contain Latin letters,
numbers, hyphens, or underscores.

Recommended examples include:

\begin{verbatim}
company_structure.pdf
motor_control.png
chapter_01.tex
\end{verbatim}

Spaces and unnecessary special characters in file names should be
avoided where possible.

\vspace{0.3cm}

% ================================================================
% XXI. DOCUMENT COMPILATION AND FINAL CHECK
% ================================================================

\noindent
\textbf{XXI.\quad DOCUMENT COMPILATION AND FINAL CHECK}

\vspace{0.15cm}

Before submission, the complete \LaTeX{} project shall compile
successfully and produce the final PDF report.

The final document shall not contain unresolved references such as:

\begin{verbatim}
??
[?]
\end{verbatim}

Before generating the final PDF, students shall check:

\begin{itemize}[leftmargin=0.8cm, itemsep=2pt]
    \item the Table of Contents;
    \item page numbering;
    \item chapter, section, and subsection numbering;
    \item figure numbering;
    \item table numbering;
    \item equation numbering;
    \item cross-references;
    \item citations;
    \item bibliography;
    \item font display;
    \item mathematical symbols;
    \item image quality;
    \item completeness of all required forms.
\end{itemize}

The document may be compiled using TeX Live, MiKTeX, or another
compatible \LaTeX{} environment. Regardless of the compilation
environment, the resulting PDF shall conform to the official Faculty
template.

\vspace{0.3cm}

% ================================================================
% XXII. CONTENT THAT STUDENTS MAY EDIT
% ================================================================

\noindent
\textbf{XXII.\quad CONTENT THAT STUDENTS MAY EDIT}

\vspace{0.15cm}

Students may edit or replace:

\begin{itemize}[leftmargin=0.8cm, itemsep=2pt]
    \item student name;
    \item class;
    \item internship organization;
    \item supervisor name;
    \item Abstract;
    \item List of Symbols and Abbreviations;
    \item chapter titles and chapter content;
    \item figures and photographs;
    \item tables;
    \item equations;
    \item citations and references.
\end{itemize}

Students shall not modify the main document formatting settings in the
preamble, including the document class, page geometry, line spacing,
heading formatting, table-of-contents formatting, font configuration,
page layout, or cover-page macros unless specifically instructed.

\vspace{0.3cm}

% ================================================================
% XXIII. FINAL SUBMISSION REQUIREMENTS
% ================================================================

\noindent
\textbf{XXIII.\quad FINAL SUBMISSION REQUIREMENTS}

\vspace{0.15cm}

Before submission, the final report shall:

\begin{itemize}[leftmargin=0.8cm, itemsep=2pt]
    \item compile without errors that affect the report content or presentation;
    \item contain all required forms and report sections;
    \item contain correct student, class, internship organization, and supervisor information;
    \item contain an updated Table of Contents, List of Figures, and List of Tables;
    \item contain no unused example content from the template;
    \item contain no placeholder text such as \textit{Title of Chapter 1},
    \textit{Section 1.1}, or other sample entries;
    \item contain all required citations and references;
    \item contain no unresolved cross-references.
\end{itemize}

If the Faculty requires submission of the \LaTeX{} source, the complete
project shall include all \verb|.tex| files, the bibliography database,
figures, and any other files required to reproduce the submitted PDF.

\vspace{0.4cm}

\noindent
\textbf{Important:}
Before submitting the final internship report, students shall remove or
comment out the following command in the main document:

\begin{verbatim}

% =================================================================
% REGULATIONS FOR PREPARING THE GRADUATION INTERNSHIP REPORT
% USING LATEX
% =================================================================
%
% This file is intended to be included from the main document with:
%
%     \input{Regulations.tex}
%
% IMPORTANT:
% This section is provided for guidance only. Students shall remove
% or comment out the \input{Regulations.tex} command before submitting
% the final internship report.
%
% =================================================================

\clearpage
\thispagestyle{empty}

\begin{center}
    {\fontsize{16pt}{20pt}\selectfont\bfseries
    SPECIFIC REGULATIONS FOR PREPARING\\
    GRADUATION INTERNSHIP REPORTS USING \LaTeX}
\end{center}

\vspace{0.5cm}

{\fontsize{13pt}{17pt}\selectfont
\setstretch{1.3}
\justifying

\noindent
\textbf{Note:}
This section is provided for guidance only and is included in the main
document using the command \verb|\input{Regulations.tex}|.
Students shall remove or comment out this command before submitting the
final internship report.

\vspace{0.4cm}

% ================================================================
% I. GENERAL REQUIREMENTS
% ================================================================

\noindent
\textbf{I.\quad GENERAL REQUIREMENTS}

\vspace{0.15cm}

The Graduation Internship Report shall be prepared using the official
\LaTeX{} template provided by the Faculty. Students shall preserve the
predefined page layout, typography, line spacing, heading styles,
numbering system, figure and table formatting, and bibliography settings
unless otherwise instructed by the Faculty or the supervisor.

The report shall be written in English.

Students are responsible for the academic and technical content of the
report. The official \LaTeX{} template is responsible for standardized
document formatting.

\vspace{0.3cm}

% ================================================================
% II. REPORT STRUCTURE
% ================================================================

\noindent
\textbf{II.\quad REPORT STRUCTURE}

\vspace{0.15cm}

The report shall normally be arranged in the following order:

\begin{enumerate}[label=\arabic*., leftmargin=1.2cm, itemsep=2pt]
    \item Outer Cover;
    \item Comments of the Internship Host Organization;
    \item Supervisor's Evaluation and Graduation Internship Defense Results;
    \item Table of Contents;
    \item List of Symbols and Abbreviations, if applicable;
    \item List of Figures;
    \item List of Tables;
    \item Abstract;
    \item Main Chapters of the Internship Report;
    \item References;
    \item Appendices, if applicable.
\end{enumerate}

The preliminary pages shall not be numbered as chapters. The main report
content shall be organized into consecutively numbered chapters beginning
with Chapter 1.

\vspace{0.3cm}

% ================================================================
% III. PAGE SIZE AND MARGINS
% ================================================================

\noindent
\textbf{III.\quad PAGE SIZE AND MARGINS}

\vspace{0.15cm}

The report shall be formatted on A4 paper with the following margins:

\begin{center}
\begin{tabular}{ll}
    Left margin:   & 3.5 cm \\
    Right margin:  & 2.0 cm \\
    Top margin:    & 2.5 cm \\
    Bottom margin: & 2.5 cm
\end{tabular}
\end{center}

These settings are predefined in the official template and shall not be
modified by students.

\vspace{0.3cm}

% ================================================================
% IV. FONT, FONT SIZE, AND LINE SPACING
% ================================================================

\noindent
\textbf{IV.\quad FONT, FONT SIZE, AND LINE SPACING}

\vspace{0.15cm}

The main text shall use a 13 pt font size with 1.3-line spacing.
A Times-family typeface shall be used in accordance with the official
template.

The text shall be fully justified.

Students shall not use repeated spaces, repeated manual line breaks, or
multiple blank lines to position text. Paragraph and heading spacing shall
be controlled by the template.

\vspace{0.3cm}

% ================================================================
% V. OUTER COVER
% ================================================================

\noindent
\textbf{V.\quad OUTER COVER}

\vspace{0.15cm}

The outer cover shall follow the official format provided in the template
and shall contain the following information:

\begin{itemize}[leftmargin=0.8cm, itemsep=2pt]
    \item Thai Nguyen University of Technology;
    \item Faculty of International Training;
    \item Engineering and Technology Department;
    \item University logo and Faculty logo;
    \item the title \textbf{INTERNSHIP REPORT};
    \item Internship Organization;
    \item Author;
    \item Class;
    \item Supervisor;
    \item place and year of completion.
\end{itemize}

Students shall only replace the required information fields and shall not
modify the cover-page layout unless instructed otherwise.

\vspace{0.3cm}

% ================================================================
% VI. COMMENTS OF THE INTERNSHIP HOST ORGANIZATION
% ================================================================

\noindent
\textbf{VI.\quad COMMENTS OF THE INTERNSHIP HOST ORGANIZATION}

\vspace{0.15cm}

The page entitled \textbf{COMMENTS OF THE INTERNSHIP HOST ORGANIZATION}
is reserved for the internship host organization.

It includes:

\begin{itemize}[leftmargin=0.8cm, itemsep=2pt]
    \item the student's full name;
    \item class;
    \item comments from the host organization;
    \item score;
    \item date;
    \item signature and official seal.
\end{itemize}

Students shall not alter or complete official evaluation fields unless
specifically authorized.

\vspace{0.3cm}

% ================================================================
% VII. SUPERVISOR'S EVALUATION AND DEFENSE RESULTS
% ================================================================

\noindent
\textbf{VII.\quad SUPERVISOR'S EVALUATION AND DEFENSE RESULTS}

\vspace{0.15cm}

The section entitled
\textbf{SUPERVISOR'S EVALUATION AND GRADUATION INTERNSHIP DEFENSE RESULTS}
is reserved for the supervisor and defense committee.

The evaluation may include:

\begin{itemize}[leftmargin=0.8cm, itemsep=2pt]
    \item Internship Progress and Student's Attitude;
    \item Report Content;
    \item Report Presentation;
    \item Other Comments;
    \item Supervisor's Opinion;
    \item Defense Result;
    \item Score.
\end{itemize}

Students shall not enter evaluation scores or defense results before they
are officially determined.

\vspace{0.3cm}

% ================================================================
% VIII. CHAPTERS, SECTIONS, AND SUBSECTIONS
% ================================================================

\noindent
\textbf{VIII.\quad CHAPTERS, SECTIONS, AND SUBSECTIONS}

\vspace{0.15cm}

The report structure shall be generated using the standard \LaTeX{}
commands:

\begin{verbatim}
\chapter{...}
\section{...}
\subsection{...}
\end{verbatim}

Chapter, section, and subsection numbers shall be generated automatically
by \LaTeX{}.

Students shall not manually type chapter or section numbers and shall not
reproduce heading styles using commands such as \verb|\textbf|,
\verb|\fontsize|, or \verb|\centering| instead of the appropriate
structural command.

Excessive subdivision should be avoided. In most cases, the hierarchy
Chapter--Section--Subsection is sufficient.

\vspace{0.3cm}

% ================================================================
% IX. ORGANIZATION OF THE MAIN CONTENT
% ================================================================

\noindent
\textbf{IX.\quad ORGANIZATION OF THE MAIN CONTENT}

\vspace{0.15cm}

The specific content of the report shall depend on the internship
organization and the assigned tasks. A typical structure may include:

\begin{itemize}[leftmargin=0.8cm, itemsep=2pt]
    \item \textbf{Chapter 1. Overview of the Internship Organization};
    \item \textbf{Chapter 2. Internship Activities and Assigned Tasks};
    \item \textbf{Chapter 3. Results, Experience, and Evaluation}.
\end{itemize}

These chapter titles are recommendations only and may be adapted to the
actual internship content.

The report shall reflect work actually performed, observed, analyzed, or
learned during the internship. Unrelated theoretical material should not
be included merely to increase report length.

\vspace{0.3cm}

% ================================================================
% X. TABLE OF CONTENTS
% ================================================================

\noindent
\textbf{X.\quad TABLE OF CONTENTS}

\vspace{0.15cm}

The Table of Contents shall be generated automatically using:

\begin{verbatim}
\tableofcontents
\end{verbatim}

Students shall not create or manually edit page numbers in the Table of
Contents.

The document shall be compiled sufficiently to ensure that all page
numbers and entries are correctly updated.

\vspace{0.3cm}

% ================================================================
% XI. LIST OF SYMBOLS AND ABBREVIATIONS
% ================================================================

\noindent
\textbf{XI.\quad LIST OF SYMBOLS AND ABBREVIATIONS}

\vspace{0.15cm}

Symbols and abbreviations that are used repeatedly in the report shall be
listed in the \textbf{List of Symbols and Abbreviations}.

The list may contain the following columns:

\begin{center}
\begin{tabular}{llll}
    No. & Symbol / Abbreviation & Unit & Definition
\end{tabular}
\end{center}

Only symbols and abbreviations actually used in the final report shall be
retained. Example entries contained in the template shall be removed if
they are not relevant to the student's report.

Each abbreviation should normally be defined when it first appears in
the main text.

\vspace{0.3cm}

% ================================================================
% XII. FIGURES, DIAGRAMS, AND PHOTOGRAPHS
% ================================================================

\noindent
\textbf{XII.\quad FIGURES, DIAGRAMS, AND PHOTOGRAPHS}

\vspace{0.15cm}

Figures shall be clear, legible, and directly relevant to the report
content.

Figures shall be inserted using the standard \verb|figure| environment.

Example:

\begin{verbatim}
\begin{figure}[H]
    \centering
    \includegraphics[width=0.8\textwidth]
                    {figures/company_structure.pdf}
    \caption{Organizational structure of the company}
    \label{fig:company-structure}
\end{figure}
\end{verbatim}

Figure captions shall be placed below the figure.

Each referenced figure shall have a unique \verb|\label| and shall be
referenced in the text using, for example:

\begin{verbatim}
Figure~\ref{fig:company-structure}
\end{verbatim}

Students shall not manually type figure numbers.

Figures reproduced or adapted from external sources shall include an
appropriate citation.

Photographs taken during the internship shall be relevant and sufficiently
clear. Students shall not include confidential, restricted, or proprietary
information without permission from the internship organization.

Vector graphics or high-resolution images are recommended whenever
possible.

\vspace{0.3cm}

% ================================================================
% XIII. TABLES
% ================================================================

\noindent
\textbf{XIII.\quad TABLES}

\vspace{0.15cm}

Tables shall be created using suitable \LaTeX{} environments such as
\verb|table|, \verb|tabular|, \verb|tabularx|, or \verb|longtable|.

Table captions shall be placed above the table.

Example:

\begin{verbatim}
\begin{table}[H]
    \centering
    \caption{Internship schedule}
    \label{tab:internship-schedule}

    \begin{tabular}{lll}
        \toprule
        Week & Task & Result \\
        \midrule
        1 & Orientation & Completed \\
        2 & System study & Completed \\
        \bottomrule
    \end{tabular}
\end{table}
\end{verbatim}

Tables shall be referenced using:

\begin{verbatim}
Table~\ref{tab:internship-schedule}
\end{verbatim}

Students shall not manually type table numbers.

Tables should be created directly in \LaTeX{} rather than inserted as
screenshots whenever practical.

\vspace{0.3cm}

% ================================================================
% XIV. EQUATIONS AND TECHNICAL QUANTITIES
% ================================================================

\noindent
\textbf{XIV.\quad EQUATIONS AND TECHNICAL QUANTITIES}

\vspace{0.15cm}

Mathematical expressions shall be typeset directly using \LaTeX{}.
Images or screenshots of equations shall not be used.

Short expressions may be written in inline mathematical mode, for example:

\begin{verbatim}
$R = 0.7\,\Omega$
\end{verbatim}

Important equations shall use a numbered equation environment.

Example:

\begin{verbatim}
\begin{equation}
    P = UI
    \label{eq:power}
\end{equation}
\end{verbatim}

Equations shall be referenced using:

\begin{verbatim}
Equation~\eqref{eq:power}
\end{verbatim}

Equation numbers shall not be manually entered.

SI units shall be used wherever applicable. The \verb|siunitx| package
should be used where appropriate.

Examples:

\begin{verbatim}
\SI{220}{\volt}
\SI{50}{\hertz}
\SI{2.5}{\ampere}
\end{verbatim}

Notation and units shall be consistent throughout the report.

\vspace{0.3cm}

% ================================================================
% XV. SOURCE CODE
% ================================================================

\noindent
\textbf{XV.\quad SOURCE CODE}

\vspace{0.15cm}

If the internship involves programming, PLCs, microcontrollers,
simulation, or software development, relevant source code may be included
in the report.

Source code shall be presented using the \verb|listings| environment or
another environment defined by the official template.

Example:

\begin{verbatim}
\begin{lstlisting}[language=st]
IF Start = TRUE THEN
    Motor := TRUE;
ELSE
    Motor := FALSE;
END_IF;
\end{lstlisting}
\end{verbatim}

Only code necessary for explaining the work should appear in the main
report. Long source-code listings should be moved to an appendix or
submitted as supplementary files if required.

Source code should not normally be inserted as screenshots.

\vspace{0.3cm}

% ================================================================
% XVI. ABSTRACT
% ================================================================

\noindent
\textbf{XVI.\quad ABSTRACT}

\vspace{0.15cm}

The Abstract shall briefly summarize:

\begin{itemize}[leftmargin=0.8cm, itemsep=2pt]
    \item the internship organization;
    \item the objectives of the internship;
    \item the main tasks performed;
    \item the main results, knowledge, and skills obtained.
\end{itemize}

The Abstract should be concise and should not reproduce the full content
of the report chapter by chapter.

\vspace{0.3cm}

% ================================================================
% XVII. REFERENCES AND CITATIONS
% ================================================================

\noindent
\textbf{XVII.\quad REFERENCES AND CITATIONS}

\vspace{0.15cm}

All books, journal papers, conference papers, standards, websites,
datasheets, manuals, technical documents, figures, tables, or other
external sources used in preparing the report shall be properly cited.

References shall be managed through the bibliography system provided with
the official template.

Bibliographic entries shall normally be stored in:

\begin{verbatim}
References.bib
\end{verbatim}

References shall be cited in the report using:

\begin{verbatim}
\cite{reference-key}
\end{verbatim}

Students shall not manually type reference numbers such as ``[1]'',
``[2]'', etc.

Titles of foreign-language publications shall remain in their original
language in the bibliography.

All borrowed ideas, data, figures, tables, equations, methods, or results
shall be properly acknowledged.

\vspace{0.3cm}

% ================================================================
% XVIII. CROSS-REFERENCING
% ================================================================

\noindent
\textbf{XVIII.\quad CROSS-REFERENCING}

\vspace{0.15cm}

Numbered elements shall be cross-referenced using the standard \LaTeX{}
mechanism.

Use the following commands where appropriate:

\begin{verbatim}
\label{...}
\ref{...}
\eqref{...}
\end{verbatim}

For example:

\begin{verbatim}
As shown in Figure~\ref{fig:system}, ...
\end{verbatim}

Students shall not manually type numbers such as ``Figure 2.3'' when
those numbers can be generated automatically.

\vspace{0.3cm}

% ================================================================
% XIX. PAGE NUMBERING
% ================================================================

\noindent
\textbf{XIX.\quad PAGE NUMBERING}

\vspace{0.15cm}

Cover pages and official evaluation forms at the beginning of the report
shall not display page numbers.

The main body of the report shall use Arabic page numbers beginning from
Chapter 1.

Page numbering shall be generated automatically by the official template.
Students shall not manually enter page numbers.

\vspace{0.3cm}

% ================================================================
% XX. PROJECT FILE ORGANIZATION
% ================================================================

\noindent
\textbf{XX.\quad PROJECT FILE ORGANIZATION}

\vspace{0.15cm}

Students should maintain a clear and organized \LaTeX{} project
structure. A recommended structure is:

\begin{verbatim}
InternshipReport/
|
|-- main.tex
|-- References.bib
|
|-- figures/
|   |-- company.pdf
|   |-- system.png
|   `-- ...
|
`-- chapters/
    |-- chapter1.tex
    |-- chapter2.tex
    `-- chapter3.tex
\end{verbatim}

Figures should normally be stored in the \verb|figures| directory
specified by the official template.

File and directory names should preferably contain Latin letters,
numbers, hyphens, or underscores.

Recommended examples include:

\begin{verbatim}
company_structure.pdf
motor_control.png
chapter_01.tex
\end{verbatim}

Spaces and unnecessary special characters in file names should be
avoided where possible.

\vspace{0.3cm}

% ================================================================
% XXI. DOCUMENT COMPILATION AND FINAL CHECK
% ================================================================

\noindent
\textbf{XXI.\quad DOCUMENT COMPILATION AND FINAL CHECK}

\vspace{0.15cm}

Before submission, the complete \LaTeX{} project shall compile
successfully and produce the final PDF report.

The final document shall not contain unresolved references such as:

\begin{verbatim}
??
[?]
\end{verbatim}

Before generating the final PDF, students shall check:

\begin{itemize}[leftmargin=0.8cm, itemsep=2pt]
    \item the Table of Contents;
    \item page numbering;
    \item chapter, section, and subsection numbering;
    \item figure numbering;
    \item table numbering;
    \item equation numbering;
    \item cross-references;
    \item citations;
    \item bibliography;
    \item font display;
    \item mathematical symbols;
    \item image quality;
    \item completeness of all required forms.
\end{itemize}

The document may be compiled using TeX Live, MiKTeX, or another
compatible \LaTeX{} environment. Regardless of the compilation
environment, the resulting PDF shall conform to the official Faculty
template.

\vspace{0.3cm}

% ================================================================
% XXII. CONTENT THAT STUDENTS MAY EDIT
% ================================================================

\noindent
\textbf{XXII.\quad CONTENT THAT STUDENTS MAY EDIT}

\vspace{0.15cm}

Students may edit or replace:

\begin{itemize}[leftmargin=0.8cm, itemsep=2pt]
    \item student name;
    \item class;
    \item internship organization;
    \item supervisor name;
    \item Abstract;
    \item List of Symbols and Abbreviations;
    \item chapter titles and chapter content;
    \item figures and photographs;
    \item tables;
    \item equations;
    \item citations and references.
\end{itemize}

Students shall not modify the main document formatting settings in the
preamble, including the document class, page geometry, line spacing,
heading formatting, table-of-contents formatting, font configuration,
page layout, or cover-page macros unless specifically instructed.

\vspace{0.3cm}

% ================================================================
% XXIII. FINAL SUBMISSION REQUIREMENTS
% ================================================================

\noindent
\textbf{XXIII.\quad FINAL SUBMISSION REQUIREMENTS}

\vspace{0.15cm}

Before submission, the final report shall:

\begin{itemize}[leftmargin=0.8cm, itemsep=2pt]
    \item compile without errors that affect the report content or presentation;
    \item contain all required forms and report sections;
    \item contain correct student, class, internship organization, and supervisor information;
    \item contain an updated Table of Contents, List of Figures, and List of Tables;
    \item contain no unused example content from the template;
    \item contain no placeholder text such as \textit{Title of Chapter 1},
    \textit{Section 1.1}, or other sample entries;
    \item contain all required citations and references;
    \item contain no unresolved cross-references.
\end{itemize}

If the Faculty requires submission of the \LaTeX{} source, the complete
project shall include all \verb|.tex| files, the bibliography database,
figures, and any other files required to reproduce the submitted PDF.

\vspace{0.4cm}

\noindent
\textbf{Important:}
Before submitting the final internship report, students shall remove or
comment out the following command in the main document:

\begin{verbatim}
\input{Regulations.tex}
\end{verbatim}

For example:

\begin{verbatim}
% \input{Regulations.tex}
\end{verbatim}

This removes the guidance pages from the final submitted report while
preserving the regulations file in the \LaTeX{} project for reference.

} % End 13pt content

\clearpage

\end{verbatim}

For example:

\begin{verbatim}
% 
% =================================================================
% REGULATIONS FOR PREPARING THE GRADUATION INTERNSHIP REPORT
% USING LATEX
% =================================================================
%
% This file is intended to be included from the main document with:
%
%     \input{Regulations.tex}
%
% IMPORTANT:
% This section is provided for guidance only. Students shall remove
% or comment out the \input{Regulations.tex} command before submitting
% the final internship report.
%
% =================================================================

\clearpage
\thispagestyle{empty}

\begin{center}
    {\fontsize{16pt}{20pt}\selectfont\bfseries
    SPECIFIC REGULATIONS FOR PREPARING\\
    GRADUATION INTERNSHIP REPORTS USING \LaTeX}
\end{center}

\vspace{0.5cm}

{\fontsize{13pt}{17pt}\selectfont
\setstretch{1.3}
\justifying

\noindent
\textbf{Note:}
This section is provided for guidance only and is included in the main
document using the command \verb|\input{Regulations.tex}|.
Students shall remove or comment out this command before submitting the
final internship report.

\vspace{0.4cm}

% ================================================================
% I. GENERAL REQUIREMENTS
% ================================================================

\noindent
\textbf{I.\quad GENERAL REQUIREMENTS}

\vspace{0.15cm}

The Graduation Internship Report shall be prepared using the official
\LaTeX{} template provided by the Faculty. Students shall preserve the
predefined page layout, typography, line spacing, heading styles,
numbering system, figure and table formatting, and bibliography settings
unless otherwise instructed by the Faculty or the supervisor.

The report shall be written in English.

Students are responsible for the academic and technical content of the
report. The official \LaTeX{} template is responsible for standardized
document formatting.

\vspace{0.3cm}

% ================================================================
% II. REPORT STRUCTURE
% ================================================================

\noindent
\textbf{II.\quad REPORT STRUCTURE}

\vspace{0.15cm}

The report shall normally be arranged in the following order:

\begin{enumerate}[label=\arabic*., leftmargin=1.2cm, itemsep=2pt]
    \item Outer Cover;
    \item Comments of the Internship Host Organization;
    \item Supervisor's Evaluation and Graduation Internship Defense Results;
    \item Table of Contents;
    \item List of Symbols and Abbreviations, if applicable;
    \item List of Figures;
    \item List of Tables;
    \item Abstract;
    \item Main Chapters of the Internship Report;
    \item References;
    \item Appendices, if applicable.
\end{enumerate}

The preliminary pages shall not be numbered as chapters. The main report
content shall be organized into consecutively numbered chapters beginning
with Chapter 1.

\vspace{0.3cm}

% ================================================================
% III. PAGE SIZE AND MARGINS
% ================================================================

\noindent
\textbf{III.\quad PAGE SIZE AND MARGINS}

\vspace{0.15cm}

The report shall be formatted on A4 paper with the following margins:

\begin{center}
\begin{tabular}{ll}
    Left margin:   & 3.5 cm \\
    Right margin:  & 2.0 cm \\
    Top margin:    & 2.5 cm \\
    Bottom margin: & 2.5 cm
\end{tabular}
\end{center}

These settings are predefined in the official template and shall not be
modified by students.

\vspace{0.3cm}

% ================================================================
% IV. FONT, FONT SIZE, AND LINE SPACING
% ================================================================

\noindent
\textbf{IV.\quad FONT, FONT SIZE, AND LINE SPACING}

\vspace{0.15cm}

The main text shall use a 13 pt font size with 1.3-line spacing.
A Times-family typeface shall be used in accordance with the official
template.

The text shall be fully justified.

Students shall not use repeated spaces, repeated manual line breaks, or
multiple blank lines to position text. Paragraph and heading spacing shall
be controlled by the template.

\vspace{0.3cm}

% ================================================================
% V. OUTER COVER
% ================================================================

\noindent
\textbf{V.\quad OUTER COVER}

\vspace{0.15cm}

The outer cover shall follow the official format provided in the template
and shall contain the following information:

\begin{itemize}[leftmargin=0.8cm, itemsep=2pt]
    \item Thai Nguyen University of Technology;
    \item Faculty of International Training;
    \item Engineering and Technology Department;
    \item University logo and Faculty logo;
    \item the title \textbf{INTERNSHIP REPORT};
    \item Internship Organization;
    \item Author;
    \item Class;
    \item Supervisor;
    \item place and year of completion.
\end{itemize}

Students shall only replace the required information fields and shall not
modify the cover-page layout unless instructed otherwise.

\vspace{0.3cm}

% ================================================================
% VI. COMMENTS OF THE INTERNSHIP HOST ORGANIZATION
% ================================================================

\noindent
\textbf{VI.\quad COMMENTS OF THE INTERNSHIP HOST ORGANIZATION}

\vspace{0.15cm}

The page entitled \textbf{COMMENTS OF THE INTERNSHIP HOST ORGANIZATION}
is reserved for the internship host organization.

It includes:

\begin{itemize}[leftmargin=0.8cm, itemsep=2pt]
    \item the student's full name;
    \item class;
    \item comments from the host organization;
    \item score;
    \item date;
    \item signature and official seal.
\end{itemize}

Students shall not alter or complete official evaluation fields unless
specifically authorized.

\vspace{0.3cm}

% ================================================================
% VII. SUPERVISOR'S EVALUATION AND DEFENSE RESULTS
% ================================================================

\noindent
\textbf{VII.\quad SUPERVISOR'S EVALUATION AND DEFENSE RESULTS}

\vspace{0.15cm}

The section entitled
\textbf{SUPERVISOR'S EVALUATION AND GRADUATION INTERNSHIP DEFENSE RESULTS}
is reserved for the supervisor and defense committee.

The evaluation may include:

\begin{itemize}[leftmargin=0.8cm, itemsep=2pt]
    \item Internship Progress and Student's Attitude;
    \item Report Content;
    \item Report Presentation;
    \item Other Comments;
    \item Supervisor's Opinion;
    \item Defense Result;
    \item Score.
\end{itemize}

Students shall not enter evaluation scores or defense results before they
are officially determined.

\vspace{0.3cm}

% ================================================================
% VIII. CHAPTERS, SECTIONS, AND SUBSECTIONS
% ================================================================

\noindent
\textbf{VIII.\quad CHAPTERS, SECTIONS, AND SUBSECTIONS}

\vspace{0.15cm}

The report structure shall be generated using the standard \LaTeX{}
commands:

\begin{verbatim}
\chapter{...}
\section{...}
\subsection{...}
\end{verbatim}

Chapter, section, and subsection numbers shall be generated automatically
by \LaTeX{}.

Students shall not manually type chapter or section numbers and shall not
reproduce heading styles using commands such as \verb|\textbf|,
\verb|\fontsize|, or \verb|\centering| instead of the appropriate
structural command.

Excessive subdivision should be avoided. In most cases, the hierarchy
Chapter--Section--Subsection is sufficient.

\vspace{0.3cm}

% ================================================================
% IX. ORGANIZATION OF THE MAIN CONTENT
% ================================================================

\noindent
\textbf{IX.\quad ORGANIZATION OF THE MAIN CONTENT}

\vspace{0.15cm}

The specific content of the report shall depend on the internship
organization and the assigned tasks. A typical structure may include:

\begin{itemize}[leftmargin=0.8cm, itemsep=2pt]
    \item \textbf{Chapter 1. Overview of the Internship Organization};
    \item \textbf{Chapter 2. Internship Activities and Assigned Tasks};
    \item \textbf{Chapter 3. Results, Experience, and Evaluation}.
\end{itemize}

These chapter titles are recommendations only and may be adapted to the
actual internship content.

The report shall reflect work actually performed, observed, analyzed, or
learned during the internship. Unrelated theoretical material should not
be included merely to increase report length.

\vspace{0.3cm}

% ================================================================
% X. TABLE OF CONTENTS
% ================================================================

\noindent
\textbf{X.\quad TABLE OF CONTENTS}

\vspace{0.15cm}

The Table of Contents shall be generated automatically using:

\begin{verbatim}
\tableofcontents
\end{verbatim}

Students shall not create or manually edit page numbers in the Table of
Contents.

The document shall be compiled sufficiently to ensure that all page
numbers and entries are correctly updated.

\vspace{0.3cm}

% ================================================================
% XI. LIST OF SYMBOLS AND ABBREVIATIONS
% ================================================================

\noindent
\textbf{XI.\quad LIST OF SYMBOLS AND ABBREVIATIONS}

\vspace{0.15cm}

Symbols and abbreviations that are used repeatedly in the report shall be
listed in the \textbf{List of Symbols and Abbreviations}.

The list may contain the following columns:

\begin{center}
\begin{tabular}{llll}
    No. & Symbol / Abbreviation & Unit & Definition
\end{tabular}
\end{center}

Only symbols and abbreviations actually used in the final report shall be
retained. Example entries contained in the template shall be removed if
they are not relevant to the student's report.

Each abbreviation should normally be defined when it first appears in
the main text.

\vspace{0.3cm}

% ================================================================
% XII. FIGURES, DIAGRAMS, AND PHOTOGRAPHS
% ================================================================

\noindent
\textbf{XII.\quad FIGURES, DIAGRAMS, AND PHOTOGRAPHS}

\vspace{0.15cm}

Figures shall be clear, legible, and directly relevant to the report
content.

Figures shall be inserted using the standard \verb|figure| environment.

Example:

\begin{verbatim}
\begin{figure}[H]
    \centering
    \includegraphics[width=0.8\textwidth]
                    {figures/company_structure.pdf}
    \caption{Organizational structure of the company}
    \label{fig:company-structure}
\end{figure}
\end{verbatim}

Figure captions shall be placed below the figure.

Each referenced figure shall have a unique \verb|\label| and shall be
referenced in the text using, for example:

\begin{verbatim}
Figure~\ref{fig:company-structure}
\end{verbatim}

Students shall not manually type figure numbers.

Figures reproduced or adapted from external sources shall include an
appropriate citation.

Photographs taken during the internship shall be relevant and sufficiently
clear. Students shall not include confidential, restricted, or proprietary
information without permission from the internship organization.

Vector graphics or high-resolution images are recommended whenever
possible.

\vspace{0.3cm}

% ================================================================
% XIII. TABLES
% ================================================================

\noindent
\textbf{XIII.\quad TABLES}

\vspace{0.15cm}

Tables shall be created using suitable \LaTeX{} environments such as
\verb|table|, \verb|tabular|, \verb|tabularx|, or \verb|longtable|.

Table captions shall be placed above the table.

Example:

\begin{verbatim}
\begin{table}[H]
    \centering
    \caption{Internship schedule}
    \label{tab:internship-schedule}

    \begin{tabular}{lll}
        \toprule
        Week & Task & Result \\
        \midrule
        1 & Orientation & Completed \\
        2 & System study & Completed \\
        \bottomrule
    \end{tabular}
\end{table}
\end{verbatim}

Tables shall be referenced using:

\begin{verbatim}
Table~\ref{tab:internship-schedule}
\end{verbatim}

Students shall not manually type table numbers.

Tables should be created directly in \LaTeX{} rather than inserted as
screenshots whenever practical.

\vspace{0.3cm}

% ================================================================
% XIV. EQUATIONS AND TECHNICAL QUANTITIES
% ================================================================

\noindent
\textbf{XIV.\quad EQUATIONS AND TECHNICAL QUANTITIES}

\vspace{0.15cm}

Mathematical expressions shall be typeset directly using \LaTeX{}.
Images or screenshots of equations shall not be used.

Short expressions may be written in inline mathematical mode, for example:

\begin{verbatim}
$R = 0.7\,\Omega$
\end{verbatim}

Important equations shall use a numbered equation environment.

Example:

\begin{verbatim}
\begin{equation}
    P = UI
    \label{eq:power}
\end{equation}
\end{verbatim}

Equations shall be referenced using:

\begin{verbatim}
Equation~\eqref{eq:power}
\end{verbatim}

Equation numbers shall not be manually entered.

SI units shall be used wherever applicable. The \verb|siunitx| package
should be used where appropriate.

Examples:

\begin{verbatim}
\SI{220}{\volt}
\SI{50}{\hertz}
\SI{2.5}{\ampere}
\end{verbatim}

Notation and units shall be consistent throughout the report.

\vspace{0.3cm}

% ================================================================
% XV. SOURCE CODE
% ================================================================

\noindent
\textbf{XV.\quad SOURCE CODE}

\vspace{0.15cm}

If the internship involves programming, PLCs, microcontrollers,
simulation, or software development, relevant source code may be included
in the report.

Source code shall be presented using the \verb|listings| environment or
another environment defined by the official template.

Example:

\begin{verbatim}
\begin{lstlisting}[language=st]
IF Start = TRUE THEN
    Motor := TRUE;
ELSE
    Motor := FALSE;
END_IF;
\end{lstlisting}
\end{verbatim}

Only code necessary for explaining the work should appear in the main
report. Long source-code listings should be moved to an appendix or
submitted as supplementary files if required.

Source code should not normally be inserted as screenshots.

\vspace{0.3cm}

% ================================================================
% XVI. ABSTRACT
% ================================================================

\noindent
\textbf{XVI.\quad ABSTRACT}

\vspace{0.15cm}

The Abstract shall briefly summarize:

\begin{itemize}[leftmargin=0.8cm, itemsep=2pt]
    \item the internship organization;
    \item the objectives of the internship;
    \item the main tasks performed;
    \item the main results, knowledge, and skills obtained.
\end{itemize}

The Abstract should be concise and should not reproduce the full content
of the report chapter by chapter.

\vspace{0.3cm}

% ================================================================
% XVII. REFERENCES AND CITATIONS
% ================================================================

\noindent
\textbf{XVII.\quad REFERENCES AND CITATIONS}

\vspace{0.15cm}

All books, journal papers, conference papers, standards, websites,
datasheets, manuals, technical documents, figures, tables, or other
external sources used in preparing the report shall be properly cited.

References shall be managed through the bibliography system provided with
the official template.

Bibliographic entries shall normally be stored in:

\begin{verbatim}
References.bib
\end{verbatim}

References shall be cited in the report using:

\begin{verbatim}
\cite{reference-key}
\end{verbatim}

Students shall not manually type reference numbers such as ``[1]'',
``[2]'', etc.

Titles of foreign-language publications shall remain in their original
language in the bibliography.

All borrowed ideas, data, figures, tables, equations, methods, or results
shall be properly acknowledged.

\vspace{0.3cm}

% ================================================================
% XVIII. CROSS-REFERENCING
% ================================================================

\noindent
\textbf{XVIII.\quad CROSS-REFERENCING}

\vspace{0.15cm}

Numbered elements shall be cross-referenced using the standard \LaTeX{}
mechanism.

Use the following commands where appropriate:

\begin{verbatim}
\label{...}
\ref{...}
\eqref{...}
\end{verbatim}

For example:

\begin{verbatim}
As shown in Figure~\ref{fig:system}, ...
\end{verbatim}

Students shall not manually type numbers such as ``Figure 2.3'' when
those numbers can be generated automatically.

\vspace{0.3cm}

% ================================================================
% XIX. PAGE NUMBERING
% ================================================================

\noindent
\textbf{XIX.\quad PAGE NUMBERING}

\vspace{0.15cm}

Cover pages and official evaluation forms at the beginning of the report
shall not display page numbers.

The main body of the report shall use Arabic page numbers beginning from
Chapter 1.

Page numbering shall be generated automatically by the official template.
Students shall not manually enter page numbers.

\vspace{0.3cm}

% ================================================================
% XX. PROJECT FILE ORGANIZATION
% ================================================================

\noindent
\textbf{XX.\quad PROJECT FILE ORGANIZATION}

\vspace{0.15cm}

Students should maintain a clear and organized \LaTeX{} project
structure. A recommended structure is:

\begin{verbatim}
InternshipReport/
|
|-- main.tex
|-- References.bib
|
|-- figures/
|   |-- company.pdf
|   |-- system.png
|   `-- ...
|
`-- chapters/
    |-- chapter1.tex
    |-- chapter2.tex
    `-- chapter3.tex
\end{verbatim}

Figures should normally be stored in the \verb|figures| directory
specified by the official template.

File and directory names should preferably contain Latin letters,
numbers, hyphens, or underscores.

Recommended examples include:

\begin{verbatim}
company_structure.pdf
motor_control.png
chapter_01.tex
\end{verbatim}

Spaces and unnecessary special characters in file names should be
avoided where possible.

\vspace{0.3cm}

% ================================================================
% XXI. DOCUMENT COMPILATION AND FINAL CHECK
% ================================================================

\noindent
\textbf{XXI.\quad DOCUMENT COMPILATION AND FINAL CHECK}

\vspace{0.15cm}

Before submission, the complete \LaTeX{} project shall compile
successfully and produce the final PDF report.

The final document shall not contain unresolved references such as:

\begin{verbatim}
??
[?]
\end{verbatim}

Before generating the final PDF, students shall check:

\begin{itemize}[leftmargin=0.8cm, itemsep=2pt]
    \item the Table of Contents;
    \item page numbering;
    \item chapter, section, and subsection numbering;
    \item figure numbering;
    \item table numbering;
    \item equation numbering;
    \item cross-references;
    \item citations;
    \item bibliography;
    \item font display;
    \item mathematical symbols;
    \item image quality;
    \item completeness of all required forms.
\end{itemize}

The document may be compiled using TeX Live, MiKTeX, or another
compatible \LaTeX{} environment. Regardless of the compilation
environment, the resulting PDF shall conform to the official Faculty
template.

\vspace{0.3cm}

% ================================================================
% XXII. CONTENT THAT STUDENTS MAY EDIT
% ================================================================

\noindent
\textbf{XXII.\quad CONTENT THAT STUDENTS MAY EDIT}

\vspace{0.15cm}

Students may edit or replace:

\begin{itemize}[leftmargin=0.8cm, itemsep=2pt]
    \item student name;
    \item class;
    \item internship organization;
    \item supervisor name;
    \item Abstract;
    \item List of Symbols and Abbreviations;
    \item chapter titles and chapter content;
    \item figures and photographs;
    \item tables;
    \item equations;
    \item citations and references.
\end{itemize}

Students shall not modify the main document formatting settings in the
preamble, including the document class, page geometry, line spacing,
heading formatting, table-of-contents formatting, font configuration,
page layout, or cover-page macros unless specifically instructed.

\vspace{0.3cm}

% ================================================================
% XXIII. FINAL SUBMISSION REQUIREMENTS
% ================================================================

\noindent
\textbf{XXIII.\quad FINAL SUBMISSION REQUIREMENTS}

\vspace{0.15cm}

Before submission, the final report shall:

\begin{itemize}[leftmargin=0.8cm, itemsep=2pt]
    \item compile without errors that affect the report content or presentation;
    \item contain all required forms and report sections;
    \item contain correct student, class, internship organization, and supervisor information;
    \item contain an updated Table of Contents, List of Figures, and List of Tables;
    \item contain no unused example content from the template;
    \item contain no placeholder text such as \textit{Title of Chapter 1},
    \textit{Section 1.1}, or other sample entries;
    \item contain all required citations and references;
    \item contain no unresolved cross-references.
\end{itemize}

If the Faculty requires submission of the \LaTeX{} source, the complete
project shall include all \verb|.tex| files, the bibliography database,
figures, and any other files required to reproduce the submitted PDF.

\vspace{0.4cm}

\noindent
\textbf{Important:}
Before submitting the final internship report, students shall remove or
comment out the following command in the main document:

\begin{verbatim}
\input{Regulations.tex}
\end{verbatim}

For example:

\begin{verbatim}
% \input{Regulations.tex}
\end{verbatim}

This removes the guidance pages from the final submitted report while
preserving the regulations file in the \LaTeX{} project for reference.

} % End 13pt content

\clearpage

\end{verbatim}

This removes the guidance pages from the final submitted report while
preserving the regulations file in the \LaTeX{} project for reference.

} % End 13pt content

\clearpage

%
% IMPORTANT:
% This section is provided for guidance only. Students shall remove
% or comment out the 
% =================================================================
% REGULATIONS FOR PREPARING THE GRADUATION INTERNSHIP REPORT
% USING LATEX
% =================================================================
%
% This file is intended to be included from the main document with:
%
%     
% =================================================================
% REGULATIONS FOR PREPARING THE GRADUATION INTERNSHIP REPORT
% USING LATEX
% =================================================================
%
% This file is intended to be included from the main document with:
%
%     \input{Regulations.tex}
%
% IMPORTANT:
% This section is provided for guidance only. Students shall remove
% or comment out the \input{Regulations.tex} command before submitting
% the final internship report.
%
% =================================================================

\clearpage
\thispagestyle{empty}

\begin{center}
    {\fontsize{16pt}{20pt}\selectfont\bfseries
    SPECIFIC REGULATIONS FOR PREPARING\\
    GRADUATION INTERNSHIP REPORTS USING \LaTeX}
\end{center}

\vspace{0.5cm}

{\fontsize{13pt}{17pt}\selectfont
\setstretch{1.3}
\justifying

\noindent
\textbf{Note:}
This section is provided for guidance only and is included in the main
document using the command \verb|\input{Regulations.tex}|.
Students shall remove or comment out this command before submitting the
final internship report.

\vspace{0.4cm}

% ================================================================
% I. GENERAL REQUIREMENTS
% ================================================================

\noindent
\textbf{I.\quad GENERAL REQUIREMENTS}

\vspace{0.15cm}

The Graduation Internship Report shall be prepared using the official
\LaTeX{} template provided by the Faculty. Students shall preserve the
predefined page layout, typography, line spacing, heading styles,
numbering system, figure and table formatting, and bibliography settings
unless otherwise instructed by the Faculty or the supervisor.

The report shall be written in English.

Students are responsible for the academic and technical content of the
report. The official \LaTeX{} template is responsible for standardized
document formatting.

\vspace{0.3cm}

% ================================================================
% II. REPORT STRUCTURE
% ================================================================

\noindent
\textbf{II.\quad REPORT STRUCTURE}

\vspace{0.15cm}

The report shall normally be arranged in the following order:

\begin{enumerate}[label=\arabic*., leftmargin=1.2cm, itemsep=2pt]
    \item Outer Cover;
    \item Comments of the Internship Host Organization;
    \item Supervisor's Evaluation and Graduation Internship Defense Results;
    \item Table of Contents;
    \item List of Symbols and Abbreviations, if applicable;
    \item List of Figures;
    \item List of Tables;
    \item Abstract;
    \item Main Chapters of the Internship Report;
    \item References;
    \item Appendices, if applicable.
\end{enumerate}

The preliminary pages shall not be numbered as chapters. The main report
content shall be organized into consecutively numbered chapters beginning
with Chapter 1.

\vspace{0.3cm}

% ================================================================
% III. PAGE SIZE AND MARGINS
% ================================================================

\noindent
\textbf{III.\quad PAGE SIZE AND MARGINS}

\vspace{0.15cm}

The report shall be formatted on A4 paper with the following margins:

\begin{center}
\begin{tabular}{ll}
    Left margin:   & 3.5 cm \\
    Right margin:  & 2.0 cm \\
    Top margin:    & 2.5 cm \\
    Bottom margin: & 2.5 cm
\end{tabular}
\end{center}

These settings are predefined in the official template and shall not be
modified by students.

\vspace{0.3cm}

% ================================================================
% IV. FONT, FONT SIZE, AND LINE SPACING
% ================================================================

\noindent
\textbf{IV.\quad FONT, FONT SIZE, AND LINE SPACING}

\vspace{0.15cm}

The main text shall use a 13 pt font size with 1.3-line spacing.
A Times-family typeface shall be used in accordance with the official
template.

The text shall be fully justified.

Students shall not use repeated spaces, repeated manual line breaks, or
multiple blank lines to position text. Paragraph and heading spacing shall
be controlled by the template.

\vspace{0.3cm}

% ================================================================
% V. OUTER COVER
% ================================================================

\noindent
\textbf{V.\quad OUTER COVER}

\vspace{0.15cm}

The outer cover shall follow the official format provided in the template
and shall contain the following information:

\begin{itemize}[leftmargin=0.8cm, itemsep=2pt]
    \item Thai Nguyen University of Technology;
    \item Faculty of International Training;
    \item Engineering and Technology Department;
    \item University logo and Faculty logo;
    \item the title \textbf{INTERNSHIP REPORT};
    \item Internship Organization;
    \item Author;
    \item Class;
    \item Supervisor;
    \item place and year of completion.
\end{itemize}

Students shall only replace the required information fields and shall not
modify the cover-page layout unless instructed otherwise.

\vspace{0.3cm}

% ================================================================
% VI. COMMENTS OF THE INTERNSHIP HOST ORGANIZATION
% ================================================================

\noindent
\textbf{VI.\quad COMMENTS OF THE INTERNSHIP HOST ORGANIZATION}

\vspace{0.15cm}

The page entitled \textbf{COMMENTS OF THE INTERNSHIP HOST ORGANIZATION}
is reserved for the internship host organization.

It includes:

\begin{itemize}[leftmargin=0.8cm, itemsep=2pt]
    \item the student's full name;
    \item class;
    \item comments from the host organization;
    \item score;
    \item date;
    \item signature and official seal.
\end{itemize}

Students shall not alter or complete official evaluation fields unless
specifically authorized.

\vspace{0.3cm}

% ================================================================
% VII. SUPERVISOR'S EVALUATION AND DEFENSE RESULTS
% ================================================================

\noindent
\textbf{VII.\quad SUPERVISOR'S EVALUATION AND DEFENSE RESULTS}

\vspace{0.15cm}

The section entitled
\textbf{SUPERVISOR'S EVALUATION AND GRADUATION INTERNSHIP DEFENSE RESULTS}
is reserved for the supervisor and defense committee.

The evaluation may include:

\begin{itemize}[leftmargin=0.8cm, itemsep=2pt]
    \item Internship Progress and Student's Attitude;
    \item Report Content;
    \item Report Presentation;
    \item Other Comments;
    \item Supervisor's Opinion;
    \item Defense Result;
    \item Score.
\end{itemize}

Students shall not enter evaluation scores or defense results before they
are officially determined.

\vspace{0.3cm}

% ================================================================
% VIII. CHAPTERS, SECTIONS, AND SUBSECTIONS
% ================================================================

\noindent
\textbf{VIII.\quad CHAPTERS, SECTIONS, AND SUBSECTIONS}

\vspace{0.15cm}

The report structure shall be generated using the standard \LaTeX{}
commands:

\begin{verbatim}
\chapter{...}
\section{...}
\subsection{...}
\end{verbatim}

Chapter, section, and subsection numbers shall be generated automatically
by \LaTeX{}.

Students shall not manually type chapter or section numbers and shall not
reproduce heading styles using commands such as \verb|\textbf|,
\verb|\fontsize|, or \verb|\centering| instead of the appropriate
structural command.

Excessive subdivision should be avoided. In most cases, the hierarchy
Chapter--Section--Subsection is sufficient.

\vspace{0.3cm}

% ================================================================
% IX. ORGANIZATION OF THE MAIN CONTENT
% ================================================================

\noindent
\textbf{IX.\quad ORGANIZATION OF THE MAIN CONTENT}

\vspace{0.15cm}

The specific content of the report shall depend on the internship
organization and the assigned tasks. A typical structure may include:

\begin{itemize}[leftmargin=0.8cm, itemsep=2pt]
    \item \textbf{Chapter 1. Overview of the Internship Organization};
    \item \textbf{Chapter 2. Internship Activities and Assigned Tasks};
    \item \textbf{Chapter 3. Results, Experience, and Evaluation}.
\end{itemize}

These chapter titles are recommendations only and may be adapted to the
actual internship content.

The report shall reflect work actually performed, observed, analyzed, or
learned during the internship. Unrelated theoretical material should not
be included merely to increase report length.

\vspace{0.3cm}

% ================================================================
% X. TABLE OF CONTENTS
% ================================================================

\noindent
\textbf{X.\quad TABLE OF CONTENTS}

\vspace{0.15cm}

The Table of Contents shall be generated automatically using:

\begin{verbatim}
\tableofcontents
\end{verbatim}

Students shall not create or manually edit page numbers in the Table of
Contents.

The document shall be compiled sufficiently to ensure that all page
numbers and entries are correctly updated.

\vspace{0.3cm}

% ================================================================
% XI. LIST OF SYMBOLS AND ABBREVIATIONS
% ================================================================

\noindent
\textbf{XI.\quad LIST OF SYMBOLS AND ABBREVIATIONS}

\vspace{0.15cm}

Symbols and abbreviations that are used repeatedly in the report shall be
listed in the \textbf{List of Symbols and Abbreviations}.

The list may contain the following columns:

\begin{center}
\begin{tabular}{llll}
    No. & Symbol / Abbreviation & Unit & Definition
\end{tabular}
\end{center}

Only symbols and abbreviations actually used in the final report shall be
retained. Example entries contained in the template shall be removed if
they are not relevant to the student's report.

Each abbreviation should normally be defined when it first appears in
the main text.

\vspace{0.3cm}

% ================================================================
% XII. FIGURES, DIAGRAMS, AND PHOTOGRAPHS
% ================================================================

\noindent
\textbf{XII.\quad FIGURES, DIAGRAMS, AND PHOTOGRAPHS}

\vspace{0.15cm}

Figures shall be clear, legible, and directly relevant to the report
content.

Figures shall be inserted using the standard \verb|figure| environment.

Example:

\begin{verbatim}
\begin{figure}[H]
    \centering
    \includegraphics[width=0.8\textwidth]
                    {figures/company_structure.pdf}
    \caption{Organizational structure of the company}
    \label{fig:company-structure}
\end{figure}
\end{verbatim}

Figure captions shall be placed below the figure.

Each referenced figure shall have a unique \verb|\label| and shall be
referenced in the text using, for example:

\begin{verbatim}
Figure~\ref{fig:company-structure}
\end{verbatim}

Students shall not manually type figure numbers.

Figures reproduced or adapted from external sources shall include an
appropriate citation.

Photographs taken during the internship shall be relevant and sufficiently
clear. Students shall not include confidential, restricted, or proprietary
information without permission from the internship organization.

Vector graphics or high-resolution images are recommended whenever
possible.

\vspace{0.3cm}

% ================================================================
% XIII. TABLES
% ================================================================

\noindent
\textbf{XIII.\quad TABLES}

\vspace{0.15cm}

Tables shall be created using suitable \LaTeX{} environments such as
\verb|table|, \verb|tabular|, \verb|tabularx|, or \verb|longtable|.

Table captions shall be placed above the table.

Example:

\begin{verbatim}
\begin{table}[H]
    \centering
    \caption{Internship schedule}
    \label{tab:internship-schedule}

    \begin{tabular}{lll}
        \toprule
        Week & Task & Result \\
        \midrule
        1 & Orientation & Completed \\
        2 & System study & Completed \\
        \bottomrule
    \end{tabular}
\end{table}
\end{verbatim}

Tables shall be referenced using:

\begin{verbatim}
Table~\ref{tab:internship-schedule}
\end{verbatim}

Students shall not manually type table numbers.

Tables should be created directly in \LaTeX{} rather than inserted as
screenshots whenever practical.

\vspace{0.3cm}

% ================================================================
% XIV. EQUATIONS AND TECHNICAL QUANTITIES
% ================================================================

\noindent
\textbf{XIV.\quad EQUATIONS AND TECHNICAL QUANTITIES}

\vspace{0.15cm}

Mathematical expressions shall be typeset directly using \LaTeX{}.
Images or screenshots of equations shall not be used.

Short expressions may be written in inline mathematical mode, for example:

\begin{verbatim}
$R = 0.7\,\Omega$
\end{verbatim}

Important equations shall use a numbered equation environment.

Example:

\begin{verbatim}
\begin{equation}
    P = UI
    \label{eq:power}
\end{equation}
\end{verbatim}

Equations shall be referenced using:

\begin{verbatim}
Equation~\eqref{eq:power}
\end{verbatim}

Equation numbers shall not be manually entered.

SI units shall be used wherever applicable. The \verb|siunitx| package
should be used where appropriate.

Examples:

\begin{verbatim}
\SI{220}{\volt}
\SI{50}{\hertz}
\SI{2.5}{\ampere}
\end{verbatim}

Notation and units shall be consistent throughout the report.

\vspace{0.3cm}

% ================================================================
% XV. SOURCE CODE
% ================================================================

\noindent
\textbf{XV.\quad SOURCE CODE}

\vspace{0.15cm}

If the internship involves programming, PLCs, microcontrollers,
simulation, or software development, relevant source code may be included
in the report.

Source code shall be presented using the \verb|listings| environment or
another environment defined by the official template.

Example:

\begin{verbatim}
\begin{lstlisting}[language=st]
IF Start = TRUE THEN
    Motor := TRUE;
ELSE
    Motor := FALSE;
END_IF;
\end{lstlisting}
\end{verbatim}

Only code necessary for explaining the work should appear in the main
report. Long source-code listings should be moved to an appendix or
submitted as supplementary files if required.

Source code should not normally be inserted as screenshots.

\vspace{0.3cm}

% ================================================================
% XVI. ABSTRACT
% ================================================================

\noindent
\textbf{XVI.\quad ABSTRACT}

\vspace{0.15cm}

The Abstract shall briefly summarize:

\begin{itemize}[leftmargin=0.8cm, itemsep=2pt]
    \item the internship organization;
    \item the objectives of the internship;
    \item the main tasks performed;
    \item the main results, knowledge, and skills obtained.
\end{itemize}

The Abstract should be concise and should not reproduce the full content
of the report chapter by chapter.

\vspace{0.3cm}

% ================================================================
% XVII. REFERENCES AND CITATIONS
% ================================================================

\noindent
\textbf{XVII.\quad REFERENCES AND CITATIONS}

\vspace{0.15cm}

All books, journal papers, conference papers, standards, websites,
datasheets, manuals, technical documents, figures, tables, or other
external sources used in preparing the report shall be properly cited.

References shall be managed through the bibliography system provided with
the official template.

Bibliographic entries shall normally be stored in:

\begin{verbatim}
References.bib
\end{verbatim}

References shall be cited in the report using:

\begin{verbatim}
\cite{reference-key}
\end{verbatim}

Students shall not manually type reference numbers such as ``[1]'',
``[2]'', etc.

Titles of foreign-language publications shall remain in their original
language in the bibliography.

All borrowed ideas, data, figures, tables, equations, methods, or results
shall be properly acknowledged.

\vspace{0.3cm}

% ================================================================
% XVIII. CROSS-REFERENCING
% ================================================================

\noindent
\textbf{XVIII.\quad CROSS-REFERENCING}

\vspace{0.15cm}

Numbered elements shall be cross-referenced using the standard \LaTeX{}
mechanism.

Use the following commands where appropriate:

\begin{verbatim}
\label{...}
\ref{...}
\eqref{...}
\end{verbatim}

For example:

\begin{verbatim}
As shown in Figure~\ref{fig:system}, ...
\end{verbatim}

Students shall not manually type numbers such as ``Figure 2.3'' when
those numbers can be generated automatically.

\vspace{0.3cm}

% ================================================================
% XIX. PAGE NUMBERING
% ================================================================

\noindent
\textbf{XIX.\quad PAGE NUMBERING}

\vspace{0.15cm}

Cover pages and official evaluation forms at the beginning of the report
shall not display page numbers.

The main body of the report shall use Arabic page numbers beginning from
Chapter 1.

Page numbering shall be generated automatically by the official template.
Students shall not manually enter page numbers.

\vspace{0.3cm}

% ================================================================
% XX. PROJECT FILE ORGANIZATION
% ================================================================

\noindent
\textbf{XX.\quad PROJECT FILE ORGANIZATION}

\vspace{0.15cm}

Students should maintain a clear and organized \LaTeX{} project
structure. A recommended structure is:

\begin{verbatim}
InternshipReport/
|
|-- main.tex
|-- References.bib
|
|-- figures/
|   |-- company.pdf
|   |-- system.png
|   `-- ...
|
`-- chapters/
    |-- chapter1.tex
    |-- chapter2.tex
    `-- chapter3.tex
\end{verbatim}

Figures should normally be stored in the \verb|figures| directory
specified by the official template.

File and directory names should preferably contain Latin letters,
numbers, hyphens, or underscores.

Recommended examples include:

\begin{verbatim}
company_structure.pdf
motor_control.png
chapter_01.tex
\end{verbatim}

Spaces and unnecessary special characters in file names should be
avoided where possible.

\vspace{0.3cm}

% ================================================================
% XXI. DOCUMENT COMPILATION AND FINAL CHECK
% ================================================================

\noindent
\textbf{XXI.\quad DOCUMENT COMPILATION AND FINAL CHECK}

\vspace{0.15cm}

Before submission, the complete \LaTeX{} project shall compile
successfully and produce the final PDF report.

The final document shall not contain unresolved references such as:

\begin{verbatim}
??
[?]
\end{verbatim}

Before generating the final PDF, students shall check:

\begin{itemize}[leftmargin=0.8cm, itemsep=2pt]
    \item the Table of Contents;
    \item page numbering;
    \item chapter, section, and subsection numbering;
    \item figure numbering;
    \item table numbering;
    \item equation numbering;
    \item cross-references;
    \item citations;
    \item bibliography;
    \item font display;
    \item mathematical symbols;
    \item image quality;
    \item completeness of all required forms.
\end{itemize}

The document may be compiled using TeX Live, MiKTeX, or another
compatible \LaTeX{} environment. Regardless of the compilation
environment, the resulting PDF shall conform to the official Faculty
template.

\vspace{0.3cm}

% ================================================================
% XXII. CONTENT THAT STUDENTS MAY EDIT
% ================================================================

\noindent
\textbf{XXII.\quad CONTENT THAT STUDENTS MAY EDIT}

\vspace{0.15cm}

Students may edit or replace:

\begin{itemize}[leftmargin=0.8cm, itemsep=2pt]
    \item student name;
    \item class;
    \item internship organization;
    \item supervisor name;
    \item Abstract;
    \item List of Symbols and Abbreviations;
    \item chapter titles and chapter content;
    \item figures and photographs;
    \item tables;
    \item equations;
    \item citations and references.
\end{itemize}

Students shall not modify the main document formatting settings in the
preamble, including the document class, page geometry, line spacing,
heading formatting, table-of-contents formatting, font configuration,
page layout, or cover-page macros unless specifically instructed.

\vspace{0.3cm}

% ================================================================
% XXIII. FINAL SUBMISSION REQUIREMENTS
% ================================================================

\noindent
\textbf{XXIII.\quad FINAL SUBMISSION REQUIREMENTS}

\vspace{0.15cm}

Before submission, the final report shall:

\begin{itemize}[leftmargin=0.8cm, itemsep=2pt]
    \item compile without errors that affect the report content or presentation;
    \item contain all required forms and report sections;
    \item contain correct student, class, internship organization, and supervisor information;
    \item contain an updated Table of Contents, List of Figures, and List of Tables;
    \item contain no unused example content from the template;
    \item contain no placeholder text such as \textit{Title of Chapter 1},
    \textit{Section 1.1}, or other sample entries;
    \item contain all required citations and references;
    \item contain no unresolved cross-references.
\end{itemize}

If the Faculty requires submission of the \LaTeX{} source, the complete
project shall include all \verb|.tex| files, the bibliography database,
figures, and any other files required to reproduce the submitted PDF.

\vspace{0.4cm}

\noindent
\textbf{Important:}
Before submitting the final internship report, students shall remove or
comment out the following command in the main document:

\begin{verbatim}
\input{Regulations.tex}
\end{verbatim}

For example:

\begin{verbatim}
% \input{Regulations.tex}
\end{verbatim}

This removes the guidance pages from the final submitted report while
preserving the regulations file in the \LaTeX{} project for reference.

} % End 13pt content

\clearpage

%
% IMPORTANT:
% This section is provided for guidance only. Students shall remove
% or comment out the 
% =================================================================
% REGULATIONS FOR PREPARING THE GRADUATION INTERNSHIP REPORT
% USING LATEX
% =================================================================
%
% This file is intended to be included from the main document with:
%
%     \input{Regulations.tex}
%
% IMPORTANT:
% This section is provided for guidance only. Students shall remove
% or comment out the \input{Regulations.tex} command before submitting
% the final internship report.
%
% =================================================================

\clearpage
\thispagestyle{empty}

\begin{center}
    {\fontsize{16pt}{20pt}\selectfont\bfseries
    SPECIFIC REGULATIONS FOR PREPARING\\
    GRADUATION INTERNSHIP REPORTS USING \LaTeX}
\end{center}

\vspace{0.5cm}

{\fontsize{13pt}{17pt}\selectfont
\setstretch{1.3}
\justifying

\noindent
\textbf{Note:}
This section is provided for guidance only and is included in the main
document using the command \verb|\input{Regulations.tex}|.
Students shall remove or comment out this command before submitting the
final internship report.

\vspace{0.4cm}

% ================================================================
% I. GENERAL REQUIREMENTS
% ================================================================

\noindent
\textbf{I.\quad GENERAL REQUIREMENTS}

\vspace{0.15cm}

The Graduation Internship Report shall be prepared using the official
\LaTeX{} template provided by the Faculty. Students shall preserve the
predefined page layout, typography, line spacing, heading styles,
numbering system, figure and table formatting, and bibliography settings
unless otherwise instructed by the Faculty or the supervisor.

The report shall be written in English.

Students are responsible for the academic and technical content of the
report. The official \LaTeX{} template is responsible for standardized
document formatting.

\vspace{0.3cm}

% ================================================================
% II. REPORT STRUCTURE
% ================================================================

\noindent
\textbf{II.\quad REPORT STRUCTURE}

\vspace{0.15cm}

The report shall normally be arranged in the following order:

\begin{enumerate}[label=\arabic*., leftmargin=1.2cm, itemsep=2pt]
    \item Outer Cover;
    \item Comments of the Internship Host Organization;
    \item Supervisor's Evaluation and Graduation Internship Defense Results;
    \item Table of Contents;
    \item List of Symbols and Abbreviations, if applicable;
    \item List of Figures;
    \item List of Tables;
    \item Abstract;
    \item Main Chapters of the Internship Report;
    \item References;
    \item Appendices, if applicable.
\end{enumerate}

The preliminary pages shall not be numbered as chapters. The main report
content shall be organized into consecutively numbered chapters beginning
with Chapter 1.

\vspace{0.3cm}

% ================================================================
% III. PAGE SIZE AND MARGINS
% ================================================================

\noindent
\textbf{III.\quad PAGE SIZE AND MARGINS}

\vspace{0.15cm}

The report shall be formatted on A4 paper with the following margins:

\begin{center}
\begin{tabular}{ll}
    Left margin:   & 3.5 cm \\
    Right margin:  & 2.0 cm \\
    Top margin:    & 2.5 cm \\
    Bottom margin: & 2.5 cm
\end{tabular}
\end{center}

These settings are predefined in the official template and shall not be
modified by students.

\vspace{0.3cm}

% ================================================================
% IV. FONT, FONT SIZE, AND LINE SPACING
% ================================================================

\noindent
\textbf{IV.\quad FONT, FONT SIZE, AND LINE SPACING}

\vspace{0.15cm}

The main text shall use a 13 pt font size with 1.3-line spacing.
A Times-family typeface shall be used in accordance with the official
template.

The text shall be fully justified.

Students shall not use repeated spaces, repeated manual line breaks, or
multiple blank lines to position text. Paragraph and heading spacing shall
be controlled by the template.

\vspace{0.3cm}

% ================================================================
% V. OUTER COVER
% ================================================================

\noindent
\textbf{V.\quad OUTER COVER}

\vspace{0.15cm}

The outer cover shall follow the official format provided in the template
and shall contain the following information:

\begin{itemize}[leftmargin=0.8cm, itemsep=2pt]
    \item Thai Nguyen University of Technology;
    \item Faculty of International Training;
    \item Engineering and Technology Department;
    \item University logo and Faculty logo;
    \item the title \textbf{INTERNSHIP REPORT};
    \item Internship Organization;
    \item Author;
    \item Class;
    \item Supervisor;
    \item place and year of completion.
\end{itemize}

Students shall only replace the required information fields and shall not
modify the cover-page layout unless instructed otherwise.

\vspace{0.3cm}

% ================================================================
% VI. COMMENTS OF THE INTERNSHIP HOST ORGANIZATION
% ================================================================

\noindent
\textbf{VI.\quad COMMENTS OF THE INTERNSHIP HOST ORGANIZATION}

\vspace{0.15cm}

The page entitled \textbf{COMMENTS OF THE INTERNSHIP HOST ORGANIZATION}
is reserved for the internship host organization.

It includes:

\begin{itemize}[leftmargin=0.8cm, itemsep=2pt]
    \item the student's full name;
    \item class;
    \item comments from the host organization;
    \item score;
    \item date;
    \item signature and official seal.
\end{itemize}

Students shall not alter or complete official evaluation fields unless
specifically authorized.

\vspace{0.3cm}

% ================================================================
% VII. SUPERVISOR'S EVALUATION AND DEFENSE RESULTS
% ================================================================

\noindent
\textbf{VII.\quad SUPERVISOR'S EVALUATION AND DEFENSE RESULTS}

\vspace{0.15cm}

The section entitled
\textbf{SUPERVISOR'S EVALUATION AND GRADUATION INTERNSHIP DEFENSE RESULTS}
is reserved for the supervisor and defense committee.

The evaluation may include:

\begin{itemize}[leftmargin=0.8cm, itemsep=2pt]
    \item Internship Progress and Student's Attitude;
    \item Report Content;
    \item Report Presentation;
    \item Other Comments;
    \item Supervisor's Opinion;
    \item Defense Result;
    \item Score.
\end{itemize}

Students shall not enter evaluation scores or defense results before they
are officially determined.

\vspace{0.3cm}

% ================================================================
% VIII. CHAPTERS, SECTIONS, AND SUBSECTIONS
% ================================================================

\noindent
\textbf{VIII.\quad CHAPTERS, SECTIONS, AND SUBSECTIONS}

\vspace{0.15cm}

The report structure shall be generated using the standard \LaTeX{}
commands:

\begin{verbatim}
\chapter{...}
\section{...}
\subsection{...}
\end{verbatim}

Chapter, section, and subsection numbers shall be generated automatically
by \LaTeX{}.

Students shall not manually type chapter or section numbers and shall not
reproduce heading styles using commands such as \verb|\textbf|,
\verb|\fontsize|, or \verb|\centering| instead of the appropriate
structural command.

Excessive subdivision should be avoided. In most cases, the hierarchy
Chapter--Section--Subsection is sufficient.

\vspace{0.3cm}

% ================================================================
% IX. ORGANIZATION OF THE MAIN CONTENT
% ================================================================

\noindent
\textbf{IX.\quad ORGANIZATION OF THE MAIN CONTENT}

\vspace{0.15cm}

The specific content of the report shall depend on the internship
organization and the assigned tasks. A typical structure may include:

\begin{itemize}[leftmargin=0.8cm, itemsep=2pt]
    \item \textbf{Chapter 1. Overview of the Internship Organization};
    \item \textbf{Chapter 2. Internship Activities and Assigned Tasks};
    \item \textbf{Chapter 3. Results, Experience, and Evaluation}.
\end{itemize}

These chapter titles are recommendations only and may be adapted to the
actual internship content.

The report shall reflect work actually performed, observed, analyzed, or
learned during the internship. Unrelated theoretical material should not
be included merely to increase report length.

\vspace{0.3cm}

% ================================================================
% X. TABLE OF CONTENTS
% ================================================================

\noindent
\textbf{X.\quad TABLE OF CONTENTS}

\vspace{0.15cm}

The Table of Contents shall be generated automatically using:

\begin{verbatim}
\tableofcontents
\end{verbatim}

Students shall not create or manually edit page numbers in the Table of
Contents.

The document shall be compiled sufficiently to ensure that all page
numbers and entries are correctly updated.

\vspace{0.3cm}

% ================================================================
% XI. LIST OF SYMBOLS AND ABBREVIATIONS
% ================================================================

\noindent
\textbf{XI.\quad LIST OF SYMBOLS AND ABBREVIATIONS}

\vspace{0.15cm}

Symbols and abbreviations that are used repeatedly in the report shall be
listed in the \textbf{List of Symbols and Abbreviations}.

The list may contain the following columns:

\begin{center}
\begin{tabular}{llll}
    No. & Symbol / Abbreviation & Unit & Definition
\end{tabular}
\end{center}

Only symbols and abbreviations actually used in the final report shall be
retained. Example entries contained in the template shall be removed if
they are not relevant to the student's report.

Each abbreviation should normally be defined when it first appears in
the main text.

\vspace{0.3cm}

% ================================================================
% XII. FIGURES, DIAGRAMS, AND PHOTOGRAPHS
% ================================================================

\noindent
\textbf{XII.\quad FIGURES, DIAGRAMS, AND PHOTOGRAPHS}

\vspace{0.15cm}

Figures shall be clear, legible, and directly relevant to the report
content.

Figures shall be inserted using the standard \verb|figure| environment.

Example:

\begin{verbatim}
\begin{figure}[H]
    \centering
    \includegraphics[width=0.8\textwidth]
                    {figures/company_structure.pdf}
    \caption{Organizational structure of the company}
    \label{fig:company-structure}
\end{figure}
\end{verbatim}

Figure captions shall be placed below the figure.

Each referenced figure shall have a unique \verb|\label| and shall be
referenced in the text using, for example:

\begin{verbatim}
Figure~\ref{fig:company-structure}
\end{verbatim}

Students shall not manually type figure numbers.

Figures reproduced or adapted from external sources shall include an
appropriate citation.

Photographs taken during the internship shall be relevant and sufficiently
clear. Students shall not include confidential, restricted, or proprietary
information without permission from the internship organization.

Vector graphics or high-resolution images are recommended whenever
possible.

\vspace{0.3cm}

% ================================================================
% XIII. TABLES
% ================================================================

\noindent
\textbf{XIII.\quad TABLES}

\vspace{0.15cm}

Tables shall be created using suitable \LaTeX{} environments such as
\verb|table|, \verb|tabular|, \verb|tabularx|, or \verb|longtable|.

Table captions shall be placed above the table.

Example:

\begin{verbatim}
\begin{table}[H]
    \centering
    \caption{Internship schedule}
    \label{tab:internship-schedule}

    \begin{tabular}{lll}
        \toprule
        Week & Task & Result \\
        \midrule
        1 & Orientation & Completed \\
        2 & System study & Completed \\
        \bottomrule
    \end{tabular}
\end{table}
\end{verbatim}

Tables shall be referenced using:

\begin{verbatim}
Table~\ref{tab:internship-schedule}
\end{verbatim}

Students shall not manually type table numbers.

Tables should be created directly in \LaTeX{} rather than inserted as
screenshots whenever practical.

\vspace{0.3cm}

% ================================================================
% XIV. EQUATIONS AND TECHNICAL QUANTITIES
% ================================================================

\noindent
\textbf{XIV.\quad EQUATIONS AND TECHNICAL QUANTITIES}

\vspace{0.15cm}

Mathematical expressions shall be typeset directly using \LaTeX{}.
Images or screenshots of equations shall not be used.

Short expressions may be written in inline mathematical mode, for example:

\begin{verbatim}
$R = 0.7\,\Omega$
\end{verbatim}

Important equations shall use a numbered equation environment.

Example:

\begin{verbatim}
\begin{equation}
    P = UI
    \label{eq:power}
\end{equation}
\end{verbatim}

Equations shall be referenced using:

\begin{verbatim}
Equation~\eqref{eq:power}
\end{verbatim}

Equation numbers shall not be manually entered.

SI units shall be used wherever applicable. The \verb|siunitx| package
should be used where appropriate.

Examples:

\begin{verbatim}
\SI{220}{\volt}
\SI{50}{\hertz}
\SI{2.5}{\ampere}
\end{verbatim}

Notation and units shall be consistent throughout the report.

\vspace{0.3cm}

% ================================================================
% XV. SOURCE CODE
% ================================================================

\noindent
\textbf{XV.\quad SOURCE CODE}

\vspace{0.15cm}

If the internship involves programming, PLCs, microcontrollers,
simulation, or software development, relevant source code may be included
in the report.

Source code shall be presented using the \verb|listings| environment or
another environment defined by the official template.

Example:

\begin{verbatim}
\begin{lstlisting}[language=st]
IF Start = TRUE THEN
    Motor := TRUE;
ELSE
    Motor := FALSE;
END_IF;
\end{lstlisting}
\end{verbatim}

Only code necessary for explaining the work should appear in the main
report. Long source-code listings should be moved to an appendix or
submitted as supplementary files if required.

Source code should not normally be inserted as screenshots.

\vspace{0.3cm}

% ================================================================
% XVI. ABSTRACT
% ================================================================

\noindent
\textbf{XVI.\quad ABSTRACT}

\vspace{0.15cm}

The Abstract shall briefly summarize:

\begin{itemize}[leftmargin=0.8cm, itemsep=2pt]
    \item the internship organization;
    \item the objectives of the internship;
    \item the main tasks performed;
    \item the main results, knowledge, and skills obtained.
\end{itemize}

The Abstract should be concise and should not reproduce the full content
of the report chapter by chapter.

\vspace{0.3cm}

% ================================================================
% XVII. REFERENCES AND CITATIONS
% ================================================================

\noindent
\textbf{XVII.\quad REFERENCES AND CITATIONS}

\vspace{0.15cm}

All books, journal papers, conference papers, standards, websites,
datasheets, manuals, technical documents, figures, tables, or other
external sources used in preparing the report shall be properly cited.

References shall be managed through the bibliography system provided with
the official template.

Bibliographic entries shall normally be stored in:

\begin{verbatim}
References.bib
\end{verbatim}

References shall be cited in the report using:

\begin{verbatim}
\cite{reference-key}
\end{verbatim}

Students shall not manually type reference numbers such as ``[1]'',
``[2]'', etc.

Titles of foreign-language publications shall remain in their original
language in the bibliography.

All borrowed ideas, data, figures, tables, equations, methods, or results
shall be properly acknowledged.

\vspace{0.3cm}

% ================================================================
% XVIII. CROSS-REFERENCING
% ================================================================

\noindent
\textbf{XVIII.\quad CROSS-REFERENCING}

\vspace{0.15cm}

Numbered elements shall be cross-referenced using the standard \LaTeX{}
mechanism.

Use the following commands where appropriate:

\begin{verbatim}
\label{...}
\ref{...}
\eqref{...}
\end{verbatim}

For example:

\begin{verbatim}
As shown in Figure~\ref{fig:system}, ...
\end{verbatim}

Students shall not manually type numbers such as ``Figure 2.3'' when
those numbers can be generated automatically.

\vspace{0.3cm}

% ================================================================
% XIX. PAGE NUMBERING
% ================================================================

\noindent
\textbf{XIX.\quad PAGE NUMBERING}

\vspace{0.15cm}

Cover pages and official evaluation forms at the beginning of the report
shall not display page numbers.

The main body of the report shall use Arabic page numbers beginning from
Chapter 1.

Page numbering shall be generated automatically by the official template.
Students shall not manually enter page numbers.

\vspace{0.3cm}

% ================================================================
% XX. PROJECT FILE ORGANIZATION
% ================================================================

\noindent
\textbf{XX.\quad PROJECT FILE ORGANIZATION}

\vspace{0.15cm}

Students should maintain a clear and organized \LaTeX{} project
structure. A recommended structure is:

\begin{verbatim}
InternshipReport/
|
|-- main.tex
|-- References.bib
|
|-- figures/
|   |-- company.pdf
|   |-- system.png
|   `-- ...
|
`-- chapters/
    |-- chapter1.tex
    |-- chapter2.tex
    `-- chapter3.tex
\end{verbatim}

Figures should normally be stored in the \verb|figures| directory
specified by the official template.

File and directory names should preferably contain Latin letters,
numbers, hyphens, or underscores.

Recommended examples include:

\begin{verbatim}
company_structure.pdf
motor_control.png
chapter_01.tex
\end{verbatim}

Spaces and unnecessary special characters in file names should be
avoided where possible.

\vspace{0.3cm}

% ================================================================
% XXI. DOCUMENT COMPILATION AND FINAL CHECK
% ================================================================

\noindent
\textbf{XXI.\quad DOCUMENT COMPILATION AND FINAL CHECK}

\vspace{0.15cm}

Before submission, the complete \LaTeX{} project shall compile
successfully and produce the final PDF report.

The final document shall not contain unresolved references such as:

\begin{verbatim}
??
[?]
\end{verbatim}

Before generating the final PDF, students shall check:

\begin{itemize}[leftmargin=0.8cm, itemsep=2pt]
    \item the Table of Contents;
    \item page numbering;
    \item chapter, section, and subsection numbering;
    \item figure numbering;
    \item table numbering;
    \item equation numbering;
    \item cross-references;
    \item citations;
    \item bibliography;
    \item font display;
    \item mathematical symbols;
    \item image quality;
    \item completeness of all required forms.
\end{itemize}

The document may be compiled using TeX Live, MiKTeX, or another
compatible \LaTeX{} environment. Regardless of the compilation
environment, the resulting PDF shall conform to the official Faculty
template.

\vspace{0.3cm}

% ================================================================
% XXII. CONTENT THAT STUDENTS MAY EDIT
% ================================================================

\noindent
\textbf{XXII.\quad CONTENT THAT STUDENTS MAY EDIT}

\vspace{0.15cm}

Students may edit or replace:

\begin{itemize}[leftmargin=0.8cm, itemsep=2pt]
    \item student name;
    \item class;
    \item internship organization;
    \item supervisor name;
    \item Abstract;
    \item List of Symbols and Abbreviations;
    \item chapter titles and chapter content;
    \item figures and photographs;
    \item tables;
    \item equations;
    \item citations and references.
\end{itemize}

Students shall not modify the main document formatting settings in the
preamble, including the document class, page geometry, line spacing,
heading formatting, table-of-contents formatting, font configuration,
page layout, or cover-page macros unless specifically instructed.

\vspace{0.3cm}

% ================================================================
% XXIII. FINAL SUBMISSION REQUIREMENTS
% ================================================================

\noindent
\textbf{XXIII.\quad FINAL SUBMISSION REQUIREMENTS}

\vspace{0.15cm}

Before submission, the final report shall:

\begin{itemize}[leftmargin=0.8cm, itemsep=2pt]
    \item compile without errors that affect the report content or presentation;
    \item contain all required forms and report sections;
    \item contain correct student, class, internship organization, and supervisor information;
    \item contain an updated Table of Contents, List of Figures, and List of Tables;
    \item contain no unused example content from the template;
    \item contain no placeholder text such as \textit{Title of Chapter 1},
    \textit{Section 1.1}, or other sample entries;
    \item contain all required citations and references;
    \item contain no unresolved cross-references.
\end{itemize}

If the Faculty requires submission of the \LaTeX{} source, the complete
project shall include all \verb|.tex| files, the bibliography database,
figures, and any other files required to reproduce the submitted PDF.

\vspace{0.4cm}

\noindent
\textbf{Important:}
Before submitting the final internship report, students shall remove or
comment out the following command in the main document:

\begin{verbatim}
\input{Regulations.tex}
\end{verbatim}

For example:

\begin{verbatim}
% \input{Regulations.tex}
\end{verbatim}

This removes the guidance pages from the final submitted report while
preserving the regulations file in the \LaTeX{} project for reference.

} % End 13pt content

\clearpage
 command before submitting
% the final internship report.
%
% =================================================================

\clearpage
\thispagestyle{empty}

\begin{center}
    {\fontsize{16pt}{20pt}\selectfont\bfseries
    SPECIFIC REGULATIONS FOR PREPARING\\
    GRADUATION INTERNSHIP REPORTS USING \LaTeX}
\end{center}

\vspace{0.5cm}

{\fontsize{13pt}{17pt}\selectfont
\setstretch{1.3}
\justifying

\noindent
\textbf{Note:}
This section is provided for guidance only and is included in the main
document using the command \verb|
% =================================================================
% REGULATIONS FOR PREPARING THE GRADUATION INTERNSHIP REPORT
% USING LATEX
% =================================================================
%
% This file is intended to be included from the main document with:
%
%     \input{Regulations.tex}
%
% IMPORTANT:
% This section is provided for guidance only. Students shall remove
% or comment out the \input{Regulations.tex} command before submitting
% the final internship report.
%
% =================================================================

\clearpage
\thispagestyle{empty}

\begin{center}
    {\fontsize{16pt}{20pt}\selectfont\bfseries
    SPECIFIC REGULATIONS FOR PREPARING\\
    GRADUATION INTERNSHIP REPORTS USING \LaTeX}
\end{center}

\vspace{0.5cm}

{\fontsize{13pt}{17pt}\selectfont
\setstretch{1.3}
\justifying

\noindent
\textbf{Note:}
This section is provided for guidance only and is included in the main
document using the command \verb|\input{Regulations.tex}|.
Students shall remove or comment out this command before submitting the
final internship report.

\vspace{0.4cm}

% ================================================================
% I. GENERAL REQUIREMENTS
% ================================================================

\noindent
\textbf{I.\quad GENERAL REQUIREMENTS}

\vspace{0.15cm}

The Graduation Internship Report shall be prepared using the official
\LaTeX{} template provided by the Faculty. Students shall preserve the
predefined page layout, typography, line spacing, heading styles,
numbering system, figure and table formatting, and bibliography settings
unless otherwise instructed by the Faculty or the supervisor.

The report shall be written in English.

Students are responsible for the academic and technical content of the
report. The official \LaTeX{} template is responsible for standardized
document formatting.

\vspace{0.3cm}

% ================================================================
% II. REPORT STRUCTURE
% ================================================================

\noindent
\textbf{II.\quad REPORT STRUCTURE}

\vspace{0.15cm}

The report shall normally be arranged in the following order:

\begin{enumerate}[label=\arabic*., leftmargin=1.2cm, itemsep=2pt]
    \item Outer Cover;
    \item Comments of the Internship Host Organization;
    \item Supervisor's Evaluation and Graduation Internship Defense Results;
    \item Table of Contents;
    \item List of Symbols and Abbreviations, if applicable;
    \item List of Figures;
    \item List of Tables;
    \item Abstract;
    \item Main Chapters of the Internship Report;
    \item References;
    \item Appendices, if applicable.
\end{enumerate}

The preliminary pages shall not be numbered as chapters. The main report
content shall be organized into consecutively numbered chapters beginning
with Chapter 1.

\vspace{0.3cm}

% ================================================================
% III. PAGE SIZE AND MARGINS
% ================================================================

\noindent
\textbf{III.\quad PAGE SIZE AND MARGINS}

\vspace{0.15cm}

The report shall be formatted on A4 paper with the following margins:

\begin{center}
\begin{tabular}{ll}
    Left margin:   & 3.5 cm \\
    Right margin:  & 2.0 cm \\
    Top margin:    & 2.5 cm \\
    Bottom margin: & 2.5 cm
\end{tabular}
\end{center}

These settings are predefined in the official template and shall not be
modified by students.

\vspace{0.3cm}

% ================================================================
% IV. FONT, FONT SIZE, AND LINE SPACING
% ================================================================

\noindent
\textbf{IV.\quad FONT, FONT SIZE, AND LINE SPACING}

\vspace{0.15cm}

The main text shall use a 13 pt font size with 1.3-line spacing.
A Times-family typeface shall be used in accordance with the official
template.

The text shall be fully justified.

Students shall not use repeated spaces, repeated manual line breaks, or
multiple blank lines to position text. Paragraph and heading spacing shall
be controlled by the template.

\vspace{0.3cm}

% ================================================================
% V. OUTER COVER
% ================================================================

\noindent
\textbf{V.\quad OUTER COVER}

\vspace{0.15cm}

The outer cover shall follow the official format provided in the template
and shall contain the following information:

\begin{itemize}[leftmargin=0.8cm, itemsep=2pt]
    \item Thai Nguyen University of Technology;
    \item Faculty of International Training;
    \item Engineering and Technology Department;
    \item University logo and Faculty logo;
    \item the title \textbf{INTERNSHIP REPORT};
    \item Internship Organization;
    \item Author;
    \item Class;
    \item Supervisor;
    \item place and year of completion.
\end{itemize}

Students shall only replace the required information fields and shall not
modify the cover-page layout unless instructed otherwise.

\vspace{0.3cm}

% ================================================================
% VI. COMMENTS OF THE INTERNSHIP HOST ORGANIZATION
% ================================================================

\noindent
\textbf{VI.\quad COMMENTS OF THE INTERNSHIP HOST ORGANIZATION}

\vspace{0.15cm}

The page entitled \textbf{COMMENTS OF THE INTERNSHIP HOST ORGANIZATION}
is reserved for the internship host organization.

It includes:

\begin{itemize}[leftmargin=0.8cm, itemsep=2pt]
    \item the student's full name;
    \item class;
    \item comments from the host organization;
    \item score;
    \item date;
    \item signature and official seal.
\end{itemize}

Students shall not alter or complete official evaluation fields unless
specifically authorized.

\vspace{0.3cm}

% ================================================================
% VII. SUPERVISOR'S EVALUATION AND DEFENSE RESULTS
% ================================================================

\noindent
\textbf{VII.\quad SUPERVISOR'S EVALUATION AND DEFENSE RESULTS}

\vspace{0.15cm}

The section entitled
\textbf{SUPERVISOR'S EVALUATION AND GRADUATION INTERNSHIP DEFENSE RESULTS}
is reserved for the supervisor and defense committee.

The evaluation may include:

\begin{itemize}[leftmargin=0.8cm, itemsep=2pt]
    \item Internship Progress and Student's Attitude;
    \item Report Content;
    \item Report Presentation;
    \item Other Comments;
    \item Supervisor's Opinion;
    \item Defense Result;
    \item Score.
\end{itemize}

Students shall not enter evaluation scores or defense results before they
are officially determined.

\vspace{0.3cm}

% ================================================================
% VIII. CHAPTERS, SECTIONS, AND SUBSECTIONS
% ================================================================

\noindent
\textbf{VIII.\quad CHAPTERS, SECTIONS, AND SUBSECTIONS}

\vspace{0.15cm}

The report structure shall be generated using the standard \LaTeX{}
commands:

\begin{verbatim}
\chapter{...}
\section{...}
\subsection{...}
\end{verbatim}

Chapter, section, and subsection numbers shall be generated automatically
by \LaTeX{}.

Students shall not manually type chapter or section numbers and shall not
reproduce heading styles using commands such as \verb|\textbf|,
\verb|\fontsize|, or \verb|\centering| instead of the appropriate
structural command.

Excessive subdivision should be avoided. In most cases, the hierarchy
Chapter--Section--Subsection is sufficient.

\vspace{0.3cm}

% ================================================================
% IX. ORGANIZATION OF THE MAIN CONTENT
% ================================================================

\noindent
\textbf{IX.\quad ORGANIZATION OF THE MAIN CONTENT}

\vspace{0.15cm}

The specific content of the report shall depend on the internship
organization and the assigned tasks. A typical structure may include:

\begin{itemize}[leftmargin=0.8cm, itemsep=2pt]
    \item \textbf{Chapter 1. Overview of the Internship Organization};
    \item \textbf{Chapter 2. Internship Activities and Assigned Tasks};
    \item \textbf{Chapter 3. Results, Experience, and Evaluation}.
\end{itemize}

These chapter titles are recommendations only and may be adapted to the
actual internship content.

The report shall reflect work actually performed, observed, analyzed, or
learned during the internship. Unrelated theoretical material should not
be included merely to increase report length.

\vspace{0.3cm}

% ================================================================
% X. TABLE OF CONTENTS
% ================================================================

\noindent
\textbf{X.\quad TABLE OF CONTENTS}

\vspace{0.15cm}

The Table of Contents shall be generated automatically using:

\begin{verbatim}
\tableofcontents
\end{verbatim}

Students shall not create or manually edit page numbers in the Table of
Contents.

The document shall be compiled sufficiently to ensure that all page
numbers and entries are correctly updated.

\vspace{0.3cm}

% ================================================================
% XI. LIST OF SYMBOLS AND ABBREVIATIONS
% ================================================================

\noindent
\textbf{XI.\quad LIST OF SYMBOLS AND ABBREVIATIONS}

\vspace{0.15cm}

Symbols and abbreviations that are used repeatedly in the report shall be
listed in the \textbf{List of Symbols and Abbreviations}.

The list may contain the following columns:

\begin{center}
\begin{tabular}{llll}
    No. & Symbol / Abbreviation & Unit & Definition
\end{tabular}
\end{center}

Only symbols and abbreviations actually used in the final report shall be
retained. Example entries contained in the template shall be removed if
they are not relevant to the student's report.

Each abbreviation should normally be defined when it first appears in
the main text.

\vspace{0.3cm}

% ================================================================
% XII. FIGURES, DIAGRAMS, AND PHOTOGRAPHS
% ================================================================

\noindent
\textbf{XII.\quad FIGURES, DIAGRAMS, AND PHOTOGRAPHS}

\vspace{0.15cm}

Figures shall be clear, legible, and directly relevant to the report
content.

Figures shall be inserted using the standard \verb|figure| environment.

Example:

\begin{verbatim}
\begin{figure}[H]
    \centering
    \includegraphics[width=0.8\textwidth]
                    {figures/company_structure.pdf}
    \caption{Organizational structure of the company}
    \label{fig:company-structure}
\end{figure}
\end{verbatim}

Figure captions shall be placed below the figure.

Each referenced figure shall have a unique \verb|\label| and shall be
referenced in the text using, for example:

\begin{verbatim}
Figure~\ref{fig:company-structure}
\end{verbatim}

Students shall not manually type figure numbers.

Figures reproduced or adapted from external sources shall include an
appropriate citation.

Photographs taken during the internship shall be relevant and sufficiently
clear. Students shall not include confidential, restricted, or proprietary
information without permission from the internship organization.

Vector graphics or high-resolution images are recommended whenever
possible.

\vspace{0.3cm}

% ================================================================
% XIII. TABLES
% ================================================================

\noindent
\textbf{XIII.\quad TABLES}

\vspace{0.15cm}

Tables shall be created using suitable \LaTeX{} environments such as
\verb|table|, \verb|tabular|, \verb|tabularx|, or \verb|longtable|.

Table captions shall be placed above the table.

Example:

\begin{verbatim}
\begin{table}[H]
    \centering
    \caption{Internship schedule}
    \label{tab:internship-schedule}

    \begin{tabular}{lll}
        \toprule
        Week & Task & Result \\
        \midrule
        1 & Orientation & Completed \\
        2 & System study & Completed \\
        \bottomrule
    \end{tabular}
\end{table}
\end{verbatim}

Tables shall be referenced using:

\begin{verbatim}
Table~\ref{tab:internship-schedule}
\end{verbatim}

Students shall not manually type table numbers.

Tables should be created directly in \LaTeX{} rather than inserted as
screenshots whenever practical.

\vspace{0.3cm}

% ================================================================
% XIV. EQUATIONS AND TECHNICAL QUANTITIES
% ================================================================

\noindent
\textbf{XIV.\quad EQUATIONS AND TECHNICAL QUANTITIES}

\vspace{0.15cm}

Mathematical expressions shall be typeset directly using \LaTeX{}.
Images or screenshots of equations shall not be used.

Short expressions may be written in inline mathematical mode, for example:

\begin{verbatim}
$R = 0.7\,\Omega$
\end{verbatim}

Important equations shall use a numbered equation environment.

Example:

\begin{verbatim}
\begin{equation}
    P = UI
    \label{eq:power}
\end{equation}
\end{verbatim}

Equations shall be referenced using:

\begin{verbatim}
Equation~\eqref{eq:power}
\end{verbatim}

Equation numbers shall not be manually entered.

SI units shall be used wherever applicable. The \verb|siunitx| package
should be used where appropriate.

Examples:

\begin{verbatim}
\SI{220}{\volt}
\SI{50}{\hertz}
\SI{2.5}{\ampere}
\end{verbatim}

Notation and units shall be consistent throughout the report.

\vspace{0.3cm}

% ================================================================
% XV. SOURCE CODE
% ================================================================

\noindent
\textbf{XV.\quad SOURCE CODE}

\vspace{0.15cm}

If the internship involves programming, PLCs, microcontrollers,
simulation, or software development, relevant source code may be included
in the report.

Source code shall be presented using the \verb|listings| environment or
another environment defined by the official template.

Example:

\begin{verbatim}
\begin{lstlisting}[language=st]
IF Start = TRUE THEN
    Motor := TRUE;
ELSE
    Motor := FALSE;
END_IF;
\end{lstlisting}
\end{verbatim}

Only code necessary for explaining the work should appear in the main
report. Long source-code listings should be moved to an appendix or
submitted as supplementary files if required.

Source code should not normally be inserted as screenshots.

\vspace{0.3cm}

% ================================================================
% XVI. ABSTRACT
% ================================================================

\noindent
\textbf{XVI.\quad ABSTRACT}

\vspace{0.15cm}

The Abstract shall briefly summarize:

\begin{itemize}[leftmargin=0.8cm, itemsep=2pt]
    \item the internship organization;
    \item the objectives of the internship;
    \item the main tasks performed;
    \item the main results, knowledge, and skills obtained.
\end{itemize}

The Abstract should be concise and should not reproduce the full content
of the report chapter by chapter.

\vspace{0.3cm}

% ================================================================
% XVII. REFERENCES AND CITATIONS
% ================================================================

\noindent
\textbf{XVII.\quad REFERENCES AND CITATIONS}

\vspace{0.15cm}

All books, journal papers, conference papers, standards, websites,
datasheets, manuals, technical documents, figures, tables, or other
external sources used in preparing the report shall be properly cited.

References shall be managed through the bibliography system provided with
the official template.

Bibliographic entries shall normally be stored in:

\begin{verbatim}
References.bib
\end{verbatim}

References shall be cited in the report using:

\begin{verbatim}
\cite{reference-key}
\end{verbatim}

Students shall not manually type reference numbers such as ``[1]'',
``[2]'', etc.

Titles of foreign-language publications shall remain in their original
language in the bibliography.

All borrowed ideas, data, figures, tables, equations, methods, or results
shall be properly acknowledged.

\vspace{0.3cm}

% ================================================================
% XVIII. CROSS-REFERENCING
% ================================================================

\noindent
\textbf{XVIII.\quad CROSS-REFERENCING}

\vspace{0.15cm}

Numbered elements shall be cross-referenced using the standard \LaTeX{}
mechanism.

Use the following commands where appropriate:

\begin{verbatim}
\label{...}
\ref{...}
\eqref{...}
\end{verbatim}

For example:

\begin{verbatim}
As shown in Figure~\ref{fig:system}, ...
\end{verbatim}

Students shall not manually type numbers such as ``Figure 2.3'' when
those numbers can be generated automatically.

\vspace{0.3cm}

% ================================================================
% XIX. PAGE NUMBERING
% ================================================================

\noindent
\textbf{XIX.\quad PAGE NUMBERING}

\vspace{0.15cm}

Cover pages and official evaluation forms at the beginning of the report
shall not display page numbers.

The main body of the report shall use Arabic page numbers beginning from
Chapter 1.

Page numbering shall be generated automatically by the official template.
Students shall not manually enter page numbers.

\vspace{0.3cm}

% ================================================================
% XX. PROJECT FILE ORGANIZATION
% ================================================================

\noindent
\textbf{XX.\quad PROJECT FILE ORGANIZATION}

\vspace{0.15cm}

Students should maintain a clear and organized \LaTeX{} project
structure. A recommended structure is:

\begin{verbatim}
InternshipReport/
|
|-- main.tex
|-- References.bib
|
|-- figures/
|   |-- company.pdf
|   |-- system.png
|   `-- ...
|
`-- chapters/
    |-- chapter1.tex
    |-- chapter2.tex
    `-- chapter3.tex
\end{verbatim}

Figures should normally be stored in the \verb|figures| directory
specified by the official template.

File and directory names should preferably contain Latin letters,
numbers, hyphens, or underscores.

Recommended examples include:

\begin{verbatim}
company_structure.pdf
motor_control.png
chapter_01.tex
\end{verbatim}

Spaces and unnecessary special characters in file names should be
avoided where possible.

\vspace{0.3cm}

% ================================================================
% XXI. DOCUMENT COMPILATION AND FINAL CHECK
% ================================================================

\noindent
\textbf{XXI.\quad DOCUMENT COMPILATION AND FINAL CHECK}

\vspace{0.15cm}

Before submission, the complete \LaTeX{} project shall compile
successfully and produce the final PDF report.

The final document shall not contain unresolved references such as:

\begin{verbatim}
??
[?]
\end{verbatim}

Before generating the final PDF, students shall check:

\begin{itemize}[leftmargin=0.8cm, itemsep=2pt]
    \item the Table of Contents;
    \item page numbering;
    \item chapter, section, and subsection numbering;
    \item figure numbering;
    \item table numbering;
    \item equation numbering;
    \item cross-references;
    \item citations;
    \item bibliography;
    \item font display;
    \item mathematical symbols;
    \item image quality;
    \item completeness of all required forms.
\end{itemize}

The document may be compiled using TeX Live, MiKTeX, or another
compatible \LaTeX{} environment. Regardless of the compilation
environment, the resulting PDF shall conform to the official Faculty
template.

\vspace{0.3cm}

% ================================================================
% XXII. CONTENT THAT STUDENTS MAY EDIT
% ================================================================

\noindent
\textbf{XXII.\quad CONTENT THAT STUDENTS MAY EDIT}

\vspace{0.15cm}

Students may edit or replace:

\begin{itemize}[leftmargin=0.8cm, itemsep=2pt]
    \item student name;
    \item class;
    \item internship organization;
    \item supervisor name;
    \item Abstract;
    \item List of Symbols and Abbreviations;
    \item chapter titles and chapter content;
    \item figures and photographs;
    \item tables;
    \item equations;
    \item citations and references.
\end{itemize}

Students shall not modify the main document formatting settings in the
preamble, including the document class, page geometry, line spacing,
heading formatting, table-of-contents formatting, font configuration,
page layout, or cover-page macros unless specifically instructed.

\vspace{0.3cm}

% ================================================================
% XXIII. FINAL SUBMISSION REQUIREMENTS
% ================================================================

\noindent
\textbf{XXIII.\quad FINAL SUBMISSION REQUIREMENTS}

\vspace{0.15cm}

Before submission, the final report shall:

\begin{itemize}[leftmargin=0.8cm, itemsep=2pt]
    \item compile without errors that affect the report content or presentation;
    \item contain all required forms and report sections;
    \item contain correct student, class, internship organization, and supervisor information;
    \item contain an updated Table of Contents, List of Figures, and List of Tables;
    \item contain no unused example content from the template;
    \item contain no placeholder text such as \textit{Title of Chapter 1},
    \textit{Section 1.1}, or other sample entries;
    \item contain all required citations and references;
    \item contain no unresolved cross-references.
\end{itemize}

If the Faculty requires submission of the \LaTeX{} source, the complete
project shall include all \verb|.tex| files, the bibliography database,
figures, and any other files required to reproduce the submitted PDF.

\vspace{0.4cm}

\noindent
\textbf{Important:}
Before submitting the final internship report, students shall remove or
comment out the following command in the main document:

\begin{verbatim}
\input{Regulations.tex}
\end{verbatim}

For example:

\begin{verbatim}
% \input{Regulations.tex}
\end{verbatim}

This removes the guidance pages from the final submitted report while
preserving the regulations file in the \LaTeX{} project for reference.

} % End 13pt content

\clearpage
|.
Students shall remove or comment out this command before submitting the
final internship report.

\vspace{0.4cm}

% ================================================================
% I. GENERAL REQUIREMENTS
% ================================================================

\noindent
\textbf{I.\quad GENERAL REQUIREMENTS}

\vspace{0.15cm}

The Graduation Internship Report shall be prepared using the official
\LaTeX{} template provided by the Faculty. Students shall preserve the
predefined page layout, typography, line spacing, heading styles,
numbering system, figure and table formatting, and bibliography settings
unless otherwise instructed by the Faculty or the supervisor.

The report shall be written in English.

Students are responsible for the academic and technical content of the
report. The official \LaTeX{} template is responsible for standardized
document formatting.

\vspace{0.3cm}

% ================================================================
% II. REPORT STRUCTURE
% ================================================================

\noindent
\textbf{II.\quad REPORT STRUCTURE}

\vspace{0.15cm}

The report shall normally be arranged in the following order:

\begin{enumerate}[label=\arabic*., leftmargin=1.2cm, itemsep=2pt]
    \item Outer Cover;
    \item Comments of the Internship Host Organization;
    \item Supervisor's Evaluation and Graduation Internship Defense Results;
    \item Table of Contents;
    \item List of Symbols and Abbreviations, if applicable;
    \item List of Figures;
    \item List of Tables;
    \item Abstract;
    \item Main Chapters of the Internship Report;
    \item References;
    \item Appendices, if applicable.
\end{enumerate}

The preliminary pages shall not be numbered as chapters. The main report
content shall be organized into consecutively numbered chapters beginning
with Chapter 1.

\vspace{0.3cm}

% ================================================================
% III. PAGE SIZE AND MARGINS
% ================================================================

\noindent
\textbf{III.\quad PAGE SIZE AND MARGINS}

\vspace{0.15cm}

The report shall be formatted on A4 paper with the following margins:

\begin{center}
\begin{tabular}{ll}
    Left margin:   & 3.5 cm \\
    Right margin:  & 2.0 cm \\
    Top margin:    & 2.5 cm \\
    Bottom margin: & 2.5 cm
\end{tabular}
\end{center}

These settings are predefined in the official template and shall not be
modified by students.

\vspace{0.3cm}

% ================================================================
% IV. FONT, FONT SIZE, AND LINE SPACING
% ================================================================

\noindent
\textbf{IV.\quad FONT, FONT SIZE, AND LINE SPACING}

\vspace{0.15cm}

The main text shall use a 13 pt font size with 1.3-line spacing.
A Times-family typeface shall be used in accordance with the official
template.

The text shall be fully justified.

Students shall not use repeated spaces, repeated manual line breaks, or
multiple blank lines to position text. Paragraph and heading spacing shall
be controlled by the template.

\vspace{0.3cm}

% ================================================================
% V. OUTER COVER
% ================================================================

\noindent
\textbf{V.\quad OUTER COVER}

\vspace{0.15cm}

The outer cover shall follow the official format provided in the template
and shall contain the following information:

\begin{itemize}[leftmargin=0.8cm, itemsep=2pt]
    \item Thai Nguyen University of Technology;
    \item Faculty of International Training;
    \item Engineering and Technology Department;
    \item University logo and Faculty logo;
    \item the title \textbf{INTERNSHIP REPORT};
    \item Internship Organization;
    \item Author;
    \item Class;
    \item Supervisor;
    \item place and year of completion.
\end{itemize}

Students shall only replace the required information fields and shall not
modify the cover-page layout unless instructed otherwise.

\vspace{0.3cm}

% ================================================================
% VI. COMMENTS OF THE INTERNSHIP HOST ORGANIZATION
% ================================================================

\noindent
\textbf{VI.\quad COMMENTS OF THE INTERNSHIP HOST ORGANIZATION}

\vspace{0.15cm}

The page entitled \textbf{COMMENTS OF THE INTERNSHIP HOST ORGANIZATION}
is reserved for the internship host organization.

It includes:

\begin{itemize}[leftmargin=0.8cm, itemsep=2pt]
    \item the student's full name;
    \item class;
    \item comments from the host organization;
    \item score;
    \item date;
    \item signature and official seal.
\end{itemize}

Students shall not alter or complete official evaluation fields unless
specifically authorized.

\vspace{0.3cm}

% ================================================================
% VII. SUPERVISOR'S EVALUATION AND DEFENSE RESULTS
% ================================================================

\noindent
\textbf{VII.\quad SUPERVISOR'S EVALUATION AND DEFENSE RESULTS}

\vspace{0.15cm}

The section entitled
\textbf{SUPERVISOR'S EVALUATION AND GRADUATION INTERNSHIP DEFENSE RESULTS}
is reserved for the supervisor and defense committee.

The evaluation may include:

\begin{itemize}[leftmargin=0.8cm, itemsep=2pt]
    \item Internship Progress and Student's Attitude;
    \item Report Content;
    \item Report Presentation;
    \item Other Comments;
    \item Supervisor's Opinion;
    \item Defense Result;
    \item Score.
\end{itemize}

Students shall not enter evaluation scores or defense results before they
are officially determined.

\vspace{0.3cm}

% ================================================================
% VIII. CHAPTERS, SECTIONS, AND SUBSECTIONS
% ================================================================

\noindent
\textbf{VIII.\quad CHAPTERS, SECTIONS, AND SUBSECTIONS}

\vspace{0.15cm}

The report structure shall be generated using the standard \LaTeX{}
commands:

\begin{verbatim}
\chapter{...}
\section{...}
\subsection{...}
\end{verbatim}

Chapter, section, and subsection numbers shall be generated automatically
by \LaTeX{}.

Students shall not manually type chapter or section numbers and shall not
reproduce heading styles using commands such as \verb|\textbf|,
\verb|\fontsize|, or \verb|\centering| instead of the appropriate
structural command.

Excessive subdivision should be avoided. In most cases, the hierarchy
Chapter--Section--Subsection is sufficient.

\vspace{0.3cm}

% ================================================================
% IX. ORGANIZATION OF THE MAIN CONTENT
% ================================================================

\noindent
\textbf{IX.\quad ORGANIZATION OF THE MAIN CONTENT}

\vspace{0.15cm}

The specific content of the report shall depend on the internship
organization and the assigned tasks. A typical structure may include:

\begin{itemize}[leftmargin=0.8cm, itemsep=2pt]
    \item \textbf{Chapter 1. Overview of the Internship Organization};
    \item \textbf{Chapter 2. Internship Activities and Assigned Tasks};
    \item \textbf{Chapter 3. Results, Experience, and Evaluation}.
\end{itemize}

These chapter titles are recommendations only and may be adapted to the
actual internship content.

The report shall reflect work actually performed, observed, analyzed, or
learned during the internship. Unrelated theoretical material should not
be included merely to increase report length.

\vspace{0.3cm}

% ================================================================
% X. TABLE OF CONTENTS
% ================================================================

\noindent
\textbf{X.\quad TABLE OF CONTENTS}

\vspace{0.15cm}

The Table of Contents shall be generated automatically using:

\begin{verbatim}
\tableofcontents
\end{verbatim}

Students shall not create or manually edit page numbers in the Table of
Contents.

The document shall be compiled sufficiently to ensure that all page
numbers and entries are correctly updated.

\vspace{0.3cm}

% ================================================================
% XI. LIST OF SYMBOLS AND ABBREVIATIONS
% ================================================================

\noindent
\textbf{XI.\quad LIST OF SYMBOLS AND ABBREVIATIONS}

\vspace{0.15cm}

Symbols and abbreviations that are used repeatedly in the report shall be
listed in the \textbf{List of Symbols and Abbreviations}.

The list may contain the following columns:

\begin{center}
\begin{tabular}{llll}
    No. & Symbol / Abbreviation & Unit & Definition
\end{tabular}
\end{center}

Only symbols and abbreviations actually used in the final report shall be
retained. Example entries contained in the template shall be removed if
they are not relevant to the student's report.

Each abbreviation should normally be defined when it first appears in
the main text.

\vspace{0.3cm}

% ================================================================
% XII. FIGURES, DIAGRAMS, AND PHOTOGRAPHS
% ================================================================

\noindent
\textbf{XII.\quad FIGURES, DIAGRAMS, AND PHOTOGRAPHS}

\vspace{0.15cm}

Figures shall be clear, legible, and directly relevant to the report
content.

Figures shall be inserted using the standard \verb|figure| environment.

Example:

\begin{verbatim}
\begin{figure}[H]
    \centering
    \includegraphics[width=0.8\textwidth]
                    {figures/company_structure.pdf}
    \caption{Organizational structure of the company}
    \label{fig:company-structure}
\end{figure}
\end{verbatim}

Figure captions shall be placed below the figure.

Each referenced figure shall have a unique \verb|\label| and shall be
referenced in the text using, for example:

\begin{verbatim}
Figure~\ref{fig:company-structure}
\end{verbatim}

Students shall not manually type figure numbers.

Figures reproduced or adapted from external sources shall include an
appropriate citation.

Photographs taken during the internship shall be relevant and sufficiently
clear. Students shall not include confidential, restricted, or proprietary
information without permission from the internship organization.

Vector graphics or high-resolution images are recommended whenever
possible.

\vspace{0.3cm}

% ================================================================
% XIII. TABLES
% ================================================================

\noindent
\textbf{XIII.\quad TABLES}

\vspace{0.15cm}

Tables shall be created using suitable \LaTeX{} environments such as
\verb|table|, \verb|tabular|, \verb|tabularx|, or \verb|longtable|.

Table captions shall be placed above the table.

Example:

\begin{verbatim}
\begin{table}[H]
    \centering
    \caption{Internship schedule}
    \label{tab:internship-schedule}

    \begin{tabular}{lll}
        \toprule
        Week & Task & Result \\
        \midrule
        1 & Orientation & Completed \\
        2 & System study & Completed \\
        \bottomrule
    \end{tabular}
\end{table}
\end{verbatim}

Tables shall be referenced using:

\begin{verbatim}
Table~\ref{tab:internship-schedule}
\end{verbatim}

Students shall not manually type table numbers.

Tables should be created directly in \LaTeX{} rather than inserted as
screenshots whenever practical.

\vspace{0.3cm}

% ================================================================
% XIV. EQUATIONS AND TECHNICAL QUANTITIES
% ================================================================

\noindent
\textbf{XIV.\quad EQUATIONS AND TECHNICAL QUANTITIES}

\vspace{0.15cm}

Mathematical expressions shall be typeset directly using \LaTeX{}.
Images or screenshots of equations shall not be used.

Short expressions may be written in inline mathematical mode, for example:

\begin{verbatim}
$R = 0.7\,\Omega$
\end{verbatim}

Important equations shall use a numbered equation environment.

Example:

\begin{verbatim}
\begin{equation}
    P = UI
    \label{eq:power}
\end{equation}
\end{verbatim}

Equations shall be referenced using:

\begin{verbatim}
Equation~\eqref{eq:power}
\end{verbatim}

Equation numbers shall not be manually entered.

SI units shall be used wherever applicable. The \verb|siunitx| package
should be used where appropriate.

Examples:

\begin{verbatim}
\SI{220}{\volt}
\SI{50}{\hertz}
\SI{2.5}{\ampere}
\end{verbatim}

Notation and units shall be consistent throughout the report.

\vspace{0.3cm}

% ================================================================
% XV. SOURCE CODE
% ================================================================

\noindent
\textbf{XV.\quad SOURCE CODE}

\vspace{0.15cm}

If the internship involves programming, PLCs, microcontrollers,
simulation, or software development, relevant source code may be included
in the report.

Source code shall be presented using the \verb|listings| environment or
another environment defined by the official template.

Example:

\begin{verbatim}
\begin{lstlisting}[language=st]
IF Start = TRUE THEN
    Motor := TRUE;
ELSE
    Motor := FALSE;
END_IF;
\end{lstlisting}
\end{verbatim}

Only code necessary for explaining the work should appear in the main
report. Long source-code listings should be moved to an appendix or
submitted as supplementary files if required.

Source code should not normally be inserted as screenshots.

\vspace{0.3cm}

% ================================================================
% XVI. ABSTRACT
% ================================================================

\noindent
\textbf{XVI.\quad ABSTRACT}

\vspace{0.15cm}

The Abstract shall briefly summarize:

\begin{itemize}[leftmargin=0.8cm, itemsep=2pt]
    \item the internship organization;
    \item the objectives of the internship;
    \item the main tasks performed;
    \item the main results, knowledge, and skills obtained.
\end{itemize}

The Abstract should be concise and should not reproduce the full content
of the report chapter by chapter.

\vspace{0.3cm}

% ================================================================
% XVII. REFERENCES AND CITATIONS
% ================================================================

\noindent
\textbf{XVII.\quad REFERENCES AND CITATIONS}

\vspace{0.15cm}

All books, journal papers, conference papers, standards, websites,
datasheets, manuals, technical documents, figures, tables, or other
external sources used in preparing the report shall be properly cited.

References shall be managed through the bibliography system provided with
the official template.

Bibliographic entries shall normally be stored in:

\begin{verbatim}
References.bib
\end{verbatim}

References shall be cited in the report using:

\begin{verbatim}
\cite{reference-key}
\end{verbatim}

Students shall not manually type reference numbers such as ``[1]'',
``[2]'', etc.

Titles of foreign-language publications shall remain in their original
language in the bibliography.

All borrowed ideas, data, figures, tables, equations, methods, or results
shall be properly acknowledged.

\vspace{0.3cm}

% ================================================================
% XVIII. CROSS-REFERENCING
% ================================================================

\noindent
\textbf{XVIII.\quad CROSS-REFERENCING}

\vspace{0.15cm}

Numbered elements shall be cross-referenced using the standard \LaTeX{}
mechanism.

Use the following commands where appropriate:

\begin{verbatim}
\label{...}
\ref{...}
\eqref{...}
\end{verbatim}

For example:

\begin{verbatim}
As shown in Figure~\ref{fig:system}, ...
\end{verbatim}

Students shall not manually type numbers such as ``Figure 2.3'' when
those numbers can be generated automatically.

\vspace{0.3cm}

% ================================================================
% XIX. PAGE NUMBERING
% ================================================================

\noindent
\textbf{XIX.\quad PAGE NUMBERING}

\vspace{0.15cm}

Cover pages and official evaluation forms at the beginning of the report
shall not display page numbers.

The main body of the report shall use Arabic page numbers beginning from
Chapter 1.

Page numbering shall be generated automatically by the official template.
Students shall not manually enter page numbers.

\vspace{0.3cm}

% ================================================================
% XX. PROJECT FILE ORGANIZATION
% ================================================================

\noindent
\textbf{XX.\quad PROJECT FILE ORGANIZATION}

\vspace{0.15cm}

Students should maintain a clear and organized \LaTeX{} project
structure. A recommended structure is:

\begin{verbatim}
InternshipReport/
|
|-- main.tex
|-- References.bib
|
|-- figures/
|   |-- company.pdf
|   |-- system.png
|   `-- ...
|
`-- chapters/
    |-- chapter1.tex
    |-- chapter2.tex
    `-- chapter3.tex
\end{verbatim}

Figures should normally be stored in the \verb|figures| directory
specified by the official template.

File and directory names should preferably contain Latin letters,
numbers, hyphens, or underscores.

Recommended examples include:

\begin{verbatim}
company_structure.pdf
motor_control.png
chapter_01.tex
\end{verbatim}

Spaces and unnecessary special characters in file names should be
avoided where possible.

\vspace{0.3cm}

% ================================================================
% XXI. DOCUMENT COMPILATION AND FINAL CHECK
% ================================================================

\noindent
\textbf{XXI.\quad DOCUMENT COMPILATION AND FINAL CHECK}

\vspace{0.15cm}

Before submission, the complete \LaTeX{} project shall compile
successfully and produce the final PDF report.

The final document shall not contain unresolved references such as:

\begin{verbatim}
??
[?]
\end{verbatim}

Before generating the final PDF, students shall check:

\begin{itemize}[leftmargin=0.8cm, itemsep=2pt]
    \item the Table of Contents;
    \item page numbering;
    \item chapter, section, and subsection numbering;
    \item figure numbering;
    \item table numbering;
    \item equation numbering;
    \item cross-references;
    \item citations;
    \item bibliography;
    \item font display;
    \item mathematical symbols;
    \item image quality;
    \item completeness of all required forms.
\end{itemize}

The document may be compiled using TeX Live, MiKTeX, or another
compatible \LaTeX{} environment. Regardless of the compilation
environment, the resulting PDF shall conform to the official Faculty
template.

\vspace{0.3cm}

% ================================================================
% XXII. CONTENT THAT STUDENTS MAY EDIT
% ================================================================

\noindent
\textbf{XXII.\quad CONTENT THAT STUDENTS MAY EDIT}

\vspace{0.15cm}

Students may edit or replace:

\begin{itemize}[leftmargin=0.8cm, itemsep=2pt]
    \item student name;
    \item class;
    \item internship organization;
    \item supervisor name;
    \item Abstract;
    \item List of Symbols and Abbreviations;
    \item chapter titles and chapter content;
    \item figures and photographs;
    \item tables;
    \item equations;
    \item citations and references.
\end{itemize}

Students shall not modify the main document formatting settings in the
preamble, including the document class, page geometry, line spacing,
heading formatting, table-of-contents formatting, font configuration,
page layout, or cover-page macros unless specifically instructed.

\vspace{0.3cm}

% ================================================================
% XXIII. FINAL SUBMISSION REQUIREMENTS
% ================================================================

\noindent
\textbf{XXIII.\quad FINAL SUBMISSION REQUIREMENTS}

\vspace{0.15cm}

Before submission, the final report shall:

\begin{itemize}[leftmargin=0.8cm, itemsep=2pt]
    \item compile without errors that affect the report content or presentation;
    \item contain all required forms and report sections;
    \item contain correct student, class, internship organization, and supervisor information;
    \item contain an updated Table of Contents, List of Figures, and List of Tables;
    \item contain no unused example content from the template;
    \item contain no placeholder text such as \textit{Title of Chapter 1},
    \textit{Section 1.1}, or other sample entries;
    \item contain all required citations and references;
    \item contain no unresolved cross-references.
\end{itemize}

If the Faculty requires submission of the \LaTeX{} source, the complete
project shall include all \verb|.tex| files, the bibliography database,
figures, and any other files required to reproduce the submitted PDF.

\vspace{0.4cm}

\noindent
\textbf{Important:}
Before submitting the final internship report, students shall remove or
comment out the following command in the main document:

\begin{verbatim}

% =================================================================
% REGULATIONS FOR PREPARING THE GRADUATION INTERNSHIP REPORT
% USING LATEX
% =================================================================
%
% This file is intended to be included from the main document with:
%
%     \input{Regulations.tex}
%
% IMPORTANT:
% This section is provided for guidance only. Students shall remove
% or comment out the \input{Regulations.tex} command before submitting
% the final internship report.
%
% =================================================================

\clearpage
\thispagestyle{empty}

\begin{center}
    {\fontsize{16pt}{20pt}\selectfont\bfseries
    SPECIFIC REGULATIONS FOR PREPARING\\
    GRADUATION INTERNSHIP REPORTS USING \LaTeX}
\end{center}

\vspace{0.5cm}

{\fontsize{13pt}{17pt}\selectfont
\setstretch{1.3}
\justifying

\noindent
\textbf{Note:}
This section is provided for guidance only and is included in the main
document using the command \verb|\input{Regulations.tex}|.
Students shall remove or comment out this command before submitting the
final internship report.

\vspace{0.4cm}

% ================================================================
% I. GENERAL REQUIREMENTS
% ================================================================

\noindent
\textbf{I.\quad GENERAL REQUIREMENTS}

\vspace{0.15cm}

The Graduation Internship Report shall be prepared using the official
\LaTeX{} template provided by the Faculty. Students shall preserve the
predefined page layout, typography, line spacing, heading styles,
numbering system, figure and table formatting, and bibliography settings
unless otherwise instructed by the Faculty or the supervisor.

The report shall be written in English.

Students are responsible for the academic and technical content of the
report. The official \LaTeX{} template is responsible for standardized
document formatting.

\vspace{0.3cm}

% ================================================================
% II. REPORT STRUCTURE
% ================================================================

\noindent
\textbf{II.\quad REPORT STRUCTURE}

\vspace{0.15cm}

The report shall normally be arranged in the following order:

\begin{enumerate}[label=\arabic*., leftmargin=1.2cm, itemsep=2pt]
    \item Outer Cover;
    \item Comments of the Internship Host Organization;
    \item Supervisor's Evaluation and Graduation Internship Defense Results;
    \item Table of Contents;
    \item List of Symbols and Abbreviations, if applicable;
    \item List of Figures;
    \item List of Tables;
    \item Abstract;
    \item Main Chapters of the Internship Report;
    \item References;
    \item Appendices, if applicable.
\end{enumerate}

The preliminary pages shall not be numbered as chapters. The main report
content shall be organized into consecutively numbered chapters beginning
with Chapter 1.

\vspace{0.3cm}

% ================================================================
% III. PAGE SIZE AND MARGINS
% ================================================================

\noindent
\textbf{III.\quad PAGE SIZE AND MARGINS}

\vspace{0.15cm}

The report shall be formatted on A4 paper with the following margins:

\begin{center}
\begin{tabular}{ll}
    Left margin:   & 3.5 cm \\
    Right margin:  & 2.0 cm \\
    Top margin:    & 2.5 cm \\
    Bottom margin: & 2.5 cm
\end{tabular}
\end{center}

These settings are predefined in the official template and shall not be
modified by students.

\vspace{0.3cm}

% ================================================================
% IV. FONT, FONT SIZE, AND LINE SPACING
% ================================================================

\noindent
\textbf{IV.\quad FONT, FONT SIZE, AND LINE SPACING}

\vspace{0.15cm}

The main text shall use a 13 pt font size with 1.3-line spacing.
A Times-family typeface shall be used in accordance with the official
template.

The text shall be fully justified.

Students shall not use repeated spaces, repeated manual line breaks, or
multiple blank lines to position text. Paragraph and heading spacing shall
be controlled by the template.

\vspace{0.3cm}

% ================================================================
% V. OUTER COVER
% ================================================================

\noindent
\textbf{V.\quad OUTER COVER}

\vspace{0.15cm}

The outer cover shall follow the official format provided in the template
and shall contain the following information:

\begin{itemize}[leftmargin=0.8cm, itemsep=2pt]
    \item Thai Nguyen University of Technology;
    \item Faculty of International Training;
    \item Engineering and Technology Department;
    \item University logo and Faculty logo;
    \item the title \textbf{INTERNSHIP REPORT};
    \item Internship Organization;
    \item Author;
    \item Class;
    \item Supervisor;
    \item place and year of completion.
\end{itemize}

Students shall only replace the required information fields and shall not
modify the cover-page layout unless instructed otherwise.

\vspace{0.3cm}

% ================================================================
% VI. COMMENTS OF THE INTERNSHIP HOST ORGANIZATION
% ================================================================

\noindent
\textbf{VI.\quad COMMENTS OF THE INTERNSHIP HOST ORGANIZATION}

\vspace{0.15cm}

The page entitled \textbf{COMMENTS OF THE INTERNSHIP HOST ORGANIZATION}
is reserved for the internship host organization.

It includes:

\begin{itemize}[leftmargin=0.8cm, itemsep=2pt]
    \item the student's full name;
    \item class;
    \item comments from the host organization;
    \item score;
    \item date;
    \item signature and official seal.
\end{itemize}

Students shall not alter or complete official evaluation fields unless
specifically authorized.

\vspace{0.3cm}

% ================================================================
% VII. SUPERVISOR'S EVALUATION AND DEFENSE RESULTS
% ================================================================

\noindent
\textbf{VII.\quad SUPERVISOR'S EVALUATION AND DEFENSE RESULTS}

\vspace{0.15cm}

The section entitled
\textbf{SUPERVISOR'S EVALUATION AND GRADUATION INTERNSHIP DEFENSE RESULTS}
is reserved for the supervisor and defense committee.

The evaluation may include:

\begin{itemize}[leftmargin=0.8cm, itemsep=2pt]
    \item Internship Progress and Student's Attitude;
    \item Report Content;
    \item Report Presentation;
    \item Other Comments;
    \item Supervisor's Opinion;
    \item Defense Result;
    \item Score.
\end{itemize}

Students shall not enter evaluation scores or defense results before they
are officially determined.

\vspace{0.3cm}

% ================================================================
% VIII. CHAPTERS, SECTIONS, AND SUBSECTIONS
% ================================================================

\noindent
\textbf{VIII.\quad CHAPTERS, SECTIONS, AND SUBSECTIONS}

\vspace{0.15cm}

The report structure shall be generated using the standard \LaTeX{}
commands:

\begin{verbatim}
\chapter{...}
\section{...}
\subsection{...}
\end{verbatim}

Chapter, section, and subsection numbers shall be generated automatically
by \LaTeX{}.

Students shall not manually type chapter or section numbers and shall not
reproduce heading styles using commands such as \verb|\textbf|,
\verb|\fontsize|, or \verb|\centering| instead of the appropriate
structural command.

Excessive subdivision should be avoided. In most cases, the hierarchy
Chapter--Section--Subsection is sufficient.

\vspace{0.3cm}

% ================================================================
% IX. ORGANIZATION OF THE MAIN CONTENT
% ================================================================

\noindent
\textbf{IX.\quad ORGANIZATION OF THE MAIN CONTENT}

\vspace{0.15cm}

The specific content of the report shall depend on the internship
organization and the assigned tasks. A typical structure may include:

\begin{itemize}[leftmargin=0.8cm, itemsep=2pt]
    \item \textbf{Chapter 1. Overview of the Internship Organization};
    \item \textbf{Chapter 2. Internship Activities and Assigned Tasks};
    \item \textbf{Chapter 3. Results, Experience, and Evaluation}.
\end{itemize}

These chapter titles are recommendations only and may be adapted to the
actual internship content.

The report shall reflect work actually performed, observed, analyzed, or
learned during the internship. Unrelated theoretical material should not
be included merely to increase report length.

\vspace{0.3cm}

% ================================================================
% X. TABLE OF CONTENTS
% ================================================================

\noindent
\textbf{X.\quad TABLE OF CONTENTS}

\vspace{0.15cm}

The Table of Contents shall be generated automatically using:

\begin{verbatim}
\tableofcontents
\end{verbatim}

Students shall not create or manually edit page numbers in the Table of
Contents.

The document shall be compiled sufficiently to ensure that all page
numbers and entries are correctly updated.

\vspace{0.3cm}

% ================================================================
% XI. LIST OF SYMBOLS AND ABBREVIATIONS
% ================================================================

\noindent
\textbf{XI.\quad LIST OF SYMBOLS AND ABBREVIATIONS}

\vspace{0.15cm}

Symbols and abbreviations that are used repeatedly in the report shall be
listed in the \textbf{List of Symbols and Abbreviations}.

The list may contain the following columns:

\begin{center}
\begin{tabular}{llll}
    No. & Symbol / Abbreviation & Unit & Definition
\end{tabular}
\end{center}

Only symbols and abbreviations actually used in the final report shall be
retained. Example entries contained in the template shall be removed if
they are not relevant to the student's report.

Each abbreviation should normally be defined when it first appears in
the main text.

\vspace{0.3cm}

% ================================================================
% XII. FIGURES, DIAGRAMS, AND PHOTOGRAPHS
% ================================================================

\noindent
\textbf{XII.\quad FIGURES, DIAGRAMS, AND PHOTOGRAPHS}

\vspace{0.15cm}

Figures shall be clear, legible, and directly relevant to the report
content.

Figures shall be inserted using the standard \verb|figure| environment.

Example:

\begin{verbatim}
\begin{figure}[H]
    \centering
    \includegraphics[width=0.8\textwidth]
                    {figures/company_structure.pdf}
    \caption{Organizational structure of the company}
    \label{fig:company-structure}
\end{figure}
\end{verbatim}

Figure captions shall be placed below the figure.

Each referenced figure shall have a unique \verb|\label| and shall be
referenced in the text using, for example:

\begin{verbatim}
Figure~\ref{fig:company-structure}
\end{verbatim}

Students shall not manually type figure numbers.

Figures reproduced or adapted from external sources shall include an
appropriate citation.

Photographs taken during the internship shall be relevant and sufficiently
clear. Students shall not include confidential, restricted, or proprietary
information without permission from the internship organization.

Vector graphics or high-resolution images are recommended whenever
possible.

\vspace{0.3cm}

% ================================================================
% XIII. TABLES
% ================================================================

\noindent
\textbf{XIII.\quad TABLES}

\vspace{0.15cm}

Tables shall be created using suitable \LaTeX{} environments such as
\verb|table|, \verb|tabular|, \verb|tabularx|, or \verb|longtable|.

Table captions shall be placed above the table.

Example:

\begin{verbatim}
\begin{table}[H]
    \centering
    \caption{Internship schedule}
    \label{tab:internship-schedule}

    \begin{tabular}{lll}
        \toprule
        Week & Task & Result \\
        \midrule
        1 & Orientation & Completed \\
        2 & System study & Completed \\
        \bottomrule
    \end{tabular}
\end{table}
\end{verbatim}

Tables shall be referenced using:

\begin{verbatim}
Table~\ref{tab:internship-schedule}
\end{verbatim}

Students shall not manually type table numbers.

Tables should be created directly in \LaTeX{} rather than inserted as
screenshots whenever practical.

\vspace{0.3cm}

% ================================================================
% XIV. EQUATIONS AND TECHNICAL QUANTITIES
% ================================================================

\noindent
\textbf{XIV.\quad EQUATIONS AND TECHNICAL QUANTITIES}

\vspace{0.15cm}

Mathematical expressions shall be typeset directly using \LaTeX{}.
Images or screenshots of equations shall not be used.

Short expressions may be written in inline mathematical mode, for example:

\begin{verbatim}
$R = 0.7\,\Omega$
\end{verbatim}

Important equations shall use a numbered equation environment.

Example:

\begin{verbatim}
\begin{equation}
    P = UI
    \label{eq:power}
\end{equation}
\end{verbatim}

Equations shall be referenced using:

\begin{verbatim}
Equation~\eqref{eq:power}
\end{verbatim}

Equation numbers shall not be manually entered.

SI units shall be used wherever applicable. The \verb|siunitx| package
should be used where appropriate.

Examples:

\begin{verbatim}
\SI{220}{\volt}
\SI{50}{\hertz}
\SI{2.5}{\ampere}
\end{verbatim}

Notation and units shall be consistent throughout the report.

\vspace{0.3cm}

% ================================================================
% XV. SOURCE CODE
% ================================================================

\noindent
\textbf{XV.\quad SOURCE CODE}

\vspace{0.15cm}

If the internship involves programming, PLCs, microcontrollers,
simulation, or software development, relevant source code may be included
in the report.

Source code shall be presented using the \verb|listings| environment or
another environment defined by the official template.

Example:

\begin{verbatim}
\begin{lstlisting}[language=st]
IF Start = TRUE THEN
    Motor := TRUE;
ELSE
    Motor := FALSE;
END_IF;
\end{lstlisting}
\end{verbatim}

Only code necessary for explaining the work should appear in the main
report. Long source-code listings should be moved to an appendix or
submitted as supplementary files if required.

Source code should not normally be inserted as screenshots.

\vspace{0.3cm}

% ================================================================
% XVI. ABSTRACT
% ================================================================

\noindent
\textbf{XVI.\quad ABSTRACT}

\vspace{0.15cm}

The Abstract shall briefly summarize:

\begin{itemize}[leftmargin=0.8cm, itemsep=2pt]
    \item the internship organization;
    \item the objectives of the internship;
    \item the main tasks performed;
    \item the main results, knowledge, and skills obtained.
\end{itemize}

The Abstract should be concise and should not reproduce the full content
of the report chapter by chapter.

\vspace{0.3cm}

% ================================================================
% XVII. REFERENCES AND CITATIONS
% ================================================================

\noindent
\textbf{XVII.\quad REFERENCES AND CITATIONS}

\vspace{0.15cm}

All books, journal papers, conference papers, standards, websites,
datasheets, manuals, technical documents, figures, tables, or other
external sources used in preparing the report shall be properly cited.

References shall be managed through the bibliography system provided with
the official template.

Bibliographic entries shall normally be stored in:

\begin{verbatim}
References.bib
\end{verbatim}

References shall be cited in the report using:

\begin{verbatim}
\cite{reference-key}
\end{verbatim}

Students shall not manually type reference numbers such as ``[1]'',
``[2]'', etc.

Titles of foreign-language publications shall remain in their original
language in the bibliography.

All borrowed ideas, data, figures, tables, equations, methods, or results
shall be properly acknowledged.

\vspace{0.3cm}

% ================================================================
% XVIII. CROSS-REFERENCING
% ================================================================

\noindent
\textbf{XVIII.\quad CROSS-REFERENCING}

\vspace{0.15cm}

Numbered elements shall be cross-referenced using the standard \LaTeX{}
mechanism.

Use the following commands where appropriate:

\begin{verbatim}
\label{...}
\ref{...}
\eqref{...}
\end{verbatim}

For example:

\begin{verbatim}
As shown in Figure~\ref{fig:system}, ...
\end{verbatim}

Students shall not manually type numbers such as ``Figure 2.3'' when
those numbers can be generated automatically.

\vspace{0.3cm}

% ================================================================
% XIX. PAGE NUMBERING
% ================================================================

\noindent
\textbf{XIX.\quad PAGE NUMBERING}

\vspace{0.15cm}

Cover pages and official evaluation forms at the beginning of the report
shall not display page numbers.

The main body of the report shall use Arabic page numbers beginning from
Chapter 1.

Page numbering shall be generated automatically by the official template.
Students shall not manually enter page numbers.

\vspace{0.3cm}

% ================================================================
% XX. PROJECT FILE ORGANIZATION
% ================================================================

\noindent
\textbf{XX.\quad PROJECT FILE ORGANIZATION}

\vspace{0.15cm}

Students should maintain a clear and organized \LaTeX{} project
structure. A recommended structure is:

\begin{verbatim}
InternshipReport/
|
|-- main.tex
|-- References.bib
|
|-- figures/
|   |-- company.pdf
|   |-- system.png
|   `-- ...
|
`-- chapters/
    |-- chapter1.tex
    |-- chapter2.tex
    `-- chapter3.tex
\end{verbatim}

Figures should normally be stored in the \verb|figures| directory
specified by the official template.

File and directory names should preferably contain Latin letters,
numbers, hyphens, or underscores.

Recommended examples include:

\begin{verbatim}
company_structure.pdf
motor_control.png
chapter_01.tex
\end{verbatim}

Spaces and unnecessary special characters in file names should be
avoided where possible.

\vspace{0.3cm}

% ================================================================
% XXI. DOCUMENT COMPILATION AND FINAL CHECK
% ================================================================

\noindent
\textbf{XXI.\quad DOCUMENT COMPILATION AND FINAL CHECK}

\vspace{0.15cm}

Before submission, the complete \LaTeX{} project shall compile
successfully and produce the final PDF report.

The final document shall not contain unresolved references such as:

\begin{verbatim}
??
[?]
\end{verbatim}

Before generating the final PDF, students shall check:

\begin{itemize}[leftmargin=0.8cm, itemsep=2pt]
    \item the Table of Contents;
    \item page numbering;
    \item chapter, section, and subsection numbering;
    \item figure numbering;
    \item table numbering;
    \item equation numbering;
    \item cross-references;
    \item citations;
    \item bibliography;
    \item font display;
    \item mathematical symbols;
    \item image quality;
    \item completeness of all required forms.
\end{itemize}

The document may be compiled using TeX Live, MiKTeX, or another
compatible \LaTeX{} environment. Regardless of the compilation
environment, the resulting PDF shall conform to the official Faculty
template.

\vspace{0.3cm}

% ================================================================
% XXII. CONTENT THAT STUDENTS MAY EDIT
% ================================================================

\noindent
\textbf{XXII.\quad CONTENT THAT STUDENTS MAY EDIT}

\vspace{0.15cm}

Students may edit or replace:

\begin{itemize}[leftmargin=0.8cm, itemsep=2pt]
    \item student name;
    \item class;
    \item internship organization;
    \item supervisor name;
    \item Abstract;
    \item List of Symbols and Abbreviations;
    \item chapter titles and chapter content;
    \item figures and photographs;
    \item tables;
    \item equations;
    \item citations and references.
\end{itemize}

Students shall not modify the main document formatting settings in the
preamble, including the document class, page geometry, line spacing,
heading formatting, table-of-contents formatting, font configuration,
page layout, or cover-page macros unless specifically instructed.

\vspace{0.3cm}

% ================================================================
% XXIII. FINAL SUBMISSION REQUIREMENTS
% ================================================================

\noindent
\textbf{XXIII.\quad FINAL SUBMISSION REQUIREMENTS}

\vspace{0.15cm}

Before submission, the final report shall:

\begin{itemize}[leftmargin=0.8cm, itemsep=2pt]
    \item compile without errors that affect the report content or presentation;
    \item contain all required forms and report sections;
    \item contain correct student, class, internship organization, and supervisor information;
    \item contain an updated Table of Contents, List of Figures, and List of Tables;
    \item contain no unused example content from the template;
    \item contain no placeholder text such as \textit{Title of Chapter 1},
    \textit{Section 1.1}, or other sample entries;
    \item contain all required citations and references;
    \item contain no unresolved cross-references.
\end{itemize}

If the Faculty requires submission of the \LaTeX{} source, the complete
project shall include all \verb|.tex| files, the bibliography database,
figures, and any other files required to reproduce the submitted PDF.

\vspace{0.4cm}

\noindent
\textbf{Important:}
Before submitting the final internship report, students shall remove or
comment out the following command in the main document:

\begin{verbatim}
\input{Regulations.tex}
\end{verbatim}

For example:

\begin{verbatim}
% \input{Regulations.tex}
\end{verbatim}

This removes the guidance pages from the final submitted report while
preserving the regulations file in the \LaTeX{} project for reference.

} % End 13pt content

\clearpage

\end{verbatim}

For example:

\begin{verbatim}
% 
% =================================================================
% REGULATIONS FOR PREPARING THE GRADUATION INTERNSHIP REPORT
% USING LATEX
% =================================================================
%
% This file is intended to be included from the main document with:
%
%     \input{Regulations.tex}
%
% IMPORTANT:
% This section is provided for guidance only. Students shall remove
% or comment out the \input{Regulations.tex} command before submitting
% the final internship report.
%
% =================================================================

\clearpage
\thispagestyle{empty}

\begin{center}
    {\fontsize{16pt}{20pt}\selectfont\bfseries
    SPECIFIC REGULATIONS FOR PREPARING\\
    GRADUATION INTERNSHIP REPORTS USING \LaTeX}
\end{center}

\vspace{0.5cm}

{\fontsize{13pt}{17pt}\selectfont
\setstretch{1.3}
\justifying

\noindent
\textbf{Note:}
This section is provided for guidance only and is included in the main
document using the command \verb|\input{Regulations.tex}|.
Students shall remove or comment out this command before submitting the
final internship report.

\vspace{0.4cm}

% ================================================================
% I. GENERAL REQUIREMENTS
% ================================================================

\noindent
\textbf{I.\quad GENERAL REQUIREMENTS}

\vspace{0.15cm}

The Graduation Internship Report shall be prepared using the official
\LaTeX{} template provided by the Faculty. Students shall preserve the
predefined page layout, typography, line spacing, heading styles,
numbering system, figure and table formatting, and bibliography settings
unless otherwise instructed by the Faculty or the supervisor.

The report shall be written in English.

Students are responsible for the academic and technical content of the
report. The official \LaTeX{} template is responsible for standardized
document formatting.

\vspace{0.3cm}

% ================================================================
% II. REPORT STRUCTURE
% ================================================================

\noindent
\textbf{II.\quad REPORT STRUCTURE}

\vspace{0.15cm}

The report shall normally be arranged in the following order:

\begin{enumerate}[label=\arabic*., leftmargin=1.2cm, itemsep=2pt]
    \item Outer Cover;
    \item Comments of the Internship Host Organization;
    \item Supervisor's Evaluation and Graduation Internship Defense Results;
    \item Table of Contents;
    \item List of Symbols and Abbreviations, if applicable;
    \item List of Figures;
    \item List of Tables;
    \item Abstract;
    \item Main Chapters of the Internship Report;
    \item References;
    \item Appendices, if applicable.
\end{enumerate}

The preliminary pages shall not be numbered as chapters. The main report
content shall be organized into consecutively numbered chapters beginning
with Chapter 1.

\vspace{0.3cm}

% ================================================================
% III. PAGE SIZE AND MARGINS
% ================================================================

\noindent
\textbf{III.\quad PAGE SIZE AND MARGINS}

\vspace{0.15cm}

The report shall be formatted on A4 paper with the following margins:

\begin{center}
\begin{tabular}{ll}
    Left margin:   & 3.5 cm \\
    Right margin:  & 2.0 cm \\
    Top margin:    & 2.5 cm \\
    Bottom margin: & 2.5 cm
\end{tabular}
\end{center}

These settings are predefined in the official template and shall not be
modified by students.

\vspace{0.3cm}

% ================================================================
% IV. FONT, FONT SIZE, AND LINE SPACING
% ================================================================

\noindent
\textbf{IV.\quad FONT, FONT SIZE, AND LINE SPACING}

\vspace{0.15cm}

The main text shall use a 13 pt font size with 1.3-line spacing.
A Times-family typeface shall be used in accordance with the official
template.

The text shall be fully justified.

Students shall not use repeated spaces, repeated manual line breaks, or
multiple blank lines to position text. Paragraph and heading spacing shall
be controlled by the template.

\vspace{0.3cm}

% ================================================================
% V. OUTER COVER
% ================================================================

\noindent
\textbf{V.\quad OUTER COVER}

\vspace{0.15cm}

The outer cover shall follow the official format provided in the template
and shall contain the following information:

\begin{itemize}[leftmargin=0.8cm, itemsep=2pt]
    \item Thai Nguyen University of Technology;
    \item Faculty of International Training;
    \item Engineering and Technology Department;
    \item University logo and Faculty logo;
    \item the title \textbf{INTERNSHIP REPORT};
    \item Internship Organization;
    \item Author;
    \item Class;
    \item Supervisor;
    \item place and year of completion.
\end{itemize}

Students shall only replace the required information fields and shall not
modify the cover-page layout unless instructed otherwise.

\vspace{0.3cm}

% ================================================================
% VI. COMMENTS OF THE INTERNSHIP HOST ORGANIZATION
% ================================================================

\noindent
\textbf{VI.\quad COMMENTS OF THE INTERNSHIP HOST ORGANIZATION}

\vspace{0.15cm}

The page entitled \textbf{COMMENTS OF THE INTERNSHIP HOST ORGANIZATION}
is reserved for the internship host organization.

It includes:

\begin{itemize}[leftmargin=0.8cm, itemsep=2pt]
    \item the student's full name;
    \item class;
    \item comments from the host organization;
    \item score;
    \item date;
    \item signature and official seal.
\end{itemize}

Students shall not alter or complete official evaluation fields unless
specifically authorized.

\vspace{0.3cm}

% ================================================================
% VII. SUPERVISOR'S EVALUATION AND DEFENSE RESULTS
% ================================================================

\noindent
\textbf{VII.\quad SUPERVISOR'S EVALUATION AND DEFENSE RESULTS}

\vspace{0.15cm}

The section entitled
\textbf{SUPERVISOR'S EVALUATION AND GRADUATION INTERNSHIP DEFENSE RESULTS}
is reserved for the supervisor and defense committee.

The evaluation may include:

\begin{itemize}[leftmargin=0.8cm, itemsep=2pt]
    \item Internship Progress and Student's Attitude;
    \item Report Content;
    \item Report Presentation;
    \item Other Comments;
    \item Supervisor's Opinion;
    \item Defense Result;
    \item Score.
\end{itemize}

Students shall not enter evaluation scores or defense results before they
are officially determined.

\vspace{0.3cm}

% ================================================================
% VIII. CHAPTERS, SECTIONS, AND SUBSECTIONS
% ================================================================

\noindent
\textbf{VIII.\quad CHAPTERS, SECTIONS, AND SUBSECTIONS}

\vspace{0.15cm}

The report structure shall be generated using the standard \LaTeX{}
commands:

\begin{verbatim}
\chapter{...}
\section{...}
\subsection{...}
\end{verbatim}

Chapter, section, and subsection numbers shall be generated automatically
by \LaTeX{}.

Students shall not manually type chapter or section numbers and shall not
reproduce heading styles using commands such as \verb|\textbf|,
\verb|\fontsize|, or \verb|\centering| instead of the appropriate
structural command.

Excessive subdivision should be avoided. In most cases, the hierarchy
Chapter--Section--Subsection is sufficient.

\vspace{0.3cm}

% ================================================================
% IX. ORGANIZATION OF THE MAIN CONTENT
% ================================================================

\noindent
\textbf{IX.\quad ORGANIZATION OF THE MAIN CONTENT}

\vspace{0.15cm}

The specific content of the report shall depend on the internship
organization and the assigned tasks. A typical structure may include:

\begin{itemize}[leftmargin=0.8cm, itemsep=2pt]
    \item \textbf{Chapter 1. Overview of the Internship Organization};
    \item \textbf{Chapter 2. Internship Activities and Assigned Tasks};
    \item \textbf{Chapter 3. Results, Experience, and Evaluation}.
\end{itemize}

These chapter titles are recommendations only and may be adapted to the
actual internship content.

The report shall reflect work actually performed, observed, analyzed, or
learned during the internship. Unrelated theoretical material should not
be included merely to increase report length.

\vspace{0.3cm}

% ================================================================
% X. TABLE OF CONTENTS
% ================================================================

\noindent
\textbf{X.\quad TABLE OF CONTENTS}

\vspace{0.15cm}

The Table of Contents shall be generated automatically using:

\begin{verbatim}
\tableofcontents
\end{verbatim}

Students shall not create or manually edit page numbers in the Table of
Contents.

The document shall be compiled sufficiently to ensure that all page
numbers and entries are correctly updated.

\vspace{0.3cm}

% ================================================================
% XI. LIST OF SYMBOLS AND ABBREVIATIONS
% ================================================================

\noindent
\textbf{XI.\quad LIST OF SYMBOLS AND ABBREVIATIONS}

\vspace{0.15cm}

Symbols and abbreviations that are used repeatedly in the report shall be
listed in the \textbf{List of Symbols and Abbreviations}.

The list may contain the following columns:

\begin{center}
\begin{tabular}{llll}
    No. & Symbol / Abbreviation & Unit & Definition
\end{tabular}
\end{center}

Only symbols and abbreviations actually used in the final report shall be
retained. Example entries contained in the template shall be removed if
they are not relevant to the student's report.

Each abbreviation should normally be defined when it first appears in
the main text.

\vspace{0.3cm}

% ================================================================
% XII. FIGURES, DIAGRAMS, AND PHOTOGRAPHS
% ================================================================

\noindent
\textbf{XII.\quad FIGURES, DIAGRAMS, AND PHOTOGRAPHS}

\vspace{0.15cm}

Figures shall be clear, legible, and directly relevant to the report
content.

Figures shall be inserted using the standard \verb|figure| environment.

Example:

\begin{verbatim}
\begin{figure}[H]
    \centering
    \includegraphics[width=0.8\textwidth]
                    {figures/company_structure.pdf}
    \caption{Organizational structure of the company}
    \label{fig:company-structure}
\end{figure}
\end{verbatim}

Figure captions shall be placed below the figure.

Each referenced figure shall have a unique \verb|\label| and shall be
referenced in the text using, for example:

\begin{verbatim}
Figure~\ref{fig:company-structure}
\end{verbatim}

Students shall not manually type figure numbers.

Figures reproduced or adapted from external sources shall include an
appropriate citation.

Photographs taken during the internship shall be relevant and sufficiently
clear. Students shall not include confidential, restricted, or proprietary
information without permission from the internship organization.

Vector graphics or high-resolution images are recommended whenever
possible.

\vspace{0.3cm}

% ================================================================
% XIII. TABLES
% ================================================================

\noindent
\textbf{XIII.\quad TABLES}

\vspace{0.15cm}

Tables shall be created using suitable \LaTeX{} environments such as
\verb|table|, \verb|tabular|, \verb|tabularx|, or \verb|longtable|.

Table captions shall be placed above the table.

Example:

\begin{verbatim}
\begin{table}[H]
    \centering
    \caption{Internship schedule}
    \label{tab:internship-schedule}

    \begin{tabular}{lll}
        \toprule
        Week & Task & Result \\
        \midrule
        1 & Orientation & Completed \\
        2 & System study & Completed \\
        \bottomrule
    \end{tabular}
\end{table}
\end{verbatim}

Tables shall be referenced using:

\begin{verbatim}
Table~\ref{tab:internship-schedule}
\end{verbatim}

Students shall not manually type table numbers.

Tables should be created directly in \LaTeX{} rather than inserted as
screenshots whenever practical.

\vspace{0.3cm}

% ================================================================
% XIV. EQUATIONS AND TECHNICAL QUANTITIES
% ================================================================

\noindent
\textbf{XIV.\quad EQUATIONS AND TECHNICAL QUANTITIES}

\vspace{0.15cm}

Mathematical expressions shall be typeset directly using \LaTeX{}.
Images or screenshots of equations shall not be used.

Short expressions may be written in inline mathematical mode, for example:

\begin{verbatim}
$R = 0.7\,\Omega$
\end{verbatim}

Important equations shall use a numbered equation environment.

Example:

\begin{verbatim}
\begin{equation}
    P = UI
    \label{eq:power}
\end{equation}
\end{verbatim}

Equations shall be referenced using:

\begin{verbatim}
Equation~\eqref{eq:power}
\end{verbatim}

Equation numbers shall not be manually entered.

SI units shall be used wherever applicable. The \verb|siunitx| package
should be used where appropriate.

Examples:

\begin{verbatim}
\SI{220}{\volt}
\SI{50}{\hertz}
\SI{2.5}{\ampere}
\end{verbatim}

Notation and units shall be consistent throughout the report.

\vspace{0.3cm}

% ================================================================
% XV. SOURCE CODE
% ================================================================

\noindent
\textbf{XV.\quad SOURCE CODE}

\vspace{0.15cm}

If the internship involves programming, PLCs, microcontrollers,
simulation, or software development, relevant source code may be included
in the report.

Source code shall be presented using the \verb|listings| environment or
another environment defined by the official template.

Example:

\begin{verbatim}
\begin{lstlisting}[language=st]
IF Start = TRUE THEN
    Motor := TRUE;
ELSE
    Motor := FALSE;
END_IF;
\end{lstlisting}
\end{verbatim}

Only code necessary for explaining the work should appear in the main
report. Long source-code listings should be moved to an appendix or
submitted as supplementary files if required.

Source code should not normally be inserted as screenshots.

\vspace{0.3cm}

% ================================================================
% XVI. ABSTRACT
% ================================================================

\noindent
\textbf{XVI.\quad ABSTRACT}

\vspace{0.15cm}

The Abstract shall briefly summarize:

\begin{itemize}[leftmargin=0.8cm, itemsep=2pt]
    \item the internship organization;
    \item the objectives of the internship;
    \item the main tasks performed;
    \item the main results, knowledge, and skills obtained.
\end{itemize}

The Abstract should be concise and should not reproduce the full content
of the report chapter by chapter.

\vspace{0.3cm}

% ================================================================
% XVII. REFERENCES AND CITATIONS
% ================================================================

\noindent
\textbf{XVII.\quad REFERENCES AND CITATIONS}

\vspace{0.15cm}

All books, journal papers, conference papers, standards, websites,
datasheets, manuals, technical documents, figures, tables, or other
external sources used in preparing the report shall be properly cited.

References shall be managed through the bibliography system provided with
the official template.

Bibliographic entries shall normally be stored in:

\begin{verbatim}
References.bib
\end{verbatim}

References shall be cited in the report using:

\begin{verbatim}
\cite{reference-key}
\end{verbatim}

Students shall not manually type reference numbers such as ``[1]'',
``[2]'', etc.

Titles of foreign-language publications shall remain in their original
language in the bibliography.

All borrowed ideas, data, figures, tables, equations, methods, or results
shall be properly acknowledged.

\vspace{0.3cm}

% ================================================================
% XVIII. CROSS-REFERENCING
% ================================================================

\noindent
\textbf{XVIII.\quad CROSS-REFERENCING}

\vspace{0.15cm}

Numbered elements shall be cross-referenced using the standard \LaTeX{}
mechanism.

Use the following commands where appropriate:

\begin{verbatim}
\label{...}
\ref{...}
\eqref{...}
\end{verbatim}

For example:

\begin{verbatim}
As shown in Figure~\ref{fig:system}, ...
\end{verbatim}

Students shall not manually type numbers such as ``Figure 2.3'' when
those numbers can be generated automatically.

\vspace{0.3cm}

% ================================================================
% XIX. PAGE NUMBERING
% ================================================================

\noindent
\textbf{XIX.\quad PAGE NUMBERING}

\vspace{0.15cm}

Cover pages and official evaluation forms at the beginning of the report
shall not display page numbers.

The main body of the report shall use Arabic page numbers beginning from
Chapter 1.

Page numbering shall be generated automatically by the official template.
Students shall not manually enter page numbers.

\vspace{0.3cm}

% ================================================================
% XX. PROJECT FILE ORGANIZATION
% ================================================================

\noindent
\textbf{XX.\quad PROJECT FILE ORGANIZATION}

\vspace{0.15cm}

Students should maintain a clear and organized \LaTeX{} project
structure. A recommended structure is:

\begin{verbatim}
InternshipReport/
|
|-- main.tex
|-- References.bib
|
|-- figures/
|   |-- company.pdf
|   |-- system.png
|   `-- ...
|
`-- chapters/
    |-- chapter1.tex
    |-- chapter2.tex
    `-- chapter3.tex
\end{verbatim}

Figures should normally be stored in the \verb|figures| directory
specified by the official template.

File and directory names should preferably contain Latin letters,
numbers, hyphens, or underscores.

Recommended examples include:

\begin{verbatim}
company_structure.pdf
motor_control.png
chapter_01.tex
\end{verbatim}

Spaces and unnecessary special characters in file names should be
avoided where possible.

\vspace{0.3cm}

% ================================================================
% XXI. DOCUMENT COMPILATION AND FINAL CHECK
% ================================================================

\noindent
\textbf{XXI.\quad DOCUMENT COMPILATION AND FINAL CHECK}

\vspace{0.15cm}

Before submission, the complete \LaTeX{} project shall compile
successfully and produce the final PDF report.

The final document shall not contain unresolved references such as:

\begin{verbatim}
??
[?]
\end{verbatim}

Before generating the final PDF, students shall check:

\begin{itemize}[leftmargin=0.8cm, itemsep=2pt]
    \item the Table of Contents;
    \item page numbering;
    \item chapter, section, and subsection numbering;
    \item figure numbering;
    \item table numbering;
    \item equation numbering;
    \item cross-references;
    \item citations;
    \item bibliography;
    \item font display;
    \item mathematical symbols;
    \item image quality;
    \item completeness of all required forms.
\end{itemize}

The document may be compiled using TeX Live, MiKTeX, or another
compatible \LaTeX{} environment. Regardless of the compilation
environment, the resulting PDF shall conform to the official Faculty
template.

\vspace{0.3cm}

% ================================================================
% XXII. CONTENT THAT STUDENTS MAY EDIT
% ================================================================

\noindent
\textbf{XXII.\quad CONTENT THAT STUDENTS MAY EDIT}

\vspace{0.15cm}

Students may edit or replace:

\begin{itemize}[leftmargin=0.8cm, itemsep=2pt]
    \item student name;
    \item class;
    \item internship organization;
    \item supervisor name;
    \item Abstract;
    \item List of Symbols and Abbreviations;
    \item chapter titles and chapter content;
    \item figures and photographs;
    \item tables;
    \item equations;
    \item citations and references.
\end{itemize}

Students shall not modify the main document formatting settings in the
preamble, including the document class, page geometry, line spacing,
heading formatting, table-of-contents formatting, font configuration,
page layout, or cover-page macros unless specifically instructed.

\vspace{0.3cm}

% ================================================================
% XXIII. FINAL SUBMISSION REQUIREMENTS
% ================================================================

\noindent
\textbf{XXIII.\quad FINAL SUBMISSION REQUIREMENTS}

\vspace{0.15cm}

Before submission, the final report shall:

\begin{itemize}[leftmargin=0.8cm, itemsep=2pt]
    \item compile without errors that affect the report content or presentation;
    \item contain all required forms and report sections;
    \item contain correct student, class, internship organization, and supervisor information;
    \item contain an updated Table of Contents, List of Figures, and List of Tables;
    \item contain no unused example content from the template;
    \item contain no placeholder text such as \textit{Title of Chapter 1},
    \textit{Section 1.1}, or other sample entries;
    \item contain all required citations and references;
    \item contain no unresolved cross-references.
\end{itemize}

If the Faculty requires submission of the \LaTeX{} source, the complete
project shall include all \verb|.tex| files, the bibliography database,
figures, and any other files required to reproduce the submitted PDF.

\vspace{0.4cm}

\noindent
\textbf{Important:}
Before submitting the final internship report, students shall remove or
comment out the following command in the main document:

\begin{verbatim}
\input{Regulations.tex}
\end{verbatim}

For example:

\begin{verbatim}
% \input{Regulations.tex}
\end{verbatim}

This removes the guidance pages from the final submitted report while
preserving the regulations file in the \LaTeX{} project for reference.

} % End 13pt content

\clearpage

\end{verbatim}

This removes the guidance pages from the final submitted report while
preserving the regulations file in the \LaTeX{} project for reference.

} % End 13pt content

\clearpage
 command before submitting
% the final internship report.
%
% =================================================================

\clearpage
\thispagestyle{empty}

\begin{center}
    {\fontsize{16pt}{20pt}\selectfont\bfseries
    SPECIFIC REGULATIONS FOR PREPARING\\
    GRADUATION INTERNSHIP REPORTS USING \LaTeX}
\end{center}

\vspace{0.5cm}

{\fontsize{13pt}{17pt}\selectfont
\setstretch{1.3}
\justifying

\noindent
\textbf{Note:}
This section is provided for guidance only and is included in the main
document using the command \verb|
% =================================================================
% REGULATIONS FOR PREPARING THE GRADUATION INTERNSHIP REPORT
% USING LATEX
% =================================================================
%
% This file is intended to be included from the main document with:
%
%     
% =================================================================
% REGULATIONS FOR PREPARING THE GRADUATION INTERNSHIP REPORT
% USING LATEX
% =================================================================
%
% This file is intended to be included from the main document with:
%
%     \input{Regulations.tex}
%
% IMPORTANT:
% This section is provided for guidance only. Students shall remove
% or comment out the \input{Regulations.tex} command before submitting
% the final internship report.
%
% =================================================================

\clearpage
\thispagestyle{empty}

\begin{center}
    {\fontsize{16pt}{20pt}\selectfont\bfseries
    SPECIFIC REGULATIONS FOR PREPARING\\
    GRADUATION INTERNSHIP REPORTS USING \LaTeX}
\end{center}

\vspace{0.5cm}

{\fontsize{13pt}{17pt}\selectfont
\setstretch{1.3}
\justifying

\noindent
\textbf{Note:}
This section is provided for guidance only and is included in the main
document using the command \verb|\input{Regulations.tex}|.
Students shall remove or comment out this command before submitting the
final internship report.

\vspace{0.4cm}

% ================================================================
% I. GENERAL REQUIREMENTS
% ================================================================

\noindent
\textbf{I.\quad GENERAL REQUIREMENTS}

\vspace{0.15cm}

The Graduation Internship Report shall be prepared using the official
\LaTeX{} template provided by the Faculty. Students shall preserve the
predefined page layout, typography, line spacing, heading styles,
numbering system, figure and table formatting, and bibliography settings
unless otherwise instructed by the Faculty or the supervisor.

The report shall be written in English.

Students are responsible for the academic and technical content of the
report. The official \LaTeX{} template is responsible for standardized
document formatting.

\vspace{0.3cm}

% ================================================================
% II. REPORT STRUCTURE
% ================================================================

\noindent
\textbf{II.\quad REPORT STRUCTURE}

\vspace{0.15cm}

The report shall normally be arranged in the following order:

\begin{enumerate}[label=\arabic*., leftmargin=1.2cm, itemsep=2pt]
    \item Outer Cover;
    \item Comments of the Internship Host Organization;
    \item Supervisor's Evaluation and Graduation Internship Defense Results;
    \item Table of Contents;
    \item List of Symbols and Abbreviations, if applicable;
    \item List of Figures;
    \item List of Tables;
    \item Abstract;
    \item Main Chapters of the Internship Report;
    \item References;
    \item Appendices, if applicable.
\end{enumerate}

The preliminary pages shall not be numbered as chapters. The main report
content shall be organized into consecutively numbered chapters beginning
with Chapter 1.

\vspace{0.3cm}

% ================================================================
% III. PAGE SIZE AND MARGINS
% ================================================================

\noindent
\textbf{III.\quad PAGE SIZE AND MARGINS}

\vspace{0.15cm}

The report shall be formatted on A4 paper with the following margins:

\begin{center}
\begin{tabular}{ll}
    Left margin:   & 3.5 cm \\
    Right margin:  & 2.0 cm \\
    Top margin:    & 2.5 cm \\
    Bottom margin: & 2.5 cm
\end{tabular}
\end{center}

These settings are predefined in the official template and shall not be
modified by students.

\vspace{0.3cm}

% ================================================================
% IV. FONT, FONT SIZE, AND LINE SPACING
% ================================================================

\noindent
\textbf{IV.\quad FONT, FONT SIZE, AND LINE SPACING}

\vspace{0.15cm}

The main text shall use a 13 pt font size with 1.3-line spacing.
A Times-family typeface shall be used in accordance with the official
template.

The text shall be fully justified.

Students shall not use repeated spaces, repeated manual line breaks, or
multiple blank lines to position text. Paragraph and heading spacing shall
be controlled by the template.

\vspace{0.3cm}

% ================================================================
% V. OUTER COVER
% ================================================================

\noindent
\textbf{V.\quad OUTER COVER}

\vspace{0.15cm}

The outer cover shall follow the official format provided in the template
and shall contain the following information:

\begin{itemize}[leftmargin=0.8cm, itemsep=2pt]
    \item Thai Nguyen University of Technology;
    \item Faculty of International Training;
    \item Engineering and Technology Department;
    \item University logo and Faculty logo;
    \item the title \textbf{INTERNSHIP REPORT};
    \item Internship Organization;
    \item Author;
    \item Class;
    \item Supervisor;
    \item place and year of completion.
\end{itemize}

Students shall only replace the required information fields and shall not
modify the cover-page layout unless instructed otherwise.

\vspace{0.3cm}

% ================================================================
% VI. COMMENTS OF THE INTERNSHIP HOST ORGANIZATION
% ================================================================

\noindent
\textbf{VI.\quad COMMENTS OF THE INTERNSHIP HOST ORGANIZATION}

\vspace{0.15cm}

The page entitled \textbf{COMMENTS OF THE INTERNSHIP HOST ORGANIZATION}
is reserved for the internship host organization.

It includes:

\begin{itemize}[leftmargin=0.8cm, itemsep=2pt]
    \item the student's full name;
    \item class;
    \item comments from the host organization;
    \item score;
    \item date;
    \item signature and official seal.
\end{itemize}

Students shall not alter or complete official evaluation fields unless
specifically authorized.

\vspace{0.3cm}

% ================================================================
% VII. SUPERVISOR'S EVALUATION AND DEFENSE RESULTS
% ================================================================

\noindent
\textbf{VII.\quad SUPERVISOR'S EVALUATION AND DEFENSE RESULTS}

\vspace{0.15cm}

The section entitled
\textbf{SUPERVISOR'S EVALUATION AND GRADUATION INTERNSHIP DEFENSE RESULTS}
is reserved for the supervisor and defense committee.

The evaluation may include:

\begin{itemize}[leftmargin=0.8cm, itemsep=2pt]
    \item Internship Progress and Student's Attitude;
    \item Report Content;
    \item Report Presentation;
    \item Other Comments;
    \item Supervisor's Opinion;
    \item Defense Result;
    \item Score.
\end{itemize}

Students shall not enter evaluation scores or defense results before they
are officially determined.

\vspace{0.3cm}

% ================================================================
% VIII. CHAPTERS, SECTIONS, AND SUBSECTIONS
% ================================================================

\noindent
\textbf{VIII.\quad CHAPTERS, SECTIONS, AND SUBSECTIONS}

\vspace{0.15cm}

The report structure shall be generated using the standard \LaTeX{}
commands:

\begin{verbatim}
\chapter{...}
\section{...}
\subsection{...}
\end{verbatim}

Chapter, section, and subsection numbers shall be generated automatically
by \LaTeX{}.

Students shall not manually type chapter or section numbers and shall not
reproduce heading styles using commands such as \verb|\textbf|,
\verb|\fontsize|, or \verb|\centering| instead of the appropriate
structural command.

Excessive subdivision should be avoided. In most cases, the hierarchy
Chapter--Section--Subsection is sufficient.

\vspace{0.3cm}

% ================================================================
% IX. ORGANIZATION OF THE MAIN CONTENT
% ================================================================

\noindent
\textbf{IX.\quad ORGANIZATION OF THE MAIN CONTENT}

\vspace{0.15cm}

The specific content of the report shall depend on the internship
organization and the assigned tasks. A typical structure may include:

\begin{itemize}[leftmargin=0.8cm, itemsep=2pt]
    \item \textbf{Chapter 1. Overview of the Internship Organization};
    \item \textbf{Chapter 2. Internship Activities and Assigned Tasks};
    \item \textbf{Chapter 3. Results, Experience, and Evaluation}.
\end{itemize}

These chapter titles are recommendations only and may be adapted to the
actual internship content.

The report shall reflect work actually performed, observed, analyzed, or
learned during the internship. Unrelated theoretical material should not
be included merely to increase report length.

\vspace{0.3cm}

% ================================================================
% X. TABLE OF CONTENTS
% ================================================================

\noindent
\textbf{X.\quad TABLE OF CONTENTS}

\vspace{0.15cm}

The Table of Contents shall be generated automatically using:

\begin{verbatim}
\tableofcontents
\end{verbatim}

Students shall not create or manually edit page numbers in the Table of
Contents.

The document shall be compiled sufficiently to ensure that all page
numbers and entries are correctly updated.

\vspace{0.3cm}

% ================================================================
% XI. LIST OF SYMBOLS AND ABBREVIATIONS
% ================================================================

\noindent
\textbf{XI.\quad LIST OF SYMBOLS AND ABBREVIATIONS}

\vspace{0.15cm}

Symbols and abbreviations that are used repeatedly in the report shall be
listed in the \textbf{List of Symbols and Abbreviations}.

The list may contain the following columns:

\begin{center}
\begin{tabular}{llll}
    No. & Symbol / Abbreviation & Unit & Definition
\end{tabular}
\end{center}

Only symbols and abbreviations actually used in the final report shall be
retained. Example entries contained in the template shall be removed if
they are not relevant to the student's report.

Each abbreviation should normally be defined when it first appears in
the main text.

\vspace{0.3cm}

% ================================================================
% XII. FIGURES, DIAGRAMS, AND PHOTOGRAPHS
% ================================================================

\noindent
\textbf{XII.\quad FIGURES, DIAGRAMS, AND PHOTOGRAPHS}

\vspace{0.15cm}

Figures shall be clear, legible, and directly relevant to the report
content.

Figures shall be inserted using the standard \verb|figure| environment.

Example:

\begin{verbatim}
\begin{figure}[H]
    \centering
    \includegraphics[width=0.8\textwidth]
                    {figures/company_structure.pdf}
    \caption{Organizational structure of the company}
    \label{fig:company-structure}
\end{figure}
\end{verbatim}

Figure captions shall be placed below the figure.

Each referenced figure shall have a unique \verb|\label| and shall be
referenced in the text using, for example:

\begin{verbatim}
Figure~\ref{fig:company-structure}
\end{verbatim}

Students shall not manually type figure numbers.

Figures reproduced or adapted from external sources shall include an
appropriate citation.

Photographs taken during the internship shall be relevant and sufficiently
clear. Students shall not include confidential, restricted, or proprietary
information without permission from the internship organization.

Vector graphics or high-resolution images are recommended whenever
possible.

\vspace{0.3cm}

% ================================================================
% XIII. TABLES
% ================================================================

\noindent
\textbf{XIII.\quad TABLES}

\vspace{0.15cm}

Tables shall be created using suitable \LaTeX{} environments such as
\verb|table|, \verb|tabular|, \verb|tabularx|, or \verb|longtable|.

Table captions shall be placed above the table.

Example:

\begin{verbatim}
\begin{table}[H]
    \centering
    \caption{Internship schedule}
    \label{tab:internship-schedule}

    \begin{tabular}{lll}
        \toprule
        Week & Task & Result \\
        \midrule
        1 & Orientation & Completed \\
        2 & System study & Completed \\
        \bottomrule
    \end{tabular}
\end{table}
\end{verbatim}

Tables shall be referenced using:

\begin{verbatim}
Table~\ref{tab:internship-schedule}
\end{verbatim}

Students shall not manually type table numbers.

Tables should be created directly in \LaTeX{} rather than inserted as
screenshots whenever practical.

\vspace{0.3cm}

% ================================================================
% XIV. EQUATIONS AND TECHNICAL QUANTITIES
% ================================================================

\noindent
\textbf{XIV.\quad EQUATIONS AND TECHNICAL QUANTITIES}

\vspace{0.15cm}

Mathematical expressions shall be typeset directly using \LaTeX{}.
Images or screenshots of equations shall not be used.

Short expressions may be written in inline mathematical mode, for example:

\begin{verbatim}
$R = 0.7\,\Omega$
\end{verbatim}

Important equations shall use a numbered equation environment.

Example:

\begin{verbatim}
\begin{equation}
    P = UI
    \label{eq:power}
\end{equation}
\end{verbatim}

Equations shall be referenced using:

\begin{verbatim}
Equation~\eqref{eq:power}
\end{verbatim}

Equation numbers shall not be manually entered.

SI units shall be used wherever applicable. The \verb|siunitx| package
should be used where appropriate.

Examples:

\begin{verbatim}
\SI{220}{\volt}
\SI{50}{\hertz}
\SI{2.5}{\ampere}
\end{verbatim}

Notation and units shall be consistent throughout the report.

\vspace{0.3cm}

% ================================================================
% XV. SOURCE CODE
% ================================================================

\noindent
\textbf{XV.\quad SOURCE CODE}

\vspace{0.15cm}

If the internship involves programming, PLCs, microcontrollers,
simulation, or software development, relevant source code may be included
in the report.

Source code shall be presented using the \verb|listings| environment or
another environment defined by the official template.

Example:

\begin{verbatim}
\begin{lstlisting}[language=st]
IF Start = TRUE THEN
    Motor := TRUE;
ELSE
    Motor := FALSE;
END_IF;
\end{lstlisting}
\end{verbatim}

Only code necessary for explaining the work should appear in the main
report. Long source-code listings should be moved to an appendix or
submitted as supplementary files if required.

Source code should not normally be inserted as screenshots.

\vspace{0.3cm}

% ================================================================
% XVI. ABSTRACT
% ================================================================

\noindent
\textbf{XVI.\quad ABSTRACT}

\vspace{0.15cm}

The Abstract shall briefly summarize:

\begin{itemize}[leftmargin=0.8cm, itemsep=2pt]
    \item the internship organization;
    \item the objectives of the internship;
    \item the main tasks performed;
    \item the main results, knowledge, and skills obtained.
\end{itemize}

The Abstract should be concise and should not reproduce the full content
of the report chapter by chapter.

\vspace{0.3cm}

% ================================================================
% XVII. REFERENCES AND CITATIONS
% ================================================================

\noindent
\textbf{XVII.\quad REFERENCES AND CITATIONS}

\vspace{0.15cm}

All books, journal papers, conference papers, standards, websites,
datasheets, manuals, technical documents, figures, tables, or other
external sources used in preparing the report shall be properly cited.

References shall be managed through the bibliography system provided with
the official template.

Bibliographic entries shall normally be stored in:

\begin{verbatim}
References.bib
\end{verbatim}

References shall be cited in the report using:

\begin{verbatim}
\cite{reference-key}
\end{verbatim}

Students shall not manually type reference numbers such as ``[1]'',
``[2]'', etc.

Titles of foreign-language publications shall remain in their original
language in the bibliography.

All borrowed ideas, data, figures, tables, equations, methods, or results
shall be properly acknowledged.

\vspace{0.3cm}

% ================================================================
% XVIII. CROSS-REFERENCING
% ================================================================

\noindent
\textbf{XVIII.\quad CROSS-REFERENCING}

\vspace{0.15cm}

Numbered elements shall be cross-referenced using the standard \LaTeX{}
mechanism.

Use the following commands where appropriate:

\begin{verbatim}
\label{...}
\ref{...}
\eqref{...}
\end{verbatim}

For example:

\begin{verbatim}
As shown in Figure~\ref{fig:system}, ...
\end{verbatim}

Students shall not manually type numbers such as ``Figure 2.3'' when
those numbers can be generated automatically.

\vspace{0.3cm}

% ================================================================
% XIX. PAGE NUMBERING
% ================================================================

\noindent
\textbf{XIX.\quad PAGE NUMBERING}

\vspace{0.15cm}

Cover pages and official evaluation forms at the beginning of the report
shall not display page numbers.

The main body of the report shall use Arabic page numbers beginning from
Chapter 1.

Page numbering shall be generated automatically by the official template.
Students shall not manually enter page numbers.

\vspace{0.3cm}

% ================================================================
% XX. PROJECT FILE ORGANIZATION
% ================================================================

\noindent
\textbf{XX.\quad PROJECT FILE ORGANIZATION}

\vspace{0.15cm}

Students should maintain a clear and organized \LaTeX{} project
structure. A recommended structure is:

\begin{verbatim}
InternshipReport/
|
|-- main.tex
|-- References.bib
|
|-- figures/
|   |-- company.pdf
|   |-- system.png
|   `-- ...
|
`-- chapters/
    |-- chapter1.tex
    |-- chapter2.tex
    `-- chapter3.tex
\end{verbatim}

Figures should normally be stored in the \verb|figures| directory
specified by the official template.

File and directory names should preferably contain Latin letters,
numbers, hyphens, or underscores.

Recommended examples include:

\begin{verbatim}
company_structure.pdf
motor_control.png
chapter_01.tex
\end{verbatim}

Spaces and unnecessary special characters in file names should be
avoided where possible.

\vspace{0.3cm}

% ================================================================
% XXI. DOCUMENT COMPILATION AND FINAL CHECK
% ================================================================

\noindent
\textbf{XXI.\quad DOCUMENT COMPILATION AND FINAL CHECK}

\vspace{0.15cm}

Before submission, the complete \LaTeX{} project shall compile
successfully and produce the final PDF report.

The final document shall not contain unresolved references such as:

\begin{verbatim}
??
[?]
\end{verbatim}

Before generating the final PDF, students shall check:

\begin{itemize}[leftmargin=0.8cm, itemsep=2pt]
    \item the Table of Contents;
    \item page numbering;
    \item chapter, section, and subsection numbering;
    \item figure numbering;
    \item table numbering;
    \item equation numbering;
    \item cross-references;
    \item citations;
    \item bibliography;
    \item font display;
    \item mathematical symbols;
    \item image quality;
    \item completeness of all required forms.
\end{itemize}

The document may be compiled using TeX Live, MiKTeX, or another
compatible \LaTeX{} environment. Regardless of the compilation
environment, the resulting PDF shall conform to the official Faculty
template.

\vspace{0.3cm}

% ================================================================
% XXII. CONTENT THAT STUDENTS MAY EDIT
% ================================================================

\noindent
\textbf{XXII.\quad CONTENT THAT STUDENTS MAY EDIT}

\vspace{0.15cm}

Students may edit or replace:

\begin{itemize}[leftmargin=0.8cm, itemsep=2pt]
    \item student name;
    \item class;
    \item internship organization;
    \item supervisor name;
    \item Abstract;
    \item List of Symbols and Abbreviations;
    \item chapter titles and chapter content;
    \item figures and photographs;
    \item tables;
    \item equations;
    \item citations and references.
\end{itemize}

Students shall not modify the main document formatting settings in the
preamble, including the document class, page geometry, line spacing,
heading formatting, table-of-contents formatting, font configuration,
page layout, or cover-page macros unless specifically instructed.

\vspace{0.3cm}

% ================================================================
% XXIII. FINAL SUBMISSION REQUIREMENTS
% ================================================================

\noindent
\textbf{XXIII.\quad FINAL SUBMISSION REQUIREMENTS}

\vspace{0.15cm}

Before submission, the final report shall:

\begin{itemize}[leftmargin=0.8cm, itemsep=2pt]
    \item compile without errors that affect the report content or presentation;
    \item contain all required forms and report sections;
    \item contain correct student, class, internship organization, and supervisor information;
    \item contain an updated Table of Contents, List of Figures, and List of Tables;
    \item contain no unused example content from the template;
    \item contain no placeholder text such as \textit{Title of Chapter 1},
    \textit{Section 1.1}, or other sample entries;
    \item contain all required citations and references;
    \item contain no unresolved cross-references.
\end{itemize}

If the Faculty requires submission of the \LaTeX{} source, the complete
project shall include all \verb|.tex| files, the bibliography database,
figures, and any other files required to reproduce the submitted PDF.

\vspace{0.4cm}

\noindent
\textbf{Important:}
Before submitting the final internship report, students shall remove or
comment out the following command in the main document:

\begin{verbatim}
\input{Regulations.tex}
\end{verbatim}

For example:

\begin{verbatim}
% \input{Regulations.tex}
\end{verbatim}

This removes the guidance pages from the final submitted report while
preserving the regulations file in the \LaTeX{} project for reference.

} % End 13pt content

\clearpage

%
% IMPORTANT:
% This section is provided for guidance only. Students shall remove
% or comment out the 
% =================================================================
% REGULATIONS FOR PREPARING THE GRADUATION INTERNSHIP REPORT
% USING LATEX
% =================================================================
%
% This file is intended to be included from the main document with:
%
%     \input{Regulations.tex}
%
% IMPORTANT:
% This section is provided for guidance only. Students shall remove
% or comment out the \input{Regulations.tex} command before submitting
% the final internship report.
%
% =================================================================

\clearpage
\thispagestyle{empty}

\begin{center}
    {\fontsize{16pt}{20pt}\selectfont\bfseries
    SPECIFIC REGULATIONS FOR PREPARING\\
    GRADUATION INTERNSHIP REPORTS USING \LaTeX}
\end{center}

\vspace{0.5cm}

{\fontsize{13pt}{17pt}\selectfont
\setstretch{1.3}
\justifying

\noindent
\textbf{Note:}
This section is provided for guidance only and is included in the main
document using the command \verb|\input{Regulations.tex}|.
Students shall remove or comment out this command before submitting the
final internship report.

\vspace{0.4cm}

% ================================================================
% I. GENERAL REQUIREMENTS
% ================================================================

\noindent
\textbf{I.\quad GENERAL REQUIREMENTS}

\vspace{0.15cm}

The Graduation Internship Report shall be prepared using the official
\LaTeX{} template provided by the Faculty. Students shall preserve the
predefined page layout, typography, line spacing, heading styles,
numbering system, figure and table formatting, and bibliography settings
unless otherwise instructed by the Faculty or the supervisor.

The report shall be written in English.

Students are responsible for the academic and technical content of the
report. The official \LaTeX{} template is responsible for standardized
document formatting.

\vspace{0.3cm}

% ================================================================
% II. REPORT STRUCTURE
% ================================================================

\noindent
\textbf{II.\quad REPORT STRUCTURE}

\vspace{0.15cm}

The report shall normally be arranged in the following order:

\begin{enumerate}[label=\arabic*., leftmargin=1.2cm, itemsep=2pt]
    \item Outer Cover;
    \item Comments of the Internship Host Organization;
    \item Supervisor's Evaluation and Graduation Internship Defense Results;
    \item Table of Contents;
    \item List of Symbols and Abbreviations, if applicable;
    \item List of Figures;
    \item List of Tables;
    \item Abstract;
    \item Main Chapters of the Internship Report;
    \item References;
    \item Appendices, if applicable.
\end{enumerate}

The preliminary pages shall not be numbered as chapters. The main report
content shall be organized into consecutively numbered chapters beginning
with Chapter 1.

\vspace{0.3cm}

% ================================================================
% III. PAGE SIZE AND MARGINS
% ================================================================

\noindent
\textbf{III.\quad PAGE SIZE AND MARGINS}

\vspace{0.15cm}

The report shall be formatted on A4 paper with the following margins:

\begin{center}
\begin{tabular}{ll}
    Left margin:   & 3.5 cm \\
    Right margin:  & 2.0 cm \\
    Top margin:    & 2.5 cm \\
    Bottom margin: & 2.5 cm
\end{tabular}
\end{center}

These settings are predefined in the official template and shall not be
modified by students.

\vspace{0.3cm}

% ================================================================
% IV. FONT, FONT SIZE, AND LINE SPACING
% ================================================================

\noindent
\textbf{IV.\quad FONT, FONT SIZE, AND LINE SPACING}

\vspace{0.15cm}

The main text shall use a 13 pt font size with 1.3-line spacing.
A Times-family typeface shall be used in accordance with the official
template.

The text shall be fully justified.

Students shall not use repeated spaces, repeated manual line breaks, or
multiple blank lines to position text. Paragraph and heading spacing shall
be controlled by the template.

\vspace{0.3cm}

% ================================================================
% V. OUTER COVER
% ================================================================

\noindent
\textbf{V.\quad OUTER COVER}

\vspace{0.15cm}

The outer cover shall follow the official format provided in the template
and shall contain the following information:

\begin{itemize}[leftmargin=0.8cm, itemsep=2pt]
    \item Thai Nguyen University of Technology;
    \item Faculty of International Training;
    \item Engineering and Technology Department;
    \item University logo and Faculty logo;
    \item the title \textbf{INTERNSHIP REPORT};
    \item Internship Organization;
    \item Author;
    \item Class;
    \item Supervisor;
    \item place and year of completion.
\end{itemize}

Students shall only replace the required information fields and shall not
modify the cover-page layout unless instructed otherwise.

\vspace{0.3cm}

% ================================================================
% VI. COMMENTS OF THE INTERNSHIP HOST ORGANIZATION
% ================================================================

\noindent
\textbf{VI.\quad COMMENTS OF THE INTERNSHIP HOST ORGANIZATION}

\vspace{0.15cm}

The page entitled \textbf{COMMENTS OF THE INTERNSHIP HOST ORGANIZATION}
is reserved for the internship host organization.

It includes:

\begin{itemize}[leftmargin=0.8cm, itemsep=2pt]
    \item the student's full name;
    \item class;
    \item comments from the host organization;
    \item score;
    \item date;
    \item signature and official seal.
\end{itemize}

Students shall not alter or complete official evaluation fields unless
specifically authorized.

\vspace{0.3cm}

% ================================================================
% VII. SUPERVISOR'S EVALUATION AND DEFENSE RESULTS
% ================================================================

\noindent
\textbf{VII.\quad SUPERVISOR'S EVALUATION AND DEFENSE RESULTS}

\vspace{0.15cm}

The section entitled
\textbf{SUPERVISOR'S EVALUATION AND GRADUATION INTERNSHIP DEFENSE RESULTS}
is reserved for the supervisor and defense committee.

The evaluation may include:

\begin{itemize}[leftmargin=0.8cm, itemsep=2pt]
    \item Internship Progress and Student's Attitude;
    \item Report Content;
    \item Report Presentation;
    \item Other Comments;
    \item Supervisor's Opinion;
    \item Defense Result;
    \item Score.
\end{itemize}

Students shall not enter evaluation scores or defense results before they
are officially determined.

\vspace{0.3cm}

% ================================================================
% VIII. CHAPTERS, SECTIONS, AND SUBSECTIONS
% ================================================================

\noindent
\textbf{VIII.\quad CHAPTERS, SECTIONS, AND SUBSECTIONS}

\vspace{0.15cm}

The report structure shall be generated using the standard \LaTeX{}
commands:

\begin{verbatim}
\chapter{...}
\section{...}
\subsection{...}
\end{verbatim}

Chapter, section, and subsection numbers shall be generated automatically
by \LaTeX{}.

Students shall not manually type chapter or section numbers and shall not
reproduce heading styles using commands such as \verb|\textbf|,
\verb|\fontsize|, or \verb|\centering| instead of the appropriate
structural command.

Excessive subdivision should be avoided. In most cases, the hierarchy
Chapter--Section--Subsection is sufficient.

\vspace{0.3cm}

% ================================================================
% IX. ORGANIZATION OF THE MAIN CONTENT
% ================================================================

\noindent
\textbf{IX.\quad ORGANIZATION OF THE MAIN CONTENT}

\vspace{0.15cm}

The specific content of the report shall depend on the internship
organization and the assigned tasks. A typical structure may include:

\begin{itemize}[leftmargin=0.8cm, itemsep=2pt]
    \item \textbf{Chapter 1. Overview of the Internship Organization};
    \item \textbf{Chapter 2. Internship Activities and Assigned Tasks};
    \item \textbf{Chapter 3. Results, Experience, and Evaluation}.
\end{itemize}

These chapter titles are recommendations only and may be adapted to the
actual internship content.

The report shall reflect work actually performed, observed, analyzed, or
learned during the internship. Unrelated theoretical material should not
be included merely to increase report length.

\vspace{0.3cm}

% ================================================================
% X. TABLE OF CONTENTS
% ================================================================

\noindent
\textbf{X.\quad TABLE OF CONTENTS}

\vspace{0.15cm}

The Table of Contents shall be generated automatically using:

\begin{verbatim}
\tableofcontents
\end{verbatim}

Students shall not create or manually edit page numbers in the Table of
Contents.

The document shall be compiled sufficiently to ensure that all page
numbers and entries are correctly updated.

\vspace{0.3cm}

% ================================================================
% XI. LIST OF SYMBOLS AND ABBREVIATIONS
% ================================================================

\noindent
\textbf{XI.\quad LIST OF SYMBOLS AND ABBREVIATIONS}

\vspace{0.15cm}

Symbols and abbreviations that are used repeatedly in the report shall be
listed in the \textbf{List of Symbols and Abbreviations}.

The list may contain the following columns:

\begin{center}
\begin{tabular}{llll}
    No. & Symbol / Abbreviation & Unit & Definition
\end{tabular}
\end{center}

Only symbols and abbreviations actually used in the final report shall be
retained. Example entries contained in the template shall be removed if
they are not relevant to the student's report.

Each abbreviation should normally be defined when it first appears in
the main text.

\vspace{0.3cm}

% ================================================================
% XII. FIGURES, DIAGRAMS, AND PHOTOGRAPHS
% ================================================================

\noindent
\textbf{XII.\quad FIGURES, DIAGRAMS, AND PHOTOGRAPHS}

\vspace{0.15cm}

Figures shall be clear, legible, and directly relevant to the report
content.

Figures shall be inserted using the standard \verb|figure| environment.

Example:

\begin{verbatim}
\begin{figure}[H]
    \centering
    \includegraphics[width=0.8\textwidth]
                    {figures/company_structure.pdf}
    \caption{Organizational structure of the company}
    \label{fig:company-structure}
\end{figure}
\end{verbatim}

Figure captions shall be placed below the figure.

Each referenced figure shall have a unique \verb|\label| and shall be
referenced in the text using, for example:

\begin{verbatim}
Figure~\ref{fig:company-structure}
\end{verbatim}

Students shall not manually type figure numbers.

Figures reproduced or adapted from external sources shall include an
appropriate citation.

Photographs taken during the internship shall be relevant and sufficiently
clear. Students shall not include confidential, restricted, or proprietary
information without permission from the internship organization.

Vector graphics or high-resolution images are recommended whenever
possible.

\vspace{0.3cm}

% ================================================================
% XIII. TABLES
% ================================================================

\noindent
\textbf{XIII.\quad TABLES}

\vspace{0.15cm}

Tables shall be created using suitable \LaTeX{} environments such as
\verb|table|, \verb|tabular|, \verb|tabularx|, or \verb|longtable|.

Table captions shall be placed above the table.

Example:

\begin{verbatim}
\begin{table}[H]
    \centering
    \caption{Internship schedule}
    \label{tab:internship-schedule}

    \begin{tabular}{lll}
        \toprule
        Week & Task & Result \\
        \midrule
        1 & Orientation & Completed \\
        2 & System study & Completed \\
        \bottomrule
    \end{tabular}
\end{table}
\end{verbatim}

Tables shall be referenced using:

\begin{verbatim}
Table~\ref{tab:internship-schedule}
\end{verbatim}

Students shall not manually type table numbers.

Tables should be created directly in \LaTeX{} rather than inserted as
screenshots whenever practical.

\vspace{0.3cm}

% ================================================================
% XIV. EQUATIONS AND TECHNICAL QUANTITIES
% ================================================================

\noindent
\textbf{XIV.\quad EQUATIONS AND TECHNICAL QUANTITIES}

\vspace{0.15cm}

Mathematical expressions shall be typeset directly using \LaTeX{}.
Images or screenshots of equations shall not be used.

Short expressions may be written in inline mathematical mode, for example:

\begin{verbatim}
$R = 0.7\,\Omega$
\end{verbatim}

Important equations shall use a numbered equation environment.

Example:

\begin{verbatim}
\begin{equation}
    P = UI
    \label{eq:power}
\end{equation}
\end{verbatim}

Equations shall be referenced using:

\begin{verbatim}
Equation~\eqref{eq:power}
\end{verbatim}

Equation numbers shall not be manually entered.

SI units shall be used wherever applicable. The \verb|siunitx| package
should be used where appropriate.

Examples:

\begin{verbatim}
\SI{220}{\volt}
\SI{50}{\hertz}
\SI{2.5}{\ampere}
\end{verbatim}

Notation and units shall be consistent throughout the report.

\vspace{0.3cm}

% ================================================================
% XV. SOURCE CODE
% ================================================================

\noindent
\textbf{XV.\quad SOURCE CODE}

\vspace{0.15cm}

If the internship involves programming, PLCs, microcontrollers,
simulation, or software development, relevant source code may be included
in the report.

Source code shall be presented using the \verb|listings| environment or
another environment defined by the official template.

Example:

\begin{verbatim}
\begin{lstlisting}[language=st]
IF Start = TRUE THEN
    Motor := TRUE;
ELSE
    Motor := FALSE;
END_IF;
\end{lstlisting}
\end{verbatim}

Only code necessary for explaining the work should appear in the main
report. Long source-code listings should be moved to an appendix or
submitted as supplementary files if required.

Source code should not normally be inserted as screenshots.

\vspace{0.3cm}

% ================================================================
% XVI. ABSTRACT
% ================================================================

\noindent
\textbf{XVI.\quad ABSTRACT}

\vspace{0.15cm}

The Abstract shall briefly summarize:

\begin{itemize}[leftmargin=0.8cm, itemsep=2pt]
    \item the internship organization;
    \item the objectives of the internship;
    \item the main tasks performed;
    \item the main results, knowledge, and skills obtained.
\end{itemize}

The Abstract should be concise and should not reproduce the full content
of the report chapter by chapter.

\vspace{0.3cm}

% ================================================================
% XVII. REFERENCES AND CITATIONS
% ================================================================

\noindent
\textbf{XVII.\quad REFERENCES AND CITATIONS}

\vspace{0.15cm}

All books, journal papers, conference papers, standards, websites,
datasheets, manuals, technical documents, figures, tables, or other
external sources used in preparing the report shall be properly cited.

References shall be managed through the bibliography system provided with
the official template.

Bibliographic entries shall normally be stored in:

\begin{verbatim}
References.bib
\end{verbatim}

References shall be cited in the report using:

\begin{verbatim}
\cite{reference-key}
\end{verbatim}

Students shall not manually type reference numbers such as ``[1]'',
``[2]'', etc.

Titles of foreign-language publications shall remain in their original
language in the bibliography.

All borrowed ideas, data, figures, tables, equations, methods, or results
shall be properly acknowledged.

\vspace{0.3cm}

% ================================================================
% XVIII. CROSS-REFERENCING
% ================================================================

\noindent
\textbf{XVIII.\quad CROSS-REFERENCING}

\vspace{0.15cm}

Numbered elements shall be cross-referenced using the standard \LaTeX{}
mechanism.

Use the following commands where appropriate:

\begin{verbatim}
\label{...}
\ref{...}
\eqref{...}
\end{verbatim}

For example:

\begin{verbatim}
As shown in Figure~\ref{fig:system}, ...
\end{verbatim}

Students shall not manually type numbers such as ``Figure 2.3'' when
those numbers can be generated automatically.

\vspace{0.3cm}

% ================================================================
% XIX. PAGE NUMBERING
% ================================================================

\noindent
\textbf{XIX.\quad PAGE NUMBERING}

\vspace{0.15cm}

Cover pages and official evaluation forms at the beginning of the report
shall not display page numbers.

The main body of the report shall use Arabic page numbers beginning from
Chapter 1.

Page numbering shall be generated automatically by the official template.
Students shall not manually enter page numbers.

\vspace{0.3cm}

% ================================================================
% XX. PROJECT FILE ORGANIZATION
% ================================================================

\noindent
\textbf{XX.\quad PROJECT FILE ORGANIZATION}

\vspace{0.15cm}

Students should maintain a clear and organized \LaTeX{} project
structure. A recommended structure is:

\begin{verbatim}
InternshipReport/
|
|-- main.tex
|-- References.bib
|
|-- figures/
|   |-- company.pdf
|   |-- system.png
|   `-- ...
|
`-- chapters/
    |-- chapter1.tex
    |-- chapter2.tex
    `-- chapter3.tex
\end{verbatim}

Figures should normally be stored in the \verb|figures| directory
specified by the official template.

File and directory names should preferably contain Latin letters,
numbers, hyphens, or underscores.

Recommended examples include:

\begin{verbatim}
company_structure.pdf
motor_control.png
chapter_01.tex
\end{verbatim}

Spaces and unnecessary special characters in file names should be
avoided where possible.

\vspace{0.3cm}

% ================================================================
% XXI. DOCUMENT COMPILATION AND FINAL CHECK
% ================================================================

\noindent
\textbf{XXI.\quad DOCUMENT COMPILATION AND FINAL CHECK}

\vspace{0.15cm}

Before submission, the complete \LaTeX{} project shall compile
successfully and produce the final PDF report.

The final document shall not contain unresolved references such as:

\begin{verbatim}
??
[?]
\end{verbatim}

Before generating the final PDF, students shall check:

\begin{itemize}[leftmargin=0.8cm, itemsep=2pt]
    \item the Table of Contents;
    \item page numbering;
    \item chapter, section, and subsection numbering;
    \item figure numbering;
    \item table numbering;
    \item equation numbering;
    \item cross-references;
    \item citations;
    \item bibliography;
    \item font display;
    \item mathematical symbols;
    \item image quality;
    \item completeness of all required forms.
\end{itemize}

The document may be compiled using TeX Live, MiKTeX, or another
compatible \LaTeX{} environment. Regardless of the compilation
environment, the resulting PDF shall conform to the official Faculty
template.

\vspace{0.3cm}

% ================================================================
% XXII. CONTENT THAT STUDENTS MAY EDIT
% ================================================================

\noindent
\textbf{XXII.\quad CONTENT THAT STUDENTS MAY EDIT}

\vspace{0.15cm}

Students may edit or replace:

\begin{itemize}[leftmargin=0.8cm, itemsep=2pt]
    \item student name;
    \item class;
    \item internship organization;
    \item supervisor name;
    \item Abstract;
    \item List of Symbols and Abbreviations;
    \item chapter titles and chapter content;
    \item figures and photographs;
    \item tables;
    \item equations;
    \item citations and references.
\end{itemize}

Students shall not modify the main document formatting settings in the
preamble, including the document class, page geometry, line spacing,
heading formatting, table-of-contents formatting, font configuration,
page layout, or cover-page macros unless specifically instructed.

\vspace{0.3cm}

% ================================================================
% XXIII. FINAL SUBMISSION REQUIREMENTS
% ================================================================

\noindent
\textbf{XXIII.\quad FINAL SUBMISSION REQUIREMENTS}

\vspace{0.15cm}

Before submission, the final report shall:

\begin{itemize}[leftmargin=0.8cm, itemsep=2pt]
    \item compile without errors that affect the report content or presentation;
    \item contain all required forms and report sections;
    \item contain correct student, class, internship organization, and supervisor information;
    \item contain an updated Table of Contents, List of Figures, and List of Tables;
    \item contain no unused example content from the template;
    \item contain no placeholder text such as \textit{Title of Chapter 1},
    \textit{Section 1.1}, or other sample entries;
    \item contain all required citations and references;
    \item contain no unresolved cross-references.
\end{itemize}

If the Faculty requires submission of the \LaTeX{} source, the complete
project shall include all \verb|.tex| files, the bibliography database,
figures, and any other files required to reproduce the submitted PDF.

\vspace{0.4cm}

\noindent
\textbf{Important:}
Before submitting the final internship report, students shall remove or
comment out the following command in the main document:

\begin{verbatim}
\input{Regulations.tex}
\end{verbatim}

For example:

\begin{verbatim}
% \input{Regulations.tex}
\end{verbatim}

This removes the guidance pages from the final submitted report while
preserving the regulations file in the \LaTeX{} project for reference.

} % End 13pt content

\clearpage
 command before submitting
% the final internship report.
%
% =================================================================

\clearpage
\thispagestyle{empty}

\begin{center}
    {\fontsize{16pt}{20pt}\selectfont\bfseries
    SPECIFIC REGULATIONS FOR PREPARING\\
    GRADUATION INTERNSHIP REPORTS USING \LaTeX}
\end{center}

\vspace{0.5cm}

{\fontsize{13pt}{17pt}\selectfont
\setstretch{1.3}
\justifying

\noindent
\textbf{Note:}
This section is provided for guidance only and is included in the main
document using the command \verb|
% =================================================================
% REGULATIONS FOR PREPARING THE GRADUATION INTERNSHIP REPORT
% USING LATEX
% =================================================================
%
% This file is intended to be included from the main document with:
%
%     \input{Regulations.tex}
%
% IMPORTANT:
% This section is provided for guidance only. Students shall remove
% or comment out the \input{Regulations.tex} command before submitting
% the final internship report.
%
% =================================================================

\clearpage
\thispagestyle{empty}

\begin{center}
    {\fontsize{16pt}{20pt}\selectfont\bfseries
    SPECIFIC REGULATIONS FOR PREPARING\\
    GRADUATION INTERNSHIP REPORTS USING \LaTeX}
\end{center}

\vspace{0.5cm}

{\fontsize{13pt}{17pt}\selectfont
\setstretch{1.3}
\justifying

\noindent
\textbf{Note:}
This section is provided for guidance only and is included in the main
document using the command \verb|\input{Regulations.tex}|.
Students shall remove or comment out this command before submitting the
final internship report.

\vspace{0.4cm}

% ================================================================
% I. GENERAL REQUIREMENTS
% ================================================================

\noindent
\textbf{I.\quad GENERAL REQUIREMENTS}

\vspace{0.15cm}

The Graduation Internship Report shall be prepared using the official
\LaTeX{} template provided by the Faculty. Students shall preserve the
predefined page layout, typography, line spacing, heading styles,
numbering system, figure and table formatting, and bibliography settings
unless otherwise instructed by the Faculty or the supervisor.

The report shall be written in English.

Students are responsible for the academic and technical content of the
report. The official \LaTeX{} template is responsible for standardized
document formatting.

\vspace{0.3cm}

% ================================================================
% II. REPORT STRUCTURE
% ================================================================

\noindent
\textbf{II.\quad REPORT STRUCTURE}

\vspace{0.15cm}

The report shall normally be arranged in the following order:

\begin{enumerate}[label=\arabic*., leftmargin=1.2cm, itemsep=2pt]
    \item Outer Cover;
    \item Comments of the Internship Host Organization;
    \item Supervisor's Evaluation and Graduation Internship Defense Results;
    \item Table of Contents;
    \item List of Symbols and Abbreviations, if applicable;
    \item List of Figures;
    \item List of Tables;
    \item Abstract;
    \item Main Chapters of the Internship Report;
    \item References;
    \item Appendices, if applicable.
\end{enumerate}

The preliminary pages shall not be numbered as chapters. The main report
content shall be organized into consecutively numbered chapters beginning
with Chapter 1.

\vspace{0.3cm}

% ================================================================
% III. PAGE SIZE AND MARGINS
% ================================================================

\noindent
\textbf{III.\quad PAGE SIZE AND MARGINS}

\vspace{0.15cm}

The report shall be formatted on A4 paper with the following margins:

\begin{center}
\begin{tabular}{ll}
    Left margin:   & 3.5 cm \\
    Right margin:  & 2.0 cm \\
    Top margin:    & 2.5 cm \\
    Bottom margin: & 2.5 cm
\end{tabular}
\end{center}

These settings are predefined in the official template and shall not be
modified by students.

\vspace{0.3cm}

% ================================================================
% IV. FONT, FONT SIZE, AND LINE SPACING
% ================================================================

\noindent
\textbf{IV.\quad FONT, FONT SIZE, AND LINE SPACING}

\vspace{0.15cm}

The main text shall use a 13 pt font size with 1.3-line spacing.
A Times-family typeface shall be used in accordance with the official
template.

The text shall be fully justified.

Students shall not use repeated spaces, repeated manual line breaks, or
multiple blank lines to position text. Paragraph and heading spacing shall
be controlled by the template.

\vspace{0.3cm}

% ================================================================
% V. OUTER COVER
% ================================================================

\noindent
\textbf{V.\quad OUTER COVER}

\vspace{0.15cm}

The outer cover shall follow the official format provided in the template
and shall contain the following information:

\begin{itemize}[leftmargin=0.8cm, itemsep=2pt]
    \item Thai Nguyen University of Technology;
    \item Faculty of International Training;
    \item Engineering and Technology Department;
    \item University logo and Faculty logo;
    \item the title \textbf{INTERNSHIP REPORT};
    \item Internship Organization;
    \item Author;
    \item Class;
    \item Supervisor;
    \item place and year of completion.
\end{itemize}

Students shall only replace the required information fields and shall not
modify the cover-page layout unless instructed otherwise.

\vspace{0.3cm}

% ================================================================
% VI. COMMENTS OF THE INTERNSHIP HOST ORGANIZATION
% ================================================================

\noindent
\textbf{VI.\quad COMMENTS OF THE INTERNSHIP HOST ORGANIZATION}

\vspace{0.15cm}

The page entitled \textbf{COMMENTS OF THE INTERNSHIP HOST ORGANIZATION}
is reserved for the internship host organization.

It includes:

\begin{itemize}[leftmargin=0.8cm, itemsep=2pt]
    \item the student's full name;
    \item class;
    \item comments from the host organization;
    \item score;
    \item date;
    \item signature and official seal.
\end{itemize}

Students shall not alter or complete official evaluation fields unless
specifically authorized.

\vspace{0.3cm}

% ================================================================
% VII. SUPERVISOR'S EVALUATION AND DEFENSE RESULTS
% ================================================================

\noindent
\textbf{VII.\quad SUPERVISOR'S EVALUATION AND DEFENSE RESULTS}

\vspace{0.15cm}

The section entitled
\textbf{SUPERVISOR'S EVALUATION AND GRADUATION INTERNSHIP DEFENSE RESULTS}
is reserved for the supervisor and defense committee.

The evaluation may include:

\begin{itemize}[leftmargin=0.8cm, itemsep=2pt]
    \item Internship Progress and Student's Attitude;
    \item Report Content;
    \item Report Presentation;
    \item Other Comments;
    \item Supervisor's Opinion;
    \item Defense Result;
    \item Score.
\end{itemize}

Students shall not enter evaluation scores or defense results before they
are officially determined.

\vspace{0.3cm}

% ================================================================
% VIII. CHAPTERS, SECTIONS, AND SUBSECTIONS
% ================================================================

\noindent
\textbf{VIII.\quad CHAPTERS, SECTIONS, AND SUBSECTIONS}

\vspace{0.15cm}

The report structure shall be generated using the standard \LaTeX{}
commands:

\begin{verbatim}
\chapter{...}
\section{...}
\subsection{...}
\end{verbatim}

Chapter, section, and subsection numbers shall be generated automatically
by \LaTeX{}.

Students shall not manually type chapter or section numbers and shall not
reproduce heading styles using commands such as \verb|\textbf|,
\verb|\fontsize|, or \verb|\centering| instead of the appropriate
structural command.

Excessive subdivision should be avoided. In most cases, the hierarchy
Chapter--Section--Subsection is sufficient.

\vspace{0.3cm}

% ================================================================
% IX. ORGANIZATION OF THE MAIN CONTENT
% ================================================================

\noindent
\textbf{IX.\quad ORGANIZATION OF THE MAIN CONTENT}

\vspace{0.15cm}

The specific content of the report shall depend on the internship
organization and the assigned tasks. A typical structure may include:

\begin{itemize}[leftmargin=0.8cm, itemsep=2pt]
    \item \textbf{Chapter 1. Overview of the Internship Organization};
    \item \textbf{Chapter 2. Internship Activities and Assigned Tasks};
    \item \textbf{Chapter 3. Results, Experience, and Evaluation}.
\end{itemize}

These chapter titles are recommendations only and may be adapted to the
actual internship content.

The report shall reflect work actually performed, observed, analyzed, or
learned during the internship. Unrelated theoretical material should not
be included merely to increase report length.

\vspace{0.3cm}

% ================================================================
% X. TABLE OF CONTENTS
% ================================================================

\noindent
\textbf{X.\quad TABLE OF CONTENTS}

\vspace{0.15cm}

The Table of Contents shall be generated automatically using:

\begin{verbatim}
\tableofcontents
\end{verbatim}

Students shall not create or manually edit page numbers in the Table of
Contents.

The document shall be compiled sufficiently to ensure that all page
numbers and entries are correctly updated.

\vspace{0.3cm}

% ================================================================
% XI. LIST OF SYMBOLS AND ABBREVIATIONS
% ================================================================

\noindent
\textbf{XI.\quad LIST OF SYMBOLS AND ABBREVIATIONS}

\vspace{0.15cm}

Symbols and abbreviations that are used repeatedly in the report shall be
listed in the \textbf{List of Symbols and Abbreviations}.

The list may contain the following columns:

\begin{center}
\begin{tabular}{llll}
    No. & Symbol / Abbreviation & Unit & Definition
\end{tabular}
\end{center}

Only symbols and abbreviations actually used in the final report shall be
retained. Example entries contained in the template shall be removed if
they are not relevant to the student's report.

Each abbreviation should normally be defined when it first appears in
the main text.

\vspace{0.3cm}

% ================================================================
% XII. FIGURES, DIAGRAMS, AND PHOTOGRAPHS
% ================================================================

\noindent
\textbf{XII.\quad FIGURES, DIAGRAMS, AND PHOTOGRAPHS}

\vspace{0.15cm}

Figures shall be clear, legible, and directly relevant to the report
content.

Figures shall be inserted using the standard \verb|figure| environment.

Example:

\begin{verbatim}
\begin{figure}[H]
    \centering
    \includegraphics[width=0.8\textwidth]
                    {figures/company_structure.pdf}
    \caption{Organizational structure of the company}
    \label{fig:company-structure}
\end{figure}
\end{verbatim}

Figure captions shall be placed below the figure.

Each referenced figure shall have a unique \verb|\label| and shall be
referenced in the text using, for example:

\begin{verbatim}
Figure~\ref{fig:company-structure}
\end{verbatim}

Students shall not manually type figure numbers.

Figures reproduced or adapted from external sources shall include an
appropriate citation.

Photographs taken during the internship shall be relevant and sufficiently
clear. Students shall not include confidential, restricted, or proprietary
information without permission from the internship organization.

Vector graphics or high-resolution images are recommended whenever
possible.

\vspace{0.3cm}

% ================================================================
% XIII. TABLES
% ================================================================

\noindent
\textbf{XIII.\quad TABLES}

\vspace{0.15cm}

Tables shall be created using suitable \LaTeX{} environments such as
\verb|table|, \verb|tabular|, \verb|tabularx|, or \verb|longtable|.

Table captions shall be placed above the table.

Example:

\begin{verbatim}
\begin{table}[H]
    \centering
    \caption{Internship schedule}
    \label{tab:internship-schedule}

    \begin{tabular}{lll}
        \toprule
        Week & Task & Result \\
        \midrule
        1 & Orientation & Completed \\
        2 & System study & Completed \\
        \bottomrule
    \end{tabular}
\end{table}
\end{verbatim}

Tables shall be referenced using:

\begin{verbatim}
Table~\ref{tab:internship-schedule}
\end{verbatim}

Students shall not manually type table numbers.

Tables should be created directly in \LaTeX{} rather than inserted as
screenshots whenever practical.

\vspace{0.3cm}

% ================================================================
% XIV. EQUATIONS AND TECHNICAL QUANTITIES
% ================================================================

\noindent
\textbf{XIV.\quad EQUATIONS AND TECHNICAL QUANTITIES}

\vspace{0.15cm}

Mathematical expressions shall be typeset directly using \LaTeX{}.
Images or screenshots of equations shall not be used.

Short expressions may be written in inline mathematical mode, for example:

\begin{verbatim}
$R = 0.7\,\Omega$
\end{verbatim}

Important equations shall use a numbered equation environment.

Example:

\begin{verbatim}
\begin{equation}
    P = UI
    \label{eq:power}
\end{equation}
\end{verbatim}

Equations shall be referenced using:

\begin{verbatim}
Equation~\eqref{eq:power}
\end{verbatim}

Equation numbers shall not be manually entered.

SI units shall be used wherever applicable. The \verb|siunitx| package
should be used where appropriate.

Examples:

\begin{verbatim}
\SI{220}{\volt}
\SI{50}{\hertz}
\SI{2.5}{\ampere}
\end{verbatim}

Notation and units shall be consistent throughout the report.

\vspace{0.3cm}

% ================================================================
% XV. SOURCE CODE
% ================================================================

\noindent
\textbf{XV.\quad SOURCE CODE}

\vspace{0.15cm}

If the internship involves programming, PLCs, microcontrollers,
simulation, or software development, relevant source code may be included
in the report.

Source code shall be presented using the \verb|listings| environment or
another environment defined by the official template.

Example:

\begin{verbatim}
\begin{lstlisting}[language=st]
IF Start = TRUE THEN
    Motor := TRUE;
ELSE
    Motor := FALSE;
END_IF;
\end{lstlisting}
\end{verbatim}

Only code necessary for explaining the work should appear in the main
report. Long source-code listings should be moved to an appendix or
submitted as supplementary files if required.

Source code should not normally be inserted as screenshots.

\vspace{0.3cm}

% ================================================================
% XVI. ABSTRACT
% ================================================================

\noindent
\textbf{XVI.\quad ABSTRACT}

\vspace{0.15cm}

The Abstract shall briefly summarize:

\begin{itemize}[leftmargin=0.8cm, itemsep=2pt]
    \item the internship organization;
    \item the objectives of the internship;
    \item the main tasks performed;
    \item the main results, knowledge, and skills obtained.
\end{itemize}

The Abstract should be concise and should not reproduce the full content
of the report chapter by chapter.

\vspace{0.3cm}

% ================================================================
% XVII. REFERENCES AND CITATIONS
% ================================================================

\noindent
\textbf{XVII.\quad REFERENCES AND CITATIONS}

\vspace{0.15cm}

All books, journal papers, conference papers, standards, websites,
datasheets, manuals, technical documents, figures, tables, or other
external sources used in preparing the report shall be properly cited.

References shall be managed through the bibliography system provided with
the official template.

Bibliographic entries shall normally be stored in:

\begin{verbatim}
References.bib
\end{verbatim}

References shall be cited in the report using:

\begin{verbatim}
\cite{reference-key}
\end{verbatim}

Students shall not manually type reference numbers such as ``[1]'',
``[2]'', etc.

Titles of foreign-language publications shall remain in their original
language in the bibliography.

All borrowed ideas, data, figures, tables, equations, methods, or results
shall be properly acknowledged.

\vspace{0.3cm}

% ================================================================
% XVIII. CROSS-REFERENCING
% ================================================================

\noindent
\textbf{XVIII.\quad CROSS-REFERENCING}

\vspace{0.15cm}

Numbered elements shall be cross-referenced using the standard \LaTeX{}
mechanism.

Use the following commands where appropriate:

\begin{verbatim}
\label{...}
\ref{...}
\eqref{...}
\end{verbatim}

For example:

\begin{verbatim}
As shown in Figure~\ref{fig:system}, ...
\end{verbatim}

Students shall not manually type numbers such as ``Figure 2.3'' when
those numbers can be generated automatically.

\vspace{0.3cm}

% ================================================================
% XIX. PAGE NUMBERING
% ================================================================

\noindent
\textbf{XIX.\quad PAGE NUMBERING}

\vspace{0.15cm}

Cover pages and official evaluation forms at the beginning of the report
shall not display page numbers.

The main body of the report shall use Arabic page numbers beginning from
Chapter 1.

Page numbering shall be generated automatically by the official template.
Students shall not manually enter page numbers.

\vspace{0.3cm}

% ================================================================
% XX. PROJECT FILE ORGANIZATION
% ================================================================

\noindent
\textbf{XX.\quad PROJECT FILE ORGANIZATION}

\vspace{0.15cm}

Students should maintain a clear and organized \LaTeX{} project
structure. A recommended structure is:

\begin{verbatim}
InternshipReport/
|
|-- main.tex
|-- References.bib
|
|-- figures/
|   |-- company.pdf
|   |-- system.png
|   `-- ...
|
`-- chapters/
    |-- chapter1.tex
    |-- chapter2.tex
    `-- chapter3.tex
\end{verbatim}

Figures should normally be stored in the \verb|figures| directory
specified by the official template.

File and directory names should preferably contain Latin letters,
numbers, hyphens, or underscores.

Recommended examples include:

\begin{verbatim}
company_structure.pdf
motor_control.png
chapter_01.tex
\end{verbatim}

Spaces and unnecessary special characters in file names should be
avoided where possible.

\vspace{0.3cm}

% ================================================================
% XXI. DOCUMENT COMPILATION AND FINAL CHECK
% ================================================================

\noindent
\textbf{XXI.\quad DOCUMENT COMPILATION AND FINAL CHECK}

\vspace{0.15cm}

Before submission, the complete \LaTeX{} project shall compile
successfully and produce the final PDF report.

The final document shall not contain unresolved references such as:

\begin{verbatim}
??
[?]
\end{verbatim}

Before generating the final PDF, students shall check:

\begin{itemize}[leftmargin=0.8cm, itemsep=2pt]
    \item the Table of Contents;
    \item page numbering;
    \item chapter, section, and subsection numbering;
    \item figure numbering;
    \item table numbering;
    \item equation numbering;
    \item cross-references;
    \item citations;
    \item bibliography;
    \item font display;
    \item mathematical symbols;
    \item image quality;
    \item completeness of all required forms.
\end{itemize}

The document may be compiled using TeX Live, MiKTeX, or another
compatible \LaTeX{} environment. Regardless of the compilation
environment, the resulting PDF shall conform to the official Faculty
template.

\vspace{0.3cm}

% ================================================================
% XXII. CONTENT THAT STUDENTS MAY EDIT
% ================================================================

\noindent
\textbf{XXII.\quad CONTENT THAT STUDENTS MAY EDIT}

\vspace{0.15cm}

Students may edit or replace:

\begin{itemize}[leftmargin=0.8cm, itemsep=2pt]
    \item student name;
    \item class;
    \item internship organization;
    \item supervisor name;
    \item Abstract;
    \item List of Symbols and Abbreviations;
    \item chapter titles and chapter content;
    \item figures and photographs;
    \item tables;
    \item equations;
    \item citations and references.
\end{itemize}

Students shall not modify the main document formatting settings in the
preamble, including the document class, page geometry, line spacing,
heading formatting, table-of-contents formatting, font configuration,
page layout, or cover-page macros unless specifically instructed.

\vspace{0.3cm}

% ================================================================
% XXIII. FINAL SUBMISSION REQUIREMENTS
% ================================================================

\noindent
\textbf{XXIII.\quad FINAL SUBMISSION REQUIREMENTS}

\vspace{0.15cm}

Before submission, the final report shall:

\begin{itemize}[leftmargin=0.8cm, itemsep=2pt]
    \item compile without errors that affect the report content or presentation;
    \item contain all required forms and report sections;
    \item contain correct student, class, internship organization, and supervisor information;
    \item contain an updated Table of Contents, List of Figures, and List of Tables;
    \item contain no unused example content from the template;
    \item contain no placeholder text such as \textit{Title of Chapter 1},
    \textit{Section 1.1}, or other sample entries;
    \item contain all required citations and references;
    \item contain no unresolved cross-references.
\end{itemize}

If the Faculty requires submission of the \LaTeX{} source, the complete
project shall include all \verb|.tex| files, the bibliography database,
figures, and any other files required to reproduce the submitted PDF.

\vspace{0.4cm}

\noindent
\textbf{Important:}
Before submitting the final internship report, students shall remove or
comment out the following command in the main document:

\begin{verbatim}
\input{Regulations.tex}
\end{verbatim}

For example:

\begin{verbatim}
% \input{Regulations.tex}
\end{verbatim}

This removes the guidance pages from the final submitted report while
preserving the regulations file in the \LaTeX{} project for reference.

} % End 13pt content

\clearpage
|.
Students shall remove or comment out this command before submitting the
final internship report.

\vspace{0.4cm}

% ================================================================
% I. GENERAL REQUIREMENTS
% ================================================================

\noindent
\textbf{I.\quad GENERAL REQUIREMENTS}

\vspace{0.15cm}

The Graduation Internship Report shall be prepared using the official
\LaTeX{} template provided by the Faculty. Students shall preserve the
predefined page layout, typography, line spacing, heading styles,
numbering system, figure and table formatting, and bibliography settings
unless otherwise instructed by the Faculty or the supervisor.

The report shall be written in English.

Students are responsible for the academic and technical content of the
report. The official \LaTeX{} template is responsible for standardized
document formatting.

\vspace{0.3cm}

% ================================================================
% II. REPORT STRUCTURE
% ================================================================

\noindent
\textbf{II.\quad REPORT STRUCTURE}

\vspace{0.15cm}

The report shall normally be arranged in the following order:

\begin{enumerate}[label=\arabic*., leftmargin=1.2cm, itemsep=2pt]
    \item Outer Cover;
    \item Comments of the Internship Host Organization;
    \item Supervisor's Evaluation and Graduation Internship Defense Results;
    \item Table of Contents;
    \item List of Symbols and Abbreviations, if applicable;
    \item List of Figures;
    \item List of Tables;
    \item Abstract;
    \item Main Chapters of the Internship Report;
    \item References;
    \item Appendices, if applicable.
\end{enumerate}

The preliminary pages shall not be numbered as chapters. The main report
content shall be organized into consecutively numbered chapters beginning
with Chapter 1.

\vspace{0.3cm}

% ================================================================
% III. PAGE SIZE AND MARGINS
% ================================================================

\noindent
\textbf{III.\quad PAGE SIZE AND MARGINS}

\vspace{0.15cm}

The report shall be formatted on A4 paper with the following margins:

\begin{center}
\begin{tabular}{ll}
    Left margin:   & 3.5 cm \\
    Right margin:  & 2.0 cm \\
    Top margin:    & 2.5 cm \\
    Bottom margin: & 2.5 cm
\end{tabular}
\end{center}

These settings are predefined in the official template and shall not be
modified by students.

\vspace{0.3cm}

% ================================================================
% IV. FONT, FONT SIZE, AND LINE SPACING
% ================================================================

\noindent
\textbf{IV.\quad FONT, FONT SIZE, AND LINE SPACING}

\vspace{0.15cm}

The main text shall use a 13 pt font size with 1.3-line spacing.
A Times-family typeface shall be used in accordance with the official
template.

The text shall be fully justified.

Students shall not use repeated spaces, repeated manual line breaks, or
multiple blank lines to position text. Paragraph and heading spacing shall
be controlled by the template.

\vspace{0.3cm}

% ================================================================
% V. OUTER COVER
% ================================================================

\noindent
\textbf{V.\quad OUTER COVER}

\vspace{0.15cm}

The outer cover shall follow the official format provided in the template
and shall contain the following information:

\begin{itemize}[leftmargin=0.8cm, itemsep=2pt]
    \item Thai Nguyen University of Technology;
    \item Faculty of International Training;
    \item Engineering and Technology Department;
    \item University logo and Faculty logo;
    \item the title \textbf{INTERNSHIP REPORT};
    \item Internship Organization;
    \item Author;
    \item Class;
    \item Supervisor;
    \item place and year of completion.
\end{itemize}

Students shall only replace the required information fields and shall not
modify the cover-page layout unless instructed otherwise.

\vspace{0.3cm}

% ================================================================
% VI. COMMENTS OF THE INTERNSHIP HOST ORGANIZATION
% ================================================================

\noindent
\textbf{VI.\quad COMMENTS OF THE INTERNSHIP HOST ORGANIZATION}

\vspace{0.15cm}

The page entitled \textbf{COMMENTS OF THE INTERNSHIP HOST ORGANIZATION}
is reserved for the internship host organization.

It includes:

\begin{itemize}[leftmargin=0.8cm, itemsep=2pt]
    \item the student's full name;
    \item class;
    \item comments from the host organization;
    \item score;
    \item date;
    \item signature and official seal.
\end{itemize}

Students shall not alter or complete official evaluation fields unless
specifically authorized.

\vspace{0.3cm}

% ================================================================
% VII. SUPERVISOR'S EVALUATION AND DEFENSE RESULTS
% ================================================================

\noindent
\textbf{VII.\quad SUPERVISOR'S EVALUATION AND DEFENSE RESULTS}

\vspace{0.15cm}

The section entitled
\textbf{SUPERVISOR'S EVALUATION AND GRADUATION INTERNSHIP DEFENSE RESULTS}
is reserved for the supervisor and defense committee.

The evaluation may include:

\begin{itemize}[leftmargin=0.8cm, itemsep=2pt]
    \item Internship Progress and Student's Attitude;
    \item Report Content;
    \item Report Presentation;
    \item Other Comments;
    \item Supervisor's Opinion;
    \item Defense Result;
    \item Score.
\end{itemize}

Students shall not enter evaluation scores or defense results before they
are officially determined.

\vspace{0.3cm}

% ================================================================
% VIII. CHAPTERS, SECTIONS, AND SUBSECTIONS
% ================================================================

\noindent
\textbf{VIII.\quad CHAPTERS, SECTIONS, AND SUBSECTIONS}

\vspace{0.15cm}

The report structure shall be generated using the standard \LaTeX{}
commands:

\begin{verbatim}
\chapter{...}
\section{...}
\subsection{...}
\end{verbatim}

Chapter, section, and subsection numbers shall be generated automatically
by \LaTeX{}.

Students shall not manually type chapter or section numbers and shall not
reproduce heading styles using commands such as \verb|\textbf|,
\verb|\fontsize|, or \verb|\centering| instead of the appropriate
structural command.

Excessive subdivision should be avoided. In most cases, the hierarchy
Chapter--Section--Subsection is sufficient.

\vspace{0.3cm}

% ================================================================
% IX. ORGANIZATION OF THE MAIN CONTENT
% ================================================================

\noindent
\textbf{IX.\quad ORGANIZATION OF THE MAIN CONTENT}

\vspace{0.15cm}

The specific content of the report shall depend on the internship
organization and the assigned tasks. A typical structure may include:

\begin{itemize}[leftmargin=0.8cm, itemsep=2pt]
    \item \textbf{Chapter 1. Overview of the Internship Organization};
    \item \textbf{Chapter 2. Internship Activities and Assigned Tasks};
    \item \textbf{Chapter 3. Results, Experience, and Evaluation}.
\end{itemize}

These chapter titles are recommendations only and may be adapted to the
actual internship content.

The report shall reflect work actually performed, observed, analyzed, or
learned during the internship. Unrelated theoretical material should not
be included merely to increase report length.

\vspace{0.3cm}

% ================================================================
% X. TABLE OF CONTENTS
% ================================================================

\noindent
\textbf{X.\quad TABLE OF CONTENTS}

\vspace{0.15cm}

The Table of Contents shall be generated automatically using:

\begin{verbatim}
\tableofcontents
\end{verbatim}

Students shall not create or manually edit page numbers in the Table of
Contents.

The document shall be compiled sufficiently to ensure that all page
numbers and entries are correctly updated.

\vspace{0.3cm}

% ================================================================
% XI. LIST OF SYMBOLS AND ABBREVIATIONS
% ================================================================

\noindent
\textbf{XI.\quad LIST OF SYMBOLS AND ABBREVIATIONS}

\vspace{0.15cm}

Symbols and abbreviations that are used repeatedly in the report shall be
listed in the \textbf{List of Symbols and Abbreviations}.

The list may contain the following columns:

\begin{center}
\begin{tabular}{llll}
    No. & Symbol / Abbreviation & Unit & Definition
\end{tabular}
\end{center}

Only symbols and abbreviations actually used in the final report shall be
retained. Example entries contained in the template shall be removed if
they are not relevant to the student's report.

Each abbreviation should normally be defined when it first appears in
the main text.

\vspace{0.3cm}

% ================================================================
% XII. FIGURES, DIAGRAMS, AND PHOTOGRAPHS
% ================================================================

\noindent
\textbf{XII.\quad FIGURES, DIAGRAMS, AND PHOTOGRAPHS}

\vspace{0.15cm}

Figures shall be clear, legible, and directly relevant to the report
content.

Figures shall be inserted using the standard \verb|figure| environment.

Example:

\begin{verbatim}
\begin{figure}[H]
    \centering
    \includegraphics[width=0.8\textwidth]
                    {figures/company_structure.pdf}
    \caption{Organizational structure of the company}
    \label{fig:company-structure}
\end{figure}
\end{verbatim}

Figure captions shall be placed below the figure.

Each referenced figure shall have a unique \verb|\label| and shall be
referenced in the text using, for example:

\begin{verbatim}
Figure~\ref{fig:company-structure}
\end{verbatim}

Students shall not manually type figure numbers.

Figures reproduced or adapted from external sources shall include an
appropriate citation.

Photographs taken during the internship shall be relevant and sufficiently
clear. Students shall not include confidential, restricted, or proprietary
information without permission from the internship organization.

Vector graphics or high-resolution images are recommended whenever
possible.

\vspace{0.3cm}

% ================================================================
% XIII. TABLES
% ================================================================

\noindent
\textbf{XIII.\quad TABLES}

\vspace{0.15cm}

Tables shall be created using suitable \LaTeX{} environments such as
\verb|table|, \verb|tabular|, \verb|tabularx|, or \verb|longtable|.

Table captions shall be placed above the table.

Example:

\begin{verbatim}
\begin{table}[H]
    \centering
    \caption{Internship schedule}
    \label{tab:internship-schedule}

    \begin{tabular}{lll}
        \toprule
        Week & Task & Result \\
        \midrule
        1 & Orientation & Completed \\
        2 & System study & Completed \\
        \bottomrule
    \end{tabular}
\end{table}
\end{verbatim}

Tables shall be referenced using:

\begin{verbatim}
Table~\ref{tab:internship-schedule}
\end{verbatim}

Students shall not manually type table numbers.

Tables should be created directly in \LaTeX{} rather than inserted as
screenshots whenever practical.

\vspace{0.3cm}

% ================================================================
% XIV. EQUATIONS AND TECHNICAL QUANTITIES
% ================================================================

\noindent
\textbf{XIV.\quad EQUATIONS AND TECHNICAL QUANTITIES}

\vspace{0.15cm}

Mathematical expressions shall be typeset directly using \LaTeX{}.
Images or screenshots of equations shall not be used.

Short expressions may be written in inline mathematical mode, for example:

\begin{verbatim}
$R = 0.7\,\Omega$
\end{verbatim}

Important equations shall use a numbered equation environment.

Example:

\begin{verbatim}
\begin{equation}
    P = UI
    \label{eq:power}
\end{equation}
\end{verbatim}

Equations shall be referenced using:

\begin{verbatim}
Equation~\eqref{eq:power}
\end{verbatim}

Equation numbers shall not be manually entered.

SI units shall be used wherever applicable. The \verb|siunitx| package
should be used where appropriate.

Examples:

\begin{verbatim}
\SI{220}{\volt}
\SI{50}{\hertz}
\SI{2.5}{\ampere}
\end{verbatim}

Notation and units shall be consistent throughout the report.

\vspace{0.3cm}

% ================================================================
% XV. SOURCE CODE
% ================================================================

\noindent
\textbf{XV.\quad SOURCE CODE}

\vspace{0.15cm}

If the internship involves programming, PLCs, microcontrollers,
simulation, or software development, relevant source code may be included
in the report.

Source code shall be presented using the \verb|listings| environment or
another environment defined by the official template.

Example:

\begin{verbatim}
\begin{lstlisting}[language=st]
IF Start = TRUE THEN
    Motor := TRUE;
ELSE
    Motor := FALSE;
END_IF;
\end{lstlisting}
\end{verbatim}

Only code necessary for explaining the work should appear in the main
report. Long source-code listings should be moved to an appendix or
submitted as supplementary files if required.

Source code should not normally be inserted as screenshots.

\vspace{0.3cm}

% ================================================================
% XVI. ABSTRACT
% ================================================================

\noindent
\textbf{XVI.\quad ABSTRACT}

\vspace{0.15cm}

The Abstract shall briefly summarize:

\begin{itemize}[leftmargin=0.8cm, itemsep=2pt]
    \item the internship organization;
    \item the objectives of the internship;
    \item the main tasks performed;
    \item the main results, knowledge, and skills obtained.
\end{itemize}

The Abstract should be concise and should not reproduce the full content
of the report chapter by chapter.

\vspace{0.3cm}

% ================================================================
% XVII. REFERENCES AND CITATIONS
% ================================================================

\noindent
\textbf{XVII.\quad REFERENCES AND CITATIONS}

\vspace{0.15cm}

All books, journal papers, conference papers, standards, websites,
datasheets, manuals, technical documents, figures, tables, or other
external sources used in preparing the report shall be properly cited.

References shall be managed through the bibliography system provided with
the official template.

Bibliographic entries shall normally be stored in:

\begin{verbatim}
References.bib
\end{verbatim}

References shall be cited in the report using:

\begin{verbatim}
\cite{reference-key}
\end{verbatim}

Students shall not manually type reference numbers such as ``[1]'',
``[2]'', etc.

Titles of foreign-language publications shall remain in their original
language in the bibliography.

All borrowed ideas, data, figures, tables, equations, methods, or results
shall be properly acknowledged.

\vspace{0.3cm}

% ================================================================
% XVIII. CROSS-REFERENCING
% ================================================================

\noindent
\textbf{XVIII.\quad CROSS-REFERENCING}

\vspace{0.15cm}

Numbered elements shall be cross-referenced using the standard \LaTeX{}
mechanism.

Use the following commands where appropriate:

\begin{verbatim}
\label{...}
\ref{...}
\eqref{...}
\end{verbatim}

For example:

\begin{verbatim}
As shown in Figure~\ref{fig:system}, ...
\end{verbatim}

Students shall not manually type numbers such as ``Figure 2.3'' when
those numbers can be generated automatically.

\vspace{0.3cm}

% ================================================================
% XIX. PAGE NUMBERING
% ================================================================

\noindent
\textbf{XIX.\quad PAGE NUMBERING}

\vspace{0.15cm}

Cover pages and official evaluation forms at the beginning of the report
shall not display page numbers.

The main body of the report shall use Arabic page numbers beginning from
Chapter 1.

Page numbering shall be generated automatically by the official template.
Students shall not manually enter page numbers.

\vspace{0.3cm}

% ================================================================
% XX. PROJECT FILE ORGANIZATION
% ================================================================

\noindent
\textbf{XX.\quad PROJECT FILE ORGANIZATION}

\vspace{0.15cm}

Students should maintain a clear and organized \LaTeX{} project
structure. A recommended structure is:

\begin{verbatim}
InternshipReport/
|
|-- main.tex
|-- References.bib
|
|-- figures/
|   |-- company.pdf
|   |-- system.png
|   `-- ...
|
`-- chapters/
    |-- chapter1.tex
    |-- chapter2.tex
    `-- chapter3.tex
\end{verbatim}

Figures should normally be stored in the \verb|figures| directory
specified by the official template.

File and directory names should preferably contain Latin letters,
numbers, hyphens, or underscores.

Recommended examples include:

\begin{verbatim}
company_structure.pdf
motor_control.png
chapter_01.tex
\end{verbatim}

Spaces and unnecessary special characters in file names should be
avoided where possible.

\vspace{0.3cm}

% ================================================================
% XXI. DOCUMENT COMPILATION AND FINAL CHECK
% ================================================================

\noindent
\textbf{XXI.\quad DOCUMENT COMPILATION AND FINAL CHECK}

\vspace{0.15cm}

Before submission, the complete \LaTeX{} project shall compile
successfully and produce the final PDF report.

The final document shall not contain unresolved references such as:

\begin{verbatim}
??
[?]
\end{verbatim}

Before generating the final PDF, students shall check:

\begin{itemize}[leftmargin=0.8cm, itemsep=2pt]
    \item the Table of Contents;
    \item page numbering;
    \item chapter, section, and subsection numbering;
    \item figure numbering;
    \item table numbering;
    \item equation numbering;
    \item cross-references;
    \item citations;
    \item bibliography;
    \item font display;
    \item mathematical symbols;
    \item image quality;
    \item completeness of all required forms.
\end{itemize}

The document may be compiled using TeX Live, MiKTeX, or another
compatible \LaTeX{} environment. Regardless of the compilation
environment, the resulting PDF shall conform to the official Faculty
template.

\vspace{0.3cm}

% ================================================================
% XXII. CONTENT THAT STUDENTS MAY EDIT
% ================================================================

\noindent
\textbf{XXII.\quad CONTENT THAT STUDENTS MAY EDIT}

\vspace{0.15cm}

Students may edit or replace:

\begin{itemize}[leftmargin=0.8cm, itemsep=2pt]
    \item student name;
    \item class;
    \item internship organization;
    \item supervisor name;
    \item Abstract;
    \item List of Symbols and Abbreviations;
    \item chapter titles and chapter content;
    \item figures and photographs;
    \item tables;
    \item equations;
    \item citations and references.
\end{itemize}

Students shall not modify the main document formatting settings in the
preamble, including the document class, page geometry, line spacing,
heading formatting, table-of-contents formatting, font configuration,
page layout, or cover-page macros unless specifically instructed.

\vspace{0.3cm}

% ================================================================
% XXIII. FINAL SUBMISSION REQUIREMENTS
% ================================================================

\noindent
\textbf{XXIII.\quad FINAL SUBMISSION REQUIREMENTS}

\vspace{0.15cm}

Before submission, the final report shall:

\begin{itemize}[leftmargin=0.8cm, itemsep=2pt]
    \item compile without errors that affect the report content or presentation;
    \item contain all required forms and report sections;
    \item contain correct student, class, internship organization, and supervisor information;
    \item contain an updated Table of Contents, List of Figures, and List of Tables;
    \item contain no unused example content from the template;
    \item contain no placeholder text such as \textit{Title of Chapter 1},
    \textit{Section 1.1}, or other sample entries;
    \item contain all required citations and references;
    \item contain no unresolved cross-references.
\end{itemize}

If the Faculty requires submission of the \LaTeX{} source, the complete
project shall include all \verb|.tex| files, the bibliography database,
figures, and any other files required to reproduce the submitted PDF.

\vspace{0.4cm}

\noindent
\textbf{Important:}
Before submitting the final internship report, students shall remove or
comment out the following command in the main document:

\begin{verbatim}

% =================================================================
% REGULATIONS FOR PREPARING THE GRADUATION INTERNSHIP REPORT
% USING LATEX
% =================================================================
%
% This file is intended to be included from the main document with:
%
%     \input{Regulations.tex}
%
% IMPORTANT:
% This section is provided for guidance only. Students shall remove
% or comment out the \input{Regulations.tex} command before submitting
% the final internship report.
%
% =================================================================

\clearpage
\thispagestyle{empty}

\begin{center}
    {\fontsize{16pt}{20pt}\selectfont\bfseries
    SPECIFIC REGULATIONS FOR PREPARING\\
    GRADUATION INTERNSHIP REPORTS USING \LaTeX}
\end{center}

\vspace{0.5cm}

{\fontsize{13pt}{17pt}\selectfont
\setstretch{1.3}
\justifying

\noindent
\textbf{Note:}
This section is provided for guidance only and is included in the main
document using the command \verb|\input{Regulations.tex}|.
Students shall remove or comment out this command before submitting the
final internship report.

\vspace{0.4cm}

% ================================================================
% I. GENERAL REQUIREMENTS
% ================================================================

\noindent
\textbf{I.\quad GENERAL REQUIREMENTS}

\vspace{0.15cm}

The Graduation Internship Report shall be prepared using the official
\LaTeX{} template provided by the Faculty. Students shall preserve the
predefined page layout, typography, line spacing, heading styles,
numbering system, figure and table formatting, and bibliography settings
unless otherwise instructed by the Faculty or the supervisor.

The report shall be written in English.

Students are responsible for the academic and technical content of the
report. The official \LaTeX{} template is responsible for standardized
document formatting.

\vspace{0.3cm}

% ================================================================
% II. REPORT STRUCTURE
% ================================================================

\noindent
\textbf{II.\quad REPORT STRUCTURE}

\vspace{0.15cm}

The report shall normally be arranged in the following order:

\begin{enumerate}[label=\arabic*., leftmargin=1.2cm, itemsep=2pt]
    \item Outer Cover;
    \item Comments of the Internship Host Organization;
    \item Supervisor's Evaluation and Graduation Internship Defense Results;
    \item Table of Contents;
    \item List of Symbols and Abbreviations, if applicable;
    \item List of Figures;
    \item List of Tables;
    \item Abstract;
    \item Main Chapters of the Internship Report;
    \item References;
    \item Appendices, if applicable.
\end{enumerate}

The preliminary pages shall not be numbered as chapters. The main report
content shall be organized into consecutively numbered chapters beginning
with Chapter 1.

\vspace{0.3cm}

% ================================================================
% III. PAGE SIZE AND MARGINS
% ================================================================

\noindent
\textbf{III.\quad PAGE SIZE AND MARGINS}

\vspace{0.15cm}

The report shall be formatted on A4 paper with the following margins:

\begin{center}
\begin{tabular}{ll}
    Left margin:   & 3.5 cm \\
    Right margin:  & 2.0 cm \\
    Top margin:    & 2.5 cm \\
    Bottom margin: & 2.5 cm
\end{tabular}
\end{center}

These settings are predefined in the official template and shall not be
modified by students.

\vspace{0.3cm}

% ================================================================
% IV. FONT, FONT SIZE, AND LINE SPACING
% ================================================================

\noindent
\textbf{IV.\quad FONT, FONT SIZE, AND LINE SPACING}

\vspace{0.15cm}

The main text shall use a 13 pt font size with 1.3-line spacing.
A Times-family typeface shall be used in accordance with the official
template.

The text shall be fully justified.

Students shall not use repeated spaces, repeated manual line breaks, or
multiple blank lines to position text. Paragraph and heading spacing shall
be controlled by the template.

\vspace{0.3cm}

% ================================================================
% V. OUTER COVER
% ================================================================

\noindent
\textbf{V.\quad OUTER COVER}

\vspace{0.15cm}

The outer cover shall follow the official format provided in the template
and shall contain the following information:

\begin{itemize}[leftmargin=0.8cm, itemsep=2pt]
    \item Thai Nguyen University of Technology;
    \item Faculty of International Training;
    \item Engineering and Technology Department;
    \item University logo and Faculty logo;
    \item the title \textbf{INTERNSHIP REPORT};
    \item Internship Organization;
    \item Author;
    \item Class;
    \item Supervisor;
    \item place and year of completion.
\end{itemize}

Students shall only replace the required information fields and shall not
modify the cover-page layout unless instructed otherwise.

\vspace{0.3cm}

% ================================================================
% VI. COMMENTS OF THE INTERNSHIP HOST ORGANIZATION
% ================================================================

\noindent
\textbf{VI.\quad COMMENTS OF THE INTERNSHIP HOST ORGANIZATION}

\vspace{0.15cm}

The page entitled \textbf{COMMENTS OF THE INTERNSHIP HOST ORGANIZATION}
is reserved for the internship host organization.

It includes:

\begin{itemize}[leftmargin=0.8cm, itemsep=2pt]
    \item the student's full name;
    \item class;
    \item comments from the host organization;
    \item score;
    \item date;
    \item signature and official seal.
\end{itemize}

Students shall not alter or complete official evaluation fields unless
specifically authorized.

\vspace{0.3cm}

% ================================================================
% VII. SUPERVISOR'S EVALUATION AND DEFENSE RESULTS
% ================================================================

\noindent
\textbf{VII.\quad SUPERVISOR'S EVALUATION AND DEFENSE RESULTS}

\vspace{0.15cm}

The section entitled
\textbf{SUPERVISOR'S EVALUATION AND GRADUATION INTERNSHIP DEFENSE RESULTS}
is reserved for the supervisor and defense committee.

The evaluation may include:

\begin{itemize}[leftmargin=0.8cm, itemsep=2pt]
    \item Internship Progress and Student's Attitude;
    \item Report Content;
    \item Report Presentation;
    \item Other Comments;
    \item Supervisor's Opinion;
    \item Defense Result;
    \item Score.
\end{itemize}

Students shall not enter evaluation scores or defense results before they
are officially determined.

\vspace{0.3cm}

% ================================================================
% VIII. CHAPTERS, SECTIONS, AND SUBSECTIONS
% ================================================================

\noindent
\textbf{VIII.\quad CHAPTERS, SECTIONS, AND SUBSECTIONS}

\vspace{0.15cm}

The report structure shall be generated using the standard \LaTeX{}
commands:

\begin{verbatim}
\chapter{...}
\section{...}
\subsection{...}
\end{verbatim}

Chapter, section, and subsection numbers shall be generated automatically
by \LaTeX{}.

Students shall not manually type chapter or section numbers and shall not
reproduce heading styles using commands such as \verb|\textbf|,
\verb|\fontsize|, or \verb|\centering| instead of the appropriate
structural command.

Excessive subdivision should be avoided. In most cases, the hierarchy
Chapter--Section--Subsection is sufficient.

\vspace{0.3cm}

% ================================================================
% IX. ORGANIZATION OF THE MAIN CONTENT
% ================================================================

\noindent
\textbf{IX.\quad ORGANIZATION OF THE MAIN CONTENT}

\vspace{0.15cm}

The specific content of the report shall depend on the internship
organization and the assigned tasks. A typical structure may include:

\begin{itemize}[leftmargin=0.8cm, itemsep=2pt]
    \item \textbf{Chapter 1. Overview of the Internship Organization};
    \item \textbf{Chapter 2. Internship Activities and Assigned Tasks};
    \item \textbf{Chapter 3. Results, Experience, and Evaluation}.
\end{itemize}

These chapter titles are recommendations only and may be adapted to the
actual internship content.

The report shall reflect work actually performed, observed, analyzed, or
learned during the internship. Unrelated theoretical material should not
be included merely to increase report length.

\vspace{0.3cm}

% ================================================================
% X. TABLE OF CONTENTS
% ================================================================

\noindent
\textbf{X.\quad TABLE OF CONTENTS}

\vspace{0.15cm}

The Table of Contents shall be generated automatically using:

\begin{verbatim}
\tableofcontents
\end{verbatim}

Students shall not create or manually edit page numbers in the Table of
Contents.

The document shall be compiled sufficiently to ensure that all page
numbers and entries are correctly updated.

\vspace{0.3cm}

% ================================================================
% XI. LIST OF SYMBOLS AND ABBREVIATIONS
% ================================================================

\noindent
\textbf{XI.\quad LIST OF SYMBOLS AND ABBREVIATIONS}

\vspace{0.15cm}

Symbols and abbreviations that are used repeatedly in the report shall be
listed in the \textbf{List of Symbols and Abbreviations}.

The list may contain the following columns:

\begin{center}
\begin{tabular}{llll}
    No. & Symbol / Abbreviation & Unit & Definition
\end{tabular}
\end{center}

Only symbols and abbreviations actually used in the final report shall be
retained. Example entries contained in the template shall be removed if
they are not relevant to the student's report.

Each abbreviation should normally be defined when it first appears in
the main text.

\vspace{0.3cm}

% ================================================================
% XII. FIGURES, DIAGRAMS, AND PHOTOGRAPHS
% ================================================================

\noindent
\textbf{XII.\quad FIGURES, DIAGRAMS, AND PHOTOGRAPHS}

\vspace{0.15cm}

Figures shall be clear, legible, and directly relevant to the report
content.

Figures shall be inserted using the standard \verb|figure| environment.

Example:

\begin{verbatim}
\begin{figure}[H]
    \centering
    \includegraphics[width=0.8\textwidth]
                    {figures/company_structure.pdf}
    \caption{Organizational structure of the company}
    \label{fig:company-structure}
\end{figure}
\end{verbatim}

Figure captions shall be placed below the figure.

Each referenced figure shall have a unique \verb|\label| and shall be
referenced in the text using, for example:

\begin{verbatim}
Figure~\ref{fig:company-structure}
\end{verbatim}

Students shall not manually type figure numbers.

Figures reproduced or adapted from external sources shall include an
appropriate citation.

Photographs taken during the internship shall be relevant and sufficiently
clear. Students shall not include confidential, restricted, or proprietary
information without permission from the internship organization.

Vector graphics or high-resolution images are recommended whenever
possible.

\vspace{0.3cm}

% ================================================================
% XIII. TABLES
% ================================================================

\noindent
\textbf{XIII.\quad TABLES}

\vspace{0.15cm}

Tables shall be created using suitable \LaTeX{} environments such as
\verb|table|, \verb|tabular|, \verb|tabularx|, or \verb|longtable|.

Table captions shall be placed above the table.

Example:

\begin{verbatim}
\begin{table}[H]
    \centering
    \caption{Internship schedule}
    \label{tab:internship-schedule}

    \begin{tabular}{lll}
        \toprule
        Week & Task & Result \\
        \midrule
        1 & Orientation & Completed \\
        2 & System study & Completed \\
        \bottomrule
    \end{tabular}
\end{table}
\end{verbatim}

Tables shall be referenced using:

\begin{verbatim}
Table~\ref{tab:internship-schedule}
\end{verbatim}

Students shall not manually type table numbers.

Tables should be created directly in \LaTeX{} rather than inserted as
screenshots whenever practical.

\vspace{0.3cm}

% ================================================================
% XIV. EQUATIONS AND TECHNICAL QUANTITIES
% ================================================================

\noindent
\textbf{XIV.\quad EQUATIONS AND TECHNICAL QUANTITIES}

\vspace{0.15cm}

Mathematical expressions shall be typeset directly using \LaTeX{}.
Images or screenshots of equations shall not be used.

Short expressions may be written in inline mathematical mode, for example:

\begin{verbatim}
$R = 0.7\,\Omega$
\end{verbatim}

Important equations shall use a numbered equation environment.

Example:

\begin{verbatim}
\begin{equation}
    P = UI
    \label{eq:power}
\end{equation}
\end{verbatim}

Equations shall be referenced using:

\begin{verbatim}
Equation~\eqref{eq:power}
\end{verbatim}

Equation numbers shall not be manually entered.

SI units shall be used wherever applicable. The \verb|siunitx| package
should be used where appropriate.

Examples:

\begin{verbatim}
\SI{220}{\volt}
\SI{50}{\hertz}
\SI{2.5}{\ampere}
\end{verbatim}

Notation and units shall be consistent throughout the report.

\vspace{0.3cm}

% ================================================================
% XV. SOURCE CODE
% ================================================================

\noindent
\textbf{XV.\quad SOURCE CODE}

\vspace{0.15cm}

If the internship involves programming, PLCs, microcontrollers,
simulation, or software development, relevant source code may be included
in the report.

Source code shall be presented using the \verb|listings| environment or
another environment defined by the official template.

Example:

\begin{verbatim}
\begin{lstlisting}[language=st]
IF Start = TRUE THEN
    Motor := TRUE;
ELSE
    Motor := FALSE;
END_IF;
\end{lstlisting}
\end{verbatim}

Only code necessary for explaining the work should appear in the main
report. Long source-code listings should be moved to an appendix or
submitted as supplementary files if required.

Source code should not normally be inserted as screenshots.

\vspace{0.3cm}

% ================================================================
% XVI. ABSTRACT
% ================================================================

\noindent
\textbf{XVI.\quad ABSTRACT}

\vspace{0.15cm}

The Abstract shall briefly summarize:

\begin{itemize}[leftmargin=0.8cm, itemsep=2pt]
    \item the internship organization;
    \item the objectives of the internship;
    \item the main tasks performed;
    \item the main results, knowledge, and skills obtained.
\end{itemize}

The Abstract should be concise and should not reproduce the full content
of the report chapter by chapter.

\vspace{0.3cm}

% ================================================================
% XVII. REFERENCES AND CITATIONS
% ================================================================

\noindent
\textbf{XVII.\quad REFERENCES AND CITATIONS}

\vspace{0.15cm}

All books, journal papers, conference papers, standards, websites,
datasheets, manuals, technical documents, figures, tables, or other
external sources used in preparing the report shall be properly cited.

References shall be managed through the bibliography system provided with
the official template.

Bibliographic entries shall normally be stored in:

\begin{verbatim}
References.bib
\end{verbatim}

References shall be cited in the report using:

\begin{verbatim}
\cite{reference-key}
\end{verbatim}

Students shall not manually type reference numbers such as ``[1]'',
``[2]'', etc.

Titles of foreign-language publications shall remain in their original
language in the bibliography.

All borrowed ideas, data, figures, tables, equations, methods, or results
shall be properly acknowledged.

\vspace{0.3cm}

% ================================================================
% XVIII. CROSS-REFERENCING
% ================================================================

\noindent
\textbf{XVIII.\quad CROSS-REFERENCING}

\vspace{0.15cm}

Numbered elements shall be cross-referenced using the standard \LaTeX{}
mechanism.

Use the following commands where appropriate:

\begin{verbatim}
\label{...}
\ref{...}
\eqref{...}
\end{verbatim}

For example:

\begin{verbatim}
As shown in Figure~\ref{fig:system}, ...
\end{verbatim}

Students shall not manually type numbers such as ``Figure 2.3'' when
those numbers can be generated automatically.

\vspace{0.3cm}

% ================================================================
% XIX. PAGE NUMBERING
% ================================================================

\noindent
\textbf{XIX.\quad PAGE NUMBERING}

\vspace{0.15cm}

Cover pages and official evaluation forms at the beginning of the report
shall not display page numbers.

The main body of the report shall use Arabic page numbers beginning from
Chapter 1.

Page numbering shall be generated automatically by the official template.
Students shall not manually enter page numbers.

\vspace{0.3cm}

% ================================================================
% XX. PROJECT FILE ORGANIZATION
% ================================================================

\noindent
\textbf{XX.\quad PROJECT FILE ORGANIZATION}

\vspace{0.15cm}

Students should maintain a clear and organized \LaTeX{} project
structure. A recommended structure is:

\begin{verbatim}
InternshipReport/
|
|-- main.tex
|-- References.bib
|
|-- figures/
|   |-- company.pdf
|   |-- system.png
|   `-- ...
|
`-- chapters/
    |-- chapter1.tex
    |-- chapter2.tex
    `-- chapter3.tex
\end{verbatim}

Figures should normally be stored in the \verb|figures| directory
specified by the official template.

File and directory names should preferably contain Latin letters,
numbers, hyphens, or underscores.

Recommended examples include:

\begin{verbatim}
company_structure.pdf
motor_control.png
chapter_01.tex
\end{verbatim}

Spaces and unnecessary special characters in file names should be
avoided where possible.

\vspace{0.3cm}

% ================================================================
% XXI. DOCUMENT COMPILATION AND FINAL CHECK
% ================================================================

\noindent
\textbf{XXI.\quad DOCUMENT COMPILATION AND FINAL CHECK}

\vspace{0.15cm}

Before submission, the complete \LaTeX{} project shall compile
successfully and produce the final PDF report.

The final document shall not contain unresolved references such as:

\begin{verbatim}
??
[?]
\end{verbatim}

Before generating the final PDF, students shall check:

\begin{itemize}[leftmargin=0.8cm, itemsep=2pt]
    \item the Table of Contents;
    \item page numbering;
    \item chapter, section, and subsection numbering;
    \item figure numbering;
    \item table numbering;
    \item equation numbering;
    \item cross-references;
    \item citations;
    \item bibliography;
    \item font display;
    \item mathematical symbols;
    \item image quality;
    \item completeness of all required forms.
\end{itemize}

The document may be compiled using TeX Live, MiKTeX, or another
compatible \LaTeX{} environment. Regardless of the compilation
environment, the resulting PDF shall conform to the official Faculty
template.

\vspace{0.3cm}

% ================================================================
% XXII. CONTENT THAT STUDENTS MAY EDIT
% ================================================================

\noindent
\textbf{XXII.\quad CONTENT THAT STUDENTS MAY EDIT}

\vspace{0.15cm}

Students may edit or replace:

\begin{itemize}[leftmargin=0.8cm, itemsep=2pt]
    \item student name;
    \item class;
    \item internship organization;
    \item supervisor name;
    \item Abstract;
    \item List of Symbols and Abbreviations;
    \item chapter titles and chapter content;
    \item figures and photographs;
    \item tables;
    \item equations;
    \item citations and references.
\end{itemize}

Students shall not modify the main document formatting settings in the
preamble, including the document class, page geometry, line spacing,
heading formatting, table-of-contents formatting, font configuration,
page layout, or cover-page macros unless specifically instructed.

\vspace{0.3cm}

% ================================================================
% XXIII. FINAL SUBMISSION REQUIREMENTS
% ================================================================

\noindent
\textbf{XXIII.\quad FINAL SUBMISSION REQUIREMENTS}

\vspace{0.15cm}

Before submission, the final report shall:

\begin{itemize}[leftmargin=0.8cm, itemsep=2pt]
    \item compile without errors that affect the report content or presentation;
    \item contain all required forms and report sections;
    \item contain correct student, class, internship organization, and supervisor information;
    \item contain an updated Table of Contents, List of Figures, and List of Tables;
    \item contain no unused example content from the template;
    \item contain no placeholder text such as \textit{Title of Chapter 1},
    \textit{Section 1.1}, or other sample entries;
    \item contain all required citations and references;
    \item contain no unresolved cross-references.
\end{itemize}

If the Faculty requires submission of the \LaTeX{} source, the complete
project shall include all \verb|.tex| files, the bibliography database,
figures, and any other files required to reproduce the submitted PDF.

\vspace{0.4cm}

\noindent
\textbf{Important:}
Before submitting the final internship report, students shall remove or
comment out the following command in the main document:

\begin{verbatim}
\input{Regulations.tex}
\end{verbatim}

For example:

\begin{verbatim}
% \input{Regulations.tex}
\end{verbatim}

This removes the guidance pages from the final submitted report while
preserving the regulations file in the \LaTeX{} project for reference.

} % End 13pt content

\clearpage

\end{verbatim}

For example:

\begin{verbatim}
% 
% =================================================================
% REGULATIONS FOR PREPARING THE GRADUATION INTERNSHIP REPORT
% USING LATEX
% =================================================================
%
% This file is intended to be included from the main document with:
%
%     \input{Regulations.tex}
%
% IMPORTANT:
% This section is provided for guidance only. Students shall remove
% or comment out the \input{Regulations.tex} command before submitting
% the final internship report.
%
% =================================================================

\clearpage
\thispagestyle{empty}

\begin{center}
    {\fontsize{16pt}{20pt}\selectfont\bfseries
    SPECIFIC REGULATIONS FOR PREPARING\\
    GRADUATION INTERNSHIP REPORTS USING \LaTeX}
\end{center}

\vspace{0.5cm}

{\fontsize{13pt}{17pt}\selectfont
\setstretch{1.3}
\justifying

\noindent
\textbf{Note:}
This section is provided for guidance only and is included in the main
document using the command \verb|\input{Regulations.tex}|.
Students shall remove or comment out this command before submitting the
final internship report.

\vspace{0.4cm}

% ================================================================
% I. GENERAL REQUIREMENTS
% ================================================================

\noindent
\textbf{I.\quad GENERAL REQUIREMENTS}

\vspace{0.15cm}

The Graduation Internship Report shall be prepared using the official
\LaTeX{} template provided by the Faculty. Students shall preserve the
predefined page layout, typography, line spacing, heading styles,
numbering system, figure and table formatting, and bibliography settings
unless otherwise instructed by the Faculty or the supervisor.

The report shall be written in English.

Students are responsible for the academic and technical content of the
report. The official \LaTeX{} template is responsible for standardized
document formatting.

\vspace{0.3cm}

% ================================================================
% II. REPORT STRUCTURE
% ================================================================

\noindent
\textbf{II.\quad REPORT STRUCTURE}

\vspace{0.15cm}

The report shall normally be arranged in the following order:

\begin{enumerate}[label=\arabic*., leftmargin=1.2cm, itemsep=2pt]
    \item Outer Cover;
    \item Comments of the Internship Host Organization;
    \item Supervisor's Evaluation and Graduation Internship Defense Results;
    \item Table of Contents;
    \item List of Symbols and Abbreviations, if applicable;
    \item List of Figures;
    \item List of Tables;
    \item Abstract;
    \item Main Chapters of the Internship Report;
    \item References;
    \item Appendices, if applicable.
\end{enumerate}

The preliminary pages shall not be numbered as chapters. The main report
content shall be organized into consecutively numbered chapters beginning
with Chapter 1.

\vspace{0.3cm}

% ================================================================
% III. PAGE SIZE AND MARGINS
% ================================================================

\noindent
\textbf{III.\quad PAGE SIZE AND MARGINS}

\vspace{0.15cm}

The report shall be formatted on A4 paper with the following margins:

\begin{center}
\begin{tabular}{ll}
    Left margin:   & 3.5 cm \\
    Right margin:  & 2.0 cm \\
    Top margin:    & 2.5 cm \\
    Bottom margin: & 2.5 cm
\end{tabular}
\end{center}

These settings are predefined in the official template and shall not be
modified by students.

\vspace{0.3cm}

% ================================================================
% IV. FONT, FONT SIZE, AND LINE SPACING
% ================================================================

\noindent
\textbf{IV.\quad FONT, FONT SIZE, AND LINE SPACING}

\vspace{0.15cm}

The main text shall use a 13 pt font size with 1.3-line spacing.
A Times-family typeface shall be used in accordance with the official
template.

The text shall be fully justified.

Students shall not use repeated spaces, repeated manual line breaks, or
multiple blank lines to position text. Paragraph and heading spacing shall
be controlled by the template.

\vspace{0.3cm}

% ================================================================
% V. OUTER COVER
% ================================================================

\noindent
\textbf{V.\quad OUTER COVER}

\vspace{0.15cm}

The outer cover shall follow the official format provided in the template
and shall contain the following information:

\begin{itemize}[leftmargin=0.8cm, itemsep=2pt]
    \item Thai Nguyen University of Technology;
    \item Faculty of International Training;
    \item Engineering and Technology Department;
    \item University logo and Faculty logo;
    \item the title \textbf{INTERNSHIP REPORT};
    \item Internship Organization;
    \item Author;
    \item Class;
    \item Supervisor;
    \item place and year of completion.
\end{itemize}

Students shall only replace the required information fields and shall not
modify the cover-page layout unless instructed otherwise.

\vspace{0.3cm}

% ================================================================
% VI. COMMENTS OF THE INTERNSHIP HOST ORGANIZATION
% ================================================================

\noindent
\textbf{VI.\quad COMMENTS OF THE INTERNSHIP HOST ORGANIZATION}

\vspace{0.15cm}

The page entitled \textbf{COMMENTS OF THE INTERNSHIP HOST ORGANIZATION}
is reserved for the internship host organization.

It includes:

\begin{itemize}[leftmargin=0.8cm, itemsep=2pt]
    \item the student's full name;
    \item class;
    \item comments from the host organization;
    \item score;
    \item date;
    \item signature and official seal.
\end{itemize}

Students shall not alter or complete official evaluation fields unless
specifically authorized.

\vspace{0.3cm}

% ================================================================
% VII. SUPERVISOR'S EVALUATION AND DEFENSE RESULTS
% ================================================================

\noindent
\textbf{VII.\quad SUPERVISOR'S EVALUATION AND DEFENSE RESULTS}

\vspace{0.15cm}

The section entitled
\textbf{SUPERVISOR'S EVALUATION AND GRADUATION INTERNSHIP DEFENSE RESULTS}
is reserved for the supervisor and defense committee.

The evaluation may include:

\begin{itemize}[leftmargin=0.8cm, itemsep=2pt]
    \item Internship Progress and Student's Attitude;
    \item Report Content;
    \item Report Presentation;
    \item Other Comments;
    \item Supervisor's Opinion;
    \item Defense Result;
    \item Score.
\end{itemize}

Students shall not enter evaluation scores or defense results before they
are officially determined.

\vspace{0.3cm}

% ================================================================
% VIII. CHAPTERS, SECTIONS, AND SUBSECTIONS
% ================================================================

\noindent
\textbf{VIII.\quad CHAPTERS, SECTIONS, AND SUBSECTIONS}

\vspace{0.15cm}

The report structure shall be generated using the standard \LaTeX{}
commands:

\begin{verbatim}
\chapter{...}
\section{...}
\subsection{...}
\end{verbatim}

Chapter, section, and subsection numbers shall be generated automatically
by \LaTeX{}.

Students shall not manually type chapter or section numbers and shall not
reproduce heading styles using commands such as \verb|\textbf|,
\verb|\fontsize|, or \verb|\centering| instead of the appropriate
structural command.

Excessive subdivision should be avoided. In most cases, the hierarchy
Chapter--Section--Subsection is sufficient.

\vspace{0.3cm}

% ================================================================
% IX. ORGANIZATION OF THE MAIN CONTENT
% ================================================================

\noindent
\textbf{IX.\quad ORGANIZATION OF THE MAIN CONTENT}

\vspace{0.15cm}

The specific content of the report shall depend on the internship
organization and the assigned tasks. A typical structure may include:

\begin{itemize}[leftmargin=0.8cm, itemsep=2pt]
    \item \textbf{Chapter 1. Overview of the Internship Organization};
    \item \textbf{Chapter 2. Internship Activities and Assigned Tasks};
    \item \textbf{Chapter 3. Results, Experience, and Evaluation}.
\end{itemize}

These chapter titles are recommendations only and may be adapted to the
actual internship content.

The report shall reflect work actually performed, observed, analyzed, or
learned during the internship. Unrelated theoretical material should not
be included merely to increase report length.

\vspace{0.3cm}

% ================================================================
% X. TABLE OF CONTENTS
% ================================================================

\noindent
\textbf{X.\quad TABLE OF CONTENTS}

\vspace{0.15cm}

The Table of Contents shall be generated automatically using:

\begin{verbatim}
\tableofcontents
\end{verbatim}

Students shall not create or manually edit page numbers in the Table of
Contents.

The document shall be compiled sufficiently to ensure that all page
numbers and entries are correctly updated.

\vspace{0.3cm}

% ================================================================
% XI. LIST OF SYMBOLS AND ABBREVIATIONS
% ================================================================

\noindent
\textbf{XI.\quad LIST OF SYMBOLS AND ABBREVIATIONS}

\vspace{0.15cm}

Symbols and abbreviations that are used repeatedly in the report shall be
listed in the \textbf{List of Symbols and Abbreviations}.

The list may contain the following columns:

\begin{center}
\begin{tabular}{llll}
    No. & Symbol / Abbreviation & Unit & Definition
\end{tabular}
\end{center}

Only symbols and abbreviations actually used in the final report shall be
retained. Example entries contained in the template shall be removed if
they are not relevant to the student's report.

Each abbreviation should normally be defined when it first appears in
the main text.

\vspace{0.3cm}

% ================================================================
% XII. FIGURES, DIAGRAMS, AND PHOTOGRAPHS
% ================================================================

\noindent
\textbf{XII.\quad FIGURES, DIAGRAMS, AND PHOTOGRAPHS}

\vspace{0.15cm}

Figures shall be clear, legible, and directly relevant to the report
content.

Figures shall be inserted using the standard \verb|figure| environment.

Example:

\begin{verbatim}
\begin{figure}[H]
    \centering
    \includegraphics[width=0.8\textwidth]
                    {figures/company_structure.pdf}
    \caption{Organizational structure of the company}
    \label{fig:company-structure}
\end{figure}
\end{verbatim}

Figure captions shall be placed below the figure.

Each referenced figure shall have a unique \verb|\label| and shall be
referenced in the text using, for example:

\begin{verbatim}
Figure~\ref{fig:company-structure}
\end{verbatim}

Students shall not manually type figure numbers.

Figures reproduced or adapted from external sources shall include an
appropriate citation.

Photographs taken during the internship shall be relevant and sufficiently
clear. Students shall not include confidential, restricted, or proprietary
information without permission from the internship organization.

Vector graphics or high-resolution images are recommended whenever
possible.

\vspace{0.3cm}

% ================================================================
% XIII. TABLES
% ================================================================

\noindent
\textbf{XIII.\quad TABLES}

\vspace{0.15cm}

Tables shall be created using suitable \LaTeX{} environments such as
\verb|table|, \verb|tabular|, \verb|tabularx|, or \verb|longtable|.

Table captions shall be placed above the table.

Example:

\begin{verbatim}
\begin{table}[H]
    \centering
    \caption{Internship schedule}
    \label{tab:internship-schedule}

    \begin{tabular}{lll}
        \toprule
        Week & Task & Result \\
        \midrule
        1 & Orientation & Completed \\
        2 & System study & Completed \\
        \bottomrule
    \end{tabular}
\end{table}
\end{verbatim}

Tables shall be referenced using:

\begin{verbatim}
Table~\ref{tab:internship-schedule}
\end{verbatim}

Students shall not manually type table numbers.

Tables should be created directly in \LaTeX{} rather than inserted as
screenshots whenever practical.

\vspace{0.3cm}

% ================================================================
% XIV. EQUATIONS AND TECHNICAL QUANTITIES
% ================================================================

\noindent
\textbf{XIV.\quad EQUATIONS AND TECHNICAL QUANTITIES}

\vspace{0.15cm}

Mathematical expressions shall be typeset directly using \LaTeX{}.
Images or screenshots of equations shall not be used.

Short expressions may be written in inline mathematical mode, for example:

\begin{verbatim}
$R = 0.7\,\Omega$
\end{verbatim}

Important equations shall use a numbered equation environment.

Example:

\begin{verbatim}
\begin{equation}
    P = UI
    \label{eq:power}
\end{equation}
\end{verbatim}

Equations shall be referenced using:

\begin{verbatim}
Equation~\eqref{eq:power}
\end{verbatim}

Equation numbers shall not be manually entered.

SI units shall be used wherever applicable. The \verb|siunitx| package
should be used where appropriate.

Examples:

\begin{verbatim}
\SI{220}{\volt}
\SI{50}{\hertz}
\SI{2.5}{\ampere}
\end{verbatim}

Notation and units shall be consistent throughout the report.

\vspace{0.3cm}

% ================================================================
% XV. SOURCE CODE
% ================================================================

\noindent
\textbf{XV.\quad SOURCE CODE}

\vspace{0.15cm}

If the internship involves programming, PLCs, microcontrollers,
simulation, or software development, relevant source code may be included
in the report.

Source code shall be presented using the \verb|listings| environment or
another environment defined by the official template.

Example:

\begin{verbatim}
\begin{lstlisting}[language=st]
IF Start = TRUE THEN
    Motor := TRUE;
ELSE
    Motor := FALSE;
END_IF;
\end{lstlisting}
\end{verbatim}

Only code necessary for explaining the work should appear in the main
report. Long source-code listings should be moved to an appendix or
submitted as supplementary files if required.

Source code should not normally be inserted as screenshots.

\vspace{0.3cm}

% ================================================================
% XVI. ABSTRACT
% ================================================================

\noindent
\textbf{XVI.\quad ABSTRACT}

\vspace{0.15cm}

The Abstract shall briefly summarize:

\begin{itemize}[leftmargin=0.8cm, itemsep=2pt]
    \item the internship organization;
    \item the objectives of the internship;
    \item the main tasks performed;
    \item the main results, knowledge, and skills obtained.
\end{itemize}

The Abstract should be concise and should not reproduce the full content
of the report chapter by chapter.

\vspace{0.3cm}

% ================================================================
% XVII. REFERENCES AND CITATIONS
% ================================================================

\noindent
\textbf{XVII.\quad REFERENCES AND CITATIONS}

\vspace{0.15cm}

All books, journal papers, conference papers, standards, websites,
datasheets, manuals, technical documents, figures, tables, or other
external sources used in preparing the report shall be properly cited.

References shall be managed through the bibliography system provided with
the official template.

Bibliographic entries shall normally be stored in:

\begin{verbatim}
References.bib
\end{verbatim}

References shall be cited in the report using:

\begin{verbatim}
\cite{reference-key}
\end{verbatim}

Students shall not manually type reference numbers such as ``[1]'',
``[2]'', etc.

Titles of foreign-language publications shall remain in their original
language in the bibliography.

All borrowed ideas, data, figures, tables, equations, methods, or results
shall be properly acknowledged.

\vspace{0.3cm}

% ================================================================
% XVIII. CROSS-REFERENCING
% ================================================================

\noindent
\textbf{XVIII.\quad CROSS-REFERENCING}

\vspace{0.15cm}

Numbered elements shall be cross-referenced using the standard \LaTeX{}
mechanism.

Use the following commands where appropriate:

\begin{verbatim}
\label{...}
\ref{...}
\eqref{...}
\end{verbatim}

For example:

\begin{verbatim}
As shown in Figure~\ref{fig:system}, ...
\end{verbatim}

Students shall not manually type numbers such as ``Figure 2.3'' when
those numbers can be generated automatically.

\vspace{0.3cm}

% ================================================================
% XIX. PAGE NUMBERING
% ================================================================

\noindent
\textbf{XIX.\quad PAGE NUMBERING}

\vspace{0.15cm}

Cover pages and official evaluation forms at the beginning of the report
shall not display page numbers.

The main body of the report shall use Arabic page numbers beginning from
Chapter 1.

Page numbering shall be generated automatically by the official template.
Students shall not manually enter page numbers.

\vspace{0.3cm}

% ================================================================
% XX. PROJECT FILE ORGANIZATION
% ================================================================

\noindent
\textbf{XX.\quad PROJECT FILE ORGANIZATION}

\vspace{0.15cm}

Students should maintain a clear and organized \LaTeX{} project
structure. A recommended structure is:

\begin{verbatim}
InternshipReport/
|
|-- main.tex
|-- References.bib
|
|-- figures/
|   |-- company.pdf
|   |-- system.png
|   `-- ...
|
`-- chapters/
    |-- chapter1.tex
    |-- chapter2.tex
    `-- chapter3.tex
\end{verbatim}

Figures should normally be stored in the \verb|figures| directory
specified by the official template.

File and directory names should preferably contain Latin letters,
numbers, hyphens, or underscores.

Recommended examples include:

\begin{verbatim}
company_structure.pdf
motor_control.png
chapter_01.tex
\end{verbatim}

Spaces and unnecessary special characters in file names should be
avoided where possible.

\vspace{0.3cm}

% ================================================================
% XXI. DOCUMENT COMPILATION AND FINAL CHECK
% ================================================================

\noindent
\textbf{XXI.\quad DOCUMENT COMPILATION AND FINAL CHECK}

\vspace{0.15cm}

Before submission, the complete \LaTeX{} project shall compile
successfully and produce the final PDF report.

The final document shall not contain unresolved references such as:

\begin{verbatim}
??
[?]
\end{verbatim}

Before generating the final PDF, students shall check:

\begin{itemize}[leftmargin=0.8cm, itemsep=2pt]
    \item the Table of Contents;
    \item page numbering;
    \item chapter, section, and subsection numbering;
    \item figure numbering;
    \item table numbering;
    \item equation numbering;
    \item cross-references;
    \item citations;
    \item bibliography;
    \item font display;
    \item mathematical symbols;
    \item image quality;
    \item completeness of all required forms.
\end{itemize}

The document may be compiled using TeX Live, MiKTeX, or another
compatible \LaTeX{} environment. Regardless of the compilation
environment, the resulting PDF shall conform to the official Faculty
template.

\vspace{0.3cm}

% ================================================================
% XXII. CONTENT THAT STUDENTS MAY EDIT
% ================================================================

\noindent
\textbf{XXII.\quad CONTENT THAT STUDENTS MAY EDIT}

\vspace{0.15cm}

Students may edit or replace:

\begin{itemize}[leftmargin=0.8cm, itemsep=2pt]
    \item student name;
    \item class;
    \item internship organization;
    \item supervisor name;
    \item Abstract;
    \item List of Symbols and Abbreviations;
    \item chapter titles and chapter content;
    \item figures and photographs;
    \item tables;
    \item equations;
    \item citations and references.
\end{itemize}

Students shall not modify the main document formatting settings in the
preamble, including the document class, page geometry, line spacing,
heading formatting, table-of-contents formatting, font configuration,
page layout, or cover-page macros unless specifically instructed.

\vspace{0.3cm}

% ================================================================
% XXIII. FINAL SUBMISSION REQUIREMENTS
% ================================================================

\noindent
\textbf{XXIII.\quad FINAL SUBMISSION REQUIREMENTS}

\vspace{0.15cm}

Before submission, the final report shall:

\begin{itemize}[leftmargin=0.8cm, itemsep=2pt]
    \item compile without errors that affect the report content or presentation;
    \item contain all required forms and report sections;
    \item contain correct student, class, internship organization, and supervisor information;
    \item contain an updated Table of Contents, List of Figures, and List of Tables;
    \item contain no unused example content from the template;
    \item contain no placeholder text such as \textit{Title of Chapter 1},
    \textit{Section 1.1}, or other sample entries;
    \item contain all required citations and references;
    \item contain no unresolved cross-references.
\end{itemize}

If the Faculty requires submission of the \LaTeX{} source, the complete
project shall include all \verb|.tex| files, the bibliography database,
figures, and any other files required to reproduce the submitted PDF.

\vspace{0.4cm}

\noindent
\textbf{Important:}
Before submitting the final internship report, students shall remove or
comment out the following command in the main document:

\begin{verbatim}
\input{Regulations.tex}
\end{verbatim}

For example:

\begin{verbatim}
% \input{Regulations.tex}
\end{verbatim}

This removes the guidance pages from the final submitted report while
preserving the regulations file in the \LaTeX{} project for reference.

} % End 13pt content

\clearpage

\end{verbatim}

This removes the guidance pages from the final submitted report while
preserving the regulations file in the \LaTeX{} project for reference.

} % End 13pt content

\clearpage
|.
Students shall remove or comment out this command before submitting the
final internship report.

\vspace{0.4cm}

% ================================================================
% I. GENERAL REQUIREMENTS
% ================================================================

\noindent
\textbf{I.\quad GENERAL REQUIREMENTS}

\vspace{0.15cm}

The Graduation Internship Report shall be prepared using the official
\LaTeX{} template provided by the Faculty. Students shall preserve the
predefined page layout, typography, line spacing, heading styles,
numbering system, figure and table formatting, and bibliography settings
unless otherwise instructed by the Faculty or the supervisor.

The report shall be written in English.

Students are responsible for the academic and technical content of the
report. The official \LaTeX{} template is responsible for standardized
document formatting.

\vspace{0.3cm}

% ================================================================
% II. REPORT STRUCTURE
% ================================================================

\noindent
\textbf{II.\quad REPORT STRUCTURE}

\vspace{0.15cm}

The report shall normally be arranged in the following order:

\begin{enumerate}[label=\arabic*., leftmargin=1.2cm, itemsep=2pt]
    \item Outer Cover;
    \item Comments of the Internship Host Organization;
    \item Supervisor's Evaluation and Graduation Internship Defense Results;
    \item Table of Contents;
    \item List of Symbols and Abbreviations, if applicable;
    \item List of Figures;
    \item List of Tables;
    \item Abstract;
    \item Main Chapters of the Internship Report;
    \item References;
    \item Appendices, if applicable.
\end{enumerate}

The preliminary pages shall not be numbered as chapters. The main report
content shall be organized into consecutively numbered chapters beginning
with Chapter 1.

\vspace{0.3cm}

% ================================================================
% III. PAGE SIZE AND MARGINS
% ================================================================

\noindent
\textbf{III.\quad PAGE SIZE AND MARGINS}

\vspace{0.15cm}

The report shall be formatted on A4 paper with the following margins:

\begin{center}
\begin{tabular}{ll}
    Left margin:   & 3.5 cm \\
    Right margin:  & 2.0 cm \\
    Top margin:    & 2.5 cm \\
    Bottom margin: & 2.5 cm
\end{tabular}
\end{center}

These settings are predefined in the official template and shall not be
modified by students.

\vspace{0.3cm}

% ================================================================
% IV. FONT, FONT SIZE, AND LINE SPACING
% ================================================================

\noindent
\textbf{IV.\quad FONT, FONT SIZE, AND LINE SPACING}

\vspace{0.15cm}

The main text shall use a 13 pt font size with 1.3-line spacing.
A Times-family typeface shall be used in accordance with the official
template.

The text shall be fully justified.

Students shall not use repeated spaces, repeated manual line breaks, or
multiple blank lines to position text. Paragraph and heading spacing shall
be controlled by the template.

\vspace{0.3cm}

% ================================================================
% V. OUTER COVER
% ================================================================

\noindent
\textbf{V.\quad OUTER COVER}

\vspace{0.15cm}

The outer cover shall follow the official format provided in the template
and shall contain the following information:

\begin{itemize}[leftmargin=0.8cm, itemsep=2pt]
    \item Thai Nguyen University of Technology;
    \item Faculty of International Training;
    \item Engineering and Technology Department;
    \item University logo and Faculty logo;
    \item the title \textbf{INTERNSHIP REPORT};
    \item Internship Organization;
    \item Author;
    \item Class;
    \item Supervisor;
    \item place and year of completion.
\end{itemize}

Students shall only replace the required information fields and shall not
modify the cover-page layout unless instructed otherwise.

\vspace{0.3cm}

% ================================================================
% VI. COMMENTS OF THE INTERNSHIP HOST ORGANIZATION
% ================================================================

\noindent
\textbf{VI.\quad COMMENTS OF THE INTERNSHIP HOST ORGANIZATION}

\vspace{0.15cm}

The page entitled \textbf{COMMENTS OF THE INTERNSHIP HOST ORGANIZATION}
is reserved for the internship host organization.

It includes:

\begin{itemize}[leftmargin=0.8cm, itemsep=2pt]
    \item the student's full name;
    \item class;
    \item comments from the host organization;
    \item score;
    \item date;
    \item signature and official seal.
\end{itemize}

Students shall not alter or complete official evaluation fields unless
specifically authorized.

\vspace{0.3cm}

% ================================================================
% VII. SUPERVISOR'S EVALUATION AND DEFENSE RESULTS
% ================================================================

\noindent
\textbf{VII.\quad SUPERVISOR'S EVALUATION AND DEFENSE RESULTS}

\vspace{0.15cm}

The section entitled
\textbf{SUPERVISOR'S EVALUATION AND GRADUATION INTERNSHIP DEFENSE RESULTS}
is reserved for the supervisor and defense committee.

The evaluation may include:

\begin{itemize}[leftmargin=0.8cm, itemsep=2pt]
    \item Internship Progress and Student's Attitude;
    \item Report Content;
    \item Report Presentation;
    \item Other Comments;
    \item Supervisor's Opinion;
    \item Defense Result;
    \item Score.
\end{itemize}

Students shall not enter evaluation scores or defense results before they
are officially determined.

\vspace{0.3cm}

% ================================================================
% VIII. CHAPTERS, SECTIONS, AND SUBSECTIONS
% ================================================================

\noindent
\textbf{VIII.\quad CHAPTERS, SECTIONS, AND SUBSECTIONS}

\vspace{0.15cm}

The report structure shall be generated using the standard \LaTeX{}
commands:

\begin{verbatim}
\chapter{...}
\section{...}
\subsection{...}
\end{verbatim}

Chapter, section, and subsection numbers shall be generated automatically
by \LaTeX{}.

Students shall not manually type chapter or section numbers and shall not
reproduce heading styles using commands such as \verb|\textbf|,
\verb|\fontsize|, or \verb|\centering| instead of the appropriate
structural command.

Excessive subdivision should be avoided. In most cases, the hierarchy
Chapter--Section--Subsection is sufficient.

\vspace{0.3cm}

% ================================================================
% IX. ORGANIZATION OF THE MAIN CONTENT
% ================================================================

\noindent
\textbf{IX.\quad ORGANIZATION OF THE MAIN CONTENT}

\vspace{0.15cm}

The specific content of the report shall depend on the internship
organization and the assigned tasks. A typical structure may include:

\begin{itemize}[leftmargin=0.8cm, itemsep=2pt]
    \item \textbf{Chapter 1. Overview of the Internship Organization};
    \item \textbf{Chapter 2. Internship Activities and Assigned Tasks};
    \item \textbf{Chapter 3. Results, Experience, and Evaluation}.
\end{itemize}

These chapter titles are recommendations only and may be adapted to the
actual internship content.

The report shall reflect work actually performed, observed, analyzed, or
learned during the internship. Unrelated theoretical material should not
be included merely to increase report length.

\vspace{0.3cm}

% ================================================================
% X. TABLE OF CONTENTS
% ================================================================

\noindent
\textbf{X.\quad TABLE OF CONTENTS}

\vspace{0.15cm}

The Table of Contents shall be generated automatically using:

\begin{verbatim}
\tableofcontents
\end{verbatim}

Students shall not create or manually edit page numbers in the Table of
Contents.

The document shall be compiled sufficiently to ensure that all page
numbers and entries are correctly updated.

\vspace{0.3cm}

% ================================================================
% XI. LIST OF SYMBOLS AND ABBREVIATIONS
% ================================================================

\noindent
\textbf{XI.\quad LIST OF SYMBOLS AND ABBREVIATIONS}

\vspace{0.15cm}

Symbols and abbreviations that are used repeatedly in the report shall be
listed in the \textbf{List of Symbols and Abbreviations}.

The list may contain the following columns:

\begin{center}
\begin{tabular}{llll}
    No. & Symbol / Abbreviation & Unit & Definition
\end{tabular}
\end{center}

Only symbols and abbreviations actually used in the final report shall be
retained. Example entries contained in the template shall be removed if
they are not relevant to the student's report.

Each abbreviation should normally be defined when it first appears in
the main text.

\vspace{0.3cm}

% ================================================================
% XII. FIGURES, DIAGRAMS, AND PHOTOGRAPHS
% ================================================================

\noindent
\textbf{XII.\quad FIGURES, DIAGRAMS, AND PHOTOGRAPHS}

\vspace{0.15cm}

Figures shall be clear, legible, and directly relevant to the report
content.

Figures shall be inserted using the standard \verb|figure| environment.

Example:

\begin{verbatim}
\begin{figure}[H]
    \centering
    \includegraphics[width=0.8\textwidth]
                    {figures/company_structure.pdf}
    \caption{Organizational structure of the company}
    \label{fig:company-structure}
\end{figure}
\end{verbatim}

Figure captions shall be placed below the figure.

Each referenced figure shall have a unique \verb|\label| and shall be
referenced in the text using, for example:

\begin{verbatim}
Figure~\ref{fig:company-structure}
\end{verbatim}

Students shall not manually type figure numbers.

Figures reproduced or adapted from external sources shall include an
appropriate citation.

Photographs taken during the internship shall be relevant and sufficiently
clear. Students shall not include confidential, restricted, or proprietary
information without permission from the internship organization.

Vector graphics or high-resolution images are recommended whenever
possible.

\vspace{0.3cm}

% ================================================================
% XIII. TABLES
% ================================================================

\noindent
\textbf{XIII.\quad TABLES}

\vspace{0.15cm}

Tables shall be created using suitable \LaTeX{} environments such as
\verb|table|, \verb|tabular|, \verb|tabularx|, or \verb|longtable|.

Table captions shall be placed above the table.

Example:

\begin{verbatim}
\begin{table}[H]
    \centering
    \caption{Internship schedule}
    \label{tab:internship-schedule}

    \begin{tabular}{lll}
        \toprule
        Week & Task & Result \\
        \midrule
        1 & Orientation & Completed \\
        2 & System study & Completed \\
        \bottomrule
    \end{tabular}
\end{table}
\end{verbatim}

Tables shall be referenced using:

\begin{verbatim}
Table~\ref{tab:internship-schedule}
\end{verbatim}

Students shall not manually type table numbers.

Tables should be created directly in \LaTeX{} rather than inserted as
screenshots whenever practical.

\vspace{0.3cm}

% ================================================================
% XIV. EQUATIONS AND TECHNICAL QUANTITIES
% ================================================================

\noindent
\textbf{XIV.\quad EQUATIONS AND TECHNICAL QUANTITIES}

\vspace{0.15cm}

Mathematical expressions shall be typeset directly using \LaTeX{}.
Images or screenshots of equations shall not be used.

Short expressions may be written in inline mathematical mode, for example:

\begin{verbatim}
$R = 0.7\,\Omega$
\end{verbatim}

Important equations shall use a numbered equation environment.

Example:

\begin{verbatim}
\begin{equation}
    P = UI
    \label{eq:power}
\end{equation}
\end{verbatim}

Equations shall be referenced using:

\begin{verbatim}
Equation~\eqref{eq:power}
\end{verbatim}

Equation numbers shall not be manually entered.

SI units shall be used wherever applicable. The \verb|siunitx| package
should be used where appropriate.

Examples:

\begin{verbatim}
\SI{220}{\volt}
\SI{50}{\hertz}
\SI{2.5}{\ampere}
\end{verbatim}

Notation and units shall be consistent throughout the report.

\vspace{0.3cm}

% ================================================================
% XV. SOURCE CODE
% ================================================================

\noindent
\textbf{XV.\quad SOURCE CODE}

\vspace{0.15cm}

If the internship involves programming, PLCs, microcontrollers,
simulation, or software development, relevant source code may be included
in the report.

Source code shall be presented using the \verb|listings| environment or
another environment defined by the official template.

Example:

\begin{verbatim}
\begin{lstlisting}[language=st]
IF Start = TRUE THEN
    Motor := TRUE;
ELSE
    Motor := FALSE;
END_IF;
\end{lstlisting}
\end{verbatim}

Only code necessary for explaining the work should appear in the main
report. Long source-code listings should be moved to an appendix or
submitted as supplementary files if required.

Source code should not normally be inserted as screenshots.

\vspace{0.3cm}

% ================================================================
% XVI. ABSTRACT
% ================================================================

\noindent
\textbf{XVI.\quad ABSTRACT}

\vspace{0.15cm}

The Abstract shall briefly summarize:

\begin{itemize}[leftmargin=0.8cm, itemsep=2pt]
    \item the internship organization;
    \item the objectives of the internship;
    \item the main tasks performed;
    \item the main results, knowledge, and skills obtained.
\end{itemize}

The Abstract should be concise and should not reproduce the full content
of the report chapter by chapter.

\vspace{0.3cm}

% ================================================================
% XVII. REFERENCES AND CITATIONS
% ================================================================

\noindent
\textbf{XVII.\quad REFERENCES AND CITATIONS}

\vspace{0.15cm}

All books, journal papers, conference papers, standards, websites,
datasheets, manuals, technical documents, figures, tables, or other
external sources used in preparing the report shall be properly cited.

References shall be managed through the bibliography system provided with
the official template.

Bibliographic entries shall normally be stored in:

\begin{verbatim}
References.bib
\end{verbatim}

References shall be cited in the report using:

\begin{verbatim}
\cite{reference-key}
\end{verbatim}

Students shall not manually type reference numbers such as ``[1]'',
``[2]'', etc.

Titles of foreign-language publications shall remain in their original
language in the bibliography.

All borrowed ideas, data, figures, tables, equations, methods, or results
shall be properly acknowledged.

\vspace{0.3cm}

% ================================================================
% XVIII. CROSS-REFERENCING
% ================================================================

\noindent
\textbf{XVIII.\quad CROSS-REFERENCING}

\vspace{0.15cm}

Numbered elements shall be cross-referenced using the standard \LaTeX{}
mechanism.

Use the following commands where appropriate:

\begin{verbatim}
\label{...}
\ref{...}
\eqref{...}
\end{verbatim}

For example:

\begin{verbatim}
As shown in Figure~\ref{fig:system}, ...
\end{verbatim}

Students shall not manually type numbers such as ``Figure 2.3'' when
those numbers can be generated automatically.

\vspace{0.3cm}

% ================================================================
% XIX. PAGE NUMBERING
% ================================================================

\noindent
\textbf{XIX.\quad PAGE NUMBERING}

\vspace{0.15cm}

Cover pages and official evaluation forms at the beginning of the report
shall not display page numbers.

The main body of the report shall use Arabic page numbers beginning from
Chapter 1.

Page numbering shall be generated automatically by the official template.
Students shall not manually enter page numbers.

\vspace{0.3cm}

% ================================================================
% XX. PROJECT FILE ORGANIZATION
% ================================================================

\noindent
\textbf{XX.\quad PROJECT FILE ORGANIZATION}

\vspace{0.15cm}

Students should maintain a clear and organized \LaTeX{} project
structure. A recommended structure is:

\begin{verbatim}
InternshipReport/
|
|-- main.tex
|-- References.bib
|
|-- figures/
|   |-- company.pdf
|   |-- system.png
|   `-- ...
|
`-- chapters/
    |-- chapter1.tex
    |-- chapter2.tex
    `-- chapter3.tex
\end{verbatim}

Figures should normally be stored in the \verb|figures| directory
specified by the official template.

File and directory names should preferably contain Latin letters,
numbers, hyphens, or underscores.

Recommended examples include:

\begin{verbatim}
company_structure.pdf
motor_control.png
chapter_01.tex
\end{verbatim}

Spaces and unnecessary special characters in file names should be
avoided where possible.

\vspace{0.3cm}

% ================================================================
% XXI. DOCUMENT COMPILATION AND FINAL CHECK
% ================================================================

\noindent
\textbf{XXI.\quad DOCUMENT COMPILATION AND FINAL CHECK}

\vspace{0.15cm}

Before submission, the complete \LaTeX{} project shall compile
successfully and produce the final PDF report.

The final document shall not contain unresolved references such as:

\begin{verbatim}
??
[?]
\end{verbatim}

Before generating the final PDF, students shall check:

\begin{itemize}[leftmargin=0.8cm, itemsep=2pt]
    \item the Table of Contents;
    \item page numbering;
    \item chapter, section, and subsection numbering;
    \item figure numbering;
    \item table numbering;
    \item equation numbering;
    \item cross-references;
    \item citations;
    \item bibliography;
    \item font display;
    \item mathematical symbols;
    \item image quality;
    \item completeness of all required forms.
\end{itemize}

The document may be compiled using TeX Live, MiKTeX, or another
compatible \LaTeX{} environment. Regardless of the compilation
environment, the resulting PDF shall conform to the official Faculty
template.

\vspace{0.3cm}

% ================================================================
% XXII. CONTENT THAT STUDENTS MAY EDIT
% ================================================================

\noindent
\textbf{XXII.\quad CONTENT THAT STUDENTS MAY EDIT}

\vspace{0.15cm}

Students may edit or replace:

\begin{itemize}[leftmargin=0.8cm, itemsep=2pt]
    \item student name;
    \item class;
    \item internship organization;
    \item supervisor name;
    \item Abstract;
    \item List of Symbols and Abbreviations;
    \item chapter titles and chapter content;
    \item figures and photographs;
    \item tables;
    \item equations;
    \item citations and references.
\end{itemize}

Students shall not modify the main document formatting settings in the
preamble, including the document class, page geometry, line spacing,
heading formatting, table-of-contents formatting, font configuration,
page layout, or cover-page macros unless specifically instructed.

\vspace{0.3cm}

% ================================================================
% XXIII. FINAL SUBMISSION REQUIREMENTS
% ================================================================

\noindent
\textbf{XXIII.\quad FINAL SUBMISSION REQUIREMENTS}

\vspace{0.15cm}

Before submission, the final report shall:

\begin{itemize}[leftmargin=0.8cm, itemsep=2pt]
    \item compile without errors that affect the report content or presentation;
    \item contain all required forms and report sections;
    \item contain correct student, class, internship organization, and supervisor information;
    \item contain an updated Table of Contents, List of Figures, and List of Tables;
    \item contain no unused example content from the template;
    \item contain no placeholder text such as \textit{Title of Chapter 1},
    \textit{Section 1.1}, or other sample entries;
    \item contain all required citations and references;
    \item contain no unresolved cross-references.
\end{itemize}

If the Faculty requires submission of the \LaTeX{} source, the complete
project shall include all \verb|.tex| files, the bibliography database,
figures, and any other files required to reproduce the submitted PDF.

\vspace{0.4cm}

\noindent
\textbf{Important:}
Before submitting the final internship report, students shall remove or
comment out the following command in the main document:

\begin{verbatim}

% =================================================================
% REGULATIONS FOR PREPARING THE GRADUATION INTERNSHIP REPORT
% USING LATEX
% =================================================================
%
% This file is intended to be included from the main document with:
%
%     
% =================================================================
% REGULATIONS FOR PREPARING THE GRADUATION INTERNSHIP REPORT
% USING LATEX
% =================================================================
%
% This file is intended to be included from the main document with:
%
%     \input{Regulations.tex}
%
% IMPORTANT:
% This section is provided for guidance only. Students shall remove
% or comment out the \input{Regulations.tex} command before submitting
% the final internship report.
%
% =================================================================

\clearpage
\thispagestyle{empty}

\begin{center}
    {\fontsize{16pt}{20pt}\selectfont\bfseries
    SPECIFIC REGULATIONS FOR PREPARING\\
    GRADUATION INTERNSHIP REPORTS USING \LaTeX}
\end{center}

\vspace{0.5cm}

{\fontsize{13pt}{17pt}\selectfont
\setstretch{1.3}
\justifying

\noindent
\textbf{Note:}
This section is provided for guidance only and is included in the main
document using the command \verb|\input{Regulations.tex}|.
Students shall remove or comment out this command before submitting the
final internship report.

\vspace{0.4cm}

% ================================================================
% I. GENERAL REQUIREMENTS
% ================================================================

\noindent
\textbf{I.\quad GENERAL REQUIREMENTS}

\vspace{0.15cm}

The Graduation Internship Report shall be prepared using the official
\LaTeX{} template provided by the Faculty. Students shall preserve the
predefined page layout, typography, line spacing, heading styles,
numbering system, figure and table formatting, and bibliography settings
unless otherwise instructed by the Faculty or the supervisor.

The report shall be written in English.

Students are responsible for the academic and technical content of the
report. The official \LaTeX{} template is responsible for standardized
document formatting.

\vspace{0.3cm}

% ================================================================
% II. REPORT STRUCTURE
% ================================================================

\noindent
\textbf{II.\quad REPORT STRUCTURE}

\vspace{0.15cm}

The report shall normally be arranged in the following order:

\begin{enumerate}[label=\arabic*., leftmargin=1.2cm, itemsep=2pt]
    \item Outer Cover;
    \item Comments of the Internship Host Organization;
    \item Supervisor's Evaluation and Graduation Internship Defense Results;
    \item Table of Contents;
    \item List of Symbols and Abbreviations, if applicable;
    \item List of Figures;
    \item List of Tables;
    \item Abstract;
    \item Main Chapters of the Internship Report;
    \item References;
    \item Appendices, if applicable.
\end{enumerate}

The preliminary pages shall not be numbered as chapters. The main report
content shall be organized into consecutively numbered chapters beginning
with Chapter 1.

\vspace{0.3cm}

% ================================================================
% III. PAGE SIZE AND MARGINS
% ================================================================

\noindent
\textbf{III.\quad PAGE SIZE AND MARGINS}

\vspace{0.15cm}

The report shall be formatted on A4 paper with the following margins:

\begin{center}
\begin{tabular}{ll}
    Left margin:   & 3.5 cm \\
    Right margin:  & 2.0 cm \\
    Top margin:    & 2.5 cm \\
    Bottom margin: & 2.5 cm
\end{tabular}
\end{center}

These settings are predefined in the official template and shall not be
modified by students.

\vspace{0.3cm}

% ================================================================
% IV. FONT, FONT SIZE, AND LINE SPACING
% ================================================================

\noindent
\textbf{IV.\quad FONT, FONT SIZE, AND LINE SPACING}

\vspace{0.15cm}

The main text shall use a 13 pt font size with 1.3-line spacing.
A Times-family typeface shall be used in accordance with the official
template.

The text shall be fully justified.

Students shall not use repeated spaces, repeated manual line breaks, or
multiple blank lines to position text. Paragraph and heading spacing shall
be controlled by the template.

\vspace{0.3cm}

% ================================================================
% V. OUTER COVER
% ================================================================

\noindent
\textbf{V.\quad OUTER COVER}

\vspace{0.15cm}

The outer cover shall follow the official format provided in the template
and shall contain the following information:

\begin{itemize}[leftmargin=0.8cm, itemsep=2pt]
    \item Thai Nguyen University of Technology;
    \item Faculty of International Training;
    \item Engineering and Technology Department;
    \item University logo and Faculty logo;
    \item the title \textbf{INTERNSHIP REPORT};
    \item Internship Organization;
    \item Author;
    \item Class;
    \item Supervisor;
    \item place and year of completion.
\end{itemize}

Students shall only replace the required information fields and shall not
modify the cover-page layout unless instructed otherwise.

\vspace{0.3cm}

% ================================================================
% VI. COMMENTS OF THE INTERNSHIP HOST ORGANIZATION
% ================================================================

\noindent
\textbf{VI.\quad COMMENTS OF THE INTERNSHIP HOST ORGANIZATION}

\vspace{0.15cm}

The page entitled \textbf{COMMENTS OF THE INTERNSHIP HOST ORGANIZATION}
is reserved for the internship host organization.

It includes:

\begin{itemize}[leftmargin=0.8cm, itemsep=2pt]
    \item the student's full name;
    \item class;
    \item comments from the host organization;
    \item score;
    \item date;
    \item signature and official seal.
\end{itemize}

Students shall not alter or complete official evaluation fields unless
specifically authorized.

\vspace{0.3cm}

% ================================================================
% VII. SUPERVISOR'S EVALUATION AND DEFENSE RESULTS
% ================================================================

\noindent
\textbf{VII.\quad SUPERVISOR'S EVALUATION AND DEFENSE RESULTS}

\vspace{0.15cm}

The section entitled
\textbf{SUPERVISOR'S EVALUATION AND GRADUATION INTERNSHIP DEFENSE RESULTS}
is reserved for the supervisor and defense committee.

The evaluation may include:

\begin{itemize}[leftmargin=0.8cm, itemsep=2pt]
    \item Internship Progress and Student's Attitude;
    \item Report Content;
    \item Report Presentation;
    \item Other Comments;
    \item Supervisor's Opinion;
    \item Defense Result;
    \item Score.
\end{itemize}

Students shall not enter evaluation scores or defense results before they
are officially determined.

\vspace{0.3cm}

% ================================================================
% VIII. CHAPTERS, SECTIONS, AND SUBSECTIONS
% ================================================================

\noindent
\textbf{VIII.\quad CHAPTERS, SECTIONS, AND SUBSECTIONS}

\vspace{0.15cm}

The report structure shall be generated using the standard \LaTeX{}
commands:

\begin{verbatim}
\chapter{...}
\section{...}
\subsection{...}
\end{verbatim}

Chapter, section, and subsection numbers shall be generated automatically
by \LaTeX{}.

Students shall not manually type chapter or section numbers and shall not
reproduce heading styles using commands such as \verb|\textbf|,
\verb|\fontsize|, or \verb|\centering| instead of the appropriate
structural command.

Excessive subdivision should be avoided. In most cases, the hierarchy
Chapter--Section--Subsection is sufficient.

\vspace{0.3cm}

% ================================================================
% IX. ORGANIZATION OF THE MAIN CONTENT
% ================================================================

\noindent
\textbf{IX.\quad ORGANIZATION OF THE MAIN CONTENT}

\vspace{0.15cm}

The specific content of the report shall depend on the internship
organization and the assigned tasks. A typical structure may include:

\begin{itemize}[leftmargin=0.8cm, itemsep=2pt]
    \item \textbf{Chapter 1. Overview of the Internship Organization};
    \item \textbf{Chapter 2. Internship Activities and Assigned Tasks};
    \item \textbf{Chapter 3. Results, Experience, and Evaluation}.
\end{itemize}

These chapter titles are recommendations only and may be adapted to the
actual internship content.

The report shall reflect work actually performed, observed, analyzed, or
learned during the internship. Unrelated theoretical material should not
be included merely to increase report length.

\vspace{0.3cm}

% ================================================================
% X. TABLE OF CONTENTS
% ================================================================

\noindent
\textbf{X.\quad TABLE OF CONTENTS}

\vspace{0.15cm}

The Table of Contents shall be generated automatically using:

\begin{verbatim}
\tableofcontents
\end{verbatim}

Students shall not create or manually edit page numbers in the Table of
Contents.

The document shall be compiled sufficiently to ensure that all page
numbers and entries are correctly updated.

\vspace{0.3cm}

% ================================================================
% XI. LIST OF SYMBOLS AND ABBREVIATIONS
% ================================================================

\noindent
\textbf{XI.\quad LIST OF SYMBOLS AND ABBREVIATIONS}

\vspace{0.15cm}

Symbols and abbreviations that are used repeatedly in the report shall be
listed in the \textbf{List of Symbols and Abbreviations}.

The list may contain the following columns:

\begin{center}
\begin{tabular}{llll}
    No. & Symbol / Abbreviation & Unit & Definition
\end{tabular}
\end{center}

Only symbols and abbreviations actually used in the final report shall be
retained. Example entries contained in the template shall be removed if
they are not relevant to the student's report.

Each abbreviation should normally be defined when it first appears in
the main text.

\vspace{0.3cm}

% ================================================================
% XII. FIGURES, DIAGRAMS, AND PHOTOGRAPHS
% ================================================================

\noindent
\textbf{XII.\quad FIGURES, DIAGRAMS, AND PHOTOGRAPHS}

\vspace{0.15cm}

Figures shall be clear, legible, and directly relevant to the report
content.

Figures shall be inserted using the standard \verb|figure| environment.

Example:

\begin{verbatim}
\begin{figure}[H]
    \centering
    \includegraphics[width=0.8\textwidth]
                    {figures/company_structure.pdf}
    \caption{Organizational structure of the company}
    \label{fig:company-structure}
\end{figure}
\end{verbatim}

Figure captions shall be placed below the figure.

Each referenced figure shall have a unique \verb|\label| and shall be
referenced in the text using, for example:

\begin{verbatim}
Figure~\ref{fig:company-structure}
\end{verbatim}

Students shall not manually type figure numbers.

Figures reproduced or adapted from external sources shall include an
appropriate citation.

Photographs taken during the internship shall be relevant and sufficiently
clear. Students shall not include confidential, restricted, or proprietary
information without permission from the internship organization.

Vector graphics or high-resolution images are recommended whenever
possible.

\vspace{0.3cm}

% ================================================================
% XIII. TABLES
% ================================================================

\noindent
\textbf{XIII.\quad TABLES}

\vspace{0.15cm}

Tables shall be created using suitable \LaTeX{} environments such as
\verb|table|, \verb|tabular|, \verb|tabularx|, or \verb|longtable|.

Table captions shall be placed above the table.

Example:

\begin{verbatim}
\begin{table}[H]
    \centering
    \caption{Internship schedule}
    \label{tab:internship-schedule}

    \begin{tabular}{lll}
        \toprule
        Week & Task & Result \\
        \midrule
        1 & Orientation & Completed \\
        2 & System study & Completed \\
        \bottomrule
    \end{tabular}
\end{table}
\end{verbatim}

Tables shall be referenced using:

\begin{verbatim}
Table~\ref{tab:internship-schedule}
\end{verbatim}

Students shall not manually type table numbers.

Tables should be created directly in \LaTeX{} rather than inserted as
screenshots whenever practical.

\vspace{0.3cm}

% ================================================================
% XIV. EQUATIONS AND TECHNICAL QUANTITIES
% ================================================================

\noindent
\textbf{XIV.\quad EQUATIONS AND TECHNICAL QUANTITIES}

\vspace{0.15cm}

Mathematical expressions shall be typeset directly using \LaTeX{}.
Images or screenshots of equations shall not be used.

Short expressions may be written in inline mathematical mode, for example:

\begin{verbatim}
$R = 0.7\,\Omega$
\end{verbatim}

Important equations shall use a numbered equation environment.

Example:

\begin{verbatim}
\begin{equation}
    P = UI
    \label{eq:power}
\end{equation}
\end{verbatim}

Equations shall be referenced using:

\begin{verbatim}
Equation~\eqref{eq:power}
\end{verbatim}

Equation numbers shall not be manually entered.

SI units shall be used wherever applicable. The \verb|siunitx| package
should be used where appropriate.

Examples:

\begin{verbatim}
\SI{220}{\volt}
\SI{50}{\hertz}
\SI{2.5}{\ampere}
\end{verbatim}

Notation and units shall be consistent throughout the report.

\vspace{0.3cm}

% ================================================================
% XV. SOURCE CODE
% ================================================================

\noindent
\textbf{XV.\quad SOURCE CODE}

\vspace{0.15cm}

If the internship involves programming, PLCs, microcontrollers,
simulation, or software development, relevant source code may be included
in the report.

Source code shall be presented using the \verb|listings| environment or
another environment defined by the official template.

Example:

\begin{verbatim}
\begin{lstlisting}[language=st]
IF Start = TRUE THEN
    Motor := TRUE;
ELSE
    Motor := FALSE;
END_IF;
\end{lstlisting}
\end{verbatim}

Only code necessary for explaining the work should appear in the main
report. Long source-code listings should be moved to an appendix or
submitted as supplementary files if required.

Source code should not normally be inserted as screenshots.

\vspace{0.3cm}

% ================================================================
% XVI. ABSTRACT
% ================================================================

\noindent
\textbf{XVI.\quad ABSTRACT}

\vspace{0.15cm}

The Abstract shall briefly summarize:

\begin{itemize}[leftmargin=0.8cm, itemsep=2pt]
    \item the internship organization;
    \item the objectives of the internship;
    \item the main tasks performed;
    \item the main results, knowledge, and skills obtained.
\end{itemize}

The Abstract should be concise and should not reproduce the full content
of the report chapter by chapter.

\vspace{0.3cm}

% ================================================================
% XVII. REFERENCES AND CITATIONS
% ================================================================

\noindent
\textbf{XVII.\quad REFERENCES AND CITATIONS}

\vspace{0.15cm}

All books, journal papers, conference papers, standards, websites,
datasheets, manuals, technical documents, figures, tables, or other
external sources used in preparing the report shall be properly cited.

References shall be managed through the bibliography system provided with
the official template.

Bibliographic entries shall normally be stored in:

\begin{verbatim}
References.bib
\end{verbatim}

References shall be cited in the report using:

\begin{verbatim}
\cite{reference-key}
\end{verbatim}

Students shall not manually type reference numbers such as ``[1]'',
``[2]'', etc.

Titles of foreign-language publications shall remain in their original
language in the bibliography.

All borrowed ideas, data, figures, tables, equations, methods, or results
shall be properly acknowledged.

\vspace{0.3cm}

% ================================================================
% XVIII. CROSS-REFERENCING
% ================================================================

\noindent
\textbf{XVIII.\quad CROSS-REFERENCING}

\vspace{0.15cm}

Numbered elements shall be cross-referenced using the standard \LaTeX{}
mechanism.

Use the following commands where appropriate:

\begin{verbatim}
\label{...}
\ref{...}
\eqref{...}
\end{verbatim}

For example:

\begin{verbatim}
As shown in Figure~\ref{fig:system}, ...
\end{verbatim}

Students shall not manually type numbers such as ``Figure 2.3'' when
those numbers can be generated automatically.

\vspace{0.3cm}

% ================================================================
% XIX. PAGE NUMBERING
% ================================================================

\noindent
\textbf{XIX.\quad PAGE NUMBERING}

\vspace{0.15cm}

Cover pages and official evaluation forms at the beginning of the report
shall not display page numbers.

The main body of the report shall use Arabic page numbers beginning from
Chapter 1.

Page numbering shall be generated automatically by the official template.
Students shall not manually enter page numbers.

\vspace{0.3cm}

% ================================================================
% XX. PROJECT FILE ORGANIZATION
% ================================================================

\noindent
\textbf{XX.\quad PROJECT FILE ORGANIZATION}

\vspace{0.15cm}

Students should maintain a clear and organized \LaTeX{} project
structure. A recommended structure is:

\begin{verbatim}
InternshipReport/
|
|-- main.tex
|-- References.bib
|
|-- figures/
|   |-- company.pdf
|   |-- system.png
|   `-- ...
|
`-- chapters/
    |-- chapter1.tex
    |-- chapter2.tex
    `-- chapter3.tex
\end{verbatim}

Figures should normally be stored in the \verb|figures| directory
specified by the official template.

File and directory names should preferably contain Latin letters,
numbers, hyphens, or underscores.

Recommended examples include:

\begin{verbatim}
company_structure.pdf
motor_control.png
chapter_01.tex
\end{verbatim}

Spaces and unnecessary special characters in file names should be
avoided where possible.

\vspace{0.3cm}

% ================================================================
% XXI. DOCUMENT COMPILATION AND FINAL CHECK
% ================================================================

\noindent
\textbf{XXI.\quad DOCUMENT COMPILATION AND FINAL CHECK}

\vspace{0.15cm}

Before submission, the complete \LaTeX{} project shall compile
successfully and produce the final PDF report.

The final document shall not contain unresolved references such as:

\begin{verbatim}
??
[?]
\end{verbatim}

Before generating the final PDF, students shall check:

\begin{itemize}[leftmargin=0.8cm, itemsep=2pt]
    \item the Table of Contents;
    \item page numbering;
    \item chapter, section, and subsection numbering;
    \item figure numbering;
    \item table numbering;
    \item equation numbering;
    \item cross-references;
    \item citations;
    \item bibliography;
    \item font display;
    \item mathematical symbols;
    \item image quality;
    \item completeness of all required forms.
\end{itemize}

The document may be compiled using TeX Live, MiKTeX, or another
compatible \LaTeX{} environment. Regardless of the compilation
environment, the resulting PDF shall conform to the official Faculty
template.

\vspace{0.3cm}

% ================================================================
% XXII. CONTENT THAT STUDENTS MAY EDIT
% ================================================================

\noindent
\textbf{XXII.\quad CONTENT THAT STUDENTS MAY EDIT}

\vspace{0.15cm}

Students may edit or replace:

\begin{itemize}[leftmargin=0.8cm, itemsep=2pt]
    \item student name;
    \item class;
    \item internship organization;
    \item supervisor name;
    \item Abstract;
    \item List of Symbols and Abbreviations;
    \item chapter titles and chapter content;
    \item figures and photographs;
    \item tables;
    \item equations;
    \item citations and references.
\end{itemize}

Students shall not modify the main document formatting settings in the
preamble, including the document class, page geometry, line spacing,
heading formatting, table-of-contents formatting, font configuration,
page layout, or cover-page macros unless specifically instructed.

\vspace{0.3cm}

% ================================================================
% XXIII. FINAL SUBMISSION REQUIREMENTS
% ================================================================

\noindent
\textbf{XXIII.\quad FINAL SUBMISSION REQUIREMENTS}

\vspace{0.15cm}

Before submission, the final report shall:

\begin{itemize}[leftmargin=0.8cm, itemsep=2pt]
    \item compile without errors that affect the report content or presentation;
    \item contain all required forms and report sections;
    \item contain correct student, class, internship organization, and supervisor information;
    \item contain an updated Table of Contents, List of Figures, and List of Tables;
    \item contain no unused example content from the template;
    \item contain no placeholder text such as \textit{Title of Chapter 1},
    \textit{Section 1.1}, or other sample entries;
    \item contain all required citations and references;
    \item contain no unresolved cross-references.
\end{itemize}

If the Faculty requires submission of the \LaTeX{} source, the complete
project shall include all \verb|.tex| files, the bibliography database,
figures, and any other files required to reproduce the submitted PDF.

\vspace{0.4cm}

\noindent
\textbf{Important:}
Before submitting the final internship report, students shall remove or
comment out the following command in the main document:

\begin{verbatim}
\input{Regulations.tex}
\end{verbatim}

For example:

\begin{verbatim}
% \input{Regulations.tex}
\end{verbatim}

This removes the guidance pages from the final submitted report while
preserving the regulations file in the \LaTeX{} project for reference.

} % End 13pt content

\clearpage

%
% IMPORTANT:
% This section is provided for guidance only. Students shall remove
% or comment out the 
% =================================================================
% REGULATIONS FOR PREPARING THE GRADUATION INTERNSHIP REPORT
% USING LATEX
% =================================================================
%
% This file is intended to be included from the main document with:
%
%     \input{Regulations.tex}
%
% IMPORTANT:
% This section is provided for guidance only. Students shall remove
% or comment out the \input{Regulations.tex} command before submitting
% the final internship report.
%
% =================================================================

\clearpage
\thispagestyle{empty}

\begin{center}
    {\fontsize{16pt}{20pt}\selectfont\bfseries
    SPECIFIC REGULATIONS FOR PREPARING\\
    GRADUATION INTERNSHIP REPORTS USING \LaTeX}
\end{center}

\vspace{0.5cm}

{\fontsize{13pt}{17pt}\selectfont
\setstretch{1.3}
\justifying

\noindent
\textbf{Note:}
This section is provided for guidance only and is included in the main
document using the command \verb|\input{Regulations.tex}|.
Students shall remove or comment out this command before submitting the
final internship report.

\vspace{0.4cm}

% ================================================================
% I. GENERAL REQUIREMENTS
% ================================================================

\noindent
\textbf{I.\quad GENERAL REQUIREMENTS}

\vspace{0.15cm}

The Graduation Internship Report shall be prepared using the official
\LaTeX{} template provided by the Faculty. Students shall preserve the
predefined page layout, typography, line spacing, heading styles,
numbering system, figure and table formatting, and bibliography settings
unless otherwise instructed by the Faculty or the supervisor.

The report shall be written in English.

Students are responsible for the academic and technical content of the
report. The official \LaTeX{} template is responsible for standardized
document formatting.

\vspace{0.3cm}

% ================================================================
% II. REPORT STRUCTURE
% ================================================================

\noindent
\textbf{II.\quad REPORT STRUCTURE}

\vspace{0.15cm}

The report shall normally be arranged in the following order:

\begin{enumerate}[label=\arabic*., leftmargin=1.2cm, itemsep=2pt]
    \item Outer Cover;
    \item Comments of the Internship Host Organization;
    \item Supervisor's Evaluation and Graduation Internship Defense Results;
    \item Table of Contents;
    \item List of Symbols and Abbreviations, if applicable;
    \item List of Figures;
    \item List of Tables;
    \item Abstract;
    \item Main Chapters of the Internship Report;
    \item References;
    \item Appendices, if applicable.
\end{enumerate}

The preliminary pages shall not be numbered as chapters. The main report
content shall be organized into consecutively numbered chapters beginning
with Chapter 1.

\vspace{0.3cm}

% ================================================================
% III. PAGE SIZE AND MARGINS
% ================================================================

\noindent
\textbf{III.\quad PAGE SIZE AND MARGINS}

\vspace{0.15cm}

The report shall be formatted on A4 paper with the following margins:

\begin{center}
\begin{tabular}{ll}
    Left margin:   & 3.5 cm \\
    Right margin:  & 2.0 cm \\
    Top margin:    & 2.5 cm \\
    Bottom margin: & 2.5 cm
\end{tabular}
\end{center}

These settings are predefined in the official template and shall not be
modified by students.

\vspace{0.3cm}

% ================================================================
% IV. FONT, FONT SIZE, AND LINE SPACING
% ================================================================

\noindent
\textbf{IV.\quad FONT, FONT SIZE, AND LINE SPACING}

\vspace{0.15cm}

The main text shall use a 13 pt font size with 1.3-line spacing.
A Times-family typeface shall be used in accordance with the official
template.

The text shall be fully justified.

Students shall not use repeated spaces, repeated manual line breaks, or
multiple blank lines to position text. Paragraph and heading spacing shall
be controlled by the template.

\vspace{0.3cm}

% ================================================================
% V. OUTER COVER
% ================================================================

\noindent
\textbf{V.\quad OUTER COVER}

\vspace{0.15cm}

The outer cover shall follow the official format provided in the template
and shall contain the following information:

\begin{itemize}[leftmargin=0.8cm, itemsep=2pt]
    \item Thai Nguyen University of Technology;
    \item Faculty of International Training;
    \item Engineering and Technology Department;
    \item University logo and Faculty logo;
    \item the title \textbf{INTERNSHIP REPORT};
    \item Internship Organization;
    \item Author;
    \item Class;
    \item Supervisor;
    \item place and year of completion.
\end{itemize}

Students shall only replace the required information fields and shall not
modify the cover-page layout unless instructed otherwise.

\vspace{0.3cm}

% ================================================================
% VI. COMMENTS OF THE INTERNSHIP HOST ORGANIZATION
% ================================================================

\noindent
\textbf{VI.\quad COMMENTS OF THE INTERNSHIP HOST ORGANIZATION}

\vspace{0.15cm}

The page entitled \textbf{COMMENTS OF THE INTERNSHIP HOST ORGANIZATION}
is reserved for the internship host organization.

It includes:

\begin{itemize}[leftmargin=0.8cm, itemsep=2pt]
    \item the student's full name;
    \item class;
    \item comments from the host organization;
    \item score;
    \item date;
    \item signature and official seal.
\end{itemize}

Students shall not alter or complete official evaluation fields unless
specifically authorized.

\vspace{0.3cm}

% ================================================================
% VII. SUPERVISOR'S EVALUATION AND DEFENSE RESULTS
% ================================================================

\noindent
\textbf{VII.\quad SUPERVISOR'S EVALUATION AND DEFENSE RESULTS}

\vspace{0.15cm}

The section entitled
\textbf{SUPERVISOR'S EVALUATION AND GRADUATION INTERNSHIP DEFENSE RESULTS}
is reserved for the supervisor and defense committee.

The evaluation may include:

\begin{itemize}[leftmargin=0.8cm, itemsep=2pt]
    \item Internship Progress and Student's Attitude;
    \item Report Content;
    \item Report Presentation;
    \item Other Comments;
    \item Supervisor's Opinion;
    \item Defense Result;
    \item Score.
\end{itemize}

Students shall not enter evaluation scores or defense results before they
are officially determined.

\vspace{0.3cm}

% ================================================================
% VIII. CHAPTERS, SECTIONS, AND SUBSECTIONS
% ================================================================

\noindent
\textbf{VIII.\quad CHAPTERS, SECTIONS, AND SUBSECTIONS}

\vspace{0.15cm}

The report structure shall be generated using the standard \LaTeX{}
commands:

\begin{verbatim}
\chapter{...}
\section{...}
\subsection{...}
\end{verbatim}

Chapter, section, and subsection numbers shall be generated automatically
by \LaTeX{}.

Students shall not manually type chapter or section numbers and shall not
reproduce heading styles using commands such as \verb|\textbf|,
\verb|\fontsize|, or \verb|\centering| instead of the appropriate
structural command.

Excessive subdivision should be avoided. In most cases, the hierarchy
Chapter--Section--Subsection is sufficient.

\vspace{0.3cm}

% ================================================================
% IX. ORGANIZATION OF THE MAIN CONTENT
% ================================================================

\noindent
\textbf{IX.\quad ORGANIZATION OF THE MAIN CONTENT}

\vspace{0.15cm}

The specific content of the report shall depend on the internship
organization and the assigned tasks. A typical structure may include:

\begin{itemize}[leftmargin=0.8cm, itemsep=2pt]
    \item \textbf{Chapter 1. Overview of the Internship Organization};
    \item \textbf{Chapter 2. Internship Activities and Assigned Tasks};
    \item \textbf{Chapter 3. Results, Experience, and Evaluation}.
\end{itemize}

These chapter titles are recommendations only and may be adapted to the
actual internship content.

The report shall reflect work actually performed, observed, analyzed, or
learned during the internship. Unrelated theoretical material should not
be included merely to increase report length.

\vspace{0.3cm}

% ================================================================
% X. TABLE OF CONTENTS
% ================================================================

\noindent
\textbf{X.\quad TABLE OF CONTENTS}

\vspace{0.15cm}

The Table of Contents shall be generated automatically using:

\begin{verbatim}
\tableofcontents
\end{verbatim}

Students shall not create or manually edit page numbers in the Table of
Contents.

The document shall be compiled sufficiently to ensure that all page
numbers and entries are correctly updated.

\vspace{0.3cm}

% ================================================================
% XI. LIST OF SYMBOLS AND ABBREVIATIONS
% ================================================================

\noindent
\textbf{XI.\quad LIST OF SYMBOLS AND ABBREVIATIONS}

\vspace{0.15cm}

Symbols and abbreviations that are used repeatedly in the report shall be
listed in the \textbf{List of Symbols and Abbreviations}.

The list may contain the following columns:

\begin{center}
\begin{tabular}{llll}
    No. & Symbol / Abbreviation & Unit & Definition
\end{tabular}
\end{center}

Only symbols and abbreviations actually used in the final report shall be
retained. Example entries contained in the template shall be removed if
they are not relevant to the student's report.

Each abbreviation should normally be defined when it first appears in
the main text.

\vspace{0.3cm}

% ================================================================
% XII. FIGURES, DIAGRAMS, AND PHOTOGRAPHS
% ================================================================

\noindent
\textbf{XII.\quad FIGURES, DIAGRAMS, AND PHOTOGRAPHS}

\vspace{0.15cm}

Figures shall be clear, legible, and directly relevant to the report
content.

Figures shall be inserted using the standard \verb|figure| environment.

Example:

\begin{verbatim}
\begin{figure}[H]
    \centering
    \includegraphics[width=0.8\textwidth]
                    {figures/company_structure.pdf}
    \caption{Organizational structure of the company}
    \label{fig:company-structure}
\end{figure}
\end{verbatim}

Figure captions shall be placed below the figure.

Each referenced figure shall have a unique \verb|\label| and shall be
referenced in the text using, for example:

\begin{verbatim}
Figure~\ref{fig:company-structure}
\end{verbatim}

Students shall not manually type figure numbers.

Figures reproduced or adapted from external sources shall include an
appropriate citation.

Photographs taken during the internship shall be relevant and sufficiently
clear. Students shall not include confidential, restricted, or proprietary
information without permission from the internship organization.

Vector graphics or high-resolution images are recommended whenever
possible.

\vspace{0.3cm}

% ================================================================
% XIII. TABLES
% ================================================================

\noindent
\textbf{XIII.\quad TABLES}

\vspace{0.15cm}

Tables shall be created using suitable \LaTeX{} environments such as
\verb|table|, \verb|tabular|, \verb|tabularx|, or \verb|longtable|.

Table captions shall be placed above the table.

Example:

\begin{verbatim}
\begin{table}[H]
    \centering
    \caption{Internship schedule}
    \label{tab:internship-schedule}

    \begin{tabular}{lll}
        \toprule
        Week & Task & Result \\
        \midrule
        1 & Orientation & Completed \\
        2 & System study & Completed \\
        \bottomrule
    \end{tabular}
\end{table}
\end{verbatim}

Tables shall be referenced using:

\begin{verbatim}
Table~\ref{tab:internship-schedule}
\end{verbatim}

Students shall not manually type table numbers.

Tables should be created directly in \LaTeX{} rather than inserted as
screenshots whenever practical.

\vspace{0.3cm}

% ================================================================
% XIV. EQUATIONS AND TECHNICAL QUANTITIES
% ================================================================

\noindent
\textbf{XIV.\quad EQUATIONS AND TECHNICAL QUANTITIES}

\vspace{0.15cm}

Mathematical expressions shall be typeset directly using \LaTeX{}.
Images or screenshots of equations shall not be used.

Short expressions may be written in inline mathematical mode, for example:

\begin{verbatim}
$R = 0.7\,\Omega$
\end{verbatim}

Important equations shall use a numbered equation environment.

Example:

\begin{verbatim}
\begin{equation}
    P = UI
    \label{eq:power}
\end{equation}
\end{verbatim}

Equations shall be referenced using:

\begin{verbatim}
Equation~\eqref{eq:power}
\end{verbatim}

Equation numbers shall not be manually entered.

SI units shall be used wherever applicable. The \verb|siunitx| package
should be used where appropriate.

Examples:

\begin{verbatim}
\SI{220}{\volt}
\SI{50}{\hertz}
\SI{2.5}{\ampere}
\end{verbatim}

Notation and units shall be consistent throughout the report.

\vspace{0.3cm}

% ================================================================
% XV. SOURCE CODE
% ================================================================

\noindent
\textbf{XV.\quad SOURCE CODE}

\vspace{0.15cm}

If the internship involves programming, PLCs, microcontrollers,
simulation, or software development, relevant source code may be included
in the report.

Source code shall be presented using the \verb|listings| environment or
another environment defined by the official template.

Example:

\begin{verbatim}
\begin{lstlisting}[language=st]
IF Start = TRUE THEN
    Motor := TRUE;
ELSE
    Motor := FALSE;
END_IF;
\end{lstlisting}
\end{verbatim}

Only code necessary for explaining the work should appear in the main
report. Long source-code listings should be moved to an appendix or
submitted as supplementary files if required.

Source code should not normally be inserted as screenshots.

\vspace{0.3cm}

% ================================================================
% XVI. ABSTRACT
% ================================================================

\noindent
\textbf{XVI.\quad ABSTRACT}

\vspace{0.15cm}

The Abstract shall briefly summarize:

\begin{itemize}[leftmargin=0.8cm, itemsep=2pt]
    \item the internship organization;
    \item the objectives of the internship;
    \item the main tasks performed;
    \item the main results, knowledge, and skills obtained.
\end{itemize}

The Abstract should be concise and should not reproduce the full content
of the report chapter by chapter.

\vspace{0.3cm}

% ================================================================
% XVII. REFERENCES AND CITATIONS
% ================================================================

\noindent
\textbf{XVII.\quad REFERENCES AND CITATIONS}

\vspace{0.15cm}

All books, journal papers, conference papers, standards, websites,
datasheets, manuals, technical documents, figures, tables, or other
external sources used in preparing the report shall be properly cited.

References shall be managed through the bibliography system provided with
the official template.

Bibliographic entries shall normally be stored in:

\begin{verbatim}
References.bib
\end{verbatim}

References shall be cited in the report using:

\begin{verbatim}
\cite{reference-key}
\end{verbatim}

Students shall not manually type reference numbers such as ``[1]'',
``[2]'', etc.

Titles of foreign-language publications shall remain in their original
language in the bibliography.

All borrowed ideas, data, figures, tables, equations, methods, or results
shall be properly acknowledged.

\vspace{0.3cm}

% ================================================================
% XVIII. CROSS-REFERENCING
% ================================================================

\noindent
\textbf{XVIII.\quad CROSS-REFERENCING}

\vspace{0.15cm}

Numbered elements shall be cross-referenced using the standard \LaTeX{}
mechanism.

Use the following commands where appropriate:

\begin{verbatim}
\label{...}
\ref{...}
\eqref{...}
\end{verbatim}

For example:

\begin{verbatim}
As shown in Figure~\ref{fig:system}, ...
\end{verbatim}

Students shall not manually type numbers such as ``Figure 2.3'' when
those numbers can be generated automatically.

\vspace{0.3cm}

% ================================================================
% XIX. PAGE NUMBERING
% ================================================================

\noindent
\textbf{XIX.\quad PAGE NUMBERING}

\vspace{0.15cm}

Cover pages and official evaluation forms at the beginning of the report
shall not display page numbers.

The main body of the report shall use Arabic page numbers beginning from
Chapter 1.

Page numbering shall be generated automatically by the official template.
Students shall not manually enter page numbers.

\vspace{0.3cm}

% ================================================================
% XX. PROJECT FILE ORGANIZATION
% ================================================================

\noindent
\textbf{XX.\quad PROJECT FILE ORGANIZATION}

\vspace{0.15cm}

Students should maintain a clear and organized \LaTeX{} project
structure. A recommended structure is:

\begin{verbatim}
InternshipReport/
|
|-- main.tex
|-- References.bib
|
|-- figures/
|   |-- company.pdf
|   |-- system.png
|   `-- ...
|
`-- chapters/
    |-- chapter1.tex
    |-- chapter2.tex
    `-- chapter3.tex
\end{verbatim}

Figures should normally be stored in the \verb|figures| directory
specified by the official template.

File and directory names should preferably contain Latin letters,
numbers, hyphens, or underscores.

Recommended examples include:

\begin{verbatim}
company_structure.pdf
motor_control.png
chapter_01.tex
\end{verbatim}

Spaces and unnecessary special characters in file names should be
avoided where possible.

\vspace{0.3cm}

% ================================================================
% XXI. DOCUMENT COMPILATION AND FINAL CHECK
% ================================================================

\noindent
\textbf{XXI.\quad DOCUMENT COMPILATION AND FINAL CHECK}

\vspace{0.15cm}

Before submission, the complete \LaTeX{} project shall compile
successfully and produce the final PDF report.

The final document shall not contain unresolved references such as:

\begin{verbatim}
??
[?]
\end{verbatim}

Before generating the final PDF, students shall check:

\begin{itemize}[leftmargin=0.8cm, itemsep=2pt]
    \item the Table of Contents;
    \item page numbering;
    \item chapter, section, and subsection numbering;
    \item figure numbering;
    \item table numbering;
    \item equation numbering;
    \item cross-references;
    \item citations;
    \item bibliography;
    \item font display;
    \item mathematical symbols;
    \item image quality;
    \item completeness of all required forms.
\end{itemize}

The document may be compiled using TeX Live, MiKTeX, or another
compatible \LaTeX{} environment. Regardless of the compilation
environment, the resulting PDF shall conform to the official Faculty
template.

\vspace{0.3cm}

% ================================================================
% XXII. CONTENT THAT STUDENTS MAY EDIT
% ================================================================

\noindent
\textbf{XXII.\quad CONTENT THAT STUDENTS MAY EDIT}

\vspace{0.15cm}

Students may edit or replace:

\begin{itemize}[leftmargin=0.8cm, itemsep=2pt]
    \item student name;
    \item class;
    \item internship organization;
    \item supervisor name;
    \item Abstract;
    \item List of Symbols and Abbreviations;
    \item chapter titles and chapter content;
    \item figures and photographs;
    \item tables;
    \item equations;
    \item citations and references.
\end{itemize}

Students shall not modify the main document formatting settings in the
preamble, including the document class, page geometry, line spacing,
heading formatting, table-of-contents formatting, font configuration,
page layout, or cover-page macros unless specifically instructed.

\vspace{0.3cm}

% ================================================================
% XXIII. FINAL SUBMISSION REQUIREMENTS
% ================================================================

\noindent
\textbf{XXIII.\quad FINAL SUBMISSION REQUIREMENTS}

\vspace{0.15cm}

Before submission, the final report shall:

\begin{itemize}[leftmargin=0.8cm, itemsep=2pt]
    \item compile without errors that affect the report content or presentation;
    \item contain all required forms and report sections;
    \item contain correct student, class, internship organization, and supervisor information;
    \item contain an updated Table of Contents, List of Figures, and List of Tables;
    \item contain no unused example content from the template;
    \item contain no placeholder text such as \textit{Title of Chapter 1},
    \textit{Section 1.1}, or other sample entries;
    \item contain all required citations and references;
    \item contain no unresolved cross-references.
\end{itemize}

If the Faculty requires submission of the \LaTeX{} source, the complete
project shall include all \verb|.tex| files, the bibliography database,
figures, and any other files required to reproduce the submitted PDF.

\vspace{0.4cm}

\noindent
\textbf{Important:}
Before submitting the final internship report, students shall remove or
comment out the following command in the main document:

\begin{verbatim}
\input{Regulations.tex}
\end{verbatim}

For example:

\begin{verbatim}
% \input{Regulations.tex}
\end{verbatim}

This removes the guidance pages from the final submitted report while
preserving the regulations file in the \LaTeX{} project for reference.

} % End 13pt content

\clearpage
 command before submitting
% the final internship report.
%
% =================================================================

\clearpage
\thispagestyle{empty}

\begin{center}
    {\fontsize{16pt}{20pt}\selectfont\bfseries
    SPECIFIC REGULATIONS FOR PREPARING\\
    GRADUATION INTERNSHIP REPORTS USING \LaTeX}
\end{center}

\vspace{0.5cm}

{\fontsize{13pt}{17pt}\selectfont
\setstretch{1.3}
\justifying

\noindent
\textbf{Note:}
This section is provided for guidance only and is included in the main
document using the command \verb|
% =================================================================
% REGULATIONS FOR PREPARING THE GRADUATION INTERNSHIP REPORT
% USING LATEX
% =================================================================
%
% This file is intended to be included from the main document with:
%
%     \input{Regulations.tex}
%
% IMPORTANT:
% This section is provided for guidance only. Students shall remove
% or comment out the \input{Regulations.tex} command before submitting
% the final internship report.
%
% =================================================================

\clearpage
\thispagestyle{empty}

\begin{center}
    {\fontsize{16pt}{20pt}\selectfont\bfseries
    SPECIFIC REGULATIONS FOR PREPARING\\
    GRADUATION INTERNSHIP REPORTS USING \LaTeX}
\end{center}

\vspace{0.5cm}

{\fontsize{13pt}{17pt}\selectfont
\setstretch{1.3}
\justifying

\noindent
\textbf{Note:}
This section is provided for guidance only and is included in the main
document using the command \verb|\input{Regulations.tex}|.
Students shall remove or comment out this command before submitting the
final internship report.

\vspace{0.4cm}

% ================================================================
% I. GENERAL REQUIREMENTS
% ================================================================

\noindent
\textbf{I.\quad GENERAL REQUIREMENTS}

\vspace{0.15cm}

The Graduation Internship Report shall be prepared using the official
\LaTeX{} template provided by the Faculty. Students shall preserve the
predefined page layout, typography, line spacing, heading styles,
numbering system, figure and table formatting, and bibliography settings
unless otherwise instructed by the Faculty or the supervisor.

The report shall be written in English.

Students are responsible for the academic and technical content of the
report. The official \LaTeX{} template is responsible for standardized
document formatting.

\vspace{0.3cm}

% ================================================================
% II. REPORT STRUCTURE
% ================================================================

\noindent
\textbf{II.\quad REPORT STRUCTURE}

\vspace{0.15cm}

The report shall normally be arranged in the following order:

\begin{enumerate}[label=\arabic*., leftmargin=1.2cm, itemsep=2pt]
    \item Outer Cover;
    \item Comments of the Internship Host Organization;
    \item Supervisor's Evaluation and Graduation Internship Defense Results;
    \item Table of Contents;
    \item List of Symbols and Abbreviations, if applicable;
    \item List of Figures;
    \item List of Tables;
    \item Abstract;
    \item Main Chapters of the Internship Report;
    \item References;
    \item Appendices, if applicable.
\end{enumerate}

The preliminary pages shall not be numbered as chapters. The main report
content shall be organized into consecutively numbered chapters beginning
with Chapter 1.

\vspace{0.3cm}

% ================================================================
% III. PAGE SIZE AND MARGINS
% ================================================================

\noindent
\textbf{III.\quad PAGE SIZE AND MARGINS}

\vspace{0.15cm}

The report shall be formatted on A4 paper with the following margins:

\begin{center}
\begin{tabular}{ll}
    Left margin:   & 3.5 cm \\
    Right margin:  & 2.0 cm \\
    Top margin:    & 2.5 cm \\
    Bottom margin: & 2.5 cm
\end{tabular}
\end{center}

These settings are predefined in the official template and shall not be
modified by students.

\vspace{0.3cm}

% ================================================================
% IV. FONT, FONT SIZE, AND LINE SPACING
% ================================================================

\noindent
\textbf{IV.\quad FONT, FONT SIZE, AND LINE SPACING}

\vspace{0.15cm}

The main text shall use a 13 pt font size with 1.3-line spacing.
A Times-family typeface shall be used in accordance with the official
template.

The text shall be fully justified.

Students shall not use repeated spaces, repeated manual line breaks, or
multiple blank lines to position text. Paragraph and heading spacing shall
be controlled by the template.

\vspace{0.3cm}

% ================================================================
% V. OUTER COVER
% ================================================================

\noindent
\textbf{V.\quad OUTER COVER}

\vspace{0.15cm}

The outer cover shall follow the official format provided in the template
and shall contain the following information:

\begin{itemize}[leftmargin=0.8cm, itemsep=2pt]
    \item Thai Nguyen University of Technology;
    \item Faculty of International Training;
    \item Engineering and Technology Department;
    \item University logo and Faculty logo;
    \item the title \textbf{INTERNSHIP REPORT};
    \item Internship Organization;
    \item Author;
    \item Class;
    \item Supervisor;
    \item place and year of completion.
\end{itemize}

Students shall only replace the required information fields and shall not
modify the cover-page layout unless instructed otherwise.

\vspace{0.3cm}

% ================================================================
% VI. COMMENTS OF THE INTERNSHIP HOST ORGANIZATION
% ================================================================

\noindent
\textbf{VI.\quad COMMENTS OF THE INTERNSHIP HOST ORGANIZATION}

\vspace{0.15cm}

The page entitled \textbf{COMMENTS OF THE INTERNSHIP HOST ORGANIZATION}
is reserved for the internship host organization.

It includes:

\begin{itemize}[leftmargin=0.8cm, itemsep=2pt]
    \item the student's full name;
    \item class;
    \item comments from the host organization;
    \item score;
    \item date;
    \item signature and official seal.
\end{itemize}

Students shall not alter or complete official evaluation fields unless
specifically authorized.

\vspace{0.3cm}

% ================================================================
% VII. SUPERVISOR'S EVALUATION AND DEFENSE RESULTS
% ================================================================

\noindent
\textbf{VII.\quad SUPERVISOR'S EVALUATION AND DEFENSE RESULTS}

\vspace{0.15cm}

The section entitled
\textbf{SUPERVISOR'S EVALUATION AND GRADUATION INTERNSHIP DEFENSE RESULTS}
is reserved for the supervisor and defense committee.

The evaluation may include:

\begin{itemize}[leftmargin=0.8cm, itemsep=2pt]
    \item Internship Progress and Student's Attitude;
    \item Report Content;
    \item Report Presentation;
    \item Other Comments;
    \item Supervisor's Opinion;
    \item Defense Result;
    \item Score.
\end{itemize}

Students shall not enter evaluation scores or defense results before they
are officially determined.

\vspace{0.3cm}

% ================================================================
% VIII. CHAPTERS, SECTIONS, AND SUBSECTIONS
% ================================================================

\noindent
\textbf{VIII.\quad CHAPTERS, SECTIONS, AND SUBSECTIONS}

\vspace{0.15cm}

The report structure shall be generated using the standard \LaTeX{}
commands:

\begin{verbatim}
\chapter{...}
\section{...}
\subsection{...}
\end{verbatim}

Chapter, section, and subsection numbers shall be generated automatically
by \LaTeX{}.

Students shall not manually type chapter or section numbers and shall not
reproduce heading styles using commands such as \verb|\textbf|,
\verb|\fontsize|, or \verb|\centering| instead of the appropriate
structural command.

Excessive subdivision should be avoided. In most cases, the hierarchy
Chapter--Section--Subsection is sufficient.

\vspace{0.3cm}

% ================================================================
% IX. ORGANIZATION OF THE MAIN CONTENT
% ================================================================

\noindent
\textbf{IX.\quad ORGANIZATION OF THE MAIN CONTENT}

\vspace{0.15cm}

The specific content of the report shall depend on the internship
organization and the assigned tasks. A typical structure may include:

\begin{itemize}[leftmargin=0.8cm, itemsep=2pt]
    \item \textbf{Chapter 1. Overview of the Internship Organization};
    \item \textbf{Chapter 2. Internship Activities and Assigned Tasks};
    \item \textbf{Chapter 3. Results, Experience, and Evaluation}.
\end{itemize}

These chapter titles are recommendations only and may be adapted to the
actual internship content.

The report shall reflect work actually performed, observed, analyzed, or
learned during the internship. Unrelated theoretical material should not
be included merely to increase report length.

\vspace{0.3cm}

% ================================================================
% X. TABLE OF CONTENTS
% ================================================================

\noindent
\textbf{X.\quad TABLE OF CONTENTS}

\vspace{0.15cm}

The Table of Contents shall be generated automatically using:

\begin{verbatim}
\tableofcontents
\end{verbatim}

Students shall not create or manually edit page numbers in the Table of
Contents.

The document shall be compiled sufficiently to ensure that all page
numbers and entries are correctly updated.

\vspace{0.3cm}

% ================================================================
% XI. LIST OF SYMBOLS AND ABBREVIATIONS
% ================================================================

\noindent
\textbf{XI.\quad LIST OF SYMBOLS AND ABBREVIATIONS}

\vspace{0.15cm}

Symbols and abbreviations that are used repeatedly in the report shall be
listed in the \textbf{List of Symbols and Abbreviations}.

The list may contain the following columns:

\begin{center}
\begin{tabular}{llll}
    No. & Symbol / Abbreviation & Unit & Definition
\end{tabular}
\end{center}

Only symbols and abbreviations actually used in the final report shall be
retained. Example entries contained in the template shall be removed if
they are not relevant to the student's report.

Each abbreviation should normally be defined when it first appears in
the main text.

\vspace{0.3cm}

% ================================================================
% XII. FIGURES, DIAGRAMS, AND PHOTOGRAPHS
% ================================================================

\noindent
\textbf{XII.\quad FIGURES, DIAGRAMS, AND PHOTOGRAPHS}

\vspace{0.15cm}

Figures shall be clear, legible, and directly relevant to the report
content.

Figures shall be inserted using the standard \verb|figure| environment.

Example:

\begin{verbatim}
\begin{figure}[H]
    \centering
    \includegraphics[width=0.8\textwidth]
                    {figures/company_structure.pdf}
    \caption{Organizational structure of the company}
    \label{fig:company-structure}
\end{figure}
\end{verbatim}

Figure captions shall be placed below the figure.

Each referenced figure shall have a unique \verb|\label| and shall be
referenced in the text using, for example:

\begin{verbatim}
Figure~\ref{fig:company-structure}
\end{verbatim}

Students shall not manually type figure numbers.

Figures reproduced or adapted from external sources shall include an
appropriate citation.

Photographs taken during the internship shall be relevant and sufficiently
clear. Students shall not include confidential, restricted, or proprietary
information without permission from the internship organization.

Vector graphics or high-resolution images are recommended whenever
possible.

\vspace{0.3cm}

% ================================================================
% XIII. TABLES
% ================================================================

\noindent
\textbf{XIII.\quad TABLES}

\vspace{0.15cm}

Tables shall be created using suitable \LaTeX{} environments such as
\verb|table|, \verb|tabular|, \verb|tabularx|, or \verb|longtable|.

Table captions shall be placed above the table.

Example:

\begin{verbatim}
\begin{table}[H]
    \centering
    \caption{Internship schedule}
    \label{tab:internship-schedule}

    \begin{tabular}{lll}
        \toprule
        Week & Task & Result \\
        \midrule
        1 & Orientation & Completed \\
        2 & System study & Completed \\
        \bottomrule
    \end{tabular}
\end{table}
\end{verbatim}

Tables shall be referenced using:

\begin{verbatim}
Table~\ref{tab:internship-schedule}
\end{verbatim}

Students shall not manually type table numbers.

Tables should be created directly in \LaTeX{} rather than inserted as
screenshots whenever practical.

\vspace{0.3cm}

% ================================================================
% XIV. EQUATIONS AND TECHNICAL QUANTITIES
% ================================================================

\noindent
\textbf{XIV.\quad EQUATIONS AND TECHNICAL QUANTITIES}

\vspace{0.15cm}

Mathematical expressions shall be typeset directly using \LaTeX{}.
Images or screenshots of equations shall not be used.

Short expressions may be written in inline mathematical mode, for example:

\begin{verbatim}
$R = 0.7\,\Omega$
\end{verbatim}

Important equations shall use a numbered equation environment.

Example:

\begin{verbatim}
\begin{equation}
    P = UI
    \label{eq:power}
\end{equation}
\end{verbatim}

Equations shall be referenced using:

\begin{verbatim}
Equation~\eqref{eq:power}
\end{verbatim}

Equation numbers shall not be manually entered.

SI units shall be used wherever applicable. The \verb|siunitx| package
should be used where appropriate.

Examples:

\begin{verbatim}
\SI{220}{\volt}
\SI{50}{\hertz}
\SI{2.5}{\ampere}
\end{verbatim}

Notation and units shall be consistent throughout the report.

\vspace{0.3cm}

% ================================================================
% XV. SOURCE CODE
% ================================================================

\noindent
\textbf{XV.\quad SOURCE CODE}

\vspace{0.15cm}

If the internship involves programming, PLCs, microcontrollers,
simulation, or software development, relevant source code may be included
in the report.

Source code shall be presented using the \verb|listings| environment or
another environment defined by the official template.

Example:

\begin{verbatim}
\begin{lstlisting}[language=st]
IF Start = TRUE THEN
    Motor := TRUE;
ELSE
    Motor := FALSE;
END_IF;
\end{lstlisting}
\end{verbatim}

Only code necessary for explaining the work should appear in the main
report. Long source-code listings should be moved to an appendix or
submitted as supplementary files if required.

Source code should not normally be inserted as screenshots.

\vspace{0.3cm}

% ================================================================
% XVI. ABSTRACT
% ================================================================

\noindent
\textbf{XVI.\quad ABSTRACT}

\vspace{0.15cm}

The Abstract shall briefly summarize:

\begin{itemize}[leftmargin=0.8cm, itemsep=2pt]
    \item the internship organization;
    \item the objectives of the internship;
    \item the main tasks performed;
    \item the main results, knowledge, and skills obtained.
\end{itemize}

The Abstract should be concise and should not reproduce the full content
of the report chapter by chapter.

\vspace{0.3cm}

% ================================================================
% XVII. REFERENCES AND CITATIONS
% ================================================================

\noindent
\textbf{XVII.\quad REFERENCES AND CITATIONS}

\vspace{0.15cm}

All books, journal papers, conference papers, standards, websites,
datasheets, manuals, technical documents, figures, tables, or other
external sources used in preparing the report shall be properly cited.

References shall be managed through the bibliography system provided with
the official template.

Bibliographic entries shall normally be stored in:

\begin{verbatim}
References.bib
\end{verbatim}

References shall be cited in the report using:

\begin{verbatim}
\cite{reference-key}
\end{verbatim}

Students shall not manually type reference numbers such as ``[1]'',
``[2]'', etc.

Titles of foreign-language publications shall remain in their original
language in the bibliography.

All borrowed ideas, data, figures, tables, equations, methods, or results
shall be properly acknowledged.

\vspace{0.3cm}

% ================================================================
% XVIII. CROSS-REFERENCING
% ================================================================

\noindent
\textbf{XVIII.\quad CROSS-REFERENCING}

\vspace{0.15cm}

Numbered elements shall be cross-referenced using the standard \LaTeX{}
mechanism.

Use the following commands where appropriate:

\begin{verbatim}
\label{...}
\ref{...}
\eqref{...}
\end{verbatim}

For example:

\begin{verbatim}
As shown in Figure~\ref{fig:system}, ...
\end{verbatim}

Students shall not manually type numbers such as ``Figure 2.3'' when
those numbers can be generated automatically.

\vspace{0.3cm}

% ================================================================
% XIX. PAGE NUMBERING
% ================================================================

\noindent
\textbf{XIX.\quad PAGE NUMBERING}

\vspace{0.15cm}

Cover pages and official evaluation forms at the beginning of the report
shall not display page numbers.

The main body of the report shall use Arabic page numbers beginning from
Chapter 1.

Page numbering shall be generated automatically by the official template.
Students shall not manually enter page numbers.

\vspace{0.3cm}

% ================================================================
% XX. PROJECT FILE ORGANIZATION
% ================================================================

\noindent
\textbf{XX.\quad PROJECT FILE ORGANIZATION}

\vspace{0.15cm}

Students should maintain a clear and organized \LaTeX{} project
structure. A recommended structure is:

\begin{verbatim}
InternshipReport/
|
|-- main.tex
|-- References.bib
|
|-- figures/
|   |-- company.pdf
|   |-- system.png
|   `-- ...
|
`-- chapters/
    |-- chapter1.tex
    |-- chapter2.tex
    `-- chapter3.tex
\end{verbatim}

Figures should normally be stored in the \verb|figures| directory
specified by the official template.

File and directory names should preferably contain Latin letters,
numbers, hyphens, or underscores.

Recommended examples include:

\begin{verbatim}
company_structure.pdf
motor_control.png
chapter_01.tex
\end{verbatim}

Spaces and unnecessary special characters in file names should be
avoided where possible.

\vspace{0.3cm}

% ================================================================
% XXI. DOCUMENT COMPILATION AND FINAL CHECK
% ================================================================

\noindent
\textbf{XXI.\quad DOCUMENT COMPILATION AND FINAL CHECK}

\vspace{0.15cm}

Before submission, the complete \LaTeX{} project shall compile
successfully and produce the final PDF report.

The final document shall not contain unresolved references such as:

\begin{verbatim}
??
[?]
\end{verbatim}

Before generating the final PDF, students shall check:

\begin{itemize}[leftmargin=0.8cm, itemsep=2pt]
    \item the Table of Contents;
    \item page numbering;
    \item chapter, section, and subsection numbering;
    \item figure numbering;
    \item table numbering;
    \item equation numbering;
    \item cross-references;
    \item citations;
    \item bibliography;
    \item font display;
    \item mathematical symbols;
    \item image quality;
    \item completeness of all required forms.
\end{itemize}

The document may be compiled using TeX Live, MiKTeX, or another
compatible \LaTeX{} environment. Regardless of the compilation
environment, the resulting PDF shall conform to the official Faculty
template.

\vspace{0.3cm}

% ================================================================
% XXII. CONTENT THAT STUDENTS MAY EDIT
% ================================================================

\noindent
\textbf{XXII.\quad CONTENT THAT STUDENTS MAY EDIT}

\vspace{0.15cm}

Students may edit or replace:

\begin{itemize}[leftmargin=0.8cm, itemsep=2pt]
    \item student name;
    \item class;
    \item internship organization;
    \item supervisor name;
    \item Abstract;
    \item List of Symbols and Abbreviations;
    \item chapter titles and chapter content;
    \item figures and photographs;
    \item tables;
    \item equations;
    \item citations and references.
\end{itemize}

Students shall not modify the main document formatting settings in the
preamble, including the document class, page geometry, line spacing,
heading formatting, table-of-contents formatting, font configuration,
page layout, or cover-page macros unless specifically instructed.

\vspace{0.3cm}

% ================================================================
% XXIII. FINAL SUBMISSION REQUIREMENTS
% ================================================================

\noindent
\textbf{XXIII.\quad FINAL SUBMISSION REQUIREMENTS}

\vspace{0.15cm}

Before submission, the final report shall:

\begin{itemize}[leftmargin=0.8cm, itemsep=2pt]
    \item compile without errors that affect the report content or presentation;
    \item contain all required forms and report sections;
    \item contain correct student, class, internship organization, and supervisor information;
    \item contain an updated Table of Contents, List of Figures, and List of Tables;
    \item contain no unused example content from the template;
    \item contain no placeholder text such as \textit{Title of Chapter 1},
    \textit{Section 1.1}, or other sample entries;
    \item contain all required citations and references;
    \item contain no unresolved cross-references.
\end{itemize}

If the Faculty requires submission of the \LaTeX{} source, the complete
project shall include all \verb|.tex| files, the bibliography database,
figures, and any other files required to reproduce the submitted PDF.

\vspace{0.4cm}

\noindent
\textbf{Important:}
Before submitting the final internship report, students shall remove or
comment out the following command in the main document:

\begin{verbatim}
\input{Regulations.tex}
\end{verbatim}

For example:

\begin{verbatim}
% \input{Regulations.tex}
\end{verbatim}

This removes the guidance pages from the final submitted report while
preserving the regulations file in the \LaTeX{} project for reference.

} % End 13pt content

\clearpage
|.
Students shall remove or comment out this command before submitting the
final internship report.

\vspace{0.4cm}

% ================================================================
% I. GENERAL REQUIREMENTS
% ================================================================

\noindent
\textbf{I.\quad GENERAL REQUIREMENTS}

\vspace{0.15cm}

The Graduation Internship Report shall be prepared using the official
\LaTeX{} template provided by the Faculty. Students shall preserve the
predefined page layout, typography, line spacing, heading styles,
numbering system, figure and table formatting, and bibliography settings
unless otherwise instructed by the Faculty or the supervisor.

The report shall be written in English.

Students are responsible for the academic and technical content of the
report. The official \LaTeX{} template is responsible for standardized
document formatting.

\vspace{0.3cm}

% ================================================================
% II. REPORT STRUCTURE
% ================================================================

\noindent
\textbf{II.\quad REPORT STRUCTURE}

\vspace{0.15cm}

The report shall normally be arranged in the following order:

\begin{enumerate}[label=\arabic*., leftmargin=1.2cm, itemsep=2pt]
    \item Outer Cover;
    \item Comments of the Internship Host Organization;
    \item Supervisor's Evaluation and Graduation Internship Defense Results;
    \item Table of Contents;
    \item List of Symbols and Abbreviations, if applicable;
    \item List of Figures;
    \item List of Tables;
    \item Abstract;
    \item Main Chapters of the Internship Report;
    \item References;
    \item Appendices, if applicable.
\end{enumerate}

The preliminary pages shall not be numbered as chapters. The main report
content shall be organized into consecutively numbered chapters beginning
with Chapter 1.

\vspace{0.3cm}

% ================================================================
% III. PAGE SIZE AND MARGINS
% ================================================================

\noindent
\textbf{III.\quad PAGE SIZE AND MARGINS}

\vspace{0.15cm}

The report shall be formatted on A4 paper with the following margins:

\begin{center}
\begin{tabular}{ll}
    Left margin:   & 3.5 cm \\
    Right margin:  & 2.0 cm \\
    Top margin:    & 2.5 cm \\
    Bottom margin: & 2.5 cm
\end{tabular}
\end{center}

These settings are predefined in the official template and shall not be
modified by students.

\vspace{0.3cm}

% ================================================================
% IV. FONT, FONT SIZE, AND LINE SPACING
% ================================================================

\noindent
\textbf{IV.\quad FONT, FONT SIZE, AND LINE SPACING}

\vspace{0.15cm}

The main text shall use a 13 pt font size with 1.3-line spacing.
A Times-family typeface shall be used in accordance with the official
template.

The text shall be fully justified.

Students shall not use repeated spaces, repeated manual line breaks, or
multiple blank lines to position text. Paragraph and heading spacing shall
be controlled by the template.

\vspace{0.3cm}

% ================================================================
% V. OUTER COVER
% ================================================================

\noindent
\textbf{V.\quad OUTER COVER}

\vspace{0.15cm}

The outer cover shall follow the official format provided in the template
and shall contain the following information:

\begin{itemize}[leftmargin=0.8cm, itemsep=2pt]
    \item Thai Nguyen University of Technology;
    \item Faculty of International Training;
    \item Engineering and Technology Department;
    \item University logo and Faculty logo;
    \item the title \textbf{INTERNSHIP REPORT};
    \item Internship Organization;
    \item Author;
    \item Class;
    \item Supervisor;
    \item place and year of completion.
\end{itemize}

Students shall only replace the required information fields and shall not
modify the cover-page layout unless instructed otherwise.

\vspace{0.3cm}

% ================================================================
% VI. COMMENTS OF THE INTERNSHIP HOST ORGANIZATION
% ================================================================

\noindent
\textbf{VI.\quad COMMENTS OF THE INTERNSHIP HOST ORGANIZATION}

\vspace{0.15cm}

The page entitled \textbf{COMMENTS OF THE INTERNSHIP HOST ORGANIZATION}
is reserved for the internship host organization.

It includes:

\begin{itemize}[leftmargin=0.8cm, itemsep=2pt]
    \item the student's full name;
    \item class;
    \item comments from the host organization;
    \item score;
    \item date;
    \item signature and official seal.
\end{itemize}

Students shall not alter or complete official evaluation fields unless
specifically authorized.

\vspace{0.3cm}

% ================================================================
% VII. SUPERVISOR'S EVALUATION AND DEFENSE RESULTS
% ================================================================

\noindent
\textbf{VII.\quad SUPERVISOR'S EVALUATION AND DEFENSE RESULTS}

\vspace{0.15cm}

The section entitled
\textbf{SUPERVISOR'S EVALUATION AND GRADUATION INTERNSHIP DEFENSE RESULTS}
is reserved for the supervisor and defense committee.

The evaluation may include:

\begin{itemize}[leftmargin=0.8cm, itemsep=2pt]
    \item Internship Progress and Student's Attitude;
    \item Report Content;
    \item Report Presentation;
    \item Other Comments;
    \item Supervisor's Opinion;
    \item Defense Result;
    \item Score.
\end{itemize}

Students shall not enter evaluation scores or defense results before they
are officially determined.

\vspace{0.3cm}

% ================================================================
% VIII. CHAPTERS, SECTIONS, AND SUBSECTIONS
% ================================================================

\noindent
\textbf{VIII.\quad CHAPTERS, SECTIONS, AND SUBSECTIONS}

\vspace{0.15cm}

The report structure shall be generated using the standard \LaTeX{}
commands:

\begin{verbatim}
\chapter{...}
\section{...}
\subsection{...}
\end{verbatim}

Chapter, section, and subsection numbers shall be generated automatically
by \LaTeX{}.

Students shall not manually type chapter or section numbers and shall not
reproduce heading styles using commands such as \verb|\textbf|,
\verb|\fontsize|, or \verb|\centering| instead of the appropriate
structural command.

Excessive subdivision should be avoided. In most cases, the hierarchy
Chapter--Section--Subsection is sufficient.

\vspace{0.3cm}

% ================================================================
% IX. ORGANIZATION OF THE MAIN CONTENT
% ================================================================

\noindent
\textbf{IX.\quad ORGANIZATION OF THE MAIN CONTENT}

\vspace{0.15cm}

The specific content of the report shall depend on the internship
organization and the assigned tasks. A typical structure may include:

\begin{itemize}[leftmargin=0.8cm, itemsep=2pt]
    \item \textbf{Chapter 1. Overview of the Internship Organization};
    \item \textbf{Chapter 2. Internship Activities and Assigned Tasks};
    \item \textbf{Chapter 3. Results, Experience, and Evaluation}.
\end{itemize}

These chapter titles are recommendations only and may be adapted to the
actual internship content.

The report shall reflect work actually performed, observed, analyzed, or
learned during the internship. Unrelated theoretical material should not
be included merely to increase report length.

\vspace{0.3cm}

% ================================================================
% X. TABLE OF CONTENTS
% ================================================================

\noindent
\textbf{X.\quad TABLE OF CONTENTS}

\vspace{0.15cm}

The Table of Contents shall be generated automatically using:

\begin{verbatim}
\tableofcontents
\end{verbatim}

Students shall not create or manually edit page numbers in the Table of
Contents.

The document shall be compiled sufficiently to ensure that all page
numbers and entries are correctly updated.

\vspace{0.3cm}

% ================================================================
% XI. LIST OF SYMBOLS AND ABBREVIATIONS
% ================================================================

\noindent
\textbf{XI.\quad LIST OF SYMBOLS AND ABBREVIATIONS}

\vspace{0.15cm}

Symbols and abbreviations that are used repeatedly in the report shall be
listed in the \textbf{List of Symbols and Abbreviations}.

The list may contain the following columns:

\begin{center}
\begin{tabular}{llll}
    No. & Symbol / Abbreviation & Unit & Definition
\end{tabular}
\end{center}

Only symbols and abbreviations actually used in the final report shall be
retained. Example entries contained in the template shall be removed if
they are not relevant to the student's report.

Each abbreviation should normally be defined when it first appears in
the main text.

\vspace{0.3cm}

% ================================================================
% XII. FIGURES, DIAGRAMS, AND PHOTOGRAPHS
% ================================================================

\noindent
\textbf{XII.\quad FIGURES, DIAGRAMS, AND PHOTOGRAPHS}

\vspace{0.15cm}

Figures shall be clear, legible, and directly relevant to the report
content.

Figures shall be inserted using the standard \verb|figure| environment.

Example:

\begin{verbatim}
\begin{figure}[H]
    \centering
    \includegraphics[width=0.8\textwidth]
                    {figures/company_structure.pdf}
    \caption{Organizational structure of the company}
    \label{fig:company-structure}
\end{figure}
\end{verbatim}

Figure captions shall be placed below the figure.

Each referenced figure shall have a unique \verb|\label| and shall be
referenced in the text using, for example:

\begin{verbatim}
Figure~\ref{fig:company-structure}
\end{verbatim}

Students shall not manually type figure numbers.

Figures reproduced or adapted from external sources shall include an
appropriate citation.

Photographs taken during the internship shall be relevant and sufficiently
clear. Students shall not include confidential, restricted, or proprietary
information without permission from the internship organization.

Vector graphics or high-resolution images are recommended whenever
possible.

\vspace{0.3cm}

% ================================================================
% XIII. TABLES
% ================================================================

\noindent
\textbf{XIII.\quad TABLES}

\vspace{0.15cm}

Tables shall be created using suitable \LaTeX{} environments such as
\verb|table|, \verb|tabular|, \verb|tabularx|, or \verb|longtable|.

Table captions shall be placed above the table.

Example:

\begin{verbatim}
\begin{table}[H]
    \centering
    \caption{Internship schedule}
    \label{tab:internship-schedule}

    \begin{tabular}{lll}
        \toprule
        Week & Task & Result \\
        \midrule
        1 & Orientation & Completed \\
        2 & System study & Completed \\
        \bottomrule
    \end{tabular}
\end{table}
\end{verbatim}

Tables shall be referenced using:

\begin{verbatim}
Table~\ref{tab:internship-schedule}
\end{verbatim}

Students shall not manually type table numbers.

Tables should be created directly in \LaTeX{} rather than inserted as
screenshots whenever practical.

\vspace{0.3cm}

% ================================================================
% XIV. EQUATIONS AND TECHNICAL QUANTITIES
% ================================================================

\noindent
\textbf{XIV.\quad EQUATIONS AND TECHNICAL QUANTITIES}

\vspace{0.15cm}

Mathematical expressions shall be typeset directly using \LaTeX{}.
Images or screenshots of equations shall not be used.

Short expressions may be written in inline mathematical mode, for example:

\begin{verbatim}
$R = 0.7\,\Omega$
\end{verbatim}

Important equations shall use a numbered equation environment.

Example:

\begin{verbatim}
\begin{equation}
    P = UI
    \label{eq:power}
\end{equation}
\end{verbatim}

Equations shall be referenced using:

\begin{verbatim}
Equation~\eqref{eq:power}
\end{verbatim}

Equation numbers shall not be manually entered.

SI units shall be used wherever applicable. The \verb|siunitx| package
should be used where appropriate.

Examples:

\begin{verbatim}
\SI{220}{\volt}
\SI{50}{\hertz}
\SI{2.5}{\ampere}
\end{verbatim}

Notation and units shall be consistent throughout the report.

\vspace{0.3cm}

% ================================================================
% XV. SOURCE CODE
% ================================================================

\noindent
\textbf{XV.\quad SOURCE CODE}

\vspace{0.15cm}

If the internship involves programming, PLCs, microcontrollers,
simulation, or software development, relevant source code may be included
in the report.

Source code shall be presented using the \verb|listings| environment or
another environment defined by the official template.

Example:

\begin{verbatim}
\begin{lstlisting}[language=st]
IF Start = TRUE THEN
    Motor := TRUE;
ELSE
    Motor := FALSE;
END_IF;
\end{lstlisting}
\end{verbatim}

Only code necessary for explaining the work should appear in the main
report. Long source-code listings should be moved to an appendix or
submitted as supplementary files if required.

Source code should not normally be inserted as screenshots.

\vspace{0.3cm}

% ================================================================
% XVI. ABSTRACT
% ================================================================

\noindent
\textbf{XVI.\quad ABSTRACT}

\vspace{0.15cm}

The Abstract shall briefly summarize:

\begin{itemize}[leftmargin=0.8cm, itemsep=2pt]
    \item the internship organization;
    \item the objectives of the internship;
    \item the main tasks performed;
    \item the main results, knowledge, and skills obtained.
\end{itemize}

The Abstract should be concise and should not reproduce the full content
of the report chapter by chapter.

\vspace{0.3cm}

% ================================================================
% XVII. REFERENCES AND CITATIONS
% ================================================================

\noindent
\textbf{XVII.\quad REFERENCES AND CITATIONS}

\vspace{0.15cm}

All books, journal papers, conference papers, standards, websites,
datasheets, manuals, technical documents, figures, tables, or other
external sources used in preparing the report shall be properly cited.

References shall be managed through the bibliography system provided with
the official template.

Bibliographic entries shall normally be stored in:

\begin{verbatim}
References.bib
\end{verbatim}

References shall be cited in the report using:

\begin{verbatim}
\cite{reference-key}
\end{verbatim}

Students shall not manually type reference numbers such as ``[1]'',
``[2]'', etc.

Titles of foreign-language publications shall remain in their original
language in the bibliography.

All borrowed ideas, data, figures, tables, equations, methods, or results
shall be properly acknowledged.

\vspace{0.3cm}

% ================================================================
% XVIII. CROSS-REFERENCING
% ================================================================

\noindent
\textbf{XVIII.\quad CROSS-REFERENCING}

\vspace{0.15cm}

Numbered elements shall be cross-referenced using the standard \LaTeX{}
mechanism.

Use the following commands where appropriate:

\begin{verbatim}
\label{...}
\ref{...}
\eqref{...}
\end{verbatim}

For example:

\begin{verbatim}
As shown in Figure~\ref{fig:system}, ...
\end{verbatim}

Students shall not manually type numbers such as ``Figure 2.3'' when
those numbers can be generated automatically.

\vspace{0.3cm}

% ================================================================
% XIX. PAGE NUMBERING
% ================================================================

\noindent
\textbf{XIX.\quad PAGE NUMBERING}

\vspace{0.15cm}

Cover pages and official evaluation forms at the beginning of the report
shall not display page numbers.

The main body of the report shall use Arabic page numbers beginning from
Chapter 1.

Page numbering shall be generated automatically by the official template.
Students shall not manually enter page numbers.

\vspace{0.3cm}

% ================================================================
% XX. PROJECT FILE ORGANIZATION
% ================================================================

\noindent
\textbf{XX.\quad PROJECT FILE ORGANIZATION}

\vspace{0.15cm}

Students should maintain a clear and organized \LaTeX{} project
structure. A recommended structure is:

\begin{verbatim}
InternshipReport/
|
|-- main.tex
|-- References.bib
|
|-- figures/
|   |-- company.pdf
|   |-- system.png
|   `-- ...
|
`-- chapters/
    |-- chapter1.tex
    |-- chapter2.tex
    `-- chapter3.tex
\end{verbatim}

Figures should normally be stored in the \verb|figures| directory
specified by the official template.

File and directory names should preferably contain Latin letters,
numbers, hyphens, or underscores.

Recommended examples include:

\begin{verbatim}
company_structure.pdf
motor_control.png
chapter_01.tex
\end{verbatim}

Spaces and unnecessary special characters in file names should be
avoided where possible.

\vspace{0.3cm}

% ================================================================
% XXI. DOCUMENT COMPILATION AND FINAL CHECK
% ================================================================

\noindent
\textbf{XXI.\quad DOCUMENT COMPILATION AND FINAL CHECK}

\vspace{0.15cm}

Before submission, the complete \LaTeX{} project shall compile
successfully and produce the final PDF report.

The final document shall not contain unresolved references such as:

\begin{verbatim}
??
[?]
\end{verbatim}

Before generating the final PDF, students shall check:

\begin{itemize}[leftmargin=0.8cm, itemsep=2pt]
    \item the Table of Contents;
    \item page numbering;
    \item chapter, section, and subsection numbering;
    \item figure numbering;
    \item table numbering;
    \item equation numbering;
    \item cross-references;
    \item citations;
    \item bibliography;
    \item font display;
    \item mathematical symbols;
    \item image quality;
    \item completeness of all required forms.
\end{itemize}

The document may be compiled using TeX Live, MiKTeX, or another
compatible \LaTeX{} environment. Regardless of the compilation
environment, the resulting PDF shall conform to the official Faculty
template.

\vspace{0.3cm}

% ================================================================
% XXII. CONTENT THAT STUDENTS MAY EDIT
% ================================================================

\noindent
\textbf{XXII.\quad CONTENT THAT STUDENTS MAY EDIT}

\vspace{0.15cm}

Students may edit or replace:

\begin{itemize}[leftmargin=0.8cm, itemsep=2pt]
    \item student name;
    \item class;
    \item internship organization;
    \item supervisor name;
    \item Abstract;
    \item List of Symbols and Abbreviations;
    \item chapter titles and chapter content;
    \item figures and photographs;
    \item tables;
    \item equations;
    \item citations and references.
\end{itemize}

Students shall not modify the main document formatting settings in the
preamble, including the document class, page geometry, line spacing,
heading formatting, table-of-contents formatting, font configuration,
page layout, or cover-page macros unless specifically instructed.

\vspace{0.3cm}

% ================================================================
% XXIII. FINAL SUBMISSION REQUIREMENTS
% ================================================================

\noindent
\textbf{XXIII.\quad FINAL SUBMISSION REQUIREMENTS}

\vspace{0.15cm}

Before submission, the final report shall:

\begin{itemize}[leftmargin=0.8cm, itemsep=2pt]
    \item compile without errors that affect the report content or presentation;
    \item contain all required forms and report sections;
    \item contain correct student, class, internship organization, and supervisor information;
    \item contain an updated Table of Contents, List of Figures, and List of Tables;
    \item contain no unused example content from the template;
    \item contain no placeholder text such as \textit{Title of Chapter 1},
    \textit{Section 1.1}, or other sample entries;
    \item contain all required citations and references;
    \item contain no unresolved cross-references.
\end{itemize}

If the Faculty requires submission of the \LaTeX{} source, the complete
project shall include all \verb|.tex| files, the bibliography database,
figures, and any other files required to reproduce the submitted PDF.

\vspace{0.4cm}

\noindent
\textbf{Important:}
Before submitting the final internship report, students shall remove or
comment out the following command in the main document:

\begin{verbatim}

% =================================================================
% REGULATIONS FOR PREPARING THE GRADUATION INTERNSHIP REPORT
% USING LATEX
% =================================================================
%
% This file is intended to be included from the main document with:
%
%     \input{Regulations.tex}
%
% IMPORTANT:
% This section is provided for guidance only. Students shall remove
% or comment out the \input{Regulations.tex} command before submitting
% the final internship report.
%
% =================================================================

\clearpage
\thispagestyle{empty}

\begin{center}
    {\fontsize{16pt}{20pt}\selectfont\bfseries
    SPECIFIC REGULATIONS FOR PREPARING\\
    GRADUATION INTERNSHIP REPORTS USING \LaTeX}
\end{center}

\vspace{0.5cm}

{\fontsize{13pt}{17pt}\selectfont
\setstretch{1.3}
\justifying

\noindent
\textbf{Note:}
This section is provided for guidance only and is included in the main
document using the command \verb|\input{Regulations.tex}|.
Students shall remove or comment out this command before submitting the
final internship report.

\vspace{0.4cm}

% ================================================================
% I. GENERAL REQUIREMENTS
% ================================================================

\noindent
\textbf{I.\quad GENERAL REQUIREMENTS}

\vspace{0.15cm}

The Graduation Internship Report shall be prepared using the official
\LaTeX{} template provided by the Faculty. Students shall preserve the
predefined page layout, typography, line spacing, heading styles,
numbering system, figure and table formatting, and bibliography settings
unless otherwise instructed by the Faculty or the supervisor.

The report shall be written in English.

Students are responsible for the academic and technical content of the
report. The official \LaTeX{} template is responsible for standardized
document formatting.

\vspace{0.3cm}

% ================================================================
% II. REPORT STRUCTURE
% ================================================================

\noindent
\textbf{II.\quad REPORT STRUCTURE}

\vspace{0.15cm}

The report shall normally be arranged in the following order:

\begin{enumerate}[label=\arabic*., leftmargin=1.2cm, itemsep=2pt]
    \item Outer Cover;
    \item Comments of the Internship Host Organization;
    \item Supervisor's Evaluation and Graduation Internship Defense Results;
    \item Table of Contents;
    \item List of Symbols and Abbreviations, if applicable;
    \item List of Figures;
    \item List of Tables;
    \item Abstract;
    \item Main Chapters of the Internship Report;
    \item References;
    \item Appendices, if applicable.
\end{enumerate}

The preliminary pages shall not be numbered as chapters. The main report
content shall be organized into consecutively numbered chapters beginning
with Chapter 1.

\vspace{0.3cm}

% ================================================================
% III. PAGE SIZE AND MARGINS
% ================================================================

\noindent
\textbf{III.\quad PAGE SIZE AND MARGINS}

\vspace{0.15cm}

The report shall be formatted on A4 paper with the following margins:

\begin{center}
\begin{tabular}{ll}
    Left margin:   & 3.5 cm \\
    Right margin:  & 2.0 cm \\
    Top margin:    & 2.5 cm \\
    Bottom margin: & 2.5 cm
\end{tabular}
\end{center}

These settings are predefined in the official template and shall not be
modified by students.

\vspace{0.3cm}

% ================================================================
% IV. FONT, FONT SIZE, AND LINE SPACING
% ================================================================

\noindent
\textbf{IV.\quad FONT, FONT SIZE, AND LINE SPACING}

\vspace{0.15cm}

The main text shall use a 13 pt font size with 1.3-line spacing.
A Times-family typeface shall be used in accordance with the official
template.

The text shall be fully justified.

Students shall not use repeated spaces, repeated manual line breaks, or
multiple blank lines to position text. Paragraph and heading spacing shall
be controlled by the template.

\vspace{0.3cm}

% ================================================================
% V. OUTER COVER
% ================================================================

\noindent
\textbf{V.\quad OUTER COVER}

\vspace{0.15cm}

The outer cover shall follow the official format provided in the template
and shall contain the following information:

\begin{itemize}[leftmargin=0.8cm, itemsep=2pt]
    \item Thai Nguyen University of Technology;
    \item Faculty of International Training;
    \item Engineering and Technology Department;
    \item University logo and Faculty logo;
    \item the title \textbf{INTERNSHIP REPORT};
    \item Internship Organization;
    \item Author;
    \item Class;
    \item Supervisor;
    \item place and year of completion.
\end{itemize}

Students shall only replace the required information fields and shall not
modify the cover-page layout unless instructed otherwise.

\vspace{0.3cm}

% ================================================================
% VI. COMMENTS OF THE INTERNSHIP HOST ORGANIZATION
% ================================================================

\noindent
\textbf{VI.\quad COMMENTS OF THE INTERNSHIP HOST ORGANIZATION}

\vspace{0.15cm}

The page entitled \textbf{COMMENTS OF THE INTERNSHIP HOST ORGANIZATION}
is reserved for the internship host organization.

It includes:

\begin{itemize}[leftmargin=0.8cm, itemsep=2pt]
    \item the student's full name;
    \item class;
    \item comments from the host organization;
    \item score;
    \item date;
    \item signature and official seal.
\end{itemize}

Students shall not alter or complete official evaluation fields unless
specifically authorized.

\vspace{0.3cm}

% ================================================================
% VII. SUPERVISOR'S EVALUATION AND DEFENSE RESULTS
% ================================================================

\noindent
\textbf{VII.\quad SUPERVISOR'S EVALUATION AND DEFENSE RESULTS}

\vspace{0.15cm}

The section entitled
\textbf{SUPERVISOR'S EVALUATION AND GRADUATION INTERNSHIP DEFENSE RESULTS}
is reserved for the supervisor and defense committee.

The evaluation may include:

\begin{itemize}[leftmargin=0.8cm, itemsep=2pt]
    \item Internship Progress and Student's Attitude;
    \item Report Content;
    \item Report Presentation;
    \item Other Comments;
    \item Supervisor's Opinion;
    \item Defense Result;
    \item Score.
\end{itemize}

Students shall not enter evaluation scores or defense results before they
are officially determined.

\vspace{0.3cm}

% ================================================================
% VIII. CHAPTERS, SECTIONS, AND SUBSECTIONS
% ================================================================

\noindent
\textbf{VIII.\quad CHAPTERS, SECTIONS, AND SUBSECTIONS}

\vspace{0.15cm}

The report structure shall be generated using the standard \LaTeX{}
commands:

\begin{verbatim}
\chapter{...}
\section{...}
\subsection{...}
\end{verbatim}

Chapter, section, and subsection numbers shall be generated automatically
by \LaTeX{}.

Students shall not manually type chapter or section numbers and shall not
reproduce heading styles using commands such as \verb|\textbf|,
\verb|\fontsize|, or \verb|\centering| instead of the appropriate
structural command.

Excessive subdivision should be avoided. In most cases, the hierarchy
Chapter--Section--Subsection is sufficient.

\vspace{0.3cm}

% ================================================================
% IX. ORGANIZATION OF THE MAIN CONTENT
% ================================================================

\noindent
\textbf{IX.\quad ORGANIZATION OF THE MAIN CONTENT}

\vspace{0.15cm}

The specific content of the report shall depend on the internship
organization and the assigned tasks. A typical structure may include:

\begin{itemize}[leftmargin=0.8cm, itemsep=2pt]
    \item \textbf{Chapter 1. Overview of the Internship Organization};
    \item \textbf{Chapter 2. Internship Activities and Assigned Tasks};
    \item \textbf{Chapter 3. Results, Experience, and Evaluation}.
\end{itemize}

These chapter titles are recommendations only and may be adapted to the
actual internship content.

The report shall reflect work actually performed, observed, analyzed, or
learned during the internship. Unrelated theoretical material should not
be included merely to increase report length.

\vspace{0.3cm}

% ================================================================
% X. TABLE OF CONTENTS
% ================================================================

\noindent
\textbf{X.\quad TABLE OF CONTENTS}

\vspace{0.15cm}

The Table of Contents shall be generated automatically using:

\begin{verbatim}
\tableofcontents
\end{verbatim}

Students shall not create or manually edit page numbers in the Table of
Contents.

The document shall be compiled sufficiently to ensure that all page
numbers and entries are correctly updated.

\vspace{0.3cm}

% ================================================================
% XI. LIST OF SYMBOLS AND ABBREVIATIONS
% ================================================================

\noindent
\textbf{XI.\quad LIST OF SYMBOLS AND ABBREVIATIONS}

\vspace{0.15cm}

Symbols and abbreviations that are used repeatedly in the report shall be
listed in the \textbf{List of Symbols and Abbreviations}.

The list may contain the following columns:

\begin{center}
\begin{tabular}{llll}
    No. & Symbol / Abbreviation & Unit & Definition
\end{tabular}
\end{center}

Only symbols and abbreviations actually used in the final report shall be
retained. Example entries contained in the template shall be removed if
they are not relevant to the student's report.

Each abbreviation should normally be defined when it first appears in
the main text.

\vspace{0.3cm}

% ================================================================
% XII. FIGURES, DIAGRAMS, AND PHOTOGRAPHS
% ================================================================

\noindent
\textbf{XII.\quad FIGURES, DIAGRAMS, AND PHOTOGRAPHS}

\vspace{0.15cm}

Figures shall be clear, legible, and directly relevant to the report
content.

Figures shall be inserted using the standard \verb|figure| environment.

Example:

\begin{verbatim}
\begin{figure}[H]
    \centering
    \includegraphics[width=0.8\textwidth]
                    {figures/company_structure.pdf}
    \caption{Organizational structure of the company}
    \label{fig:company-structure}
\end{figure}
\end{verbatim}

Figure captions shall be placed below the figure.

Each referenced figure shall have a unique \verb|\label| and shall be
referenced in the text using, for example:

\begin{verbatim}
Figure~\ref{fig:company-structure}
\end{verbatim}

Students shall not manually type figure numbers.

Figures reproduced or adapted from external sources shall include an
appropriate citation.

Photographs taken during the internship shall be relevant and sufficiently
clear. Students shall not include confidential, restricted, or proprietary
information without permission from the internship organization.

Vector graphics or high-resolution images are recommended whenever
possible.

\vspace{0.3cm}

% ================================================================
% XIII. TABLES
% ================================================================

\noindent
\textbf{XIII.\quad TABLES}

\vspace{0.15cm}

Tables shall be created using suitable \LaTeX{} environments such as
\verb|table|, \verb|tabular|, \verb|tabularx|, or \verb|longtable|.

Table captions shall be placed above the table.

Example:

\begin{verbatim}
\begin{table}[H]
    \centering
    \caption{Internship schedule}
    \label{tab:internship-schedule}

    \begin{tabular}{lll}
        \toprule
        Week & Task & Result \\
        \midrule
        1 & Orientation & Completed \\
        2 & System study & Completed \\
        \bottomrule
    \end{tabular}
\end{table}
\end{verbatim}

Tables shall be referenced using:

\begin{verbatim}
Table~\ref{tab:internship-schedule}
\end{verbatim}

Students shall not manually type table numbers.

Tables should be created directly in \LaTeX{} rather than inserted as
screenshots whenever practical.

\vspace{0.3cm}

% ================================================================
% XIV. EQUATIONS AND TECHNICAL QUANTITIES
% ================================================================

\noindent
\textbf{XIV.\quad EQUATIONS AND TECHNICAL QUANTITIES}

\vspace{0.15cm}

Mathematical expressions shall be typeset directly using \LaTeX{}.
Images or screenshots of equations shall not be used.

Short expressions may be written in inline mathematical mode, for example:

\begin{verbatim}
$R = 0.7\,\Omega$
\end{verbatim}

Important equations shall use a numbered equation environment.

Example:

\begin{verbatim}
\begin{equation}
    P = UI
    \label{eq:power}
\end{equation}
\end{verbatim}

Equations shall be referenced using:

\begin{verbatim}
Equation~\eqref{eq:power}
\end{verbatim}

Equation numbers shall not be manually entered.

SI units shall be used wherever applicable. The \verb|siunitx| package
should be used where appropriate.

Examples:

\begin{verbatim}
\SI{220}{\volt}
\SI{50}{\hertz}
\SI{2.5}{\ampere}
\end{verbatim}

Notation and units shall be consistent throughout the report.

\vspace{0.3cm}

% ================================================================
% XV. SOURCE CODE
% ================================================================

\noindent
\textbf{XV.\quad SOURCE CODE}

\vspace{0.15cm}

If the internship involves programming, PLCs, microcontrollers,
simulation, or software development, relevant source code may be included
in the report.

Source code shall be presented using the \verb|listings| environment or
another environment defined by the official template.

Example:

\begin{verbatim}
\begin{lstlisting}[language=st]
IF Start = TRUE THEN
    Motor := TRUE;
ELSE
    Motor := FALSE;
END_IF;
\end{lstlisting}
\end{verbatim}

Only code necessary for explaining the work should appear in the main
report. Long source-code listings should be moved to an appendix or
submitted as supplementary files if required.

Source code should not normally be inserted as screenshots.

\vspace{0.3cm}

% ================================================================
% XVI. ABSTRACT
% ================================================================

\noindent
\textbf{XVI.\quad ABSTRACT}

\vspace{0.15cm}

The Abstract shall briefly summarize:

\begin{itemize}[leftmargin=0.8cm, itemsep=2pt]
    \item the internship organization;
    \item the objectives of the internship;
    \item the main tasks performed;
    \item the main results, knowledge, and skills obtained.
\end{itemize}

The Abstract should be concise and should not reproduce the full content
of the report chapter by chapter.

\vspace{0.3cm}

% ================================================================
% XVII. REFERENCES AND CITATIONS
% ================================================================

\noindent
\textbf{XVII.\quad REFERENCES AND CITATIONS}

\vspace{0.15cm}

All books, journal papers, conference papers, standards, websites,
datasheets, manuals, technical documents, figures, tables, or other
external sources used in preparing the report shall be properly cited.

References shall be managed through the bibliography system provided with
the official template.

Bibliographic entries shall normally be stored in:

\begin{verbatim}
References.bib
\end{verbatim}

References shall be cited in the report using:

\begin{verbatim}
\cite{reference-key}
\end{verbatim}

Students shall not manually type reference numbers such as ``[1]'',
``[2]'', etc.

Titles of foreign-language publications shall remain in their original
language in the bibliography.

All borrowed ideas, data, figures, tables, equations, methods, or results
shall be properly acknowledged.

\vspace{0.3cm}

% ================================================================
% XVIII. CROSS-REFERENCING
% ================================================================

\noindent
\textbf{XVIII.\quad CROSS-REFERENCING}

\vspace{0.15cm}

Numbered elements shall be cross-referenced using the standard \LaTeX{}
mechanism.

Use the following commands where appropriate:

\begin{verbatim}
\label{...}
\ref{...}
\eqref{...}
\end{verbatim}

For example:

\begin{verbatim}
As shown in Figure~\ref{fig:system}, ...
\end{verbatim}

Students shall not manually type numbers such as ``Figure 2.3'' when
those numbers can be generated automatically.

\vspace{0.3cm}

% ================================================================
% XIX. PAGE NUMBERING
% ================================================================

\noindent
\textbf{XIX.\quad PAGE NUMBERING}

\vspace{0.15cm}

Cover pages and official evaluation forms at the beginning of the report
shall not display page numbers.

The main body of the report shall use Arabic page numbers beginning from
Chapter 1.

Page numbering shall be generated automatically by the official template.
Students shall not manually enter page numbers.

\vspace{0.3cm}

% ================================================================
% XX. PROJECT FILE ORGANIZATION
% ================================================================

\noindent
\textbf{XX.\quad PROJECT FILE ORGANIZATION}

\vspace{0.15cm}

Students should maintain a clear and organized \LaTeX{} project
structure. A recommended structure is:

\begin{verbatim}
InternshipReport/
|
|-- main.tex
|-- References.bib
|
|-- figures/
|   |-- company.pdf
|   |-- system.png
|   `-- ...
|
`-- chapters/
    |-- chapter1.tex
    |-- chapter2.tex
    `-- chapter3.tex
\end{verbatim}

Figures should normally be stored in the \verb|figures| directory
specified by the official template.

File and directory names should preferably contain Latin letters,
numbers, hyphens, or underscores.

Recommended examples include:

\begin{verbatim}
company_structure.pdf
motor_control.png
chapter_01.tex
\end{verbatim}

Spaces and unnecessary special characters in file names should be
avoided where possible.

\vspace{0.3cm}

% ================================================================
% XXI. DOCUMENT COMPILATION AND FINAL CHECK
% ================================================================

\noindent
\textbf{XXI.\quad DOCUMENT COMPILATION AND FINAL CHECK}

\vspace{0.15cm}

Before submission, the complete \LaTeX{} project shall compile
successfully and produce the final PDF report.

The final document shall not contain unresolved references such as:

\begin{verbatim}
??
[?]
\end{verbatim}

Before generating the final PDF, students shall check:

\begin{itemize}[leftmargin=0.8cm, itemsep=2pt]
    \item the Table of Contents;
    \item page numbering;
    \item chapter, section, and subsection numbering;
    \item figure numbering;
    \item table numbering;
    \item equation numbering;
    \item cross-references;
    \item citations;
    \item bibliography;
    \item font display;
    \item mathematical symbols;
    \item image quality;
    \item completeness of all required forms.
\end{itemize}

The document may be compiled using TeX Live, MiKTeX, or another
compatible \LaTeX{} environment. Regardless of the compilation
environment, the resulting PDF shall conform to the official Faculty
template.

\vspace{0.3cm}

% ================================================================
% XXII. CONTENT THAT STUDENTS MAY EDIT
% ================================================================

\noindent
\textbf{XXII.\quad CONTENT THAT STUDENTS MAY EDIT}

\vspace{0.15cm}

Students may edit or replace:

\begin{itemize}[leftmargin=0.8cm, itemsep=2pt]
    \item student name;
    \item class;
    \item internship organization;
    \item supervisor name;
    \item Abstract;
    \item List of Symbols and Abbreviations;
    \item chapter titles and chapter content;
    \item figures and photographs;
    \item tables;
    \item equations;
    \item citations and references.
\end{itemize}

Students shall not modify the main document formatting settings in the
preamble, including the document class, page geometry, line spacing,
heading formatting, table-of-contents formatting, font configuration,
page layout, or cover-page macros unless specifically instructed.

\vspace{0.3cm}

% ================================================================
% XXIII. FINAL SUBMISSION REQUIREMENTS
% ================================================================

\noindent
\textbf{XXIII.\quad FINAL SUBMISSION REQUIREMENTS}

\vspace{0.15cm}

Before submission, the final report shall:

\begin{itemize}[leftmargin=0.8cm, itemsep=2pt]
    \item compile without errors that affect the report content or presentation;
    \item contain all required forms and report sections;
    \item contain correct student, class, internship organization, and supervisor information;
    \item contain an updated Table of Contents, List of Figures, and List of Tables;
    \item contain no unused example content from the template;
    \item contain no placeholder text such as \textit{Title of Chapter 1},
    \textit{Section 1.1}, or other sample entries;
    \item contain all required citations and references;
    \item contain no unresolved cross-references.
\end{itemize}

If the Faculty requires submission of the \LaTeX{} source, the complete
project shall include all \verb|.tex| files, the bibliography database,
figures, and any other files required to reproduce the submitted PDF.

\vspace{0.4cm}

\noindent
\textbf{Important:}
Before submitting the final internship report, students shall remove or
comment out the following command in the main document:

\begin{verbatim}
\input{Regulations.tex}
\end{verbatim}

For example:

\begin{verbatim}
% \input{Regulations.tex}
\end{verbatim}

This removes the guidance pages from the final submitted report while
preserving the regulations file in the \LaTeX{} project for reference.

} % End 13pt content

\clearpage

\end{verbatim}

For example:

\begin{verbatim}
% 
% =================================================================
% REGULATIONS FOR PREPARING THE GRADUATION INTERNSHIP REPORT
% USING LATEX
% =================================================================
%
% This file is intended to be included from the main document with:
%
%     \input{Regulations.tex}
%
% IMPORTANT:
% This section is provided for guidance only. Students shall remove
% or comment out the \input{Regulations.tex} command before submitting
% the final internship report.
%
% =================================================================

\clearpage
\thispagestyle{empty}

\begin{center}
    {\fontsize{16pt}{20pt}\selectfont\bfseries
    SPECIFIC REGULATIONS FOR PREPARING\\
    GRADUATION INTERNSHIP REPORTS USING \LaTeX}
\end{center}

\vspace{0.5cm}

{\fontsize{13pt}{17pt}\selectfont
\setstretch{1.3}
\justifying

\noindent
\textbf{Note:}
This section is provided for guidance only and is included in the main
document using the command \verb|\input{Regulations.tex}|.
Students shall remove or comment out this command before submitting the
final internship report.

\vspace{0.4cm}

% ================================================================
% I. GENERAL REQUIREMENTS
% ================================================================

\noindent
\textbf{I.\quad GENERAL REQUIREMENTS}

\vspace{0.15cm}

The Graduation Internship Report shall be prepared using the official
\LaTeX{} template provided by the Faculty. Students shall preserve the
predefined page layout, typography, line spacing, heading styles,
numbering system, figure and table formatting, and bibliography settings
unless otherwise instructed by the Faculty or the supervisor.

The report shall be written in English.

Students are responsible for the academic and technical content of the
report. The official \LaTeX{} template is responsible for standardized
document formatting.

\vspace{0.3cm}

% ================================================================
% II. REPORT STRUCTURE
% ================================================================

\noindent
\textbf{II.\quad REPORT STRUCTURE}

\vspace{0.15cm}

The report shall normally be arranged in the following order:

\begin{enumerate}[label=\arabic*., leftmargin=1.2cm, itemsep=2pt]
    \item Outer Cover;
    \item Comments of the Internship Host Organization;
    \item Supervisor's Evaluation and Graduation Internship Defense Results;
    \item Table of Contents;
    \item List of Symbols and Abbreviations, if applicable;
    \item List of Figures;
    \item List of Tables;
    \item Abstract;
    \item Main Chapters of the Internship Report;
    \item References;
    \item Appendices, if applicable.
\end{enumerate}

The preliminary pages shall not be numbered as chapters. The main report
content shall be organized into consecutively numbered chapters beginning
with Chapter 1.

\vspace{0.3cm}

% ================================================================
% III. PAGE SIZE AND MARGINS
% ================================================================

\noindent
\textbf{III.\quad PAGE SIZE AND MARGINS}

\vspace{0.15cm}

The report shall be formatted on A4 paper with the following margins:

\begin{center}
\begin{tabular}{ll}
    Left margin:   & 3.5 cm \\
    Right margin:  & 2.0 cm \\
    Top margin:    & 2.5 cm \\
    Bottom margin: & 2.5 cm
\end{tabular}
\end{center}

These settings are predefined in the official template and shall not be
modified by students.

\vspace{0.3cm}

% ================================================================
% IV. FONT, FONT SIZE, AND LINE SPACING
% ================================================================

\noindent
\textbf{IV.\quad FONT, FONT SIZE, AND LINE SPACING}

\vspace{0.15cm}

The main text shall use a 13 pt font size with 1.3-line spacing.
A Times-family typeface shall be used in accordance with the official
template.

The text shall be fully justified.

Students shall not use repeated spaces, repeated manual line breaks, or
multiple blank lines to position text. Paragraph and heading spacing shall
be controlled by the template.

\vspace{0.3cm}

% ================================================================
% V. OUTER COVER
% ================================================================

\noindent
\textbf{V.\quad OUTER COVER}

\vspace{0.15cm}

The outer cover shall follow the official format provided in the template
and shall contain the following information:

\begin{itemize}[leftmargin=0.8cm, itemsep=2pt]
    \item Thai Nguyen University of Technology;
    \item Faculty of International Training;
    \item Engineering and Technology Department;
    \item University logo and Faculty logo;
    \item the title \textbf{INTERNSHIP REPORT};
    \item Internship Organization;
    \item Author;
    \item Class;
    \item Supervisor;
    \item place and year of completion.
\end{itemize}

Students shall only replace the required information fields and shall not
modify the cover-page layout unless instructed otherwise.

\vspace{0.3cm}

% ================================================================
% VI. COMMENTS OF THE INTERNSHIP HOST ORGANIZATION
% ================================================================

\noindent
\textbf{VI.\quad COMMENTS OF THE INTERNSHIP HOST ORGANIZATION}

\vspace{0.15cm}

The page entitled \textbf{COMMENTS OF THE INTERNSHIP HOST ORGANIZATION}
is reserved for the internship host organization.

It includes:

\begin{itemize}[leftmargin=0.8cm, itemsep=2pt]
    \item the student's full name;
    \item class;
    \item comments from the host organization;
    \item score;
    \item date;
    \item signature and official seal.
\end{itemize}

Students shall not alter or complete official evaluation fields unless
specifically authorized.

\vspace{0.3cm}

% ================================================================
% VII. SUPERVISOR'S EVALUATION AND DEFENSE RESULTS
% ================================================================

\noindent
\textbf{VII.\quad SUPERVISOR'S EVALUATION AND DEFENSE RESULTS}

\vspace{0.15cm}

The section entitled
\textbf{SUPERVISOR'S EVALUATION AND GRADUATION INTERNSHIP DEFENSE RESULTS}
is reserved for the supervisor and defense committee.

The evaluation may include:

\begin{itemize}[leftmargin=0.8cm, itemsep=2pt]
    \item Internship Progress and Student's Attitude;
    \item Report Content;
    \item Report Presentation;
    \item Other Comments;
    \item Supervisor's Opinion;
    \item Defense Result;
    \item Score.
\end{itemize}

Students shall not enter evaluation scores or defense results before they
are officially determined.

\vspace{0.3cm}

% ================================================================
% VIII. CHAPTERS, SECTIONS, AND SUBSECTIONS
% ================================================================

\noindent
\textbf{VIII.\quad CHAPTERS, SECTIONS, AND SUBSECTIONS}

\vspace{0.15cm}

The report structure shall be generated using the standard \LaTeX{}
commands:

\begin{verbatim}
\chapter{...}
\section{...}
\subsection{...}
\end{verbatim}

Chapter, section, and subsection numbers shall be generated automatically
by \LaTeX{}.

Students shall not manually type chapter or section numbers and shall not
reproduce heading styles using commands such as \verb|\textbf|,
\verb|\fontsize|, or \verb|\centering| instead of the appropriate
structural command.

Excessive subdivision should be avoided. In most cases, the hierarchy
Chapter--Section--Subsection is sufficient.

\vspace{0.3cm}

% ================================================================
% IX. ORGANIZATION OF THE MAIN CONTENT
% ================================================================

\noindent
\textbf{IX.\quad ORGANIZATION OF THE MAIN CONTENT}

\vspace{0.15cm}

The specific content of the report shall depend on the internship
organization and the assigned tasks. A typical structure may include:

\begin{itemize}[leftmargin=0.8cm, itemsep=2pt]
    \item \textbf{Chapter 1. Overview of the Internship Organization};
    \item \textbf{Chapter 2. Internship Activities and Assigned Tasks};
    \item \textbf{Chapter 3. Results, Experience, and Evaluation}.
\end{itemize}

These chapter titles are recommendations only and may be adapted to the
actual internship content.

The report shall reflect work actually performed, observed, analyzed, or
learned during the internship. Unrelated theoretical material should not
be included merely to increase report length.

\vspace{0.3cm}

% ================================================================
% X. TABLE OF CONTENTS
% ================================================================

\noindent
\textbf{X.\quad TABLE OF CONTENTS}

\vspace{0.15cm}

The Table of Contents shall be generated automatically using:

\begin{verbatim}
\tableofcontents
\end{verbatim}

Students shall not create or manually edit page numbers in the Table of
Contents.

The document shall be compiled sufficiently to ensure that all page
numbers and entries are correctly updated.

\vspace{0.3cm}

% ================================================================
% XI. LIST OF SYMBOLS AND ABBREVIATIONS
% ================================================================

\noindent
\textbf{XI.\quad LIST OF SYMBOLS AND ABBREVIATIONS}

\vspace{0.15cm}

Symbols and abbreviations that are used repeatedly in the report shall be
listed in the \textbf{List of Symbols and Abbreviations}.

The list may contain the following columns:

\begin{center}
\begin{tabular}{llll}
    No. & Symbol / Abbreviation & Unit & Definition
\end{tabular}
\end{center}

Only symbols and abbreviations actually used in the final report shall be
retained. Example entries contained in the template shall be removed if
they are not relevant to the student's report.

Each abbreviation should normally be defined when it first appears in
the main text.

\vspace{0.3cm}

% ================================================================
% XII. FIGURES, DIAGRAMS, AND PHOTOGRAPHS
% ================================================================

\noindent
\textbf{XII.\quad FIGURES, DIAGRAMS, AND PHOTOGRAPHS}

\vspace{0.15cm}

Figures shall be clear, legible, and directly relevant to the report
content.

Figures shall be inserted using the standard \verb|figure| environment.

Example:

\begin{verbatim}
\begin{figure}[H]
    \centering
    \includegraphics[width=0.8\textwidth]
                    {figures/company_structure.pdf}
    \caption{Organizational structure of the company}
    \label{fig:company-structure}
\end{figure}
\end{verbatim}

Figure captions shall be placed below the figure.

Each referenced figure shall have a unique \verb|\label| and shall be
referenced in the text using, for example:

\begin{verbatim}
Figure~\ref{fig:company-structure}
\end{verbatim}

Students shall not manually type figure numbers.

Figures reproduced or adapted from external sources shall include an
appropriate citation.

Photographs taken during the internship shall be relevant and sufficiently
clear. Students shall not include confidential, restricted, or proprietary
information without permission from the internship organization.

Vector graphics or high-resolution images are recommended whenever
possible.

\vspace{0.3cm}

% ================================================================
% XIII. TABLES
% ================================================================

\noindent
\textbf{XIII.\quad TABLES}

\vspace{0.15cm}

Tables shall be created using suitable \LaTeX{} environments such as
\verb|table|, \verb|tabular|, \verb|tabularx|, or \verb|longtable|.

Table captions shall be placed above the table.

Example:

\begin{verbatim}
\begin{table}[H]
    \centering
    \caption{Internship schedule}
    \label{tab:internship-schedule}

    \begin{tabular}{lll}
        \toprule
        Week & Task & Result \\
        \midrule
        1 & Orientation & Completed \\
        2 & System study & Completed \\
        \bottomrule
    \end{tabular}
\end{table}
\end{verbatim}

Tables shall be referenced using:

\begin{verbatim}
Table~\ref{tab:internship-schedule}
\end{verbatim}

Students shall not manually type table numbers.

Tables should be created directly in \LaTeX{} rather than inserted as
screenshots whenever practical.

\vspace{0.3cm}

% ================================================================
% XIV. EQUATIONS AND TECHNICAL QUANTITIES
% ================================================================

\noindent
\textbf{XIV.\quad EQUATIONS AND TECHNICAL QUANTITIES}

\vspace{0.15cm}

Mathematical expressions shall be typeset directly using \LaTeX{}.
Images or screenshots of equations shall not be used.

Short expressions may be written in inline mathematical mode, for example:

\begin{verbatim}
$R = 0.7\,\Omega$
\end{verbatim}

Important equations shall use a numbered equation environment.

Example:

\begin{verbatim}
\begin{equation}
    P = UI
    \label{eq:power}
\end{equation}
\end{verbatim}

Equations shall be referenced using:

\begin{verbatim}
Equation~\eqref{eq:power}
\end{verbatim}

Equation numbers shall not be manually entered.

SI units shall be used wherever applicable. The \verb|siunitx| package
should be used where appropriate.

Examples:

\begin{verbatim}
\SI{220}{\volt}
\SI{50}{\hertz}
\SI{2.5}{\ampere}
\end{verbatim}

Notation and units shall be consistent throughout the report.

\vspace{0.3cm}

% ================================================================
% XV. SOURCE CODE
% ================================================================

\noindent
\textbf{XV.\quad SOURCE CODE}

\vspace{0.15cm}

If the internship involves programming, PLCs, microcontrollers,
simulation, or software development, relevant source code may be included
in the report.

Source code shall be presented using the \verb|listings| environment or
another environment defined by the official template.

Example:

\begin{verbatim}
\begin{lstlisting}[language=st]
IF Start = TRUE THEN
    Motor := TRUE;
ELSE
    Motor := FALSE;
END_IF;
\end{lstlisting}
\end{verbatim}

Only code necessary for explaining the work should appear in the main
report. Long source-code listings should be moved to an appendix or
submitted as supplementary files if required.

Source code should not normally be inserted as screenshots.

\vspace{0.3cm}

% ================================================================
% XVI. ABSTRACT
% ================================================================

\noindent
\textbf{XVI.\quad ABSTRACT}

\vspace{0.15cm}

The Abstract shall briefly summarize:

\begin{itemize}[leftmargin=0.8cm, itemsep=2pt]
    \item the internship organization;
    \item the objectives of the internship;
    \item the main tasks performed;
    \item the main results, knowledge, and skills obtained.
\end{itemize}

The Abstract should be concise and should not reproduce the full content
of the report chapter by chapter.

\vspace{0.3cm}

% ================================================================
% XVII. REFERENCES AND CITATIONS
% ================================================================

\noindent
\textbf{XVII.\quad REFERENCES AND CITATIONS}

\vspace{0.15cm}

All books, journal papers, conference papers, standards, websites,
datasheets, manuals, technical documents, figures, tables, or other
external sources used in preparing the report shall be properly cited.

References shall be managed through the bibliography system provided with
the official template.

Bibliographic entries shall normally be stored in:

\begin{verbatim}
References.bib
\end{verbatim}

References shall be cited in the report using:

\begin{verbatim}
\cite{reference-key}
\end{verbatim}

Students shall not manually type reference numbers such as ``[1]'',
``[2]'', etc.

Titles of foreign-language publications shall remain in their original
language in the bibliography.

All borrowed ideas, data, figures, tables, equations, methods, or results
shall be properly acknowledged.

\vspace{0.3cm}

% ================================================================
% XVIII. CROSS-REFERENCING
% ================================================================

\noindent
\textbf{XVIII.\quad CROSS-REFERENCING}

\vspace{0.15cm}

Numbered elements shall be cross-referenced using the standard \LaTeX{}
mechanism.

Use the following commands where appropriate:

\begin{verbatim}
\label{...}
\ref{...}
\eqref{...}
\end{verbatim}

For example:

\begin{verbatim}
As shown in Figure~\ref{fig:system}, ...
\end{verbatim}

Students shall not manually type numbers such as ``Figure 2.3'' when
those numbers can be generated automatically.

\vspace{0.3cm}

% ================================================================
% XIX. PAGE NUMBERING
% ================================================================

\noindent
\textbf{XIX.\quad PAGE NUMBERING}

\vspace{0.15cm}

Cover pages and official evaluation forms at the beginning of the report
shall not display page numbers.

The main body of the report shall use Arabic page numbers beginning from
Chapter 1.

Page numbering shall be generated automatically by the official template.
Students shall not manually enter page numbers.

\vspace{0.3cm}

% ================================================================
% XX. PROJECT FILE ORGANIZATION
% ================================================================

\noindent
\textbf{XX.\quad PROJECT FILE ORGANIZATION}

\vspace{0.15cm}

Students should maintain a clear and organized \LaTeX{} project
structure. A recommended structure is:

\begin{verbatim}
InternshipReport/
|
|-- main.tex
|-- References.bib
|
|-- figures/
|   |-- company.pdf
|   |-- system.png
|   `-- ...
|
`-- chapters/
    |-- chapter1.tex
    |-- chapter2.tex
    `-- chapter3.tex
\end{verbatim}

Figures should normally be stored in the \verb|figures| directory
specified by the official template.

File and directory names should preferably contain Latin letters,
numbers, hyphens, or underscores.

Recommended examples include:

\begin{verbatim}
company_structure.pdf
motor_control.png
chapter_01.tex
\end{verbatim}

Spaces and unnecessary special characters in file names should be
avoided where possible.

\vspace{0.3cm}

% ================================================================
% XXI. DOCUMENT COMPILATION AND FINAL CHECK
% ================================================================

\noindent
\textbf{XXI.\quad DOCUMENT COMPILATION AND FINAL CHECK}

\vspace{0.15cm}

Before submission, the complete \LaTeX{} project shall compile
successfully and produce the final PDF report.

The final document shall not contain unresolved references such as:

\begin{verbatim}
??
[?]
\end{verbatim}

Before generating the final PDF, students shall check:

\begin{itemize}[leftmargin=0.8cm, itemsep=2pt]
    \item the Table of Contents;
    \item page numbering;
    \item chapter, section, and subsection numbering;
    \item figure numbering;
    \item table numbering;
    \item equation numbering;
    \item cross-references;
    \item citations;
    \item bibliography;
    \item font display;
    \item mathematical symbols;
    \item image quality;
    \item completeness of all required forms.
\end{itemize}

The document may be compiled using TeX Live, MiKTeX, or another
compatible \LaTeX{} environment. Regardless of the compilation
environment, the resulting PDF shall conform to the official Faculty
template.

\vspace{0.3cm}

% ================================================================
% XXII. CONTENT THAT STUDENTS MAY EDIT
% ================================================================

\noindent
\textbf{XXII.\quad CONTENT THAT STUDENTS MAY EDIT}

\vspace{0.15cm}

Students may edit or replace:

\begin{itemize}[leftmargin=0.8cm, itemsep=2pt]
    \item student name;
    \item class;
    \item internship organization;
    \item supervisor name;
    \item Abstract;
    \item List of Symbols and Abbreviations;
    \item chapter titles and chapter content;
    \item figures and photographs;
    \item tables;
    \item equations;
    \item citations and references.
\end{itemize}

Students shall not modify the main document formatting settings in the
preamble, including the document class, page geometry, line spacing,
heading formatting, table-of-contents formatting, font configuration,
page layout, or cover-page macros unless specifically instructed.

\vspace{0.3cm}

% ================================================================
% XXIII. FINAL SUBMISSION REQUIREMENTS
% ================================================================

\noindent
\textbf{XXIII.\quad FINAL SUBMISSION REQUIREMENTS}

\vspace{0.15cm}

Before submission, the final report shall:

\begin{itemize}[leftmargin=0.8cm, itemsep=2pt]
    \item compile without errors that affect the report content or presentation;
    \item contain all required forms and report sections;
    \item contain correct student, class, internship organization, and supervisor information;
    \item contain an updated Table of Contents, List of Figures, and List of Tables;
    \item contain no unused example content from the template;
    \item contain no placeholder text such as \textit{Title of Chapter 1},
    \textit{Section 1.1}, or other sample entries;
    \item contain all required citations and references;
    \item contain no unresolved cross-references.
\end{itemize}

If the Faculty requires submission of the \LaTeX{} source, the complete
project shall include all \verb|.tex| files, the bibliography database,
figures, and any other files required to reproduce the submitted PDF.

\vspace{0.4cm}

\noindent
\textbf{Important:}
Before submitting the final internship report, students shall remove or
comment out the following command in the main document:

\begin{verbatim}
\input{Regulations.tex}
\end{verbatim}

For example:

\begin{verbatim}
% \input{Regulations.tex}
\end{verbatim}

This removes the guidance pages from the final submitted report while
preserving the regulations file in the \LaTeX{} project for reference.

} % End 13pt content

\clearpage

\end{verbatim}

This removes the guidance pages from the final submitted report while
preserving the regulations file in the \LaTeX{} project for reference.

} % End 13pt content

\clearpage

\end{verbatim}

For example:

\begin{verbatim}
% 
% =================================================================
% REGULATIONS FOR PREPARING THE GRADUATION INTERNSHIP REPORT
% USING LATEX
% =================================================================
%
% This file is intended to be included from the main document with:
%
%     
% =================================================================
% REGULATIONS FOR PREPARING THE GRADUATION INTERNSHIP REPORT
% USING LATEX
% =================================================================
%
% This file is intended to be included from the main document with:
%
%     \input{Regulations.tex}
%
% IMPORTANT:
% This section is provided for guidance only. Students shall remove
% or comment out the \input{Regulations.tex} command before submitting
% the final internship report.
%
% =================================================================

\clearpage
\thispagestyle{empty}

\begin{center}
    {\fontsize{16pt}{20pt}\selectfont\bfseries
    SPECIFIC REGULATIONS FOR PREPARING\\
    GRADUATION INTERNSHIP REPORTS USING \LaTeX}
\end{center}

\vspace{0.5cm}

{\fontsize{13pt}{17pt}\selectfont
\setstretch{1.3}
\justifying

\noindent
\textbf{Note:}
This section is provided for guidance only and is included in the main
document using the command \verb|\input{Regulations.tex}|.
Students shall remove or comment out this command before submitting the
final internship report.

\vspace{0.4cm}

% ================================================================
% I. GENERAL REQUIREMENTS
% ================================================================

\noindent
\textbf{I.\quad GENERAL REQUIREMENTS}

\vspace{0.15cm}

The Graduation Internship Report shall be prepared using the official
\LaTeX{} template provided by the Faculty. Students shall preserve the
predefined page layout, typography, line spacing, heading styles,
numbering system, figure and table formatting, and bibliography settings
unless otherwise instructed by the Faculty or the supervisor.

The report shall be written in English.

Students are responsible for the academic and technical content of the
report. The official \LaTeX{} template is responsible for standardized
document formatting.

\vspace{0.3cm}

% ================================================================
% II. REPORT STRUCTURE
% ================================================================

\noindent
\textbf{II.\quad REPORT STRUCTURE}

\vspace{0.15cm}

The report shall normally be arranged in the following order:

\begin{enumerate}[label=\arabic*., leftmargin=1.2cm, itemsep=2pt]
    \item Outer Cover;
    \item Comments of the Internship Host Organization;
    \item Supervisor's Evaluation and Graduation Internship Defense Results;
    \item Table of Contents;
    \item List of Symbols and Abbreviations, if applicable;
    \item List of Figures;
    \item List of Tables;
    \item Abstract;
    \item Main Chapters of the Internship Report;
    \item References;
    \item Appendices, if applicable.
\end{enumerate}

The preliminary pages shall not be numbered as chapters. The main report
content shall be organized into consecutively numbered chapters beginning
with Chapter 1.

\vspace{0.3cm}

% ================================================================
% III. PAGE SIZE AND MARGINS
% ================================================================

\noindent
\textbf{III.\quad PAGE SIZE AND MARGINS}

\vspace{0.15cm}

The report shall be formatted on A4 paper with the following margins:

\begin{center}
\begin{tabular}{ll}
    Left margin:   & 3.5 cm \\
    Right margin:  & 2.0 cm \\
    Top margin:    & 2.5 cm \\
    Bottom margin: & 2.5 cm
\end{tabular}
\end{center}

These settings are predefined in the official template and shall not be
modified by students.

\vspace{0.3cm}

% ================================================================
% IV. FONT, FONT SIZE, AND LINE SPACING
% ================================================================

\noindent
\textbf{IV.\quad FONT, FONT SIZE, AND LINE SPACING}

\vspace{0.15cm}

The main text shall use a 13 pt font size with 1.3-line spacing.
A Times-family typeface shall be used in accordance with the official
template.

The text shall be fully justified.

Students shall not use repeated spaces, repeated manual line breaks, or
multiple blank lines to position text. Paragraph and heading spacing shall
be controlled by the template.

\vspace{0.3cm}

% ================================================================
% V. OUTER COVER
% ================================================================

\noindent
\textbf{V.\quad OUTER COVER}

\vspace{0.15cm}

The outer cover shall follow the official format provided in the template
and shall contain the following information:

\begin{itemize}[leftmargin=0.8cm, itemsep=2pt]
    \item Thai Nguyen University of Technology;
    \item Faculty of International Training;
    \item Engineering and Technology Department;
    \item University logo and Faculty logo;
    \item the title \textbf{INTERNSHIP REPORT};
    \item Internship Organization;
    \item Author;
    \item Class;
    \item Supervisor;
    \item place and year of completion.
\end{itemize}

Students shall only replace the required information fields and shall not
modify the cover-page layout unless instructed otherwise.

\vspace{0.3cm}

% ================================================================
% VI. COMMENTS OF THE INTERNSHIP HOST ORGANIZATION
% ================================================================

\noindent
\textbf{VI.\quad COMMENTS OF THE INTERNSHIP HOST ORGANIZATION}

\vspace{0.15cm}

The page entitled \textbf{COMMENTS OF THE INTERNSHIP HOST ORGANIZATION}
is reserved for the internship host organization.

It includes:

\begin{itemize}[leftmargin=0.8cm, itemsep=2pt]
    \item the student's full name;
    \item class;
    \item comments from the host organization;
    \item score;
    \item date;
    \item signature and official seal.
\end{itemize}

Students shall not alter or complete official evaluation fields unless
specifically authorized.

\vspace{0.3cm}

% ================================================================
% VII. SUPERVISOR'S EVALUATION AND DEFENSE RESULTS
% ================================================================

\noindent
\textbf{VII.\quad SUPERVISOR'S EVALUATION AND DEFENSE RESULTS}

\vspace{0.15cm}

The section entitled
\textbf{SUPERVISOR'S EVALUATION AND GRADUATION INTERNSHIP DEFENSE RESULTS}
is reserved for the supervisor and defense committee.

The evaluation may include:

\begin{itemize}[leftmargin=0.8cm, itemsep=2pt]
    \item Internship Progress and Student's Attitude;
    \item Report Content;
    \item Report Presentation;
    \item Other Comments;
    \item Supervisor's Opinion;
    \item Defense Result;
    \item Score.
\end{itemize}

Students shall not enter evaluation scores or defense results before they
are officially determined.

\vspace{0.3cm}

% ================================================================
% VIII. CHAPTERS, SECTIONS, AND SUBSECTIONS
% ================================================================

\noindent
\textbf{VIII.\quad CHAPTERS, SECTIONS, AND SUBSECTIONS}

\vspace{0.15cm}

The report structure shall be generated using the standard \LaTeX{}
commands:

\begin{verbatim}
\chapter{...}
\section{...}
\subsection{...}
\end{verbatim}

Chapter, section, and subsection numbers shall be generated automatically
by \LaTeX{}.

Students shall not manually type chapter or section numbers and shall not
reproduce heading styles using commands such as \verb|\textbf|,
\verb|\fontsize|, or \verb|\centering| instead of the appropriate
structural command.

Excessive subdivision should be avoided. In most cases, the hierarchy
Chapter--Section--Subsection is sufficient.

\vspace{0.3cm}

% ================================================================
% IX. ORGANIZATION OF THE MAIN CONTENT
% ================================================================

\noindent
\textbf{IX.\quad ORGANIZATION OF THE MAIN CONTENT}

\vspace{0.15cm}

The specific content of the report shall depend on the internship
organization and the assigned tasks. A typical structure may include:

\begin{itemize}[leftmargin=0.8cm, itemsep=2pt]
    \item \textbf{Chapter 1. Overview of the Internship Organization};
    \item \textbf{Chapter 2. Internship Activities and Assigned Tasks};
    \item \textbf{Chapter 3. Results, Experience, and Evaluation}.
\end{itemize}

These chapter titles are recommendations only and may be adapted to the
actual internship content.

The report shall reflect work actually performed, observed, analyzed, or
learned during the internship. Unrelated theoretical material should not
be included merely to increase report length.

\vspace{0.3cm}

% ================================================================
% X. TABLE OF CONTENTS
% ================================================================

\noindent
\textbf{X.\quad TABLE OF CONTENTS}

\vspace{0.15cm}

The Table of Contents shall be generated automatically using:

\begin{verbatim}
\tableofcontents
\end{verbatim}

Students shall not create or manually edit page numbers in the Table of
Contents.

The document shall be compiled sufficiently to ensure that all page
numbers and entries are correctly updated.

\vspace{0.3cm}

% ================================================================
% XI. LIST OF SYMBOLS AND ABBREVIATIONS
% ================================================================

\noindent
\textbf{XI.\quad LIST OF SYMBOLS AND ABBREVIATIONS}

\vspace{0.15cm}

Symbols and abbreviations that are used repeatedly in the report shall be
listed in the \textbf{List of Symbols and Abbreviations}.

The list may contain the following columns:

\begin{center}
\begin{tabular}{llll}
    No. & Symbol / Abbreviation & Unit & Definition
\end{tabular}
\end{center}

Only symbols and abbreviations actually used in the final report shall be
retained. Example entries contained in the template shall be removed if
they are not relevant to the student's report.

Each abbreviation should normally be defined when it first appears in
the main text.

\vspace{0.3cm}

% ================================================================
% XII. FIGURES, DIAGRAMS, AND PHOTOGRAPHS
% ================================================================

\noindent
\textbf{XII.\quad FIGURES, DIAGRAMS, AND PHOTOGRAPHS}

\vspace{0.15cm}

Figures shall be clear, legible, and directly relevant to the report
content.

Figures shall be inserted using the standard \verb|figure| environment.

Example:

\begin{verbatim}
\begin{figure}[H]
    \centering
    \includegraphics[width=0.8\textwidth]
                    {figures/company_structure.pdf}
    \caption{Organizational structure of the company}
    \label{fig:company-structure}
\end{figure}
\end{verbatim}

Figure captions shall be placed below the figure.

Each referenced figure shall have a unique \verb|\label| and shall be
referenced in the text using, for example:

\begin{verbatim}
Figure~\ref{fig:company-structure}
\end{verbatim}

Students shall not manually type figure numbers.

Figures reproduced or adapted from external sources shall include an
appropriate citation.

Photographs taken during the internship shall be relevant and sufficiently
clear. Students shall not include confidential, restricted, or proprietary
information without permission from the internship organization.

Vector graphics or high-resolution images are recommended whenever
possible.

\vspace{0.3cm}

% ================================================================
% XIII. TABLES
% ================================================================

\noindent
\textbf{XIII.\quad TABLES}

\vspace{0.15cm}

Tables shall be created using suitable \LaTeX{} environments such as
\verb|table|, \verb|tabular|, \verb|tabularx|, or \verb|longtable|.

Table captions shall be placed above the table.

Example:

\begin{verbatim}
\begin{table}[H]
    \centering
    \caption{Internship schedule}
    \label{tab:internship-schedule}

    \begin{tabular}{lll}
        \toprule
        Week & Task & Result \\
        \midrule
        1 & Orientation & Completed \\
        2 & System study & Completed \\
        \bottomrule
    \end{tabular}
\end{table}
\end{verbatim}

Tables shall be referenced using:

\begin{verbatim}
Table~\ref{tab:internship-schedule}
\end{verbatim}

Students shall not manually type table numbers.

Tables should be created directly in \LaTeX{} rather than inserted as
screenshots whenever practical.

\vspace{0.3cm}

% ================================================================
% XIV. EQUATIONS AND TECHNICAL QUANTITIES
% ================================================================

\noindent
\textbf{XIV.\quad EQUATIONS AND TECHNICAL QUANTITIES}

\vspace{0.15cm}

Mathematical expressions shall be typeset directly using \LaTeX{}.
Images or screenshots of equations shall not be used.

Short expressions may be written in inline mathematical mode, for example:

\begin{verbatim}
$R = 0.7\,\Omega$
\end{verbatim}

Important equations shall use a numbered equation environment.

Example:

\begin{verbatim}
\begin{equation}
    P = UI
    \label{eq:power}
\end{equation}
\end{verbatim}

Equations shall be referenced using:

\begin{verbatim}
Equation~\eqref{eq:power}
\end{verbatim}

Equation numbers shall not be manually entered.

SI units shall be used wherever applicable. The \verb|siunitx| package
should be used where appropriate.

Examples:

\begin{verbatim}
\SI{220}{\volt}
\SI{50}{\hertz}
\SI{2.5}{\ampere}
\end{verbatim}

Notation and units shall be consistent throughout the report.

\vspace{0.3cm}

% ================================================================
% XV. SOURCE CODE
% ================================================================

\noindent
\textbf{XV.\quad SOURCE CODE}

\vspace{0.15cm}

If the internship involves programming, PLCs, microcontrollers,
simulation, or software development, relevant source code may be included
in the report.

Source code shall be presented using the \verb|listings| environment or
another environment defined by the official template.

Example:

\begin{verbatim}
\begin{lstlisting}[language=st]
IF Start = TRUE THEN
    Motor := TRUE;
ELSE
    Motor := FALSE;
END_IF;
\end{lstlisting}
\end{verbatim}

Only code necessary for explaining the work should appear in the main
report. Long source-code listings should be moved to an appendix or
submitted as supplementary files if required.

Source code should not normally be inserted as screenshots.

\vspace{0.3cm}

% ================================================================
% XVI. ABSTRACT
% ================================================================

\noindent
\textbf{XVI.\quad ABSTRACT}

\vspace{0.15cm}

The Abstract shall briefly summarize:

\begin{itemize}[leftmargin=0.8cm, itemsep=2pt]
    \item the internship organization;
    \item the objectives of the internship;
    \item the main tasks performed;
    \item the main results, knowledge, and skills obtained.
\end{itemize}

The Abstract should be concise and should not reproduce the full content
of the report chapter by chapter.

\vspace{0.3cm}

% ================================================================
% XVII. REFERENCES AND CITATIONS
% ================================================================

\noindent
\textbf{XVII.\quad REFERENCES AND CITATIONS}

\vspace{0.15cm}

All books, journal papers, conference papers, standards, websites,
datasheets, manuals, technical documents, figures, tables, or other
external sources used in preparing the report shall be properly cited.

References shall be managed through the bibliography system provided with
the official template.

Bibliographic entries shall normally be stored in:

\begin{verbatim}
References.bib
\end{verbatim}

References shall be cited in the report using:

\begin{verbatim}
\cite{reference-key}
\end{verbatim}

Students shall not manually type reference numbers such as ``[1]'',
``[2]'', etc.

Titles of foreign-language publications shall remain in their original
language in the bibliography.

All borrowed ideas, data, figures, tables, equations, methods, or results
shall be properly acknowledged.

\vspace{0.3cm}

% ================================================================
% XVIII. CROSS-REFERENCING
% ================================================================

\noindent
\textbf{XVIII.\quad CROSS-REFERENCING}

\vspace{0.15cm}

Numbered elements shall be cross-referenced using the standard \LaTeX{}
mechanism.

Use the following commands where appropriate:

\begin{verbatim}
\label{...}
\ref{...}
\eqref{...}
\end{verbatim}

For example:

\begin{verbatim}
As shown in Figure~\ref{fig:system}, ...
\end{verbatim}

Students shall not manually type numbers such as ``Figure 2.3'' when
those numbers can be generated automatically.

\vspace{0.3cm}

% ================================================================
% XIX. PAGE NUMBERING
% ================================================================

\noindent
\textbf{XIX.\quad PAGE NUMBERING}

\vspace{0.15cm}

Cover pages and official evaluation forms at the beginning of the report
shall not display page numbers.

The main body of the report shall use Arabic page numbers beginning from
Chapter 1.

Page numbering shall be generated automatically by the official template.
Students shall not manually enter page numbers.

\vspace{0.3cm}

% ================================================================
% XX. PROJECT FILE ORGANIZATION
% ================================================================

\noindent
\textbf{XX.\quad PROJECT FILE ORGANIZATION}

\vspace{0.15cm}

Students should maintain a clear and organized \LaTeX{} project
structure. A recommended structure is:

\begin{verbatim}
InternshipReport/
|
|-- main.tex
|-- References.bib
|
|-- figures/
|   |-- company.pdf
|   |-- system.png
|   `-- ...
|
`-- chapters/
    |-- chapter1.tex
    |-- chapter2.tex
    `-- chapter3.tex
\end{verbatim}

Figures should normally be stored in the \verb|figures| directory
specified by the official template.

File and directory names should preferably contain Latin letters,
numbers, hyphens, or underscores.

Recommended examples include:

\begin{verbatim}
company_structure.pdf
motor_control.png
chapter_01.tex
\end{verbatim}

Spaces and unnecessary special characters in file names should be
avoided where possible.

\vspace{0.3cm}

% ================================================================
% XXI. DOCUMENT COMPILATION AND FINAL CHECK
% ================================================================

\noindent
\textbf{XXI.\quad DOCUMENT COMPILATION AND FINAL CHECK}

\vspace{0.15cm}

Before submission, the complete \LaTeX{} project shall compile
successfully and produce the final PDF report.

The final document shall not contain unresolved references such as:

\begin{verbatim}
??
[?]
\end{verbatim}

Before generating the final PDF, students shall check:

\begin{itemize}[leftmargin=0.8cm, itemsep=2pt]
    \item the Table of Contents;
    \item page numbering;
    \item chapter, section, and subsection numbering;
    \item figure numbering;
    \item table numbering;
    \item equation numbering;
    \item cross-references;
    \item citations;
    \item bibliography;
    \item font display;
    \item mathematical symbols;
    \item image quality;
    \item completeness of all required forms.
\end{itemize}

The document may be compiled using TeX Live, MiKTeX, or another
compatible \LaTeX{} environment. Regardless of the compilation
environment, the resulting PDF shall conform to the official Faculty
template.

\vspace{0.3cm}

% ================================================================
% XXII. CONTENT THAT STUDENTS MAY EDIT
% ================================================================

\noindent
\textbf{XXII.\quad CONTENT THAT STUDENTS MAY EDIT}

\vspace{0.15cm}

Students may edit or replace:

\begin{itemize}[leftmargin=0.8cm, itemsep=2pt]
    \item student name;
    \item class;
    \item internship organization;
    \item supervisor name;
    \item Abstract;
    \item List of Symbols and Abbreviations;
    \item chapter titles and chapter content;
    \item figures and photographs;
    \item tables;
    \item equations;
    \item citations and references.
\end{itemize}

Students shall not modify the main document formatting settings in the
preamble, including the document class, page geometry, line spacing,
heading formatting, table-of-contents formatting, font configuration,
page layout, or cover-page macros unless specifically instructed.

\vspace{0.3cm}

% ================================================================
% XXIII. FINAL SUBMISSION REQUIREMENTS
% ================================================================

\noindent
\textbf{XXIII.\quad FINAL SUBMISSION REQUIREMENTS}

\vspace{0.15cm}

Before submission, the final report shall:

\begin{itemize}[leftmargin=0.8cm, itemsep=2pt]
    \item compile without errors that affect the report content or presentation;
    \item contain all required forms and report sections;
    \item contain correct student, class, internship organization, and supervisor information;
    \item contain an updated Table of Contents, List of Figures, and List of Tables;
    \item contain no unused example content from the template;
    \item contain no placeholder text such as \textit{Title of Chapter 1},
    \textit{Section 1.1}, or other sample entries;
    \item contain all required citations and references;
    \item contain no unresolved cross-references.
\end{itemize}

If the Faculty requires submission of the \LaTeX{} source, the complete
project shall include all \verb|.tex| files, the bibliography database,
figures, and any other files required to reproduce the submitted PDF.

\vspace{0.4cm}

\noindent
\textbf{Important:}
Before submitting the final internship report, students shall remove or
comment out the following command in the main document:

\begin{verbatim}
\input{Regulations.tex}
\end{verbatim}

For example:

\begin{verbatim}
% \input{Regulations.tex}
\end{verbatim}

This removes the guidance pages from the final submitted report while
preserving the regulations file in the \LaTeX{} project for reference.

} % End 13pt content

\clearpage

%
% IMPORTANT:
% This section is provided for guidance only. Students shall remove
% or comment out the 
% =================================================================
% REGULATIONS FOR PREPARING THE GRADUATION INTERNSHIP REPORT
% USING LATEX
% =================================================================
%
% This file is intended to be included from the main document with:
%
%     \input{Regulations.tex}
%
% IMPORTANT:
% This section is provided for guidance only. Students shall remove
% or comment out the \input{Regulations.tex} command before submitting
% the final internship report.
%
% =================================================================

\clearpage
\thispagestyle{empty}

\begin{center}
    {\fontsize{16pt}{20pt}\selectfont\bfseries
    SPECIFIC REGULATIONS FOR PREPARING\\
    GRADUATION INTERNSHIP REPORTS USING \LaTeX}
\end{center}

\vspace{0.5cm}

{\fontsize{13pt}{17pt}\selectfont
\setstretch{1.3}
\justifying

\noindent
\textbf{Note:}
This section is provided for guidance only and is included in the main
document using the command \verb|\input{Regulations.tex}|.
Students shall remove or comment out this command before submitting the
final internship report.

\vspace{0.4cm}

% ================================================================
% I. GENERAL REQUIREMENTS
% ================================================================

\noindent
\textbf{I.\quad GENERAL REQUIREMENTS}

\vspace{0.15cm}

The Graduation Internship Report shall be prepared using the official
\LaTeX{} template provided by the Faculty. Students shall preserve the
predefined page layout, typography, line spacing, heading styles,
numbering system, figure and table formatting, and bibliography settings
unless otherwise instructed by the Faculty or the supervisor.

The report shall be written in English.

Students are responsible for the academic and technical content of the
report. The official \LaTeX{} template is responsible for standardized
document formatting.

\vspace{0.3cm}

% ================================================================
% II. REPORT STRUCTURE
% ================================================================

\noindent
\textbf{II.\quad REPORT STRUCTURE}

\vspace{0.15cm}

The report shall normally be arranged in the following order:

\begin{enumerate}[label=\arabic*., leftmargin=1.2cm, itemsep=2pt]
    \item Outer Cover;
    \item Comments of the Internship Host Organization;
    \item Supervisor's Evaluation and Graduation Internship Defense Results;
    \item Table of Contents;
    \item List of Symbols and Abbreviations, if applicable;
    \item List of Figures;
    \item List of Tables;
    \item Abstract;
    \item Main Chapters of the Internship Report;
    \item References;
    \item Appendices, if applicable.
\end{enumerate}

The preliminary pages shall not be numbered as chapters. The main report
content shall be organized into consecutively numbered chapters beginning
with Chapter 1.

\vspace{0.3cm}

% ================================================================
% III. PAGE SIZE AND MARGINS
% ================================================================

\noindent
\textbf{III.\quad PAGE SIZE AND MARGINS}

\vspace{0.15cm}

The report shall be formatted on A4 paper with the following margins:

\begin{center}
\begin{tabular}{ll}
    Left margin:   & 3.5 cm \\
    Right margin:  & 2.0 cm \\
    Top margin:    & 2.5 cm \\
    Bottom margin: & 2.5 cm
\end{tabular}
\end{center}

These settings are predefined in the official template and shall not be
modified by students.

\vspace{0.3cm}

% ================================================================
% IV. FONT, FONT SIZE, AND LINE SPACING
% ================================================================

\noindent
\textbf{IV.\quad FONT, FONT SIZE, AND LINE SPACING}

\vspace{0.15cm}

The main text shall use a 13 pt font size with 1.3-line spacing.
A Times-family typeface shall be used in accordance with the official
template.

The text shall be fully justified.

Students shall not use repeated spaces, repeated manual line breaks, or
multiple blank lines to position text. Paragraph and heading spacing shall
be controlled by the template.

\vspace{0.3cm}

% ================================================================
% V. OUTER COVER
% ================================================================

\noindent
\textbf{V.\quad OUTER COVER}

\vspace{0.15cm}

The outer cover shall follow the official format provided in the template
and shall contain the following information:

\begin{itemize}[leftmargin=0.8cm, itemsep=2pt]
    \item Thai Nguyen University of Technology;
    \item Faculty of International Training;
    \item Engineering and Technology Department;
    \item University logo and Faculty logo;
    \item the title \textbf{INTERNSHIP REPORT};
    \item Internship Organization;
    \item Author;
    \item Class;
    \item Supervisor;
    \item place and year of completion.
\end{itemize}

Students shall only replace the required information fields and shall not
modify the cover-page layout unless instructed otherwise.

\vspace{0.3cm}

% ================================================================
% VI. COMMENTS OF THE INTERNSHIP HOST ORGANIZATION
% ================================================================

\noindent
\textbf{VI.\quad COMMENTS OF THE INTERNSHIP HOST ORGANIZATION}

\vspace{0.15cm}

The page entitled \textbf{COMMENTS OF THE INTERNSHIP HOST ORGANIZATION}
is reserved for the internship host organization.

It includes:

\begin{itemize}[leftmargin=0.8cm, itemsep=2pt]
    \item the student's full name;
    \item class;
    \item comments from the host organization;
    \item score;
    \item date;
    \item signature and official seal.
\end{itemize}

Students shall not alter or complete official evaluation fields unless
specifically authorized.

\vspace{0.3cm}

% ================================================================
% VII. SUPERVISOR'S EVALUATION AND DEFENSE RESULTS
% ================================================================

\noindent
\textbf{VII.\quad SUPERVISOR'S EVALUATION AND DEFENSE RESULTS}

\vspace{0.15cm}

The section entitled
\textbf{SUPERVISOR'S EVALUATION AND GRADUATION INTERNSHIP DEFENSE RESULTS}
is reserved for the supervisor and defense committee.

The evaluation may include:

\begin{itemize}[leftmargin=0.8cm, itemsep=2pt]
    \item Internship Progress and Student's Attitude;
    \item Report Content;
    \item Report Presentation;
    \item Other Comments;
    \item Supervisor's Opinion;
    \item Defense Result;
    \item Score.
\end{itemize}

Students shall not enter evaluation scores or defense results before they
are officially determined.

\vspace{0.3cm}

% ================================================================
% VIII. CHAPTERS, SECTIONS, AND SUBSECTIONS
% ================================================================

\noindent
\textbf{VIII.\quad CHAPTERS, SECTIONS, AND SUBSECTIONS}

\vspace{0.15cm}

The report structure shall be generated using the standard \LaTeX{}
commands:

\begin{verbatim}
\chapter{...}
\section{...}
\subsection{...}
\end{verbatim}

Chapter, section, and subsection numbers shall be generated automatically
by \LaTeX{}.

Students shall not manually type chapter or section numbers and shall not
reproduce heading styles using commands such as \verb|\textbf|,
\verb|\fontsize|, or \verb|\centering| instead of the appropriate
structural command.

Excessive subdivision should be avoided. In most cases, the hierarchy
Chapter--Section--Subsection is sufficient.

\vspace{0.3cm}

% ================================================================
% IX. ORGANIZATION OF THE MAIN CONTENT
% ================================================================

\noindent
\textbf{IX.\quad ORGANIZATION OF THE MAIN CONTENT}

\vspace{0.15cm}

The specific content of the report shall depend on the internship
organization and the assigned tasks. A typical structure may include:

\begin{itemize}[leftmargin=0.8cm, itemsep=2pt]
    \item \textbf{Chapter 1. Overview of the Internship Organization};
    \item \textbf{Chapter 2. Internship Activities and Assigned Tasks};
    \item \textbf{Chapter 3. Results, Experience, and Evaluation}.
\end{itemize}

These chapter titles are recommendations only and may be adapted to the
actual internship content.

The report shall reflect work actually performed, observed, analyzed, or
learned during the internship. Unrelated theoretical material should not
be included merely to increase report length.

\vspace{0.3cm}

% ================================================================
% X. TABLE OF CONTENTS
% ================================================================

\noindent
\textbf{X.\quad TABLE OF CONTENTS}

\vspace{0.15cm}

The Table of Contents shall be generated automatically using:

\begin{verbatim}
\tableofcontents
\end{verbatim}

Students shall not create or manually edit page numbers in the Table of
Contents.

The document shall be compiled sufficiently to ensure that all page
numbers and entries are correctly updated.

\vspace{0.3cm}

% ================================================================
% XI. LIST OF SYMBOLS AND ABBREVIATIONS
% ================================================================

\noindent
\textbf{XI.\quad LIST OF SYMBOLS AND ABBREVIATIONS}

\vspace{0.15cm}

Symbols and abbreviations that are used repeatedly in the report shall be
listed in the \textbf{List of Symbols and Abbreviations}.

The list may contain the following columns:

\begin{center}
\begin{tabular}{llll}
    No. & Symbol / Abbreviation & Unit & Definition
\end{tabular}
\end{center}

Only symbols and abbreviations actually used in the final report shall be
retained. Example entries contained in the template shall be removed if
they are not relevant to the student's report.

Each abbreviation should normally be defined when it first appears in
the main text.

\vspace{0.3cm}

% ================================================================
% XII. FIGURES, DIAGRAMS, AND PHOTOGRAPHS
% ================================================================

\noindent
\textbf{XII.\quad FIGURES, DIAGRAMS, AND PHOTOGRAPHS}

\vspace{0.15cm}

Figures shall be clear, legible, and directly relevant to the report
content.

Figures shall be inserted using the standard \verb|figure| environment.

Example:

\begin{verbatim}
\begin{figure}[H]
    \centering
    \includegraphics[width=0.8\textwidth]
                    {figures/company_structure.pdf}
    \caption{Organizational structure of the company}
    \label{fig:company-structure}
\end{figure}
\end{verbatim}

Figure captions shall be placed below the figure.

Each referenced figure shall have a unique \verb|\label| and shall be
referenced in the text using, for example:

\begin{verbatim}
Figure~\ref{fig:company-structure}
\end{verbatim}

Students shall not manually type figure numbers.

Figures reproduced or adapted from external sources shall include an
appropriate citation.

Photographs taken during the internship shall be relevant and sufficiently
clear. Students shall not include confidential, restricted, or proprietary
information without permission from the internship organization.

Vector graphics or high-resolution images are recommended whenever
possible.

\vspace{0.3cm}

% ================================================================
% XIII. TABLES
% ================================================================

\noindent
\textbf{XIII.\quad TABLES}

\vspace{0.15cm}

Tables shall be created using suitable \LaTeX{} environments such as
\verb|table|, \verb|tabular|, \verb|tabularx|, or \verb|longtable|.

Table captions shall be placed above the table.

Example:

\begin{verbatim}
\begin{table}[H]
    \centering
    \caption{Internship schedule}
    \label{tab:internship-schedule}

    \begin{tabular}{lll}
        \toprule
        Week & Task & Result \\
        \midrule
        1 & Orientation & Completed \\
        2 & System study & Completed \\
        \bottomrule
    \end{tabular}
\end{table}
\end{verbatim}

Tables shall be referenced using:

\begin{verbatim}
Table~\ref{tab:internship-schedule}
\end{verbatim}

Students shall not manually type table numbers.

Tables should be created directly in \LaTeX{} rather than inserted as
screenshots whenever practical.

\vspace{0.3cm}

% ================================================================
% XIV. EQUATIONS AND TECHNICAL QUANTITIES
% ================================================================

\noindent
\textbf{XIV.\quad EQUATIONS AND TECHNICAL QUANTITIES}

\vspace{0.15cm}

Mathematical expressions shall be typeset directly using \LaTeX{}.
Images or screenshots of equations shall not be used.

Short expressions may be written in inline mathematical mode, for example:

\begin{verbatim}
$R = 0.7\,\Omega$
\end{verbatim}

Important equations shall use a numbered equation environment.

Example:

\begin{verbatim}
\begin{equation}
    P = UI
    \label{eq:power}
\end{equation}
\end{verbatim}

Equations shall be referenced using:

\begin{verbatim}
Equation~\eqref{eq:power}
\end{verbatim}

Equation numbers shall not be manually entered.

SI units shall be used wherever applicable. The \verb|siunitx| package
should be used where appropriate.

Examples:

\begin{verbatim}
\SI{220}{\volt}
\SI{50}{\hertz}
\SI{2.5}{\ampere}
\end{verbatim}

Notation and units shall be consistent throughout the report.

\vspace{0.3cm}

% ================================================================
% XV. SOURCE CODE
% ================================================================

\noindent
\textbf{XV.\quad SOURCE CODE}

\vspace{0.15cm}

If the internship involves programming, PLCs, microcontrollers,
simulation, or software development, relevant source code may be included
in the report.

Source code shall be presented using the \verb|listings| environment or
another environment defined by the official template.

Example:

\begin{verbatim}
\begin{lstlisting}[language=st]
IF Start = TRUE THEN
    Motor := TRUE;
ELSE
    Motor := FALSE;
END_IF;
\end{lstlisting}
\end{verbatim}

Only code necessary for explaining the work should appear in the main
report. Long source-code listings should be moved to an appendix or
submitted as supplementary files if required.

Source code should not normally be inserted as screenshots.

\vspace{0.3cm}

% ================================================================
% XVI. ABSTRACT
% ================================================================

\noindent
\textbf{XVI.\quad ABSTRACT}

\vspace{0.15cm}

The Abstract shall briefly summarize:

\begin{itemize}[leftmargin=0.8cm, itemsep=2pt]
    \item the internship organization;
    \item the objectives of the internship;
    \item the main tasks performed;
    \item the main results, knowledge, and skills obtained.
\end{itemize}

The Abstract should be concise and should not reproduce the full content
of the report chapter by chapter.

\vspace{0.3cm}

% ================================================================
% XVII. REFERENCES AND CITATIONS
% ================================================================

\noindent
\textbf{XVII.\quad REFERENCES AND CITATIONS}

\vspace{0.15cm}

All books, journal papers, conference papers, standards, websites,
datasheets, manuals, technical documents, figures, tables, or other
external sources used in preparing the report shall be properly cited.

References shall be managed through the bibliography system provided with
the official template.

Bibliographic entries shall normally be stored in:

\begin{verbatim}
References.bib
\end{verbatim}

References shall be cited in the report using:

\begin{verbatim}
\cite{reference-key}
\end{verbatim}

Students shall not manually type reference numbers such as ``[1]'',
``[2]'', etc.

Titles of foreign-language publications shall remain in their original
language in the bibliography.

All borrowed ideas, data, figures, tables, equations, methods, or results
shall be properly acknowledged.

\vspace{0.3cm}

% ================================================================
% XVIII. CROSS-REFERENCING
% ================================================================

\noindent
\textbf{XVIII.\quad CROSS-REFERENCING}

\vspace{0.15cm}

Numbered elements shall be cross-referenced using the standard \LaTeX{}
mechanism.

Use the following commands where appropriate:

\begin{verbatim}
\label{...}
\ref{...}
\eqref{...}
\end{verbatim}

For example:

\begin{verbatim}
As shown in Figure~\ref{fig:system}, ...
\end{verbatim}

Students shall not manually type numbers such as ``Figure 2.3'' when
those numbers can be generated automatically.

\vspace{0.3cm}

% ================================================================
% XIX. PAGE NUMBERING
% ================================================================

\noindent
\textbf{XIX.\quad PAGE NUMBERING}

\vspace{0.15cm}

Cover pages and official evaluation forms at the beginning of the report
shall not display page numbers.

The main body of the report shall use Arabic page numbers beginning from
Chapter 1.

Page numbering shall be generated automatically by the official template.
Students shall not manually enter page numbers.

\vspace{0.3cm}

% ================================================================
% XX. PROJECT FILE ORGANIZATION
% ================================================================

\noindent
\textbf{XX.\quad PROJECT FILE ORGANIZATION}

\vspace{0.15cm}

Students should maintain a clear and organized \LaTeX{} project
structure. A recommended structure is:

\begin{verbatim}
InternshipReport/
|
|-- main.tex
|-- References.bib
|
|-- figures/
|   |-- company.pdf
|   |-- system.png
|   `-- ...
|
`-- chapters/
    |-- chapter1.tex
    |-- chapter2.tex
    `-- chapter3.tex
\end{verbatim}

Figures should normally be stored in the \verb|figures| directory
specified by the official template.

File and directory names should preferably contain Latin letters,
numbers, hyphens, or underscores.

Recommended examples include:

\begin{verbatim}
company_structure.pdf
motor_control.png
chapter_01.tex
\end{verbatim}

Spaces and unnecessary special characters in file names should be
avoided where possible.

\vspace{0.3cm}

% ================================================================
% XXI. DOCUMENT COMPILATION AND FINAL CHECK
% ================================================================

\noindent
\textbf{XXI.\quad DOCUMENT COMPILATION AND FINAL CHECK}

\vspace{0.15cm}

Before submission, the complete \LaTeX{} project shall compile
successfully and produce the final PDF report.

The final document shall not contain unresolved references such as:

\begin{verbatim}
??
[?]
\end{verbatim}

Before generating the final PDF, students shall check:

\begin{itemize}[leftmargin=0.8cm, itemsep=2pt]
    \item the Table of Contents;
    \item page numbering;
    \item chapter, section, and subsection numbering;
    \item figure numbering;
    \item table numbering;
    \item equation numbering;
    \item cross-references;
    \item citations;
    \item bibliography;
    \item font display;
    \item mathematical symbols;
    \item image quality;
    \item completeness of all required forms.
\end{itemize}

The document may be compiled using TeX Live, MiKTeX, or another
compatible \LaTeX{} environment. Regardless of the compilation
environment, the resulting PDF shall conform to the official Faculty
template.

\vspace{0.3cm}

% ================================================================
% XXII. CONTENT THAT STUDENTS MAY EDIT
% ================================================================

\noindent
\textbf{XXII.\quad CONTENT THAT STUDENTS MAY EDIT}

\vspace{0.15cm}

Students may edit or replace:

\begin{itemize}[leftmargin=0.8cm, itemsep=2pt]
    \item student name;
    \item class;
    \item internship organization;
    \item supervisor name;
    \item Abstract;
    \item List of Symbols and Abbreviations;
    \item chapter titles and chapter content;
    \item figures and photographs;
    \item tables;
    \item equations;
    \item citations and references.
\end{itemize}

Students shall not modify the main document formatting settings in the
preamble, including the document class, page geometry, line spacing,
heading formatting, table-of-contents formatting, font configuration,
page layout, or cover-page macros unless specifically instructed.

\vspace{0.3cm}

% ================================================================
% XXIII. FINAL SUBMISSION REQUIREMENTS
% ================================================================

\noindent
\textbf{XXIII.\quad FINAL SUBMISSION REQUIREMENTS}

\vspace{0.15cm}

Before submission, the final report shall:

\begin{itemize}[leftmargin=0.8cm, itemsep=2pt]
    \item compile without errors that affect the report content or presentation;
    \item contain all required forms and report sections;
    \item contain correct student, class, internship organization, and supervisor information;
    \item contain an updated Table of Contents, List of Figures, and List of Tables;
    \item contain no unused example content from the template;
    \item contain no placeholder text such as \textit{Title of Chapter 1},
    \textit{Section 1.1}, or other sample entries;
    \item contain all required citations and references;
    \item contain no unresolved cross-references.
\end{itemize}

If the Faculty requires submission of the \LaTeX{} source, the complete
project shall include all \verb|.tex| files, the bibliography database,
figures, and any other files required to reproduce the submitted PDF.

\vspace{0.4cm}

\noindent
\textbf{Important:}
Before submitting the final internship report, students shall remove or
comment out the following command in the main document:

\begin{verbatim}
\input{Regulations.tex}
\end{verbatim}

For example:

\begin{verbatim}
% \input{Regulations.tex}
\end{verbatim}

This removes the guidance pages from the final submitted report while
preserving the regulations file in the \LaTeX{} project for reference.

} % End 13pt content

\clearpage
 command before submitting
% the final internship report.
%
% =================================================================

\clearpage
\thispagestyle{empty}

\begin{center}
    {\fontsize{16pt}{20pt}\selectfont\bfseries
    SPECIFIC REGULATIONS FOR PREPARING\\
    GRADUATION INTERNSHIP REPORTS USING \LaTeX}
\end{center}

\vspace{0.5cm}

{\fontsize{13pt}{17pt}\selectfont
\setstretch{1.3}
\justifying

\noindent
\textbf{Note:}
This section is provided for guidance only and is included in the main
document using the command \verb|
% =================================================================
% REGULATIONS FOR PREPARING THE GRADUATION INTERNSHIP REPORT
% USING LATEX
% =================================================================
%
% This file is intended to be included from the main document with:
%
%     \input{Regulations.tex}
%
% IMPORTANT:
% This section is provided for guidance only. Students shall remove
% or comment out the \input{Regulations.tex} command before submitting
% the final internship report.
%
% =================================================================

\clearpage
\thispagestyle{empty}

\begin{center}
    {\fontsize{16pt}{20pt}\selectfont\bfseries
    SPECIFIC REGULATIONS FOR PREPARING\\
    GRADUATION INTERNSHIP REPORTS USING \LaTeX}
\end{center}

\vspace{0.5cm}

{\fontsize{13pt}{17pt}\selectfont
\setstretch{1.3}
\justifying

\noindent
\textbf{Note:}
This section is provided for guidance only and is included in the main
document using the command \verb|\input{Regulations.tex}|.
Students shall remove or comment out this command before submitting the
final internship report.

\vspace{0.4cm}

% ================================================================
% I. GENERAL REQUIREMENTS
% ================================================================

\noindent
\textbf{I.\quad GENERAL REQUIREMENTS}

\vspace{0.15cm}

The Graduation Internship Report shall be prepared using the official
\LaTeX{} template provided by the Faculty. Students shall preserve the
predefined page layout, typography, line spacing, heading styles,
numbering system, figure and table formatting, and bibliography settings
unless otherwise instructed by the Faculty or the supervisor.

The report shall be written in English.

Students are responsible for the academic and technical content of the
report. The official \LaTeX{} template is responsible for standardized
document formatting.

\vspace{0.3cm}

% ================================================================
% II. REPORT STRUCTURE
% ================================================================

\noindent
\textbf{II.\quad REPORT STRUCTURE}

\vspace{0.15cm}

The report shall normally be arranged in the following order:

\begin{enumerate}[label=\arabic*., leftmargin=1.2cm, itemsep=2pt]
    \item Outer Cover;
    \item Comments of the Internship Host Organization;
    \item Supervisor's Evaluation and Graduation Internship Defense Results;
    \item Table of Contents;
    \item List of Symbols and Abbreviations, if applicable;
    \item List of Figures;
    \item List of Tables;
    \item Abstract;
    \item Main Chapters of the Internship Report;
    \item References;
    \item Appendices, if applicable.
\end{enumerate}

The preliminary pages shall not be numbered as chapters. The main report
content shall be organized into consecutively numbered chapters beginning
with Chapter 1.

\vspace{0.3cm}

% ================================================================
% III. PAGE SIZE AND MARGINS
% ================================================================

\noindent
\textbf{III.\quad PAGE SIZE AND MARGINS}

\vspace{0.15cm}

The report shall be formatted on A4 paper with the following margins:

\begin{center}
\begin{tabular}{ll}
    Left margin:   & 3.5 cm \\
    Right margin:  & 2.0 cm \\
    Top margin:    & 2.5 cm \\
    Bottom margin: & 2.5 cm
\end{tabular}
\end{center}

These settings are predefined in the official template and shall not be
modified by students.

\vspace{0.3cm}

% ================================================================
% IV. FONT, FONT SIZE, AND LINE SPACING
% ================================================================

\noindent
\textbf{IV.\quad FONT, FONT SIZE, AND LINE SPACING}

\vspace{0.15cm}

The main text shall use a 13 pt font size with 1.3-line spacing.
A Times-family typeface shall be used in accordance with the official
template.

The text shall be fully justified.

Students shall not use repeated spaces, repeated manual line breaks, or
multiple blank lines to position text. Paragraph and heading spacing shall
be controlled by the template.

\vspace{0.3cm}

% ================================================================
% V. OUTER COVER
% ================================================================

\noindent
\textbf{V.\quad OUTER COVER}

\vspace{0.15cm}

The outer cover shall follow the official format provided in the template
and shall contain the following information:

\begin{itemize}[leftmargin=0.8cm, itemsep=2pt]
    \item Thai Nguyen University of Technology;
    \item Faculty of International Training;
    \item Engineering and Technology Department;
    \item University logo and Faculty logo;
    \item the title \textbf{INTERNSHIP REPORT};
    \item Internship Organization;
    \item Author;
    \item Class;
    \item Supervisor;
    \item place and year of completion.
\end{itemize}

Students shall only replace the required information fields and shall not
modify the cover-page layout unless instructed otherwise.

\vspace{0.3cm}

% ================================================================
% VI. COMMENTS OF THE INTERNSHIP HOST ORGANIZATION
% ================================================================

\noindent
\textbf{VI.\quad COMMENTS OF THE INTERNSHIP HOST ORGANIZATION}

\vspace{0.15cm}

The page entitled \textbf{COMMENTS OF THE INTERNSHIP HOST ORGANIZATION}
is reserved for the internship host organization.

It includes:

\begin{itemize}[leftmargin=0.8cm, itemsep=2pt]
    \item the student's full name;
    \item class;
    \item comments from the host organization;
    \item score;
    \item date;
    \item signature and official seal.
\end{itemize}

Students shall not alter or complete official evaluation fields unless
specifically authorized.

\vspace{0.3cm}

% ================================================================
% VII. SUPERVISOR'S EVALUATION AND DEFENSE RESULTS
% ================================================================

\noindent
\textbf{VII.\quad SUPERVISOR'S EVALUATION AND DEFENSE RESULTS}

\vspace{0.15cm}

The section entitled
\textbf{SUPERVISOR'S EVALUATION AND GRADUATION INTERNSHIP DEFENSE RESULTS}
is reserved for the supervisor and defense committee.

The evaluation may include:

\begin{itemize}[leftmargin=0.8cm, itemsep=2pt]
    \item Internship Progress and Student's Attitude;
    \item Report Content;
    \item Report Presentation;
    \item Other Comments;
    \item Supervisor's Opinion;
    \item Defense Result;
    \item Score.
\end{itemize}

Students shall not enter evaluation scores or defense results before they
are officially determined.

\vspace{0.3cm}

% ================================================================
% VIII. CHAPTERS, SECTIONS, AND SUBSECTIONS
% ================================================================

\noindent
\textbf{VIII.\quad CHAPTERS, SECTIONS, AND SUBSECTIONS}

\vspace{0.15cm}

The report structure shall be generated using the standard \LaTeX{}
commands:

\begin{verbatim}
\chapter{...}
\section{...}
\subsection{...}
\end{verbatim}

Chapter, section, and subsection numbers shall be generated automatically
by \LaTeX{}.

Students shall not manually type chapter or section numbers and shall not
reproduce heading styles using commands such as \verb|\textbf|,
\verb|\fontsize|, or \verb|\centering| instead of the appropriate
structural command.

Excessive subdivision should be avoided. In most cases, the hierarchy
Chapter--Section--Subsection is sufficient.

\vspace{0.3cm}

% ================================================================
% IX. ORGANIZATION OF THE MAIN CONTENT
% ================================================================

\noindent
\textbf{IX.\quad ORGANIZATION OF THE MAIN CONTENT}

\vspace{0.15cm}

The specific content of the report shall depend on the internship
organization and the assigned tasks. A typical structure may include:

\begin{itemize}[leftmargin=0.8cm, itemsep=2pt]
    \item \textbf{Chapter 1. Overview of the Internship Organization};
    \item \textbf{Chapter 2. Internship Activities and Assigned Tasks};
    \item \textbf{Chapter 3. Results, Experience, and Evaluation}.
\end{itemize}

These chapter titles are recommendations only and may be adapted to the
actual internship content.

The report shall reflect work actually performed, observed, analyzed, or
learned during the internship. Unrelated theoretical material should not
be included merely to increase report length.

\vspace{0.3cm}

% ================================================================
% X. TABLE OF CONTENTS
% ================================================================

\noindent
\textbf{X.\quad TABLE OF CONTENTS}

\vspace{0.15cm}

The Table of Contents shall be generated automatically using:

\begin{verbatim}
\tableofcontents
\end{verbatim}

Students shall not create or manually edit page numbers in the Table of
Contents.

The document shall be compiled sufficiently to ensure that all page
numbers and entries are correctly updated.

\vspace{0.3cm}

% ================================================================
% XI. LIST OF SYMBOLS AND ABBREVIATIONS
% ================================================================

\noindent
\textbf{XI.\quad LIST OF SYMBOLS AND ABBREVIATIONS}

\vspace{0.15cm}

Symbols and abbreviations that are used repeatedly in the report shall be
listed in the \textbf{List of Symbols and Abbreviations}.

The list may contain the following columns:

\begin{center}
\begin{tabular}{llll}
    No. & Symbol / Abbreviation & Unit & Definition
\end{tabular}
\end{center}

Only symbols and abbreviations actually used in the final report shall be
retained. Example entries contained in the template shall be removed if
they are not relevant to the student's report.

Each abbreviation should normally be defined when it first appears in
the main text.

\vspace{0.3cm}

% ================================================================
% XII. FIGURES, DIAGRAMS, AND PHOTOGRAPHS
% ================================================================

\noindent
\textbf{XII.\quad FIGURES, DIAGRAMS, AND PHOTOGRAPHS}

\vspace{0.15cm}

Figures shall be clear, legible, and directly relevant to the report
content.

Figures shall be inserted using the standard \verb|figure| environment.

Example:

\begin{verbatim}
\begin{figure}[H]
    \centering
    \includegraphics[width=0.8\textwidth]
                    {figures/company_structure.pdf}
    \caption{Organizational structure of the company}
    \label{fig:company-structure}
\end{figure}
\end{verbatim}

Figure captions shall be placed below the figure.

Each referenced figure shall have a unique \verb|\label| and shall be
referenced in the text using, for example:

\begin{verbatim}
Figure~\ref{fig:company-structure}
\end{verbatim}

Students shall not manually type figure numbers.

Figures reproduced or adapted from external sources shall include an
appropriate citation.

Photographs taken during the internship shall be relevant and sufficiently
clear. Students shall not include confidential, restricted, or proprietary
information without permission from the internship organization.

Vector graphics or high-resolution images are recommended whenever
possible.

\vspace{0.3cm}

% ================================================================
% XIII. TABLES
% ================================================================

\noindent
\textbf{XIII.\quad TABLES}

\vspace{0.15cm}

Tables shall be created using suitable \LaTeX{} environments such as
\verb|table|, \verb|tabular|, \verb|tabularx|, or \verb|longtable|.

Table captions shall be placed above the table.

Example:

\begin{verbatim}
\begin{table}[H]
    \centering
    \caption{Internship schedule}
    \label{tab:internship-schedule}

    \begin{tabular}{lll}
        \toprule
        Week & Task & Result \\
        \midrule
        1 & Orientation & Completed \\
        2 & System study & Completed \\
        \bottomrule
    \end{tabular}
\end{table}
\end{verbatim}

Tables shall be referenced using:

\begin{verbatim}
Table~\ref{tab:internship-schedule}
\end{verbatim}

Students shall not manually type table numbers.

Tables should be created directly in \LaTeX{} rather than inserted as
screenshots whenever practical.

\vspace{0.3cm}

% ================================================================
% XIV. EQUATIONS AND TECHNICAL QUANTITIES
% ================================================================

\noindent
\textbf{XIV.\quad EQUATIONS AND TECHNICAL QUANTITIES}

\vspace{0.15cm}

Mathematical expressions shall be typeset directly using \LaTeX{}.
Images or screenshots of equations shall not be used.

Short expressions may be written in inline mathematical mode, for example:

\begin{verbatim}
$R = 0.7\,\Omega$
\end{verbatim}

Important equations shall use a numbered equation environment.

Example:

\begin{verbatim}
\begin{equation}
    P = UI
    \label{eq:power}
\end{equation}
\end{verbatim}

Equations shall be referenced using:

\begin{verbatim}
Equation~\eqref{eq:power}
\end{verbatim}

Equation numbers shall not be manually entered.

SI units shall be used wherever applicable. The \verb|siunitx| package
should be used where appropriate.

Examples:

\begin{verbatim}
\SI{220}{\volt}
\SI{50}{\hertz}
\SI{2.5}{\ampere}
\end{verbatim}

Notation and units shall be consistent throughout the report.

\vspace{0.3cm}

% ================================================================
% XV. SOURCE CODE
% ================================================================

\noindent
\textbf{XV.\quad SOURCE CODE}

\vspace{0.15cm}

If the internship involves programming, PLCs, microcontrollers,
simulation, or software development, relevant source code may be included
in the report.

Source code shall be presented using the \verb|listings| environment or
another environment defined by the official template.

Example:

\begin{verbatim}
\begin{lstlisting}[language=st]
IF Start = TRUE THEN
    Motor := TRUE;
ELSE
    Motor := FALSE;
END_IF;
\end{lstlisting}
\end{verbatim}

Only code necessary for explaining the work should appear in the main
report. Long source-code listings should be moved to an appendix or
submitted as supplementary files if required.

Source code should not normally be inserted as screenshots.

\vspace{0.3cm}

% ================================================================
% XVI. ABSTRACT
% ================================================================

\noindent
\textbf{XVI.\quad ABSTRACT}

\vspace{0.15cm}

The Abstract shall briefly summarize:

\begin{itemize}[leftmargin=0.8cm, itemsep=2pt]
    \item the internship organization;
    \item the objectives of the internship;
    \item the main tasks performed;
    \item the main results, knowledge, and skills obtained.
\end{itemize}

The Abstract should be concise and should not reproduce the full content
of the report chapter by chapter.

\vspace{0.3cm}

% ================================================================
% XVII. REFERENCES AND CITATIONS
% ================================================================

\noindent
\textbf{XVII.\quad REFERENCES AND CITATIONS}

\vspace{0.15cm}

All books, journal papers, conference papers, standards, websites,
datasheets, manuals, technical documents, figures, tables, or other
external sources used in preparing the report shall be properly cited.

References shall be managed through the bibliography system provided with
the official template.

Bibliographic entries shall normally be stored in:

\begin{verbatim}
References.bib
\end{verbatim}

References shall be cited in the report using:

\begin{verbatim}
\cite{reference-key}
\end{verbatim}

Students shall not manually type reference numbers such as ``[1]'',
``[2]'', etc.

Titles of foreign-language publications shall remain in their original
language in the bibliography.

All borrowed ideas, data, figures, tables, equations, methods, or results
shall be properly acknowledged.

\vspace{0.3cm}

% ================================================================
% XVIII. CROSS-REFERENCING
% ================================================================

\noindent
\textbf{XVIII.\quad CROSS-REFERENCING}

\vspace{0.15cm}

Numbered elements shall be cross-referenced using the standard \LaTeX{}
mechanism.

Use the following commands where appropriate:

\begin{verbatim}
\label{...}
\ref{...}
\eqref{...}
\end{verbatim}

For example:

\begin{verbatim}
As shown in Figure~\ref{fig:system}, ...
\end{verbatim}

Students shall not manually type numbers such as ``Figure 2.3'' when
those numbers can be generated automatically.

\vspace{0.3cm}

% ================================================================
% XIX. PAGE NUMBERING
% ================================================================

\noindent
\textbf{XIX.\quad PAGE NUMBERING}

\vspace{0.15cm}

Cover pages and official evaluation forms at the beginning of the report
shall not display page numbers.

The main body of the report shall use Arabic page numbers beginning from
Chapter 1.

Page numbering shall be generated automatically by the official template.
Students shall not manually enter page numbers.

\vspace{0.3cm}

% ================================================================
% XX. PROJECT FILE ORGANIZATION
% ================================================================

\noindent
\textbf{XX.\quad PROJECT FILE ORGANIZATION}

\vspace{0.15cm}

Students should maintain a clear and organized \LaTeX{} project
structure. A recommended structure is:

\begin{verbatim}
InternshipReport/
|
|-- main.tex
|-- References.bib
|
|-- figures/
|   |-- company.pdf
|   |-- system.png
|   `-- ...
|
`-- chapters/
    |-- chapter1.tex
    |-- chapter2.tex
    `-- chapter3.tex
\end{verbatim}

Figures should normally be stored in the \verb|figures| directory
specified by the official template.

File and directory names should preferably contain Latin letters,
numbers, hyphens, or underscores.

Recommended examples include:

\begin{verbatim}
company_structure.pdf
motor_control.png
chapter_01.tex
\end{verbatim}

Spaces and unnecessary special characters in file names should be
avoided where possible.

\vspace{0.3cm}

% ================================================================
% XXI. DOCUMENT COMPILATION AND FINAL CHECK
% ================================================================

\noindent
\textbf{XXI.\quad DOCUMENT COMPILATION AND FINAL CHECK}

\vspace{0.15cm}

Before submission, the complete \LaTeX{} project shall compile
successfully and produce the final PDF report.

The final document shall not contain unresolved references such as:

\begin{verbatim}
??
[?]
\end{verbatim}

Before generating the final PDF, students shall check:

\begin{itemize}[leftmargin=0.8cm, itemsep=2pt]
    \item the Table of Contents;
    \item page numbering;
    \item chapter, section, and subsection numbering;
    \item figure numbering;
    \item table numbering;
    \item equation numbering;
    \item cross-references;
    \item citations;
    \item bibliography;
    \item font display;
    \item mathematical symbols;
    \item image quality;
    \item completeness of all required forms.
\end{itemize}

The document may be compiled using TeX Live, MiKTeX, or another
compatible \LaTeX{} environment. Regardless of the compilation
environment, the resulting PDF shall conform to the official Faculty
template.

\vspace{0.3cm}

% ================================================================
% XXII. CONTENT THAT STUDENTS MAY EDIT
% ================================================================

\noindent
\textbf{XXII.\quad CONTENT THAT STUDENTS MAY EDIT}

\vspace{0.15cm}

Students may edit or replace:

\begin{itemize}[leftmargin=0.8cm, itemsep=2pt]
    \item student name;
    \item class;
    \item internship organization;
    \item supervisor name;
    \item Abstract;
    \item List of Symbols and Abbreviations;
    \item chapter titles and chapter content;
    \item figures and photographs;
    \item tables;
    \item equations;
    \item citations and references.
\end{itemize}

Students shall not modify the main document formatting settings in the
preamble, including the document class, page geometry, line spacing,
heading formatting, table-of-contents formatting, font configuration,
page layout, or cover-page macros unless specifically instructed.

\vspace{0.3cm}

% ================================================================
% XXIII. FINAL SUBMISSION REQUIREMENTS
% ================================================================

\noindent
\textbf{XXIII.\quad FINAL SUBMISSION REQUIREMENTS}

\vspace{0.15cm}

Before submission, the final report shall:

\begin{itemize}[leftmargin=0.8cm, itemsep=2pt]
    \item compile without errors that affect the report content or presentation;
    \item contain all required forms and report sections;
    \item contain correct student, class, internship organization, and supervisor information;
    \item contain an updated Table of Contents, List of Figures, and List of Tables;
    \item contain no unused example content from the template;
    \item contain no placeholder text such as \textit{Title of Chapter 1},
    \textit{Section 1.1}, or other sample entries;
    \item contain all required citations and references;
    \item contain no unresolved cross-references.
\end{itemize}

If the Faculty requires submission of the \LaTeX{} source, the complete
project shall include all \verb|.tex| files, the bibliography database,
figures, and any other files required to reproduce the submitted PDF.

\vspace{0.4cm}

\noindent
\textbf{Important:}
Before submitting the final internship report, students shall remove or
comment out the following command in the main document:

\begin{verbatim}
\input{Regulations.tex}
\end{verbatim}

For example:

\begin{verbatim}
% \input{Regulations.tex}
\end{verbatim}

This removes the guidance pages from the final submitted report while
preserving the regulations file in the \LaTeX{} project for reference.

} % End 13pt content

\clearpage
|.
Students shall remove or comment out this command before submitting the
final internship report.

\vspace{0.4cm}

% ================================================================
% I. GENERAL REQUIREMENTS
% ================================================================

\noindent
\textbf{I.\quad GENERAL REQUIREMENTS}

\vspace{0.15cm}

The Graduation Internship Report shall be prepared using the official
\LaTeX{} template provided by the Faculty. Students shall preserve the
predefined page layout, typography, line spacing, heading styles,
numbering system, figure and table formatting, and bibliography settings
unless otherwise instructed by the Faculty or the supervisor.

The report shall be written in English.

Students are responsible for the academic and technical content of the
report. The official \LaTeX{} template is responsible for standardized
document formatting.

\vspace{0.3cm}

% ================================================================
% II. REPORT STRUCTURE
% ================================================================

\noindent
\textbf{II.\quad REPORT STRUCTURE}

\vspace{0.15cm}

The report shall normally be arranged in the following order:

\begin{enumerate}[label=\arabic*., leftmargin=1.2cm, itemsep=2pt]
    \item Outer Cover;
    \item Comments of the Internship Host Organization;
    \item Supervisor's Evaluation and Graduation Internship Defense Results;
    \item Table of Contents;
    \item List of Symbols and Abbreviations, if applicable;
    \item List of Figures;
    \item List of Tables;
    \item Abstract;
    \item Main Chapters of the Internship Report;
    \item References;
    \item Appendices, if applicable.
\end{enumerate}

The preliminary pages shall not be numbered as chapters. The main report
content shall be organized into consecutively numbered chapters beginning
with Chapter 1.

\vspace{0.3cm}

% ================================================================
% III. PAGE SIZE AND MARGINS
% ================================================================

\noindent
\textbf{III.\quad PAGE SIZE AND MARGINS}

\vspace{0.15cm}

The report shall be formatted on A4 paper with the following margins:

\begin{center}
\begin{tabular}{ll}
    Left margin:   & 3.5 cm \\
    Right margin:  & 2.0 cm \\
    Top margin:    & 2.5 cm \\
    Bottom margin: & 2.5 cm
\end{tabular}
\end{center}

These settings are predefined in the official template and shall not be
modified by students.

\vspace{0.3cm}

% ================================================================
% IV. FONT, FONT SIZE, AND LINE SPACING
% ================================================================

\noindent
\textbf{IV.\quad FONT, FONT SIZE, AND LINE SPACING}

\vspace{0.15cm}

The main text shall use a 13 pt font size with 1.3-line spacing.
A Times-family typeface shall be used in accordance with the official
template.

The text shall be fully justified.

Students shall not use repeated spaces, repeated manual line breaks, or
multiple blank lines to position text. Paragraph and heading spacing shall
be controlled by the template.

\vspace{0.3cm}

% ================================================================
% V. OUTER COVER
% ================================================================

\noindent
\textbf{V.\quad OUTER COVER}

\vspace{0.15cm}

The outer cover shall follow the official format provided in the template
and shall contain the following information:

\begin{itemize}[leftmargin=0.8cm, itemsep=2pt]
    \item Thai Nguyen University of Technology;
    \item Faculty of International Training;
    \item Engineering and Technology Department;
    \item University logo and Faculty logo;
    \item the title \textbf{INTERNSHIP REPORT};
    \item Internship Organization;
    \item Author;
    \item Class;
    \item Supervisor;
    \item place and year of completion.
\end{itemize}

Students shall only replace the required information fields and shall not
modify the cover-page layout unless instructed otherwise.

\vspace{0.3cm}

% ================================================================
% VI. COMMENTS OF THE INTERNSHIP HOST ORGANIZATION
% ================================================================

\noindent
\textbf{VI.\quad COMMENTS OF THE INTERNSHIP HOST ORGANIZATION}

\vspace{0.15cm}

The page entitled \textbf{COMMENTS OF THE INTERNSHIP HOST ORGANIZATION}
is reserved for the internship host organization.

It includes:

\begin{itemize}[leftmargin=0.8cm, itemsep=2pt]
    \item the student's full name;
    \item class;
    \item comments from the host organization;
    \item score;
    \item date;
    \item signature and official seal.
\end{itemize}

Students shall not alter or complete official evaluation fields unless
specifically authorized.

\vspace{0.3cm}

% ================================================================
% VII. SUPERVISOR'S EVALUATION AND DEFENSE RESULTS
% ================================================================

\noindent
\textbf{VII.\quad SUPERVISOR'S EVALUATION AND DEFENSE RESULTS}

\vspace{0.15cm}

The section entitled
\textbf{SUPERVISOR'S EVALUATION AND GRADUATION INTERNSHIP DEFENSE RESULTS}
is reserved for the supervisor and defense committee.

The evaluation may include:

\begin{itemize}[leftmargin=0.8cm, itemsep=2pt]
    \item Internship Progress and Student's Attitude;
    \item Report Content;
    \item Report Presentation;
    \item Other Comments;
    \item Supervisor's Opinion;
    \item Defense Result;
    \item Score.
\end{itemize}

Students shall not enter evaluation scores or defense results before they
are officially determined.

\vspace{0.3cm}

% ================================================================
% VIII. CHAPTERS, SECTIONS, AND SUBSECTIONS
% ================================================================

\noindent
\textbf{VIII.\quad CHAPTERS, SECTIONS, AND SUBSECTIONS}

\vspace{0.15cm}

The report structure shall be generated using the standard \LaTeX{}
commands:

\begin{verbatim}
\chapter{...}
\section{...}
\subsection{...}
\end{verbatim}

Chapter, section, and subsection numbers shall be generated automatically
by \LaTeX{}.

Students shall not manually type chapter or section numbers and shall not
reproduce heading styles using commands such as \verb|\textbf|,
\verb|\fontsize|, or \verb|\centering| instead of the appropriate
structural command.

Excessive subdivision should be avoided. In most cases, the hierarchy
Chapter--Section--Subsection is sufficient.

\vspace{0.3cm}

% ================================================================
% IX. ORGANIZATION OF THE MAIN CONTENT
% ================================================================

\noindent
\textbf{IX.\quad ORGANIZATION OF THE MAIN CONTENT}

\vspace{0.15cm}

The specific content of the report shall depend on the internship
organization and the assigned tasks. A typical structure may include:

\begin{itemize}[leftmargin=0.8cm, itemsep=2pt]
    \item \textbf{Chapter 1. Overview of the Internship Organization};
    \item \textbf{Chapter 2. Internship Activities and Assigned Tasks};
    \item \textbf{Chapter 3. Results, Experience, and Evaluation}.
\end{itemize}

These chapter titles are recommendations only and may be adapted to the
actual internship content.

The report shall reflect work actually performed, observed, analyzed, or
learned during the internship. Unrelated theoretical material should not
be included merely to increase report length.

\vspace{0.3cm}

% ================================================================
% X. TABLE OF CONTENTS
% ================================================================

\noindent
\textbf{X.\quad TABLE OF CONTENTS}

\vspace{0.15cm}

The Table of Contents shall be generated automatically using:

\begin{verbatim}
\tableofcontents
\end{verbatim}

Students shall not create or manually edit page numbers in the Table of
Contents.

The document shall be compiled sufficiently to ensure that all page
numbers and entries are correctly updated.

\vspace{0.3cm}

% ================================================================
% XI. LIST OF SYMBOLS AND ABBREVIATIONS
% ================================================================

\noindent
\textbf{XI.\quad LIST OF SYMBOLS AND ABBREVIATIONS}

\vspace{0.15cm}

Symbols and abbreviations that are used repeatedly in the report shall be
listed in the \textbf{List of Symbols and Abbreviations}.

The list may contain the following columns:

\begin{center}
\begin{tabular}{llll}
    No. & Symbol / Abbreviation & Unit & Definition
\end{tabular}
\end{center}

Only symbols and abbreviations actually used in the final report shall be
retained. Example entries contained in the template shall be removed if
they are not relevant to the student's report.

Each abbreviation should normally be defined when it first appears in
the main text.

\vspace{0.3cm}

% ================================================================
% XII. FIGURES, DIAGRAMS, AND PHOTOGRAPHS
% ================================================================

\noindent
\textbf{XII.\quad FIGURES, DIAGRAMS, AND PHOTOGRAPHS}

\vspace{0.15cm}

Figures shall be clear, legible, and directly relevant to the report
content.

Figures shall be inserted using the standard \verb|figure| environment.

Example:

\begin{verbatim}
\begin{figure}[H]
    \centering
    \includegraphics[width=0.8\textwidth]
                    {figures/company_structure.pdf}
    \caption{Organizational structure of the company}
    \label{fig:company-structure}
\end{figure}
\end{verbatim}

Figure captions shall be placed below the figure.

Each referenced figure shall have a unique \verb|\label| and shall be
referenced in the text using, for example:

\begin{verbatim}
Figure~\ref{fig:company-structure}
\end{verbatim}

Students shall not manually type figure numbers.

Figures reproduced or adapted from external sources shall include an
appropriate citation.

Photographs taken during the internship shall be relevant and sufficiently
clear. Students shall not include confidential, restricted, or proprietary
information without permission from the internship organization.

Vector graphics or high-resolution images are recommended whenever
possible.

\vspace{0.3cm}

% ================================================================
% XIII. TABLES
% ================================================================

\noindent
\textbf{XIII.\quad TABLES}

\vspace{0.15cm}

Tables shall be created using suitable \LaTeX{} environments such as
\verb|table|, \verb|tabular|, \verb|tabularx|, or \verb|longtable|.

Table captions shall be placed above the table.

Example:

\begin{verbatim}
\begin{table}[H]
    \centering
    \caption{Internship schedule}
    \label{tab:internship-schedule}

    \begin{tabular}{lll}
        \toprule
        Week & Task & Result \\
        \midrule
        1 & Orientation & Completed \\
        2 & System study & Completed \\
        \bottomrule
    \end{tabular}
\end{table}
\end{verbatim}

Tables shall be referenced using:

\begin{verbatim}
Table~\ref{tab:internship-schedule}
\end{verbatim}

Students shall not manually type table numbers.

Tables should be created directly in \LaTeX{} rather than inserted as
screenshots whenever practical.

\vspace{0.3cm}

% ================================================================
% XIV. EQUATIONS AND TECHNICAL QUANTITIES
% ================================================================

\noindent
\textbf{XIV.\quad EQUATIONS AND TECHNICAL QUANTITIES}

\vspace{0.15cm}

Mathematical expressions shall be typeset directly using \LaTeX{}.
Images or screenshots of equations shall not be used.

Short expressions may be written in inline mathematical mode, for example:

\begin{verbatim}
$R = 0.7\,\Omega$
\end{verbatim}

Important equations shall use a numbered equation environment.

Example:

\begin{verbatim}
\begin{equation}
    P = UI
    \label{eq:power}
\end{equation}
\end{verbatim}

Equations shall be referenced using:

\begin{verbatim}
Equation~\eqref{eq:power}
\end{verbatim}

Equation numbers shall not be manually entered.

SI units shall be used wherever applicable. The \verb|siunitx| package
should be used where appropriate.

Examples:

\begin{verbatim}
\SI{220}{\volt}
\SI{50}{\hertz}
\SI{2.5}{\ampere}
\end{verbatim}

Notation and units shall be consistent throughout the report.

\vspace{0.3cm}

% ================================================================
% XV. SOURCE CODE
% ================================================================

\noindent
\textbf{XV.\quad SOURCE CODE}

\vspace{0.15cm}

If the internship involves programming, PLCs, microcontrollers,
simulation, or software development, relevant source code may be included
in the report.

Source code shall be presented using the \verb|listings| environment or
another environment defined by the official template.

Example:

\begin{verbatim}
\begin{lstlisting}[language=st]
IF Start = TRUE THEN
    Motor := TRUE;
ELSE
    Motor := FALSE;
END_IF;
\end{lstlisting}
\end{verbatim}

Only code necessary for explaining the work should appear in the main
report. Long source-code listings should be moved to an appendix or
submitted as supplementary files if required.

Source code should not normally be inserted as screenshots.

\vspace{0.3cm}

% ================================================================
% XVI. ABSTRACT
% ================================================================

\noindent
\textbf{XVI.\quad ABSTRACT}

\vspace{0.15cm}

The Abstract shall briefly summarize:

\begin{itemize}[leftmargin=0.8cm, itemsep=2pt]
    \item the internship organization;
    \item the objectives of the internship;
    \item the main tasks performed;
    \item the main results, knowledge, and skills obtained.
\end{itemize}

The Abstract should be concise and should not reproduce the full content
of the report chapter by chapter.

\vspace{0.3cm}

% ================================================================
% XVII. REFERENCES AND CITATIONS
% ================================================================

\noindent
\textbf{XVII.\quad REFERENCES AND CITATIONS}

\vspace{0.15cm}

All books, journal papers, conference papers, standards, websites,
datasheets, manuals, technical documents, figures, tables, or other
external sources used in preparing the report shall be properly cited.

References shall be managed through the bibliography system provided with
the official template.

Bibliographic entries shall normally be stored in:

\begin{verbatim}
References.bib
\end{verbatim}

References shall be cited in the report using:

\begin{verbatim}
\cite{reference-key}
\end{verbatim}

Students shall not manually type reference numbers such as ``[1]'',
``[2]'', etc.

Titles of foreign-language publications shall remain in their original
language in the bibliography.

All borrowed ideas, data, figures, tables, equations, methods, or results
shall be properly acknowledged.

\vspace{0.3cm}

% ================================================================
% XVIII. CROSS-REFERENCING
% ================================================================

\noindent
\textbf{XVIII.\quad CROSS-REFERENCING}

\vspace{0.15cm}

Numbered elements shall be cross-referenced using the standard \LaTeX{}
mechanism.

Use the following commands where appropriate:

\begin{verbatim}
\label{...}
\ref{...}
\eqref{...}
\end{verbatim}

For example:

\begin{verbatim}
As shown in Figure~\ref{fig:system}, ...
\end{verbatim}

Students shall not manually type numbers such as ``Figure 2.3'' when
those numbers can be generated automatically.

\vspace{0.3cm}

% ================================================================
% XIX. PAGE NUMBERING
% ================================================================

\noindent
\textbf{XIX.\quad PAGE NUMBERING}

\vspace{0.15cm}

Cover pages and official evaluation forms at the beginning of the report
shall not display page numbers.

The main body of the report shall use Arabic page numbers beginning from
Chapter 1.

Page numbering shall be generated automatically by the official template.
Students shall not manually enter page numbers.

\vspace{0.3cm}

% ================================================================
% XX. PROJECT FILE ORGANIZATION
% ================================================================

\noindent
\textbf{XX.\quad PROJECT FILE ORGANIZATION}

\vspace{0.15cm}

Students should maintain a clear and organized \LaTeX{} project
structure. A recommended structure is:

\begin{verbatim}
InternshipReport/
|
|-- main.tex
|-- References.bib
|
|-- figures/
|   |-- company.pdf
|   |-- system.png
|   `-- ...
|
`-- chapters/
    |-- chapter1.tex
    |-- chapter2.tex
    `-- chapter3.tex
\end{verbatim}

Figures should normally be stored in the \verb|figures| directory
specified by the official template.

File and directory names should preferably contain Latin letters,
numbers, hyphens, or underscores.

Recommended examples include:

\begin{verbatim}
company_structure.pdf
motor_control.png
chapter_01.tex
\end{verbatim}

Spaces and unnecessary special characters in file names should be
avoided where possible.

\vspace{0.3cm}

% ================================================================
% XXI. DOCUMENT COMPILATION AND FINAL CHECK
% ================================================================

\noindent
\textbf{XXI.\quad DOCUMENT COMPILATION AND FINAL CHECK}

\vspace{0.15cm}

Before submission, the complete \LaTeX{} project shall compile
successfully and produce the final PDF report.

The final document shall not contain unresolved references such as:

\begin{verbatim}
??
[?]
\end{verbatim}

Before generating the final PDF, students shall check:

\begin{itemize}[leftmargin=0.8cm, itemsep=2pt]
    \item the Table of Contents;
    \item page numbering;
    \item chapter, section, and subsection numbering;
    \item figure numbering;
    \item table numbering;
    \item equation numbering;
    \item cross-references;
    \item citations;
    \item bibliography;
    \item font display;
    \item mathematical symbols;
    \item image quality;
    \item completeness of all required forms.
\end{itemize}

The document may be compiled using TeX Live, MiKTeX, or another
compatible \LaTeX{} environment. Regardless of the compilation
environment, the resulting PDF shall conform to the official Faculty
template.

\vspace{0.3cm}

% ================================================================
% XXII. CONTENT THAT STUDENTS MAY EDIT
% ================================================================

\noindent
\textbf{XXII.\quad CONTENT THAT STUDENTS MAY EDIT}

\vspace{0.15cm}

Students may edit or replace:

\begin{itemize}[leftmargin=0.8cm, itemsep=2pt]
    \item student name;
    \item class;
    \item internship organization;
    \item supervisor name;
    \item Abstract;
    \item List of Symbols and Abbreviations;
    \item chapter titles and chapter content;
    \item figures and photographs;
    \item tables;
    \item equations;
    \item citations and references.
\end{itemize}

Students shall not modify the main document formatting settings in the
preamble, including the document class, page geometry, line spacing,
heading formatting, table-of-contents formatting, font configuration,
page layout, or cover-page macros unless specifically instructed.

\vspace{0.3cm}

% ================================================================
% XXIII. FINAL SUBMISSION REQUIREMENTS
% ================================================================

\noindent
\textbf{XXIII.\quad FINAL SUBMISSION REQUIREMENTS}

\vspace{0.15cm}

Before submission, the final report shall:

\begin{itemize}[leftmargin=0.8cm, itemsep=2pt]
    \item compile without errors that affect the report content or presentation;
    \item contain all required forms and report sections;
    \item contain correct student, class, internship organization, and supervisor information;
    \item contain an updated Table of Contents, List of Figures, and List of Tables;
    \item contain no unused example content from the template;
    \item contain no placeholder text such as \textit{Title of Chapter 1},
    \textit{Section 1.1}, or other sample entries;
    \item contain all required citations and references;
    \item contain no unresolved cross-references.
\end{itemize}

If the Faculty requires submission of the \LaTeX{} source, the complete
project shall include all \verb|.tex| files, the bibliography database,
figures, and any other files required to reproduce the submitted PDF.

\vspace{0.4cm}

\noindent
\textbf{Important:}
Before submitting the final internship report, students shall remove or
comment out the following command in the main document:

\begin{verbatim}

% =================================================================
% REGULATIONS FOR PREPARING THE GRADUATION INTERNSHIP REPORT
% USING LATEX
% =================================================================
%
% This file is intended to be included from the main document with:
%
%     \input{Regulations.tex}
%
% IMPORTANT:
% This section is provided for guidance only. Students shall remove
% or comment out the \input{Regulations.tex} command before submitting
% the final internship report.
%
% =================================================================

\clearpage
\thispagestyle{empty}

\begin{center}
    {\fontsize{16pt}{20pt}\selectfont\bfseries
    SPECIFIC REGULATIONS FOR PREPARING\\
    GRADUATION INTERNSHIP REPORTS USING \LaTeX}
\end{center}

\vspace{0.5cm}

{\fontsize{13pt}{17pt}\selectfont
\setstretch{1.3}
\justifying

\noindent
\textbf{Note:}
This section is provided for guidance only and is included in the main
document using the command \verb|\input{Regulations.tex}|.
Students shall remove or comment out this command before submitting the
final internship report.

\vspace{0.4cm}

% ================================================================
% I. GENERAL REQUIREMENTS
% ================================================================

\noindent
\textbf{I.\quad GENERAL REQUIREMENTS}

\vspace{0.15cm}

The Graduation Internship Report shall be prepared using the official
\LaTeX{} template provided by the Faculty. Students shall preserve the
predefined page layout, typography, line spacing, heading styles,
numbering system, figure and table formatting, and bibliography settings
unless otherwise instructed by the Faculty or the supervisor.

The report shall be written in English.

Students are responsible for the academic and technical content of the
report. The official \LaTeX{} template is responsible for standardized
document formatting.

\vspace{0.3cm}

% ================================================================
% II. REPORT STRUCTURE
% ================================================================

\noindent
\textbf{II.\quad REPORT STRUCTURE}

\vspace{0.15cm}

The report shall normally be arranged in the following order:

\begin{enumerate}[label=\arabic*., leftmargin=1.2cm, itemsep=2pt]
    \item Outer Cover;
    \item Comments of the Internship Host Organization;
    \item Supervisor's Evaluation and Graduation Internship Defense Results;
    \item Table of Contents;
    \item List of Symbols and Abbreviations, if applicable;
    \item List of Figures;
    \item List of Tables;
    \item Abstract;
    \item Main Chapters of the Internship Report;
    \item References;
    \item Appendices, if applicable.
\end{enumerate}

The preliminary pages shall not be numbered as chapters. The main report
content shall be organized into consecutively numbered chapters beginning
with Chapter 1.

\vspace{0.3cm}

% ================================================================
% III. PAGE SIZE AND MARGINS
% ================================================================

\noindent
\textbf{III.\quad PAGE SIZE AND MARGINS}

\vspace{0.15cm}

The report shall be formatted on A4 paper with the following margins:

\begin{center}
\begin{tabular}{ll}
    Left margin:   & 3.5 cm \\
    Right margin:  & 2.0 cm \\
    Top margin:    & 2.5 cm \\
    Bottom margin: & 2.5 cm
\end{tabular}
\end{center}

These settings are predefined in the official template and shall not be
modified by students.

\vspace{0.3cm}

% ================================================================
% IV. FONT, FONT SIZE, AND LINE SPACING
% ================================================================

\noindent
\textbf{IV.\quad FONT, FONT SIZE, AND LINE SPACING}

\vspace{0.15cm}

The main text shall use a 13 pt font size with 1.3-line spacing.
A Times-family typeface shall be used in accordance with the official
template.

The text shall be fully justified.

Students shall not use repeated spaces, repeated manual line breaks, or
multiple blank lines to position text. Paragraph and heading spacing shall
be controlled by the template.

\vspace{0.3cm}

% ================================================================
% V. OUTER COVER
% ================================================================

\noindent
\textbf{V.\quad OUTER COVER}

\vspace{0.15cm}

The outer cover shall follow the official format provided in the template
and shall contain the following information:

\begin{itemize}[leftmargin=0.8cm, itemsep=2pt]
    \item Thai Nguyen University of Technology;
    \item Faculty of International Training;
    \item Engineering and Technology Department;
    \item University logo and Faculty logo;
    \item the title \textbf{INTERNSHIP REPORT};
    \item Internship Organization;
    \item Author;
    \item Class;
    \item Supervisor;
    \item place and year of completion.
\end{itemize}

Students shall only replace the required information fields and shall not
modify the cover-page layout unless instructed otherwise.

\vspace{0.3cm}

% ================================================================
% VI. COMMENTS OF THE INTERNSHIP HOST ORGANIZATION
% ================================================================

\noindent
\textbf{VI.\quad COMMENTS OF THE INTERNSHIP HOST ORGANIZATION}

\vspace{0.15cm}

The page entitled \textbf{COMMENTS OF THE INTERNSHIP HOST ORGANIZATION}
is reserved for the internship host organization.

It includes:

\begin{itemize}[leftmargin=0.8cm, itemsep=2pt]
    \item the student's full name;
    \item class;
    \item comments from the host organization;
    \item score;
    \item date;
    \item signature and official seal.
\end{itemize}

Students shall not alter or complete official evaluation fields unless
specifically authorized.

\vspace{0.3cm}

% ================================================================
% VII. SUPERVISOR'S EVALUATION AND DEFENSE RESULTS
% ================================================================

\noindent
\textbf{VII.\quad SUPERVISOR'S EVALUATION AND DEFENSE RESULTS}

\vspace{0.15cm}

The section entitled
\textbf{SUPERVISOR'S EVALUATION AND GRADUATION INTERNSHIP DEFENSE RESULTS}
is reserved for the supervisor and defense committee.

The evaluation may include:

\begin{itemize}[leftmargin=0.8cm, itemsep=2pt]
    \item Internship Progress and Student's Attitude;
    \item Report Content;
    \item Report Presentation;
    \item Other Comments;
    \item Supervisor's Opinion;
    \item Defense Result;
    \item Score.
\end{itemize}

Students shall not enter evaluation scores or defense results before they
are officially determined.

\vspace{0.3cm}

% ================================================================
% VIII. CHAPTERS, SECTIONS, AND SUBSECTIONS
% ================================================================

\noindent
\textbf{VIII.\quad CHAPTERS, SECTIONS, AND SUBSECTIONS}

\vspace{0.15cm}

The report structure shall be generated using the standard \LaTeX{}
commands:

\begin{verbatim}
\chapter{...}
\section{...}
\subsection{...}
\end{verbatim}

Chapter, section, and subsection numbers shall be generated automatically
by \LaTeX{}.

Students shall not manually type chapter or section numbers and shall not
reproduce heading styles using commands such as \verb|\textbf|,
\verb|\fontsize|, or \verb|\centering| instead of the appropriate
structural command.

Excessive subdivision should be avoided. In most cases, the hierarchy
Chapter--Section--Subsection is sufficient.

\vspace{0.3cm}

% ================================================================
% IX. ORGANIZATION OF THE MAIN CONTENT
% ================================================================

\noindent
\textbf{IX.\quad ORGANIZATION OF THE MAIN CONTENT}

\vspace{0.15cm}

The specific content of the report shall depend on the internship
organization and the assigned tasks. A typical structure may include:

\begin{itemize}[leftmargin=0.8cm, itemsep=2pt]
    \item \textbf{Chapter 1. Overview of the Internship Organization};
    \item \textbf{Chapter 2. Internship Activities and Assigned Tasks};
    \item \textbf{Chapter 3. Results, Experience, and Evaluation}.
\end{itemize}

These chapter titles are recommendations only and may be adapted to the
actual internship content.

The report shall reflect work actually performed, observed, analyzed, or
learned during the internship. Unrelated theoretical material should not
be included merely to increase report length.

\vspace{0.3cm}

% ================================================================
% X. TABLE OF CONTENTS
% ================================================================

\noindent
\textbf{X.\quad TABLE OF CONTENTS}

\vspace{0.15cm}

The Table of Contents shall be generated automatically using:

\begin{verbatim}
\tableofcontents
\end{verbatim}

Students shall not create or manually edit page numbers in the Table of
Contents.

The document shall be compiled sufficiently to ensure that all page
numbers and entries are correctly updated.

\vspace{0.3cm}

% ================================================================
% XI. LIST OF SYMBOLS AND ABBREVIATIONS
% ================================================================

\noindent
\textbf{XI.\quad LIST OF SYMBOLS AND ABBREVIATIONS}

\vspace{0.15cm}

Symbols and abbreviations that are used repeatedly in the report shall be
listed in the \textbf{List of Symbols and Abbreviations}.

The list may contain the following columns:

\begin{center}
\begin{tabular}{llll}
    No. & Symbol / Abbreviation & Unit & Definition
\end{tabular}
\end{center}

Only symbols and abbreviations actually used in the final report shall be
retained. Example entries contained in the template shall be removed if
they are not relevant to the student's report.

Each abbreviation should normally be defined when it first appears in
the main text.

\vspace{0.3cm}

% ================================================================
% XII. FIGURES, DIAGRAMS, AND PHOTOGRAPHS
% ================================================================

\noindent
\textbf{XII.\quad FIGURES, DIAGRAMS, AND PHOTOGRAPHS}

\vspace{0.15cm}

Figures shall be clear, legible, and directly relevant to the report
content.

Figures shall be inserted using the standard \verb|figure| environment.

Example:

\begin{verbatim}
\begin{figure}[H]
    \centering
    \includegraphics[width=0.8\textwidth]
                    {figures/company_structure.pdf}
    \caption{Organizational structure of the company}
    \label{fig:company-structure}
\end{figure}
\end{verbatim}

Figure captions shall be placed below the figure.

Each referenced figure shall have a unique \verb|\label| and shall be
referenced in the text using, for example:

\begin{verbatim}
Figure~\ref{fig:company-structure}
\end{verbatim}

Students shall not manually type figure numbers.

Figures reproduced or adapted from external sources shall include an
appropriate citation.

Photographs taken during the internship shall be relevant and sufficiently
clear. Students shall not include confidential, restricted, or proprietary
information without permission from the internship organization.

Vector graphics or high-resolution images are recommended whenever
possible.

\vspace{0.3cm}

% ================================================================
% XIII. TABLES
% ================================================================

\noindent
\textbf{XIII.\quad TABLES}

\vspace{0.15cm}

Tables shall be created using suitable \LaTeX{} environments such as
\verb|table|, \verb|tabular|, \verb|tabularx|, or \verb|longtable|.

Table captions shall be placed above the table.

Example:

\begin{verbatim}
\begin{table}[H]
    \centering
    \caption{Internship schedule}
    \label{tab:internship-schedule}

    \begin{tabular}{lll}
        \toprule
        Week & Task & Result \\
        \midrule
        1 & Orientation & Completed \\
        2 & System study & Completed \\
        \bottomrule
    \end{tabular}
\end{table}
\end{verbatim}

Tables shall be referenced using:

\begin{verbatim}
Table~\ref{tab:internship-schedule}
\end{verbatim}

Students shall not manually type table numbers.

Tables should be created directly in \LaTeX{} rather than inserted as
screenshots whenever practical.

\vspace{0.3cm}

% ================================================================
% XIV. EQUATIONS AND TECHNICAL QUANTITIES
% ================================================================

\noindent
\textbf{XIV.\quad EQUATIONS AND TECHNICAL QUANTITIES}

\vspace{0.15cm}

Mathematical expressions shall be typeset directly using \LaTeX{}.
Images or screenshots of equations shall not be used.

Short expressions may be written in inline mathematical mode, for example:

\begin{verbatim}
$R = 0.7\,\Omega$
\end{verbatim}

Important equations shall use a numbered equation environment.

Example:

\begin{verbatim}
\begin{equation}
    P = UI
    \label{eq:power}
\end{equation}
\end{verbatim}

Equations shall be referenced using:

\begin{verbatim}
Equation~\eqref{eq:power}
\end{verbatim}

Equation numbers shall not be manually entered.

SI units shall be used wherever applicable. The \verb|siunitx| package
should be used where appropriate.

Examples:

\begin{verbatim}
\SI{220}{\volt}
\SI{50}{\hertz}
\SI{2.5}{\ampere}
\end{verbatim}

Notation and units shall be consistent throughout the report.

\vspace{0.3cm}

% ================================================================
% XV. SOURCE CODE
% ================================================================

\noindent
\textbf{XV.\quad SOURCE CODE}

\vspace{0.15cm}

If the internship involves programming, PLCs, microcontrollers,
simulation, or software development, relevant source code may be included
in the report.

Source code shall be presented using the \verb|listings| environment or
another environment defined by the official template.

Example:

\begin{verbatim}
\begin{lstlisting}[language=st]
IF Start = TRUE THEN
    Motor := TRUE;
ELSE
    Motor := FALSE;
END_IF;
\end{lstlisting}
\end{verbatim}

Only code necessary for explaining the work should appear in the main
report. Long source-code listings should be moved to an appendix or
submitted as supplementary files if required.

Source code should not normally be inserted as screenshots.

\vspace{0.3cm}

% ================================================================
% XVI. ABSTRACT
% ================================================================

\noindent
\textbf{XVI.\quad ABSTRACT}

\vspace{0.15cm}

The Abstract shall briefly summarize:

\begin{itemize}[leftmargin=0.8cm, itemsep=2pt]
    \item the internship organization;
    \item the objectives of the internship;
    \item the main tasks performed;
    \item the main results, knowledge, and skills obtained.
\end{itemize}

The Abstract should be concise and should not reproduce the full content
of the report chapter by chapter.

\vspace{0.3cm}

% ================================================================
% XVII. REFERENCES AND CITATIONS
% ================================================================

\noindent
\textbf{XVII.\quad REFERENCES AND CITATIONS}

\vspace{0.15cm}

All books, journal papers, conference papers, standards, websites,
datasheets, manuals, technical documents, figures, tables, or other
external sources used in preparing the report shall be properly cited.

References shall be managed through the bibliography system provided with
the official template.

Bibliographic entries shall normally be stored in:

\begin{verbatim}
References.bib
\end{verbatim}

References shall be cited in the report using:

\begin{verbatim}
\cite{reference-key}
\end{verbatim}

Students shall not manually type reference numbers such as ``[1]'',
``[2]'', etc.

Titles of foreign-language publications shall remain in their original
language in the bibliography.

All borrowed ideas, data, figures, tables, equations, methods, or results
shall be properly acknowledged.

\vspace{0.3cm}

% ================================================================
% XVIII. CROSS-REFERENCING
% ================================================================

\noindent
\textbf{XVIII.\quad CROSS-REFERENCING}

\vspace{0.15cm}

Numbered elements shall be cross-referenced using the standard \LaTeX{}
mechanism.

Use the following commands where appropriate:

\begin{verbatim}
\label{...}
\ref{...}
\eqref{...}
\end{verbatim}

For example:

\begin{verbatim}
As shown in Figure~\ref{fig:system}, ...
\end{verbatim}

Students shall not manually type numbers such as ``Figure 2.3'' when
those numbers can be generated automatically.

\vspace{0.3cm}

% ================================================================
% XIX. PAGE NUMBERING
% ================================================================

\noindent
\textbf{XIX.\quad PAGE NUMBERING}

\vspace{0.15cm}

Cover pages and official evaluation forms at the beginning of the report
shall not display page numbers.

The main body of the report shall use Arabic page numbers beginning from
Chapter 1.

Page numbering shall be generated automatically by the official template.
Students shall not manually enter page numbers.

\vspace{0.3cm}

% ================================================================
% XX. PROJECT FILE ORGANIZATION
% ================================================================

\noindent
\textbf{XX.\quad PROJECT FILE ORGANIZATION}

\vspace{0.15cm}

Students should maintain a clear and organized \LaTeX{} project
structure. A recommended structure is:

\begin{verbatim}
InternshipReport/
|
|-- main.tex
|-- References.bib
|
|-- figures/
|   |-- company.pdf
|   |-- system.png
|   `-- ...
|
`-- chapters/
    |-- chapter1.tex
    |-- chapter2.tex
    `-- chapter3.tex
\end{verbatim}

Figures should normally be stored in the \verb|figures| directory
specified by the official template.

File and directory names should preferably contain Latin letters,
numbers, hyphens, or underscores.

Recommended examples include:

\begin{verbatim}
company_structure.pdf
motor_control.png
chapter_01.tex
\end{verbatim}

Spaces and unnecessary special characters in file names should be
avoided where possible.

\vspace{0.3cm}

% ================================================================
% XXI. DOCUMENT COMPILATION AND FINAL CHECK
% ================================================================

\noindent
\textbf{XXI.\quad DOCUMENT COMPILATION AND FINAL CHECK}

\vspace{0.15cm}

Before submission, the complete \LaTeX{} project shall compile
successfully and produce the final PDF report.

The final document shall not contain unresolved references such as:

\begin{verbatim}
??
[?]
\end{verbatim}

Before generating the final PDF, students shall check:

\begin{itemize}[leftmargin=0.8cm, itemsep=2pt]
    \item the Table of Contents;
    \item page numbering;
    \item chapter, section, and subsection numbering;
    \item figure numbering;
    \item table numbering;
    \item equation numbering;
    \item cross-references;
    \item citations;
    \item bibliography;
    \item font display;
    \item mathematical symbols;
    \item image quality;
    \item completeness of all required forms.
\end{itemize}

The document may be compiled using TeX Live, MiKTeX, or another
compatible \LaTeX{} environment. Regardless of the compilation
environment, the resulting PDF shall conform to the official Faculty
template.

\vspace{0.3cm}

% ================================================================
% XXII. CONTENT THAT STUDENTS MAY EDIT
% ================================================================

\noindent
\textbf{XXII.\quad CONTENT THAT STUDENTS MAY EDIT}

\vspace{0.15cm}

Students may edit or replace:

\begin{itemize}[leftmargin=0.8cm, itemsep=2pt]
    \item student name;
    \item class;
    \item internship organization;
    \item supervisor name;
    \item Abstract;
    \item List of Symbols and Abbreviations;
    \item chapter titles and chapter content;
    \item figures and photographs;
    \item tables;
    \item equations;
    \item citations and references.
\end{itemize}

Students shall not modify the main document formatting settings in the
preamble, including the document class, page geometry, line spacing,
heading formatting, table-of-contents formatting, font configuration,
page layout, or cover-page macros unless specifically instructed.

\vspace{0.3cm}

% ================================================================
% XXIII. FINAL SUBMISSION REQUIREMENTS
% ================================================================

\noindent
\textbf{XXIII.\quad FINAL SUBMISSION REQUIREMENTS}

\vspace{0.15cm}

Before submission, the final report shall:

\begin{itemize}[leftmargin=0.8cm, itemsep=2pt]
    \item compile without errors that affect the report content or presentation;
    \item contain all required forms and report sections;
    \item contain correct student, class, internship organization, and supervisor information;
    \item contain an updated Table of Contents, List of Figures, and List of Tables;
    \item contain no unused example content from the template;
    \item contain no placeholder text such as \textit{Title of Chapter 1},
    \textit{Section 1.1}, or other sample entries;
    \item contain all required citations and references;
    \item contain no unresolved cross-references.
\end{itemize}

If the Faculty requires submission of the \LaTeX{} source, the complete
project shall include all \verb|.tex| files, the bibliography database,
figures, and any other files required to reproduce the submitted PDF.

\vspace{0.4cm}

\noindent
\textbf{Important:}
Before submitting the final internship report, students shall remove or
comment out the following command in the main document:

\begin{verbatim}
\input{Regulations.tex}
\end{verbatim}

For example:

\begin{verbatim}
% \input{Regulations.tex}
\end{verbatim}

This removes the guidance pages from the final submitted report while
preserving the regulations file in the \LaTeX{} project for reference.

} % End 13pt content

\clearpage

\end{verbatim}

For example:

\begin{verbatim}
% 
% =================================================================
% REGULATIONS FOR PREPARING THE GRADUATION INTERNSHIP REPORT
% USING LATEX
% =================================================================
%
% This file is intended to be included from the main document with:
%
%     \input{Regulations.tex}
%
% IMPORTANT:
% This section is provided for guidance only. Students shall remove
% or comment out the \input{Regulations.tex} command before submitting
% the final internship report.
%
% =================================================================

\clearpage
\thispagestyle{empty}

\begin{center}
    {\fontsize{16pt}{20pt}\selectfont\bfseries
    SPECIFIC REGULATIONS FOR PREPARING\\
    GRADUATION INTERNSHIP REPORTS USING \LaTeX}
\end{center}

\vspace{0.5cm}

{\fontsize{13pt}{17pt}\selectfont
\setstretch{1.3}
\justifying

\noindent
\textbf{Note:}
This section is provided for guidance only and is included in the main
document using the command \verb|\input{Regulations.tex}|.
Students shall remove or comment out this command before submitting the
final internship report.

\vspace{0.4cm}

% ================================================================
% I. GENERAL REQUIREMENTS
% ================================================================

\noindent
\textbf{I.\quad GENERAL REQUIREMENTS}

\vspace{0.15cm}

The Graduation Internship Report shall be prepared using the official
\LaTeX{} template provided by the Faculty. Students shall preserve the
predefined page layout, typography, line spacing, heading styles,
numbering system, figure and table formatting, and bibliography settings
unless otherwise instructed by the Faculty or the supervisor.

The report shall be written in English.

Students are responsible for the academic and technical content of the
report. The official \LaTeX{} template is responsible for standardized
document formatting.

\vspace{0.3cm}

% ================================================================
% II. REPORT STRUCTURE
% ================================================================

\noindent
\textbf{II.\quad REPORT STRUCTURE}

\vspace{0.15cm}

The report shall normally be arranged in the following order:

\begin{enumerate}[label=\arabic*., leftmargin=1.2cm, itemsep=2pt]
    \item Outer Cover;
    \item Comments of the Internship Host Organization;
    \item Supervisor's Evaluation and Graduation Internship Defense Results;
    \item Table of Contents;
    \item List of Symbols and Abbreviations, if applicable;
    \item List of Figures;
    \item List of Tables;
    \item Abstract;
    \item Main Chapters of the Internship Report;
    \item References;
    \item Appendices, if applicable.
\end{enumerate}

The preliminary pages shall not be numbered as chapters. The main report
content shall be organized into consecutively numbered chapters beginning
with Chapter 1.

\vspace{0.3cm}

% ================================================================
% III. PAGE SIZE AND MARGINS
% ================================================================

\noindent
\textbf{III.\quad PAGE SIZE AND MARGINS}

\vspace{0.15cm}

The report shall be formatted on A4 paper with the following margins:

\begin{center}
\begin{tabular}{ll}
    Left margin:   & 3.5 cm \\
    Right margin:  & 2.0 cm \\
    Top margin:    & 2.5 cm \\
    Bottom margin: & 2.5 cm
\end{tabular}
\end{center}

These settings are predefined in the official template and shall not be
modified by students.

\vspace{0.3cm}

% ================================================================
% IV. FONT, FONT SIZE, AND LINE SPACING
% ================================================================

\noindent
\textbf{IV.\quad FONT, FONT SIZE, AND LINE SPACING}

\vspace{0.15cm}

The main text shall use a 13 pt font size with 1.3-line spacing.
A Times-family typeface shall be used in accordance with the official
template.

The text shall be fully justified.

Students shall not use repeated spaces, repeated manual line breaks, or
multiple blank lines to position text. Paragraph and heading spacing shall
be controlled by the template.

\vspace{0.3cm}

% ================================================================
% V. OUTER COVER
% ================================================================

\noindent
\textbf{V.\quad OUTER COVER}

\vspace{0.15cm}

The outer cover shall follow the official format provided in the template
and shall contain the following information:

\begin{itemize}[leftmargin=0.8cm, itemsep=2pt]
    \item Thai Nguyen University of Technology;
    \item Faculty of International Training;
    \item Engineering and Technology Department;
    \item University logo and Faculty logo;
    \item the title \textbf{INTERNSHIP REPORT};
    \item Internship Organization;
    \item Author;
    \item Class;
    \item Supervisor;
    \item place and year of completion.
\end{itemize}

Students shall only replace the required information fields and shall not
modify the cover-page layout unless instructed otherwise.

\vspace{0.3cm}

% ================================================================
% VI. COMMENTS OF THE INTERNSHIP HOST ORGANIZATION
% ================================================================

\noindent
\textbf{VI.\quad COMMENTS OF THE INTERNSHIP HOST ORGANIZATION}

\vspace{0.15cm}

The page entitled \textbf{COMMENTS OF THE INTERNSHIP HOST ORGANIZATION}
is reserved for the internship host organization.

It includes:

\begin{itemize}[leftmargin=0.8cm, itemsep=2pt]
    \item the student's full name;
    \item class;
    \item comments from the host organization;
    \item score;
    \item date;
    \item signature and official seal.
\end{itemize}

Students shall not alter or complete official evaluation fields unless
specifically authorized.

\vspace{0.3cm}

% ================================================================
% VII. SUPERVISOR'S EVALUATION AND DEFENSE RESULTS
% ================================================================

\noindent
\textbf{VII.\quad SUPERVISOR'S EVALUATION AND DEFENSE RESULTS}

\vspace{0.15cm}

The section entitled
\textbf{SUPERVISOR'S EVALUATION AND GRADUATION INTERNSHIP DEFENSE RESULTS}
is reserved for the supervisor and defense committee.

The evaluation may include:

\begin{itemize}[leftmargin=0.8cm, itemsep=2pt]
    \item Internship Progress and Student's Attitude;
    \item Report Content;
    \item Report Presentation;
    \item Other Comments;
    \item Supervisor's Opinion;
    \item Defense Result;
    \item Score.
\end{itemize}

Students shall not enter evaluation scores or defense results before they
are officially determined.

\vspace{0.3cm}

% ================================================================
% VIII. CHAPTERS, SECTIONS, AND SUBSECTIONS
% ================================================================

\noindent
\textbf{VIII.\quad CHAPTERS, SECTIONS, AND SUBSECTIONS}

\vspace{0.15cm}

The report structure shall be generated using the standard \LaTeX{}
commands:

\begin{verbatim}
\chapter{...}
\section{...}
\subsection{...}
\end{verbatim}

Chapter, section, and subsection numbers shall be generated automatically
by \LaTeX{}.

Students shall not manually type chapter or section numbers and shall not
reproduce heading styles using commands such as \verb|\textbf|,
\verb|\fontsize|, or \verb|\centering| instead of the appropriate
structural command.

Excessive subdivision should be avoided. In most cases, the hierarchy
Chapter--Section--Subsection is sufficient.

\vspace{0.3cm}

% ================================================================
% IX. ORGANIZATION OF THE MAIN CONTENT
% ================================================================

\noindent
\textbf{IX.\quad ORGANIZATION OF THE MAIN CONTENT}

\vspace{0.15cm}

The specific content of the report shall depend on the internship
organization and the assigned tasks. A typical structure may include:

\begin{itemize}[leftmargin=0.8cm, itemsep=2pt]
    \item \textbf{Chapter 1. Overview of the Internship Organization};
    \item \textbf{Chapter 2. Internship Activities and Assigned Tasks};
    \item \textbf{Chapter 3. Results, Experience, and Evaluation}.
\end{itemize}

These chapter titles are recommendations only and may be adapted to the
actual internship content.

The report shall reflect work actually performed, observed, analyzed, or
learned during the internship. Unrelated theoretical material should not
be included merely to increase report length.

\vspace{0.3cm}

% ================================================================
% X. TABLE OF CONTENTS
% ================================================================

\noindent
\textbf{X.\quad TABLE OF CONTENTS}

\vspace{0.15cm}

The Table of Contents shall be generated automatically using:

\begin{verbatim}
\tableofcontents
\end{verbatim}

Students shall not create or manually edit page numbers in the Table of
Contents.

The document shall be compiled sufficiently to ensure that all page
numbers and entries are correctly updated.

\vspace{0.3cm}

% ================================================================
% XI. LIST OF SYMBOLS AND ABBREVIATIONS
% ================================================================

\noindent
\textbf{XI.\quad LIST OF SYMBOLS AND ABBREVIATIONS}

\vspace{0.15cm}

Symbols and abbreviations that are used repeatedly in the report shall be
listed in the \textbf{List of Symbols and Abbreviations}.

The list may contain the following columns:

\begin{center}
\begin{tabular}{llll}
    No. & Symbol / Abbreviation & Unit & Definition
\end{tabular}
\end{center}

Only symbols and abbreviations actually used in the final report shall be
retained. Example entries contained in the template shall be removed if
they are not relevant to the student's report.

Each abbreviation should normally be defined when it first appears in
the main text.

\vspace{0.3cm}

% ================================================================
% XII. FIGURES, DIAGRAMS, AND PHOTOGRAPHS
% ================================================================

\noindent
\textbf{XII.\quad FIGURES, DIAGRAMS, AND PHOTOGRAPHS}

\vspace{0.15cm}

Figures shall be clear, legible, and directly relevant to the report
content.

Figures shall be inserted using the standard \verb|figure| environment.

Example:

\begin{verbatim}
\begin{figure}[H]
    \centering
    \includegraphics[width=0.8\textwidth]
                    {figures/company_structure.pdf}
    \caption{Organizational structure of the company}
    \label{fig:company-structure}
\end{figure}
\end{verbatim}

Figure captions shall be placed below the figure.

Each referenced figure shall have a unique \verb|\label| and shall be
referenced in the text using, for example:

\begin{verbatim}
Figure~\ref{fig:company-structure}
\end{verbatim}

Students shall not manually type figure numbers.

Figures reproduced or adapted from external sources shall include an
appropriate citation.

Photographs taken during the internship shall be relevant and sufficiently
clear. Students shall not include confidential, restricted, or proprietary
information without permission from the internship organization.

Vector graphics or high-resolution images are recommended whenever
possible.

\vspace{0.3cm}

% ================================================================
% XIII. TABLES
% ================================================================

\noindent
\textbf{XIII.\quad TABLES}

\vspace{0.15cm}

Tables shall be created using suitable \LaTeX{} environments such as
\verb|table|, \verb|tabular|, \verb|tabularx|, or \verb|longtable|.

Table captions shall be placed above the table.

Example:

\begin{verbatim}
\begin{table}[H]
    \centering
    \caption{Internship schedule}
    \label{tab:internship-schedule}

    \begin{tabular}{lll}
        \toprule
        Week & Task & Result \\
        \midrule
        1 & Orientation & Completed \\
        2 & System study & Completed \\
        \bottomrule
    \end{tabular}
\end{table}
\end{verbatim}

Tables shall be referenced using:

\begin{verbatim}
Table~\ref{tab:internship-schedule}
\end{verbatim}

Students shall not manually type table numbers.

Tables should be created directly in \LaTeX{} rather than inserted as
screenshots whenever practical.

\vspace{0.3cm}

% ================================================================
% XIV. EQUATIONS AND TECHNICAL QUANTITIES
% ================================================================

\noindent
\textbf{XIV.\quad EQUATIONS AND TECHNICAL QUANTITIES}

\vspace{0.15cm}

Mathematical expressions shall be typeset directly using \LaTeX{}.
Images or screenshots of equations shall not be used.

Short expressions may be written in inline mathematical mode, for example:

\begin{verbatim}
$R = 0.7\,\Omega$
\end{verbatim}

Important equations shall use a numbered equation environment.

Example:

\begin{verbatim}
\begin{equation}
    P = UI
    \label{eq:power}
\end{equation}
\end{verbatim}

Equations shall be referenced using:

\begin{verbatim}
Equation~\eqref{eq:power}
\end{verbatim}

Equation numbers shall not be manually entered.

SI units shall be used wherever applicable. The \verb|siunitx| package
should be used where appropriate.

Examples:

\begin{verbatim}
\SI{220}{\volt}
\SI{50}{\hertz}
\SI{2.5}{\ampere}
\end{verbatim}

Notation and units shall be consistent throughout the report.

\vspace{0.3cm}

% ================================================================
% XV. SOURCE CODE
% ================================================================

\noindent
\textbf{XV.\quad SOURCE CODE}

\vspace{0.15cm}

If the internship involves programming, PLCs, microcontrollers,
simulation, or software development, relevant source code may be included
in the report.

Source code shall be presented using the \verb|listings| environment or
another environment defined by the official template.

Example:

\begin{verbatim}
\begin{lstlisting}[language=st]
IF Start = TRUE THEN
    Motor := TRUE;
ELSE
    Motor := FALSE;
END_IF;
\end{lstlisting}
\end{verbatim}

Only code necessary for explaining the work should appear in the main
report. Long source-code listings should be moved to an appendix or
submitted as supplementary files if required.

Source code should not normally be inserted as screenshots.

\vspace{0.3cm}

% ================================================================
% XVI. ABSTRACT
% ================================================================

\noindent
\textbf{XVI.\quad ABSTRACT}

\vspace{0.15cm}

The Abstract shall briefly summarize:

\begin{itemize}[leftmargin=0.8cm, itemsep=2pt]
    \item the internship organization;
    \item the objectives of the internship;
    \item the main tasks performed;
    \item the main results, knowledge, and skills obtained.
\end{itemize}

The Abstract should be concise and should not reproduce the full content
of the report chapter by chapter.

\vspace{0.3cm}

% ================================================================
% XVII. REFERENCES AND CITATIONS
% ================================================================

\noindent
\textbf{XVII.\quad REFERENCES AND CITATIONS}

\vspace{0.15cm}

All books, journal papers, conference papers, standards, websites,
datasheets, manuals, technical documents, figures, tables, or other
external sources used in preparing the report shall be properly cited.

References shall be managed through the bibliography system provided with
the official template.

Bibliographic entries shall normally be stored in:

\begin{verbatim}
References.bib
\end{verbatim}

References shall be cited in the report using:

\begin{verbatim}
\cite{reference-key}
\end{verbatim}

Students shall not manually type reference numbers such as ``[1]'',
``[2]'', etc.

Titles of foreign-language publications shall remain in their original
language in the bibliography.

All borrowed ideas, data, figures, tables, equations, methods, or results
shall be properly acknowledged.

\vspace{0.3cm}

% ================================================================
% XVIII. CROSS-REFERENCING
% ================================================================

\noindent
\textbf{XVIII.\quad CROSS-REFERENCING}

\vspace{0.15cm}

Numbered elements shall be cross-referenced using the standard \LaTeX{}
mechanism.

Use the following commands where appropriate:

\begin{verbatim}
\label{...}
\ref{...}
\eqref{...}
\end{verbatim}

For example:

\begin{verbatim}
As shown in Figure~\ref{fig:system}, ...
\end{verbatim}

Students shall not manually type numbers such as ``Figure 2.3'' when
those numbers can be generated automatically.

\vspace{0.3cm}

% ================================================================
% XIX. PAGE NUMBERING
% ================================================================

\noindent
\textbf{XIX.\quad PAGE NUMBERING}

\vspace{0.15cm}

Cover pages and official evaluation forms at the beginning of the report
shall not display page numbers.

The main body of the report shall use Arabic page numbers beginning from
Chapter 1.

Page numbering shall be generated automatically by the official template.
Students shall not manually enter page numbers.

\vspace{0.3cm}

% ================================================================
% XX. PROJECT FILE ORGANIZATION
% ================================================================

\noindent
\textbf{XX.\quad PROJECT FILE ORGANIZATION}

\vspace{0.15cm}

Students should maintain a clear and organized \LaTeX{} project
structure. A recommended structure is:

\begin{verbatim}
InternshipReport/
|
|-- main.tex
|-- References.bib
|
|-- figures/
|   |-- company.pdf
|   |-- system.png
|   `-- ...
|
`-- chapters/
    |-- chapter1.tex
    |-- chapter2.tex
    `-- chapter3.tex
\end{verbatim}

Figures should normally be stored in the \verb|figures| directory
specified by the official template.

File and directory names should preferably contain Latin letters,
numbers, hyphens, or underscores.

Recommended examples include:

\begin{verbatim}
company_structure.pdf
motor_control.png
chapter_01.tex
\end{verbatim}

Spaces and unnecessary special characters in file names should be
avoided where possible.

\vspace{0.3cm}

% ================================================================
% XXI. DOCUMENT COMPILATION AND FINAL CHECK
% ================================================================

\noindent
\textbf{XXI.\quad DOCUMENT COMPILATION AND FINAL CHECK}

\vspace{0.15cm}

Before submission, the complete \LaTeX{} project shall compile
successfully and produce the final PDF report.

The final document shall not contain unresolved references such as:

\begin{verbatim}
??
[?]
\end{verbatim}

Before generating the final PDF, students shall check:

\begin{itemize}[leftmargin=0.8cm, itemsep=2pt]
    \item the Table of Contents;
    \item page numbering;
    \item chapter, section, and subsection numbering;
    \item figure numbering;
    \item table numbering;
    \item equation numbering;
    \item cross-references;
    \item citations;
    \item bibliography;
    \item font display;
    \item mathematical symbols;
    \item image quality;
    \item completeness of all required forms.
\end{itemize}

The document may be compiled using TeX Live, MiKTeX, or another
compatible \LaTeX{} environment. Regardless of the compilation
environment, the resulting PDF shall conform to the official Faculty
template.

\vspace{0.3cm}

% ================================================================
% XXII. CONTENT THAT STUDENTS MAY EDIT
% ================================================================

\noindent
\textbf{XXII.\quad CONTENT THAT STUDENTS MAY EDIT}

\vspace{0.15cm}

Students may edit or replace:

\begin{itemize}[leftmargin=0.8cm, itemsep=2pt]
    \item student name;
    \item class;
    \item internship organization;
    \item supervisor name;
    \item Abstract;
    \item List of Symbols and Abbreviations;
    \item chapter titles and chapter content;
    \item figures and photographs;
    \item tables;
    \item equations;
    \item citations and references.
\end{itemize}

Students shall not modify the main document formatting settings in the
preamble, including the document class, page geometry, line spacing,
heading formatting, table-of-contents formatting, font configuration,
page layout, or cover-page macros unless specifically instructed.

\vspace{0.3cm}

% ================================================================
% XXIII. FINAL SUBMISSION REQUIREMENTS
% ================================================================

\noindent
\textbf{XXIII.\quad FINAL SUBMISSION REQUIREMENTS}

\vspace{0.15cm}

Before submission, the final report shall:

\begin{itemize}[leftmargin=0.8cm, itemsep=2pt]
    \item compile without errors that affect the report content or presentation;
    \item contain all required forms and report sections;
    \item contain correct student, class, internship organization, and supervisor information;
    \item contain an updated Table of Contents, List of Figures, and List of Tables;
    \item contain no unused example content from the template;
    \item contain no placeholder text such as \textit{Title of Chapter 1},
    \textit{Section 1.1}, or other sample entries;
    \item contain all required citations and references;
    \item contain no unresolved cross-references.
\end{itemize}

If the Faculty requires submission of the \LaTeX{} source, the complete
project shall include all \verb|.tex| files, the bibliography database,
figures, and any other files required to reproduce the submitted PDF.

\vspace{0.4cm}

\noindent
\textbf{Important:}
Before submitting the final internship report, students shall remove or
comment out the following command in the main document:

\begin{verbatim}
\input{Regulations.tex}
\end{verbatim}

For example:

\begin{verbatim}
% \input{Regulations.tex}
\end{verbatim}

This removes the guidance pages from the final submitted report while
preserving the regulations file in the \LaTeX{} project for reference.

} % End 13pt content

\clearpage

\end{verbatim}

This removes the guidance pages from the final submitted report while
preserving the regulations file in the \LaTeX{} project for reference.

} % End 13pt content

\clearpage

\end{verbatim}

This removes the guidance pages from the final submitted report while
preserving the regulations file in the \LaTeX{} project for reference.

} % End 13pt content

\clearpage

\end{verbatim}

This removes the guidance pages from the final submitted report while
preserving the regulations file in the \LaTeX{} project for reference.

} % End 13pt content

\clearpage
